% Auto-generated from data/segments.json + layout.json (+ transforms) — do not edit by hand.
% volume: 28Khu11
% mode: printing
% segments: 2518
% transform_rules: 158
% toc_marks: 550

% Volume layout defaults (layout.json ``layout``).
\csromanlayoutapply{1}{1.5}{21.6pt}{8pt}{8pt}{65pt}{2.5em}%

\setcsromanpitaka{สุตฺตนฺตปิฏก}
\setcsromannikaya{ขุทฺทกนิกาย}
\csromanheader{2}{\textbf{ขุทฺทกนิกาย}}
\csromanmatikaanchor{mtk.0}
\csromantocmarknopage{chapter}{มิลินฺทปญฺหปาฬิ}{mtk.0}
\bookmark[dest={mtk.0},level=0]{มิลินฺทปญฺหปาฬิ}
\gambhira{\textbf{มิลินฺทปญฺห}\-\textbf{ปาฬิ}}
\namakkaram{นโม ตสฺส ภควโต อรหโต สมฺมา\-สมฺพุทฺธสฺส.}
\proseitem{1}{\csromanfolio{1}มิลินฺโท นาม โส ราชา, สาคลายํ ปุรุตฺตเม.}

\setlength{\csromangathapagewidth}{0pt}
\setlength{\csromangathawakwidth}{0pt}
\csromangathameasuregroup{อุปคญฺฉิ นาคเสนํ, \\ อาสชฺช ราชา จิตฺรกถิํ, \\ อปุจฺฉิ นิปุเณ ปเญฺห, \\ ปุจฺฉา วิสชฺชนา\textsuperscript{1} เจว, \\ หทยงฺคมา กณฺณสุขา, \\ อภิธมฺมวินโยคาฬฺหา, \\ นาคเสนกถา จิตฺรา, \\ ตตฺถ ญาณํ ปณิธาย, \\ สุณาถ นิปุเณ ปเญฺห,}{\csromangathabat{อุปคญฺฉิ นาคเสนํ,}{คงฺคา จ\textsuperscript{1} ยถา สาครํ.} \\ \csromangathabat{อาสชฺช ราชา จิตฺรกถิํ,}{อุกฺกาธารํ ตโมนุทํ.} \\ \csromangathabat{อปุจฺฉิ นิปุเณ ปเญฺห,}{ฐานาฏฺฐานคเต ปุถู.} \\ \csromangathabat{ปุจฺฉา วิสชฺชนา\textsuperscript{1} เจว,}{คมฺภีรตฺถูปนิสฺสิตา.} \\ \csromangathabat{หทยงฺคมา กณฺณสุขา,}{อพฺภุตา โลมหํสนา.} \\ \csromangathabat{อภิธมฺมวินโยคาฬฺหา,}{สุตฺตชาลสมตฺติตา.} \\ \csromangathabat{นาคเสนกถา จิตฺรา,}{โอปมฺเมหิ นเยหิ จ.} \\ \csromangathabat{ตตฺถ ญาณํ ปณิธาย,}{หาสยิตฺวาน มานสํ.} \\ \csromangathabat{สุณาถ นิปุเณ ปเญฺห,}{กงฺขาฏฺฐานวิทาลเนติ.}}
\csromangathasetleft{อุปคญฺฉิ นาคเสนํ, \\ อาสชฺช ราชา จิตฺรกถิํ, \\ อปุจฺฉิ นิปุเณ ปเญฺห, \\ ปุจฺฉา วิสชฺชนา\textsuperscript{1} เจว, \\ หทยงฺคมา กณฺณสุขา, \\ อภิธมฺมวินโยคาฬฺหา, \\ นาคเสนกถา จิตฺรา, \\ ตตฺถ ญาณํ ปณิธาย, \\ สุณาถ นิปุเณ ปเญฺห,}
\csromangathagroup{\csromangathastanza{\csromangathabat{อุปคญฺฉิ นาคเสนํ,}{คงฺคา จ\csromansharedfootnote{90dbfd120f9a986f}{คงฺคาว (สี, อิ)} ยถา สาครํ.} \\ \csromangathabat{อาสชฺช ราชา จิตฺรกถิํ,}{อุกฺกาธารํ ตโมนุทํ.}}\csromangathastanza{\csromangathabat{อปุจฺฉิ นิปุเณ ปเญฺห,}{ฐานาฏฺฐานคเต ปุถู.} \\ \csromangathabat{ปุจฺฉา วิสชฺชนา\csromansharedfootnote{e3f20fa65a817f75}{วิสฺสชฺชนา (สี, อิ)} เจว,}{คมฺภีรตฺถูปนิสฺสิตา.}}\csromangathastanza{\csromangathabat{หทยงฺคมา กณฺณสุขา,}{อพฺภุตา โลมหํสนา.} \\ \csromangathabat{อภิธมฺมวินโยคาฬฺหา,}{สุตฺตชาลสมตฺติตา.}}\csromangathastanza{\csromangathabat{นาคเสนกถา จิตฺรา,}{โอปมฺเมหิ นเยหิ จ.} \\ \csromangathabat{ตตฺถ ญาณํ ปณิธาย,}{หาสยิตฺวาน มานสํ.} \\ \csromangathabat{สุณาถ นิปุเณ ปเญฺห,}{กงฺขาฏฺฐานวิทาลเนติ.}}}

\proseitem{2}{ตํ ยถานุสูยเต\csromandash{}อตฺถิ โยนกานํ นานา\-ปุฏเภทนํ สาคลํ นาม นครํ นทีปพฺพต\-โสภิตํ รมณีย\-ภูมิปฺปเทส\-ภาคํ อารามุ\-ยฺยาโน\-ปวน\-ตฬาก\-โปกฺขรณิ\-สมฺปนฺนํ นทีปพฺพต\-วนรามเณยฺยกํ สุตวนฺต\-นิมฺมิตํ นิหต\-ปจฺจตฺถิกํ\csromansharedfootnote{aa3c4e5433ceabec}{นิปฺปจฺจตฺถิกํ (ก)} ปจฺจามิตฺตา\-นุป\-ปีฬิตํ วิวิธ\-วิจิตฺร\-ทฬฺหมฏฺฏาลโกฏฺฐกํ วรปวร\-โคปุร\csromansharedfootnote{4f50de5133d7e94a}{ปวรปจุรโคปุร (สี)} โตรณํ คมฺภีร\-ปริขา\-ปณฺฑร\-ปาการปริกฺขิตฺต\-นฺเตปุรํ \csromanfolio{2}สุวิภตฺต\-วีถิ\-จจฺจร\-จตุกฺก\-สิงฺฆาฏกํ สุปฺป\-สาริตา\-เนก\-วิธ\-วรภณฺฑ\-ปริปูริต\-นฺตราปณํ วิวิธ\-ทานคฺค\-สต\-สมุปโสภิตํ\csromansharedfootnote{fa8acdef14f731df}{สตตมุปโสภิตํ (ก)} หิมคิริสิขร\-สงฺกาส\-วร\-ภวน\-สตสหสฺส\-ปฺปฏิมณฺฑิตํ คชหย\-รถ\-ปตฺติ\-สมากุลํ อภิรูป\-นร\-นาริ\-คณา\-นุจริตํ อากิณฺณ\-ชน\-มนุสฺสํ ปุถุขตฺติย\-พฺราหฺมณ\-เวสฺส\-สุทฺทํ วิวิธ\-สมณพฺราหฺมณ\-สภาชน\csromansharedfootnote{c737e4a9b3e52418}{สมภาชน (ก), สมฺมาภาชน (ก)} สงฺฆฏิตํ พหุวิธ\-วิชฺชา\-วนฺต\csromansharedfootnote{40b382c5fe1855d6}{วิชฺชาธร (ก)} นรจิร\csromansharedfootnote{9c4148f2768408a5}{นรวิร (สี, อิ)} นิเสวิตํ กาสิกโกฏุมฺพริกาทินานาวิธวตฺถาปณสมฺปนฺนํ สุปฺป\-สาริต\-รุจิร\-พหุวิธ\-ปุปฺผ\-คนฺธาปณํ คนฺธคนฺธิตํ อาสีส\-นีย\-พหุรตน\-ปริปูริตํ ทิสามุข\-สุปฺป\-สาริตาปณํ สิงฺคาร\-วาณิช\-คณา\-นุจริตํ กหาปณ\-รชต\-สุวณฺณ\-กํส\-ปตฺถร\-ปริปูรํ ปชฺโชต\-มาน\-นิธิ\-นิเกตํ ปหูตธนธญฺญ\-วิตฺตูปกรณํ ปริปุณฺณ\-โกส\-โกฏฺฐาคารํ พหฺวนฺนปานํ พหุวิธ\-ขชฺชโภชฺช\-เลยฺย\-เปยฺย\-สายนียํ อุตฺตรกุรุ\-สงฺกาสํ สมฺปนฺนสสฺสํ อาฬกมนฺทา วิย เทวปุรํ.}

\prose{เอตฺถ ฐตฺวา เตสํ ปุพฺพกมฺมํ กเถตพฺพํ, กเถนฺเตน จ ฉธา วิภชิตฺวา กเถตพฺพํ. เสยฺยถีทํ\csromandash{}ปุพฺพโยโค มิลินฺทปญฺหํ ลกฺขณปญฺหํ เมณฺฑกปญฺหํ อนุมานปญฺหํ โอปมฺมกถาปญฺหนฺติ.}

\prose{ตตฺถ มิลินฺทปโญฺห ลกฺขณปโญฺห, วิมติจฺเฉทนปโญฺหติ ทุวิโธ. เมณฺฑกปโญฺหปิ มหาวคฺโค, โยคิ\-กถา\-ป\-โญฺหติ ทุวิโธ.}

\csromanmatikaanchor{mtk.1}
\csromanmatikaanchor{mtk.2}
\csromantocmarknopage{chapter}{1. พาหิรกถา}{mtk.1}
\bookmark[dest={mtk.1},level=0]{1. พาหิรกถา}
\csromantocmarkref{section}{ปุพฺพโยคาทิ}{mtk.2}
\bookmark[dest={mtk.2},level=1]{ปุพฺพโยคาทิ}
\csromanheader{1}{\textbf{ปุพฺพโยคาทิ}}
\proseitem{3}{\csromanfolio{3}\textbf{ปุพฺพโยโค}ติ เตสํ ปุพฺพกมฺมํ. อตีเต กิร กสฺสปสฺส ภควโต สาสเน วตฺตมาเน คงฺคาย สมีเป เอกสฺมิํ อาวาเส มหา\-ภิกฺขุสํโฆ ปฏิวสติ, ตตฺถ วตฺต\-สีลสมฺปนฺนา ภิกฺขู ปาโตว อุฏฺฐาย ยฏฺฐิ\-สมฺมชฺชนิโย\csromansharedfootnote{9071b4cb2c4eb34b}{ยฏฺฐิสมฺมุญฺชนิโย (สี, อิ)} อาทาย พุทฺธคุเณ อาวชฺเชนฺตา องฺคณํ สมฺมชฺชิตฺวา กจวรพฺยูหํ กโรนฺติ. อเถโก ภิกฺขุ เอกํ สามเณรํ “เอหิ สามเณร, อิมํ กจวรํ ฉฑฺเฑหี”ติ อาห, โส อสุณนฺโต วิย คจฺฉติ, โส ทุติยมฺปิ. ตติยมฺปิ อามนฺติยมาโน อสุณนฺโต วิย คจฺฉเตว. ตโต โส ภิกฺขุ “ทุพฺพโจ วตายํ สามเณโร”ติ กุทฺโธ สมฺมชฺชนิ\-ทณฺเฑน ปหารํ อทาสิ. ตโต โส โรทนฺโต ภเยน กจวรํ ฉฑฺเฑนฺโต “อิมินา กจวร\-ฉฑฺฑน\-ปุญฺญกมฺเมน ยาวาหํ นิพฺพานํ ปาปุณามิ\csromansharedfootnote{e0e1cb7890c0f451}{น ปาปุณามิ (สฺยา)}, เอตฺถนฺตเร นิพฺพตฺต\-นิพฺพตฺต\-ฏฺฐาเน มชฺฌนฺหิก\-สูริโย\csromansharedfootnote{eed3549d63969832}{สุริโย (สี, อิ)} วิย มเหสกฺโข มหาเตโช ภเวยฺยนฺ”ติ ปฐมํ ปตฺถนํ ปฏฺฐเปสิ. กจวรํ ฉฑฺเฑตฺวา นหานตฺถาย คงฺคาติตฺถํ คโต คงฺคาย อูมิเวคํ คคฺครายมานํ ทิสฺวา “ยาวาหํ นิพฺพานํ ปาปุณามิ\csromansharedfootnote{e0e1cb7890c0f451}{น ปาปุณามิ (สฺยา)}, เอตฺถนฺตเร นิพฺพตฺต\-นิพฺพตฺต\-ฏฺฐาเน อยํ อูมิเวโค วิย ฐานุปฺปตฺติตปฏิภาโน ภเวยฺยํ อกฺขย\-ปฏิภาโน”ติ ทุติยมฺปิ ปตฺถนํ ปฏฺฐเปสิ.}

\prose{โสปิ ภิกฺขุ สมฺมชฺชนิ\-สาลาย สมฺมชฺชนิํ ฐเปตฺวา นหานตฺถาย คงฺคาติตฺถํ คจฺฉนฺโต สามเณรสฺส ปตฺถนํ สุตฺวา “เอส มยา ปโยชิโตปิ ตาว เอวํ ปตฺเถติ, มยฺหํ กิํ น สมิชฺฌิสฺสตี”ติ จินฺเตตฺวา “ยาวาหํ นิพฺพานํ ปาปุณามิ\csromansharedfootnote{e0e1cb7890c0f451}{น ปาปุณามิ (สฺยา)}, เอตฺถนฺตเร นิพฺพตฺต\-นิพฺพตฺต\-ฏฺฐาเน อยํ คงฺคา- อูมิเวโค วิย อกฺขย\-ปฏิภาโน ภเวยฺยํ, อิมินา ปุจฺฉิต\-ปุจฺฉิตํ สพฺพํ ปญฺห\-ปฏิภานํ วิชเฏตุํ นิพฺเพเฐตุํ สมตฺโถ ภเวยฺยนฺ”ติ ปตฺถนํ ปฏฺฐเปสิ.}

\prose{เต อุโภปิ เทเวสุ จ มนุสฺเสสุ จ สํสรนฺตา เอกํ พุทฺธนฺตรํ เขเปสุํ. อถ อมฺหากํ ภควตาปิ ยถา โมคฺคลิปุตฺต\-ติสฺสตฺเถโร ทิสฺสติ, เอวเมเตปิ ทิสฺสนฺติ มม ปรินิพฺพานโต ปญฺจวสฺสสเต อติกฺกนฺเต เอเต อุปฺปชฺชิสฺสนฺติ, ยํ มยา สุขุมํ กตฺวา เทสิตํ ธมฺมวินยํ, ตํ เอเต ปญฺหปุจฺฉนโอปมฺมยุตฺติวเสน นิชฺชฏํ นิคฺคุมฺพํ กตฺวา วิภชิสฺสนฺตีติ นิทฺทิฏฺฐา.}

\proseitem{4}{\csromanfolio{4}เตสุ สามเณโร ชมฺพุทีเป สาคลนคเร มิลินฺโท นาม ราชา อโหสิ ปณฺฑิโต พฺยตฺโต เมธาวี ปฏิพโล อตีตา\-นาคต\-ปจฺจุปฺปนฺนานํ มนฺต\-โยค\-วิธาน\-กิริยานํ\csromansharedfootnote{144d37e3b9b314bb}{สมนฺตโยค... (สี, อิ)}, กรณกาเล นิสมฺมการี โหติ, พหูนิ จสฺส สตฺถานิ อุคฺคหิตานิ โหนฺติ. เสยฺยถีทํ, สุติ สมฺมุติ สงฺขฺยา โยคา นีติ วิเสสิกา คณิกา คนฺธพฺพา ติกิจฺฉา จตุพฺเพทา\csromansharedfootnote{ee078a738e43bceb}{ธนุพฺเพทา (สี, อิ)} ปุราณา อิติหาสา โชติสา มายา เกตุ\csromansharedfootnote{18d1cf1fc40b5cf9}{เหตุ (สี, อิ)} มนฺตนา ยุทฺธา ฉนฺทสา พุทฺธวจเนน\csromansharedfootnote{0a1b5d11fba31734}{ฉนฺทสามุทฺทวจเนน (สี, อิ)} เอกูนวีสติ, วิตณฺฑวาที\csromansharedfootnote{fbccac1f2a38ca77}{วาที (สี, อิ)} ทุราสโท ทุปฺปสโห ปุถุติตฺถกรานํ อคฺคมกฺขายติ, สกล\-ชมฺพุทีเป มิลินฺเทน รญฺญา สโม โกจิ นาโหสิ ยทิทํ ถาเมน ชเวน สูเรน ปญฺญาย, อฑฺโฒ มหทฺธโน มหาโภโค อนนฺต\-พล\-วาหโน.}

\proseitem{5}{อเถกทิวสํ มิลินฺโท ราชา อนนฺต\-พลวาหนํ จตุรงฺคินิํ พลคฺค\-เสนาพฺยูหํ ทสฺสน\-กมฺยตาย นครา นิกฺขมิตฺวา พหินคเร เสนงฺค\-ทสฺสนํ กตฺวา\csromansharedfootnote{80c831fb6b7e0b50}{เสนาคณนํ กาเรตฺวา (สี, อิ)} สาเรตฺวา โส ราชา ภสฺส\-ปฺปวาทโก โลกายตวิตณฺฑชนสลฺลาปปฺลว\csromansharedfootnote{20f30176421a60f3}{ปวตฺต (สี, อิ)} จิตฺตโกตูหโล วิสารโท วิชมฺภโก สูริยํ โอโลเกตฺวา อมจฺเจ อามนฺเตสิ “พหุ ภเณ ตาว ทิวสาวเสโส กิํ กริสฺสาม, อิทาเนว นครํ ปวิสิตฺวา อตฺถิ โกจิ ปณฺฑิโต สมโณ วา พฺราหฺมโณ วา สํฆี คณี คณาจริโย อปิ อรหนฺตํ สมฺมา\-สมฺพุทฺธํ ปฏิชานมาโน, โย มยา สทฺธิํ สลฺลปิตุํ สกฺโกติ กงฺขํ ปฏิวิเนตุํ, ตํ อุปสงฺกมิตฺวา ปญฺหํ ปุจฺฉิสฺสาม, กงฺขํ ปฏิวินยิสฺสามา”ติ.}

\prose{เอวํ วุตฺเต ปญฺจสตา โยนกา ราชานํ มิลินฺทํ เอตทโวจุํ “อตฺถิ มหาราช ฉ สตฺถาโร ปูรโณ กสฺสโป มกฺขลิโคสาโล นิคณฺโฐ นาฏปุตฺโต\csromansharedfootnote{be303b1eabc46dd2}{นาถปุตฺโต (สี, อิ)} สญฺชโย เพลฏฺฐปุตฺโต อชิโต เกสกมฺพโล ปกุโธ กจฺจายโน, เต สํฆิโน คณิโน คณาจริยกา ญาตา ยสสฺสิโน ติตฺถกรา สาธุสมฺมตา พหุชนสฺส, คจฺฉ ตฺวํ มหาราช, เต ปญฺหํ ปุจฺฉสฺสุ, กงฺขํ ปฏิวินยสฺสู”ติ.}

\proseitem{6}{อถ โข มิลินฺโท ราชา ปญฺจหิ โยนกสเตหิ ปริวุโต ภทฺรวาหนํ รถวรมา\-รุยฺห เยน ปูรโณ กสฺสโป เตนุปสงฺกมิ, \csromanfolio{5}อุปสงฺกมิตฺวา ปูรเณน กสฺสเปน สทฺธิํ สมฺโมทิ, สมฺโมทนียํ กถํ สารณียํ วีติสาเรตฺวา เอกมนฺตํ นิสีทิ, เอกมนฺตํ นิสินฺโน โข มิลินฺโท ราชา ปูรณํ กสฺสปํ เอตทโวจ “โก ภนฺเต กสฺสป โลกํ ปาเลตี”ติ. ปถวี\csromansharedfootnote{832d2c542b555d38}{ปฐวี (สี, สฺยา, อิ)} มหาราช โลกํ ปาเลตีติ. ยทิ ภนฺเต กสฺสป ปถวี โลกํ ปาเลติ, อถ กสฺมา อวีจินิรยํ คจฺฉนฺตา สตฺตา ปถวิํ อติกฺกมิตฺวา คจฺฉนฺตีติ. เอวํ วุตฺเต ปูรโณ กสฺสโป เนว สกฺขิ โอคิลิตุํ, โน สกฺขิ อุคฺคิลิตุํ, อโธมุโข ปตฺตกฺขนฺโธ ตุณฺหีภูโต ปชฺฌายนฺโต นิสีทิ.}

\proseitem{7}{อถ โข มิลินฺโท ราชา มกฺขลิํ โคสาลํ เอตทโวจ “อตฺถิ ภนฺเต โคสาล กุสลากุสลานิ กมฺมานิ, อตฺถิ สุกต\-ทุกฺกฏานํ กมฺมานํ ผลํ วิปาโก”ติ. นตฺถิ มหาราช กุสลากุสลานิ กมฺมานิ, นตฺถิ สุกต\-ทุกฺกฏานํ กมฺมานํ ผลํ วิปาโก, เย เต มหาราช อิธ โลเก ขตฺติยา, เต ปรโลกํ คนฺตฺวาปิ ปุน ขตฺติยาว ภวิสฺสนฺติ, เย เต พฺราหฺมณา เวสฺสา สุทฺทา จณฺฑาลา ปุกฺกุสา, เต ปรโลกํ คนฺตฺวาปิ ปุน พฺราหฺมณา เวสฺสา สุทฺทา จณฺฑาลา ปุกฺกุสาว ภวิสฺสนฺติ, กิํ กุสลากุสเลหิ กมฺเมหีติ. ยทิ ภนฺเต โคสาล อิธ โลเก ขตฺติยา พฺราหฺมณา เวสฺสา สุทฺทา จณฺฑาลา ปุกฺกุสา, เต ปรโลกํ คนฺตฺวาปิ ปุน ขตฺติยา พฺราหฺมณา เวสฺสา สุทฺทา จณฺฑาลา ปุกฺกุสาว ภวิสฺสนฺติ, นตฺถิ กุสลากุสเลหิ กมฺเมหิ กรณียํ. เตน หิ ภนฺเต โคสาล เย เต อิธ โลเก หตฺถจฺฉินฺนา, เต ปรโลกํ คนฺตฺวาปิ ปุน หตฺถจฺฉินฺนาว ภวิสฺสนฺติ. เย ปาทจฺฉินฺนา, เต ปาทจฺฉินฺนาว ภวิสฺสนฺติ. เย หตฺถ\-ปาท\-จฺฉินฺนา, เต หตฺถ\-ปาท\-จฺฉินฺนาว ภวิสฺสนฺติ. เย กณฺณจฺฉินฺนา, เต กณฺณจฺฉินฺนาว ภวิสฺสนฺติ. เย นาสจฺฉินฺนา, เต นาสจฺฉินฺนาว ภวิสฺสนฺติ. เย กณฺณ\-นาส\-จฺฉินฺนา, เต กณฺณ\-นาส\-จฺฉินฺนาว ภวิสฺสนฺตีติ. เอวํ วุตฺเต โคสาโล ตุณฺหี อโหสิ.}

\prose{อถ โข มิลินฺทสฺส รญฺโญ เอตทโหสิ “ตุจฺโฉ วต โภ ชมฺพุทีโป, ปลาโป วต โภ ชมฺพุทีโป, นตฺถิ โกจิ สมโณ วา พฺราหฺมโณ วา, โย มยา สทฺธิํ สลฺลปิตุํ สกฺโกติ กงฺขํ ปฏิวิเนตุนฺ”ติ.}

\prose{\csromanfolio{6}อถ โข มิลินฺโท ราชา อมจฺเจ อามนฺเตสิ “รมณียา วต โภ โทสินา รตฺติ, กํ นุ ขฺวชฺช สมณํ วา พฺราหฺมณํ วา อุปสงฺกเมยฺยาม ปญฺหํ ปุจฺฉิตุํ, โก มยา สทฺธิํ สลฺลปิตุํ สกฺโกติ กงฺขํ ปฏิวิเนตุนฺ”ติ. เอวํ วุตฺเต อมจฺจา ตุณฺหีภูตา รญฺโญ มุขํ โอโลกยมานา อฏฺฐํสุ.}

\prose{เตน โข ปน สมเยน สาคลนครํ ทฺวาทส วสฺสานิ สุญฺญํ อโหสิ สมณพฺราหฺมณ\-คหปติปณฺฑิเตหิ, ยตฺถ สมณพฺราหฺมณ\-คหปติ\-ปณฺฑิตา ปฏิวสนฺตีติ สุณาติ, ตตฺถ คนฺตฺวา ราชา เต ปญฺหํ ปุจฺฉติ, เต สพฺเพปิ ปญฺห\-วิสชฺช\-เนน ราชานํ อาราเธตุํ อสกฺโกนฺตา เยน วา เตน วา ปกฺกมนฺติ. เย อญฺญํ ทิสํ น ปกฺกมนฺติ, เต สพฺเพ ตุณฺหีภูตา อจฺฉนฺติ. ภิกฺขู ปน เยภุยฺเยน หิมวนฺตเมว คจฺฉนฺติ.}

\proseitem{8}{เตน โข ปน สมเยน โกฏิสตา อรหนฺโต หิมวนฺเต ปพฺพเต รกฺขิตตเล ปฏิวสนฺติ. อถ โข อายสฺมา อสฺสคุตฺโต ทิพฺพาย โสตธาตุยา มิลินฺทสฺส รญฺโญ วจนํ สุตฺวา ยุคนฺธร\-มตฺถเก ภิกฺขุสํฆํ สนฺนิปาเตตฺวา ภิกฺขู ปุจฺฉิ “อตฺถาวุโส โกจิ ภิกฺขุ ปฏิพโล มิลินฺเทน รญฺญา สทฺธิํ สลฺลปิตุํ กงฺขํ ปฏิวิเนตุนฺ”ติ.}

\prose{เอวํ วุตฺเต โกฏิสตา อรหนฺโต ตุณฺหี อเหสุํ. ทุติยมฺปิ. ตติยมฺปิ ปุฏฺฐา ตุณฺหี อเหสุํ. อถ โข อายสฺมา อสฺสคุตฺโต ภิกฺขุสํฆํ เอตทโวจ “อตฺถาวุโส ตาวติํส\-ภวเน เวชยนฺตสฺส ปาจีนโต \textbf{เกตุมตี} นาม วิมานํ, ตตฺถ \textbf{มหาเสโน} นาม เทวปุตฺโต ปฏิวสติ, โส ปฏิพโล เตน มิลินฺเทน รญฺญา สทฺธิํ สลฺลปิตุํ กงฺขํ ปฏิวิเนตุนฺ”ติ.}

\prose{อถ โข โกฏิสตา อรหนฺโต ยุคนฺธร\-ปพฺพเต อนฺตรหิตา ตาวติํส\-ภวเน ปาตุรเหสุํ. อทฺทสา โข สกฺโก เทวานมินฺโท เต ภิกฺขู ทูรโตว อาคจฺฉนฺเต, ทิสฺวาน เยนายสฺมา อสฺสคุตฺโต เตนุปสงฺกมิ, อุปสงฺกมิตฺวา อายสฺมนฺตํ อสฺสคุตฺตํ อภิวาเทตฺวา เอกมนฺตํ อฏฺฐาสิ, เอกมนฺตํ ฐิโต โข สกฺโก เทวานมินฺโท อายสฺมนฺตํ อสฺสคุตฺตํ เอตทโวจ “มหา โข ภนฺเต ภิกฺขุสํโฆ อนุปฺปตฺโต, อหํ สํฆสฺส อารามิโก, เกนตฺโถ, กิํ มยา กรณียนฺ”ติ.}

\prose{อถ โข อายสฺมา อสฺสคุตฺโต สกฺกํ เทวานมินฺทํ เอตทโวจ “อยํ โข มหาราช ชมฺพุทีเป สาคลนคเร \textbf{มิลินฺโท} นาม ราชา วิตณฺฑวาที \csromanfolio{7}ทุราสโท ทุปฺปสโห ปุถุติตฺถกรานํ อคฺคมกฺขายติ, โส ภิกฺขุสํฆํ อุปสงฺกมิตฺวา ทิฏฺฐิวาเทน ปญฺหํ ปุจฺฉิตฺวา ภิกฺขุสํฆํ วิเหเฐตี”ติ.}

\prose{อถ โข สกฺโก เทวานมินฺโท อายสฺมนฺตํ อสฺสคุตฺตํ เอตทโวจ “อยํ โข ภนฺเต มิลินฺโท ราชา อิโต จุโต มนุสฺเสสุ อุปฺปนฺโน, เอโส โข ภนฺเต เกตุมติวิมาเน มหาเสโน นาม เทวปุตฺโต ปฏิวสติ, โส ปฏิพโล เตน มิลินฺเทน รญฺญา สทฺธิํ สลฺลปิตุํ กงฺขํ ปฏิวิเนตุํ, ตํ เทวปุตฺตํ ยาจิสฺสาม มนุสฺสโลกู\-ปปตฺติยา”ติ.}

\prose{อถ โข สกฺโก เทวานมินฺโท ภิกฺขุสํฆํ ปุรกฺขตฺวา เกตุมติ\-วิมานํ ปวิสิตฺวา มหาเสนํ เทวปุตฺตํ อาลิงฺคิตฺวา เอตทโวจ “ยาจติ ตํ มาริส ภิกฺขุสํโฆ มนุสฺสโลกู\-ปปตฺติยา”ติ. น เม ภนฺเต มนุสฺสโลเกน\-ตฺโถ กมฺมพหุเลน, ติพฺโพ มนุสฺสโลโก, อิเธวาหํ ภนฺเต เทวโลเก อุปรูปรูปปตฺติโก หุตฺวา ปรินิพฺพายิสฺสามี”ติ. ทุติยมฺปิ. ตติยมฺปิ โข สกฺเกน เทวานมินฺเทน ยาจิโต มหาเสโน เทวปุตฺโต เอวมาห “น เม ภนฺเต มนุสฺสโลเกน\-ตฺโถ กมฺมพหุเลน, ติพฺโพ มนุสฺสโลโก, อิเธวาหํ ภนฺเต เทวโลเก อุปรูปรูปปตฺติโก หุตฺวา ปรินิพฺพายิสฺสามี”ติ.}

\prose{อถ โข อายสฺมา อสฺสคุตฺโต มหาเสนํ เทวปุตฺตํ เอตทโวจ “อิธ มยํ มาริส สเทวกํ โลกํ อนุวิโลกยมานา อญฺญตฺร ตยา มิลินฺทสฺส รญฺโญ วาทํ ภินฺทิตฺวา สาสนํ ปคฺคเหตุํ สมตฺถํ อญฺญํ กญฺจิ น ปสฺสาม, ยาจติ ตํ มาริส ภิกฺขุสํโฆ, สาธุ สปฺปุริส มนุสฺสโลเก นิพฺพตฺติตฺวา ทสพลสฺส สาสนํ ปคฺคณฺหาหี”ติ. เอวํ วุตฺเต มหาเสโน เทวปุตฺโต “อหํ กิร มิลินฺทสฺส รญฺโญ วาทํ ภินฺทิตฺวา พุทฺธสาสนํ ปคฺคเหตุํ สมตฺโถ ภวิสฺสามี”ติ หฏฺฐปหฏฺโฐ อุทคฺคุทคฺโค หุตฺวา “สาธุ ภนฺเต มนุสฺสโลเก อุปฺปชฺชิสฺสามี”ติ ปฏิญฺญํ อทาสิ.}

\proseitem{9}{อถ โข เต ภิกฺขู เทวโลเก ตํ กรณียํ ตีเรตฺวา เทเวสุ ตาวติํเสสุ อนฺตรหิตา หิมวนฺเต ปพฺพเต รกฺขิตตเล ปาตุรเหสุํ.}

\prose{อถ โข อายสฺมา อสฺสคุตฺโต ภิกฺขุสํฆํ เอตทโวจ “อตฺถาวุโส อิมสฺมิํ ภิกฺขุสํเฆ โกจิ ภิกฺขุ สนฺนิปาตํ อนาคโต”ติ. เอวํ วุตฺเต อญฺญตโร ภิกฺขุ อายสฺมนฺตํ อสฺสคุตฺตํ เอตทโวจ “อตฺถิ \csromanfolio{8}ภนฺเต อายสฺมา \textbf{โรหโณ} อิโต สตฺตเม ทิวเส หิมวนฺตํ ปพฺพตํ ปวิสิตฺวา นิโรธํ สมาปนฺโน, ตสฺส สนฺติเก ทูตํ ปาเหถา”ติ. อายสฺมาปิ โรหโณ ตงฺขณญฺเญว นิโรธา วุฏฺฐาย “สํโฆ มํ ปฏิมาเนตี”ติ หิมวนฺเต ปพฺพเต อนฺตรหิโต รกฺขิตตเล โกฏิสตานํ อรหนฺตานํ ปุรโต ปาตุรโหสิ.}

\prose{อถ โข อายสฺมา อสฺสคุตฺโต อายสฺมนฺตํ โรหณํ เอตทโวจ “กิํ นุ โข อาวุโส โรหณ พุทฺธสาสเน ภิชฺชนฺเต\csromansharedfootnote{771a728a7cb3fa19}{ปลุชฺชนฺเต (สี, อิ)} น ปสฺสสิ สํฆสฺส กรณียานี”ติ. อมนสิกาโร เม ภนฺเต อโหสีติ.}

\prose{เตน หาวุโส โรหณ ทณฺฑกมฺมํ กโรหีติ. กิํ ภนฺเต กโรมีติ. อตฺถาวุโส โรหณ หิมวนฺต\-ปพฺพ\-ตปสฺเส \textbf{คชงฺคลํ}\csromansharedfootnote{6a3fd9326fe99ee9}{กชงฺคลํ (สี, อิ)} นาม พฺราหฺมณคาโม, ตตฺถ \textbf{โสณุตฺตโร} นาม พฺราหฺมโณ ปฏิวสติ, ตสฺส ปุตฺโต อุปฺปชฺชิสฺสติ \textbf{นาคเสโน} นาม ทารโก, เตน หิ ตฺวํ อาวุโส โรหณ ทสมา\-สาธิกานิ สตฺต วสฺสานิ ตํ กุลํ ปิณฺฑาย ปวิสิตฺวา นาคเสนํ ทารกํ นีหริตฺวา ปพฺพาเชหิ, ปพฺพชิเตว ตสฺมิํ ทณฺฑกมฺมโต มุจฺจิสฺสสีติ. อายสฺมาปิ โข โรหโณ “สาธู”ติ สมฺปฏิจฺฉิ.}

\proseitem{10}{มหาเสโนปิ โข เทวปุตฺโต เทวโลกา จวิตฺวา โสณุตฺตร\-พฺราหฺมณสฺส ภริยาย กุจฺฉิสฺมิํ ปฏิสนฺธิํ อคฺคเหสิ, สห ปฏิสนฺธิ\-คฺคหณา ตโย อจฺฉริยา อพฺภุตา ธมฺมา ปาตุรเหสุํ, อาวุธภณฺฑานิ ปชฺชลิํสุ, อคฺคสสฺสํ อภินิปฺผนฺนํ, มหาเมโฆ อภิปฺปวสฺสิ. อายสฺมาปิ โข โรหโณ ตสฺส ปฏิสนฺธิ\-คฺคหณโต ปฏฺฐาย ทสมา\-สาธิกานิ สตฺต วสฺสานิ ตํ กุลํ ปิณฺฑาย ปวิสนฺโต เอกทิวสมฺปิ กฏจฺฉุมตฺตํ ภตฺตํ วา อุฬุงฺกมตฺตํ ยาคุํ วา อภิวาทนํ วา อญฺชลิกมฺมํ วา สามีจิกมฺมํ วา นาลตฺถ, อถ โข อกฺโกสญฺเญว ปริภาสญฺเญว ปฏิลภติ “อติจฺฉถ ภนฺเต”ติ วจนมตฺตมฺปิ วตฺตา นาม นาโหสิ, ทสมา\-สาธิกานํ ปน สตฺตนฺนํ วสฺสานํ อจฺจเยน เอกทิวสํ “อติจฺฉถ ภนฺเต”ติ วจนมตฺตํ อลตฺถ. ตํ ทิวสเมว พฺราหฺมโณปิ พหิ กมฺมนฺตา อาคจฺฉนฺโต ปฏิปเถ เถรํ ทิสฺวา “กิํ โภ ปพฺพชิต อมฺหากํ เคหํ อคมิตฺถา”ติ อาห. อาม พฺราหฺมณ อคมมฺหาติ. อปิ กิญฺจิ ลภิตฺถาติ. อาม พฺราหฺมณ ลภิมฺหาติ. โส อนตฺตมโน เคหํ คนฺตฺวา ปุจฺฉิ \csromanfolio{9}“ตสฺส ปพฺพชิตสฺส กิญฺจิ อทตฺถา”ติ. น กิญฺจิ อทมฺหาติ. พฺราหฺมโณ ทุติยทิวเส ฆรทฺวาเร เยว นิสีทิ “อชฺช ปพฺพชิตํ มุสาวาเทน นิคฺคเหสฺสามี”ติ. เถโร ทุติยทิวเส พฺราหฺมณสฺส ฆรทฺวารํ สมฺปตฺโต.}

\prose{พฺราหฺมโณ เถรํ ทิสฺวาว เอวมาห “ตุเมฺห หิยฺโย อมฺหากํ เคเห กิญฺจิ อลภิตฺวาว “ลภิมฺหา”ติ อโวจุตฺถ, วฏฺฏติ นุ โข ตุมฺหากํ มุสาวาโท”ติ. เถโร อาห “มยํ พฺราหฺมณ ตุมฺหากํ เคเห ( )\csromansharedfootnote{7bafc82681165472}{(ปวิสนฺตา) (ก)} ทสมา\-สาธิกานิ สตฺต วสฺสานิ ‘อติจฺฉถา’ติ วจนมตฺตมฺปิ อลภิตฺวา หิยฺโย ‘อติจฺฉถา’ติ วจนมตฺตํ ลภิมฺหา, อเถตํ วาจาปฏิสนฺธารํ\csromansharedfootnote{e5cfc26636ee27a3}{ปฏิสนฺถารํ (สี, อิ)} อุปาทาย เอวมโวจุมฺหา”ติ.}

\prose{พฺราหฺมโณ จินฺเตสิ “อิเม วาจาปฏิสนฺธารมตฺตมฺปิ ลภิตฺวา ชนมชฺเฌ ‘ลภิมฺหา’ติ ปสํสนฺติ, อญฺญํ กิญฺจิ ขาทนียํ วา โภชนียํ วา ลภิตฺวา กสฺมา นปฺปสํสนฺตี”ติ ปสีทิตฺวา อตฺตโน อตฺถาย ปฏิยาทิต\-ภตฺตโต กฏจฺฉุภิกฺขํ, ตทุปิยญฺจ พฺยญฺชนํ ทาเปตฺวา “อิมํ ภิกฺขํ สพฺพกาลํ ตุเมฺห ลภิสฺสถา”ติ อาห.}

\prose{โส ปุนทิวสโต ปภุติ อุปส\-งฺกม\-นฺตสฺส เถรสฺส อุปสมํ ทิสฺวา ภิยฺโยโส มตฺตาย ปสีทิตฺวา เถรํ นิจฺจกาลํ อตฺตโน ฆเร ภตฺต\-วิสฺสคฺค\-กรณ\-ตฺถาย ยาจิ. เถโร ตุณฺหีภาเวน อธิวาเสตฺวา ทิวเส ทิวเส ภตฺตกิจฺจํ กตฺวา คจฺฉนฺโต โถกํ โถกํ พุทฺธวจนํ กเถตฺวา คจฺฉติ. สาปิ โข พฺราหฺมณี ทสมาส\-จฺจเยน ปุตฺตํ วิชายิ, “นาคเสโน”ติสฺส นามมกํสุ, โส อนุกฺกเมน วฑฺฒนฺโต สตฺตวสฺสิโก ชาโต.}

\proseitem{11}{อถ โข นาคเสนสฺส ทารกสฺส ปิตา นาคเสนํ ทารกํ เอตทโวจ “อิมสฺมิํ โข ตาต นาคเสน พฺราหฺมณกุเล สิกฺขานิ สิกฺเขยฺยาสี”ติ. กตมานิ ตาต อิมสฺมิํ พฺราหฺมณกุเล สิกฺขานิ นามาติ. ตโย โข ตาต นาคเสน เวทา \textbf{สิกฺขานิ} นาม, อวเสสานิ สิปฺปานิ \textbf{สิปฺปํ} นามาติ. เตน หิ ตาต สิกฺขิสฺสามีติ.}

\prose{อถ โข โสณุตฺตโร พฺราหฺมโณ อาจริย\-พฺราหฺมณสฺส อาจริยภาคํ สหสฺสํ ทตฺวา อนฺโตปาสาเท เอกสฺมิํ คพฺเภ เอกโต มญฺจกํ ปญฺญเปตฺวา อาจริย\-พฺราหฺมณํ เอตทโวจ “สชฺฌาเปหิ โข ตฺวํ พฺราหฺมณ \csromanfolio{10}อิมํ ทารกํ มนฺตานี”ติ. เตน หิ ตาต ทารก อุคฺคณฺหาหิ มนฺตานีติ. อาจริย\-พฺราหฺมโณ สชฺฌายติ นาคเสนสฺส ทารกสฺส เอเกเนว อุทฺเทเสน ตโย เวทา หทยงฺคตา วาจุคฺคตา สูปธาริตา สุววตฺถาปิตา สุมนสิกตา อเหสุํ, สกิเมว จกฺขุํ อุทปาทิ ตีสุ เวเทสุ สนิฆณฺฑุ\-เกฏุเภสุ\csromansharedfootnote{af3a8e01b98d9bba}{สนิฆณฺฏุเกฏุเภสุ (ก)} สากฺขร\-ปฺปเภเทสุ อิติหาส\-ปญฺจเมสุ ปทโก เวยฺยากรโณ โลกายต\-มหา\-ปุริส\-ลกฺขเณสุ อนวโย อโหสิ.}

\prose{อถ โข นาคเสโน ทารโก ปิตรํ เอตทโวจ “อตฺถิ นุ โข ตาต อิมสฺมิํ พฺราหฺมณกุเล อิโต อุตฺตริมฺปิ สิกฺขิตพฺพานิ, อุทาหุ เอตฺตกาเนวา”ติ. นตฺถิ ตาต นาคเสน อิมสฺมิํ พฺราหฺมณกุเล อิโต อุตฺตริํ สิกฺขิตพฺพานิ, เอตฺตกาเนว สิกฺขิตพฺพานีติ.}

\prose{อถ โข นาคเสโน ทารโก อาจริยสฺส อนุโยคํ ทตฺวา ปาสาทา โอรุยฺห ปุพฺพวาสนาย โจทิตหทโย รโหคโต ปฏิสลฺลีโน อตฺตโน สิปฺปสฺส อาทิ\-มชฺฌ\-ปริโยสานํ โอโลเกนฺโต อาทิมฺหิ วา มชฺเฌ วา ปริโยสาเน วา อปฺปมตฺตกมฺปิ สารํ อทิสฺวา “ตุจฺฉา วต โภ อิเม เวทา, ปลาปา วต โภ อิเม เวทา อสารา นิสฺสารา”ติ วิปฺปฏิสารี อนตฺตมโน อโหสิ.}

\proseitem{12}{เตน โข ปน สมเยน อายสฺมา โรหโณ วตฺตนิเย เสนาสเน นิสินฺโน นาคเสนสฺส ทารกสฺส เจตสา เจโต\-ปริวิตกฺกม\-ญฺญาย นิวาเสตฺวา ปตฺต\-จีวรมา\-ทาย วตฺตนิเย เสนาสเน อนฺตรหิโต คชงฺคล\-พฺราหฺมณคามสฺส ปุรโต ปาตุรโหสิ. อทฺทสา โข นาคเสโน ทารโก อตฺตโน ทฺวารโกฏฺฐเก ฐิโต อายสฺมนฺตํ โรหณํ ทูรโตว อาคจฺฉนฺตํ, ทิสฺวาน อตฺตมโน อุทคฺโค ปมุทิโต ปีติโสมนสฺส\-ชาโต “อปฺเปว นามายํ ปพฺพชิโต กญฺจิ สารํ ชาเนยฺยา”ติ เยนายสฺมา โรหโณ เตนุปสงฺกมิ, อุปสงฺกมิตฺวา อายสฺมนฺตํ โรหณํ เอตทโวจ “โก นุ โข ตฺวํ มาริส เอทิโส ภณฺฑุ\-กาสาว\-วสโน”ติ. ปพฺพชิโต นามาหํ ทารกาติ. เกน ตฺวํ มาริส ปพฺพชิโต นามาสีติ. ปาปกานิ มลานิ ปพฺพาเชติ\csromansharedfootnote{280b70d8d2919c5b}{ปาปกานํ มลานํ ปพฺพาเชตุํ ปพฺพชิโต (สี, อิ)}, ตสฺมาหํ ทารก ปพฺพชิโต นามาติ. กิํการณา มาริส เกสา เต น ยถา \csromanfolio{11}อญฺเญสนฺติ. โสฬสิเม ทารก ปลิโพเธ ทิสฺวา เกสมสฺสุํ โอหาเรตฺวา ปพฺพชิโต. กตเม โสฬส, อลงฺการ\-ปลิโพโธ มณฺฑน\-ปลิโพโธ เตลมกฺขน\-ปลิโพโธ โธวน\-ปลิโพโธ มาลาปลิโพโธ คนฺธ\-ปลิโพโธ วาสน\-ปลิโพโธ หรีฏกปลิโพโธ อามลก\-ปลิโพโธ รงฺคปลิโพโธ พนฺธน\-ปลิโพโธ โกจฺฉ\-ปลิโพโธ กปฺปก\-ปลิโพโธ วิชฏน\-ปลิโพโธ อูกาปลิโพโธ, เกเสสุ วิลูเนสุ โสจนฺติ กิลมนฺติ ปริเทวนฺติ อุรตฺตาฬิํ กนฺทนฺติ สมฺโมหํ อาปชฺชนฺติ, อิเมสุ โข ทารก โสฬสสุ ปลิโพเธสุ ปลิคุณฺฐิตา มนุสฺสา สพฺพานิ อติสุขุมานิ สิปฺปานิ นาเสนฺตีติ. กิํการณา มาริส วตฺถานิปิ เต น ยถา อญฺเญสนฺติ. กามนิสฺสิตานิ โข ทารก วตฺถานิ, กามนิสฺสิตานิ คิหิพฺยญฺชน\-ภณฺฑานิ\csromansharedfootnote{011971c8d468f005}{กมนียานิ คิหิพฺยญฺชนานิ (สี, อิ)}, ยานิ กานิจิ โข ภยานิ วตฺถโต อุปฺปชฺชนฺติ, ตานิ กาสาว\-วสนสฺส น โหนฺติ, ตสฺมา วตฺถานิปิ เม น ยถา อญฺเญสนฺติ. ชานาสิ โข ตฺวํ มาริส สิปฺปานิ นามาติ. อาม ทารก, ชานามหํ สิปฺปานิ, ยํ โลเก อุตฺตมํ มนฺตํ, ตมฺปิ ชานามีติ. มยฺหมฺปิ ตํ มาริส ทาตุํ สกฺกาติ. อาม ทารก สกฺกาติ. เตน หิ เม เทหีติ. อกาโล โข ทารก, อนฺตรฆรํ ปิณฺฑาย ปวิฏฺฐมฺหาติ.}

\prose{อถ โข นาคเสโน ทารโก อายสฺมโต โรหณสฺส หตฺถโต ปตฺตํ คเหตฺวา ฆรํ ปเวเสตฺวา ปณีเตน ขาทนีเยน โภชนีเยน สหตฺถา สนฺตปฺเปตฺวา สมฺปวาเรตฺวา อายสฺมนฺตํ โรหณํ ภุตฺตาวิํ โอนีต\-ปตฺต\-ปาณิํ เอตทโวจ “เทหิ เม ทานิ มาริส มนฺตนฺ”ติ. “ยทา โข ตฺวํ ทารก นิปฺปลิโพโธ หุตฺวา มาตาปิตโร อนุชานาเปตฺวา มยา คหิตํ ปพฺพชิตเวสํ คณฺหิสฺสสิ, ตทา ทสฺสามี”ติ อาห.}

\prose{อถ โข นาคเสโน ทารโก มาตาปิตโร อุปสงฺกมิตฺวา อาห “อมฺมตาตา อยํ ปพฺพชิโต ‘ยํ โลเก อุตฺตมํ มนฺตํ, ตํ ชานามี’ติ วทติ, น จ อตฺตโน สนฺติเก อปพฺพชิตสฺส เทติ, อหํ เอตสฺส สนฺติเก ปพฺพชิตฺวา ตํ อุตฺตมํ มนฺตํ อุคฺคณฺหิสฺสามี”ติ. อถสฺส มาตาปิตโร “ปพฺพชิตฺวาปิ โน ปุตฺโต มนฺตํ คณฺหตุ, คเหตฺวา ปุน อาคจฺฉิสฺสตี”ติ มญฺญมานา “คณฺห ปุตฺตา”ติ อนุชานิํสุ.}

\proseitem{13}{\csromanfolio{12}อถ โข อายสฺมา โรหโณ นาคเสนํ ทารกํ อาทาย เยน วตฺตนิยํ เสนาสนํ, เยน วิชมฺภวตฺถุ เตนุปสงฺกมิ, อุปสงฺกมิตฺวา วิชมฺภ\-วตฺถุสฺมิํ เสนาสเน เอกรตฺตํ วสิตฺวา เยน รกฺขิตตลํ เตนุปสงฺกมิ, อุปสงฺกมิตฺวา โกฏิสตานํ อรหนฺตานํ มชฺเฌ นาคเสนํ ทารกํ ปพฺพาเชสิ. ปพฺพชิโต จ ปนายสฺมา นาคเสโน อายสฺมนฺตํ โรหณํ เอตทโวจ “คหิโต เม ภนฺเต ตว เวโส, เทถ เม ทานิ มนฺตนฺ”ติ. อถ โข อายสฺมา โรหโณ “กิมฺหิ นุ โขหํ นาคเสนํ วิเนยฺยํ ปฐมํ วินเย วา สุตฺตนฺเต วา อภิธมฺเม วา”ติ จินฺเตตฺวา ปณฺฑิโต โข อยํ นาคเสโน, สกฺโกติ สุเขเนว อภิธมฺมํ ปริยาปุณิตุนฺติ ปฐมํ อภิธมฺเม วิเนสิ.}

\prose{อายสฺมา จ นาคเสโน “กุสลา ธมฺมา, อกุสลา ธมฺมา, อพฺยากตา ธมฺมา”ติ ติกทุก\-ปฏิมณฺฑิตํ ธมฺมสงฺคณี\-ปกรณํ, ขนฺธ\-วิภงฺคา\-ทิ\-อฏฺฐารส วิภงฺค\-ปฏิมณฺฑิตํ วิภงฺค\-ปฺปกรณํ, “สงฺคโห อสงฺคโห”ติอาทินา จุทฺทสวิเธน วิภตฺตํ ธาตุกถา\-ปกรณํ, “ขนฺธ\-ปญฺญตฺติ อายตนปญฺญตฺตี”ติอาทินา ฉพฺพิเธน วิภตฺตํ ปุคฺคลปญฺญตฺติ\-ปฺปกรณํ, สกวาเท ปญฺจ\-สุตฺตสตานิ ปรวาเท ปญฺจ\-สุตฺตสตานีติ สุตฺตสหสฺสํ สโมธาเนตฺวา วิภตฺตํ กถาวตฺถุ\-ปฺปกรณํ, “มูลยมกํ ขนฺธยมกนฺ”ติ- อาทินา ทสวิเธน วิภตฺตํ ยมก\-ปฺปกรณํ, “เหตุปจฺจโย อารมฺมณ\-ปจฺจโย”ติอาทินา จตุวีสติวิเธน วิภตฺตํ ปฏฺฐานปฺปกรณนฺติ สพฺพํ ตํ อภิธมฺม\-ปิฏกํ เอเกเนว สชฺฌาเยน ปคุณํ กตฺวา “ติฏฺฐถ ภนฺเต, น ปุน โอสาเรถ, เอตฺตเกเนวาหํ สชฺฌายิสฺสามี”ติ อาห.}

\proseitem{14}{อถ โข อายสฺมา นาคเสโน เยน โกฏิสตา อรหนฺโต เตนุปสงฺกมิ, อุปสงฺกมิตฺวา โกฏิสเต อรหนฺเต เอตทโวจ “อหํ โข ภนฺเต ‘กุสลา ธมฺมา, อกุสลา ธมฺมา, อพฺยากตา ธมฺมา’ติ อิเมสุ ตีสุ ปเทสุ ปกฺขิปิตฺวา สพฺพํ ตํ อภิธมฺม\-ปิฏกํ วิตฺถาเรน โอสาเรสฺสามี”ติ. สาธุ นาคเสน โอสาเรหีติ.}

\prose{อถ โข อายสฺมา นาคเสโน สตฺต มาสานิ สตฺต ปกรณานิ วิตฺถาเรน โอสาเรสิ, ปถวี อุนฺนทิ, เทวตา สาธุการมทํสุ, พฺรหฺมาโน อปฺโผเฏสุํ, ทิพฺพานิ จนฺทนจุณฺณานิ ทิพฺพานิ จ มนฺทารว\-ปุปฺผานิ อภิปฺปวสฺสิํสุ.}

\proseitem{15}{\csromanfolio{13}อถ โข โกฏิสตา อรหนฺโต อายสฺมนฺตํ นาคเสนํ ปริปุณฺณ\-วีสติวสฺสํ รกฺขิตตเล อุปสมฺปาเทสุํ. อุปสมฺปนฺโน จ ปนายสฺมา นาคเสโน ตสฺสา รตฺติยา อจฺจเยน ปุพฺพณฺหสมยํ นิวาเสตฺวา ปตฺต\-จีวรมา\-ทาย อุปชฺฌาเยน สทฺธิํ คามํ ปิณฺฑาย ปวิสนฺโต เอวรูปํ ปริวิตกฺกํ อุปฺปาเทสิ “ตุจฺโฉ วต เม อุปชฺฌาโย, พาโล วต เม อุปชฺฌาโย, ฐเปตฺวา อวเสสํ พุทฺธวจนํ ปฐมํ มํ อภิธมฺเม วิเนสี”ติ.}

\prose{อถ โข อายสฺมา โรหโณ อายสฺมโต นาคเสนสฺส เจตสา เจโต\-ปริวิตกฺกม\-ญฺญาย อายสฺมนฺตํ นาคเสนํ เอตทโวจ “อนนุจฺฉวิกํ โข นาคเสน ปริวิตกฺกํ วิตกฺเกสิ, น โข ปเนตํ นาคเสน ตวานุจฺฉวิกนฺ”ติ.}

\prose{อถ โข อายสฺมโต นาคเสนสฺส เอตทโหสิ “อจฺฉริยํ วต โภ, อพฺภุตํ วต โภ, ยตฺร หิ นาม เม อุปชฺฌาโย เจตสา เจโต\-ปริวิตกฺกํ ชานิสฺสติ, ปณฺฑิโต วต เม อุปชฺฌาโย, ยนฺนูนาหํ อุปชฺฌายํ ขมาเปยฺยนฺ”ติ. อถ โข อายสฺมา นาคเสโน อายสฺมนฺตํ โรหณํ เอตทโวจ “ขมถ เม ภนฺเต, น ปุน เอวรูปํ วิตกฺเกสฺสามี”ติ.}

\prose{อถ โข อายสฺมา โรหโณ อายสฺมนฺตํ นาคเสนํ เอตทโวจ “น โข ตฺยาหํ นาคเสน เอตฺตาวตา ขมามิ, อตฺถิ โข นาคเสน สาคลํ นาม นครํ, ตตฺถ มิลินฺโท นาม ราชา รชฺชํ กาเรติ, โส ทิฏฺฐิวาเทน ปญฺหํ ปุจฺฉิตฺวา ภิกฺขุสํฆํ วิเหเฐติ, สเจ ตฺวํ ตตฺถ คนฺตฺวา ตํ ราชานํ ทเมตฺวา พุทฺธสาสเน ปสาเทสฺสสิ, เอวาหํ ตํ ขมิสฺสามี”ติ.}

\prose{ติฏฺฐตุ ภนฺเต เอโก มิลินฺโท ราชา, สเจ ภนฺเต สกล\-ชมฺพุทีเป สพฺเพ ราชาโน อาคนฺตฺวา มํ ปญฺหํ ปุจฺเฉยฺยุํ, สพฺพํ ตํ วิสชฺเชตฺวา สมฺปทาเลสฺสามิ, “ขมถ เม ภนฺเต”ติ วตฺวา, “น ขมามี”ติ วุตฺเต “เตน หิ ภนฺเต อิมํ เตมาสํ กสฺส สนฺติเก วสิสฺสามี”ติ อาห. อยํ โข นาคเสน อายสฺมา อสฺสคุตฺโต วตฺตนิเย เสนาสเน วิหรติ, คจฺฉ ตฺวํ นาคเสน เยนายสฺมา อสฺสคุตฺโต เตนุปสงฺกม, อุปสงฺกมิตฺวา มม วจเนน อายสฺมโต อสฺสคุตฺตสฺส ปาเท สิรสา วนฺท, เอวญฺจ นํ วเทหิ “อุปชฺฌาโย เม ภนฺเต ตุมฺหากํ ปาเท สิรสา วนฺทติ, อปฺปาพาธํ อปฺปาตงฺกํ ลหุฏฺฐานํ พลํ ผาสุวิหารํ ปุจฺฉติ, อุปชฺฌาโย เม ภนฺเต อิมํ เตมาสํ ตุมฺหากํ สนฺติเก วสิตุํ มํ \csromanfolio{14}ปหิณี”ติ, “โกนาโม เต อุปชฺฌาโย”ติ จ วุตฺเต “โรหณตฺเถโร นาม ภนฺเต”ติ วเทยฺยาสิ, “อหํ โกนาโม”ติ วุตฺเต เอวํ วเทยฺยาสิ “มม อุปชฺฌาโย ภนฺเต ตุมฺหากํ นามํ ชานาตี”ติ. “เอวํ ภนฺเต”ติ โข อายสฺมา นาคเสโน อายสฺมนฺตํ โรหณํ อภิวาเทตฺวา ปทกฺขิณํ กตฺวา ปตฺต\-จีวรมา\-ทาย อนุปุพฺเพน จาริกํ จรมาโน เยน วตฺตนิยํ เสนาสนํ, เยนายสฺมา อสฺสคุตฺโต เตนุปสงฺกมิ, อุปสงฺกมิตฺวา อายสฺมนฺตํ อสฺสคุตฺตํ อภิวาเทตฺวา เอกมนฺตํ อฏฺฐาสิ, เอกมนฺตํ ฐิโต โข อายสฺมา นาคเสโน อายสฺมนฺตํ อสฺสคุตฺตํ เอตทโวจ “อุปชฺฌาโย เม ภนฺเต ตุมฺหากํ ปาเท สิรสา วนฺทติ, เอวญฺจ วเทติ อปฺปาพาธํ อปฺปาตงฺกํ ลหุฏฺฐานํ พลํ ผาสุวิหารํ ปุจฺฉติ, อุปชฺฌาโย เม ภนฺเต อิมํ เตมาสํ ตุมฺหากํ สนฺติเก วสิตุํ มํ ปหิณี”ติ.}

\prose{อถ โข อายสฺมา อสฺสคุตฺโต อายสฺมนฺตํ นาคเสนํ เอตทโวจ “ตฺวํ กินฺนาโมสี”ติ. อหํ ภนฺเต นาคเสโน นามาติ. โกนาโม เต อุปชฺฌาโยติ. อุปชฺฌาโย เม ภนฺเต โรหโณ นามาติ. อหํ โกนาโมติ. อุปชฺฌาโย เม ภนฺเต ตุมฺหากํ นามํ ชานาตีติ.}

\prose{สาธุ นาคเสน ปตฺตจีวรํ ปฏิสาเมหีติ. “สาธุ ภนฺเต”ติ ปตฺตจีวรํ ปฏิสาเมตฺวา ปุนทิวเส ปริเวณํ สมฺมชฺชิตฺวา มุโขทกํ ทนฺตโปณํ อุปฏฺฐเปสิ. เถโร สมฺมชฺชิต\-ฏฺฐานํ ปฏิสมฺมชฺชิ, ตํ อุทกํ ฉฑฺเฑตฺวา อญฺญํ อุทกํ อาหริ, ตญฺจ ทนฺตกฏฺฐํ อปเนตฺวา อญฺญํ ทนฺตกฏฺฐํ คณฺหิ, น อาลาปสลฺลาปํ อกาสิ, เอวํ สตฺต ทิวสานิ กตฺวา สตฺตเม ทิวเส ปุน ปุจฺฉิตฺวา ปุน เตน ตเถว วุตฺเต วสฺสวาสํ อนุชานิ.}

\proseitem{16}{เตน โข ปน สมเยน เอกา มหาอุปาสิกา อายสฺมนฺตํ อสฺสคุตฺตํ ติํสมตฺตานิ วสฺสานิ อุปฏฺฐาสิ. อถ โข สา มหาอุปาสิกา เตมาสจฺจเยน เยนายสฺมา อสฺสคุตฺโต เตนุปสงฺกมิ, อุปสงฺกมิตฺวา อายสฺมนฺตํ อสฺสคุตฺตํ เอตทโวจ “อตฺถิ นุ โข ตาต ตุมฺหากํ สนฺติเก อญฺโญ ภิกฺขู”ติ. อตฺถิ มหาอุปาสิเก อมฺหากํ สนฺติเก นาคเสโน นาม ภิกฺขูติ. เตน หิ ตาต อสฺสคุตฺต อธิวาเสหิ นาคเสเนน สทฺธิํ สฺวาตนาย ภตฺตนฺติ. อธิวาเสสิ โข อายสฺมา อสฺสคุตฺโต ตุณฺหีภาเวน.}

\prose{\csromanfolio{15}อถ โข อายสฺมา อสฺสคุตฺโต ตสฺสา รตฺติยา อจฺจเยน ปุพฺพณฺหสมยํ นิวาเสตฺวา ปตฺต\-จีวรมา\-ทาย อายสฺมตา นาคเสเนน สทฺธิํ ปจฺฉาสมเณน เยน มหาอุปาสิกาย นิเวสนํ เตนุปสงฺกมิ, อุปสงฺกมิตฺวา ปญฺญตฺเต อาสเน นิสีทิ. อถ โข สา มหาอุปาสิกา อายสฺมนฺตํ อสฺสคุตฺตํ อายสฺมนฺตญฺจ นาคเสนํ ปณีเตน ขาทนีเยน โภชนีเยน สหตฺถา สนฺตปฺเปสิ สมฺปวาเรสิ. อถ โข อายสฺมา อสฺสคุตฺโต ภุตฺตาวิํ โอนีต\-ปตฺต\-ปาณิํ อายสฺมนฺตํ นาคเสนํ เอตทโวจ “ตฺวํ นาคเสน มหาอุปาสิกาย อนุโมทนํ กโรหี”ติ อิทํ วตฺวา อุฏฺฐายาสนา ปกฺกามิ.}

\prose{อถ โข สา มหาอุปาสิกา อายสฺมนฺตํ นาคเสนํ เอตทโวจ “มหลฺลิกา โขหํ ตาต นาคเสน คมฺภีราย ธมฺมกถาย มยฺหํ อนุโมทนํ กโรหี”ติ. อถ โข อายสฺมา นาคเสโน ตสฺสา มหาอุปาสิกาย คมฺภีราย ธมฺมกถาย โลกุตฺตราย สุญฺญตปฺปฏิ สํยุตฺตาย อนุโมทนํ อกาสิ. อถ โข ตสฺสา มหาอุปาสิกาย ตสฺมิํ เยว อาสเน วิรชํ วีตมลํ ธมฺมจกฺขุํ อุทปาทิ “ยํ กิญฺจิ สมุทย\-ธมฺมํ, สพฺพํ ตํ นิโรธธมฺมนฺ”ติ. อายสฺมาปิ โข นาคเสโน ตสฺสา มหาอุปาสิกาย อนุโมทนํ กตฺวา อตฺตนา เทสิตํ ธมฺมํ ปจฺจเวกฺขนฺโต วิปสฺสนํ ปฏฺฐเปตฺวา ตสฺมิํ เยว อาสเนนิสินฺโน โสตาปตฺติผเล ปติฏฺฐาสิ.}

\prose{อถ โข อายสฺมา อสฺสคุตฺโต มณฺฑลมาเฬ นิสินฺโน ทฺวินฺนมฺปิ ธมฺมจกฺขุ\-ปฏิลาภํ ญตฺวา สาธุการํ ปวตฺเตสิ “สาธุ สาธุ นาคเสน, เอเกน กณฺฑปฺปหาเรน เทฺว มหากายา ปทาลิตา”ติ, อเนกานิ จ เทวตา\-สหสฺสานิ สาธุการํ ปวตฺเตสุํ.}

\proseitem{17}{อถ โข อายสฺมา นาคเสโน อุฏฺฐายาสนา เยนายสฺมา อสฺสคุตฺโต เตนุปสงฺกมิ, อุปสงฺกมิตฺวา อายสฺมนฺตํ อสฺสคุตฺตํ อภิวาเทตฺวา เอกมนฺตํ นิสีทิ, เอกมนฺตํ นิสินฺนํ โข อายสฺมนฺตํ นาคเสนํ อายสฺมา อสฺสคุตฺโต เอตทโวจ “คจฺฉ ตฺวํ นาคเสน ปาฏลิปุตฺตํ, ปาฏลิปุตฺต\-นคเร อโสการาเม อายสฺมา ธมฺมรกฺขิโต ปฏิวสติ, ตสฺส สนฺติเก พุทฺธวจนํ ปริยาปุณาหี”ติ. กีว ทูโร ภนฺเต อิโต ปาฏลิปุตฺตนาครนฺติ. โยชนสตานิ โข นาคเสนาติ. ทูโร โข ภนฺเต \csromanfolio{16}มคฺโค, อนฺตรามคฺเค ภิกฺขา ทุลฺลภา, กถาหํ คมิสฺสามีติ. คจฺฉ ตฺวํ นาคเสน อนฺตรามคฺเค ปิณฺฑปาตํ ลภิสฺสสิ สาลีนํ โอทนํ วิคตกาฬกํ อเนกสูปํ อเนก\-พฺยญฺชนนฺติ. “เอวํ ภนฺเต”ติ โข อายสฺมา นาคเสโน อายสฺมนฺตํ อสฺสคุตฺตํ อภิวาเทตฺวา ปทกฺขิณํ กตฺวา ปตฺต\-จีวรมา\-ทาย เยน ปาฏลิปุตฺตํ เตน จาริกํ ปกฺกามิ.}

\proseitem{18}{เตน โข ปน สมเยน ปาฏลิปุตฺตโก เสฏฺฐิ ปญฺจหิ สกฏสเตหิ ปาฏลิปุตฺต\-คามิ\-มคฺคํ ปฏิปนฺโน โหติ. อทฺทสา โข ปาฏลิปุตฺตโก เสฏฺฐิ อายสฺมนฺตํ นาคเสนํ ทูรโตว อาคจฺฉนฺตํ, ทิสฺวาน เยนายสฺมา นาคเสโน เตนุปสงฺกมิ, อุปสงฺกมิตฺวา อายสฺมนฺตํ นาคเสนํ อภิวาเทตฺวา “กุหิํ คจฺฉสิ ตาตา”ติ อาห. ปาฏลิปุตฺตํ คหปตีติ. สาธุ ตาต, มยมฺปิ ปาฏลิปุตฺตํ คจฺฉาม, อเมฺหหิ สทฺธิํ สุขํ คจฺฉถาติ.}

\prose{อถ โข ปาฏลิปุตฺตโก เสฏฺฐิ อายสฺมโต นาคเสนสฺส อิริยาปเถ ปสีทิตฺวา อายสฺมนฺตํ นาคเสนํ ปณีเตน ขาทนีเยน โภชนีเยน สหตฺถา สนฺตปฺเปตฺวา สมฺปวาเรตฺวา อายสฺมนฺตํ นาคเสนํ ภุตฺตาวิํ โอนีต\-ปตฺต\-ปาณิํ อญฺญตรํ นีจํ อาสนํ คเหตฺวา เอกมนฺตํ นิสีทิ, เอกมนฺตํ นิสินฺโน โข ปาฏลิปุตฺตโก เสฏฺฐิ อายสฺมนฺตํ นาคเสนํ เอตทโวจ “กินฺนาโมสิ ตฺวํ ตาตา”ติ. อหํ คหปติ นาคเสโน นามาติ. ชานาสิ โข ตฺวํ ตาต พุทฺธวจนํ นามาติ. ชานามิ โขหํ คหปติ อภิธมฺม\-ปทานีติ. ลาภา โน ตาต, สุลทฺธํ โน ตาต, อหมฺปิ โข ตาต อาภิธมฺมิโก, ตฺวมฺปิ อาภิธมฺมิโก, ภณ ตาต อภิธมฺม\-ปทานีติ. อถ โข อายสฺมา นาคเสโน ปาฏลิ\-ปุตฺต\-กสฺส เสฏฺฐิสฺส อภิธมฺมํ เทเสสิ, เทเสนฺเต เยว ปาฏลิ\-ปุตฺต\-กสฺส เสฏฺฐิสฺส วิรชํ วีตมลํ ธมฺมจกฺขุํ อุทปาทิ “ยํ กิญฺจิ สมุทย\-ธมฺมํ, สพฺพํ ตํ นิโรธธมฺมนฺ”ติ.}

\prose{อถ โข ปาฏลิปุตฺตโก เสฏฺฐิ ปญฺจมตฺตานิ สกฏสตานิ ปุรโต อุยฺโยเชตฺวา สยํ ปจฺฉโต คจฺฉนฺโต ปาฏลิปุตฺตสฺส อวิทูเร เทฺวธาปเถ ฐตฺวา อายสฺมนฺตํ นาคเสนํ เอตทโวจ “อยํ โข ตาต นาคเสน อโสการามสฺส มคฺโค, อิทํ โข ตาต อมฺหากํ กมฺพลรตนํ โสฬสหตฺถํ อายาเมน, อฏฺฐหตฺถํ วิตฺถาเรน, ปฏิคฺคณฺหาหิ โข ตาต \csromanfolio{17}อิทํ กมฺพลรตนํ อนุกมฺปํ อุปาทายา”ติ. ปฏิคฺคเหสิ โข อายสฺมา นาคเสโน ตํ กมฺพลรตนํ อนุกมฺปํ อุปาทาย. อถ โข ปาฏลิปุตฺตโก เสฏฺฐิ อตฺตมโน อุทคฺโค ปมุทิโต ปีติโสมนสฺส\-ชาโต อายสฺมนฺตํ นาคเสนํ อภิวาเทตฺวา ปทกฺขิณํ กตฺวา ปกฺกามิ.}

\proseitem{19}{อถ โข อายสฺมา นาคเสโน เยน อโสการาโม เยนายสฺมา ธมฺมรกฺขิโต เตนุปสงฺกมิ, อุปสงฺกมิตฺวา อายสฺมนฺตํ ธมฺม\-รกฺขิตํ อภิวาเทตฺวา อตฺตโน อาคตการณํ กเถตฺวา อายสฺมโต ธมฺม\-รกฺขิตสฺส สนฺติเก เตปิฏกํ พุทฺธวจนํ เอเกเนว อุทฺเทเสน ตีหิ มาเสหิ พฺยญฺชนโส ปริยาปุณิตฺวา ปุน ตีหิ มาเสหิ อตฺถโส มนสากาสิ.}

\prose{อถ โข อายสฺมา ธมฺมรกฺขิโต อายสฺมนฺตํ นาคเสนํ เอตทโวจ “เสยฺยถาปิ นาคเสน โคปาลโก คาโว รกฺขติ, อญฺเญ โครสํ ปริภุญฺชนฺติ. เอวเมว โข ตฺวํ นาคเสน เตปิฏกํ พุทฺธวจนํ ธาเรนฺโตปิ น ภาคี สามญฺญสฺสา”ติ. โหตุ ภนฺเต อลํ เอตฺตเกนาติ. เตเนว ทิวสภาเคน เตน รตฺติภาเคน สห ปฏิสมฺภิทาหิ อรหตฺตํ ปาปุณิ, สห สจฺจ\-ปฺปฏิเวเธน อายสฺมโต นาคเสนสฺส สพฺเพ เทวา สาธุการมทํสุ, ปถวี อุนฺนทิ, พฺรหฺมาโน อปฺโผเฏสุํ, ทิพฺพานิ จนฺทนจุณฺณานิ ทิพฺพานิ จ มนฺทารว\-ปุปฺผานิ อภิปฺปวสฺสิํสุ.}

\proseitem{20}{เตน โข ปน สมเยน โกฏิสตา อรหนฺโต หิมวนฺเต ปพฺพเต รกฺขิตตเล สนฺนิปติตฺวา อายสฺมโต นาคเสนสฺส สนฺติเก ทูตํ ปาเหสุํ “อาคจฺฉตุ นาคเสโน, ทสฺสนกามา มยํ นาคเสนนฺ”ติ. อถ โข อายสฺมา นาคเสโน ทูตสฺส วจนํ สุตฺวา อโสการาเม อนฺตรหิโต หิมวนฺเต ปพฺพเต รกฺขิตตเล โกฏิสตานํ อรหนฺตานํ ปุรโต ปาตุรโหสิ.}

\prose{อถ โข โกฏิสตา อรหนฺโต อายสฺมนฺตํ นาคเสนํ เอตทโวจุํ “เอโส โข นาคเสน มิลินฺโท ราชา ภิกฺขุสํฆํ วิเหเฐติ วาท\-ปฺปฏิวาเทน ปญฺหปุจฺฉาย, สาธุ นาคเสน คจฺฉ ตฺวํ มิลินฺทํ ราชานํ ทเมหี”ติ. ติฏฺฐตุ ภนฺเต เอโก มิลินฺโท ราชา, สเจ ภนฺเต สกล\-ชมฺพุทีเป ราชาโน อาคนฺตฺวา มํ ปญฺหํ ปุจฺเฉยฺยุํ, สพฺพํ ตํ วิสชฺเชตฺวา สมฺปทาเลสฺสามิ, \csromanfolio{18}คจฺฉถ โว ภนฺเต อจฺฉมฺภิตา สาคล\-นครนฺติ. อถ โข เถรา ภิกฺขู สาคลนครํ กาสาว\-ปฺปชฺโชตํ อิสิวาตปฏิวาตํ อกํสุ.}

\proseitem{21}{เตน โข ปน สมเยน อายสฺมา อายุปาโล สงฺเขฺยยฺย\-ปริเวเณ ปฏิวสติ. อถ โข มิลินฺโท ราชา อมจฺเจ เอตทโวจ “รมณียา วต โภ โทสินา รตฺติ, กนฺนุ ขฺวชฺช สมณํ วา พฺราหฺมณํ วา อุปสงฺกเมยฺยาม สากจฺฉาย ปญฺห\-ปุจฺฉนาย, โก มยา สทฺธิํ สลฺลปิตุํ อุสฺสหติ กงฺขํ ปฏิวิเนตุนฺ”ติ. เอวํ วุตฺเต ปญฺจสตา โยนกา ราชานํ มิลินฺทํ เอตทโวจุํ “อตฺถิ มหาราช อายุปาโล นาม เถโร เตปิฏโก พหุสฺสุโต อาคตาคโม, โส เอตรหิ สงฺเขฺยยฺย\-ปริเวเณ ปฏิวสติ, คจฺฉ ตฺวํ มหาราช อายสฺมนฺตํ อายุปาลํ ปญฺหํ ปุจฺฉสฺสู”ติ. เตน หิ ภเณ ภทนฺตสฺส อาโรเจถาติ.}

\prose{อถ โข เนมิตฺติโก อายสฺมโต อายุปาลสฺส สนฺติเก ทูตํ ปาเหสิ “ราชา ภนฺเต มิลินฺโท อายสฺมนฺตํ อายุปาลํ ทสฺสนกาโม”ติ. อายสฺมาปิ โข อายุปาโล เอวมาห “เตน หิ อาคจฺฉตู”ติ. อถ โข มิลินฺโท ราชา ปญฺจมตฺเตหิ โยนกสเตหิ ปริวุโต รถวรมา\-รุยฺห เยน สงฺเขฺยยฺย\-ปริเวณํ เยนายสฺมา อายุปาโล เตนุปสงฺกมิ, อุปสงฺกมิตฺวา อายสฺมตา อายุปาเลน สทฺธิํ สมฺโมทิ, สมฺโมทนียํ กถํ สารณียํ วีติสาเรตฺวา เอกมนฺตํ นิสีทิ, เอกมนฺตํ นิสินฺโน โข มิลินฺโท ราชา อายสฺมนฺตํ อายุปาลํ เอตทโวจ “กิมตฺถิยา ภนฺเต อายุปาล ตุมฺหากํ ปพฺพชฺชา, โก จ ตุมฺหากํ ปรมตฺโถ”ติ. เถโร อาห “ธมฺมจริย\-สมจริย\-ตฺถา โข มหาราช ปพฺพชฺชา, สามญฺญผลํ โข ปน อมฺหากํ ปรมตฺโถ”ติ. อตฺถิ ปน ภนฺเต โกจิ คิหีปิ ธมฺมจารี สมจารีติ. อาม มหาราช อตฺถิ คิหีปิ ธมฺมจารี สมจารี, ภควติ โข มหาราช พาราณสิยํ อิสิปตเน มิคทาเย ธมฺมจกฺกํ ปวตฺเตนฺเต อฏฺฐารสนฺนํ พฺรหฺมโกฏีนํ ธมฺมา\-ภิสมโย อโหสิ, เทวตานํ ปน ธมฺมา\-ภิสมโย คณนปถํ วีติวตฺโต, สพฺเพเต คิหิภูตา, น ปพฺพชิตา.}

\prose{ปุน จปรํ มหาราช ภควตา โข มหาสมย\-สุตฺตนฺเต เทสิยมาเน, มหามงฺคล\-สุตฺตนฺเต เทสิยมาเน, สมจิตฺต\-ปริยาย\-สุตฺตนฺเต เทสิยมาเน, ราหุโลวาท\-สุตฺตนฺเต เทสิยมาเน, ปราภว\-สุตฺตนฺเต เทสิยมาเน คณนปถํ วีติวตฺตานํ เทวตานํ ธมฺมา\-ภิสมโย อโหสิ, \csromanfolio{19}สพฺเพเต คิหิภูตา, น ปพฺพชิตาติ. เตน หิ ภนฺเต อายุปาล นิรตฺถิกา ตุมฺหากํ ปพฺพชฺชา, ปุพฺเพ กตสฺส ปาปกมฺมสฺส นิสฺสนฺเทน สมณา สกฺยปุตฺติยา ปพฺพชนฺติ ธุตงฺคานิ จ ปริหรนฺติ, เย โข เต ภนฺเต อายุปาล ภิกฺขู เอกาสนิกา, นูน เต ปุพฺเพ ปเรสํ โภคหารกา โจรา, เต ปเรสํ โภเค อจฺฉินฺทิตฺวา ตสฺส กมฺมสฺส นิสฺสนฺเทน เอตรหิ เอกาสนิกา ภวนฺติ, น ลภนฺติ กาเลน กาลํ ปริภุญฺชิตุํ, นตฺถิ เตสํ สีลํ, นตฺถิ ตโป, นตฺถิ พฺรหฺมจริยํ, เย โข ปน เต ภนฺเต อายุปาล ภิกฺขู อพฺโภกาสิกา, นูน เต ปุพฺเพ คามฆาตกา โจรา, เต ปเรสํ เคหานิ วินาเสตฺวา ตสฺส กมฺมสฺส นิสฺสนฺเทน เอตรหิ อพฺโภกาสิกา ภวนฺติ, น ลภนฺติ เสนาสนานิ ปริภุญฺชิตุํ, นตฺถิ เตสํ สีลํ, นตฺถิ ตโป, นตฺถิ พฺรหฺมจริยํ, เย โข ปน เต ภนฺเต อายุปาล ภิกฺขู เนสชฺชิกา, นูน เต ปุพฺเพ ปนฺถทูสกา โจรา, เต ปเรสํ ปถิเก ชเน คเหตฺวา พนฺธิตฺวา นิสีทาเปตฺวา ตสฺส กมฺมสฺส นิสฺสนฺเทน เอตรหิ เนสชฺชิกา ภวนฺติ, น ลภนฺติ เสยฺยํ กปฺเปตุํ, นตฺถิ เตสํ สีลํ, นตฺถิ ตโป, นตฺถิ พฺรหฺม\-จริยนฺติ อาห.}

\prose{เอวํ วุตฺเต อายสฺมา อายุปาโล ตุณฺหี อโหสิ, น กิญฺจิ ปฏิภาสิ. อถ โข ปญฺจสตา โยนกา ราชานํ มิลินฺทํ เอตทโวจุํ “ปณฺฑิโต มหาราช เถโร, อปิ จ โข อวิสารโท น กิญฺจิ ปฏิภาสตี”ติ.}

\prose{อถ โข มิลินฺโท ราชา อายสฺมนฺตํ อายุปาลํ ตุณฺหีภูตํ ทิสฺวา อปฺโผเฏตฺวา อุกฺกุฏฺฐิํ กตฺวา โยนเก เอตทโว จ “ตุจฺโฉ วต โภ ชมฺพุทีโป, ปลาโป วต โภ ชมฺพุทีโป. นตฺถิ โกจิ สมโณ วา พฺราหฺมโณ วา, โย มยา สทฺธิํ สลฺลปิตุํ อุสฺสหติ กงฺขํ ปฏิวิเนตุนฺ”ติ.}

\proseitem{22}{อถ โข มิลินฺทสฺส รญฺโญ สพฺพํ ตํ ปริสํ อนุวิโลเกนฺตสฺส อภีเต อมงฺกุภูเต โยนเก ทิสฺวา เอตทโหสิ “นิสฺสํสยํ อตฺถิ มญฺเญ อญฺโญ โกจิ ปณฺฑิโต ภิกฺขุ, โย มยา สทฺธิํ สลฺลปิตุํ อุสฺสหติ, เยนิเม โยนกา น มงฺกุภูตา”ติ. อถ โข มิลินฺโท ราชา โยนเก เอตทโวจ “อตฺถิ ภเณ อญฺโญ โกจิ ปณฺฑิโต ภิกฺขุ, โย มยา สทฺธิํ สลฺลปิตุํ อุสฺสหติ กงฺขํ ปฏิวิเนตุนฺ”ติ.}

\prose{เตน โข ปน สมเยน อายสฺมา นาคเสโน สมณ\-คณ\-ปริวุโต สํฆี คณี คณาจริโย ญาโต ยสสฺสี สาธุสมฺมโต \csromanfolio{20}พหุชนสฺส ปณฺฑิโต พฺยตฺโต เมธาวี นิปุโณ วิญฺญู วิภาวี วินีโต วิสารโท พหุสฺสุโต เตปิฏโก เวทคู ปภินฺน\-พุทฺธิมา อาคตาคโม ปภินฺนปฏิสมฺภิโท นวงฺค\-สตฺถุสาสเน ปริยตฺติธโร ปารมิปฺปตฺโต ชินวจเน ธมฺมตฺถเทสนา ปฏิเวธกุสโล อกฺขย\-วิจิตฺร\-ปฏิภาโน จิตฺรกถี กลฺยาณ\-วากฺกรโณ ทุราสโท ทุปฺปสโห ทุรุตฺตโร ทุราวรโณ ทุนฺนิวารโย, สาคโร วิย อกฺโขโภ, คิริราชา วิย นิจฺจโล, รณญฺชโห ตโมนุโท ปภงฺกโร มหากถี ปรคณิคณมถโน ปรติตฺถิย\-มทฺทโน ภิกฺขูนํ ภิกฺขุนีนํ อุปาสกานํ อุปาสิกานํ ราชูนํ ราช\-มหามตฺตานํ สกฺกโต ครุกโต มานิโต ปูชิโต อปจิโต ลาภี จีวร\-ปิณฺฑปาต\-เสนาสน\-คิลานปฺปจฺจย\-เภสชฺช\-ปริกฺขารานํ ลาภคฺค\-ยสคฺค\-ปฺปตฺโต วุทฺธานํ วิญฺญูนํ โสตาวธาเนน สมนฺนาคตานํ สนฺทสฺเสนฺโต นวงฺคํ ชินสาสน\-รตนํ, อุปทิสนฺโต ธมฺมมคฺคํ, ธาเรนฺโต ธมฺม\-ปฺปชฺโชตํ, อุสฺสาเปนฺโต ธมฺมยูปํ, ยชนฺโต ธมฺมยาคํ, ปคฺคณฺหนฺโต ธมฺมทฺธชํ, อุสฺสาเปนฺโต ธมฺมเกตุํ, ธเมนฺโต\csromansharedfootnote{da714b86a297363c}{อุปฺปฬาเสนฺโต (สี, อิ)} ธมฺมสงฺขํ, อาหนนฺโต ธมฺมเภริํ, นทนฺโต สีหนาทํ, คชฺชนฺโต อินฺทคชฺชิตํ, มธุร\-คิร\-คชฺชิเตน ญาณวรวิชฺชุชาลปริเวฐิเตน กรุณา\-ชล\-ภริเตน มหตา ธมฺมามต\-เมเฆน สกลโลกมภิตปฺปยนฺโต คามนิคม\-ราชธานีสุ จาริกํ จรมาโน อนุปุพฺเพน สาคลนครํ อนุปฺปตฺโต โหติ. ตตฺร สุทํ อายสฺมา นาคเสโน อสีติยา ภิกฺขุ\-สหสฺเสหิ สทฺธิํ สงฺเขฺยยฺย\-ปริเวเณ ปฏิวสติ. เตนาหุ โปราณา\csromandash{}}

\setlength{\csromangathapagewidth}{0pt}
\setlength{\csromangathawakwidth}{0pt}
\csromangathameasuregroup{“พหุสฺสุโต จิตฺรกถี, \\ สามยิโก จ กุสโล, \\ เต จ เตปิฏกา ภิกฺขู, \\ จตุเนกายิกา เจว, \\ คมฺภีรปญฺโญ เมธาวี, \\ อุตฺตมตฺถํ อนุปฺปตฺโต, \\ เตหิ ภิกฺขูหิ ปริวุโต, \\ จรนฺโต คามนิคมํ, \\ สงฺเขฺยยฺยปริเวณสฺมิํ, \\ กเถติ โส มนุสฺโสหิ,}{\csromangathabat{“พหุสฺสุโต จิตฺรกถี,}{นิปุโณ จ วิสารโท.} \\ \csromangathabat{สามยิโก จ กุสโล,}{ปฏิภาเน จ โกวิโท.} \\ \csromangathabat{เต จ เตปิฏกา ภิกฺขู,}{ปญฺจเนกายิกาปิ จ.} \\ \csromangathabat{จตุเนกายิกา เจว,}{นาคเสนํ ปุรกฺขรุํ.} \\ \csromangathabat{คมฺภีรปญฺโญ เมธาวี,}{มคฺคามคฺคสฺส โกวิโท.} \\ \csromangathabat{อุตฺตมตฺถํ อนุปฺปตฺโต,}{นาคเสโน วิสารโท.} \\ \csromangathabat{เตหิ ภิกฺขูหิ ปริวุโต,}{นิปุเณหิ สจฺจวาทิภิ.} \\ \csromangathabat{จรนฺโต คามนิคมํ,}{สาคลํ อุปสงฺกมิ.} \\ \csromangathabat{สงฺเขฺยยฺยปริเวณสฺมิํ,}{นาคเสโน ตทา วสิ.} \\ \csromangathabat{กเถติ โส มนุสฺโสหิ,}{ปพฺพเต เกสรี ยถา”ติ.}}
\csromangathasetleft{“พหุสฺสุโต จิตฺรกถี, \\ สามยิโก จ กุสโล, \\ เต จ เตปิฏกา ภิกฺขู, \\ จตุเนกายิกา เจว, \\ คมฺภีรปญฺโญ เมธาวี, \\ อุตฺตมตฺถํ อนุปฺปตฺโต, \\ เตหิ ภิกฺขูหิ ปริวุโต, \\ จรนฺโต คามนิคมํ, \\ สงฺเขฺยยฺยปริเวณสฺมิํ, \\ กเถติ โส มนุสฺโสหิ,}
\csromangathagroup{\csromangathastanza{\csromangathabat{“พหุสฺสุโต จิตฺรกถี,}{นิปุโณ จ วิสารโท.} \\ \csromangathabat{สามยิโก จ กุสโล,}{ปฏิภาเน จ โกวิโท.}}\csromangathastanza{\csromangathabat{เต จ เตปิฏกา ภิกฺขู,}{ปญฺจเนกายิกาปิ จ.} \\ \csromangathabat{จตุเนกายิกา เจว,}{นาคเสนํ ปุรกฺขรุํ.}}\csromangathastanza{\csromangathabat{คมฺภีรปญฺโญ เมธาวี,}{มคฺคามคฺคสฺส โกวิโท.} \\ \csromangathabat{อุตฺตมตฺถํ อนุปฺปตฺโต,}{นาคเสโน วิสารโท.}}\csromangathastanza{\csromangathabat{เตหิ ภิกฺขูหิ ปริวุโต,}{นิปุเณหิ สจฺจวาทิภิ.} \\ \csromangathabat{จรนฺโต คามนิคมํ,}{สาคลํ อุปสงฺกมิ.}}\csromangathastanza{\csromanfolio{21}\csromangathabat{สงฺเขฺยยฺยปริเวณสฺมิํ,}{นาคเสโน ตทา วสิ.} \\ \csromangathabat{กเถติ โส มนุสฺโสหิ,}{ปพฺพเต เกสรี ยถา”ติ.}}}

\proseitem{23}{อถ โข เทวมนฺติโย ราชานํ มิลินฺทํ เอตทโวจ “อาคเมหิ ตฺวํ มหาราช, อตฺถิ มหาราช นาคเสโน นาม เถโร ปณฺฑิโต พฺยตฺโต เมธาวี วินีโต วิสารโท พหุสฺสุโต จิตฺรกถี กลฺยาณ\-ปฏิภาโน อตฺถ\-ธมฺม\-นิรุตฺติ\-ปฏิภานปฏิสมฺภิทาสุ ปารมิปฺปตฺโต, โส เอตรหิ สงฺเขฺยยฺย\-ปริเวเณ ปฏิวสติ, คจฺฉ ตฺวํ มหาราช อายสฺมนฺตํ นาคเสนํ ปญฺหํ ปุจฺฉสฺสุ, อุสฺสหติ โส ตยา สทฺธิํ สลฺลปิตุํ กงฺขํ ปฏิวิเนตุนฺ”ติ. อถ โข มิลินฺทสฺส รญฺโญ สหสา “นาคเสโน”ติ สทฺทํ สุตฺวาว อหุเทว ภยํ, อหุเทว ฉมฺภิตตฺตํ, อหุเทว โลมหํโส. อถ โข มิลินฺโท ราชา เทวมนฺติยํ เอตทโวจ “อุสฺสหติ โภ นาคเสโน ภิกฺขุ มยา สทฺธิํ สลฺลปิตุนฺ”ติ. อุสฺสหติ มหาราช อปิ อินฺท\-ยมวรุณ\-กุเวร\-ปชาปติ สุยาม สนฺตุสิตโลกปาเลหิปิ ปิตุปิตามเหน มหาพฺรหฺมุนาปิ สทฺธิํ สลฺลปิตุํ, กิมงฺคํ ปน มนุสฺสภูเตนา”ติ.}

\prose{อถ โข มิลินฺโท ราชา เทวมนฺติยํ เอตทโวจ “เตน หิ ตฺวํ เทวมนฺติย ภทนฺตสฺส สนฺติเก ทูตํ เปเสหีติ. “เอวํ เทวา”ติ โข เทวมนฺติโย อายสฺมโต นาคเสนสฺส สนฺติเก ทูตํ ปาเหสิ “ราชา ภนฺเต มิลินฺโท อายสฺมนฺตํ ทสฺสนกาโม”ติ. อายสฺมาปิ โข นาคเสโน เอวมาห “เตน หิ อาคจฺฉตู”ติ.}

\prose{อถ โข มิลินฺโท ราชา ปญฺจมตฺเตหิ โยนกสเตหิ ปริวุโต รถวรมา\-รุยฺห มหตา พลกาเยน สทฺธิํ เยน สงฺเขฺยยฺย\-ปริเวณํ เยนายสฺมา นาคเสโน เตนุปสงฺกมิ. เตน โข ปน สมเยน อายสฺมา นาคเสโน อสีติยา ภิกฺขุ\-สหสฺเสหิ สทฺธิํ มณฺฑลมาเฬ นิสินฺโน โหติ. อทฺทสา โข มิลินฺโท ราชา อายสฺมโต นาคเสนสฺส ปริสํ ทูรโตว, ทิสฺวาน เทวมนฺติยํ เอตทโวจ “กสฺเสสา เทวมนฺติย มหตี ปริสา”ติ. อายสฺมโต โข มหาราช นาคเสนสฺส ปริสาติ.}

\prose{อถ โข มิลินฺทสฺส รญฺโญ อายสฺมโต นาคเสนสฺส ปริสํ ทูรโตว ทิสฺวา อหุเทว ภยํ, อหุเทว ฉมฺภิตตฺตํ, อหุเทว โลมหํโส. อถ โข มิลินฺโท ราชา ขคฺค\-ปริวาริโต วิย คโช, \csromanfolio{22}ครุฬ\-ปริวาริโต วิย นาโค, อชคร\-ปริวาริโต วิย โกตฺถุโก\csromansharedfootnote{e577bce7386b5b68}{โกฏฺฐุโก (ก)}, มหิํส\-ปริวุโต วิย อจฺโฉ, นาคานุพทฺโธ วิย มณฺฑูโก, สทฺทูลา\-นุพทฺโธ วิย มิโค, อหิตุณฺฑิก\-สมาคโต\csromansharedfootnote{7ba36ee17d1d1b87}{อหิคุณฺฐิกสมาคโต (สี, อิ)} วิย ปนฺนโค, มชฺชาร\-สมาคโต วิย อุนฺทูโร, ภูตเวชฺช\-สมาคโต วิย ปิสาโจ, ราหุมุขคโต วิย จนฺโท, ปนฺนโค วิย เปฬนฺตรคโต, สกุโณ วิย ปญฺชร\-นฺตร\-คโต, มจฺโฉ วิย ชาลนฺตรคโต, วาฬ\-วนม\-นุปฺปวิฏฺโฐ วิย ปุริโส, เวสฺสวณา\-ปราธิโก วิย ยกฺโข, ปริกฺขีณายุโก วิย เทวปุตฺโต ภีโต อุพฺพิคฺโค อุตฺรสฺโต สํวิคฺโค โลมหฏฺฐชาโต วิมโน ทุมฺมโน ภนฺตจิตฺโต วิปริณต\-มานโส “มา มํ อยํ ปริชโน ปริภวี”ติ สติํ\csromansharedfootnote{cb388bed69a8c1f8}{ธีติํ (สี, อิ)} อุปฏฺฐเปตฺวา เทวมนฺติยํ เอตทโวจ\csromandash{}“มา โข ตฺวํ เทวมนฺติย อายสฺมนฺตํ นาคเสนํ มยฺหํ อาจิกฺเขยฺยาสิ, อนกฺขาตญฺเญวาหํ นาคเสนํ ชานิสฺสามี”ติ. สาธุ มหาราช ตฺวญฺเญว ชานาหีติ.}

\prose{เตน โข ปน สมเยน อายสฺมา นาคเสโน ตสฺสา ภิกฺขุ\-ปริสาย ปุรโต จตฺตาลีสาย ภิกฺขุ\-สหสฺสานํ นวกตโร โหติ ปจฺฉโต จตฺตาลีสาย ภิกฺขุ\-สหสฺสานํ วุฑฺฒตโร.}

\prose{อถ โข มิลินฺโท ราชา สพฺพํ ตํ ภิกฺขุสํฆํ ปุรโต จ ปจฺฉโต จ มชฺฌโต จ อนุวิโลเกนฺโต อทฺทสา โข อายสฺมนฺตํ นาคเสนํ ทูรโตว ภิกฺขุ\-สํฆสฺส มชฺเฌ นิสินฺนํ เกสรสีหํ วิย วิคต\-ภยเภรวํ วิคต\-โลมหํสํ วิคตภย\-สารชฺชํ, ทิสฺวาน อากาเรเนว อญฺญาสิ “เอโส โข เอตฺถ นาคเสโน”ติ.}

\prose{อถ โข มิลินฺโท ราชา เทวมนฺติยํ เอตทโวจ “เอโส โข เทวมนฺติย อายสฺมา นาคเสโน”ติ. อาม มหราช เอโส โข นาคเสโน, สุฏฺฐุ โข ตฺวํ มหาราช นาคเสนํ อญฺญาสีติ. ตโต ราชา ตุฏฺโฐ อโหสิ “อนกฺขาโตว มยา นาคเสโน อญฺญาโต”ติ. อถ โข มิลินฺทสฺส รญฺโญ อายสฺมนฺตํ นาคเสนํ ทิสฺวาว อหุเทว ภยํ, อหุเทว ฉมฺภิตตฺตํ, อหุเทว โลมหํโส.}

\prose{\csromanfolio{23}เตนาหุ\csromandash{}}

\setlength{\csromangathapagewidth}{0pt}
\setlength{\csromangathawakwidth}{0pt}
\csromangathameasuregroup{“จรเณน จ สมฺปนฺนํ, \\ ทิสฺวา ราชา นาคเสนํ, \\ กถิตา\textsuperscript{1} มยา พหู ทิฏฺฐา, \\ น ตาทิสํ ภยํ อาสิ, \\ นิสฺสํสยํ ปราชโย, \\ ชโย จ นาคเสนสฺส,}{\csromangathabat{“จรเณน จ สมฺปนฺนํ,}{สุทนฺตํ อุตฺตเม ทเม.} \\ \csromangathabat{ทิสฺวา ราชา นาคเสนํ,}{อิทํ วจนมพฺรวิ.} \\ \csromangathabat{กถิตา\textsuperscript{1} มยา พหู ทิฏฺฐา,}{สากจฺฉา โอสฏา พหู.} \\ \csromangathabat{น ตาทิสํ ภยํ อาสิ,}{อชฺช ตาโส ยถา มม.} \\ \csromangathabat{นิสฺสํสยํ ปราชโย,}{มม อชฺช ภวิสฺสติ.} \\ \csromangathabat{ชโย จ นาคเสนสฺส,}{ยถา จิตฺตํ น สณฺฐิตนฺ”ติ.}}
\csromangathasetleft{“จรเณน จ สมฺปนฺนํ, \\ ทิสฺวา ราชา นาคเสนํ, \\ กถิตา\textsuperscript{1} มยา พหู ทิฏฺฐา, \\ น ตาทิสํ ภยํ อาสิ, \\ นิสฺสํสยํ ปราชโย, \\ ชโย จ นาคเสนสฺส,}
\csromangathagroupwithcloser{\csromangathastanza{\csromangathabat{“จรเณน จ สมฺปนฺนํ,}{สุทนฺตํ อุตฺตเม ทเม.} \\ \csromangathabat{ทิสฺวา ราชา นาคเสนํ,}{อิทํ วจนมพฺรวิ.}}\csromangathastanza{\csromangathabat{กถิตา\csromansharedfootnote{671ed94da9879954}{กถิกา (สี)} มยา พหู ทิฏฺฐา,}{สากจฺฉา โอสฏา พหู.} \\ \csromangathabat{น ตาทิสํ ภยํ อาสิ,}{อชฺช ตาโส ยถา มม.}}\csromangathastanza{\csromangathabat{นิสฺสํสยํ ปราชโย,}{มม อชฺช ภวิสฺสติ.} \\ \csromangathabat{ชโย จ นาคเสนสฺส,}{ยถา จิตฺตํ น สณฺฐิตนฺ”ติ.}}}{\textbf{พาหิรกถา นิฏฺฐิตา.}}
\csromanmatikaanchor{mtk.3}
\csromantocmarknopage{chapter}{2. มิลินฺทปญฺห}{mtk.3}
\bookmark[dest={mtk.3},level=0]{2. มิลินฺทปญฺห}
\chapterheadpageread{2. \textbf{มิลินฺทปญฺห}}
\csromanmatikaanchor{mtk.4}
\csromanmatikaanchor{mtk.5}
\csromantocmarknopage{section}{1. มหาวคฺค}{mtk.4}
\bookmark[dest={mtk.4},level=1]{1. มหาวคฺค}
\csromantocmarkref{subsection}{1. ปญฺญตฺติปญฺห}{mtk.5}
\bookmark[dest={mtk.5},level=2]{1. ปญฺญตฺติปญฺห}
\csromanmarkh{1. มหาวคฺค}\csromanmarkhii{1. ปญฺญตฺติปญฺห}\csromanheadernomark{2}{1. \textbf{มหาวคฺค 1. ปญฺญตฺติปญฺห}}
\proseitem{1}{\csromanfolio{24}อถ โข มิลินฺโท ราชา เยนายสฺมา นาคเสโน เตนุปสงฺกมิ, อุปสงฺกมิตฺวา อายสฺมตา นาคเสเนน สทฺธิํ สมฺโมทิ, สมฺโมทนียํ กถํ สารณียํ วีติสาเรตฺวา เอกมนฺตํ นิสีทิ. อายสฺมาปิ โข นาคเสโน ปฏิสมฺโมทนีเยเนว\csromansharedfootnote{3ff8212a50f22289}{ปฏิสมฺโมทิ, เตเนว (สี)} มิลินฺทสฺส รญฺโญ จิตฺตํ อาราเธสิ. อถ โข มิลินฺโท ราชา อายสฺมนฺตํ นาคเสนํ เอตทโวจ “กถํ ภทนฺโต ญายติ, กินฺนาโมสิ ภนฺเต”ติ. “นาคเสโน”ติ โข อหํ มหาราช ญายามิ, “นาคเสโน”ติ โข มํ มหาราช สพฺรหฺมจารี สมุทาจรนฺติ, อปิ จ มาตาปิตโร นามํ กโรนฺติ “นาคเสโน”ติ วา “สูรเสโน”ติ วา “วีรเสโน”ติ วา “สีหเสโน”ติ วา, อปิ จ โข มหาราช สงฺขา สมญฺญา ปญฺญตฺติ โวหาโร นามมตฺตํ ยทิทํ นาคเสโนติ, น เตตฺถ ปุคฺคโล อุปลพฺภตีติ.}

\prose{อถ โข มิลินฺโท ราชา เอวมาห “สุณนฺตุ เม โภนฺโต ปญฺจสตา โยนกา อสีติสหสฺสา จ ภิกฺขู, อยํ นาคเสโน เอวมาห ‘น เหตฺถ ปุคฺคโล อุปลพฺภตี’ติ, กลฺลํ นุโข\csromansharedfootnote{ca12a5f72c7f51e2}{นุ โข เทฺว นิปาตา \csromandash{} ม.พ.ป.} ตทภินนฺทิตุนฺ”ติ. อถ โข มิลินฺโท ราชา อายสฺมนฺตํ นาคเสนํ เอตทโวจ\csromandash{}สเจ ภนฺเต นาคเสน ปุคฺคโล นูปลพฺภติ, โก จรหิ ตุมฺหากํ จีวร\-ปิณฺฑปาต\-เสนาสน\-คิลานปฺปจฺจย\-เภสชฺช\-ปริกฺขารํ เทติ, โก ตํ ปริภุญฺชติ, โก สีลํ รกฺขติ, โก ภาวนม\-นุยุญฺชติ, โก มคฺค\-ผลนิพฺพานานิ สจฺฉิกโรติ, โก ปาณํ หนติ, โก อทินฺนํ อาทิยติ, โก กาเมสุ\-มิจฺฉาจารํ จรติ, โก มุสา ภณติ, โก มชฺชํ ปิวติ, โก ปญฺจานนฺตริย\-กมฺมํ กโรติ, ตสฺมา นตฺถิ กุสลํ, นตฺถิ อกุสลํ, นตฺถิ กุสลา\-กุสลานํ กมฺมานํ กตฺตา วา กาเรตา วา, นตฺถิ สุกต\-ทุกฺกฏานํ กมฺมานํ ผลํ วิปาโก, สเจ ภนฺเต นาคเสน โย ตุเมฺห มาเรติ, นตฺถิ ตสฺสาปิ ปาณาติปาโต, ตุมฺหากมฺปิ ภนฺเต นาคเสน นตฺถิ อาจริโย, นตฺถิ อุปชฺฌาโย, นตฺถิ อุปสมฺปทา. “นาคเสโนติ มํ มหาราช สพฺรหฺมจารี สมุทาจรนฺตี”ติ ยํ วเทสิ, กตโม เอตฺถ \csromanfolio{25}นาคเสโน. กินฺนุ โข ภนฺเต เกสา นาคเสโนติ. น หิ มหาราชาติ. โลมา นาคเสโนติ. น หิ มหาราชาติ. นขา ฯเปฯ. ทนฺตา. ตโจ. มํสํ. นฺหารุ. อฏฺฐิ. อฏฺฐิมิญฺชํ. วกฺกํ. หทยํ. ยกนํ. กิโลมกํ. ปิหกํ. ปปฺผาสํ. อนฺตํ. อนฺตคุณํ. อุทริยํ. กรีสํ. ปิตฺตํ. เสมฺหํ. ปุพฺโพ. โลหิตํ. เสโท. เมโท. อสฺสุ. วสา. เขโฬ. สิงฺฆาณิกา. ลสิกา. มุตฺตํ. ป. มตฺถเก มตฺถลุงฺคํ นาคเสโนติ. น หิ มหาราชาติ. กินฺนุ โข ภนฺเต รูปํ นาคเสโนติ. น หิ มหาราชติ. เวทนา นาคเสโนติ. น หิ มหาราชาติ. สญฺญา นาคเสโนติ. น หิ มหาราชาติ. สงฺขารา นาคเสโนติ. น หิ มหาราชาติ. วิญฺญาณํ นาคเสโนติ. น หิ มหาราชาติ. กิํ ปน ภนฺเต รูปเวทนา\-สญฺญา\-สงฺขาร\-วิญฺญาณํ นาคเสโนติ. น หิ มหาราชาติ. กิํ ปน ภนฺเต อญฺญตฺร รูปเวทนา\-สญฺญา\-สงฺขาร\-วิญฺญาณํ นาคเสโนติ. น หิ มหาราชาติ. ตมหํ ภนฺเต ปุจฺฉนฺโต ปุจฺฉนฺโต น ปสฺสามิ นาคเสนํ, นาคเสนสทฺโท เยว นุ โข ภนฺเต นาคเสโนติ. น หิ มหาราชาติ. โก ปเนตฺถ นาคเสโน, อลิกํ ตฺวํ ภนฺเต ภาสสิ มุสาวาทํ, นตฺถิ นาคเสโนติ.}

\prose{อถ โข อายสฺมา นาคเสโน มิลินฺทํ ราชานํ เอตทโวจ “ตฺวํ โขสิ มหาราช ขตฺติยสุขุมาโล อจฺจนฺต\-สุขุมาโล, ตสฺส เต มหาราช มชฺฌนฺหิกสมยํ ตตฺตาย ภูมิยา อุณฺหาย วาลิกาย ขราย สกฺขร\-กถลิกาย\csromansharedfootnote{e73dacf1ba3e79ee}{ขรา สกฺขรกฐลวาลิกา (สี, อิ)} มทฺทิตฺวา ปาเทนา\-คจฺฉนฺตสฺส ปาทา รุชฺชนฺติ, กาโย กิลมติ, จิตฺตํ อุปหญฺญติ, ทุกฺข\-สหคตํ กายวิญฺญาณํ อุปฺปชฺชติ, กินฺนุ โข ตฺวํ ปาเทนาคโตสิ, อุทาหุ วาหเนนาติ. นาหํ ภนฺเต ปาเทนาคจฺฉามิ, รเถนาหํ อาคโตสฺมีติ. สเจ ตฺวํ มหาราช รเถนาคโตสิ, รถํ เม อาโรเจหิ, กินฺนุ โข มหาราช อีสา รโถติ. น หิ ภนฺเตติ. อกฺโข รโถติ. น หิ ภนฺเตติ. จกฺกานิ รโถติ. น หิ ภนฺเตติ. รถปญฺชรํ รโถติ. น หิ ภนฺเตติ. รถทณฺฑโก รโถติ. น หิ ภนฺเตติ. ยุคํ รโถติ. น หิ ภนฺเตติ. รสฺมิโย รโถติ. น หิ ภนฺเตติ. ปโตทลฏฺฐิ รโถติ. น หิ ภนฺเตติ. กิํ นุ โข มหาราช อีสา- อกฺขจกฺกรถปญฺชรรถทณฺฑยุครสฺมิปโตทา รโถติ. น หิ ภนฺเตติ. กิํ ปน \csromanfolio{26}มหาราช อญฺญตฺร อีสา\-อกฺข\-จกฺก\-รถปญฺชร\-รถ\-ทณฺฑ\-ยุค\-รสฺมิ\-ปโตทา รโถติ. น หิ ภนฺเตติ. ตมหํ มหาราช ปุจฺฉนฺโต ปุจฺฉนฺโต น ปสฺสามิ รถํ, รถสทฺโท เยว นุ โข มหาราช รโถติ. น หิ ภนฺเตติ. โก ปเนตฺถ รโถ, อลิกํ ตฺวํ มหาราช ภาสสิ มุสาวาทํ, นตฺถิ รโถ, ตฺวํสิ มหาราช สกล\-ชมฺพุทีเป อคฺคราชา, กสฺส ปน ตฺวํ ภายิตฺวา มุสาวาทํ ภาสสิ, สุณนฺตุ เม โภนฺโต ปญฺจสตา โยนกา อสีติสหสฺสา จ ภิกฺขู, อยํ มิลินฺโท ราชา เอวมาห “รเถนาหํ อาคโตสฺมีติ, สเจ ตฺวํ มหาราช รเถนาคโตสิ, ‘รถํ เม อาโรเจหี’ติ วุตฺโต สมาโน รถํ น สมฺปาเทติ, กลฺลํ นุ โข ตทภินนฺทิตุนฺ”ติ. เอวํ วุตฺเต ปญฺจสตา โยนากา อายสฺมโต นาคเสนสฺส สาธุการํ ทตฺวา มิลินฺทํ ราชานํ เอตทโวจุํ “อิทานิ โข ตฺวํ มหาราช สกฺโกนฺโต ภาสสฺสู”ติ.}

\proseitem{2}{อถ โข มิลินฺโท ราชา อายสฺมนฺตํ นาคเสนํ เอตทโวจ “นาหํ ภนฺเต นาคเสน มุสา ภณามิ, อีสญฺจ ปฏิจฺจ อกฺขญฺจ ปฏิจฺจ จกฺกานิ จ ปฏิจฺจ รถปญฺชรญฺจ ปฏิจฺจ รถทณฺฑ\-กญฺจ ปฏิจฺจ ‘รโถ’ติ สงฺขา สมญฺญา ปญฺญตฺติ โวหาโร นามมตฺตํ ปวตฺตตี”ติ.}

\prose{สาธุ โข ตฺวํ มหาราช รถํ ชานาสิ, เอวเมว โข มหาราช มยฺหมฺปิ เกเส จ ปฏิจฺจ โลเม จ ปฏิจฺจ ฯเปฯ มตฺถเก มตฺถลุงฺคญฺจ ปฏิจฺจ รูปญฺจ ปฏิจฺจ เวทนญฺจ ปฏิจฺจ สญฺญญฺจ ปฏิจฺจ สงฺขาเร จ ปฏิจฺจ วิญฺญาณญฺจ ปฏิจฺจ “นาคเสโน”ติ สงฺขา สมญฺญา ปญฺญตฺติ โวหาโร นามมตฺตํ ปวตฺตติ, ปรมตฺถโต ปเนตฺถ ปุคฺคโล นูปลพฺภติ, ภาสิตมฺเปตํ มหาราช วชิราย ภิกฺขุนิยา ภควโต สมฺมุขา\csromandash{}}

\setlength{\csromangathapagewidth}{0pt}
\setlength{\csromangathawakwidth}{0pt}
\csromangathameasuregroup{“ยถา หิ องฺคสมฺภารา, \\ เอวํ ขนฺเธสุ สนฺเตสุ,}{\csromangathabat{“ยถา หิ องฺคสมฺภารา,}{โหติ สทฺโท รโถ อิติ.} \\ \csromangathabat{เอวํ ขนฺเธสุ สนฺเตสุ,}{โหติ ‘สตฺโต’ติ สมฺมุตี”ติ\textsuperscript{1}.}}
\csromangathasetleft{“ยถา หิ องฺคสมฺภารา, \\ เอวํ ขนฺเธสุ สนฺเตสุ,}
\csromangathagroup{\csromangathastanza{\csromangathabat{“ยถา หิ องฺคสมฺภารา,}{โหติ สทฺโท รโถ อิติ.} \\ \csromangathabat{เอวํ ขนฺเธสุ สนฺเตสุ,}{โหติ ‘สตฺโต’ติ สมฺมุตี”ติ\csromansharedfootnote{d3fc529a68302519}{ปสฺส สํ 1. 137; ขุ 7. 347 ปิฏฺเฐสุ.}.}}}

\prose{อจฺฉริยํ ภนฺเต นาคเสน, อพฺภุตํ ภนฺเต นาคเสน, อติจิตฺรานิ ปญฺห\-ปฏิภานานิ วิสชฺชิตานิ, ยทิ พุทฺโธ ติฏฺเฐยฺย สาธุการํ ทเทยฺย “สาธุ สาธุ นาคเสน, อติจิตฺรานิ ปญฺห\-ปฏิภานานิ วิสชฺชิตานี”ติ.}

\vspace{\tipitakaparskip}
\csromancenter{ปญฺญตฺติปโญฺห ปฐโม.}
\csromanmatikaanchor{mtk.6}
\csromantocmarkref{subsection}{2. วสฺสคณนปญฺห}{mtk.6}
\bookmark[dest={mtk.6},level=2]{2. วสฺสคณนปญฺห}
\csromanmarkh{2. มิลินฺทปญฺห}\csromanmarkhii{2. วสฺส\-คณน\-ปญฺห}\csromanheadernomark{2}{2. มิลินฺทปญฺห \textbf{2. วสฺส}\-\textbf{คณน}\-\textbf{ปญฺห}}
\proseitem{2}{\csromanfolio{27}กติวสฺโสสิ ตฺวํ ภนฺเต นาคเสนาติ. สตฺตวสฺโสหํ มหาราชาติ. เก เต ภนฺเต สตฺต, ตฺวํ วา สตฺต, คณนา วา สตฺตาติ.}

\prosewithcloserruled{เตน โข ปน สมเยน มิลินฺทสฺส รญฺโญ สพฺพา\-ภรณ\-ปฏิมณฺฑิตสฺส อลงฺกต\-ปฏิยตฺตสฺส ปถวิยํ ฉายา ทิสฺสติ, อุทกมณิเก จ ฉายา ทิสฺสติ. อถ โข อายสฺมา นาคเสโน มิลินฺทํ ราชานํ เอตทโวจ “อยํ เต มหาราช ฉายา ปถวิยํ อุทกมณิเก จ ทิสฺสติ, กิํ ปน มหาราช ตฺวํ วา ราชา, ฉายา วา ราชาติ. อหํ ภนฺเต นาคเสน ราชา, นายํ ฉายา ราชา, มํ ปน นิสฺสาย ฉายา ปวตฺตตีติ. เอวเมว โข มหาราช วสฺสานํ คณนา สตฺต, น ปนาหํ สตฺต, มํ ปน นิสฺสาย สตฺต ปวตฺตติ, ฉายูปมํ มหาราชาติ. อจฺฉริยํ ภนฺเต นาคเสน, อพฺภุตํ ภนฺเต นาคเสน, อติจิตฺรานิ ปญฺห\-ปฏิภานานิ วิสชฺชิตานีติ.}{วสฺส\-คณน\-ปโญฺห ทุติโย.}
\csromanmatikaanchor{mtk.7}
\csromantocmarkref{subsection}{3. วีมํสนปญฺห}{mtk.7}
\bookmark[dest={mtk.7},level=2]{3. วีมํสนปญฺห}
\csromanheader{2}{3. \textbf{วีมํสนปญฺห}}
\proseitem{3}{ราชา อาห “ภนฺเต นาคเสน สลฺลปิสฺสสิ มยา สทฺธินฺ”ติ. สเจ ตฺวํ มหาราช ปณฺฑิตวาทํ\csromansharedfootnote{06305475a6ec91fa}{ปณฺฑิตวาทา (สี, อิ)} สลฺลปิสฺสสิ สลฺลปิสฺสามิ, สเจ ปน ราชวาทํ สลฺลปิสฺสสิ น สลฺลปิสฺสามีติ. กถํ ภนฺเต นาคเสน ปณฺฑิตา สลฺลปนฺตีติ. ปณฺฑิตานํ โข มหาราช สลฺลาเป อาเวฐนมฺปิ กยิรติ, นิพฺเพฐนมฺปิ กยิรติ, นิคฺคโหปิ กยิรติ, ปฏิกมฺมมฺปิ กยิรติ, วิสฺสาโสปิ\csromansharedfootnote{62d0ba9d9bbe27d4}{วิเสโสปิ (สี, อิ)} กยิรติ, ปฏิวิสฺสาโสปิ กยิรติ, น จ เตน ปณฺฑิตา กุปฺปนฺติ, เอวํ โข มหาราช ปณฺฑิตา สลฺลปนฺตีติ. กถํ ปน ภนฺเต ราชาโน สลฺลปนฺตีติ. ราชาโน โข มหาราช สลฺลาเป เอกํ วตฺถุํ ปฏิชานนฺติ, โย ตํ วตฺถุํ วิโลเมติ, ตสฺส ทณฺฑํ อาณาเปนฺติ “อิมสฺส ทณฺฑํ ปเณถา”ติ, เอวํ โข มหาราช ราชาโน สลฺลปนฺตีติ. ปณฺฑิตวาทาหํ ภนฺเต สลฺลปิสฺสามิ, โน ราชวาทํ, วิสฺสฏฺโฐ ภทนฺโต สลฺลปตุ ยถา ภิกฺขุนา วา สามเณเรน วา อุปาสเกน วา \csromanfolio{28}อารามิเกน วา สทฺธิํ สลฺลปติ, เอวํ วิสฺสฏฺโฐ ภทนฺโต สลฺลปตุ มา ภายตูติ. “สุฏฺฐุ มหาราชา”ติ เถโร อพฺภานุโมทิ.}

\prosewithcloserruled{ราชา อาห “ภนฺเต นาคเสน ปุจฺฉิสฺสามี”ติ. ปุจฺฉ มหาราชาติ. ปุจฺฉิโตสิ เม ภนฺเตติ. วิสชฺชิตํ มหาราชาติ. กิํ ปน ภนฺเต ตยา วิสชฺชิตนฺติ. กิํ ปน มหาราช ตยา ปุจฺฉิตนฺติ.}{วีมํสนปโญฺห ตติโย.}
\csromanmatikaanchor{mtk.8}
\csromantocmarkref{subsection}{4. อนนฺตกายปญฺห}{mtk.8}
\bookmark[dest={mtk.8},level=2]{4. อนนฺตกายปญฺห}
\csromanheader{2}{4. \textbf{อนนฺตกาย}\-\textbf{ปญฺห}}
\proseitem{3}{อถ โข มิลินฺทสฺส รญฺโญ เอตทโหสิ “ปณฺฑิโต โข อยํ ภิกฺขุ ปฏิพโล มยา สทฺธิํ สลฺลปิตุํ, พหุกานิ จ เม ฐานานิ ปุจฺฉิ ตพฺพานิ ภวิสฺสนฺติ, ยาว อปุจฺฉิตานิ เยว ตานิ ฐานานิ ภวิสฺสนฺติ, อถ สูริโย อตฺถํ คมิสฺสติ, ยนฺนูนาหํ เสฺว อนฺเตปุเร สลฺลเปยฺยนฺ”ติ. อถ โข ราชา เทวมนฺติยํ เอตทโวจ\csromandash{}เตน หิ ตฺวํ เทวมนฺติย ภทนฺตสฺส อาโรเจยฺยาสิ “เสฺว อนฺเตปุเร รญฺญา สทฺธิํ สลฺลาโป ภวิสฺสตี”ติ. อิทํ วตฺวา มิลินฺโท ราชา อุฏฺฐายาสนา เถรํ นาคเสนํ อาปุจฺฉิตฺวา รถํ อภิรูหิตฺวา “นาคเสโน นาคเสโน”ติ สชฺฌายํ กโรนฺโต ปกฺกามิ.}

\prose{อถ โข เทวมนฺติโย อายสฺมนฺตํ นาคเสนํ เอตทโวจ “ราชา ภนฺเต มิลินฺโท เอวมาห\csromandash{}เสฺว อนฺเตปุเร รญฺญา สทฺธิํ สลฺลาโป ภวิสฺสตี”ติ. “สุฏฺฐู”ติ เถโร อพฺภานุโมทิ. อถ โข ตสฺสา รตฺติยา อจฺจเยน เทวมนฺติโย จ อนนฺตกาโย จ มงฺกุโร จ สพฺพทินฺโน จ เยน มิลินฺโท ราชา เตนุ\-ปสงฺกมิํสุ, อุปสงฺกมิตฺวา ราชานํ มิลินฺทํ เอตทโวจุํ “อาคจฺฉตุ มหาราช ภทนฺโต นาคเสโน”ติ. อาม อาคจฺฉตูติ. กิตฺตเกหิ ภิกฺขูหิ สทฺธิํ อาคจฺฉตูติ. ยตฺตเก ภิกฺขู อิจฺฉติ, ตตฺตเกหิ ภิกฺขูหิ สทฺธิํ อาคจฺฉตูติ.}

\prose{อถ โข สพฺพทินฺโน อาห “อาคจฺฉตุ มหาราช ทสหิ ภิกฺขูหิ สทฺธินฺ”ติ, ทุติยมฺปิ โข ราชา อาห “ยตฺตเก ภิกฺขู อิจฺฉติ, ตตฺตเกหิ ภิกฺขูหิ สทฺธิํ อาคจฺฉตู”ติ. ทุติยมฺปิ โข สพฺพทินฺโน อาห “อาคจฺฉตุ มหาราช ทสหิ ภิกฺขูหิ สทฺธินฺ”ติ. ตติยมฺปิ โข ราชา อาห \csromanfolio{29}“ยตฺตเก ภิกฺขู อิจฺฉติ, ตตฺตเกหิ ภิกฺขูหิ สทฺธิํ อาคจฺฉตู”ติ. ตติยมฺปิ โข สพฺพทินฺโน อาห “อาคจฺฉตุ มหาราช ทสหิ ภิกฺขูหิ สทฺธินฺ”ติ. สพฺโพ ปนายํ สกฺกาโร ปฏิยาทิโต, อหํ ภณามิ “ยตฺตเก ภิกฺขู อิจฺฉติ, ตตฺตเกหิ ภิกฺขูหิ สทฺธิํ อาคจฺฉตู”ติ. อยํ ภเณ สพฺพทินฺโน อญฺญถา ภณติ, กินฺนุ มยํ นปฺปฏิพลา ภิกฺขูนํ โภชนํ ทาตุนฺติ. เอวํ วุตฺเต สพฺพทินฺโน มงฺกุ อโหสิ.}

\proseitem{4}{อถ โข เทวมนฺติโย จ อนนฺตกาโย จ มงฺกุโร จ เยนายสฺมา นาคเสโน เตนุ\-ปสงฺกมิํสุ, อุปสงฺกมิตฺวา อายสฺมนฺตํ นาคเสนํ เอตทโวจุํ ราชา ภนฺเต มิลินฺโท เอวมาห “ยตฺตเก ภิกฺขู อิจฺฉติ, ตตฺตเกหิ ภิกฺขูหิ สทฺธิํ อาคจฺฉตู”ติ. อถ โข อายสฺมา นาคเสโน ปุพฺพณฺหสมยํ นิวาเสตฺวา ปตฺต\-จีวรมา\-ทาย อสีติยา ภิกฺขุ\-สหสฺเสหิ สทฺธิํ สาคลํ ปาวิสิ.}

\prose{อถ โข อนนฺตกาโย อายสฺมนฺตํ นาคเสนํ นิสฺสาย คจฺฉนฺโต อายสฺมนฺตํ นาคเสนํ เอตทโวจ “ภนฺเต นาคเสน ยํ ปเนตํ พฺรูสิ “นาคเสโน”ติ, กตโม เอตฺถ นาคเสโนติ. เถโร อาห “โก ปเนตฺถ “นาคเสโน”ติ มญฺญสี”ติ. โย โส ภนฺเต อพฺภนฺตเร วาโต ชีโว ปวิสติ จ นิกฺขมติ จ, โส “นาคเสโน”ติ มญฺญามีติ. ยทิ ปเนโส วาโต นิกฺขมิตฺวา นปฺปวิเสยฺย, ปวิสิตฺวา น นิกฺขเมยฺย, ชีเวยฺย นุ โข โส ปุริโสติ. น หิ ภนฺเตติ. เย ปนิเม สงฺขธมกา สงฺขํ ธเมนฺติ, เตสํ วาโต ปุน ปวิสตีติ. น หิ ภนฺเตติ. เย ปนิเม วํสธมกา วํสํ ธเมนฺติ, เตสํ วาโต ปุน ปวิสตีติ. น หิ ภนฺเตติ. เย ปนิเม สิงฺคธมกา สิงฺคํ ธเมนฺติ, เตสํ วาโต ปุน ปวิสตีติ. น หิ ภนฺเตติ. อถ กิสฺส ปน เต น มรนฺตีติ. นาหํ ปฏิพโล ตยา วาทินา สทฺธิํ สลฺลปิตุํ, สาธุ ภนฺเต อตฺถํ ชปฺเปหีติ. เนโส ชีโว, อสฺสาสปสฺสาสา นาเมเต กายสงฺขาราติ เถโร อภิธมฺมกถํ กเถสิ. อถ อนนฺตกาโย อุปาสกตฺตํ ปฏิเวเทสีติ.}

\vspace{\tipitakaparskip}
\csromancenter{อนนฺตกาย\-ปโญฺห จตุตฺโถ.}
\csromanmatikaanchor{mtk.9}
\csromantocmarkref{subsection}{5. ปพฺพชฺชปญฺห}{mtk.9}
\bookmark[dest={mtk.9},level=2]{5. ปพฺพชฺชปญฺห}
\csromanheader{2}{5. \textbf{ปพฺพชฺชปญฺห}}
\proseitem{5}{\csromanfolio{30}อถ โข อายสฺมา นาคเสโน เยน มิลินฺทสฺส รญฺโญ นิเวสนํ เตนุปสงฺกมิ, อุปสงฺกมิตฺวา ปญฺญตฺเต อาสเน นิสีทิ. อถ โข มิลินฺโท ราชา อายสฺมนฺตํ นาคเสนํ สปริสํ ปณีเตน ขาทนีเยน โภชนีเยน สหตฺถา สนฺตปฺเปตฺวา สมฺปวาเรตฺวา เอกเมกํ ภิกฺขุํ เอกเมเกน ทุสฺสยุเคน อจฺฉาเทตฺวา อายสฺมนฺตํ นาคเสนํ ติจีวเรน อจฺฉาเทตฺวา อายสฺมนฺตํ นาคเสนํ เอตทโวจ “ภนฺเต นาคเสน ทสหิ ภิกฺขูหิ สทฺธิํ อิธ นิสีทถ, อวเสสา คจฺฉนฺตู”ติ.}

\prose{อถ โข มิลินฺโท ราชา อายสฺมนฺตํ นาคเสนํ ภุตฺตาวิํ โอนีต\-ปตฺต\-ปาณิํ วิทิตฺวา อญฺญตรํ นีจํ อาสนํ คเหตฺวา เอกมนฺตํ นิสีทิ, เอกมนฺตํ นิสินฺโน โข มิลินฺโท ราชา อายสฺมนฺตํ นาคเสนํ เอตทโวจ “ภนฺเต นาคเสน กิมฺหิ โหติ กถาสลฺลาโป”ติ. อตฺเถน มยํ มหาราช อตฺถิกา, อตฺเถ โหตุ กถาสลฺลาโปติ.}

\prose{ราชา อาห “กิมตฺถิยา ภนฺเต นาคเสน ตุมฺหากํ ปพฺพชฺชา, โก จ ตุมฺหากํ ปรมตฺโถ”ติ. เถโร อาห “กินฺติ มหาราช อิทํ ทุกฺขํ นิรุชฺเฌยฺย, อญฺญญฺจ ทุกฺขํ น อุปฺปชฺเชยฺยาติ. เอตทตฺถา มหาราช อมฺหากํ ปพฺพชฺชา, อนุปาทา ปรินิพฺพานํ โข ปน อมฺหากํ ปรมตฺโถ”ติ.}

\prose{กิํ ปน ภนฺเต นาคเสน สพฺเพ เอตทตฺถาย ปพฺพชนฺตีติ. น หิ มหาราช เกจิ เอตทตฺถาย ปพฺพชนฺติ, เกจิ ราชาภินีตา\csromansharedfootnote{0df0bdf018b7811e}{ราชภีติตา (สี)} ปพฺพชนฺติ, เกจิ โจราภินีตา\csromansharedfootnote{a2380bdb6b6c9c90}{โจรภีติตา (สี)} ปพฺพชนฺติ, เกจิ อิณฏฺฏา ปพฺพชนฺติ, เกจิ อาชีวิกตฺถาย ปพฺพชนฺติ, เย ปน สมฺมา ปพฺพชนฺติ, เต เอตทตฺถาย ปพฺพชนฺตีติ.}

\prose{ตฺวํ ปน ภนฺเต เอตทตฺถาย ปพฺพชิโตสีติ. อหํ โข มหาราช ทหรโก สนฺโต ปพฺพชิโต, น ชานามิ อิมสฺส นามตฺถาย ปพฺพชามีติ, อปิ จ โข เม เอวํ อโหสิ “ปณฺฑิตา อิเม สมณา สกฺยปุตฺติยา, เต มํ สิกฺขาเปสฺสนฺตี”ติ, สฺวาหํ เตหิ สิกฺขาปิโต ชานามิ จ ปสฺสามิ จ “อิมสฺส นามตฺถาย ปพฺพชฺชา”ติ.}

\prose{กลฺโลสิ ภนฺเต นาคเสนาติ.}

\vspace{\tipitakaparskip}
\csromancenter{ปพฺพชฺชปโญฺห ปญฺจโม.}
\csromanmatikaanchor{mtk.10}
\csromantocmarkref{subsection}{6. ปฏิสนฺธิปญฺห}{mtk.10}
\bookmark[dest={mtk.10},level=2]{6. ปฏิสนฺธิปญฺห}
\csromanmarkh{2. มิลินฺทปญฺห}\csromanmarkhii{6. ปฏิสนฺธิ\-ปญฺห}\csromanheadernomark{2}{2. มิลินฺทปญฺห \textbf{6. ปฏิสนฺธิ}\-\textbf{ปญฺห}}
\proseitem{6}{\csromanfolio{31}ราชา อาห “ภนฺเต นาคเสน อตฺถิ โกจิ มโต น ปฏิสนฺทหตี”ติ. เถโร อาห “โกจิ ปฏิสนฺทหติ, โกจิ น ปฏิสนฺทหตี”ติ. โก ปฏิสนฺทหติ, โก น ปฏิสนฺทหตีติ. สกิเลโส มหาราช ปฏิสนฺทหติ, นิกฺกิเลโส น ปฏิสนฺทหตีติ. ตฺวํ ปน ภนฺเต นาคเสน ปฏิสนฺทหิสฺสสีติ. สเจ มหาราช สอุปาทาโน ภวิสฺสามิ ปฏิสนฺทหิสฺสามิ, สเจ อนุปาทาโน ภวิสฺสามิ น ปฏิสนฺทหิสฺสามีติ.}

\prosewithcloserruled{กลฺโลสิ ภนฺเต นาคเสนาติ.}{ปฏิสนฺธิ\-ปโญฺห ฉฏฺโฐ.}
\csromanmatikaanchor{mtk.11}
\csromantocmarkref{subsection}{7. โยนิโสมนสิการปญฺห}{mtk.11}
\bookmark[dest={mtk.11},level=2]{7. โยนิโสมนสิการปญฺห}
\csromanheader{2}{7. \textbf{โยนิโสมนสิการ}\-\textbf{ปญฺห}}
\proseitem{7}{ราชา อาห “ภนฺเต นาคเสน โย น ปฏิสนฺทหติ, นนุ โส โยนิโส มนสิกาเรน น ปฏิสนฺทหตี”ติ. โยนิโส จ มหาราช มนสิกาเรน ปญฺญาย จ อญฺเญหิ จ กุสเลหิ ธมฺเมหีติ. นนุ ภนฺเต โยนิโส มนสิกาโร เยว ปญฺญาติ. น หิ มหาราช อญฺโญ มนสิกาโร, อญฺญา ปญฺญา, อิเมสํ โข มหาราช อเชฬก โคณมหิํ สโอฏฺฐคทฺรภานมฺปิ มนสิกาโร อตฺถิ, ปญฺญา ปน เตสํ นตฺถีติ.}

\prosewithcloserruled{กลฺโลสิ ภนฺเต นาคเสนาติ.}{โยนิโส มนสิการ\-ปโญฺห สตฺตโม.}
\csromanmatikaanchor{mtk.12}
\csromantocmarkref{subsection}{8. มนสิการลกฺขณปญฺห}{mtk.12}
\bookmark[dest={mtk.12},level=2]{8. มนสิการลกฺขณปญฺห}
\csromanheader{2}{8. \textbf{มนสิการ}\-\textbf{ลกฺขณ}\-\textbf{ปญฺห}}
\proseitem{8}{ราชา อาห “กิํลกฺขโณ ภนฺเต นาคเสน มนสิกาโร, กิํลกฺขณา ปญฺญา”ติ. อูหนลกฺขโณ โข มหาราช มนสิกาโร, เฉทน\-ลกฺขณา ปญฺญาติ.}

\prose{กถํ อูหนลกฺขโณ มนสิกาโร, กถํ เฉทน\-ลกฺขณา ปญฺญา, โอปมฺมํ กโรหีติ. ชานาสิ ตฺวํ มหาราช ยวลาวเกติ. อาม ภนฺเต \csromanfolio{32}ชานามีติ. กถํ มหาราช ยวลาวกา ยวํ ลุนนฺตีติ. วาเมน ภนฺเต หตฺเถน ยวกลาปํ คเหตฺวา ทกฺขิเณน หตฺเถน ทาตฺตํ คเหตฺวา ทาตฺเตน ฉินฺทนฺตีติ.}

\prose{ยถา มหาราช ยวลาวโก วาเมน หตฺเถน ยวกลาปํ คเหตฺวา ทกฺขิเณน หตฺเถน ทาตฺตํ คเหตฺวา ยวํ ฉินฺทติ, เอวเมว โข มหาราช โยคาวจโร มนสิกาเรน มานสํ คเหตฺวา ปญฺญาย กิเลเส ฉินฺทติ, เอวํ โข มหาราช อูหนลกฺขโณ มนสิกาโร, เอวํ เฉทน\-ลกฺขณา ปญฺญาติ.}

\prosewithcloserruled{กลฺโลสิ ภนฺเต นาคเสนาติ.}{มนสิการ\-ลกฺขณ\-ปโญฺห อฏฺฐโม.}
\csromanmatikaanchor{mtk.13}
\csromantocmarkref{subsection}{9. สีลลกฺขณปญฺห}{mtk.13}
\bookmark[dest={mtk.13},level=2]{9. สีลลกฺขณปญฺห}
\csromanheader{2}{9. \textbf{สีลลกฺขณ}\-\textbf{ปญฺห}}
\proseitem{9}{ราชา อาห “ภนฺเต นาคเสน ยํ ปเนตํ พฺรูสิ ‘อญฺเญหิ จ กุสเลหิ ธมฺเมหี’ติ, กตเม เต กุสลา ธมฺมา”ติ. สีลํ มหาราช สทฺธา วีริยํ สติ สมาธิ, อิเม เต กุสลา ธมฺมาติ. กิํลกฺขณํ ภนฺเต สีลนฺติ. ปติฏฺฐาน\-ลกฺขณํ มหาราช สีลํ สพฺเพสํ กุสลานํ ธมฺมานํ, อินฺทฺริยพลโพชฺฌงฺคมคฺคงฺคสติปฏฺฐานสมฺมปฺปธาน- อิทฺธิปาทฌานวิโมกฺขสมาธิสมาสมาปตฺตีนํ สีลํ ปติฏฺฐํ, สีเล ปติฏฺฐิโต โข มหาราช โยคาวจโร สีลํ นิสฺสาย สีเล ปติฏฺฐาย ปญฺจินฺทฺริยานิ ภาเวติ สทฺธินฺทฺริยํ วีริยินฺทฺริยํ สตินฺทฺริยํ สมาธินฺทฺริยํ ปญฺญินฺทฺริยนฺติ, สพฺเพ กุสลา ธมฺมา น ปริหายนฺตีติ. โอปมฺมํ กโรหีติ. ยถา มหาราช เย เกจิ พีชคาม\-ภูตคามา วุฑฺฒิํ วิรูฬฺหิํ เวปุลฺลํ อาปชฺชนฺติ, สพฺเพ เต ปถวิํ นิสฺสาย ปถวิยํ ปติฏฺฐาย วุฑฺฒิํ วิรูฬฺหิํ เวปุลฺลํ อาปชฺชนฺติ. เอวเมว โข มหาราช โยคาวจโร สีลํ นิสฺสาย สีเล ปติฏฺฐาย ปญฺจินฺทฺริยานิ ภาเวติ สทฺธินฺทฺริยํ วีริยินฺทฺริยํ สตินฺทฺริยํ สมาธินฺทฺริยํ ปญฺญินฺทฺริยนฺติ.}

\prose{ภิยฺโย โอปมฺมํ กโรหีติ. ยถา มหาราช เย เกจิ พลกรณียา กมฺมนฺตา กยิรนฺติ, สพฺเพ เต ปถวิํ นิสฺสาย ปถวิยํ ปติฏฺฐาย กยิรนฺติ. เอวเมว โข มหาราช โยคาวจโร สีลํ นิสฺสาย สีเล ปติฏฺฐาย \csromanfolio{33}ปญฺจินฺทฺริยานิ ภาเวติ สทฺธินฺทฺริยํ วีริยินฺทฺริยํ สตินฺทฺริยํ สมาธินฺทฺริยํ ปญฺญินฺทฺริยนฺติ.}

\prose{ภิยฺโย โอปมฺมํ กโรหีติ. ยถา มหาราช นครวฑฺฒกี นครํ มาเปตุกาโม ปฐมํ นครฏฺฐานํ โสธาเปตฺวา ขาณุกณฺฏกํ อปกฑฺฒาเปตฺวา ภูมิํ สมํ การาเปตฺวา ตโต อปรภาเค วีถิจตุกฺก\-สิงฺฆาฏกา\-ทิ\-ปริจฺเฉเทน วิภชิตฺวา นครํ มาเปติ. เอวเมว โข มหาราช โยคาวจโร สีลํ นิสฺสาย สีเล ปติฏฺฐาย ปญฺจินฺทฺริยานิ ภาเวติ สทฺธินฺทฺริยํ วีริยินฺทฺริยํ สตินฺทฺริยํ สมาธินฺทฺริยํ ปญฺญินฺทฺริยนฺติ.}

\prose{ภิยฺโย โอปมฺมํ กโรหีติ. ยถา มหาราช ลงฺฆโก สิปฺปํ ทสฺเสตุกาโม ปถวิํ ขณาเปตฺวา สกฺขร\-กถลํ อปกฑฺฒาเปตฺวา ภูมิํ สมํ การาเปตฺวา มุทุกาย ภูมิยา สิปฺปํ ทสฺเสติ. เอวเมว โข มหาราช โยคาวจโร สีลํ นิสฺสาย สีเล ปติฏฺฐาย ปญฺจินฺทฺริยานิ ภาเวติ สทฺธินฺทฺริยํ วีริยินฺทฺริยํ สตินฺทฺริยํ สมาธินฺทฺริยํ ปญฺญินฺทฺริยนฺติ. ภาสิตมฺเปตํ มหาราช ภควตา\csromandash{}}

\setlength{\csromangathapagewidth}{0pt}
\setlength{\csromangathawakwidth}{0pt}
\csromangathameasuregroup{“สีเล ปติฏฺฐาย นโร สปญฺโญ, \\ อาตาปี นิปโก ภิกฺขุ,}{\csromangathabat{“สีเล ปติฏฺฐาย นโร สปญฺโญ,}{จิตฺตํ ปญฺญญฺจ ภาวยํ.} \\ \csromangathabat{อาตาปี นิปโก ภิกฺขุ,}{โส อิมํ วิชฏเย ชฏนฺ”ติ\textsuperscript{1}.} \\ “อยํ ปติฏฺฐา ธรณีว ปาณินํ, \\ อิทญฺจ มูลํ กุสลาภิวุฑฺฒิยา. \\ มุขญฺจิทํ สพฺพชินานุสาสเน, \\ โย สีลกฺขนฺโธ วรปาติโมกฺขิโย”ติ.}
\csromangathameasurewak{“อยํ ปติฏฺฐา ธรณีว ปาณินํ, \\ อิทญฺจ มูลํ กุสลาภิวุฑฺฒิยา. \\ มุขญฺจิทํ สพฺพชินานุสาสเน, \\ โย สีลกฺขนฺโธ วรปาติโมกฺขิโย”ติ.}
\csromangathasetleft{“สีเล ปติฏฺฐาย นโร สปญฺโญ, \\ อาตาปี นิปโก ภิกฺขุ,}
\csromangathagroup{\csromangathastanza{\csromangathabat{“สีเล ปติฏฺฐาย นโร สปญฺโญ,}{จิตฺตํ ปญฺญญฺจ ภาวยํ.} \\ \csromangathabat{อาตาปี นิปโก ภิกฺขุ,}{โส อิมํ วิชฏเย ชฏนฺ”ติ\csromansharedfootnote{0c352e03ce84fecb}{ปสฺส สํ 1. 13 ปิฏฺเฐ.}.}}\csromangathastanza{\csromangathawak{“อยํ ปติฏฺฐา ธรณีว ปาณินํ, \\ อิทญฺจ มูลํ กุสลาภิวุฑฺฒิยา. \\ มุขญฺจิทํ สพฺพชินานุสาสเน, \\ โย สีลกฺขนฺโธ วรปาติโมกฺขิโย”ติ.}}}

\prosewithcloserruled{กลฺโลสิ ภนฺเต นาคเสนาติ.}{สีลลกฺขณ\-ปโญฺห นวโม.}
\csromanmatikaanchor{mtk.14}
\csromantocmarkref{subsection}{10. สมฺปสาทนลกฺขณสทฺธาปญฺห}{mtk.14}
\bookmark[dest={mtk.14},level=2]{10. สมฺปสาทนลกฺขณสทฺธาปญฺห}
\csromanheader{2}{10. \textbf{สมฺปสาทน}\-\textbf{ลกฺขณ}\-\textbf{สทฺธา}\-\textbf{ปญฺห}}
\proseitem{10}{ราชา อาห “ภนฺเต นาคเสน กิํลกฺขณา สทฺธา”ติ. สมฺปสาทน\-ลกฺขณา จ มหาราช สทฺธา สมฺปกฺขนฺทนลกฺขณา จาติ. กถํ ภนฺเต สมฺปสาทน\-ลกฺขณา สทฺธาติ. สทฺธา โข มหาราช อุปฺปชฺชมานา \csromanfolio{34}นีวรเณ วิกฺขมฺเภติ, วินีวรณํ จิตฺตํ โหติ อจฺฉํ วิปฺปสนฺนํ อนาวิลํ. เอวํ โข มหาราช สมฺปสาทน\-ลกฺขณา สทฺธาติ.}

\prose{โอปมฺมํ กโรหีติ. ยถา มหาราช ราชา จกฺกวตฺตี จตุรงฺคินิยา เสนาย สทฺธิํ อทฺธาน\-มคฺค\-ปฺปฏิปนฺโน ปริตฺตํ อุทกํ ตเรยฺย, ตํ อุทกํ หตฺถีหิ จ อสฺเสหิ จ รเถหิ จ ปตฺตีหิ จ ขุภิตํ ภเวยฺย อาวิลํ ลุฬิตํ กลลีภูตํ. อุตฺติณฺโณ จ ราชา จกฺกวตฺตี มนุสฺเส อาณาเปยฺย “ปานียํ ภเณ อาหรถ, ปิวิสฺสามี”ติ, รญฺโญ จ อุทกปฺปสาทโก มณิ ภเวยฺย. “เอวํ เทวา”ติ โข เต มนุสฺสา รญฺโญ จกฺกวตฺติสฺส ปฏิสฺสุตฺวา ตํ อุทก\-ปฺปสาทกํ มณิํ อุทเก ปกฺขิเปยฺยุํ, ตสฺมิํ อุทเก ปกฺขิตฺตมตฺเต สงฺข\-เสวาล\-ปณกํ วิคจฺเฉยฺย, กทฺทโม จ สนฺนิสีเทยฺย, อจฺฉํ ภเวยฺย อุทกํ วิปฺปสนฺนํ อนาวิลํ. ตโต รญฺโญ จกฺกวตฺติสฺส ปานียํ อุปนาเมยฺยุํ “ปิวตุ เทว ปานียนฺ”ติ.}

\prose{ยถา มหาราช อุทกํ, เอวํ จิตฺตํ ทฏฺฐพฺพํ, ยถา เต มนุสฺสา, เอวํ โยคาวจโร ทฏฺฐพฺโพ, ยถา สงฺข\-เสวาล\-ปณกํ กทฺทโม จ, เอวํ กิเลสา ทฏฺฐพฺพา. ยถา อุทกปฺปสาทโก มณิ, เอวํ สทฺธา ทฏฺฐพฺพา, ยถา อุทกปฺปสาทเก มณิมฺหิ อุทเก ปกฺขิตฺตมตฺเต สงฺข\-เสวาล\-ปณกํ วิคจฺเฉยฺย, กทฺทโม จ สนฺนิสีเทยฺย, อจฺฉํ ภเวยฺย อุทกํ วิปฺปสนฺนํ อนาวิลํ, เอวเมว โข มหาราช สทฺธา อุปฺปชฺชมานา นีวรเณ วิกฺขมฺเภติ, วินีวรณํ จิตฺตํ โหติ อจฺฉํ วิปฺปสนฺนํ อนาวิลํ, เอวํ โข มหาราช สมฺปสาทน\-ลกฺขณา สทฺธาติ.}

\prosewithcloserruled{กลฺโลสิ ภนฺเต นาคเสนาติ.}{สมฺปสาทน\-ลกฺขณ\-สทฺธา\-ปโญฺห ทสโม.}
\csromanmatikaanchor{mtk.15}
\csromantocmarkref{subsection}{11. สมฺปกฺขนฺทนลกฺขณสทฺธาปญฺห}{mtk.15}
\bookmark[dest={mtk.15},level=2]{11. สมฺปกฺขนฺทนลกฺขณสทฺธาปญฺห}
\csromanheader{2}{11. \textbf{สมฺปกฺขนฺทนลกฺขณสทฺธาปญฺห}}
\proseitem{11}{กถํ ภนฺเต สมฺปกฺขนฺทนลกฺขณา สทฺธาติ. ยถา มหาราช โยคาวจโร อญฺเญสํ จิตฺตํ วิมุตฺตํ ปสฺสิตฺวา โสตาปตฺติผเล วา สกทาคามิ\-ผเล วา อนาคามิผเล วา อรหตฺเต วา สมฺปกฺขนฺทติ โยคํ กโรติ อปฺปตฺตสฺส ปตฺติยา อนธิคตสฺส อธิคมาย อสจฺฉิกตสฺส สจฺฉิกิริยาย. เอวํ โข มหาราช สมฺปกฺขนฺทนลกฺขณา สทฺธาติ.}

\prose{\csromanfolio{35}โอปมฺมํ กโรหีติ. ยถา มหาราช อุปริปพฺพเต มหาเมโฆ อภิปฺปวสฺเสยฺย, ตํ อุทกํ ยถานินฺนํ ปวตฺตมานํ ปพฺพต\-กนฺทร\-ปทร\-สาขา ปริปูเรตฺวา นทิํ ปริปูเรยฺย, สา อุภโต กูลานิ สํวิสฺสนฺทนฺตี คจฺเฉยฺย, อถ มหาชนกาโย อาคนฺตฺวา ตสฺสา นทิยา อุตฺตานตํ วา คมฺภีรตํ วา อชานนฺโต ภีโต วิตฺถโต ตีเร ติฏฺเฐยฺย, อถญฺญตโร ปุริโส อาคนฺตฺวา อตฺตโน ถามญฺจ พลญฺจ สมฺปสฺสนฺโต คาฬฺหํ กจฺฉํ พนฺธิตฺวา ปกฺขนฺทิตฺวา ตเรยฺย, ตํ ติณฺณํ ปสฺสิตฺวา มหาชนกาโยปิ ตเรยฺย. เอวเมว โข มหาราช โยคาวจโร อญฺเญสํ จิตฺตํ วิมุตฺตํ ปสฺสิตฺวา โสตาปตฺติผเล วา สกทาคามิ\-ผเล วา อนาคามิผเล วา อรหตฺเต วา สมฺปกฺขนฺทติ โยคํ กโรติ อปฺปตฺตสฺส ปตฺติยา อนธิคตสฺส อธิคมาย อสจฺฉิกตสฺส สจฺฉิกิริยาย. เอวํ โข มหาราช สมฺปกฺขนฺทนลกฺขณา สทฺธาติ. ภาสิตมฺเปตํ มหาราช ภควตา สํยุตฺตนิกาย\-วเร\csromandash{}}

\setlength{\csromangathapagewidth}{0pt}
\setlength{\csromangathawakwidth}{0pt}
\csromangathameasuregroup{“สทฺธาย ตรตี โอฆํ, \\ วีริเยน ทุกฺขมจฺเจติ,}{\csromangathabat{“สทฺธาย ตรตี โอฆํ,}{อปฺปมาเทน อณฺณวํ.} \\ \csromangathabat{วีริเยน ทุกฺขมจฺเจติ,}{ปญฺญาย ปริสุชฺฌตี”ติ\textsuperscript{1}.}}
\csromangathasetleft{“สทฺธาย ตรตี โอฆํ, \\ วีริเยน ทุกฺขมจฺเจติ,}
\csromangathagroup{\csromangathastanza{\csromangathabat{“สทฺธาย ตรตี โอฆํ,}{อปฺปมาเทน อณฺณวํ.} \\ \csromangathabat{วีริเยน ทุกฺขมจฺเจติ,}{ปญฺญาย ปริสุชฺฌตี”ติ\csromansharedfootnote{2fba8541fd73da29}{ปสฺส สํ 1. 217 ปิฏฺเฐ.}.}}}

\prosewithcloserruled{กลฺโลสิ ภนฺเต นาคเสนาติ.}{สมฺปกฺขนฺทนลกฺขณสทฺธาปโญฺห เอกาทสโม.}
\csromanmatikaanchor{mtk.16}
\csromantocmarkref{subsection}{12. วีริยลกฺขณปญฺห}{mtk.16}
\bookmark[dest={mtk.16},level=2]{12. วีริยลกฺขณปญฺห}
\csromanheader{2}{12. \textbf{วีริย}\-\textbf{ลกฺขณ}\-\textbf{ปญฺห}}
\proseitem{12}{ราชา อาห “ภนฺเต นาคเสน กิํลกฺขณํ วีริยนฺ”ติ. อุปตฺถมฺภน\-ลกฺขณํ มหาราช วีริยํ, วีริยู\-ปตฺถมฺภิตา สพฺเพ กุสลา ธมฺมา น ปริหายนฺตีติ.}

\prose{โอปมฺมํ กโรหีติ. ยถา มหาราช ปุริโส เคเห ปตนฺเต อญฺเญน ทารุนา อุปตฺถมฺเภยฺย, อุปตฺถมฺภิตํ สนฺตํ เอวํ ตํ เคหํ น ปเตยฺย. เอวเมว โข มหาราช อุปตฺถมฺภน\-ลกฺขณํ วีริยํ, วีริยู\-ปตฺถมฺภิตา สพฺเพ กุสลา ธมฺมา น ปริหายนฺตีติ. ภิยฺโย โอปมฺมํ กโรหีติ. ยถา มหาราช ปริตฺตกํ เสนํ มหตี เสนา ภญฺเชยฺย, ตโต ราชา อญฺญมญฺญํ อนุสฺสาเรยฺย อนุเปเสยฺย \csromanfolio{36}อตฺตโน ปริตฺตกาย เสนาย พลํ อนุปทํ ทเทยฺย, ตาย สทฺธิํ ปริตฺตกา เสนา มหติํ เสนํ ภญฺเชยฺย. เอวเมว โข มหาราช อุปตฺถมฺภน\-ลกฺขณํ วีริยํ, วีริยู\-ปตฺถมฺภิตา สพฺเพ กุสลา ธมฺมา น ปริหายนฺติ. ภาสิตมฺเปตํ มหาราช ภควตา\csromandash{}“วีริยวา โข ภิกฺขเว อริยสาวโก อกุสลํ ปชหติ, กุสลํ ภาเวติ. สาวชฺชํ ปชหติ, อนวชฺชํ ภาเวติ. สุทฺธมตฺตานํ ปริหรตี”ติ *.}

\prosewithcloserruled{กลฺโลสิ ภนฺเต นาคเสนาติ.}{วีริย\-ลกฺขณ\-ปโญฺห ทฺวาทสโม.}
\csromanmatikaanchor{mtk.17}
\csromantocmarkref{subsection}{13. สติลกฺขณปญฺห}{mtk.17}
\bookmark[dest={mtk.17},level=2]{13. สติลกฺขณปญฺห}
\csromanheader{2}{13. \textbf{สติลกฺขณ}\-\textbf{ปญฺห}}
\proseitem{13}{ราชา อาห “ภนฺเต นาคเสน กิํลกฺขณา สตี”ติ. อปิลาปน\-ลกฺขณา มหาราช สติ, อุปคฺคณฺหน\-ลกฺขณา จาติ. กถํ ภนฺเต อปิลาปน\-ลกฺขณา สตีติ. สติ มหาราช อุปฺปชฺชมานา กุสลากุสล\-สาวชฺชานวชฺช\-หีนปฺปณีต\-กณฺหสุกฺกสปฺปฏิภาค\-ธมฺเม อปิลาเปติ “อิเม จตฺตาโร สติปฏฺฐานา, อิเม จตฺตาโร สมฺมปฺปธานา, อิเม จตฺตาโร อิทฺธิปาทา, อิมานิ ปญฺจินฺทฺริยานิ, อิมานิ ปญฺจ พลานิ, อิเม สตฺต โพชฺฌงฺคา, อยํ อริโย อฏฺฐงฺคิโก มคฺโค, อยํ สมโถ, อยํ วิปสฺสนา, อยํ วิชฺชา, อยํ วิมุตฺตี”ติ. ตโต โยคาวจโร เสวิตพฺเพ ธมฺเม เสวติ, อเสวิตพฺเพ ธมฺเม น เสวติ. ภชิตพฺเพ ธมฺเม ภชติ อภชิตพฺเพ ธมฺเม น ภชติ. เอวํ โข มหาราช อปิลาปน\-ลกฺขณา สตีติ.}

\prose{โอปมฺมํ กโรหีติ. ยถา มหาราช รญฺโญ จกฺกวตฺติสฺส ภณฺฑาคาริโก ราชานํ จกฺกวตฺติํ สายํ ปาตํ ยสํ สราเปติ “เอตฺตกา เทว เต หตฺถี, เอตฺตกา อสฺสา, เอตฺตกา รถา, เอตฺตกา ปตฺตี, เอตฺตกํ หิรญฺญํ, เอตฺตกํ สุวณฺณํ, เอตฺตกํ สาปเตยฺยํ, ตํ เทโว สรตู”ติ รญฺโญ สาปเตยฺยํ อปิลาเปติ. เอวเมว โข มหาราช สติ อุปฺปชฺชมานา กุสลากุสล\-สาวชฺชานวชฺช\-หีนปฺปณีต\-กณฺหสุกฺกสปฺปฏิภาค\-ธมฺเม อปิลาเปติ “อิเม จตฺตาโร สติปฏฺฐานา, อิเม จตฺตาโร สมฺมปฺปธานา, อิเม จตฺตาโร อิทฺธิปาทา, อิมานิ ปญฺจินฺทฺริยานิ, อิมานิ ปญฺจ พลานิ, อิเม สตฺต โพชฺฌงฺคา, อยํ \csromanfolio{37}อริโย อฏฺฐงฺคิโก มคฺโค, อยํ สมโถ, อยํ วิปสฺสนา, อยํ วิชฺชา, อยํ วิมุตฺตี”ติ. ตโต โยคาวจโร เสวิตพฺเพ ธมฺเม เสวติ, อเสวิตพฺเพ ธมฺเม น เสวติ. ภชิตพฺเพ ธมฺเม ภชติ, อภชิตพฺเพ ธมฺเม น ภชติ. เอวํ โข มหาราช อปิลาปน\-ลกฺขณา สตีติ.}

\prose{กถํ ภนฺเต อุปคฺคณฺหน\-ลกฺขณา สตีติ. สติ มหาราช อุปฺปชฺชมานา หิตาหิตานํ ธมฺมานํ คติโย สมเนฺวติ “อิเม ธมฺมา หิตา, อิเม ธมฺมา อหิตา. อิเม ธมฺมา อุปการา, อิเม ธมฺมา อนุปการา”ติ. ตโต โยคาวจโร อหิเต ธมฺเม อปนุเทติ, หิเต ธมฺเม อุปคฺคณฺหาติ. อนุปกาเร ธมฺเม อปนุเทติ, อุปกาเร ธมฺเม อุปคฺคณฺหาติ. เอวํ โข มหาราช อุปคฺคณฺหน\-ลกฺขณา สตีติ.}

\prose{โอปมฺมํ กโรหีติ. ยถา มหาราช รญฺโญ จกฺกวตฺติสฺส ปริณายก\-รตนํ รญฺโญ หิตาหิเต ชานาติ “อิเม รญฺโญ หิตา, อิเม อหิตา. อิเม อุปการา, อิเม อนุปการา”ติ. ตโต อหิเต อปนุเทติ, หิเต อุปคฺคณฺหาติ. อนุปกาเร อปนุเทติ, อุปกาเร อุปคฺคณฺหาติ. เอวเมว โข มหาราช สติ อุปฺปชฺชมานา หิตาหิตานํ ธมฺมานํ คติโย สมเนฺวติ “อิเม ธมฺมา หิตา, อิเม ธมฺมา อหิตา. อิเม ธมฺมา อุปการา, อิเม ธมฺมา อนุปการา”ติ. ตโต โยคาวจโร อหิเต ธมฺเม อปนุเทติ, หิเต ธมฺเม อุปคฺคณฺหาติ. อนุปกาเร ธมฺเม อปนุเทติ, อุปกาเร ธมฺเม อุปคฺคณฺหาติ. เอวํ โข มหาราช อุปคฺคณฺหน\-ลกฺขณา สติ. ภาสิตมฺเปตํ มหาราช ภควตา\csromandash{}“สติญฺจ ขฺวาหํ ภิกฺขเว สพฺพตฺถิกํ วทามี”ติ *.}

\prosewithcloserruled{กลฺโลสิ ภนฺเต นาคเสนาติ.}{สติลกฺขณ\-ปโญฺห เตรสโม.}
\csromanmatikaanchor{mtk.18}
\csromantocmarkref{subsection}{14. สมาธิปญฺห}{mtk.18}
\bookmark[dest={mtk.18},level=2]{14. สมาธิปญฺห}
\csromanheader{2}{14. \textbf{สมาธิปญฺห}}
\proseitem{14}{ราชา อาห “ภนฺเต นาคเสน กิํลกฺขโณ สมาธี”ติ. ปมุข\-ลกฺขโณ มหาราช สมาธิ, เย เกจิ กุสลา ธมฺมา, สพฺเพ เต สมาธิปมุขา โหนฺติ สมาธินินฺนา สมาธิโปณา สมาธิ\-ปพฺภาราติ.}

\prose{\csromanfolio{38}โอปมฺมํ กโรหีติ. ยถา มหาราช กูฏาคารสฺส ยา กาจิ โคปานสิโย, สพฺพา ตา กูฏงฺคมา โหนฺติ กูฏนินฺนา กูฏสโมสรณา, กูฏํ ตาสํ อคฺคมกฺขายติ. เอวเมว โข มหาราช เย เกจิ กุสลา ธมฺมา, สพฺเพ เต สมาธิปมุขา โหนฺติ สมาธินินฺนา สมาธิโปณา สมาธิ\-ปพฺภาราติ.}

\prose{ภิยฺโย โอปมฺมํ กโรหีติ. ยถา มหาราช โกจิ ราชา จตุรงฺคินิยา เสนาย สทฺธิํ สงฺคามํ โอตเรยฺย, สพฺพาว เสนา หตฺถี จ อสฺสา จ รถา จ ปตฺตี จ ตปฺปมุขา\csromansharedfootnote{cf751a2b8a290bce}{ตมฺปมุขา (สฺยา, ก)} ภเวยฺยุํ ตนฺนินฺนา ตปฺโปณา ตปฺปพฺภารา ตํ เยว\csromansharedfootnote{cf7dc7e6d0b8f44d}{ตํเยว นิคฺคหีตสนฺธิ \csromandash{} ม.พ.ป.} อนุปริยาเยยฺยุํ. เอวเมว โข มหาราช เย เกจิ กุสลา ธมฺมา, สพฺเพ เต สมาธิปมุขา โหนฺติ สมาธินินฺนา สมาธิโปณา สมาธิ\-ปพฺภารา. เอวํ โข มหาราช ปมุข\-ลกฺขโณ สมาธิ. ภาสิตมฺเปตํ มหาราช ภควตา\csromandash{}“สมาธิํ ภิกฺขเว ภาเวถ, สมาหิโต ภิกฺขเว ภิกฺขุ ยถาภูตํ ปชานาตี”ติ *.}

\prosewithcloserruled{กลฺโลสิ ภนฺเต นาคเสนาติ.}{สมาธิปโญฺห จุทฺทสโม.}
\csromanmatikaanchor{mtk.19}
\csromantocmarkref{subsection}{15. ปญฺญาลกฺขณปญฺห}{mtk.19}
\bookmark[dest={mtk.19},level=2]{15. ปญฺญาลกฺขณปญฺห}
\csromanheader{2}{15. \textbf{ปญฺญา}\-\textbf{ลกฺขณ}\-\textbf{ปญฺห}}
\proseitem{15}{ราชา อาห “ภนฺเต นาคเสน กิํ ลกฺขณา ปญฺญา”ติ. ปุพฺเพว โข มหาราช มยา วุตฺตํ “เฉทน\-ลกฺขณา ปญฺญา”ติ, อปิ จ โอภาสน\-ลกฺขณา ปญฺญาติ. กถํ ภนฺเต โอภาสน\-ลกฺขณา ปญฺญาติ ปญฺญา มหาราช อุปฺปชฺชมานา อวิชฺช\-นฺธการํ วิธเมติ, วิชฺโชภาสํ ชเนติ, ญาณาโลกํ วิทํเสติ, อริยสจฺจานิ ปากฏานิ กโรติ. ตโต โยคาวจโร “อนิจฺจนฺ”ติ วา “ทุกฺขนฺ”ติ วา “อนตฺตา”ติ วา สมฺมปฺปญฺญาย ปสฺสตีติ.}

\prose{โอปมฺมํ กโรหีติ. ยถา มหาราช ปุริโส อนฺธกาเร เคเห ปทีปํ ปเวเสยฺย, ปวิฏฺโฐ ปทีโป อนฺธการํ วิธเมติ, โอภาสํ ชเนติ, อาโลกํ วิทํเสติ, รูปานิ ปากฏานิ กโรติ. เอวเมว โข มหาราช \csromanfolio{39}ปญฺญา อุปฺปชฺชมานา อวิชฺช\-นฺธการํ วิธเมติ, วิชฺโชภาสํ ชเนติ, ญาณาโลกํ วิทํเสติ, อริยสจฺจานิ ปากฏานิ กโรติ. ตโต โยคาวจโร “อนิจฺจนฺ”ติ วา “ทุกฺขนฺ”ติ วา “อนตฺตา”ติ วา สมฺมปฺปญฺญาย ปสฺสติ. เอวํ โข มหาราช โอภาสน\-ลกฺขณา ปญฺญาติ.}

\prosewithcloserruled{กลฺโลสิ ภนฺเต นาคเสนาติ.}{ปญฺญา\-ลกฺขณ\-ปโญฺห ปนฺนรสโม.}
\csromanmatikaanchor{mtk.20}
\csromantocmarkref{subsection}{16. นานาธมฺมานํ เอกกิจฺจอภินิปฺผาทนปญฺห}{mtk.20}
\bookmark[dest={mtk.20},level=2]{16. นานาธมฺมานํ เอกกิจฺจอภินิปฺผาทนปญฺห}
\csromanheader{2}{16. นานาธมฺมานํ เอกกิจฺจอภินิปฺผาทนปญฺห}
\proseitem{16}{ราชา อาห “ภนฺเต นาคเสน อิเม ธมฺมา นานา สนฺตา เอกํ อตฺถํ อภินิ\-ปฺผาเทนฺตี”ติ. อาม มหาราช อิเม ธมฺมา นานา สนฺตา เอกํ อตฺถํ อภินิ\-ปฺผาเทนฺติ, กิเลเส หนนฺตีติ.}

\prose{กถํ ภนฺเต อิเม ธมฺมา นานา สนฺตา เอกํ อตฺถํ อภินิ\-ปฺผาเทนฺติ, กิเลเส หนนฺติ. โอปมฺมํ กโรหีติ. ยถา มหาราช เสนา นานา สนฺตา หตฺถี จ อสฺสา จ รถา จ ปตฺตี จ เอกํ อตฺถํ อภินิ\-ปฺผาเทนฺติ, สงฺคาเม ปรเสนํ อภิวิชินนฺติ. เอวเมว โข มหาราช อิเม ธมฺมา นานา สนฺตา เอกํ อตฺถํ อภินิ\-ปฺผาเทนฺติ, กิเลเส หนนฺตีติ.}

\setlength{\csromangathapagewidth}{0pt}
\setlength{\csromangathawakwidth}{0pt}
\csromangathameasuregroup{}{กลฺโลสิ ภนฺเต นาคเสนาติ. \\ นานาธมฺมานํ เอกกิจฺจอภินิปฺผาทนปโญฺห โสฬสโม.}
\csromangathameasurewak{กลฺโลสิ ภนฺเต นาคเสนาติ. \\ นานาธมฺมานํ เอกกิจฺจอภินิปฺผาทนปโญฺห โสฬสโม.}
\setlength{\csromangathaleftcol}{0pt}
\csromangathagroupwithclosermidruled{\csromangathastanza{\csromangathawak{\csromangathaitemline{16}{กลฺโลสิ ภนฺเต นาคเสนาติ.} \\ \csromangathacontline{นานาธมฺมานํ เอกกิจฺจอภินิปฺผาทนปโญฺห โสฬสโม.}}}}{\textbf{มหาวคฺโค ปฐโม.}}
\csromancenter{อิมสฺมิํ วคฺเค โสฬส ปญฺหา.}
\csromanmatikaanchor{mtk.21}
\csromantocmarknopage{section}{2. อทฺธานวคฺค}{mtk.21}
\bookmark[dest={mtk.21},level=1]{2. อทฺธานวคฺค}
\csromanheader{1}{2. \textbf{อทฺธานวคฺค}}
\csromanmatikaanchor{mtk.22}
\csromantocmarkref{subsection}{1. ธมฺมสนฺตติปญฺห}{mtk.22}
\bookmark[dest={mtk.22},level=2]{1. ธมฺมสนฺตติปญฺห}
\csromanheader{2}{1. \textbf{ธมฺม}\-\textbf{สนฺตติ}\-\textbf{ปญฺห}}
\proseitem{1}{\csromanfolio{40}ราชา อาห “ภนฺเต นาคเสน โย อุปฺปชฺชติ, โส เอว โส, อุทาหุ อญฺโญ”ติ. เถโร อาห “น จ โส, น จ อญฺโญ”ติ.}

\prose{โอปมฺมํ กโรหีติ. ตํ กิํ มญฺญสิ มหาราช ยทา ตฺวํ ทหโร ตรุโณ มนฺโท อุตฺตานเสยฺยโก อโหสิ, โส เยว\csromansharedfootnote{cb63cef12ad01f03}{โสเยว สนฺธิ \csromandash{} ม.พ.ป.} ตฺวํ เอตรหิ มหนฺโตติ. น หิ ภนฺเต อญฺโญ โส ทหโร ตรุโณ มนฺโท อุตฺตานเสยฺยโก อโหสิ, อญฺโญ อหํ เอตรหิ มหนฺโตติ. เอวํ สนฺเต โข มหาราช มาตาติปิ น ภวิสฺสติ, ปิตาติปิ น ภวิสฺสติ, อาจริโยติปิ น ภวิสฺสติ, สิปฺปวาติปิ น ภวิสฺสติ, สีลวาติปิ น ภวิสฺสติ, ปญฺญวาติปิ น ภวิสฺสติ, กินฺนุ โข มหาราช อญฺญา เอว กลลสฺส มาตา, อญฺญา อพฺพุทสฺส มาตา, อญฺญา เปสิยา มาตา, อญฺญา ฆนสฺส มาตา, อญฺญา ขุทฺทกสฺส มาตา, อญฺญา มหนฺตสฺส มาตา, อญฺโญ สิปฺปํ สิกฺขติ, อญฺโญ สิกฺขิโต ภวติ, อญฺโญ ปาปกมฺมํ กโรติ, อญฺญสฺส หตฺถปาทา ฉิชฺชนฺตีติ. น หิ ภนฺเต. ตฺวํ ปน ภนฺเต เอวํ วุตฺเต กิํ วเทยฺยาสีติ. เถโร อาห “อหญฺเญว โข มหาราช ทหโร อโหสิํ ตรุโณ มนฺโท อุตฺตานเสยฺยโก, อหญฺเญว เอตรหิ มหนฺโต, อิมเมว กายํ นิสฺสาย สพฺเพ เต เอกสงฺคหิตา”ติ.}

\prose{ภิยฺโย โอปมฺมํ กโรหีติ. ยถา มหาราช โกจิเทว ปุริโส ปทีปํ ปทีเปยฺย, กิํ โส สพฺพรตฺติํ ปทีเปยฺยาติ. อาม ภนฺเต สพฺพรตฺติํ ปทีเปยฺยาติ. กินฺนุ โข มหาราช ยา ปุริเม ยาเม อจฺจิ, สา มชฺฌิเม ยาเม อจฺจีติ. น หิ ภนฺเตติ. ยา มชฺฌิเม ยาเม อจฺจิ, สา ปจฺฉิเม ยาเม อจฺจีติ. น หิ ภนฺเตติ. กินฺนุ โข มหาราช อญฺโญ โส อโหสิ ปุริเม ยาเม ปทีโป, อญฺโญ มชฺฌิเม ยาเม ปทีโป, อญฺโญ ปจฺฉิเม ยาเม ปทีโปติ. น หิ ภนฺเต, ตํ เยว\csromansharedfootnote{cf7dc7e6d0b8f44d}{ตํเยว นิคฺคหีตสนฺธิ \csromandash{} ม.พ.ป.} นิสฺสาย สพฺพรตฺติํ ปทีปิโตติ. เอวเมว โข มหาราช ธมฺมา สนฺตติ สนฺทหติ, อญฺโญ อุปฺปชฺชติ, อญฺโญ นิรุชฺฌติ, อปุพฺพํ อจริมํ วิย สนฺทหติ, เตน น จ โส, น จ อญฺโญ, ปุริมวิญฺญาเณ ปจฺฉิม\-วิญฺญาณํ สงฺคหํ คจฺฉตีติ.}

\prose{\csromanfolio{41}ภิยฺโย โอปมฺมํ กโรหีติ. ยถา มหาราช ขีรํ ทุยฺหมานํ กาลนฺตเรน ทธิ ปริวตฺเตยฺย, ทธิโต นวนีตํ, นวนีตโต ฆตํ ปริวตฺเตยฺย, โย นุ โข มหาราช เอวํ วเทยฺย “ยํ เยว ขีรํ ตํ เยว\csromansharedfootnote{cf7dc7e6d0b8f44d}{ตํเยว นิคฺคหีตสนฺธิ \csromandash{} ม.พ.ป.} ทธิ, ยํ เยว ทธิ ตํ เยว\csromansharedfootnote{cf7dc7e6d0b8f44d}{ตํเยว นิคฺคหีตสนฺธิ \csromandash{} ม.พ.ป.} นวนีตํ, ยํ เยว นวนีตํ ตํ เยว\csromansharedfootnote{cf7dc7e6d0b8f44d}{ตํเยว นิคฺคหีตสนฺธิ \csromandash{} ม.พ.ป.} ฆตนฺ”ติ, สมฺมา นุ โข โส มหาราช วทมาโน วเทยฺยาติ. น หิ ภนฺเต, ตํเยว นิสฺสาย สมฺภูตนฺติ. เอวเมว โข มหาราช ธมฺมสนฺตติ สนฺทหติ, อญฺโญ อุปฺปชฺชติ, อญฺโญ นิรุชฺฌติ, อปุพฺพํ อจริมํ วิย สนฺทหติ, เตน น จ โส, น จ อญฺโญ, ปุริมวิญฺญาเณ ปจฺฉิม\-วิญฺญาณํ สงฺคหํ คจฺฉตีติ.}

\prosewithcloserruled{กลฺโลสิ ภนฺเต นาคเสนาติ.}{ธมฺม\-สนฺตติ\-ปโญฺห ปฐโม.}
\csromanmatikaanchor{mtk.23}
\csromantocmarkref{subsection}{2. ปฏิสนฺทหนปญฺห}{mtk.23}
\bookmark[dest={mtk.23},level=2]{2. ปฏิสนฺทหนปญฺห}
\csromanheader{2}{2. \textbf{ปฏิสนฺทห}\-\textbf{น}\-\textbf{ปญฺห}}
\proseitem{2}{ราชา อาห “ภนฺเต นาคเสน โย น ปฏิสนฺทหติ, ชานาติ โส “น ปฏิสนฺทหิสฺสามี”ติ. อาม มหาราช โย น ปฏิสนฺทหติ, ชานาติ โส “น ปฏิสนฺทหิสฺสามี”ติ. กถํ ภนฺเต ชานาตีติ. โย เหตุ โย ปจฺจโย มหาราช ปฏิสนฺทห\-นาย, ตสฺส เหตุสฺส ตสฺส ปจฺจยสฺส อุปรมา ชานาติ โส “น ปฏิสนฺทหิสฺสามี”ติ.}

\prose{โอปมฺมํ กโรหีติ. ยถา มหาราช กสฺสโก คหปติโก กสิตฺวา จ วปิตฺวา จ ธญฺญาคารํ ปริปูเรยฺย, โส อปเรน สมเยน เนว กสฺเสยฺย น วปฺเปยฺย, ยถา\-สมฺภตญฺจ ธญฺญํ ปริภุญฺเชยฺย วา วิสชฺเชยฺย วา ยถา ปจฺจยํ วา กเรยฺย, ชาเนยฺย โส มหาราช กสฺสโก คหปติโก “น เม ธญฺญาคารํ ปริปูเรสฺสตี”ติ. อาม ภนฺเต ชาเนยฺยาติ. กถํ ชาเนยฺยาติ. โย เหตุ โย ปจฺจโย ธญฺญาคารสฺส ปริปูรณาย, ตสฺส เหตุสฺส ตสฺส ปจฺจยสฺส อุปรมา ชานาติ “น เม ธญฺญาคารํ ปริปูเรสฺสตี”ติ. เอวเมว โข มหาราช โย เหตุ โย ปจฺจโย ปฏิสนฺทห\-นาย, ตสฺส เหตุสฺส ตสฺส ปจฺจยสฺส อุปรมา ชานาติ โส “น ปฏิสนฺทหิสฺสามี”ติ.}

\prose{กลฺโลสิ ภนฺเต นาคเสนาติ.}

\vspace{\tipitakaparskip}
\csromancenter{ปฏิสนฺทห\-น\-ปโญฺห ทุติโย.}
\csromanmatikaanchor{mtk.24}
\csromantocmarkref{subsection}{3. ญาณปญฺญาปญฺห}{mtk.24}
\bookmark[dest={mtk.24},level=2]{3. ญาณปญฺญาปญฺห}
\csromanheader{2}{3. \textbf{ญาณปญฺญาปญฺห}}
\proseitem{3}{\csromanfolio{42}ราชา อาห “ภนฺเต นาคเสน ยสฺส ญาณํ อุปฺปนฺนํ, ตสฺส ปญฺญา อุปฺปนฺนา”ติ. อาม มหาราช ยสฺส ญาณํ อุปฺปนฺนํ, ตสฺส ปญฺญา อุปฺปนฺนาติ. กิํ ภนฺเต ยญฺเญว ญาณํ สา เยว ปญฺญาติ. อาม มหาราช ยญฺเญว ญาณํ สา เยว ปญฺญาติ. ยสฺส ปน ภนฺเต ตญฺเญว ญาณํ สา เยว ปญฺญา อุปฺปนฺนา, กิํ สมฺมุเยฺหยฺย โส, อุทาหุ น สมฺมุเยฺหยฺยาติ. กตฺถจิ มหาราช สมฺมุเยฺหยฺย, กตฺถจิ น สมฺมุเยฺหยฺยาติ. กุหิํ ภนฺเต สมฺมุเยฺหยฺยาติ. อญฺญาตปุพฺเพสุ วา มหาราช สิปฺปฏฺฐาเนสุ, อคตปุพฺพาย วา ทิสาย, อสฺสุตปุพฺพาย วา นามปญฺญตฺติยา สมฺมุเยฺหยฺยติ. กุหิํ น สมฺมุเยฺหยฺยาติ. ยํ โข ปน มหาราช ตาย ปญฺญาย กตํ “อนิจฺจนฺ”ติ วา “ทุกฺขนฺ”ติ วา “อนตฺตา”ติ วา, ตหิํ น สมฺมุเยฺหยฺยาติ. โมโห ปนสฺส ภนฺเต กุหิํ คจฺฉตีติ. โมโห โข มหาราช ญาเณ อุปฺปนฺนมตฺเต ตตฺเถว นิรุชฺฌตีติ.}

\prose{โอปมฺมํ กโรหีติ. ยถา มหาราช โกจิเทว ปุริโส อนฺธการเคเห ปทีปํ อาโรเปยฺย, ตโต อนฺธกาโร นิรุชฺเฌยฺย, อาโลโก ปาตุภเวยฺย. เอวเมว โข มหาราช ญาเณ อุปฺปนฺนมตฺเต โมโห ตตฺเถว นิรุชฺฌตีติ.}

\prose{ปญฺญา ปน ภนฺเต กุหิํ คจฺฉตีติ, ปญฺญาปิ โข มหาราช สกิจฺจยํ กตฺวา ตตฺเถว นิรุชฺฌติ, ยํ ปน ตาย ปญฺญาย กตํ “อนิจฺจนฺ”ติ วา “ทุกฺขนฺ”ติ วา “อนตฺตา”ติ วา, ตํ น นิรุชฺฌตีติ.}

\prose{ภนฺเต นาคเสน ยํ ปเนตํ พฺรูสิ “ปญฺญา สกิจฺจยํ กตฺวา ตตฺเถว นิรุชฺฌติ, ยํ ปน ตาย ปญฺญาย กตํ “อนิจฺจนฺ”ติ วา “ทุกฺขนฺ”ติ วา “อนตฺตา”ติ วา, ตํ น นิรุชฺฌตี”ติ, ตสฺส โอปมฺมํ กโรหีติ. ยถา มหาราช โย โกจิ ปุริโส รตฺติํ เลขํ เปเสตุกาโม เลขกํ ปกฺโกสาเปตฺวา ปทีปํ อาโรเปตฺวา เลขํ ลิขาเปยฺย, ลิขิเต ปน เลเข ปทีปํ วิชฺฌาเปยฺย, วิชฺฌาปิเตปิ ปทีเป เลขํ น วินสฺเสยฺย. เอวเมว โข มหาราช ปญฺญา สกิจฺจยํ กตฺวา ตตฺเถว นิรุชฺฌติ, ยํ ปน ตาย ปญฺญาย กตํ “อนิจฺจนฺ”ติ วา “ทุกฺขนฺ”ติ วา “อนตฺตา”ติ วา, ตํ น นิรุชฺฌตีติ.}

\prose{\csromanfolio{43}ภิยฺโย โอปมฺมํ กโรหีติ. ยถา มหาราช ปุรตฺถิเมสุ ชนปเทสุ มนุสฺสา อนุฆรํ ปญฺจ ปญฺจ อุทกฆฏกานิ ฐเปนฺติ อาลิมฺปนํ วิชฺฌาเปตุํ, ฆเร ปทิตฺเต ตานิ ปญฺจ อุทกฆฏกานิ ฆรสฺสูปริ ขิปนฺติ, ตโต อคฺคิ วิชฺฌายติ, กิํ นุ โข มหาราช เตสํ มนุสฺสานํ เอวํ โหติ “ปุน เตหิ ฆเฏหิ ฆฏกิจฺจํ กริสฺสามาติ. น หิ ภนฺเต อลํ เตหิ ฆเฏหิ, กิํ เตหิ ฆเฏหีติ. ยถา มหาราช ปญฺจ อุทกฆฏกานิ, เอวํ ปญฺจินฺทฺริยานิ ทฏฺฐพฺพานิ สทฺธินฺทฺริยํ วีริยินฺทฺริยํ สตินฺทฺริยํ สมาธินฺทฺริยํ ปญฺญินฺทฺริยํ. ยถา เต มนุสฺสา, เอวํ โยคาวจโร ทฏฺฐพฺโพ. ยถา อคฺคิ, เอวํ กิเลสา ทฏฺฐพฺพา. ยถา ปญฺจหิ อุทกฆฏเกหิ อคฺคิ วิชฺฌาปียติ, เอวํ ปญฺจินฺทฺริเยหิ กิเลสา วิชฺฌาปิยนฺติ, วิชฺฌาปิตาปิ กิเลสา น ปุน สมฺภวนฺติ. เอวเมว โข มหาราช ปญฺญา สกิจฺจยํ กตฺวา ตตฺเถว นิรุชฺฌติ, ยํ ปน ตาย ปญฺญาย กตํ “อนิจฺจนฺ”ติ วา “ทุกฺขนฺ”ติ วา “อนตฺตา”ติ วา, ตํ น นิรุชฺฌตีติ.}

\prose{ภิยฺโย โอปมฺมํ กโรหีติ. ยถา มหาราช เวชฺโช ปญฺจ\-มูลเภสชฺชานิ คเหตฺวา คิลานกํ อุปสงฺกมิตฺวา ตานิ ปญฺจ\-มูลเภสชฺชานิ ปิสิตฺวา\csromansharedfootnote{3b655640083f44dd}{ปิํสิตฺวา (สี, อิ)} คิลานกํ ปาเยยฺย, เตหิ จ โทสา นิทฺธเมยฺยุํ, กิํ นุ โข มหาราช ตสฺส เวชฺชสฺส เอวํ โหติ “ปุน เตหิ ปญฺจ\-มูลเภสชฺเชหิ เภสชฺชกิจฺจํ กริสฺสามี”ติ. น หิ ภนฺเต อลํ เตหิ ปญฺจ\-มูลเภสชฺเชหิ, กิํ เตหิ ปญฺจ\-มูลเภสชฺเชหีติ. ยถา มหาราช ปญฺจ\-มูลเภสชฺชานิ, เอวํ ปญฺจินฺทฺริยานิ ทฏฺฐพฺพานิ สทฺธินฺทฺริยํ วีริยินฺทฺริยํ สตินฺทฺริยํ สมาธินฺทฺริยํ ปญฺญินฺทฺริยํ. ยถา เวชฺโช, เอวํ โยคาวจโร ทฏฺฐพฺโพ. ยถา พฺยาธิ, เอวํ กิเลสา ทฏฺฐพฺพา. ยถา พฺยาธิโต ปุริโส, เอวํ ปุถุชฺชโน ทฏฺฐพฺโพ. ยถา ปญฺจ\-มูลเภสชฺเชหิ คิลานสฺส โทสา นิทฺธนฺตา, โทเส นิทฺธนฺเต คิลาโน อโรโค โหติ, เอวํ ปญฺจินฺทฺริเยหิ กิเลสา นิทฺธมียนฺติ, นิทฺธมิตา จ กิเลสา น ปุน สมฺภวนฺติ. เอวเมว โข มหาราช ปญฺญา สกิจฺจยํ กตฺวา ตตฺเถว นิรุชฺฌติ, ยํ ปน ตาย ปญฺญาย กตํ “อนิจฺจนฺ”ติ วา “ทุกฺขนฺ”ติ วา “อนตฺตา”ติ วา, ตํ น นิรุชฺฌตีติ.}

\prose{ภิยฺโย โอปมฺมํ กโรหีติ. ยถา มหาราช สงฺคามาวจโร โยโธ ปญฺจ กณฺฑานิ คเหตฺวา สงฺคามํ โอตเรยฺย ปรเสนํ วิเชตุํ, โส สงฺคามคโต ตานิ ปญฺจ กณฺฑานิ ขิเปยฺย, เตหิ จ ปรเสนา}

\csromanmatikaanchor{mtk.25}
\csromantocmarkref{subsection}{4. ปฏิสนฺทหนปุคฺคลเวทิยปญฺห}{mtk.25}
\bookmark[dest={mtk.25},level=2]{4. ปฏิสนฺทหนปุคฺคลเวทิยปญฺห}
\prose{\csromanfolio{44}ภิชฺเชยฺย, กินฺนุ โข มหาราช ตสฺส สงฺคามา\-วจรสฺส โยธสฺส เอวํ โหติ “ปุน เตหิ กณฺเฑหิ กณฺฑกิจฺจํ กริสฺสามี”ติ. น หิ ภนฺเต อลํ เตหิ กณฺเฑหิ, กิํ เตหิ กณฺเฑหีติ. ยถา มหาราช ปญฺจ กณฺฑานิ, เอวํ ปญฺจินฺทฺริยานิ ทฏฺฐพฺพานิ สทฺธินฺทฺริยํ วีริยินฺทฺริยํ สตินฺทฺริยํ สมาธินฺทฺริยํ ปญฺญินฺทฺริยํ. ยถา มหาราช สงฺคามาวจโร โยโธ, เอวํ โยคาวจโร ทฏฺฐพฺโพ. ยถา ปรเสน เอวํ กิเลสา ทฏฺฐพฺพา. ยถา ปญฺจหิ กณฺเฑหิ ปรเสนา ภิชฺชติ, เอวํ ปญฺจินฺทฺริเยหิ กิเลสา ภิชฺชนฺติ, ภคฺคา จ กิเลสา น ปุน สมฺภวนฺติ. เอวเมว โข มหาราช ปญฺญา สกิจฺจยํ กตฺวา ตตฺเถว นิรุชฺฌติ, ยํ ปน ตาย ปญฺญาย กตํ “อนิจฺจนฺ”ติ วา “ทุกฺขนฺ”ติ วา “อนตฺตา”ติ วา, ตํ น นิรุชฺฌตีติ.}

\prosewithcloserruled{กลฺโลสิ ภนฺเต นาคเสนาติ.}{ญาณปญฺญาปโญฺห ตติโย.}
\csromanheader{2}{4. \textbf{ปฏิสนฺทหนปุคฺคลเวทิยนปญฺห}}
\proseitem{4}{ราชา อาห “ภนฺเต นาคเสน โย น ปฏิสนฺทหติ, เวเทติ โส กิญฺจิ ทุกฺขํ เวทนนฺ”ติ. เถโร อาห “กิญฺจิ เวเทติ, กิญฺจิ น เวเทตี”ติ. กิํ เวเทติ, กิํ น เวเทตีติ. กายิกํ มหาราช เวทนํ เวเทติ, เจตสิกํ เวทนํ น เวเทตีติ. กถํ ภนฺเต กายิกํ เวทนํ เวเทติ, กถํ เจตสิกํ เวทนํ น เวเทตีติ. โย เหตุ โย ปจฺจโย กายิกาย ทุกฺขเวทนาย อุปฺปตฺติยา, ตสฺส เหตุสฺส ตสฺส ปจฺจยสฺส อนุปรมา กายิกํ ทุกฺขเวทนํ เวเทติ, โย เหตุ โย ปจฺจโย เจตสิกาย ทุกฺขเวทนาย อุปฺปตฺติยา, ตสฺส เหตุสฺส ตสฺส ปจฺจยสฺส อุปรมา เจตสิกํ ทุกฺขเวทนํ น เวเทติ. ภาสิตมฺเปตํ มหาราช ภควตา\csromandash{} “โส เอกํ เวทนํ เวเทติ กายิกํ น เจตสิกนฺ”ติ.}

\prose{ภนฺเต นาคเสน โย ทุกฺขํ เวทนํ เวเทติ, กสฺมา โส น ปรินิพฺพายตีติ. นตฺถิ มหาราช อรหโต อนุนโย วา ปฏิโฆ วา, น จ อรหนฺโต อปกฺกํ ปาเตนฺติ ปริปากํ อาคเมนฺติ ปณฺฑิตา. ภาสิตมฺเปตํ มหาราช เถเรน สาริปุตฺเตน ธมฺม\-เสนาปตินา\csromandash{}}

\setlength{\csromangathapagewidth}{0pt}
\setlength{\csromangathawakwidth}{0pt}
\csromangathameasuregroup{“นาภินนฺทามิ มรณํ, \\ กาลญฺจ ปฏิกงฺขามิ, \\ นาภินนฺทามิ มรณํ, \\ กาลญฺจ ปฏิกงฺขามิ,}{\csromangathabat{“นาภินนฺทามิ มรณํ,}{นาภินนฺทามิ ชีวิตํ.} \\ \csromangathabat{กาลญฺจ ปฏิกงฺขามิ,}{นิพฺพิสํ ภตโก ยถา.} \\ \csromangathabat{นาภินนฺทามิ มรณํ,}{นาภินนฺทามิ ชีวิตํ.} \\ \csromangathabat{กาลญฺจ ปฏิกงฺขามิ,}{สมฺปชาโน ปติสฺสโต”ติ\textsuperscript{1}.}}
\csromangathasetleft{“นาภินนฺทามิ มรณํ, \\ กาลญฺจ ปฏิกงฺขามิ, \\ นาภินนฺทามิ มรณํ, \\ กาลญฺจ ปฏิกงฺขามิ,}
\csromangathagroup{\csromangathastanza{\csromanfolio{45}\csromangathabat{“นาภินนฺทามิ มรณํ,}{นาภินนฺทามิ ชีวิตํ.} \\ \csromangathabat{กาลญฺจ ปฏิกงฺขามิ,}{นิพฺพิสํ ภตโก ยถา.}}\csromangathastanza{\csromangathabat{นาภินนฺทามิ มรณํ,}{นาภินนฺทามิ ชีวิตํ.} \\ \csromangathabat{กาลญฺจ ปฏิกงฺขามิ,}{สมฺปชาโน ปติสฺสโต”ติ\csromansharedfootnote{5146d3ed92062e2c}{ปสฺส ขุ 2. 314 ปิฏฺเฐ.}.}}}

\prosewithcloserruled{กลฺโลสิ ภนฺเต นาคเสนาติ.}{ปฏิสนฺทหนปุคฺคลเวทิยนปโญฺห จตุตฺโถ.}
\csromanmatikaanchor{mtk.26}
\csromantocmarkref{subsection}{5. เวทนาปญฺห}{mtk.26}
\bookmark[dest={mtk.26},level=2]{5. เวทนาปญฺห}
\csromanheader{2}{5. \textbf{เวทนาปญฺห}}
\proseitem{5}{ราชา อาห “ภนฺเต นาคเสน สุขา เวทนา กุสลา วา อกุสลา วา อพฺยากตา วา”ติ. สิยา มหาราช กุสลา, สิยา อกุสลา, สิยา อพฺยากตาติ. ยทิ ภนฺเต กุสลา น ทุกฺขา, ยทิ ทุกฺขา น กุสลา, กุสลํ ทุกฺขนฺติ นุปฺปชฺชตีติ. ตํ กิํ มญฺญสิ มหาราช, อิธ ปุริสสฺส หตฺเถ ตตฺถํ อโยคุฬํ นิกฺขิเปยฺย, ทุติเย หตฺเถ สีตํ หิมปิณฺฑํ นิกฺขิเปยฺย, กิํ นุ โข มหาราช อุโภปิ เต ทเหยฺยุนฺติ. อาม ภนฺเต อุโภปิ เต ทเหยฺยุนฺติ. กิํ นุ โข เต มหาราช อุโภปิ อุณฺหาติ. น หิ ภนฺเตติ. กิํ ปน เต มหาราช อุโภปิ สีตลาติ. น หิ ภนฺเตติ. อาชานาหิ นิคฺคหํ ยทิ ตตฺตํ ทหติ, น จ เต อุโภปิ อุณฺหา, เตน นุปฺปชฺชติ. ยทิ สีตลํ ทหติ, น จ เต อุโภปิ สีตลา, เตน นุปฺปชฺชติ. กิสฺส ปน เต มหาราช อุโภปิ ทหนฺติ, น จ เต อุโภปิ อุณฺหา, น จ เต อุโภปิ สีตลา, เอกํ อุณฺหํ, เอกํ สีตลํ, อุโภปิ เต ทหนฺติ, เตน นุปฺปชฺชตีติ. นาหํ ปฏิพโล ตยา วาทินา สทฺธิํ สลฺลปิตุํ, สาธุ อตฺถํ ชปฺเปหีติ. ตโต เถโร อภิธมฺม\-สํยุตฺตาย กถาย ราชานํ มิลินฺทํ สญฺญาเปสิ\csromandash{}}

\prose{ฉยิมานิ มหาราช เคหนิสฺสิตานิ โสมนสฺสานิ, ฉ เนกฺขมฺม\-นิสฺสิตานิ โสมนสฺสานิ, ฉ เคหนิสฺสิตานิ โทมนสฺสานิ, ฉ เนกฺขมฺม\-นิสฺสิตานิ โทมนสฺสานิ, ฉ เคหนิสฺสิตา อุเปกฺขา, ฉ เนกฺขมฺม\-นิสฺสิตา อุเปกฺขาติ, อิมานิ ฉ ฉกฺกานิ, อตีตาปิ ฉตฺติํสวิธา เวทนา, \csromanfolio{46}อนาคตาปิ ฉตฺติํสวิธา เวทนา, ปจฺจุปฺปนฺนาปิ ฉตฺติํสวิธา เวทนา, ตเทกชฺฌํ อภิสญฺญูหิตฺวา อภิสมฺปิณฺเฑตฺวา อฏฺฐสตํ เวทนา โหนฺตีติ.}

\prosewithcloserruled{กลฺโลสิ ภนฺเต นาคเสนาติ.}{เวทนาปโญฺห ปญฺจโม.}
\csromanmatikaanchor{mtk.27}
\csromantocmarkref{subsection}{6. นามรูปเอกตฺตนานตฺตปญฺห}{mtk.27}
\bookmark[dest={mtk.27},level=2]{6. นามรูปเอกตฺตนานตฺตปญฺห}
\csromanheader{2}{6. นามรูปเอกตฺตนานตฺตปญฺห}
\proseitem{5}{ราชา อาห “ภนฺเต นาคเสน โก ปฏิสนฺทหตี”ติ. เถโร อาห “นามรูปํ โข มหาราช ปฏิสนฺทหตี”ติ. กิํ อิมํ เยว นามรูปํ ปฏิสนฺทหตีติ. น โข มหาราช อิมํ เยว นามรูปํ ปฏิสนฺทหติ, อิมินา ปน มหาราช นามรูเปน กมฺมํ กโรติ โสภนํ วา ปาปกํ วา, เตน กมฺเมน อญฺญํ นามรูปํ ปฏิสนฺทหตีติ. ยทิ ภนฺเต น อิมํ เยว นามรูปํ ปฏิสนฺทหติ, นนุ โส มุตฺโต ภวิสฺสติ ปาปเกหิ กมฺเมหีติ. เถโร อาห “ยทิ น ปฏิสนฺทเหยฺย, มุตฺโต ภเวยฺย ปาปเกหิ กมฺเมหิ. ยสฺมา จ โข มหาราช ปฏิสนฺทหติ, ตสฺมา น มุตฺโต ปาปเกหิ กมฺเมหี”ติ.}

\prose{โอปมฺมํ กโรหีติ. ยถา มหาราช โกจิเทว ปุริโส อญฺญตรสฺส ปุริสสฺส อมฺพํ อวหเรยฺย, ตเมนํ อมฺพสามิโก คเหตฺวา รญฺโญ ทสฺเสยฺย “อิมินา เทว ปุริเสน มยฺหํ อมฺพา อวหฏา”ติ, โส เอวํ วเทยฺย “นาหํ เทว อิมสฺส อมฺเพ อวหรามิ, อญฺเญ เต อมฺพา, เย อิมินา โรปิตา, อญฺเญ เต อมฺพา, เย มยา อวหฏา, นาหํ ทณฺฑปฺปตฺโต”ติ. กินฺนุ โข โส มหาราช ปุริโส ทณฺฑปฺปตฺโต ภเวยฺยาติ. อาม ภนฺเต ทณฺฑปฺปตฺโต ภเวยฺยาติ. เกน การเณนาติ. กิญฺจาปิ โส เอวํ วเทยฺย, ปุริมํ ภนฺเต อมฺพํ อปฺปจฺจกฺขาย ปจฺฉิเมน อมฺเพน โส ปุริโส ทณฺฑปฺปตฺโต ภเวยฺยาติ. เอวเมว โข มหาราช อิมินา นามรูเปน กมฺมํ กโรติ โสภนํ วา ปาปกํ วา, เตน กมฺเมน อญฺญํ นามรูปํ ปฏิสนฺทหติ, ตสฺมา น มุตฺโต ปาปเกหิ กมฺเมหีติ.}

\prose{ภิยฺโย โอปมฺมํ กโรหีติ. ยถา มหาราช โกจิเทว ปุริโส อญฺญตรสฺส ปุริสสฺส สาลิํ อวหเรยฺย ฯเปฯ. อุจฺฉุํ อวหเรยฺย ฯเปฯ. ยถา มหาราช โกจิ ปุริโส เหมนฺตกาเล อคฺคิํ ชาเลตฺวา \csromanfolio{47}วิสิพฺเพตฺวา\csromansharedfootnote{63bfb3128dbb020e}{วิสีเวตฺวา (สี, อิ)} อวิชฺฌาเปตฺวา ปกฺกเมยฺย, อถ โข โส อคฺคิ อญฺญตรสฺส ปุริสสฺส เขตฺตํ ฑเหยฺย\csromansharedfootnote{38198bd309afd8d1}{อุปฑเหยฺย (ก)}, ตเมนํ เขตฺตสามิโก คเหตฺวา รญฺโญ ทสฺเสยฺย “อิมินา เทว ปุริเสน มยฺหํ เขตฺตํ ทฑฺฒนฺ”ติ. โส เอวํ วเทยฺย “นาหํ เทว อิมสฺส เขตฺตํ ฌาเปมิ, อญฺโญ โส อคฺคิ, โย มยา อวิชฺฌาปิโต, อญฺโญ โส อคฺคิ, เยนิมสฺส เขตฺตํ ทฑฺฒํ, นาหํ ทณฺฑปฺปตฺโตติ. กิํ นุ โข โส มหาราช ปุริโส ทณฺฑปฺปตฺโต ภเวยฺยาติ. อาม ภนฺเต ทณฺฑปฺปตฺโต ภเวยฺยาติ. เกน การเณนาติ. กิญฺจาปิ โส เอวํ วเทยฺย, ปุริมํ ภนฺเต อคฺคิํ อปฺปจฺจกฺขาย ปจฺฉิเมน อคฺคินา โส ปุริโส ทณฺฑปฺปตฺโต ภเวยฺยาติ. เอวเมว โข มหาราช อิมินา นามรูเปน กมฺมํ กโรติ โสภนํ วา ปาปกํ วา, เตน กมฺเมน อญฺญํ นามรูปํ ปฏิสนฺทหติ, ตสฺมา น มุตฺโต ปาปเกหิ กมฺเมหีติ.}

\proseitem{6}{ภิยฺโย โอปมฺมํ กโรหีติ. ยถา มหาราช โกจิเทว ปุริโส ปทีปํ อาทาย ปาสาทํ อภิรูหิตฺวา ภุญฺเชยฺย, ปทีโป ฌายมาโน ติณํ ฌาเปยฺย, ติณํ ฌายมานํ ฆรํ ฌาเปยฺย, ฆรํ ฌายมานํ คามํ ฌาเปยฺย, คามชโน ตํ ปุริสํ คเหตฺวา เอวํ วเทยฺย “กิสฺส ตฺวํ โภ ปุริส คามํ ฌาเปสี”ติ, โส เอวํ วเทยฺย “นาหํ โภ คามํ ฌาเปมิ, อญฺโญ โส ปทีปคฺคิ, ยสฺสาหํ อาโลเกน ภุญฺชิํ, อญฺโญ โส อคฺคิ, เยน คาโม ฌาปิโต”ติ, เต วิวทมานา ตว สนฺติเก อาคจฺเฉยฺยุํ, กสฺส ตฺวํ มหาราช อฏฺฏํ\csromansharedfootnote{380cd2f410186f22}{อตฺถํ (สี, อิ)} ธาเรยฺยาสีติ. คามชนสฺส ภนฺเตติ. กิํ การณาติ. กิญฺจาปิ โส เอวํ วเทยฺย, อปิ จ ตโต เอว โส อคฺคิ นิพฺพตฺโตติ. เอวเมว โข มหาราช กิญฺจาปิ อญฺญํ มารณนฺติกํ นามรูปํ, อญฺญํ ปฏิสนฺธิสฺมิํ นามรูปํ, อปิ จ ตโต เยว ตํ นิพฺพตฺตํ, ตสฺมา น มุตฺโต ปาปเกหิ กมฺเมหีติ.}

\prose{ภิยฺโย โอปมฺมํ กโรหีติ. ยถา มหาราช โกจิเทว ปุริโส ทหริํ ทาริกํ วาเรตฺวา สุงฺกํ ทตฺวา ปกฺกเมยฺย, สา อปเรน สมเยน มหตี อสฺส วยปฺปตฺตา, ตโต อญฺโญ ปุริโส สุงฺกํ ทตฺวา วิวาหํ กเรยฺย, อิตโร อาคนฺตฺวา เอวํ วเทยฺย “กิสฺส ปน เม ตฺวํ อมฺโภ ปุริส ภริยํ เนสี”ติ. โส เอวํ วเทยฺย “นาหํ ตว ภริยํ เนมิ, อญฺญา สา ทาริกา ทหรี ตรุณี, \csromanfolio{48}ยา ตยา วาริตา จ ทินฺนสุงฺกา จ, อญฺญายํ ทาริกา มหตี วยปฺปตฺตา มยา วาริตา จ ทินฺนสุงฺกา จา”ติ, เต วิวทมานา ตว สนฺติเก อาคจฺเฉยฺยุํ. กสฺส ตฺวํ มหาราช อฏฺฏํ ธาเรยฺยาสีติ. ปุริมสฺส ภนฺเตติ. กิํ การณาติ. กิญฺจาปิ โส เอวํ วเทยฺย, อปิ จ ตโต เยว สา มหตี นิพฺพตฺตาติ. เอวเมว โข มหาราช กิญฺจาปิ อญฺญํ มารณนฺติกํ นามรูปํ, อญฺญํ ปฏิสนฺธิสฺมิํ นามรูปํ, อปิ จ ตโต เยว ตํ นิพฺพตฺตํ, ตสฺมา น ปริมุตฺโต ปาปเกหิ กมฺเมหีติ.}

\proseitem{2}{ภิยฺโย โอปมฺมํ กโรหีติ. ยถา มหาราช โกจิเทว ปุริโส โคปาลกสฺส หตฺถโต ขีรฆฏํ กิณิตฺวา ตสฺเสว หตฺเถ นิกฺขิปิตฺวา ปกฺกเมยฺย “เสฺว คเหตฺวา คมิสฺสามี”ติ, ตํ อปรชฺชุ ทธิ สมฺปชฺเชยฺย, โส อาคนฺตฺวา เอวํ วเทยฺย “เทหิ เม ขีรฆฏนฺ”ติ. โส ทธิํ ทสฺเสยฺย. อิตโร เอวํ วเทยฺย “นาหํ ตว หตฺถโต ทธิํ กิณามิ, เทหิ เม ขีรฆฏนฺ”ติ. โส เอวํ วเทยฺย “อชานโต เต ขีรํ ทธิภูตนฺ”ติ เต วิวทมานา ตว สนฺติเก อาคจฺเฉยฺยุํ, กสฺส ตฺวํ มหาราช อฏฺฏํ ธาเรยฺยาสีติ. โคปาลกสฺส ภนฺเตติ. กิํ การณาติ. กิญฺจาปิ โส เอวํ วเทยฺย, อปิ จ ตโต เยว ตํ นิพฺพตฺตนฺติ. เอวเมว โข มหาราช กิญฺจาปิ อญฺญํ มารณนฺติกํ นามรูปํ, อญฺญํ ปฏิสนฺธิสฺมิํ นามรูปํ, อปิ จ ตโต เยว ตํ นิพฺพตฺตํ, ตสฺมา น ปริมุตฺโต ปาปเกหิ กมฺเมหีติ.}

\prose{กลฺโลสิ ภนฺเต นาคเสนาติ.}

\prose{นามรูปเอกตฺตนานตฺตปโญฺห ฉฏฺโฐ.}
\csromansectionrule

\csromanmatikaanchor{mtk.28}
\csromantocmarkref{subsection}{7. เถรปฏิสนฺทหนาปฏิสนฺทหนปญฺห}{mtk.28}
\bookmark[dest={mtk.28},level=2]{7. เถรปฏิสนฺทหนาปฏิสนฺทหนปญฺห}
\csromanheader{2}{7. \textbf{เถรปฏิสนฺทหนาปฏิสนฺทหนปญฺห}}
\proseitem{7}{ราชา อาห “ภนฺเต นาคเสน ตฺวํ ปน ปฏิสนฺทหิสฺสสี”ติ. อลํ มหาราช กิํ เต เตน ปุจฺฉิเตน, นนุ มยา ปฏิกจฺเจว อกฺขาตํ “สเจ มหาราช สอุปาทาโน ภวิสฺสามิ, ปฏิสนฺทหิสฺสามิ, สเจ อนุปาทาโน ภวิสฺสามิ, น ปฏิสนฺทหิสฺสามี”ติ.}

\prose{โอปมฺมํ กโรหีติ. ยถา มหาราช โกจิเทว ปุริโส รญฺโญ อธิการํ กเรยฺย, ราชา ตุฏฺโฐ อธิการํ ทเทยฺย, โส เตน อธิกาเรน ปญฺจหิ กามคุเณหิ สมปฺปิโต สมงฺคิภูโต ปริจเรยฺย, \csromanfolio{49}โส เจ ชนสฺส อาโรเจยฺย “น เม ราชา กิญฺจิ ปฏิกโรตี”ติ. กินฺนุ โข โส มหาราช ปุริโส ยุตฺตการี ภเวยฺยาติ. น หิ ภนฺเตติ. เอวเมว โข มหาราช กิํ เต เตน ปุจฺฉิเตน, นนุ มยา ปฏิกจฺเจว อกฺขาตํ “สเจ ส- อุปาทาโน ภวิสฺสามิ ปฏิสนฺทหิสฺสามิ, สเจ อนุปาทาโน ภวิสฺสามิ, น ปฏิสนฺทหิสฺสามี”ติ.}

\prosewithcloserruled{กลฺโลสิ ภนฺเต นาคเสนาติ.}{เถรปฏิสนฺทหนาปฏิสนฺทหนปโญฺห สตฺตโม.}
\csromanmatikaanchor{mtk.29}
\csromantocmarkref{subsection}{8. นามรูปปฏิสนฺทหนปญฺห}{mtk.29}
\bookmark[dest={mtk.29},level=2]{8. นามรูปปฏิสนฺทหนปญฺห}
\csromanheader{2}{8. \textbf{นามรูป}\-\textbf{ปฏิสนฺทห}\-\textbf{น}\-\textbf{ปญฺห}}
\proseitem{8}{ราชา อาห “ภนฺเต นาคเสน ยํ ปเนตํ พฺรูสิ ‘นามรูปนฺ’ติ, ตตฺถ กตมํ นามํ, กตมํ รูปนฺ”ติ. ยํ ตตฺถ มหาราช โอฬาริกํ, เอตํ รูปํ, เย ตตฺถ สุขุมา จิตฺตเจตสิกา ธมฺมา, เอตํ นามนฺติ. ภนฺเต นาคเสน เกน การเณน นามํ เยว น ปฏิสนฺทหติ, รูปํ เยว วาติ. อญฺญมญฺญู\-ปนิสฺสิตา มหาราช เอเต ธมฺมา เอกโตว อุปฺปชฺชนฺตีติ.}

\prose{โอปมฺมํ กโรหีติ. ยถา มหาราช กุกฺกุฏิยา กลลํ น ภเวยฺย, อณฺฑมฺปิ น ภเวยฺย. ยญฺจ ตตฺถ กลลํ, ยญฺจ อณฺฑํ, อุโภเปเต อญฺญมญฺญู\-ปนิสฺสิตา, เอกโตว เนสํ อุปฺปตฺติ โหติ. เอวเมว โข มหาราช ยทิ ตตฺถ นามํ น ภเวยฺย, รูปมฺปิ น ภเวยฺย, ยญฺเจว ตตฺถ นามํ, ยญฺเจว รูปํ, อุโภเปเต อญฺญมญฺญู\-ปนิสฺสิตา, เอกโตว เนสํ อุปฺปตฺติ โหติ. เอวเมตํ ทีฆมทฺธานํ สนฺธาวิตนฺติ.}

\prosewithcloserruled{กลฺโลสิ ภนฺเต นาคเสนาติ.}{นามรูป\-ปฏิสนฺทห\-น\-ปโญฺห อฏฺฐโม.}
\csromanmatikaanchor{mtk.30}
\csromantocmarkref{subsection}{9. อทฺธานปญฺห}{mtk.30}
\bookmark[dest={mtk.30},level=2]{9. อทฺธานปญฺห}
\csromanheader{2}{9. \textbf{อทฺธานปญฺห}}
\proseitem{9}{ราชา อาห “ภนฺเต นาคเสน ยํ ปเนตํ พฺรูสิ ‘ทีฆมทฺธานนฺ’ติ, กิเมตํ อทฺธานํ นามา”ติ. อตีโต มหาราช อทฺธา, อนาคโต อทฺธา, ปจฺจุปฺปนฺโน อทฺธาติ. กิํ ปน ภนฺเต สพฺเพ อทฺธา อตฺถีติ. โกจิ \csromanfolio{50}มหาราช อทฺธา อตฺถิ, โกจิ นตฺถีติ. กตโม ปน ภนฺเต อตฺถิ, กตโม นตฺถีติ. เย เต มหาราช สงฺขารา อตีตา วิคตา นิรุทฺธา วิปริณตา, โส อทฺธา นตฺถิ. เย ธมฺมา วิปากา, เย จ วิปาก\-ธมฺม\-ธมฺมา, เย จ อญฺญตฺร ปฏิสนฺธิํ เทนฺติ, โส อทฺธา อตฺถิ. เย สตฺตา กาลงฺกตา อญฺญตฺร อุปฺปนฺนา, โส จ อทฺธา อตฺถิ. เย สตฺตา กาลงฺกตา อญฺญตฺร อนุปฺปนฺนา, โส อทฺธา นตฺถิ. เย จ สตฺตา ปรินิพฺพุตา, โส จ อทฺธา นตฺถิ ปรินิพฺพุต\-ตฺตาติ.}

\prose{กลฺโลสิ ภนฺเต นาคเสนาติ.}

\vspace{\tipitakaparskip}
\csromancenter{อทฺธานปโญฺห นวโม.}
\nitthitammidruled{\textbf{อทฺธานวคฺโค ทุติโย.}}
\csromancenter{อิมสฺมิํ วคฺเค นว ปญฺหา.}
\csromanmatikaanchor{mtk.31}
\csromantocmarknopage{section}{3. วิจารวคฺค}{mtk.31}
\bookmark[dest={mtk.31},level=1]{3. วิจารวคฺค}
\csromanheader{1}{3. \textbf{วิจารวคฺค}}
\csromanmatikaanchor{mtk.32}
\csromantocmarkref{subsection}{1. อทฺธานมูลปญฺห}{mtk.32}
\bookmark[dest={mtk.32},level=2]{1. อทฺธานมูลปญฺห}
\csromanheader{2}{1. \textbf{อทฺธาน}\-\textbf{มูล}\-\textbf{ปญฺห}}
\proseitem{1}{\csromanfolio{51}ราชา อาห “ภนฺเต นาคเสน อตีตสฺส อทฺธานสฺส กิํ มูลํ, อนาคตสฺส อทฺธานสฺส กิํ มูลํ, ปจฺจุปฺปนฺนสฺส อทฺธานสฺส กิํ มูลนฺ”ติ. อตีตสฺส จ มหาราช อทฺธานสฺส อนาคตสฺส จ อทฺธานสฺส ปจฺจุปฺปนฺนสฺส จ อทฺธานสฺส อวิชฺชา มูลํ, อวิชฺชาปจฺจยา สงฺขารา, สงฺขาร\-ปจฺจยา วิญฺญาณํ, วิญฺญาณปจฺจยา นามรูปํ, นามรูป\-ปจฺจยา สฬายตนํ, สฬายตน\-ปจฺจยา ผสฺโส, ผสฺสปจฺจยา เวทนา, เวทนาปจฺจยา ตณฺหา, ตณฺหาปจฺจยา อุปาทานํ, อุปาทานปจฺจยา ภโว, ภวปจฺจยา ชาติ, ชาติปจฺจยา ชรามรณํ โสก\-ปริเทว\-ทุกฺข\-โทมนสฺสุ\-ปายาสา สมฺภวนฺติ. เอวเมตสฺส เกวลสฺส ทุกฺข\-กฺขนฺธสฺส อทฺธานสฺส\csromansharedfootnote{752335962175f101}{อทฺธานสฺส (อิ, ก)} ปุริมา โกฏิ น ปญฺญายตีติ.}

\prosewithcloserruled{กลฺโลสิ ภนฺเต นาคเสนาติ.}{อทฺธาน\-มูล\-ปโญฺห ปฐโม.}
\csromanmatikaanchor{mtk.33}
\csromantocmarkref{subsection}{2. ปุริมโกฏิปญฺห}{mtk.33}
\bookmark[dest={mtk.33},level=2]{2. ปุริมโกฏิปญฺห}
\csromanheader{2}{2. \textbf{ปุริม}\-\textbf{โกฏิ}\-\textbf{ปญฺห}}
\proseitem{2}{ราชา อาห “ภนฺเต นาคเสน ยํ ปเนตํ พฺรูสิ ‘ปุริมา โกฏิ น ปญฺญายตี’ติ, ตสฺส โอปมฺมํ กโรหี”ติ. ยถา มหาราช ปุริโส ปริตฺตํ\csromansharedfootnote{1652a26ef181c304}{ปริปกฺกํ (ก)} พีชํ ปถวิยํ นิกฺขิเปยฺย, ตโต องฺกุโร อุฏฺฐหิตฺวา อนุปุพฺเพน วุฑฺฒิํ วิรูฬฺหิํ เวปุลฺลํ อาปชฺชิตฺวา ผลํ ทเทยฺย. ตโต พีชํ คเหตฺวา ปุน โรเปยฺย, ตโตปิ องฺกุโร อุฏฺฐหิตฺวา อนุปุพฺเพน วุฑฺฒิํ วิรูฬฺหิํ เวปุลฺลํ อาปชฺชิตฺวา ผลํ ทเทยฺย. เอวเมติสฺสา สนฺตติยา อตฺถิ อนฺโตติ. นตฺถิ ภนฺเตติ. เอวเมว โข มหาราช อทฺธานสฺสาปิ ปุริมา โกฏิ น ปญฺญายตีติ.}

\prose{ภิยฺโย โอปมฺมํ กโรหีติ. ยถา มหาราช กุกฺกุฏิยา อณฺฑํ ภเวยฺย, อณฺฑโต กุกฺกุฏี กุกฺกุฏิยา อณฺฑนฺติ. เอวเมติสฺสา สนฺตติยา อตฺถิ อนฺโตติ. นตฺถิ ภนฺเตติ. เอวเมว โข มหาราช อทฺธานสฺสาปิ ปุริมา โกฏิ น ปญฺญายตีติ.}

\prose{\csromanfolio{52}ภิยฺโย โอปมฺมํ กโรหีติ. เถโร ปถวิยา จกฺกํ ลิขิตฺวา มิลินฺทํ ราชานํ เอตทโวจ “อตฺถิ มหาราช อิมสฺส จกฺกสฺส อนฺโต”ติ. นตฺถิ ภนฺเตติ. เอวเมว โข มหาราช อิมานิ จกฺกานิ วุตฺตานิ ภควตา “จกฺขุญฺจ ปฏิจฺจ รูเป จ อุปฺปชฺชติ จกฺขุวิญฺญาณํ, ติณฺณํ สงฺคติ ผสฺโส, ผสฺสปจฺจยา เวทนา, เวทนาปจฺจยา ตณฺหา, ตณฺหาปจฺจยา อุปาทานํ, อุปาทานปจฺจยา กมฺมํ, กมฺมโต ปุน จกฺขุํ ชายตี”ติ. เอวเมติสฺสา สนฺตติยา อตฺถิ อนฺโตติ. นตฺถิ ภนฺเตติ.}

\prose{โสตญฺจ ปฏิจฺจ สทฺเท จ ฯเปฯ. มนญฺจ ปฏิจฺจ ธมฺเม จ อุปฺปชฺชติ มโนวิญฺญาณํ, ติณฺณํ สงฺคติ ผสฺโส, ผสฺสปจฺจยา เวทนา, เวทนาปจฺจยา ตณฺหา, ตณฺหาปจฺจยา อุปาทานํ, อุปาทานปจฺจยา กมฺมํ, กมฺมโต ปุน มโน ชายตีติ. เอวเมติสฺสา สนฺตติยา อตฺถิ อนฺโตติ. นตฺถิ ภนฺเตติ. เอวเมว โข มหาราช อทฺธานสฺสาปิ ปุริมา โกฏิ น ปญฺญายตีติ.}

\prosewithcloserruled{กลฺโลสิ ภนฺเต นาคเสนาติ.}{ปุริม\-โกฏิ\-ปโญฺห ทุติโย.}
\csromanmatikaanchor{mtk.34}
\csromantocmarkref{subsection}{3. โกฏิปญฺญายนปญฺห}{mtk.34}
\bookmark[dest={mtk.34},level=2]{3. โกฏิปญฺญายนปญฺห}
\csromanheader{2}{3. \textbf{โกฏิปญฺญายนปญฺห}}
\proseitem{3}{ราชา อาห “ภนฺเต นาคเสน ยํ ปเนตํ พฺรูสิ ‘ปุริมา โกฏิ น ปญฺญายตี’ติ, กตมา จ สา ปุริมา โกฏี”ติ. โย โข มหาราช อตีโต อทฺธา, เอสา ปุริมา โกฏีติ. ภนฺเต นาคเสน ยํ ปเนตํ พฺรูสิ “ปุริมา โกฏิ น ปญฺญายตี”ติ, กิํ ปน ภนฺเต สพฺพาปิ ปุริมา โกฏิ น ปญฺญายตีติ. กาจิ มหาราช ปญฺญายติ, กาจิ น ปญฺญายตีติ. กตมา ภนฺเต ปญฺญายติ, กตมา น ปญฺญายตีติ. อิโต ปุพฺเพ มหาราช สพฺเพน สพฺพํ สพฺพถา สพฺพํ อวิชฺชา นาโหสีติ เอสา ปุริมา โกฏิ น ปญฺญายติ, ยํ อหุตฺวา สมฺโภติ, หุตฺวา ปฏิวิคจฺฉติ, เอสา ปุริมา โกฏิ ปญฺญายตีติ.}

\prose{ภนฺเต นาคเสน ยํ อหุตฺวา สมฺโภติ, หุตฺวา ปฏิวิคจฺฉติ, นนุ ตํ อุภโต ฉินฺนํ อตฺถํ คจฺฉตีติ. ยทิ มหาราช อุภโต ฉินฺนํ อตฺถํ คจฺฉติ, \csromanfolio{53}\textbf{อุภโต ฉินฺนา สกฺกา วฑฺเฒตุนฺติ.}\textbf{ อาม สาปิ สกฺกา วฑฺเฒตุนฺติ.}\textbf{ นาหํ} ภนฺเต เอตํ ปุจฺฉามิ โกฏิโต สกฺกา วฑฺเฒตุนฺติ. อาม สกฺกา วฑฺเฒตุนฺติ.}

\prose{โอปมฺมํ กโรหีติ. เถโร ตสฺส รุกฺขูปมํ อกาสิ, ขนฺธา จ เกวลสฺส ทุกฺข\-กฺขนฺธสฺส พีชานีติ.}

\prosewithcloserruled{กลฺโลสิ ภนฺเต นาคเสนาติ.}{โกฏิปญฺญายนปโญฺห ตติโย.}
\csromanmatikaanchor{mtk.35}
\csromantocmarkref{subsection}{4. สงฺขารชายมานปญฺห}{mtk.35}
\bookmark[dest={mtk.35},level=2]{4. สงฺขารชายมานปญฺห}
\csromanheader{2}{4. \textbf{สงฺขาร}\-\textbf{ชายมาน}\-\textbf{ปญฺห}}
\proseitem{3}{ราชา อาห “ภนฺเต นาคเสน อตฺถิ เกจิ สงฺขารา, เย ชายนฺตี”ติ. อาม มหาราช อตฺถิ สงฺขารา, เย ชายนฺตีติ. กตเม เต ภนฺเตติ. จกฺขุสฺมิญฺจ โข มหาราช สติ รูเปสุ จ จกฺขุวิญฺญาณํ โหติ, จกฺขุวิญฺญาเณ สติ จกฺขุ\-สมฺผสฺโส โหติ, จกฺขุ\-สมฺผสฺเส สติ เวทนา โหติ, เวทนาย สติ ตณฺหา โหติ, ตณฺหาย สติ อุปาทานํ โหติ, อุปาทาเน สติ ภโว โหติ, ภเว สติ ชาติ โหติ, ชาติยา สติ ชรามรณํ โสก\-ปริเทว\-ทุกฺข\-โทมนสฺสุ\-ปายาสา สมฺภวนฺติ, เอวเมตสฺส เกวลสฺส ทุกฺข\-กฺขนฺธสฺส สมุทโห โหติ. จกฺขุสฺมิญฺจ โข มหาราช อสติ รูเปสุ จ อสติ จกฺขุ วิญฺญาณํ น โหติ, จกฺขุวิญฺญาเณ อสติ จกฺขุ\-สมฺผสฺโส น โหติ, จกฺขุ\-สมฺผสฺเส อสติ เวทนา น โหติ, เวทนาย อสติ ตณฺหา น โหติ, ตณฺหาย อสติ อุปาทานํ น โหติ, อุปาทาเน อสติ ภโว น โหติ, ภเว อสติ ชาติ น โหติ, ชาติยา อสติ ชารามรณํ โสก\-ปริเทว\-ทุกฺข\-โทมนสฺสุ\-ปายาสา น โหนฺติ, เอวเมตสฺส เกวลสฺส ทุกฺข\-กฺขนฺธสฺส นิโรโธ โหตีติ.}

\prose{กลฺโลสิ ภนฺเต นาคเสนาติ.}

\vspace{\tipitakaparskip}
\csromancenter{สงฺขาร\-ชายมาน\-ปโญฺห จตุตฺโถ.}
\csromanmatikaanchor{mtk.36}
\csromantocmarkref{subsection}{5. ภวนฺตสงฺขารชายมานปญฺห}{mtk.36}
\bookmark[dest={mtk.36},level=2]{5. ภวนฺตสงฺขารชายมานปญฺห}
\csromanheader{2}{5. \textbf{ภวนฺต}\-\textbf{สงฺขาร}\-\textbf{ชายมาน}\-\textbf{ปญฺห}}
\proseitem{5}{\csromanfolio{54}ราชา อาห “ภนฺเต นาคเสน อตฺถิ เกจิ สงฺขารา, เย อภวนฺตา ชายนฺตี”ติ. นตฺถิ มหาราช เกจิ สงฺขารา, เย อภวนฺตา ชายนฺติ, ภวนฺตา เยว โข มหาราช สงฺขารา ชายนฺตีติ.}

\prose{โอปมฺมํ กโรหีติ. ตํ กิํ มญฺญสิ มหาราช, อิทํ เคหํ อภวนฺตํ ชาตํ, ยตฺถ ตฺวํ นิสินฺโนสีติ. นตฺถิ กิญฺจิ ภนฺเต อิธ อภวนฺตํ ชาตํ, ภวนฺตํ เยว ชาตํ, อิมานิ โข ภนฺเต ทารูนิ วเน อเหสุํ, อยญฺจ มตฺติกา ปถวิยํ อโหสิ, อิตฺถีนญฺจ ปุริสานญฺจ ตชฺเชน วายาเมน เอวมิทํ เคหํ นิพฺพตฺตนฺติ. เอวเมว โข มหาราช นตฺถิ เกจิ สงฺขารา, เย อภวนฺตา ชายนฺติ, ภวนฺตา เยว สงฺขารา ชายนฺตีติ.}

\prose{ภิยฺโย โอปมฺมํ กโรหีติ. ยถา มหาราช เย เกจิ พีชคาม\-ภูตคามา ปถวิยํ นิกฺขิตฺตา อนุปุพฺเพน วุฑฺฒิํ วิรูฬฺหิํ เวปุลฺลํ อาปชฺชมานา ปุปฺผานิ จ ผลานิ จ ทเทยฺยุํ, น เต รุกฺขา อภวนฺตา ชาตา, ภวนฺตา เยว เต รุกฺขา ชาตา. เอวเมว โข มหาราช นตฺถิ เกจิ สงฺขารา, เย อภวนฺตา ชายนฺติ, ภวนฺตา เยว เต สงฺขารา ชายนฺตีติ.}

\prose{ภิยฺโย โอปมฺมํ กโรหีติ. ยถา มหาราช กุมฺภกาโร ปถวิยา มตฺติกํ อุทฺธริตฺวา นานาภาชนานิ กโรติ, น ตานิ ภาชนานิ อภวนฺตานิ ชาตานิ, ภวนฺตานิ เยว ชาตานิ. เอวเมว โข มหาราช นตฺถิ เกจิ สงฺขารา, เย อภวนฺตา ชายนฺติ, ภวนฺตา เยว สงฺขารา ชายนฺตีติ.}

\prose{ภิยฺโย โอปมฺมํ กโรหีติ. ยถา มหาราช วีณาย ปตฺตํ น สิยา, จมฺมํ น สิยา, โทณิ น สิยา, ทณฺโฑ น สิยา, อุปวีโณ น สิยา, ตนฺติโย น สิยุํ, โกโณ น สิยา, ปุริสสฺส จ ตชฺโช วายาโม น สิยา, ชาเยยฺย สทฺโทติ. น หิ ภนฺเตติ. ยโต จ โข มหาราช วีณาย ปตฺตํ สิยา, จมฺมํ สิยา, โทณิ สิยา, ทณฺโฑ สิยา, อุปวีโณ สิยา, ตนฺติโย สิยุํ, โกโณ สิยา, ปุริสสฺส จ ตชฺโช วายาโม สิยา, ชาเยยฺย สทฺโทติ. อาม ภนฺเต ชาเยยฺยาติ. เอวเมว โข มหาราช นตฺถิ เกจิ สงฺขารา, เย อภวนฺตา ชายนฺติ, ภวนฺตา เยว โข สงฺขารา ชายนฺตีติ.}

\prose{ภิยฺโย โอปมฺมํ กโรหีติ. ยถา มหาราช อรณิ น สิยา, อรณิโปตโก น สิยา, อรณิโยตฺตกํ น สิยา, อุตฺตรารณิ \csromanfolio{55}น สิยา, โจฬกํ น สิยา, ปุริสสฺส จ ตชฺโช วายาโม น สิยา, ชาเยยฺย โส อคฺคีติ. น หิ ภนฺเตติ. ยโต จ โข มหาราช อรณิ สิยา, อรณิโปตโก สิยา, อรณิโยตฺตกํ สิยา, อุตฺตรารณิ สิยา, โจฬกํ สิยา, ปุริสสฺส จ ตชฺโช วายโม สิยา, ชาเยยฺย โส อคฺคีติ. อาม ภนฺเต ชาเยยฺยาติ. เอวเมว โข มหาราช นตฺถิ เกจิ สงฺขารา, เย อภวนฺตา ชายนฺติ, ภวนฺตา เยว โข สงฺขารา ชายนฺตีติ.}

\prose{ภิยฺโย โอปมฺมํ กโรหีติ. ยถา มหาราช มณิ น สิยา, อาตโป น สิยา, โคมยํ น สิยา, ชาเยยฺย โส อคฺคีติ. น หิ ภนฺเตติ. ยโต จ โข มหาราช มณิ สิยา, อาตโป สิยา, โคมยํ สิยา, ชาเยยฺย โส อคฺคีติ. อาม ภนฺเต ชาเยยฺยาติ. เอวเมว โข มหาราช นตฺถิ เกจิ สงฺขารา, เย อภวนฺตา ชายนฺติ, ภวนฺตา เยว โข สงฺขารา ชายนฺตีติ.}

\prose{ภิยฺโย โอปมฺมํ กโรหีติ. ยถา มหาราช อาทาโส น สิยา, อาภา น สิยา, มุขํ น สิยา, ชาเยยฺย อตฺตาติ. น หิ ภนฺเตติ. ยโต จ โข มหาราช อาทาโส สิยา, อาภา สิยา, มุขํ สิยา, ชาเยยฺย อตฺตาติ. อาม ภนฺเต ชาเยยฺยาติ. เอวเมว โข มหาราช นตฺถิ เกจิ สงฺขารา, เย อภวนฺตา ชายนฺติ, ภวนฺตา เยว โข สงฺขารา ชายนฺตีติ.}

\prosewithcloserruled{กลฺโลสิ ภนฺเต นาคเสนาติ.}{ภวนฺต\-สงฺขาร\-ชายมาน\-ปโญฺห ปญฺจโม.}
\csromanmatikaanchor{mtk.37}
\csromantocmarkref{subsection}{6. เวทคูปญฺห}{mtk.37}
\bookmark[dest={mtk.37},level=2]{6. เวทคูปญฺห}
\csromanheader{2}{6. \textbf{เวทคูปญฺห}}
\proseitem{5}{ราชา อาห “ภนฺเต นาคเสน เวทคู อุปลพฺภตี”ติ. โก ปเนส มหาราช เวทคู นามาติ. โย ภนฺเต อพฺภนฺตเร ชีโว จกฺขุนา รูปํ ปสฺสติ, โสเตน สทฺทํ สุณาติ, ฆาเนน คนฺธํ ฆายติ, ชิวฺหาย รสํ สายติ, กาเยน โผฏฺฐพฺพํ ผุสติ, มนสา ธมฺมํ วิชานาติ, ยถา มยํ อิธ ปาสาเท นิสินฺนา เยน เยน วาตปาเนน อิจฺเฉยฺยาม ปสฺสิตุํ, เตน เตน วาตปาเนน ปสฺเสยฺยาม, ปุรตฺถิเมนปิ วาตปาเนน ปสฺเสยฺยาม, \csromanfolio{56}ปจฺฉิเมนปิ วาตปาเนน ปสฺเสยฺยาม, อุตฺตเรนปิ วาตปาเนน ปสฺเสยฺยาม, ทกฺขิเณนปิ วาตปาเนน ปสฺเสยฺยาม. เอวเมว โข ภนฺเต อยํ อพฺภนฺตเร ชีโว เยน เยน ทฺวาเรน อิจฺฉติ ปสฺสิตุํ, เตน เตน ทฺวาเรน ปสฺสตีติ.}

\proseitem{6}{เถโร อาห “ปญฺจทฺวารํ มหาราช ภณิสฺสามิ, ตํ สุโณหิ, สาธุกํ มนสิกโรหิ, ยทิ อพฺภนฺตเร ชีโว จกฺขุนา รูปํ ปสฺสติ, ยถา มยํ อิธ ปาสาเท นิสินฺนา เยน เยน วาตปาเนน อิจฺเฉยฺยาม ปสฺสิตุํ, เตน เตน วาตปาเนน รูปํ เยว ปสฺเสยฺยาม, ปุรตฺถิเมนปิ วาตปาเนน รูปํ เยว ปสฺเสยฺยาม, ปจฺฉิเมนปิ วาตปาเนน รูปํ เยว ปสฺเสยฺยาม, อุตฺตเรนปิ วาตปาเนน รูปํ เยว ปสฺเสยฺยาม, ทกฺขิเณนปิ วาตปาเนน รูปํ เยว ปสฺเสยฺยาม, เอวเมเตน อพฺภนฺตเร ชีเวน โสเตนปิ รูปํ เยว ปสฺสิตพฺพํ, ฆาเนนปิ รูปํ เยว ปสฺสิตพฺพํ, ชิวฺหายปิ รูปํ เยว ปสฺสิตพฺพํ, กาเยนปิ รูปํ เยว ปสฺสิตพฺพํ, มนสาปิ รูปํ เยว ปสฺสิตพฺพํ. จกฺขุนาปิ สทฺโท เยว โสตพฺโพ, ฆาเนนปิ สทฺโท เยว โสตพฺโพ, ชิวฺหายปิ สทฺโท เยว โสตพฺโพ, กาเยนปิ สทฺโท เยว โสตพฺโพ, มนสาปิ สทฺโท เยว โสตพฺโพ. จกฺขุนาปิ คนฺโธ เยว ฆายิตพฺโพ, โสเตนปิ คนฺโธ เยว ฆายิตพฺโพ, ชิวฺหายปิ คนฺโธ เยว ฆายิตพฺโพ, กาเยนปิ คนฺโธ เยว ฆายิตพฺโพ, มนสาปิ คนฺโธ เยว ฆายิตพฺโพ. จกฺขุนาปิ รโส เยว สายิตพฺโพ, โสเตนปิ รโส เยว สายิตพฺโพ, ฆาเนนปิ รโส เยว สายิตพฺโพ, กาเยนปิ รโส เยว สายิตพฺโพ, มนสาปิ รโส เยว สายิตพฺโพ. จกฺขุนาปิ โผฏฺฐพฺพํ เยว ผุสิตพฺพํ, โสเตนปิ โผฏฺฐพฺพํ เยว ผุสิตพฺพํ, ฆาเนนปิ โผฏฺฐพฺพํ เยว ผุสิตพฺพํ, ชิวฺหายปิ โผฏฺฐพฺพํ เยว ผุสิตพฺพํ, มนสาปิ โผฏฺฐพฺพํ เยว ผุสิตพฺพํ. จกฺขุนาปิ ธมฺมํ เยว วิชานิตพฺพํ, โสเตนปิ ธมฺมํ เยว วิชานิตพฺพํ, ฆาเนนปิ ธมฺมํ เยว วิชานิตพฺพํ, ชิวฺหายปิ ธมฺมํ เยว วิชานิตพฺพํ, กาเยนปิ ธมฺมํ เยว วิชานิตพฺพนฺติ. น หิ ภนฺเตติ.}

\prose{น โข เต มหาราช ยุชฺชติ ปุริเมน วา ปจฺฉิมํ, ปจฺฉิเมน วา ปุริมํ, ยถา วา ปน มหาราช มยํ อิธ ปาสาเท นิสินฺนา อิเมสุ ชาลวาตปาเนสุ อุคฺฆาฏิเตสุ มหนฺเตน อากาเสน พหิมุขา สุฏฺฐุตรํ รูปํ ปสฺสาม, เอวเมเตน อพฺภนฺตเร ชีเวนาปิ จกฺขุทฺวาเรสุ อุคฺฆาฏิเตสุ มหนฺเตน อากาเสน สุฏฺฐุตรํ รูปํ ปสฺสิตพฺพํ, โสเตสุ อุคฺฆาฏิเตสุ. \csromanfolio{57}ฆาเน อุคฺฆาฏิเต. ชิวฺหาย อุคฺฆาฏิตาย. กาเย อุคฺฆาฏิเต มหนฺเตน อากาเสน สุฏฺฐุตรํ สทฺโท โสตพฺโพ, คนฺโธ ฆายิตพฺโพ, รโส สายิตพฺโพ, โผฏฺฐพฺโพ ผุสิตพฺโพติ. น หิ ภนฺเตติ.}

\prose{น โข เต มหาราช ยุชฺชติ ปุริเมน วา ปจฺฉิมํ, ปจฺฉิเมน วา ปุริมํ, ยถา วา ปน มหาราช อยํ ทินฺโน นิกฺขมิตฺวา พหิทฺวาร\-โกฏฺฐเก ติฏฺเฐยฺย, ชานาสิ ตฺวํ มหาราช “อยํ ทินฺโน นิกฺขมิตฺวา พหิทฺวาร\-โกฏฺฐเก ฐิโต”ติ. อาม ภนฺเต ชานามีติ. ยถา วา ปน มหาราช อยํ ทินฺโน อนฺโต ปวิสิตฺวา ตว ปุรโต ติฏฺเฐยฺย, ชานาสิ ตฺวํ มหาราช “อยํ ทินฺโน อนฺโต ปวิสิตฺวา มม ปุรโต ฐิโต”ติ. อาม ภนฺเต ชานามีติ. เอวเมว โข มหาราช อพฺภนฺตเร โส ชีโว ชิวฺหาย รเส นิกฺขิตฺเต ชาเนยฺย อมฺพิลตฺตํ วา ลวณตฺตํ วา ติตฺตกตฺตํ วา กฏุกตฺตํ วา กสายตฺตํ วา มธุรตฺตํ วาติ. อาม ภนฺเต ชาเนยฺยาติ. เต รเส อนฺโต ปวิฏฺเฐ ชาเนยฺย อมฺพิลตฺตํ วา ลวณตฺตํ วา ติตฺตกตฺตํ วา กฏุกตฺตํ วากสายตฺตํ วา มธุรตฺตํ วาติ. น หิ ภนฺเตติ.}

\prose{น โข เต มหาราช ยุชฺชติ ปุริเมน วา ปจฺฉิมํ, ปจฺฉิเมน วา ปุริมํ, ยถา มหาราช โกจิเทว ปุริโส มธุฆฏสตํ อาหราเปตฺวา มธุโทณิํ ปูราเปตฺวา ปุริสสฺส มุขํ ปิทหิตฺวา\csromansharedfootnote{ebec54fdf08f896d}{ปิทหิตฺวาว (ก)} มธุโทณิยา ปกฺขิเปยฺย, ชาเนยฺย มหาราช โส ปุริโส มธุํ สมฺปนฺนํ วา น สมฺปนฺนํ วาติ. น หิ ภนฺเตติ. เกน การเณนาติ. น หิ ตสฺส ภนฺเต มุเข มธุ ปวิฏฺฐนฺติ.}

\prose{น โข เต มหาราช ยุชฺชติ ปุริเมน วา ปจฺฉิมํ, ปจฺฉิเมน วา ปุริมนฺติ. นาหํ ปฏิพโล ตยา วาทินา สทฺธิํ สลฺลปิตุํ, สาธุ ภนฺเต อตฺถํ ชปฺเปหีติ.}

\prose{เถโร อภิธมฺม\-สํยุตฺตาย กถาย ราชานํ มิลินฺทํ สญฺญาเปสิ\csromandash{} อิธ มหาราช จกฺขุญฺจ ปฏิจฺจ รูเป จ อุปฺปชฺชติ จกฺขุวิญฺญาณํ, ตํสหชาตา ผสฺโส เวทนา สญฺญา เจตนา เอกคฺคตา ชีวิตินฺทฺริยํ มนสิกาโรติ เอวเมเต ธมฺมา ปจฺจยโต ชายนฺติ, น เหตฺถ เวทคู อุปลพฺภติ, โสตญฺจ ปฏิจฺจ สทฺเท จ ฯเปฯ มนญฺจ ปฏิจฺจ ธมฺเม จ อุปฺปชฺชติ มโนวิญฺญาณํ, ตํสหชาตา ผสฺโส เวทนา สญฺญา เจตนา เอกคฺคตา ชีวิตินฺทฺริยํ \csromanfolio{58}มนสิกาโรติ เอวเมเต ธมฺมา ปจฺจยโต ชายนฺติ, น เหตฺถ เวทคู อุปลพฺภตีติ.}

\proseitemwithcloserruled{2}{กลฺโลสิ ภนฺเต นาคเสนาติ.}{เวทคูปโญฺห ฉฏฺโฐ.}
\csromanmatikaanchor{mtk.38}
\csromantocmarkref{subsection}{7. จกฺขุวิญฺญาณาทิปญฺห}{mtk.38}
\bookmark[dest={mtk.38},level=2]{7. จกฺขุวิญฺญาณาทิปญฺห}
\csromanheader{2}{7. \textbf{จกฺขุวิญฺญาณา}\-\textbf{ทิ}\-\textbf{ปญฺห}}
\proseitem{7}{ราชา อาห “ภนฺเต นาคเสน ยตฺถ จกฺขุวิญฺญาณํ อุปฺปชฺชติ, ตตฺถ มโนวิญฺญาณมฺปิ อุปฺปชฺชตี”ติ. อาม มหาราช ยตฺถ จกฺขุวิญฺญาณํ อุปฺปชฺชติ, ตตฺถ มโนวิญฺญาณมฺปิ อุปฺปชฺชตีติ.}

\prose{กินฺนุ โข ภนฺเต นาคเสน ปฐมํ จกฺขุวิญฺญาณํ อุปฺปชฺชติ, ปจฺฉา มโนวิญฺญาณํ, อุทาหุ มโนวิญฺญาณํ ปฐมํ อุปฺปชฺชติ, ปจฺฉา จกฺขุ\-วิญฺญาณนฺติ. ปฐมํ มหาราช จกฺขุวิญฺญาณํ อุปฺปชฺชติ, ปจฺฉา มโนวิญฺญาณนฺติ.}

\prose{กินฺนุ โข ภนฺเต นาคเสน จกฺขุวิญฺญาณํ มโนวิญฺญาณํ อาณาเปติ “ยตฺถาหํ อุปฺปชฺชามิ, ตฺวมฺปิ ตตฺถ อุปฺปชฺชาหี”ติ, อุทาหุ มโนวิญฺญาณํ จกฺขุวิญฺญาณํ อาณาเปติ “ยตฺถ ตฺวํ อุปฺปชฺชิสฺสสิ, อหมฺปิ ตตฺถ อุปฺปชฺชิสฺสามี”ติ. น หิ มหาราช อนาลาโป เตสํ อญฺญมญฺเญหีติ.}

\prose{กถํ ภนฺเต นาคเสน ยตฺถ จกฺขุวิญฺญาณํ อุปฺปชฺชติ, ตตฺถ มโนวิญฺญาณมฺปิ อุปฺปชฺชตีติ. นินฺนตฺตา จ มหาราช ทฺวารตฺตา จ จิณฺณตฺตา จ สมุทาจริตตฺตา จาติ.}

\prose{กถํ ภนฺเต นาคเสน นินฺนตฺตา ยตฺถ จกฺขุวิญฺญาณํ อุปฺปชฺชตี, ตตฺถ มโนวิญฺญาณมฺปิ อุปฺปชฺชติ.}

\prose{โอปมฺมํ กโรหีติ. ตํ กิํ มญฺญสิ มหาราช, เทเว วสฺสนฺเต กตเมน อุทกํ คจฺเฉยฺยาติ. เยน ภนฺเต นินฺนํ, เตน คจฺเฉยฺยาติ. อถาปเรน สมเยน เทโว วสฺเสยฺย, กตเมน ตํ อุทกํ คจฺเฉยฺยาติ. เยน ภนฺเต ปุริมํ อุทกํ คตํ, ตมฺปิ เตน คจฺเฉยฺยาติ.}

\prose{กินฺนุ โข มหาราช ปุริมํ อุทกํ ปจฺฉิมํ อุทกํ อาณาเปติ “เยนาหํ คจฺฉามิ, ตฺวมฺปิ เตน คจฺฉาหี”ติ, ปจฺฉิมํ วา อุทกํ ปุริมํ อุทกํ อาณาเปติ “เยน ตฺวํ คจฺฉิสฺสสิ, อหมฺปิ เตน คจฺฉิสฺสามี”ติ. น หิ ภนฺเต, \csromanfolio{59}อนาลาโป เตสํ อญฺญมญฺเญหิ, นินฺนตฺตา คจฺฉนฺตีติ. เอวเมว โข มหาราช นินฺนตฺตา ยตฺถ จกฺขุวิญฺญาณํ อุปฺปชฺชติ, ตตฺถ มโนวิญฺญาณมฺปิ อุปฺปชฺชติ, น จกฺขุวิญฺญาณํ มโนวิญฺญาณํ อาณาเปติ “ยตฺถาหํ อุปฺปชฺชามิ, ตฺวมฺปิ ตตฺถ อุปฺปชฺชาหี”ติ, นาปิ มโนวิญฺญาณํ จกฺขุวิญฺญาณํ อาณาเปติ “ยตฺถ ตฺวํ อุปฺปชฺชิสฺสสิ, อหมฺปิ ตตฺถ อุปฺปชฺชิสฺสามี”ติ, อนาลาโป เตสํ อญฺญมญฺเญหิ, นินฺนตฺตา อุปฺปชฺชนฺตีติ.}

\prose{กถํ ภนฺเต นาคเสน ทฺวารตฺตา ยตฺถ จกฺขุวิญฺญาณํ อุปฺปชฺชติ, ตตฺถ มโนวิญฺญาณมฺปิ อุปฺปชฺชติ.}

\prose{โอปมฺมํ กโรหีติ. ตํ กิํ มญฺญสิ มหาราช, รญฺโญ ปจฺจนฺติมํ นครํ อสฺส ทฬฺห\-ปาการ\-โตรณํ เอกทฺวารํ, ตโต ปุริโส นิกฺขมิตุกาโม ภเวยฺย, กตเมน นิกฺขเมยฺยาติ. ทฺวาเรน ภนฺเต นิกฺขเมยฺยาติ. อถาปโร ปุริโส นิกฺขมิตุกาโม ภเวยฺย, กตเมน โส นิกฺขเมยฺยาติ. เยน ภนฺเต ปุริโม ปุริโส นิกฺขนฺโต, โสปิ เตน นิกฺขเมยฺยาติ.}

\prose{กินฺนุ โข มหาราช ปุริโม ปุริโส ปจฺฉิมํ ปุริสํ อาณาเปติ “เยนาหํ คจฺฉามิ, ตฺวมฺปิ เตน คจฺฉาหี”ติ, ปจฺฉิโม วา ปุริโส ปุริมํ ปุริสํ อาณาเปติ “เยน ตฺวํ คจฺฉิสฺสสิ, อหมฺปิ เตน คจฺฉิสฺสามี”ติ. น หิ ภนฺเต, อนาลาโป เตสํ อญฺญมญฺเญหิ, ทฺวารตฺตา คจฺฉนฺตีติ. เอวเมว โข มหาราช ทฺวารตฺตา ยตฺถ จกฺขุวิญฺญาณํ อุปฺปชฺชติ, ตตฺถ มโนวิญฺญาณมฺปิ อุปฺปชฺชติ, น จกฺขุวิญฺญาณํ มโนวิญฺญาณํ อาณาเปติ “ยตฺถาหํ อุปฺปชฺชามิ, ตฺวมฺปิ ตตฺถ อุปฺปชฺชาหี”ติ, นาปิ มโนวิญฺญาณํ จกฺขุวิญฺญาณํ อาณาเปติ “ยตฺถ ตฺวํ อุปฺปชฺชิสฺสสิ, อหมฺปิ ตตฺถ อุปฺปชฺชิสฺสามี”ติ, อนาลาโป เตสํ อญฺญมญฺเญหิ, ทฺวารตฺตา อุปฺปชฺชนฺตีติ.}

\prose{กถํ ภนฺเต นาคเสน จิณฺณตฺตา ยตฺถ จกฺขุวิญฺญาณํ อุปฺปชฺชติ, ตตฺถ มโนวิญฺญาณมฺปิ อุปฺปชฺชติ.}

\prose{โอปมฺมํ กโรหีติ. ตํ กิํ มญฺญสิ มหาราช, ปฐมํ เอกํ สกฏํ คจฺเฉยฺย, อถ ทุติยํ สกฏํ กตเมน คจฺเฉยฺยาติ. เยน ภนฺเต ปุริมํ สกฏํ คตํ, ตมฺปิ เตน คจฺเฉยฺยาติ.}

\prose{กินฺนุ โข มหาราช ปุริมํ สกฏํ ปจฺฉิมํ สกฏํ อาณาเปติ “เยนาหํ คจฺฉามิ, ตฺวมฺปิ เตน คจฺฉาหี”ติ, ปจฺฉิมํ วา สกฏํ ปุริมํ \csromanfolio{60}สกฏํ อาณาเปติ “เยน ตฺวํ คจฺฉิสฺสสิ, อหมฺปิ เตน คจฺฉิสฺสามี”ติ. น หิ ภนฺเต, อนาลาโป เตสํ อญฺญมญฺเญหิ, จิณฺณตฺตา คจฺฉนฺตีติ. เอวเมว โข มหาราช จิณฺณตฺตา ยตฺถ จกฺขุวิญฺญาณํ อุปฺปชฺชติ, ตตฺถ มโนวิญฺญาณมฺปิ อุปฺปชฺชติ, น จกฺขุวิญฺญาณํ มโนวิญฺญาณํ อาณาเปติ “ยตฺถาหํ อุปฺปชฺชามิ, ตฺวมฺปิ ตตฺถ อุปฺปชฺชาหี”ติ, นาปิ มโนวิญฺญาณํ จกฺขุวิญฺญาณํ อาณาเปติ “ยตฺถ ตฺวํ อุปฺปชฺชิสฺสสิ, อหมฺปิ ตตฺถ อุปฺปชฺชิสฺสามี”ติ, อนาลาโป เตสํ อญฺญมญฺเญหิ, จิณฺณตฺตา อุปฺปชฺชนฺตีติ.}

\proseitem{2}{กถํ ภนฺเต นาคเสน สมุทาจริตตฺตา ยตฺถ จกฺขุวิญฺญาณํ อุปฺปชฺชติ, ตตฺถ มโนวิญฺญาณมฺปิ อุปฺปชฺชติ, โอปมฺมํ กโรหีติ. ยถา มหาราช มุทฺทา\-คณนา\-สงฺขฺยา\-เลขา\-สิปฺปฏฺฐาเนสุ อาทิกมฺมิกสฺส ทนฺธายนา ภวติ, อถาปเรน สมเยน นิสมฺม\-กิริยาย สมุทาจริตตฺตา อทนฺธายนา ภวติ. เอวเมว โข มหาราช สมุทาจริตตฺตา ยตฺถ จกฺขุวิญฺญาณํ อุปฺปชฺชติ, ตตฺถ มโนวิญฺญาณมฺปิ อุปฺปชฺชติ, น จกฺขุวิญฺญาณํ มโนวิญฺญาณํ อาณาเปติ “ยตฺถาหํ อุปฺปชฺชามิ, ตฺวมฺปิ ตตฺถ อุปฺปชฺชาหี”ติ, นาปิ มโนวิญฺญาณํ จกฺขุวิญฺญาณํ อาณาเปติ “ยตฺถ ตฺวํ อุปฺปชฺชิสฺสสิ, อหมฺปิ ตตฺถ อุปฺปชฺชิสฺสามี”ติ, อนาลาโป เตสํ อญฺญมญฺเญหิ, สมุทาจริตตฺตา อุปฺปชฺชนฺตีติ.}

\prose{ภนฺเต นาคเสน ยตฺถ โสตวิญฺญาณํ อุปฺปชฺชติ, ตตฺถ มโนวิญฺญาณมฺปิ อุปฺปชฺชตีติ ฯเปฯ ยตฺถ ฆานวิญฺญาณํ อุปฺปชฺชติ. ยตฺถ ชิวฺหาวิญฺญาณํ อุปฺปชฺชติ. ยตฺถ กายวิญฺญาณํ อุปฺปชฺชติ, ตตฺถ มโนวิญฺญาณมฺปิ อุปฺปชฺชตีติ. อาม มหาราช ยตฺถ กายวิญฺญาณํ อุปฺปชฺชติ, ตตฺถ มโนวิญฺญาณมฺปิ อุปฺปชฺชตีติ.}

\prose{กินฺนุ โข ภนฺเต นาคเสน ปฐมํ กายวิญฺญาณํ อุปฺปชฺชติ, ปจฺฉา มโนวิญฺญาณํ, อุทาหุ มโนวิญฺญาณํ ปฐมํ อุปฺปชฺชติ, ปจฺฉา กายวิญฺญาณนฺติ. กายวิญฺญาณํ มหาราช ปฐมํ อุปฺปชฺชติ, ปจฺฉา มโนวิญฺญาณนฺติ.}

\prose{กินฺนุ โข ภนฺเต นาคเสน ฯเปฯ อนาลาโป เตสํ อญฺญมญฺเญหิ, สมุทาจริตตฺตา อุปฺปชฺชนฺตีติ.}

\prose{กลฺโลสิ ภนฺเต นาคเสนาติ.}

\vspace{\tipitakaparskip}
\csromancenter{จกฺขุวิญฺญาณา\-ทิ\-ปโญฺห สตฺตโม.}
\csromanmatikaanchor{mtk.39}
\csromantocmarkref{subsection}{8. ผสฺสลกฺขณปญฺห}{mtk.39}
\bookmark[dest={mtk.39},level=2]{8. ผสฺสลกฺขณปญฺห}
\csromanmarkh{2. มิลินฺทปญฺห}\csromanmarkhii{8. ผสฺส\-ลกฺขณ\-ปญฺห}\csromanheadernomark{2}{2. มิลินฺทปญฺห \textbf{8. ผสฺส}\-\textbf{ลกฺขณ}\-\textbf{ปญฺห}}
\proseitem{8}{\csromanfolio{61}ราชา อาห “ภนฺเต นาคเสน ยตฺถ มโนวิญฺญาณํ อุปฺปชฺชติ, ผสฺโสปิ เวทนาปิ ตตฺถ อุปฺปชฺชตี”ติ. อาม มหาราช ยตฺถ มโนวิญฺญาณํ อุปฺปชฺชติ, ผสฺโสปิ ตตฺถ อุปฺปชฺชติ, เวทนาปิ ตตฺถ อุปฺปชฺชติ, สญฺญาปิ ตตฺถ อุปฺปชฺชติ, เจตนาปิ ตตฺถ อุปฺปชฺชติ, วิตกฺโกปิ ตตฺถ อุปฺปชฺชติ, วิจาโรปิ ตตฺถ อุปฺปชฺชติ, สพฺเพปิ ผสฺสปฺปมุขา ธมฺมา ตตฺถ อุปฺปชฺชนฺตีติ.}

\prose{ภนฺเต นาคเสน กิํลกฺขโณ ผสฺโสติ. ผุสน\-ลกฺขโณ มหาราช ผสฺโสติ.}

\prose{โอปมฺมํ กโรหีติ. ยถา มหาราช เทฺว เมณฺฑา ยุชฺเฌยฺยุํ, เตสุ ยถา เอโก เมณฺโฑ, เอวํ จกฺขุ ทฏฺฐพฺพํ. ยถา ทุติโย เมณฺโฑ, เอวํ รูปํ ทฏฺฐพฺพํ. ยถา เตสํ สนฺนิปาโต, เอวํ ผสฺโส ทฏฺฐพฺโพติ.}

\prose{ภิยฺโย โอปมฺมํ กโรหีติ. ยถา มหาราช เทฺว ปาณี วชฺเชยฺยุํ, เตสุ ยถา เอโก ปาณิ, เอวํ จกฺขุ ทฏฺฐพฺพํ. ยถา ทุติโย ปาณิ, เอวํ รูปํ ทฏฺฐพฺพํ. ยถา เตสํ สนฺนิปาโต, เอวํ ผสฺโส ทฏฺฐพฺโพติ.}

\prose{ภิยฺโย โอปมฺมํ กโรหีติ. ยถา มหาราช เทฺว สมฺมา วชฺเชยฺยุํ, เตสุ ยถา เอโก สมฺโม, เอวํ จกฺขุ ทฏฺฐพฺพํ. ยถา ทุติโย สมฺโม, เอวํ รูปํ ทฏฺฐพฺพํ. ยถา เตสํ สนฺนิปาโต, เอวํ ผสฺโส ทฏฺฐพฺโพติ.}

\prosewithcloserruled{กลฺโลสิ ภนฺเต นาคเสนาติ.}{ผสฺส\-ลกฺขณ\-ปโญฺห อฏฺฐโม.}
\csromanmatikaanchor{mtk.40}
\csromantocmarkref{subsection}{9. เวทนาลกฺขณปญฺห}{mtk.40}
\bookmark[dest={mtk.40},level=2]{9. เวทนาลกฺขณปญฺห}
\csromanheader{2}{9. \textbf{เวทนา}\-\textbf{ลกฺขณ}\-\textbf{ปญฺห}}
\proseitem{9}{ภนฺเต นาคเสน กิํลกฺขณา เวทนาติ. เวทยิต\-ลกฺขณา มหาราช เวทนา อนุภวน\-ลกฺขณา จาติ.}

\prose{โอปมฺมํ กโรหีติ. ยถา มหาราช โกจิเทว ปุริโส รญฺโญ อธิการํ กเรยฺย, ตสฺส ราชา ตุฏฺโฐ อธิการํ ทเทยฺย, โส เตน อธิกาเรน ปญฺจหิ กามคุเณหิ สมปฺปิโต สมงฺคิภูโต ปริจเรยฺย, ตสฺส เอวมสฺส\csromandash{}มยา โข ปุพฺเพ รญฺโญ อธิกาโร กโต, ตสฺส \csromanfolio{62}เม ราชา ตุฏฺโฐ อธิการํ อทาสิ, สฺวาหํ ตโตนิทานํ อิมํ เอวรูปํ เวทนํ เวทยามีติ.}

\prose{ยถา วา ปน มหาราช โกจิเทว ปุริโส กุสลํ กมฺมํ กตฺวา กายสฺส เภทา ปรํ มรณา สุคติํ สคฺคํ โลกํ อุปปชฺเชยฺย, โส จ ตตฺถ ทิพฺเพหิ ปญฺจหิ กามคุเณหิ สมปฺปิโต สมงฺคิภูโต ปริจเรยฺย, ตสฺส เอวมสฺส\csromandash{}สฺวาหํ โข ปุพฺเพ กุสลํ กมฺมํ อกาสิํ, โสหํ ตโตนิทานํ อิมํ เอวรูปํ เวทนํ เวทยามีติ, เอวํ โข มหาราช เวทยิต\-ลกฺขณา เวทนา อนุภวน\-ลกฺขณา จาติ.}

\prosewithcloserruled{กลฺโลสิ ภนฺเต นาคเสนาติ.}{เวทนา\-ลกฺขณ\-ปโญฺห นวโม.}
\csromanmatikaanchor{mtk.41}
\csromantocmarkref{subsection}{10. สญฺญาลกฺขณปญฺห}{mtk.41}
\bookmark[dest={mtk.41},level=2]{10. สญฺญาลกฺขณปญฺห}
\csromanheader{2}{10. \textbf{สญฺญา}\-\textbf{ลกฺขณ}\-\textbf{ปญฺห}}
\proseitem{10}{ภนฺเต นาคเสน กิํลกฺขณา สญฺญาติ. สญฺชานน\-ลกฺขณา มหาราช สญฺญา. กิํ สญฺชานาติ, นีลมฺปิ สญฺชานาติ, ปีตมฺปิ สญฺชานาติ, โลหิตมฺปิ สญฺชานาติ, โอทาตมฺปิ สญฺชานาติ, มญฺชิฏฺฐมฺปิ\csromansharedfootnote{c938f1b2bb97cfd5}{มญฺเชฏฺฐมฺปิ (สี, อิ)} สญฺชานาติ, เอวํ โข มหาราช สญฺชานน\-ลกฺขณา สญฺญาติ.}

\prose{โอปมฺมํ กโรหีติ. ยถา มหาราช รญฺโญ ภณฺฑาคาริโก ภณฺฑาคารํ ปวิสิตฺวา นีลปีตโลหิโตทาต มญฺชิฏฺฐานิ\csromansharedfootnote{3482d74e9550b151}{มญฺเชฏฺฐานิ (สี, อิ)} ราชโภคานิ รูปานิ ปสฺสิตฺวา สญฺชานาติ. เอวํ โข มหาราช สญฺชานน\-ลกฺขณา สญฺญาติ.}

\prosewithcloserruled{กลฺโลสิ ภนฺเต นาคเสนาติ.}{สญฺญา\-ลกฺขณ\-ปโญฺห ทสโม.}
\csromanmatikaanchor{mtk.42}
\csromantocmarkref{subsection}{11. เจตนาลกฺขณปญฺห}{mtk.42}
\bookmark[dest={mtk.42},level=2]{11. เจตนาลกฺขณปญฺห}
\csromanheader{2}{11. \textbf{เจตนา}\-\textbf{ลกฺขณ}\-\textbf{ปญฺห}}
\proseitem{11}{ภนฺเต นาคเสน กิํลกฺขณา เจตนาติ. เจตยิต\-ลกฺขณา มหาราช เจตนา อภิสงฺขรณ\-ลกฺขณา จาติ.}

\prose{\csromanfolio{63}โอปมฺมํ กโรหีติ. ยถา มหาราช โกจิเทว ปุริโส วิสํ อภิส\-งฺขริ\-ตฺวา อตฺตนา จ ปิเวยฺย, ปเร จ ปาเยยฺย, โส อตฺตนาปิ ทุกฺขิโต ภเวยฺย, ปเรปิ ทุกฺขิตา ภเวยฺยุํ. เอวเมว โข มหาราช อิเธกจฺโจ ปุคฺคโล อกุสลํ กมฺมํ เจตนาย เจตยิตฺวา กายสฺส เภทา ปรํ มรณา อปายํ ทุคฺคติํ วินิปาตํ นิรยํ อุปปชฺเชยฺย, เยปิ ตสฺส อนุสิกฺขนฺติ, เตปิ กายสฺส เภทา ปรํ มรณา อปายํ ทุคฺคติํ วินิปาตํ นิรยํ อุปปชฺชนฺติ.}

\prose{ยถา วา ปน มหาราช โกจิเทว ปุริโส สปฺปิ\-นวนีต\-เตล\-มธุ\-ผาณิตํ เอกชฺฌํ อภิส\-งฺขริ\-ตฺวา อตฺตนา จ ปิเวยฺย, ปเร จ ปาเยยฺย, โส อตฺตนา สุขิโต ภเวยฺย, ปเรปิ สุขิตา ภเวยฺยุํ. เอวเมว โข มหาราช อิเธกจฺโจ ปุคฺคโล กุสลํ กมฺมํ เจตนาย เจตยิตฺวา กายสฺส เภทา ปรํ มรณา สุคติํ สคฺคํ โลกํ อุปปชฺชติ, เยปิ ตสฺส อนุสิกฺขนฺติ, เตปิ กายสฺส เภทา ปรํ มรณา สุคติํ สคฺคํ โลกํ อุปปชฺชนฺติ. เอวํ โข มหาราช เจตยิต\-ลกฺขณา เจตนา อภิสงฺขรณ\-ลกฺขณา จาติ.}

\prosewithcloserruled{กลฺโลสิ ภนฺเต นาคเสนาติ.}{เจตนา\-ลกฺขณ\-ปโญฺห เอกาทสโม.}
\csromanmatikaanchor{mtk.43}
\csromantocmarkref{subsection}{12. วิญฺญาณลกฺขณปญฺห}{mtk.43}
\bookmark[dest={mtk.43},level=2]{12. วิญฺญาณลกฺขณปญฺห}
\csromanheader{2}{12. \textbf{วิญฺญาณ}\-\textbf{ลกฺขณ}\-\textbf{ปญฺห}}
\proseitem{12}{ภนฺเต นาคเสน กิํลกฺขณํ วิญฺญาณนฺติ. วิชานน\-ลกฺขณํ มหาราช วิญฺญาณนฺติ.}

\prose{โอปมฺมํ กโรหีติ. ยถา มหาราช นครคุตฺติโก มชฺเฌ นคร\-สิงฺฆาฏเก นิสินฺโน ปสฺเสยฺย ปุรตฺถิม\-ทิสโต ปุริสํ อาคจฺฉนฺตํ, ปสฺเสยฺย ทกฺขิณทิสโต ปุริสํ อาคจฺฉนฺตํ, ปสฺเสยฺย ปจฺฉิมทิสโต ปุริสํ อาคจฺฉนฺตํ, ปสฺเสยฺย อุตฺตรทิสโต ปุริสํ อาคจฺฉนฺตํ. เอวเมว โข มหาราช ยญฺจ ปุริโส จกฺขุนา รูปํ ปสฺสติ, ตํ วิญฺญาเณน วิชานาติ. ยญฺจ โสเตน สทฺทํ สุณาติ, ตํ วิญฺญาเณน วิชานาติ. ยญฺจ ฆาเนน คนฺธํ ฆายติ, ตํ วิญฺญาเณน วิชานาติ. ยญฺจ ชิวฺหาย รสํ สายติ, ตํ \csromanfolio{64}วิญฺญาเณน วิชานาติ. ยญฺจ กาเยน โผฏฺฐพฺพํ ผุสติ, ตํ วิญฺญาเณน วิชานาติ. ยญฺจ มนสา ธมฺมํ วิชานาติ, ตํ วิญฺญาเณน วิชานาติ. เอวํ โข มหาราช วิชานน\-ลกฺขณํ วิญฺญาณนฺติ.}

\prosewithcloserruled{กลฺโลสิ ภนฺเต นาคเสนาติ.}{วิญฺญาณ\-ลกฺขณ\-ปโญฺห ทฺวาทสโม.}
\csromanmatikaanchor{mtk.44}
\csromantocmarkref{subsection}{13. วิตกฺกลกฺขณปญฺห}{mtk.44}
\bookmark[dest={mtk.44},level=2]{13. วิตกฺกลกฺขณปญฺห}
\csromanheader{2}{13. \textbf{วิตกฺก}\-\textbf{ลกฺขณ}\-\textbf{ปญฺห}}
\proseitem{13}{ภนฺเต นาคเสน กิํลกฺขโณ วิตกฺโกติ. อปฺปนาลกฺขโณ มหาราช วิตกฺโกติ.}

\prose{โอปมฺมํ กโรหีติ. ยถา มหาราช วฑฺฒกี สุปริกมฺมกตํ ทารุํ สนฺธิสฺมิํ อปฺเปติ, เอวเมว โข มหาราช อปฺปนาลกฺขโณ วิตกฺโกติ.}

\prosewithcloserruled{กลฺโลสิ ภนฺเต นาคเสนาติ.}{วิตกฺก\-ลกฺขณ\-ปโญฺห เตรสโม.}
\csromanmatikaanchor{mtk.45}
\csromantocmarkref{subsection}{14. วิจารลกฺขณปญฺห}{mtk.45}
\bookmark[dest={mtk.45},level=2]{14. วิจารลกฺขณปญฺห}
\csromanheader{2}{14. \textbf{วิจาร}\-\textbf{ลกฺขณ}\-\textbf{ปญฺห}}
\proseitem{14}{ภนฺเต นาคเสน กิํลกฺขโณ วิจาโรติ. อนุมชฺชน\-ลกฺขโณ มหาราช วิจาโรติ.}

\prose{โอปมฺมํ กโรหีติ. ยถา มหาราช กํสถาลํ อาโกฏิตํ ปจฺฉา อนุรวติ อนุสนฺทหติ\csromansharedfootnote{c9184a8b53a6e199}{อนุสทฺทายติ (ก)}, ยถา มหาราช อาโกฏนา, เอวํ วิตกฺโก ทฏฺฐพฺโพ. ยถา อนุรวนา\csromansharedfootnote{5dd015e0d0b417bb}{อนุมชฺชนา (ก)}, เอวํ วิจาโร ทฏฺฐพฺโพติ.}

\prose{กลฺโลสิ ภนฺเต นาคเสนาติ.}

\vspace{\tipitakaparskip}
\csromancenter{วิจาร\-ลกฺขณ\-ปโญฺห จุทฺทสโม.}
\nitthitammidruled{\textbf{วิจารวคฺโค ตติโย.}}
\csromancenter{อิมสฺมิํ วคฺเค จุทฺทส ปญฺหา.}
\csromanmatikaanchor{mtk.46}
\csromantocmarknopage{section}{4. นิพฺพานวคฺค}{mtk.46}
\bookmark[dest={mtk.46},level=1]{4. นิพฺพานวคฺค}
\csromanheader{1}{4. \textbf{นิพฺพานวคฺค}}
\csromanmatikaanchor{mtk.47}
\csromantocmarkref{subsection}{1. ผสฺสาทิวินิพฺภุชนปญฺห}{mtk.47}
\bookmark[dest={mtk.47},level=2]{1. ผสฺสาทิวินิพฺภุชนปญฺห}
\csromanheader{2}{1. \textbf{ผสฺสา}\-\textbf{ทิ}\-\textbf{วินิพฺภุช}\-\textbf{น}\-\textbf{ปญฺห}}
\proseitem{1}{\csromanfolio{65}ราชา อาห “ภนฺเต นาคเสน สกฺกา อิเมสํ ธมฺมานํ เอกโต\-ภาว\-คตานํ วินิพฺภุชิตฺวา วินิพฺภุชิตฺวา นานากรณํ ปญฺญาเปตุํ ‘อยํ ผสฺโส, อยํ เวทนา, อยํ สญฺญา, อยํ เจตนา, อิทํ วิญฺญาณํ, อยํ วิตกฺโก, อยํ วิจาโร’ติ”. น สกฺกา มหาราช อิเมสํ ธมฺมานํ เอกโต\-ภาว\-คตานํ วินิพฺภุชิตฺวา วินิพฺภุชิตฺวา นานากรณํ ปญฺญาเปตุํ “อยํ ผสฺโส, อยํ เวทนา, อยํ สญฺญา, อยํ เจตนา, อิทํ วิญฺญาณํ, อยํ วิตกฺโก, อยํ วิจาโร”ติ.}

\prose{โอปมฺมํ กโรหีติ. ยถา มหาราช รญฺโญ สูโท อรสํ วา รสํ วา\csromansharedfootnote{ffc7901ec4854063}{ยูสํ วา รสํ วา (สี, สฺยา, อิ)} กเรยฺย, โส ตตฺถ ทธิมฺปิ ปกฺขิเปยฺย, โลณมฺปิ ปกฺขิเปยฺย, สิงฺคิ เวรมฺปิ ปกฺขิเปยฺย, ชีรกมฺปิ ปกฺขิเปยฺย, มริจมฺปิ ปกฺขิเปยฺย, อญฺญานิปิ ปการานิ ปกฺขิเปยฺย, ตเมนํ ราชา เอวํ วเทยฺย, “ทธิสฺส เม รสํ อาหร, โลณสฺสเม รสํ อาหร, สิงฺคิเวรสฺส เม รสํ อาหร, ชีรกสฺส เม รสํ อาหร, มริจสฺส เม รสํ อาหร, สพฺเพสํ เม ปกฺขิตฺตานํ รสํ อาหรา”ติ. สกฺกา นุ โข มหาราช เตสํ รสานํ เอกโต\-ภาว\-คตานํ วินิพฺภุชิตฺวา วินิพฺภุชิตฺวา รสํ อาหริตุํ อมฺพิลตฺตํ วา ลวณตฺตํ วา ติตฺตกตฺตํ วา กฏุกตฺตํ วา กสายตฺตํ วา มธุรตฺตํ วาติ. น หิ ภนฺเต สกฺกา เตสํ รสานํ เอกโต\-ภาว\-คตานํ วินิพฺภุชิตฺวา วินิพฺภุชิตฺวา รสํ อาหริตุํ อมฺพิลตฺตํ วา ลวณตฺตํ วา ติตฺตกตฺตํ วา กฏุกตฺตํ วา กสายตฺตํ วา มธุรตฺตํ วา, อปิ จ โข ปน สเกน สเกน ลกฺขเณน อุปฏฺฐหนฺตีติ. เอวเมว โข มหาราช น สกฺกา อิเมสํ ธมฺมานํ เอกโต\-ภาว\-คตานํ วินิพฺภุชิตฺวา วินิพฺภุชิตฺวา นานากรณํ ปญฺญาเปตุํ “อยํ ผสฺโส, อยํ เวทนา, อยํ สญฺญา, อยํ เจตนา, อิทํ วิญฺญาณํ, อยํ วิตกฺโก, อยํ วิจาโร”ติ, อปิ จ โข ปน สเกน สเกน ลกฺขเณน อุปฏฺฐหนฺตีติ.}

\prose{กลฺโลสิ ภนฺเต นาคเสนาติ.}

\vspace{\tipitakaparskip}
\csromancenter{ผสฺสา\-ทิ\-วินิพฺภุช\-น\-ปโญฺห ปฐโม.}
\csromanmatikaanchor{mtk.48}
\csromantocmarkref{subsection}{2. นาคเสนปญฺห}{mtk.48}
\bookmark[dest={mtk.48},level=2]{2. นาคเสนปญฺห}
\csromanheader{2}{2. \textbf{นาคเสนปญฺห}}
\proseitem{2}{\csromanfolio{66}เถโร อาห “โลณํ มหาราช จกฺขุวิญฺเญยฺยนฺ”ติ. อาม ภนฺเต จกฺขุ\-วิญฺเญยฺยนฺติ. สุฏฺฐุ โข มหาราช ชานาหีติ. กิํ ปน ภนฺเต ชิวฺหา\-วิญฺเญยฺยนฺติ. อาม มหาราช ชิวฺหา\-วิญฺเญยฺยนฺติ. กิํ ปน ภนฺเต สพฺพํ โลณํ ชิวฺหาย วิชานาตีติ. อาม มหาราช สพฺพํ โลณํ ชิวฺหาย วิชานาติ.}

\prose{ยทิ ภนฺเต สพฺพํ โลณํ ชิวฺหาย วิชานาติ, กิสฺส ปน ตํ สกเฏหิ พลีพทฺทา\csromansharedfootnote{c69caa16d1a3357a}{พลิพทฺทา (สี, อิ)} อาหรนฺติ, นนุ โลณเมว อาหริตพฺพนฺติ. น สกฺกา มหาราช โลณเมว อาหริตุํ เอกโตภาวคตา เอเต ธมฺมา โคจร\-นานตฺต\-คตา โลณํ ครุภาโว จาติ. สกฺกา ปน มหาราช โลณํ ตุลาย ตุลยิตุนฺติ. อาม ภนฺเต สกฺกาติ. น สกฺกา มหาราช โลณํ ตุลาย ตุลยิตุํ, ครุภาโว ตุลาย ตุลิยตีติ.}

\prosewithcloserruled{กลฺโลสิ ภนฺเต นาคเสนาติ.}{นาคเสนปโญฺห ทุติโย.}
\csromanmatikaanchor{mtk.49}
\csromantocmarkref{subsection}{3. ปญฺจายตนกมฺมนิพฺพตฺตปญฺห}{mtk.49}
\bookmark[dest={mtk.49},level=2]{3. ปญฺจายตนกมฺมนิพฺพตฺตปญฺห}
\csromanheader{2}{3. \textbf{ปญฺจายตน}\-\textbf{กมฺมนิพฺพตฺต}\-\textbf{ปญฺห}}
\proseitem{3}{ราชา อาห “ภนฺเต นาคเสน ยานิมานิ ปญฺจายตนานิ, กินฺนุ ตานิ นานากมฺเมหิ นิพฺพตฺตานิ, อุทาหุ เอเกน กมฺเมนา”ติ. นานากมฺเมหิ มหาราช นิพฺพตฺตานิ, น เอเกน กมฺเมนาติ.}

\prose{โอปมฺมํ กโรหีติ. ตํ กิํ มญฺญสิ มหาราช, เอกสฺมิํ เขตฺเต นานาพีชานิ วปฺเปยฺยุํ, เตสํ นานาพีชานํ นานาผลานิ นิพฺพตฺเตยฺยุนฺติ. อาม ภนฺเต นิพฺพตฺเตยฺยุนฺติ. เอวเมว โข มหาราช ยานิ ยานิ ปญฺจายตนานิ, ตานิ ตานิ นานากมฺเมหิ นิพฺพตฺตานิ, น เอเกน กมฺเมนาติ.}

\prose{กลฺโลสิ ภนฺเต นาคเสนาติ.}

\vspace{\tipitakaparskip}
\csromancenter{ปญฺจายตน\-กมฺมนิพฺพตฺต\-ปโญฺห ตติโย.}
\csromanmatikaanchor{mtk.50}
\csromantocmarkref{subsection}{4. กมฺมนานากรณปญฺห}{mtk.50}
\bookmark[dest={mtk.50},level=2]{4. กมฺมนานากรณปญฺห}
\csromanmarkh{2. มิลินฺทปญฺห}\csromanmarkhii{4. กมฺม\-นานากรณ\-ปญฺห}\csromanheadernomark{2}{2. มิลินฺทปญฺห \textbf{4. กมฺม}\-\textbf{นานากรณ}\-\textbf{ปญฺห}}
\proseitem{4}{\csromanfolio{67}ราชา อาห “ภนฺเต นาคเสน เกน การเณน มนุสฺสา น สพฺเพ สมกา, อญฺเญ อปฺปายุกา, อญฺเญ ทีฆายุกา, อญฺเญ พหฺวาพาธา อญฺเญ อปฺปาพาธา, อญฺเญ ทุพฺพณฺณา, อญฺเญ วณฺณวนฺโต, อญฺเญ อปฺเปสกฺขา, อญฺเญ มเหสกฺขา, อญฺเญ อปฺปโภคา, อญฺเญ มหาโภคา, อญฺเญ นีจกุลีนา, อญฺเญ มหากุลีนา, อญฺเญ ทุปฺปญฺญา, อญฺเญ ปญฺญวนฺโต”ติ.}

\prose{เถโร อาห “กิสฺส ปน มหาราช รุกฺขา น สพฺเพ สมกา, อญฺเญ อมฺพิลา, อญฺเญ ลวณา, อญฺเญ ติตฺตกา, อญฺเญ กฏุกา, อญฺเญ กสาวา, อญฺเญ มธุรา”ติ. มญฺญามิ ภนฺเต พีชานํ นานากรเณนาติ. เอวเมว โข มหาราช กมฺมานํ นานากรเณน มนุสฺสา น สพฺเพ สมกา, อญฺเญ อปฺปายุกา, อญฺเญ ทีฆายุกา, อญฺเญ พหฺวาพาธา, อญฺเญ อปฺปาพาธา, อญฺเญ ทุพฺพณฺณา, อญฺเญ วณฺณวนฺโต, อญฺเญ อปฺเปสกฺขา, อญฺเญ มเหสกฺขา, อญฺเญ อปฺปโภคา, อญฺเญ มหาโภคา, อญฺเญ นีจกุลีนา, อญฺเญ มหากุลีนา, อญฺเญ ทุปฺปญฺญา, อญฺเญ ปญฺญวนฺโต, ภาสิตมฺเปตํ มหาราช ภควตา\csromandash{}“กมฺมสฺสกา มาณว สตฺตา กมฺมทายาทา กมฺมโยนี กมฺมพนฺธู กมฺม\-ปฺปฏิสรณา, กมฺมํ สตฺเต วิภชติ ยทิทํ หีนปฺปณีตตายา”ติ *.}

\prosewithcloserruled{กลฺโลสิ ภนฺเต นาคเสนาติ.}{กมฺม\-นานากรณ\-ปโญฺห จตุตฺโถ.}
\csromanmatikaanchor{mtk.51}
\csromantocmarkref{subsection}{5. วายามกรณปญฺห}{mtk.51}
\bookmark[dest={mtk.51},level=2]{5. วายามกรณปญฺห}
\csromanheader{2}{5. \textbf{วายาม}\-\textbf{กรณ}\-\textbf{ปญฺห}}
\proseitem{5}{ราชา อาห “ภนฺเต นาคเสน ตุเมฺห ภณถ ‘กินฺติ อิมํ ทุกฺขํ นิรุชฺเฌยฺย, อญฺญญฺจ ทุกฺขํ นุปฺปชฺเชยฺยา’ติ”. เอตทตฺถา มหาราช อมฺหากํ ปพฺพชฺชาติ. กิํ ปฏิกจฺเจว วายมิเตน, นนุ สมฺปตฺเต กาเล วายมิตพฺพนฺติ. เถโร อาห “สมฺปตฺเต กาเล มหาราช วายาโม อกิจฺจกโร ภวติ, ปฏิกจฺเจว วายาโม กิจฺจกโร ภวตี”ติ.}

\prose{โอปมฺมํ กโรหีติ. ตํ กิํ มญฺญสิ มหาราช, ยทา ตฺวํ ปิปาสิโต ภเวยฺยาสิ, ตทา ตฺวํ อุทปานํ ขณาเปยฺยาสิ, ตฬากํ ขณาเปยฺยาสิ “ปานียํ ปิวิสฺสามี”ติ. น หิ ภนฺเตติ. เอวเมว โข มหาราช สมฺปตฺเต \csromanfolio{68}กาเล วายาโม อกิจฺจกโร ภวติ, ปฏิกจฺเจว วายาโม กิจฺจกโร ภวตีติ.}

\prose{ภิยฺโย โอปมฺมํ กโรหีติ. ตํ กิํ มญฺญสิ มหาราช, ยทา ตฺวํ พุภุกฺขิโต ภเวยฺยาสิ, ตทา ตฺวํ เขตฺตํ กสาเปยฺยาสิ, สาลิํ โรปาเปยฺยาสิ, ธญฺญํ อติหราเปยฺยาสิ “ภตฺตํ ภุญฺชิสฺสามี”ติ. น หิ ภนฺเตติ. เอวเมว โข มหาราช สมฺปตฺเต กาเล วายาโม อกิจฺจกโร ภวติ, ปฏิกจฺเจว วายาโม กิจฺจกโร ภวตีติ.}

\prose{ภิยฺโย โอปมฺมํ กโรหีติ. ตํ กิํ มญฺญสิ มหาราช, ยทา เต สงฺคาโม ปจฺจุปฏฺฐิโต ภเวยฺย, ตทา ตฺวํ ปริขํ ขณาเปยฺยาสิ, ปาการํ การาเปยฺยาสิ, โคปุรํ การาเปยฺยาสิ, อฏฺฏาลกํ การาเปยฺยาสิ, ธญฺญํ อติหราเปยฺยาสิ, ตทา ตฺวํ หตฺถิสฺมิํ สิกฺเขยฺยาสิ, อสฺสสฺมิํ สิกฺเขยฺยาสิ, รถสฺมิํ สิกฺเขยฺยาสิ, ธนุสฺมิํ สิกฺเขยฺยาสิ, ถรุสฺมิํ สิกฺเขยฺยาสีติ. น หิ ภนฺเตติ. เอวเมว โข มหาราช สมฺปตฺเต กาเล วายาโม อกิจฺจกโร ภวติ, ปฏิกจฺเจว วายาโม กิจฺจกโร ภวติ, ภาสิตมฺเปตํ มหาราช ภควตา\csromandash{}}

\setlength{\csromangathapagewidth}{0pt}
\setlength{\csromangathawakwidth}{0pt}
\csromangathameasuregroup{“ปฏิกจฺเจว ตํ กยิรา, \\ น สากฏิกจินฺตาย, \\ ยถ สากฏิโก มฏฺฐํ\textsuperscript{1}, \\ วิสมํ มคฺคมารุยฺห, \\ เอวํ ธมฺมา อปกฺกมฺม, \\ มนฺโท มจฺจุ มุขํ ปตฺโต,}{\csromangathabat{“ปฏิกจฺเจว ตํ กยิรา,}{ยํ ชญฺญา หิตมตฺตโน.} \\ \csromangathabat{น สากฏิกจินฺตาย,}{มนฺตา ธีโร ปรกฺกเม.} \\ \csromangathabat{ยถ สากฏิโก มฏฺฐํ\textsuperscript{1},}{สมํ หิตฺวา มหาปถํ.} \\ \csromangathabat{วิสมํ มคฺคมารุยฺห,}{อกฺขจฺฉินฺโนว ฌายติ.} \\ \csromangathabat{เอวํ ธมฺมา อปกฺกมฺม,}{อธมฺมมนุวตฺติย.} \\ \csromangathabat{มนฺโท มจฺจุ มุขํ ปตฺโต,}{อกฺขจฺฉินฺโนว ฌายตี”ติ\textsuperscript{1}.}}
\csromangathasetleft{“ปฏิกจฺเจว ตํ กยิรา, \\ น สากฏิกจินฺตาย, \\ ยถ สากฏิโก มฏฺฐํ\textsuperscript{1}, \\ วิสมํ มคฺคมารุยฺห, \\ เอวํ ธมฺมา อปกฺกมฺม, \\ มนฺโท มจฺจุ มุขํ ปตฺโต,}
\csromangathagroup{\csromangathastanza{\csromangathabat{“ปฏิกจฺเจว ตํ กยิรา,}{ยํ ชญฺญา หิตมตฺตโน.} \\ \csromangathabat{น สากฏิกจินฺตาย,}{มนฺตา ธีโร ปรกฺกเม.}}\csromangathastanza{\csromangathabat{ยถ สากฏิโก มฏฺฐํ\csromansharedfootnote{3d047f6c41463169}{นาม (สี, อิ, ก) สํ 1. 56 ปิฏฺเฐ ปสฺสิตพฺพํ.},}{สมํ หิตฺวา มหาปถํ.} \\ \csromangathabat{วิสมํ มคฺคมารุยฺห,}{อกฺขจฺฉินฺโนว ฌายติ.}}\csromangathastanza{\csromangathabat{เอวํ ธมฺมา อปกฺกมฺม,}{อธมฺมมนุวตฺติย.} \\ \csromangathabat{มนฺโท มจฺจุ มุขํ ปตฺโต,}{อกฺขจฺฉินฺโนว ฌายตี”ติ\csromansharedfootnote{35fe6c2350ae53d2}{โสจตีติ (สพฺพตฺถ)}.}}}

\prosewithcloserruled{กลฺโลสิ ภนฺเต นาคเสนาติ.}{วายาม\-กรณ\-ปโญฺห ปญฺจโม.}
\csromanmatikaanchor{mtk.52}
\csromantocmarkref{subsection}{6. เนรยิกคฺคิอุณฺหภาวปญฺห}{mtk.52}
\bookmark[dest={mtk.52},level=2]{6. เนรยิกคฺคิอุณฺหภาวปญฺห}
\csromanheader{2}{6. เนรยิกคฺคิ\-อุณฺห\-ภาว\-ปญฺห}
\proseitem{5}{ราชา อาห “ภนฺเต นาคเสน ตุเมฺห ภณถ ‘ปากติกอคฺคิโต เนรยิโก อคฺคิ มหาภิตาป\-ตโร โหติ, ขุทฺทโกปิ \csromanfolio{69}ปาสาโณ ปากติเก อคฺคิมฺหิ ปกฺขิตฺโต ทิวสมฺปิ ปจฺจมาโน\csromansharedfootnote{8ec659811437eb4c}{ธมามาโน (สี, อิ)} น วิลยํ คจฺฉติ, กูฏาคาร\-มตฺโตปิ ปาสาโณ เนรยิกคฺคิมฺหิ ปกฺขิตฺโต ขเณน วิลยํ คจฺฉตี’ติ”, เอตํ วจนํ น สทฺทหามิ, เอวญฺจ ปน วเทถ “เย จ ตตฺถ อุปฺปนฺนา สตฺตา, เต อเนกานิปิ วสฺสสหสฺสานิ นิรเย ปจฺจมานา น วิลยํ คจฺฉนฺตีติ, ตมฺปิ วจนํ น สทฺทหามีติ.}

\proseitem{6}{เถโร อาห “ตํ กิํ มญฺญสิ มหาราช, ยา ตา สนฺติ มกรินิโยปิ สุสุมารินิโยปิ กจฺฉปินิโยปิ โมรินิโยปิ กโปตินิโยปิ, กินฺนุ ตา กกฺขฬานิ ปาสาณานิ สกฺขราโย จ ขาทนฺตี”ติ. อาม ภนฺเต ขาทนฺตีติ. กิํ ปน ตานิ ตาสํ กุจฺฉิยํ โกฏฺฐ\-พฺภนฺตรคตานิ วิลยํ คจฺฉนฺตีติ. อาม ภนฺเต วิลยํ คจฺฉนฺตีติ. โย ปน ตาสํ กุจฺฉิยํ คพฺโภ, โสปิ วิลยํ คจฺฉตีติ. น หิ ภนฺเตติ. เกน การเณนาติ. มญฺญามิ ภนฺเต กมฺมาธิกเตน น วิลยํ คจฺฉตีติ. เอวเมว โข มหาราช กมฺมาธิกเตน เนรยิกา สตฺตา อเนกานิปิ วสฺสสหสฺสานิ นิรเย ปจฺจมานาน วิลยํ คจฺฉนฺติ. ภาสิตมฺเปตํ มหาราช ภควตา\csromandash{} *“โส น ตาว กาลํ กโรติ, ยาว น ตํ ปาปกมฺมํ พฺยนฺตีโหตี”ติ.}

\prose{ภิยฺโย โอปมฺมํ กโรหีติ. ตํ กิํ มญฺญสิ มหาราช, ยา ตา สนฺติ สีหินิโยปิ พฺยคฺฆินิโยปิ ทีปินิโยปิ กุกฺกุรินิโยปิ. กินฺนุ ตา กกฺขฬานิ อฏฺฐิกานิ มํสานิ ขาทนฺตีติ, อาม ภนฺเต ขาทนฺตีติ. กิํ ปน ตานิ ตาสํ กุจฺฉิยํ โกฏฺฐ\-พฺภนฺตรคตานิ วิลยํ คจฺฉนฺตีติ. อาม ภนฺเต วิลยํ คจฺฉนฺตีติ. โย ปน ตาสํ กุจฺฉิยํ คพฺโภ, โสปิ วิลยํ คจฺฉตีติ. น หิ ภนฺเตติ. เกน การเณนาติ. มญฺญามิ ภนฺเต กมฺมาธิกเตน น วิลยํ คจฺฉตีติ. เอวเมว โข มหาราช กมฺมาธิกเตน เนรยิกา สตฺตา อเนกานิปิ วสฺสสหสฺสานิ นิรเย ปจฺจมานา น วิลยํ คจฺฉนฺตีติ.}

\prose{ภิยฺโย โอปมฺมํ กโรหีติ. ตํ กิํ มญฺญสิ มหาราช ยา ตา สนฺติ โยนกสุขุมาลินิโยปิ ขตฺติยสุขุมาลินิโยปิ พฺราหฺมณสุขุมาลินิโยปิ คหปติสุขุมาลินิโยปิ. กินฺนุ ตา กกฺขฬานิ ขชฺชกานิ มํสานิ ขาทนฺตีติ. อาม ภนฺเต ขาทนฺตีติ. กิํ ปน ตานิ ตาสํ กุจฺฉิยํ โกฏฺฐ\-พฺภนฺตรคตานิ วิลยํ คจฺฉนฺตีติ. อาม ภนฺเต วิลยํ คจฺฉนฺตีติ. \csromanfolio{70}โย ปน ตาสํ กุจฺฉิยํ คพฺโภ โสปิ วิลยํ คจฺฉตีติ. น หิ ภนฺเตติ. เกน การเณนาติ. มญฺญามิ ภนฺเต กมฺมาธิกเตน น วิลยํ คจฺฉตีติ. เอวเมว โข มหาราช กมฺมาธิกเตน เนรยิกา สตฺตา อเนกานิปิ วสฺสสหสฺสานิ นิรเย ปจฺจมานา น วิลยํ คจฺฉนฺติ. ภาสิตมฺเปตํ มหาราช ภควตา\csromandash{}* “โส น ตาว กาลํ กโรติ, ยาว น ตํ ปาปกมฺมํ พฺยนฺตีโหตี”ติ.}

\proseitem{2}{กลฺโลสิ ภนฺเต นาคเสนาติ.}

\prose{เนรยิกคฺคิ\-อุณฺห\-ภาว\-ปโญฺห ฉฏฺโฐ.}
\csromansectionrule

\csromanmatikaanchor{mtk.53}
\csromantocmarkref{subsection}{7. ปถวิสนฺธารกปญฺห}{mtk.53}
\bookmark[dest={mtk.53},level=2]{7. ปถวิสนฺธารกปญฺห}
\csromanheader{2}{7. \textbf{ปถวิ}\-\textbf{สนฺธารก}\-\textbf{ปญฺห}}
\proseitem{7}{ราชา อาห “ภนฺเต นาคเสน ตุเมฺห ภณถ ‘อยํ มหาปถวี อุทเก ปติฏฺฐิตา, อุทกํ วาเต ปติฏฺฐิตํ, วาโต อากาเส ปติฏฺฐิโต’ติ, เอตมฺปิ วจนํ น สทฺทหามี”ติ. เถโร ธมฺมกรเกน\csromansharedfootnote{6df7f2d97dc1ea6d}{ธมฺมกรเณน (ก)} อุทกํ คเหตฺวา ราชานํ มิลินฺทํ สญฺญาเปสิ “ยถา มหาราช อิมํ อุทกํ วาเตน อาธาริตํ, เอวํ ตมฺปิ อุทกํ วาเตน อาธาริตนฺ”ติ.}

\prosewithcloserruled{กลฺโลสิ ภนฺเต นาคเสนาติ.}{ปถวิ\-สนฺธารก\-ปโญฺห สตฺตโม.}
\csromanmatikaanchor{mtk.54}
\csromantocmarkref{subsection}{8. นิโรธนิพฺพานปญฺห}{mtk.54}
\bookmark[dest={mtk.54},level=2]{8. นิโรธนิพฺพานปญฺห}
\csromanheader{2}{8. \textbf{นิโรธ}\-\textbf{นิพฺพาน}\-\textbf{ปญฺห}}
\proseitem{8}{ราชา อาห “ภนฺเต นาคเสน นิโรโธ นิพฺพานนฺ”ติ. อาม มหาราช นิโรโธ นิพฺพานนฺติ. กถํ ภนฺเต นาคเสน นิโรโธ นิพฺพานนฺติ. สพฺเพ พาลปุถุชฺชนา โข มหาราช อชฺฌตฺติก\-พาหิเร อายตเน อภินนฺทนฺติ อภิวทนฺติ อชฺโฌสาย ติฏฺฐนฺติ, เต เตน โสเตน วุยฺหนฺติ, น ปริมุจฺจนฺติ ชาติยา ชราย มรเณน โสเกน ปริเทเวน ทุกฺเขหิ โทมนสฺเสหิ อุปายาเสหิ น ปริมุจฺจนฺติ ทุกฺขสฺมาติ วทามิ. สุตวา จ โข มหาราช อริยสาวโก อชฺฌตฺติก\-พาหิเร \csromanfolio{71}อายตเน นาภินนฺทติ นาภิวทติ นาชฺโฌสาย ติฏฺฐติ, ตสฺส ตํ อนภินนฺทโต อนภิวทโต อนชฺโฌสาย ติฏฺฐโต ตณฺหา นิรุชฺฌติ, ตณฺหานิโรธา อุปาทานนิโรโธ, อุปาทานนิโรธา ภวนิโรโธ, ภวนิโรธา ชาตินิโรโธ, ชาตินิโรธา ชรามรณํ โสก\-ปริเทว\-ทุกฺข\-โทมนสฺสุ\-ปายาสา นิรุชฺฌนฺติ, เอวเมตสฺส เกวลสฺส ทุกฺข\-กฺขนฺธสฺส นิโรโธ โหติ, เอวํ โข มหาราช นิโรโธ นิพฺพานนฺติ.}

\prosewithcloserruled{กลฺโลสิ ภนฺเต นาคเสนาติ.}{นิโรธ\-นิพฺพาน\-ปโญฺห อฏฺฐโม.}
\csromanmatikaanchor{mtk.55}
\csromantocmarkref{subsection}{9. นิพฺพานลภนปญฺห}{mtk.55}
\bookmark[dest={mtk.55},level=2]{9. นิพฺพานลภนปญฺห}
\csromanheader{2}{9. \textbf{นิพฺพาน}\-\textbf{ลภน}\-\textbf{ปญฺห}}
\proseitem{9}{ราชา อาห “ภนฺเต นาคเสน สพฺเพว ลภนฺติ นิพฺพานนฺ”ติ. น โข มหาราช สพฺเพว ลภนฺติ นิพฺพานํ, อปิ จ โข มหาราช โย สมฺมา ปฏิปนฺโน อภิญฺเญยฺเย ธมฺเม อภิชานาติ, ปริญฺเญยฺเย ธมฺเม ปริชานาติ, ปหาตพฺเพ ธมฺเม ปชหติ, ภาเวตพฺเพ ธมฺเม ภาเวติ, สจฺฉิกาตพฺเพ ธมฺเม สจฺฉิกโรติ, โส ลภติ นิพฺพานนฺติ.}

\prosewithcloserruled{กลฺโลสิ ภนฺเต นาคเสนาติ.}{นิพฺพาน\-ลภน\-ปโญฺห นวโม.}
\csromanmatikaanchor{mtk.56}
\csromantocmarkref{subsection}{10. นิพฺพานสุขชานนปญฺห}{mtk.56}
\bookmark[dest={mtk.56},level=2]{10. นิพฺพานสุขชานนปญฺห}
\csromanheader{2}{10. \textbf{นิพฺพานสุข}\-\textbf{ชานน}\-\textbf{ปญฺห}}
\proseitem{10}{ราชา อาห “ภนฺเต นาคเสน โย น ลภติ นิพฺพานํ, ชานาติ โส “สุขํ นิพฺพานนฺ”ติ. อาม มหาราช โย น ลภติ นิพฺพานํ, ชานาติ โส “สุขํ นิพฺพานนฺ”ติ. กถํ ภนฺเต นาคเสน อลภนฺโต ชานาติ “สุขํ นิพฺพานนฺ”ติ. ตํ กิํ มญฺญสิ มหาราช, เยสํ นจฺฉินฺนา \csromanfolio{72}หตฺถปาทา, ชาเนยฺยุํ เต มหาราช “ทุกฺขํ หตฺถปาทจฺเฉทนนฺ”ติ. อาม ภนฺเต ชาเนยฺยุนฺติ. กถํ ชานายฺยุนฺติ. อญฺเญสํ ภนฺเต \textbf{ฉินฺน}\-\textbf{หตฺถปาทานํ ปริเทวิต}\-\textbf{สทฺทํ สุตฺวา ชานนฺติ “ทุกฺขํ} หตฺถปาทจฺเฉทนนฺ”ติ. เอวเมว โข มหาราช เยสํ ทิฏฺฐํ นิพฺพานํ, เตสํ สทฺทํ สุตฺวา ชานาติ “สุขํ นิพฺพานนฺ”ติ.}

\prose{กลฺโลสิ ภนฺเต นาคเสนาติ.}

\vspace{\tipitakaparskip}
\csromancenter{นิพฺพานสุข\-ชานน\-ปโญฺห ทสโม.}
\nitthitammidruled{\textbf{นิพฺพานวคฺโค จตุตฺโถ.}}
\csromancenter{อิมสฺมิํ วคฺเค ทส ปญฺหา.}
\chapterheadpageread{5. \textbf{พุทฺธวคฺค}}
\csromanmatikaanchor{mtk.57}
\csromanmatikaanchor{mtk.58}
\csromantocmarknopage{section}{5. พุทฺธวคฺค}{mtk.57}
\bookmark[dest={mtk.57},level=1]{5. พุทฺธวคฺค}
\csromantocmarkref{subsection}{1. พุทฺธสฺส อตฺถินตฺถิภาวปญฺห}{mtk.58}
\bookmark[dest={mtk.58},level=2]{1. พุทฺธสฺส อตฺถินตฺถิภาวปญฺห}
\csromanheader{2}{1. \textbf{พุทฺธสฺส อตฺถินตฺถิภาวปญฺห}}
\proseitem{1}{\csromanfolio{73}ราชา อาห “ภนฺเต นาคเสน พุทฺโธ ตยา ทิฏฺโฐ”ติ. น หิ มหาราชาติ. อถ เต อาจริเยหิ พุทฺโธ ทิฏฺโฐติ. น หิ มหาราชาติ. เตน หิ ภนฺเต นาคเสน นตฺถิ พุทฺโธติ. กิํ ปน มหาราช หิมวติ อูหา นที ตยา ทิฏฺฐาติ. น หิ ภนฺเตติ. อถ เต ปิตรา อูหา นที ทิฏฺฐาติ. น หิ ภนฺเตติ. เตน หิ มหาราช นตฺถิ อูหา นทีติ. อตฺถิ ภนฺเต กิญฺจาปิ มยา อูหา นที น ทิฏฺฐา, ปิตราปิ เม อูหา นที น ทิฏฺฐา, อปิ จ อตฺถิ อูหา นทีติ. เอวเมว โข มหาราช กิญฺจาปิ มยา ภควา น ทิฏฺโฐ, อาจริเยหิปิ เม ภควา น ทิฏฺโฐ, อปิ จ อตฺถิ ภควาติ.}

\prosewithcloserruled{กลฺโลสิ ภนฺเต นาคเสนาติ.}{พุทฺธสฺส อตฺถิ\-นตฺถิภาว\-ปโญฺห ปฐโม.}
\csromanmatikaanchor{mtk.59}
\csromantocmarkref{subsection}{2. พุทฺธสฺส อนุตฺตรภาวปญฺห}{mtk.59}
\bookmark[dest={mtk.59},level=2]{2. พุทฺธสฺส อนุตฺตรภาวปญฺห}
\csromanheader{2}{2. \textbf{พุทฺธสฺส อนุตฺตรภาวปญฺห}}
\proseitem{2}{ราชา อาห “ภนฺเต นาคเสน พุทฺโธ อนุตฺตโร”ติ. อาม มหาราช ภควา อนุตฺตโรติ. กถํ ภนฺเต นาคเสน อทิฏฺฐปุพฺพํ ชานาสิ “พุทฺโธ อนุตฺตโร”ติ. ตํ กิํ มญฺญสิ มหาราช, เยหิ อทิฏฺฐปุพฺโพ มหาสมุทฺโท, ชาเนยฺยุํ เต มหาราช มหนฺโต โข มหาสมุทฺโท คมฺภีโร อปฺปเมยฺโย ทุปฺปริโยคาโห, ยตฺถิมา ปญฺจ มหานทิโย สตตํ สมิตํ อปฺเปนฺติ, เสยฺยถีทํ, คงฺคา ยมุนา อจิรวตี สรภู มหี, เนว ตสฺส อูนตฺตํ วา ปูรตฺตํ วา ปญฺญายตีติ. อาม ภนฺเต ชาเนยฺยุนฺติ. เอวเมว โข มหาราช สาวเก มหนฺเต ปรินิพฺพุเต ปสฺสิตฺวา ชานามิ “ภควา อนุตฺตโร”ติ.}

\prose{กลฺโลสิ ภนฺเต นาคเสนาติ.}

\vspace{\tipitakaparskip}
\csromancenter{พุทฺธสฺส อนุตฺตร\-ภาว\-ปโญฺห ทุติโย.}
\csromanmatikaanchor{mtk.60}
\csromantocmarkref{subsection}{3. พุทฺธสฺส อนุตฺตรภาวชานนปญฺห}{mtk.60}
\bookmark[dest={mtk.60},level=2]{3. พุทฺธสฺส อนุตฺตรภาวชานนปญฺห}
\csromanheader{2}{3. \textbf{พุทฺธสฺส อนุตฺตรภาวชานนปญฺห}}
\proseitem{3}{\csromanfolio{74}ราชา อาห “ภนฺเต นาคเสน สกฺกา ชานิตุํ “พุทฺโธ อนุตฺตโร”ติ. อาม มหาราช สกฺกา ชานิตุํ “ภควา อนุตฺตโร”ติ. กถํ ภนฺเต นาคเสน สกฺกา ชานิตุํ “พุทฺโธ อนุตฺตโร”ติ. ภูตปุพฺพํ มหาราช ติสฺสตฺเถโร นาม เลขาจริโย อโหสิ, พหูนิ วสฺสานิ อพฺภตีตานิ กาลงฺกตสฺส กถํ โส ญายตีติ. เลเขน ภนฺเตติ. เอวเมว โข มหาราช โย ธมฺมํ ปสฺสติ, โส ภควนฺตํ ปสฺสติ, ธมฺโม หิ มหาราช ภควตา เทสิโตติ.}

\prosewithcloserruled{กลฺโลสิ ภนฺเต นาคเสนาติ.}{พุทฺธสฺส อนุตฺตร\-ภาว\-ชานน\-ปโญฺห ตติโย.}
\csromanmatikaanchor{mtk.61}
\csromantocmarkref{subsection}{4. ธมฺมทิฏฺฐปญฺห}{mtk.61}
\bookmark[dest={mtk.61},level=2]{4. ธมฺมทิฏฺฐปญฺห}
\csromanheader{2}{4. \textbf{ธมฺมทิฏฺฐ}\-\textbf{ปญฺห}}
\proseitem{4}{ราชา อาห “ภนฺเต นาคเสน ธมฺโม ตยา ทิฏฺโฐ”ติ. พุทฺธเนตฺติยา โข มหาราช พุทฺธ\-ปญฺญตฺติยา ยาวชีวํ สาวเกหิ วตฺติตพฺพนฺติ.}

\prosewithcloserruled{กลฺโลสิ ภนฺเต นาคเสนาติ.}{ธมฺมทิฏฺฐ\-ปโญฺห จตุตฺโถ.}
\csromanmatikaanchor{mtk.62}
\csromantocmarkref{subsection}{5. อสงฺกมนปฏิสนฺทหนปญฺห}{mtk.62}
\bookmark[dest={mtk.62},level=2]{5. อสงฺกมนปฏิสนฺทหนปญฺห}
\csromanheader{2}{5. \textbf{อสงฺกมนปฏิสนฺทหนปญฺห}}
\proseitem{5}{ราชา อาห “ภนฺเต นาคเสน น จ สงฺกมติ ปฏิสนฺทหติ จา”ติ. อาม มหาราช น จ สงฺกมติ ปฏิสนฺทหติ จาติ. กถํ ภนฺเต นาคเสน น จ สงฺกมติ ปฏิสนฺทหติ จ, โอปมฺมํ กโรหีติ. ยถา มหาราช โกจิเทว ปุริโส ปทีปโต ปทีปํ ปทีเปยฺย, กินฺนุ โข โส มหาราช ปทีโป ปทีปมฺหา สงฺกนฺโตติ. น หิ ภนฺเตติ. เอวเมว โข มหาราช น จ สงฺกมติ ปฏิสนฺทหติ จาติ.}

\prose{ภิยฺโย โอปมฺมํ กโรหีติ. อภิชานาสิ นุ ตฺวํ มหาราช ทหรโก สนฺโต สิโลกา\-จริยสฺส สนฺติเก กิญฺจิ สิโลกํ คหิตนฺติ. อาม \csromanfolio{75}ภนฺเตติ. กิํ นุ โข มหาราช โส สิโลโก อาจริยมฺหา สงฺกนฺโตติ. น หิ ภนฺเตติ. เอวเมว โข มหาราช น จ สงฺกมติ ปฏิสนฺทหติ จาติ.}

\prosewithcloserruled{กลฺโลสิ ภนฺเต นาคเสนาติ.}{อสงฺกมนปฏิสนฺทหนปโญฺห ปญฺจโม.}
\csromanmatikaanchor{mtk.63}
\csromantocmarkref{subsection}{6. เวทคูปญฺห}{mtk.63}
\bookmark[dest={mtk.63},level=2]{6. เวทคูปญฺห}
\csromanheader{2}{6. \textbf{เวทคูปญฺห}}
\proseitem{5}{ราชา อาห “ภนฺเต นาคเสน เวทคู อุปลพฺภตี”ติ. เถโร อาห “ปรมตฺเถน โข มหาราช เวทคู นุปลพฺภตีติ.}

\prosewithcloserruled{กลฺโลสิ ภนฺเต นาคเสนาติ.}{เวทคูปโญฺห ฉฏฺโฐ.}
\csromanmatikaanchor{mtk.64}
\csromantocmarkref{subsection}{7. อญฺญกายสงฺกมนปญฺห}{mtk.64}
\bookmark[dest={mtk.64},level=2]{7. อญฺญกายสงฺกมนปญฺห}
\csromanheader{2}{7. \textbf{อญฺญ}\-\textbf{กาย}\-\textbf{สงฺกมน}\-\textbf{ปญฺห}}
\proseitem{7}{ราชา อาห “ภนฺเต นาคเสน อตฺถิ โกจิ สตฺโต, โย อิมมฺหา กายา อญฺญํ กายํ สงฺกมตี”ติ. น หิ มหาราชาติ. ยทิ ภนฺเต นาคเสน อิมมฺหา กายา อญฺญํ กายํ สงฺกมนฺโต นตฺถิ, นนุ มุตฺโต ภวิสฺสติ ปาปเกหิ กมฺเมหีติ. อาม มหาราช, ยทิ น ปฏิสนฺทเหยฺย, มุตฺโต ภวิสฺสติ ปาปเกหิ กมฺเมหีติ, ยสฺมา จ โข มหาราช ปฏิสนฺทหติ, ตสฺมา น ปริมุตฺโต ปาปเกหิ กมฺเมหีติ.}

\prose{โอปมฺมํ กโรหีติ. ยถา มหาราช โกจิเทว ปุริโส อญฺญตรสฺส ปุริสสฺส อมฺพํ อวหเรยฺย, กิํ โส ทณฺฑปฺปตฺโต ภเวยฺยาติ. อาม ภนฺเต ทณฺฑปฺปตฺโต ภเวยฺยาติ. น โข โส มหาราช ตานิ อมฺพานิ อวหริ, ยานิ เตน โรปิตานิ, กสฺมา ทณฺฑปฺปตฺโต ภเวยฺยาติ. ตานิ ภนฺเต อมฺพานิ นิสฺสาย ชาตานิ, ตสฺมา ทณฺฑปฺปตฺโต ภเวยฺยาติ. เอวเมว โข มหาราช อิมินา นามรูเปน กมฺมํ กโรติ โสภนํ วา \csromanfolio{76}อโสภนํ วา, เตน กมฺเมน อญฺญํ นามรูปํ ปฏิสนฺทหติ, ตสฺมา น ปริมุตฺโต ปาปเกหิ กมฺเมหีติ.}

\prosewithcloserruled{กลฺโลสิ ภนฺเต นาคเสนาติ.}{อญฺญ\-กาย\-สงฺกมน\-ปโญฺห สตฺตโม.}
\csromanmatikaanchor{mtk.65}
\csromantocmarkref{subsection}{8. กมฺมผลอตฺถิภาวปญฺห}{mtk.65}
\bookmark[dest={mtk.65},level=2]{8. กมฺมผลอตฺถิภาวปญฺห}
\csromanheader{2}{8. กมฺมผล\-อตฺถิภาว\-ปญฺห}
\proseitem{8}{ราชา อาห “ภนฺเต นาคเสน อิมินา นามรูเปน กมฺมํ กตํ กุสลํ วา อกุสลํ วา, กุหิํ ตานิ กมฺมานิ ติฏฺฐนฺตี”ติ. อนุพนฺเธยฺยุํ โข มหาราช ตานิ กมฺมานิ ฉายาว อนปายินีติ\csromansharedfootnote{69a6b70bfe1eaa59}{อนุปายินีติ (ก)}. สกฺกา ปน ภนฺเต ตานิ กมฺมานิ ทสฺเสตุํ “อิธ วา อิธ วา ตานิ กมฺมานิ ติฏฺฐนฺตี”ติ. น สกฺกา มหาราช ตานิ กมฺมานิ ทสฺเสตุํ “อิธ วา อิธ วา ตานิ กมฺมานิ ติฏฺฐนฺตี”ติ.}

\prose{โอปมฺมํ กโรหีติ. ตํ กิํ มญฺญสิ มหาราช, ยานิมานิ รุกฺขานิ อนิพฺพตฺตผลานิ, สกฺกา เตสํ ผลานิ ทสฺเสตุํ “อิธ วา อิธ วา ตานิ ผลานิ ติฏฺฐนฺตี”ติ. น หิ ภนฺเตติ. เอวเมว โข มหาราช อพฺโพจฺฉินฺนาย สนฺตติยา น สกฺกา ตานิ กมฺมานิ ทสฺเสตุํ “อิธ วา อิธ วา ตานิ กมฺมานิ ติฏฺฐนฺตี”ติ.}

\prose{กลฺโลสิ ภนฺเต นาคเสนาติ.}

\prose{กมฺมผล\-อตฺถิภาว\-ปโญฺห อฏฺฐโม.}
\csromansectionrule

\csromanmatikaanchor{mtk.66}
\csromantocmarkref{subsection}{9. อุปฺปชฺชติชานนปญฺห}{mtk.66}
\bookmark[dest={mtk.66},level=2]{9. อุปฺปชฺชติชานนปญฺห}
\csromanheader{2}{9. \textbf{อุปฺปชฺชติ}\-\textbf{ชานน}\-\textbf{ปญฺห}}
\proseitem{9}{ราชา อาห “ภนฺเต นาคเสน โย อุปฺปชฺชติ, ชานาติ โส “อุปฺปชฺชิสฺสามี”ติ. อาม มหาราช โย อุปฺปชฺชติ ชานาติ โส “อุปฺปชฺชิสฺสามี”ติ. โอปมฺมํ กโรหีติ. ยถา มหาราช กสฺสโก คหปติโก พีชานิ ปถวิยํ นิกฺขิปิตฺวา สมฺมา เทเว วสฺสนฺเต ชานาติ \csromanfolio{77}“ธญฺญํ นิพฺพตฺติสฺสตี”ติ. อาม ภนฺเต ชาเนยฺยาติ. เอวเมว โข มหาราช โย อุปฺปชฺชติ, ชานาติ โส “อุปฺปชฺชิสฺสามี”ติ.}

\prosewithcloserruled{กลฺโลสิ ภนฺเต นาคเสนาติ.}{อุปฺปชฺชติ\-ชานน\-ปโญฺห นวโม.}
\csromanmatikaanchor{mtk.67}
\csromantocmarkref{subsection}{10. พุทฺธนิทสฺสนปญฺห}{mtk.67}
\bookmark[dest={mtk.67},level=2]{10. พุทฺธนิทสฺสนปญฺห}
\csromanheader{2}{10. \textbf{พุทฺธ}\-\textbf{นิทสฺสน}\-\textbf{ปญฺห}}
\proseitem{10}{ราชา อาห “ภนฺเต นาคเสน พุทฺโธ อตฺถี”ติ. อาม มหาราช ภควา อตฺถีติ. สกฺกา ปน ภนฺเต นาคเสน พุทฺโธ นิทสฺเสตุํ “อิธ วา อิธ วา”ติ. ปรินิพฺพุโต มหาราช ภควา อนุปาทิเสสาย นิพฺพานธาตุยา, น สกฺกา ภควา นิทสฺเสตุํ “อิธ วา อิธ วา”ติ.}

\prose{โอปมฺมํ กโรหีติ. ตํ กิํ มญฺญสิ มหาราช, มหโต อคฺคิ\-กฺขนฺธสฺส ชลมานสฺส ยา อจฺจิ อตฺถงฺคตา, สกฺกา สา อจฺจิ ทสฺเสตุํ “อิธ วา อิธ วา”ติ. น หิ ภนฺเต, นิรุทฺธา สา อจฺจิ อปฺปญฺญตฺติํ คตาติ. เอวเมว โข มหาราช ภควา อนุปาทิเสสาย นิพฺพานธาตุยา ปรินิพฺพุโต อตฺถงฺคโต, น สกฺกา ภควา นิทสฺเสตุํ “อิธ วา อิธ วา”ติ, ธมฺมกาเยน ปน โข มหาราช สกฺกา ภควา นิทสฺเสตุํ, ธมฺโม หิ มหาราช ภควตา เทสิโตติ.}

\prose{กลฺโลสิ ภนฺเต นาคเสนาติ.}

\vspace{\tipitakaparskip}
\csromancenter{พุทฺธ\-นิทสฺสน\-ปโญฺห ทสโม.}
\nitthitammidruled{\textbf{พุทฺธวคฺโค ปญฺจโม.}}
\csromancenter{อิมสฺมิํ วคฺเค ทส ปญฺหา.}
\csromanmatikaanchor{mtk.68}
\csromantocmarknopage{section}{6. สติวคฺค}{mtk.68}
\bookmark[dest={mtk.68},level=1]{6. สติวคฺค}
\csromanheader{1}{6. \textbf{สติวคฺค}}
\csromanmatikaanchor{mtk.69}
\csromantocmarkref{subsection}{1. กายปิยายนปญฺห}{mtk.69}
\bookmark[dest={mtk.69},level=2]{1. กายปิยายนปญฺห}
\csromanheader{2}{1. \textbf{กาย}\-\textbf{ปิยายน}\-\textbf{ปญฺห}}
\proseitem{1}{\csromanfolio{78}ราชา อาห “ภนฺเต นาคเสน ปิโย ปพฺพชิตานํ กาโย”ติ. น โข มหาราช ปิโย ปพฺพชิตานํ กาโยติ. อถ กิสฺส นุ โข ภนฺเต เกลายถ มมายถาติ. กิํ ปน เต มหาราช กทาจิ กรหจิ สงฺคามคตสฺส กณฺฑปฺปหาโร โหตีติ. อาม ภนฺเต โหตีติ. กินฺนุ โข มหาราช โส วโณ อาเลเปน จ อาลิมฺปียติ เตเลน จ มกฺขียติ สุขุเมน จ โจฬปฏฺเฏน ปลิเวฐียตีติ. อาม ภนฺเต อาเลเปน จ อาลิมฺปียติ เตเลน จ มกฺขียติ สุขุเมน จ โจฬปฏฺเฏน ปลิเวฐียตีติ. กินฺนุ โข มหาราช ปิโย เต วโณ, เตน อาเลเปน จ อาลิมฺปียติ เตเลน จ มกฺขียติ สุขุเมน จ โจฬปฏฺเฏน ปลิเวฐียตีติ. น เม ภนฺเต ปิโย วโณ, อปิ จ มํสสฺส รุหนตฺถาย อาเลเปน จ อาลิมฺปียติ เตเลน จ มกฺขียติ สุขุเมน จ โจฬปฏฺเฏน ปลิเวฐียตีติ. เอวเมว โข มหาราช อปฺปิโย ปพฺพชิตานํ กาโย, อถ จ ปพฺพชิตา อนชฺโฌสิตา กายํ ปริหรนฺติ พฺรหฺมจริยา\-นุคฺคหาย, อปิ จ โข มหาราช วณูปโม กาโย วุตฺโต ภควตา, เตน ปพฺพชิตา วณมิว กายํ ปริหรนฺติ อนชฺโฌสิตา. ภาสิตมฺเปตํ มหาราช ภควตา\csromandash{}}

\setlength{\csromangathapagewidth}{0pt}
\setlength{\csromangathawakwidth}{0pt}
\csromangathameasuregroup{\textsuperscript{*}“อลฺลจมฺมปฺปฏิจฺฉนฺโน, \\ สมนฺตโต ปคฺฆรติ,}{\csromangathabat{\textsuperscript{*}“อลฺลจมฺมปฺปฏิจฺฉนฺโน,}{นวทฺวาโร มหาวโณ.} \\ \csromangathabat{สมนฺตโต ปคฺฆรติ,}{อสุจิปูติคนฺธิโย”ติ.}}
\csromangathasetleft{\textsuperscript{*}“อลฺลจมฺมปฺปฏิจฺฉนฺโน, \\ สมนฺตโต ปคฺฆรติ,}
\csromangathagroup{\csromangathastanza{\csromangathabat{\csromansymbolfootnote{*}{วิสุทฺธิ 1. 190 ปิฏฺเฐปิ.}“อลฺลจมฺมปฺปฏิจฺฉนฺโน,}{นวทฺวาโร มหาวโณ.} \\ \csromangathabat{สมนฺตโต ปคฺฆรติ,}{อสุจิปูติคนฺธิโย”ติ.}}}

\prosewithcloserruled{กลฺโลสิ ภนฺเต นาคเสนาติ.}{กาย\-ปิยายน\-ปโญฺห ปฐโม.}
\csromanmatikaanchor{mtk.70}
\csromantocmarkref{subsection}{2. สพฺพญฺญูภาวปญฺห}{mtk.70}
\bookmark[dest={mtk.70},level=2]{2. สพฺพญฺญูภาวปญฺห}
\csromanheader{2}{2. \textbf{สพฺพญฺญู}\-\textbf{ภาว}\-\textbf{ปญฺห}}
\proseitem{2}{ราชา อาห “ภนฺเต นาคเสน พุทฺโธ สพฺพญฺญู สพฺพทสฺสาวี”ติ. อาม มหาราช ภควา สพฺพญฺญู สพฺพทสฺสาวีติ. อถ กิสฺส นุ โข ภนฺเต นาคเสน สาวกานํ อนุปุพฺเพน สิกฺขาปทํ ปญฺญเปสีติ. อตฺถิ ปน \csromanfolio{79}เต มหาราช โกจิ เวชฺโช, โย อิมิสฺสํ ปถวิยํ สพฺพ\-เภสชฺชานิ ชานาตีติ. อาม ภนฺเต อตฺถีติ. กินฺนุ โข มหาราช โส เวชฺโช คิลานกํ สมฺปตฺเต กาเล เภสชฺชํ ปาเยติ, อุทาหุ อสมฺปตฺเต กาเลติ. สมฺปตฺเต กาเล ภนฺเต คิลานกํ เภสชฺชํ ปาเยติ, โน อสมฺปตฺเต กาเลติ. เอวเมว โข มหาราช ภควา สพฺพญฺญู สพฺพทสฺสาวี น อสมฺปตฺเต กาเล สาวกานํ สิกฺขาปทํ ปญฺญาเปติ, สมฺปตฺเต กาเล สาวกานํ สิกฺขาปทํ ปญฺญาเปติ ยาวชีวํ อนติกฺกมนียนฺติ.}

\prosewithcloserruled{กลฺโลสิ ภนฺเต นาคเสนาติ.}{สพฺพญฺญู\-ภาว\-ปโญฺห ทุติโย.}
\csromanmatikaanchor{mtk.71}
\csromantocmarkref{subsection}{3. มหาปุริสลกฺขณปญฺห}{mtk.71}
\bookmark[dest={mtk.71},level=2]{3. มหาปุริสลกฺขณปญฺห}
\csromanheader{2}{3. \textbf{มหาปุริสลกฺขณ}\-\textbf{ปญฺห}}
\proseitem{3}{ราชา อาห “ภนฺเต นาคเสน พุทฺโธ ทฺวตฺติํส\-มหา\-ปุริส\-ลกฺขเณหิ สมนฺนาคโต อสีติยา จ อนุพฺยญฺชเนหิ ปริรญฺชิโต สุวณฺณวณฺโณ กญฺจน\-สนฺนิภ\-ตฺตโจ พฺยามปฺปโภ”ติ. อาม มหาราช ภควา ทฺวตฺติํส\-มหา\-ปุริส\-ลกฺขเณหิ สมนฺนาคโต อสีติยา จ อนุพฺยญฺชเนหิ ปริรญฺชิโต สุวณฺณวณฺโณ กญฺจน\-สนฺนิภ\-ตฺตโจ พฺยามปฺปโภติ.}

\prose{กิํ ปนสฺส ภนฺเต มาตาปิตโรปิ ทฺวตฺติํส\-มหา\-ปุริส\-ลกฺขเณหิ สมนฺนาคตา อสีติยา จ อนุพฺยญฺชเนหิ ปริรญฺชิตา สุวณฺณวณฺณา กญฺจน\-สนฺนิภ\-ตฺตจา พฺยามปฺปภาติ. โน จสฺส มหาราช มาตาปิตโร ทฺวตฺติํส\-มหา\-ปุริส\-ลกฺขเณหิ สมนฺนาคตา อสีติยา จ อนุพฺยญฺชเนหิ ปริรญฺชิตา สุวณฺณวณฺณา กญฺจน\-สนฺนิภ\-ตฺตจา พฺยามปฺปภาติ.}

\prose{เอวํ สนฺเต โข ภนฺเต นาคเสน น อุปฺปชฺชติ พุทฺโธ ทฺวตฺติํส\-มหา\-ปุริส\-ลกฺขเณหิ สมนฺนาคโต อสีติยา จ อนุพฺยญฺชเนหิ ปริรญฺชิโต สุวณฺณวณฺโณ กญฺจน\-สนฺนิภ\-ตฺตโจ พฺยามปฺปโภติ, อปิ จ มาตุสทิโส วา ปุตฺโต โหติ มาตุปกฺโข วา, ปิตุสทิโส วา ปุตฺโต โหติ ปิตุปกฺโข วาติ. เถโร อาห “อตฺถิ ปน มหาราช กิญฺจิ ปทุมํ สตปตฺตนฺ”ติ. อาม ภนฺเต อตฺถีติ. ตสฺส ปน กุหิํ สมฺภโวติ. กทฺทเม ชายติ อุทเก อาสียตีติ. กินฺนุ โข มหาราช ปทุมํ กทฺทเมน สทิสํ วณฺเณน วา คนฺเธน วา รเสน วาติ. น หิ ภนฺเตติ. อถ อุทเกน \csromanfolio{80}วา คนฺเธน วา รเสน วาติ. น หิ ภนฺเตติ. เอวเมว โข มหาราช ภควา ทฺวตฺติํส\-มหา\-ปุริส\-ลกฺขเณหิ สมนฺนาคโต อสีติยา จ อนุพฺยญฺชเนหิ ปริรญฺชิโต สุวณฺณวณฺโณ กญฺจน\-สนฺนิภ\-ตฺตโจ พฺยามปฺปโภ, โน จสฺส มาตาปิตโร ทฺวตฺติํส\-มหา\-ปุริส\-ลกฺขเณหิ สมนฺนาคตา อสีติยา จ อนุพฺยญฺชเนหิ ปริรญฺชิตา สุวณฺณวณฺณา กญฺจน\-สนฺนิภ\-ตฺตจา พฺยามปฺปภาติ.}

\prosewithcloserruled{กลฺโลสิ ภนฺเต นาคเสนาติ.}{มหาปุริสลกฺขณ\-ปโญฺห ตติโย.}
\csromanmatikaanchor{mtk.72}
\csromantocmarkref{subsection}{4. ภควโต พฺรหฺมจาริปญฺห}{mtk.72}
\bookmark[dest={mtk.72},level=2]{4. ภควโต พฺรหฺมจาริปญฺห}
\csromanheader{2}{4. \textbf{ภควโต พฺรหฺมจาริปญฺห}}
\proseitem{3}{ราชา อาห “ภนฺเต นาคเสน พุทฺโธ พฺรหฺมจารี”ติ. อาม มหาราช ภควา พฺรหฺมจารีติ. เตน หิ ภนฺเต นาคเสน พุทฺโธ พฺรหฺมุโน สิสฺโสติ. อตฺถิ ปน เต มหาราช หตฺถิปาโมกฺโขติ. อตฺถิ ภนฺเตติ. กินฺนุ โข มหาราช โส หตฺถี กทาจิ กรหจิ โกญฺจนาทํ นทตีติ. อาม ภนฺเต นทตีติ. เตน หิ มหาราช โส หตฺถี โกญฺจสกุณสฺส สิสฺโสติ. น หิ ภนฺเตติ. กิํ ปน มหาราช พฺรหฺมา สพุทฺธิโก อพุทฺธิโกติ. สพุทฺธิโก ภนฺเตติ. เตน หิ มหาราช พฺรหฺมา ภควโต สิสฺโสติ.}

\prosewithcloserruled{กลฺโลสิ ภนฺเต นาคเสนาติ.}{ภควโต พฺรหฺมจาริ\-ปโญฺห จตุตฺโถ.}
\csromanmatikaanchor{mtk.73}
\csromantocmarkref{subsection}{5. ภควโต อุปสมฺปทาปญฺห}{mtk.73}
\bookmark[dest={mtk.73},level=2]{5. ภควโต อุปสมฺปทาปญฺห}
\csromanheader{2}{5. \textbf{ภควโต อุปสมฺปทาปญฺห}}
\proseitem{3}{ราชา อาห “ภนฺเต นาคเสน อุปสมฺปทา สุนฺทรา”ติ. อาม มหาราช อุปสมฺปทา สุนฺทราติ. อตฺถิ ปน ภนฺเต พุทฺธสฺส อุปสมฺปทา, อุทาหุ นตฺถีติ. อุปสมฺปนฺโน โข มหาราช ภควา โพธิรุกฺข\-มูเล สห สพฺพญฺญุต\-ญาเณน, นตฺถิ ภควโต อุปสมฺปทา อญฺเญหิ ทินฺนา, ยถา สาวกานํ มหาราช ภควา สิกฺขาปทํ ปญฺญเปติ ยาวชีวํ อนติกฺกมนียนฺติ.}

\prose{กลฺโลสิ ภนฺเต นาคเสนาติ.}

\vspace{\tipitakaparskip}
\csromancenter{ภควโต อุปสมฺปทา\-ปโญฺห ปญฺจโม.}
\csromanmatikaanchor{mtk.74}
\csromantocmarkref{subsection}{6. อสฺสุเภสชฺชาเภสชฺชปญฺห}{mtk.74}
\bookmark[dest={mtk.74},level=2]{6. อสฺสุเภสชฺชาเภสชฺชปญฺห}
\csromanmarkh{2. มิลินฺทปญฺห}\csromanmarkhii{6. อสฺสุ\-เภสชฺชา\-เภสชฺช\-ปญฺห}\csromanheadernomark{2}{2. มิลินฺทปญฺห \textbf{6. อสฺสุ}\-\textbf{เภสชฺชา}\-\textbf{เภสชฺช}\-\textbf{ปญฺห}}
\proseitem{6}{\csromanfolio{81}ราชา อาห “ภนฺเต นาคเสน โย จ มาตริ มตาย โรทติ, โย จ ธมฺมเปเมน โรทติ, อุภินฺนํ เตสํ โรทนฺตานํ กสฺส อสฺสุ เภสชฺชํ, กสฺส น เภสชฺชนฺ”ติ. เอกสฺส โข มหาราช อสฺสุ ราคโทสโมเหหิ สมลํ อุณฺหํ, เอกสฺส ปีติโสมนสฺเสน วิมลํ สีตลํ, ยํ โข มหาราช สีตลํ, ตํ เภสชฺชํ, ยํ อุณฺหํ, ตํ น เภสชฺชนฺติ.}

\prosewithcloserruled{กลฺโลสิ ภนฺเต นาคเสนาติ.}{อสฺสุ\-เภสชฺชา\-เภสชฺช\-ปโญฺห ฉฏฺโฐ.}
\csromanmatikaanchor{mtk.75}
\csromantocmarkref{subsection}{7. สราควีตราคนานากรณปญฺห}{mtk.75}
\bookmark[dest={mtk.75},level=2]{7. สราควีตราคนานากรณปญฺห}
\csromanheader{2}{7. \textbf{สราค}\-\textbf{วีตราค}\-\textbf{นานากรณ}\-\textbf{ปญฺห}}
\proseitem{7}{ราชา อาห “ภนฺเต นาคเสน กิํ นานากรณํ สราคสฺส จ วีตราคสฺส จา”ติ. เอโก โข มหาราช อชฺโฌสิโต, เอโก อนชฺโฌสิโตติ. กิํ เอตํ ภนฺเต อชฺโฌสิโต อนชฺโฌสิโต นามาติ, เอโก โข มหาราช อตฺถิโก, เอโก อนตฺถิโกติ. ปสฺสามหํ ภนฺเต เอวรูปํ โย จ สราโค, โย จ วีตราโค, สพฺโพเปโส โสภนํ เยว อิจฺฉติ ขาทนียํ วา โภชนียํ วา, น โกจิ ปาปกํ อิจฺฉตีติ. อวีตราโค โข มหาราช รสปฏิสํเวที จ รสราค\-ปฏิสํเวที จ โภชนํ ภุญฺชติ, วีตราโค ปน รสปฏิสํเวที โภชนํ ภุญฺชติ, โน จ โข รสราค\-ปฏิสํเวทีติ.}

\prosewithcloserruled{กลฺโลสิ ภนฺเต นาคเสนาติ.}{สราค\-วีตราค\-นานากรณ\-ปโญฺห สตฺตโม.}
\csromanmatikaanchor{mtk.76}
\csromantocmarkref{subsection}{8. ปญฺญาปติฏฺฐานปญฺห}{mtk.76}
\bookmark[dest={mtk.76},level=2]{8. ปญฺญาปติฏฺฐานปญฺห}
\csromanheader{2}{8. \textbf{ปญฺญา}\-\textbf{ปติฏฺฐาน}\-\textbf{ปญฺห}}
\proseitem{8}{ราชา อาห “ภนฺเต นาคเสน ปญฺญา กุหิํ ปฏิวสตี”ติ. น กตฺถจิ มหาราชาติ. เตน หิ ภนฺเต นาคเสน นตฺถิ ปญฺญาติ. วาโต มหาราช กุหิํ ปฏิวสตีติ. น กตฺถจิ ภนฺเตติ. เตน หิ มหาราช นตฺถิ วาโตติ.}

\prose{กลฺโลสิ ภนฺเต นาคเสนาติ.}

\vspace{\tipitakaparskip}
\csromancenter{ปญฺญา\-ปติฏฺฐาน\-ปโญฺห อฏฺฐโม.}
\csromanmatikaanchor{mtk.77}
\csromantocmarkref{subsection}{9. สํสารปญฺห}{mtk.77}
\bookmark[dest={mtk.77},level=2]{9. สํสารปญฺห}
\csromanheader{2}{9. \textbf{สํสารปญฺห}}
\proseitem{9}{\csromanfolio{82}ราชา อาห “ภนฺเต นาคเสน ยํ ปเนตํ พฺรูสิ ‘สํสาโร’ติ, กตโม โส สํสาโร”ติ. อิธ มหาราช ชาโต อิเธว มรติ, อิธ มโต อญฺญตฺร อุปฺปชฺชติ, ตหิํ ชาโต ตหิํ เยว มรติ, ตหิํ มโต อญฺญตฺร อุปฺปชฺชติ, เอวํ โข มหาราช สํสาโร โหตีติ.}

\prose{โอปมฺมํ กโรหีติ, ยถา มหาราช โกจิเทว ปุริโส ปกฺกํ อมฺพํ ขาทิตฺวา อฏฺฐิํ โรเปยฺย, ตโต มหนฺโต อมฺพรุกฺโข นิพฺพตฺติตฺวา ผลานิ ทเทยฺย, อถ โส ปุริโส ตโตปิ ปกฺกํ อมฺพํ ขาทิตฺวา อฏฺฐิํ โรเปยฺย, ตโตปิ มหนฺโต อมฺพรุกฺโข นิพฺพตฺติตฺวา ผลานิ ทเทยฺย, เอวเมเตสํ รุกฺขานํ โกฏิ น ปญฺญายติ, เอวเมว โข มหาราช อิธ ชาโต อิเธว มรติ, อิธ มโต อญฺญตฺร อุปฺปชฺชติ, ตหิํ ชาโต ตหิํ เยว มรติ, ตหิํ มโต อญฺญตฺร อุปฺปชฺชติ, เอวํ โข มหาราช สํสาโร โหตีติ.}

\prosewithcloserruled{กลฺโลสิ ภนฺเต นาคเสนาติ.}{สํสารปโญฺห นวโม.}
\csromanmatikaanchor{mtk.78}
\csromantocmarkref{subsection}{10. จิรกตสรณปญฺห}{mtk.78}
\bookmark[dest={mtk.78},level=2]{10. จิรกตสรณปญฺห}
\csromanheader{2}{10. \textbf{จิรกต}\-\textbf{สรณ}\-\textbf{ปญฺห}}
\proseitem{10}{ราชา อาห “ภนฺเต นาคเสน เกน อตีตํ จิรกตํ สรตี”ติ. สติยา มหาราชาติ. นนุ ภนฺเต นาคเสน จิตฺเตน สรติ โน สติยาติ. อภิชานาสิ นุ ตฺวํ มหาราช กิญฺจิเทว กรณียํ กตฺวา ปมุฏฺฐนฺติ. อาม ภนฺเตติ. กินฺนุ โข ตฺวํ มหาราช ตสฺมิํ สมเย อจิตฺตโก อโหสีติ. น หิ ภนฺเต, สติ ตสฺมิํ สมเย นาโหสีติ. อถ กสฺมา ตฺวํ มหาราช เอวมาห “จิตฺเตน สรติ, โน สติยา”ติ.}

\prose{กลฺโลสิ ภนฺเต นาคเสนาติ.}

\vspace{\tipitakaparskip}
\csromancenter{จิรกต\-สรณ\-ปโญฺห ทสโม.}
\csromanmatikaanchor{mtk.79}
\csromantocmarkref{subsection}{11. อภิชานนฺตสติปญฺห}{mtk.79}
\bookmark[dest={mtk.79},level=2]{11. อภิชานนฺตสติปญฺห}
\csromanmarkh{2. มิลินฺทปญฺห}\csromanmarkhii{11. อภิชานนฺต\-สติ\-ปญฺห}\csromanheadernomark{2}{2. มิลินฺทปญฺห \textbf{11. อภิชานนฺต}\-\textbf{สติ}\-\textbf{ปญฺห}}
\proseitem{11}{\csromanfolio{83}ราชา อาห “ภนฺเต นาคเสน สพฺพา สติ อภิชานนฺตี อุปฺปชฺชติ, อุทาหุ กฏุมิกาว สตี”ติ. อภิชานนฺตีปิ มหาราช กฏุมิกาปิ สตีติ. เอวญฺหิ โข ภนฺเต นาคเสน สพฺพา สติ อภิชานนฺตี, นตฺถิ กฏุมิกา สตีติ. ยทิ นตฺถิ มหาราช กฏุมิกา สติ, นตฺถิ กิญฺจิ สิปฺปิกานํ กมฺมายตเนหิ วา สิปฺปายตเนหิ วา วิชฺชาฏฺฐาเนหิ วา กรณียํ, นิรตฺถกา อาจริยา, ยสฺมา จ โข มหาราช อตฺถิ กฏุมิกา สติ, ตสฺมา อตฺถิ กมฺมายตเนหิ วา สิปฺปายตเนหิ วา วิชฺชาฏฺฐาเนหิ วา กรณียํ, อตฺโถ จ อาจริเยหีติ.}

\prose{กลฺโลสิ ภนฺเต นาคเสนาติ.}

\vspace{\tipitakaparskip}
\csromancenter{อภิชานนฺต\-สติ\-ปโญฺห เอกาทสโม.}
\nitthitammidruled{\textbf{สติวคฺโค ฉฏฺโฐ.}}
\csromancenter{อิมสฺมิํ วคฺเค เอกาทส ปญฺหา.}
\csromanmatikaanchor{mtk.80}
\csromantocmarknopage{section}{7. อรูปธมฺมววตฺถานวคฺค}{mtk.80}
\bookmark[dest={mtk.80},level=1]{7. อรูปธมฺมววตฺถานวคฺค}
\csromanheader{1}{7. \textbf{อรูป}\-\textbf{ธมฺมววตฺถาน}\-\textbf{วคฺค}}
\csromanmatikaanchor{mtk.81}
\csromantocmarkref{subsection}{1. สติอุปฺปชฺชนปญฺห}{mtk.81}
\bookmark[dest={mtk.81},level=2]{1. สติอุปฺปชฺชนปญฺห}
\prose{\csromanfolio{84}1. สติอุปฺปชฺชน\-ปญฺห}

\proseitem{1}{ราชา อาห “ภนฺเต นาคเสน กติหากาเรหิ สติ อุปฺปชฺชตี”ติ. สตฺตรสหา\-กาเรหิ มหาราช สติ อุปฺปชฺชตีติ. กตเมหิ สตฺตรสหา\-กาเรหีติ. อภิชานโตปิ มหาราช สติ อุปฺปชฺชติ, กฏุมิกายปิ สติ อุปฺปชฺชติ, โอฬาริกวิญฺญาณโตปิ สติ อุปฺปชฺชติ, หิตวิญฺญาณโตปิ สติ อุปฺปชฺชติ, อหิตวิญฺญาณโตปิ สติ อุปฺปชฺชติ, สภาคนิมิตฺตโตปิ สติ อุปฺปชฺชติ, วิสภาคนิมิตฺตโตปิ สติ อุปฺปชฺชติ, กถาภิญฺญาณโตปิ สติ อุปฺปชฺชติ, ลกฺขณโตปิ สติ อุปฺปชฺชติ, สารณโตปิ สติ อุปฺปชฺชติ, มุทฺทาโตปิ สติ อุปฺปชฺชติ, คณนาโตปิ สติ อุปฺปชฺชติ, ธารณโตปิ สติ อุปฺปชฺชติ, ภาวนโตปิ สติ อุปฺปชฺชติ, โปตฺถกนิพนฺธนโตปิ สติ อุปฺปชฺชติ, อุปนิกฺเขปโตปิ สติ อุปฺปชฺชติ, อนุภูตโตปิ สติ อุปฺปชฺชตีติ.}

\prose{กถํ อภิชานโต สติ อุปฺปชฺชติ, ยถา มหาราช อายสฺมา จ อานนฺโท ขุชฺชุตฺตรา จ อุปาสิกา, เย วา ปน อญฺเญปิ เกจิ ชาติสฺสรา ชาติํ สรนฺติ, เอวํ อภิชานโต สติ อุปฺปชฺชติ.}

\prose{กถํ กฏุมิกาย สติ อุปฺปชฺชติ, โย ปกติยา มุฏฺฐสฺสติโก, ปเร จ ตํ สราปนตฺถํ นิพนฺธนฺติ, เอวํ กฏุมิกาย สติ อุปฺปชฺชติ.}

\prose{กถํ โอฬาริก\-วิญฺญาณโต สติ อุปฺปชฺชติ, ยทา รชฺเช วา อภิสิตฺโต โหติ, โสตาปตฺติผลํ วา ปตฺโต โหติ, เอวํ โอฬาริก\-วิญฺญาณโต สติ อุปฺปชฺชติ.}

\prose{กถํ หิตวิญฺญาณโต สติ อุปฺปชฺชติ, ยมฺหิ สุขาปิโต, “อมุกสฺมิํ เอวํ สุขาปิโต”ติ สรติ, เอวํ หิตวิญฺญาณโต สติ อุปฺปชฺชติ.}

\prose{กถํ อหิตวิญฺญาณโต สติ อุปฺปชฺชติ, ยมฺหิ ทุกฺขา ปิโต, “อมุกสฺมิํ เอวํ ทุกฺขาปิโต”ติ สรติ, เอวํ อหิตวิญฺญาณโต สติ อุปฺปชฺชติ.}

\prose{กถํ สภาค\-นิมิตฺตโต สติ อุปฺปชฺชติ, สทิสํ ปุคฺคลํ ทิสฺวา มาตรํ วา ปิตรํ วา ภาตรํ วา ภคินิํ วา สรติ, โอฏฺฐํ วา โคณํ วา คทฺรภํ วา ทิสฺวา อญฺญํ ตาทิสํ โอฏฺฐํ วา โคณํ วา คทฺรภํ วา สรติ, เอวํ สภาค\-นิมิตฺตโต สติ อุปฺปชฺชติ.}

\prose{\csromanfolio{85}กถํ วิสภาค\-นิมิตฺตโต สติ อุปฺปชฺชติ, อสุกสฺส นาม วณฺโณ เอทิโส, สทฺโท เอทิโส, คนฺโธ เอทิโส, รโส เอทิโส, โผฏฺฐพฺโพ เอทิโสติ สรติ, เอวมฺปิ วิสภาค\-นิมิตฺตโตปิ สติ อุปฺปชฺชติ.}

\prose{กถํ กถา\-ภิญฺญาณโต สติ อุปฺปชฺชติ, โย ปกติยา มุฏฺฐสฺสติโก โหติ, ตํ ปเร สราเปนฺติ, เตน โส สรติ, เอวํ กถา\-ภิญฺญาณโต สติ อุปฺปชฺชติ.}

\prose{กถํ ลกฺขณโต สติ อุปฺปชฺชติ, โย ปกติยา พลีพทฺทานํ องฺเคน ชานาติ, ลกฺขเณน ชานาติ, เอวํ ลกฺขณโต สติ อุปฺปชฺชติ.}

\prose{กถํ สารณโต สติ อุปฺปชฺชติ, โย ปกติยา มุฏฺฐสฺสติโก โหติ, โย ตํ “สราหิ โภ, สราหิ โภ”ติ ปุนปฺปุนํ สราเปติ, เอวํ สารณโต สติ อุปฺปชฺชติ.}

\prose{กถํ มุทฺทาโต สติ อุปฺปชฺชติ, ลิปิยา สิกฺขิตตฺตา ชานาติ “อิมสฺส อกฺขรสฺส อนนฺตรํ อิมํ อกฺขรํ กาตพฺพนฺ”ติ, เอวํ มุทฺทาโต สติ อุปฺปชฺชติ.}

\prose{กถํ คณนาโต สติ อุปฺปชฺชติ, คณนาย สิกฺขิตตฺตา คณกา พหุมฺปิ คเณนฺติ, เอวํ คณนาโต สติ อุปฺปชฺชติ.}

\prose{กถํ ธารณโต สติ อุปฺปชฺชติ, ธารณาย สิกฺขิตตฺตา ธารณกา พหุมฺปิ ธาเรนฺติ, เอวํ ธารณโต สติ อุปฺปชฺชติ.}

\prose{กถํ ภาวนาโต สติ อุปฺปชฺชติ, อิธ ภิกฺขุ อเนกวิหิตํ ปุพฺเพนิวาสํ อนุสฺสรติ, เสยฺยถีทํ, เอกมฺปิ ชาติํ เทฺวปิ ชาติโย ฯเปฯ อิติ สาการํ สอุทฺเทสํ ปุพฺเพนิวาสํ อนุสฺสรติ, เอวํ ภาวนาโต สติ อุปฺปชฺชติ.}

\prose{กถํ โปตฺถก\-นิพนฺธนโต สติ อุปฺปชฺชติ, ราชาโน อนุสาสนิยํ อสฺสรนฺตา\csromansharedfootnote{d5da96e617b46066}{อนุสฺสรนฺตา (สพฺพตฺถ)} เอตํ โปตฺถกํ อาหรถาติ, เตน โปตฺถเกน อนุสฺสรนฺติ, เอวํ โปตฺถก\-นิพนฺธนโต สติ อุปฺปชฺชติ.}

\prose{กถํ อุปนิกฺเขปโต สติ อุปฺปชฺชติ, อุปนิกฺขิตฺตํ ภณฺฑํ ทิสฺวา สรติ, เอวํ อุปนิกฺเขปโต สติ อุปฺปชฺชติ.}

\prose{กถํ อนุภูตโต สติ อุปฺปชฺชติ, ทิฏฺฐตฺตา รูปํ สรติ, สุตตฺตา สทฺทํ สรติ, ฆายิตตฺตา คนฺธํ สรติ, สายิตตฺตา รสํ สรติ, ผุฏฺฐตฺตา}

\prose{\csromanfolio{86}โผฏฺฐพฺพํ สรติ, วิญฺญาตตฺตา ธมฺมํ สรติ, เอวํ อนุภูตโต สติ อุปฺปชฺชติ. อิเมหิ โข มหาราช สตฺตรสหา\-กาเรหิ สติ อุปฺปชฺชตีติ.}

\prose{กลฺโลสิ ภนฺเต นาคเสนาติ.}

\prose{สติอุปฺปชฺชน\-ปโญฺห ปฐโม.}
\csromansectionrule

\csromanmatikaanchor{mtk.82}
\csromantocmarkref{subsection}{2. พุทฺธคุณสติปฏิลาภปญฺห}{mtk.82}
\bookmark[dest={mtk.82},level=2]{2. พุทฺธคุณสติปฏิลาภปญฺห}
\csromanheader{2}{2. \textbf{พุทฺธคุณ}\-\textbf{สติ}\-\textbf{ปฏิลาภ}\-\textbf{ปญฺห}}
\proseitem{2}{ราชา อาห “ภนฺเต นาคเสน ตุเมฺห เอตํ ภณถ ‘โย วสฺสสตํ อกุสลํ กเรยฺย, มรณกาเล จ เอกํ พุทฺธคุณํ สติํ ปฏิลเภยฺย. โส เทเวสุ อุปฺปชฺเชยฺยา’ติ เอตํ น สทฺทหามิ, เอวญฺจ ปน วเทถ ‘เอเกน ปาณาติปาเตน นิรเย อุปฺปชฺเชยฺยา’ติ เอตมฺปิ น สทฺทหามี”ติ.}

\prose{ตํ กิํ มญฺญสิ มหาราช, ขุทฺทโกปิ ปาสาโณ วินา นาวาย อุทเก อุปฺปิลเวยฺยาติ. น หิ ภนฺเตติ. กิํ นุ โข มหาราช วาหสตมฺปิ ปาสาณานํนาวาย อาโรปิตํ อุทเก อุปฺปิลเวยฺยาติ. อาม ภนฺเตติ. ยถา มหาราช นาวา, เอวํ กุสลานิ กมฺมานิ ทฏฺฐพฺพานีติ.}

\prosewithcloserruled{กลฺโลสิ ภนฺเต นาคเสนาติ.}{พุทฺธคุณ\-สติ\-ปฏิลาภ\-ปโญฺห ทุติโย.}
\csromanmatikaanchor{mtk.83}
\csromantocmarkref{subsection}{3. ทุกฺขปฺปหานวายมปญฺห}{mtk.83}
\bookmark[dest={mtk.83},level=2]{3. ทุกฺขปฺปหานวายมปญฺห}
\csromanheader{2}{3. \textbf{ทุกฺข}\-\textbf{ปฺปหาน}\-\textbf{วายม}\-\textbf{ปญฺห}}
\proseitem{3}{ราชา อาห “ภนฺเต นาคเสน กิํ ตุเมฺห อตีตสฺส ทุกฺขสฺส ปหานาย วายมถา”ติ. น หิ มหาราชาติ. กิํ ปน ภนฺเต อนาคตสฺส ทุกฺขสฺส ปหานาย วายมถาติ. น หิ มหาราชาติ. กิํ ปน ปจฺจุปฺปนฺนสฺส ทุกฺขสฺส ปหานาย วายมถาติ. น หิ มหาราชาติ. ยทิ ตุเมฺห น อตีตสฺส ทุกฺขสฺส ปหานาย วายมถ, น อนาคตสฺส ทุกฺขสฺส ปหานาย วายมถ, น ปจฺจุปฺปนฺนสฺส ทุกฺขสฺส ปหานาย วายมถ, อถ กิมตฺถาย วายมถาติ. เถโร อาห “กินฺติ มหาราช อิทญฺจ ทุกฺขํ นิรุชฺเฌยฺย, อญฺญญฺจ ทุกฺขํ นุปฺปชฺเชยฺยา”ติ เอตทตฺถาย วายมามาติ.}

\prose{\csromanfolio{87}อตฺถิ ปน เต ภนฺเต นาคเสน อนาคตํ ทุกฺขนฺติ. นตฺถิ\csromansharedfootnote{231f91e455b85226}{อภิ 4. 417 ปิฏฺเฐ ปสฺสิตพฺพํ.} มหาราชาติ. ตุเมฺห โข ภนฺเต นาคเสน อติปณฺฑิตา, เย ตุเมฺห อสนฺตานํ อนาคตานํ ทุกฺขานํ ปหานาย วายมถาติ. อตฺถิ ปน เต มหาราช เกจิ ปฏิราชาโน ปจฺจตฺถิกา ปจฺจามิตฺตา ปจฺจุปฏฺฐิตา โหนฺตีติ. อาม ภนฺเต อตฺถีติ. กิํ นุ โข มหาราช ตทา ตุเมฺห ปริขํ ขณาเปยฺยาถ, ปาการํ จินาเปยฺยาถ โคปุรํ การาเปยฺยาถ, อฏฺฏาลกํ การาเปยฺยาถ, ธญฺญํ อติหราเปยฺยาถาติ. น หิ ภนฺเต, ปฏิกจฺเจว ตํ ปฏิยตฺตํ โหตีติ. กิํ ตุเมฺห มหาราช ตทา หตฺถิสฺมิํ สิกฺเขยฺยาถ, อสฺสสฺมิํ สิกฺเขยฺยาถ, รถสฺมิํ สิกฺเขยฺยาถ, ธนุสฺมิํ สิกฺเขยฺยาถ, ถรุสฺมิํ สิกฺเขยฺยาถาติ. น หิ ภนฺเต, ปฏิกจฺเจว ตํ สิกฺขิตํ โหตีติ. กิสฺสตฺถายาติ. อนาคตานํ ภนฺเต ภยานํ ปฏิพาหน\-ตฺถายาติ. กิํ นุ โข มหาราช อตฺถิ อนาคตํ ภยนฺติ. นตฺถิ ภนฺเตติ. ตุเมฺห จ โข มหาราช อติปณฺฑิตา, เย ตุเมฺห อสนฺตานํ อนาคตานํ ภยานํ ปฏิพาหน\-ตฺถาย ปฏิยาเทถาติ.}

\prose{ภิยฺโย โอปมฺมํ กโรหีติ. ตํ กิํ มญฺญสิ มหาราช, ยทา ตฺวํ ปิปาสิโต ภเวยฺยาสิ, ตทา ตฺวํ อุทปานํ ขณาเปยฺยาสิ, โปกฺขรณิํ ขณาเปยฺยาสิ, ตฬากํ ขณาเปยฺยาสิ “ปานียํ ปิวิสฺสามี”ติ. น หิ ภนฺเต, ปฏิกจฺเจว ตํ ปฏิยตฺตํ โหตีติ. กิสฺสตฺถายาติ. อนาคตานํ ภนฺเต ปิปาสานํ ปฏิพาหน\-ตฺถาย ปฏิยตฺตํ โหตีติ. อตฺถิ ปน มหาราช อนาคตา ปิปาสาติ. นตฺถิ ภนฺเตติ. ตุเมฺห โข มหาราช อติปณฺฑิตา, เย ตุเมฺห อสนฺตานํ อนาคตานํ ปิปาสานํ ปฏิพาหน\-ตฺถาย ตํ ปฏิยาเทถาติ.}

\prose{ภิยฺโย โอปมฺมํ กโรหีติ. ตํ กิํ มญฺญสิ มหาราช, ยทา ตฺวํ พุภุกฺขิโต ภเวยฺยาสิ, ตทา ตฺวํ เขตฺตํ กสาเปยฺยาสิ, สาลิํ วปาเปยฺยาสิ “ภตฺตํ ภุญฺชิสฺสามี”ติ. น หิ ภนฺเต, ปฏิกจฺเจว ตํ ปฏิยตฺตํ โหตีติ. กิสฺสตฺถายาติ. อนาคตานํ ภนฺเต พุภุกฺขานํ ปฏิพาหน\-ตฺถายาติ. อตฺถิ ปน มหาราช อนาคตา พุภุกฺขาติ. นตฺถิ ภนฺเตติ.}

\prose{\csromanfolio{88}ตุเมฺห โข มหาราช อติปณฺฑิตา, เย ตุเมฺห อสนฺตานํ อนาคตานํ พุภุกฺขานํ ปฏิพาหน\-ตฺถาย ปฏิยาเทถาติ.}

\prosewithcloserruled{กลฺโลสิ ภนฺเต นาคเสนาติ.}{ทุกฺข\-ปฺปหาน\-วายม\-ปโญฺห ตติโย.}
\csromanmatikaanchor{mtk.84}
\csromantocmarkref{subsection}{4. พฺรหฺมโลกปญฺห}{mtk.84}
\bookmark[dest={mtk.84},level=2]{4. พฺรหฺมโลกปญฺห}
\csromanheader{2}{4. \textbf{พฺรหฺมโลก}\-\textbf{ปญฺห}}
\proseitem{4}{ราชา อาห “ภนฺเต นาคเสน กีวทูโร อิโต พฺรหฺมโลโก”ติ. ทูโร โข มหาราช อิโต พฺรหฺมโลโก กูฏาคารมตฺตา สิลา ตมฺหา ปติตา อโหรตฺเตน อฏฺฐจตฺตาลีส\-โยชนสหสฺสานิ ภสฺสมานา จตูหิ มาเสหิ ปถวิยํ ปติฏฺฐเหยฺยาติ.}

\prose{ภนฺเต นาคเสน ตุเมฺห เอวํ ภณถ “เสยฺยถาปิ พลวา ปุริโส สมิญฺชิตํ วา พาหํ ปสาเรยฺย, ปสาริตํ วา พาหํ สมิญฺเชยฺย, เอวเมว อิทฺธิมา ภิกฺขุ เจโตวสิปฺปตฺโต ชมฺพุทีเป อนฺตรหิโต พฺรหฺมโลเก ปาตุภเวยฺยา”ติ เอตํ วจนํ น สทฺทหามิ, เอวํ อติสีฆํ ตาว พหูนิ โยชนสตานิ คจฺฉิสฺสตีติ.}

\prose{เถโร อาห “กุหิํ ปน มหาราช ตว ชาตภูมี”ติ. อตฺถิ ภนฺเต อลสนฺโท นาม ทีโป, ตตฺถาหํ ชาโตติ. กีวทูโร มหาราช อิโต อลสนฺโท โหตีติ. ทฺวิมตฺตานิ ภนฺเต โยชนสตานีติ. อภิชานาสิ นุ ตฺวํ มหาราช ตตฺถ กิญฺจิเทว กรณียํ กริตฺวา สริตาติ. อาม ภนฺเต สรามีติ. ลหุํ โข ตฺวํ มหาราช คโตสิ ทฺวิมตฺตานิ โยชนสตานีติ.}

\prosewithcloserruled{กลฺโลสิ ภนฺเต นาคเสนาติ.}{พฺราหฺมโลกปโญฺห จตุตฺโถ.}
\csromanmatikaanchor{mtk.85}
\csromantocmarkref{subsection}{5. ทฺวนฺนํ โลกุปฺปนฺนานํ สมกภาวปญฺห}{mtk.85}
\bookmark[dest={mtk.85},level=2]{5. ทฺวนฺนํ โลกุปฺปนฺนานํ สมกภาวปญฺห}
\csromanheader{2}{5. \textbf{ทฺวินฺนํ โลกุปฺปนฺนานํ สมกภาวปญฺห}}
\proseitem{5}{ราชา อาห “ภนฺเต นาคเสน โย อิธ กาลงฺกโต พฺรหฺมโลเก อุปฺปชฺเชยฺย, โย จ อิธ กาลงฺกโต กสฺมีเร อุปฺปชฺเชยฺย, โก จิรตรํ โก สีฆตรนฺ”ติ. สมกํ มหาราชาติ.}

\prose{\csromanfolio{89}โอปมฺมํ กโรหีติ. กุหิํ ปน มหาราช ตว ชาตนครนฺติ. อตฺถิ ภนฺเต กลสิคาโม นาม, ตตฺถาหํ ชาโตติ. กีว ทูโร มหาราช อิโต กลสิคาโม โหตีติ. ทฺวิมตฺตานิ ภนฺเต โยชนสตานีติ. กีว ทูรํ มหาราช อิโต กสฺมีรํ โหตีติ. ทฺวาทส ภนฺเต โยชนานีติ. อิงฺฆ ตฺวํ มหาราช กลสิคามํ จินฺเตหีติ. จินฺติโต ภนฺเตติ. อิงฺฆ ตฺวํ มหาราช กสฺมีรํ จินฺเตหีติ. จินฺติตํ ภนฺเตติ. กตมํ นุ โข มหาราช จิเรน จินฺติตํ, กตมํ สีฆตรนฺติ. สมกํ ภนฺเตติ. เอวเมว โข มหาราช โย อิธ กาลงฺกโต พฺรหฺมโลเก อุปฺปชฺเชยฺย, โย จ อิธ กาลงฺกโต กสฺมีเร อุปฺปชฺเชยฺย, สมกํ เยว อุปฺปชฺชนฺตีติ.}

\prose{ภิยฺโย โอปมฺมํ กโรหีติ. ตํ กิํ มญฺญสิ มหาราช, เทฺว สกุณา อากาเสน คจฺเฉยฺยุํ, เตสุ เอโก อุจฺเจ รุกฺเข นิสีเทยฺย, เอโก นีเจ รุกฺเข นิสีเทยฺย, เตสํ สมกํ ปติฏฺฐิตานํ กตมสฺส ฉายา ปฐมตรํ ปถวิยํ ปติฏฺฐเหยฺย, กตมสฺส ฉายา จิเรน ปถวิยํ ปติฏฺฐเหยฺยาติ. สมกํ ภนฺเตติ. เอวเมว โข มหาราช โย อิธ กาลงฺกโต พฺรหฺมโลเก อุปฺปชฺเชยฺย, โย จ อิธ กาลงฺกโต กสฺมีเร อุปฺปชฺเชยฺย, สมกํ เยว อุปฺปชฺชนฺตีติ.}

\prose{กลฺโลสิ ภนฺเต นาคเสนาติ.}

\prose{ทฺวินฺนํ โลกุปฺปนฺนานํ สมก\-ภาว\-ปโญฺห ปญฺจโม.}
\csromansectionrule

\csromanmatikaanchor{mtk.86}
\csromantocmarkref{subsection}{6. โพชฺฌงฺคปญฺห}{mtk.86}
\bookmark[dest={mtk.86},level=2]{6. โพชฺฌงฺคปญฺห}
\csromanheader{2}{6. \textbf{โพชฺฌงฺคปญฺห}}
\proseitem{6}{ราชา อาห “กติ นุ โข ภนฺเต นาคเสน โพชฺฌงฺคา”ติ. สตฺต โข มหาราช โพชฺฌงฺคาติ. กติหิ ปน ภนฺเต โพชฺฌงฺเคหิ พุชฺฌตีติ. เอเกน โข มหาราช โพชฺฌงฺเคน พุชฺฌติ ธมฺมวิจย\-สมฺโพชฺฌงฺเคนาติ. อถ กิสฺส นุ โข ภนฺเต วุจฺจนฺติ “สตฺต โพชฺฌงฺคา”ติ. ตํ กิํ มญฺญสิ มหาราช, อสิ โกสิยา ปกฺขิตฺโต อคฺคหิโต หตฺเถน อุสฺสหติ เฉชฺชํ ฉินฺทิตุนฺติ. น หิ ภนฺเตติ. เอวเมว โข มหาราช ธมฺมวิจย\-สมฺโพชฺฌงฺเคน วินา ฉหิ โพชฺฌงฺเคหิ น พุชฺฌตีติ.}

\prose{กลฺโลสิ ภนฺเต นาคเสนาติ.}

\vspace{\tipitakaparskip}
\csromancenter{โพชฺฌงฺคปโญฺห ฉฏฺโฐ.}
\csromanmatikaanchor{mtk.87}
\csromantocmarkref{subsection}{7. ปาปปุญฺญานํ อปฺปานปฺปภาวปญฺห}{mtk.87}
\bookmark[dest={mtk.87},level=2]{7. ปาปปุญฺญานํ อปฺปานปฺปภาวปญฺห}
\csromanheader{2}{7. \textbf{ปาปปุญฺญานํ อปฺปาน}\-\textbf{ปฺปภาว}\-\textbf{ปญฺห}}
\proseitem{7}{\csromanfolio{90}ราชา อาห “ภนฺเต นาคเสน กตรํ นุ โข พหุตรํ ปุญฺญํ วา อปุญฺญํ วา”ติ. ปุญฺญํ โข มหาราช พหุตรํ, อปุญฺญํ โถกนฺติ. เกน การเณนาติ. อปุญฺญํ โข มหาราช กโรนฺโต วิปฺปฏิสารี โหติ “ปาปกมฺมํ มยา กตนฺ”ติ, เตน ปาปํ น วฑฺฒติ. ปุญฺญํ โข มหาราช กโรนฺโต อวิปฺปฏิสารี โหติ, อวิปฺปฏิสาริโน ปาโมชฺชํ ชายติ, ปมุทิตสฺส ปีติ ชายติ, ปีติมนสฺส กาโย ปสฺสมฺภติ, ปสฺสทฺธกาโย สุขํ เวเทติ, สุขิโน จิตฺตํ สมาธิยติ, สมาหิโต ยถาภูตํ ปชานาติ, เตน การเณน ปุญฺญํ วฑฺฒติ. ปุริโส โข มหาราช ฉินฺน\-หตฺถปาโท ภควโต เอกํ อุปฺปลหตฺถํ ทตฺวา เอกนวุติ\-กปฺปานิ วินิปาตํ น คจฺฉิสฺสติ, อิมินาปิ มหาราช การเณน ภณามิ “ปุญฺญํ พหุตรํ, อปุญฺญํ โถกนฺ”ติ.}

\prosewithcloserruled{กลฺโลสิ ภนฺเต นาคเสนาติ.}{ปาปปุญฺญานํ อปฺปาน\-ปฺปภาว\-ปโญฺห สตฺตโม.}
\csromanmatikaanchor{mtk.88}
\csromantocmarkref{subsection}{8. ชานนฺตาชานนฺตปาปกรณปญฺห}{mtk.88}
\bookmark[dest={mtk.88},level=2]{8. ชานนฺตาชานนฺตปาปกรณปญฺห}
\csromanheader{2}{8. \textbf{ชานนฺตา}\-\textbf{ชานนฺต}\-\textbf{ปาปกรณ}\-\textbf{ปญฺห}}
\proseitem{8}{ราชา อาห “ภนฺเต นาคเสน โย ชานนฺโต ปาปกมฺมํ กโรติ, โย อชานนฺโต ปาปกมฺมํ กโรติ, กสฺส พหุตรํ อปุญฺญนฺ”ติ. เถโร อาห “โย โข มหาราช อชานนฺโต ปาปกมฺมํ กโรติ, ตสฺส พหุตรํ อปุญฺญนฺ”ติ. เตน หิ ภนฺเต นาคเสน โย อมฺหากํ ราชปุตฺโต วา ราชมหามตฺโต วา อชานนฺโต ปาปกมฺมํ กโรติ, ตํ มยํ ทิคุณํ ทณฺเฑมาติ. ตํ กิํ มญฺญสิ มหาราช, ตตฺตํ อโยคุฬํ อาทิตฺตํ สมฺปชฺชลิตํ สโชติภูตํ เอโก ชานนฺโต คเณฺหยฺย, เอโก อชานนฺโต คเณฺหยฺย, กตโม\csromansharedfootnote{adc110b73e5652ee}{กสฺส (ก)} พลวตรํ ฑเยฺหยฺยาติ. โย โข ภนฺเต อชานนฺโต คเณฺหยฺย, โส\csromansharedfootnote{74f3b3db4495d6cc}{ตสฺส (อิ, ก)} พลวตรํ ฑเยฺหยฺยาติ. เอวเมว โข มหาราช โย อชานนฺโต ปาปกมฺมํ กโรติ, ตสฺส พหุตรํ อปุญฺญนฺติ.}

\prose{กลฺโลสิ ภนฺเต นาคเสนาติ.}

\vspace{\tipitakaparskip}
\csromancenter{ชานนฺตา\-ชานนฺต\-ปาปกรณ\-ปโญฺห อฏฺฐโม.}
\csromanmatikaanchor{mtk.89}
\csromantocmarkref{subsection}{9. อุตฺตรกุรุกาทิคมนปญฺห}{mtk.89}
\bookmark[dest={mtk.89},level=2]{9. อุตฺตรกุรุกาทิคมนปญฺห}
\csromanmarkh{2. มิลินฺทปญฺห}\csromanmarkhii{9. อุตฺตรกุรุกา\-ทิ\-คมน\-ปญฺห}\csromanheadernomark{2}{2. มิลินฺทปญฺห \textbf{9. อุตฺตรกุรุกา}\-\textbf{ทิ}\-\textbf{คมน}\-\textbf{ปญฺห}}
\csromanmatikaanchor{mtk.90}
\csromantocmarkref{subsection}{10. ธีฆฏฺฐิปญฺห}{mtk.90}
\bookmark[dest={mtk.90},level=2]{10. ธีฆฏฺฐิปญฺห}
\proseitem{9}{\csromanfolio{91}ราชา อาห “ภนฺเต นาคเสน อตฺถิ โกจิ, โย อิมินา สรีเรน อุตฺตรกุรุํ วา คจฺเฉยฺย, พฺราหฺมโลกํ วา, อญฺญํ วา ปน ทีปนฺ”ติ. อตฺถิ มหาราช โย อิมินา จาตุมฺมหาภูติเกน กาเยน อุตฺตรกุรุํ วา คจฺเฉยฺย, พฺราหฺมโลกํ วา, อญฺญํ วา ปน ทีปนฺติ.}

\prose{กถํ ภนฺเต นาคเสน อิมินา จาตุมฺมหาภูติเกน กาเยน อุตฺตรกุรุํ วา คจฺเฉยฺย, พฺรหฺมโลกํ วา, อญฺญํ วา ปน ทีปนฺติ. อภิชานาสิ นุ ตฺวํ มหาราช อิมิสฺสา ปถวิยา วิทตฺถิํ วา รตนํ วา ลงฺฆิตาติ. อาม ภนฺเต อภิชานามิ “อหํ ภนฺเต นาคเสน อฏฺฐปิ รตนิโย ลงฺเฆมี”ติ. กถํ ตฺวํ มหาราช อฏฺฐปิ รตนิโย ลงฺเฆสีติ. อหญฺหิ ภนฺเต จิตฺตํ อุปฺปาเทมิ “เอตฺถ นิปติสฺสามี”ติ สห จิตฺตุปฺปาเทน กาโย เม ลหุโก โหตีติ. เอวเมว โข มหาราช อิทฺธิมา ภิกฺขุ เจโตวสิปฺปตฺโต กายํ จิตฺเต สมาโรเปตฺวา จิตฺตวเสน เวหาสํ คจฺฉตีติ.}

\prosewithcloserruled{กลฺโลสิ ภนฺเต นาคเสนาติ.}{อุตฺตรกุรุกา\-ทิ\-คมน\-ปโญฺห นวโม.}
\csromanheader{2}{10. \textbf{ทีฆฏฺฐิปญฺห}}
\proseitem{10}{ราชา อาห “ภนฺเต นาคเสน ตุเมฺห เอวํ ภณถ ‘อฏฺฐิกานิ ทีฆานิ โยชนสติกานิปี’ติ, รุกฺโขปิ ตาว นตฺถิ โยชนสติโก, กุโต ปน อฏฺฐิกานิ ทีฆานิ โยชนสติกานิ ภวิสฺสนฺตี”ติ.}

\prose{ตํ กิํ มญฺญสิ มหาราช, สุตํ เต “มหาสมุทฺเท ปญฺจ\-โยชน\-สติกาปิ มจฺฉา อตฺถี”ติ. อาม ภนฺเต สุตนฺติ. นนุ มหาราช ปญฺจ\-โยชน\-สติกสฺส มจฺฉสฺส อฏฺฐิกานิ ทีฆานิ ภวิสฺสนฺติ โยชนสติกานิปีติ.}

\prose{กลฺโลสิ ภนฺเต นาคเสนาติ.}

\vspace{\tipitakaparskip}
\csromancenter{ทีฆฏฺฐิปโญฺห ทสโม.}
\csromanmatikaanchor{mtk.91}
\csromantocmarkref{subsection}{11. อสฺสาสปสฺสาสนิโรธปญฺห}{mtk.91}
\bookmark[dest={mtk.91},level=2]{11. อสฺสาสปสฺสาสนิโรธปญฺห}
\csromanheader{2}{11. \textbf{อสฺสาสปสฺสาส}\-\textbf{นิโรธ}\-\textbf{ปญฺห}}
\proseitem{11}{\csromanfolio{92}ราชา อาห “ภนฺเต นาคเสน ตุเมฺห เอวํ ภณถ ‘สกฺกา อสฺสาสปสฺสาเส นิโรเธตุนฺ’ติ”. อาม มหาราช สกฺกา อสฺสาสปสฺสาเส นิโรเธตุนฺติ. กถํ ภนฺเต นาคเสน สกฺกา อสฺสาสปสฺสาเส นิโรเธตุนฺติ. ตํ กิํ มญฺญสิ มหาราช, สุตปุพฺโพ เต โกจิ กากจฺฉมาโนติ. อาม ภนฺเต สุตปุพฺโพติ. กินฺนุ โข มหาราช โส สทฺโท กาเย นมิเต วิรเมยฺยาติ. อาม ภนฺเต วิรเมยฺยาติ. โส หิ นาม มหาราช สทฺโท อภาวิตกายสฺส อภาวิตสีลสฺส อภาวิต\-จิตฺตสฺส อภาวิต\-ปญฺญสฺส กาเย นมิเต วิรมิสฺสติ, กิํ ปน ภาวิตกายสฺส ภาวิตสีลสฺส ภาวิต\-จิตฺตสฺส ภาวิต\-ปญฺญสฺส จตุตฺถ\-ชฺฌานํ สมาปนฺนสฺส อสฺสาสปสฺสาสา น นิรุชฺฌิสฺสนฺตีติ.}

\prosewithcloserruled{กลฺโลสิ ภนฺเต นาคเสนาติ.}{อสฺสาสปสฺสาส\-นิโรธ\-ปโญฺห เอกาทสโม.}
\csromanmatikaanchor{mtk.92}
\csromanmatikaanchor{mtk.93}
\csromantocmarkref{subsection}{12. สมุทฺทปญฺห}{mtk.92}
\bookmark[dest={mtk.92},level=2]{12. สมุทฺทปญฺห}
\csromantocmarkref{subsection}{13. สมุทฺทเอกรสปญฺห}{mtk.93}
\bookmark[dest={mtk.93},level=2]{13. สมุทฺทเอกรสปญฺห}
\csromanheader{2}{12. \textbf{สมุทฺทปญฺห}}
\proseitem{12}{ราชา อาห “ภนฺเต นาคเสน ‘สมุทฺโท สมุทฺโท’ติ วุจฺจติ, เกน การเณน อุทกํ ‘สมุทฺโท’ติ วุจฺจตี”ติ. เถโร อาห “ยตฺตกํ มหาราช อุทกํ, ตตฺตกํ โลณํ. ยตฺตกํ โลณํ, ตตฺตกํ อุทกํ. ตสฺมา ‘สมุทฺโท’ติ วุจฺจตี”ติ.}

\prosewithcloserruled{กลฺโลสิ ภนฺเต นาคเสนาติ.}{สมุทฺทปโญฺห ทฺวาทสโม.}
\prose{13. สมุทฺทเอกรสปญฺห}

\proseitem{13}{ราชา อาห “ภนฺเต นาคเสน เกน การเณน สมุทฺโท เอกรโส โลณรโส”ติ. จิรสณฺฐิตตฺตา โข มหาราช อุทกสฺส สมุทฺโท เอกรโส โลณรโสติ.}

\prose{กลฺโลสิ ภนฺเต นาคเสนาติ.}

\prose{สมุทฺทเอกรสปโญฺห เตรสโม.}

\csromanmatikaanchor{mtk.94}
\csromantocmarkref{subsection}{14. สุขุมปญฺห}{mtk.94}
\bookmark[dest={mtk.94},level=2]{14. สุขุมปญฺห}
\csromanmarkh{2. มิลินฺทปญฺห}\csromanmarkhii{14. สุขุมปญฺห}\csromanheadernomark{2}{2. มิลินฺทปญฺห \textbf{14. สุขุมปญฺห}}
\csromanmatikaanchor{mtk.95}
\csromantocmarkref{subsection}{15. วิญฺญาณนานาตฺถปญฺห}{mtk.95}
\bookmark[dest={mtk.95},level=2]{15. วิญฺญาณนานาตฺถปญฺห}
\proseitem{14}{\csromanfolio{93}ราชา อาห “ภนฺเต นาคเสน สกฺกา สพฺพํ สุขุมํ ฉินฺทิตุนฺ”ติ. อาม มหาราช สกฺกา สพฺพํ สุขุมํ ฉินฺทิตุนฺติ. กิํ ปน ภนฺเต สพฺพํ สุขุมนฺติ. ธมฺโม โข มหาราช สพฺพสุขุโม, น โข มหาราช ธมฺมา สพฺเพ สุขุมา, “สุขุมนฺ”ติ วา “ถูลนฺ”ติ วา ธมฺมานเม\-ตม\-ธิวจนํ. ยํ กิญฺจิ ฉินฺทิตพฺพํ, สพฺพํ ตํ ปญฺญาย ฉินฺทติ, นตฺถิ ทุติยํ ปญฺญาย เฉทนนฺติ.}

\prosewithcloserruled{กลฺโลสิ ภนฺเต นาคเสนาติ.}{สุขุมปโญฺห จุทฺทสโม.}
\csromanheader{2}{15. \textbf{วิญฺญาณ}\-\textbf{นานตฺถ}\-\textbf{ปญฺห}}
\proseitem{15}{ราชา อาห “ภนฺเต นาคเสน ‘วิญฺญาณนฺ’ติ วา ‘ปญฺญา’ติ วา ‘ภูตสฺมิํ ชีโว’ติ วา อิเม ธมฺมา นานตฺถา เจว นานาพฺยญฺชนา จ, อุทาหุ เอกตฺถา พฺยญฺชนเมว นานนฺ”ติ. วิชานน\-ลกฺขณํ มหาราช วิญฺญาณํ, ปชานน\-ลกฺขณา ปญฺญา, ภูตสฺมิํ ชีโว นุปลพฺภตีติ. ยทิ ชีโว นุปลพฺภติ, อถ โก จรหิ จกฺขุนา รูปํ ปสฺสติ, โสเตน สทฺทํ สุณาติ, ฆาเนน คนฺธํ ฆายติ, ชิวฺหาย รสํ สายติ, กาเยน โผฏฺฐพฺพํ ผุสติ, มนสา ธมฺมํ วิชานาตีติ. เถโร อาห “ยทิ ชีโว จกฺขุนา รูปํ ปสฺสติ ฯเปฯ มนสา ธมฺมํ วิชานาติ, โส ชีโว จกฺขุทฺวาเรสุ อุปฺปาฏิเตสุ มหนฺเตน อากาเสน พหิมุโข สุฏฺฐุตรํ รูปํ ปสฺเสยฺย, โสเตสุ อุปฺปาฏิเตสุ, ฆาเน อุปฺปาฏิเต, ชิวฺหาย อุปฺปาฏิตาย, กาเย อุปฺปาฏิเต มหนฺเตน อากาเสน สุฏฺฐุตรํ สทฺทํ สุเณยฺย, คนฺธํ ฆาเยยฺย, รสํ สาเยยฺย, โผฏฺฐพฺพํ ผุเสยฺยา”ติ. น หิ ภนฺเตติ. เตน หิ มหาราช ภูตสฺมิํ ชีโว นุปลพฺภตีติ.}

\prose{กลฺโลสิ ภนฺเต นาคเสนาติ.}

\vspace{\tipitakaparskip}
\csromancenter{วิญฺญาณ\-นานตฺถ\-ปโญฺห ปนฺนรสโม.}
\csromanmatikaanchor{mtk.96}
\csromantocmarkref{subsection}{16. อรูปธมฺมววตฺถานทุกฺกรปญฺห}{mtk.96}
\bookmark[dest={mtk.96},level=2]{16. อรูปธมฺมววตฺถานทุกฺกรปญฺห}
\csromanheader{2}{16. \textbf{อรูป}\-\textbf{ธมฺมววตฺถาน}\-\textbf{ทุกฺกร}\-\textbf{ปญฺห}}
\proseitem{16}{\csromanfolio{94}ราชา อาห “ภนฺเต นาคเสน ทุกฺกรํ นุ โข ภควตา กตนฺ”ติ. เถโร อาห “ทุกฺกรํ มหาราช ภควตา กตนฺ”ติ. กิํ ปน ภนฺเต นาคเสน ภควตา ทุกฺกรํ กตนฺติ. ทุกฺกรํ มหาราช ภควตา กตํ อิเมสํ อรูปีนํ จิตฺต\-เจตสิกานํ ธมฺมานํ เอการมฺมเณ วตฺตมานานํ ววตฺถานํ อกฺขาตํ “อยํ ผสฺโส, อยํ เวทนา, อยํ สญฺญา, อยํ เจตนา, อิทํ จิตฺตนฺ”ติ.}

\prose{โอปมฺมํ กโรหีติ. ยถา มหาราช โกจิเทว ปุริโส นาวาย มหาสมุทฺทํ อชฺโฌคาเหตฺวา หตฺถปุเฏน อุทกํ คเหตฺวา ชิวฺหาย สายิตฺวา ชาเนยฺย นุ โข มหาราช โส ปุริโส “อิทํ คงฺคาย อุทกํ, อิทํ ยมุนาย อุทกํ, อิทํ อจิรวติยา อุทกํ, อิทํ สรภุยา อุทกํ, อิทํ มหิยา อุทกนฺ”ติ. ทุกฺกรํ ภนฺเต ชานิตุนฺติ. อิโต ทุกฺกรตรํ โข มหาราช ภควตา กตํ อิเมสํ อรูปีนํ จิตฺต\-เจตสิกานํ ธมฺมานํ เอการมฺมเณ วตฺตมานานํ ววตฺถานํ อกฺขาตํ “อยํ ผสฺโส, อยํ เวทนา, อยํ สญฺญา, อยํ เจตนา, อิทํ จิตฺตนฺ”ติ. สุฏฺฐุ ภนฺเตติ ราชา อพฺภานุโมทีติ.}

\vspace{\tipitakaparskip}
\csromancenter{อรูป\-ธมฺมววตฺถาน\-ทุกฺกร\-ปโญฺห โสฬสโม.}
\nitthitammidruled{\textbf{อรูป}\-\textbf{ธมฺมววตฺถาน}\-\textbf{วคฺโค สตฺตโม.}}
\csromancenter{อิมสฺมิํ วคฺเค โสฬส ปญฺหา.}
\csromanmatikaanchor{mtk.97}
\csromantocmarkref{subsubsection}{มิลินฺทปญฺหปุจฺฉาวิสชฺชนา}{mtk.97}
\bookmark[dest={mtk.97},level=3]{มิลินฺทปญฺหปุจฺฉาวิสชฺชนา}
\csromanheader{3}{2. มิลินฺทปญฺห \textbf{มิลินฺทปญฺห}\-\textbf{ปุจฺฉา}\-\textbf{วิส}\-\textbf{ชฺชนา}}
\prose{\csromanfolio{95}เถโร อาห “ชานาสิ โข มหาราช สมฺปติ กา เวลา”ติ. อาม ภนฺเต ชานามิ “สมฺปติ ปฐโม ยาโม อติกฺกนฺโต, มชฺฌิโม ยาโม ปวตฺตติ, อุกฺกา ปทีปียนฺติ, จตฺตาริ ปฏากานิ อาณตฺตานิ คมิสฺสนฺติ ภณฺฑโต ราชเทยฺยานี”ติ.}

\prose{โยนกา เอวมาหํสุ “กลฺโลสิ มหาราช ปณฺฑิโต เถโร”ติ. อาม ภเณ ปณฺฑิโต เถโร, เอทิโส อาจริโย ภเวยฺย มาทิโส จ อนฺเตวาสี, นจิรสฺเสว ปณฺฑิโต ธมฺมํ อาชาเนยฺยาติ. ตสฺส ปญฺห\-เวยฺยากรเณน ตุฏฺโฐ ราชา เถรํ นาคเสนํ สตสหสฺส\-คฺฆนเกน กมฺพเลน อจฺฉาเทตฺวา “ภนฺเต นาคเสน อชฺชตคฺเค เต อฏฺฐสตํ ภตฺตํ ปญฺญเปมิ, ยํ กิญฺจิ อนฺเตปุเร กปฺปิยํ, เตน จ ปวาเรมี”ติ อาห. อลํ มหาราช ชีวามีติ. ชานามิ ภนฺเต นาคเสน ชีวสิ, อปิ จ อตฺตานญฺจ รกฺข, มมญฺจ รกฺขาหีติ. กถํ อตฺตานํ รกฺขสิ, “นาคเสโน มิลินฺทํ ราชานํ ปสาเทติ, น จ กิญฺจิ อลภี”ติ ปราปวาโท\csromansharedfootnote{5b80c885a3eacf27}{ปรปฺปวาโท (ก)} อาคจฺเฉยฺยาติ, เอวํ อตฺตานํ รกฺข. กถํ มมํ รกฺขสิ, “มิลินฺโท ราชา ปสนฺโน ปสนฺนาการํ น กโรตี”ติ ปราปวาโท อาคจฺเฉยฺยาติ, เอวํ มมํ รกฺขาหีติ. ตถา โหตุ มหาราชาติ. เสยฺยถาปิ ภนฺเต สีโห มิคราชา สุวณฺณปญฺชเร ปกฺขิตฺโตปิ พหิมุโข เยว โหติ, เอวเมว โข อหํ ภนฺเต กิญฺจาปิ อคารํ อชฺฌาวสามิ พหิมุโข เยว ปน อจฺฉามิ. สเจ อหํ ภนฺเต อคารสฺมา อนาคาริยํ ปพฺพเชยฺยํ, น จิรํ ชีเวยฺยํ, พหู เม ปจฺจตฺถิกาติ.}

\prose{อถ โข อายสฺมา นาคเสโน มิลินฺทสฺส รญฺโญ ปญฺหํ วิสชฺเชตฺวา อุฏฺฐายาสนา สํฆารามํ อคมาสิ. อจิรปกฺกนฺเต จ อายสฺมนฺเต นาคเสเน มิลินฺทสฺส รญฺโญ เอตทโหสิ “กิํ มยา ปุจฺฉิตํ, กิํ ภทนฺเตน นาคเสเนน วิสชฺชิตนฺ”ติ. อถ โข มิลินฺทสฺส รญฺโญ เอตทโหสิ “สพฺพํ มยา สุปุจฺฉิตํ, สพฺพํ ภทนฺเตน นาคเสเนน สุวิสชฺชิตนฺ”ติ. อายสฺมโตปิ นาคเสนสฺส สํฆารามคตสฺส เอตทโหสิ “กิํ}

\prose{\csromanfolio{96}มิลินฺเทน รญฺญา ปุจฺฉิตํ, กิํ มยา วิสชฺชิตนฺ”ติ. อถ โข อายสฺมโต นาคเสนสฺส เอตทโหสิ “สพฺพํ มิลินฺเทน รญฺญา สุปุจฺฉิตํ, สพฺพํ มยา สุวิสชฺชิตนฺ”ติ.}

\prose{อถ โข อายสฺมา นาคเสโน ตสฺสา รตฺติยา อจฺจเยน ปุพฺพณฺหสมยํ นิวาเสตฺวา ปตฺต\-จีวรมา\-ทาย เยน มิลินฺทสฺส รญฺโญ นิเวสนํ เตนุปสงฺกมิ, อุปสงฺกมิตฺวา ปญฺญตฺเต อาสเน นิสีทิ. อถ โข มิลินฺโท ราชา อายสฺมนฺตํ นาคเสนํ อภิวาเทตฺวา เอกมนฺตํ นิสีทิ, เอกมนฺตํ นิสินฺโน โข มิลินฺโท ราชา อายสฺมนฺตํ นาคเสนํ เอตทโวจ\csromandash{}}

\prose{มา โข ภทนฺตสฺส เอวํ อโหสิ “นาคเสโน มยา ปญฺหํ ปุจฺฉิโต”ติ เตเนว โสมนสฺเสน ตํ รตฺตาวเสสํ วีตินาเมสีติ น เต เอวํ ทฏฺฐพฺพํ. ตสฺส มยฺหํ ภนฺเต ตํ รตฺตาวเสสํ เอตทโหสิ “กิํ มยา ปุจฺฉิตํ, กิํ ภทนฺเตน วิสชฺชิตนฺ”ติ, “สพฺพํ มยา สุปุจฺฉิตํ, สพฺพํ ภทนฺเตน สุวิสชฺชิตนฺ”ติ.}

\prosewithcloser{เถโรปิ เอวมาห\csromandash{}มา โข มหาราชสฺส เอวํ อโหสิ “มิลินฺทสฺส รญฺโญ มยา ปโญฺห วิสชฺชิโต”ติ เตเนว โสมนสฺเสน ตํ รตฺตาวเสสํ วีตินาเมสีติ น เต เอวํ ทฏฺฐพฺพํ. ตสฺส มยฺหํ มหาราช ตํ รตฺตาวเสสํ เอตทโหสิ “กิํ มิลินฺเทน รญฺญา ปุจฺฉิตํ, กิํ มยา วิสชฺชิตนฺ”ติ, “สพฺพํ มิลินฺเทน รญฺญา สุปุจฺฉิตํ, สพฺพํ มยา สุวิสชฺชิตนฺ”ติ อิติห เต มหานาคา อญฺญมญฺญสฺส สุภาสิตํ สมนุ\-โมทิํสูติ.}{\textbf{มิลินฺทปญฺห}\-\textbf{ปุจฺฉา}\-\textbf{วิส}\-\textbf{ชฺชนา นิฏฺฐิตา.}}
\csromanmatikaanchor{mtk.98}
\csromantocmarknopage{section}{เมณฺฑกปญฺหารมฺภกถา}{mtk.98}
\bookmark[dest={mtk.98},level=1]{เมณฺฑกปญฺหารมฺภกถา}
\csromanheader{1}{\textbf{เมณฺฑก}\-\textbf{ปญฺหา}\-\textbf{รมฺภ}\-\textbf{กถา}}
\csromanmatikaanchor{mtk.99}
\csromantocmarkref{subsection}{อฏฺฐมนฺตปริวชฺชนียฏฺฐาน}{mtk.99}
\bookmark[dest={mtk.99},level=2]{อฏฺฐมนฺตปริวชฺชนียฏฺฐาน}
\csromanheader{2}{\textbf{อฏฺฐ}\-\textbf{มนฺต}\-\textbf{ปริวชฺชนีย}\-\textbf{ฏฺฐาน}}
\setlength{\csromangathapagewidth}{0pt}
\setlength{\csromangathawakwidth}{0pt}
\csromangathameasuregroup{ภสฺสปฺปวาโท\textsuperscript{1} เวตณฺฑี, \\ มิลินฺโท ญาณเภทาย, \\ วสนฺโต ตสฺส ฉายาย, \\ ปหินฺนพุทฺธิ หุตฺวาน, \\ นวงฺคํ อนุมชฺชนฺโต, \\ อทฺทกฺขิ เมณฺฑเก ปเญฺห, \\ ปริยายภาสิตํ อตฺถิ, \\ สภาวภาสิตํ อตฺถิ, \\ เตสมตฺถํ อวิญฺญาย, \\ อนาคตมฺหิ อทฺธาเน, \\ หนฺท กถิํ ปสาเทตฺวา, \\ ตสฺส นิทฺทิฏฺฐมคฺเคน,}{\csromangathabat{ภสฺสปฺปวาโท\textsuperscript{1} เวตณฺฑี,}{อติพุทฺธิ วิจกฺขโณ.} \\ \csromangathabat{มิลินฺโท ญาณเภทาย,}{นาคเสนมุปาคมิ.} \\ \csromangathabat{วสนฺโต ตสฺส ฉายาย,}{ปริปุจฺฉํ ปุนปฺปุนํ.} \\ \csromangathabat{ปหินฺนพุทฺธิ หุตฺวาน,}{โสปิ อาสิ ติเปฏโก.} \\ \csromangathabat{นวงฺคํ อนุมชฺชนฺโต,}{รตฺติภาเค รโหคโต.} \\ \csromangathabat{อทฺทกฺขิ เมณฺฑเก ปเญฺห,}{ทุนฺนิเวเฐ สนิคฺคเห.} \\ \csromangathabat{ปริยายภาสิตํ อตฺถิ,}{อตฺถิ สนฺธายภาสิตํ.} \\ \csromangathabat{สภาวภาสิตํ อตฺถิ,}{ธมฺมราชสฺส สาสเน.} \\ \csromangathabat{เตสมตฺถํ อวิญฺญาย,}{เมณฺฑเก ชินภาสิเต.} \\ \csromangathabat{อนาคตมฺหิ อทฺธาเน,}{วิคฺคโห ตตฺถ เหสฺสติ.} \\ \csromangathabat{หนฺท กถิํ ปสาเทตฺวา,}{เฉชฺชาเปสฺสามิ เมณฺฑเก.} \\ \csromangathabat{ตสฺส นิทฺทิฏฺฐมคฺเคน,}{นิทฺทิสิสฺสนฺตฺยนาคเตติ.}}
\csromangathasetleft{ภสฺสปฺปวาโท\textsuperscript{1} เวตณฺฑี, \\ มิลินฺโท ญาณเภทาย, \\ วสนฺโต ตสฺส ฉายาย, \\ ปหินฺนพุทฺธิ หุตฺวาน, \\ นวงฺคํ อนุมชฺชนฺโต, \\ อทฺทกฺขิ เมณฺฑเก ปเญฺห, \\ ปริยายภาสิตํ อตฺถิ, \\ สภาวภาสิตํ อตฺถิ, \\ เตสมตฺถํ อวิญฺญาย, \\ อนาคตมฺหิ อทฺธาเน, \\ หนฺท กถิํ ปสาเทตฺวา, \\ ตสฺส นิทฺทิฏฺฐมคฺเคน,}
\csromangathagroup{\csromangathastanza{\csromanfolio{97}\csromangathabat{ภสฺสปฺปวาโท\csromansharedfootnote{4c58ccab34a8e97c}{ภสฺสปฺปเวที (สี, อิ)} เวตณฺฑี,}{อติพุทฺธิ วิจกฺขโณ.} \\ \csromangathabat{มิลินฺโท ญาณเภทาย,}{นาคเสนมุปาคมิ.}}\csromangathastanza{\csromangathabat{วสนฺโต ตสฺส ฉายาย,}{ปริปุจฺฉํ ปุนปฺปุนํ.} \\ \csromangathabat{ปหินฺนพุทฺธิ หุตฺวาน,}{โสปิ อาสิ ติเปฏโก.}}\csromangathastanza{\csromangathabat{นวงฺคํ อนุมชฺชนฺโต,}{รตฺติภาเค รโหคโต.} \\ \csromangathabat{อทฺทกฺขิ เมณฺฑเก ปเญฺห,}{ทุนฺนิเวเฐ สนิคฺคเห.}}\csromangathastanza{\csromangathabat{ปริยายภาสิตํ อตฺถิ,}{อตฺถิ สนฺธายภาสิตํ.} \\ \csromangathabat{สภาวภาสิตํ อตฺถิ,}{ธมฺมราชสฺส สาสเน.}}\csromangathastanza{\csromangathabat{เตสมตฺถํ อวิญฺญาย,}{เมณฺฑเก ชินภาสิเต.} \\ \csromangathabat{อนาคตมฺหิ อทฺธาเน,}{วิคฺคโห ตตฺถ เหสฺสติ.}}\csromangathastanza{\csromangathabat{หนฺท กถิํ ปสาเทตฺวา,}{เฉชฺชาเปสฺสามิ เมณฺฑเก.} \\ \csromangathabat{ตสฺส นิทฺทิฏฺฐมคฺเคน,}{นิทฺทิสิสฺสนฺตฺยนาคเตติ.}}}

\prose{อถ โข มิลินฺโท ราชา ปภาตาย รตฺติยา อุทฺธสฺเต\csromansharedfootnote{47cf1b4f5db2a130}{อุฏฺฐิเต (สฺยา), อุคฺคเต (สี, อิ)} อรุเณ สีสํ นฺหตฺวา สิรสิ อญฺชลิํ ปคฺคเหตฺวา อตีตา\-นาคต\-ปจฺจุปฺปนฺเน สมฺมาสมฺพุทฺเธ อนุสฺสริตฺวา อฏฺฐ วตฺตปทานิ สมาทิยิ “อิโต เม อนาคตานิ สตฺต ทิวสานิ อฏฺฐ คุเณ สมาทิยิตฺวา ตโป จริตพฺโพ ภวิสฺสติ, โสหํ จิณฺณตโป สมาโน อาจริยํ อาราเธตฺวา เมณฺฑเก ปเญฺห ปุจฺฉิสฺสามี”ติ. อถ โข มิลินฺโท ราชา ปกติ\-ทุสฺสยุคํ อปเนตฺวา อาภรณานิ จ โอมุญฺจิตฺวา กาสาวํ นิวาเสตฺวา มุณฺฑกปฏิสีสกํ สีเส ปฏิมุญฺจิตฺวา มุนิภาวมุ\-ปคนฺตฺวา อฏฺฐ คุเณ สมาทิยิ “อิมํ สตฺตาหํ มยา น ราชตฺโถ อนุสาสิตพฺโพ, น ราคูปสญฺหิตํ จิตฺตํ อุปฺปาเทตพฺพํ, น โทสูปสญฺหิตํ จิตฺตํ อุปฺปาเทตพฺพํ, น โมหูปสญฺหิตํ จิตฺตํ อุปฺปาเทตพฺพํ, ทาส\-กมฺมกร\-โปริเส ชเนปิ นิวาตวุตฺตินา ภวิตพฺพํ, กายิกํ}

\prose{\csromanfolio{98}วาจสิกํ อนุรกฺขิตพฺพํ, ฉปิ อายตนานิ นิรวเสสโต อนุรกฺขิตพฺพานิ, เมตฺตาภาวนาย มานสํ ปกฺขิปิตพฺพนฺ”ติ. อิเม อฏฺฐ คุเณ สมาทิยิตฺวา เตเสฺวว อฏฺฐสุ คุเณสุ มานสํ ปติฏฺฐเปตฺวา พหิ อนิกฺขมิตฺวา สตฺตาหํ วีตินาเมตฺวา อฏฺฐเม ทิวเส ปภาตาย รตฺติยา ปเคว ปาตราสํ กตฺวา โอกฺขิตฺตจกฺขุ มิตภาณี สุสณฺฐิเตน อิริยาปเถน อวิกฺขิตฺเตน จิตฺเตน หฏฺเฐน อุทคฺเคน วิปฺปสนฺเนน เถรํ นาคเสนํ อุปสงฺกมิตฺวา เถรสฺส ปาเท สิรสา วนฺทิตฺวา เอกมนฺตํ ฐิโต อิทมโวจ\csromandash{}}

\prose{อตฺถิ เม ภนฺเต นาคเสน โกจิ อตฺโถ ตุเมฺหหิ สทฺธิํ มนฺตยิตพฺโพ, น ตตฺถ อญฺโญ โกจิ ตติโย อิจฺฉิตพฺโพ, สุญฺเญ โอกาเส ปวิวิตฺเต อรญฺเญ อฏฺฐงฺคุปาคเต สมณสารุปฺเป. ตตฺถ โส ปโญฺห ปุจฺฉิตพฺโพ ภวิสฺสติ, ตตฺถ เม คุยฺหํ น กาตพฺพํ น รหสฺสกํ, อรหามหํ รหสฺสกํ สุณิตุํ สุมนฺตเน อุปคเต, อุปมายปิ โส อตฺโถ อุปปริ\-กฺขิ\-ตพฺโพ, ยถา กิํวิย, ยถา นาม ภนฺเต นาคเสน มหาปถวี นิกฺเขปํ อรหติ นิกฺเขเป อุปคเต. เอวเมว โข ภนฺเต นาคเสน อรหามหํ รหสฺสกํ สุณิตุํ สุมนฺตเน อุปคเตติ.}

\prose{ครุนา สห ปวิวิตฺต\-ปวนํ ปวิสิตฺวา อิทมโวจ\csromandash{}ภนฺเต นาคเสน อิธ ปุริเสน มนฺตยิตุกาเมน อฏฺฐ ฐานานิ ปริวชฺช\-ยิตพฺพานิ ภวนฺติ, น เตสุ ฐาเนสุ วิญฺญู ปุริโส อตฺถํ มนฺเตติ, มนฺติโตปิ อตฺโถ ปริปตติ น สมฺภวติ. กตมานิ อฏฺฐ ฐานานิ, วิสมฏฺฐานํ ปริวชฺชนียํ, สภยํ ปริวชฺชนียํ, อติวาตฏฺฐานํ ปริวชฺชนียํ, ปฏิจฺฉนฺน\-ฏฺฐานํ ปริวชฺชนียํ, เทวฏฺฐานํ ปริวชฺชนียํ, ปนฺโถ ปริวชฺชนีโย, สงฺคาโม\csromansharedfootnote{7ace64e15e7b9288}{สงฺกโม (สี, อิ)} ปริวชฺชนีโย, อุทกติตฺถํ ปริวชฺชนียํ. อิมานิ อฏฺฐ ฐานานิ ปริวชฺชนียานีติ.}

\prose{เถโร อาห “โก โทโส วิสมฏฺฐาเน, สภเย, อติวาเต, ปฏิจฺฉนฺเน, เทวฏฺฐาเน, ปนฺเถ, สงฺคาเม, อุทกติตฺเถ”ติ. วิสเม ภนฺเต นาคเสน มนฺติโต อตฺโถ วิกิรติ วิธมติ ปคฺฆรติ น สมฺภวติ, สภเย มโน สนฺตสฺสติ, สนฺตสฺสิโต น สมฺมา อตฺถํ สมนุปสฺสติ, อติวาเต สทฺโท อวิภูโต โหติ, ปฏิจฺฉนฺเน อุปสฺสุติํ ติฏฺฐนฺติ,}

\prose{\csromanfolio{99}เทวฏฺฐาเน มนฺติโต อตฺโถ ครุกํ ปริณมติ, ปนฺเถ มนฺติโต อตฺโถ ตุจฺโฉ ภวติ, สงฺคาเม จญฺจโล ภวติ, อุทกติตฺเถ ปากโฏ ภวติ. ภวตีห\csromandash{}}

\setlength{\csromangathapagewidth}{0pt}
\setlength{\csromangathawakwidth}{0pt}
\csromangathameasuregroup{“วิสมํ สภยํ อติวาโต, \\ ปนฺโถ จ สงฺคาโม ติตฺถํ,}{\csromangathabat{“วิสมํ สภยํ อติวาโต,}{ปฏิจฺฉนฺนํ เทวนิสฺสิตํ.} \\ \csromangathabat{ปนฺโถ จ สงฺคาโม ติตฺถํ,}{อฏฺเฐเต ปริวชฺชิยา”ติ.} \\ อฏฺฐ มนฺตนสฺส ปริวชฺชนียฏฺฐานานิ.}
\csromangathameasurewak{อฏฺฐ มนฺตนสฺส ปริวชฺชนียฏฺฐานานิ.}
\csromangathasetleft{“วิสมํ สภยํ อติวาโต, \\ ปนฺโถ จ สงฺคาโม ติตฺถํ,}
\csromangathagroup{\csromangathastanza{\csromangathabat{“วิสมํ สภยํ อติวาโต,}{ปฏิจฺฉนฺนํ เทวนิสฺสิตํ.} \\ \csromangathabat{ปนฺโถ จ สงฺคาโม ติตฺถํ,}{อฏฺเฐเต ปริวชฺชิยา”ติ.}}\csromangathastanza{\csromangathawak{อฏฺฐ มนฺตนสฺส ปริวชฺชนียฏฺฐานานิ.}}}

\csromanmatikaanchor{mtk.100}
\csromantocmarkref{subsection}{อฏฺฐมนฺตวินาสกปุคฺคล}{mtk.100}
\bookmark[dest={mtk.100},level=2]{อฏฺฐมนฺตวินาสกปุคฺคล}
\csromanheader{2}{\textbf{อฏฺฐ}\-\textbf{มนฺต}\-\textbf{วินาสก}\-\textbf{ปุคฺคล}}
\prose{ภนฺเต นาคเสน อฏฺฐิเม ปุคฺคลา มนฺติยมานา มนฺติตํ อตฺถํ พฺยาปาเทนฺติ. กตเม อฏฺฐ, ราคจริโต โทสจริโต โมหจริโต มานจริโต ลุทฺโธ อลโส เอกจินฺตี พาโลติ. อิเม อฏฺฐ ปุคฺคลา มนฺติตํ อตฺถํ พฺยาปาเทนฺตีติ.}

\prose{เถโร อาห “เตสํ โก โทโส”ติ. ราคจริโต ภนฺเต นาคเสน ราควเสน มนฺติตํ อตฺถํ พฺยาปาเทติ, โทสจริโต โทสวเสน มนฺติตํ อตฺถํ พฺยาปาเทติ, โมหจริโต โมหวเสน มนฺติตํ อตฺถํ พฺยาปาเทติ, มานจริโต มานวเสน มนฺติตํ อตฺถํ พฺยาปาเทติ, ลุทฺโธ โลภวเสน มนฺติตํ อตฺถํ พฺยาปาเทติ, อลโส อลสตาย มนฺติตํ อตฺถํ พฺยาปาเทติ, เอกจินฺตี เอกจินฺติตาย มนฺติตํ อตฺถํ พฺยาปาเทติ, พาโล พาลตาย มนฺติตํ อตฺถํ พฺยาปาเทติ. ภวตีห\csromandash{}}

\setlength{\csromangathapagewidth}{0pt}
\setlength{\csromangathawakwidth}{0pt}
\csromangathameasuregroup{“รตฺโต ทุฏฺโฐ จ มูโฬ จ, \\ เอกจินฺตี จ พาโล จ,}{\csromangathabat{“รตฺโต ทุฏฺโฐ จ มูโฬ จ,}{มานี ลุทฺโธ ตถาลโส.} \\ \csromangathabat{เอกจินฺตี จ พาโล จ,}{เอเต อตฺถวินาสกา”ติ.} \\ อฏฺฐ มนฺตวินาสกปุคฺคลา.}
\csromangathameasurewak{อฏฺฐ มนฺตวินาสกปุคฺคลา.}
\csromangathasetleft{“รตฺโต ทุฏฺโฐ จ มูโฬ จ, \\ เอกจินฺตี จ พาโล จ,}
\csromangathagroup{\csromangathastanza{\csromangathabat{“รตฺโต ทุฏฺโฐ จ มูโฬ จ,}{มานี ลุทฺโธ ตถาลโส.} \\ \csromangathabat{เอกจินฺตี จ พาโล จ,}{เอเต อตฺถวินาสกา”ติ.}}\csromangathastanza{\csromangathawak{อฏฺฐ มนฺตวินาสกปุคฺคลา.}}}

\csromanmatikaanchor{mtk.101}
\csromantocmarkref{subsection}{นวคุยฺหมนฺตวิธํสก}{mtk.101}
\bookmark[dest={mtk.101},level=2]{นวคุยฺหมนฺตวิธํสก}
\csromanheader{2}{\textbf{นวคุยฺหมนฺต}\-\textbf{วิธํสก}}
\prose{ภนฺเต นาคเสน นวิเม ปุคฺคลา มนฺติตํ คุยฺหํ วิวรนฺติ น ธาเรนฺติ. กตเม นว, ราคจริโต โทสจริโต โมหจริโต ภีรุโก อามิสครุโก อิตฺถี โสณฺโฑ ปณฺฑโก ทารโกติ.}

\prose{เถโร อาห “เตสํ โก โทโส”ติ. ราคจริโต ภนฺเต นาคเสน ราควเสน มนฺติตํ คุยฺหํ วิวรติ น ธาเรติ, โทสจริโต}

\csromanmatikaanchor{mtk.102}
\csromantocmarkref{subsection}{อฏฺฐปญฺญาปฏิลาภการณ}{mtk.102}
\bookmark[dest={mtk.102},level=2]{อฏฺฐปญฺญาปฏิลาภการณ}
\prose{\csromanfolio{100}ภนฺเต โทสวเสน มนฺติตํ คุยฺหํ วิวรติ น ธาเรติ, มูโฬฺห โมหวเสน มนฺติตํ คุยฺหํ วิวรติ น ธาเรติ, ภีรุโก ภยวเสน มนฺติตํ คุยฺหํ วิวรติ น ธาเรติ, อามิสครุโก อามิสเหตุ มนฺติตํ คุยฺหํ วิวรติ น ธาเรติ, อิตฺถี ปญฺญาย อิตฺตรตาย มนฺติตํ คุยฺหํ วิวรติ น ธาเรติ, โสณฺฑิโก สุราโลลตาย มนฺติตํ คุยฺหํ วิวรติ น ธาเรติ, ปณฺฑโก อเนกํสิกตาย มนฺติตํ คุยฺหํ วิวรติ น ธาเรติ, ทารโก จปลตาย มนฺติตํ คุยฺหํ วิวรติ น ธาเรติ. ภวตีห\csromandash{}}

\setlength{\csromangathapagewidth}{0pt}
\setlength{\csromangathawakwidth}{0pt}
\csromangathameasuregroup{“รตฺโต ทุฏฺโฐ จ มูโฬฺห จ, \\ อิตฺถี โสณฺโฑ ปณฺฑโก จ, \\ นเวเต ปุคฺคลา โลเก, \\ เอเตหิ มนฺติตํ คุยฺหํ,}{\csromangathabat{“รตฺโต ทุฏฺโฐ จ มูโฬฺห จ,}{ภีรุ อามิสครุโก\textsuperscript{1}.} \\ \csromangathabat{อิตฺถี โสณฺโฑ ปณฺฑโก จ,}{นวโม ภวติ ทารโก.} \\ \csromangathabat{นเวเต ปุคฺคลา โลเก,}{อิตฺตรา จลิตา จลา.} \\ \csromangathabat{เอเตหิ มนฺติตํ คุยฺหํ,}{ขิปฺปํ ภวติ ปากฏนฺ”ติ.} \\ นว คุยฺหมนฺตวิธํสกา ปุคฺคลา.}
\csromangathameasurewak{นว คุยฺหมนฺตวิธํสกา ปุคฺคลา.}
\csromangathasetleft{“รตฺโต ทุฏฺโฐ จ มูโฬฺห จ, \\ อิตฺถี โสณฺโฑ ปณฺฑโก จ, \\ นเวเต ปุคฺคลา โลเก, \\ เอเตหิ มนฺติตํ คุยฺหํ,}
\csromangathagroup{\csromangathastanza{\csromangathabat{“รตฺโต ทุฏฺโฐ จ มูโฬฺห จ,}{ภีรุ อามิสครุโก\csromansharedfootnote{1f15952fea36985e}{อามิสจกฺขุโก (สี, อิ)}.} \\ \csromangathabat{อิตฺถี โสณฺโฑ ปณฺฑโก จ,}{นวโม ภวติ ทารโก.}}\csromangathastanza{\csromangathabat{นเวเต ปุคฺคลา โลเก,}{อิตฺตรา จลิตา จลา.} \\ \csromangathabat{เอเตหิ มนฺติตํ คุยฺหํ,}{ขิปฺปํ ภวติ ปากฏนฺ”ติ.}}\csromangathastanza{\csromangathawak{นว คุยฺหมนฺตวิธํสกา ปุคฺคลา.}}}

\csromanheader{2}{\textbf{อฏฺฐ ปญฺญาปฏิลาภการณ}}
\prose{ภนฺเต นาคเสน อฏฺฐหิ การเณหิ พุทฺธิ ปริณมติ ปริปากํ คจฺฉติ. กตเมหิ อฏฺฐหิ, วยปริณาเมน พุทฺธิ ปริณมติ ปริปากํ คจฺฉติ, ยสปริณาเมน พุทฺธิ ปริณมติ ปริปากํ คจฺฉติ, ปริปุจฺฉาย พุทฺธิ ปริณมติ ปริปากํ คจฺฉติ, ติตฺถ\-สํวาเสน พุทฺธิ ปริณมติ ปริปากํ คจฺฉติ, โยนิโส มนสิกาเรน พุทฺธิ ปริณมติ ปริปากํ คจฺฉติ, สากจฺฉาย พุทฺธิ ปริณมติ ปริปากํ คจฺฉติ, สฺเนหู\-ปเสว\-เนน พุทฺธิ ปริณมติ ปริปากํ คจฺฉติ, ปติรูป\-เทส\-วาเสน พุทฺธิ ปริณมติ ปริปากํ คจฺฉติ. ภวตีห\csromandash{}}

\setlength{\csromangathapagewidth}{0pt}
\setlength{\csromangathawakwidth}{0pt}
\csromangathameasuregroup{“วเยน ยสปุจฺฉาหิ, \\ สากจฺฉา สฺเนหสํเสวา, \\ เอตานิ อฏฺฐ ฐานานิ, \\ เยสํ เอตานิ สมฺโภนฺติ,}{\csromangathabat{“วเยน ยสปุจฺฉาหิ,}{ติตฺถวาเสน โยนิโส.} \\ \csromangathabat{สากจฺฉา สฺเนหสํเสวา,}{ปติรูปวเสน จ.} \\ \csromangathabat{เอตานิ อฏฺฐ ฐานานิ,}{พุทฺธิวิสทการณา.} \\ \csromangathabat{เยสํ เอตานิ สมฺโภนฺติ,}{เตสํ พุทฺธิ ปภิชฺชตี”ติ.} \\ อฏฺฐ ปญฺญาปฏิลาภการณานิ.}
\csromangathameasurewak{อฏฺฐ ปญฺญาปฏิลาภการณานิ.}
\csromangathasetleft{“วเยน ยสปุจฺฉาหิ, \\ สากจฺฉา สฺเนหสํเสวา, \\ เอตานิ อฏฺฐ ฐานานิ, \\ เยสํ เอตานิ สมฺโภนฺติ,}
\csromangathagroup{\csromangathastanza{\csromangathabat{“วเยน ยสปุจฺฉาหิ,}{ติตฺถวาเสน โยนิโส.} \\ \csromangathabat{สากจฺฉา สฺเนหสํเสวา,}{ปติรูปวเสน จ.}}\csromangathastanza{\csromangathabat{เอตานิ อฏฺฐ ฐานานิ,}{พุทฺธิวิสทการณา.} \\ \csromangathabat{เยสํ เอตานิ สมฺโภนฺติ,}{เตสํ พุทฺธิ ปภิชฺชตี”ติ.}}\csromangathastanza{\csromangathawak{อฏฺฐ ปญฺญาปฏิลาภการณานิ.}}}

\csromanmatikaanchor{mtk.103}
\csromantocmarkref{subsection}{อาจริยคุณ}{mtk.103}
\bookmark[dest={mtk.103},level=2]{อาจริยคุณ}
\csromanheader{2}{เมณฺฑก\-ปญฺหา\-รมฺภ\-กถา \textbf{อาจริยคุณ}}
\prose{\csromanfolio{101}ภนฺเต นาคเสน อยํ ภูมิภาโค อฏฺฐ มนฺถโทสวิวชฺชิโต, อหญฺจ โลเก ปรโม มนฺติสหาโย\csromansharedfootnote{a1ce0145b57d860c}{มนฺตสหาโย (สี)}, คุยฺหม\-นุรกฺขี จาหํ ยาวาหํ ชีวิสฺสามิ ตาว คุยฺหม\-นุรกฺขิสฺสามิ, อฏฺฐหิ จ เม การเณหิ พุทฺธิ ปริณามํ คตา, ทุลฺลโภ เอตรหิ มาทิโส อนฺเตวาสี, สมฺมา ปฏิปนฺเน อนฺเตวาสิเก เย อาจริยานํ ปญฺจวีสติ อาจริยคุณา, เตหิ คุเณหิ อาจริเยน สมฺมา ปฏิปชฺชิตพฺพํ. กตเม ปญฺจวีสติ คุณา.}

\prose{อิธ ภนฺเต นาคเสน อาจริเยน อนฺเตวาสิมฺหิ สตตํ สมิตํ อารกฺขา อุปฏฺฐเปตพฺพา, อเสวนเสวนา ชานิตพฺพา, ปมตฺตา\-ปฺปมตฺตา ชานิตพฺพา, เสยฺยาวกาโส ชานิตพฺโพ, เคลญฺญํ ชานิตพฺพํ, โภชนสฺส\csromansharedfootnote{61ec5c7427e09383}{โภชนียํ (สฺยา)} ลทฺธาลทฺธํ ชานิตพฺพํ, วิเสโส ชานิตพฺโพ, ปตฺตคตํ สํวิภชิตพฺพํ, อสฺสาสิตพฺโพ “มา ภายิ, อตฺโถ เต อภิกฺกมตี”ติ, “อิมินา ปุคฺคเลน ปฏิจรตี”ติ\csromansharedfootnote{d857ca591a37d78f}{ปฏิจราหีติ (ก)} ปฏิจาโร ชานิตพฺโพ, คาเม ปฏิจาโร ชานิตพฺโพ, วิหาเร ปฏิจาโร ชานิตพฺโพ, น เตน หาโส ทโว กาตพฺโพ, เตน สห อาลาโป กาตพฺโพ\csromansharedfootnote{6b7a4c89e6e889de}{น เตน สห สลฺลาโป กาตพฺโพ (สี, อิ)}, ฉิทฺทํ ทิสฺวา อธิวาเสตพฺพํ, สกฺกจฺจการินา ภวิตพฺพํ, อขณฺฑการินา ภวิตพฺพํ, อรหสฺสการินา ภวิตพฺพํ, นิรวเสส\-การินา ภวิตพฺพํ, “ชเนมิมํ\csromansharedfootnote{ec259be8b81ce81b}{ชาเนมิมํ (สฺยา)} สิปฺเปสู”ติ ชนกจิตฺตํ อุปฏฺฐเป\-ตพฺพํ, “กถํ อยํ น ปริหาเยยฺยา”ติ วฑฺฒิจิตฺตํ อุปฏฺฐเป\-ตพฺพํ, “พลวํ อิมํ กโรมิ สิกฺขาพเลนี”ติ จิตฺตํ อุปฏฺฐเป\-ตพฺพํ, เมตฺตจิตฺตํ อุปฏฺฐเป\-ตพฺพํ, อาปทาสุ น วิชหิตพฺพํ, กรณีเย นปฺปมชฺชิตพฺพํ, ขลิเต ธมฺเมน ปคฺคเหตพฺโพติ. อิเม โข ภนฺเต ปญฺจวีสติ อาจริยสฺส อาจริยคุณา, เตหิ คุเณหิ มยิ สมฺมา ปฏิปชฺชสฺสุ, สํสโย เม ภนฺเต อุปฺปนฺโน, อตฺถิ เมณฺฑกปญฺหา ชินภาสิตา, อนาคเต อทฺธาเน ตตฺถ วิคฺคโห อุปฺปชฺชิสฺสติ, อนาคเต จ อทฺธาเน ทุลฺลภา ภวิสฺสนฺติ ตุมฺหาทิสา พุทฺธิมนฺโต, เตสุ เม ปเญฺหสุ จกฺขุํ เทหิ ปรวาทานํ นิคฺคหายาติ.}

\csromanmatikaanchor{mtk.104}
\csromantocmarkref{subsection}{อุปาสกคุณ}{mtk.104}
\bookmark[dest={mtk.104},level=2]{อุปาสกคุณ}
\csromanheader{2}{\textbf{อุปาสกคุณ}}
\prose{เถโร “สาธู”ติ สมฺปฏิจฺฉิตฺวา ทส อุปาสกสฺส อุปาสกคุเณ ปริทีเปสิ. ทส อิเม มหาราช อุปาสกสฺส อุปาสกคุณา. กตเม}

\prosewithcloser{\csromanfolio{102}ทส, อิธ มหาราช อุปาสโก สํเฆน สมาน\-สุข\-ทุกฺโข โหติ, ธมฺมา\-ธิปเตยฺโย โหติ, ยถาพลํ สํวิภาครโต โหติ, ชินสาสน\-ปริหานิํ ทิสฺวา อภิวฑฺฒิยา วายมติ. สมฺมาทิฏฺฐิโก โหติ, อปคต\-โกตูหล\-มงฺคลิโก ชีวิตเหตุปิ น อญฺญํ สตฺถารํ อุทฺทิสติ, กายิกวาจสิก\-ญฺจสฺส รกฺขิตํ โหติ, สมคฺคาราโม โหติ สมคฺครโต, อนุสูยโก โหติ, น จ กุหนวเสน สาสเน จรติ, พุทฺธํ สรณํ คโต โหติ, ธมฺมํ สรณํ คโต โหติ, สํฆํ สรณํ คโต โหติ. อิเม โข มหาราช ทส อุปาสกสฺส อุปาสกคุณา, เต สพฺเพ คุณา ตยิ สํวิชฺชนฺติ, ตํ เต ยุตฺตํ ปตฺตํ อนุจฺฉวิกํ ปติรูปํ ยํ ตฺวํ ชินสาสน\-ปริหานิํ ทิสฺวา อภิวฑฺฒิํ อิจฺฉสิ, กโรมิ เต โอกาสํ, ปุจฺฉ มํ ตฺวํ ยถาสุขนฺติ.}{\textbf{เมณฺฑก}\-\textbf{ปญฺหา}\-\textbf{รมฺภ}\-\textbf{กถา นิฏฺฐิตา.}}
\csromanmatikaanchor{mtk.105}
\csromanmatikaanchor{mtk.106}
\csromanmatikaanchor{mtk.107}
\csromantocmarknopage{section}{4. เมณฺฑกปญฺห}{mtk.105}
\bookmark[dest={mtk.105},level=1]{4. เมณฺฑกปญฺห}
\csromantocmarknopage{subsection}{1. อิทฺธิพลวคฺค}{mtk.106}
\bookmark[dest={mtk.106},level=2]{1. อิทฺธิพลวคฺค}
\csromantocmarkref{subsubsection}{1. กตาธิการสผลปญฺห}{mtk.107}
\bookmark[dest={mtk.107},level=3]{1. กตาธิการสผลปญฺห}
\csromanmarkh{1. อิทฺธิพลวคฺค}\csromanmarkhii{1. กตาธิการ\-สผล\-ปญฺห}\csromanheadernomark{3}{1. \textbf{อิทฺธิพลวคฺค 1. กตาธิการ}\-\textbf{สผล}\-\textbf{ปญฺห}}
\proseitem{1}{\csromanfolio{103}อถ โข มิลินฺโท ราชา กตาวกาโส นิปจฺจ ครุโน ปาเท สิรสิ อญฺชลิํ กตฺวา เอตทโวจ “ภนฺเต นาคเสน อิเม ติตฺถิยา เอวํ ภณนฺติ ‘ยทิ พุทฺโธ ปูชํ สาทิยติ, น ปรินิพฺพุโต พุทฺโธ สํยุตฺโต โลเกน อนฺโตภวิโก โลกสฺมิํ โลกสาธารโณ, ตสฺมา ตสฺส กโต อธิกาโร อวญฺโฌ ภวติ สผโล\csromansharedfootnote{1e6740c3c762f02d}{วญฺโฌ ภวติ อผโล (สี, อิ, ก)}. ยทิ ปรินิพฺพุโต วิสํยุตฺโต โลเกน นิสฺสโฏ สพฺพภเวหิ, ตสฺส ปูชา นุปฺปชฺชติ, ปรินิพฺพุโต น กิญฺจิ สาทิยติ, อสาทิยนฺตสฺส กโต อธิกาโร วญฺโฌ ภวติ อผโล’ติ อุภโต โกฏิโก เอโส ปโญฺห, เนโส วิสโย อปฺปตฺต\-มานสานํ, มหนฺตานํ เยเวโส วิสโย, ภินฺเทตํ ทิฏฺฐิชาลํ เอกํเส ฐปย, ตเวโส ปโญฺห อนุปฺปตฺโต, อนาคตานํ ชินปุตฺตานํ จกฺขุํ เทหิ ปรวาทนิคฺคหายา”ติ.}

\prose{เถโร อาห\csromandash{}ปรินิพฺพุโต มหาราช ภควา, น จ ภควา ปูชํ สาทิยติ, โพธิมูเล เยว ตถาคตสฺส สาทิยนา ปหีนา, กิํ ปน อนุปาทิเสสาย นิพฺพานธาตุยา ปรินิพฺพุตสฺส. ภาสิตมฺเปตํ มหาราช เถเรน สาริปุตฺเตน ธมฺม\-เสนาปตินา\csromandash{}}

\setlength{\csromangathapagewidth}{0pt}
\setlength{\csromangathawakwidth}{0pt}
\csromangathameasuregroup{“ปูชิยนฺตา\textsuperscript{1} อสมสมา, \\ น สาทิยนฺติ สกฺการํ,}{\csromangathabat{“ปูชิยนฺตา\textsuperscript{1} อสมสมา,}{สเทวมานุเสหิ เต.} \\ \csromangathabat{น สาทิยนฺติ สกฺการํ,}{พุทฺธานํ เอส ธมฺมตา”ติ.}}
\csromangathasetleft{“ปูชิยนฺตา\textsuperscript{1} อสมสมา, \\ น สาทิยนฺติ สกฺการํ,}
\csromangathagroup{\csromangathastanza{\csromangathabat{“ปูชิยนฺตา\csromansharedfootnote{3edc1101668a8e46}{ปูชิตา (สฺยา)} อสมสมา,}{สเทวมานุเสหิ เต.} \\ \csromangathabat{น สาทิยนฺติ สกฺการํ,}{พุทฺธานํ เอส ธมฺมตา”ติ.}}}

\prose{ราชา อาห “ภนฺเต นาคเสน ปุตฺโต วา ปิตุโน วณฺณํ ภาสติ, ปิตา วา ปุตฺตสฺส วณฺณํ ภาสติ, น เจตํ การณํ ปรวาทานํ นิคฺคหาย, ปสาท\-ปฺปกาสนํ นาเมตํ, อิงฺฆ เม ตฺวํ ตตฺถ การณํ สมฺมา พฺรูหิ สกวาทสฺส ปติฏฺฐาปนาย ทิฏฺฐิชาลวินิเวฐนายา”ติ.}

\prose{เถโร อาห ปรินิพฺพุโต มหาราช ภควา, น จ ภควา ปูชํ สาทิยติ, อสาทิยนฺตสฺเสว ตถาคตสฺส เทวมนุสฺสา ธาตุรตนํ วตฺถุํ กริตฺวา ตถาคตสฺส ญาณ\-รตนา\-รมฺมเณน สมฺมา\-ปฏิปตฺติํ เสวนฺตา ติสฺโส สมฺปตฺติโย ปฏิลภนฺติ.}

\prose{\csromanfolio{104}ยถา มหาราช มหติ\-มหา\-อคฺคิกฺขนฺโธ ปชฺชลิตฺวา นิพฺพาเยยฺย, อปิ นุ โข โส มหาราช มหา\-อคฺคิกฺขนฺโธ สาทิยติ ติณกฏฺฐุ\-ปาทานนฺติ. ชลมาโนปิ โส ภนฺเต มหา\-อคฺคิกฺขนฺโธ ติณกฏฺฐุ\-ปาทานํ น สาทิยติ, กิํ ปน นิพฺพุโต อุปสนฺโต อเจตโน สาทิยติ, ตสฺมิํ ปน มหาราช อคฺคิกฺขนฺเธ อุปรเต อุปสนฺเต โลเก อคฺคิ สุญฺโญ โหตีติ. น หิ ภนฺเต กฏฺฐํ อคฺคิสฺส วตฺถุ โหติ อุปาทานํ, เย เกจิ มนุสฺสา อคฺคิกามา, เต อตฺตโน ถาม\-พลวีริเยน ปจฺจตฺต\-ปุริสกาเรน กฏฺฐํ มนฺถยิตฺวา\csromansharedfootnote{5d43d9293290c088}{มทฺทิตฺวา (ก)} อคฺคิํ นิพฺพตฺเตตฺวา เตน อคฺคินา อคฺคิกรณียานิ กมฺมานิ กโรนฺตีติ. เตน หิ มหาราช ติตฺถิยานํ วจนํ มิจฺฉา ภวติ “อสาทิยนฺตสฺส กโต อธิกาโร วญฺโฌ ภวติ อผโล”ติ.}

\prose{ยถา มหาราช มหติ\-มหา\-อคฺคิกฺขนฺโธ ปชฺชลิ, เอวเมว ภควา ทสสหสฺสิยา\csromansharedfootnote{9064689a06deaa74}{ทสสหสฺสิมฺหิ (สี, อิ, ก)} โลกธาตุยา พุทฺธสิริยา ปชฺชลิ. ยถา มหาราช มหติ\-มหา\-อคฺคิกฺขนฺโธ ปชฺชลิตฺวา นิพฺพุโต, เอวเมว ภควา ทสสหสฺสิยา โลกธาตุยา พุทฺธสิริยา ปชฺชลิตฺวา อนุปาทิเสสาย นิพฺพานธาตุยา ปรินิพฺพุโต. ยถา มหาราช นิพฺพุโต อคฺคิกฺขนฺโธ ติณกฏฺฐุ\-ปาทานํ น สาทิยติ, เอวเมว โข โลกหิตสฺส สาทิยนา ปหีนา อุปสนฺตา. ยถา มหาราช มนุสฺสา นิพฺพุเต อคฺคิกฺขนฺเธ อนุปาทาเน อตฺตโน ถาม\-พลวีริเยน ปจฺจตฺต\-ปุริสกาเรน กฏฺฐํ มนฺถยิตฺวา อคฺคิํ นิพฺพตฺเตตฺวา เตน อคฺคินา อคฺคิกรณียานิ กมฺมานิ กโรนฺติ, เอวเมว โข เทวมนุสฺสา ตถาคตสฺส ปรินิพฺพุตสฺส อสาทิยนฺตสฺเสว ธาตุรตนํ วตฺถุํ กริตฺวา ตถาคตสฺส ญาณ\-รตนา\-รมฺมเณน สมฺมา\-ปฏิปตฺติํ เสวนฺตา ติสฺโส สมฺปตฺติโย ปฏิลภนฺติ, อิมินาปิ มหาราช การเณน ตถาคตสฺส ปรินิพฺพุตสฺส อสาทิยนฺตสฺเสว กโต อธิกาโร อวญฺโฌ ภวติ สผโล.}

\prose{อปรมฺปิ มหาราช อุตฺตริํ การณํ สุโณหิ เยน การเณน ตถาคตสฺส ปรินิพฺพุตสฺส อสาทิยนฺตสฺเสว กโต อธิกาโร อวญฺโฌ ภวติ สผโล, ยถา มหาราช มหติมหาวาโต วายิตฺวา อุปรเมยฺย, อปิ นุ โข โส มหาราช อุปรโต วาโต สาทิยติ ปุน นิพฺพตฺตาปนนฺติ. น หิ ภนฺเต อุปรตสฺส วาตสฺส อาโภโค \csromanfolio{105}วา มนสิกาโร วา ปุน นิพฺพตฺตาปนาย, กิํ การณํ, อเจตนา สา วาโยธาตูติ. อปิ นุ ตสฺส มหาราช อุปรตสฺส วาตสฺส วาโตติ สมญฺญา อปคจฺฉตีติ. น หิ ภนฺเต ตาลวณฺฏ\-วิธูปนานิ วาตสฺส อุปฺปตฺติยา ปจฺจยา, เย เกจิ มนุสฺสา อุณฺหาภิตตฺตา ปริฬาห\-ปริปีฬิตา, เต ตาลวณฺเฏน วา วิธูปเนน วา อตฺตโน ถาม\-พลวีริเยน ปจฺจตฺต\-ปุริสกาเรน ตํ นิพฺพตฺเตตฺวา เตน วาเตน อุณฺหํ นิพฺพาเปนฺติ ปริฬาหํ วูปสเมนฺตีติ. เตน หิ มหาราช ติตฺถิยานํ วจนํ มิจฺฉา ภวติ “อสาทิยนฺตสฺส กโต อธิกาโร วญฺโฌ ภวติ อผโล”ติ.}

\prose{ยถา มหาราช มหติมหาวาโต วายิ, เอวเมว ภควา ทสสหสฺสิยา โลกธาตุยา สีตล\-มธุร\-สนฺต\-สุขุม\-เมตฺตา\-วาเตน อุปวายิ. ยถา มหาราช มหติมหาวาโต วายิตฺวา อุปรโต, เอวเมว ภควา สีตล\-มธุร\-สนฺต\-สุขุม\-เมตฺตา\-วาเตน อุปวายิตฺวา อนุปาทิเสสาย นิพฺพานธาตุยา ปรินิพฺพุโต. ยถา มหาราช อุปรโต วาโต ปุน นิพฺพตฺตาปนํ น สาทิยติ, เอวเมว โลกหิตสฺส สาทิยนา ปหีนา อุปสนฺตา. ยถา มหาราช เต มนุสฺสา อุณฺหาภิตตฺตา ปริฬาห\-ปริปีฬิตา, เอวเมว เทวมนุสฺสา ติวิธคฺคิ\-สนฺตาป\-ปริฬาห\-ปริปีฬิตา. ยถา ตาลวณฺฏ\-วิธูปนานิ วาตสฺส นิพฺพตฺติยา ปจฺจยา โหนฺติ, เอวเมว ตถาคตสฺส ธาตุ จ ญาณรตนญฺจ ปจฺจโย โหติ ติสฺสนฺนํ สมฺปตฺตีนํ ปฏิลาภาย. ยถา มนุสฺสา อุณฺหาภิตตฺตา ปริฬาห\-ปริปีฬิตา ตาลวณฺเฏน วา วิธูปเนน วา วาตํ นิพฺพตฺเตตฺวา อุณฺหํ นิพฺพาเปนฺติ ปริฬาหํ วูปสเมนฺติ, เอวเมว เทวมนุสฺสา ตถาคตสฺส ปรินิพฺพุตสฺส อสาทิยนฺตสฺเสว ธาตุญฺจ ญาณรตนญฺจ ปูเชตฺวา กุสลํ นิพฺพตฺเตตฺวา เตน กุสเลน ติวิธคฺคิ\-สนฺตาป\-ปริฬาหํ นิพฺพาเปนฺติ วูปสเมนฺติ. อิมินาปิ มหาราช การเณน ตถาคตสฺส ปรินิพฺพุตสฺส อสาทิยนฺตสฺเสว กโต อธิกาโร อวญฺโฌ ภวติ สผโลติ.}

\prose{อปรมฺปิ มหาราช อุตฺตริํ การณํ สุโณหิ ปรวาทานํ นิคฺคหาย, ยถา มหาราช ปุริโส เภริํ อาโกเฏตฺวา สทฺทํ นิพฺพตฺเตยฺย, โย โส เภริสทฺโท ปุริเสน นิพฺพตฺติโต, โส สทฺโท อนฺตรธาเยยฺย, อปิ นุ โข โส มหาราช สทฺโท สาทิยติ ปุน นิพฺพตฺตาปนนฺติ. น หิ ภนฺเต อนฺตรหิโต โส สทฺโท, นตฺถิ ตสฺส ปุน อุปฺปาทาย \csromanfolio{106}อาโภโค วา มนสิกาโร วา, สกิํ นิพฺพตฺเต เภริสทฺเท อนฺตรหิเต โส เภริสทฺโท สมุจฺฉินฺโน โหติ. เภรี ปน ภนฺเต ปจฺจโย โหติ สทฺทสฺส นิพฺพตฺติยา, อถ ปุริโส ปจฺจเย สติ อตฺตเชน วายาเมน เภริํ อโกเฏตฺวา สทฺทํ นิพฺพตฺเตตีติ, เอวเมว โข มหาราช ภควา สีลสมาธิ\-ปญฺญาวิมุตฺติ\-วิมุตฺติญาณทสฺสน\-ปริภาวิตํ ธาตุรตนญฺจ ธมฺมญฺจ วินยญฺจ อนุสิฏฺฐญฺจ\csromansharedfootnote{07d5c9e7a9c02808}{อนุสตฺถิญฺจ (สี, อิ)} สตฺถารํ ฐปยิตฺวา สยํ อนุปาทิเสสาย นิพฺพานธาตุยา ปรินิพฺพุโต, น จ ปรินิพฺพุเต ภควติ สมฺปตฺติลาโภ อุปจฺฉินฺโน โหติ, ภวทุกฺข\-ปฏิปีฬิตา สตฺตา ธาตุรตนญฺจ ธมฺมญฺจ วินยญฺจ อนุสิฏฺฐญฺจ ปจฺจยํ กริตฺวา สมฺปตฺติกามา สมฺปตฺติโย ปฏิลภนฺติ, อิมินาปิ มหาราช การเณน ตถาคตสฺส ปรินิพฺพุตสฺส อสาทิยนฺตสฺเสว กโต อธิกาโร อวญฺโฌ ภวติ สผโลติ.}

\proseitem{4}{ทิฏฺฐญฺเจตํ มหาราช ภควตา อนาคตมทฺธานํ. กถิตญฺจ ภณิตญฺจ อาจิกฺขิตญฺจ\csromansymbolfootnote{*}{วิ 4. 480; ที 2. 134 ปิฏฺเฐสุ.}“สิยา โข ปนานนฺท ตุมฺหากํ เอวมสฺส อตีตสตฺถุกํ ปาวจนํ นตฺถิ โน สตฺถาติ, น โข ปเนตํ อานนฺท เอวํ ทฏฺฐพฺพํ, โย โว อานนฺท มยา ธมฺโม จ วินโย จ เทสิโต ปญฺญตฺโต, โส โว มมจฺจเยน สตฺถา”ติ. ปรินิพฺพุตสฺส ตถาคตสฺส อสาทิยนฺตสฺส กโต อธิกาโร วญฺโฌ ภวติ อผโลติ, ตํ เตสํ ติตฺถิยานํ วจนํ มิจฺฉา อภูตํ วิตถํ อลิกํ วิรุทฺธํ วิปรีตํ ทุกฺขทายกํ ทุกฺขวิปากํ อปาย\-คมนียนฺติ.}

\prose{อปรมฺปิ มหาราช อุตฺตริํ การณํ สุโณหิ เยน การเณน ตถาคตสฺส ปรินิพฺพุตสฺส อสาทิยนฺตสฺเสว กโต อธิกาโร อวญฺโฌ ภวติ สผโล, สาทิยติ นุ โข มหาราช อยํ มหาปถวี “สพฺพพีชานิ มยิ สํวิรุหนฺตู”ติ. น หิ ภนฺเตติ. กิสฺส ปน ตานิ มหาราช พีชานิ อสาทิยนฺติยา มหาปถวิยา สํวิรุหิตฺวา ทฬฺห มูลชฏาปติฏฺฐิตา ขนฺธสารสาขาปริวิตฺถิณฺณา ปุปฺผ\-ผล\-ธรา โหนฺตีติ. อสาทิยนฺตีปิ ภนฺเต มหาปถวี เตสํ พีชานํ วตฺถุ โหติ ปจฺจยํ เทติ วิรุหนาย, ตานิ พีชานิ ตํ วตฺถุํ นิสฺสาย เตน ปจฺจเยน สํวิรุหิตฺวา ทฬมูลชฏาปติฏฺฐิตา ขนฺธสารสาขาปริวิตฺถิณฺณา ปุปฺผ\-ผล\-ธรา โหนฺตีติ. \csromanfolio{107}เตน หิ มหาราช ติตฺถิยา สเก วาเท นฏฺฐา โหนฺติ หตา วิรุทฺธา, สเจ เต ภณนฺติ “อสาทิยนฺตสฺส กโต อธิกาโร วญฺโฌ ภวติ อผโล”ติ.}

\prose{ยถา มหาราช มหาปถวี, เอวํ ตถาคโต อรหํ สมฺมาสมฺพุทฺโธ. ยถา มหาราช มหาปถวี น กิญฺจิ สาทิยติ, เอวํ ตถาคโต น กิญฺจิ สาทิยติ. ยถา มหาราช ตานิ พีชานิ ปถวิํ นิสฺสาย สํวิรุหิตฺวา ทฬฺหมูล\-ชฏา\-ปติฏฺฐิตา ขนฺธสารสาขาปริวิตฺถิณฺณา ปุปฺผ\-ผล\-ธรา โหนฺติ, เอวํ เทวมนุสฺสา ตถาคตสฺส ปรินิพฺพุตสฺส อสาทิยนฺตสฺเสว ธาตุญฺจ ญาณรตนญฺจ นิสฺสาย ทฬฺห\-กุสลมูล\-ปติฏฺฐิตา สมาธิกฺขนฺธธมฺมสารสีลสาขาปริวิตฺถิณฺณา วิมุตฺติ\-ปุปฺผ\-สามญฺญผล\-ธรา โหนฺติ, อิมินาปิ มหาราช การเณน ตถาคตสฺส ปรินิพฺพุกสฺส อสาทิยนฺตสฺเสว กโต อธิกาโร อวญฺโฌ ภวติ สผโลติ.}

\prose{อปรมฺปิ มหาราช อุตฺตริํ การณํ สุโณหิ เยน การเณน ตถาคตสฺส ปรินิพฺพุตสฺส อสาทิยนฺตสฺเสว กโต อธิกาโร อวญฺโฌ ภวติ สผโล, สาทิยนฺติ นุ โข มหาราช อิเม โอฏฺฐา โคณา คทฺรภา อชา ปสู มนุสฺสา อนฺโตกุจฺฉิสฺมิํ กิมิกุลานํ สมฺภวนฺติ. น หิ ภนฺเตติ. กิสฺส ปน เต มหาราช กิมโย เตสํ อสาทิยนฺตานํ อนฺโตกุจฺฉิสฺมิํ สมฺภวิตฺวา พหุปุตฺตนตฺตา เวปุลฺลตํ ปาปุณนฺตีติ. ปาปสฺส ภนฺเต กมฺมสฺส พลวตาย อสาทิยนฺตานํ เยว เตสํ สตฺตานํ อนฺโตกุจฺฉิสฺมิํ กิมโย สมฺภวิตฺวา พหุปุตฺตนตฺตา เวปุลฺลตํ ปาปุณนฺตีติ, เอวเมว โข มหาราช ตถาคตสฺส ปรินิพฺพุตสฺส อสาทิยนฺตสฺเสว ธาตุสฺส จ ญาณารมฺมณสฺส จ พลวตาย ตถาคเต กโต อธิกาโร อวญฺโฌ ภวติ สผโลติ.}

\prose{อปรมฺปิ มหาราช อุตฺตริํ การณํ สุโณหิ เยน การเณน ตถาคตสฺส ปรินิพฺพุตสฺส อสาทิยนฺตสฺเสว กโต อธิกาโร อวญฺโฌ ภวติ สผโล, สาทิยนฺติ นุ โข มหาราช อิเม มนุสฺสา อิเม อฏฺฐนวุติ โรคา กาเย นิพฺพตฺตนฺตูติ. น หิ ภนฺเตติ. กิสฺส ปน เต มหาราช โรคา อสาทิยนฺตานํ กาเย นิปตนฺตีติ. ปุพฺเพ กเตน ภนฺเต ทุจฺจริเตนาติ. ยทิ มหาราช ปุพฺเพ กตํ อกุสลํ \csromanfolio{108}อิธ เวทนียํ โหติ, เตน หิ มหาราช ปุพฺเพ กตมฺปิ อิธ กตมฺปิ กุสลากุสลํ กมฺมํ อวญฺฌํ ภวติ สผลนฺติ, อิมินาปิ มหาราช การเณน ตถาคตสฺส ปรินิพฺพุตสฺส อสาทิยนฺตสฺเสว กโต อธิกาโร อวญฺโฌ ภวติ สผโลติ.}

\prose{สุตปุพฺพํ ปน ตยา มหาราช\csromansymbolfootnote{*}{ขุ 1. 123 ปิฏฺเฐ.}นนฺทโก นาม ยกฺโข เถรํ สาริปุตฺตํ อาสาทยิตฺวา ปถวิํ ปวิฏฺโฐติ. อาม ภนฺเต สุยฺยติ, โลเก ปากโฏ เอโสติ, อปิ นุ โข มหาราช เถโร สาริปุตฺโต สาทิยิ นนฺทกสฺส ยกฺขสฺส มหาปถวิคิลนนฺติ. อุพฺพตฺติยนฺเตปิ\csromansharedfootnote{6d050c4d0d6184b9}{ปวตฺตมาเนปิ (สฺยา)} ภนฺเต สเทวเก โลเก ปตมาเนปิ ฉมายํ จนฺทิมสูริเย วิกิรนฺเตปิ สิเนรุ\-ปพฺพตราเช เถโร สาริปุตฺโต น ปรสฺส ทุกฺขํ สาทิเยยฺย, ตํ กิสฺส เหตุ, เยน เหตุนา เถโร สาริปุตฺโต กุชฺเฌยฺย วา ทุสฺเสยฺย วา, โส เหตุ เถรสฺส สาริปุตฺตสฺส สมูหโต สมุจฺฉินฺโน, เหตุโน สมุคฺฆาติต\-ตฺตา ภนฺเต เถโร สาริปุตฺโต ชีวิตหารเกปิ โกปํ น กเรยฺยาติ. ยทิ มหาราช เถโร สาริปุตฺโต นนฺทกสฺส ยกฺขสฺส ปถวิคิลนํ น สาทิยิ, กิสฺส ปน นนฺทโก ยกฺโข ปถวิํ ปวิฏฺโฐติ, อกุสลสฺส ภนฺเต กมฺมสฺส พลวตายาติ. ยทิ มหาราช อกุสลสฺส กมฺมสฺส พลวตาย นนฺทโก ยกฺโข ปถวิํ ปวิฏฺโฐ, อสาทิยนฺตสฺสาปิ กโต อปราโธ อวญฺโฌ ภวติ สผโล. เตน หิ มหาราช อกุสลสฺสปิ กมฺมสฺส พลวตาย อสาทิยนฺตสฺส กโต อธิกาโร อวญฺโฌ ภวติ สผโลติ. อิมินาปิ มหาราช การเณน ตถาคตสฺส ปรินิพฺพุตสฺส อสาทิยนฺตสฺเสว กโต อธิกาโร อวญฺโฌ ภวติ สผโลติ.}

\prosewithcloserruled{กติ นุ โข เต มหาราช มนุสฺสา, เย เอตรหิ มหาปถวิํ ปวิฏฺฐา, อตฺถิ เต ตตฺถ สวณนฺติ. อาม ภนฺเต สุยฺยตีติ, อิงฺฆ ตฺวํ มหาราช สาเวหีติ. จิญฺจมาณวิกา ภนฺเต, สุปฺปพุทฺโธ จ สกฺโก, เทวทตฺโต จ เถโร, นนฺทโก จ ยกฺโข, นนฺโท จ มาณวโกติ. สุตเมตํ ภนฺเต อิเม ปญฺจ ชนา มหาปถวิํ ปวิฏฺฐาติ. กิสฺมิํ เต มหาราช อปรทฺธาติ, ภควติ จ ภนฺเต สาวเกสุ จาติ. อปิ นุ โข \csromanfolio{109}มหาราช ภควา วา สาวกา วา สาทิยิํสุ อิเมสํ มหาปถวิปวิสนนฺติ. น หิ ภนฺเตติ, เตน หิ มหาราช ตถาคตสฺส ปรินิพฺพุตสฺส อสาทิยนฺตสฺเสว กโต อธิกาโร อวญฺโฌ ภวติ สผโลติ, สุวิญฺญาปิโต ภนฺเต นาคเสน ปโญฺห คมฺภีโร อุตฺตานี กโต, คุยฺหํ วิทํสิตํ, คณฺฐิ ภินฺโน, คหนํ อคหนํ กตํ, นฏฺฐา ปรวาทา, ภคฺคา กุทิฏฺฐี, นิปฺปภา ชาตา กุติตฺถิยา, ตฺวํ คณีวรปวรมาสชฺชานฺติ.}{กตาธิการ\-สผล\-ปโญฺห ปฐโม.}
\csromanmatikaanchor{mtk.108}
\csromantocmarkref{subsubsection}{2. สพฺพญฺญุภาวปญฺห}{mtk.108}
\bookmark[dest={mtk.108},level=3]{2. สพฺพญฺญุภาวปญฺห}
\csromanheader{3}{2. \textbf{สพฺพญฺญุ}\-\textbf{ภาว}\-\textbf{ปญฺห}}
\proseitem{4}{ภนฺเต นาคเสน พุทฺโธ สพฺพญฺญูติ. อาม มหาราช ภควา สพฺพญฺญู, น จ ภควโต สตตํ สมิตํ ญาณทสฺสนํ ปจฺจุปฏฺฐิตํ, อาวชฺชนปฏิพทฺธํ ภควโต สพฺพญฺญุต\-ญาณํ, อาวชฺชิตฺวา ยทิจฺฉกํ ชานาตีติ, เตน หิ ภนฺเต นาคเสน พุทฺโธ อสพฺพญฺญูติ. ยทิ ตสฺส ปริเยสนาย สพฺพญฺญุต\-ญาณํ โหตีติ. วาหสตํ โข มหาราช วีหีนํ อฑฺฒจูฬญฺจ วาหา วีหิสตฺต\-มฺพณานิ เทฺว จ ตุมฺพา เอก\-จฺฉรากฺขเณ ปวตฺต\-จิตฺตสฺส เอตฺตกา วีหี ลกฺขํ ฐปียมานา\csromansharedfootnote{62b2b660aa6a4441}{ฐปียมาเน (สี, อิ)} ปริกฺขยํ ปริยาทานํ คจฺเฉยฺยุํ.}

\prose{ตตฺริเม สตฺตวิธา จิตฺตา ปวตฺตนฺติ, เย เต มหาราช สราคา สโทสา สโมหา สกิเลสา อภาวิตกายา อภาวิตสีลา อภาวิตจิตฺตา อภาวิตปญฺญา, เตสํ ตํ จิตฺตํ ครุกํ อุปฺปชฺชติ ทนฺธํ ปวตฺตติ, กิํ การณา, อภาวิตตฺตา จิตฺตสฺส. ยถา มหาราช วํสนาฬสฺส วิตตสฺส วิสาลสฺส วิตฺถิณฺณสฺส สํสิพฺพิตวิสิพฺพิตสฺส สาขา\-ชฏา\-ชฏิตสฺส อากฑฺฒิ\-ยนฺตสฺส ครุกํ โหติ อาคมนํ ทนฺธํ. กิํ การณา, สํสิพฺพิตวิสิพฺพิตตฺตา สาขานํ, เอวเมว โข มหาราช เย เต สราคา สโทสา สโมหา สกิเลสา อภาวิตกายา อภาวิตสีลา อภาวิตจิตฺตา อภาวิตปญฺญา, เตสํ ตํ จิตฺตํ ครุกํ อุปฺปชฺชติ ทนฺธํ ปวตฺตติ. กิํ การณา, สํสิพฺพิตวิสิพฺพิตตฺตา กิเลเสหิ, อิทํ ปฐมํ จิตฺตํ.}

\prose{ตตฺริทํ ทุติยํ จิตฺตํ วิภตฺตมา\-ปชฺชติ\csromandash{}เย เต มหาราช โสตาปนฺนา ปิหิตาปายา ทิฏฺฐิปฺปตฺตา วิญฺญาต\-สตฺถุสาสนา, เตสํ ตํ จิตฺตํ \csromanfolio{110}ตีสุ ฐาเนสุ ลหุกํ อุปฺปชฺชติ ลหุกํ ปวตฺตติ. อุปริภูมีสุ ครุกํ อุปฺปชฺชติ ทนฺธํ ปวตฺตติ. กิํ การณา, ตีสุ ฐาเนสุ จิตฺตสฺส ปริสุทฺธตฺตา อุปริ กิเลสานํ อปฺปหีนตฺตา. ยถา มหาราช วํสนาฬสฺส ติปพฺพคณฺฐิปริสุทฺธสฺส อุปริ สาขา\-ชฏา\-ชฏิตสฺส อากฑฺฒิ\-ยนฺตสฺส ยาว ติปพฺพํ ตาว ลหุกํ เอติ, ตโต อุปริ ถทฺธํ, กิํ การณา, เหฏฺฐา ปริสุทฺธตฺตา อุปริ สาขา\-ชฏา\-ชฏิ\-ตตฺตา, เอวเมว โข มหาราช เย เต โสตาปนฺนา ปิหิตาปายา ทิฏฺฐิปฺปตฺตา วิญฺญาต\-สตฺถุสาสนา, เตสํ ตํ จิตฺตํ ตีสุ ฐาเนสุ ลหุกํ อุปฺปชฺชติ ลหุกํ ปวตฺตติ, อุปริภูมีสุ ครุกํ อุปฺปชฺชติ ทนฺธํ ปวตฺตติ. กิํ การณา, ตีสุ ฐาเนสุ จิตฺตสฺส ปริสุทฺธตฺตา อุปริ กิเลสานํ อปฺปหีนตฺตา, อิทํ ทุติยํ จิตฺตํ.}

\proseitem{2}{ตตฺริทํ ตติยํ จิตฺตํ วิภตฺตมา\-ปชฺชติ\csromandash{}เย เต มหาราช สกทาคามิโน, เยสํ ราคโทสโมหา ตนุภูตา, เตสํ ตํ จิตฺตํ ปญฺจสุ ฐาเนสุ ลหุกํ อุปฺปชฺชติ ลหุกํ ปวตฺตติ, อุปริภูมีสุ ครุกํ อุปฺปชฺชติ ทนฺธํ ปวตฺตติ. กิํ การณา, ปญฺจสุ ฐาเนสุ จิตฺตสฺส ปริสุทฺธตฺตา อุปริ กิเลสานํ อปฺปหีนตฺตา. ยถา มหาราช วํสนาฬสฺส ปญฺจ\-ปพฺพ\-คณฺฐิ\-ปริสุทฺธสฺส อุปริ สาขา\-ชฏา\-ชฏิตสฺส อากฑฺฒิ\-ยนฺตสฺส ยาว ปญฺจปพฺพํ ตาว ลหุกํ เอติ, ตโต อุปริ ถทฺธํ, กิํ การณา, เหฏฺฐา ปริสุทฺธตฺตา อุปริ สาขา\-ชฏา\-ชฏิ\-ตตฺตา, เอวเมว โข มหาราช เย เต สกทาคามิโน, เยสํ ราคโทสโมหา ตนุภูตา, เตสํ ตํ จิตฺตํ ปญฺจสุ ฐาเนสุ ลหุกํ อุปฺปชฺชติ ลหุกํ ปวตฺตติ, อุปริภูมีสุ ครุกํ อุปฺปชฺชติ ทนฺธํ ปวตฺตติ, กิํ การณา, ปญฺจสุ ฐาเนสุ จิตฺตสฺส ปริสุทฺธตฺตา อุปริ กิเลสานํ อปฺปหีนตฺตา, อิทํ ตติยํ จิตฺตํ.}

\prose{ตตฺริทํ จตุตฺถํ จิตฺตํ วิภตฺตมา\-ปชฺชติ\csromandash{}เย เต มหาราช อนาคามิโน, เยสํ ปญฺโจ\-รมฺภาคิยานิ สญฺโญชนานิ ปหีนานิ, เตสํ ตํ จิตฺตํ ทสสุ ฐาเนสุ ลหุกํ อุปฺปชฺชติ ลหุกํ ปวตฺตติ, อุปริภูมีสุ ครุกํ อุปฺปชฺชติ ทนฺธํ ปวตฺตติ, กิํ การณา, ทสสุ ฐาเนสุ จิตฺตสฺส ปริสุทฺธตฺตา อุปริ กิเลสานํ อปฺปหีนตฺตา. ยถา มหาราช วํสนาฬสฺส ทสปพฺพ\-คณฺฐิ\-ปริสุทฺธสฺส อุปริ สาขา\-ชฏา\-ชฏิตสฺส อากฑฺฒิ\-ยนฺตสฺส ยาว ทสปพฺพํ ตาว ลหุกํ เอติ, ตโต อุปริ ตทฺธํ, กิํ การณา, เหฏฺฐา ปริสุทฺธตฺตา อุปริ สาขา\-ชฏา\-ชฏิ\-ตตฺตา, เอวเมว โข มหาราช เย เต อนาคมิโน, เยสํ ปญฺโจ\-รมฺภาคิยานิ สญฺโญชนานิ ปหีนานิ, เตสํ \csromanfolio{111}ตํ จิตฺตํ ทสสุ ฐาเนสุ ลหุกํ อุปฺปชฺชติ ลหุกํ ปวตฺตติ, อุปริภูมีสุ ครุกํ อุปฺปชฺชติ ทนฺธํ ปวตฺตติ, กิํ การณา, ทสสุ ฐาเนสุ จิตฺตสฺส ปริสุทฺธตฺตา อุปริ กิเลสานํ อปฺปหีนตฺตา, อิทํ จตุตฺถํ จิตฺตํ.}

\prose{ตตฺริทํ ปญฺจมํ จิตฺตํ วิภตฺตมา\-ปชฺชติ\csromandash{}เย เต มหาราช อรหนฺโต ขีณาสวา โธตมลา วนฺตกิเลสา วุสิตวนฺโต กตกรณียา โอหิตภารา อนุปฺปตฺต\-สทตฺถา ปริกฺขีณ\-ภว\-สญฺโญชนา ปตฺต\-ปฏิสมฺภิทา สาวกภูมีสุ ปริสุทฺธา, เตสํ ตํ จิตฺตํ สาวกวิสเย ลหุกํ อุปฺปชฺชติ ลหุกํ ปวตฺตติ, ปจฺเจก\-พุทฺธภูมีสุ ครุกํ อุปฺปชฺชติ ทนฺธํ ปวตฺตติ, กิํ การณา, ปริสุทฺธตฺตา สาวกวิสเย, อปริสุทฺธตฺตา ปจฺเจก\-พุทฺธวิสเย. ยถา มหาราช วํสนาฬสฺส สพฺพ\-ปพฺพ\-คณฺฐิ\-ปริสุทฺธสฺส อากฑฺฒิ\-ยนฺตสฺส ลหุกํ โหติ อาคมนํ อทนฺธํ, กิํ การณา, สพฺพ\-ปพฺพ\-คณฺฐิ\-ปริสุทฺธตฺตา อคหนตฺตา วํสสฺส, เอวเมว โข มหาราช เย เต อรหนฺโต ขีณาสวา โธตมลา วนฺตกิเลสา วุสิตวนฺโต กตกรณียา โอหิตภารา อนุปฺปตฺต\-สทตฺถา ปริกฺขีณ\-ภว\-สญฺโญชนา ปตฺต\-ปฏิสมฺภิทา สาวกภูมีสุ ปริสุทฺธา, เตสํ ตํ จิตฺตํ สาวกวิสเย ลหุกํ อุปฺปชฺชติ ลหุกํ ปวตฺตติ, ปจฺเจก\-พุทฺธภูมีสุ ครุกํ อุปฺปชฺชติ ทนฺธํ ปวตฺตติ, กิํ การณา, ปริสุทฺธตฺตา สาวกวิสเย, อปริสุทฺธตฺตา ปจฺเจก\-พุทฺธวิสเย, อิทํ ปญฺจมํ จิตฺตํ.}

\prose{ตตฺริทํ ฉฏฺฐํ จิตฺตํ วิภตฺตมา\-ปชฺชติ\csromandash{}เย เต มหาราช ปจฺเจกพุทฺธาสยมฺภุโน อนาจริยกา เอกจาริโน ขคฺควิสาณ\-กปฺปา สกวิสเย ปริสุทฺธ\-วิมล\-จิตฺตา, เตสํ ตํ จิตฺตํ สกวิสเย ลหุกํ อุปฺปชฺชติ ลหุกํ ปวตฺตติ, สพฺพญฺญุ\-พุทฺธภูมีสุ ครุกํ อุปฺปชฺชติ ทนฺธํ ปวตฺตติ, กิํ การณา, ปริสุทฺธตฺตา สกวิสเย มหนฺตตฺตา สพฺพญฺญุ\-พุทฺธวิสยสฺส. ยถา มหาราช ปุริโส สกวิสยํ ปริตฺตํ นทิํ รตฺติมฺปิ ทิวาปิ ยทิจฺฉกํ อจฺฉมฺภิโต โอตเรยฺย, อถ ปรโต มหาสมุทฺทํ คมฺภีรํ วิตฺถตํ อคาธมปารํ ทิสฺวา ภาเยยฺย, ทนฺธาเยยฺย น วิสเหยฺย โอตริตุํ. กิํ การณา, ติณฺณตฺตา\csromansharedfootnote{df4e20f8de8f4cbb}{จิณฺณตฺตา (สี, สฺยา, อิ)} สกวิสยสฺส, มหนฺตตฺต จ มหาสมุทฺทสฺส, เอวเมว โข มหาราช เย เต ปจฺเจกพุทฺธา สยมฺภุโน อนาจริยกา เอกจาริโน ขคฺควิสาณ\-กปฺปา สกวิสเย ปริสุทฺธ\-วิมล\-จิตฺตา, \csromanfolio{112}เตสํ ตํ จิตฺตํ สกวิสเย ลหุกํ อุปฺปชฺชติ ลหุกํ ปวตฺตติ, สพฺพญฺญุ\-พุทฺธภูมีสุ ครุกํ อุปฺปชฺชติ ทนฺธํ ปวตฺตติ, กิํ การณา, ปริสุทฺธตฺตา สกวิสเย มหนฺตตฺตา สพฺพญฺญุ\-พุทฺธวิสยสฺส, อิทํ ฉฏฺฐํ จิตฺตํ.}

\proseitem{4}{ตตฺริทํ สตฺตมํ จิตฺตํ วิภตฺตมา\-ปชฺชติ\_\_เย เต มหาราช สมฺมาสมฺพุทฺธา สพฺพญฺญุโน ทสพลธรา จตุเวสารชฺช\-วิสารทา อฏฺฐารสหิ พุทฺธธมฺเมหิ สมนฺนาคตา อนนฺตชินา อนาวรณญาณา, เตสํ ตํ จิตฺตํ สพฺพตฺถ ลหุกํ อุปฺปชฺชติ ลหุกํ ปวตฺตติ, กิํ การณา, สพฺพตฺถ ปริสุทฺธตฺตา, อปิ นุ โข มหาราช นาราจสฺส สุโธตสฺส วิมลสฺส นิคฺคณฺฐิสฺส สุขุม\-ธารสฺส อชิมฺหสฺส อวงฺกสฺส อกุฏิลสฺส ทฬฺห\-จาป\-สมารูฬฺหสฺส โขมสุขุเม วา กปฺปาสสุขุเม วา กมฺพลสุขุเม วา พลว\-นิปาติตสฺส ทนฺธายิตตฺตํ วา ลคฺคนํ วา โหตีติ. น หิ ภนฺเต, กิํ การณา, สุขุมตฺตา วตฺถานํ สุโธตตฺตา นาราจสฺส นิปาตสฺส จ พลวตฺตาติ, เอวเมว โข มหาราช เย เต สมฺมาสมฺพุทฺธา สพฺพญฺญุโน ทสพลธรา จตุเวสารชฺช\-วิสารทา อฏฺฐารสหิ พุทฺธธมฺเมหิ สมนฺนาคตา อนนฺตชินา อนาวรณญาณา, เตสํ ตํ จิตฺตํ สพฺพตฺถ ลหุกํ อุปฺปชฺชติ ลหุกํ ปวตฺตติ, กิํ การณา, สพฺพตฺถ ปริสุทฺธตฺตา, อิทํ สตฺตมํ จิตฺตํ.}

\prose{ตตฺร มหาราช ยทิทํ สพฺพญฺญุ\-พุทฺธานํ จิตฺตํ, ตํ ฉนฺนมฺปิ จิตฺตานํ คณนํ อติกฺกมิตฺวา อสงฺเขฺยยฺเยน คุเณน ปริสุทฺธญฺจ ลหุกญฺจ, ยสฺมา จ ภควโต จิตฺตํ ปริสุทฺธญฺจ ลหุกญฺจ, ตสฺมา มหาราช ภควา ยมกปาฏิหีรํ ทสฺเสติ. ยมกปาฏิหีเร มหาราช ญาตพฺพํ พุทฺธานํ ภควนฺตานํ จิตฺตํ เอวํ ลหุปริวตฺตนฺติ, น ตตฺถ สกฺกา อุตฺตริํ การณํ วตฺตุํ, เตปิ มหาราช ปาฏิหีรา สพฺพญฺญุ\-พุทฺธานํ จิตฺตํ อุปาทาย คณนมฺปิ สงฺขมฺปิ กลมฺปิ กลภาคมฺปิ น อุเปนฺติ, อาวชฺชนปฏิพทฺธํ มหาราช ภควโต สพฺพญฺญุต\-ญาณํ, อาวชฺเชตฺวา ยทิจฺฉกํ ชานาติ.}

\prose{ยถา มหาราช ปุริโส หตฺเถ ฐปิตํ ยํ กิญฺจิ ทุติเย หตฺเถ ฐเปยฺย วิวเฏน มุเขน วาจํ นิจฺฉาเรยฺย, มุขคตํ โภชนํ คิเลยฺย, อุมฺมีเลตฺวา วา นิมีเลยฺย, นิมีเลตฺวา วา อุมฺมีเลยฺย, สมิญฺชิตํ วา พาหํ ปสาเรยฺย, ปสาริตํ วา พาหํ สมิญฺเชยฺย, จิรตรํ เอตํ มหาราช, ลหุตรํ ภควโต สพฺพญฺญุต\-ญาณํ, ลหุตรํ อาวชฺชนํ, อาวชฺเชตฺวา ยทิจฺฉกํ ชานาติ, อาวชฺชน\-วิกล\-มตฺตเกน น ตาวตา พุทฺธา ภควนฺโต อสพฺพญฺญุโน นาม โหนฺตีติ.}

\prose{\csromanfolio{113}อาวชฺชนมฺปิ ภนฺเต นาคเสน ปริเยสนาย กาตพฺพํ, อิงฺฆ มํ ตตฺถ การเณน สญฺญาเปหีติ. ยถา มหาราช ปุริสสฺส อฑฺฒสฺส มหทฺธนสฺส มหาโภคสฺส ปหูต\-ชาตรูป\-รชตสฺส ปหูต\-วิตฺตู\-ปกรณสฺส ปหูต\-ธน\-ธญฺญสฺส สาลิวีหิยวตณฺฑุลติลมุคฺคมาสปุพฺพณฺณาปรณฺณสปฺปิเตลนวนีตขีรท ธิมธุคุฬผาณิตา จ ขโฬปิ\-กุมฺภิ\-ปีฐร\-โกฏฺฐ\-ภาชนคตา ภเวยฺยุํ, ตสฺส จ ปุริสสฺส ปาหุนโก อาคจฺเฉยฺย ภตฺตารโห ภตฺตาภิ\-กงฺขี, ตสฺส จ เคเห ยํ รนฺธํ โภชนํ, ตํ ปรินิฏฺฐิตํ ภเวยฺย, กุมฺภิโต ตณฺฑุเล นีหริตฺวา โภชนํ รนฺเธยฺย, อปิ จ โข โส มหาราช ตาวตเกน โภชน\-เวกลฺล\-มตฺตเกน อธโน นาม กปโณ นาม ภเวยฺยาติ. น หิ ภนฺเต, จกฺกวตฺติรญฺโญ ฆเรปิ ภนฺเต อกาเล โภชน\-เวกลฺลํ โหติ, กิํ ปน คหปติ\-กสฺสาติ. เอวเมว โข มหาราช ตถาคตสฺส อาวชฺชน\-วิกล\-มตฺตกํ สพฺพญฺญุต\-ญาณํ อาวชฺเชตฺวา ยทิจฺฉกํ ชานาติ.}

\prose{ยถา วา ปน มหาราช รุกฺโข อสฺส ผลิโต โอณตวินโต ปิณฺฑิ\-ภาร\-ภริโต, น กิญฺจิ ตตฺถ ปติตํ ผลํ ภเวยฺย, อปิ นุ โข โส มหาราช รุกฺโข ตาวตเกน ปติตผล\-เวกลฺล\-มตฺตเกน อผโล นาม ภเวยฺยาติ. น หิ ภนฺเต, ปตน\-ปฏิพทฺธานิ ตานิ รุกฺขผลานิ, ปติเต ยทิจฺฉกํ ลภตีติ. เอวเมว โข มหาราช ตถาคตสฺสํ อาวชฺชนปฏิพทฺธํ สพฺพญฺญุต\-ญาณํ อาวชฺเชตฺวา ยทิจฺฉกํ ชานาตีติ.}

\prose{ภนฺเต นาคเสน อาวชฺเชตฺวา อาวชฺเชตฺวา พุทฺโธ ยทิจฺฉกํ ชานาตีติ. อาม มหาราช ภควา อาวชฺเชตฺวา อาวชฺเชตฺวา ยทิจฺฉกํ ชานาตีติ.}

\prosewithcloserruled{ยถา มหาราช จกฺกวตฺตี ราชา ยทา จกฺกรตนํ สรติ “อุเปตุ เม จกฺกรตนนฺ”ติ, สริเต จกฺกรตนํ อุเปติ, เอวเมว โข มหาราช ตถาคโต อาวชฺเชตฺวา อาวชฺเชตฺวา ยทิจฺฉกํ ชานาตีติ. ทฬฺหํ ภนฺเต นาคเสน การณํ, พุทฺโธ สพฺพญฺญู, สมฺปฏิจฺฉาม พุทฺโธ สพฺพญฺญูติ.}{พุทฺธ\-สพฺพญฺญุ\-ภาว\-ปโญฺห ทุติโย.}
\csromanmatikaanchor{mtk.109}
\csromantocmarkref{subsubsection}{3. เทวทตฺตปพฺพชฺชปญฺห}{mtk.109}
\bookmark[dest={mtk.109},level=3]{3. เทวทตฺตปพฺพชฺชปญฺห}
\csromanheader{3}{3. \textbf{เทวทตฺต}\-\textbf{ปพฺพชฺช}\-\textbf{ปญฺห}}
\proseitem{3}{ภนฺเต นาคเสน เทวทตฺโต เกน ปพฺพาชิโตติ. ฉ ยิเม มหาราช ขตฺติยกุมารา ภทฺทิโย จ อนุรุทฺโธ จ อานนฺโท จ ภคุ จ \csromanfolio{114}กิมิโล\csromansharedfootnote{538d71e1f24b7f12}{กิมฺพิโล (สี, อิ) วิ 3. 338; ม 2. 125 ปิฏฺเฐ ปสฺสิตพฺพํ.} จ เทวทตฺโต จ อุปาลิกปฺปโก สตฺตโม อภิสมฺพุทฺเธ สตฺถริ สกฺยกุลา\-นนฺทชนเน ภควนฺตํ อนุปพฺพชนฺตา นิกฺขมิํสุ, เต ภควา ปพฺพาเชสีติ. นนุ ภนฺเต เทวทตฺเตน ปพฺพชิตฺวา สํโฆ ภินฺโนติ. อาม มหาราช เทวทตฺเตน ปพฺพชิตฺวา สํโฆ ภินฺโน, น คิหี สํฆํ ภินฺทติ, น ภิกฺขุนี, น สิกฺขมานา, น สามเณโร, น สามเณรี สํฆํ ภินฺทติ, ภิกฺขุ ปกตตฺโต สมาน\-สํวาสโก สมานสีมายํ ฐิโต สํฆํ ภินฺทตีติ. สํฆเภทโก ภนฺเต ปุคฺคโล กิํ กมฺมํ ผุสตีติ. กปฺปฏฺฐิติกํ มหาราช กมฺมํ ผุสตีติ.}

\prose{กิํ ปน ภนฺเต นาคเสน พุทฺโธ ชานาติ “เทวทตฺโต ปพฺพชิตฺวา สํฆํ ภินฺทิสฺสติ, สํฆํ ภินฺทิตฺวา กปฺปํ นิรเย ปจฺจิสฺสตี”ติ. อาม มหาราช ตถาคโต ชานาติ “เทวทตฺโต ปพฺพชิตฺวา สํฆํ ภินฺทิสฺสติ, สํฆํ ภินฺทิตฺวา กปฺปํ นิรเย ปจฺจิสฺสตี”ติ. ยทิ ภนฺเต นาคเสน พุทฺโธ ชานาติ “เทวทตฺโต ปพฺพชิตฺวา สํฆํ ภินฺทิสฺสติ, สํฆํ ภินฺทิตฺวา กปฺปํ นิรเย ปจฺจิสฺสตี”ติ, เตน หิ ภนฺเต นาคเสน พุทฺโธ การุณิโก อนุกมฺปโก หิเตสี สพฺพสตฺตานํ อหิตํ อปเนตฺวา หิตมุ\-ปทหตีติ ยํ วจนํ, ตํ มิจฺฉา. ยทิ ตํ อชานิตฺวา ปพฺพาเชสิ, เตน หิ พุทฺโธ อสพฺพญฺญูติ, อยมฺปิ อุภโต โกฏิโก ปโญฺห ตวานุปฺปตฺโต, วิชเฏหิ เอตํ มหาชฏํ, ภินฺท ปราปวาทํ, อนาคเต อทฺธาเน ตยา สทิสา พุทฺธิมนฺโต ภิกฺขู ทุลฺลภา ภวิสฺสนฺติ, เอตฺถ ตว พลํ ปกาเสหีติ.}

\prose{การุณิโก มหาราช ภควา สพฺพญฺญู จ, การุญฺเญน มหาราช ภควา สพฺพญฺญุต\-ญาเณน เทวทตฺตสฺส คติํ โอโลเกนฺโต อทฺทส เทวทตฺตํ อาปายิกํ กมฺมํ\csromansharedfootnote{edcda40c0a1814a0}{อปราปริยกมฺมํ (สี, สฺยา, อิ)} อายูหิตฺวา อเนกานิ กปฺป\-โกฏิ\-สต\-สหสฺสานิ นิรเยน นิรยํ วินิปาเตน วินิปาตํ คจฺฉนฺตํ, ตํ ภควา สพฺพญฺญูตญาเณน ชานิตฺวา อิมสฺส อปริยนฺตกตํ กมฺมํ มม สาสเน ปพฺพชิตสฺส ปริยนฺตกตํ ภวิสฺสติ, ปุริมํ อุปาทาย ปริยนฺตกตํ ทุกฺขํ ภวิสฺสติ, อปพฺพชิโตปิ อยํ โมฆปุริโส กปฺปฏฺฐิยเมว กมฺมํ อายูหิสฺสตีติ การุญฺเญน เทวทตฺตํ ปพฺพาเชสีติ.}

\prose{เตน หิ ภนฺเต นาคเสน พุทฺโธ วธิตฺวา เตเลน มกฺเขติ, ปปาเต ปาเตตฺวา หตฺถํ เทติ, มาเรตฺวา ชีวิตํ ปริเยสติ, ยํ โส ปฐมํ \csromanfolio{115}ทุกฺขํ ทตฺวา ปจฺฉา สุขํ อุปทหตีติ. วเธติปิ มหาราช ตถาคโต สตฺตานํ หิตวเสน, ปาเตติปิ สตฺตานํ หิตวเสน, มาเรติปิ สตฺตานํ หิตวเสน, วธิตฺวาปิ มหาราช ตถาคโต สตฺตานํ หิตเมว อุปทหติ, ปาเตตฺวาปิ สตฺตานํ หิตเมว อุปทหติ, มาเรตฺวาปิ สตฺตานํ หิตเมว อุปทหติ. ยถา มหาราช มาตาปิตโร นาม วธิตฺวาปิ ปาตยิตฺวาปิ ปุตฺตานํ หิตเมว อุปทหนฺติ, เอวเมว โข มหาราช ตถาคโต วเธติปิ สตฺตานํ หิตวเสน, ปาเตติปิ สตฺตานํ หิตวเสน, มาเรติปิ สตฺตานํ หิตวเสน, วธิตฺวาปิ มหาราช ตถาคโต สตฺตานํ หิตเมว อุปทหติ, ปาเตตฺวาปิ สตฺตานํ หิตเมว อุปทหติ, มาเรตฺวาปิ สตฺตานํ หิตเมว อุปทหติ, เยน เยน โยเคน สตฺตานํ คุณวุฑฺฒิ โหติ, เตน เตน โยเคน สพฺพสตฺตานํ หิตเมว อุปทหติ. สเจ มหาราช เทวทตฺโต น ปพฺพาเชยฺย, คิหิภูโต สมาโน นิรยสํวตฺตนิกํ พหุํ ปาปกมฺมํ กตฺวา อเนกานิ กปฺป\-โกฏิ\-สต\-สหสฺสานิ นิรเยน นิรยํ วินิปาเตน วินิปาตํ คจฺฉนฺโต พหุํ ทุกฺขํ เวทยิสฺสติ, ตํ ภควา ชานมาโน การุญฺเญน เทวทตฺตํ ปพฺพาเชสิ, “มม สาสเน ปพฺพชิตสฺส ทุกฺขํ ปริยนฺตกตํ ภวิสฺสตี”ติ การุญฺเญน ครุกํ ทุกฺขํ ลหุกํ อกาสิ.}

\prose{ยถา วา มหาราช ธนยส\-สิริ\-ญาติพเลน พลวา ปุริโส อตฺตโน ญาติํ วา มิตฺตํ วา รญฺญา ครุกํ ทณฺฑํ ธาเรนฺตํ อตฺตโน พหุวิสฺสตฺถ\-ภาเวน สมตฺถตาย ครุกํ ทณฺฑํ ลหุกํ อกาสิ, เอวเมว โข มหาราช ภควา พหูนิ กปฺป\-โกฏิ\-สต\-สหสฺสานิ ทุกฺขํ เวทยมานํ เทวทตฺตํ ปพฺพาเชตฺวา สีลสมาธิปญฺญาวิมุตฺติ\-พล\-สมตฺถภาเวน ครุกํ ทุกฺขํ ลหุกํ อกาสิ.}

\prose{ยถา วา ปน มหาราช กุสโล ภิสกฺโก สลฺลกตฺโต ครุกํ โรคํ พลโว\-สธ\-พเลน ลหุกํ กโรติ, เอวเมว โข มหาราช พหูนิ กปฺป\-โกฏิ\-สต\-สหสฺสานิ ทุกฺขํ เวทยมานํ เทวทตฺตํ ภควา โรคญฺญุตาย ปพฺพเชตฺวา การุญฺญพโล ปตฺถทฺธ\-ธมฺโม\-สธ\-พเลน ครุกํ ทุกฺขํ ลหุกํ อกาสิ. อปิ นุ โข โส มหาราช ภควา พหุเวทนียํ เทวทตฺตํ อปฺปเวทนียํ กโรนฺโต กิญฺจิ อปุญฺญํ อาปชฺเชยฺยาติ. น กิญฺจิ ภนฺเต อปุญฺญํ อาปชฺเชยฺย อนฺตมโส คทฺทูหนมตฺตมฺปีติ. อิมมฺปิ โข มหาราช การณํ อตฺถโต สมฺปฏิจฺฉ, เยน การเณน ภควา เทวทตฺตํ ปพฺพาเชสิ.}

\proseitem{4}{\csromanfolio{116}อปรมฺปิ มหาราช อุตฺตริํ การณํ สุโณหิ, เยน การเณน ภควา เทวทตฺตํ ปพฺพาเชสิ, ยถา มหาราช โจรํ อาคุจาริํ คเหตฺวา รญฺโญ ทสฺเสยฺยุํ, อยํ โข เทว โจโร อาคุจารี, อิมสฺส ยํ อิจฺฉสิ, ตํ ทณฺฑํ ปเณหีติ. ตเมนํ ราชา เอวํ วเทยฺย “เตน หิ ภเณ อิมํ โจรํ พหินครํ นีหริตฺวา อาฆาตเน สีสํ ฉินฺทถา”ติ, “เอวํ เทวา”ติ โข เต รญฺโญ ปฏิสฺสุตฺวา ตํ พหินครํ นีหริตฺวา อาฆาตนํ นเยยฺยุํ. ตเมนํ ปสฺเสยฺย โกจิเทว ปุริโส รญฺโญ สนฺติกา ลทฺธวโร ลทฺธ\-ยส\-ธน\-โภโค อาเทยฺยวจโน พลวิ\-จฺฉิต\-การี, โส ตสฺส การุญฺญํ กตฺวา เต ปุริเส เอวํ วเทยฺย “อลํ โภ กิํ ตุมฺหากํ อิมสฺส สีสจฺเฉทเนน, เตน หิ โภ อิมสฺส หตฺถํ วา ปาทํ วา ฉินฺทิตฺวา ชีวิตํ รกฺขถ, อหเมตสฺส การณา รญฺโญ สนฺติเก ปฏิวจนํ กริสฺสามี”ติ. เต ตสฺส พลวโต วจเนน ตสฺส โจรสฺส หตฺถํ วา ปาทํ วา ฉินฺทิตฺวา ชีวิตํ รกฺเขยฺยุํ. อปิ นุ โข โส มหาราช ปุริโส เอวํ การี ตสฺส โจรสฺส กิจฺจการี อสฺสาติ. ชีวิตทายโก โส ภนฺเต ปุริโส ตสฺส โจรสฺส, ชีวิเต ทินฺเน กิํ ตสฺส อกตํ นาม อตฺถีติ. ยา ปน หตฺถปาท\-จฺเฉทเน เวทนา, โส ตาย เวทนาย กิญฺจิ อปุญฺญํ อาปชฺเชยฺยาติ. อตฺตโน กเตน โส ภนฺเต โจโร ทุกฺขเวทนํ เวทยติ, ชีวิตทายโก ปน ปุริโส น กิญฺจิ อปุญฺญํ อาปชฺเชยฺยาติ, เอวเมว โข มหาราช ภควา การุญฺเญน เทวทตฺตํ ปพฺพาเชสิ “มม สาสเน ปพฺพชิตสฺส ทุกฺขํ ปริยนฺตกตํ ภวิสฺสตี”ติ. ปริยนฺตกตญฺจ มหาราช เทวทตฺตสฺส ทุกฺขํ, เทวทตฺโต มหาราช มรณกาเล\csromandash{}}

\setlength{\csromangathapagewidth}{0pt}
\setlength{\csromangathawakwidth}{0pt}
\csromangathameasuregroup{}{“อิเมหิ อฏฺฐีหิ ต’มคฺคปุคฺคลํ, \\ เทวาติเทวํ นรทมฺมสารถิํ. \\ สมนฺตจกฺขุํ สตปุญฺญลกฺขณํ, \\ ปาเณหิ พุทฺธํ สรณํ อุเปมี”ติ.}
\csromangathameasurewak{“อิเมหิ อฏฺฐีหิ ต’มคฺคปุคฺคลํ, \\ เทวาติเทวํ นรทมฺมสารถิํ. \\ สมนฺตจกฺขุํ สตปุญฺญลกฺขณํ, \\ ปาเณหิ พุทฺธํ สรณํ อุเปมี”ติ.}
\setlength{\csromangathaleftcol}{0pt}
\csromangathagroup{\csromangathastanza{\csromangathawak{“อิเมหิ อฏฺฐีหิ ต’มคฺคปุคฺคลํ, \\ เทวาติเทวํ นรทมฺม\-สารถิํ. \\ สมนฺตจกฺขุํ สตปุญฺญ\-ลกฺขณํ, \\ ปาเณหิ พุทฺธํ สรณํ อุเปมี”ติ.}}}

\prose{ปาณุเปตํ สรณมคมาสิ. เทวทตฺโต มหาราช ฉ โกฏฺฐาเส กเต กปฺเป อติกฺกนฺเต ปฐม\-โกฏฺฐาเส สํฆํ ภินฺทิ, ปญฺจ โกฏฺฐาเส นิรเย ปจฺจิตฺวา ตโต มุจฺจิตฺวา อฏฺฐิสฺสโร นาม ปจฺเจกพุทฺโธ ภวิสฺสติ. อปิ นุ โข โส มหาราช ภควา เอวํ การี เทวทตฺตสฺส กิจฺจการี อสฺสาติ. สพฺพทโท ภนฺเต นาคเสน ตถาคโต เทวทตฺตสฺส, ยํ \csromanfolio{117}ตถาคโต เทวทตฺตํ ปจฺเจกโพธิํ ปาเปสฺสติ, กิํ ตถาคเตน เทวทตฺตสฺส อกตํ นาม อตฺถีติ. ยํ ปน มหาราช เทวทตฺโต สํฆํ ภินฺทิตฺวา นิรเย ทุกฺขเวทนํ เวทยติ, อปิ นุ โข ภควา ตโตนิทานํ กิญฺจิ อปุญฺญํ อาปชฺเชยฺยาติ. น หิ ภนฺเต, อตฺตนา กเตน ภนฺเต เทวทตฺโต กปฺปํ นิรเย ปจฺจติ, ทุกฺข\-ปริยนฺต\-การโก สตฺถา น กิญฺจิ อปุญฺญํ อาปชฺชตีติ. อิมมฺปิ โข ตฺวํ มหาราช การณํ อตฺถโต สมฺปฏิจฺฉ, เยน การเณน ภควา เทวทตฺตํ ปพฺพาเชสิ.}

\prose{อปรมฺปิ มหาราช อุตฺตริํ การณํ สุโณหิ, เยน การเณน ภควา เทวทตฺตํ ปพฺพาเชสิ, ยถา มหาราช กุสโล ภิสกฺโก สลฺลกตฺโต วาตปิตฺตเสมฺหสนฺนิปาตอุตุปริณามวิสมปริหารโอปกฺกมิโกปกฺกนฺตํ ปูติกุณป\-ทุคฺคนฺธาภิ\-สญฺฉนฺนํ อนฺโตสลฺลํ สุสิรคตํ ปุพฺพ\-รุหิร\-สมฺปุณฺณํ วณํ วูปสเมนฺโต วณมุขํ กกฺขฬ\-ติขิณ\-ขาร\-กฏุเกน เภสชฺเชน อนุลิมฺปติ ปริปจฺจนาย, ปริปจฺจิตฺวา มุทุภาวมุ\-ปคตํ สตฺเถน วิกนฺตยิตฺวา ฑหติ สลากาย, ทฑฺเฒ ขารลวณํ เทติ, เภสชฺเชน อนุลิมฺปติ วณรุหนาย พฺยาธิตสฺส โสตฺถิภาวม\-นุปฺปตฺติยา, อปิ นุ โข โส มหาราช ภิสกฺโก สลฺลกตฺโต อหิตจิตฺโต เภสชฺเชน อนุลิมฺปติ, สตฺเถน วิกนฺเตติ, ฑหติ สลากาย, ขารลวณํ เทตีติ. น หิ ภนฺเต, หิตจิตฺโต โสตฺถิกาโม ตานิ กิริยานิ กโรตีติ. ยา ปนสฺส เภสชฺช\-กิริยา\-กรเณน อุปฺปนฺนา ทุกฺขเวทนา, ตโตนิทานํ โส ภิสกฺโก สลฺลกตฺโต กิญฺจิ อปุญฺญํ อาปชฺเชยฺยาติ. หิตจิตฺโต ภนฺเต โสตฺถิกาโม ภิสกฺโก สลฺลกตฺโต ตานิ กิริยานิ กโรติ, กิํ โส ตโตนิทานํ อปุญฺญํ อาปชฺเชยฺย, สคฺคคามี โส ภนฺเต ภิสกฺโก สลฺลกตฺโตติ, เอวเมว โข มหาราช การุญฺเญน ภควา เทวทตฺตํ ปพฺพาเชสิ ทุกฺข\-ปริมุตฺติยา.}

\prose{อปรมฺปิ มหาราช อุตฺตริํ การณํ สุโณหิ, เยน การเณน ภควา เทวทตฺตํ ปพฺพาเชสิ, ยถา มหาราช ปุริโส กณฺฏเกน วิทฺโธ อสฺส, อถญฺญตโร ปุริโส ตสฺส หิตกาโม โสตฺถิกาโม ติเณฺหน กณฺฏเกน วา สตฺถมุเขน วา สมนฺตโต ฉินฺทิตฺวา ปคฺฆรนฺเตน โลหิเตน ตํ กณฺฏกํ นีหเรยฺย, อปิ นุ โข โส มหาราช ปุริโส อหิตกาโม ตํ กณฺฏกํ นีหรตีติ. น หิ ภนฺเต, หิตกาโม โส ภนฺเต ปุริโส โสตฺถิกาโม ตํ กณฺฏกํ นีหรติ. \csromanfolio{118}สเจ โส ภนฺเต ตํ กณฺฏกํ น นีหเรยฺย, มรณํ วา โส เตน ปาปุเณยฺย มรณมตฺตํ วา ทุกฺขนฺติ, เอวเมว โข มหาราช ตถาคโต การุญฺเญน เทวทตฺตํ ปพฺพาเชสิ ทุกฺข\-ปริมุตฺติยา. สเจ มหาราช ภควา เทวทตฺตํ น ปพฺพาเชยฺย, กปฺปโกฏิสตสหสฺสมฺปิ เทวทตฺโต ภวปรมฺปราย นิรเย ปจฺเจยฺยาติ.}

\prosewithcloserruled{อนุโสตคามิํ ภนฺเต นาคเสน เทวทตฺต ตถาคโต ปฏิโสตํ ปาเปสิ, วิปนฺถปฏิปนฺนํ เทวทตฺตํ ปนฺเถ ปฏิปาเทสิ, ปปาเต ปติตสฺส เทวทตฺตสฺส ปติฏฺฐํ อทาสิ, วิสมคตํ เทวทตฺตํ ตถาคโต สมํ อาโรเปสิ, อิเม จ ภนฺเต นาคเสน เหตู อิมานิ จ การณานิ น สกฺกา อญฺเญน สนฺทสฺเสตุํ อญฺญตฺร ตวาทิเสน พุทฺธิมตาติ.}{เทวทตฺต\-ปพฺพชฺช\-ปโญฺห ตติโย.}
\csromanmatikaanchor{mtk.110}
\csromantocmarkref{subsubsection}{4. ปถวิจลนปญฺห}{mtk.110}
\bookmark[dest={mtk.110},level=3]{4. ปถวิจลนปญฺห}
\csromanheader{3}{4. \textbf{ปถวิ}\-\textbf{จลน}\-\textbf{ปญฺห}}
\proseitem{4}{ภนฺเต นาคเสน ภาสิตมฺเปตํ ภควตา\csromandash{}“อฏฺฐิเม ภิกฺขเว\csromansharedfootnote{15065a1736c4f010}{อฏฺฐิเม อานนฺท (ที 2. 90; อํ 3. 132 ปิฏฺเฐสุ.)} เหตู อฏฺฐ ปจฺจยา มหโต ภูมิจาลสฺส ปาตุภาวายา”ติ. อเสสวจนํ อิทํ, นิสฺเสสวจนํ อิทํ, นิปฺปริยาย\-วจนํ อิทํ, นตฺถญฺโญ นวโม เหตุ มหโต ภูมิจาลสฺส ปาตุภาวาย. ยทิ ภนฺเต นาคเสน อญฺโญ นวโม เหตุ ภเวยฺย มหโต ภูมิจาลสฺส ปาตุภาวาย, ตมฺปิ เหตุํ ภควา กเถยฺย. ยสฺมา จ โข ภนฺเต นาคเสน นตฺถญฺโญ นวโม เหตุ มหโต ภูมิจาลสฺส ปาตุภาวาย, ตสฺมา อนาจิกฺขิโต ภควตา, อยญฺจ นวโม เหตุ ทิสฺสติ มหโต ภูมิจาลสฺส ปาตุภาวาย, ยํ เวสฺสนฺตเรน รญฺญา มหาทาเน ทียมาเน สตฺตกฺขตฺตุํ มหาปถวี กมฺปิตาติ. ยทิ ภนฺเต นาคเสน อฏฺเฐว เหตู อฏฺฐ ปจฺจยา มหโต ภูมิจาลสฺส ปาตุภาวาย, เตน หิ เวสฺสนฺตเรน รญฺญา มหาทาเน ทียมาเน สตฺตกฺขตฺตุํ มหาปถวี กมฺปิตาติ ยํ วจนํ, ตํ มิจฺฉา. ยทิ เวสฺสนฺตเรน รญฺญา มหาทาเน ทียมาเน สตฺตกฺขตฺตุํ มหาปถวี กมฺปิตา, เตน หิ อฏฺเฐว เหตู อฏฺฐ ปจฺจยา มหโต ภูมิจาลสฺส ปาตุภาวายาติ ตมฺปิ วจนํ มิจฺฉา. อยมฺปิ อุภโต โกฏิโก ปโญฺห สุขุโม \csromanfolio{119}ทุนฺนิเวฐิโย อนฺธกรโณ เจว คมฺภีโร จ, โส ตวานุปฺปตฺโต, เนโส อญฺเญน อิตฺตรปญฺเญน สกฺกา วิสชฺเชตุํ อญฺญตฺร ตวาทิเสน พุทฺธิมตาติ.}

\prose{ภาสิตมฺเปตํ มหาราช ภควตา\csromandash{}“อฏฺฐิเม ภิกฺขเว เหตู อฏฺฐ ปจฺจยา มหโต ภูมิจาลสฺส ปาตุภาวายา”ติ. ยํ เวสฺสนฺตเรน รญฺญา มหาทาเน ทียมาเน สตฺตกฺขตฺตุํ มหาปถวี กมฺปิตา, ตญฺจ ปน อกาลิกํ กทา\-จุ\-ปฺปตฺติกํ อฏฺฐหิ เหตูหิ วิปฺปมุตฺตํ, ตสฺมา อคณิตํ อฏฺฐหิ เหตูหิ.}

\prose{ยถา มหาราช โลเก ตโย เยว เมฆา คณียนฺติ วสฺสิโก เหมนฺติโก ปาวุสโกติ. ยทิ เต มุญฺจิตฺวา อญฺโญ เมโฆ ปวสฺสติ, น โส เมโฆ คณียติ สมฺมเตหิ เมเฆหิ, อกาลเมโฆเตฺวว สงฺขํ คจฺฉติ. เอวเมว โข มหาราช เวสฺสนฺตเรน รญฺญา มหาทาเน ทียมาเน ยํ สตฺตกฺขตฺตุํ มหาปถวี กมฺปิตา, อกาลิกํ เอตํ กทา\-จุ\-ปฺปตฺติกํ อฏฺฐหิ เหตูหิ วิปฺปมุตฺตํ, น ตํ คณียติ อฏฺฐหิ เหตูหิ.}

\prose{ยถา วา ปน มหาราช หิมวนฺตา ปพฺพตา ปญฺจ นทิสตานิ สนฺทนฺติ, เตสํ มหาราช ปญฺจนฺนํ นทิสตานํ ทเสว นทิโย นทิคณนาย คณียนฺติ. เสยฺยถีทํ, คงฺคา ยมุนา อจิรวตี สรภู มหี สินฺธุ สรสฺสตี เวตฺรวตี วีตํสา จนฺทภาคาติ, อวเสสา นทิโย นทิคณนาย อคณิตา, กิํ การณา, น ตา นทิโย ธุวสลิลา. เอวเมว โข มหาราช เวสฺสนฺตเรน รญฺญา มหาทาเน ทียมาเน ยํ สตฺตกฺขตฺตุํ มหาปถวี กมฺปิตา, อกาลิกํ เอตํ กทา\-จุ\-ปฺปตฺติกํ อฏฺฐหิ เหตูหิ วิปฺปมุตฺตํ, น ตํ คณียติ อฏฺฐหิ เหตูหิ.}

\prose{ยถา วา ปน มหาราช รญฺโญ สตมฺปิ ทฺวิสตมฺปิ ติสตมฺปิ อมจฺจา โหนฺติ, เตสํ ฉ เยว ชนา อมจฺจคณนาย คณียนฺติ. เสยฺยถีทํ, เสนาปติ ปุโรหิโต อกฺขทสฺโส ภณฺฑาคาริโก ฉตฺตคฺคาหโก ขคฺคคฺคาหโก. เอเต เยว อมจฺจคณนาย คณียนฺติ. กิํ การณา, ยุตฺตตฺตา ราชคุเณหิ, อวเสสา อคณิตา, สพฺเพ อมจฺจาเตฺวว สงฺขํ คจฺฉนฺตี. เอวเมวโข มหาราช เวสฺสนฺตเรน รญฺญา มหาทาเน ทียมาเน ยํ สตฺตกฺขตฺตุํ มหาปถวี กมฺปิตา, อกาลิกํ เอตํ กทา\-จุ\-ปฺปตฺติกํ อฏฺฐหิ เหตูหิ วิปฺปมุตฺตํ, น ตํ คณียติ อฏฺฐหิ เหตูหิ.}

\prose{\csromanfolio{120}สุยฺยติ นุ โข มหาราช เอตรหิ ชินสาสเน กตาธิการานํ ทิฏฺฐธมฺม\-สุขเวทนีย\-กมฺมํ, กิตฺติ จ เยสํ อพฺภุคฺคตา เทวมนุสฺเสสูติ. อาม ภนฺเต สุยฺยติ เอตรหิ ชินสาสเน กตาธิการานํ ทิฏฺฐธมฺม\-สุขเวทนีย\-กมฺมํ, กิตฺติ จ เยสํ อพฺภุคฺคตา เทวมนุสฺเสสุ สตฺต ชนาติ. เก จ เต มหาราชาติ. สุมโน จ ภนฺเต มาลากาโร, เอกสาฏโก จ พฺราหฺมโณ, ปุณฺโณ จ ภตโก, มลฺลิกา จ เทวี, โคปาลมาตา จ เทวี, สุปฺปิยา จ อุปาสิกา, ปุณฺณา จ ทาสีติ อิเม สตฺต ทิฏฺฐธมฺม\-สุขเวทนียา สตฺตา, กิตฺติ จ อิเมสํ อพฺภุคฺคตา เทวมนุสฺเสสูติ.}

\prose{อปเรปิ สุยฺยนฺติ นุโข\csromansharedfootnote{ca12a5f72c7f51e2}{นุ โข เทฺว นิปาตา \csromandash{} ม.พ.ป.} อตีเต มานุสเกเนว สรีรเทเหน ติทสภวนํ คตาติ. อาม ภนฺเต สุยฺยนฺตีติ, เก จ เต มหาราชาติ. คุตฺติโล จ คนฺธพฺโพ, สาธีโน จ ราชา, นิมิ จ ราชา, มนฺธาตา จ ราชาติ อิเม จตุโร ชนา สุยฺยนฺติ, เตเนว มานุสเกน สรีรเทเหน ติทสภวนํ คตาติ. สุจิรมฺปิ กตํ สุยฺยติ สุกต\-ทุกฺกฏนฺติ. สุตปุพฺพํ ปน ตยา มหาราช อตีเต วา อทฺธาเน วตฺตมาเน วา อทฺธาเน อิตฺถนฺนามสฺส ทาเน ทียมาเน สกิํ วา ทฺวิกฺขตฺตุํ วา ติกฺขตฺตุํ วา มหาปถวี กมฺปิตาติ. น หิ ภนฺเตติ. อตฺถิ เม มหาราช อาคโม อธิคโม ปริยตฺติ สวนํ สิกฺขาพลํ สุสฺสูสา ปริปุจฺฉา อาจริยุปาสนํ, มยาปิ น สุตปุพฺพํ “อิตฺถนฺนามสฺส ทาเน ทียมาเน สกิํ วา ทฺวิกฺขตฺตุํ วา ติกฺขตฺตุํ วา มหาปถวี กมฺปิตี”ติ ฐเปตฺวา เวสฺสนฺตรสฺส ราชวสภสฺส ทานวรํ. ภควโต จ มหาราช กสฺสปสฺส, ภควโต จ สกฺยมุนิโนติ ทฺวินฺนํ พุทฺธานํ อนฺตเร คณนปถํ วีติวตฺตา วสฺสโกฏิโย อติกฺกนฺตา, ตตฺถปิ เม สวนํ นตฺถิ “อิตฺถนฺนามสฺส ทาเน ทียมาเน สกิํ วา ทฺวิกฺขตฺตุํ วา ติกฺขตฺตุํ วา มหาปถวี กมฺปิตา”ติ. น มหาราช ตาวตเกน วีริเยน ตาวตเกน ปรกฺกเมน มหาปถวี กมฺปติ, คุณภาร\-ภริตา มหาราช สพฺพ\-โสเจยฺย\-กิริย\-คุณ\-ภาร\-ภริตา ธาเรตุํ น วิสหนฺตี มหาปถวี จลติ กมฺปติ ปเวธติ.}

\prose{ยถา มหาราช สกฏสฺส อติภาร\-ภริตสฺส นาภิโย จ เนมิโย จ ผลนฺติ อกฺโข ภิชฺชติ, เอวเมว โข มหาราช สพฺพ\-โสเจยฺย\-กิริย\-คุณ\-ภาร\-ภริตา มหาปถวี ธาเรตุํ น วิสหนฺตี จลติ กมฺปติ ปเวธติ.}

\prose{\csromanfolio{121}ยถา วา ปน มหาราช คคนํ อนิล\-ชล\-เวค\-สญฺฉาทิตํ อุสฺสนฺน\-ชล\-ภาร\-ภริตํ อติวาเตน ผุฏิตตฺตา นทติ รวติ คฬคฬายติ, เอวเมว โข มหาราช มหาปถวี รญฺโญ เวสฺสนฺตรสฺส ทานพลวิปุล- อุสฺสนฺนภารภริตา ธาเรตุํ น วิสหนฺตี จลติ กมฺปติ ปเวธติ, น หิ มหาราช รญฺโญ เวสฺสนฺตรสฺส จิตฺตํ ราควเสน ปวตฺตติ, น โทสวเสน ปวตฺตติ, น โมหวเสน ปวตฺตติ, น มานวเสน ปวตฺตติ, น ทิฏฺฐิวเสน ปวตฺตติ, น กิเลสวเสน ปวตฺตติ, น วิตกฺกวเสน ปวตฺตติ, น อรติวเสน ปวตฺตติ, อถ โข ทานวเสน พหุลํ ปวตฺตติ “กินฺติ อนาคตา ยาจกา มม สนฺติเก อาคจฺเฉยฺยุํ, อาคตา จ ยาจกา ยถากามํ ลภิตฺวา อตฺตมนา ภเวยฺยุนฺ”ติ สตตํ สมิตํ ทานํ ปติ มานสํ ฐปิตํ โหติ. รญฺโญ มหาราช เวสฺสนฺตรสฺส สตตํ สมิตํ ทสสุ ฐาเนสุ มานสํ ฐปิตํ โหติ ทเม สเม ขนฺติยํ สํวเร ยเม นิยเม อกฺโกเธ อวิหิํสายํ สจฺเจ โสเจยฺเย, รญฺโญ มหาราช เวสฺสนฺตรสฺส กาเมสนา ปหีนา, ภเวสนา ปฏิปฺปสฺสทฺธา, พฺรหฺมจริเย\-สนาย เยว อุสฺสุกฺกํ อาปนฺโน, รญฺโญ มหาราช เวสฺสนฺตรสฺส อตฺตรกฺขา\csromansharedfootnote{d8c2760546eacef2}{ปรรกฺขาย (สี, อิ)} ปหีนา, สพฺพสตฺต\-รกฺขาย อุสฺสุกฺกํ อาปนฺโน “กินฺติ อิเม สตฺตา สมคฺคา อสฺสุ อโรคา สธนา ทีฆายุกา”ติ พหุลํ เยว มานสํ ปวตฺตติ. ททมาโน จ มหาราช เวสฺสนฺตโร ราชา ตํ ทานํ น ภวสมฺปตฺติ\-เหตุ เทติ, น ธนเหตุ เทติ, น ปฏิทานเหตุ เทติ, น อุปลาปนเหตุ เทติ, น อายุเหตุ เทติ, น วณฺณเหตุ เทติ, น สุขเหตุ เทติ, น พลเหตุ เทติ, น ยสเหตุ เทติ,น ปุตฺตเหตุ เทติ, น ธีตุเหตุ เทติ, อถ โข สพฺพญฺญุตญาณ\-เหตุ สพฺพญฺญุตญาณ\-รตนสฺส การณา เอวรูเป อตุล\-วิปุลา\-นุตฺตเร ทานวเร อทาสิ, สพฺพญฺญุตํ ปตฺโต จ อิมํ คาถํ อภาสิ\csromandash{}}

\setlength{\csromangathapagewidth}{0pt}
\setlength{\csromangathawakwidth}{0pt}
\csromangathameasuregroup{\textsuperscript{*}“ชาลิํ กณฺหาชินํ ธีตํ, \\ จชมาโน น จินฺเตสิํ,}{\csromangathabat{\textsuperscript{*}“ชาลิํ กณฺหาชินํ ธีตํ,}{มทฺทิเทวิํ ปติพฺพตํ.} \\ \csromangathabat{จชมาโน น จินฺเตสิํ,}{โพธิยา เยว การณา”ติ.}}
\csromangathasetleft{\textsuperscript{*}“ชาลิํ กณฺหาชินํ ธีตํ, \\ จชมาโน น จินฺเตสิํ,}
\csromangathagroup{\csromangathastanza{\csromangathabat{\csromansymbolfootnote{*}{ขุ 4. 395 ปิฏฺเฐ จริยาปิฏเก.}“ชาลิํ กณฺหาชินํ ธีตํ,}{มทฺทิเทวิํ ปติพฺพตํ.} \\ \csromangathabat{จชมาโน น จินฺเตสิํ,}{โพธิยา เยว การณา”ติ.}}}

\prose{เวสฺสนฺตโร มหาราช ราชา อกฺโกเธน โกธํ ชินาติ, อสาธุํ สาธุนา ชินาติ, กทริยํ ทาเนน ชินาติ, อลิกวาทินํ สจฺเจน ชินาติ, สพฺพํ อกุสลํ กุสเลน ชินาติ. ตสฺส เอวํ ททมานสฺส ธมฺมา\-นุคตสฺส ธมฺม\-สีสกสฺส\csromansharedfootnote{3caf2f737737bf11}{ธมฺมาสีสกสฺส (ก)} ทาน\-นิสฺสนฺท\-พลว\csromansharedfootnote{bc5bb8d937cc7d75}{ทานนิสฺสนฺทพล (สี, อิ)} วีริย\-วิปุล\-วิปฺผาเรน เหฏฺฐา มหาวาตา}

\prose{\csromanfolio{122}สญฺจลนฺติ สณิกํ สณิกํ สกิํ สกิํ อากุลากุลา วายนฺติ โอนมนฺติ อุนฺนมนฺติ วินมนฺติ, ฉินฺน\-ปตฺต\-ปาทปา\csromansharedfootnote{08a63d1d761249af}{สีนปฺปตฺตปาทปา (สี)} ปปตนฺติ, คุมฺพํ คุมฺพํ วลาหกา คคเน สนฺธาวนฺติ, รโชสญฺจิตา วาตา ทารุณา โหนฺติ, คคนํ อุปฺปีฬิตา วาตา วายนฺติ, สหสา ธมธมา\-ยนฺติ, มหาภีโม สทฺโท นิจฺฉรติ, เตสุ วาเตสุ กุปิเตสุ อุทกํ สณิกํ สณิกํ จลติ, อุทเก จลิเต ขุพฺภนฺติ มจฺฉกจฺฉปา, ยมกยมกา อูมิโย ชายนฺติ, ชลจรา สตฺตา ตสนฺติ, ชลวีจิ ยุคนทฺโธ วตฺตติ, วีจินาโท ปวตฺตติ, โฆรา พุพฺพุฬา\csromansharedfootnote{2528dfafc69a0201}{ปุพฺพุฬา (ก)} อุฏฺฐหนฺติ, เผณมาลา ภวนฺติ, อุตฺตรติ มหาสมุทฺโท, ทิสาวิทิสํ ธาวติ อุทกํ, อุทฺธํโสต\-ปฏิโสต\-มุขา สนฺทนฺติ สลิลธารา, ตสนฺติ อสุรา ครุฬา นาคา ยกฺขา, อุพฺพิชฺชนฺติ “กิํ นุ โข, กถํ นุ โข, สาคโร วิปริวตฺตตี”ติ, คมนปถเมสนฺติ ภีตจิตฺตา, ขุภิเต ลุฬิเต ชลธาเร ปกมฺปติ มหาปถวี สนคา สสาครา, ปริวตฺตติ สิเนรุคิริ กูฏเสล\-สิขโร วินมมาโน โหติ, วิมนา โหนฺติ อหินกุล\-พิฬาร\-โกฏฺฐุก\-สูกร\-มิคปกฺขิโน, รุทนฺติ ยกฺขา อปฺเปสกฺขา, หสนฺติ ยกฺขา มเหสกฺขา กมฺปมานาย มหาปถวิยา.}

\prose{ยถา มหาราช มหติ มหาปริโยเค อุทฺธนคเต อุทกสมฺปุณฺเณ อากิณฺณตณฺฑุเล เหฏฺฐโต อคฺคิ ชลมาโน ปฐมํ ตาว ปริโยคํ สนฺตาเปติ, ปริโยโค สนฺตตฺโต อุทกํ สนฺตาเปติ, อุทกํ สนฺตตฺตํ ตณฺฑุลํ สนฺตาเปติ, ตณฺฑุลํ สนฺตตฺตํ อุมฺมุชฺชติ นิมุชฺชติ, พุพฺพุฬกชาตํ โหติ, เผณมาลา อุตฺตรติ. เอวเมว โข มหาราช เวสฺสนฺตโร ราชา ยํ โลเก ทุจฺจชํ, ตํ จชิ, ตสฺส ตํ ทุจฺจชํ จชนฺตสฺส ทานสฺส สภาว\-นิสฺสนฺเทน เหฏฺฐา มหาวาตา ธาเรตุํ น วิสหนฺตา ปริกุปฺปิํสุ\csromansharedfootnote{31f92c5f3548b905}{ปริกมฺปิํสุ (ก)}, มหาวาเตสุ ปริกุปิเตสุ\csromansharedfootnote{ec27f900d9b1fe08}{ปริขุพฺภิเตสุ (สฺยา)} อุทกํ กมฺปิ, อุทเก กมฺปิเต มหาปถวี กมฺปิ, อิติ ตทา มหาวาตา จ อุทกญฺจ มหาปถวี จาติ อิเม ตโย เอกมนา วิย อเหสุํ มหาทาน\-นิสฺสนฺเทน วิปุล\-พลวีริเยน, นตฺเถทิโส มหาราช อญฺญสฺส ทานานุภาโว, ยถา เวสฺสนฺตรสฺส รญฺโญ มหาทานา\-นุภาโว. ยถา มหาราช มหิยา พหุวิธา มณโย วิชฺชนฺติ. เสยฺยถีทํ, อินฺทนีโล มหานีโล โชติรโส เวฬุริโย อุมฺมาปุปฺโผ สิรีสปุปฺโผ มโนหโร สูริยกนฺโต จนฺทกนฺโต วชิโร ขชฺโชปนโก \csromanfolio{123}ผุสฺสราโค โลหิตงฺโค มสารคลฺโลติ, เอเต สพฺเพ อติกฺกมฺม จกฺกวตฺติมณิ อคฺคมกฺขายติ, จกฺกวตฺติมณิ มหาราช สมนฺตา โยชนํ โอภาเสติ. เอวเมว โข มหาราช ยํ กิญฺจิ มหิยา ทานํ วิชฺชติ อปิ อสทิสทานํ ปรมํ, ตํ สพฺพํ อติกฺกมฺม เวสฺสนฺตรสฺส รญฺโญ มหาทานํ อคฺคมกฺขายติ, เวสฺสนฺตรสฺส มหาราช รญฺโญ มหาทาเน ทียมาเน สตฺตกฺขตฺตุํ มหาปถวี กมฺปิตาติ.}

\proseitemwithcloserruled{4}{อจฺฉริยํ ภนฺเต นาคเสน พุทฺธานํ, อพฺภุตํ ภนฺเต นาคเสน พุทฺธานํ, ยํ ตถาคโต โพธิสตฺโต สมาโน อสโม โลเกน เอวํขนฺติ เอวํจิตฺโต เอวํอธิมุตฺติ เอวํอธิปฺปาโย, โพธิสตฺตานํ ภนฺเต นาคเสน ปรกฺกโม ทกฺขาปิโต, ปารมี จ ชินานํ ภิยฺโย โอภาสิตา, จริยํ จรโตปิ ตาว ตถาคตสฺส สเทวเก โลเก เสฏฺฐภาโว อนุทสฺสิโต. สาธุ ภนฺเต นาคเสน โถมิตํ ชินสาสนํ, โชติตา ชินปารมี, ฉินฺโน ติตฺถิยานํ วาทคณฺฐิ, ภินฺโน ปราปวาท\-กุมฺโภ\csromansharedfootnote{63d117b1e418cd93}{คุมฺโพ ตยา วิทฺทํสิโต (สฺยา)}, ปโญฺห คมฺภีโร อุตฺตานีกโต, คหนํ อคหนํ กตํ, สมฺมา ลทฺธํ ชินปุตฺตานํ นิพฺพาหนํ\csromansharedfootnote{eb97da1cbd50d03d}{นิพฺพายนํ (ก)}, เอวเมตํ คณิวรปวร ตถา สมฺปฏิจฺฉามาติ.}{ปถวิ\-จลน\-ปโญฺห จตุตฺโถ.}
\csromanmatikaanchor{mtk.111}
\csromantocmarkref{subsubsection}{5. สิวิราชจกฺขุทานปญฺห}{mtk.111}
\bookmark[dest={mtk.111},level=3]{5. สิวิราชจกฺขุทานปญฺห}
\csromanheader{3}{5. \textbf{สิวิราช}\-\textbf{จกฺขุทาน}\-\textbf{ปญฺห}}
\proseitem{4}{ภนฺเต นาคเสน ตุเมฺห เอวํ ภณถ “สิวิราเชน ยาจกสฺส จกฺขูนิ ทินฺนานิ, อนฺธสฺส สโต ปุน ทิพฺพจกฺขูนิ อุปฺปนฺนานี”ติ, เอตมฺปิ วจนํ สกสฏํ สนิคฺคหํ สโทสํ “เหตุสมุคฺฆาเต อเหตุสฺมิํ อวตฺถุสฺมิํ นตฺถิ ทิพฺพจกฺขุสฺส อุปฺปาโท”ติ\csromansymbolfootnote{*}{ขุ 1. 231; อภิ 4. 191 ปิฏฺเฐสุ อตฺถโต สมานํ.}สุตฺเต วุตฺตํ, ยทิ ภนฺเต นาคเสน สิวิราเชน ยาจกสฺส จกฺขูนิ ทินฺนานิ, เตน หิ “ปุน ทิพฺพจกฺขูนิ อุปฺปนฺนานี”ติ ยํ วจนํ, ตํ มิจฺฉา, ยทิ ทิพฺพจกฺขูนิ อุปฺปนฺนานิ, เตน หิ “สิวิราเชน ยาจกสฺส จกฺขูนิ ทินฺนานี”ติ ยํ วจนํ, ตมฺปิ มิจฺฉา, อยมฺปิ อุภโต โกฏิโก ปโญฺห คณฺฐิโตปิ คณฺฐิตโร เวฐโตปิ เวฐตโร คหนโตปิ คหนตโร, โส ตวานุปฺปตฺโต, ตตฺถ ฉนฺทมภิชเนหิ นิพฺพาหนาย ปรวาทานํ นิคฺคตายาติ.}

\proseitem{5}{\csromanfolio{124}ทินฺนานิ มหาราช สิวิราเชน ยาจกสฺส จกฺขูนิ, ตตฺถ มา วิมติํ อุปฺปาเทหิ, ปุน ทิพฺพานิ จ จกฺขูนิ อุปฺปนฺนานิ, ตตฺถาปิ มา วิมติํ ชเนหีติ. อปิ นุ โข ภนฺเต นาคเสน เหตุสมุคฺฆาเต อเหตุสฺมิํ อวตฺถุสฺมิํ ทิพฺพจกฺขุ อุปฺปชฺชตีติ. น หิ มหาราชาติ. กิํ ปน ภนฺเต เอตฺถ การณํ, เยน การเณน เหตุสมุคฺฆาเต อเหตุสฺมิํ อวตฺถุสฺมิํ ทิพฺพจกฺขุ อุปฺปชฺชติ, อิงฺฆ ตาว การเณน มํ สญฺญาเปหีติ.}

\prose{กิํ ปน มหาราช อตฺถิ โลเก สจฺจํ นาม, เยน สจฺจวาทิโน สจฺจกิริยํ กโรนฺตีติ. อาม ภนฺเต, อตฺถิ โลเก สจฺจํ นาม, สจฺเจน ภนฺเต นาคเสน สจฺจวาทิโน สจฺจกิริยํ กตฺวา เทวํ วสฺสาเปนฺติ, อคฺคิํ นิพฺพาเปนฺติ, วิสํ ปฏิหนนฺติ, อญฺญมฺปิ วิวิธํ กตฺตพฺพํ กโรนฺตีติ. เตน หิ มหาราช ยุชฺชติ สเมติ สิวิราชสฺส สจฺจพเลน ทิพฺพจกฺขูนิ อุปฺปนฺนานีติ, สจฺจพเลน มหาราช อวตฺถุสฺมิํ ทิพฺพจกฺขุ อุปฺปชฺชติ, สจฺจํ เยว ตตฺถ วตฺถุ ภวติ ทิพฺพจกฺขุสฺส อุปฺปาทาย.}

\prose{ยถา มหาราช เย เกจิ สตฺตา สจฺจม\-นุคายนฺติ “มหาเมโฆ ปวสฺสตู”ติ, เตสํ สห สจฺจม\-นุ\-คีเตน มหาเมโฆ ปวสฺสติ, อปิ นุ โข มหาราช อตฺถิ อากาเส วสฺสเหตุ สนฺนิจิโต “เยน เหตุนา มหาเมโฆ ปวสฺสตี”ติ. น หิ ภนฺเต, สจฺจํ เยว ตตฺถ เหตุ ภวติ มหโต เมฆสฺส ปวสฺสนายาติ. เอวเมว โข มหาราช นตฺถิ ตสฺส ปกติเหตุ, สจฺจํ เยเวตฺถ วตฺถุ ภวติ ทิพฺพจกฺขุสฺส อุปฺปาทายาติ.}

\prose{ยถา วา ปน มหาราช เย เกจิ สตฺตา สจฺจม\-นุคายนฺติ “ชลิต\-ปชฺชลิต\-มหาอคฺคิกฺขนฺโธ ปฏินิวตฺตตู”ติ, เตสํ สห สจฺจม\-นุ\-คีเตน ชลิต\-ปชฺชลิต\-มหาอคฺคิกฺขนฺโธ ขเณน ปฏินิวตฺตติ. อปิ นุ โข มหาราช อตฺถิ ตสฺมิํ ชลิต\-ปชฺชลิเต มหา\-อคฺคิกฺขนฺเธ เหตุ สนฺนิจิโต “เยน เหตุนา ชลิต\-ปชฺชลิต\-มหาอคฺคิกฺขนฺโธ ขเณน ปฏินิวตฺตตี”ติ. น หิ ภนฺเต, สจฺจํ เยว ตตฺถ วตฺถุ โหติ ตสฺส ชลิต\-ปชฺชลิตสฺส มหา\-อคฺคิ\-กฺขนฺธสฺส ขเณน ปฏินิวตฺต\-นายาติ. เอวเมว โข มหาราช นตฺถิ ตสฺส ปกติเหตุ, สจฺจํ เยเวตฺถ วตฺถุ ภวติ ทิพฺพจกฺขุสฺส อุปฺปาทายาติ.}

\prose{ยถา วา ปน มหาราช เย เกจิ สตฺตา สจฺจม\-นุคายนฺติ “วิสํ หลาหลํ อคทํ ภวตู”ติ. เตสํ สห สจฺจม\-นุ\-คีเตน วิสํ หลาหลํ \csromanfolio{125}ขเณน อคทํ ภวติ, อปิ นุ โข มหาราช อตฺถิ ตสฺมิํ หลาหลวิเส เหตุ สนฺนิจิโต “เยน เหตุนา วิสํ หลาหลํ ขเณน อคทํ ภวตี”ติ. น หิ ภนฺเต, สจฺจํ เยว ตตฺถ เหตุ ภวติ วิสสฺส หลาหลสฺส ขเณน ปฏิฆาตายาติ. เอวเมว โข มหาราช วินา ปกติเหตุํ สจฺจํ เยเวตฺถ วตฺถุ ภวติ ทิพฺพจกฺขุสฺส อุปฺปาทายาติ.}

\prose{จตุนฺนมฺปิ มหาราช อริยสจฺจานํ ปฏิเวธาย นตฺถญฺญํ วตฺถุ, สจฺจํ วตฺถุํ กตฺวา จตฺตาริ อริยสจฺจานิ ปฏิวิชฺฌนฺตีติ. อตฺถิ มหาราช จีนวิสเย จีนราชา, โส มหาสมุทฺเท กีฬิตุกาโม\csromansharedfootnote{042e4c84ea9b2e35}{พลิํ กาตุกาโม (สี, อิ)} จตุมาเส จตุมาเส สจฺจกิริยํ กตฺวา สห รเถน อนฺโต\-มหาสมุทฺเท โยชนํ ปวิสติ, ตสฺส ราถสีสสฺส ปุรโต ปุรโต มหา\-วาริ\-กฺขนฺโธ ปฏิกฺกมติ, นิกฺขนฺตสฺส ปุน โอตฺถรติ, อปิ นุ โข มหาราช โส มหาสมุทฺโท สเทว\-มนุสฺเสนปิ โลเกน ปกติ\-กายพเลน สกฺกา ปฏิกฺกมาเปตุนฺติ. อติปริตฺตเกปิ ภนฺเต ตฬาเก อุทกํ น สกฺกา สเทว\-มนุสฺเสนปิ โลเกน ปกติ\-กายพเลน ปฏิกฺกมาเปตุํ, กิํ ปน มหาสมุทฺเท อุทกนฺติ. อิมินาปิ มหาราช การเณน สจฺจพลํ ญาตพฺพํ “นตฺถิ ตํ ฐานํ, ยํ สจฺเจน น ปตฺตพฺพนฺ”ติ.}

\prose{นคเร มหาราช ปาฏลิปุตฺเต อโสโก ธมฺมราชา สเนคมชานปท- อมจฺจภฏพลมหามตฺเตหิ ปริวุโต คงฺคํ นทิํ\csromansharedfootnote{9ea9f810a97efa10}{คงฺคานทิํ (สี)} นวสลิล\-สมฺปุณฺณํ สมติตฺถิกํ สมฺภริตํ ปญฺจ\-โยชนสตา\-ยามํ โยชนปุถุลํ สนฺทมานํ ทิสฺวา อมจฺเจ เอวมาห “อตฺถิ โกจิ ภเณ สมตฺโถ, โย อิมํ มหาคงฺคํ ปฏิโสตํ สนฺทาเปตุนฺ”ติ. อมจฺจา อาหํสุ “ทุกฺกรํ เทวา”ติ.}

\prose{ตสฺมิํ เยว คงฺคากูเล ฐิตา พนฺธุมตี นาม คณิกา อสฺโสสิ รญฺญา กิร เอวํ วุตฺตํ “สกฺกา นุ โข อิมํ มหาคงฺคํ ปฏิโสตํ สนฺทาเปตุนฺ”ติ, สา เอวมาห “อหญฺหิ นคเร ปาฏลิปุตฺเต คณิกา รูปูปชีวินี อนฺติมชีวิกา, มม ตาว ราชา สจฺจกิริยํ ปสฺสตู”ติ, อถ สา สจฺจกิริยํ อกาสิ, สห ตสฺสา สจฺจกิริยาย ขเณน สา มหาคงฺคา คฬคฬายนฺตี ปฏิโสตํ สนฺทิตฺถ มหโต ชนกายสฺส ปสฺสโต.}

\prose{อถ ราชา คงฺคาย อาวฏฺฏ\-อูมิ\-เวค\-ชนิตํ หลาหลสทฺทํ สุตฺวา วิมฺหิโต อจฺฉริย\-พฺภุต\-ชาโต อมจฺเจ เอวมาห “กิสฺสายํ ภเณ มหาคงฺคา ปฏิโสตํ สนฺทตี”ติ. พนฺธุมตี มหาราช คณิกา ตว วจนํ \csromanfolio{126}สุตฺวา สจฺจกิริยํ อกาสิ, ตสฺสา สจฺจกิริยาย มหาคงฺคา อุทฺธํมุขา สนฺทตีติ.}

\proseitem{4}{อถ สํวิคฺคหทโย ราชา ตุริตตุริโต สยํ คนฺตฺวา ตํ คณิกํ ปุจฺฉิ “สจฺจํ กิร เช ตยา สจฺจกิริยาย อยํ คงฺคา ปฏิโสตํ สนฺทาปิตา”ติ. อาม เทวาติ. ราชา อาห “กิํ เต ตตฺถ พลํ อตฺถิ, โก วา เต วจนํ อาทิยติ อนุมฺมตฺโต, เกน ตฺวํ พเลน อิมํ มหาคงฺคํ ปฏิโสตํ สนฺทาเปสีติ. สา อาห “สจฺจพเลนาหํ มหาราช อิมํ มหาคงฺคํ ปฏิโสตํ สนฺทาเปสินฺ”ติ. ราชา อาห “กิํ เต สจฺจพลํ อตฺถิ โจริยา ธุตฺติยา อสติยา ฉินฺนิกาย ปาปิยา ภินฺนสีลาย\csromansharedfootnote{0b152e3ecf435500}{ปาปิกาย ภินฺนสีมาย (สี)} หิริอติกฺกนฺติกาย อนฺธชนปโลภิกายา”ติ. สจฺจํ มหาราช ตาทิสิกา อหํ, ตาทิสิกายปิ เม มหาราช สจฺจกิริยา อตฺถิ, ยายาหํ อิจฺฉมานา สเทวกมฺปิ โลกํ ปริวตฺเตยฺยนฺติ. ราชา อาห “กตมา ปน สา โหติ สจฺจกิริยา, อิงฺฆ มํ สาเวหี”ติ. โย เม มหาราช ธนํ เทติ ขตฺติโย วา พฺราหฺมโณ วา เวสฺโส วา สุทฺโท วา อญฺโญ วา โกจิ, เตสํ สมกํ เยว อุปฏฺฐหามิ, “ขตฺติโย”ติ วิเสโส นตฺถิ, “สุทฺโท”ติ อติมญฺญนา\csromansharedfootnote{b93aba2a130ebcc3}{อติมญฺญมาโน (ก)} นตฺถิ, อนุนย\-ปฺปฏิฆ\-วิปฺปมุตฺตา ธนสฺสามิกํ ปริจรามิ, เอสา เม เทว สจฺจกิริยา, ยายาหํ อิมํ มหาคงฺคํ ปฏิโสตํ สนฺทาเปสินฺติ.}

\prosewithcloserruled{อิติปิ มหาราช สจฺเจ ฐิตา น กิญฺจิ อตฺถํ น วินฺทนฺติ. ทินฺนานิ จ มหาราช สิวิราเชน ยาจกสฺส จกฺขูนิ, ทิพฺพจกฺขูนิ จ อุปฺปนฺนานิ, ตญฺจ สจฺจกิริยาย. ยํ ปน สุตฺเต วุตฺตํ “มํสจกฺขุสฺมิํ นฏฺเฐ อเหตุสฺมิํ อวตฺถุสฺมิํ นตฺถิ ทิพฺพจกฺขุสฺส อุปฺปาโท”ติ. ตํ ภาวนามยํ จกฺขุํ สนฺธาย วุตฺตํ, เอวเมตํ มหาราช ธาเรหีติ. สาธุ ภนฺเต นาคเสน สุนิพฺเพฐิโต ปโญฺห, สุนิทฺทิฏฺโฐ นิคฺคโห, สุมทฺทิตา ปรวาทา, เอวเมตํ ตถา สมฺปฏิจฺฉามีติ.}{สิวิราช\-จกฺขุทาน\-ปโญฺห ปญฺจโม.}
\csromanmatikaanchor{mtk.112}
\csromantocmarkref{subsubsection}{6. คพฺภาวกฺกนฺติปญฺห}{mtk.112}
\bookmark[dest={mtk.112},level=3]{6. คพฺภาวกฺกนฺติปญฺห}
\csromanheader{3}{6. \textbf{คพฺภาวกฺกนฺติ}\-\textbf{ปญฺห}}
\proseitem{4}{ภนฺเต นาคเสน ภาสิตมฺเปตํ ภควตา “ติณฺณํ โข ปน ภิกฺขเว สนฺนิปาตา คพฺภสฺส อวกฺกนฺติ\csromansharedfootnote{60b4a7758a3c1036}{คพฺภสฺสาวกฺกนฺติ (ม 1. 332 ปิฏฺเฐ.)} โหติ, อิธ มาตาปิตโร จ \csromanfolio{127}สนฺนิปติตา โหนฺติ, มาตา จ อุตุนี โหติ, คนฺธพฺโพ จ ปจฺจุปฏฺฐิโต โหติ, อิเมสํ โข ภิกฺขเว ติณฺณํ สนฺนิปาตา คพฺภสฺส อวกฺกนฺติ โหตี”ติ, อเสส\-วจนเมตํ, นิสฺเสส\-วจนเมตํ, นิปฺปริยาย\-วจนเมตํ, อรหสฺส\-วจนเมตํ, สเทว\-มนุสฺสานํ มชฺเฌ นิสีทิตฺวา ภณิตํ, อยญฺจ ทฺวินฺนํ สนฺนิปาตา คพฺภสฺส อวกฺกนฺติ ทิสฺสติ, ทุกูเลน ตาปเสน ปาริกาย ตาปสิยา อุตุนิกาเล ทกฺขิเณน หตฺถ\-งฺคุฏฺเฐน นาภิ ปรามฏฺฐา, ตสฺส เตน นาภิ\-ปรามสเนน สามกุมาโร นิพฺพตฺโต. มาตงฺเคนาปิ อิสินา พฺราหฺมณ\-กญฺญาย อุตุนิกาเล ทกฺขิเณน หตฺถ\-งฺคุฏฺเฐน นาภิ ปรามฏฺฐา, ตสฺส เตน นาภิ\-ปรามสเนน มณฺฑโพฺย นาม มาณวโก นิพฺพตฺโตติ. ยทิ ภนฺเต นาคเสน ภควตา ภณิตํ “ติณฺณํ โข ปน ภิกฺขเว สนฺนิปาตา คพฺภสฺส อวกฺกนฺติ โหตี”ติ. เตน หิ สาโม จ กุมาโร มณฺฑโพฺย จ มาณวโก อุโภปิ เต นาภิ\-ปรามสเนน นิพฺพตฺตาติ ยํ วจนํ, ตํ มิจฺฉา, ยทิ ภนฺเต ตถาคเตน ภณิตํ “สาโม จ กุมาโร มณฺฑโพฺย จ มาณวโก นาภิ\-ปรามสเนน นิพฺพตฺตา”ติ, เตน หิ “ติณฺณํ โข ปน ภิกฺขเว สนฺนิปาตา คพฺภสฺส อวกฺกนฺติ โหตี”ติ ยํ วจนํ, ตมฺปิ มิจฺฉา, อยมฺปิ อุภโต โกฏิโก ปโญฺห สุคมฺภีโร สุนิปุโณ วิสโย พุทฺธิมนฺตานํ, โส ตวานุปฺปตฺโต, ฉินฺท วิมติปถํ, ธาเรหิ ญาณวรปฺปชฺโชตนฺติ.}

\proseitem{6}{ภาสิตมฺเปตํ มหาราช ภควตา “ติณฺณํ โข ปน ภิกฺขเว สนฺนิปาตา คพฺภสฺส อวกฺกนฺติ โหติ, อิธ มาตาปิตโร จ สนฺนิปติตา โหนฺติ, มาตา จ อุตุนี โหติ, คนฺธพฺโพ จ ปจฺจุปฏฺฐิโต โหติ, เอวํ ติณฺณํ สนฺนิปาตา คพฺภสฺส อวกฺกนฺติ โหตี”ติ. ภณิตญฺจ “สาโม จ กุมาโร มณฺฑโพฺย จ มาณวโก นาภิ\-ปรามสเนน นิพฺพตฺตา”ติ. เตน หิ ภนฺเต นาคเสน เยน การเณน ปโญฺห สุวินิจฺฉิโต โหติ, เตน การเณน มํ สญฺญาเปหีติ.}

\prose{สุตปุพฺพํ ปน ตยา มหาราช สํกิจฺโจ จ กุมาโร อิสิสิงฺโค จ ตาปโส เถโร จ กุมารกสฺสโป “อิมินา นาม เต นิพฺพตฺตา”ติ. อาม ภนฺเต สุยฺยติ, อพฺภุคฺคตา เตสํ ชาติ, เทฺว มิคเธนุโย ตาว อุตุนิกาเล ทฺวินฺนํ ตาปสานํ ปสฺสาวฏฺฐานํ อาคนฺตฺวา สสมฺภวํ ปสฺสาวํ ปิวิํสุ, เตน ปสฺสาว\-สมฺภเวน สํกิจฺโจ จ กุมาโร อิสิสิงฺโค จ ตาปโส นิพฺพตฺตา. เถรสฺส อุทายิสฺส ภิกฺขุนุ\-ปสฺสยํ \csromanfolio{128}อุปคตสฺส รตฺตจิตฺเตน ภิกฺขุนิยา องฺคชาตํ อุปนิ\-ชฺฌาย\-นฺตสฺส สมฺภวํ กาสาเว มุจฺจิ, อถ โข อายสฺมา อุทายิ ตํ ภิกฺขุนิํ เอตทโวจ “คจฺฉ ภคินิ, อุทกํ อาหร อนฺตรวาสกํ โธวิสฺสามี”ติ. อาหร’ยฺย อหเมว โธวิสฺสามีติ. ตโต สา ภิกฺขุนี อุตุนิสมเย ตํ สมฺภวํ เอกเทสํ มุเขน อคฺคเหสิ, เอกเทสํ องฺคชาเต ปกฺขิปิ, เตน เถโร กุมารกสฺสโป นิพฺพตฺโตติ เอตํ ชโน อาหาติ.}

\proseitem{4}{อปิ นุโข\csromansharedfootnote{ca12a5f72c7f51e2}{นุ โข เทฺว นิปาตา \csromandash{} ม.พ.ป.} ตฺวํ มหาราช สทฺทหสิ ตํ วจนนฺติ. อาม ภนฺเต, พลวํ ตตฺถ มยํ การณํ อุปลภาม, เยน มยํ การเณน สทฺทหาม “อิมินา การเณน นิพฺพตฺตา”ติ. กิํ ปเนตฺถ มหาราช การณนฺติ. สุปริกมฺม\-กเต ภนฺเต กลเล พีชํ นิปติตฺวา ขิปฺปํ สํวิรุหตีติ. อาม มหาราชาติ. เอวเมว โข ภนฺเต สา ภิกฺขุนี อุตุนี สมานา สณฺฐิเต กลเล รุหิเร ปจฺฉินฺนเวเค ฐิตาย ธาตุยา ตํ สมฺภวํ คเหตฺวา ตสฺมิํ กลเล ปกฺขิปิ, เตน ตสฺสา คพฺโภ สณฺฐาสิ, เอวํ ตตฺถ การณํ ปจฺเจม เตสํ นิพฺพตฺติยาติ. เอวเมตํ มหาราช ตถา สมฺปฏิจฺฉามิ, โยนิปฺปเวเสน คพฺโภ สมฺภวตีติ. สมฺปฏิจฺฉสิ ปน ตฺวํ มหาราช เถรสฺส กุมาร\-กสฺสปสฺส คพฺภาวกฺกมนนฺติ. อาม ภนฺเตติ. สาธุ มหาราช ปจฺจาคโตสิ มม วิสยํ, เอกวิเธนปิ คพฺภา\-วกฺกนฺติํ กถยนฺโต มมานุพลํ ภวิสฺสสิ, อถ ยา ปน ตา เทฺว มิคเธนุโย ปสฺสาวํ ปิวิตฺวา คพฺภํ ปฏิลภิํสุ, ตาสํ ตฺวํ สทฺทหสิ คพฺภสฺสาวกฺกมนนฺติ. อาม ภนฺเต, ยํ กิญฺจิ ภุตฺตํ ปีตํ ขายิตํ เลหิตํ, สพฺพํ ตํ กลลํ โอสรติ, ฐานคตํ วุฑฺฒิมา\-ปชฺชติ. ยถา ภนฺเต นาคเสน ยา กาจิ สริตา นาม, สพฺพา ตา มหาสมุทฺทํ โอสรนฺติ, ฐานคตา วุฑฺฒิมา\-ปชฺชนฺติ. เอวเมว โข ภนฺเต นาคเสน ยํ กิญฺจิ ภุตฺตํ ปีตํ ขายิตํ เลหิตํ, สพฺพํ ตํ กลลํ โอสรติ, ฐานคตํ วุฑฺฒิมา\-ปชฺชติ, เตนาหํ การเณน สทฺทหามิ มุขคเตนปิ คพฺภสฺส อวกฺกนฺติ โหตีติ. สาธุ มหาราช คาฬฺหตรํ อุปคโตสิ มม วิสยํ, มุขปาเนนปิ ทฺวยสนฺนิปาโต ภวติ. สํกิจฺจสฺส จ มหาราช กุมารสฺส อิสิสิงฺคสฺส จ ตาปสสฺส เถรสฺส จ กุมาร\-กสฺสปสฺส คพฺภา\-วกฺกมนํ สมฺปฏิจฺฉสีติ. อาม ภนฺเต สนฺนิปาโต โอสรตีติ.}

\prose{สาโมปิ มหาราช กุมาโร มณฺฑโพฺยปิ มาณวโก ตีสุ สนฺนิปาเตสุ อนฺโตคธา, เอกรสา เยว ปุริเมน, ตตฺถ การณํ วกฺขามิ. \csromanfolio{129}ทุกูโล จ มหาราช ตาปโส ปาริกา จ ตาปสี อุโภปิ เต อรญฺญวาสา อเหสุํ ปวิเวกา\-ธิมุตฺตา อุตฺตมตฺถ\-คเวสกา, ตปเตเชน ยาว พฺรหฺมโลกํ สนฺตาเปสุํ. เตสํ ตทา สกฺโก เทวานมินฺโท สายํ ปาตํ อุปฏฺฐานํ อาคจฺฉติ, โส เตสํ ครุกต\-เมตฺต\-ตาย อุปธาเรนฺโต อทฺทส อนาคตมทฺธาเน ทฺวินฺนมฺปิ เตสํ จกฺขูนํ อนฺตรธานํ, ทิสฺวา เต เอวมาห “เอกํ เม โภนฺโต วจนํ กโรถ, สาธุ เอกํ ปุตฺตํ ชเนยฺยาถ, โส ตุมฺหากํ อุปฏฺฐาโก ภวิสฺสติ อาลมฺพโน จา”ติ. อลํ โกสิย, มา เอวํ ภณีติ. เต ตสฺส ตํ วจนํ น สมฺปฏิจฺฉิํสุ. อนุกมฺปโก อตฺถกาโม สกฺโก เทวานมินฺโท ทุติยมฺปิ. ตติยมฺปิ เต เอวมาห “เอกํ เม โภนฺโต วจนํ กโรถ, สาธุ เอกํ ปุตฺตํ ชเนยฺยาถ, โส ตุมฺหากํ อุปฏฺฐาโก ภวิสฺสติ อาลมฺพโน จา”ติ. ตติยมฺปิ เต อาหํสุ “อลํ โกสิย, มา ตฺวํ โข อเมฺห อนตฺเถ นิโยเชหิ, กทายํ กาโย น ภิชฺชิสฺสติ, ภิชฺชตุ อยํ กาโย เภทนธมฺโม, ภิชฺชนฺติยาปิ มรณิยา ปตนฺเตปิ เสลสิขเร ผลนฺเตปิ อากาเส ปตนฺเตปิ จนฺทิมสูริเย เนว มยํ โลกธมฺเมหิ มิสฺสยิสฺสาม, มา ตฺวํ อมฺหากํ สมฺมุขภาวํ อุปคจฺฉ, อุปคตสฺส เต เอโส วิสฺสาโส, อนตฺถจโร ตฺวํ มญฺเญ”ติ.}

\prose{ตโต สกฺโก เทวานมินฺโท เตสํ มนํ อลภมาโน ครุกโต ปญฺชลิโก ปุน ยาจิ “ยทิ เม วจนํ น อุสฺสหถ กาตุํ, ยทา ตาปสี อุตุนี โหติ ปุปฺผวตี, ตทา ตฺวํ ภนฺเต ทกฺขิเณน หตฺถ\-งฺคุฏฺเฐน นาภิํ ปรามเสยฺยาสิ, เตน สา คพฺภํ ลจฺฉติ, สนฺนิปาโต เยเวส คพฺภา\-วกฺกนฺติยา”ติ. “สกฺโกมหํ โกสิย ตํ วจนํ กาตุํ, น ตาวตเกน อมฺหากํ ตโป ภิชฺชติ, โหตู”ติ สมฺปฏิจฺฉิํสุ. ตาย จ ปน เวลาย เทวภวเน อตฺถิ เทวปุตฺโต อุสฺสนฺน\-กุสล\-มูโล ขีณายุโก อายุกฺขย\-ปฺปตฺโต ยทิจฺฉกํ สมตฺโถ โอกฺกมิตุํ อปิ จกฺกวตฺติกุเลปิ. อถ สกฺโก เทวานมินฺโท ตํ เทวปุตฺตํ อุปสงฺกมิตฺวา เอวมาห “เอหิ โข มาริส, สุปภาโต เต ทิวโส, อตฺถสิทฺธิ อุปคตา, ยมหํ เต อุปฏฺฐานมาคมิํ, รมณีเย เต โอกาเส วาโส ภวิสฺสติ, ปติรูเป กุเล ปฏิสนฺธิ ภวิสฺสติ, สุนฺทเรหิ มาตาปิตูหิ วฑฺเฒตพฺโพ, เอหิ เม วจนํ กโรหี”ติ ยาจิ. ทุติยมฺปิ. ตติยมฺปิ ยาจิ สิรสิ ปญฺชลิกโต.}

\prose{\csromanfolio{130}ตโต โส เทวปุตฺโต เอวมาห “กตมํ ตํ มาริส กุลํ, ยํ ตฺวํ อภิกฺขณํ กิตฺตยสิ ปุนปฺปุนนฺ”ติ. ทุกูโล จ ตาปโส ปาริกา จ ตาปสีติ. โส ตสฺส วจนํ สุตฺวา ตุฏฺโฐ สมฺปฏิจฺฉิ “สาธุ มาริส โย ตว ฉนฺโท, โส โหตุ, อากงฺขมาโน อหํ มาริส ปตฺถิเต กุเล อุปฺปชฺเชยฺยํ, กิมฺหิ กุเล อุปฺปชฺชามิ อณฺฑเช วา ชลาพุเช วา สํเสทเช วา โอปปาติเก วา”ติ. ชลาพุชาย มาริส โยนิยา อุปฺปชฺชาหีติ. อถ สกฺโก เทวานมินฺโท อุปฺปตฺติทิวสํ วิคเณตฺวา ทุกูลสฺส ตาปสสฺส อาโรเจสิ “อสุกสฺมิํ นาม ทิวเส ตาปสี อุตุนี ภวิสฺสติ ปุปฺผวตี, ตทา ตฺวํ ภนฺเต ทกฺขิเณน หตฺถ\-งฺคุฏฺเฐน นาภิํ ปรามเสยฺยาสี”ติ. ตสฺมิํ มหาราช ทิวเส ตาปสี จ อุตุนี ปุปฺผวตี อโหสิ, เทวปุตฺโต จ ตตฺถูปโค ปจฺจุปฏฺฐิโต อโหสิ, ตาปโส จ ทกฺขิเณน หตฺถ\-งฺคุฏฺเฐน ตาปสิยา นาภิํ ปรามสิ, อิติ เต ตโย สนฺนิปาตา อเหสุํ, นาภิ\-ปรามสเนน ตาปสิยา ราโค อุทปาทิ, โส ปนสฺสา ราโค นาภิ\-ปรามสนํ ปฏิจฺจ มา ตฺวํ สนฺนิปาตํ อชฺฌาจารเมว มญฺญิ, อูหสนมฺปิ\csromansharedfootnote{dd46f55047c927ce}{หสนมฺปิ (ก)} สนฺนิปาโต, อุลฺลปนมฺปิ สนฺนิปาโต, อุปนิ\-ชฺฌาย\-นมฺปิ สนฺนิปาโต, ปุพฺพภาค\-ภาวโต ราคสฺส อุปฺปาทาย อามสเนน สนฺนิปาโต ชายติ, สนฺนิปาตา โอกฺกมนํ โหตีติ.}

\prose{อนชฺฌาจาเรปิ มหาราช ปรามสเนน คพฺภาวกฺกนฺติ โหติ. ยถา มหาราช อคฺคิ ชลมาโน อปรามสโนปิ อุปคตสฺส สีตํ พฺยปหนฺติ, เอวเมว โข มหาราช อนชฺฌาจาเรปิ ปรามสเนน คพฺภาวกฺกนฺติ โหติ.}

\prose{จตุนฺนํ มหาราช วเสน สตฺตานํ คพฺภาวกฺกนฺติ โหติ กมฺมวเสน โยนิวเสน กุลวเสน อายาจนวเสน, อปิ จ สพฺเพเปเต สตฺตา กมฺมสมฺภวา กมฺมสมุฏฺฐานา.}

\prose{กถํ มหาราช กมฺมวเสน สตฺตานํ คพฺภาวกฺกนฺติ โหติ. อุสฺสนฺน\-กุสล\-มูลา มหาราช สตฺตา ยทิจฺฉกํ อุปฺปชฺชนฺติ ขตฺติย\-มหาสาล\-กุเล วา พฺราหฺมณ\-มหาสาล\-กุเล วา คหปติ\-มหาสาล\-กุเล วา เทเวสุ วา อณฺฑชาย วา โยนิยา ชลาพุชาย วา โยนิยา สํเสทชาย วา โยนิยา โอปปาติกาย วา โยนิยา. ยถา มหาราช ปุริโส \csromanfolio{131}อฑฺโฒ มหทฺธโน มหาโภโค ปหูต\-ชาตรูป\-รชโต ปหูต\-วิตฺตู\-ปกรโณ ปหูต\-ธน\-ธญฺโญ ปหูต\-ญาติ\-ปกฺโข ทาสิํ วา ทาสํ วา เขตฺตํ วา วตฺถุํ วา คามํ วา นิคมํ วา ชนปทํ วา ยํ กิญฺจิ มนสา อภิปตฺถิตํ, ยทิจฺฉกํ ทฺวิคุณติคุณมฺปิ ธนํ ทตฺวา กิณาติ, เอวเมว โข มหาราช อุสฺสนฺน\-กุสล\-มูลา สตฺตา ยทิจฺฉกํ อุปฺปชฺชนฺติ ขตฺติย\-มหาสาล\-กุเล วา พฺราหฺมณ\-มหาสาล\-กุเล วา คหปติ\-มหาสาล\-กุเล วา เทเวสุ วา อณฺฑชาย วา โยนิยา ชลาพุชาย วา โยนิยา สํเสทชาย วา โยนิยา โอปปาติกาย วา โยนิยา. เอวํ กมฺมวเสน สตฺตานํ คพฺภาวกฺกนฺติ โหติ.}

\prose{กถํ โยนิวเสน สตฺตานํ คพฺภาวกฺกนฺติ โหติ. กุกฺกุฏานํ มหาราช วาเตน คพฺภาวกฺกนฺติ โหติ. พลากานํ เมฆสทฺเทน คพฺภาวกฺกนฺติ โหติ. สพฺเพปิ เทวา อคพฺภเสยฺยกา สตฺตา เยว, เตสํ นานาวณฺเณน คพฺภาวกฺกนฺติ โหติ. ยถา มหาราช มนุสฺสา นานาวณฺเณน มหิยา จรนฺติ, เกจิ ปุรโต ปฏิจฺฉาเทนฺติ, เกจิ ปจฺฉโต ปฏิจฺฉาเทนฺติ, เกจิ นคฺคา โหนฺติ, เกจิ ภณฺฑู โหนฺติ เสตปฏธรา, เกจิ โมฬิพทฺธา โหนฺติ, เกจิ ภณฺฑู กาสาววสนา โหนฺติ, เกจิ กาสาววสนา โมฬิพทฺธา โหนฺติ, เกจิ ชฏิโน วากจีรธรา\csromansharedfootnote{d387b9fc987dbb7f}{วากจีรา (ก)} โหนฺติ, เกจิ จมฺมวสนา โหนฺติ, เกจิ รสฺมิโย นิวาเสนฺติ, สพฺเพเปเต มนุสฺสา นานาวณฺเณน มหิยา จรนฺติ. เอวเมว โข มหาราช สตฺตา เยว เต สพฺเพ, เตสํ นานาวณฺเณน คพฺภาวกฺกนฺติ โหติ. เอวํ โยนิวเสน สตฺตานํ คพฺภาวกฺกนฺติ โหติ.}

\prose{กถํ กุลวเสน สตฺตานํ คพฺภาวกฺกนฺติ โหติ. กุลํ นาม มหาราช จตฺตาริ กุลานิ อณฺฑชํ ชลาพุชํ สํเสทชํ โอปปาติกํ. ยทิ ตตฺถ คนฺธพฺโพ ยโต กุโตจิ อาคนฺตฺวา อณฺฑเช กุเล อุปฺปชฺชติ, โส ตตฺถ อณฺฑโช โหติ ฯเปฯ ชลาพุเช กุเล. สํเสทเช กุเล. โอปปาติเก กุเล อุปฺปชฺชติ, โส ตตฺถ โอปปาติโก โหติ. เตสุ เตสุ กุเลสุ ตาทิสา เยว สตฺตา สมฺภวนฺติ. ยถา มหาราช หิมวติ เนรุปพฺพตํ เย เกจิ มิคปกฺขิโน อุเปนฺติ, สพฺเพ \csromanfolio{132}เต สกวณฺณํ วิชหิตฺวา สุวณฺณวณฺณา โหนฺติ, เอวเมว โข มหาราช โย โกจิ คนฺธพฺโพ ยโต กุโตจิ อาคนฺตฺวา อณฺฑชํ โยนิํ อุปคนฺตฺวา สภาววณฺณํ วิชหิตฺวา อณฺฑโช โหติ ฯเปฯ. ชลาพุชํ. สํเสทชํ. โอปปาติกํ โยนิํ อุปคนฺตฺวา สภาววณฺณํ วิชหิตฺวา โอปปาติโก โหติ, เอวํ กุลวเสน สตฺตานํ คพฺภาวกฺกนฺติ โหติ.}

\prose{กถํ อายาจนวเสน สตฺตานํ คพฺภาวกฺกนฺติ โหติ. อิธ มหาราช กุลํ โหติ อปุตฺตกํ พหุสาปเตยฺยํ สทฺธํ ปสนฺนํ สีลวนฺตํ กลฺยาณธมฺมํ ตปนิสฺสิตํ, เทวปุตฺโต จ อุสฺสนฺน\-กุสล\-มูโล จวนธมฺโม โหติ. อถ สกฺโก เทวานมินฺโท ตสฺส กุลสฺส อนุกมฺปาย ตํ เทวปุตฺตํ อายาจติ “ปณิเธหิ มาริส อสุกสฺส กุลสฺส มเหสิยา กุจฺฉินฺ”ติ. โส ตสฺส อายาจนเหตุ ตํ กุลํ ปณิเธติ. ยถา มหาราช มนุสฺสา ปุณฺณกามา สมณํ มโนภาวนียํ อายาจิตฺวา เคหํ อุปเนนฺติ, อยํ อุปคนฺตฺวา สพฺพสฺส กุลสฺส สุขาวโห ภวิสฺสตีติ. เอวเมว โข มหาราช สกฺโก เทวานมินฺโท ตํ เทวปุตฺตํ อายาจิตฺวา ตํ กุลํ อุปเนติ. เอวํ อายาจนวเสน สตฺตานํ คพฺภาวกฺกนฺติ โหติ.}

\prose{สาโม มหาราช กุมาโร สกฺเกน เทวานมินฺเทน อายาจิโต ปาริกาย ตาปสิยา กุจฺฉิํ โอกฺกนฺโต. สาโม มหาราช กุมาโร กตปุญฺโญ, มาตาปิตโร สีลวนฺโต กลฺยาณธมฺมา, อายาจโก สกฺโก, ติณฺณํ เจโตปณิธิยา สาโม กุมาโร นิพฺพตฺโต. อิธ มหาราช นยกุสโล ปุริโส สุกฏฺเฐ อนูปเขตฺเต พีชํ โรเปยฺย, อปิ นุ ตสฺส พีชสฺส อนฺตรายํ วิวชฺเชนฺตสฺส วุฑฺฒิยา โกจิ อนฺตราโย ภเวยฺยาติ. น หิ ภนฺเต, นิรุปฆาตํ พีชํ ขิปฺปํ สํวิรุเหยฺยาติ. เอวเมว โข มหาราช สาโม กุมาโร มุตฺโต อุปฺปนฺน\-นฺตราเยหิ ติณฺณํ เจโตปณิธิยา นิพฺพตฺโต.}

\prose{อปิ นุ โข มหาราช สุตปุพฺพํ ตยา อิสีนํ มโนปโทเสน อิทฺโธ ผีโต มหาชนปโท สชโน สมุจฺฉินฺโนติ. อาม ภนฺเต สุยฺยติ. มหิยา ทณฺฑการญฺญํ\csromansharedfootnote{8c556fd87269f3b0}{ทณฺฑกีรญฺญํ (ม 2. 41 ปิฏฺเฐ.)} มชฺฌารญฺญํ กาลิงฺคารญฺญํ มาตงฺคารญฺญํ, สพฺพํ ตํ อรญฺญํ อรญฺญภูตํ, สพฺเพเปเต ชนปทา อิสีนํ มโนปโทเสน ขยํ คตาติ. \csromanfolio{133}ยทิ มหาราช เตสํ มโนปโทเสน สุสมิทฺธา ชนปทา อุจฺฉิชฺชนฺติ, อปิ นุ โข เตสํ มโนปสาเทน กิญฺจิ นิพฺพตฺเตยฺยาติ. อาม ภนฺเตติ. เตน หิ มหาราช สาโม กุมาโร ติณฺณํ พลวนฺตานํ เจโตปสาเทน นิพฺพตฺโต อิสินิมฺมิโต เทวนิมฺมิโต ปุญฺญนิมฺมิโตติ. เอวเมตํ มหาราช ธาเรหิ.}

\prosewithcloserruled{ตโยเม มหาราช เทวปุตฺตา สกฺเกน เทวานมินฺเทน อายาจิตา กุลํ อุปฺปนฺนา. กตเม ตโย, สาโม กุมาโร มหาปนาโท กุสราชา, ตโยเปเต โพธิสตฺตาติ. สุนิทฺทิฏฺฐา ภนฺเต นาคเสน คพฺภาวกฺกนฺติ, สุกถิตํ การณํ, อนฺธกาโร อาโลโก กโต, ชฏา วิชฏิตา, นิจฺฉุทฺธา ปรวาทา, เอวเมตํ ตถา สมฺปฏิจฺฉามีติ.}{คพฺภาวกฺกนฺติ\-ปโญฺห ฉฏฺโฐ.}
\csromanmatikaanchor{mtk.113}
\csromantocmarkref{subsubsection}{7. สทฺธมฺมนฺตรธานปญฺห}{mtk.113}
\bookmark[dest={mtk.113},level=3]{7. สทฺธมฺมนฺตรธานปญฺห}
\csromanheader{3}{7. \textbf{สทฺธมฺม}\-\textbf{นฺตรธาน}\-\textbf{ปญฺห}}
\proseitem{7}{ภนฺเต นาคเสน ภาสิตมฺเปตํ ภควตา “ปญฺเจว ทานิ อานนฺท วสฺสสตานิ\csromansharedfootnote{613289542f63314c}{วสฺสสหสฺสานิ (สี) ปสฺส วิ 4. 446; อํ 3. 105 ปิฏฺเฐสุ.} สทฺธมฺโม ฐสฺสตี”ติ. ปุน จ ปรินิพฺพาน\-สมเย สุภทฺเทน ปริพฺพาชเกน ปญฺหํ ปุฏฺเฐน ภควตา ภณิตํ “อิเม จ สุภทฺท\csromansharedfootnote{d6f0b29cc7dc3a61}{ที 2. 125 ปิฏฺเฐ ปสฺสิตพฺพํ.} ภิกฺขู สมฺมา วิหเรยฺยุํ, อสุญฺโญ โลโก อรหนฺเตหิ อสฺสา”ติ, อเสส\-วจนเมตํ, นิสฺเสส\-วจนเมตํ, นิปฺปริยาย\-วจนเมตํ. ยทิ ภนฺเต นาคเสน ตถาคเตน ภณิตํ “ปญฺเจว ทานิ อานนฺท วสฺสสตานิ สทฺธมฺโม ฐสฺสตี”ติ, เตน หิ “อสุญฺโญ โลโก อรหนฺเตหิ อสฺสา”ติ ยํ วจนํ, ตํ มิจฺฉา. ยทิ ตถาคเตน ภณิตํ “อสุญฺโญ โลโก อรหนฺเตหิ อสฺสา”ติ, เตน หิ “ปญฺเจว ทานิ อานนฺท วสฺสสตานิ สทฺธมฺโม ฐสฺสตี”ติ ตมฺปิ วจนํ มิจฺฉา, อยมฺปิ อุภโต โกฏิโก ปโญฺห คหนโตปิ คหนตโร พลวโตปิ พลวตโร คณฺฐิโตปิ คณฺฐิตโร, โส ตวานุปฺปตฺโต, ตตฺถ เต ญาณพล\-วิปฺผารํ ทสฺเสหิ มกโร วิย สาครพฺภนฺตรคโตติ.}

\prose{ภาสิตมฺเปตํ มหาราช ภควตา “ปญฺเจว ทานิ อานนฺท วสฺสสตานิ สทฺธมฺโม ฐสฺสตี”ติ. ปรินิพฺพาน\-สมเย จ สุภทฺทสฺส ปริพฺพาชกสฺส \csromanfolio{134}ภณิตํ”อิเม จ สุภทฺท ภิกฺขู สมฺมา วิหเรยฺยุํ, อสุญฺโญ โลโก อรหนฺเตหิ อสฺสา”ติ. ตญฺจ ปน มหาราช ภควโต วจนํ นานตฺถญฺเจว โหติ นานาพฺยญฺชนญฺจ, อยํ สาสน\-ปริจฺเฉโท, อยํ ปฏิปตฺติ ปริทีปนาติ ทูรํ วิวชฺชิตา เต อุโภ อญฺญมญฺญํ. ยถา มหาราช นภํ ปถวิโต ทูรํ วิวชฺชิตํ, นิรยํ สคฺคโต ทูรํ วิวชฺชิตํ, กุสลํ อกุสลโก ทูรํ วิวชฺชิตํ, สุขํ ทุกฺขโต ทูรํ วิวชฺชิตํ. เอวเมว โข มหาราช เต อุโภ อญฺญมญฺญํ ทูรํ วิวชฺชิตา.}

\prose{อปิ จ มหาราช มา เต ปุจฺฉา โมฆา อสฺส\csromansharedfootnote{a4af543e7d227821}{อสฺสุ (สี, สฺยา)}, รสโต เต สํสนฺทิตฺวา กถยิสฺสามิ “ปญฺเจว ทานิ อานนฺท วสฺสสตานิ สทฺธมฺโม ฐสฺสตี”ติ ยํ ภควา อาห, ตํ ขยํ ปริทีปยนฺโต เสสกํ ปริจฺฉินฺทิ, วสฺสสหสฺสํ อานนฺท สทฺธมฺโม ติฏฺเฐยฺย, สเจ ภิกฺขุนิโย น ปพฺพาเชยฺยุํ. ปญฺเจว ทานิ อานนฺท วสฺสสตานิ สทฺธมฺโม ฐสฺสตีติ. อปิ นุ โข มหาราช ภควา เอวํ วทนฺโต สทฺธมฺมสฺส อนฺตรธานํ วา วเทติ อภิสมยํ วา ปฏิกฺโกสตีติ. น หิ ภนฺเตติ. นฏฺฐํ มหาราช ปริกิตฺตยนฺโต เสสกํ ปริทีปยนฺโต ปริจฺฉินฺทิ. ยถา มหาราช ปุริโส นฏฺฐายิโก สพฺพเสสกํ คเหตฺวา ชนสฺส ปริทีเปยฺย “เอตฺตกํ เม ภณฺฑํ นฏฺฐํ, อิทํ เสสกนฺ”ติ. เอวเมว โข มหาราช ภควา นฏฺฐํ ปริทีปยนฺโต เสสกํ เทวมนุสฺสานํ กเถสิ “ปญฺเจว ทานิ อานนฺท วสฺสสตานิ สทฺธมฺโม ฐสฺสภี”ติ. ยํ ปน มหาราช ภควตา ภณิตํ “ปญฺเจว ทานิ อานนฺท วสฺสสตานิ สทฺธมฺโม ฐสฺสตี”ติ, สาสน\-ปริจฺเฉโท เอโส.}

\prose{ยํ ปน ปรินิพฺพาน\-สมเย สุภทฺทสฺส ปริพฺพาชกสฺส สมเณ ปริกิตฺตยนฺโต อาห “อิเม จ สุภทฺท ภิกฺขู สมฺมา วิหเรยฺยุํ, อสุญฺโญ โลโก อรหนฺเตหิ อสฺสา”ติ, ปฏิปตฺติ\-ปริทีปนา เอสา. ตฺวํ ปน ตํ ปริจฺเฉทญฺจ ปริทีปนญฺจ เอกรสํ กโรสิ. ยทิ ปน เต ฉนฺโท, เอกรสํ กตฺวา กถยิสฺสามิ, สาธุกํ สุโณหิ มนสิกโรหิ อวิกฺขิตฺตมานโส\csromansharedfootnote{60e1471505ab6a49}{อวิจลมานโส (สี) อวิมนมานโส (อิ, ก)}.}

\prose{อิธ มหาราช ตฬาโก ภเวยฺย นวสลิล\-สมฺปุณฺโณ สมฺมุขมุ\-ตฺตริย\-มาโน ปริจฺฉินฺโน ปริวฏุมกโต, อปริยาทิณฺเณ เยว \csromanfolio{135}ตสฺมิํ ตฬาเก อุทกูปริ มหาเมโฆ อปราปรํ อนุปฺปพนฺโธ อภิวสฺเสยฺย, อปิ นุ โข มหาราช ตสฺมิํ ตฬาเก อุทกํ ปริกฺขยํ ปริยาทานํ คจฺเฉยฺยาติ. น หิ ภนฺเตติ. เกน การเณน มหาราชาติ. เมฆสฺส ภนฺเต อนุปฺปพนฺธ\-ตายาติ. เอวเมว โข มหาราช ชินสาสน\-วร\-สทฺธมฺม\-ตฬาโก อาจาร\-สีล\-คุณ\-วตฺตปฏิปตฺติ\-วิมล\-นว\-สลิล\-สมฺปุณฺโณ อุตฺตริยมาโน ภวคฺคม\-ภิภวิตฺวา ฐิโต. ยทิ ตตฺถ พุทฺธปุตฺตา อาจาร\-สีล\-คุณ\-วตฺตปฏิปตฺติ\-เมฆ\-วสฺสํ อปราปรํ อนุปฺปพนฺธา\-เปยฺยุํ อภิวสฺสา\-เปยฺยุํ. เอวมิทํ ชินสาสน\-วร\-สทฺธมฺม\-ตฬาโก จิรํ ทีฆมทฺธานํ ติฏฺเฐยฺย, อรหนฺเตหิ โลโก อสุญฺโญ ภเวยฺย, อิมมตฺถํ ภควตา สนฺธาย ภาสิตํ “อิเม จ สุภทฺท ภิกฺขู สมฺมา วิหเรยฺยุํ, อสุญฺโญ โลโก อรหนฺเตหิ อสฺสา”ติ.}

\prose{อิธ ปน มหาราช มหติ มหา\-อคฺคิกฺขนฺเธ ชลมาเน อปราปรํ สุกฺข\-ติณกฏฺฐ\-โคมยานิ อุปสํหเรยฺยุํ, อปิ นุ โข โส มหาราช อคฺคิกฺขนฺโธ นิพฺพาเยยฺยาติ. น หิ ภนฺเต, ภิยฺโย ภิยฺโย โส อคฺคิกฺขนฺโธ ชเลยฺย, ภิยฺโย ภิยฺโย ปภาเสยฺยาติ. เอวเมว โข มหาราช ทสสหสฺสิยา\csromansharedfootnote{f502cf67557e353e}{ทสสหสฺสิมฺหิ (พหูสุ)} โลกธาตุยา ชินสาสนวรมฺปิ อาจาร\-สีล\-คุณ\-วตฺตปฏิปตฺติยา ชลติ ปภาสติ. ยทิ ปน มหาราช ตทุตฺตริํ พุทฺธปุตฺตา ปญฺจหิ ปธานิยงฺเคหิ สมนฺนาคตา สตตม\-ปฺปมตฺตา ปทเหยฺยุํ, ตีสุ สิกฺขาสุ ฉนฺทชาตา สิกฺเขยฺยุํ, จาริตฺตญฺจ สีลํ สมตฺตํ ปริปูเรยฺยุํ, เอวมิทํ ชินสาสนวรํ ภิยฺโย ภิยฺโย จิรํ ทีฆมทฺธานํ ติฏฺเฐยฺย, อสุญฺโญ โลโก อรหนฺเตหิ อสฺสาติ อิมมตฺถํ ภควตา สนฺธาย ภาสิตํ “อิเม จ สุภทฺท ภิกฺขู สมฺมา วิหเรยฺยุํ, อสุญฺโญ โลโก อรหนฺเตหิ อสฺสา”ติ.}

\prose{อิธ ปน มหาราช สินิทฺธสมสุมชฺชิตสปฺปภาสวิมลาทาสํ\csromansharedfootnote{8916a2e2b404ce5f}{สปฺปภํ สุวิมลาทาสํ (สี)} สณฺหสุขุม\-เครุก\-จุณฺเณน อปราปรํ มชฺเชยฺยุํ, อปิ นุ โข มหาราช ตสฺมิํ อาทาเส มลกทฺทม\-รโชชลฺลํ ชาเยยฺยาติ. น หิ ภนฺเต, อญฺญทตฺถุ วิมลตรํ เยว ภเวยฺยาติ. เอวเมว โข มหาราช ชินสาสนวรํ ปกตินิมฺมลํ พฺยปคตกิเลสมลรโชชลฺลํ, ยทิ ตํ พุทฺธปุตฺตา อาจาร\-สีล\-คุณ\-วตฺตปฏิปตฺติ\-สลฺเลข\-ธุตคุเณน ชินสาสนวรํ สลฺลกฺเขยฺยุํ\csromansharedfootnote{f9cff8e4575515fc}{สลฺลิกฺเขยฺยุํ (สี, อิ)}, เอวมิทํ ชินสาสนวรํ จิรํ ทีฆมทฺธานํ ติฏฺเฐยฺย, อสุญฺโญ จ โลโก อรหนฺเตหิ \csromanfolio{136}อสฺสาติ อิมมตฺถํ ภควตา สนฺธาย ภาสิตํ “อิเม จ สุภทฺท ภิกฺขู สมฺมา วิหเรยฺยุํ, อสุญฺโญ โลโก อรหนฺเตหิ อสฺสา”ติ. ปฏิปตฺติมูลกํ มหาราช สตฺถุสาสนํ ปฏิปตฺติ\-การณํ ปฏิปตฺติยา อนนฺตรหิตาย ติฏฺฐตีติ.}

\proseitem{4}{ภนฺเต นาคเสน “สทฺธมฺมนฺตรธานนฺ”ติ ยํ วเทสิ, กตมํ ตํ สทฺธมฺมนฺตรธานนฺติ. ตีณิมานิ มหาราช สาส\-นนฺตรธานานิ. กตมานิ ตีณิ, อธิคม\-นฺตรธานํ ปฏิปตฺต\-นฺตรธานํ ลิงฺค\-นฺตรธานํ, อธิคเม มหาราช อนฺตรหิเต สุปฺปฏิปนฺนสฺสาปิ ธมฺมา\-ภิสมโย น โหติ, ปฏิปตฺติยา อนฺตรหิตาย สิกฺขาปท\-ปญฺญตฺติ อนฺตรธายติ, ลิงฺคํ เยว ติฏฺฐติ, ลิงฺเค อนฺตรหิเต ปเวณุปจฺเฉโท โหติ, อิมานิ โข มหาราช ตีณิ อนฺตรธานานีติ.}

\prosewithcloserruled{สุวิญฺญาปิโต ภนฺเต นาคเสน ปโญฺห, คมฺภีโร อุตฺตานีกโต, คณฺฐิ ภินฺโน, นฏฺฐา ปรวาทา ภคฺคา, นิปฺปภา กตา, ตฺวํ คณิวร\-วสภมา\-สชฺชาติ.}{สทฺธมฺม\-นฺตรธาน\-ปโญฺห สตฺตโม.}
\csromanmatikaanchor{mtk.114}
\csromantocmarkref{subsubsection}{8. อกุสลจฺเฉทนปญฺห}{mtk.114}
\bookmark[dest={mtk.114},level=3]{8. อกุสลจฺเฉทนปญฺห}
\csromanheader{3}{8. \textbf{อกุสล}\-\textbf{จฺเฉทน}\-\textbf{ปญฺห}}
\proseitem{8}{ภนฺเต นาคเสน ตถาคโต สพฺพํ อกุสลํ ฌาเปตฺวา สพฺพญฺญุตํ ปตฺโต, อุทาหุ สาวเสเส อกุสเล สพฺพญฺญุตํ ปตฺโตติ. สพฺพํ มหาราช อกุสลํ ฌาเปตฺวา ภควา สพฺพญฺญุตํ ปตฺโต, นตฺถิ ภควโต เสสกํ อกุสลนฺติ.}

\prose{กิํ ปน ภนฺเต ทุกฺขา เวทนา ตถาคตสฺส กาเย อุปฺปนฺนปุพฺพาติ. อาม มหาราช,\csromansymbolfootnote{*}{วิ 4. 354 ปิฏฺเฐ.}ราชคเห ภควโต ปาโท สกลิกาย\csromansharedfootnote{2b6bdff839eece1c}{สกฺขลิกาย (สฺยา, ก)} ขโต, โลหิตปกฺขนฺทิกา\-พาโธ อุปฺปนฺโน, กาเย อภิสนฺเน ชีวเกน วิเรโก การิโต, วาตาพาเธ อุปฺปนฺเน อุปฏฺฐาเกน เถเรน อุโณฺหทกํ ปริยิฏฺฐนฺติ.}

\prose{\csromanfolio{137}ยทิ ภนฺเต นาคเสน ตถาคโต สพฺพํ อกุสลํ ฌาเปตฺวา สพฺพญฺญุตํ ปตฺโต, เตน หิ ภควโต ปาโท สกลิกาย ขโต, โลหิต\-ปกฺขนฺทิกา จ อาพาโธ อุปฺปนฺโนติ ยํ วจนํ, ตํ มิจฺฉา. ยทิ ตถาคตสฺส ปาโท สกลิกาย ขโต, โลหิต\-ปกฺขนฺทิกา จ อาพาโธ อุปฺปนฺโน, เตน หิ ตถาคโต สพฺพํ อกุสลํ ฌาเปตฺวา สพฺพญฺญุตํ ปตฺโตติ ตมฺปิ วจนํ มิจฺฉา, นตฺถิ ภนฺเต วินา กมฺเมน เวทยิตํ, สพฺพํ ตํ เวทยิตํ กมฺมมูลกํ, ตํ กมฺเมเนว เวทยติ, อยมฺปิ อุภโต โกฏิโก ปโญฺห ตวานุปฺปตฺโต, โส ตยา นิพฺพาหิตพฺโพติ\csromansharedfootnote{accd0ce25904992b}{นิพฺพายิตพฺโพติ (ก)}.}

\prose{น หิ มหาราช สพฺพํ ตํ เวทยิตํ กมฺมมูลกํ, อฏฺฐหิ มหาราช การเณหิ เวทยิตานิ อุปฺปชฺชนฺติ, เยหิ การเณหิ ปุถู สตฺตา เวทนา เวทิยนฺติ. กตเมหิ อฏฺฐหิ, วาต\-สมุฏฺฐานานิปิ โข มหาราช อิเธกจฺจานิ เวทยิตานิ อุปฺปชฺชนฺติ, ปิตฺต\-สมุฏฺฐานานิปิ โข มหาราช ฯเปฯ เสมฺห\-สมุฏฺฐานานิปิ โข มหาราช ฯเปฯ สนฺนิปาติกานิปิ โข มหาราช ฯเปฯ อุตุปริณาม\-ชานิปิ โข มหาราช ฯเปฯ วิสม\-ปริหาร\-ชานิปิ โข มหาราช ฯเปฯ โอปกฺกมิกานิปิ โข มหาราช ฯเปฯ กมฺม\-วิปาก\-ชานิปิ โข มหาราช อิเธกจฺจานิ เวทยิตานิ อุปฺปชฺชนฺติ, อิเมหิ โข มหาราช อฏฺฐหิ การเณหิ ปุถู สตฺตา เวทนา เวทยนฺติ. ตตฺถ เย เต ปุคฺคลา “สตฺเต กมฺมํ วิพาธตี”ติ วเทยฺยุํ, เต อิเม ปุคฺคลา สตฺตการณํ ปฏิพาหนฺติ. เตสํ ตํ วจนํ มิจฺฉาติ. ภนฺเต นาคเสน ยญฺจ วาติกํ ยญฺจ ปิตฺติกํ ยญฺจ เสมฺหิกํ ยญฺจ สนฺนิปาติกํ ยญฺจ อุตุปริณามชํ ยญฺจ วิสม\-ปริหารชํ ยญฺจ โอปกฺกมิกํ, สพฺเพเต กมฺมสมุฏฺฐานา เยว, กมฺเมเนว เต สพฺเพ สมฺภวนฺตีติ.}

\prose{ยทิ มหาราช เตปิ สพฺเพ กมฺมสมุฏฺฐานาว อาพาธา ภเวยฺยุํ, น เตสํ โกฏฺฐาสโต ลกฺขณานิ ภเวยฺยุํ. วาโต โข มหาราช กุปฺปมาโน ทสวิเธน กุปฺปติ สีเตน อุเณฺหน ชิฆจฺฉาย ปิปาสาย อติภุตฺเตน ฐาเนน ปธาเนน อาธาวเนน อุปกฺกเมน กมฺมวิปาเกน. ตตฺร เย เต นว วิธา, น เต อตีเต, น อนาคเต, วตฺตมานเก ภเว อุปฺปชฺชนฺติ, ตสฺมา น วตฺตพฺพา “กมฺมสมฺภวา สพฺพา เวทนี”ติ. ปิตฺตํ มหาราช กุปฺปมานํ ติวิเธน กุปฺปติ สีเตน อุเณฺหน วิสม\-โภชเนน. เสมฺหํ มหาราช กุปฺปมานํ ติวิเธน กุปฺปติ สีเตน อุเณฺหน อนฺนปาเนน.}

\prose{\csromanfolio{138}โย จ มหาราช วาโต ยญฺจ ปิตฺตํ ยญฺจ เสมฺหํ, เตหิ เตหิ โกเปหิ กุปฺปิตฺวา มิสฺสี หุตฺวา สกํ สกํ เวทนํ อากฑฺฒติ. อุตุปริณามชา มหาราช เวทนา อุตุปริณาเมน อุปฺปชฺชติ. วิสม\-ปริหารชา เวทนา วิสม\-ปริหาเรน อุปฺปชฺชติ. โอปกฺกมิกา มหาราช เวทนา อตฺถิ กิริยา, อตฺถิ กมฺมวิปากา, กมฺมวิปากชา เวทนา ปุพฺเพ กเตน กมฺเมน อุปฺปชฺชติ. อิติ โข มหาราช อปฺปํ กมฺมวิปากชํ, พหุตรํ อวเสสํ. ตตฺถ พาลา “สพฺพํ กมฺมวิปากชํ เยวา”ติ อติธาวนฺติ. ตํ กมฺมํ น สกฺกา วินา พุทฺธญาเณน ววตฺถานํ กาตุํ.}

\prose{ยํ ปน มหาราช ภควโต ปาโท สกลิกาย ขโต, ตํ เวทยิตํ เนว วาต\-สมุฏฺฐานํ, น ปิตฺต\-สมุฏฺฐานํ, น เสมฺห\-สมุฏฺฐานํ, น สนฺนิปาติกํ, น อุตุปริณามชํ, น วิสม\-ปริหารชํ, น กมฺมวิปากชํ, โอปกฺกมิกํ เยว. เทวทตฺโต หิ มหาราช พหูนิ ชาติ\-สต\-สหสฺสานิ ตถาคเต อาฆาตํ พนฺธิ, โส เตน อาฆาเตน มหติํ ครุํ สิลํ คเหตฺวา “มตฺถเก ปาเตสฺสามี”ติ มุญฺจิ, อถญฺเญ เทฺว เสลา อาคนฺตฺวา ตํ สิลํ ตถาคตํ อสมฺปตฺตํ เยว สมฺปฏิจฺฉิํสุ, ตาสํ ปหาเรน ปปฏิกา ภิชฺชิตฺวา ภควโต ปาเท ปติตฺวา รุหิรํ\csromansharedfootnote{8a3afbe62c8ea859}{นิปติตฺวา รุธิรํ (สฺยา)} อุปฺปาเทสิ, กมฺมวิปากโต วา มหาราช ภควโต เอสา เวทนา นิพฺพตฺตา กิริยโต วา, ตตุทฺธํ นตฺถญฺญา เวทนา.}

\prose{ยถา มหาราช เขตฺต\-ทุฏฺฐ\-ตาย วา พีชํ น สมฺภวติ พีชทุฏฺฐตาย วา. เอวเมว โข มหาราช กมฺมวิปากโต วา ภควโต เอสา เวทนา นิพฺพตฺตา กิริยโต วา, ตตุทฺธํ นตฺถญฺญา เวทนา.}

\prose{ยถา วา ปน มหาราช โกฏฺฐ\-ทุฏฺฐ\-ตาย วา โภชนํ วิสมํ ปริณมติ อาหาร\-ทุฏฺฐ\-ตาย วา, เอวเมว โข มหาราช กมฺมวิปากโต วา ภควโต เอสา เวทนา นิพฺพตฺตา กิริยโต วา, ตตุทฺธํ นตฺถญฺญา เวทนา. อปิ จ มหาราช นตฺถิ ภควโต กมฺมวิปากชา เวทนา, นตฺถิ วิสม\-ปริหารชา เวทนา, อวเสเสหิ สมุฏฺฐาเนหิ ภควโต เวทนา อุปฺปชฺชติ, ตาย จ ปน เวทนาย น สกฺกา ภควนฺตํ ชีวิตา โวโรเปตุํ.}

\prose{\csromanfolio{139}นิปตนฺติ มหาราช อิมสฺมิํ จาตุมหาภูติเก\csromansharedfootnote{e3042e26310f78e8}{จาตุมฺมหาภูติเก (สี) 2. โมลิยสิวเก (สฺยา, ก) สํ 2. 428 ปิฏฺเฐ ปสฺสิตพฺพํ.} กาเย อิฏฺฐานิฏฺฐา สุภาสุภ\-เวทนา. อิธ มหาราช อากาเส ขิตฺโต เลฑฺฑุ มหาปถวิยา นิปตติ, อปิ นุ โข โส มหาราช เลฑฺฑุ ปุพฺเพ กเตน มหาปถวิยา นิปตีติ. น หิ ภนฺเต, นตฺถิ โส ภนฺเต เหตุ มหาปถวิยา, เยน เหตุนา มหาปถวี กุสลา\-กุสลวิปากํ ปฏิสํเวเทยฺย, ปจฺจุปฺปนฺเนน ภนฺเต อกมฺมเกน เหตุนา โส เลฑฺฑุ มหาปถวิยํ นิปตติ. ยถา มหาราช มหาปถวี, เอวํ ตถาคโต ทฏฺฐพฺโพ. ยถา เลฑฺฑุ ปุพฺเพ อกเตน มหาปถวิยํ นิปตติ, เอวเมว โข มหาราช ตถาคตสฺส ปุพฺเพ อกเตน สา สกลิกา ปาเท นิปติตา.}

\prose{อิธ ปน มหาราช มนุสฺสา มหาปถวิํ ภินฺทนฺติ จ ขณนฺติ จ, อปิ นุ โข มหาราช เต มนุสฺสา ปุพฺเพ กเตน มหาปถวิํ ภินฺทนฺติ จ ขณนฺติ จาติ. น หิ ภนฺเตติ. เอวเมว โข มหาราช ยา สา สกลิกา ภควโต ปาเท นิปติตา, น สา สกลิกา ปุพฺเพ กเตน ภควโต ปาเท นิปติตา. โยปิ มหาราช ภควโต โลหิตปกฺขนฺทิกา\-พาโธ อุปฺปนฺโน, โสปิ อาพาโธ น ปุพฺเพ กเตน อุปฺปนฺโน, สนฺนิปาติเกเนว อุปฺปนฺโน, เย เกจิ มหาราช ภควโต กายิกา อาพาธา อุปฺปนฺนา, น เต กมฺมา\-ภินิพฺพตฺตา, ฉนฺนํ เอเตสํ สมุฏฺฐานานํ อญฺญตรโต นิพฺพตฺตา.}

\prose{ภาสิตมฺเปตํ มหาราช ภควตา เทวาติเทเวน สํยุตฺตนิกาย\-วร\-ลญฺฉเก โมฬิยสีวเก\csromansharedfootnote{97bfe227ef14287c}{โมลิยสิวเก (สฺยา, ก) สํ 2. 428 ปิฏฺเฐ ปสฺสิตพฺพํ.} เวยฺยากรเณ\csromandash{} “ปิตฺต\-สมุฏฺฐานานิปิ โข สีวก อิเธกจฺจานิ เวทยิตานิ อุปฺปชฺชนฺติ. สามมฺปิ โข เอตํ สีวก เวทิตพฺพํ, ยถา ปิตฺต\-สมุฏฺฐานานิปิ อิเธกจฺจานิ เวทยิตานิ อุปฺปชฺชนฺติ. โลกสฺสปิ โข เอตํ สีวก สจฺจสมฺมตํ, ยถา ปิตฺต\-สมุฏฺฐานานิปิ อิเธกจฺจานิ เวทยิตานิ อุปฺปชฺชนฺติ. ตตฺร สีวก เย เต สมณพฺราหฺมณา เอวํวาทิโน เอวํทิฏฺฐิโน “ยํ กิญฺจายํ ปุริสปุคฺคโล ปฏิสํเวเทติ สุขํ วา ทุกฺขํ วา อทุกฺขมสุขํ วา, สพฺพํ ตํ ปุพฺเพ กตเหตู”ติ. ยญฺจ สามํ ญาตํ, ตญฺจ อติธาวนฺติ, ยญฺจ โลเก สจฺจสมฺมตํ, ตญฺจ อติธาวนฺติ. ตสฺมา เตสํ สมณ\-พฺราหฺมณานํ มิจฺฉาติ วทามิ.}

\prose{\csromanfolio{140}เสมฺห\-สมุฏฺฐานานิปิ โข สีวก อิเธกจฺจานิ เวทยิตานิ อุปฺปชฺชนฺติ. วาต\-สมุฏฺฐานานิปิ โข สีวก ฯเปฯ สนฺนิปาติกานิปิ โข สีวก ฯเปฯ อุตุปริณาม\-ชานิปิ โข สีวก ฯเปฯ วิสม\-ปริหาร\-ชานิปิ โข สีวก ฯเปฯ โอปกฺกมิกานิปิ โข สีวก ฯเปฯ กมฺม\-วิปาก\-ชานิปิ โข สีวก อิเธกจฺจานิ เวทยิตานิ อุปฺปชฺชนฺติ. สามมฺปิ โข เอตํ สีวก เวทิตพฺพํ, ยถา กมฺม\-วิปาก\-ชานิปิ อิเธกจฺจานิ เวทยิตานิ อุปฺปชฺชนฺติ. โลกสฺสปิ โข เอตํ สีวก สจฺจสมฺมตํ, ยถา กมฺม\-วิปาก\-ชานิปิ อิเธกจฺจานิ เวทยิตานิ อุปฺปชฺชนฺติ. ตตฺร สีวก เย เต สมณพฺราหฺมณา เอวํวาทิโน เอวํทิฏฺฐิโน “ยํ กิญฺจายํ ปุริสปุคฺคโล ปฏิสํเวเทติ สุขํ วา ทุกฺขํ วา อทุกฺขมสุขํ วา, สพฺพํ ตํ ปุพฺเพ กตเหตู”ติ. ยญฺจ สามํ ญาตํ, ตญฺจ อติธาวนฺติ, ยญฺจ โลเก สจฺจสมฺมตํ, ตญฺจ อติธาวนฺติ. ตสฺมา เตสํ สมณ\-พฺราหฺมณานํ มิจฺฉาติ วทามี”ติ.}

\prosewithcloserruled{อิติปิ มหาราช น สพฺพา เวทนา กมฺมวิปากชา, สพฺพํ มหาราช อกุสลํ ฌาเปตฺวา ภควา สพฺพญฺญุตํ ปตฺโตติ เอวเมตํ ธาเรหีติ. สาธุ ภนฺเต นาคเสน, เอวเมตํ ตถา สมฺปฏิจฺฉามีติ.}{อกุสล\-จฺเฉทน\-ปโญฺห อฏฺฐโม.}
\csromanmatikaanchor{mtk.115}
\csromantocmarkref{subsubsection}{9. อุตฺตริกรณียปญฺห}{mtk.115}
\bookmark[dest={mtk.115},level=3]{9. อุตฺตริกรณียปญฺห}
\csromanheader{3}{9. \textbf{อุตฺตริกรณีย}\-\textbf{ปญฺห}}
\proseitem{9}{ภนฺเต นาคเสน ตุเมฺห ภณถ “ยํ กิญฺจิ กรณียํ ตถาคตสฺส, สพฺพํ ตํ โพธิยา เยว มูเล ปรินิฏฺฐิตํ, นตฺถิ ตถาคตสฺส อุตฺตริํ กรณียํ, กตสฺส วา ปติจโย”ติ, อิทญฺจ\csromansymbolfootnote{*}{วิ 1. 336 ปิฏฺเฐ.}เตมาสํ ปฏิสลฺลานํ ทิสฺสติ. ยทิ ภนฺเต นาคเสน ยํ กิญฺจิ กรณียํ ตถาคตสฺส, สพฺพํ ตํ โพธิยา เยว มูเล ปรินิฏฺฐิตํ, นตฺถิ ตถาคตสฺส อุตฺตริํ กรณียํ, กตสฺส วา ปติจโย, เตน หิ “เตมาสํ ปฏิสลฺลีโน”ติ ยํ วจนํ, ตํ มิจฺฉา. ยทิ เตมาสํ ปฏิสลฺลีโน, เตน หิ “ยํ กิญฺจิ กรณียํ, ตถาคตสฺส, สพฺพํ ตํ โพธิยา เยว มูเล ปรินิฏฺฐิตนฺ”ติ ตมฺปิ วจนํ \csromanfolio{141}มิจฺฉา. นตฺถิ กตกรณียสฺส ปฏิสลฺลานํ, สกรณียสฺเสว ปฏิสลฺลานํ, ยถา นาม พฺยาธิตสฺเสว เภสชฺเชน กรณียํ โหติ, อพฺยาธิตสฺส กิํ เภสชฺเชน. ฉาตสฺเสว โภชเนน กรณียํ โหติ, อฉาตสฺส กิํ โภชเนน. เอวเมว โข ภนฺเต นาคเสน นตฺถิ กตกรณียสฺส ปฏิสลฺลานํ, สกรณียสฺเสว ปฏิสลฺลานํ. อยมฺปิ อุภโต โกฏิโก ปโญฺห ตวานุปฺปตฺโต, โส ตยา นิพฺพาหิตพฺโพติ.}

\prose{ยํ กิญฺจิ มหาราช กรณียํ ตถาคตสฺส, สพฺพํ ตํ โพธิยา เยว มูเล ปรินิฏฺฐิตํ, นตฺถิ ตถาคตสฺส อุตฺตริํ กรณียํ, กตสฺส วา ปติจโย, ภควา จ เตมาสํ ปฏิสลฺลีโน, ปฏิสลฺลานํ โข มหาราช พหุคุณํ, สพฺเพปิ ตถาคตา ปฏิสลฺลียิตฺวา สพฺพญฺญุตํ ปตฺตา, ตํ เต สุกต\-คุณม\-นุสฺสรนฺตา ปฏิสลฺลานํ เสวนฺติ. ยถา มหาราช ปุริโส รญฺโญ สนฺติกา ลทฺธวโร ปฏิลทฺธ\-โภโค ตํ สุกต\-คุณม\-นุสฺสรนฺโต อปราปรํ รญฺโญ อุปฏฺฐานํ เอติ. เอวเมว โข มหาราช สพฺเพปิ ตถาคตา ปฏิสลฺลียิตฺวา สพฺพญฺญุตํ ปตฺตา, ตํ เต สุกต\-คุณม\-นุสฺสรนฺตา ปฏิสลฺลานํ เสวนฺติ.}

\prose{ยถา วา ปน มหาราช ปุริโส อาตุโร ทุกฺขิโต พาฬฺหคิลาโน ภิสกฺกมุปเสวิตฺวา โสตฺถิม\-นุปฺปตฺโต ตํ สุกต\-คุณม\-นุสฺสรนฺโต อปราปรํ ภิสกฺกมุ\-ปเสวติ. เอวเมว โข มหาราช สพฺเพปิ ตถาคตา ปฏิสลฺลียิตฺวา สพฺพญฺญุตํ ปตฺตา, ตํ เต สุกต\-คุณม\-นุสฺสรนฺตา ปฏิสลฺลานํ เสวนฺติ.}

\prose{อฏฺฐวีสติ โข ปนิเม มหาราช ปฏิสลฺลาน\-คุณา, เย คุเณ สมนุสฺสรนฺตา\csromansharedfootnote{bb71bc784b89a4d6}{สมนุปสฺสนฺตา (สี, อิ)} ตถาคตา ปฏิสลฺลานํ เสวนฺติ. กตเม อฏฺฐวีสติ, อิธ มหาราช ปฏิสลฺลานํ ปฏิสลฺลีย\-มานํ อตฺตานํ รกฺขติ, อายุํ วฑฺเฒติ, พลํ เทติ, วชฺชํ ปิทหติ, อยสมปเนติ, ยสมุปเนติ, อรติํ วิโนเทติ, รติมุปทหติ, ภยมปเนติ, เวสารชฺชํ กโรติ, โกสชฺชม\-ปเนติ, วีริยมภิชเนติ, ราคมปเนติ, โทสมปเนติ, โมหมปเนติ, มานํ นิหนฺติ, วิตกฺกํ ภญฺชติ, จิตฺตํ เอกคฺคํ กโรติ, มานสํ สฺเนหยติ\csromansharedfootnote{ae25c8bb0e629dae}{โสภยติ (สี)}, หาสํ ชเนติ, ครุกํ กโรติ, ลาภมุ\-ปฺปาทยติ, นมสฺสิยํ \csromanfolio{142}กโรติ, ปีติํ ปาเปติ, ปาโมชฺชํ กโรติ, สงฺขารานํ สภาวํ ทสฺสยติ, ภวปฺปฏิสนฺธิํ อุคฺฆาเฏติ, สพฺพสามญฺญํ เทติ. อิเม โข มหาราช อฏฺฐวีสติ ปฏิสลฺลาน\-คุณา, เย คุเณ สมนุสฺสรนฺตา ตถาคตา ปฏิสลฺลานํ เสวนฺติ.}

\proseitemwithcloserruled{4}{อปิ จ โข มหาราช ตถาคตา สนฺตํ สุขํ สมาปตฺติรติํ อนุภวิตุกามา ปฏิสลฺลานํ เสวนฺติ ปริโยสิต\-สงฺกปฺปา. จตูหิ โข มหาราช การเณหิ ตถาคตา ปฏิสลฺลานํ เสวนฺติ. กตเมหิ จตูหิ, วิหาร\-ผาสุตายปิ มหาราช ตถาคตา ปฏิสลฺลานํ เสวนฺติ, อนวชฺช\-คุณ\-พหุลตายปิ ตถาคตา ปฏิสลฺลานํ เสวนฺติ, อเสส- อริยวีถิโตปิ ตถาคตา ปฏิสลฺลานํ เสวนฺติ, สพฺพพุทฺธานํ ถุตโถมิตวณฺณิตปสตฺถโตปิ ตถาคตา ปฏิสลฺลานํ เสวนฺติ. อิเมหิ โข มหาราช จตูหิ การเณหิ ตถาคตา ปฏิสลฺลานํ เสวนฺติ. อิติ โข มหาราช ตถาคตา ปฏิสลฺลานํ เสวนฺติ น สกรณียตาย, น กตสฺส วา ปติจยาย, อถ โข คุณวิเสส\-ทสฺสาวิ\-ตาย ตถาคตา ปฏิสลฺลานํ เสวนฺตีติ. สาธุ ภนฺเต นาคเสน, เอวเมตํ ตถา สมฺปฏิจฺฉามีติ.}{อุตฺตริกรณีย\-ปโญฺห นวโม.}
\csromanmatikaanchor{mtk.116}
\csromantocmarkref{subsubsection}{10. อิทฺธิพลทสฺสนปญฺห}{mtk.116}
\bookmark[dest={mtk.116},level=3]{10. อิทฺธิพลทสฺสนปญฺห}
\csromanheader{3}{10. \textbf{อิทฺธิพล}\-\textbf{ทสฺสน}\-\textbf{ปญฺห}}
\proseitem{10}{ภนฺเต นาคเสน ภาสิตมฺเปตํ ภควตา\csromansymbolfootnote{*}{ที 2. 86; สํ 3. 227; อํ 3. 129; ขุ 1. 152 ปิฏฺเฐสุ.}“ตถาคตสฺส โข อานนฺท จตฺตาโร อิทฺธิปาทา ภาวิตา พหุลีกตา ยานีกตา วตฺถุกตา อนุฏฺฐิตา ปริจิตา สุสมารทฺธา, โส อากงฺขมาโน อานนฺท ตถาคโต กปฺปํ วา ติฏฺเฐยฺย กปฺปาวเสสํ วา”ติ. ปุน จ ภณิตํ\csromansymbolfootnote{+}{ที 2. 89; สํ 3. 230; อํ 3. 131; ขุ 1. 155 ปิฏฺเฐสุ.}“อิโต ติณฺณํ มาสานํ อจฺจเยน ตถาคโต ปรินิพฺพายิสฺสตี”ติ. ยทิ ภนฺเต นาคเสน ภควตา ภณิตํ “ตถาคตสฺส โข อานนฺท จตฺตาโร อิทฺธิปาทา ภาวิตา ฯเปฯ กปฺปาวเสสํ วา”ติ, เตน หิ เตมาส\-ปริจฺเฉโท มิจฺฉา. ยทิ ภนฺเต ตถาคเตน ภณิตํ “อิโต ติณฺณํ มาสานํ อจฺจเยน ตถาคโต ปรินิพฺพายิสฺสตี”ติ, เตน หิ “ตถาคตสฺส โข อานนฺท จตฺตาโร อิทฺธิปาทา ภาวิตา ฯเปฯ กปฺปาวเสสํ \csromanfolio{143}วา”ติ ตมฺปิ วจนํ มิจฺฉา. นตฺถิ ตถาคตานํ อฏฺฐาเน คชฺชิตํ. อโมฆวจนา พุทฺธา ภควนฺโต ตถวจนา อเทฺวชฺฌวจนา. อยมฺปิ อุภโต โกฏิโก ปโญฺห คมฺภีโร สุนิปุโณ ทุนฺนิชฺฌาปโย ตวานุปฺปตฺโต, ภินฺเทตํ ทิฏฺฐิชาลํ, เอกํเส ฐปย, ภินฺท ปรวาทนฺติ.}

\prose{ภาสิตมฺเปตํ มหาราช ภควตา “ตถาคตสฺส โข อานนฺท จตฺตาโร อิทฺธิปาทา ภาวิตา ฯเปฯ กปฺปาวเสสํ วา”ติ, เตมาส\-ปริจฺเฉโท จ ภณิโต, โส จ ปน กปฺโป อายุกปฺโป วุจฺจติ. น มหาราช ภควา อตฺตโน พลํ กิตฺตยมาโน เอวมาห, อิทฺธิพลํ ปน มหาราช ภควา ปริกิตฺตยมาโน เอวมาห “ตถาคตสฺส โข อานนฺท จตฺตาโร อิทฺธิปาทา ภาวิตา ฯเปฯ กปฺปาวเสสํ วา”ติ.}

\prose{ยถา มหาราช รญฺโญ อสฺสาชานีโย ภเวยฺย สีฆคติ อนิลชโว, ตสฺส ราชา ชวพลํ ปริกิตฺตยนฺโต สเนคมชานปท\-ภฏ\-พล\-พฺราหฺมณคหปติก\-อมจฺจ\-ชนมชฺเฌ เอวํ วเทยฺย “อากงฺขมาโน เม โภ อยํ หยวโร สาคร\-ชล\-ปริยนฺตํ มหิํ อนุวิจริตฺวา ขเณน อิธาคจฺเฉยฺยา”ติ, น จ ตํ ชวคติํ ตสฺสํ ปริสายํ ทสฺเสยฺย, วิชฺชติ จ โส ชโว ตสฺส, สมตฺโถ จ โส ขเณน สาคร\-ชล\-ปริยนฺตํ มหิํ อนุวิจริตุํ. เอวเมว โข มหาราช ภควา อตฺตโน อิทฺธิพลํ ปริกิตฺตยมาโน เอวมาห, ตมฺปิ เตวิชฺชานํ ฉฬภิญฺญานํ อรหนฺตานํ วิมล\-ขีณาสวานํ เทว\-มนุสฺสา\-นญฺจ มชฺเฌ นิสีทิตฺวา ภณิตํ “ตถาคตสฺส โข อานนฺท จตฺตาโร อิทฺธิปาทา ภาวิตา พหุลีกตา ยานีกตา วตฺถุกตา อนุฏฺฐิตา ปริจิตา สุสมารทฺธา, โส อากงฺขมาโน อานนฺท ตถาคโต กปฺปํ วา ติฏฺเฐยฺย กปฺปาวเสสํ วา”ติ. วิชฺชติ จ ตํ มหาราช อิทฺธิพลํ ภควโต, สมตฺโถ จ ภควา อิทฺธิพเลน กปฺปํ วา ฐาตุํ กปฺปาวเสสํ วา, น จ ภควา ตํ อิทฺธิพลํ ตสฺสํ ปริสายํ ทสฺเสติ, อนตฺถิโก มหาราช ภควา สพฺพภเวหิ, ครหิตา จ ตถาคตสฺส สพฺพภวา. ภาสิตมฺเปตํ มหาราช ภควตา\csromansymbolfootnote{*}{อํ 1. 36 ปิฏฺเฐ.}“เสยฺยถาปิ ภิกฺขเว อปฺปมตฺตโกปิ คูโถ ทุคฺคนฺโธ \csromanfolio{144}โหติ. เอวเมว โข อหํ ภิกฺขเว อปฺปมตฺตกมฺปิ ภวํ น วณฺเณมิ อนฺตมโส อจฺฉราสงฺฆาตมตฺตมฺปี”ติ. อปิ นุ โข มหาราช ภควา สพฺพภว\-คติ\-โยนิโย คูถสมํ ทิสฺวา อิทฺธิพลํ นิสฺสาย ภเวสุ ฉนฺทราคํ กเรยฺยาติ. น หิ ภนฺเตติ. เตน หิ มหาราช ภควา อิทฺธิพลํ ปริกิตฺตยมาโน เอวรูปํ พุทฺธ\-สีหนาทม\-ภินทีติ. สาธุ ภนฺเต นาคเสน, เอวเมตํ ตถา สมฺปฏิจฺฉามีติ.}

\vspace{\tipitakaparskip}
\csromancenter{อิทฺธิพล\-ทสฺสน\-ปโญฺห ทสโม.}
\nitthitammidruled{\textbf{อิทฺธิพลวคฺโค ปฐโม.}}
\csromancenter{อิมสฺมิํ วคฺเค ทส ปญฺหา.}
\csromanmatikaanchor{mtk.117}
\csromantocmarknopage{subsection}{2. อเภชฺชวคฺค}{mtk.117}
\bookmark[dest={mtk.117},level=2]{2. อเภชฺชวคฺค}
\csromanheader{2}{2. \textbf{อเภชฺชวคฺค}}
\csromanmatikaanchor{mtk.118}
\csromantocmarkref{subsubsection}{1. ขุทฺทานุขุทฺทกปญฺห}{mtk.118}
\bookmark[dest={mtk.118},level=3]{1. ขุทฺทานุขุทฺทกปญฺห}
\csromanheader{3}{1. \textbf{ขุทฺทานุขุทฺทก}\-\textbf{ปญฺห}}
\proseitem{1}{\csromanfolio{145}ภนฺเต นาคเสน ภาสิตมฺเปตํ ภควตา\csromansymbolfootnote{*}{อํ 1. 280; อภิ 4. 407 ปิฏฺเฐสุ.}“อภิญฺญายาหํ ภิกฺขเว ธมฺมํ เทเสมิ โน อนภิญฺญายา”ติ. ปุน จ วินย\-ปญฺญตฺติยา เอวํ ภณิตํ\csromansymbolfootnote{+}{วิ 4. 484; ที 2. 127 ปิฏฺเฐสุ.}“อากงฺขมาโน อานนฺท สํโฆ มมจฺจเยน ขุทฺทา\-นุขุทฺท\-กานิ สิกฺขาปทานิ สมูหนตู”ติ. กิํ นุ โข ภนฺเต นาคเสน ขุทฺทา\-นุขุทฺท\-กานิ สิกฺขาปทานิ ทุปฺปญฺญตฺตานิ, อุทาหุ อวตฺถุสฺมิํ อชานิตฺวา ปญฺญตฺตานิ, ยํ ภควา อตฺตโน อจฺจเยน ขุทฺทา\-นุขุทฺท\-กานิ สิกฺขาปทานิ สมูหนาเปติ. ยทิ ภนฺเต นาคเสน ภควตา ภณิตํ “อภิญฺญายาหํ ภิกฺขเว ธมฺมํ เทเสมิโน อนภิญฺญายา”ติ, เตน หิ “อากงฺขมาโน อานนฺท สํโฆ มมจฺจเยน ขุทฺทา\-นุขุทฺท\-กานิ สิกฺขาปทานิ สมูหนตู”ติ ยํ วจนํ, ตํ มิจฺฉา. ยทิ ตถาคเตน วินย\-ปญฺญตฺติยา เอวํ ภณิตํ “อากงฺขมาโน อานนฺท สํโฆ มมจฺจเยน ขุทฺทา\-นุขุทฺท\-กานิ สิกฺขาปทานิ สมูหนตู”ติ. เตน หิ “อภิญฺญายาหํ ภิกฺขเว ธมฺมํ เทเสมิ โน อนภิญฺญายา”ติ ตมฺปิ วจนํ มิจฺฉา. อยมฺปิ อุภโต โกฏิโก ปโญฺห สุขุโม นิปุโณ คมฺภีโร สุคมฺภีโร ทุนฺนิชฺฌาปโย, โส ตวานุปฺปตฺโต, ตตฺถ เต ญาณพล\-วิปฺผารํ ทสฺเสหีติ.}

\prose{ภาสิตมฺเปตํ มหาราช ภควตา “อภิญฺญายาหํ ภิกฺขเว ธมฺมํ เทเสมิ โน อนภิญฺญายา”ติ, วินย\-ปญฺญตฺติยาปิ เอวํ ภณิตํ “อากงฺขมาโน อานนฺท สํโฆ มมจฺจเยน ขุทฺทา\-นุขุทฺท\-กานิ สิกฺขาปทานิ สมูหนตู”ติ, ตํ ปน มหราช ตถาคโต ภิกฺขู วีมํสมาโน อาห “อุกฺกเลสฺสนฺติ\csromansharedfootnote{6f961b1d444a2f06}{อุกฺกฑฺฒิสฺสนฺติ (สี), อุสฺสกฺกิสฺสนฺติ (สฺยา)} นุ โข มม สาวกา มยา วิสฺสชฺชา\-ปีย\-มานา มมจฺจเยน ขุทฺทา\-นุขุทฺท\-กานิ สิกฺขาปทานิ, อุทาหุ อาทิยิสฺสนฺตี”ติ.}

\prose{ยถา มหาราช จกฺกวตฺตี ราชา ปุตฺเต เอวํ วเทยฺย “อยํ โข ตาตา มหาชนปโท สพฺพทิสาสุ สาคร\-ปริยนฺโต, ทุกฺกโร ตาตา ตาวตเกน พเลน ธาเรตุํ, เอถ ตุเมฺห ตาตา มมจฺจเยน ปจฺจนฺเต ปจฺจนฺเต เทเส ปชหถา”ติ. อปิ นุ โข เต มหาราช กุมารา ปิตุอจฺจเยน \csromanfolio{146}หตฺถคเต ชนปเท สพฺเพ เต ปจฺจนฺเต ปจฺเจนฺเต เทเส มุญฺเจยฺยุนฺติ. น หิ ภนฺเต, ราชโต\csromansharedfootnote{786cb3627db81d2c}{ราชาโน (สี, อิ)} ภนฺเต ลุทฺธตรา\csromansharedfootnote{2f62afe9cde89b7d}{ลทฺธตรา (ก)} กุมารา รชฺชโลเภน ตทุตฺตริํ ทิคุณติคุณํ ชนปทํ ปริคฺคเณฺหยฺยุํ\csromansharedfootnote{220a5bb395254e3a}{ปริกฑฺเฒยฺยุํ (สี, อิ)}, กิํ ปน เต หตฺถคตํ ชนปทํ มุญฺเจยฺยุนฺติ. เอวเมว โข มหาราช ตถาคโต ภิกฺขู วีมํสมาโน เอวมาห “อากงฺขมาโน อานนฺท สํโฆ มมจฺจเยน ขุทฺทา\-นุขุทฺท\-กานิ สิกฺขาปทานิ สมูหนตู”ติ. ทุกฺข\-ปริมุตฺติยา มหาราช พุทฺธปุตฺตา ธมฺมโลเภน อญฺญมฺปิ อุตฺตริํ ทิยฑฺฒ\-สิกฺขาปทสตํ โคเปยฺยุํ, กิํ ปน ปกติ ปญฺญตฺตํ สิกฺขาปทํ มุญฺเจยฺยุนฺติ.}

\prosewithcloserruled{ภนฺเต นาคเสน ยํ ภควา อาย “ขุทฺทา\-นุขุทฺท\-กานิ สิกฺขาปทานี”ติ, เอตฺถายํ ชโน สมฺมูโฬฺห วิมติชาโต อธิกโต สํสย\-ปกฺขนฺโท. กตมานิ ตานิ ขุทฺทกานิ สิกฺขาปทานิ, กตมานิ อนุขุทฺทกานิ สิกฺขาปทานีติ. ทุกฺกฏํ มหาราช ขุทฺทกํ สิกฺขาปทํ, ทุพฺภาสิตํ อนุขุทฺทกํ สิกฺขาปทํ, อิมานิ เทฺว ขุทฺทา\-นุขุทฺท\-กานิ สิกฺขาปทานิ, ปุพฺพเกหิปิ มหาราช มหาเถเรหิ เอตฺถ วิมติ อุปฺปาทิตา, เตหิปิ เอกชฺฌํ น กโต ธมฺม\-สณฺฐิติ\-ปริยาเย ภควตา เอโส ปโญฺห อุปทิฏฺโฐติ. จิรนิกฺขิตฺตํ ภนฺเต นาคเสน ชินรหสฺสํ อชฺเชตรหิ โลเก วิวฏํ ปากฏํ กตนฺติ.}{ขุทฺทานุขุทฺทก\-ปโญฺห ปฐโม.}
\csromanmatikaanchor{mtk.119}
\csromantocmarkref{subsubsection}{2. อพฺยากรณียปญฺห}{mtk.119}
\bookmark[dest={mtk.119},level=3]{2. อพฺยากรณียปญฺห}
\csromanheader{3}{2. \textbf{อพฺยากรณียปญฺห}}
\proseitem{1}{ภนฺเต นาคเสน ภาสิตมฺเปตํ ภควตา\csromansymbolfootnote{*}{ที 2. 85; สํ 3. 133 ปิฏฺเฐสุ.}“นตฺถานนฺท ตถาคตสฺส ธมฺเมสุ อาจริยมุฏฺฐี”ติ, ปุน จ เถเรน มาลุกฺยปุตฺเตน\csromansharedfootnote{bacb72088b3a415c}{มาลุงฺกฺยปุตฺเตน (สี, สฺยา, อิ) ม 2. 89; สํ 2. 294 ปิฏฺเฐ; อํ 1. 571 ปิฏฺเฐ จ ปสฺสิตพฺพํ.} ปญฺหํ ปุฏฺโฐ น พฺยากาสิ, เอโส โข ภนฺเต นาคเสน ปโญฺห ทฺวยนฺโต\csromansharedfootnote{5d2298ca57036465}{ทฺวยโต (สี)} เอกนฺตนิสฺสิโต ภวิสฺสติ อชานเนน วา คุยฺหกรเณน วา. ยทิ ภนฺเต นาคเสน ภควตา ภณิตํ “นตฺถานนฺท ตถาคตสฺส ธมฺเมสุ อาจริยมุฏฺฐี”ติ, เตน หิ เถรสฺส มาลุกฺย\-ปุตฺตสฺส อชานนฺเตน \csromanfolio{147}น อพฺยากตํ. ยทิ อชานนฺเตน น พฺยากตํ, เตน หิ อตฺถิ ตถาคตสฺส ธมฺเมสุ อาจริยมุฏฺฐิ. อยมฺปิ อุภโต โกฏิโก ปโญฺห ตวานุปฺปตฺโต, โส ตยา นิพฺพาหิตพฺโพติ.}

\proseitem{2}{ภาสิตมฺเปตํ มหาราช ภควตา “นตฺถานนฺท ตถาคตสฺส ธมฺเมสุ อาจริยมุฏฺฐี”ติ, อพฺยากโต จ เถเรน มาลุกฺยปุตฺเตน ปุจฺฉิโต ปโญฺห, ตญฺจ ปน น อชานนฺเตน น คุยฺหกรเณน. จตฺตาริมานิ มหาราช ปญฺห\-พฺยากรณานิ. กตมานิ จตฺตาริ, เอกํส\-พฺยากรณีโย ปโญฺห วิภชฺช\-พฺยากรณีโย ปโญฺห ปฏิปุจฺฉา\-พฺยากรณีโย ปโญฺห ฐปนีโย ปโญฺหติ.}

\prose{กตโม จ มหาราช เอกํส\-พฺยากรณีโย ปโญฺห, “รูปํ อนิจฺจนฺ”ติ เอกํส\-พฺยากรณีโย ปโญฺห, “เวทนา อนิจฺจา”ติ ฯเปฯ “สญฺญา อนิจฺจา”ติ ฯเปฯ “สงฺขารา อนิจฺจา”ติ ฯเปฯ “วิญฺญาณํ อนิจฺจนฺ”ติ เอกํส\-พฺยากรณีโย ปโญฺห, อยํ เอกํส\-พฺยากรณีโย ปโญฺห.}

\prose{กตโม วิภชฺช\-พฺยากรณีโย ปโญฺห, “อนิจฺจํ ปน รูปนฺ”ติ วิภชฺช\-พฺยากรณีโย ปโญฺห, “อนิจฺจา ปน เวทนา”ติ ฯเปฯ “อนิจฺจา ปน สญฺญา”ติ ฯเปฯ “อนิจฺจา ปน สงฺขารา”ติ ฯเปฯ “อนิจฺจํ ปน วิญฺญาณนฺ”ติ วิภชฺช\-พฺยากรณีโย ปโญฺห, อยํ วิภชฺช\-พฺยากรณีโย ปโญฺห.}

\prose{กตโม ปฏิปุจฺฉา\-พฺยากรณีโย ปโญฺห, “กิํ นุ โข จกฺขุนา สพฺพํ วิชานาตี”ติ อยํ ปฏิปุจฺฉา\-พฺยากรณีโย ปโญฺห.}

\prose{กตโม ฐปนีโย ปโญฺห, “สสฺสโต โลโก”ติ ฐปนีโย ปโญฺห, “อสสฺสโต โลโก”ติ. “อนฺตวา โลโก”ติ. “อนนฺตวา โลโก”ติ. “อนฺตวา จ อนนฺตวา จ โลโก”ติ. “เนวนฺตวา นานนฺตวา โลโก”ติ. “ตํ ชีวํ ตํ สรีรนฺ”ติ. “อญฺญํ ชีวํ อญฺญํ สรีรนฺ”ติ. “โหติ ตถาคโต ปรํ มรณา”ติ. “น โหติ ตถาคโต ปรํ มรณา”ติ. “โหติ จ น จ โหติ ตถาคโต ปรํ มรณา”ติ. “เนว โหติ น น โหติ ตถาคโต ปรํ มรณา”ติ ฐปนีโย ปโญฺห, อยํ ฐปนีโย ปโญฺห.}

\prosewithcloserruled{ภควา มหาราช เถรสฺส มาลุกฺย\-ปุตฺตสฺส ตํ ฐปนียํ ปญฺหํ น พฺยากาสิ. โส ปน ปโญฺห กิํ การณา ฐปนีโย. น ตสฺส ทีปนาย \csromanfolio{148}เหตุ วา การณํ วา อตฺถิ, ตสฺมา โส ปโญฺห ฐปนีโย. นตฺถิ พุทฺธานํ ภควนฺตานํ อการณม\-เหตุกํ คิรมุทีรณนฺติ. สาธุ ภนฺเต นาคเสน, เอวเมตํ ตถา สมฺปฏิจฺฉามีติ.}{อพฺยากรณียปโญฺห ทุติโย.}
\csromanmatikaanchor{mtk.120}
\csromantocmarkref{subsubsection}{3. มจฺจุภายนาภายนปญฺห}{mtk.120}
\bookmark[dest={mtk.120},level=3]{3. มจฺจุภายนาภายนปญฺห}
\csromanheader{3}{3. \textbf{มจฺจุภา}\-\textbf{ยนา}\-\textbf{ภายน}\-\textbf{ปญฺห}}
\proseitem{3}{ภนฺเต นาคเสน ภาสิตมฺเปตํ ภควตา\csromansymbolfootnote{*}{ขุ 1. 32 ปิฏฺเฐ ธมฺมปเท.}“สพฺเพ ตสนฺติ ทณฺฑสฺส, สพฺเพ ภายนฺติ มจฺจุโน”ติ, ปุน ภณิตํ “อรหา สพฺพ\-ภยม\-ติกฺกนฺโต”ติ. กินฺนุ โข ภนฺเต นาคเสน อรหา ทณฺฑภยา ตสติ, นิรเย วา เนรยิกา สตฺตา ชลิตา กุถิตา ตตฺตา สนฺตตฺตา ตมฺหา ชลิต\-คฺคิ\-ชาลกา มหานิรยา จวมานา มจฺจุโน ภายนฺติ. ยทิ ภนฺเต นาคเสน ภควตา ภณิตํ “สพฺเพ ตสนฺติ ทณฺฑสฺส, สพฺเพ ภายนฺติ มจฺจุโน”ติ, เตน หิ “อรหา สพฺพ\-ภยม\-ติกฺกนฺโต”ติ ยํ วจนํ, ตํ มิจฺฉา. ยทิ ภควตา ภณิตํ “อรหา สพฺพ\-ภยม\-ติกฺกนฺโต”ติ, เตน หิ “สพฺเพ ตสนฺติ ทณฺฑสฺส, สพฺเพ ภายนฺติ มจฺจุโน”ติ ตมฺปิ วจนํ มิจฺฉา. อยํ อุภโต โกฏิโก ปโญฺห ตวานุปฺปตฺโต, โส ตยา นิพฺพาหิตพฺโพติ.}

\prose{เนตํ มหาราช วจนํ ภควตา อรหนฺเต อุปาทาย ภณิตํ “สพฺเพ ตสนฺติ ทณฺฑสฺส, สพฺเพ ภายนฺติ มจฺจุโน”ติ. ฐปิโต อรหา ตสฺมิํ วตฺถุสฺมิํ, สมูหโต ภยเหตุ อรหโต. เย เต มหาราช สตฺตา สกิเลสา, เยสญฺจ อธิมตฺตา อตฺตานุทิฏฺฐิ, เย จ สุขทุกฺเขสุ อุนฺนตาวนตา, เต อุปาทาย ภควตา ภณิตํ “สพฺเพ ตสนฺติ ทณฺฑสฺส, สพฺเพ ภายนฺติ มจฺจุโน”ติ. อรหโต มหาราช สพฺพคติ อุปจฺฉินฺนา, โยนิ วิทฺธํสิตา, ปฏิสนฺธิอุปหตา, ภคฺคา ผาสุกา, สมูหตา สพฺพภวาลยา, สมุจฺฉินฺนา สพฺพสงฺขารา, หตํ กุสลากุสลํ, วิหตา อวิชฺชา, อพีชํ วิญฺญาณํ กตํ, ทฑฺฒา สพฺพกิเลสา, อติวตฺตา โลกธมฺมา, ตสฺมา อรหา น ตสติ สพฺพภเยหิ.}

\prose{อิธ มหาราช รญฺโญ จตฺตาโร มหามตฺตา ภเวยฺยุํ อนุรกฺขา ลทฺธยสา วิสฺสาสิกา ฐปิตา มหติ อิสฺสริเย ฐาเน. อถ ราชา \csromanfolio{149}กิสฺมิญฺจิ เทว กรณีเย สมุปฺปนฺเน ยาวตา สกวิชิเต สพฺพชนสฺส อาณาเปยฺย “สพฺเพว เม พลิํ กโรนฺตุ, สาเธถ ตุเมฺห จตฺตาโร มหามตฺตา ตํ กรณียนฺ”ติ. อปิ นุ โข มหาราช เตสํ จตุนฺนํ มหามตฺตานํ พลิภยา สนฺตาโส อุปฺปชฺเชยฺยาติ. น หิ ภนฺเตติ. เกน การเณน มหาราชาติ. ฐปิตา เต ภนฺเต รญฺญา อุตฺตมฏฺฐาเน, นตฺถิ เตสํ พลิ, สมติกฺกนฺต\-พลิโน เต, อวเสเส อุปาทาย รญฺญา อาณาปิตํ “สพฺเพว เม พลิํ กโรนฺตู”ติ. เอวเมว โข มหาราช เนตํ วจนํ ภควตา อรหนฺเต อุปาทาย ภณิตํ, ฐปิโต อรหา ตสฺมิํ วตฺถุสฺมิํ, สมูหโต ภยเหตุ อรหโต, เย เต มหาราช สตฺตา สกิเลสา, เยสญฺจ อธิมตฺตา อตฺตานุทิฏฺฐิ, เย จ สุขทุกฺเขสุ อุนฺนตาวนตา, เต อุปาทาย ภควตา ภณิตํ “สพฺเพ ตสนฺติ ทณฺฑสฺส, สพฺเพ ภายนฺติ มจฺจุโน”ติ. ตสฺมา อรหา น ตสติ สพฺพภเยหีติ.}

\prose{เนตํ ภนฺเต นาคเสน วจนํ สาวเสสํ, นิรวเสส\-วจนเมตํ “สพฺเพ”ติ. ตตฺถ เม อุตฺตริํ การณํ พฺรูหิ ตํ วจนํ ปติฏฺฐาเป\-ตุนฺติ.}

\prose{อิธ มหาราช คาเม คามสฺสามิโก อาณาปกํ อาณาเปยฺย “เอหิ โภ อาณาปก ยาวตา คาเม คามิกา, เต สพฺเพ สีฆํ มม สนฺติเก สนฺนิปาเตหี”ติ. โส “สาธุ สามี”ติ สมฺปฏิจฺฉิตฺวา คามมชฺเฌ ฐตฺวา ติกฺขตฺตุํ สทฺทมนุสฺสาเวยฺย “ยาวตา คาเม คามิกา, เต สพฺเพ สีฆสีฆํ สามิโน สนฺติเก สนฺนิปตนฺตู”ติ. ตโต เต คามิกา อาณาปกสฺส วจเนน ตุริตตุริตา สนฺนิปติตฺวา คามสฺสามิกสฺส อาโรเจนฺติ “สนฺนิปติตา สามิ สพฺเพ คามิกา, ยํ เต กรณียํ ตํ กโรหี”ติ. อิติ โส มหาราช คามสฺสามิโก กุฏิปุริเส สนฺนิปาเตนฺโต สพฺเพ คามิเก อาณาเปติ, เต จ อาณตฺตา น สพฺเพ สนฺนิปตนฺติ, กุฏิปุริสา เยว สนฺนิปตนฺติ, “เอตฺตกา เยว เม คามิกา”ติ คามสฺสามิโก จ ตถา สมฺปฏิจฺฉติ, อญฺเญ พหุตรา อนาคตา อิตฺถิปุริสา ทาสิทาสา ภตกา กมฺมกรา คามิกา คิลานา โคมหิํสา อเชฬกา สุวานา, เย อนาคตา, สพฺเพ เต อคณิตา, กุฏิปุริเส เยว อุปาทาย อาณาปิตตฺตา “สพฺเพ สนฺนิปตนฺตู”ติ. เอวเมว โข มหาราช เนตํ วจนํ ภควตา อรหนฺเต อุปาทาย ภณิตํ, ฐปิโต อรหา ตสฺมิํ วตฺถุสฺมิํ, สมูหโต ภยเหตุ อรหโต, เย เต มหาราช สตฺตา สกิเลสา, เยสญฺจ อธิมตฺตา อตฺตานุทิฏฺฐิ, \csromanfolio{150}เย จ สุขทุกฺเขสุ อุนฺนตาวนตา, เต อุปาทาย ภควตา ภณิตํ “สพฺเพ ตสนฺติ ทณฺฑสฺส, สพฺเพ ภายนฺติ มจฺจุโน”ติ. ตสฺมา อรหา น ตสติ สพฺพภเยหิ.}

\proseitem{4}{อตฺถิ มหาราช สาวเสสํ วจนํ สาวเสโส อตฺโถ, อตฺถิ สาวเสสํ วจนํ นิรวเสโส อตฺโถ, อตฺถิ นิรวเสสํ วจนํ สาวเสโส อตฺโถ, อตฺถิ นิรวเสสํ วจนํ นิรวเสโส อตฺโถ. เตน เตน อตฺโถ สมฺปฏิจฺฉิตพฺโพ.}

\prose{ปญฺจวิเธหิ มหาราช การเณหิ อตฺโถ สมฺปฏิจฺฉิตพฺโพ อาหจฺจปเทน รเสน อาจริยวํเสน\csromansharedfootnote{2dd7ae06cff966dc}{อาจริยวํสตาย (อิ, ก)} อธิปฺปายา การณุ\-ตฺตริ\-ยตาย. เอตฺถ หิ \textbf{อาหจฺจปท}นฺติ สุตฺตํ อธิปฺเปตํ. \textbf{รโส}ติ สุตฺตานุโลมํ. \textbf{อาจริยวํโส}ติ อาจริยวาโท. \textbf{อธิปฺปาโย}ติ อตฺตโน มติ. \textbf{การณุ}\-\textbf{ตฺตริ}\-\textbf{ยตา}ติ อิเมหิ จตูหิ สเมนฺตํ\csromansharedfootnote{8585f396c7ddad2a}{สเมตํ (สี)} การณํ. อิเมหิ โข มหาราช ปญฺจหิ การเณหิ อตฺโถ สมฺปฏิจฺฉิตพฺโพ. เอวเมโส ปโญฺห สุวินิจฺฉิโต โหตีติ.}

\prose{โหตุ ภนฺเต นาคเสน, ตถา ตํ สมฺปฏิจฺฉามิ. ฐปิโต โหตุ อรหา ตสฺมิํ วตฺถุสฺมิํ, ตสนฺตุ อวเสสา สตฺตา, นิรเย ปน เนรยิกา สตฺตา ทุกฺขา ติพฺพา กฏุกา เวทนา เวทยมานา ชลิต\-ปชฺชลิต\-สพฺพงฺคปจฺจงฺคา รุณฺณ\-การุญฺญ\-กนฺทิต\-ปริเทวิต\-ลาลปฺปิ\-ต\-มุขา อสยฺห\-ติพฺพ\-ทุกฺขาภิ\-ภูตา อตาณา อสรณา อสรณีภูตา อนปฺปโสกาตุรา อนฺติม\-ปจฺฉิม\-คติกา เอกนฺต\-โสก\-ปรายณา อุณฺห\-ติขิณ\-จณฺฑ\-ขร\-ตปน\-เตชวนฺโต ภีมภยชนกนินาทมหาสทฺทา สํสิพฺพิต\-ฉพฺพิธ\-ชาลา\-มาลา\-กุลา สมนฺตา สตโยชนา\-นุ\-ผรณ\-จฺจิ\-เวคา กทริยา ตปนา มหานิรยา จวมานา มจฺจุโน ภายนฺตีติ. อาม มหาราชาติ.}

\prose{นนุ ภนฺเต นาคเสน นิรโย เอกนฺต\-ทุกฺขเวทนีโย, กิสฺส ปน เต เนรยิกา สตฺตา เอกนฺต\-ทุกฺขเวทนียา นิรยา จวมานา มจฺจุโน ภายนฺติ, กิสฺส นิรเย รมนฺตีติ. น เต มหาราช เนรยิกา สตฺตา นิรเย รมนฺติ, มุญฺจิตุกามาว เต นิรยา. มรณสฺเสว โส\csromansharedfootnote{421b4652dd7d5ad4}{มรณสฺเสโส (สี, อิ)} มหาราช อานุภาโว, เยน เตสํ สนฺตาโส อุปฺปชฺชตีติ. เอตํ โข ภนฺเต นาคเสน น สทฺทหามิ, ยํ มุจฺจิตุกามานํ จุติยา สนฺตาโส อุปฺปชฺชตีติ, \csromanfolio{151}หาสนียํ ภนฺเต นาคเสน ตํ ฐานํ, ยํ เต ปตฺถิตํ ลภนฺติ, การเณน มํ สญฺญาเปหีติ.}

\prose{มรณนฺติ โข มหาราช เอตํ อทิฏฺฐสจฺจานํ ตาส\-นีย\-ฏฺฐานํ, เอตฺถายํ ชโน ตสติ จ อุพฺพิชฺชติ จ. โย จ มหาราช กณฺหสปฺปสฺส ภายติ, โส มรณสฺส ภายนฺโต กณฺหสปฺปสฺส ภายติ. โย จ หตฺถิสฺส ภายติ ฯเปฯ. สีหสฺส. พฺยคฺฆสฺส. ทีปิสฺส. อจฺฉสฺส. ตรจฺฉสฺส. มหิํสสฺส. ควยสฺส. อคฺคิสฺส. อุทกสฺส. ขาณุกสฺส. กณฺฏกสฺส ภายติ. โย จ สตฺติยา ภายติ, โส มรณสฺส ภายนฺโต สตฺติยา ภายติ. มรณสฺเสว โส\csromansharedfootnote{421b4652dd7d5ad4}{มรณสฺเสโส (สี, อิ)} มหาราช สรส\-สภาว\-เตโช\csromansharedfootnote{6f4a2af43a7d1b3d}{สรสภาวเตโช (สี, อิ)}, ตสฺส สรส\-สภาว\-เตเชน สกิเลสา สตฺตา มรณสฺส ตสนฺติ ภายนฺติ, มุจฺจิตุกามาปิ มหาราช เนรยิกา สตฺตา มรณสฺส ตสนฺติ ภายนฺติ.}

\prose{อิธ มหาราช ปุริสสฺส กาเย เมโท คณฺฐิ อุปฺปชฺเชยฺย, โส เตน โรเคน ทุกฺขิโต อุปทฺทวา ปริมุจฺจิตุกาโม ภิสกฺกํ สลฺลกตฺตํ อามนฺตาเปยฺย, ตสฺส วจนํ โส ภิสกฺโก สลฺลกตฺโต สมฺปฏิจฺฉิตฺวา ตสฺส โรคสฺส อุทฺธรณาย อุปกรณํ อุปฏฺฐาเปยฺย, สตฺถกํ ติขิณํ กเรยฺย, ยมกสลากา\csromansharedfootnote{cf4fb60eab688a2e}{ทหนสลากํ (ก)} อคฺคิมฺหิ ปกฺขิเปยฺย, ขารลวณํ นิสทาย ปิสาเปยฺย, อปิ นุ โข มหาราช ตสฺส อาตุรสฺส ติขิณ\-สตฺถก\-จฺเฉทเนน ยมก\-สลากา\-ทหเนน ขาร\-โลณ\-ปฺปเวสเนน ตาโส อุปฺปชฺเชยฺยาติ. อาม ภนฺเตติ. อิติ มหาราช ตสฺส อาตุรสฺส โรคา มุจฺจิตุกามสฺสาปิ เวทนาภยา สนฺตาโส อุปฺปชฺชติ. เอวเมว โข มหาราช นิรยา มุจฺจิตุ\-กามานมฺปิ เนรยิกานํ สตฺตานํ มรณภยา สนฺตาโส อุปฺปชฺชติ.}

\prose{อิธ มหาราช ปุริโส อิสฺสรา\-ปราธิโก พทฺโธ สงฺขลิก\-พนฺธเนน คพฺเภ ปกฺขิตฺโต ปริมุจฺจิตุกาโม อสฺส, ตเมนํ โส อิสฺสโร เมเจตุกาโม ปกฺโกสาเปยฺย. อปิ นุ โข มหาราช ตสฺส อิสฺสรา\-ปราธิกสฺส ปุริสสฺส “กตโทโส อหนฺ”ติ ชานนฺตสฺส อิสฺสร\-ทสฺสเนน สนฺตาโส อุปฺปชฺเชยฺยาติ. อาม ภนฺเตติ. อิติ มหาราช ตสฺส อิสฺสรา\-ปราธิกสฺส ปุริสสฺส ปริมุจฺจิตุกามสฺสาปิ อิสฺสรภยา สนฺตาโส อุปฺปชฺชติ. เอวเมว โข มหาราช นิรยา มุจฺจิตุ\-กามานมฺปิ เนรยิกานํ สตฺตานํ มรณภยา สนฺตาโส อุปฺปชฺชตีติ.}

\prosewithcloserruled{\csromanfolio{152}อปรมฺปิ ภนฺเต อุตฺตริํ การณํ พฺรูหิ, เยนาหํ การเณน โอกปฺเปยฺยนฺติ. อิธ มหาราช ปุริโส ทฏฺฐวิเสน อาสีวิเสน ทฏฺโฐ ภเวยฺย, โส เตน วิสวิกาเรน ปเตยฺย อุปฺปเตยฺย วฏฺเฏยฺย ปวฏฺเฏยฺย, อถญฺญตโร ปุริโส พลวนฺเตน มนฺตปเทน ตํ ทฏฺฐวิสํ อาสีวิสํ อาเนตฺวา ตํ ทฏฺฐวิสํ ปจฺจาจมาเปยฺย, อปิ นุ โข มหาราช ตสฺส วิสคตสฺส ปุริสสฺส ตสฺมิํ ทฏฺฐวิเส สปฺเป โสตฺถิเหตุ อุปคจฺฉนฺเต สนฺตาโส อุปฺปชฺเชยฺยาติ. อาม ภนฺเตติ. อิติ มหาราช ตถารูเป อหิมฺหิ โสตฺถิเหตุปิ อุปคจฺฉนฺเต ตสฺส สนฺตาโส อุปฺปชฺชติ. เอวเมว โข มหาราช นิรยา มุจฺจิตุ\-กามานมฺปิ เนรยิกานํ สตฺตานํ มรณภยา สนฺตาโส อุปฺปชฺชติ. อนิฏฺฐํ มหาราช สพฺพสตฺตานํ มรณํ, ตสฺมา เนรยิกา สตฺตา นิรยา ปริมุจฺจิตุกามาปิ มจฺจุโน ภายนฺตีติ. สาธุ ภนฺเต นาคเสน, เอวเมตํ ตถา สมฺปฏิจฺฉามีติ.}{มจฺจุภา\-ยนา\-ภายน\-ปโญฺห ตติโย.}
\csromanmatikaanchor{mtk.121}
\csromantocmarkref{subsubsection}{4. มจฺจุปาสมุตฺติปญฺห}{mtk.121}
\bookmark[dest={mtk.121},level=3]{4. มจฺจุปาสมุตฺติปญฺห}
\csromanheader{3}{4. \textbf{มจฺจุปาส}\-\textbf{มุตฺติ}\-\textbf{ปญฺห}}
\prose{4. ภนฺเต นาคเสน ภาสิตมฺเปตํ ภควตา\csromandash{}}

\setlength{\csromangathapagewidth}{0pt}
\setlength{\csromangathawakwidth}{0pt}
\csromangathameasuregroup{}{\textsuperscript{*}“น อนฺตลิกฺเข น สมุทฺทมชฺเฌ, \\ น ปพฺพตานํ วิวรํ ปวิสฺส. \\ น วิชฺชตี โส ชคติปฺปเทโส, \\ ยตฺถฏฺฐิโต มุจฺเจยฺย มจฺจุปาสา”ติ.}
\csromangathameasurewak{\textsuperscript{*}“น อนฺตลิกฺเข น สมุทฺทมชฺเฌ, \\ น ปพฺพตานํ วิวรํ ปวิสฺส. \\ น วิชฺชตี โส ชคติปฺปเทโส, \\ ยตฺถฏฺฐิโต มุจฺเจยฺย มจฺจุปาสา”ติ.}
\setlength{\csromangathaleftcol}{0pt}
\csromangathagroup{\csromangathastanza{\csromangathawak{\csromansymbolfootnote{*}{ขุ 1. 32 ปิฏฺเฐ ธมฺมปเท.}“น อนฺตลิกฺเข น สมุทฺทมชฺเฌ, \\ น ปพฺพตานํ วิวรํ ปวิสฺส. \\ น วิชฺชตี โส ชคติปฺปเทโส, \\ ยตฺถฏฺฐิโต มุจฺเจยฺย มจฺจุปาสา”ติ.}}}

\prose{ปุน ภควตา ปริตฺตา จ อุทฺทิฏฺฐา. เสยฺยถีทํ, รตนสุตฺตํ เมตฺตสุตฺตํ ขนฺธ\-ปริตฺตํ โมรปริตฺตํ ธชคฺค\-ปริตฺตํ อาฏานาฏิย\-ปริตฺตํ องฺคุลิมาล\-ปริตฺตํ. ยทิ ภนฺเต นาคเสน อากาสคโตปิ สมุทฺทมชฺฌคโตปิ ปาสาทกุฏิเลณคุหาปพฺภารทริพิลคิริ วิวรปพฺพตนฺตรคโตปิ น มุจฺจติ มจฺจุปาสา, เตน หิ ปริตฺตกมฺมํ มิจฺฉา. ยทิ ปริตฺต\-กรเณน มจฺจุปาสา ปริมุตฺติ ภวติ, เตน หิ “น อนฺต ลิกฺเข ฯเปฯ มจฺจุปาสา”ติ ตมฺปิ วจนํ มิจฺฉา. อยมฺปิ อุภโต โกฏิโก ปโญฺห คณฺฐิโตปิ คณฺฐิตโร ตวานุปฺปตฺโต, โส ตยา นิพฺพาหิตพฺโพติ.}

\prose{\csromanfolio{153}ภาสิตมฺเปตํ มหาราช ภควตา “น อนฺตลิกฺเข ฯเปฯ มจฺจุปาสา”ติ, ปริตฺตา จ ภควตา อุทฺทิฏฺฐา, ตญฺจ ปน สาวเสสา\-ยุกสฺส วยสมฺปนฺนสฺส อเปต\-กมฺมาวรณสฺส, นตฺถิ มหาราช ขีณายุกสฺส ฐิติยา กิริยา วา อุปกฺกโม วา.}

\prose{ยถา มหาราช มตสฺส รุกฺขสฺส สุกฺขสฺส โกฬาปสฺส นิสฺเนหสฺส อุปรุทฺธ\-ชีวิตสฺส คตา\-ยุสงฺขารสฺส กุมฺภ\-สหสฺเสนปิ อุทเก อากิรนฺเต อลฺลตฺตํ วา ปลฺลวิ\-ต\-หริต\-ภาโว วา น ภเวยฺย. เอวเมว โข มหาราช เภสชฺช\-ปริตฺต\-กมฺเมน นตฺถิ ขีณายุกสฺส ฐิติยา กิริยา วา อุปกฺกโม วา, ยานิ ตานิ มหาราช มหิยา โอสธานิ เภสชฺชานิ, ตานิปิ ขีณายุกสฺส อกิจฺจกรานิ ภวนฺติ. สาวเสสายุกํ มหาราช วยสมฺปนฺนํ อเปต\-กมฺมาวรณํ ปริตฺตํ รกฺขติ โคเปติ, ตสฺสตฺถาย ภควตา ปริตฺตา อุทฺทิฏฺฐา.}

\prose{ยถา มหาราช กสฺสโก ปริปกฺเก ธญฺเญ มเต สสฺสนาเฬ อุทก\-ปฺปเวสนํ วาเรยฺย, ยํ ปน สสฺสํ ตรุณํ เมฆสนฺนิภํ วยสมฺปนฺนํ, ตํ อุทกวฑฺฒิยา วฑฺฒติ. เอวเมว โข มหาราช ขีณายุกสฺส เภสชฺช\-ปริตฺต\-กิริยา ฐปิตา ปฏิกฺขิตฺตา, เย ปน เต มนุสฺสา สาวเสสายุกา วยสมฺปนฺนา, เตสํ อตฺถาย ปริตฺต\-เภสชฺชานิ ภณิตานิ, เต ปริตฺต\-เภสชฺเชหิ วฑฺฒนฺตีติ.}

\prose{ยทิ ภนฺเต นาคเสน ขีณายุโก มรติ, สาวเสสายุโก ชีวติ, เตน หิ ปริตฺต\-เภสชฺชานิ นิรตฺถกานิ โหนฺตีติ. ทิฏฺฐปุพฺโพ ปน ตยา มหาราช โกจิ โรโค เภสชฺเชหิ ปฏินิวตฺติโตติ. อาม ภนฺเต อเนกสตานิ ทิฏฺฐานีติ. เตน หิ มหาราช “ปริตฺต\-เภสชฺช\-กิริยา นิรตฺถกา”ติ ยํ วจนํ, ตํ มิจฺฉา ภวตีติ.}

\prose{ทิสฺสนฺติ ภนฺเต นาคเสน เวชฺชานํ อุปกฺกมา เภสชฺช\-ปานา\-นุ\-เลปา, เตน เตสํ อุปกฺกเมน โรโค ปฏินิวตฺตตีติ. ปริตฺตานมฺปิ มหาราช ปวตฺตียมานานํ สทฺโท สุยฺยติ, ชิวฺหา สุกฺขติ, หทยํ พฺยาวฏฺฏติ, กณฺโฐ อาตุรติ. เตน เตสํ ปวตฺเตน สพฺเพ พฺยาธโย วูปสมนฺติ, สพฺพา อีติโย อปคจฺฉนฺตีติ.}

\prose{ทิฏฺฐปุพฺโพ ปน ตยา มหาราช โกจิ อหินา ทฏฺโฐ มนฺตปเทน วิสํ ปาตียมาโน วิสํ จิกฺขสฺสนฺโต อุทฺธมโธ อาจมยมาโนติ. อาม}

\prose{\csromanfolio{154}ภนฺเต อชฺเชตรหิปิ ตํ โลเก วตฺตตีติ. เตน หิ มหาราช “ปริตฺต\-เภสชฺช\-กิริยา นิรตฺถกา”ติ ยํ วจนํ, ตํ มิจฺฉา ภวติ. กตปริตฺตญฺหิ มหาราช ปุริสํ ฑํสิตุกาโม อหิ น ฑํสติ, วิวฏํ มุขํ ปิทหติ, โจรานํ อุกฺขิตฺตลคุฬมฺปิ น สมฺภวติ, เต ลคุฬํ มุญฺจิตฺวา เปมํ กโรนฺติ, กุปิโตปิ หตฺถินาโค สมาคนฺตฺวา อุปรมติ, ปชฺชลิตมหาอคฺคิกฺขนฺโธปิ อุปคนฺตฺวา นิพฺพายติ, วิสํ หลาหลมฺปิ ขายิตํ อคทํ สมฺปชฺชติ, อาหารตฺถํ วา ผรติ, วธกา หนฺตุกามา อุปคนฺตฺวา ทาสภูตา สมฺปชฺชนฺติ, อกฺกนฺโตปิ ปาโส น สํวรติ\csromansharedfootnote{f8dcd834c8dc3d4e}{น สํจรติ (สี)}.}

\prose{สุตปุพฺพํ ปน ตยา มหาราช “โมรสฺส กตปริตฺตสฺส สตฺต\-วสฺสสตานิ ลุทฺทโก นาสกฺขิ ปาสํ อุปเนตุํ, อกตปริตฺตสฺส ตํ เยว\csromansharedfootnote{cf7dc7e6d0b8f44d}{ตํเยว นิคฺคหีตสนฺธิ \csromandash{} ม.พ.ป.} ทิวสํ ปาสํ อุปเนสี”ติ. อาม ภนฺเต สุยฺยติ, อพฺภุคฺคโต โส สทฺโท สเทวเก โลเกติ. เตน หิ มหาราช “ปริตฺต\-เภสชฺช\-กิริยา นิรตฺถกา”ติ ยํ วจนํ, ตํ มิจฺฉา ภวติ.}

\prose{สุตปุพฺพํ ปน ตยา มหาราช “ทานโว ภริยํ ปริรกฺขนฺโต สมุคฺเค ปกฺขิปิตฺวา คิลิตฺวา กุจฺฉินา ปริหรติ, อเถโก วิชฺชาธโร ตสฺส ทานวสฺส มุเขน ปวิสิตฺวา ตาย สทฺธิํ อภิรมติ, ยทา โส ทานโว อญฺญาสิ, อถ สมุคฺคํ วมิตฺวา วิวริ, สห สมุคฺเค วิวเฏ วิชฺชาธโร ยถากามํ\csromansharedfootnote{3294a5a611729648}{เยน กามํ (ก)} ปกฺกามี”ติ. อาม ภนฺเต สุยฺยติ, อพฺภุคฺคโต โสปิ สทฺโท สเทวเก โลเกติ. นนุ โส มหาราช วิชฺชาธโร ปริตฺตพเลน\csromansharedfootnote{e7be8b32f4759898}{มนฺตพเลน (?)} คหณา มุตฺโตติ. อาม ภนฺเตติ. เตน หิ มหาราช อตฺถิ ปริตฺตพลํ.}

\prose{สุตปุพฺพํ ปน ตยา มหาราช “อปโรปิ วิชฺชาธโร พาราณสิรญฺโญ อนฺเตปุเร มเหสิยา สทฺธิํ สมฺปทุฏฺโฐ\csromansharedfootnote{f8f27140b66515a9}{สํสฏฺโฐ (สี)} คหณปฺปตฺโต สมาโน ขเณน อทสฺสนํ คโต มนฺตพเลนา”ติ. อาม ภนฺเต สุยฺยตีติ. นนุ โส มหาราช วิชฺชาธโร ปริตฺตพเลน คหณา มุตฺโตติ. อาม ภนฺเตติ. เตน หิ มหาราช อตฺถิ ปริตฺตพลนฺติ.}

\prose{ภนฺเต นาคเสน กิํ สพฺเพ เยว ปริตฺตํ รกฺขตีติ. เอกจฺเจ มหาราช รกฺขติ, เอกจฺเจ น รกฺขตีติ. เตน หิ ภนฺเต นาคเสน ปริตฺตํ น สพฺพตฺถิกนฺติ. อปิ นุ โข มหาราช โภชนํ สพฺเพสํ ชีวิตํ รกฺขตีติ. เอกจฺเจ \csromanfolio{155}ภนฺเต รกฺขติ, เอกจฺเจ น รกฺขตีติ. กิํ การณาติ. ยโต ภนฺเต เอกจฺเจ ตํ เยว\csromansharedfootnote{cf7dc7e6d0b8f44d}{ตํเยว นิคฺคหีตสนฺธิ \csromandash{} ม.พ.ป.} โภชนํ อติภุญฺชิตฺวา วิสูจิกาย มรนฺตีติ. เตน หิ มหาราช โภชนํ น สพฺเพสํ ชีวิตํ รกฺขตีติ. ทฺวีหิ ภนฺเต นาคเสน การเณหิ โภชนํ ชีวิตํ หรติ อติภุตฺเตน วา อุสฺมา\-ทุพฺพลตาย วา, อายุททํ ภนฺเต นาคเสน โภชนํ ทุรุปจาเรน ชีวิตํ หรตีติ. เอวเมว โข มหาราช ปริตฺตํ เอกจฺเจ รกฺขติ, เอกจฺเจ น รกฺขติ.}

\proseitemwithcloserruled{4}{ตีหิ มหาราช การเณหิ ปริตฺตํ น รกฺขติ กมฺมาวรเณน, กิเลสาวรเณน, อสทฺทหนตาย. สตฺตา\-นุรกฺขณํ มหาราช ปริตฺตํ อตฺตนา กเตน อารกฺขํ ชหติ, ยถา มหาราช มาตา ปุตฺตํ กุจฺฉิคตํ โปเสติ, หิเตน อุปจาเรน ชเนติ, ชนยิตฺวา อสุจิ\-มล\-สิงฺฆาณิกม\-ปเนตฺวา อุตฺตม\-วร\-สุคนฺธํ อุปลิมฺปติ, โส อปเรน สมเยน ปเรสํ ปุตฺเต อกฺโกสนฺเต วา ปหรนฺเต วา ปหารํ เทติ. เต ตสฺส กุชฺฌิตฺวา ปริสาย อากฑฺฒิตฺวา ตํ คเหตฺวา สามิโน อุปเนนฺติ, ยทิ ปน ตสฺสา ปุตฺโต อปรทฺโธ โหติ เวลาติวตฺโต. อถ นํ สามิโน มนุสฺสา อากฑฺฒยมานา ทณฺฑ\-มุคฺคร\-ชาณุ\-มุฏฺฐีหิ ตาเฬนฺติ โปเถนฺติ, อปิ นุ โข มหาราช ตสฺส มาตา ลภติ อากฑฺฒน\-ปริกฑฺฒนํ คาหํ สามิโน อุปนยนํ กาตุนฺติ. น หิ ภนฺเตติ. เกน การเณน มหาราชาติ. อตฺตโน ภนฺเต อปราเธนาติ. เอวเมว โข มหาราช สตฺตานํ อารกฺขํ ปริตฺตํ อตฺตโน อปราเธน วญฺจฺยํ กโรตีติ\csromansharedfootnote{cff1ea245306d4f4}{กาเรตีติ (สี)}. สาธุ ภนฺเต นาคเสน สุวินิจฺฉิโต ปโญฺห, คหนํ อคหนํ กตํ, อนฺธกาโร อาโลโก กโต, วินิเวฐิตํ ทิฏฺฐิชาลํ, ตฺวํ คณิวร\-ปวรมา\-สชฺชาติ.}{มจฺจุปาส\-มุตฺติ\-ปโญฺห จตุตฺโถ.}
\csromanmatikaanchor{mtk.122}
\csromantocmarkref{subsubsection}{5. พุทฺธลาภนฺตรายปญฺห}{mtk.122}
\bookmark[dest={mtk.122},level=3]{5. พุทฺธลาภนฺตรายปญฺห}
\csromanheader{3}{5. \textbf{พุทฺธ}\-\textbf{ลาภ}\-\textbf{นฺตราย}\-\textbf{ปญฺห}}
\proseitem{4}{ภนฺเต นาคเสน ตุเมฺห ภณถ\csromansymbolfootnote{*}{ขุ 1. 89, 128 ปิฏฺเฐสุ.}“ลาภี ตถาคโต จีวร\-ปิณฺฑปาต\-เสนาสน\-คิลานปฺปจฺจย\-เภสชฺช\-ปริกฺขารานนฺ”ติ. ปุน จ ตถาคโต\csromansymbolfootnote{+}{สํ 1. 115 ปิฏฺเฐ มารสํยุตฺเต.}ปญฺจสาลํ พฺราหฺมณคามํ ปิณฺฑาย ปวิสิตฺวา กิญฺจิเทว อลภิตฺวา \csromanfolio{156}ยถาโธเตน ปตฺเตน นิกฺขนฺโตติ. ยทิ ภนฺเต นาคเสน ตถาคโต ลาภี จีวร\-ปิณฺฑปาต\-เสนาสน\-คิลานปฺปจฺจย\-เภสชฺช\-ปริกฺขารานํ, เตน หิ ปญฺจสาลํ พฺราหฺมณคามํ ปิณฺฑาย ปวิสิตฺวา กิญฺจิเทว อลภิตฺวา ยถาโธเตน ปตฺเตน นิกฺขนฺโตติ ยํ วจนํ, ตํ มิจฺฉา. ยทิ ปญฺจสาลํ พฺราหฺมณคามํ ปิณฺฑาย ปวิสิตฺวา กิญฺจิเทว อลภิตฺวา ยถาโธเตน ปตฺเตน นิกฺขนฺโต, เตน หิ ลาภี ตถาคโต จีวร\-ปิณฺฑปาต\-เสนาสน\-คิลานปฺปจฺจย\-เภสชฺช\-ปริกฺขารานนฺติ ตมฺปิ วจนํ มิจฺฉา. อยมฺปิ อุภโต โกฏิโก ปโญฺห สุมหนฺโต ทุนฺนิพฺเพโฐ ตวานุปฺปตฺโต, โส ตยา นิพฺพาหิตพฺโพติ.}

\proseitem{5}{ลาภี มหาราช ตถาคโต จีวร\-ปิณฺฑปาต\-เสนาสนคิลาน- ปฺปจฺจยเภสชฺชปริกฺขารานํ, ปญฺจสาลญฺจ พฺราหฺมณคามํ ปิณฺฑาย ปวิสิตฺวา กิญฺจิเทว อลภิตฺวา ยถาโธเตน ปตฺเตน นิกฺขนฺโต, ตญฺจ ปน มารสฺส ปาปิมโต การณาติ. เตน หิ ภนฺเต นาคเสน ภควโต คณนปถํ วีติวตฺตกปฺเป\csromansharedfootnote{f68012fa2bb1e888}{คณนปถวีติวตฺเต กปฺเป (สี)} อภิสงฺขตํ กุสลํ กินฺติ นิฏฺฐิตํ, อธุนุฏฺฐิเตน มาเรน ปาปิมตา ตสฺส กุสลสฺส พลเวคํ\csromansharedfootnote{6ea8b59c0d8247d1}{ตํ กุสลพลเวควิปฺผารํ (สี)} กินฺติ ปิหิตํ, เตน หิ ภนฺเต นาคเสน ตสฺมิํ วตฺถุสฺมิํ ทฺวีสุ ฐาเนสุ อุปวาโท อาคจฺฉติ, กุสลโตปิ อกุสลํ พลวตรํ โหติ, พุทฺธพลโตปิ มารพลํ พลวตรํ โหตีติ, เตน หิ รุกฺขสฺส มูลโตปิ อคฺคํ ภารตรํ โหติ, คุณสมฺปริกิณฺณโตปิ ปาปิยํ พลวตรํ โหตีติ. น มหาราช ตาวตเกน กุสลโตปิ อกุสลํ พลวตรํ นาม โหติ, น พุทฺธพลโตปิ มารพลํ พลวตรํ นาม โหติ. อปิ เจตฺถ การณํ อิจฺฉิตพฺพํ.}

\prose{ยถา มหาราช ปุริโส รญฺโญ จกฺกวตฺติสฺส มธุํ วา มธุปิณฺฑิกํ วา อญฺญํ วา อุปายนํ อภิหเรยฺย, ตเมนํ รญฺโญ ทฺวารปาโล เอวํ วเทยฺย “อกาโล โภ อยํ รญฺโญ ทสฺสนาย, เตน หิ โภ ตว อุปายนํ คเหตฺวา สีฆสีฆํ ปฏินิวตฺต, ปุเร ตว ราชา ทณฺฑํ ธาเรสฺสตี”ติ\csromansharedfootnote{34e37eec77334f64}{มา เต ราชา ทณฺฑํ ปาเปยฺยาติ (สี)}. ตโต โส ปุริโส ทณฺฑภยา ตสิโต อุพฺพิคฺโค ตํ อุปายนํ อาทาย สีฆสีฆํ ปฏินิวตฺเตยฺย, อปิ นุ โข โส มหาราช ราชา จกฺกวตฺตี ตาวตเกน อุปายน\-วิกล\-มตฺตเกน ทฺวารปาลโต ทุพฺพลตโร นาม \csromanfolio{157}โหติ, อญฺญํ วา ปน กิญฺจิ อุปายนํ น ลเภยฺยาติ. น หิ ภนฺเต, อิสฺสาปกโต โส ภนฺเต ทฺวารปาโล อุปายนํ นิวาเรสิ, อญฺเญน ปน ทฺวาเรน สตสหสฺส\-คุณมฺปิ รญฺโญ อุปายนํ อุเปตีติ. เอวเมว โข มหาราช อิสฺสาปกโต มาโร ปาปิมา ปญฺจสาลเก พฺราหฺมณคหปติเก อนฺวาวิสิ, อญฺญานิ ปน อเนกานิ เทวตา\-สตสหสฺสานิ อมตํ ทิพฺพํ โอชํ คเหตฺวา อุปคตานิ “ภควโต กาเย โอชํ โอทหิสฺสามี”ติ ภควนฺตํ นมสฺสมานานิ ปญฺชลิกานิ ฐิตานีติ.}

\prose{โหตุ ภนฺเต นาคเสน สุลภา ภควโต จตฺตาโร ปจฺจยา โลเก อุตฺตม\-ปุริสสฺส, ยาจิโตว ภควา เทวมนุสฺเสหิ จตฺตาโร ปจฺจเย ปริภุญฺชติ, อปิ จ โข ปน มารสฺส โย อธิปฺปาโย, โส ตาวตเกน สิทฺโธ, ยํ โส ภควโต โภชนสฺส อนฺตรายมกาสิ. เอตฺถ เม ภนฺเต กงฺขา น ฉิชฺชติ, วิมติชาโตหํ ตตฺถ สํสย\-ปกฺขนฺโท. น เม ตตฺถ มานสํ ปกฺขนฺทติ, ยํ ตถาคตสฺส อรหโต สมฺมา\-สมฺพุทฺธสฺส สเทวเก โลเก อคฺคปุคฺคล\-วรสฺส กุสล\-วร\-ปุญฺญ\-สมฺภวสฺส อสมสมสฺส อนุปมสฺส อปฺปฏิสมสฺส ฉวกํ ลามกํ ปริตฺตํ ปาปํ อนริยํ วิปนฺนํ มาโร ลาภ\-นฺตรายมกาสีติ.}

\prose{จตฺตาโร โข มหาราช อนฺตรายา อทิฏฺฐนฺตราโย อุทฺทิสฺส กตนฺตราโย อุปกฺขฏ\-นฺตราโย ปริโภคนฺตราโยติ. ตตฺถ กตโม อทิฏฺฐนฺตราโย, อโนทิสฺส อทสฺสเนน อนภิสงฺขตํ โกจิ อนฺตรายํ กโรติ “กิํ ปรสฺส ทินฺเนนา”ติ, อยํ อทิฏฺฐนฺตราโย นาม.}

\prose{กตโม อุทฺทิสฺส กตนฺตราโย, อิเธกจฺจํ ปุคฺคลํ อุปทิสิตฺวา อุทฺทิสฺส โภชนํ ปฏิยตฺตํ โหติ, ตํ โกจิ อนฺตรายํ กโรติ, อยํ อุทฺทิสฺส กตนฺตราโย นาม.}

\prose{กตโม อุปกฺขฏ\-นฺตราโย, อิธ ยํ กิญฺจิ อุปกฺขฏํ โหติ อปฺปฏิคฺคหิตํ, ตตฺถ โกจิ อนฺตรายํ กโรติ, อยํ อุปกฺขฏ\-นฺตราโย นาม.}

\prose{กตโม ปริโภค\-นฺตราโย, อิธ ยํ กิญฺจิ ปริโภคํ, ตตฺถ โกจิ อนฺตรายํ กโรติ, อยํ ปริโภค\-นฺตราโย นาม. อิเม โข มหาราช จตฺตาโร อนฺตรายา.}

\prose{ยํ ปน มาโร ปาปิมา ปญฺจสาลเก พฺราหฺมณคหปติเก อนฺวาวิสิ, ตํ เนว ภควโต ปริโภคํ น อุปกฺขฏํ น อุทฺทิสฺสกตํ, อนาคตํ อสมฺปตฺตํ \csromanfolio{158}อทสฺสเนน อนฺตรายํ กตํ, ตํ ปน เนกสฺส ภควโต เยว, อถ โข เย เต เตน สมเยน นิกฺขนฺตา อพฺภาคตา, สพฺเพปิ เต ตํ ทิวสํ โภชนํ น ลภิํสุ, นาหํ ตํ มหาราช ปสฺสามิ สเทวเก โลเก สมารเก สพฺรหฺมเก สสฺสมณ\-พฺราหฺมณิยา ปชาย สเทว\-มนุสฺสาย, โย ตสฺส ภควโต อุทฺทิสฺส กตํ อุปกฺขฏํ ปริโภคํ อนฺตรายํ กเรยฺย. สเจ โกจิ อิสฺสาย อุทฺทิสฺส กตํ อุปกฺขฏํ ปริโภคํ อนฺตรายํ กเรยฺย, ผเลยฺย ตสฺส มุทฺธา สตธา วา สหสฺสธา วา.}

\proseitem{4}{จตฺตาโรเม มหาราช ตถาคตสฺส เกนจิ อนาวรณียา คุณา. กตเม จตฺตาโร, ลาโภ มหาราช ภควโต อุทฺทิสฺส กโต อุปกฺขโฏ น สกฺกา เกนจิ อนฺตรายํ กาตุํ, สรีรานุคตา มหาราช ภควโต พฺยามปฺปภา น สกฺกา เกนจิ อนฺตรายํ กาตุํ, สพฺพญฺญุตํ มหาราช ภควโต ญาณรตนํ น สกฺกา เกนจิ อนฺตรายํ กาตุํ, ชีวิตํ มหาราช ภควโต น สกฺกา เกนจิ อนฺตรายํ กาตุํ, อิเม โข มหาราช จตฺตาโร ตถาคตสฺส เกนจิ อนาวรณียา คุณา, สพฺเพเปเต มหาราช คุณา เอกรสา อโรคา อกุปฺปา อปรูปกฺกมา อผุสานิ กิริยานิ. อทสฺสเนน มหาราช มาโร ปาปิมา นิลียิตฺวา ปญฺจสาลเก พฺราหฺมณคหปติเก อนฺวาวิสิ.}

\prose{ยถา มหาราช รญฺโญ ปจฺจนฺเต เทเส วิสเม อทสฺสเนน นิลียิตฺวา โจรา ปนฺถํ ทูเสนฺติ. ยทิ ปน ราชา เต โจเร ปสฺเสยฺย, อปิ นุ โข เต โจรา โสตฺถิํ ลเภยฺยุนฺติ. น หิ ภนฺเต, ผรสุนา ผาลาเปยฺย สตธา วา สหสฺสธา วาติ. เอวเมว โข มหาราช อทสฺสเนน มาโร ปาปิมา นิลียิตฺวา ปญฺจสาลเก พฺราหฺมณคหปติเก อนฺวาวิสิ.}

\csromanmatikaanchor{mtk.123}
\csromantocmarkref{subsubsection}{6. อปญฺญปญฺห}{mtk.123}
\bookmark[dest={mtk.123},level=3]{6. อปญฺญปญฺห}
\prosewithcloserruled{ยถา วา ปน มหาราช อิตฺถี สปติกา อทสฺสเนน นิลียิตฺวา ปรปุริสํ เสวติ, เอวเมว โข มหาราช อทสฺสเนน มาโร ปาปิมา นิลียิตฺวา ปญฺจสาลเก พฺราหฺมณคหปติเก อนฺวาวิสิ. ยทิ มหาราช อิตฺถี สามิกสฺส สมฺมุขา ปรปุริสํ เสวติ, อปิ นุ โข สา อิตฺถี โสตฺถิํ ลเภยฺยาติ. น หิ ภนฺเต, หเนยฺยาปิ ตํ ภนฺเต สามิโก วเธยฺยาปิ พนฺเธยฺยาปิ ทาสิตฺตํ วา อุปเนยฺยาติ. เอวเมว โข มหาราช อทสฺสเนน มาโร ปาปิมา นิลียิตฺวา ปญฺจสาลเก พฺราหฺมณคหปติเก อนฺวาวิสิ. ยทิ มหาราช มาโร ปาปิมา ภควโต อุทฺทิสฺส กตํ อุปกฺขฏํ ปริโภคํ \csromanfolio{159}อนฺตรายํ กเรยฺย, ผเลยฺย ตสฺส มุทฺธา สตธา วา สหสฺสธา วาติ. เอวเมตํ ภนฺเต นาคเสน โจริกาย กตํ มาเรน ปาปิมตา, นิลียิตฺวา มาโร ปาปิมา ปญฺจสาลเก พฺราหฺมณคหปติเก อนฺวาวิสิ. สเจ โส ภนฺเต มาโร ปาปิมา ภควโต อุทฺทิสฺส กตํ อุปกฺขฏํ ปริโภคํ อนฺตรายํ กเรยฺย, มุทฺธา วาสฺส ผเลยฺย สตธา วา สหสฺสธา วา, กาโย วาสฺส ภุสมุฏฺฐิ วิย วิกิเรยฺย, สาธุ ภนฺเต นาคเสน, เอวเมตํ ตถา สมฺปฏิจฺฉามีติ.}{พุทฺธ\-ลาภ\-นฺตราย\-ปโญฺห ปญฺจโม.}
\csromanheader{2}{6. \textbf{อปุญฺญปญฺห}}
\proseitem{4}{ภนฺเต นาคเสน ตุเมฺห ภณถ\csromansymbolfootnote{+}{เหฏฺฐา 90 ปิฏฺเฐ.}“โย อชานนฺโต ปาณาติปาตํ กโรติ, โส พลวตรํ อปุญฺญํ ปสวตี”ติ. ปุน จ ภควตา วินย\-ปญฺญตฺติยา ภณิตํ\csromansymbolfootnote{*}{วิ 1. 48 ปิฏฺฐาทีสุ.}“อนาปตฺติ อชานนฺตสฺสา”ติ. ยทิ ภนฺเต นาคเสน อชานิตฺวา ปาณาติปาตํ กโรนฺโต พลวตรํ อปุญฺญํ ปสวติ, เตน หิ “อนาปตฺติ อชานนฺตสฺสา”ติ ยํ วจนํ, ตํ มิจฺฉา. ยทิ อนาปตฺติ อชานนฺตสฺส, เตน หิ “อชานิตฺวา ปาณาติปาตํ กโรนฺโต พลวตรํ อปุญฺญํ ปสวตี”ติ ตมฺปิ วจนํ มิจฺฉา, อยมฺปิ อุภโต โกฏิโก ปโญฺห ทุรุตฺตโร ทุรติกฺกโม ตวานุปฺปตฺโต, โส ตยา นิพฺพาหิตพฺโพติ.}

\prose{ภาสิตมฺเปตํ มหาราช ภควตา “โย อชานนฺโต ปาณาติปาตํ กโรติ, โส พลวตรํ อปุญฺญํ ปสวตี”ติ. ปุน จ วินย\-ปญฺญตฺติยา ภควตา ภณิตํ “อนาปตฺติ อชานนฺตสฺสา”ติ. ตตฺถ อตฺถนฺตรํ อตฺถิ. กตมํ อตฺถนฺตรํ, อตฺถิ มหาราช อาปตฺติ สญฺญาวิโมกฺขา, อตฺถิ อาปตฺติ โนสญฺญา\-วิโมกฺขา, ยายํ มหาราช อาปตฺติ สญฺญาวิโมกฺขา, ตํ อาปตฺติํ อารพฺภ ภควตา ภณิตํ “อนาปตฺติ อชานนฺตสฺสา”ติ. สาธุ ภนฺเต นาคเสน, เอวเมตํ ตถา สมฺปฏิจฺฉามีติ.}

\vspace{\tipitakaparskip}
\csromancenter{อปุญฺญปโญฺห ฉฏฺโฐ.}
\csromanheader{2}{7. \textbf{ภิกฺขุสํฆปริหรณปญฺห}}
\csromanmatikaanchor{mtk.124}
\csromantocmarkref{subsubsection}{7. ภิกฺขุสํฆปริหารณปญฺห}{mtk.124}
\bookmark[dest={mtk.124},level=3]{7. ภิกฺขุสํฆปริหารณปญฺห}
\proseitem{7}{\csromanfolio{160}ภนฺเต นาคเสน ภาสิตมฺเปตํ ภควตา\csromansymbolfootnote{*}{ที 2. 85; สํ 3. 133 ปิฏฺเฐสุ.}“ตถาคตสฺส โข อานนฺท น เอวํ โหติ ‘อหํ ภิกฺขุสํฆํ ปริหริสฺสามี’ติ วา, ‘มมุทฺเทสิโก ภิกฺขุสํโฆ’ติ วา”ติ. ปุน จ เมตฺเตยฺยสฺส ภควโต สภาวคุณํ ปริทีปย\-มาเนน ภควตา เอวํ ภณิตํ\csromansymbolfootnote{+}{ที 3. 63 ปิฏฺเฐ.}“โส อเนกสหสฺสํ ภิกฺขุสํฆํ ปริหริสฺสติ, เสยฺยถาปิ อหํ เอตรหิ อเนกสตํ ภิกฺขุสํฆํ ปริหรามี”ติ. ยทิ ภนฺเต นาคเสน ภควตา ภณิตํ “ตถาคตสฺส โข อานนฺท น เอวํ โหติ ‘อหํ ภิกฺขุสํฆํ ปริหริสฺสามี’ติ วา, ‘มมุทฺเทสิโก ภิกฺขุสํโฆ’ติ วา”ติ, เตน หิ อเนกสตํ ภิกฺขุสํฆํ ปริหรามีติ ยํ วจนํ, ตํ มิจฺฉา. ยทิ ตถาคเตน ภณิตํ “โส อเนกสหสฺสํ ภิกฺขุสํฆํ ปริหริสฺสติ, เสยฺยถาปิ อหํ เอตรหิ อเนกสตํ ภิกฺขุสํฆํ ปริหรามี”ติ, เตน หิ ตถาคตสฺส โข อานนฺท น เอวํ โหติ “อหํ ภิกฺขุสํฆํ ปริหริสฺสามี”ติ วา, “มมุทฺเทสิโก ภิกฺขุสํโฆ”ติ วาติ ตมฺปิ ตมฺปิ วจนํ มิจฺฉา, อยมฺปิ อุภโต โกฏิโก ปโญฺห ตวานุปฺปตฺโต, โส ตยา นิพฺพาหิตพฺโพติ.}

\prose{ภาสิตมฺเปตํ มหาราช ภควตา “ตถาคตสฺส โข อานนฺท น เอวํ โหติ ‘อหํ ภิกฺขุสํฆํ ปริหริสฺสามี’ติ วา, ‘มมุทฺเทสิโก ภิกฺขุ สํโฆ’ติ วา”ติ. ปุน จ เมตฺตยฺยสฺสาปิ ภควโต สภาวคุณํ ปริทีปย\-มาเนน ภควตา ภณิตํ “โส อเนกสหสฺสํ ภิกฺขุสํฆํ ปริหริสฺสติ, เสยฺยถาปิ อหํ เอตรหิ อเนกสตํ ภิกฺขุสํฆํ ปริหรามี”ติ, เอตสฺมิญฺจ มหาราช ปเญฺห เอโก อตฺโถ สาวเสโส, เอโก อตฺโถ นิรวเสโส, น มหาราช ตถาคโต ปริสาย อนุคามิโก, ปริสา ปน ตถาคตสฺส อนุคามิกา, สมฺมุติ มหาราช เอสา “อหนฺ”ติ “มมา”ติ, น ปรมตฺโถ เอโส, วิคตํ มหาราช ตถาคตสฺส เปมํ, วิคโต สิเนโห, “มยฺหนฺ”ติปิ ตถาคตสฺส คหณํ นตฺถิ, อุปาทาย ปน อวสฺสโย โหติ.}

\prose{ยถา มหาราช ปถวี ภูมฏฺฐานํ สตฺตานํ ปติฏฺฐา โหติ อุปสฺสยํ, ปถวิฏฺฐา เจเต สตฺตา, น จ มหาปถวิยา “มเยฺหเต”ติ อเปกฺขา โหติ, เอวเมว โข มหาราช ตถาคโต สพฺพสตฺตานํ ปติฏฺฐา \csromanfolio{161}โหติ อุปสฺสยํ, ตถาคตฏฺฐา\csromansharedfootnote{6dac9fb385261835}{ตถาคตปติฏฺฐา เอว (สี)} เจเต สตฺตา, น จ ตถาคตสฺส “มเยฺหเต”ติ อเปกฺขา โหติ. ยถา วา ปน มหาราช มหติ\-มหาเมโฆ อภิวสฺสนฺโต ติณรุกฺข\-ปสุ\-มนุสฺสานํ วุฑฺฒิํ เทติ สนฺตติํ อนุปาเลติ. วุฏฺฐูปชีวิโน เจเต สตฺตา สพฺเพ, น จ มหาเมฆสฺส “มเยฺหเต”ติ อเปกฺขา โหติ. เอวเมว โข มหาราช ตถาคโต สพฺพสตฺตานํ กุสลธมฺเม ชเนติ อนุปาเลติ, สตฺถูปชีวิโน เจเต สตฺตา สพฺเพ, น จ ตถาคตสฺส “มเยฺหเต”ติ อเปกฺขา โหติ. ตํ กิสฺส เหตุ, อตฺตานุทิฏฺฐิยา ปหีนตฺตาติ.}

\prosewithcloserruled{สาธุ ภนฺเต นาคเสน, สุนิพฺเพฐิโต ปโญฺห พหุวิเธหิ การเณหิ, คมฺภีโร อุตฺตานีกโต, คณฺฐิ ภินฺโน, คหนํ อคหนํ กตํ, อนฺธกาโร อาโลโก กโต, ภคฺคา ปรวาทา, ชินปุตฺตานํ จกฺขุํ อุปฺปาทิตนฺติ.}{ภิกฺขุสํฆปริหรณปโญฺห สตฺตโม.}
\csromanmatikaanchor{mtk.125}
\csromantocmarkref{subsubsection}{8. อเภชฺชปริสปญฺห}{mtk.125}
\bookmark[dest={mtk.125},level=3]{8. อเภชฺชปริสปญฺห}
\csromanheader{3}{8. \textbf{อเภชฺชปริส}\-\textbf{ปญฺห}}
\proseitem{8}{ภนฺเต นาคเสน ตุเมฺห ภณถ\csromansymbolfootnote{*}{ที 3. 140 ปิฏฺเฐ.}“ตถาคโต อเภชฺชปริโส”ติ, ปุน จ ภณถ\csromansymbolfootnote{+}{วิ 4. 362 ปิฏฺเฐ.}“เทวทตฺเตน เอกปฺปหารํ ปญฺจ ภิกฺขุสตานิ ภินฺนานี”ติ, ยทิ ภนฺเต นาคเสน ตถาคโต อเภชฺชปริโส, เตน หิ เทวทตฺเตน เอกปฺปหารํ ปญฺจ ภิกฺขุสตานิ ภินฺนานีติ ยํ วจนํ, ตํ มิจฺฉา. ยทิ เทวทตฺเตน เอกปฺปหารํ ปญฺจ ภิกฺขุสตานิ ภินฺนานิ, เตน หิ ตถาคโต อเภชฺช\-ปริโสติ ตมฺปิ วจนํ มิจฺฉา. อยมฺปิ อุภโต โกฏิโก ปโญฺห ตวานุปฺปตฺโต, คมฺภีโร ทุนฺนิเวฐิโย, คณฺฐิโตปิ คณฺฐิตโร, เอตฺถายํ ชโน อาวโฏ นิวุโต โอวุโต ปิหิโต ปริโยนทฺโธ, เอตฺถ ตว ญาณพลํ ทสฺเสหิ ปรวาเทสูติ.}

\prose{อเภชฺชปริโส มหาราช ตถาคโต, เทวทตฺเตน จ เอกปฺปหารํ ปญฺจ ภิกฺขุสตานิ ภินฺนานิ, ตญฺจ ปน เภทกสฺส พเลน, เภทเก วิชฺชมาเน นตฺถิ มหาราช อเภชฺชํ นาม. เภทเก สติ มาตาปิ ปุตฺเตน ภิชฺชติ, ปุตฺโตปิ มาตรา ภิชฺชติ, ปิตาปิ ปุตฺเตน ภิชฺชติ, ปุตฺโตปิ ปิตรา ภิชฺชติ, ภาตาปิ ภคินิยา ภิชฺชติ, ภคินีปิ ภาตรา ภิชฺชติ, \csromanfolio{162}สหาโยปิ สหาเยน ภิชฺชติ, นาวาปิ นานาทารุ\-สงฺฆฏิตา อูมิเวค\-สมฺปหาเรน ภิชฺชติ, รุกฺโขปิ มธุกปฺป\-สมฺปนฺนผโล อนิล\-พล\-เวคา\-ภิหโต ภิชฺชติ, สุวณฺณมฺปิ ชาติมนฺตํ\csromansharedfootnote{ee59a21bd5de73a0}{ชาตรูปมฺปิ (สี)} โลเหน ภิชฺชติ. อปิ จ มหาราช เนโส อธิปฺปาโย วิญฺญูนํ, เนสา พุทฺธานํ อธิมุตฺติ, เนโส ปณฺฏิตานํ ฉนฺโท “ตถาคโต เภชฺชปริโส”ติ. อปิ เจตฺถ การณํ อตฺถิ, เยน การเณน ตถาคโต วุจฺจติ “อเภชฺชปริโส”ติ. กตมํ เอตฺถ การณํ, ตถาคตสฺส มหาราช กเตน อทาเนน วา อปฺปิยวจเนน วา อนตฺถ\-จริยาย วา อสมานตฺตตาย วา ยโต กุโตจิ จริยํ จรนฺตสฺสปิ ปริสา ภินฺนาติ น สุตปุพฺพํ, เตน การเณน ตถาคโต วุจฺจติ “อเภชฺชปริโส”ติ. ตยาเปตํ มหาราช ญาตพฺพํ “อตฺถิ กิญฺจิ นวงฺเค พุทฺธวจเน สุตฺตาคตํ, อิมินา นาม การเณน โพธิสตฺตสฺส กเตน ตถาคตสฺส ปริสา ภินฺนา”ติ. นตฺถิ ภนฺเต, โน เจตํ โลเก ทิสฺสติ โนปิ สุยฺยติ, สาธุ ภนฺเต นาคเสน, เอวเมตํ ตถา สมฺปฏิจฺฉามีติ.}

\vspace{\tipitakaparskip}
\csromancenter{อเภชฺชปริส\-ปโญฺห อฏฺฐโม.}
\nitthitammidruled{\textbf{อเภชฺชวคฺโค ทุติโย.}}
\csromancenter{อิมสฺมิํ วคฺเค อฏฺฐ ปญฺหา.}
\csromanmatikaanchor{mtk.126}
\csromantocmarknopage{subsection}{3. ปณามิตวคฺค}{mtk.126}
\bookmark[dest={mtk.126},level=2]{3. ปณามิตวคฺค}
\csromanheader{2}{3. \textbf{ปณามิตวคฺค}}
\csromanmatikaanchor{mtk.127}
\csromantocmarkref{subsubsection}{1. เสฏฺฐธมฺมปญฺห}{mtk.127}
\bookmark[dest={mtk.127},level=3]{1. เสฏฺฐธมฺมปญฺห}
\csromanheader{3}{1. \textbf{เสฏฺฐ}\-\textbf{ธมฺม}\-\textbf{ปญฺห}}
\proseitem{1}{\csromanfolio{163}ภนฺเต นาคเสน ภาสิตมฺเปตํ ภควตา\csromansymbolfootnote{*}{ที 3. 68 ปิฏฺเฐ.}“ธมฺโม หิ วาเสฏฺฐ เสฏฺโฐ ชเนตสฺมิํ ทิฏฺเฐ เจว ธมฺเม อภิสมฺปราเย จา”ติ. ปุน จ “อุปาสโก คิหี โสตาปนฺโน ปิหิตาปาโย ทิฏฺฐิปฺปตฺโต วิญฺญาตสาสโน ภิกฺขุํ วา สามเณรํ วา ปุถุชฺชนํ อภิวาเทติ ปจฺจุฏฺเฐตี”ติ. ยทิ ภนฺเต นาคเสน ภควตา ภณิตํ “ธมฺโม หิ วาเสฏฺฐ เสฏฺโฐ ชเนตสฺมิํ ทิฏฺเฐ เจว ธมฺเม อภิสมฺปราเย จา”ติ, เตน หิ “อุปาสโก คิหี โสตาปนฺโน ปิหิตาปาโย ทิฏฺฐิปฺปตฺโต วิญฺญาตสาสโน ภิกฺขุํ วา สามเณรํ วา ปุถุชฺชนํ อภิวาเทติ ปจฺจุฏฺเฐตี”ติ ยํ วจนํ, ตํ มิจฺฉา. ยทิ อุปาสโก คิหี โสตาปนฺโน ปิหิตาปาโย ทิฏฺฐิปฺปตฺโต วิญฺญาตสาสโน ภิกฺขุํ วา สามเณรํ วา ปุถุชฺชนํ อภิวาเทติ ปจฺจุฏฺเฐติ, เตน หิ “ธมฺโม หิ วาเสฏฺฐ เสฏฺโฐ ชเนตสฺมิํ ทิฏฺเฐ เจว ธมฺเม อภิสมฺปราเย จา”ติ ตมฺปิ วจนํ มิจฺฉา. อยมฺปิ อุภโต โกฏิโก ปโญฺห ตวานุปฺปตฺโต, โส ตยา นิพฺพาหิตพฺโพติ.}

\prose{ภาสิตมฺเปตํ มหาราช ภควตา “ธมฺโม หิ วาเสฏฺฐ เสฏฺโฐ ชเนตสฺมิํ ทิฏฺเฐ เจว ธมฺเม อภิสมฺปราเย จา”ติ, อุปาสโก จ คิหี โสตาปนฺโน ปิหิตาปาโย ทิฏฺฐิปฺปตฺโต วิญฺญาตสาสโน ภิกฺขุํ วา สามเณรํ วา ปุถุชฺชนํ อภิวาเทติ ปจฺจุฏฺเฐติ. ตตฺถ ปน การณํ อตฺถิ, กตมํ ตํ การณํ.}

\prose{วีสติ โข ปนิเม มหาราช สมณสฺส สมณกรณา ธมฺมา เทฺว จ ลิงฺคานิ, เยหิ สมโณ อภิวาทน\-ปจฺจุฏฺฐาน\-สมฺมานน\-ปูชนารโห โหติ, กตเม วีสติ สมณสฺส สมณกรณา ธมฺมา เทฺว จ ลิงฺคานิ, เสฏฺโฐ\csromansharedfootnote{f825ffe4680a9839}{เสฏฺฐภูมิสโย (สี, สฺยา), เสฏฺโฐ ยโม (อิ)} ธมฺมาราโม, อคฺโค นิยโม, จาโร วิหาโร สํยโม สํวโร ขนฺติ โสรจฺจํ เอกตฺตจริยา เอกตฺตาภิรติ ปฏิสลฺลานํ หิริโอตฺตปฺปํ วีริยํ อปฺปมาโท สิกฺขา\-สมาทานํ\csromansharedfootnote{a3813f53906400f5}{สิกฺขาปธานํ (สี, สฺยา), สุกฺกาวทานํ (ก)} อุทฺเทโส ปริปุจฺฉา สีลาทิอภิรติ นิราลยตา สิกฺขา\-ปทปาริปูริตา, กาสาวธารณํ, ภณฺฑุภาโว. อิเม \csromanfolio{164}โข มหาราช วีสติ สมณสฺส สมณกรณา ธมฺมา เทฺว จ ลิงฺคานิ. เอเต คุเณ ภิกฺขุ สมาทาย วตฺตติ, โส เตสํ ธมฺมานํ อนูนตฺตา ปริปุณฺณตฺตา สมฺปนฺนตฺตา สมนฺนาคตตฺตา อเสกฺขภูมิํ อรหนฺตภูมิํ โอกฺกมติ, เสฏฺฐํ ภูมนฺตรํ โอกฺกมติ, อรหตฺตาสนฺนคโตติ อรหติ อุปาสโก โสตาปนฺโน ภิกฺขุํ ปุถุชฺชนํ อภิวาเทตุํ ปจฺจุฏฺฐาตุํ.}

\prose{“ขีณาสเวหิ โส สามญฺญํ อุปคโต, นตฺถิ เม โส สมโย”ติ\csromansharedfootnote{6a727279aa027da9}{ตํ สามญฺญนฺ”ติ (?)} อรหติ อุปาสโก โสตาปนฺโน ภิกฺขุํ ปุถุชฺชนํ อภิวาเทตุํ ปจฺจุฏฺฐาตุํ.}

\prose{“อคฺคปริสํ โส อุปคโต, นาหํ ตํ ฐานํ อุปคโต”ติ อรหติ อุปาสโก โสตาปนฺโน ภิกฺขุํ ปุถุชฺชนํ อภิวาเทตุํ ปจฺจุฏฺฐาตุํ.}

\prose{“ลภติ โส ปาติโมกฺขุ\-ทฺเทสํ โสตุํ, นาหํ ตํ ลภามิ โสตุนฺ”ติ อรหติ อุปาสโก โสตาปนฺโน ภิกฺขุํ ปุถุชฺชนํ อภิวาเทตุํ ปจฺจุฏฺฐาตุํ.}

\prose{“โส อญฺเญ ปพฺพาเชติ อุปสมฺปาเทติ ชินสาสนํ วฑฺเฒติ, อหเมตํ น ลภามิ กาตุนฺ”ติ อรหติ อุปาสโก โสตาปนฺโน ภิกฺขุํ ปุถุชฺชนํ อภิวาเทตุํ ปจฺจุฏฺฐาตุํ.}

\prose{“อปฺปมาเณสุ โส สิกฺขาปเทสุ สมตฺตการี, นาหํ เตสุ วตฺตามี”ติ อรหติ อุปาสโก โสตาปนฺโน ภิกฺขุํ ปุถุชฺชนํ อภิวาเทตุํ ปจฺจุฏฺฐาตุํ.}

\prose{“อุปคโต โส สมณลิงฺคํ, พุทฺธาธิปฺปาเย ฐิโต, เตนาหํ ลิงฺเคน ทูรมปคโต”ติ อรหติ อุปาสโก โสตาปนฺโน ภิกฺขุํ ปุถุชฺชนํ อภิวาเทตุํ ปจฺจุฏฺฐาตุํ.}

\prose{“ปรูฬฺห\-กจฺฉ\-โลโม โส อนญฺชิต\-อมณฺฑิโต อนุลิตฺต\-สีลคนฺโธ, อหํ ปน มณฺฑนวิภูสนาภิรโต”ติ อรหติ อุปาสโก โสตาปนฺโน ภิกฺขุํ ปุถุชฺชนํ อภิวาเทตุํ ปจฺจุฏฺฐาตุํ.}

\prose{อปิ จ มหาราช “เย เต วีสติ สมณกรณา ธมฺมา เทฺว จ ลิงฺคานิ, สพฺเพเปเต ธมฺมา ภิกฺขุสฺส สํวิชฺชนฺติ, โส เยว\csromansharedfootnote{cb63cef12ad01f03}{โสเยว สนฺธิ \csromandash{} ม.พ.ป.} เต ธมฺเม ธาเรติ, อญฺเญปิ ตตฺถ สิกฺขาเปติ, โส เม อาคโม สิกฺขาปนญฺจ นตฺถี”ติ \csromanfolio{165}อรหติ อุปาสโก โสตาปนฺโน ภิกฺขุํ ปุถุชฺชนํ อภิวาเทตุํ ปจฺจุฏฺฐาตุํ.}

\prose{ยถา มหาราช ราชกุมาโร ปุโรหิตสฺส สนฺติเก วิชฺชํ อธียติ, ขตฺติย\-ธมฺมํ สิกฺขติ, โส อปเรน สมเยน อภิสิตฺโต อาจริยํ อภิวาเทติ ปจฺจุฏฺเฐติ “สิกฺขาปโก เม อยนฺ”ติ, เอวเมว โข มหาราช “ภิกฺขุ สิกฺขาปโก วํสธโร”ติ อรหติ อุปาสโก โสตาปนฺโน ภิกฺขุํ ปุถุชฺชนํ อภิวาเทตุํ ปจฺจุฏฺฐาตุํ.}

\prosewithcloserruled{อปิ จ มหาราช อิมินาเปตํ ปริยาเยน ชานาหิ ภิกฺขุภูมิยา มหนฺตตํ อสม\-วิปุล\-ภาวํ. ยทิ มหาราช อุปาสโก โสตาปนฺโน อรหตฺตํ สจฺฉิ กโรติ,\csromansymbolfootnote{+}{อุปริ 257 ปิฏฺเฐ.}เทฺวว ตสฺส คติโย ภวนฺติ อนญฺญา ตสฺมิํ เยว ทิวเส ปรินิพฺพาเยยฺย วา, ภิกฺขุภาวํ วา อุปคจฺเฉยฺย. อจลา หิ สา มหาราช ปพฺพชฺชา, มหตี อจฺจุคฺคตา, ยทิทํ ภิกฺขุภูมีติ. ญาณคโต ภนฺเต นาคเสน ปโญฺห สุนิพฺเพฐิโต พลวตา อติพุทฺธินา ตยา, น ยิมํ ปญฺหํ สมตฺโถ อญฺโญ เอวํ วินิเวเฐตุํ อญฺญตฺร ตวาทิเสน พุทฺธิมตาติ.}{เสฏฺฐ\-ธมฺม\-ปโญฺห ปฐโม.}
\csromanmatikaanchor{mtk.128}
\csromantocmarkref{subsubsection}{2. สพฺพสตฺตหิตผรณปญฺห}{mtk.128}
\bookmark[dest={mtk.128},level=3]{2. สพฺพสตฺตหิตผรณปญฺห}
\csromanheader{3}{2. \textbf{สพฺพสตฺต}\-\textbf{หิต}\-\textbf{ผรณ}\-\textbf{ปญฺห}}
\proseitem{1}{ภนฺเต นาคเสน ตุเมฺห ภณถ “ตถาคโต สพฺพสตฺตานํ อหิตมปเนตฺวา หิตมุปทหตี”ติ. ปุน จ ภณถ\csromansymbolfootnote{*}{อํ 2. 495 ปิฏฺเฐ.}\textbf{อคฺคิกฺขนฺธู}\-\textbf{ปเม} ธมฺมปริยาเย ภญฺญมาเน “สฏฺฐิมตฺตานํ ภิกฺขูนํ อุณฺหํ โลหิตํ มุขโต อุคฺคตนฺ”ติ. อคฺคิกฺขนฺธูปมํ ภนฺเต ธมฺม\-ปริยายํ เทเสนฺเตน ตถาคเตน สฏฺฐิมตฺตานํ ภิกฺขูนํ หิตมปเนตฺวา อหิตมุ\-ปทหิตํ. ยทิ ภนฺเต นาคเสน ตถาคโต สพฺพสตฺตานํ อหิตมปเนตฺวา หิตมุปทหติ, เตน หิ อคฺคิกฺขนฺธู\-ปเม ธมฺมปริยาเย ภญฺญมาเน สฏฺฐิมตฺตานํ ภิกฺขูนํ อุณฺหํ โลหิตํ มุขโต อุคฺคตนฺติ ยํ วจนํ, ตํ มิจฺฉา. ยทิ อคฺคิกฺขนฺธู\-ปเม ธมฺมปริยาเย ภญฺญมาเน สฏฺฐิมตฺตานํ ภิกฺขูนํ อุณฺหํ โลหิตํ มุขโต อุคฺคตํ, เตน หิ ตถาคโต สพฺพสตฺตานํ อหิตมปเนตฺวา หิตมุ\-ปทหตีติ ตมฺปิ วจนํ มิจฺฉา. อยมฺปิ อุภโต โกฏิโก ปโญฺห ตวานุปฺปตฺโต, โส ตยา นิพฺพาหิตพฺโพติ.}

\proseitem{2}{\csromanfolio{166}ตถาคโต มหาราช สพฺพสตฺตานํ อหิตมปเนตฺวา หิตมุปทหติ, อคฺคิกฺขนฺธู\-ปเม ธมฺมปริยาเย ภญฺญมาเน สฏฺฐิมตฺตานํ ภิกฺขูนํ อุณฺหํ โลหิตํ มุขโต อุคฺคตํ, ตญฺจ ปน น ตถาคตสฺส กเตน, เตสํ เยว\csromansharedfootnote{9a3af51fcf41c94d}{เตสํเยว นิคฺคหีตสนฺธิ \csromandash{} ม.พ.ป.} อตฺตโน กเตนาติ.}

\prose{ยทิ ภนฺเต นาคเสน ตถาคโต อคฺคิกฺขนฺธูปมํ ธมฺม\-ปริยายํ น ภาเสยฺย, อปิ นุ เตสํ อุณฺหํ โลหิตํ มุขโต อุคฺคจฺเฉยฺยาติ. น หิ มหาราช, มิจฺฉา\-ปฏิปนฺนานํ เตสํ ภควโต ธมฺม\-ปริยายํ สุตฺวา ปริฬาโห กาเย อุปฺปชฺชิ, เตน เตสํ ปริฬาเหน อุณฺหํ โลหิตํ มุขโต อุคฺคตนฺติ. เตน หิ ภนฺเต นาคเสน ตถาคตสฺเสว กเตน เตสํ อุณฺหํ โลหิตํ มุขโต อุคฺคตํ, ตถาคโต เยว ตตฺถ อธิกาโร เตสํ นาสนาย, ยถา นาม ภนฺเต นาคเสน อหิ วมฺมิกํ ปวิเสยฺย, อถญฺญตโร ปํสุกาโม ปุริโส วมฺมิกํ ภินฺทิตฺวา ปํสุํ หเรยฺย, ตสฺส ปํสุหรเณน วมฺมิกสฺส สุสิรํ ปิทเหยฺย, อถ ตตฺเถว โส อสฺสาสํ อลภมาโน มเรยฺย, นนุ โส ภนฺเต อหิ ตสฺส ปุริสสฺส กเตน มรณปฺปตฺโตติ. อาม มหาราชาติ. เอวเมว โข ภนฺเต นาคเสน ตถาคโต เยว ตตฺถ อธิกาโร เตสํ นาสนายาติ.}

\prose{ตถาคโต มหาราช ธมฺมํ เทสยมาโน อนุนย\-ปฺปฏิฆํ น กโรติ, อนุนย\-ปฺปฏิฆ\-วิปฺปมุตฺโต ธมฺมํ เทเสติ, เอวํ ธมฺเม เทสียมาเน เย ตตฺถ สมฺมาปฏิปนฺนา, เต พุชฺฌนฺติ. เย ปน มิจฺฉา\-ปฏิปนฺนา, เต ปตนฺติ. ยถา มหาราช ปุริสสฺส อมฺพํ วา ชมฺพุํ วา มธุกํ วา จาลยมานสฺส ยานิ ตตฺถ ผลานิ สารานิ ทฬฺห\-พนฺธนานิ, ตานิ ตตฺเถว อจฺจุตานิ ติฏฺฐนฺติ, ยานิ ตตฺถ ผลานิ ปูติวณฺฏ\-มูลานิ ทุพฺพล\-พนฺธนานิ, ตานิ ปตนฺติ. เอวเมว โข มหาราช ตถาคโต ธมฺมํ เทสยมาโน อนุนย\-ปฺปฏิฆํ น กโรติ, อนุนย\-ปฺปฏิฆ\-วิปฺปมุตฺโต ธมฺมํ เทเสติ, เอวํ ธมฺเม เทสียมาเน เย ตตฺถ สมฺมาปฏิปนฺนา, เต พุชฺฌนฺติ. เย ปน มิจฺฉา\-ปฏิปนฺนา, เต ปตนฺติ.}

\prose{ยถา วา ปน มหาราช กสฺสโก ธญฺญํ โรเปตุกาโม เขตฺตํ กสติ, ตสฺส กสนฺตสฺส อเนก\-สตสหสฺสานิ ติณานิ มรนฺติ. เอวเมว โข มหาราช ตถาคโต ปริปกฺก\-มานเส สตฺเต โพเธนฺโต\csromansharedfootnote{f0f71d9fada5b689}{โพเธตุํ (สี)} \csromanfolio{167}อนุนย\-ปฺปฏิฆ\-วิปฺปมุตฺโต ธมฺมํ เทเสติ, เอวํ ธมฺเม เทสียมาเน เย ตตฺถ สมฺมาปฏิปนฺนา, เต พุชฺฌนฺติ. เย ปน มิจฺฉา\-ปฏิปนฺนา, เต ติณานิ วิย มรนฺติ.}

\prose{ยถา วา ปน มหาราช มนุสฺสา รสเหตุ ยนฺเตน อุจฺฉุํ ปีฬยนฺติ, เตสํ อุจฺฉุํ ปีฬยมานานํ เย ตตฺถ ยนฺตมุขคตา กิมโย, เต ปีฬิยนฺติ. เอวเมว โข มหาราช ตถาคโต ปริปกฺก\-มานเส สตฺเต โพเธนฺโต ธมฺมยนฺตมภิปีฬยติ\csromansharedfootnote{0e48b81dc00cc432}{ธมฺมยนฺตมติปีฬยติ (ก)}, เย ตตฺถ มิจฺฉา\-ปฏิปนฺนา, เต กิมี วิย มรนฺตีติ.}

\prose{นนุ ภนฺเต นาคเสน เต ภิกฺขู ตาย ธมฺมเทสนาย ปติตาติ. อปิ นุ โข มหาราช ตจฺฉโก รุกฺขํ ตจฺฉนฺโต\csromansharedfootnote{16b7a415cc572dfb}{รกฺขนฺโต (สี, อิ)} อุชุกํ ปริสุทฺธํ กโรตีติ. น หิ ภนฺเต, วชฺชนียํ อปเนตฺวา ตจฺฉโก รุกฺขํ อุชุกํ ปริสุทฺธํ กโรตีติ. เอวเมว โข มหาราช ตถาคโต ปริสํ รกฺขนฺโต น สกฺโกติ โพธเนยฺเย\csromansharedfootnote{c1ade8911adfe380}{อโพธนีเย (สฺยา) 4. อตฺตเชน ผเลน (สี)} สตฺเต โพเธตุํ, มิจฺฉา\-ปฏิปนฺเน ปน สตฺเต อปเนตฺวา โพธเนยฺเย สตฺเต โพเธติ, อตฺตกเตน ปน เต มหาราช มิจฺฉา\-ปฏิปนฺนา ปตนฺติ.}

\prose{ยถา มหาราช กทลี เวฬุ อสฺสตรี อตฺตเชน\csromansharedfootnote{c4108b76501a80cf}{อตฺตเชน ผเลน (สี)} หญฺญติ, เอวเมว โข มหาราช เย เต มิจฺฉา\-ปฏิปนฺนา, เต อตฺตกเตน หญฺญนฺติ ปตนฺติ.}

\prose{ยถา มหาราช โจรา อตฺตกเตน จกฺขุปฺปาฏนํ สูลาโรปนํ สีสจฺเฉทนํ ปาปุณนฺติ, เอวเมว โข มหาราช เย เต มิจฺฉา\-ปฏิปนฺนา, เต อตฺตกเตน หญฺญนฺติ ปตนฺติ\csromansharedfootnote{0cc980ae109b8560}{ชินสาสนา ปตนฺติ (สี, อิ)}, เยสํ มหาราช สฏฺฐิมตฺตานํ ภิกฺขูนํ อุณฺหํ โลหิตํ มุขโต อุคฺคตํ, เตสํ ตํ เนว ภควโต กเตน, น ปเรสํ กเตน, อถ โข อตฺตโน เยว กเตน.}

\prosewithcloserruled{ยถา มหาราช ปุริโส สพฺพชนสฺส อมตํ ทเทยฺย, เต ตํ อมตํ อสิตฺวา อโรคา ทีฆายุกา สพฺพีติโต\csromansharedfootnote{e59fe1a9bf51d7c5}{สพฺพีติยา (สี)} ปริมุจฺเจยฺยุํ, อถญฺญตโร ปุริโส ทุรุปจาเรน ตํ อสิตฺวา มรณํ ปาปุเณยฺย, อปิ นุ โข โส มหาราช อมตทายโก ปุริโส ตโตนิทานํ กิญฺจิ อปุญฺญํ อาปชฺเชยฺยาติ. น หิ ภนฺเตติ. เอวเมว โข มหาราช ตถาคโต ทสสหสฺสิยา โลกธาตุยา เทวมนุสฺสานํ อมตํ \csromanfolio{168}ธมฺมทานํ เทติ, เย เต สตฺตา ภพฺพา, เต ธมฺมามเตน พุชฺฌนฺติ. เย ปน เต สตฺตา อภพฺพา, เต ธมฺมามเตน หญฺญนฺติ ปตนฺติ. โภชนํ มหาราช สพฺพสตฺตานํ ชีวิตํ รกฺขติ, ตเมกจฺเจ ภุญฺชิตฺวา วิสูจิกาย มรนฺติ, อปิ นุ โข โส มหาราช โภชนทายโก ปุริโส ตโตนิทานํ กิญฺจิ อปุญฺญํ อาปชฺเชยฺยาติ. น หิ ภนฺเตติ. เอวเมวโข มหาราช ตถาคโต ทสสหสฺสิยา โลกธาตุยา เทวมนุสฺสานํ อมตํ ธมฺมทานํ เทติ, เย เต สตฺตา ภพฺพา, เต ธมฺมามเตน พุชฺฌนฺติ. เย ปน เต สตฺตา อภพฺพา, เต ธมฺมามเตน หญฺญนฺติ ปตนฺตีติ. สาธุ ภนฺเต นาคเสน, เอวเมตํ ตถา สมฺปฏิจฺฉามีติ.}{สพฺพสตฺต\-หิต\-ผรณ\-ปโญฺห ทุติโย.}
\csromanmatikaanchor{mtk.129}
\csromantocmarkref{subsubsection}{3. วตฺถคุยฺหนิทสฺสนปญฺห}{mtk.129}
\bookmark[dest={mtk.129},level=3]{3. วตฺถคุยฺหนิทสฺสนปญฺห}
\csromanheader{3}{3. \textbf{วตฺถคุยฺห}\-\textbf{นิทสฺสน}\-\textbf{ปญฺห}}
\prose{3. ภนฺเต นาคเสน ภาสิตมฺเปตํ ตถาคเตน\csromandash{}}

\setlength{\csromangathapagewidth}{0pt}
\setlength{\csromangathawakwidth}{0pt}
\csromangathameasuregroup{“กาเยน สํวโร สาธุ\textsuperscript{1}, \\ มนสา สํวโร สาธุ,}{\csromangathabat{“กาเยน สํวโร สาธุ\textsuperscript{1},}{สาธุ วาจาย สํวโร.} \\ \csromangathabat{มนสา สํวโร สาธุ,}{สาธุ สพฺพตฺถ สํวโร”ติ.}}
\csromangathasetleft{“กาเยน สํวโร สาธุ\textsuperscript{1}, \\ มนสา สํวโร สาธุ,}
\csromangathagroup{\csromangathastanza{\csromangathabat{“กาเยน สํวโร สาธุ\csromansharedfootnote{c9292d08ada5a5eb}{ขุ 1. 65 ปิฏฺเฐ ธมฺมปเท.},}{สาธุ วาจาย สํวโร.} \\ \csromangathabat{มนสา สํวโร สาธุ,}{สาธุ สพฺพตฺถ สํวโร”ติ.}}}

\prose{ปุน จ ตถาคโต จตุนฺนํ ปริสานํ มชฺเฌ นิสีทิตฺวา ปุรโต เทวมนุสฺสานํ\csromansymbolfootnote{*}{ม 2. 351 ปิฏฺเฐ.}เสลสฺส พฺราหฺมณสฺส โกโสหิตํ วตฺถคุยฺหํ ทสฺเสสิ. ยทิ ภนฺเต นาคเสน ภควตา ภณิตํ “กาเยน สํวโร สาธู”ติ, เตน หิ เสลสฺส พฺราหฺมณสฺส โกโสหิตํ วตฺถคุยฺหํ ทสฺเสสีติ ยํ วจนํ, ตํ มิจฺฉา. ยทิ เสลสฺส พฺราหฺมณสฺส โกโสหิตํ วตฺถคุยฺหํ ทสฺเสสิ, เตน หิ “กาเยน สํวโร สาธู”ติ ตมฺปิ วจนํ มิจฺฉา. อยมฺปิ อุภโต โกฏิโก ปโญฺห ตวานุปฺปตฺโต, โส ตยา นิพฺพาหิตพฺโพติ.}

\prose{ภาสิตมฺเปตํ มหาราช ภควตา “กาเยน สํวโร สาธู”ติ, เสลสฺส จ พฺราหฺมณสฺส โกโสหิตํ วตฺถคุยฺหํ ทสฺสิตํ. ยสฺส โข มหาราช ตถาคเต กงฺขา อุปฺปนฺนา, ตสฺส โพธนตฺถาย ภควา อิทฺธิยา ตปฺปฏิภาคํ กายํ ทสฺเสติ, โส เยว\csromansharedfootnote{cb63cef12ad01f03}{โสเยว สนฺธิ \csromandash{} ม.พ.ป.} ตํ ปาฏิหาริยํ ปสฺสตีติ.}

\prose{\csromanfolio{169}โก ปเนตํ ภนฺเต นาคเสน สทฺทหิสฺสติ, ยํ ปริสคโต เอโก เยว ตํ คุยฺหํ ปสฺสติ, อวเสสา ตตฺเถว วสนฺตา น ปสฺสนฺตีติ. อิงฺฆ เม ตฺวํ ตตฺถ การณํ อุปทิส, การเณน มํ สญฺญาเปหีติ. ทิฏฺฐปุพฺโพ ปน ตยา มหาราช โกจิ พฺยาธิโต ปุริโส ปริกิณฺโณ ญาติมิตฺเตหีติ. อาม ภนฺเตติ. อปิ นุ โข สา มหาราช ปริสา ปสฺสเตตํ เวทนํ, ยาย โส ปุริโส เวทนาย เวทยตีติ. น หิ ภนฺเต, อตฺตนา เยว โส ภนฺเต ปุริโส เวทยตีติ. เอวเมว โข มหาราช ยสฺเสว ตถาคเต กงฺขา อุปฺปนฺนา, ตสฺเสว ตถาคโต โพธนตฺถาย อิทฺธิยา ตปฺปฏิภาคํ กายํ ทสฺเสติ, โส เยว\csromansharedfootnote{cb63cef12ad01f03}{โสเยว สนฺธิ \csromandash{} ม.พ.ป.} ตํ ปาฏิหาริยํ ปสฺสติ.}

\prose{ยถา วา ปน มหาราช กญฺจิเทว ปุริสํ ภูโต อาวิเสยฺย, อปิ นุ โข สา มหาราช ปริสา ปสฺสติ ตํ ภูตาคมนนฺติ. น หิ ภนฺเต, โส เยว\csromansharedfootnote{cb63cef12ad01f03}{โสเยว สนฺธิ \csromandash{} ม.พ.ป.} อาตุโร ตสฺส ภูตสฺส อาคมนํ ปสฺสตีติ. เอวเมว โข มหาราช ยสฺเสว ตถาคเต กงฺขา อุปฺปนฺนา, ตสฺเสว ตถาคโต โพธนตฺถาย อิทฺธิยา ตปฺปฏิภาคํ กายํ ทสฺเสติ, โส เยว\csromansharedfootnote{cb63cef12ad01f03}{โสเยว สนฺธิ \csromandash{} ม.พ.ป.} ตํ ปาฏิหาริยํ ปสฺสตีติ.}

\prose{ทุกฺกรํ ภนฺเต นาคเสน ภควตา กตํ, ยํ เอกสฺสปิ อทสฺสนียํ, ตํ ทสฺเสนฺเตนาติ. น มหาราช ภควา คุยฺหํ ทสฺเสสิ, อิทฺธิยา ปน ฉายํ ทสฺเสสีติ. ฉายายปิ ภนฺเต ทิฏฺฐาย ทิฏฺฐํ เยว โหติ คุยฺหํ, ยํ ทิสฺวา นิฏฺฐํ คโตติ. ทุกฺกรญฺจาปิ มหาราช ตถาคโต กโรติ โพธเนยฺเย สตฺเต โพเธตุํ, ยทิ มหาราช ตถาคโต กิริยํ หาเปยฺย, โพธเนยฺยา สตฺตา น พุชฺเฌยฺยุํ. ยสฺมา จ โข มหาราช โยคญฺญู ตถาคโต โพธเนยฺเย สตฺเต โพเธตุํ, ตสฺมา ตถาคโต เยน เยน โยเคน โพธเนยฺยา พุชฺฌนฺติ, เตน เตน โยเคน โพธเนยฺเย โพเธติ.}

\prose{ยถา มหาราช ภิสกฺโก สลฺลกตฺโต เยน เยน เภสชฺเชน อาตุโร อโรโค โหติ, เตน เตน เภสชฺเชน อาตุรํ อุปสงฺกมติ, วมนียํ วเมติ, วิเรจนียํ วิเรเจติ, อนุเลปนียํ อนุลิมฺเปติ, อนุวาสนียํ อนุวาเสติ. เอวเมว โข มหาราช ตถาคโต เยน เยน โยเคน โพธเนยฺยา สตฺตา พุชฺฌนฺติ, เตน เตน โยเคน โพเธติ.}

\prose{\csromanfolio{170}ยถา วา ปน มหาราช อิตฺถี มูฬฺหคพฺภา ภิสกฺกสฺส อทสฺสนียํ คุยฺหํ ทสฺเสติ, เอวเมว โข มหาราช ตถาคโต โพธเนยฺเย สตฺเต โพเธตุํ อทสฺสนียํ คุยฺหํ อิทฺธิยา ฉายํ ทสฺเสสิ, นตฺถิ มหาราช อทสฺสนีโย นาม โอกาโส ปุคฺคลํ อุปาทาย, ยทิ มหาราช โกจิ ภควโต หทยํ ทิสฺวา พุชฺเฌยฺย, ตสฺสปิ ภควา โยเคน หทยํ ทสฺเสยฺย, โยคญฺญู มหาราช ตถาคโต เทสนากุสโล.}

\prose{นนุ มหาราช ตถาคโต เถรสฺส\csromansymbolfootnote{*}{ขุ 1. 101 ปิฏฺเฐ อุทาเน.}นนฺทสฺส อธิมุตฺติํ ชานิตฺวา ตํ เทวภวนํ เนตฺวา เทวกญฺญาโย ทสฺเสสิ “อิมินายํ กุลปุตฺโต พุชฺฌิสฺสตี”ติ, เตน จ โส กุลปุตฺโต พุชฺฌิ. อิติ โข มหาราช ตถาคโต อเนก\-ปริยาเยน สุภนิมิตฺตํ หีเฬนฺโต ครหนฺโต ชิคุจฺฉนฺโต ตสฺส โพธนเหตุ กกุฏปาทินิโย อจฺฉราโย ทสฺเสสิ. เอวมฺปิ ตถาคโต โยคญฺญู เทสนากุสโล.}

\prose{ปุน จปรํ มหาราช ตถาคโต เถรสฺส\csromansymbolfootnote{+}{ขุ 3. 64 ปิฏฺเฐ อปทาเน.}จูฬปนฺถกสฺส ภาตรา นิกฺกฑฺฒิตสฺส ทุกฺขิตสฺส ทุมฺมนสฺส อุปคนฺตฺวา สุขุมํ โจฬขณฺฑํ อทาสิ “อิมินายํ กุลปุตฺโต พุชฺฌิสฺสตี”ติ, โส จ กุลปุตฺโต เตน การเณน ชินสาสเน วสีภาวํ ปาปุณิ. เอวมฺปิ มหาราช ตถาคโต โยคญฺญู เทสนากุสโล.}

\prose{ปุน จปรํ มหาราช ตถาคโต พฺราหฺมณสฺส โมฆราชสฺส ยาว ตติยํ ปญฺหํ ปุฏฺโฐ น พฺยากาสิ “เอวมิมสฺส กุลปุตฺตสฺส มาโน อุปสมิสฺสติ, มานูปสมา อภิสมโย ภวิสฺสตี”ติ, เตน จ ตสฺส กุลปุตฺตสฺส มาโน อุปสมิ, มานูปสมา โส พฺราหฺมโณ ฉสุ อภิญฺญาสุ วสีภาวํ ปาปุณิ. เอวมฺปิ มหาราช ตถาคโต โยคญฺญู เทสนากุสโลติ.}

\prose{สาธุ ภนฺเต นาคเสน, สุนิพฺเพฐิโต ปโญฺห พหุวิเธหิ การเณหิ, คหนํ อคหนํ กตํ, อนฺธกาโร อาโลโก กโต, คณฺฐิ ภินฺโน, ภคฺคา ปรวาทา, ชินปุตฺตานํ จกฺขุํ ตยา อุปฺปาทิตํ, นิปฺปฏิภานา ติตฺถิยา, ตฺวํ คณิวร\-ปวรมา\-สชฺชาติ.}

\vspace{\tipitakaparskip}
\csromancenter{วตฺถคุยฺห\-นิทสฺสน\-ปโญฺห ตติโย.}
\csromanmatikaanchor{mtk.130}
\csromantocmarkref{subsubsection}{4. ผรุสวาจาภาวปญฺห}{mtk.130}
\bookmark[dest={mtk.130},level=3]{4. ผรุสวาจาภาวปญฺห}
\csromanmarkh{4. เมณฺฑกปญฺห}\csromanmarkhii{4. ผรุสวาจา\-ภาว\-ปญฺห}\csromanheadernomark{3}{4. เมณฺฑกปญฺห \textbf{4. ผรุสวาจา}\-\textbf{ภาว}\-\textbf{ปญฺห}}
\proseitem{4}{\csromanfolio{171}ภนฺเต นาคเสน ภาสิตมฺเปตํ เถเรน สาริปุตฺเตน ธมฺม\-เสนาปตินา\csromansymbolfootnote{*}{ที 3. 182; อํ 2. 459 ปิฏฺเฐสุ.}“ปริสุทฺธ\-วจี\-สมาจาโร อาวุโส ตถาคโต, นตฺถิ ตถาคตสฺส วจีทุจฺจริตํ, ยํ ตถาคโต รกฺเขยฺย ‘มา เม อิทํ ปโร อญฺญาสี’ติ”. ปุน จ ตถาคโต เถรสฺส สุทินฺนสฺส กลนฺท\-ปุตฺตสฺส อปราเธ ปาราชิกํ ปญฺญเปนฺโต ผรุสาหิ วาจาหิ\csromansymbolfootnote{+}{วิ 1. 23 ปิฏฺเฐ.}โมฆปุริส\-วาเทน สมุทาจริ, เตน จ โส เถโร โมฆปุริส\-วาเทน มงฺกุ\-จิตฺตวเสน รุนฺธิตตฺตา วิปฺปฏิสารี นาสกฺขิ อริยมคฺคํ ปฏิวิชฺฌิตุํ. ยทิ ภนฺเต นาคเสน ปริสุทฺธ\-วจี\-สมาจาโร ตถาคโต, นตฺถิ ตถาคตสฺส วจีทุจฺจริตํ, เตน หิ ตถาคเตน เถรสฺส สุทินฺนสฺส กลนฺท\-ปุตฺตสฺส อปราเธ โมฆปุริส\-วาเทน สมุทาจิณฺณนฺติ ยํ วจนํ, ตํ มิจฺฉา. ยทิ ภควตา เถรสฺส สุทินฺนสฺส กลนฺท\-ปุตฺตสฺส อปราเธ โมฆปุริส\-วาเทน สมุทาจิณฺณํ, เตน หิ ปริสุทฺธ\-วจี\-สมาจาโร ตถาคโต, นตฺถิ ตถาคตสฺส วจีทุจฺจริตนฺติ ตมฺปิ วจนํ มิจฺฉา. อยมฺปิ อุภโต โกฏิโก ปโญฺห ตวานุปฺปตฺโต, โส ตยา นิพฺพาหิตพฺโพติ.}

\prose{ภาสิตมฺเปตํ มหาราช เถเรน สาริปุตฺเตน ธมฺม\-เสนาปตินา “ปริสุทฺธ\-วจี\-สมาจาโร อาวุโส ตถาคโต, นตฺถิ ตถาคตสฺส วจีทุจฺจริตํ, ยํ ตถาคโต รกฺเขยฺย ‘มา เม อิทํ ปโร อญฺญาสี’ติ”. อายสฺมโต จ สุทินฺนสฺส กลนฺท\-ปุตฺตสฺส อปราเธ ปาราชิกํ ปญฺญเปนฺเตน ภควตา โมฆปุริส\-วาเทน สมุทาจิณฺณํ, ตญฺจ ปน อทุฏฺฐจิตฺเตน อสารมฺเภน ยาถาว\-ลกฺขเณน. กิญฺจ ตตฺถ ยาถาว\-ลกฺขณํ, ยสฺส มหาราช ปุคฺคลสฺส อิมสฺมิํ อตฺตภาเว จตุสจฺจาภิสมโย น โหติ, ตสฺส ปุริสตฺตนํ โมฆํ อญฺญํ กยิรมานํ อญฺเญน สมฺภวติ, เตน วุจฺจติ “โมฆปุริโส”ติ. อิติ มหาราช ภควตา อายสฺมโต สุทินฺนสฺส กลนฺท\-ปุตฺตสฺส สตาววจเนน สมุทาจิณฺณํ, โน อภูตวาเทนาติ.}

\prose{สภาวมฺปิ ภนฺเต นาคเสน โย อกฺโกสนฺโต ภณติ, ตสฺส มยํ กหาปณํ ทณฺฑํ ธาเรม, อปราโธ เยว โส วตฺถุํ นิสฺสาย วิสุํ โวหารํ อาจรนฺโต อกฺโกสตีติ. อตฺถิ ปน มหาราช \csromanfolio{172}สุตปุพฺพํ ตยา ขลิตสฺส อภิวาทนํ วา ปจฺจุฏฺฐานํ วา สกฺการํ วา อุปายนา\-นุปฺปทานํ วาติ. น หิ ภนฺเต, ยโต กุโตจิ ยตฺถ กตฺถจิ ขลิโต, โส ปริภาสนา\-รโห โหติ ตชฺชนารโห, อุตฺตมงฺคมฺปิสฺส ฉินฺทนฺติ หนนฺติปิ พนฺธนฺติปิ ฆาเตนฺติปิ ฌาเปนฺติปีติ\csromansharedfootnote{dd9f33c79a523d1c}{ชาเปนฺติปีติ (สี, อิ)}. เตน หิ มหาราช ภควตา กิริยา เยว กตา, โน อกิริยาติ.}

\prose{กิริยมฺปิ ภนฺเต นาคเสน กุรุมาเนน ปติรูเปน กาตพฺพํ อนุจฺฉวิเกน, สวเนนปิ ภนฺเต นาคเสน ตถาคตสฺส สเทวโก โลโก โอตฺตปฺปติ หิริยติ ภิยฺโย ทสฺสเนน ตตุตฺตริํ อุปสงฺกมเนน ปยิรุปาส\-เนนาติ. อปิ นุ โข มหาราช ติกิจฺฉโก อภิสนฺเน กาเย กุปิเต โทเส สิเนหนียานิ เภสชฺชานิ เทตีติ. น หิ ภนฺเต ติณฺหานิ เลขนียานิ เภสชฺชานิ เทตีติ. เอวเมว โข มหาราช ตถาคโต สพฺพ\-กิเลสพฺยาธิ\-วูปสมาย อนุสิฏฺฐิํ เทติ, ผรุสาปิ มหาราช ตถาคตสฺส วาจา สตฺเต สิเนหยติ, มุทุเก กโรติ. ยถา มหาราช อุณฺหมฺปิ อุทกํ ยํ กิญฺจิ สิเนหนียํ สิเนหยติ, มุทุกํ กโรติ, เอวเมว โข มหาราช ผรุสาปิ ตถาคตสฺส วาจา อตฺถวตี โหติ กรุณาสหคตา. ยถา มหาราช ปิตุวจนํ ปุตฺตานํ อตฺถวนฺตํ โหติ กรุณา\-สหคตํ, เอวเมว โข มหาราช ผรุสาปิ ตถาคตสฺส วาจา อตฺถวตี โหติ กรุณาสหคตา. ผรุสาปิ มหาราช ตถาคตสฺส วาจา สตฺตานํ กิเลสปฺปหานา\csromansharedfootnote{25b226086d06f2f8}{กิเลสปฺปหานาย (สี)} โหติ. ยถา มหาราช ทุคฺคนฺธมฺปิ โคมุตฺตํ ปีตํ วิรสมฺปิ อคทํ ขายิตํ สตฺตานํ พฺยาธิํ หนติ, เอวเมว โข มหาราช ผรุสาปิ ตถาคตสฺส วาจา อตฺถวตี กรุณาสหคตา. ยถา มหาราช มหนฺโตปิ ตูลปุญฺโช\csromansharedfootnote{a2056901cacdee0c}{ตูลปิจุ (สี, สฺยา)} ปรสฺส กาเย นิปติตฺวา รุชํ น กโรติ, เอวเมว โข มหาราช ผรุสาปิ ตถาคตสฺส วาจา น กสฺสจิ ทุกฺขํ อุปฺปาเทตีติ. สุวินิจฺฉิโต ภนฺเต นาคเสน ปโญฺห พหูหิ การเณหิ, สาธุ ภนฺเต นาคเสน, เอวเมตํ ตถา สมฺปฏิจฺฉามีติ.}

\vspace{\tipitakaparskip}
\csromancenter{ผรุสวาจา\-ภาว\-ปโญฺห จตุตฺโถ.}
\csromanmatikaanchor{mtk.131}
\csromantocmarkref{subsubsection}{5. รุกฺขอเจตนาภาวปญฺห}{mtk.131}
\bookmark[dest={mtk.131},level=3]{5. รุกฺขอเจตนาภาวปญฺห}
\csromanmarkh{4. เมณฺฑกปญฺห}\csromanmarkhii{5. รุกฺข\-อเจตนา\-ภาว\-ปญฺห}\csromanheadernomark{3}{4. เมณฺฑกปญฺห 5. รุกฺข\-อเจตนา\-ภาว\-ปญฺห}
\prose{\csromanfolio{173}5. ภนฺเต นาคเสน ภาสิตมฺเปตํ ตถาคเตน\csromandash{}}

\setlength{\csromangathapagewidth}{0pt}
\setlength{\csromangathawakwidth}{0pt}
\csromangathameasuregroup{}{“อเจตนํ พฺราหฺมณ อสฺสุณนฺตํ, \\ ชาโน อชานนฺตมิมํ ปลาสํ. \\ อารทฺธวีริโย ธุวํ อปฺปมตฺโต, \\ สุขเสยฺยํ ปุจฺฉสิ กิสฺส เหตู”ติ\textsuperscript{1}.}
\csromangathameasurewak{“อเจตนํ พฺราหฺมณ อสฺสุณนฺตํ, \\ ชาโน อชานนฺตมิมํ ปลาสํ. \\ อารทฺธวีริโย ธุวํ อปฺปมตฺโต, \\ สุขเสยฺยํ ปุจฺฉสิ กิสฺส เหตู”ติ\textsuperscript{1}.}
\setlength{\csromangathaleftcol}{0pt}
\csromangathagroup{\csromangathastanza{\csromangathawak{“อเจตนํ พฺราหฺมณ อสฺสุณนฺตํ, \\ ชาโน อชานนฺตมิมํ ปลาสํ. \\ อารทฺธวีริโย ธุวํ อปฺปมตฺโต, \\ สุขเสยฺยํ ปุจฺฉสิ กิสฺส เหตู”ติ\csromansharedfootnote{1a5701235e34b4e6}{ขุ 5. 96 ปิฏฺเฐ.}.}}}

\prose{ปุน จ ภณิตํ\csromandash{}}

\setlength{\csromangathapagewidth}{0pt}
\setlength{\csromangathawakwidth}{0pt}
\csromangathameasuregroup{“อิติ ผนฺทนรุกฺโขปิ, \\ มยฺหมฺปิ วจนํ อตฺถิ,}{\csromangathabat{“อิติ ผนฺทนรุกฺโขปิ,}{ตาวเท อชฺฌภาสถ.} \\ \csromangathabat{มยฺหมฺปิ วจนํ อตฺถิ,}{ภารทฺวาช สุโณหิ เม”ติ\textsuperscript{1}.}}
\csromangathasetleft{“อิติ ผนฺทนรุกฺโขปิ, \\ มยฺหมฺปิ วจนํ อตฺถิ,}
\csromangathagroup{\csromangathastanza{\csromangathabat{“อิติ ผนฺทนรุกฺโขปิ,}{ตาวเท อชฺฌภาสถ.} \\ \csromangathabat{มยฺหมฺปิ วจนํ อตฺถิ,}{ภารทฺวาช สุโณหิ เม”ติ\csromansharedfootnote{ef0f0e0feb4d1a39}{ขุ 5. 259 ปิฏฺเฐ.}.}}}

\prose{ยทิ ภนฺเต นาคเสน รุกฺโข อเจตโน, เตน หิ ผนฺทเนน รุกฺเขน ภารทฺวาเชน สห สลฺลปิตนฺติ ยํ วจนํ, ตํ มิจฺฉา. ยทิ ผนฺทเนน รุกฺเขน ภารทฺวาเชน สทฺธิํ สลฺลปิตํ, เตน หิ รุกฺโข อเจตโนติ ตมฺปิ วจนํ มิจฺฉา. อยมฺปิ อุภโต โกฏิโก ปโญฺห ตวานุปฺปตฺโต, โส ตยา นิพฺพาหิตพฺโพติ.}

\prose{ภาสิตมฺเปตํ มหาราช ภควตา “รุกฺโข อเจตโน”ติ, ผนฺทเนน จ รุกฺเขนํ ภารทฺวาเชน สทฺธิํ สลฺลปิตํ, ตญฺจ ปน วจนํ โลกสมญฺญาย ภณิตํ, นตฺถิ มหาราช อเจตนสฺส รุกฺขสฺส สลฺลาโป นาม, อปิ จ มหาราช ตสฺมิํ รุกฺเข อธิวตฺถาย เทวตาเยตํ อธิวจนํ รุกฺโขติ, รุกฺโข สลฺลปตีติ เจสา โลกปณฺณตฺติ, ยถา มหาราช สกฏํ ธญฺญสฺส ปริปูริตํ ธญฺญสกฏนฺติ ชโน โวหรติ, น จ ตํ ธญฺญมยํ สกฏํ, รุกฺขมยํ สกฏํ, ตสฺมิํ สกเฏ ธญฺญสฺส ปน อากิริตตฺตา ธญฺญสกฏนฺติ ชโน โวหรติ, เอวเมว โข มหาราช น รุกฺโข สลฺลปติ, รุกฺโข อเจตโน, ยา ปน ตสฺมิํ รุกฺเข อธิวตฺถา เทวตา, ตสฺสา เยว\csromansharedfootnote{19cf0a39ee311ce1}{ตสฺสาเยว สนฺธิ \csromandash{} ม.พ.ป.} ตํ อธิวจนํ รุกฺโขติ, รุกฺโข สลฺลปตีติ เจสา โลกปณฺณตฺติ.}

\prose{ยถา วา ปน มหาราช ทธิํ มนฺถยมาโน ตกฺกํ มนฺเถมีติ โวหรติ, น ตํ ตกฺกํ, ยํ โส มนฺเถติ, ทมิํ เยว โส มนฺเถนฺโต ตกฺกํ มนฺเถมีติ โวหรติ, เอวเมว โข มหาราช น รุกฺโข สลฺลปติ, รุกฺโข}

\prose{\csromanfolio{174}อเจตโน. ยา ปน ตสฺมิํ รุกฺเข อธิวตฺถา เทวตา, ตสฺสาเยว ตํ อธิวจนํ รุกฺโขติ, รุกฺโข สลฺลปตีติ เจสา โลกปณฺณตฺติ.}

\prose{ยถา วา ปน มหาราช อสนฺตํ สาเธตุกาโม สนฺตํ สาเธมีติ โวหรติ, อสิทฺธํ สิทฺธนฺติ โวหรติ, เอวเมสา โลกสมญฺญา, เอวเมว โข มหาราช น รุกฺโข สลฺลปติ, รุกฺโข อเจตโน. ยา ปน ตสฺมิํ รุกฺเข อธิวตฺถา เทวตา, ตสฺสาเยว ตํ อธิวจนํ รุกฺโขติ, รุกฺโข สลฺลปตีติ เจสา โลกปณฺณตฺติ, ยาย มหาราช โลกสมญฺญาย ชโน โวหรติ, ตถาคโตปิ ตาเยว โลกสมญฺญาย สตฺตานํ ธมฺมํ เทเสตีติ. สาธุ ภนฺเต นาคเสน, เอวเมตํ ตถา สมฺปฏิจฺฉามีติ.}

\prose{รุกฺข\-อเจตนา\-ภาว\-ปโญฺห ปญฺจโม.}
\csromansectionrule

\csromanmatikaanchor{mtk.132}
\csromantocmarkref{subsubsection}{6. ปิณฺฑปาตมหปฺผลปญฺห}{mtk.132}
\bookmark[dest={mtk.132},level=3]{6. ปิณฺฑปาตมหปฺผลปญฺห}
\csromanheader{3}{6. \textbf{ปิณฺฑปาต}\-\textbf{มหปฺผล}\-\textbf{ปญฺห}}
\proseitem{6}{ภนฺเต นาคเสน ภาสิตมฺเปตํ ธมฺม\-สงฺคีติการเกหิ เถเรหิ\_\_}

\setlength{\csromangathapagewidth}{0pt}
\setlength{\csromangathawakwidth}{0pt}
\csromangathameasuregroup{“จุนฺทสฺส ภตฺตํ ภุญฺชิตฺวา, \\ อาพาธํ สมฺผุสี ธีโร,}{\csromangathabat{“จุนฺทสฺส ภตฺตํ ภุญฺชิตฺวา,}{กมฺมารสฺสาติ เม สุตํ.} \\ \csromangathabat{อาพาธํ สมฺผุสี ธีโร,}{ปพาฬฺหํ มารณนฺติกนฺ”ติ\textsuperscript{1}.}}
\csromangathasetleft{“จุนฺทสฺส ภตฺตํ ภุญฺชิตฺวา, \\ อาพาธํ สมฺผุสี ธีโร,}
\csromangathagroup{\csromangathastanza{\csromangathabat{“จุนฺทสฺส ภตฺตํ ภุญฺชิตฺวา,}{กมฺมารสฺสาติ เม สุตํ.} \\ \csromangathabat{อาพาธํ สมฺผุสี ธีโร,}{ปพาฬฺหํ มารณนฺติกนฺ”ติ\csromansharedfootnote{e9bcb89e4db071d9}{ที 2. 106 ปิฏฺเฐ.}.}}}

\prose{ปุน จ ภควตา ภณิตํ\csromansymbolfootnote{*}{ที 2. 112 ปิฏฺเฐ.}“เทฺวเม อานนฺท ปิณฺฑปาตา สมสมผลา สมวิปากา อติวิย อญฺเญหิ ปิณฺฑปาเตหิ มหปฺผลตรา จ มหา\-นิสํส\-ตรา จ, กตเม เทฺว, ยญฺจ ปิณฺฑปาตํ ปริภุญฺชิตฺวา ตถาคโต อนุตฺตรํ สมฺมาสมฺโพธิํ อภิสมฺพุชฺฌิ, ยญฺจ ปิณฺฑปาตํ ปริภุญฺชิตฺวา ตถาคโต อนุปาทิเสสาย นิพฺพานธาตุยา ปรินิพฺพายติ, อิเม เทฺว ปิณฺฑปาตา สมสมผลา สมวิปากา, อติวิย อญฺเญหิ ปิณฺฑปาเตหิ มหปฺผลตรา จ มหา\-นิสํส\-ตรา จา”ติ. ยทิ ภนฺเต นาคเสน ภควโต จุนฺทสฺส ภตฺตํ ภุตฺตาวิสฺส\csromansharedfootnote{7d758a69dfb66f88}{ภุญฺชิตฺวา (สี)} ขโร อาพาโธ อุปฺปนฺโน, ปพาฬฺหา จ เวทนา ปวตฺตา มารณนฺติกา, เตน หิ “เทฺวเม อานนฺท ปิณฺฑปาตา สมสมผลา สมวิปากา อติวิย อญฺเญหิ ปิณฺฑปาเตหิ มหปฺผลตารา จ มหา\-นิสํส\-ตรา จา”ติ ยํ วจนํ, ตํ มิจฺฉา, ยทิ เทฺวเม ปิณฺฑปาตา สมสมผลา สมวิปากา อติวิย อญฺเญหิ ปิณฺฑปาเตหิ มหปฺผลตรา \csromanfolio{175}จ มหา\-นิสํส\-ตรา จ, เตน หิ ภควโต จุนฺทสฺส ภตฺตํ ภุตฺตาวิสฺส\csromansharedfootnote{7d758a69dfb66f88}{ภุญฺชิตฺวา (สี)} ขโร อาพาโธ อุปฺปนฺโน, ปพาฬฺหา จ เวทนา ปวตฺตา มารณนฺติกาติ ตมฺปิ วจนํ มิจฺฉา. กินฺนุ โข ภนฺเต นาคเสน โส ปิณฺฑปาโต วิสคตตาย มหปฺผโล, โรคุ\-ปฺปาท\-กตาย มหปฺผโล, อายุ\-วินา\-สกตาย มหปฺผโล, ภควโต ชีวิต\-หรณ\-ตาย มหปฺผโล. ตตฺถ เม การณํ พฺรูหิ ปรวาทานํ นิคฺคหาย, เอตฺถายํ ชโน สมฺมูโฬฺห โลภวเสน อติพหุํ ขายิเตน โลหิต\-ปกฺขนฺทิกา อุปฺปนฺนาติ. อยมฺปิ อุภโต โกฏิโก ปโญฺห ตวานุปฺปตฺโต, โส ตยา นิพฺพาหิตพฺโพติ.}

\prose{ภาสิตมฺเปตํ มหาราช ธมฺม\-สงฺคีติการเกหิ เถเรหิ\csromandash{}}

\setlength{\csromangathapagewidth}{0pt}
\setlength{\csromangathawakwidth}{0pt}
\csromangathameasuregroup{“จุนฺทสฺส ภตฺตํ ภุญฺชิตฺวา, \\ อาพาธํ สมฺผุสี ธีโร,}{\csromangathabat{“จุนฺทสฺส ภตฺตํ ภุญฺชิตฺวา,}{กมฺมารสฺสาติ เม สุตํ.} \\ \csromangathabat{อาพาธํ สมฺผุสี ธีโร,}{ปพาฬฺหํ มารณนฺติกนฺ”ติ.}}
\csromangathasetleft{“จุนฺทสฺส ภตฺตํ ภุญฺชิตฺวา, \\ อาพาธํ สมฺผุสี ธีโร,}
\csromangathagroup{\csromangathastanza{\csromangathabat{“จุนฺทสฺส ภตฺตํ ภุญฺชิตฺวา,}{กมฺมารสฺสาติ เม สุตํ.} \\ \csromangathabat{อาพาธํ สมฺผุสี ธีโร,}{ปพาฬฺหํ มารณนฺติกนฺ”ติ.}}}

\prose{ภควตา จ ภณิตํ “เทฺวเม อานนฺท ปิณฺฑปาตา สมสมผลา สมวิปากา อติวิย อญฺเญหิ ปิณฺฑปาเตหิ มหปฺผลตรา จ มหา\-นิสํส\-ตรา จ, กตเม เทฺว, ยญฺจ ปิณฺฑปาตํ ปริภุญฺชิตฺวา ตถาคโต อนุตฺตรํ สมฺมาสมฺโพธิํ อภิสมฺพุชฺฌิ, ยญฺจ ปิณฺฑปาตํ ปริภุญฺชิตฺวา ตถาคโต อนุปาทิเสสาย นิพฺพานธาตุยา ปรินิพฺพายติ\csromansharedfootnote{1b3b4f4f096f929d}{ปรินิพฺพายิ (สี)}, อิเม เทฺว ปิณฺฑปาตา สมสมผลา สมวิปากา, อติวิย อญฺเญหิ ปิณฺฑปาเตหิ มหปฺผลตรา จ มหา\-นิสํส\-ตรา จา”ติ.}

\prose{โส ปน ปิณฺฑปาโต พหุคุโณ อเนกานิสํโส. เทวตา มหาราช หฏฺฐา ปสนฺนมานสา “อยํ ภควโต ปจฺฉิโม ปิณฺฑปาโต”ติ ทิพฺพํ โอชํ สูกรมทฺทเว อากิริํสุ. ตญฺจ ปน สมฺมาปากํ ลหุปากํ\csromansharedfootnote{ffb85ab6c0acb243}{พหุปากํ (สี)} มนุญฺญํ พหุรสํ ชฏฺฐร\-คฺคิ\-เตชสฺส หิตํ. น มหาราช ตโตนิทานํ ภควโต โกจิ อนุปฺปนฺโน โรโค อุปฺปนฺโน, อปิ จ มหาราช ภควโต ปกติทุพฺพเล สรีเร ขีเณ อายุสงฺขาเร อุปฺปนฺโน โรโค ภิยฺโย อภิวฑฺฒิ.}

\prose{ยถา มหาราช ปกติยา ชลมาโน อคฺคิ อญฺญสฺมิํ อุปาทาเน ทินฺเน ภิยฺโย ปชฺชลติ, เอวเมว โข มหาราช ภควโต ปกติทุพฺพเล สรีเร ขีเณ อายุสงฺขาเร อุปฺปนฺโน โรโค ภิยฺโย อภิวฑฺฒิ.}

\prose{\csromanfolio{176}ยถา วา ปน มหาราช โสโต ปกติยา สนฺทมาโน อภิวุฏฺเฐ มหาเมเฆ ภิยฺโย มโหโฆ อุทกวาหโก โหติ, เอวเมว โข มหาราช ภควโต ปกติทุพฺพเล สรีเร ขีเณ อายุสงฺขาเร อุปฺปนฺโน โรโค ภิยฺโย อภิวฑฺฒิ.}

\prose{ยถา วา ปน มหาราช ปกติยา อภิสนฺนธาตุ กุจฺฉิ อญฺญสฺมิํ อชฺโฌหริเต ภิยฺโย อายเมยฺย\csromansharedfootnote{f88aca6bbb4b50ec}{อามเยยฺย (สี)}, เอวเมว โข มหาราช ภควโต ปกติทุพฺพเล สรีเร ขีเณ อายุสงฺขาเร อุปฺปนฺโน โรโค ภิยฺโย อภิวฑฺฒิ, นตฺถิ มหาราช ตสฺมิํ ปิณฺฑปาเต โทโส, น จ ตสฺส สกฺกา โทสํ อาโรเปตุนฺติ.}

\prose{ภนฺเต นาคเสน เกน การเณน เต เทฺว ปิณฺฑปาตา สมสมผลา สมวิปากา อติวิย อญฺเญหิ ปิณฺฑปาเตหิ มหปฺผลตรา จ มหา\-นิสํส\-ตรา จาติ. ธมฺมา\-นุมชฺชน\-สมาปตฺติวเสน มหาราช เต เทฺว ปิณฺฑปาตา สมสมผลา สมวิปากา อติวิย อญฺเญหิ ปิณฺฑปาเตหิ มหปฺผลตรา จ มหา\-นิสํส\-ตรา จาติ.}

\prose{ภนฺเต นาคเสน กตเมสํ ธมฺมานํ อนุมชฺชน\-สมาปตฺติวเสน เต เทฺว ปิณฺฑปาตา สมสมผลา สมวิปากา อติวิย อญฺเญหิ ปิณฺฑปาเตหิ มหปฺผลตรา จ มหา\-นิสํส\-ตรา จาติ. นวนฺนํ มหาราช อนุปุพฺพวิหาร\-สมาปตฺตีนํ อนุโลม\-ปฺปฏิโลม\-สมาปชฺชน\-วเสน เต เทฺว ปิณฺฑปาตา สมสมผลา สมวิปากา อติวิย อญฺเญหิ ปิณฺฑปาเตหิ มหปฺผลตรา จ มหา\-นิสํส\-ตรา จาติ.}

\prose{ภนฺเต นาคเสน ทฺวีสุ เยว ทิวเสสุ อธิมตฺตํ ตถาคโต นวา\-นุปุพฺพวิหารสมาปตฺติโย อนุโลม\-ปฺปฏิโลมํ สมาปชฺชีติ. อาม มหาราชาติ. อจฺฉริยํ ภนฺเต นาคเสน อพฺภุตํ ภนฺเต นาคเสน. ยํ อิมสฺมิํ พุทฺธกฺเขตฺเต อสทิสํ ปรมทานํ, ตมฺปิ อิเมหิ ทฺวีหิ ปิณฺฑปาเตหิ อคณิตํ. อจฺฉริยํ ภนฺเต นาคเสน, อพฺภุตํ ภนฺเต นาคเสน. ยาว มหนฺตา นวา\-นุปุพฺพวิหารสมาปตฺติโย, ยตฺร หิ นาม นวา\-นุปุพฺพวิหาร\-สมาปตฺติวเสน ทานํ มหปฺผลตรํ โหติ มหา\-นิสํส\-ตรญฺจ. สาธุ ภนฺเต นาคเสน, เอวเมตํ ตถา สมฺปฏิจฺฉามีติ.}

\vspace{\tipitakaparskip}
\csromancenter{ปิณฺฑปาต\-มหปฺผล\-ปโญฺห ฉฏฺโฐ.}
\csromanmatikaanchor{mtk.133}
\csromantocmarkref{subsubsection}{7. พุทฺธปูชนปญฺห}{mtk.133}
\bookmark[dest={mtk.133},level=3]{7. พุทฺธปูชนปญฺห}
\csromanmarkh{4. เมณฺฑกปญฺห}\csromanmarkhii{7. พุทฺธ\-ปูชน\-ปญฺห}\csromanheadernomark{3}{4. เมณฺฑกปญฺห \textbf{7. พุทฺธ}\-\textbf{ปูชน}\-\textbf{ปญฺห}}
\proseitem{7}{\csromanfolio{177}ภนฺเต นาคเสน ภาสิตมฺเปตํ ตถาคเตน\csromansymbolfootnote{*}{ที 2. 117 ปิฏฺเฐ.}“อพฺยาวฏา ตุเมฺห อานนฺท โหถ ตถาคตสฺส สรีรปูชายา”ติ. ปุน จ ภณิตํ\csromandash{}}

\prose{\csromansymbolfootnote{+}{ขุ 2. 113 ปิฏฺเฐ วิมานวตฺถุมฺหิ.}“ปูเชถ นํ ปูชนิยสฺส ธาตุํ.}

\prose{เอวํ กิร สคฺคมิโต คมิสฺสถา”ติ.}

\prose{ยทิ ภนฺเต นาคเสน ตถาคเตน ภณิตํ “อพฺยาวฏา ตุเมฺห อานนฺท โหถ ตถาคตสฺส สรีรปูชายา”ติ, เตน หิ “ปูเชถ นํ ปูชนิยสฺส ธาตุํ, เอวํ กรา สคฺคมิโต คมิสฺสถา”ติ ยํ วจนํ, ตํ มิจฺฉา. ยทิ ตถาคเตน ภณิตํ “ปูเชถ นํ ปูชนิยสฺส ธาตุํ, เอวํ กรา สคฺคมิโต คมิสฺสถา”ติ, เตน หิ “อพฺยาวฏา ตุเมฺห อานนฺท โหถ ตถาคตสฺส สรีรปูชายา”ติ ตมฺปิ วจนํ มิจฺฉา. อยมฺปิ อุภโต โกฏิโก ปโญฺห ตวานุปฺปตฺโต, โส ตยา นิพฺพาหิตพฺโพติ.}

\prose{ภาสิตมฺเปตํ มหาราช ภควตา “อพฺยาวฏา ตุเมฺห อานนฺท โหถ ตถาคตสฺส สรีรปูชายา”ติ, ปุน จ ภณิตํ “ปูเชถ นํ ปูชนิยสฺส ธาตุํ, เอวํ กรา สคฺคมิโต คมิสฺสถา”ติ, ตญฺจ ปน น สพฺเพสํ ชินปุตฺตานํ เยว อารพฺภ ภณิตํ “อพฺยาวฏา ตุเมฺห อานนฺท โหถ ตถาคตสฺส สรีรปูชายา”ติ. อกมฺมํ เหตํ มหาราช ชินปุตฺตานํ ยทิทํ ปูชา, สมฺมสนํ สงฺขารานํ, โยนิโส มนสิกาโร, สติปฏฺฐานา\-นุปสฺสนา, อารมฺมณ\-สาร\-คฺคาโห, กิเลสยุทฺธํ, สทตฺถม\-นุยุญฺชนา, เอตํ ชินปุตฺตานํ กรณียํ, อวเสสานํ เทวมนุสฺสานํ ปูชา กรณียา.}

\prose{ยถา มหาราช มหิยา ราชปุตฺตานํ หตฺถิ- อสฺสรถธนุถรุเลขมุทฺทา- สิกฺขาขคฺคมนฺตสุติสมฺมุติยุทฺธยุชฺฌาปนกิริยา กรณียา, อวเสสานํ ปุถุเวสฺส\-สุทฺทานํ กสิ วณิชฺชา โครกฺขา กรณียา, เอวเมว โข มหาราช อกมฺมํ เหตํ ชินปุตฺตานํ ยทิทํ ปูชา, สมฺมสนํ สงฺขารานํ, โยนิโส มนสิกาโร, สติปฏฺฐานา\-นุปสฺสนา, อารมฺมณ\-สาร\-คฺคาโห, กิเลสยุทฺธํ, สทตฺถม\-นุยุญฺชนา, เอตํ ชินปุตฺตานํ กรณียํ, อวเสสานํ เทวมนุสฺสานํ ปูชา กรณียา.}

\prosewithcloserruled{ยถา วา ปน มหาราช พฺราหฺมณ\-มาณวกานํ อิรุเวทํ ยชุเวทํ สามเวทํ อถพฺพณเวทํ ลกฺขณํ อิติหาสํ ปุราณํ นิฆณฺฑุ เกฏุภํ อกฺขร\-ปฺปเภทํ ปทํ เวยฺยากรณํ ภาสมคฺคํ อุปฺปาตํ สุปินํ นิมิตฺตํ ฉฬงฺคํ \csromanfolio{178}จนฺทคฺคาหํ สูริยคฺคาหํ สุกฺก\-ราหุ\-จริตํ อุฬุคฺคห\-ยุทฺธํ\csromansharedfootnote{ea7dbc16b537f45c}{โอฬุคฺคหยุทฺธํ (ก)} เทว\-ทุนฺทุภิ\-สฺสรํ โอกฺกนฺติ อุกฺกาปาตํ ภูมิกมฺมํ\csromansharedfootnote{908e3e878eba01d6}{ภูมิกมฺปํ (สี, อิ)} ทิสาทาหํ ภุมฺม\-นฺตลิกฺขํ โชติสํ โลกายติกํ สาจกฺกํ มิคจกฺกํ อนฺตรจกฺกํ มิสฺสกุปฺปาทํ สกุณ\-รุต\-รวิตํ\csromansharedfootnote{c1587eebb14dbcfc}{สกุณรุตํ (สี)} สิกฺขา กรณียา, อวเสสานํ ปุถุเวสฺส\-สุทฺทานํ กสิ วณิชฺชา โครกฺขา กรณียา, เอวเมว โข มหาราช อกมฺมํ เหตํ ชินปุตฺตานํ ยทิทํ ปูชา, สมฺมสนํ สงฺขารานํ, โยนิโส มนสิกาโร, สติปฏฺฐานา\-นุปสฺสนา, อารมฺมณ\-สาร\-คฺคาโห, กิเลสยุทฺธํ, สทตฺถม\-นุยุญฺชนา, เอตํ ชินปุตฺตานํ กรณียํ, อวเสสานํ เทวมนุสฺสานํ ปูชา กรณียา, ตสฺมา มหาราช ตถาคโต “มา อิเม อกมฺเม ยุญฺชนฺตุ, กมฺเม อิเม ยุญฺชนฺตู”ติ อาห “อพฺยาวฏา ตุเมฺห อานนฺท โหถ ตถาคตสฺส สรีรปูชายา”ติ. ยเทตํ มหาราช ตถาคโต น ภเณยฺย, ปตฺตจีวรมฺปิ อตฺตโน ปริยาทาเปตฺวา ภิกฺขู พุทฺธปูชํ เยว กเรยฺยุนฺติ. สาธุ ภนฺเต นาคเสน, เอวเมตํ ตถา สมฺปฏิจฺฉามีติ.}{พุทฺธ\-ปูชน\-ปโญฺห สตฺตโม.}
\csromanmatikaanchor{mtk.134}
\csromantocmarkref{subsubsection}{8. ปาทสกลิกาหตปญฺห}{mtk.134}
\bookmark[dest={mtk.134},level=3]{8. ปาทสกลิกาหตปญฺห}
\csromanheader{3}{8. \textbf{ปาท}\-\textbf{สกลิกา}\-\textbf{หต}\-\textbf{ปญฺห}}
\proseitem{8}{ภนฺเต นาคเสน ตุเมฺห ภณถ “ภควโต คจฺฉนฺตสฺส อยํ อเจตนา มหาปถวี นินฺนํ อุนฺนมติ, อุนฺนตํ โอนมตี”ติ, ปุน จ ภณถ\csromansymbolfootnote{*}{วิ 4. 354 ปิฏฺเฐ.}“ภควโต ปาโท สกลิกาย ขโต”ติ. ยา สา สกลิกา ภควโต ปาเท ปติตา, กิสฺส ปน สา สกลิกา ภควโต ปาทา น นิวตฺตา. ยทิ ภนฺเต นาคเสน ภควโต คจฺฉนฺตสฺส อยํ อเจตนา มหาปถวี นินฺนํ อุนฺนมติ, อุนฺนตํ โอนมติ, เตน หิ “ภควโต ปาโท สกลิกาย ขโต”ติ ยํ วจนํ, ตํ มิจฺฉา. ยทิ ภควโต ปาโท สกลิกาย ขโต, เตน หิ “ภควโต คจฺฉนฺตสฺส อยํ อเจตนา มหาปถวี นินฺนํ อุนฺนมติ อุนฺนตํ โอนมตี”ติ ตมฺปิ วจนํ มิจฺฉา. อยมฺปิ อุภโต โกฏิโก ปโญฺห ตวานุปฺปตฺโต, โส ตยา นิพฺพาหิตพฺโพติ.}

\prose{\csromanfolio{179}สจฺจํ มหาราช อตฺเถตํ ภควโต คจฺฉนฺตสฺส อยํ อเจตนา มหาปถวี นินฺนํ อุนฺนมติ อุนฺนตํ โอนมติ, ภควโต จ ปาโท สกลิกาย ขโต, น จ ปน สา สกลิกา อตฺตโน ธมฺมตาย ปติตา, เทวทตฺตสฺส อุปกฺกเมน ปติตา, เทวทตฺโต มหาราช พหูนิ ชาติ\-สต\-สหสฺสานิ ภควติ อาฆาตํ พนฺธิ, โส เตน อาฆาเตน “มหนฺตํ กูฏาคาร\-ปฺปมาณํ ปาสาณํ ภควโต อุปริ ปาเตสฺสามี”ติ มุญฺจิ. อถ เทฺว เสลา ปถวิโต อุฏฺฐหิตฺวา ตํ ปาสาณํ สมฺปฏิจฺฉิํสุ, อถ เนสํ สมฺปหาเรน ปาสาณโต ปปฏิกา ภิชฺชิตฺวา เยน วา เตน วา ปตนฺตี ภควโต ปาเท ปติตาติ.}

\prose{ยถา จ ภนฺเต นาคเสน เทฺว เสลา ปาสาณํ สมฺปฏิจฺฉิํสุ, ตเถว ปปฏิกาปิ สมฺปฏิจฺฉิตพฺพาติ. สมฺปฏิจฺฉิตมฺปิ มหาราช อิเธกจฺจํ ปคฺฆรติ ปสวติ น ฐานมุ\-ปคจฺฉติ, ยถา มหาราช อุทกํ ปาณินา คหิตํ องฺคุล\-นฺตริกาหิ ปคฺฆรติ ปสวติ น ฐานมุ\-ปคจฺฉติ, ขีรํ ตกฺกํ มธุํ สปฺปิ เตลํ มจฺฉรสํ มํสรสํ ปาณินา คหิตํ องฺคุล\-นฺตริกาหิ ปคฺฆรติ ปสวติ น ฐานมุ\-ปคจฺฉติ, เอวเมว โข มหาราช สมฺปฏิจฺฉ\-นตฺถํ อุปคตานํ ทฺวินฺนํ เสลานํ สมฺปหาเรน ปาสาณโต ปปฏิกา ภิชฺชิตฺวา เยน วา เตน วา ปตนฺตี ภควโต ปาเท ปติตา.}

\prose{ยถา วา ปน มหาราช สณฺหสุขุม\-อณุ\-รชสมํ ปุฬินํ มุฏฺฐินา คหิตํ อางฺคุลนฺตริกาหิ ปคฺฆรติ ปสวติ น ฐานมุ\-ปคจฺฉติ, เอวเมว โข มหาราช สมฺปฏิจฺฉ\-นตฺถํ อุปคตานํ ทฺวินฺนํ เสลานํ สมฺปหาเรน ปาสาณโต ปปฏิกา ภิชฺชิตฺวา เยน วา เตน วา ปตนฺตี ภควโต ปาเท ปติตา.}

\prose{ยถา วา ปน มหาราช กพโฬ มุเขน คหิโต อิเธกจฺจสฺส มุขโต มุจฺจิตฺวา ปคฺฆรติ ปสวติ น ฐานมุ\-ปคจฺฉติ, เอวเมว โข มหาราช สมฺปฏิจฺฉ\-นตฺถํ อุปคตานํ ทฺวินฺนํ เสลานํ สมฺปหาเรน ปาสาณโต ปปฏิกา ภิชฺชิตฺวา เยน วา เตน วา ปตนฺตี ภควโต ปาเท ปติตาติ.}

\prose{โหตุ ภนฺเต นาคเสน, เสเลหิ ปาสาโณ สมฺปฏิจฺฉิโต, อถ ปปฏิกายปิ อปจิติ กาตพฺพา ยเถว มหาปถวิยาติ. ทฺวาทสิเม มหาราช อปจิติํ น กโรนฺติ, กตเม ทฺวาทส, รตฺโต ราควเสน อปจิติํ น กโรติ, ทุฏฺโฐ โทสวเสน, มูโฬฺห โมหวเสน}

\prose{\csromanfolio{180}อุนฺนโต มานวเสน, นิคฺคุโณ อวิเสสตาย, อติถทฺโธ อนิเสธนตาย, หีโน หีนสภา\-วตาย, วจนกโร อนิสฺสรตาย, ปาโป กทริยตาย, ทุกฺขาปิโต ปฏิทุกฺขาปนตาย, ลุทฺโธ โลภา\-ภิภู\-ตตาย, อายูหิโต อตฺถ\-สาธน\-ตาย\csromansharedfootnote{4ab8cdbf40f973de}{อตฺถสาธเนน (สฺยา, อิ, ก)} อปจิติํ น กโรติ, อิเม โข มหาราช ทฺวาทส อปจิติํ น กโรนฺติ. สา จ ปน ปปฏิกา ปาสาณ\-สมฺปหาเรน ภิชฺชิตฺวา อนิมิตฺตกต\-ทิสา เยน วา เตน วา ปตมานา ภควโต ปาเท ปติตา.}

\prose{ยถา วา ปน มหาราช สณฺหสุขุม\-อณุ\-รโช อนิล\-พล\-สมาหโต อนิมิตฺตกตทิโส เยน วา เตน วา อภิกิรติ, เอวเมว โข มหาราช สา ปปฏิกา ปาสาณ\-สมฺปหาเรน ภิชฺชิตฺวา อนิมิตฺตกต\-ทิสา เยน วา เตน วา ปตมานา ภควโต ปาเท ปติตา. ยทิ ปน มหาราช สา ปปฏิกา ปาสาณโต วิสุํ น ภเวยฺย, ตมฺปิ เต เสลา ปาสาณ\-ปปฏิกํ อุปฺปติตฺวา คเณฺหยฺยุํ. เอสา ปน มหาราช ปปฏิกา น ภูมฏฺฐา น อากาสฏฺฐา, ปาสาณ\-สมฺปหาร\-เวเคน ภิชฺชิตฺวา อนิมิตฺตกต\-ทิสา เยน วา เตน วา ปตมานา ภควโต ปาเท ปติตา.}

\prosewithcloserruled{ยถา วา ปน มหาราช วาตมณฺฏลิกาย อุกฺขิตฺตํ ปุราณปณฺณํ อนิมิตฺตกตทิสํ เยน วา เตน วา ปตติ, เอวเมว โข มหาราช เอสา ปปฏิกา ปาสาณ\-สมฺปหาร\-เวเคน อนิมิตฺตกต\-ทิสา เยน วา เตน วา ปตมานา ภควโต ปาเท ปติตา. อปิ จ มหาราช อกตญฺญุสฺส กทริยสฺส เทวทตฺตสฺส ทุกฺขา\-นุภวนาย ปปฏิกา ภควโต ปาเท ปติตาติ. สาธุ ภนฺโต นาคเสน, เอวเมตํ ตถา สมฺปฏิจฺฉามีติ.}{ปาท\-สกลิกา\-หต\-ปโญฺห อฏฺฐโม.}
\csromanmatikaanchor{mtk.135}
\csromantocmarkref{subsubsection}{9. อคฺคคฺคสมณปญฺห}{mtk.135}
\bookmark[dest={mtk.135},level=3]{9. อคฺคคฺคสมณปญฺห}
\csromanheader{3}{9. \textbf{อคฺคคฺค}\-\textbf{สมณ}\-\textbf{ปญฺห}}
\proseitem{9}{ภนฺเต นาคเสน ภาสิตมฺเปตํ ภควตา\csromansymbolfootnote{*}{ม 1. 353 ปิฏฺเฐ.}“อาสวานํ ขยา สมโณ โหตี”ติ. ปุน จ ภณิตํ\csromandash{}}

\setlength{\csromangathapagewidth}{0pt}
\setlength{\csromangathawakwidth}{0pt}
\csromangathameasuregroup{}{\textsuperscript{+}“จตุพฺภิ ธมฺเมหิ สมงฺคิภูตํ, \\ ตํ เว นรํ สมณํ อาหุ โลเก”ติ.}
\csromangathameasurewak{\textsuperscript{+}“จตุพฺภิ ธมฺเมหิ สมงฺคิภูตํ, \\ ตํ เว นรํ สมณํ อาหุ โลเก”ติ.}
\setlength{\csromangathaleftcol}{0pt}
\csromangathagroup{\csromangathastanza{\csromangathawak{\csromansymbolfootnote{+}{ขุ 5. 204 ปิฏฺเฐ ชาตเก.}“จตุพฺภิ ธมฺเมหิ สมงฺคิภูตํ, \\ ตํ เว นรํ สมณํ อาหุ โลเก”ติ.}}}

\prose{\csromanfolio{181}ตตฺริเม จตฺตาโร ธมฺมา ขนฺติ อปฺปาหารตา รติวิปฺปหานํ อากิญฺจญฺญํ. สพฺพานิ ปเนตานิ อปริกฺขีณา\-สวสฺส สกิเลสสฺเสว โหนฺติ. ยทิ ภนฺเต นาคเสน อาสวานํ ขยา สมโณ โหติ, เตน หิ “จตุพฺภิ ธมฺเมหิ สมงฺคิภูตํ, ตํ เว นรํ สมณํ อาหุ โลเก”ติ ยํ วจนํ ตํ มิจฺฉา. ยทิ จตุพฺภิ ธมฺเมหิ สมงฺคิภูโต สมโณ โหติ, เตน หิ “อาสวานํ ขยา สมโณ โหตี”ติ ตมฺปิ วจนํ มิจฺฉา, อยมฺปิ อุภโต โกฏิโก ปโญฺห ตวานุปฺปตฺโต, โส ตยา นิพฺพาหิตพฺโพติ.}

\prose{ภาสิตมฺเปตํ มหาราช ภควตา “อาสวานํ ขยา สมโณ โหตี”ติ. ปุน จ ภณิตํ “จตุพฺภิ ธมฺเมหิ สมงฺคิภูตํ, ตํ เว นรํ สมณํ อาหุ โลเก”ติ. ตทิทํ มหาราช วจนํ เตสํ เตสํ ปุคฺคลานํ คุณวเสน ภณิตํ “จตุพฺภิ ธมฺเมหิ สมงฺคิภูตํ, ตํ เว นรํ สมณํ อาหุ โลเก”ติ, อิทํ ปน นิรวเสส\-วจนํ “อาสวานํ ขยา สมโณ โหตี”ติ.}

\prose{อปิ จ มหาราช เย เกจิ กิเลสูปสมาย ปฏิปนฺนา, เต สพฺเพ อุปาทายุปาทาย สมโณ ขีณาสโว อคฺคมกฺขายติ. ยถา มหาราช ยานิ กานิจิ ชลช\-ถลช\-ปุปฺผานิ, วสฺสิกํ เตสํ อคฺคมกฺขายติ, อวเสสานิ ยานิ กานิจิ วิวิธานิ ปุปฺผชาตานิ, สพฺพานิ ตานิ ปุปฺผานิ เยว, อุปาทายุปาทาย ปน วสฺสิกํ เยว ปุปฺผํ ชนสฺส ปตฺถิตํ ปิหยิตํ. เอวเมว โข มหาราช เย เกจิ กิเลสูปสมาย ปฏิปนฺนา, เต สพฺเพ อุปาทายุปาทาย สมโณ ขีณาสโว อคฺคมกฺขายติ.}

\prose{ยถา วา ปน มหาราช สพฺพธญฺญานํ สาลิ อคฺคมกฺขายติ, ยา กาจิ อวเสสา วิวิธา ธญฺญชาติโย, ตา สพฺพา อุปาทายุปาทาย โภชนานิ สรีรยาปนาย, สาลิ เยว เตสํ อคฺคมกฺขายติ. เอวเมว โข มหาราช เย เกจิ กิเลสูปสมาย ปฏิปนฺนา, เต สพฺเพ อุปาทายุปาทาย สมโณ ขีณาสโว อคฺคมกฺขายตีติ. สาธุ ภนฺเต นาคเสน, เอวเมตํ ตถา สมฺปฏิจฺฉามีติ.}

\vspace{\tipitakaparskip}
\csromancenter{อคฺคคฺค\-สมณ\-ปโญฺห นวโม.}
\csromanmatikaanchor{mtk.136}
\csromantocmarkref{subsubsection}{10. วณฺณภณนปญฺห}{mtk.136}
\bookmark[dest={mtk.136},level=3]{10. วณฺณภณนปญฺห}
\csromanheader{3}{10. \textbf{วณฺณภณน}\-\textbf{ปญฺห}}
\proseitem{10}{\csromanfolio{182}ภนฺเต นาคเสน ภาสิตมฺเปตํ ภควตา\csromansymbolfootnote{*}{ที 1. 3 ปิฏฺเฐ.}“มมํ วา ภิกฺขเว ปเร วณฺณํ ภาเสยฺยุํ, ธมฺมสฺส วา, สํฆสฺส วา วณฺณํ ภาเสยฺยุํ, ตตฺร ตุเมฺหหิ น อานนฺโท, น โสมนสฺสํ, น เจตโส อุปฺปิลาวิตตฺตํ กรณียนฺ”ติ, ปุน จ ตถาคโต เสลสฺส พฺราหฺมณสฺส ยถาภุจฺเจ วณฺเณ ภญฺญมาเน อานนฺทิโต สุมโน อุปฺปิลาวิโต ภิยฺโย อุตฺตริํ สกคุณํ ปกิตฺเตสิ\csromandash{}}

\setlength{\csromangathapagewidth}{0pt}
\setlength{\csromangathawakwidth}{0pt}
\csromangathameasuregroup{“ราชาหมสฺมิ เสลาติ, \\ ธมฺเมน จกฺกํ วตฺเตมิ,}{\csromangathabat{“ราชาหมสฺมิ เสลาติ,}{ธมฺมราชา อนุตฺตโร.} \\ \csromangathabat{ธมฺเมน จกฺกํ วตฺเตมิ,}{จกฺกํ อปฺปฏิวตฺติยนฺ”ติ\textsuperscript{1}.}}
\csromangathasetleft{“ราชาหมสฺมิ เสลาติ, \\ ธมฺเมน จกฺกํ วตฺเตมิ,}
\csromangathagroup{\csromangathastanza{\csromangathabat{“ราชาหมสฺมิ เสลาติ,}{ธมฺมราชา อนุตฺตโร.} \\ \csromangathabat{ธมฺเมน จกฺกํ วตฺเตมิ,}{จกฺกํ อปฺปฏิวตฺติยนฺ”ติ\csromansharedfootnote{d075cdc6c10be333}{ม 2. 352; ขุ 1. 368 ปิฏฺเฐสุ.}.}}}

\prose{ยทิ ภนฺเต นาคเสน ภควตา ภณิตํ “มมํ วา ภิกฺขเว ปเร วณฺณํ ภาเสยฺยุํ, ธมฺมสฺส วา สํฆสฺส วา วณฺณํ ภาเสยฺยุํ, ตตฺร ตุเมฺหหิ น อานนฺโท, น โสมนสฺสํ, น เจตโส อุปฺปิลาวิตตฺตํ กรณียนฺ”ติ, เตน หิ เสลสฺส พฺราหฺมณสฺส ยถาภุจฺเจ วณฺเณ ภญฺญมาเน อานนฺทิโต สุมโน อุปฺปิลาวิโต ภิยฺโย อุตฺตริํ สกคุณํ ปกิตฺเตสีติ ยํ วจนํ, ตํ มิจฺฉา. ยทิ เสลสฺส พฺราหฺมณสฺส ยถาภุจฺเจ วณฺเณ ภญฺญมาเน อานนฺทิโต สุมโน อุปฺปิลาวิโต ภิยฺโย อุตฺตริํ สกคุณํ ปกิตฺเตสิ, เตน หิ “มมํ วา ภิกฺขเว ปเร วณฺณํ ภาเสยฺยุํ, ธมฺมสฺส วา สํฆสฺส วา วณฺณํ ภาเสยฺยุํ, ตตฺร ตุเมฺหหิ น อานนฺโท, น โสมนสฺสํ, น เจตโส อุปฺปิลาวิตตฺตํ กรณียนฺ”ติ ตมฺปิ วจนํ มิจฺฉา, อยมฺปิ อุภโต โกฏิโก ปโญฺห ตวานุปฺปตฺโต, โส ตยา นิพฺพาหิตพฺโพติ.}

\prose{ภาสิตมฺเปตํ มหาราช ภควตา “มมํ วา ภิกฺขเว ปเร วณฺณํ ภาเสยฺยุํ, ธมฺมสฺส วา สํฆสฺส วา วณฺณํ ภาเสยฺยุํ, ตตฺร ตุเมฺหหิ น อานนฺโท, น โสมนสฺสํ, น เจตโส อุปฺปิลาวิตตฺตํ กรณียนฺ”ติ. เสลสฺส จ พฺราหฺมณสฺส ยถาภุจฺเจ วณฺเณ ภญฺญมาเน ภิยฺโย อุตฺตริํ สกคุณํ ปกิตฺติตํ\csromandash{}}

\setlength{\csromangathapagewidth}{0pt}
\setlength{\csromangathawakwidth}{0pt}
\csromangathameasuregroup{“ราชาหมสฺมิ เสลาติ, \\ ธมฺเมน จกฺกํ วตฺเตมิ,}{\csromangathabat{“ราชาหมสฺมิ เสลาติ,}{ธมฺมราชา อนุตฺตโร.} \\ \csromangathabat{ธมฺเมน จกฺกํ วตฺเตมิ,}{จกฺกํ อปฺปฏิวตฺติยนฺ”ติ.}}
\csromangathasetleft{“ราชาหมสฺมิ เสลาติ, \\ ธมฺเมน จกฺกํ วตฺเตมิ,}
\csromangathagroup{\csromangathastanza{\csromangathabat{“ราชาหมสฺมิ เสลาติ,}{ธมฺมราชา อนุตฺตโร.} \\ \csromangathabat{ธมฺเมน จกฺกํ วตฺเตมิ,}{จกฺกํ อปฺปฏิวตฺติยนฺ”ติ.}}}

\prosewithcloserruled{ปฐมํ มหาราช ภควตา ธมฺมสฺส สภาว\-สรสลกฺขณํ สภาวํ อวิตถํ ภูตํ ตจฺฉํ ตถตฺถํ, ปริทีปย\-มาเนน ภณิตํ “มมํ วา ภิกฺขเว \csromanfolio{183}ปเร วณฺณํ ภาเสยฺยุํ, ธมฺมสฺส วา สํฆสฺส วา วณฺณํ ภาเสยฺยุํ, ตตฺร ตุเมฺหหิ น อานนฺโท, น โสมนสฺสํ, นเจตโส อุปฺปิลาวิตตฺตํ กรณียนฺ”ติ. ยํ ปน ภควตา เสลสฺส พฺราหฺมณสฺส ยถาภุจฺเจ วณฺเณ ภญฺญมาเน ภิยฺโย อุตฺตริํ สกคุณํ ปกิตฺติตํ “ราชาหมสฺมิ เสลาติ, ธมฺมราชา อนุตฺตโร”ติ ตํ น ลาภเหตุ, น ยสเหตุ, น อตฺตเหตุ, น ปกฺขเหตุ, น อนฺเตวาสิ\-กมฺยตาย, อถ โข อนุกมฺปาย การุญฺเญน หิตวเสน เอวํ อิมสฺส ธมฺมา\-ภิสมโย ภิวิสฺสติ ติณฺณญฺจ มาณวก\-สตานนฺติ, เอวํ ภิยฺโย อุตฺตริํ สกคุณํ ภณิตํ “ราชาหมสฺมิ เสลาติ, ธมฺมราชา อนุตฺตโร”ติ. สาธุ ภนฺเต นาคเสน, เอวเมตํ ตถา สมฺปฏิจฺฉามีติ.}{วณฺณภณน\-ปโญฺห ทสโม.}
\csromanmatikaanchor{mtk.137}
\csromantocmarkref{subsubsection}{11. อหิํสานิคฺคหปญฺห}{mtk.137}
\bookmark[dest={mtk.137},level=3]{11. อหิํสานิคฺคหปญฺห}
\csromanheader{3}{11. \textbf{อหิํสา}\-\textbf{นิคฺคห}\-\textbf{ปญฺห}}
\proseitem{11}{ภนฺเต นาคเสน ภาสิตมฺเปตํ ภควตา\csromansymbolfootnote{*}{ขุ 5. 220 ปิฏฺเฐ ชาตเก.}“อหิํสาย จร โลเก, ปิโย โหหิสิ มํมิวา”ติ. ปุน จ ภณิตํ\csromansymbolfootnote{+}{ขุ 6. 3 ปิฏฺเฐ ชาตเก.}“นิคฺคเณฺห นิคฺคหารหํ, ปคฺคเณฺห ปคฺคหารหนฺ”ติ. นิคฺคโห นาม ภนฺเต นาคเสน หตฺถจฺเฉโท ปาทจฺเฉโท วโธ พนฺทนํ การณา มารณํ สนฺตติ\-วิโกปนํ, น เอตํ วจนํ ภควโต ยุตฺตํ, น จ ภควา อรหติ เอตํ วจนํ วตฺถุํ. ยทิ ภนฺเต นาคเสน ภควตา ภณิตํ “อหิํสยํ ปรํ โลเก, ปิโย โหหิสิ มามโก”ติ, เตน หิ “นิคฺคเณฺห นิคฺคหารหํ, ปคฺคเณฺห ปคฺคหารหนฺ”ติ ยํ วจนํ, ตํ มิจฺฉา. ยทิ ตถาคเตน ภณิตํ “นิคฺคเณฺห นิคฺคหารหํ, ปคฺคเณฺห ปคฺคหารหนฺ”ติ, เตน หิ “อหิํสยํ ปรํ โลเก, ปิโย โหหิสิ มามโก”ติ ตมฺปิ วจนํ มิจฺฉา. อยมฺปิ อุภโต โกฏิโก ปโญฺห ตวานุปฺปตฺโต, โส ตยา นิพฺพาหิตพฺโพติ.}

\prose{ภาสิตมฺเปตํ มหาราช ภควตา “อหิํสยํ ปรํ โลเก, ปิโย โหหิสิ มามโก”ติ, ภณิตญฺจ “นิคฺคเณฺห นิคฺคหารหํ, ปคฺคเณฺห \csromanfolio{184}ปคฺคหารหนฺ”ติ. “อหิํสยํ ปรํ โลเก, ปิโย โหหิสิ มามโก”ติ สพฺเพสํ มหาราช ตถาคตานํ อนุมตํ เอตํ, เอสา อนุสิฏฺฐิ, เอสา ธมฺมเทสนา, ธมฺโม หิ มหาราช อหิํสาลกฺขโณ, สภาววจนํ เอตํ. ยํ ปน มหาราช ตถาคโต อาห “นิคฺคเณฺห นิคฺคหารหํ, ปคฺคเณฺห ปคฺคหารหนฺ”ติ, ภาสา เอสา, อุทฺธตํ มหาราช จิตฺตํ นิคฺคเหตพฺพํ, ลีนํ จิตฺตํ ปคฺคเหตพฺพํ. อกุสลํ จิตฺตํ นิคฺคเหตพฺพํ, กุสลํ จิตฺตํ ปคฺคเหตพฺพํ. อโยนิโส มนสิกาโร นิคฺคเหตพฺโพ, โยนิโส มนสิกาโร ปคฺคเหตพฺโพ. มิจฺฉา\-ปฏิปนฺโน นิคฺคเหตพฺโพ, สมฺมาปฏิปนฺโน ปคฺคเหตพฺโพ. อนริโย นิคฺคเหตพฺโพ อริโย ปคฺคเหตพฺโพ. โจโร นิคฺคเหตพฺโพ, อโจโร ปคฺคเหตพฺโพติ.}

\prose{โหตุ ภนฺเต นาคเสน, อิทานิ ตฺวํ ปจฺจาคโตสิ มม วิสยํ, ยมหํ ปุจฺฉามิ, โส เม อตฺโถ อุปคโต. โจโร ปน ภนฺเต นาคเสน นิคฺคณฺหนฺเตน กถํ นิคฺคเหตพฺโพติ. โจโร มหาราช นิคฺคณฺหนฺเตน เอวํ นิคฺคเหตพฺโพ, ปริภาสนีโย ปริภาสิตพฺโพ, ทณฺฑนีโย ทณฺเฑตพฺโพ, ปพฺพาชนีโย ปพฺพาเชตพฺโพ, พนฺธนีโย พนฺธิตพฺโพ, ฆาตนีโย ฆาเตตพฺโพติ. ยํ ปน ภนฺเต นาคเสน โจรานํ ฆาตนํ, ตํ ตถาคตานํ อนุมตนฺติ. น หิ มหาราชาติ. กิสฺส ปน โจโร อนุสาสนีโย อนุมโต ตถาคตานนฺติ. โย โส มหาราช ฆาตียติ, น โส ตถาคตานํ อนุมติยา ฆาตียติ, สยํกเตน โส ฆาตียติ, อปิ จ ธมฺมา\-นุสิฏฺฐิยา อนุสาสียติ, สกฺกา ปน มหาราช ตยา ปุริสํ อการกํ อนปราธํ วีถิยํ จรนฺตํ คเหตฺวา ฆาตยิตุนฺติ. น สกฺกา ภนฺเตติ. เกน การเณน มหาราชาติ. อการกตฺตา ภนฺเตติ. เอวเมว โข มหาราช น โจโร ตถาคตานํ อนุมติยา หญฺญติ, สยํกเตน โส หญฺญติ, กิํ ปเนตฺถ อนุสาสโก กิญฺจิ โทสํ อาปชฺชตีติ. น หิ ภนฺเตติ. เตน หิ มหาราช ตถาคตานํ อนุสิฏฺฐิ สมฺมานุสิฏฺฐิ โหตีติ. สาธุ ภนฺเต นาคเสน เอวเมตํ ตถา สมฺปฏิจฺฉามีติ.}

\vspace{\tipitakaparskip}
\csromancenter{อหิํสา\-นิคฺคห\-ปโญฺห เอกาทสโม.}
\csromanmatikaanchor{mtk.138}
\csromantocmarkref{subsubsection}{12. ภิกฺขุปณามิตปญฺห}{mtk.138}
\bookmark[dest={mtk.138},level=3]{12. ภิกฺขุปณามิตปญฺห}
\csromanmarkh{4. เมณฺฑกปญฺห}\csromanmarkhii{12. ภิกฺขุ\-ปณามิต\-ปญฺห}\csromanheadernomark{3}{4. เมณฺฑกปญฺห \textbf{12. ภิกฺขุ}\-\textbf{ปณามิต}\-\textbf{ปญฺห}}
\proseitem{12}{\csromanfolio{185}ภนฺเต นาคเสน ภาสิตมฺเปตํ ภควตา\csromansymbolfootnote{*}{ขุ 1. 282 ปิฏฺเฐ สุตฺตนิปาเต.}“อกฺโกธโน วิคตขิโลหมสฺมี”ติ, ปุน จ ตถาคโต เถเร สาริปุตฺต\-โมคฺคลฺลาเน\csromansymbolfootnote{+}{ม 2. 120 ปิฏฺเฐ.}สปริเส ปณาเมสิ, กิํ นุ โข ภนฺเต นาคเสน ตถาคโต กุปิโต ปริสํ ปณาเมสิ, อุทาหุ ตุฏฺโฐ ปณาเมสิ, เอตํ ตาว ชานาหิ อิมํ นามาติ. ยทิ ภนฺเต นาคเสน กุปิโต ปริสํ ปณาเมสิ, เตน หิ ตถาคตสฺส โกโธ อปฺปฏิวตฺติโต, ยทิ ตุฏฺโฐ ปณาเมสิ, เตน หิ อวตฺถุสฺมิํ อชานนฺเตน ปณามิตา, อยมฺปิ อุภโต โกฏิโก ปโญฺห ตวานุปฺปตฺโต, โส ตยา นิพฺพาหิตพฺโพติ.}

\prose{ภาสิตมฺเปตํ มหาราช ภควตา “อกฺโกธโน วิคตขิโลหมสฺมี”ติ, ปณามิตา จ เถรา สาริปุตฺต\-โมคฺคลฺลานา สปริสา, ตญฺจ ปน น โกเปน. อิธ มหาราช โกจิเทว ปุริโส มหาปถวิยา มูเล วา ขาณุเก วา ปาสาเณ วา กฐเล วา วิสเม วา ภูมิภาเค ขลิตฺวา ปตติ, อปิ นุ โข มหาราช มหาปถวี กุปิตา ตํ ปาเตตีติ. น หิ ภนฺเต, นตฺถิ มหาปถวิยา โกโป วา ปสาโท วา, อนุนย\-ปฺปฏิฆ\-วิปฺปมุตฺตา มหาปถวี, สยเมว โส อลโส ขลิตฺวา ปติโตติ, เอวเมว โข มหาราช นตฺถิ ตถาคตานํ โกโป วา ปสาโท วา, อนุนย\-ปฺปฏิฆ\-วิปฺปมุตฺตา ตถาคตา อรหนฺโต สมฺมาสมฺพุทฺธา, อถ โข สยํ กเตเนว เต อตฺตโน อปราเธน ปณามิตา.}

\prose{อิธ ปน มหาราช มหาสมุทฺโท น มเตน กุณเปน สํวสติ, ยํ โหติ มหาสมุทฺเท มตํ กุณปํ, ตํ ขิปฺปเมว นิจฺฉุภติ ถลํ อุสฺสาเรติ, อปิ นุ โข มหาราช มหาสมุทฺโท กุปิโต ตํ กุณปํ นิจฺฉุภตีติ. น หิ ภนฺเต, นตฺถิ มหาสมุทฺทสฺส โกโป วา ปสาโท วา, อนุนย\-ปฺปฏิฆ\-วิปฺปมุตฺโต มหาสมุทฺโทติ. เอวเมว โข มหาราช นตฺถิ ตถาคตานํ โกโป วา ปสาโท วา, อนุนย\-ปฺปฏิฆ\-วิปฺปมุตฺตา ตถาคตา อรหนฺโต สมฺมาสมฺพุทฺธา, อถ โข สยํ กเตเนว เต อตฺตโน อปราเธน ปณามิตา.}

\prose{ยถา มหาราช ปถวิยา ขลิโต ปตียติ, เอวํ ชินสาสนวเร ขลิโต ปณามียติ. ยถา มหาราช สมุทฺเท มตํ กุณปํ \csromanfolio{186}นิจฺฉุภียติ, เอวํ ชินสาสนวเร ขลิโต ปณามียติ. ยํ ปน เต มหาราช ตถาคโต ปณาเมสิ, เตสํ อตฺถกาโม หิตกาโม สุขกาโม วิสุทฺธิกาโม “เอวํ อิเม ชาติ\-ชรา\-พฺยาธิ\-มรเณหิ ปริมุจฺจิสฺสนฺตี”ติ ปณาเมสีติ. สาธุ ภนฺเต นาคเสน, เอวเมตํ ตถา สมฺปฏิจฺฉามีติ.}

\vspace{\tipitakaparskip}
\csromancenter{ภิกฺขุ\-ปณามิต\-ปโญฺห ทฺวาทสโม.}
\nitthitammidruled{\textbf{ปณามิตวคฺโค ตติโย.}}
\csromancenter{อิมสฺมิํ วคฺเค ทฺวาทสฺส ปญฺหา.}
\csromanmatikaanchor{mtk.139}
\csromantocmarknopage{subsection}{4. สพฺพญฺญุตญาณวคฺค}{mtk.139}
\bookmark[dest={mtk.139},level=2]{4. สพฺพญฺญุตญาณวคฺค}
\csromanheader{2}{4. \textbf{สพฺพญฺญุตญาณ}\-\textbf{วคฺค}}
\csromanmatikaanchor{mtk.140}
\csromantocmarkref{subsubsection}{1. อิทฺธิกมฺมวิปากปญฺห}{mtk.140}
\bookmark[dest={mtk.140},level=3]{1. อิทฺธิกมฺมวิปากปญฺห}
\csromanheader{3}{1. \textbf{อิทฺธิ}\-\textbf{กมฺมวิปาก}\-\textbf{ปญฺห}}
\proseitem{1}{\csromanfolio{187}ภนฺเต นาคเสน ภาสิตมฺเปตํ ภควตา\csromansymbolfootnote{*}{อํ 1. 23 ปิฏฺเฐ.}“เอตทคฺคํ ภิกฺขเว มม สาวกานํ ภิกฺขูนํ อิทฺธิมนฺตานํ ยทิทํ มหาโมคฺคลฺลาโน”ติ. ปุน จ กิร โส ลคุเฬหิ ปริโปถิโต ภินฺนสีโส สญฺจุณฺณิตฏฺฐิมํสธมนิฉินฺนปริคตฺโต\csromansharedfootnote{53b8773699ec4376}{ธมนิมชฺชปริกตฺโต (สี, อิ), ธมฺมนิมิญฺชปริคตฺโต (สฺยา)} ปรินิพฺพุโต. ยทิ ภนฺเต นาคเสน เถโร มหาโมคฺคลฺลาโน อิทฺธิยา โกฏิํ คโต, เตน หิ ลคุเฬหิ โปถิโต ปรินิพฺพุโตติ ยํ วจนํ, ตํ มิจฺฉา. ยทิ ลคุเฬหิ ปริโปถิโต ปรินิพฺพุโต, เตน หิ อิทฺธิยา โกฏิํ คโตติ ตมฺปิ วจนํ มิจฺฉา. กิํ น สมตฺโถ อิทฺธิยา อตฺตโน อุปฆาตํ อปนยิตุํ, สเทวกสฺสปิ โลกสฺส ปฏิสรณํ ภวิตุํ อรโหติ. อยมฺปิ อุภโต โกฏิโก ปโญฺห ตวานุปฺปตฺโต, โส ตยา นิพฺพาหิตพฺโพติ.}

\prose{ภาสิตมฺเปตํ มหาราช ภควตา “เอตทคฺคํ ภิกฺขเว มม สาวกานํ ภิกฺขุนํ อิทฺธิมนฺตานํ ยทิทํ มหาโนคฺคลฺลาโน”ติ. อายสฺมา จ มหาโมคฺคลฺลาโน ลคุฬหโต ปรินิพฺพุโต, ตญฺจ ปน กมฺมาธิคฺคหิเตนาติ.}

\prose{นนุ ภนฺเต นาคเสน อิทฺธิมโต อิทฺธิวิสโยปิ กมฺมวิปาโกปิ เทฺว อจินฺติยา, อจินฺติเยน อจินฺติยํ อปนยิตพฺพํ. ยถา นาม ภนฺเต เกจิ ผลกามา กปิตฺเถน กปิตฺถํ โปเถนฺติ, อมฺเพน อมฺพํ โปเถนฺติ, เอวเมว โข ภนฺเต นาคเสน อจินฺติเยน อจินฺติยํ โปถยิตฺวา อปเนตพฺพนฺติ. อจินฺติยานมฺปิ มหาราช เอกํ อธิมตฺตํ พลวตรํ, ยถา มหาราช มหิยา ราชาโน โหนฺติ สมชจฺจา, สมชจฺจานมฺปิ เตสํ เอโก สพฺเพ อภิภวิตฺวา อาณํ ปวตฺเตติ, เอวเมว โข มหาราช เตสํ อจินฺติยานํ กมฺมวิปากํ เยว อธิมตฺตํ พลวตรํ, กมฺมวิปากํ เยว สพฺเพ อภิภวิย อาณํ ปวตฺเตติ, กมฺมา\-ธิคฺคหิตสฺส อวเสสา กิริยา โอกาสํ น ลภนฺติ.}

\prose{\csromanfolio{188}อิธ ปน มหาราช โกจิ ปุริโส กิสฺมิญฺจิเทว ปกรเณ อปรชฺฌติ, น ตสฺส มาตา วา ปิตา วา ภคินี วา ภาตโร วา สขี วา สหายกา วา\csromansharedfootnote{1066cb9e72826a9f}{ภคินิภาตโร วา สขิสหายกา วา (สี, อิ, ก)} ตายนฺติ, อถ โข ราชา เยว ตตฺถ อภิภวิย อาณํ ปวตฺเตติ, กิํ ตตฺถ การณํ, อปราธิกตา. เอวเมว โข มหาราช เตสํ อจินฺติยานํ กมฺมวิปากํ เยว อธิมตฺตํ พลวตรํ, กมฺมวิปากํ เยว สพฺเพ อภิภวิย อาณํ ปวตฺเตติ, กมฺมา\-ธิคฺคหิตสฺส อวเสสา กิริยา โอกาสํ น ลภนฺติ.}

\prosewithcloserruled{ยถา วา ปน มหาราช มหิยา ทวฑาเห สมุฏฺฐิเต ฆฏสหสฺสมฺปิ อุทกํ น สกฺโกติ นิพฺพาเปตุํ, อถ โข อคฺคิ เยว ตตฺถ อภิภวิย อาณํ ปวตฺเตติ, กิํ ตตฺถ การณํ, พลวตา เตชสฺส. เอวเมว โข มหาราช เตสํ อจินฺนิยานํ กมฺมวิปากํ เยว อธิมตฺตํ พลวตรํ, กมฺมวิปากํ เยว สพฺเพ อภิภวิย อาณํ ปวตฺเตติ, กมฺมา\-ธิคฺคหิตสฺส อวเสสา กิริยา โอกาสํ น ลภนฺติ, ตสฺมา มหาราช อายสฺมโต มหา\-โมคฺคลานสฺส กมฺมา\-ธิคฺคหิตสฺส ลคุเฬหิ โปถิยมานสฺส อิทฺธิยา สมนฺนาหาโร นาโหสีติ. สาธุ ภนฺเต นาคเสน, เอวเมตํ ตถา สมฺปฏิจฺฉามีติ.}{อิทฺธิ\-กมฺมวิปาก\-ปโญฺห ปฐโม.}
\csromanmatikaanchor{mtk.141}
\csromantocmarkref{subsubsection}{2. ธมฺมวินยปฏิจฺฉนฺนาปฏิจฺฉนฺนปญฺห}{mtk.141}
\bookmark[dest={mtk.141},level=3]{2. ธมฺมวินยปฏิจฺฉนฺนาปฏิจฺฉนฺนปญฺห}
\csromanheader{3}{2. \textbf{ธมฺมวินย}\-\textbf{ปฏิจฺฉนฺนา}\-\textbf{ปฏิจฺฉนฺน}\-\textbf{ปญฺห}}
\proseitem{2}{ภนฺเต นาคเสน ภาสิตมฺเปตํ ภควตา\csromansymbolfootnote{*}{วิ 5. 218; อํ 1. 216 ปิฏฺเฐ.}“ตถาคต\-ปฺปเวทิโต ภิกฺขเว ธมฺมวินโย วิวโฏ วิโรจติ โน ปฏิจฺฉนฺโน”ติ.\csromansymbolfootnote{+}{วิ 3. 155 ปิฏฺเฐ.}ปุน จ ปาติโมกฺขุ\-ทฺเทโส เกวลญฺจ วินยปิฏกํ ปิหิตํ ปฏิจฺฉนฺนํ. ยทิ ภนฺเต นาคเสน ชินสาสเน ยุตฺตํ วา ปตฺตํ วา สมยํ ลเภถ, วินยปณฺณตฺถิ วิวฏา โสเภยฺย, เกน การเณน, เกวลํ ตตฺถ สิกฺขา สํยโม นิยโม สีลคุณ\-อาจาร\-ปณฺณตฺติ อตฺถรโส ธมฺมรโส วิมุตฺติรโส. ยทิ ภนฺเต นาคเสน ภควตา ภณิตํ “ตถาคต\-ปฺปเวทิโต ภิกฺขเว ธมฺมวินโย วิวโฏ วิโรจติ โน ปฏิจฺฉนฺโน”ติ, เตน หิ “ปาติโมกฺขุ\-ทฺเทโส เกวลญฺจ วินยปิฏกํ ปิหิตํ ปฏิจฺฉนฺนนฺ”ติ ยํ วจนํ, ตํ มิจฺฉา. ยทิ ปาติโมกฺขุ\-ทฺเทโส เกวลญฺจ วินยปิฏกํ ปิหิตํ ปฏิจฺฉนฺนํ, เตน หิ \csromanfolio{189}“ตถาคต\-ปฺปเวทิโต ภิกฺขเว ธมฺมวินโย วิวโฏ วิโรจติ โน ปฏิจฺฉนฺโน”ติ ตมฺปิ วจนํ มิจฺฉา. อยมฺปิ อุภโต โกฏิโก ปโญฺห ตวานุปฺปตฺโต, โส ตยา นิพฺพาหิตพฺโพติ.}

\prose{ภาสิตมฺเปตํ มหาราช ภควตา “ตถาคต\-ปฺปเวทิโต ภิกฺขเว ธมฺมวินโย วิวโฏ วิโรจติ โน ปฏิจฺฉนฺโน”ติ. ปุน จ ปาติโมกฺขุ\-ทฺเทโส เกวลญฺจ วินยปิฏกํ ปิหิตํ ปฏิจฺฉนฺนํ, ตญฺจ ปน น สพฺเพสํ, สีมํ กตฺวา ปิหิตํ.}

\prose{ติวิเธน มหาราช ภควตา ปาติโมกฺขุ\-ทฺเทโส สีมํ กตฺวา ปิหิโต, ปุพฺพกานํ ตถาคตานํ วํสวเสน ปาติโมกฺขุ\-ทฺเทโส สีมํ กตฺวา ปิหิโต, ธมฺมสฺส ครุกตฺตา ปิหิโต, ภิกฺขุภูมิยา ครุกตฺตา ปิหิโต.}

\prose{กถํ ปุพฺพกานํ ตถาคตานํ วํสเวเสน ปาติโมกฺขุ\-ทฺเทโส สีมํ กตฺวา ปิหิโต, เอโส วํโส มหาราช สพฺเพสํ ปุพฺพกานํ ตถาคตานํ ยทิทํ ภิกฺขุมชฺเฌ ปาติโมกฺขุ\-ทฺเทโส อวเสสานํ ปิหิโต. ยถา มหาราช ขตฺติยานํ ขตฺติยมายา ขตฺติเยสุ เยว จรติ, เอวเมตํ ขตฺติยานํ โลกสฺส ปเวณี อวเสสานํ ปิหิตา. เอวเมว โข มหาราช เอโส วํโส สพฺเพสํ ปุพฺพกานํ ตถาคตานํ ยทิทํ ภิกฺขุมชฺเฌ ปาติโมกฺขุ\-ทฺเทโส อวเสสานํ ปิหิโต.}

\prose{ยถา วา ปน มหาราช มหิยา คณา วตฺตนฺติ, เสยฺยถีทํ, มลฺลา อโตณา ปพฺพตา ธมฺมคิริยา พฺรหฺมคิริยา นฏกา นจฺจกา ลงฺฆกา ปิสาจา มณิภทฺทา ปุณฺณพทฺธา จนฺทิมสูริยา สิริเทวตา กาลิเทวตา, สิวา วสุเทวา ฆนิกา อสิปาสา ภทฺทิปุตฺตาติ, เตสํ เตสํ รหสฺสํ เตสุ เตสุ คเณสุ เยว จรติ, อวเสสานํ ปิหิตํ. เอวเมว โข มหาราช เอโส วํโส สพฺเพสํ ปุพฺพกานํ ตถาคตานํ ยทิทํ ภิกฺขุมชฺเฌ ปาติโมกฺขุ\-ทฺเทโส อวเสสานํ ปิหิโต. เอวํ ปุพฺพกานํ ตถาคตานํ วํสวเสน ปาติโมกฺขุ\-ทฺเทโส สีมํ กตฺวา ปิหิโต.}

\prose{กถํ ธมฺมสฺส ครุกตฺตา ปาติโมกฺขุ\-ทฺเทโส สีมํ กตฺวา ปิหิโต, ธมฺโม มหาราช ครุโก ภาริโย, ตตฺถ สมฺมตฺตการี อญฺญํ อาราเธติ, ตํ ตตฺถ ปรมฺปรา\-สมฺมตฺต\-การิตาย ปาปุณาติ, น ตํ ตตฺถ ปรมฺปรา\-สมฺมตฺต\-การิตาย ปาปุณาติ, มา จายํ สารธมฺโม \csromanfolio{190}วรธมฺโม อสมฺมตฺตการีนํ หตฺถคโต โอญฺญาโต อวญฺญาโต หีฬิโต ขีฬิโต ครหิโต ภวตุ, มา จายํ สารธมฺโม วรธมฺโม ทุชฺชนคโต โอญฺญาโต อวญฺญาโต หีฬิโต ขีฬิโต ครหิโต ภวตูติ. เอวํ ธมฺมสฺส ครุกตฺตา ปาติโมกฺขุ\-ทฺเทโส สีมํ กตฺวา ปิหิโต.}

\proseitem{4}{ยถา มหาราช สาร\-วร\-ปวร\-อภิชาต\-ชาติมนฺต\-รตฺต\-โลหิตจนฺทนํ นาม สวร\-ปุรม\-นุคตํ โอญฺญาตํ อวญฺญาตํ หีฬิตํ ขีฬิตํ ครหิตํ ภวติ, เอวเมว โข มหาราช มา จายํ สารธมฺโม วรธมฺโม ปรมฺปราอสมฺมตฺตการีนํ หตฺถคโต โอญฺญาโต อวญฺญาโต หีฬิโต ขีฬิโต ครหิโต ภวตุ, มา จายํ สารธมฺโม วรธมฺโม ทุชฺชนคโต โอญฺญาโต อวญฺญาโต หีฬิโต ขีฬิโต ครหิโต ภวตูติ. เอวํธมฺมสฺส ครุกตฺตา ปาติโมกฺขุ\-ทฺเทโส สีมํ กตฺวา ปิหิโต.}

\prosewithcloserruled{กถํ ภิกฺขุภูมิยา ครุกตฺตา ปาติโมกฺขุ\-ทฺเทโส สีมํ กตฺวา ปิหิโต, ภิกฺขุภาโว โข มหาราช โลเก อตุลิโย อปฺปมาโณ อนคฺฆิโย, น สกฺกา เกนจิ อคฺฆาเปตุํ ตุเลตุํ ปริเมตุํ, มายํ เอวรูเป ภิกฺขุภาเว ฐิโต โลเกน สมสโม ภวตูติ ภิกฺขูนํ เยว อนฺตเร ปาติโมกฺขุ\-ทฺเทโส จรติ. ยถา มหาราช โลเก วรปวร\-ภณฺฑํ วตฺถํ วา อตฺถรณํ วา คชตุรงฺค\-รถ\-สุวณฺณรชต\-มณิมุตฺตา\-อิตฺถิรตนาทีนิ วา วิชิต\-กมฺม\-สูรา วา\csromansharedfootnote{b27947f070f14c5f}{นิชฺชิตกมฺมสูรา วา (สี, อิ)} สพฺเพ เต ราชานมุ\-ปคจฺฉนฺติ, เอวเมว โข มหาราช ยาวตา โลเก\csromansharedfootnote{68383ea03011cfb5}{โลเก สิกฺขา (สี, อิ)} สุคตา\-คม\-ปริยตฺติ\-อาจาร\-สํยม\-สีลสํวร\-คุณา, สพฺเพ เต ภิกฺขุสํฆมุปคตา ภวนฺติ. เอวํ ภิกฺขุภูมิยา ครุกตฺตา ปาติโมกฺขุ\-ทฺเทโส สีมํ กตฺวา ปิหิโตติ. สาธุ ภนฺเต นาคเสน, เอวเมตํ ตถา สมฺปฏิจฺฉามีติ.}{ธมฺมวินยปฏิจฺฉานฺนาปฏิจฺฉนฺนปโญฺห ทุติโย.}
\csromanmatikaanchor{mtk.142}
\csromantocmarkref{subsubsection}{3. มุสาวาทครุลหุภาวปญฺห}{mtk.142}
\bookmark[dest={mtk.142},level=3]{3. มุสาวาทครุลหุภาวปญฺห}
\csromanheader{3}{3. \textbf{มุสาวาท}\-\textbf{ครุ}\-\textbf{ลหุ}\-\textbf{ภาว}\-\textbf{ปญฺห}}
\proseitem{3}{ภนฺเต นาคเสน ภาสิตมฺเปตํ ภควตา “สมฺปชาน\-มุสาวาเท ปาราชิโก โหตี”ติ. ปุน จ ภณิตํ “สมฺปชาน\-มุสาวาเท ลหุกํ \csromanfolio{191}อาปตฺติํ อาปชฺชติ เอกสฺส สนฺติเก เทสนาวตฺถุกนฺ”ติ. ภนฺเต นาคเสน โก \textbf{ปเนตฺถ วิเสโส, กิํ การณํ, ยญฺเจเกน มุสาวาเทน อุจฺฉิชฺชติ, ยญฺเจเกน} \textbf{มุสาวาเทน สเตกิจฺโฉ โหติ. ยทิ ภนฺเต นาคเสน ภควตา ภณิตํ} “\textbf{สมฺปชาน}\-\textbf{มุสาวาเท ปาราชิโก โหตี”ติ, เตน หิ “สมฺปชาน}\-\textbf{มุสาวาเท ลหุกํ} อาปตฺติํ อาปชฺชติ เอกสฺส สนฺติเก เทสนาวตฺถุกนฺ”ติ ยํ วจนํ, ตํ มิจฺฉา. ยทิ ตถาคเตน ภณิตํ “สมฺปชาน\-มุสาวาเท ลหุกํ อาปตฺติํ อาปชฺชติ เอกสฺส สนฺติเก เทสนาวตฺถุกนฺ”ติ, เตน หิ “สมฺปชามุสาวาเท ปาราชิโก โหตี”ติ ตมฺปิ วจนํ มิจฺฉา. อยมฺปิ อุภโต โกฏิโก ปโญฺห ตวานุปฺปตฺโต, โส ตยา นิพฺพาหิตพฺโพติ.}

\prose{ภาสิตมฺเปตํ มหาราช ภควตา “สมฺปชานมุสวาเท ปราชิโก โหตี”ติ. ภณิตญฺจ “สมฺปชาน\-มุสาวาเท ลหุกํ อาปตฺติํ อาปชฺชติ เอกสฺส สนฺติเก เทสนาวตฺถุกนฺ”ติ, ตญฺจ ปน วตฺถุวเสน ครุกลหุกํ โหติ. ตํ กิํ มญฺญสิ มหาราช อิธ โกจิ ปุริโส ปรสฺส ปาณินา ปหารํ ทเทยฺย, ตสฺส ตุเมฺห กิํ ทณฺฑํ ธาเรถาติ. ยทิ โส ภนฺเต อาห “นกฺขมามี”ติ, ตสฺส มยํ อกฺขมมาเน กหาปณํ หราเปมาติ. อิธ ปน มหาราช โส เยว\csromansharedfootnote{cb63cef12ad01f03}{โสเยว สนฺธิ \csromandash{} ม.พ.ป.} ปุริโส ตว ปาณินา ปหารํ ทเทยฺย, ตสฺส ปน โก ทณฺโฑติ. หตฺถมฺปิสฺส ภนฺเต เฉทาเปยฺยาม, ปาทมฺปิ เฉทาเปยฺยาม, ยาว สีสํ กฬีรจฺเฉชฺชํ เฉทาเปยฺยาม, สพฺพมฺปิ ตํ เคหํ วิลุมฺปาเปยฺยาม, อุภโตปกฺเข\csromansharedfootnote{f49e6dec70c74f88}{อุภโตปสฺเส (สี, อิ, ก)} ยาว สตฺตมํ กุลํ สมุคฺฆาตา\-เปยฺยามาติ. โก ปเนตฺถ มหาราช วิเสโส, กิํ การณํ, ยํ เอกสฺส ปาณิปฺปหาเร สุขุโม กหาปโณ ทณฺโฑ, ยํ ตว ปาณิปฺปหาเร หตฺถจฺเฉชฺชํ ปาทจฺเฉชฺชํ ยาว กฬีรจฺเฉชฺชํ สพฺพเคหาทานํ อุภโตปกฺเข\csromansharedfootnote{f49e6dec70c74f88}{อุภโตปสฺเส (สี, อิ, ก)} ยาว สตฺตมกุลา สมุคฺฆาโตติ. มนุสฺสนฺตเรน ภนฺเตติ. เอวเมว โข มหาราช สมฺปชาน\-มุสาวาโท วตฺถุวเสน ครุกลหุโก โหตีติ. สาธุ ภนฺเต นาคเสน, เอวเมตํ ตถา สมฺปฏิจฺฉามีติ.}

\vspace{\tipitakaparskip}
\csromancenter{มุสาวาท\-ครุ\-ลหุ\-ภาว\-ปโญฺห ตติโย.}
\csromanmatikaanchor{mtk.143}
\csromantocmarkref{subsubsection}{4. โพธิสตฺตธมฺมตาปญฺห}{mtk.143}
\bookmark[dest={mtk.143},level=3]{4. โพธิสตฺตธมฺมตาปญฺห}
\csromanheader{3}{4. \textbf{โพธิสตฺต}\-\textbf{ธมฺมตาปญฺห}}
\proseitem{4}{\csromanfolio{192}ภนฺเต นาคเสน ภาสิตมฺเปตํ ภควตา\csromansymbolfootnote{*}{ที 2. 10 ปิฏฺเฐ.}\textbf{ธมฺมตา}\-\textbf{ธมฺมปริยาเย} “ปุพฺเพว โพธิสตฺตานํ มาตาปิตโร นิยตา โหนฺติ, โพธิ นิยตา โหติ, อคฺคสาวกา นิยตา โหนฺติ, ปุตฺโต นิยโต โหติ, อุปฏฺฐาโก นิยโต โหตี”ติ. ปุน จ ตุเมฺห ภณถ “ตุสิเต กาเย ฐิโต โพธิสตฺโต อฏฺฐ มหาวิโลกนานิ วิโลเกติ, กาลํ วิโลเกติ, ทีปํ วิโลเกติ, เทสํ วิโลเกติ, กุลํ วิโลเกติ, ชเนตฺติํ วิโลเกติ, อายุํ วิโลเกติ, มาสํ วิโลเกติ, เนกฺขมฺมํ วิโลเกตี”ติ. ภนฺเต นาคเสน อปริปกฺเก ญาเณ พุชฺฌนํ นตฺถิ, ปริปกฺเก ญาเณ น สกฺกา นิเมสนฺตรมฺปิ อาคเมตุํ, อนติกฺกมนียํ ปริปกฺก\-มานสํ. กสฺมา โพธิสตฺโต กาลํ วิโลเกติ “กมฺหิ กาเล อุปฺปชฺชามี”ติ. อปริปกฺเก ญาเณ พุชฺฌนํ นตฺถิ, ปริปกฺเก ญาเณ น สกฺกา นิเมสนฺตรมฺปิ อาคเมตุํ, กสฺมา โพธิสตฺโต กุลํ วิโลเกติ “กมฺหิ กุเล อุปฺปชฺชามี”ติ. ยทิ ภนฺเต นาคเสน ปุพฺเพว โพธิสตฺตสฺส มาตาปิตโร นิยตา, เตน หิ “กุลํ วิโลเกตี”ติ ยํ วจนํ, ตํ มิจฺฉา. ยทิ กุลํ วิโลเกติ, เตน หิ “ปุพฺเพว โพธิสตฺตสฺส มาตาปิตโร นิยตา”ติ ตมฺปิ วจนํ มิจฺฉา. อยมฺปิ อุภโต โกฏิโก ปโญฺห ตวานุปฺปตฺโต, โส ตยา นิพฺพาหิตพฺโพติ.}

\prose{นิยตา มหาราช ปุพฺเพว โพธิสตฺตสฺส มาตาปิตโร, กุสญฺจ โพธิสตฺโต วิโลเกติ. กินฺติ ปน กุลํ วิโลเกติ “เย เม มาตาปิตโร, เต ขตฺติยา อุทาหุ พฺราหฺมณา”ติ. เอวํ กุลํ วิโลเกติ.}

\prosewithcloserruled{อฏฺฐนฺนํ มหาราช ปุพฺเพว อนาคตํ โอโลเกตพฺพํ โหติ. กตเมสํ อฏฺฐนฺนํ, วาณิชสฺส มหาราช ปุพฺเพว วิกฺกยภณฺฑํ โอโลเกตพฺพํ โหติ, หตฺถินาคสฺส ปุพฺเพว โสณฺฑาย อนาคโต มคฺโค โอโลเกตพฺโพ โหติ, สากฏิกสฺส ปุพฺเพว อนาคตํ ติตฺถํ โอโลเกตพฺพํ โหติ, นิยามกสฺส ปุพฺเพว อนาคตํ ตีรํ โอโลเกตฺวา นาวา เปเสตพฺพา โหติ, ภิสกฺกสฺส ปุพฺเพว อายุํ โอโลเกตฺวา อาตุโร อุปสงฺกมิตพฺโพ โหติ, อุตฺตรเสตุสฺส ปุพฺเพว ถิรา\-ถิรภาวํ ชานิตฺวา อภิรุหิตพฺพํ โหติ, ภิกฺขุสฺส ปุพฺเพว อนาคตํ กาลํ ปจฺจเวกฺขิตฺวา โภชนํ ภุญฺชิตพฺพํ โหติ, โพธิสตฺตานํ ปุพฺเพว กุลํ โอโลเกตพฺพํ \csromanfolio{193}โหติ “ขตฺติยกุลํ วา พฺราหฺมณกุลํ วา”ติ. อิเมสํ โข มหาราช อฏฺฐนฺนํ ปุพฺเพว อนาคตํ โอโลเกตพฺพํ โหตีติ. สาธุ ภนฺเต นาคเสน, เอวเมตํ ตถา สมฺปฏิจฺฉามีติ.}{โพธิสตฺต\-ธมฺมตาปโญฺห จตุตฺโถ.}
\csromanmatikaanchor{mtk.144}
\csromantocmarkref{subsubsection}{5. อตฺตนิปาตนปญฺห}{mtk.144}
\bookmark[dest={mtk.144},level=3]{5. อตฺตนิปาตนปญฺห}
\csromanheader{3}{5. \textbf{อตฺต}\-\textbf{นิปาตน}\-\textbf{ปญฺห}}
\proseitem{4}{ภนฺเต นาคเสน ภาสิตมฺเปตํ ภควตา\csromansymbolfootnote{*}{วิ 1. 104 ปิฏฺเฐ.}“น ภิกฺขเว อตฺตานํ ปาเตตพฺพํ, โย ปาเตยฺย, ยถาธมฺโม กาเรตพฺโพ”ติ. ปุน จ ตุเมฺห ภณถ “ยตฺถ กตฺถจิ ภควา สาวกานํ ธมฺมํ เทสยมาโน อเนก\-ปริยาเยน ชาติยา ชราย พฺยาธิโน มรณสฺส สมุจฺเฉทาย ธมฺมํ เทเสติ, โย หิ โกจิ ชาติ\-ชรา\-พฺยาธิ\-มรณํ สมติกฺกมติ, ตํ ปรมาย ปสํสาย ปสํสตี”ติ. ยทิ ภนฺเต นาคเสน ภควตา ภณิตํ “น ภิกฺขเว อตฺตานํ ปาเตตพฺพํ, โย ปาเตยฺย, ยถาธมฺโม กาเรตพฺโพ”ติ, เตน หิ “ชาติยา ชราย พฺยาธิโน มรณสฺส สมุจฺเฉทาย ธมฺมํ เทเสตี”ติ ยํ วจนํ, ตํ มิจฺฉา. ยทิ ชาติยา ชราย พฺยาธิโน มรณสฺส สมุจฺเฉทาย ธมฺมํ เทเสติ, เตน หิ “น ภิกฺขเว อตฺตานํ ปาเตตพฺพํ, โย ปาเตยฺย, ยถาธมฺโม กาเรตพฺโพ”ติ ตมฺปิ วจนํ มิจฺฉา. อยมฺปิ อุภโต โกฏิโก ปโญฺห ตวานุปฺปตฺโต, โส ตยา นิพฺพาหิตพฺโพติ.}

\prose{ภาสิตมฺเปตํ มหาราช ภควตา “น ภิกฺขเว อตฺตานํ ปาเตตพฺพํ โย ปาเตยฺย, ยถาธมฺโม กาเรตพฺโพ”ติ. ยตฺถ กตฺถจิ ภควตา สาวกานํ ธมฺมํ เทสยมาเนน จ อเนก\-ปริยาเยน ชาติยา ชราย พฺยาธิโน มรณสฺส สมุจฺเฉทาย ธมฺโม เทสิโต, ตตฺถ ปน การณํ อตฺถิ, เยน ภควา การเณน ปฏิกฺขิปิ สมาทเปสิ จาติ.}

\prose{กิํ ปเนตฺถ ภนฺเต นาคเสน การณํ, เยน ภควา การเณน ปฏิกฺขิปิ สมาทเปสิ จาติ, สีลวา มหาราช สีลสมฺปนฺโน อคทสโม สตฺตานํ กิเลส\-วิส\-วินาสเน, โอสธสโม สตฺตานํ กิเลสพฺยาธิ\-วูปสเม, อุทกสโม สตฺตานํ กิเลส\-รโชชลฺลา\-ปหรเณ, มณิรตนสโม สตฺตานํ สพฺพสมฺปตฺติ\-ทาเน, นาวาสโม สตฺตานํ จตุโรฆ\-ปาร\-คมเน, สตฺถวาหสโม สตฺตานํ ชาติ\-กนฺตารตา\-รเณ, วาตสโม \csromanfolio{194}สตฺตานํ ติวิธคฺคิ\-สนฺตาป\-นิพฺพาปเน, มหาเมฆสโม สตฺตานํ มานส\-ปริปูรเณ, อาจริยสโม สตฺตานํ กุสล\-สิกฺขาปเน, สุเทสกสโม สตฺตานํ เขม\-ปถมา\-จิกฺขเณ. เอวรูโป มหาราช พหุคุโณ อเนกคุโณ อปฺปมาณคุโณ คุณราสิ คุณปุญฺโช สตฺตานํ วฑฺฒิกโร สีลวา “มา วินสฺสี”ติ สตฺตานํ อนุกมฺปาย ภควา สิกฺขาปทํ ปญฺญเปสิ “น ภิกฺขเว อตฺตานํ ปาเตตพฺพํ, โย ปาเตยฺย, ยถาธมฺโม กาเรตพฺโพ”ติ. อิทเมตฺถ มหาราช การณํ, เยน การเณน ภควา ปฏิกฺขิปิ. ภาสิตมฺเปตํ มหาราช เถเรน กุมาร\-กสฺสเปน วิจิตฺร\-กถิเกน ปายาสิ\-ราชญฺญสฺส ปรโลกํ ทีปยมาเนน\csromansymbolfootnote{*}{ที 2. 264 ปิฏฺเฐ.}“ยถา ยถา โข ราชญฺญ สมณพฺราหฺมณา สีลวนฺโต กลฺยาณธมฺมา จิรํ ทีฆมทฺธานํ ติฏฺฐนฺติ, ตถา ตถา พหุํ ปุญฺญํ ปสวนฺติ, พหุชนหิตาย จ ปฏิปชฺชนฺติ พหุชน\-สุขาย โลกานุกมฺปาย อตฺถาย หิตาย สุขาย เทวมนุสฺสานนฺ”ติ.}

\proseitem{5}{เกน ปน การเณน ภควา สมาทเปสิ, ชาติปิ มหาราช ทุกฺขา, ชราปิ ทุกฺขา, พฺยาธิปิ ทุกฺโข, มรณมฺปิ ทุกฺขํ, โสโกปิ ทุกฺโข, ปริเทโวปิ ทุกฺโข, ทุกฺขมฺปิ ทุกฺขํ, โทมนสฺสมฺปิ ทุกฺขํ, อุปายาโสปิ ทุกฺโข, อปฺปิเยหิ สมฺปโยโคปิ ทุกฺโข, ปิเยหิ วิปฺปโยโคปิ ทุกฺโข, มาตุมรณมฺปิ ทุกฺขํ. ปิตุมรณมฺปิ ทุกฺขํ, ภาตุมรณมฺปิ ทุกฺขํ, ภคินิ\-มรณมฺปิ ทุกฺขํ, ปุตฺตมรณมฺปิ ทุกฺขํ ทารมรณมฺปิ ทุกฺขํ, ทาสมรณมฺปิ ทุกฺขํ\csromansharedfootnote{6dde0ec79e190dc5}{อิทํ วากฺยํ สีอิโปตฺถเกสุ นตฺถิ.}, ญาติมรณมฺปิ ทุกฺขํ, ญาติพฺยสนมฺปิ ทุกฺขํ, โรคพฺยสนมฺปิ ทุกฺขํ, โภค\-พฺยสนมฺปิ ทุกฺขํ, สีลพฺยสนมฺปิ ทุกฺขํ, ทิฏฺฐิ\-พฺยสนมฺปิ ทุกฺขํ, ราชภยมฺปิ ทุกฺขํ, โจรภยมฺปิ ทิกฺขํ, เวริภยมฺปิ ทุกฺขํ, ทุพฺภิกฺข\-ภยมฺปิ ทุกฺขํ, อคฺคิภยมฺปิ ทุกฺขํ, อุทกภยมฺปิ ทุกฺขํ, อูมิภยมฺปิ ทุกฺขํ, อาวฏฺฏภยมฺปิ ทุกฺขํ, กุมฺภีล\-ภยมฺปิ ทุกฺขํ, สุสุกาภยมฺปิ ทุกฺขํ, อตฺตานุวาท\-ภยมฺปิ ทุกฺขํ, ปรานุวาท\-ภยมฺปิ ทุกฺขํ, ทณฺฑภยมฺปิ ทุกฺขํ, ทุคฺคติ\-ภยมฺปิ ทุกฺขํ, ปริสาสารชฺชภยมฺปิ ทุกฺขํ, อาชีวกภยมฺปิ ทุกฺขํ, มรณภยมฺปิ ทิกฺขํ, เวตฺเตหิ ตาฬนมฺปิ ทุกฺขํ, กสาหิ ตาฬนมฺปิ ทุกฺขํ, อทฺธทณฺฑเกหิ ตาฬนมฺปิ ทุกฺขํ, หตฺถ\-จฺเฉทนมฺปิ ทุกฺขํ, ปาทจฺเฉทนมฺปิ ทุกฺขํ, หตฺถปาทจฺเฉทนมฺปิ ทุกฺขํ, กณฺณจฺเฉทนมฺปิ ทุกฺขํ, นาสจฺเฉทนมฺปิ ทุกฺขํ, กณฺณนาสจฺเฉทนมฺปิ ทุกฺขํ, พิลงฺค\-ถาลิกมฺปิ ทุกฺขํ, สงฺข\-มุณฺฑิกมฺปิ ทุกฺขํ, ราหุมุขมฺปิ ทุกฺขํ, โชติมาลิกมฺปิ ทุกฺขํ, หตฺถ\-ปชฺโชติกมฺปิ ทุกฺขํ, เอรก\-วตฺติกมฺปิ ทุกฺขํ, จีรก\-วาสิกมฺปิ ทุกฺขํ, เอเณยฺยากมฺปิ \csromanfolio{195}ทุกฺขํ, พฬิส\-มํสิกมฺปิ ทุกฺขํ, กหาปณิกมฺปิ ทุกฺขํ, ขารา\-ปตจฺฉิกมฺปิ ทุกฺขํ, ปลิฆ\-ปริวตฺติกมฺปิ ทุกฺขํ, ปลาล\-ปีฐกมฺปิ ทุกฺขํ, ตตฺเตน เตเลน โอสิญฺจนมฺปิ ทุกฺขํ, สุนเขหิ ขาทาปนมฺปิ ทุกฺขํ, ชีวสูลาโรปนมฺปิ ทุกฺขํ, อสินา สีสจฺเฉทนมฺปิ ทุกฺขํ, เอวรูปานิ มหาราช พหุวิธานิ อเนกวิธานิ ทุกฺขานิ สํสารคโต อนุภวติ.}

\prosewithcloserruled{ยถา มหาราช หิมวนฺต\-ปพฺพเต อภิวุฏฺฐํ อุทกํ คงฺคาย นทิยา ปาสาณ สกฺขร ขร มรุมฺพ อาวฏฺฏ คคฺคลก อูมิกวงฺก จทิก อาวรณ\-นีวรณ\-มูลก\-สาขาสุ ปริโยตฺถรติ, เอวเมว โข มหาราช เอวรูปานิ พหุวิธานิ อเนกวิธานิ ทุกฺขานิ สํสารคโต อนุภวติ. ปวตฺตํ มหาราช ทุกฺขํ, อปฺปวตฺตํ สุขํ. อปฺปวตฺตสฺส คุณํ ปวตฺเต\csromansharedfootnote{3c42c8ca091f97fe}{ปวตฺตสฺส (ก)} จ ภยํ ทีปยมาโน มหาราช ภควา อปฺปวตฺตสฺส สจฺฉิกิริยาย ชาติชราพฺยาธิมรณ\-สมติกฺกมาย สมาทเปสิ, อิทเมตฺถ มหาราช การณํ, เยน การเณน ภควา สมาทเปสีติ. สาธุ ภนฺเต นาคเสน สุนิพฺเพฐิโต ปโญฺห, สุกถิตํ การณํ, เอวเมตํ ตถา สมฺปฏิจฺฉามีติ.}{อตฺต\-นิปาตน\-ปโญฺห ปญฺจโม.}
\csromanmatikaanchor{mtk.145}
\csromantocmarkref{subsubsection}{6. เมตฺตาภาวนานิสํสปญฺห}{mtk.145}
\bookmark[dest={mtk.145},level=3]{6. เมตฺตาภาวนานิสํสปญฺห}
\csromanheader{3}{6. \textbf{เมตฺตาภาวนา}\-\textbf{นิสํส}\-\textbf{ปญฺห}}
\proseitem{5}{ภนฺเต นาคเสน ภาสิตมฺเปตํ ภควตา “เมตฺตาย ภิกฺขเว เจโตวิมุตฺติยา อาเสวิตาย ภาวิตาย พหุลีกตาย ยานีกตาย วตฺถุกตาย อนุฏฺฐิตาย ปริจิตาย สุสมารทฺธาย เอกาทสานิสํสา ปาฏิกงฺขา. กตเม เอกาทส, สุขํ สุปติ, สุขํ ปฏิพุชฺฌติ, น ปาปกํ สุปินํ ปสฺสติ, มนุสฺสานํ ปิโย โหติ, อมนุสฺสานํ ปิโย โหติ, เทวตา รกฺขนฺติ, นาสฺส อคฺคิ วา วิสํ วา สตฺถํ วา กมติ\csromansharedfootnote{ccdfbee696f56291}{สตฺถํ กมติ (สฺยา) อํ 3. 542; วิ 5. 249; ขุ 9. 314 ปิฏฺเฐสุ ปสฺสิตพฺพํ.}, ตุวฏํ จิตฺตํ สมาธิยติ, มุขวณฺโณ วิปฺปสีทติ, อสมฺมูโฬฺห กาลํ กโรติ, อุตฺตริํ อปฺปฏิวิชฺฌนฺโต พฺรหฺมโลกูปโค โหตี”ติ. ปุน จ ตุเมฺห ภณถ\csromansymbolfootnote{*}{ขุ 6. 177 ปิฏฺเฐ.}“สาโม กุมาโร เมตฺตาวิหารี มิคสํเฆน ปริวุโต ปวเน วิจรนฺโต ปีฬิยกฺเขน รญฺญา วิทฺโธ วิสปีเตน สลฺเลน ตตฺเถว มุจฺฉิโต ปติโต”ติ.}

\proseitem{6}{\csromanfolio{196}ยทิ ภนฺเต นาคเสน ภควตา ภณิตํ “เมตฺตาย ภิกฺขเว ฯเปฯ พฺรหฺมโลกูปโค โหตี”ติ, เตน หิ “สาโม กุมาโร เมตฺตาวิหารี มิคสํเฆน ปริวุโต ปวเน วิจรนฺโต ปีฬิยกฺเขน รญฺญา วิทฺโธ วิสปีเตน สลฺเลน ตตฺเถว มุจฺฉิโต ปติโต”ติ ยํ วจนํ, ตํ มิจฺฉา. ยทิ สาโม กุมาโร เมตฺตาวิหารี มิคสํเฆน ปริวุโต ปวเน วิจรนฺโต ปีฬิยกฺเขน รญฺญา วิทฺโธ วิสปีเตน สลฺเลน ตตฺเถว มุจฺฉิโต ปติโต, เตน หิ “เมตฺตาย ภิกฺขเว ฯเปฯ สตฺถํ วา กมตี”ติ ตมฺปิ วจนํ มิจฺฉา. อยมฺปิ อุภโต โกฏิโก ปโญฺห สุนิปุโณ ปริสโณฺห สุขุโม คมฺภีโร, อปิ สุนิปุณานํ มนุชานํ คตฺเต เสทํ โมเจยฺย, โส ตวานุปฺปตฺโต, วิชเฏหิ ตํ มหา\-ชฏา\-ชฏิตํ, อนาคตานํ ชินปุตฺตานํ จกฺขุํ เทหิ นิพฺพาหนายาติ.}

\prose{ภาสิตมฺเปตํ มหาราช ภควตา “เมตฺตาย ภิกฺขเว ฯเปฯ สตฺถํ วา กมตี”ติ. สาโม จ กุมาโร เมตฺตาวิหารี มิคสํเฆน ปริวุโต ปวเน วิจรนฺโต ปีฬิยกฺเขน รญฺญา วิทฺโธ วิสปีเตน สลฺเลน ตตฺเถว มุจฺฉิโต ปติโต, ตตฺถ ปน มหาราช การณํ อตฺถิ, กตมํ ตตฺถ การณํ, เนเต มหาราช คุณา ปุคฺคลสฺส, เมตฺตาภาวนาเยเต คุณา, สาโม มหาราช กุมาโร ฆฏํ อุกฺขิปนฺโต ตสฺมิํ ขเณ เมตฺตาภาวนาย\csromansymbolfootnote{*}{ขุ 6. 174 ปิฏฺเฐ.}ปมตฺโต อโหสิ.}

\prose{ยสฺมิํ มหาราช ขเณ ปุคฺคโล เมตฺตํ สมาปนฺโน โหติ, น ตสฺส ปุคฺคลสฺส ตสฺมิํ ขเณ อคฺคิ วา วิสํ วา สตฺถํ วา กมติ, ตสฺส เย เกจิ อหิตกามา อุปคนฺตฺวา ตํ น ปสฺสนฺติ, น ตสฺมิํ โอกาสํ ลภนฺติ, เนเต มหาราช คุณา ปุคฺคลสฺส, เมตฺตาภาวนาเยเต คุณา. อิธ มหาราช ปุริโส สงฺคามสูโร อเภชฺช\-กวจ\-ชาลิกํ สนฺนยฺหิตฺวา สงฺคามํ โอตเรยฺย, ตสฺส สรา ขิตฺตา อุปคนฺตฺวา ปตนฺติ วิกิรนฺติ, น ตสฺมิํ โอกาสํ ลภนฺติ, เนโส มหาราช คุโณ สงฺคามสูรสฺส, อเภชฺช\-กวจ\-ชาลิกาเยโส คุโณ, ยสฺส สรา ขิตฺตา อุปคนฺตฺวา ปตนฺติ วิกิรนฺติ. เอวเมว โข มหาราช เนเต คุณา ปุคฺคลสฺส เมตฺตาภาวนาเยเต คุณา.}

\prose{ยสฺมิํ มหาราช ขเณ ปุคฺคโล เมตฺตํ สมาปนฺโน โหติ, น ตสฺส ปุคฺคลสฺส ตสฺมิํ ขเณ อคฺคิ วา วิสํ วา สตฺถํ วา กมติ, ตสฺส เย เกจิ อหิตกามา อุปคนฺตฺวา ตํ น ปสฺสนฺติ, ตสฺมิํ โอกาสํ น ลภนฺติ. \csromanfolio{197}เนเต มหาราช คุณา ปุคฺคลสฺส, เมตฺตาภาวนาเยเต คุณา. อิธ ปน มหาราช ปุริโส ทิพฺพํ อนฺตรธานํ มูลํ หตฺเถ กเรยฺย, ยาว ตํ มูลํ ตสฺส หตฺถคตํ โหติ, ตาว น อญฺโญ โกจิ ปกติมนุสฺโส ตํ ปุริสํ ปสฺสติ, เนโส มหาราช คุโณ ปุริสสฺส, มูลสฺเสโส คุโณ อนฺตรธานสฺส, ยํ โส ปกติ\-มนุสฺสานํ จกฺขุปเถ น ทิสฺสติ. เอวเมว โข มหาราช เนเต คุณา ปุคฺคลสฺส, เมตฺตาภาวนาเยเต คุณา.}

\prose{ยสฺมิํ มหาราช ขเณ ปุคฺคโล เมตฺตํ สมาปนฺโน โหติ, น ตสฺส ปุคฺคลสฺส ตสฺมิํ ขเณ อคฺคิ วา วิสํ วา สตฺถํ วา กมติ, ตสฺส เย เกจิ อหิตกามา อุปคนฺตฺวา ตํ น ปสฺสนฺติ, น ตสฺมิํ โอกาสํ ลภนฺติ, เนเต มหาราช คุณา ปุคฺคลสฺส, เมตฺตาภาวนาเยเต คุณา. ยถา วา ปน มหาราช ปุริสํ สุกตํ มหาเลณม\-นุปฺปวิฏฺฐํ มหติ\-มหาเมโฆ อภิวสฺสนฺโต น สกฺโกติ เตมยิตุํ, เนโส มหาราช คุโณ ปุริสสฺส, มหาเลณสฺเสโส คุโณ, ยํ มหาเมโฆ อภิวสฺสมาโน น ตํ เตเมติ. เอวเมว โข มหาราช เนเต คุณา ปุคฺคลสฺส, เมตฺตาภาวนาเยเต คุณา.}

\prosewithcloserruled{ยสฺมิํ มหาราช ขเณ ปุคฺคโล เมตฺตํ สมาปนฺโน โหติ, น ตสฺส ปุคฺคลสฺส ตสฺมิํ ขเณ อคฺคิ วา วิสํ วา สตฺถํ วา กมติ, ตสฺส เย เกจิ อหิตกามา อุปคนฺตฺวา ตํ น ปสฺสนฺติ, น ตสฺส สกฺโก นฺติ อหิตํ กาตุํ, เนเต มหาราช คุณา ปุคฺคลสฺส, เมตฺตาภาวนาเยเต คุณาติ. อจฺฉริยํ ภนฺเต นาคเสน, อพฺภุตํ ภนฺเต นาคเสน, สพฺพปาป\-นิวารณา เมตฺตาภาวนาติ. สพฺพ\-กุสล\-คุณา\-วหา มหาราช เมตฺตาภาวนา หิตานมฺปิ อหิตานมฺปิ, เย เต สตฺตา วิญฺญาณพทฺธา, สพฺเพสํ มหานิสํสา เมตฺตาภาวนา สํวิภชิตพฺพาติ.}{เมตฺตาภาวนา\-นิสํส\-ปโญฺห ฉฏฺโฐ.}
\csromanmatikaanchor{mtk.146}
\csromantocmarkref{subsubsection}{7. กุสลากุสลสมวิสมปญฺห}{mtk.146}
\bookmark[dest={mtk.146},level=3]{7. กุสลากุสลสมวิสมปญฺห}
\csromanheader{3}{7. \textbf{กุสลากุสล}\-\textbf{สมวิสม}\-\textbf{ปญฺห}}
\proseitem{7}{ภนฺเต นาคเสน “กุสล\-การิสฺสปิ อกุสล\-การิสฺสปิ วิปาโก สมสโม, อุทาหุ โกจิ วิเสโส อตฺถี”ติ. อตฺถิ มหาราช \csromanfolio{198}กุสลสฺส จ อกุสลสฺส จ วิเสโส, กุสลํ มหาราช สุขวิปากํ สคฺค\-สํวตฺตนิกํ, อกุสลํ ทุกฺขวิปากํ นิรยสํวตฺตนิกนฺติ.}

\prose{ภนฺเต นาคเสน ตุเมฺห ภณถ “เทวทตฺโต เอกนฺตกโณฺห, เอกนฺตกเณฺหหิ ธมฺเมหิ สมนฺนาคโต, โพธิสตฺโต เอกนฺตสุกฺโก, เอกนฺตสุกฺเกหิ ธมฺเมหิ สมนฺนาคโต”ติ. ปุน จ เทวทตฺโต ภเว ภเว ยเสน จ ปกฺเขน จ โพธิสตฺเตน สมสโม โหติ, กทาจิ อธิกตโร วา. ยทา เทวทตฺโต นคเร พาราณสิยํ พฺรหฺมทตฺตสฺส รญฺโญ ปุโรหิตปุตฺโต อโหสิ, ตทา โพธิสตฺโต ฉวกจณฺฑาโล อโหสิ วิชฺชาธโร, วิชฺชํ ปริชปฺปิตฺวา อกาเล อมฺพผลานิ นิพฺพตฺเตสิ, เอตฺถ ตาว โพธิสตฺโต เทวทตฺตโต ชาติยา นิหีโน ยเสน จ นิหีโน.}

\prose{ปุน จปรํ ยทา เทวทตฺโต ราชา อโหสิ มหา มหีปติ สพฺพกาม\-สมงฺคี, ตทา โพธิสตฺโต ตสฺสูปโภโค อโหสิ หตฺถินาโค สพฺพ\-ลกฺขณสมฺปนฺโน, ตสฺส จารุ\-คติ\-วิลาสํ อสหมาโน ราชา วธมิจฺฉนฺโต หตฺถาจริยํ เอวมโวจ “อสิกฺขิโต เต อาจริย หตฺถินาโค, ตสฺส อากาสคมนํ นาม การณํ กโรหี”ติ, ตตฺถปิ ตาว โพธิสตฺโต เทวทตฺตโต ชาติยา นิหีโน ลามโก ติรจฺฉานคโต.}

\prose{ปุน จปรํ ยทา เทวทตฺโต มนุสฺโส อโหสิ ปวเน นฏฺฐายิโก, ตทา โพธิสตฺโต มหาปถวี นาม มกฺกโฏ อโหสิ, เอตฺถปิ ตาว ทิสฺสติ วิเสโส มนุสฺสสฺส จ ติรจฺฉาน\-คตสฺส จ, ตตฺถปิ ตาว โพธิสตฺโต เทวทตฺตโต ชาติยา นิหีโน.}

\prose{ปุน จปรํ ยทา เทวทตฺโต มนุสฺโส อโหสิ โสณุตฺตโร นาม เนสาโท พลวา พลวตโร นาคพโล, ตทา โพธิสตฺโต ฉทฺทนฺโต นาม นาคราชา อโหสิ, ตทา โส ลุทฺทโก ตํ หตฺถินาคํ ฆาเตสิ, ตตฺถปิ ตาว เทวทตฺโตว อธิกตโร.}

\prose{ปุน จปรํ ยทา เทวทตฺโต มนุสฺโส อโหสิ วนจรโก อนิเกตวาสี, ตทา โพธิสตฺโต สกุโณ อโหสิ ติตฺติโร มนฺตชฺฌายี, ตทาปิ โส วนจรโก ตํ สกุณํ ฆาเตสิ, ตตฺถปิ ตาว เทวทตฺโตว ชาติยา อธิกตโร.}

\prose{\csromanfolio{199}ปุน จปรํ ยทา เทวทตฺโต กลาพุ นาม กาสิราชา\csromansharedfootnote{f1976408f2ad41ec}{กาสิกราชา (ก)} อโหสิ, ตทา โพธิสตฺโต ตาปโส อโหสิ ขนฺติวาที, ตทา โส ราชา ตสฺส ตาปสสฺส กุทฺโธ หตฺถปาเท วํสกฬีเร วิย เฉทาเปสิ, ตตฺถปิ ตาว เทวทตฺโต เยว อธิกตโร ชาติยา จ ยเสน จ.}

\prose{ปุน จปรํ ยทา เทวทตฺโต มนุสฺโส อโหสิ วนจโร, ตทา โพธิสตฺโต นนฺทิโย นาม วานรินฺโท อโหสิ, ตทาปิ โส วนจโร ตํ วานรินฺทํ ฆาเตสิ สทฺธิํ มาตรา กนิฏฺฐ\-ภาติเกน จ, ตตฺถปิ ตาว เทวทตฺโต เยว อธิกตโร ชาติยา.}

\prose{ปุน จปรํ ยทา เทวทตฺโต มนุสฺโส อโหสิ อเจลโก การมฺภิโย นาม, ตทา โพธิสตฺโต ปณฺฑรโก นาม นาคราชา อโหสิ, ตตฺถปิ ตาว เทวทตฺโต เยว อธิกตโร ชาติยา.}

\prose{ปุน จปรํ ยทา เทวทตฺโต มนุสฺโส อโหสิ ปวเน ชฏิลโก, ตทา โพธิสตฺโต ตจฺฉโก นาม มหาสูกโร อโหสิ, ตตฺถปิ ตาว เทวทตฺโต เยว ชาติยา อธิกโร.}

\prose{ปุน จปรํ ยทา เทวทตฺโต เจตีสุ สูรปริจโร นาม ราชา อโหสิ อุปริ ปุริสมตฺเต คคเน เวหาสงฺคโม, ตทา โพธิสตฺโต กปิโล นาม พฺราหฺมโณ อโหสิ, ตตฺถปิ ตาว เทวทตฺโต เยว อธิกตโร ชาติยา จ ยเสน จ.}

\prose{ปุน จปรํ ยทา เทวทตฺโต มนุสฺโส อโหสิ สาโน นาม, ตทา โพธิสตฺโต รุรุ นาม มิคราชา อโหสิ, ตตฺถปิ ตาว เทวทตฺโต เยว ชาติยา อธิกตโร.}

\prose{ปุน จปรํ ยทา เทวทตฺโต มนุสฺโส อโหสิ ลุทฺทโก ปวนจโร, ตทา โพธิสตฺโต หตฺถินาโค อโหสิ, โส ลุทฺทโก ตสฺส หตฺถินาคสฺส สตฺตกฺขตฺตุํ ทนฺเต ฉินฺทิตฺวา หริ, ตตฺถปิ ตาว เทวทตฺโต เยว โยนิยา อธิกตโร.}

\prose{ปุน จปรํ ยทา เทวทตฺโต สิงฺคาโล อโหสิ ขตฺติยธมฺโม, โส ยาวตา ชมฺพุทีเป ปเทสราชาโน เต สพฺเพ อนุยุตฺเต อกาสิ,}

\prose{\csromanfolio{200}ตทา โพธิสตฺโต วิธุโร นาม ปณฺฑิโต อโหสิ, ตตฺถปิ ตาว เทวทตฺโต เยว ยเสน อธิกตโร.}

\prose{ปุน จปรํ ยทา เทวทตฺโต หตฺถินาโค หุตฺวา ลฏุกิกาย สกุณิกาย ปุตฺตเก ฆาเตสิ, ตทา โพธิสตฺโตปิ หตฺถินาโค อโหสิ ยูถปติ, ตตฺถ ตาว อุโภปิ เต สมสมา อเหสุํ.}

\prose{ปุน จปรํ ยทา เทวทตฺโต ยกฺโข อโหสิ อธมฺโม นาม, ตทา โพธิสตฺโตปิ ยกฺโข อโหสิ ธมฺโม นาม, ตตฺถปิ ตาว อุโภปิ สมสมา อเหสุํ.}

\prose{ปุน จปรํ ยทา เทวทตฺโต นาวิโก อโหสิ ปญฺจนฺนํ กุลสตานํ อิสฺสโร, ตทา โพธิสตฺโตปิ นาวิโก อโหสิ ปญฺจนฺนํ กุลสตานํ อิสฺสโร, ตตฺถปิ ตาว อุโภปิ สมสมา อเหสุํ.}

\prose{ปุน จปรํ ยทา เทวทตฺโต สตฺถวาโห อโหสิ ปญฺจนฺนํ สกฏสตานํ อิสฺสโร, ตทา โพธิสตฺโตปิ สตฺถวาโห อโหสิ ปญฺจนฺนํ สกฏสตานํ อิสฺสโร, ตตฺถปิ ตาว อุโภปิ สมสมา อเหสุํ.}

\prose{ปุน จปรํ ยทา เทวทตฺโต สาโข นาม มิคราชา อโหสิ, ตทา โพธิสตฺโตปิ นิโคฺรโธ นาม มิคราชา อโหสิ, ตตฺถปิ ตาว อุโภปิ สมสมา อเหสุํ.}

\prose{ปุน จปรํ ยทา เทวทตฺโต สาโข นาม เสนาปติ อโหสิ, ตทา โพธิสตฺโตปิ นิโคฺรโธ นาม ราชา อโหสิ, ตตฺถปิ ตาว อุโภปิ สมสมา อเหสุํ.}

\prose{ปุน จปรํ ยทา เทวทตฺโต ขณฺฑหาโล นาม พฺราหฺมโณ อโหสิ, ตทา โพธิสตฺโต จนฺโท นาม ราชกุมาโร อโหสิ, ตทา โส ขณฺฑหาโล เยว อธิกตโร.}

\prose{ปุน จปรํ ยทา เทวทตฺโต พฺรหฺมทตฺโต นาม ราชา อโหสิ, ตทา โพธิสตฺโต ตสฺส ปุตฺโต มหาปทุโม นาม กุมาโร อโหสิ, ตทา โส ราชา สกปุตฺตํ โจรปปาเต ขิปาเปสิ, ยโต กุโตจิ ปิตาว ปุตฺตานํ อธิกตโร โหติ วิสิฏฺโฐติ, ตตฺถปิ ตาว เทวทตฺโต เยว อธิกตโร.}

\prose{\csromanfolio{201}ปุน จปรํ ยทา เทวทตฺโต มหาปาตาโป นาม ราช อโหสิ, ตทา โพธิสตฺโต ตสฺส ปุตฺโต ธมฺมปาโล นาม กุมาโร อโหสิ, ตทา โส ราชา สกปุตฺตสฺส หตฺถปาเท สีสญฺจ เฉทาเปสิ, ตตฺถปิ ตาว เทวทตฺโต เยว อุตฺตโร อธิกตโร.}

\prose{อชฺเชตรหิ อุโภปิ สกฺยกุเล ชายิํสุ, โพธิสตฺโต พุทฺโธ อโหสิ สพฺพญฺญู โลกนายโก, เทวทตฺโต ตสฺส เทวาติเทวสฺส สาสเน ปพฺพชิตฺวา อิทฺธํ นิพฺพตฺเตตฺวา พุทฺธาลยํ อกาสิ. กินฺนุ โข ภนฺเต นาคเสน ยํ มยา ภณิตํ, ตํ สพฺพํ ตตฺถํ อุทาหุ วิตถนฺติ.}

\prose{ยํ ตฺวํ มหาราช พหุวิธํ การณํ โอสาเรสิ, สพฺพํ ตํ ตเถว โน อญฺญถาติ. ยทิ ภนฺเต นาคเสน กโณฺหปิ สุกฺโกปิ สมสมคติกา โหนฺติ, เตน หิ กุสลมฺปิ อกุสลมฺปิ สมสม\-วิปากํ โหตีติ. น หิ มหาราช กุสลมฺปิ อกุสลมฺปิ สมสม\-วิปากํ โหติ, น หิ มหาราช เทวทตฺโต สพฺพชเนหิ ปฏิวิรุทฺโธ, โพธิสตฺเตเนว ปฏิวิรุทฺโธ. โย ตสฺส โพธิสตฺเตน ปฏิวิรุทฺโธ, โส ตสฺมิํ ตสฺมิํ เยว ภเว ปจฺจติ ผลํ เทติ. เทวทตฺโตปิ มหาราช อิสฺสริเย ฐิโต ชนปเทสุ อารกฺขํ เทติ, เสตุํ สภํ ปุญฺญสาลํ กาเรติ, สมณ\-พฺราหฺมณานํ กปณทฺธิก\-วณิพฺพกานํ นาถานาถานํ ยถาปณิหิตํ ทานํ เทติ. ตสฺส โส วิปาเกน ภเว ภเว สมฺปตฺติโย ปฏิลภติ. กสฺเสตํ มหาราช สกฺกา วตฺตุํ วินา ทาเนน ทเมน สํยเมน อุโปสถ\-กมฺเมน สมฺปตฺติํ อนุภวิสฺสตีติ.}

\prose{ยํ ปน ตฺวํ มหาราช เอวํ วเทสิ “เทวทตฺโต จ โพธิสตฺโต จ เอกโต อนุปริ\-วตฺตนฺตี”ติ, โส น ชาติสตสฺส อจฺจเยน สมาคโม อโหสิ, น ชาติสหสฺสสฺส อจฺจเยน, น ชาติ\-สต\-สหสฺสสฺส อจฺจเยน, กทาจิ กรหจิ พหูนํ อโหรตฺตานํ อจฺจเยน สมาคโม อโหสิ. ยํ ปเนตํ มหาราช ภควตา\csromansymbolfootnote{*}{สํ 3. 397 ปิฏฺเฐ.}กาณกจฺฉโปปมํ อุปทสฺสิตํ มนุสฺสตฺต\-ปฺปฏิลาภาย, ตถูปมํ มหาราช อิเมสํ สมาคมํ ธาเรหิ.}

\prose{น มหาราช โพธิสตฺตสฺส เทวทตฺเตเนว สทฺธิํ สมาคโม อโหสิ, เถโรปิ มหาราช สาริปุตฺโต อเนเกสุ ชาติ\-สต\-สหสฺเสสุ โพธิสตฺตสฺส ปิตา อโหสิ, มหาปิตา อโหสิ, จูฬปิตา}

\prose{\csromanfolio{202}อโหสิ, ภาตา อโหสิ, ปุตฺโต อโหสิ, ภาคิเนยฺโย อโหสิ, มิตฺโต อโหสิ.}

\prosewithcloserruled{โพธิสตฺโตปิ มหาราช อเนเกสุ ชาติ\-สต\-สหสฺเสสุ เถรสฺส สาริปุตฺตสฺส ปิตา อโหสิ, มหาปิตา อโหสิ, จูฬปิตา อโหสิ, ภาตา อโหสิ, ปุตฺโต อโหสิ, ภาคิเนยฺโย อโหสิ, มิตฺโต อโหสิ, สพฺเพปิ มหาราช สตฺตนิกาย\-ปริยาปนฺนา สํสารโสตม\-นุคตา สํสารโสเตน วุยฺหนฺตา อปฺปิเยหิปิ ปิเยหิปิ สมาคจฺฉนฺติ. ยถา มหาราช อุทกํ โสเตน วุยฺหมานํ สุจิ- อสุจิกลฺยาณปาปเกน สมาคจฺฉติ, เอวเมว โข มหาราช สพฺเพปิ สตฺตนิกาย\-ปริยาปนฺนา สํสารโสตม\-นุคตา สํสารโสเตน วุยฺหนฺตา อปฺปิเยหิปิ ปิเยหิปิ สมาคจฺฉนฺติ.\csromansymbolfootnote{+}{ขุ 5. 229 ปิฏฺเฐ.}เทวทตฺโต มหาราช ยกฺโข สมาโน อตฺตนา อธมฺโม ปเร อธมฺเม นิโยเชตฺวา สตฺตปญฺญาส\-วสฺสโกฏิโย สฏฺฐิ จ วสฺส\-สตสหสฺสานิ มหานิรเย ปจฺจิ, โพธิสตฺโตปิ มหาราช ยกฺโข สมาโน อตฺตนา ธมฺโม ปเร ธมฺเม นิโยเชตฺวา สตฺตปญฺญาส\-วสฺสโกฏิโย สฏฺฐิ จ วสฺส\-สตสหสฺสานิ สคฺเค โมทิ สพฺพกาม\-สมงฺคี, อปิ จ มหาราช เทวทตฺโต อิมสฺมิํ ภเว พุทฺธํ อนาสาทนียมาสาทยิตฺวา สมคฺคญฺจ สํฆํ ภินฺทิตฺวา ปถวิํ ปาวิสิ, ตถาคโต พุชฺฌิตฺวา สพฺพธมฺเม ปรินิพฺพุโต อุปธิ\-สงฺขเยติ. สาธุ ภนฺเต นาคเสน, เอวเมตํ ตถา สมฺปฏิจฺฉามีติ.}{กุสลากุสล\-สมวิสม\-ปโญฺห สตฺตโม.}
\csromanmatikaanchor{mtk.147}
\csromantocmarkref{subsubsection}{8. อมราเทวีปญฺห}{mtk.147}
\bookmark[dest={mtk.147},level=3]{8. อมราเทวีปญฺห}
\csromanheader{3}{8. \textbf{อมราเทวีปญฺห}}
\prose{8. ภนฺเต นาคเสน ภาสิตมฺเปตํ ภควตา\csromandash{}}

\setlength{\csromangathapagewidth}{0pt}
\setlength{\csromangathawakwidth}{0pt}
\csromangathameasuregroup{}{“สเจ ลเภถ ขณํ วา รโห วา, \\ นิมนฺตกํ\textsuperscript{1} วาปิ ลเภถ ตาทิสํ. \\ สพฺพาว\textsuperscript{1} อิตฺถี กยิรุํ\textsuperscript{1} นุ ปาปํ, \\ อญฺญํ อลทฺธา ปีฐสปฺปินา สทฺธินฺ”ติ.}
\csromangathameasurewak{“สเจ ลเภถ ขณํ วา รโห วา, \\ นิมนฺตกํ\textsuperscript{1} วาปิ ลเภถ ตาทิสํ. \\ สพฺพาว\textsuperscript{1} อิตฺถี กยิรุํ\textsuperscript{1} นุ ปาปํ, \\ อญฺญํ อลทฺธา ปีฐสปฺปินา สทฺธินฺ”ติ.}
\setlength{\csromangathaleftcol}{0pt}
\csromangathagroup{\csromangathastanza{\csromangathawak{“สเจ ลเภถ ขณํ วา รโห วา, \\ นิมนฺตกํ\csromansharedfootnote{784a45c58e8e9f78}{นิวาตกํ (ขุ 6. 123 ปิฏฺเฐ กุณาลชาตเก)} วาปิ ลเภถ ตาทิสํ. \\ สพฺพาว\csromansharedfootnote{313d2b77f9f8c4d3}{สพฺพาปิ (สี, อิ)} อิตฺถี กยิรุํ\csromansharedfootnote{6cae72d9479c77c5}{กเรยฺยุํ (สี, อิ, ก)} นุ ปาปํ, \\ อญฺญํ อลทฺธา ปีฐสปฺปินา สทฺธินฺ”ติ.}}}

\prose{ปุน จ กถียติ\csromansymbolfootnote{*}{ชาตกฏฺฐ 6. 210 ปิฏฺเฐ.}“มโหสธสฺส ภริยา อมรา นาม อิตฺถี คามเก ฐปิตา ปวุตฺถปติกา รโห นิสินฺนา วิวิตฺตา ราชปฺปฏิสมํ สามิกํ กริตฺวา \csromanfolio{203}สหสฺเสน นิมนฺตียมานา ปาปํ นากาสี”ติ. ยทิ ภนฺเต นาคเสน ภควตา ภณิตํ “สเจ ฯเปฯ สทฺธินฺ”ติ. เตน หิ “มโหสธสฺส ภริยา ฯเปฯ นากาสี”ติ ยํ วจนํ, ตํ มิจฺฉา. ยทิ มโหสธสฺส ภริยา ฯเปฯ นากาสิ, เตน หิ “สเจ ฯเปฯ สทฺธินฺ”ติ ตมฺปิ วจนํ มิจฺฉา. อยมฺปิ อุภโต โกฏิโก ปโญฺห ตวานุปฺปตฺโต, โส ตยา นิพฺพาหิตพฺโพติ.}

\proseitem{7}{ภาสิตมฺเปตํ มหาราช ภควตา “สเจ ฯเปฯ สทฺธินฺ”ติ. กถียติ จ “มโหสธสฺส ภริยา ฯเปฯ นากาสี”ติ. กเรยฺย สา มหาราช อิตฺถี สหสฺสํ ลภมานา ตาทิเสน ปุริเสน สทฺธิํ ปาปกมฺมํ, น สา กเรยฺย สเจ ขณํ วา รโห วา นิมนฺตกํ วาปิ ตาทิสํ ลเภยฺย, วิจินนฺตี สา มหาราช อมรา อิตฺถี น อทฺทส ขณํ วา รโห วา นิมนฺตกํ วาปิ ตาทิสํ.}

\prose{อิธ โลเก ครหภยา ขณํ น ปสฺสิ, ปรโลเก นิรยภยา ขณํ น ปสฺสิ, กฏุกวิปากํ ปาปนฺติ ขณํ น ปสฺสิ, ปิยํ อมุญฺจิตุกามา ขณํ น ปสฺสิ, สามิกสฺส ครุกตาย ขณํ น ปสฺสิ, ธมฺมํ อปจายนฺตี ขณํ น ปสฺสิ, อนริยํ ครหนฺตี ขณํ น ปสฺสิ, กิริยํ อภินฺทิตุกามา ขณํ น ปสฺสิ. เอวรูเปหิ พหูหิ การเณหิ ขณํ น ปสฺสิ.}

\prose{รโหปิ สา โลเก วิจินิตฺวา อปสฺสนฺตี ปาปํ นากาสิ. สเจ สา มนุสฺเสหิ รโห ลเภยฺย, อถ อมนุสฺเสหิ รโห น ลเภยฺย. สเจ อมนุสฺเสหิ รโห ลเภยฺย, อถ ปรจิตฺต\-วิทูหิ ปพฺพชิเตหิ รโห น ลเภยฺย. สเจ ปรจิตฺต\-วิทูหิ ปพฺพชิเตหิ รโห ลเภยฺย, อถ ปรจิตฺตวิทูนีหิ เทวตาหิ รโห น ลเภยฺย. สเจ ปรจิตฺตวิทูนีหิ เทวตาหิ รโห ลเภยฺย, อตฺตนาว ปาเปหิ รโห น ลเภยฺย. สเจ อตฺตนาว ปาเปหิ รโห ลเภยฺย, อถ อธมฺเมน รโห น ลเภยฺย. เอวรูเปหิ พหุวิเธหิ การเณหิ รโห อลภิตฺวา ปาปํ นากาสิ.}

\prosewithcloserruled{นิมนฺตกมฺปิ สา โลเก วิจินิตฺวา ตาทิสํ อลภนฺตี ปาปํ นากาสิ. มโหสโธ มหาราช ปณฺฑิโต อฏฺฐวีสติยา องฺเคหิ สมนฺนาคโต. กตเมหิ อฏฺฐวีสติยา องฺเคหิ สมนฺนาคโต, มโหสโธ มหาราช สูโร หิริมา โอตฺตปฺปี สปกฺโข มิตฺตสมฺปนฺโน ขโม สีลวา สจฺจวาที โสเจยฺย\-สมฺปนฺโน อกฺโกธโน อนติมานี อนุสูยโก วีริยวา อายูหโก สงฺคาหโก สํวิภาคี สขิโล นิวาตวุตฺติ สโณฺห อสโฐ อมายาวี อติพุทฺธิสมฺปนฺโน กิตฺติมา วิชฺชาสมฺปนฺโน หิเตสี \csromanfolio{204}อุปนิสฺสิตานํ ปตฺถิโต สพฺพชนสฺส ธนวา ยสวา. มโหสโธ มหาราช ปณฺฑิโต อิเมหิ อฏฺฐวีสติยา องฺเคหิ สมนฺนาคโต, สา อญฺญํ ตาทิสํ นิมนฺตกํ อลภิตฺวา ปาปํ นากาสีติ. สาธุ ภนฺเต นาคเสน, เอวเมตํ ตถา สมฺปฏิจฺฉามีติ.}{อมราเทวีปโญฺห อฏฺฐโม.}
\csromanmatikaanchor{mtk.148}
\csromantocmarkref{subsubsection}{9. อรหนฺตอภายนปญฺห}{mtk.148}
\bookmark[dest={mtk.148},level=3]{9. อรหนฺตอภายนปญฺห}
\csromanheader{3}{9. อรหนฺต\-อภายน\-ปญฺห}
\proseitem{9}{ภนฺเต นาคเสน ภาสิตมฺเปตํ ภควตา “วิคตภย\-สนฺตาสา อรหนฺโต”ติ.\csromansymbolfootnote{*}{วิ 4. 356 ปิฏฺเฐ.}ปุน จ นาคเร ราชคเห ธนปาลกํ หตฺถิํ ภควติ โอปตนฺตํ ทิสฺวา ปญฺจ ขีณาสว\-สตานิ ปริจฺจชิตฺวา ชินวรํ ปกฺกนฺตานิ ทิสาวิทิสํ เอกํ ฐเปตฺวา เถรํ อานนฺทํ. กินฺนุ โข ภนฺเต นาคเสน เต อรหนฺโต ภยา ปกฺกนฺตา, ปญฺญายิสฺสติ สเกน กมฺเมนาติ ทสพลํ ปาเตตุกามา ปกฺกนฺตา, อุทาหุ ตถาคตสฺส อตุลํ วิปุลมสมํ ปาฏิหาริยํ ทฏฺฐุกามา ปกฺกนฺตา. ยทิ ภนฺเต นาคเสน ภควตา ภณิตํ “วิคตภย\-สนฺตาสา อรหนฺโต”ติ, เตน หิ “นคเร ฯเปฯ อานนฺทนฺ”ติ ยํ วจนํ ตํ มิจฺฉา. ยทิ นาคเร ราชคเห ธนปาลกํ หตฺถํ ภควติ โอปตนฺตํ ทิสฺวา ปญฺจ ขีณาสว\-สตานิ ปริจฺจชิตฺวา ชินวรํ ปกฺกนฺตานิ ทิสาวิทิสํ เอกํ ฐเปตฺวา เถรํ อานนฺทํ, เตน หิ “วิคตภย\-สนฺตาสา อรหนฺโต”ติ ตมฺปิ วจนํ มิจฺฉา. อยมฺปิ อุภโต โกฏิโก ปโญฺห ตวานุปฺปตฺโต, โส ตยา นิพฺพาหิตพฺโพติ.}

\prose{ภาสิตมฺเปตํ มหาราช ภควตา “วิคตภย\-สนฺตาสา อรหนฺโต”ติ, นคเร ราชคเห ธนปาลกํ หตฺถิํ ภควติ โอปตนฺตํ ทิสฺวา ปญฺจ ขีณาสว\-สตานิ ปริจฺจชิตฺวา ชินวรํ ปกฺกนฺตานิ ทิสาวิทิสํ เอกํ ฐเปตฺวา เถรํ อานนฺทํ, ตญฺจ ปน น ภยา, นาปิ ภควนฺตํ ปาเตตุกามตาย.}

\prose{เยน ปน มหาราช เหตุนา อรหนฺโต ภาเยยฺยุํ วา ตาเสยฺยุํ วา, โส เหตุ อรหนฺตานํ สมุจฺฉินฺโน, ตสฺมา วิคตภย\-สนฺตาสา อรหนฺโต, ภายติ นุ มหาราช มหาปถวี ขณนฺเตปิ ภินฺทนฺเตปิ ธาเรนฺเตปิ สมุทฺทปพฺพตคิริสิขเรติ. น หิ ภนฺเตติ. เกน การเณน มหาราชาติ. นตฺถิ ภนฺเต มหาปถวิยา โส เหตุ, เยน เหตุนา \csromanfolio{205}มหาปถวี ภาเยยฺย วา ตาเสยฺย วาติ, เอวเมว โข มหาราช นตฺถิ อรห- อินฺตานํ โส เหตุ, เยน เหตุนา อรหนฺโต ภาเยยฺยุํ วา ตาเสยฺยุํ วา.}

\prose{ภายติ นุ มหาราช คิริสิขรํ ฉินฺทนฺเต วา ภินฺทนฺเต วา ปตนฺเต วา อคฺคินา ทหนฺเต วาติ. น หิ ภนฺเตติ. เกน การเณน มหาราชาติ. นตฺติ ภนฺเต คิริสิขรสฺส โส เหตุ, เยน เหตุนา คิริสิขรํ ภาเยยฺย วา ตาเสยฺย วาติ, เอวเมว โข มหาราช นตฺถิ อรหนฺตานํ โส เหตุ, เยน เหตุนา อรหนฺโต ภาเยยฺยุํ วา ตาเสยฺยุํ วา.}

\prose{ยทิปิ มหาราช โลกธาตุ\-สตสหสฺเสสุ เย เกจิ สตฺตนิกาย\-ปริยาปนฺนา สพฺเพปิ เต สตฺติหตฺถา เอกํ อรหนฺตํ อุปธาวิตฺวา ตาเสยฺยุํ, น ภเวยฺย อรหโต จิตฺตสฺส กิญฺจิ อญฺญถตฺตํ, กิํ การณํ, อฏฺฐานม\-นวกาสตาย.}

\prose{อปิ จ มหาราช เตสํ ขีณาสวานํ เอวํ เจโต\-ปริวิตกฺโก อโหสิ “อชฺช นรวรปวเร ชินวร\-วสเภ นครวรม\-นุปฺปวิฏฺเฐ วีถิยา ธนปาลโก หตฺถี อาปติสฺสติ, อสํส\-ยมติ\-เทวเทวํ อุปฏฺฐาโก น ปริจฺจชิสฺสติ, ยทิ มยํ สพฺเพปิ ภควนฺตํ น ปริจฺจชิสฺสาม, อานนฺทสฺส คุโณ ปากโฏ น ภวิสฺสติ, น เหว จ ตถาคตํ สมุปคมิสฺสติ หตฺถินาโค, หนฺท มยํ อาปคจฺฉาม, เอวมิทํ มหโต ชนกายสฺส กิเลส\-พนฺธนโมกฺโข ภวิสฺสติ, อานนฺทสฺส จ คุโณ ปากโฏ ภวิสฺสตี”ติ. เอวํ เต อรหนฺโต อานิสํสํ ทิสฺวา ทิสาวิทิสํ ปกฺกนฺตาติ. สุวิภตฺโต ภนฺเต นาคเสน ปโญฺห, เอวเมตํ นตฺถิ อรหนฺตานํ ภยํ วา สนฺตาโส วา, อานิสํสํ ทิสฺวา อรหนฺโต ปกฺกนฺตา ทิสาวิทิสนฺติ.}

\prose{อรหนฺต\-อภายน\-ปโญฺห นวโม.}
\csromansectionrule

\csromanmatikaanchor{mtk.149}
\csromantocmarkref{subsubsection}{10. พุทฺธสพฺพญฺญุภาวปญฺห}{mtk.149}
\bookmark[dest={mtk.149},level=3]{10. พุทฺธสพฺพญฺญุภาวปญฺห}
\csromanheader{3}{10. \textbf{พุทฺธ}\-\textbf{สพฺพญฺญุ}\-\textbf{ภาว}\-\textbf{ปญฺห}}
\proseitem{10}{ภนฺเต นาคเสน ตุเมฺห ภณถ\csromansymbolfootnote{*}{ขุ 7. 139; ขุ 8. 176; ขุ 9. 376 ปิฏฺเฐสุ.}“ตถาคโต สพฺพญฺญู”ติ. ปุน จ ภณถ\csromansymbolfootnote{+}{ม 2. 120 ปิฏฺฐาทีสุ.}“ตถาคเตน สาริปุตฺต\-โมคฺคลฺลาน\-ปฺปมุเข ภิกฺขุสํเฆ \csromanfolio{206}ปณามิเต จาตุเมยฺยกา จ สกฺยา พฺรหฺมา จ สหมฺปติ พีชูปมญฺจ วจฺฉ\-ตรุณูปมญฺจ อุปทสฺเสตฺวา ภควนฺตํ ปสาเทสุํ ขมาเปสุํ นิชฺฌตฺตํ อกํสู”ติ. กินฺนุ โข ภนฺเต นาคเสน อญฺญาตา ตา อุปมา ตถาคตสฺส, ยาหิ ตถาคโต อุปมาหิ โอรโต ขมิโต อุปสนฺโต นิชฺฌตฺตํ คโต. ยทิ ภนฺเต นาคเสน ตถาคตสฺส ตา อุปมา อญฺญาตา, เตน หิ พุทฺโธ อสพฺพญฺญู, ยทิ ญาตา, เตน หิ โอกสฺส ปสยฺห วีมํสาเปกฺโข ปณาเมสิ, เตน หิ ตสฺส อการุญฺญตา สมฺภวติ. อยมฺปิ อุภโต โกฏิโก ปโญฺห ตวานุปฺปตฺโต, โส ตยา นิพฺพาหิตพฺโพติ.}

\prose{สพฺพญฺญุ มหาราช ตถาคโต, ตาหิ จ อุปมาหิ ภควา ปสนฺโน โอรโต ขมิโต อุปสนฺโต นิชฺฌตฺตํ คโต, ธมฺมสฺสามี มหาราช ตถาคโต, ตถาคตปฺปเวทิเตเหว เต โอปมฺเมหิ ตถาคตํ อาราเธสุํ โตเสสุํ ปสาเทสุํ, เตสญฺจ ตถาคโต ปสนฺโน “สาธู”ติ อพฺภานุโมทิ.}

\prose{ยถา มหาราช อิตฺถี สามิกสฺส สนฺตเกเนว ธเนน สามิกํ อาราเธติ โตเสติ ปสาเทติ, ตญฺจ สามิโก “สาธู”ติ อพฺภานุโมทติ, เอวเมว โข มหาราช จาตุเมยฺยกา จ สกฺยา พฺรหฺมา จ สหมฺปติ ตถาคตปฺปเวทิเตเหว โอปมฺเมหิ ตถาคตํ อาราเธสุํ โตเสสุํ ปสาเทสุํ, เตสญฺจ ตถาคโต ปสนฺโน “สาธู”ติ อพฺภานุโมทิ.}

\prose{ยถา วา ปน มหาราช กปฺปโก รญฺโญ สนฺตเกเนว สุวณฺณ\-ผณเกน รญฺโญ อุตฺตมงฺคํ ปสาธยมาโน ราชานํ อาราเธติ โตเสติ ปสาเทติ, ตสฺส จ ราชา ปสนฺโน “สาธู”ติ อพฺภานุโมทติ, ยถิจฺฉิตม\-นุปฺปเทติ, เอวเมว โข มหาราช จาตุเมยฺยกา จ สกฺยา พฺรหฺมา จ สหมฺปติ ตถาคตปฺปเวทิเตเหว โอปมฺเมหิ ตถาคตํ อาราเธสุํ โตเสสุํ ปสาเทสุํ, เตสญฺจ ตถาคโต ปสนฺโน “สาธุ”ติ อพฺภานุโมทิ.}

\prose{ยถา วา ปน มหาราช สทฺธิวิหาริโก อุปชฺฌายาภตํ ปิณฺฑปาตํ คเหตฺวา อุปชฺฌายสฺส อุปนาเมนฺโต อุปชฺฌายํ อาราเธติ โตเสติ ปสาเทติ, ตญฺจ อุปชฺฌาโย ปสนฺโน “สาธู”ติ อพฺภานุโมทติ, \csromanfolio{207}เอวเมว โข มหาราช จาตุเมยฺยกา จ สกฺยา พฺรหฺมา จ สหมฺปติ ตถาคตปฺปเวทิเตเหว โอปมฺเมหิ ตถาคตํ อาราเธสุํ โตเสสุํ ปสาเทสุํ, เตสญฺจ ตถาคโต ปสนฺโน “สาธู”ติ อพฺภานุโมทิตฺวา สพฺพทุกฺข\-ปริมุตฺติยา ธมฺมํ เทเสสีติ. สาธุ ภนฺเต นาคเสน, เอวเมตํ ตถา สมฺปฏิจฺฉามีติ.}

\vspace{\tipitakaparskip}
\csromancenter{พุทฺธ\-สพฺพญฺญุ\-ภาว\-ปโญฺห ทสโม.}
\nitthitammidruled{\textbf{สพฺพญฺญุตญาณ}\-\textbf{วคฺโค จตุตฺโถ.}}
\csromancenter{อิมสฺมิํ วคฺเค ทส ปญฺหา.}
\csromanmatikaanchor{mtk.150}
\csromantocmarknopage{subsection}{5. สนฺถววคฺค}{mtk.150}
\bookmark[dest={mtk.150},level=2]{5. สนฺถววคฺค}
\csromanheader{2}{5. \textbf{สนฺถววคฺค}}
\csromanmatikaanchor{mtk.151}
\csromantocmarkref{subsubsection}{1. สนฺถวปญฺห}{mtk.151}
\bookmark[dest={mtk.151},level=3]{1. สนฺถวปญฺห}
\csromanheader{3}{1. \textbf{สนฺถวปญฺห}}
\prose{\csromanfolio{208}1. ภนฺเต นาคเสน ภาสิตมฺเปตํ ภควตา\csromandash{}}

\setlength{\csromangathapagewidth}{0pt}
\setlength{\csromangathawakwidth}{0pt}
\csromangathameasuregroup{\textsuperscript{*}“สนฺถวโต ภยํ ชาตํ, \\ อนิเกตมสนฺถวํ,}{\csromangathabat{\textsuperscript{*}“สนฺถวโต ภยํ ชาตํ,}{นิเกตา ชายเต รโช.} \\ \csromangathabat{อนิเกตมสนฺถวํ,}{เอตํ เว มุนิทสฺสนนฺ”ติ.}}
\csromangathasetleft{\textsuperscript{*}“สนฺถวโต ภยํ ชาตํ, \\ อนิเกตมสนฺถวํ,}
\csromangathagroup{\csromangathastanza{\csromangathabat{\csromansymbolfootnote{*}{ขุ 1. 309 ปิฏฺเฐ สุตฺตนิปาเต.}“สนฺถวโต ภยํ ชาตํ,}{นิเกตา ชายเต รโช.} \\ \csromangathabat{อนิเกตม\-สนฺถวํ,}{เอตํ เว มุนิทสฺสนนฺ”ติ.}}}

\prose{ปุน จ ภควตา ภณิตํ\csromansymbolfootnote{+}{วิ 4. 314 ปิฏฺเฐ.}“วิหาเร การเย รมฺเม, วาสเยตฺถ พหุสฺสุเต”ติ. ยทิ ภนฺเต นาคเสน ตถาคเตน ภณิตํ “สนฺถวโต ภยํ ชาตํ, นิเกตา ชายเต รโช. อนิเกตม\-สนฺถวํ, เอตํ เว มุนิทสฺสนนฺ”ติ. เตน หิ “วิหาเร การเย รมฺเม, วาสเยตฺถ พหุสฺสุเต”ติ ยํ วจนํ, ตํ มิจฺฉา. ยทิ ตถาคเตน ภณิตํ “วิหาเร การเย รมฺเม, วาสเยตฺถ พหุสฺสุเต”ติ, เตน หิ “สนฺถวโต ภยํ ชาตํ, นิเกตา ชายเต รโช. อนิเกตม\-สนฺถวํ, เอตํ เว มุนิทสฺสนนฺ”ติ ตมฺปิ วจนํ มิจฺฉา. อยมฺปิ อุภโต โกฏิโก ปโญฺห ตวานุปฺปตฺโต, โส ตยา นิพฺพาหิตพฺโพติ.}

\prose{ภาสิตมฺเปตํ มหาราช ภควตา “สนฺถวโต ภยํ ชาตํ, นิเกตา ชายเต รโช. อนิเกตม\-สนฺถวํ, เอตํ เว มุนิทสฺสนนฺ”ติ, ภณิตญฺจ “วิหาเร การเย รมฺเม, วาสเยตฺถ พหุสฺสุเต”ติ. ยํ มหาราช ภควตา ภณิตํ “สนฺถวโต ภยํ ชาตํ, นิเกตา ชายเต รโช. อนิเกตม\-สนฺถวํ, เอตํ เว มุนิทสฺสนนฺ”ติ, ตํ สภาววจนํ อเสสวจนํ นิสฺเสสวจนํ นิปฺปริยาย\-วจนํ สมณา\-นุจฺฉวํ สมณสารุปฺปํ สมณ\-ปฺปติรูปํ สมณารหํ สมณโคจรํ สมณ\-ปฺปฏิปทา สมณ\-ปฺปฏิปตฺติ. ยถา มหาราช อารญฺญโก มิโค อรญฺเญ ปวเน จรมาโน นิรลโย อนิเกโต ยถิจฺฉกํ สยติ, เอวเมว โข มหารช ภิกฺขุนา “สนฺถวโต ภยํ ชาตํ, นิเกตา ชายเต รโช. อนิเกตม\-สนฺถวํ, เอตํ เว มุนิทสฺสนนฺ”ติ จินฺเตตพฺพํ.}

\prose{ยํ ปน มหาราช ภควตา ภณิตํ “วิหาเร การเย รมฺเม, วาสเยตฺถ พหุสฺสุเต”ติ, ตํ เทฺว อตฺถวเส สมฺปสฺสามาเนน ภควตา ภณิตํ. กตเม เทฺว, วิหารทานํ นาม สพฺพพุทฺเธหิ วณฺณิตํ อนุมตํ \csromanfolio{209}โถมิตํ ปสตฺถํ, ตํ เต วิหารทานํ ทตฺวา ชาติชรามรณา ปริมุจฺจิสฺสนฺตีติ. อยํ ตาว ปฐโม อานิสํโส วิหารทาเน.}

\proseitemwithcloserruled{10}{ปุน จปรํ วิหเร วิชฺชมาเน ภิกฺขุนิโย พฺยตฺตสงฺเกตา ภวิสฺสนฺติ, สุลภํ ทสฺสนํ ทสฺสนกามานํ, อนิเกเต ทุทฺทสฺสนา ภวิสฺสนฺตีติ. อยํ ทุติโย อานิสํโส วิหารทาเน. อิเม เทฺว อตฺถวเส สมฺปสฺสมาเนน ภควตา ภณิตํ “วิหาเร การเย รมฺเม, วาสเยตฺถ พหุสฺสุเต”ติ, น ตตฺถ พุทฺธปุตฺเตน อาลโย กรณีโย นิเกเตติ. สาธุ ภนฺเต นาคเสน, เอวเมตํ ตถา สมฺปฏิจฺฉามีติ.}{สนฺถวปโญฺห ปฐโม.}
\csromanmatikaanchor{mtk.152}
\csromantocmarkref{subsubsection}{2. อุทรสํยตปญฺห}{mtk.152}
\bookmark[dest={mtk.152},level=3]{2. อุทรสํยตปญฺห}
\csromanheader{3}{2. \textbf{อุทร}\-\textbf{สํยต}\-\textbf{ปญฺห}}
\prose{2. ภนฺเต นาคเสน ภาสิตมฺเปตํ ภควตา\csromandash{}}

\setlength{\csromangathapagewidth}{0pt}
\setlength{\csromangathawakwidth}{0pt}
\csromangathameasuregroup{\textsuperscript{*}“อุตฺติฏฺเฐ นปฺปมชฺเชยฺย,}{\csromangathabat{\textsuperscript{*}“อุตฺติฏฺเฐ นปฺปมชฺเชยฺย,}{อุทเร สํยโต สิยา”ติ.}}
\csromangathasetleft{\textsuperscript{*}“อุตฺติฏฺเฐ นปฺปมชฺเชยฺย,}
\csromangathagroup{\csromangathastanza{\csromangathabat{\csromansymbolfootnote{*}{ขุ 1. 38 ปิฏฺเฐ ธมฺมปเท.}“อุตฺติฏฺเฐ นปฺปมชฺเชยฺย,}{อุทเร สํยโต สิยา”ติ.}}}

\prose{ปุน จ ภควตา ภณิตํ\csromansymbolfootnote{+}{ม 2. 199 ปิฏฺเฐ.}“อหํ โข ปนุทายิ อปฺเปกทา อิมินา ปตฺเตน สมติตฺติกมฺปิ ภุญฺชามิ, ภิยฺโยปิ ภุญฺชามี”ติ. ยทิ ภนฺเต นาคเสน ภควตา ภณิตํ “อุตฺติฏฺเฐ นปฺปมชฺเชยฺย, อุทเร สํยโต สิยา”ติ, เตน หิ “อหํ โข ปนุทายิ อปฺเปกทา อิมินา ปตฺเตน สมติตฺถิ\-กมฺปิ ภุญฺชามิ, ภิยฺโยปิ ภุญฺชามี”ติ ยํ วจนํ, ตํ มิจฺฉา. ยทิ ตถาคเตน ภิณิตํ “อหํ โข ปนุทายิ อปฺเปกทา อิมินา ปตฺเตน สมติตฺถิ\-กมฺปิ ภุญฺชามี”ติ, เตน หิ “อุตฺติฏฺเฐ นปฺปมชฺเชยฺย, อุทเร สํยโต สิยา”ติ ตมฺปิ วจนํ มิจฺฉา. อยมฺปิ อุภโต โกฏิโก ปโญฺห ตวานุปฺปตฺโต, โส ตยา นิพฺพาหิตพฺโพติ.}

\prose{ภาสิตมฺเปตํ มหาราช ภควตา “อุตฺติฏฺเฐ นปฺปมชฺเชยฺย, อุทเร สํยโต สิยา”ติ, ภณิตญฺจ “อหํ โข ปนุทายิ อปฺเปกทา อิมินา ปตฺเตน สมติตฺติกมฺปิ ภุญฺชามิ, ภิยฺโยปิ ภุญฺชามี”ติ. ยํ มหาราช ภควตา ภณิตํ “อุตฺติฏฺเฐ นปฺปมชฺเชยฺย, อุทเร สํยโต สิยา”ติ, ตํ สภาววจนํ อเสสวจนํ นิสฺเสสวจนํ นิปฺปริยาย\-วจนํ ภูตวจนํ ตจฺฉวจนํ ยาถาววจนํ อวิปรีต\-วจนํ อิสิวจนํ มุนิวจนํ ภควนฺตาวจนํ}

\prose{\csromanfolio{210}อรหนฺต\-วจนํ ปจฺเจก\-พุทฺธวจนํ ชินวจนํ สพฺพญฺญุ\-วจนํ ตถาคตสฺส อรหโต สมฺมา\-สมฺพุทฺธสฺส วจนํ.}

\prose{อุทเร อสํยโต มหาราช ปาณมฺปิ หนติ, อทินฺนมฺปิ อาทิยติ, ปรทารมฺปิ คจฺฉติ, มุสาปิ ภณติ, มชฺชมฺปิ ปิวติ, มาตรมฺปิ ชีวิตา โวโรเปติ, ปิตรมฺปิ ชีวิตา โวโรเปติ, อรหนฺตมฺปิ ชีวิตา โวโรเปติ, สํฆมฺปิ ภินฺทติ, ทุฏฺเฐน จิตฺเตน ตถาคตสฺส โลหิตมฺปิ อุปฺปาเทติ. นนุ มหาราช เทวทตฺโต อุทเร อสํยโต สํฆํ ภินฺทิตฺวา กปฺปฏฺฐิยํ กมฺมํ อายูหิ\csromansharedfootnote{f4049db0db82cb28}{อายูหติ (ก)}. เอวรูปานิ มหาราช อญฺญานิปิ พหุวิธานิ การณานิ ทิสฺวา ภควตา ภณิตํ “อุตฺติฏฺเฐ นปฺปมชฺเชยฺย, อุทเร สํยโต สิยา”ติ.}

\prose{อุทเร สํยโต มหาราช จตุสจฺจาภิสมยํ อภิสเมติ, จตฺตาริ สามญฺญผลานิ สจฺฉิกโรติ, จตูสุ ปฏิสมฺภิทาสุ อฏฺฐสุ สมาปตฺตีสุ ฉสุ อภิญฺญาสุ วสีภาวํ ปาปุณาติ, เกวลญฺจ สมณธมฺมํ ปูเรติ. นนุ มหาราช สุกโปตโก อุทเร สํยโต หุตฺวา ยาว ตาวติํส\-ภวนํ กมฺเปตฺวา สกฺกํ เทวานมินฺทํ อุปฏฺฐานมุ\-ปเนสิ, เอวรูปานิ มหาราช อญฺญานิปิ พหุวิธานิ การณานิ ทิสฺวา ภควตา ภณิตํ “อุตฺติฏฺเฐ นปฺปมชฺเชยฺย, อุทเร สํยโต สิยา”ติ.}

\prose{ยํ ปน มหาราช ภควตา ภณิตํ “อหํ โข ปนุทายิ อปฺเปกทา อิมินา ปตฺเตน สมติตฺติกมฺปิ ภุญฺชามิ, ภิยฺโยปิ ภุญฺชามี”ติ, ตํ กตกิจฺเจน นิฏฺฐิต\-กิริเยน สิทฺธตฺเถน วุสิตโวสาเนน นิราวรเณน สพฺพญฺญุนา สยมฺภุนา ตถาคเตน อตฺตานํ อุปาทาย ภณิตํ.}

\prose{ยถา มหาราช วนฺตสฺส วิริตฺตสฺส อนุวาสิตสฺส อาตุรสฺส สปฺปายกิริยา อิจฺฉิตพฺพา โหติ, เอวเมว โข มหาราช สกิเลสสฺส อทิฏฺฐสจฺจสฺส อุทเร สํยโม กรณีโย โหติ. ยถา มหาราช มณิรตนสฺส สปฺปภาสสฺส ชาติมนฺตสฺส อภิชาติ\-ปริสุทฺธสฺส มชฺชน\-นิฆํสน\-ปริโสธเนน กรณียํ น โหติ, เอวเมว โข มหาราช ตถาคตสฺส พุทฺธวิสเย ปารมิํ คตสฺส กิริยากรเณสุ อาวรณํ น โหตีติ. สาธุ ภนฺเต นาคเสน, เอวเมตํ ตถา สมฺปฏิจฺฉามีติ.}

\vspace{\tipitakaparskip}
\csromancenter{อุทร\-สํยต\-ปโญฺห ทุติโย.}
\csromanmatikaanchor{mtk.153}
\csromantocmarkref{subsubsection}{3. พุทฺธอปฺปาพาธปญฺห}{mtk.153}
\bookmark[dest={mtk.153},level=3]{3. พุทฺธอปฺปาพาธปญฺห}
\csromanmarkh{4. เมณฺฑกปญฺห}\csromanmarkhii{3. พุทฺธ\-อปฺปาพาธ\-ปญฺห}\csromanheadernomark{3}{4. เมณฺฑกปญฺห 3. พุทฺธ\-อปฺปาพาธ\-ปญฺห}
\proseitem{3}{\csromanfolio{211}ภนฺเต นาคเสน ภาสิตมฺเปตํ ภควตา\csromansymbolfootnote{+}{ขุ 1. 264 ปิฏฺเฐ อิติวุตฺตเก.}“อหมสฺมิ ภิกฺขเว พฺราหฺมโณ ยาจโยโค สทา ปยตปาณิ อนฺติม\-เทห\-ธโร อนุตฺตโร ภิสกฺโก สลฺลกตฺโต”ติ. ปุน จ ภณิตํ ภควตา\csromansymbolfootnote{*}{อํ 1. 25 ปิฏฺเฐ.}“เอตทคฺคํ ภิกฺขเว มม สาวกานํ ภิกฺขูนํ อปฺปาพาธานํ ยทิทํ พากุโล”ติ. ภควโต จ สรีเร พหุกฺขตฺตุํ อาพาโธ อุปฺปนฺโน ทิสฺสติ. ยทิ ภนฺเต นาคเสน ตถาคโต อนุตฺตโร, เตน หิ “เอตทคฺคํ ฯเปฯ พากุโล”ติ ยํ วจนํ, ตํ มิจฺฉา. ยทิ เถโร พากุโล อปฺปาพาธานํ อคฺโค, เตน หิ “อหมสฺมิ ฯเปฯ สลฺลกตฺโต”ติ ตมฺปิ วจนํ มิจฺฉา. อยมฺปิ อุภโต โกฏิโก ปโญฺห ตวานุปฺปตฺโต, โส ตยา นิพฺพาหิตพฺโพติ.}

\prose{ภาสิตมฺเปตํ มหาราช ภควตา “อหมสฺมิ ฯเปฯ สลฺลกตฺโต”ติ, ภณิตญฺจ “เอตทคฺคํ ฯเปฯ พากุโล”ติ, ตญฺจ ปน พาหิรานํ อาคมานํ อธิคมานํ ปริยตฺตีนํ อตฺตนิ วิชฺชมานตํ สนฺธาย ภาสิตํ.}

\prose{สนฺติ โข ปน มหาราช ภควโต สาวกา ฐานจงฺกมิกา, เต ฐาเนน จงฺกเมน ทิวารตฺติํ วีตินาเมนฺติ, ภควา ปน มหาราช ฐาเนน จงฺกเมน นิสชฺชาย สยเนน ทิวารตฺติํ วีตินาเมติ, เย เต มหาราช ภิกฺขู ฐานจงฺกมิกา, เต เตน องฺเคน อติเรกา.}

\prose{สนฺติ โข ปน มหาราช ภควโต สาวกา เอกาสนิกา, เต ชีวิตเหตุปิ ทุติยํ โภชนํ น ภุญฺชนฺติ, ภควา ปน มหาราช ทุติยมฺปิ ยาว ตติยมฺปิ โภชนํ ภุญฺชติ, เย เต มหาราช ภิกฺขู เอกาสนิกา, เต เตน องฺเคน อติเรกา, อเนกวิธานิ มหาราช ตานิ การณานิ เตสํ เตสํ ตํ ตํ สนฺธาย ภณิตานิ. ภควา ปน มหาราช อนุตฺตโร สีเลน สมาธินา ปญฺญาย วิมุตฺติยา วิมุตฺติ\-ญาณ\-ทสฺสเนน ทสหิ จ พเลหิ จตูหิ เวสารชฺเชหิ อฏฺฐารสหิ พุทฺธธมฺเมหิ ฉหิ อสาธารเณหิ ญาเณหิ, เกวเล จ พุทฺธวิสเย ตํ สนฺธาย ภณิตํ “อหมสฺมิ ฯเปฯ สลฺลกตฺโต”ติ.}

\prose{อิธ มหาราช มนุสฺเสสุ เอโก ชาติมา โหติ, เอโก ธนวา, เอโก วิชวา, เอโก สิปฺปวา, เอโก สูโร, เอโก วิจกฺขโณ, สพฺเพเปเต อภิภวิย ราชา เยว เตสํ อุตฺตโม โหติ, เอวเมว โข มหาราช ภควา สพฺพสตฺตานํ อคฺโค เชฏฺโฐ เสฏฺโฐ.}

\prose{\csromanfolio{212}ยํ ปน อายสฺมา พากุโล อปฺปาพาโธ อโหสิ, ตํ อภินีหาร\-วเสน, โส หิ มหาราช อโนมทสฺสิสฺส\csromansharedfootnote{e2f30f70ce36d753}{ขุ 3. 385 ปิฏฺเฐ อปทาเน.} ภควโต อุทรวาตาพาเธ อุปฺปนฺเน วิปสฺสิสฺส จ ภควโต อฏฺฐสฏฺฐิยา จ ภิกฺขุ\-สตสหสฺสานํ ติณปุปฺผก\-โรเค อุปฺปนฺเน สยํ ตาปโส สมาโน นานาเภสชฺเชหิ ตํ พฺยาธิํ อปเนตฺวา อปฺปาพาธตํ ปตฺโต, ภณิโต จ “เอตทคฺคํ ฯเปฯ พากุโล”ติ.}

\prose{ภควโต มหาราช พฺยาธิมฺหิ อุปฺปชฺชนฺเตปิ อนุปฺปชฺชนฺเตปิ ธุตงฺคํ อาทิยนฺเตปิ อนาทิยนฺเตปิ นตฺถิ ภควตา สทิโส โกจิ สตฺโต. ภาสิตมฺเปตํ มหาราช ภควตา เทวาติเทเวน\csromansymbolfootnote{*}{สํ 3. 38; อํ 1. 343; อํ 2. 29; อํ 3. 273; ขุ 1. 255 ปิฏฺเฐสุปิ.}สํยุตฺตนิกาย\-วร\-ลญฺฉเก\csromandash{} “ยาวตา ภิกฺขเว สตฺตา อปทา วา ทฺวิปทา วา จตุปฺปทา วา พหุปฺปทา วา รูปิโน วา อรูปิโน วา สญฺญิโน วา อสญฺญิโน วา เนว\-สญฺญี\-นา\-สญฺญิโน วา, ตถาคโต เตสํ อคฺคมกฺขายติ อรหํ สมฺมาสมฺพุทฺโธ”ติ. สาธุ ภนฺเต นาคเสน, เอวเมตํ ตถา สมฺปฏิจฺฉามีติ.}

\prose{พุทฺธ\-อปฺปาพาธ\-ปโญฺห ตติโย.}
\csromansectionrule

\csromanmatikaanchor{mtk.154}
\csromantocmarkref{subsubsection}{4. มคฺคุปฺปาทนปญฺห}{mtk.154}
\bookmark[dest={mtk.154},level=3]{4. มคฺคุปฺปาทนปญฺห}
\csromanheader{3}{4. \textbf{มคฺคุ}\-\textbf{ปฺปาทน}\-\textbf{ปญฺห}}
\proseitem{4}{ภนฺเต นาคเสน ภาสิตมฺเปตํ ภควตา\csromansymbolfootnote{+}{ม 3. 59; สํ 2. 54; ขุ 7. 138; ขุ 8. 175; ขุ 9. 375 ปิฏฺเฐสุ. ++ สํ 1. 329 ปิฏฺเฐ.}“ตถาคโต ภิกฺขเว อรหํ สมฺมาสมฺพุทฺโธ อนุปฺปนฺนสฺส มคฺคสฺส อุปฺปาเทตา”ติ. ปุน จ ภณิตํ ++ “อทฺทสํ ขฺวาหํ ภิกฺขเว ปุราณํ มคฺคํ ปุราณํ อญฺชสํ ปุพฺพเกหิ สมฺมา\-สมฺพุทฺเธหิ อนุยาตนฺ”ติ. ยทิ ภนฺเต นาคเสน ตถาคโต อนุปฺปนฺนสฺส มคฺคสฺส อุปฺปาเทตา, เตน หิ “อทฺทสํ ขฺวาหํ ภิกฺขเว ปุราณํ มคฺคํ ปุราณํ อญฺชสํ ปุพฺพเกหิ สมฺมา\-สมฺพุทฺเธหิ อนุยาตนฺ”ติ ยํ วจนํ, ตํ มิจฺฉา. ยทิ ตถาคเตน ภณิตํ “อทฺทสํ ขฺวาหํ ภิกฺขเว ปุราณํ มคฺคํ ปุราณํ อญฺชสํ ปุพฺพเกหิ สมฺมา\-สมฺพุทฺเธหิ อนุยาตนฺ”ติ, เตน หิ “ตถาคโต ภิกฺขเว อรหํ สมฺมาสมฺพุทฺโธ อนุปฺปนฺนสฺส มคฺคสฺส อุปฺปาเทตา”ติ ตมฺปิ วจนํ มิจฺฉา. อยมฺปิ อุภโต โกฏิโก ปโญฺห ตวานุปฺปตฺโต, โส ตยา นิพฺพาหิตพฺโพติ.}

\prose{\csromanfolio{213}ภาสิตมฺเปตํ มหาราช ภควตา “ตถาคโต ภิกฺขเว อรหํ สมฺมาสมฺพุทฺโธ อนุปฺปนฺนสฺส มคฺคสฺส อุปฺปาเทตา”ติ. ภณิตญฺจ “อทฺทสํ ขฺวาหํ ภิกฺขเว ปุราณํ มคฺคํ ปุราณํ อญฺชสํ ปุพฺพเกหิ สมฺมา\-สมฺพุทฺเธหิ อนุยาตนฺ”ติ, ตํ ทฺวยมฺปิ สภาว\-วจนเมว, ปุพฺพกานํ มหาราช ตถาคตานํ อนฺตรธาเนน อสติ อนุสาสเก มคฺโค อนฺตรธายิ, ตํ\csromansharedfootnote{42ec0c679dcd910f}{โส ตํ (สี, อิ, ก)} ตถาคโต มคฺคํ ลุคฺคํ ปลุคฺคํ คูฬฺหํ ปิหิตํ ปฏิจฺฉนฺนํ อสญฺจรณํ ปญฺญาจกฺขุนา สมฺปสฺสมาโน\csromansharedfootnote{53d5429aadec7f42}{สมฺมสมาโน (สี, อิ)} อทฺทส ปุพฺพเกหิ สมฺมา\-สมฺพุทฺเธหิ อนุยาตํ, ตํการณา อาห “อทฺทสํ ขฺวาหํ ภิกฺขเว ปุราณํ มคฺคํ ปุราณํ อญฺชสํ ปุพฺพเกหิ สมฺมา\-สมฺพุทฺเธหิ อนุยาตนฺ”ติ.}

\prose{ปุพฺพกานํ มหาราช ตถาคตานํ อนฺตรธาเนน อสติ อนุสาสเก ลุคฺคํ ปลุคฺคํ คูฬฺหํ ปิหิตํ ปฏิจฺฉนฺนํ อสญฺจรณํ มคฺคํ ยํ ทานิ ตถาคโต สญฺจรณํ อกาสิ, ตํการณา อาห “ตถาคโต ภิกฺขเว อรหํ สมฺมาสมฺพุทฺโธ อนุปฺปนฺนสฺส มคฺคสฺส อุปฺปาเทตา”ติ.}

\prose{อิธ มหาราช รญฺโญ จกฺกวตฺติสฺส อนฺตรธาเนน มณิรตนํ คิริสิขนฺตเร นิลียติ, อปรสฺส จกฺกวตฺติสฺส สมฺมา\-ปฏิปตฺติยา อุปคจฺฉติ, อปิ นุ โข ตํ มหารช มณิรตนํ ตสฺส ปกตนฺติ. น หิ ภนฺเต, ปากติกํ เยว ตํ มณิรตนํ, เตน ปน นิพฺพตฺติตนฺติ\csromansharedfootnote{87e44a26ad5574ed}{นิพฺพตฺตนฺติ (สี, อิ)}, เอวเมว โข มหาราช ปากติกํ ปุพฺพเกหิ ตถาคเตหิ อนุจิณฺณํ อฏฺฐงฺคิกํ สิวํ มคฺคํ อสติ อนุสาสเก ลุคฺคํ ปลุคฺคํ คูฬฺหํ ปิหิตํ ปฏิจฺฉนฺนํ อสญฺจรณํ ภควา ปญฺญาจกฺขุนา สมฺปสฺสมาโน อุปฺปาเทสิ, สญฺจรณํ อกาสิ, ตํการณา อาห “ตถาคโต ภิกฺขเว อรหํ สมฺมาสมฺพุทฺโธ อนุปฺปนฺนสฺส มคฺคสฺส อุปฺปาเทตี”ติ.}

\prose{ยถา วา ปน มหาราช สนฺตํ เยว ปุตฺตํ โยนิยา ชนยิตฺวา มาตา “ชนิกา”ติ วุจฺจติ, เอวเมว โข มหาราช ตถาคโต สนฺตํ เยว มคฺคํ ลุคฺคํ ปลุคฺคํ คูฬฺหํ ปิหิตํ ปฏิจฺฉนฺนํ อสญฺจรณํ ปญฺญาจกฺขุนา สมฺปสฺสมาโน อุปฺปาเทสิ, สญฺจรณํ อกาสิ, ตํการณา อาห “ตถาคโต ภิกฺขเว อรหํ สมฺมาสมฺพุทฺโธ อนุปฺปนฺนสฺส มคฺคสฺส อุปฺปาเทตา”ติ.}

\prose{ยถา วา ปน มหาราช โกจิ ปุริโส ยํ กิญฺจิ นฏฺฐํ ปสฺสติ, “เตน ตํ ภณฺฑํ นิพฺพตฺติตนฺ”ติ ชโน โวหรติ, เอวเมว โข มหาราช ตถาคโต สนฺตํ เยว มคฺคํ ลุคฺคํ ปลุคฺคํ คูฬฺหํ ปิหิตํ ปฏิจฺฉนฺนํ}

\prose{\csromanfolio{214}อสญฺจรณํ ปญฺญาจกฺขุนา สมฺปสฺสมาโน อุปฺปาเทสิ, สญฺจรณํ อกาสิ, ตํการณา อาห “ตถาคโต ภิกฺขเว อรหํ สมฺมาสมฺพุทฺโธ อนุปฺปนฺนสฺส มคฺคสฺส อุปฺปาเทตา”ติ.}

\prosewithcloserruled{ยถา วา ปน มหาราช โกจิ ปุริโส วนํ โสเธตฺวา ภูมิํ นีหรติ, “ตสฺส สา ภูมี”ติ ชโน โวหรติ, น เจสา ภูมิ เตน ปวตฺติตา, ตํ ภูมิํ การณํ กตฺวา ภูมิสามิโก นาม โหติ, เอวเมว โข มหาราช ตถาคโต สนฺตํ เยว มคฺคํ ลุคฺคํ ปลุคฺคํ คูฬฺหํ ปิหิตํ ปฏิจฺฉนฺนํ อสญฺจรณํ ปญฺญาย สมฺปสฺสมาโน อุปฺปาเทสิ, สญฺจรณํ อกาสิ, ตํการณา อาห “ตถาคโต ภิกฺขเว อรหํ สมฺมาสมฺพุทฺโธ อนุปฺปนฺนสฺส มคฺคสฺส อุปฺปาเทตา”ติ. สาธุ ภนฺเต นาคเสน, เอวเมตํ ตถา สมฺปฏิจฺฉามีติ.}{มคฺคุ\-ปฺปาทน\-ปโญฺห จตุตฺโถ.}
\csromanmatikaanchor{mtk.155}
\csromantocmarkref{subsubsection}{5. พุทฺธอวิเหฐกปญฺห}{mtk.155}
\bookmark[dest={mtk.155},level=3]{5. พุทฺธอวิเหฐกปญฺห}
\csromanheader{3}{5. พุทฺธ\-อวิเหฐก\-ปญฺห}
\proseitem{5}{ภนฺเต นาคเสน ภาสิตมฺเปตํ ภควตา\csromansymbolfootnote{*}{ที 3. 135 ปิฏฺเฐ.}“ปุพฺเพ วาหํ มนุสฺสภูโต สมาโน สตฺตานํ อวิเหฐก\-ชาติโก อโหสินฺ”ติ. ปุน จ ภณิตํ\csromansymbolfootnote{+}{ขุ 5. 193 ปิฏฺเฐ.}“โลมสกสฺสโป นาม อิสิ สมาโน อเนกสเต ปาเณ ฆาตยิตฺวา วาชเปยฺยํ มหายญฺญํ ยชี”ติ. ยทิ ภนฺเต นาคเสน ภควตา ภณิตํ “ปุพฺเพ วาหํ มนุสฺสภูโต สมาโน สตฺตานํ อวิเหฐก\-ชาติโก อโหสินฺ”ติ, เตน หิ “โลมส\-กสฺสเปน อิสินา อเนกสเต ปาเณ ฆาตยิตฺวา วาชเปยฺยํ มหายญฺญํ ยชิตนฺ”ติ ยํ วจนํ, ตํ มิจฺฉา. ยทิ โลมส\-กสฺสเปน อิสินา อเนกสเต ปาเณ ฆาตยิตฺวา วาชเปยฺยํ มหายญฺญํ ยชิตํ, เตน หิ “ปุพฺเพ วาหํ มนุสฺสภูโต สมาโน สตฺตานํ อวิเหฐก\-ชาติโก อโหสินฺ”ติ ตมฺปิ วจนํ มิจฺฉา. อยมฺปิ อุภโต โกฏิโก ปโญฺห ตวานุปฺปตฺโต, โส ตยา นิพฺพาหิตพฺโพติ.}

\prose{ภาสิตมฺเปตํ มหาราช ภควตา “ปุพฺเพ วาหํ มนุสฺสภูโต สมาโน สตฺตานํ อวิเหฐก\-ชาติโก อโหสินฺ”ติ. โลมส\-กสฺสเปน อิสินา อเนกสเต ปาเณ ฆาตยิตฺวา วาชเปยฺยํ มหายญฺญํ ยชิตํ, ตญฺจ ปน ราควเสน วิสญฺญินา, โน สเจตเนนาติ.}

\prose{\csromanfolio{215}อฏฺฐิเม ภนฺเต นาคเสน ปุคฺคลา ปาณํ หนนฺติ. กตเม อฏฺฐ, รตฺโต ราควเสน ปาณํ หนติ, ทุฏฺโฐ โทสวเสน ปาณํ หนติ, มูโฬฺห โมหวเสน ปาณํ หนติ, มานี มานวเสน ปาณํ หนติ, ลุทฺโธ โลภวเสน ปาณํ หนติ, อกิญฺจโน ชีวิกตฺถาย ปาณํ หนติ, พาโล หสฺสวเสน\csromansharedfootnote{c693c03f9e9c4dcd}{อญฺญาณวเสน (ก, สี)} ปาณํ หนติ, ราชา วินยนวเสน ปาณํ หนติ, อิเม โข ภนฺเต นาคเสน อฏฺฐ ปุคฺคลา ปาณํ หนนฺติ. ปากติกํ เยว ภนฺเต นาคเสน โพธิสตฺเตน กตนฺติ. น มหาราช ปากติกํ โพธิสตฺเตน กตํ, ยทิ มหาราช โพธิสตฺโต ปกติภาเวน โอนเมยฺย มหายญฺญํ ยชิตุํ, น ยิมํ คาถํ ภเณยฺย\csromandash{}}

\setlength{\csromangathapagewidth}{0pt}
\setlength{\csromangathawakwidth}{0pt}
\csromangathameasuregroup{\textsuperscript{*}“สสมุทฺทปริยายํ, \\ น อิจฺเฉ สห นินฺทาย,}{\csromangathabat{\textsuperscript{*}“สสมุทฺทปริยายํ,}{มหิํ สาครกุณฺฑลํ.} \\ \csromangathabat{น อิจฺเฉ สห นินฺทาย,}{เอวํ เสยฺห\textsuperscript{1} วิชานหี”ติ.}}
\csromangathasetleft{\textsuperscript{*}“สสมุทฺทปริยายํ, \\ น อิจฺเฉ สห นินฺทาย,}
\csromangathagroup{\csromangathastanza{\csromangathabat{\csromansymbolfootnote{*}{ขุ 5. 193 ปิฏฺเฐ.}“สสมุทฺท\-ปริยายํ,}{มหิํ สาครกุณฺฑลํ.} \\ \csromangathabat{น อิจฺเฉ สห นินฺทาย,}{เอวํ เสยฺห\csromansharedfootnote{96da2dc4fe8bac25}{สยฺห (สี, อิ)} วิชานหี”ติ.}}}

\prose{เอวํวาที มหาราช โพธิสตฺโต สห ทสฺสเนน จนฺทวติยา ราชกญฺญาย วิสญฺญี อโหสิ ขิตฺตจิตฺโต รตฺโต วิสญฺญิภูโต อากุลากุโล ตุริตตุริโต เตน วิกฺขิตฺต\-ภนฺต\-ลุฬิต\-จิตฺเตน มหติ\-มหา\-ปสุ\-ฆาต\-คล\-รุหิร\-สญฺจยํ วาชเปยฺยํ มหายญฺญํ ยชิ.}

\prose{ยถา มหาราช อุมฺมตฺตโก ขิตฺตจิตฺโต ชลิตมฺปิ ชาตเวทํ อกฺกมติ, กุปิตมฺปิ อาสีวิสํ คณฺหาติ, มตฺตมฺปิ หตฺถิํ อุเปติ, สมุทฺทมฺปิ อตีรทสฺสิํ ปกฺขนฺทติ, จนฺทนิกมฺปิ โอฬิคลฺลมฺปิ โอมทฺทติ, กณฺฏกาธานมฺปิ อภิรุหติ, ปปาเตปิ ปตติ, อสุจิมฺปิ ภกฺเขติ, นคฺโคปิ รถิยา จรติ, อญฺญมฺปิ พหุวิธํ อกิริยํ กโรติ. เอวเมว โข มหาราช โพธิสตฺโต สห ทสฺสเนน จนฺทวติยา ราชกญฺญย วิสญฺญี อโหสิ ขิตฺตจิตฺโต รตฺโต วิสญฺญิภูโต อากุลากุโล ตุริตตุริโต, เตน วิกฺขิตฺต\-ภนฺต\-ลุฬิต\-จิตฺเตน มหติ\-มหา\-ปสุ\-ฆาต\-คล\-รุหิร\-สญฺจยํ วาชเปยฺยํ มหายญฺญํ ยชิ.}

\prose{ขิตฺตจิตฺเตน มหาราช กตํ ปาปํ ทิฏฺฐธมฺเมปิ น มหาสาวชฺชํ โหติ, สมฺปราเย วิปาเกนปิ โน ตถา. อิธ มหาราช โกจิ อุมฺมตฺตโก วชฺฌมา\-ปชฺเชยฺย, ตสฺส ตุเมฺห กิํ ทณฺฑํ ธาเรถาติ. โก ภนฺเต อุมฺมตฺตกสฺส ทณฺโฑ ภวิสฺสติ, ตํ มยํ โปถาเปตฺวา นีหราเปม, เอโสว ตสฺส ทณฺโฑติ. อิติ โข มหาราช อุมฺมตฺตกสฺส อปราเธ ทณฺโฑปิ น}

\prose{\csromanfolio{216}ภวติ, ตสฺมา อุมฺมตฺตกสฺส กเตปิ น โทโส ภวติ สเตกิจฺโฉ. เอวเมว โข มหาราช โลมสกสฺสโป อิสิ สห ทสฺสเนน จนฺทวติยา ราชกญฺญาย วิสญฺญี อโหสิ ขิตฺตจิตฺโต รตฺโต วิสญฺญิภูโต วิสฏปยาโต อากุลากุโล ตุริตตุริโต, เตน วิกฺขิตฺต\-ภนฺต\-ลุฬิต\-จิตฺเตน มหติ\-มหา\-ปสุ\-ฆาต\-คล\-รุหิร\-สญฺจยํ วาชเปยฺยํ มหายญฺญํ ยชิ. ยทา จ ปน ปกติจิตฺโต อโหสิ ปฏิลทฺธ\-สฺสติ, ตทา ปุนเทว ปพฺพชิตฺวา ปญฺจาภิญฺญาโย นิพฺพตฺเตตฺวา พฺรหฺมโลกูปโค อโหสีติ. สาธุ ภนฺเต นาคเสน, เอวเมตํ ตถา สมฺปฏิจฺฉามีติ.}

\prose{พุทฺธ\-อวิเหฐก\-ปโญฺห ปญฺจโม.}
\csromansectionrule

\csromanmatikaanchor{mtk.156}
\csromantocmarkref{subsubsection}{6. ฉทฺทนฺตโชติปาลารพฺภปญฺห}{mtk.156}
\bookmark[dest={mtk.156},level=3]{6. ฉทฺทนฺตโชติปาลารพฺภปญฺห}
\csromanheader{3}{6. \textbf{ฉทฺทนฺต}\-\textbf{โชติปาลา}\-\textbf{รพฺภ}\-\textbf{ปญฺห}}
\prose{6. ภนฺเต นาคเสน ภาสิตมฺเปตํ ภควตา ฉทฺทนฺโต นาคราชา\csromandash{}}

\setlength{\csromangathapagewidth}{0pt}
\setlength{\csromangathawakwidth}{0pt}
\csromangathameasuregroup{}{\textsuperscript{*}“วธิสฺสเมตนฺติ ปรามสนฺโต, \\ กาสาวมทฺทกฺขิ ธชํ อิสีนํ. \\ ทุกฺเขน ผุฏฺฐสฺสุทปาทิ สญฺญา, \\ อรหทฺธโช สพฺภิ อวชฺฌรูโป”ติ.}
\csromangathameasurewak{\textsuperscript{*}“วธิสฺสเมตนฺติ ปรามสนฺโต, \\ กาสาวมทฺทกฺขิ ธชํ อิสีนํ. \\ ทุกฺเขน ผุฏฺฐสฺสุทปาทิ สญฺญา, \\ อรหทฺธโช สพฺภิ อวชฺฌรูโป”ติ.}
\setlength{\csromangathaleftcol}{0pt}
\csromangathagroup{\csromangathastanza{\csromangathawak{\csromansymbolfootnote{*}{ขุ 5. 374 ปิฏฺเฐ.}“วธิสฺสเมตนฺติ ปรามสนฺโต, \\ กาสาวม\-ทฺทกฺขิ ธชํ อิสีนํ. \\ ทุกฺเขน ผุฏฺฐสฺสุทปาทิ สญฺญา, \\ อรหทฺธโช สพฺภิ อวชฺฌรูโป”ติ.}}}

\prose{ปุน จ ภณิตํ\csromansymbolfootnote{+}{ม 2. 236 ปิฏฺเฐ.}“โชติปาลมาณโว สมาโน กสฺสปํ ภควนฺตํ อรหนฺตํ สมฺมา\-สมฺพุทฺธํ มุณฺฑกวาเทน สมณกวาเทน อสพฺภาหิ ผรุสาหิ วาจาหิ อกฺโกสิ ปริภาสี”ติ. ยทิ ภนฺเต นาคเสน โพธิสตฺโต ติรจฺฉานคโต สมาโน กาสาวํ อภิปูชยิ, เตน หิ “โชติปาเลน มาณเวน กสฺสโป ภควา อรหํ สมฺมาสมฺพุทฺโธ มุณฺฑกวาเทน สมณกวาเทน อสพฺภาหิ ผรุสาหิ วาจาหิ อกฺกุฏฺโฐ ปริภาสิโต”ติ ยํ วจนํ, ตํ มิจฺฉา. ยทิ โชติปาเลน มาณเวน กสฺสโป ภควา อรหํ สมฺมาสมฺพุทฺโธ มุณฺฑกวาเทน สมณกวาเทน อสพฺภาหิ ผรุสาหิ วาจาหิ อกฺกุฏฺโฐ ปริภาสิโต, เตน หิ “ฉทฺทนฺเตน นาคราเชน กาสาวํ ปูชิตนฺ”ติ ตมฺปิ วจนํ มิจฺฉา. ยทิ ติรจฺฉาน\-คเตน โพธิสตฺเตน กกฺขฬ\-ขร\-กฏุก\-เวทนํ เวทยมาเนน ลุทฺทเกน นิวตฺถํ กาสาวํ ปูชิตํ, กิํ มนุสฺสภูโต สมาโน ปริปกฺกญาโณ ปริปกฺกาย โพธิยา กสฺสปํ ภควนฺตํ อรหนฺตํ สมฺมา\-สมฺพุทฺธํ ทสพลํ โลกนายกํ \csromanfolio{217}อุทิโตทิตํ ชลิต\-พฺยาโม\-ภาสํ ปวรุตฺตมํ ปวร\-รุจิร\-กาสิก\-กาสาวม\-ภิปารุตํ ทิสฺวา น ปูชยิ. อยมฺปิ อุภโต โกฏิโก ปโญฺห ตวานุปฺปตฺโต, โส ตยา นิพฺพาหิตพฺโพติ.}

\proseitem{5}{ภาสิตมฺเปตํ มหาราช ภควตา ฉทฺทนฺโต นาคราชา “วธิสฺสเมตนฺติ ฯเปฯ อวชฺฌรูโป”ติ. โชติปาเลน จ มาณเวน กสฺสโป ภควา อรหํ สมฺมาสมฺพุทฺโธ มุณฺฑกวาเทน สมณกวาเทน อสพฺภาหิ ผรุสาหิ วาจาหิ อกฺกุฏฺโฐ ปริภาสิโต, ตญฺจ ปน ชาติวเสน กุลวเสน. โชติปาโล มหาราช มาณโว อสฺสทฺเธ อปฺปสนฺเน กุเล ปจฺจาชาโต, ตสฺส มาตาปิตโร ภคินิภาตโร ทาสิทาส\-เจฏก\-ปริวารก\-มนุสฺสา พฺรหฺมเทวตา พฺรหฺมครุกา, เต “พฺรหฺมณา เอว อุตฺตมา ปวรา”ติ อวเสเส ปพฺพชิเต ครหนฺติ ชิคุจฺฉนฺติ, เตสํ ตํ วจนํ สุตฺวา โชติปาโล มาณโว ฆฏิกาเรน กุมฺภกาเรน สตฺถารํ ทสฺสนาย ปกฺโกสิโต เอวมาห “กิํ ปน เตน มุณฺฑเกน สมณเกน ทิฏฺเฐนา”ติ.}

\prose{ยถา มหาราช อมตํ วิสมาสชฺช ติตฺตกํ โหติ, ยถา จ สีโตทกํ อคฺคิมาสชฺช อุณฺหํ โหติ, เอวเมว โข มหาราช โชติปาโล มาณโว อสฺสทฺเธ อปฺปสนฺเน กุเล ปจฺจาชาโต, โส กุลวเสน อนฺโธ หุตฺวา\csromansharedfootnote{888a90ef5589369f}{โส กุลชาติวเสน อนฺโธ ภวิตฺวา (สฺยา)} ตถาคตํ อกฺโกสิ ปริภาสิ.}

\prose{ยถา มหาราช ชลิต\-ปชฺชลิโต มหา\-อคฺคิกฺขนฺโธ สปฺปภาโส อุทกมาสชฺช อุปหต\-ปฺปภา\-เตโช สีตโล กาฬโก ภวติ ปริปกฺก\-นิคฺคุณฺฑิ\-ผล\-สทิโส, เอวเมว โข มหาราช โชติปาโล มาณโว ปุญฺญวา สทฺโธ ญาณ\-วิปุล\-สปฺปภาโส อสฺสทฺเธ อปฺปสนฺเน กุเล ปจฺจาชาโต, โส กุลวเสน อนฺโธ หุตฺวา ตถาคตํ อกฺโกสิ ปริภาสิ, อุปคนฺตฺวา จ พุทฺธคุณ\-มญฺญาย เจฏกภูโต วิย อโหสิ, ชินสาสเน ปพฺพชิตฺวา อภิญฺญา จ สมาปตฺติโย จ นิพฺพตฺเตตฺวา พฺรหฺมโลกูปโค อโหสีติ. สาธุ ภนฺเต นาคเสน, เอวเมตํ ตถา สมฺปฏิจฺฉามีติ.}

\vspace{\tipitakaparskip}
\csromancenter{ฉทฺทนฺต\-โชติปาลา\-รพฺภ\-ปโญฺห ฉฏฺโฐ.}
\csromanmatikaanchor{mtk.157}
\csromantocmarkref{subsubsection}{7. ฆฏิการปญฺห}{mtk.157}
\bookmark[dest={mtk.157},level=3]{7. ฆฏิการปญฺห}
\csromanheader{3}{7. \textbf{ฆฏิการปญฺห}}
\proseitem{7}{\csromanfolio{218}ภนฺเต นาคเสน ภสิตมฺเปตํ ภควตา\csromansymbolfootnote{*}{ม 2. 244 ปิฏฺเฐ.}“ฆฏิการสฺส กุมฺภการสฺส อาเวสนํ สพฺพํ เตมาสํ อากาสจฺฉทนํ อฏฺฐาสิ, น เทโวติวสฺสี”ติ. ปุน จ ภณิตํ “กสฺสปสฺส ตถาคตสฺส\csromansharedfootnote{dde66810241c6cbe}{ภควโต อรหโต สมฺมาสมฺพุทฺธสฺส (ม 2. 243 ปิฏฺเฐ)} กุฏิ โอวสฺสตี”ติ. กิสฺส ปน ภนฺเต นาคเสน ตถาคตสฺส เอวมุ\-สฺสนฺนกุสลมูลสฺส\csromansharedfootnote{9bf11c92aa70b15e}{เอวรูปสฺส อุสฺสนฺนกุสลมูลสฺส (ก)} กุฏิ โอวสฺสติ, ตถาคตสฺส นาม โส อานุภาโว อิจฺฉิตพฺโพ. ยทิ ภนฺเต นาคเสน ฆฏิการสฺส กุมฺภการสฺส อาเวสนํ อโนวสฺสํ อากาสจฺฉทนํ อโหสิ, เตน หิ “ตถาคตสฺส กุฏิ โอวสฺสตี”ติ ยํ วจนํ, ตํ มิจฺฉา. ยทิ ตถาคตสฺส กุฏิ โอวสฺสติ, เตน หิ “ฆฏิการสฺส กุมฺภการสฺส อาเวสนํ อโนวสฺสกํ อโหสิ อากาสจฺฉทนนฺ”ติ ตมฺปิ วจนํ มิจฺฉา. อยมฺปิ อุภโต โกฏิโก ปโญฺห ตวานุปฺปตฺโต, โส ตยา นิพฺพาหิตพฺโพติ.}

\prose{ภาสิตมฺเปตํ มหาราช ภควตา “ฆฏิการสฺส กุมฺภการสฺส อาเวสนํ สพฺพํ เตมาสํ อากาสจฺฉทนํ อฏฺฐาสิ, น เทโวติวสฺสี”ติ. ภณิตญฺจ “กสฺสปสฺส ตถาคตสฺส\csromansharedfootnote{dde66810241c6cbe}{ภควโต อรหโต สมฺมาสมฺพุทฺธสฺส (ม 2. 243 ปิฏฺเฐ)} กุฏิ โอวสฺสตี”ติ. ฆฏิกาโร มหาราช กุมฺภกาโร สีลวา กลฺยาณธมฺโม อุสฺสนฺน\-กุสล\-มูโล อนฺเธ ชิณฺเณ มาตาปิตโร โปเสติ, ตสฺส อสมฺมุขา อนาปุจฺฉาเยวสฺส ฆเร ติณํ หริตฺวา ภควโต กุฏิํ ฉาเทสุํ, โส เตน ติณหรเณน อกมฺปิตํ อสญฺจลิตํ สุสณฺฐิตํ วิปุลมสมํ ปีติํ ปฏิลภติ, ภิยฺโย โสมนสฺสญฺจ อตุลํ อุปฺปาเทสิ “อโห วต เม ภควา โลกุตฺตโม สุวิสฺสตฺโถ”ติ, เตน ตสฺส ทิฏฺฐธมฺมิโก วิปาโก นิพฺพตฺโต. น หิ มหาราช ตถาคโต ตาวตเกน วิกาเรน จลติ.}

\prose{ยถา มหาราช สิเนรุ คิริราชา อเนกสตสหสฺส\-วาต\-สมฺปหาเรนปิ น กมฺปติ น จลติ, มโหทธิ วรปฺปวร\-สาคโร อเนกสตนหุตมหาคงฺคาสตสหสฺเสหิปิ น ปูรติ น วิการมา\-ปชฺชติ, เอวเมว โข มหาราช ตถาคโต น ตาวตเกน วิกาเรน จลติ.}

\prosewithcloserruled{ยํ ปน มหาราช ตถาคตสฺส กุฏิ โอวสฺสติ, ตํ มหโต ชนกายสฺส อนุกมฺปาย. เทฺวเม มหาราช อตฺถวเส สมฺปสฺสมานา \csromanfolio{219}ตถาคตา สยํ นิมฺมิตํ ปจฺจยํ นปฺปฏิเสวนฺติ, “อยํ อคฺค\-ทกฺขิเณยฺโย สตฺถา”ติ ภควโต ปจฺจยํ ทตฺวา เทวมนุสฺสา สพฺพทุคฺคติโต ปริมุจฺจิสฺสนฺตีติ, ปาฏิหีรํ ทสฺเสตฺวา วุตฺติํ ปริเยสนฺตีติ “มา อญฺเญ อุปวเทยฺยุนฺ”ติ. อิเม เทฺว อตฺถวเส สมฺปสฺสมานา ตถาคตา สยํ นิมฺมิตํ ปจฺจยํ นปฺปฏิเสวนฺติ. ยทิ มหาราช สกฺโก วา ตํ กุฏิํ อโนวสฺสํ กเรยฺย พฺรหฺมา วา สยํ วา, สาวชฺชํ ภเวยฺย ตํ เยว\csromansharedfootnote{cf7dc7e6d0b8f44d}{ตํเยว นิคฺคหีตสนฺธิ \csromandash{} ม.พ.ป.} กรณํ\csromansharedfootnote{a1c9df95b888e7ea}{การณํ (สี, อิ)} สโทสํ สนิคฺคหํ, อิเม วิภูตํ\csromansharedfootnote{981b6383f0bf16c6}{วิภูสํ (สี, อิ)} กตฺวา โลกํ สมฺโมเหนฺติ อธิกตํ กโรนฺตีติ, ตสฺมา ตํ กรณํ วชฺชนียํ. น มหาราช ตถาคตา วตฺถุํ ยาจนฺติ, ตาย อวตฺถุยา\-จ\-นาย อปริภาสิยา ภวนฺตีติ. สาธุ ภนฺเต นาคเสน, เอวเมตํ ตถา สมฺปฏิจฺฉามีติ.}{ฆฏิการปโญฺห สตฺตโม.}
\csromanmatikaanchor{mtk.158}
\csromantocmarkref{subsubsection}{8. พฺราหฺมณราชวาทปญฺห}{mtk.158}
\bookmark[dest={mtk.158},level=3]{8. พฺราหฺมณราชวาทปญฺห}
\csromanheader{3}{8. \textbf{พฺราหฺมณราชวาทปญฺห}}
\proseitem{8}{ภนฺเต นาคเสน ภาสิตมฺเปตํ ตถาคเตน\csromansymbolfootnote{+}{ขุ 1. 264 ปิฏฺเฐ อิติวุตฺตเก.}“อหมสฺมิ ภิกฺขเว พฺราหฺมโณ ยาจโยโค”ติ. ปุน จ ภณิตํ\csromansymbolfootnote{*}{ม 2. 352; ขุ 1. 368 ปิฏฺเฐสุ.}“ราชาหมสฺมิ เสลา”ติ. ยทิ ภนฺเต นาคเสน ภควตา ภณิตํ “อหมสฺมิ ภิกฺขเว พฺราหฺมโณ ยาจโยโค”ติ, เตน หิ “ราชาหมสฺมิ เสลา”ติ, ยํ วจนํ, ตํ มิจฺฉา. ยทิ ตถาคเตน ภณิตํ “ราชาหมสฺมิ เสลา”ติ, เตน หิ “อหมสฺมิ ภิกฺขเว พฺราหฺมโณ ยาจโยโค”ติ ตมฺปิ วจนํ มิจฺฉา. ขตฺติโย วา หิ ภเวยฺย พฺราหฺมโณ วา, นตฺถิ เอกาย ชาติยา เทฺว วณฺณา นาม, อยมฺปิ อุภโต โกฏิโก ปโญฺห ตวานุปฺปตฺโต, โส ตยา นิพฺพาหิตพฺโพติ.}

\prose{ภาสิตมฺเปตํ มหาราช ภควตา “อหมสฺมิ ภิกฺขเว พฺราหฺมโณ ยาจโยโค”ติ, ปุน จ ภณิตํ “ราชหมสฺมิ เสลา”ติ, ตตฺถ การณํ อตฺถิ, เยน การเณน ตถาคโต พฺราหฺมโณ จ ราชา จ โหตีติ.}

\prose{กิํ ปน ตํ ภนฺเต นาคเสน การณํ, เยน การเณน ตถาคโต พฺราหฺมโณ จ ราชา จ โหติ, สพฺเพ มหาราช \csromanfolio{220}ปาปกา อกุสลา ธมฺมา ตถาคตสฺส พาหิตา ปหีนา อปคตา พฺยปคตา อุจฺฉินฺนา ขีณา ขยํ ปตฺตา นิพฺพุตา อุปสนฺตา, ตสฺมา ตถาคโต “พฺราหฺมโณ”ติ วุจฺจติ.}

\prose{พฺราหฺมโณ นาม สํสย\-มเนกํสํ วิมติปถํ วีติวตฺโต, ภควาปิ มหาราช สํสย\-มเนกํสํ วิมติปถํ วีติวตฺโต, เตน การเณน ตถาคโต “พฺราหฺมโณ”ติ วุจฺจติ.}

\prose{พฺราหฺมโณ นาม สพฺพภว\-คติ\-โยนิ\-นิสฺสโฏ มลรช\-คต\-วิปฺปมุตฺโต อสหาโย, ภควาปิ มหาราช สพฺพภว\-คติ\-โยนิ\-นิสฺสโฏ มลรช\-คต\-วิปฺปมุตฺโต อสหาโย, เตน การเณน ตถาคโต “พฺราหฺมโณ”ติ วุจฺจติ.}

\prose{พฺราหฺมโณ นาม อคฺค\-เสฏฺฐ\-วร\-ปวร\-ทิพฺพวิหาร\-พหุโล, ภควาปิ มหาราช อคฺค\-เสฏฺฐ\-วร\-ปวร\-ทิพฺพวิหาร\-พหุโล, เตนาปิ การเณน ตถาคโต “พฺราหฺมโณ”ติ วุจฺจติ.}

\prose{พฺราหฺมโณ นาม อชฺฌยน อชฺฌาปน ทาน ปฺปฏิคฺคหณ ทม สํยม\-นิยม\-ปุพฺพม\-นุสิฏฺฐิ ปเวณิ วํส ธรโณ, ภควาปิ มหาราช อชฺฌยน อชฺฌาปนทานปฺปฏิคฺคหณ ทม สํยม นิยม ปุพฺพชินาจิณฺณ อนุสิฏฺฐิ ปเวณิ วํส ธรโณ, เตนาปิ การเณน ตถาคโต “พฺราหฺมโณ”ติ วุจฺจติ.}

\prose{พฺราหฺมโณ นาม พฺรหา\-สุขวิหาร\-ชฺฌาน\-ฌายี, ภควาปิ มหาราช พฺรหา\-สุขวิหาร\-ชฺฌาน\-ฌายี, เตนาปิ การเณน ตถาคโต “พฺราหฺมโณ”ติ วุจฺจติ.}

\prose{พฺราหฺมโณ นาม สพฺพ\-ภวาภว\-คตีสุ อภิชาติ\-วตฺติตม\-นุจริตํ ชานาติ, ภควาปิ มหาราช สพฺพ\-ภวาภว\-คตีสุ อภิชาติ\-วตฺติตม\-นุจริตํ ชานาติ, เตนาปิ การเณน ตถาคโต “พฺราหฺมโณ”ติ วุจฺจติ.}

\prose{พฺราหฺมโณติ มหาราช ภควโต เนตํ นามํ มาตรา กตํ, น ปิตรา กตํ, น ภาตรา กตํ, น ภคินิยา กตํ, น มิตฺตามจฺเจหิ กตํ, น ญาติสาโลหิเตหิ กตํ, นสมณ\-พฺราหฺมเณหิ กตํ, น เทวตาหิ กตํ, วิโมกฺข\-นฺติกเมตํ พุทฺธานํ ภควนฺตานํ นามํ โพธิยา เยว มูเล มารเสนํ วิธมิตฺวา อตีตา\-นาคต\-ปจฺจุปฺปนฺเน ปาปเก \csromanfolio{221}อกุสเล ธมฺเม พาเหตฺวา สห สพฺพญฺญุต\-ญาณสฺส ปฏิลาภา ปฏิลทฺธ\-ปาตุภูต\-สมุปฺปนฺน\-มตฺเต สจฺฉิกา ปญฺญตฺติ ยทิทํ พฺราหฺมโณติ, เตน การเณน ตถาคโต วุจฺจติ “พฺราหฺมโณ”ติ.}

\prose{เกน ปน ภนฺเต นาคเสน การเณน ตถาคโต วุจฺจติ “ราชา”ติ. ราชา นาม มหาราช โย โกจิ รชฺชํ กาเรติ โลกมนุสาสติ, ภควาปิ มหาราช ทสสหสฺสิยา โลกธาตุยา ธมฺเมน รชฺชํ กาเรติ, สเทวกํ โลกํ สมารกํ สพฺรหฺมกํ สสฺสมณ\-พฺราหฺมณิํ ปชํ สเทวมนุสฺสํ อนุสาสติ, เตนาปิ การเณน ตถาคโต วุจฺจติ “ราชา”ติ.}

\prose{ราชา นาม มหาราช สพฺพชน\-มนุสฺเส อภิภวิตฺวา นนฺทยนฺโต ญาติสํฆํ, โสจยนฺโต อมิตฺตสํฆํ, มหติ\-มหายส\-สิริหรํ ถิรสาร\-ทณฺฑํ อนูน\-สต\-สลากา\-ลงฺกตํ อุสฺสาเปติ ปณฺฑร\-วิมล\-เสตจฺฉตฺตํ, ภควาปิ มหาราช โสจยนฺโต มารเสนํ มิจฺฉา\-ปฏิปนฺนํ, นนฺทยนฺโต เทวมนุสฺเส สมฺมาปฏิปนฺเน ทสสหสฺสิยา โลกธาตุยา มหติ\-มหายส\-สิริหรํ ขนฺติ\-ถิร\-สาร\-ทณฺฑํ ญาณ\-วร\-สต\-สลากา\-ลงฺกตํ อุสฺสาเปติ อคฺค\-วรวิมุตฺติ\-ปณฺฑร\-วิมล\-เสตจฺฉตฺตํ, เตนาปิ การเณน ตถาคโต วุจฺจติ “ราชา”ติ.}

\prose{ราชา นาม อุปคตสมฺปตฺตชานานํ พหูนม\-ภิวนฺท\-นีโย ภวติ, ภควาปิ มหาราช อุปคต\-สมฺปตฺต\-เทวมนุสฺสานํ พหูนม\-ภิวนฺท\-นีโย, เตนาปิ การเณน ตถาคโต วุจฺจติ “ราชา”ติ.}

\prose{ราชา นาม ยสฺส กสฺสจิ อาราธกสฺส ปสีทิตฺวา วริตํ วรํ ทตฺวา กาเมน ตปฺปยติ, ภควาปิ มหาราช ยสฺส กสฺสจิ กาเยน วาจาย มนสา อาราธกสฺส ปสีทิตฺวา วริตํ วรมนุตฺตรํ สพฺพทุกฺข\-ปริมุตฺติํ ทตฺวา อเสส\-กาม\-วเรน จ ตปฺปยติ, เตนาปิ การเณน ตถาคโต วุจฺจติ “ราชา”ติ.}

\prose{ราชา นาม อาณํ วีติกฺกมนฺตํ วิครหติ ฌาเปติ\csromansharedfootnote{48d98c65eb690af5}{ชาเปติ (สี, อิ)} ธํเสติ, ภควโตปิ มหาราช สาสนวเร อาณํ อติกฺกมนฺโต อลชฺชี มงฺกุภาเวน โอญฺญาโต หีฬิโต ครหิโต ภวิตฺวา วชฺชติ ชินสาสน\-วรมฺหา, เตนาปิ การเณน ตถาคโต วุจฺจติ “ราชา”ติ.}

\proseitemwithcloserruled{4}{\csromanfolio{222}ราชา นาม ปุพฺพกานํ ธมฺมิกานํ ราชูนํ ปเวณิม\-นุสิฏฺฐิยา ธมฺมาธมฺมมนุทีปยิตฺวา ธมฺเมน รชฺชํ การยมาโน ปิหยิโต ปิโย ปตฺถิโต ภวติ ชนมนุสฺสานํ, จิรํ ราชกุลวํสํ ฐปยติ ธมฺมคุณ\-พเลน, ภควาปิ มหาราช ปุพฺพกานํ สยมฺภูนํ ปเวณิม\-นุสิฏฺฐิยา ธมฺมาธมฺมมนุทีปยิตฺวา ธมฺเมน โลกม\-นุสาสมาโน ปิหยิโต ปิโย ปตฺถิโต เทวมนุสฺสานํ จิรํ สาสนํ ปวตฺเตติ ธมฺมคุณ\-พเลน, เตนาปิ การเณน ตถาคโต วุจฺจติ “ราชา”ติ. เอวม\-เนกวิธํ มหาราช การณํ, เยน การเณน ตถาคโต พฺราหฺมโณปิ ภเวยฺย ราชาปิ ภเวยฺย, สุนิปุโณ ภิกฺขุ กปฺปมฺปิ โน นํ สมฺปาเทยฺย, กิํ อติพหุํ ภณิเตน, สํขิตฺตํ สมฺปฏิจฺฉิตพฺพนฺติ. สาธุ ภนฺเต นาคเสน, เอวเมตํ ตถา สมฺปฏิจฺฉามีติ.}{พฺราหฺมณราชวาทปโญฺห อฏฺฐโม.}
\csromanmatikaanchor{mtk.159}
\csromantocmarkref{subsubsection}{9. คาถาภิคีตโภชนกถาปญฺห}{mtk.159}
\bookmark[dest={mtk.159},level=3]{9. คาถาภิคีตโภชนกถาปญฺห}
\csromanheader{3}{9. \textbf{คาถาภิคีต}\-\textbf{โภชน}\-\textbf{กถา}\-\textbf{ปญฺห}}
\prose{9. ภนฺเต นาคเสน ภาสิตมฺเปตํ ภควตา\csromandash{}}

\setlength{\csromangathapagewidth}{0pt}
\setlength{\csromangathawakwidth}{0pt}
\csromangathameasuregroup{}{“คาถาภิคีตํ เม อโภชเนยฺยํ\textsuperscript{1}, \\ สมฺปสฺสตํ พฺราหฺมณ เนส ธมฺโม. \\ คาถาภิคีตํ ปนุทนฺติ พุทฺธา, \\ ธมฺเม สตี พฺราหฺมณ วุตฺติเรสา”ติ.}
\csromangathameasurewak{“คาถาภิคีตํ เม อโภชเนยฺยํ\textsuperscript{1}, \\ สมฺปสฺสตํ พฺราหฺมณ เนส ธมฺโม. \\ คาถาภิคีตํ ปนุทนฺติ พุทฺธา, \\ ธมฺเม สตี พฺราหฺมณ วุตฺติเรสา”ติ.}
\setlength{\csromangathaleftcol}{0pt}
\csromangathagroup{\csromangathastanza{\csromangathawak{“คาถาภิคีตํ เม อโภชเนยฺยํ\csromansharedfootnote{b6cf1e9304078d27}{อโภชนียํ (ก) สํ 1. 175; ขุ 1. 292, 349 ปิฏฺเฐสุ.}, \\ สมฺปสฺสตํ พฺราหฺมณ เนส ธมฺโม. \\ คาถาภิคีตํ ปนุทนฺติ พุทฺธา, \\ ธมฺเม สตี พฺราหฺมณ วุตฺติเรสา”ติ.}}}

\prose{ปุน จ ภควา ปริสาย ธมฺมํ เทเสนฺโต กเถนฺโต อนุปุพฺพิกถํ ปฐมํ ตาว ทานกถํ กเถติ, ปจฺฉา สีลกถํ, ตสฺส ภควโต สพฺพโลกิ\-สฺสรสฺส ภาสิตํ สุตฺวา เทวมนุสฺสา อภิส\-งฺขริ\-ตฺวา ทานํ เทนฺติ, ตสฺส ตํ อุยฺโยชิตํ ทานํ สาวกา ปริภุญฺชนฺติ. ยทิ ภนฺเต นาคเสน ภควตา ภณิตํ “คาถาภิคีตํ เม อโภชเนยฺยนฺ”ติ, เตน หิ “ภควา ทานกถํ ปฐมํ กเถตี”ติ ยํ วจนํ, ตํ 0มิจฺฉา. ยทิ ทานกถํ ปฐมํ กเถติ, เตน หิ “คาถาภิคีตํ เม อโภชเนยฺยนฺ”ติ ตมฺปิ วจนํ มิจฺฉา. กิํ การณํ, โย โส ภนฺเต ทกฺขิเณยฺโย คิหีนํ ปิณฺฑปาต\-ทานสฺส วิปากํ กเถติ, ตสฺส เต ธมฺมกถํ สุตฺวา ปสนฺนจิตฺตา \csromanfolio{223}อปราปรํ ทานํ เทนฺติ, เย ตํ ทานํ ปริภุญฺชนฺติ, สพฺเพ เต คาถาภิคีตํ ปริภุญฺชนฺติ. อยมฺปิ อุภโต โกฏิโก ปโญฺห นิปุโณ คมฺภีโร ตวานุปฺปตฺโต, โส ตยา นิพฺพาหิตพฺโพติ.}

\proseitem{4}{ภาสิตมฺเปตํ มหาราช ภควตา “คาถาภิคีตํ เม อโภชเนยฺยํ, สมฺปสฺสตํ พฺราหฺมณ เนส ธมฺโม. คาถาภิคีตํ ปนุทนฺติ พุทฺธา, ธมฺเม สตี พฺราหฺมณ วุตฺติเรสา”ติ, กเถติ จ ภควา ปฐมํ ทานกถํ ตญฺจ ปน กิริยํ สพฺเพสํ ตถาคตานํ ปฐมํ ทานกถาย, ตตฺถ จิตฺตํ อภิรมาเปตฺวา ปจฺฉา สีเล นิโยเชนฺติ. ยถา มหาราช มนุสฺสา ตรุณ\-ทารกานํ ปฐมํ ตาว กีฬาภณฺฑกานิ เทนฺติ. เสยฺยถีทํ, วงฺกกํ ฆฏิกํ จิงฺคุลกํ ปตฺตาฬฺหกํ รถกํ ธนุกํ, ปจฺฉา เต สเก สเก กมฺเม นิโยเชนฺติ. เอวเมว โข มหาราช ตถาคโต ปฐมํ ทานกถาย จิตฺตํ อภิรมาเปตฺวา ปจฺฉา สีเล นิโยเชติ.}

\prose{ยถา วา ปน มหาราช ภิสกฺโก นาม อาตุรานํ ปฐมํ ตาว จตูหปญฺจาหํ เตลํ ปาเยติ พลกรณาย สิเนหนาย, ปจฺฉา วิเรเจติ. เอวเมว โข มหาราช ตถาคโต ปฐมํ ตาว ทานกถาย จิตฺตํ อภิรมาเปตฺวา ปจฺฉา สีเล นิโยเชติ. ทายกานํ มหาราช ทานปตีนํ จิตฺตํ มุทุกํ โหติ มทฺทวํ สินิทฺธํ, เตน เต ทาน\-เสตุ\-สงฺกเมน ทานนาวาย สํสารสาคร\-ปารม\-นุคจฺฉนฺติ, ตสฺมา เตสํ ปฐมํ กมฺม\-ภูมิม\-นุสาสติ, น จ เกนจิ\csromansharedfootnote{f52a20201a74ca59}{เตน (สี, อิ)} วิญฺญตฺติมา\-ปชฺชตีติ.}

\prose{ภนฺเต นาคเสน “วิญฺญตฺตินฺ”ติ ยํ วเทสิ, กติ ปน ตา วิญฺญตฺติโยติ. เทฺวมา มหาราช วิญฺญตฺติโย กายวิญฺญตฺติ วจีวิญฺญตฺติ จาติ. ตตฺถ อตฺถิ กายวิญฺญตฺติ สาวชฺชา, อตฺถิ อนวชฺชา. อตฺถิ วจีวิญฺญตฺติ สาวชฺชา, อตฺถิ อนวชฺชา.}

\prose{กตมา กายวิญฺญตฺติ สาวชฺชา, อิเธกจฺโจ ภิกฺขุ กุลานิ อุปคนฺตฺวา อโนกาเส ฐิโต ฐานํ ภญฺชติ, อยํ กายวิญฺญตฺติ สาวชฺชา. ตาย จ วิญฺญาปิตํ อริยา น ปริภุญฺชนฺติ, โส จ ปุคฺคโล อริยานํ สมเย โอญฺญาโต โหติ หีฬิโต ขีฬิโต ครหิโต ปริภูโต อจิตฺตีกโต, ภินฺนาชีโวเตฺวว สงฺขํ คจฺฉติ.}

\prose{\csromanfolio{224}ปุน จปรํ มหาราช อิเธกจฺโจ ภิกฺขุ กุลานิ อุปคนฺตฺวา อโนกาเส ฐิโต คลํ ปณาเมตฺวา โมรเปกฺขิตํ เปกฺขติ “เอวํ อิเม ปสฺสนฺตี”ติ, เตน จ เต ปสฺสนฺติ. อยมฺปิ กายวิญฺญตฺติ สาวชฺชา. ตาย จ วิญฺญาปิตํ อริยา น ปริภุญฺชนฺติ, โส จ ปุคฺคโล อริยานํ สมเย โอญฺญาโต โหติ หีฬิโต ขีฬิโต ครหิโต ปริภูโต อจิตฺตีกโต, ภินฺนาชีโวเตฺวว สงฺขํ คจฺฉติ.}

\prose{ปุน จปรํ มหาราช อิเธกจฺโจ ภิกฺขุ หนุกาย วา ภมุกาย วา องฺคุฏฺเฐน วา วิญฺญาเปติ, อยมฺปิ กายวิญฺญตฺติ สาวชฺชา, ตาย จ วิญฺญาปิตํ อริยา น ปริภุญฺชนฺติ, โส จ ปุคฺคโล อริยานํ สมเย โอญฺญาโต โหติ หีฬิโต ขีฬิโต ครหิโต ปริภูโต อจิตฺตีกโต, ภินฺนาชีโวเตฺวว สงฺขํ คจฺฉติ.}

\prose{กตมา กายวิญฺญตฺติ อนวชฺชา, อิธ ภิกฺขุ กุลานิ อุปคนฺตฺวา สโต สมาหิโต สมฺปชาโน ฐาเนปิ อฏฺฐาเนปิ ยถานุสิฏฺฐิํ คนฺตฺวา ฐาเน ติฏฺฐติ, ทาตุกาเมสุ ติฏฺฐติ, อทาตุกาเมสุ ปกฺกมติ. อยํ กายวิญฺญตฺติ อนวชฺชา, ตาย จ วิญฺญาปิตํ อริยา ปริภุญฺชนฺติ, โส จ ปุคฺคโล อริยานํ สมเย วณฺณิโต โหติ ถุโต ปสตฺโถ สลฺเลขิตาจาโร, ปริสุทฺธาชีโวเตฺวว สงฺขํ คจฺฉติ. ภาสิตมฺเปตํ มหาราช ภควตา เทวาติเทเวน\csromandash{}}

\setlength{\csromangathapagewidth}{0pt}
\setlength{\csromangathawakwidth}{0pt}
\csromangathameasuregroup{“น เว ยาจนฺติ สปฺปญฺญา, \\ อุทฺทิสฺส อริยา ติฏฺฐนฺติ,}{\csromangathabat{“น เว ยาจนฺติ สปฺปญฺญา,}{ธีโร จ เวทิตุมรหติ\textsuperscript{1}.} \\ \csromangathabat{อุทฺทิสฺส อริยา ติฏฺฐนฺติ,}{เอสา อริยาน ยาจนา”ติ.}}
\csromangathasetleft{“น เว ยาจนฺติ สปฺปญฺญา, \\ อุทฺทิสฺส อริยา ติฏฺฐนฺติ,}
\csromangathagroup{\csromangathastanza{\csromangathabat{“น เว ยาจนฺติ สปฺปญฺญา,}{ธีโร จ เวทิตุมรหติ\csromansharedfootnote{bdd9637ad8688b0c}{อริยา ครหนฺติ ยาจนํ (สี, อิ) ขุ 5. 162 ปิฏฺเฐ.}.} \\ \csromangathabat{อุทฺทิสฺส อริยา ติฏฺฐนฺติ,}{เอสา อริยาน ยาจนา”ติ.}}}

\prose{กตมา วจีวิญฺญตฺติ สาวชฺชา, อิธ มหาราช ภิกฺขุ วาจาย พหุวิธํ วิญฺญาเปติ จีวร\-ปิณฺฑปาต\-เสนาสน\-คิลานปฺปจฺจย\-เภสชฺช\-ปริกฺขารํ, อยํ วจีวิญฺญตฺติ สาวชฺชา, ตาย จ วิญฺญาปิตํ อริยา น ปริภุญฺชนฺติ, โส จ ปุคฺคโล อริยานํ สมเย โอญฺญาโต โหติ หีฬิโต ขีฬิโต ครหิโต ปริภูโต อจิตฺตีกโต, ภินฺนาชีโวเตฺวว สงฺขํ คจฺฉติ.}

\prose{ปุน จปรํ มหาราช อิเธกจฺโจ ภิกฺขุ ปเรสํ สาเวนฺโต เอวํ ภณติ “อิมินา เม อตฺโถ”ติ, ตาย จ วาจาย ปเรสํ สาวิตาย ตสฺส ลาโภ อุปฺปชฺชติ, อยมฺปิ วจีวิญฺญตฺติ สาวชฺชา, ตาย จ วิญฺญาปิตํ \csromanfolio{225}อริยา น ปริภุญฺชนฺติ, โส จ ปุคฺคโล อริยานํ สมเย โอญฺญาโต โหติ หีฬิโต ขีฬิโต ครหิโต ปริภูโต อจิตฺตีกโต, ภินฺนาชีโวเตฺวว สงฺขํ คจฺฉติ.}

\prose{ปุน จปรํ มหาราช อิเธกจฺโจ ภิกฺขุ วจีวิปฺผาเรน ปริสาย สาเวติ “เอวญฺจ เอวญฺจ ภิกฺขูนํ ทาตพฺพนฺ”ติ, ตญฺจ เต วจนํ สุตฺวา ปริกิตฺติตํ อภิหรนฺติ, อยมฺปิ วจีวิญฺญตฺติ สาวชฺชา, ตาย จ วิญฺญาปิตํ อริยา น ปริภุญฺชนฺติ, โส จ ปุคฺคโล อริยานํ สมเย โอญฺญาโต โหติ หีฬิโต ขีฬิโต ครหิโต ปริภูโต อจิตฺตีกโต, ภินฺนาชีโวเตฺวว สงฺขํ คจฺฉติ.}

\prose{นนุ มหาราช เถโรปิ สาริปุตฺโต อตฺถงฺคเต สูริเย รตฺติภาเค คิลาโน สมาโน เถเรน มหา\-โมคฺคลฺลาเนน เภสชฺชํ ปุจฺฉียมาโน วาจํ ภินฺทิ, ตสฺส เตน วจีเภเทน เภสชฺชํ อุปฺปชฺชิ, อถ เถโร สาริปุตฺโต “วจีเภเทน เม อิมํ เภสชฺชํ อุปฺปนฺนํ, มา เม อาชีโว ภิชฺชี”ติ อาชีว\-เภท\-ภยา ตํ เภสชฺชํ ปชหิ น อุปชีวิ, เอวมฺปิ วจีวิญฺญตฺติ สาวชฺชา, ตาย จ วิญฺญาปิตํ อริยา น ปริภุญฺชนฺติ, โส จ ปุคฺคโล อริยานํ สมเย โอญฺญาโต โหติ หีฬิโต ขีฬิโต ครหิโต ปริภูโต อจิตฺตีกโต, ภินฺนาชีโวเตฺวว สงฺขํ คจฺฉติ.}

\prose{กตมา วจีวิญฺญตฺติ อนวชฺชา, อิธ มหาราช ภิกฺขุ สติ ปจฺจเย เภสชฺชํ วิญฺญาเปติ ญาติปวาริเตสุ กุเลสุ, อยํ วจีวิญฺญตฺติ อนวชฺชา, ตาย จ วิญฺญาปิตํ อริยา ปริภุญฺชนฺติ, โส จ ปุคฺคโล อริยานํ สมเย วณฺณิโต โหติ โถมิโต ปสตฺโถ, ปริสุทฺธาชีโวเตฺวว สงฺขํ คจฺฉติ, อนุมโต ตถาคเตหิ อรหนฺเตหิ สมฺมา\-สมฺพุทฺเธหิ.}

\prose{ยํ ปน มหาราช ตถาคโต กสิภารทฺวาชสฺส พฺรหฺมณสฺส โภชนํ ปชหิ\csromansharedfootnote{ddc9d2a9e7cf601a}{ปชหติ (ก)}, ตํ อาเวฐน\-วินิ\-เวฐน\-กฑฺฒน\-นิคฺคห\-ปฺปฏิกมฺเมน นิพฺพตฺติ, ตสฺมา ตถาคโต ตํ ปิณฺฑปาตํ ปฏิกฺขิปิ น อุปชีวีติ.}

\prose{สพฺพกาลํ ภนฺเต นาคเสน ตถาคเต ภุญฺชมาเน เทวตา ทิพฺพํ โอชํ ปตฺเต อากิรนฺติ, อุทาหุ “สูกรมทฺทเว จ มธุปายาเส จา”ติ ทฺวีสุ เยว ปิณฺฑปาเตสุ อากิริํสูติ. สพฺพกาลํ มหาราช ตถาคเต ภุญฺชมาเน เทวตา ทิพฺพํ โอชํ คเหตฺวา อุปติฏฺฐิตฺวา อุทฺธฏุทฺธเฏ อาโลเป อากิรนฺติ.}

\prosewithcloserruled{\csromanfolio{226}ยถา มหาราช รญฺโญ สูโท รญฺโญ ภุญฺชนฺตสฺส สูปํ คเหตฺวา อุปติฏฺฐิตฺวา กพเฬ กพเฬ สูปํ อากิรติ, เอวเมว โข มหาราช สพฺพกาลํ ตถาคเต ภุญฺชมาเน เทวตา ทิพฺพํ โอชํ คเหตฺวา อุปติฏฺฐิตฺวา อุทฺธฏุทฺธเฏ อาโลเป ทิพฺพํ โอชํ อากิรนฺติ. เวรญฺชายมฺปิ มหาราช ตถาคตสฺส สุกฺข\-ยว\-ปุลเก\csromansharedfootnote{0466f41ac3fbf4be}{สุกฺขยวมูลเก (ก)} ภุญฺชมานสฺส เทวตา ทิพฺเพน โอเชน เตมยิตฺวา เตมยิตฺวา อุปสํหริํสุ, เตน ตถาคตสฺส กาโย อุปจิโต อโหสีติ. ลาภา วต ภนฺเต นาคเสน ตาสํ เทวตานํ, ยา ตถาคตสฺส สรีร\-ปฺปฏิชคฺคเน สตตํ สมิตํ อุสฺสุกฺกมาปนฺนา. สาธุ ภนฺเต นาคเสน, เอวเมตํ ตถา สมฺปฏิจฺฉามีติ.}{คาถาภิคีต\-โภชน\-กถา\-ปโญฺห นวโม.}
\csromanmatikaanchor{mtk.160}
\csromantocmarkref{subsubsection}{10. ธมฺมเทสนาย อปฺโปสฺสุกปญฺห}{mtk.160}
\bookmark[dest={mtk.160},level=3]{10. ธมฺมเทสนาย อปฺโปสฺสุกปญฺห}
\csromanheader{3}{10. \textbf{ธมฺมเทสนาย อปฺโปสฺสุกฺกปญฺห}}
\proseitem{10}{ภนฺเต นาคเสน ตุเมฺห ภณถ “ตถาคเตน จตูหิ จ อสงฺเขฺยยฺเยหิ กปฺปานํ สตสหสฺเสน จ เอตฺถนฺตเร สพฺพญฺญุต\-ญาณํ ปริปาจิตํ มหโต ชนกายสฺส สมุทฺธรณายา”ติ. ปุน จ สพฺพญฺญุตํ ปตฺตสฺส\csromansymbolfootnote{*}{วิ 3. 7; ที 2. 31; ม 1. 224; ม 2. 291; สํ 1. 138 ปิฏฺเฐสุ.}อปฺโปสฺสุกฺกตาย จิตฺตํ นมิ, โน ธมฺมเทสนายาติ.}

\prose{ยถา นาม ภนฺเต นาคเสน อิสฺสาโส วา อิสฺสาสนฺเตวาสี วา พหุเก ทิวเส สงฺคามตฺถาย อุปาสนํ สิกฺขิตฺวา สมฺปตฺเต มหายุทฺเธ โอสกฺเกยฺย, เอวเมว โข ภนฺเต นาคเสน ตถาคเตน จตูหิ จ อสงฺเขฺยยฺเยหิ กปฺปานํ สตสหสฺเสน จ เอตฺถนฺตเร สพฺพญฺญุต\-ญาณํ ปริปาเจตฺวา มหโต ชนกายสฺส สมุทฺธรณาย สพฺพญฺญุตํ ปตฺเตน ธมฺมเทสนาย โอสกฺกิตํ.}

\prose{ยถา วา ปน ภนฺเต ตาคเสน มลฺโล วา มลฺลนฺเตวาสี วา พหุเก ทิวเส นิพฺพุทฺธํ สิกฺขิตฺวา สมฺปตฺเต มลฺลยุทฺเธ โอสกฺเกยฺย, เอวเมว โข ภนฺเต นาคเสน ตถาคเตน จหูหิ จ อสงฺเขฺยยฺเยหิ กปฺปานํ สตสหสฺเสน จ เอตฺถนฺตเร สพฺพญฺญุต\-ญาณํ ปริปาเจตฺวา มหโต ชนกายสฺส สมุทฺธรณาย สพฺพญฺญุตํ ปตฺเตน ธมฺมเทสนาย โอสกฺกิตํ.}

\prose{\csromanfolio{227}กินฺนุ โข ภนฺเต นาคเสน ตถาคเตน ภยา โอสกฺกิตํ, อุทาหุ อปากฏตาย โอสกฺกิตํ, อุทาหุ ทุพฺพลตาย โอสกฺกิตํ, อุทาหุ อสพฺพญฺญุตาย โอสกฺกิตํ, กิํ ตตฺถ การณํ, อิงฺฆ เม ตฺวํ การณํ พฺรูหิ กงฺขา\-วิตรณาย, ยทิ ภนฺเต นาคเสน ตถาคเตน จตูหิ จ อสงฺเขฺยยฺเยหิ กปฺปานํ สตสหสฺเสน จ เอตฺถนฺตเร สพฺพญฺญุต\-ญาณํ ปริปาจิตํ มหโต ชนกายสฺส สมุทฺธรณาย, เตน หิ “สพฺพญฺญุตํ ปตฺตสฺส อปฺโปสฺสุกฺกตาย จิตฺตํ นมิ, โน ธมฺมเทสนายา”ติ ยํ วจนํ, ตํ มิจฺฉา. ยทิ สพฺพญฺญุตํ ปตฺตสฺส อปฺโปสฺสุกฺกตาย จิตฺตํ นมิ, โน ธมฺมเทสนาย, เตน หิ “ตถาคเตน จหูหิ จ อสงฺเขฺยยฺเยหิ กปฺปานํ สตสหสฺเสน จ เอตฺถนฺตเร สพฺพญฺญุต\-ญาณํ ปริปาจิตํ มหโต ชนกายสฺส สมุทฺธรณายา”ติ ตมฺปิ วจนํ มิจฺฉา. อยมฺปิ อุภโต โกฏิโก ปโญฺห คมฺภีโร ทุนฺนิพฺเพโฐ ตวานุปฺปตฺโต, โส ตยา นิพฺพาหิตพฺโพติ.}

\prose{ปริปาจิตญฺจ มหาราช ตถาคเตน จหูหิ จ อสงฺเขฺยยฺเยหิ กปฺปานํ สตสหสฺเสน จ เอตฺถนฺตเร สพฺพญฺญุต\-ญาณํ มหโต ชนกายสฺส สมุทฺธรณาย, ปตฺต\-สพฺพญฺญุ\-ตสฺส จ อปฺโปสฺสุกฺกตาย จิตฺตํ นมิ, โน ธมฺมเทสนาย. ตญฺจ ปน ธมฺมสฺส คมฺภีรนิปุณทุทฺทสทุรนุโพธสุขุมทุปฺปฏิเวธตํ สตฺตานญฺจ อาลยารามตํ สกฺกาย\-ทิฏฺฐิยา ทฬฺห\-สุคฺคหิต\-ตญฺจ ทิสฺวา “กินฺนุ โข, กถํ นุ โข”ติ อปฺโปสฺสุกฺกตาย จิตฺตํ นมิ, โน ธมฺมเทสนาย, สตฺตานํ ปฏิเวธ\-จินฺตน\-มานสํ เยเวตํ.}

\prose{ยถา มหาราช ภิสกฺโก สลฺลกตฺโต อเนก\-พฺยาธิ\-ปริปีฬิตํ นรํ อุปสงฺกมิตฺวา เอวํ จินฺตยติ “เกน นุ โข อุปกฺกเมน กตเมน วา เภสชฺเชน อิมสฺส พฺยาธิ วูปสเมยฺยา”ติ, เอวเมว โข มหาราช ตถาคตสฺส สพฺพ\-กิเลสพฺยาธิ\-ปริปีฬิตํ ชนํ ธมฺมสฺส จ คมฺภีรนิปุณทุทฺทสทุรนุโพธสุขุมทุปฺปฏิเวธตํ ทิสฺวา “กินฺนุ โข, กถํ นุ โข”ติ อปฺโปสฺสุกฺกตาย จิตฺตํ นมิ, โน ธมฺมเทสนาย, สตฺตานํ ปฏิเวธ\-จินฺตน\-มานสํ เยเวตํ.}

\prose{ยถา มหาราช รญฺโญ ขตฺติยสฺส มุทฺธา\-วสิตฺตสฺส โทวาริก- อนีกฏฺฐปาริสชฺชเนคมภฏพล\csromansharedfootnote{5294446c6739fca7}{พลตฺถ (สี, อิ)} อมจฺจ\-ราชญฺญ\-ราชู\-ปชีวิเน ชเน ทิสฺวา เอวํ}

\prose{\csromanfolio{228}จิตฺตมุ\-ปฺปชฺเชยฺย “กินฺนุ โข, กถํ นุ โข อิเม สงฺคณฺหิสฺสามี”ติ, เอวเมว โข มหาราช ตถาคตสฺส ธมฺมสฺส คมฺภีรนิปุณทุทฺทสทุรนุโพธสุขุมทุปฺปฏิเวธตํ สตฺตานญฺจ อาลยารามตํ สกฺกาย\-ทิฏฺฐิยา ทฬฺห\-สุคฺคหิต\-ตญฺจ ทิสฺวา “กินฺนุ โข, กถํ นุ โข”ติ อปฺโปสฺสุกฺกตาย จิตฺตํ นมิ, โน ธมฺมเทสนาย, สตฺตานํ ปฏิเวธ\-จินฺตน\-มานสํ เยเวตํ.}

\prose{อปิ จ มหาราช สพฺเพสํ ตถาคตานํ ธมฺมตา เอสา, ยํ พฺรหฺมุนา อายาจิตา ธมฺมํ เทเสนฺติ, ตตฺถ ปน กิํ การณํ, เย เตน สมเยน มนุสฺสา ตาปส\-ปริพฺพาชกา สมณพฺราหฺมณา, สพฺเพเต พฺรหฺมเทวตา โหนฺติ พฺรหฺมครุกา พฺรหฺมปรายณา, ตสฺมา ตสฺส พลวโต ยสวโต ญาตสฺส ปญฺญาตสฺส อุตฺตรสฺส อจฺจุคฺคตสฺส โอนมเนน สเทวโก โลโก โอนมิสฺสติ โอกปฺเปสฺสติ อธิมุจฺจิสฺสตีติ อิมินา จ มหาราช การเณน ตถาคตา พฺรหฺมุนา อายาจิตา ธมฺมํ เทเสนฺติ.}

\prosewithcloserruled{ยถา มหาราช โกจิ ราชา วา ราชมหามตฺโต วา ยสฺส โอนมติ อปจิติํ กโรติ, พลวตรสฺส ตสฺส โอนมเนน อวเสสา ชนตา โอนมติ อปจิติํ กโรติ, เอวเมว โข มหาราช พฺรหฺเม โอนมิเต ตถาคตานํ สเทวโก โลโก โอนมิสฺสติ, ปูชิตปูชโก มหาราช โลโก, ตสฺมา โส พฺรหฺมา สพฺเพสํ ตถาคตานํ อายาจติ ธมฺมเทนาย, เตน จ การเณน ตถาคตา พฺรหฺมุนา อายาจิตา ธมฺมํ เทเสนฺตีติ. สาธุ ภนฺเต นาคเสน สุนิพฺเพฐิโต ปโญฺห, อติภทฺรกํ เวยฺยากรณํ, เอวเมตํ ตถา สมฺปฏิจฺฉามีติ.}{ธมฺมเทสนาย อปฺโปสฺสุกฺก\-ปโญฺห ทสโม.}
\csromanmatikaanchor{mtk.161}
\csromantocmarkref{subsubsection}{11. อาจริยานาจริยปญฺห}{mtk.161}
\bookmark[dest={mtk.161},level=3]{11. อาจริยานาจริยปญฺห}
\csromanheader{3}{11. \textbf{อาจริยา}\-\textbf{นาจริย}\-\textbf{ปญฺห}}
\prose{11. ภนฺเต นาคเสน ภาสิตมฺเปตํ ภควตา\csromandash{}}

\setlength{\csromangathapagewidth}{0pt}
\setlength{\csromangathawakwidth}{0pt}
\csromangathameasuregroup{“น เม อาจริโย อตฺถิ, \\ สเทวกสฺมิํ โลกสฺมิํ,}{\csromangathabat{“น เม อาจริโย อตฺถิ,}{สทิโส เม น วิชฺชติ.} \\ \csromangathabat{สเทวกสฺมิํ โลกสฺมิํ,}{นตฺถิ เม ปฏิปุคฺคโล”ติ\textsuperscript{1}.}}
\csromangathasetleft{“น เม อาจริโย อตฺถิ, \\ สเทวกสฺมิํ โลกสฺมิํ,}
\csromangathagroup{\csromangathastanza{\csromangathabat{“น เม อาจริโย อตฺถิ,}{สทิโส เม น วิชฺชติ.} \\ \csromangathabat{สเทวกสฺมิํ โลกสฺมิํ,}{นตฺถิ เม ปฏิปุคฺคโล”ติ\csromansharedfootnote{0a3c745b212b6769}{วิ 3. 12; ม 1. 227; ม 2. 295 ปิฏฺเฐสุ.}.}}}

\prose{ปุน จ ภณิตํ\csromansymbolfootnote{*}{ม 1. 221; ม 2. 281 ปิฏฺเฐสุ.}“อิติ โข ภิกฺขเว อาฬาโร กาลาโม อาจริโย เม สมาโน อนฺเตวาสิํ มํ สมานํ อตฺตนา สมสมํ ฐเปสิ, \csromanfolio{229}อุฬาราย จ มํ ปูชาย ปูเชสี”ติ. ยทิ ภนฺเต นาคเสน ตถาคเตน ภณิตํ “น เม อาจริโย อตฺถิ, สทิโส เม น วิชฺชติ. สเทวกสฺมิํ โลกสฺมิํ, นตฺถิ เม ปฏิปุคฺคโล”ติ, เตน หิ “อิติ โข ภิกฺขเว อาฬาโร กาลาโม อาจริโย เม สมาโน อนฺเตวาสิํ มํ สมานํ อตฺตนา สมสมํ ฐเปสี”ติ ยํ วจนํ, ตํ มิจฺฉา. ยทิ ตถาคเตน ภณิตํ “อิติ โข ภิกฺขเว อาฬาโร กาลาโม อาจริโย เม สมาโน อนฺเตวาสิํ มํ สมานํ อตฺตนา สมสมํ ฐเปสี”ติ, เตน หิ “น เม อาจริโย อตฺถิ, สทิโส เม น วิชฺชติ. สเทวกสฺมิํ โลกสฺมิํ, นตฺถิ เม ปฏิปุคฺคโล”ติ ตมฺปิ วจนํ มิจฺฉา. อยมฺปิ อุภโต โกฏิโก ปโญฺห ตวานุปฺปตฺโต, โส ตยา นิพฺพาหิตพฺโพติ.}

\proseitem{10}{ภาสิตมฺเปตํ มหาราช ตถาคเตน “น เม อาจริโย อตฺถิ, สทิโส เม น วิชฺชติ. สเทวกสฺมิํ โลกสฺมิํ, นตฺถิ เม ปฏิปุคฺคโล”ติ, ภณิตญฺจ “อิติ โข ภิกฺขเว อาฬาโร กาลาโม อาจริโย เม สมาโน อนฺเตวาสิํ มํ สมานํ อตฺตนา สมสมํ ฐเปสิ, อุฬาราย จ มํ ปูชาย ปูเชสี”ติ.}

\prose{ตญฺจ ปน วจนํ ปุพฺเพว สมฺโพธา อนภิสมฺพุทฺธสฺส โพธิสตฺตสฺเสว สโต อาจริยภาวํ สนฺธาย ภาสิตํ.}

\prose{“ปญฺจิเม มหาราช ปุพฺเพว สมฺโพธา อนภิสมฺพุทฺธสฺส โพธิสตฺตสฺส สโต อาจริยา, เยหิ อนุสิฏฺโฐ โพธิสตฺโต ตตฺถ ตตฺถ ทิวสํ วีตินาเมสิ, กตเม ปญฺจ, เย เต มหาราช อฏฺฐ พฺราหฺมณา ชาตมตฺเต โพธิสตฺเต ลกฺขณานิ ปริคฺคณฺหิํสุ, เสยฺยถีทํ, ราโม ธโช ลกฺขโณ มนฺตี ยญฺโญ สุยาโม สุโภโช สุทตฺโตติ. เต ตสฺส โสตฺถิํ ปเวทยิตฺวา รกฺขากมฺมํ อกํสุ, เต จ ปฐมํ อาจริยา.}

\prose{ปุน จปรํ มหาราช โพธิสตฺตสฺส ปิตา สุทฺโธทโน ราชา ยํ เตน สมเยน อภิชาตํ อุทิจฺจ\-ชาติมนฺตํ ปทกํ เวยฺยากรณํ ฉฬงฺควนฺตํ สพฺพมิตฺตํ นาม พฺราหฺมณํ อุปเนตฺวา โสวณฺเณน ภิงฺคาเรน\csromansharedfootnote{08178526da25de7c}{ภิงฺกาเรน (สี, อิ)} อุทกํ โอโณเชตฺวา “อิมํ กุมารํ สิกฺขาเปหี”ติ อทาสิ, อยํ ทุติโย อาจริโย.}

\proseitem{4}{\csromanfolio{230}ปุน จปรํ มหาราช ยา สา เทวตา โพธิสตฺตํ สํเวเชสิ, ยสฺสา วจนํ สุตฺวา โพธิสตฺโต สํวิคฺโค อุพฺพิคฺโค ตสฺมิํ เยว ขเณ เนกฺขมฺมํ นิกฺขมิตฺวา ปพฺพชิ, อยํ ตติโย อาจริโย.}

\prose{ปุน จปรํ มหาราช อาฬาโร กาลาโม อากิญฺจญฺญา\-ยตนสฺส ปริกมฺมํ อาจิกฺขิ\csromansharedfootnote{09528490c12d308c}{อาจิกฺขติ (ก)}, อยํ จตุตฺโถ อาจริโย.}

\prose{ปุน จปรํ มหาราช อุทโก รามปุตฺโต เนว\-สญฺญา\-นา\-สญฺญา\-ยตนสฺส ปริกมฺมํ อาจิกฺขิ\csromansharedfootnote{09528490c12d308c}{อาจิกฺขติ (ก)}, อยํ ปญฺจโม อาจริโย. อิเม โข มหาราช ปุพฺเพว สมฺโพธา อนภิสมฺพุทฺธสฺส โพธิสตฺตสฺส สโต ปญฺจ อาจริยา, เต จ ปน อาจริยา โลกิเย ธมฺเม, อิมสฺมิญฺจ ปน มหาราช โลกุตฺตเร ธมฺเม สพฺพญฺญุตญาณ\-ปฺปฏิเวธาย นตฺถิ ตถาคตสฺส อนุตฺตโร อนุสาสโก, สยมฺภู มหาราช ตถาคโต อนาจริยโก, ตสฺมา การณา ตถาคเตน ภณิตํ “น เม อาจริโย อตฺถิ, สทิโส เม น วิชฺชติ. สเทวกสฺมิํ โลกสฺมิํ, นตฺถิ เม ปฏิปุคฺคโล”ติ. สาธุ ภนฺเต นาคเสน, เอวเมตํ ตถา สมฺปฏิจฺฉามีติ.}

\vspace{\tipitakaparskip}
\csromancenter{อาจริยา\-นาจริย\-ปโญฺห เอกาทสโม.}
\nitthitammidruled{สนฺถววคฺโค ปญฺจโม.}
\csromancenter{อิมสฺมิํ วคฺเค เอกาทส ปญฺหา.}
\nitthitam{\textbf{เมณฺฑกปโญฺห นิฏฺฐิโต.}}
\csromanmatikaanchor{mtk.162}
\csromanmatikaanchor{mtk.163}
\csromanmatikaanchor{mtk.164}
\csromantocmarknopage{section}{5. อนุมานปญฺห}{mtk.162}
\bookmark[dest={mtk.162},level=1]{5. อนุมานปญฺห}
\csromantocmarknopage{subsection}{1. พุทฺธวคฺค}{mtk.163}
\bookmark[dest={mtk.163},level=2]{1. พุทฺธวคฺค}
\csromantocmarkref{subsubsection}{1. ทฺวินฺนํ พุทฺธานํ อนุปฺปชฺชมานปญฺห}{mtk.164}
\bookmark[dest={mtk.164},level=3]{1. ทฺวินฺนํ พุทฺธานํ อนุปฺปชฺชมานปญฺห}
\csromanmarkh{1. พุทฺธวคฺค}\csromanmarkhii{1. ทฺวินฺนํ พุทฺธานํ อนุปฺปชฺชมานปญฺห}\csromanheadernomark{3}{1. \textbf{พุทฺธวคฺค 1. ทฺวินฺนํ พุทฺธานํ อนุปฺปชฺชมานปญฺห}}
\proseitem{1}{\csromanfolio{231}ภนฺเต นาคเสน ภาสิตมฺเปตํ ภควตา\csromansymbolfootnote{*}{ที 3. 95; ม 3. 110; อํ 1. 29; อภิ 2. 348; ขุ 10. 78 ปิฏฺเฐสุ.}“อฏฺฐานเมตํ ภิกฺขเว อนวกาโส, ยํ เอกิสฺสา โลกธาตุยา เทฺว อรหนฺโต สมฺมาสมฺพุทฺธา อปุพฺพํ อจริมํ อุปฺปชฺเชยฺยุํ, เนตํ ฐานํ วิชฺชตี”ติ. เทเสนฺตา จ ภนฺเต นาคเสน สพฺเพปิ ตถาคตา สตฺตติํส โพธิปกฺขิยธมฺเม เทเสนฺติ, กถยมานา จ จตฺตาริ อริยสจฺจานิ กเถนฺติ, สิกฺขาเปนฺตา จ ตีสุ สิกฺขาสุ สิกฺขาเปนฺติ, อนุสาสมานา จ อปฺปมาท\-ปฺปฏิปตฺติยํ อนุสาสนฺติ. ยทิ ภนฺเต นาคเสน สพฺเพสมฺปิ ตถาคตานํ เอกา เทสนา เอกา กถา เอกา สิกฺขา เอกา อนุสิฏฺฐิ, เกน การเณน เทฺว ตถาคตา เอกกฺขเณ นุปฺปชฺชนฺติ. เอเกนปิ ตาว พุทฺธุปฺปาเทน อยํ โลโก โอภาสชาโต, ยทิ ทุติโย พุทฺโธ ภเวยฺย, ทฺวินฺนํ ปภาย อยํ โลโก ภิยฺโยโส\-มตฺตาย\csromansharedfootnote{ab67f8da3c3dc98c}{ภิยฺโยโส มตฺตาย เทฺว ปทานิ \csromandash{} ม.พ.ป.} โอภาสชาโต ภเวยฺย, โอวทมานา จ เทฺว ตถาคตา สุขํ โอวเทยฺยุํ, อนุสาสมานา จ สุขํ อนุสาเสยฺยุํ, ตตฺถ เม การณํ พฺรูหิ, ยถาหํ นิสฺสํสโย ภเวยฺยนฺติ.}

\prose{อยํ มหาราช ทสสหสฺสี โลกธาตุ เอกพุทฺธธารณี, เอกสฺเสว ตถาคตสฺส คุณํ ธาเรติ, ยทิ ทุติโย พุทฺโธ อุปฺปชฺเชยฺย, นายํ ทสสหสฺสี โลกธาตุ ธาเรยฺย, จเลยฺย กมฺเปยฺย นเมยฺย โอนเมยฺย วินเมยฺย วิกิเรยฺย วิธเมยฺย วิทฺธํเสยฺย, น ฐานมุปคจฺเฉยฺย.}

\prose{ยถา มหาราช นาวา เอก\-ปุริส\-สนฺธารณี\csromansharedfootnote{093339aeb3d290e6}{เอกปุริสสนฺตารณี (สี, อิ)} ภเวยฺย, เอกสฺมิํ ปุริเส อภิรูเฬ สา นาวา สมุปาทิกา\csromansharedfootnote{4236672493c651bd}{สมุทกา (ก)} ภเวยฺย, อถ ทุติโย ปุริโส อาคจฺเฉยฺย ตาทิโส อายุนา วณฺเณน วเยน ปมาเณน กิสถูเลน สพฺพ\-งฺค\-ปจฺจงฺเคน, โส ตํ นาวํ อภิรูเหยฺย, อปิ นุ สา มหาราช นาวา ทฺวินฺนมฺปิ ธาเรยฺยาติ. น หิ ภนฺเต, จเลยฺย กมฺเปยฺย นเมยฺย โอนเมยฺย วินเมยฺย วิกิเรยฺย วิธเมยฺย วิทฺธํเสยฺย, น ฐานมุปคจฺเฉยฺย, โอสีเทยฺย อุทเกติ. เอวเมว โข มหาราช อยํ ทสสหสฺสี \csromanfolio{232}โลกธาตุ เอกพุทฺธธารณี, เอกสฺเสว ตถาคตสฺส คุณํ ธาเรติ, ยทิ ทุติโย พุทฺโธ อุปฺปชฺเชยฺย, นายํ ทสสหสฺสี โลกธาตุ ธาเรยฺย, จเลยฺย กมฺเปยฺย นเมยฺย โอนเมยฺย วินเมยฺย วิกิเรยฺย วิธเมยฺย วิทฺธํเสยฺย, น ฐานมุปคจฺเฉยฺย.}

\prose{ยถา วา ปน มหาราช ปุริโส ยาวทตฺถํ โภชนํ ภุญฺเชยฺย ฉาเทนฺตํ ยาว กณฺฐมภิปูรยิตฺวา, โส ธาโต ปีณิโต ปริปุณฺโณ นิรนฺตโร ตนฺทิกโต อโนนมิ\-ต\-ทณฺฑ\-ชาโต ปุนเทว ตตฺตกํ โภชนํ ภุญฺเชยฺย, อปิ นุ โข โส มหาราช ปุริโส สุขิโต ภเวยฺยาติ. น หิ ภนฺเต, สกิํ ภุตฺโตว มเรยฺยาติ\csromansharedfootnote{b0f5da45c39adf9c}{ภุตฺโต วเมยฺยาติ (ก)}. เอวเมว โข มหาราช อยํ ทสสหสฺสี โลกธาตุ เอกพุทฺธธารณี, เอกสฺเสว ตถาคตสฺส คุณํ ธาเรติ, ยทิ ทุติโย พุทฺโธ อุปฺปชฺเชยฺย, นายํ ทสสหสฺสี โลกธาตุ ธาเรยฺย, จเลยฺย กมฺเปยฺย นเมยฺย โอนเมยฺย วินเมยฺย วิกิเรยฺย วิธเมยฺย วิทฺธํเสยฺย, น ฐานมุปคจฺเฉยฺยาติ.}

\prose{กิํ นุ โข ภนฺเต นาคเสน อติธมฺม\-ภาเรน ปถวี จลตีติ. อิธ มหาราช เทฺว สกฏา รตน\-ปริปูริตา ภเวยฺยุํ ยาว มุขสมา, เอกสฺมา สกฏโต รตนํ คเหตฺวา เอกสฺมิํ สกเฏ อากิเรยฺยุํ, อปิ นุ โข ตํ มหาราช สกฏํ ทฺวินฺนมฺปิ สกฏานํ รตนํ ธาเรยฺยาติ. น หิ ภนฺเต, นาภิปิ ตสฺส ผเลยฺย, อราปิ ตสฺส ภิชฺเชยฺยุํ, เนมิปิ ตสฺส โอปเตยฺย, อกฺโขปิ ตสฺส ภิชฺเชยฺยาติ. กิํ นุ โข มหาราช อติรตน\-ภาเรน สกฏํ ภิชฺชตีติ. กิํ นุ โข มหาราช อติรตน\-ภาเรน สกฏํ ภิชฺชตีติ. อาม ภนฺเตติ. เอวเมว โข มหาราช อติธมฺม\-ภาเรน ปถวี จลติ.}

\prose{อปิ จ มหาราช อิมํ การณํ พุทฺธพล\-ปริทีปนาย โอสาริตํ, อญฺญมฺปิ ตตฺถ อภิรูปํ การณํ สุโณหิ, เยน การเณน เทฺว สมฺมาสมฺพุทฺธา เอกกฺขเณ นุปฺปชฺชนฺติ. ยทิ มหาราช เทฺว สมฺมาสมฺพุทฺธา เอกกฺขเณ อุปฺปชฺเชยฺยุํ, เตสํ ปริสาย วิวาโท อุปฺปชฺเชยฺย “ตุมฺหากํ พุทฺโธ, อมฺหากํ พุทฺโธ”ติ, อุภโต ปกฺขชาตา ภเวยฺยุํ, ยถา มหาราช ทฺวินฺนํ พลวามจฺจานํ ปริสาย วิวาโท อุปฺปชฺเชยฺย “ตุมฺหากํ อมจฺโจ, อมฺหากํ อมจฺโจ”ติ, อุภโต ปกฺขชาตา โหนฺติ, เอวเมว โข มหาราช ยทิ เทฺว สมฺมาสมฺพุทฺธา เอกกฺขเณ อุปฺปชฺเชยฺยุํ, เตสํ ปริสาย วิวาโท อุปฺปชฺเชยฺย \csromanfolio{233}“ตุมฺหากํ พุทฺโธ, อมฺหากํ พุทฺโธ”ติ, อุภโต ปกฺขชาตา ภเวยฺยุํ, อิทํ ตาว มหาราช เอกํ การณํ, เยน การเณน เทฺว สมฺมาสมฺพุทฺธา เอกกฺขเณ นุปฺปชฺชนฺติ.}

\prose{อปรมฺปิ มหาราช อุตฺตริํ การณํ สุโณหิ, เยน การเณน เทฺว สมฺมาสมฺพุทฺธา เอกกฺขเณ นุปฺปชฺชนฺติ, ยทิ มหาราช เทฺว สมฺมาสมฺพุทฺธา เอกกฺขเณ อุปฺปชฺเชยฺยุํ, “อคฺโค พุทฺโธ”ติ ยํ วจนํ, ตํ มิจฺฉา ภเวยฺย, “เชฏฺโฐ พุทฺโธ”ติ ยํ วจนํ, ตํ มิจฺฉา ภเวยฺย, “เสฏฺโฐ พุทฺโธ”ติ, “วิสิฏฺโฐ พุทฺโธ”ติ, “อุตฺตโม พุทฺโธ”ติ, “ปวโร พุทฺโธ”ติ, “อสโม พุทฺโธ”ติ, “อสมสโม พุทฺโธ”ติ, “อปฺปฏิโม พุทฺโธ”ติ, “อปฺปฏิภาโค พุทฺโธ”ติ, “อปฺปฏิปุคฺคโล พุทฺโธ”ติ ยํ วจนํ, ตํ มิจฺฉา ภเวยฺย, อิทมฺปิ โข ตฺวํ มหาราช การณํ อตฺถโต สมฺปฏิจฺฉ, เยน การเณน เทฺว สมฺมาสมฺพุทฺธา เอกกฺขเณ นุปฺปชฺชนฺติ.}

\prose{อปิ จ โข มหาราช พุทฺธานํ ภควนฺตานํ สภาวปกติ เอสายํ, เอโก เยว พุทฺโธ โลเก อุปฺปชฺชติ, กสฺมา การณา, มหนฺตตาย สพฺพญฺญุ\-พุทฺธคุณานํ, อญฺญมฺปิ มหาราช ยํ โลเก มหนฺตํ, ตํ เอกํ เยว โหติ. ปถวี มหาราช มหนฺตี. สา เอกา เยว. สาคโร มหนฺโต, โส เอโก เยว. สิเนรุ คิริราชา มหนฺโต, โส เอโก เยว. อากาโส มหนฺโต, โส เอโก เยว. สกฺโก มหนฺโต, โส เอโก เยว. มาโร มหนฺโต, โส เอโก เยว. มหาพฺรหฺมา มหนฺโต, โส เอโก เยว. ตถาคโต อรหํ สมฺมาสมฺพุทฺโธ มหนฺโต, โส เอโก เยว โลกสฺมิํ. ยตฺถ เต อุปฺปชฺชนฺติ, ตตฺถ อญฺญสฺส โอกาโส น โหติ, ตสฺมา มหาราช ตถาคโต อรหํ สมฺมาสมฺพุทฺโธ เอโก เยว โลกสฺมิํ อุปฺปชฺชตีติ.}

\prose{สุกถิโต ภนฺเต นาคเสน ปโญฺห โอปมฺเมหิ การเณหิ, อนิปุโณเปตํ สุตฺวา อตฺตมโน ภเวยฺย, กิํ ปน มาทิโส มหาปญฺโญ. สาธุ ภนฺเต นาคเสน, เอวเมตํ ตถา สมฺปฏิจฺฉามีติ.}

\vspace{\tipitakaparskip}
\csromancenter{ทฺวินฺนํ พุทฺธานํ อนุปฺปชฺชนปโญฺห ปฐโม.}
\csromanmatikaanchor{mtk.165}
\csromantocmarkref{subsubsection}{2. โคตมิวตฺถทานปญฺห}{mtk.165}
\bookmark[dest={mtk.165},level=3]{2. โคตมิวตฺถทานปญฺห}
\csromanheader{3}{2. \textbf{โคตมิ}\-\textbf{วตฺถทาน}\-\textbf{ปญฺห}}
\proseitem{2}{\csromanfolio{234}ภนฺเต นาคเสน ภาสิตมฺเปตํ ภควตา มาตุจฺฉาย มหาปชาปติยา โคตมิยา วสฺสิก\-สาฏิกาย ทียมานาย\csromansymbolfootnote{*}{ม 3. 296 ปิฏฺเฐ.}“สํเฆ โคตมิ เทหิ, สํเฆ เต ทินฺเน อหญฺเจว ปูชิโต ภวิสฺสามิ สํโฆ จา”ติ. กิํ นุ โข ภนฺเต นาคเสน ตถาคโต สํฆรตนโต น ภาริโก น ครุโก น ทกฺขิเณยฺโย, ยํ ตถาคโต สกาย มาตุจฺฉาย สยํ ปิญฺชิตํ\csromansharedfootnote{fcaef59533ceefc3}{ปิจฺฉิตํ (สี, อิ)} สยํ ลุญฺจิตํ สยํ โปถิตํ สยํ กนฺติตํ สยํ วายิตํ วสฺสิกสาฏิกํ อตฺตโน ทียมานํ สํฆสฺส ทาเปสิ. ยทิ ภนฺเต นาคเสน ตถาคโต สํฆรตนโต อุตฺตโร ภเวยฺย อธิโก วา วิสิฏฺโฐ วา, “มยิ ทินฺเน มหปฺผลํ ภวิสฺสตี”ติ น ตถาคโต มาตุจฺฉาย สยํ ปิญฺชิตํ สยํ ลุญฺจิตํ สยํ โปถิตํ วสฺสิกสาฏิกํ สํเฆ ทาเปยฺย, ยสฺมา จ โข ภนฺเต นาคเสน ตถาคโต อตฺตานํ น ปตฺถยติ\csromansharedfootnote{b74cb3fe0c09aa5e}{น ปตฺตียติ (สี, อิ)} น อุปนิสฺสยติ, ตสฺมา ตถาคโต มาตุจฺฉาย ตํ วสฺสิกสาฏิกํ สํฆสฺส ทาเปสีติ.}

\prose{ภาสิตมฺเปตํ มหาราช ภควตา มาตุจฺฉาย มหาปชาติยา โคตมิยา วสฺสิก\-สาฏิกาย ทียมานาย “สํเฆ โคตมิ เทหิ, สํเฆ เต ทินฺเน อหญฺเจว ปูชิโต ภวิสฺสามิ สํโฆ จา”ติ. ตํ ปน น อตฺตโน ปติมานนสฺส อวิปากตาย น อทกฺขิเณยฺยตาย, อปิ จ โข มหาราช หิตตฺถาย อนุกมฺปาย อนาคตมทฺธานํ สํโฆ มมจฺจเยน จิตฺตีกโต ภวิสฺสตีติ วิชฺชมาเน เยว คุเณ ปริกิตฺตยนฺโต เอวมาห “สํเฆ โคตมิ เทหิ, สํเฆ เต ทินฺเน อหญฺเจว ปูชิโต ภวิสฺสามิ สํโฆ จา”ติ.}

\prose{ยถา มหาราช ปิตา ธรมาโน เยว อมจฺจภฏพลโทวาริก- อนีกฏฺฐปาริสชฺชชนมชฺเฌ รญฺโญ สนฺติเก ปุตฺตสฺส วิชฺชมานํ เยว คุณํ ปกิตฺเตติ\csromansharedfootnote{5f659748ced3b732}{ปริกิตฺเตติ (ก)} “อิธ ฐปิโต อนาคตมทฺธานํ ชนมชฺเฌ ปูชิโต ภวิสฺสตี”ติ. เอวเมว โข มหาราช ตถาคโต หิตตฺถาย อนุกมฺปาย อนาคตมทฺธานํ สํโฆ มมจฺจเยน จิตฺตีกโต ภวิสฺสตีติ วิชฺชมาเน เยว คุเณ ปกิตฺตยนฺโต เอวมาห “สํเฆ โคตมิ เทหิ, สํเฆ เต ทินฺเน อหญฺเจว ปูชิโต ภวิสฺสามิ สํโฆ จา”ติ.}

\prose{\csromanfolio{235}น โข มหาราช ตาวตเกน วสฺสิกสาฏิกา\-นุปฺปทาน\-มตฺตเกน สํโฆ ตถาคตโต อธิโก นาม โหติ วิสิฏฺโฐ วา. ยถา มหาราช มาตาปิตโร ปุตฺตานํ อุจฺฉาเทนฺติ ปริมทฺทนฺติ นหาเปนฺติ สมฺพาเหนฺติ, อปิ นุ โข มหาราช ตาวตเกน อุจฺฉาทน\-ปริมทฺทน\-นหาปน\-สมฺพาหน\-มตฺตเกน “ปุตฺโต มาตาปิตูหิ อธิโก นาม โหติ วิสิฏฺโฐ วา”ติ. น หิ ภนฺเต, อกามกรณียา ภนฺเต ปุตฺตา มาตาปิตูนํ, ตสฺมา มาตาปิตโร ปุตฺตานํ อุจฺฉาทน\-ปริมทฺทน\-นหาปน\-สมฺพาหนํ กโรนฺตีติ. เอวเมว โข มหาราช น ตาวตเกน วสฺสิกสาฏิกา\-นุปฺปทาน\-มตฺตเกน สํโฆ ตถาคตโต อธิโก นาม โหติ วิสิฏฺโฐ วาติ. อปิ จ ตถาคโต อกามกรณียํ กโรนฺโต มาตุจฺฉาย ตํ วสฺสิกสาฏิกํ สํฆสฺส ทาเปสิ.}

\prose{ยถา วา ปน มหาราช โกจิเทว ปุริโส รญฺโญ อุปายนํ อาหเรยฺย, ตํ ราชา อุปายนํ อญฺญตรสฺส ภฏสฺส วา พลสฺส วา เสนาปติสฺส วา ปุโรหิตสฺส วา ทเทยฺย. อปิ นุ โข โส มหาราช ปุริโส ตาวตเกน อุปายน\-ปฏิลาภ\-มตฺตเกน รญฺญา อธิโก นาม โหติ วิสิฏฺโฐ วาติ. น หิ ภนฺเต, ราชภตฺติโก ภนฺเต โส ปุริโส ราชูปชีวี, ตฏฺฐาเน ฐเปนฺโต ราชา อุปายนํ เทตีติ. เอวเมว โข มหาราช น ตาวตเกน วสฺสิกสาฏิกา\-นุปฺปทาน\-มตฺตเกน สํโฆ ตถาคตโต อธิโก นาม โหติ วิสิฏฺโฐ วา, อถ โข ตถาคต\-ภตฺติโก ตถาคตู\-ปชีวี. ตฏฺฐาเน ฐเปนฺโต ตถาคโต สํฆสฺส วสฺสิกสาฏิกํ ทาเปสิ.}

\prose{อปิ จ มหาราช ตถาคตสฺส เอวํ อโหสิ “สภาวปฏิปูชนีโย สํโฆ, มม สนฺตเกน สํฆํ ปฏิปูเชสฺสามี”ติ สํฆสฺส วสฺสิกสาฏิกํ ทาเปสิ, น มหาราช ตถาคโต อตฺตโน เยว ปฏิปูชนํ วณฺเณติ, อถ โข เย โลเก ปฏิปูชนารหา, เตสมฺปิ ตถาคโต ปฏิปูชนํ วณฺเณติ.}

\prose{ภาสิตมฺเปตํ มหาราช ภควตา เทวาติเทเวน มชฺฌิมนิกาย\-วร\-ลญฺฉเก ธมฺมทายาท\-ธมฺมปริยาเย อปฺปิจฺฉ\-ปฺปฏิปตฺติํ ปกิตฺตยมาเนน “อสุ เยว เม ปุริโม ภิกฺขุ ปุชฺชตโร จ ปาสํสตโร จา”ติ *. นตฺถิ มหาราช ภเวสุ โกจิ สตฺโต ตถาคตโต ทกฺขิเณยฺโย}

\prose{\csromanfolio{236}วา อุตฺตโร วา อธิโก วา วิสิฏฺโฐ วา, ตถาคโตว อุตฺตโร อธิโก วิสิฏฺโฐ.}

\prose{ภาสิตมฺเปตํ มหาราช สํยุตฺตนิกาย\-วเร มาณวคามิเกน เทวปุตฺเตน ภควโต ปุรโต ฐตฺวา เทวมนุสฺส\-มชฺเฌ\csromandash{}}

\setlength{\csromangathapagewidth}{0pt}
\setlength{\csromangathawakwidth}{0pt}
\csromangathameasuregroup{“วิปุโล ราชคหียานํ\textsuperscript{1}, \\ เสโต หิมวตํ เสฏฺโฐ, \\ สมุทฺโท อุทธินํ เสฏฺโฐ, \\ สเทวกสฺส โลกสฺส,}{\csromangathabat{“วิปุโล ราชคหียานํ\textsuperscript{1},}{คิริ เสฏฺโฐ ปวุจฺจติ.} \\ \csromangathabat{เสโต หิมวตํ เสฏฺโฐ,}{อาทิจฺโจ อฆคามินํ.} \\ \csromangathabat{สมุทฺโท อุทธินํ เสฏฺโฐ,}{นกฺขตฺตานญฺจ จนฺทิมา.} \\ \csromangathabat{สเทวกสฺส โลกสฺส,}{พุทฺโธ อคฺโค ปวุจฺจตี”ติ.}}
\csromangathasetleft{“วิปุโล ราชคหียานํ\textsuperscript{1}, \\ เสโต หิมวตํ เสฏฺโฐ, \\ สมุทฺโท อุทธินํ เสฏฺโฐ, \\ สเทวกสฺส โลกสฺส,}
\csromangathagroup{\csromangathastanza{\csromangathabat{“วิปุโล ราชคหียานํ\csromansharedfootnote{ec5b4d6b087e06e3}{ราชคหิกานํ (ก) สํ 1. 67 ปิฏฺเฐ.},}{คิริ เสฏฺโฐ ปวุจฺจติ.} \\ \csromangathabat{เสโต หิมวตํ เสฏฺโฐ,}{อาทิจฺโจ อฆคามินํ.}}\csromangathastanza{\csromangathabat{สมุทฺโท อุทธินํ เสฏฺโฐ,}{นกฺขตฺตานญฺจ จนฺทิมา.} \\ \csromangathabat{สเทวกสฺส โลกสฺส,}{พุทฺโธ อคฺโค ปวุจฺจตี”ติ.}}}

\prose{ตา โข ปเนตา มหาราช มาณวคามิเกน เทวปุตฺเตน คาถา สุคีตา น ทุคฺคีตา, สุภาสิตา น ทุพฺภาสิตา, อนุมตา จ ภควตา, นนุ มหาราช เถเรนปิ สาริปุตฺเตน ธมฺม\-เสนาปตินา ภณิตํ\csromandash{}}

\prose{“เอโก มโนปสาโท.}

\prose{สรณคมนม\-ญฺชลิ\-ปณาโม วา.}

\setlength{\csromangathapagewidth}{0pt}
\setlength{\csromangathawakwidth}{0pt}
\csromangathameasuregroup{}{อุสฺสหเต ตารยิตุํ, \\ มารพลนิสูทเน พุทฺเธ”ติ.}
\csromangathameasurewak{อุสฺสหเต ตารยิตุํ, \\ มารพลนิสูทเน พุทฺเธ”ติ.}
\setlength{\csromangathaleftcol}{0pt}
\csromangathagroup{\csromangathastanza{\csromangathawak{อุสฺสหเต ตารยิตุํ, \\ มารพลนิสูทเน พุทฺเธ”ติ.}}}

\prosewithcloserruled{ภควตา จ ภณิตํ เทวาติเทเวน\csromansymbolfootnote{*}{อํ 1. 21 ปิฏฺเฐ.}“เอกปุคฺคโล ภิกฺขเว โลเก อุปฺปชฺชมาโน อุปฺปชฺชติ พหุชนหิตาย พหุชน\-สุขาย โลกานุกมฺปาย อตฺถาย หิตาย สุขาย เทวมนุสฺสานํ, กตโม เอกปุคฺคโล, ตถาคโต อรหํ สมฺมาสมฺพุทฺโธ ฯเปฯ เทวมนุสฺสานนฺ”ติ. สาธุ ภนฺเต นาคเสน, เอวเมตํ ตถา สมฺปฏิจฺฉามีติ.}{โคตมิ\-วตฺถทาน\-ปโญฺห ทุติโย.}
\csromanmatikaanchor{mtk.166}
\csromantocmarkref{subsubsection}{3. คิหิปพฺพชิตสมฺมาปฏิปตฺติปญฺห}{mtk.166}
\bookmark[dest={mtk.166},level=3]{3. คิหิปพฺพชิตสมฺมาปฏิปตฺติปญฺห}
\csromanheader{3}{3. \textbf{คิหิปพฺพชิต}\-\textbf{สมฺมาปฏิปตฺติ}\-\textbf{ปญฺห}}
\prose{(3) ภนฺเต นาคเสน ภาสิตมฺเปตํ ภควตา\csromansymbolfootnote{+}{อํ 1. 70 ปิฏฺเฐ.}“คิหิโน วาหํ ภิกฺขเว ปพฺพชิตสฺส วา สมฺมา\-ปฏิปตฺติํ วณฺเณมิ, คิหี วา ภิกฺขเว ปพฺพชิโต วา สมฺมาปฏิปนฺโน สมฺมาปฏิปตฺตา\-ธิกรณ\-เหตุ อาราธโก โหติ \csromanfolio{237}ญายํ ธมฺมํ กุสลนฺ”ติ. ยทิ ภนฺเต นาคเสน คิหี โอทาตวสโน กามโภคี ปุตฺตทาร\-สมฺพาธ\-สยนํ อชฺฌาวสนฺโต กาสิกจนฺทนํ ปจฺจนุโภนฺโต มาลา\-คนฺธ\-วิเลปนํ ธาเรนฺโต ชาตรูป\-รชตํ สาทิยนฺโต มณิกุณฺฑล\csromansharedfootnote{dd7898a215063485}{มณิกนก (สี, อิ)} วิจิตฺต\-โมฬิพทฺโธ สมฺมาปฏิปนฺโน อาราธโก โหติ ญายํ ธมฺมํ กุสลํ, ปพฺพชิโตปิ ภณฺฑุ\-กาสาววตฺถ\-วสโน ปรปิณฺฑม\-ชฺฌุปคโต จตูสุ สีลกฺขนฺเธสุ สมฺมา\-ปริปูรการี ทิยฑฺเฒสุ สิกฺขาปท\-สเตสุ สมาทาย วตฺตนฺโต เตรสสุ ธุตคุเณสุ อนวเสสํ วตฺตนฺโต สมฺมาปฏิปนฺโน อาราธโก โหติ ญายํ ธมฺมํ กุสลํ. ตตฺถ ภนฺเต โก วิเสโส คิหิโน วา ปพฺพชิตสฺส วา, อผลํ โหติ ตโปกมฺมํ, นิรตฺถกา ปพฺพชฺชา. วญฺฌา สิกฺขาปท\-โคปนา, โมฆํ ธุตคุณ\-สมาทานํ, กิํ ตตฺถ ทุกฺขม\-นุจิณฺเณน, นนุ นาม สุเขเนว สุขํ อธิคนฺตพฺพนฺติ.}

\proseitem{2}{ภาสิตมฺเปตํ มหาราช ภควตา “คิหิโน วาหํ ภิกฺขเว ปพฺพชิตสฺส วา สมฺมา\-ปฏิปตฺติํ วณฺเณมิ, คิหี วา ภิกฺขเว ปพฺพชิโต วา สมฺมาปฏิปนฺโน สมฺมาปฏิปตฺตา\-ธิกรณ\-เหตุ อาราธโก โหติ ญายํ ธมฺมํ กุสลนฺ”ติ. เอวเมตํ มหาราช สมฺมา\-ปฏิปนฺโนว เสฏฺโฐ, ปพฺพชิโตปิ มหาราช “ปพฺพชิโตมฺหี”ติ น สมฺมา ปฏิปชฺเชยฺย, อถ โข โส อารกาว สามญฺญา, อารกาว พฺรหฺมญฺญา, ปเคว คิหี โอทาตวสโน. คิหีปิ มหาราช สมฺมาปฏิปนฺโน อาราธโก โหติ ญายํ ธมฺมํ กุสลํ, ปพฺพชิโตปิ มหาราช สมฺมาปฏิปนฺโน อาราธโก โหติ ญายํ ธมฺมํ กุสลํ.}

\prose{อปิ จ โข มหาราช ปพฺพชิโตว สามญฺญสฺส อิสฺสโร อธิปติ, ปพฺพชฺชา มหาราช พหุคุณา อเนกคุณา อปฺปมาณคุณา, น สกฺกา ปพฺพชฺชาย คุณํ ปริมาณํ กาตุํ.}

\prose{ยถา มหาราช กามททสฺส มณิรตนสฺส น สกฺกา ธเนน อคฺโฆ ปริมาณํ กาตุํ “เอตฺตกํ มณิรตนสฺส มูลนฺ”ติ, เอวเมว โข มหาราช ปพฺพชฺชา พหุคุณา อเนกคุณา อปฺปมาณคุณา, น สกฺกา ปพฺพชฺชาย คุณํ ปริมาณํ กาตุํ.}

\prose{ยถา วา ปน มหาราช มหาสมุทฺเท อูมิโย น สกฺกา ปริมาณํ กาตุํ “เอตฺตกา มหาสมุทฺเท อูมิโย”ติ, เอวเมว โข มหาราช ปพฺพชฺชา \csromanfolio{238}พหุคุณา อเนกคุณา อปฺปมาณคุณา, น สกฺกา ปพฺพชฺชาย คุณํ ปริมาณํ กาตุํ.}

\proseitemwithcloserruled{5}{ปพฺพชิตสฺส มหาราช ยํ กิญฺจิ กรณียํ, สพฺพํ ตํ ขิปฺปเมว สมิชฺฌติ โน จิรรตฺตาย. กิํ การณา, ปพฺพชิโต มหาราช อปฺปิจฺโฉ โหติ สนฺตุฏฺโฐ ปวิวิตฺโต อสํสฏฺโฐ อารทฺธวีริโย นิราลโย อนิเกโต ปริปุณฺณสีโล สลฺเลขิตาจาโร ธุตปฺปฏิปตฺติ\-กุสโล โหติ, ตํการาณา ปพฺพชิตสฺส ยํ กิญฺจิ กรณียํ, สพฺพํ ตํ ขิปฺปเมว สมิชฺฌติ โน จิรรตฺตาย. ยถา มหาราช นิคฺคณฺฐิสมสุโธตอุชุวิมลนาราโจ สุสชฺชิโต สมฺมา วหติ, เอวเมว โข มหาราช ปพฺพชิตสฺส ยํ กิญฺจิ กรณียํ, สพฺพํ ตํ ขิปฺปเมว สมิชฺฌติ โน จิรรตฺตายาติ. สาธุ ภนฺเต นาคเสน, เอวเมตํ ตถา สมฺปฏิจฺฉามีติ.}{คิหิปพฺพชิต\-สมฺมาปฏิปตฺติ\-ปโญฺห ตติโย.}
\csromanmatikaanchor{mtk.167}
\csromantocmarkref{subsubsection}{4. ปฏิปทาโทสปญฺห}{mtk.167}
\bookmark[dest={mtk.167},level=3]{4. ปฏิปทาโทสปญฺห}
\csromanheader{3}{4. \textbf{ปฏิปทา}\-\textbf{โทส}\-\textbf{ปญฺห}}
\proseitem{5}{ภนฺเต นาคเสน ยทา โพธิสตฺโต ทุกฺกรการิกํ อกาสิ, เนตาทิโส อญฺญตฺร อารมฺโภ อโหสิ นิกฺกโม กิเลสยุทฺธํ มจฺจุเสนํ วิธมนํ อาหารปริคฺคโห ทุกฺกรการิกา, เอวรูเป ปรกฺกเม กิญฺจิ อสฺสาทํ อลภิตฺวา ตเมว จิตฺตํ ปริหาเปตฺวา เอวมโวจ “น โข ปนาหํ อิมาย กฏุกาย ทุกฺกร\-การิกาย อธิคจฺฉามิ อุตฺตริ\-มนุสฺสธมฺมํ อลม\-ริย\-ญาณ\-ทสฺสน\-วิเสสํ, สิยา นุ โข อญฺโญ มคฺโค โพธายา”ติ, ตโต นิพฺพินฺทิตฺวา อญฺเญน มคฺเคน สพฺพญฺญุตํ ปตฺโต, ปุน ตาย ปฏิปทาย สาวเก อนุสาสติ สมาทเปติ.}

\setlength{\csromangathapagewidth}{0pt}
\setlength{\csromangathawakwidth}{0pt}
\csromangathameasuregroup{“อารมฺภถ นิกฺขมถ, \\ ธุนาถ มจฺจุโน เสนํ,}{\csromangathabat{“อารมฺภถ นิกฺขมถ,}{ยุญฺชถ พุทฺธสาสเน.} \\ \csromangathabat{ธุนาถ มจฺจุโน เสนํ,}{นฬาคารํว กุญฺชโร”ติ\textsuperscript{1}.}}
\csromangathasetleft{“อารมฺภถ นิกฺขมถ, \\ ธุนาถ มจฺจุโน เสนํ,}
\csromangathagroup{\csromangathastanza{\csromangathabat{“อารมฺภถ นิกฺขมถ,}{ยุญฺชถ พุทฺธสาสเน.} \\ \csromangathabat{ธุนาถ มจฺจุโน เสนํ,}{นฬาคารํว กุญฺชโร”ติ\csromansharedfootnote{f55da3e3c73adef8}{สํ 1. 158; ขุ 10. 35, 217 ปิฏฺเฐสุ.}.}}}

\prose{เกน นุ โข ภนฺเต นาคเสน การเณน ตถาคโต ยาย ปฏิปทาย อตฺตนา นิพฺพินฺโน วิรตฺตรูโป, ตตฺถ สาวเก อนุสาสติ สมาทเปตีติ.}

\prose{\csromanfolio{239}ตทาปิ มหาราช เอตรหิปิ สา เยว ปฏิปทา, ตํ เยว\csromansharedfootnote{cf7dc7e6d0b8f44d}{ตํเยว นิคฺคหีตสนฺธิ \csromandash{} ม.พ.ป.} ปฏิปทํ ปฏิปชฺชิตฺวา โพธิสตฺโต สพฺพญฺญุตํ ปตฺโต. อปิ จ มหาราช โพธิสตฺโต อติวีริยํ กโรนฺโต นิรวเสสโต อาหารํ อุปรุนฺธิ, ตสฺส อาหารูปโรเธน จิตฺตทุพฺพลฺยํ อุปฺปชฺชิ, โส เตน ทุพฺพเลฺยน นาสกฺขิ สพฺพญฺญุตํ ปาปุณิตุํ, โส มตฺตมตฺตํ กพฬีการา\-หารํ เสวนฺโต ตาเยว ปฏิปทาย นจิรสฺเสว สพฺพญฺญุตํ ปาปุณิ. สา เยว มหาราช ปฏิปทา สพฺเพสํ ตถาคตานํ สพฺพญฺญุตญาณ\-ปฺปฏิลาภาย.}

\prose{ยถา มหาราช สพฺเพสํ สตฺตานํ อาหาโร อุปตฺถมฺโภ, อาหารูปนิสฺสิตา สพฺเพ สตฺตา สุขํ อนุภวนฺติ, เอวเมว โข มหาราช สา เยว ปฏิปทา สพฺเพสํ ตถาคตานํ สพฺพญฺญุตญาณ\-ปฺปฏิลาภาย, เนโส มหาราช โทโส อารมฺภสฺส, น นิกฺกมสฺส, น กิเลส\-ยุทฺธสฺส, เยน ตถาคโต ตสฺมิํ สมเย น ปาปุณิ สพฺพญฺญุต\-ญาณํ, อถ โข อาหารู\-ปโรธสฺเส\-เวโส โทโส, สทา ปฏิยตฺตา เยวสา ปฏิปทา.}

\prose{ยถา มหาราช ปุริโส อทฺธานํ อติเวเคน คจฺเฉยฺย, เตน โส ปกฺขหโต วา ภเวยฺย ปีฐสปฺปี วา อสญฺจโร ปถวิตเล. อปิ นุ โข มหาราช มหาปถวิยา โทโส อตฺถิ, เยน โส ปุริโส ปกฺขหโต อโหสีติ. น หิ ภนฺเต, สทา ปฏิยตฺตา ภนฺเต มหาปถวี, กุโต ตสฺสา โทโส, วายามสฺเสเวโส โทโส, เยน โส ปุริโส ปกฺขหโต อโหสีติ. เอวเมว โข มหาราช เนโส โทโส อารมฺภสฺส, น นิกฺกมสฺส, น กิเลส\-ยุทฺธสฺส, เยน ตถาคโต ตสฺมิํ สมเย น ปาปุณิ สพฺพญฺญุต\-ญาณํ, อถ โข อาหารู\-ปโรธสฺเส\-เวโส โทโส สทา ปฏิยตฺตา เยเวสา ปฏิปทา.}

\prose{ยถา วา ปน มหาราช ปุริโส กิลิฏฺฐํ สาฏกํ นิวาเสยฺย, น โส ตํ โธวาเปยฺย, เนโส โทโส อุทกสฺส, สทา ปฏิยตฺตํ อุทกํ, ปุริสสฺเสเวโส โทโส, เอวเมว โข มหาราช เนโส โทโส อารมฺภสฺส, น นิกฺกมสฺส, น กุเลสยุทฺธสฺส, เยน ตถาคโต ตสฺมิํ สมเย น ปาปุณิ สพฺพญฺญุต\-ญาณํ, อถ โข อาหารู\-ปโรธสฺเส\-เวโส โทโส, สทา ปฏิยตฺตา เยเวสา ปฏิปทา, ตสฺมา ตถาคโต ตาเยว ปฏิปทาย สาวเก อนุสาสติ สมาทเปติ, เอวํ โข มหาราช}

\prosewithcloserruled{\csromanfolio{240}สทา ปฏิยตฺตา อนวชฺชา สา ปฏิปทาติ. สาธุ ภนฺเต นาคเสน, เอวเมตํ ตถา สมฺปฏิจฺฉามีติ.}{ปฏิปทา\-โทส\-ปโญฺห จตุตฺโถ.}
\csromanmatikaanchor{mtk.168}
\csromantocmarkref{subsubsection}{5. หีนายาวตฺตนปญฺห}{mtk.168}
\bookmark[dest={mtk.168},level=3]{5. หีนายาวตฺตนปญฺห}
\csromanheader{3}{5. \textbf{หีนายา}\-\textbf{วตฺตน}\-\textbf{ปญฺห}}
\proseitem{5}{ภนฺเต นาคเสน มหนฺตํ อิทํ ตถาคต\-สาสนํ สารํ วรํ เสฏฺฐํ ปวรํ อนุปมํ ปริสุทฺธํ วิมลํ ปณฺฑรํ อนวชฺชํ, น ยุตฺตํ คิหิํ ตาวตกํ ปพฺพเชตุํ, คิหี เยว\csromansharedfootnote{8e6af8e5121a9a02}{คิหิํ เยว (สี, อิ)} เอกสฺมิํ ผเล วิเนตฺวา ยทา อปุนราวตฺตี โหติ, ตทา โส ปพฺพาเชตพฺโพ. กิํ การณา, อิเม ทุชฺชนา ตาว ตตฺถ สาสเน วิสุทฺเธ ปพฺพชิตฺวา ปฏินิวตฺติตฺวา หีนายาวตฺตนฺติ, เตสํ ปจฺจาคมเนน อยํ มหาชโน เอวํ วิจินฺเตติ “ตุจฺฉกํ วต โภ เอตํ สมณสฺส โคตมสฺส สาสนํ ภวิสฺสติ, ยํ อิเม ปฏินิวตฺตนฺตี”ติ. อิทเมตฺถ การณนฺติ.}

\prose{ยถา มหาราช ตฬาโก ภเวยฺย สมฺปุณฺณสุจิวิมลสีตลสลิโล, อถ โย โกจิ กิลิฏฺโฐ มลกทฺท\-มคโต ตํ ตฬากํ คนฺตฺวา อนหายิตฺวา กิลิฏฺโฐว ปฏินิวตฺเตยฺย, ตตฺถ มหาราช กตมํ ชโน ครเหยฺย กิลิฏฺฐํ วา ตฬากํ วาติ. กิลิฏฺฐํ ภนฺเต ชโน ครเหยฺย “อยํ ตฬากํ คนฺตฺวา อนหายิตฺวา กิลิฏฺโฐว ปฏินิวตฺโต, กิํ อิมํ อนหายิตุกามํ ตฬาโก สยํ นหาเปสฺสติ, โก โทโส ตฬากสฺสา”ติ. เอวเมว โข มหาราช ตถาคโต วิมุตฺติ\-วร\-สลิล\-สมฺปุณฺณํ สทฺธมฺมวร\-ตฬากํ มาเปสิ “เย เกจิ กิเลสมล\-กิลิฏฺฐา สเจตนา พุธา, เต อิธ นหายิตฺวา สพฺพกิเลเส ปวาหยิสฺสนฺตี”ติ. ยทิ โกจิ ตํ สทฺธมฺมวร\-ตฬากํ คนฺตฺวา อนหายิตฺวา สกิเลโสว ปฏินิวตฺติตฺวา หีนายาวตฺตติ. ตํ เยว\csromansharedfootnote{cf7dc7e6d0b8f44d}{ตํเยว นิคฺคหีตสนฺธิ \csromandash{} ม.พ.ป.} ชโน ครหิสฺสติ “อยํ ชินสาสเน ปพฺพชิตฺวา ตตฺถ ปติฏฺฐํ อลภิตฺวา หีนายาวตฺโต, กิํ อิมํ อปฺปฏิปชฺชนฺตํ ชินสาสนํ สยํ โพเธสฺสติ, โก โทโส ชินสาสนสฺสา”ติ.}

\prose{\csromanfolio{241}ยถา วา ปน มหาราช ปุริโส ปรมพฺยาธิโต โรคุ\-ปฺปตฺติ\-กุสลํ อโมฆ\-ธุว\-สิทฺธ\-กมฺมํ ภิสกฺกํ สลฺลกตฺตํ ทิสฺวา อติกิจฺฉาเปตฺวา สพฺยาธิโกว ปฏินิวตฺเตยฺย, ตตฺถ กตมํ ชโน ครเหยฺย อาตุรํ วา ภิสกฺกํ วาติ. อาตุรํ ภนฺเต ชโน ครเหยฺย “อยํ โรคุ\-ปฺปตฺติ\-กุสลํ อโมฆ\-ธุว\-สิทฺธ\-กมฺมํ ภิสกฺกํ สลฺลกตฺตํ ทิสฺวา อติกิจฺฉาเปตฺวา สพฺยาธิโกว ปฏินิวตฺโต, กิํ อิมํ อติกิจฺฉา\-เปนฺตํ ภิสกฺโก สยํ ติกิจฺฉิสฺสติ, โก โทโส ภิสกฺกสฺสา”ติ. เอวเมว โข มหาราช ตถาคโต อนฺโต\-สาสน\-สมุคฺเค เกวลํ สกล\-กิเลสพฺยาธิ\-วูปสมนสมตฺถํ อมโตสธํ ปกฺขิปิ, “เย เกจิ กิเลสพฺยาธิ\-ปีฬิตา สเจตนา พุธา, เต อิมํ อมโตสธํ ปิวิตฺวา สพฺพ\-กิเลสพฺยาธิํ วูปสเมสฺสนฺตี”ติ. ยทิ โกจิ ตํ อมโตสธํ อปิวิตฺวา สกิเลโสว ปฏินิวตฺติตฺวา หีนายาวตฺตติ, ตํ เยว\csromansharedfootnote{cf7dc7e6d0b8f44d}{ตํเยว นิคฺคหีตสนฺธิ \csromandash{} ม.พ.ป.} ชโน ครหิสฺสติ “อยํ ชินสาสเน ปพฺพชิตฺวา ตตฺถ ปติฏฺฐํ อลภิตฺวา หีนายาวตฺโต, กิํ อิมํ อปฺปฏิปชฺชนฺตํ ชินสาสนํ สยํ โพเธสฺสติ, โก โทโส ชินสาสนสฺสา”ติ.}

\prose{ยถา วา ปน มหาราช ฉาโต ปุริโส มหติ\-มหาปุญฺญ\-ภตฺต\-ปริเวสนํ คนฺตฺวา ตํ ภตฺตํ อภุญฺชิตฺวา ฉาโตว ปฏินิวตฺเตยฺย, ตตฺถ กตมํ ชโน ครเหยฺย ฉาตํ วา ปุญฺญภตฺตํ วาติ. ฉาตํ ภนฺเต ชโน ครเหยฺย “อยํ ขุทาปีฬิโต ปุญฺญภตฺตํ ปฏิลภิตฺวา อภุญฺชิตฺวา ฉาโตว ปฏินิวตฺโต, กิํ อิมสฺส อภุญฺชนฺตสฺส โภชนํ สยํ มุขํ ปวิสิสฺสติ, โก โทโส โภชนสฺสา”ติ. เอวเมว โข มหาราช ตถาคโต อนฺโต\-สาสน\-สมุคฺเค ปรมปวรํ สนฺตํ สิวํ ปณีตํ อมตํ ปรมมธุรํ กายคตาสติ\-โภชนํ ฐเปสิ “เย เกจิ กิเลส\-ฉาตชฺฌตฺตา ตณฺหา\-ปเรต\-มานสา สเจตนา พุธา, เต อิมํ โภชนํ ภุญฺชิตฺวา กาม\-รูปารูปภเวสุ สพฺพํ ตณฺหมปเนสฺสนฺตี”ติ. ยทิ โกจิ ตํ โภชนํ อภุญฺชิตฺวา ตณฺหาสิโตว ปฏินิวตฺติตฺวา หีนายาวตฺตติ, ตญฺเญว ชโน ครหิสฺสติ “อยํ ชินสาสเน ปพฺพชิตฺวา ตตฺถ ปติฏฺฐํ อลภิตฺวา หีนายาวตฺโต, กิํ อิมํ อปฺปฏิปชฺชนฺตํ ชินสาสนํ สยํ โพเธสฺสสิ, โก โทโส ชินสาสนสฺสา”ติ.}

\prose{ยทิ มหาราช ตถาคโต คิหิํ เยว เอกสฺมิํ ผเล วินีตํ ปพฺพาเชยฺย, น นามายํ ปพฺพชฺชา กิเลส\-ปฺปหานาย วิสุทฺธิยา วา,}

\prose{\csromanfolio{242}นตฺถิ ปพฺพชฺชาย กรณียํ. ยถา มหาราช ปุริโส อเนกสเตน กมฺเมน ตฬากํ ขณาเปตฺวา ปริสาย เอวม\-นุสฺสาเวยฺย “มา เม โภนฺโต เกจิ สํกิลิฏฺฐา อิมํ ตฬากํ โอตรถ, ปวาหิต\-รโชชลฺลา ปริสุทฺธา วิมลมฏฺฐา อิมํ ตฬากํ โอตรถา”ติ. อปิ นุ โข มหาราช เตสํ ปวาหิต\-รโชชลฺลานํ ปริสุทฺธานํ วิมล\-มฏฺฐานํ เตน ตฬาเกน กรณียํ ภเวยฺยาติ. น หิ ภนฺเต, ยสฺสตฺถาย เต ตํ ตฬากํ อุปคจฺเฉยฺยุํ, ตํ อญฺญเตฺรว เตสํ กตํ กรณียํ, กิํ เตสํ เตน ตฬาเกนาติ. เอวเมว โข มหาราช ยทิ ตถาคโต คิหิํ เยว เอกสฺมิํ ผเล วินีตํ ปพฺพาเชยฺย, ตตฺเถว เตสํ กตํ กรณียํ, กิํ เตสํ ปพฺพชฺชาย.}

\prose{ยถา วา ปน มหาราช สภาวอิสิภตฺติโก สุตมนฺตปท\-ธโร อตกฺกิโก โรคุ\-ปฺปตฺติ\-กุสโล อโมฆธุวสิทฺธกมฺโม ภิสกฺโก สลฺลกตฺโต สพฺพโรคู\-ปสม\-เภสชฺชํ สนฺนิปาเตตฺวา ปริสาย เอวม\-นุสฺสาเวยฺย “มา โข โภนฺโต เกจิ สพฺยาธิกา มม สนฺติเก อุปคจฺฉถ, อพฺยาธิกา อโรคา มม สนฺติเก อุปคจฺฉถา”ติ. อปิ นุ โข มหาราช เตสํ อพฺยาธิกานํ อโรคานํ ปริปุณฺณานํ อุทคฺคานํ เตน ภิสกฺเกน กรณียํ ภเวยฺยาติ. น หิ ภนฺเต, ยสฺสตฺถาย เต ตํ ภิสกฺกํ สลฺลกตฺตํ อุปคจฺเฉยฺยุํ, ตํ อญฺญเตฺรว เตสํ กตํ กรณียํ, กิํ เตสํ เตน ภิสกฺเกนาติ. เอวเมว โข มหาราช ยทิ ตถาคโต คิหิํ เยว เอกสฺมิํ ผเล วินีตํ ปพฺพาเชยฺย, ตตฺเถว เตสํ กตํ กรณียํ, กิํ เตสํ ปพฺพชฺชาย.}

\prose{ยถา วา ปน มหาราช โกจิ ปุริโส อเนก\-ถาลิปากสตํ โภชนํ ปฏิยาทาเปตฺวา ปริสาย เอวม\-นุสฺสาเวยฺย “มา เม โภนฺโต เกจิ ฉาตา อิมํ ปริเวสนํ อุปคจฺฉถ, สุภุตฺตา ติตฺตา สุหิตา ธาตา ปีณิตา ปริปุณฺณา อิมํ ปริเวสนํ อุปคจฺฉถา”ติ. อปิ นุ โข มหาราช เตสํ ภุตฺตาวีนํ ติตฺตานํ สุหิตานํ ธาตานํ ปีณิตานํ ปริปุณฺณานํ เตน โภชเนน กรณียํ ภเวยฺยาติ. น หิ ภนฺเต, ยสฺสตฺถาย เต ตํ ปริเวสนํ อุปคจฺเฉยฺยุํ, ตํ อญฺญเตฺรว เตสํ กตํ กรณียํ, กิํ เตสํ ตาย ปริเวสนายาติ. เอวเมว โข มหาราช ยทิ ตถาคโต คิหิํ เยว เอกสฺมิํ ผเล วินีตํ ปพฺพาเชยฺย, ตตฺเถว เตสํ กตํ กรณียํ, กิํ เตสํ ปพฺพชฺชาย.}

\prose{\csromanfolio{243}อปิ จ มหาราช เย หีนายาวตฺตนฺติ, เต ชินสาสนสฺส ปญฺจ อตุลิเย คุเณ ทสฺเสนฺติ. กตเม ปญฺจ, ภูมิมหนฺตภาวํ ทสฺเสนฺติ, ปริสุทฺธ\-วิมล\-ภาวํ ทสฺเสนฺติ, ปาเปหิ อสํวาสิยภาวํ ทสฺเสนฺติ, ทุปฺปฏิเวธ\-ภาวํ ทสฺเสนฺติ, พหุสํวร\-รกฺขิย\-ภาวํ ทสฺเสนฺติ.}

\prose{กถํ ภูมิมหนฺตภาวํ ทสฺเสนฺติ, ยถา มหาราช ปุริโส อธโน หีนชจฺโจ นิพฺพิเสโส พุทฺธิปริหีโน มหารชฺชํ ปฏิลภิตฺวา น จิรสฺเสว ปริปตติ ปริธํสติ ปริหายติ ยสโต, น สกฺโกติ อิสฺสริยํ สนฺธาเรตุํ. กิํ การณํ, มหนฺตตฺตา อิสฺสริยสฺส. เอวเมว โข มหาราช เย เกจิ นิพฺพิเสสา อกตปุญฺญา พุทฺธิปริหีนา ชินสาสเน ปพฺพชนฺติ, เต ตํ ปพฺพชฺชํ ปวรุตฺตมํ สนฺธาเรตุํ อวิสหนฺตา น จิรสฺเสว ชินสาสนา ปริปติตฺวา ปริธํสิตฺวา ปริหายิตฺวา หีนายาวตฺตนฺติ, น สกฺโกนฺติ ชินสาสนํ สนฺธาเรตุํ. กิํ การณํ, มหนฺตตฺตา ชินสาสน\-ภูมิยา. เอวํ ภูมิมหนฺตภาวํ ทสฺเสนฺติ.}

\prose{กถํ ปริสุทฺธ\-วิมล\-ภาวํ ทสฺเสนฺติ, ยถา มหาราช วาริ โปกฺขรปตฺเต วิกิรติ วิธมติ วิธํเสติ, น ฐานมุ\-ปคจฺฉติ นูปลิมฺปติ. กิํ การณํ, ปริสุทฺธ\-วิมล\-ตฺตา ปทุมสฺส. เอวเมว โข มหาราช เย เกจิ สฐา กูฏา วงฺกา กุฏิลา วิสมทิฏฺฐิโน ชินสาสเน ปพฺพชนฺติ, เต ปริสุทฺธ\-วิมล\-นิกฺกณฺฏก\-ปณฺฑร\-วร\-ปฺป\-วร\-สาสนโต น จิรสฺเสว วิกิริตฺวา วิธมิตฺวา วิธํเสตฺวา อสณฺฐหิตฺวา อนุปลิมฺปิตฺวา หีนายาวตฺตนฺติ. กิํ การณํ, ปริสุทฺธ\-วิมล\-ตฺตา ชินสาสนสฺส. เอวํ ปริสุทฺธ\-วิมล\-ภาวํ ทสฺเสนฺติ.}

\prose{กถํ ปาเปหิ อสํวาสิยภาวํ ทสฺเสนฺติ, ยถา มหาราช มหาสมุทฺโท น มเตน กุณเปน สํวสติ, ยํ โหติ มหาสมุทฺเท มตํ กุณปํ, ตํ ขิปฺปเมว ตีรํ อุปเนติ ถลํ วา อุสฺสาเรติ. กิํ การณํ, มหาภูตานํ ภวนตฺตา มหาสมุทฺทสฺส. เอวเมว โข มหาราช เย เกจิ ปาปกา อสํวุตา อหิริกา อกิริยา โอสนฺนวีริยา กุสีตา กิลิฏฺฐา ทุชฺชนา มนุสฺสา ชินสาสเน ปพฺพชนฺติ, เต น จิรสฺเสว ชินสาสนโต อรหนฺต\-วิมล\-ขีณาสว\-มหาภูต\-ภวนโต นิกฺขมิตฺวา อสํวสิตฺวา หีนายาวตฺตนฺติ. กิํ การณํ, ปาเปหิ อสํวาสิยตฺตา ชินสาสนสฺส. เอวํปาเปหิ อสํวาสิยภาวํ ทสฺเสนฺติ.}

\prose{\csromanfolio{244}กถํ ทุปฺปฏิเวธ\-ภาวํ ทสฺเสนฺติ, ยถา มหาราช เย เกจิ อเฉกา อสิกฺขิตา อสิปฺปิโน มติวิปฺปหีนา อิสฺสาสา\csromansharedfootnote{29b8402928a356dc}{อิสฺสตฺถา (สี, สฺยา, อิ)} วาลคฺคเวธํ อวิสหนฺตา วิคฬนฺติ ปกฺกมนฺติ. กิํ การณํ, สณฺหสุขุมทุปฺปฏิเวธตฺตา วาลคฺคสฺส. เอวเมว โข มหาราช เย เกจิ ทุปฺปญฺญา ชฬา เอฬมูคา มูฬฺหา ทนฺธคติกา ชนา ชินสาสเน ปพฺพชนฺติ, เต ตํ ปรม\-สณฺหสุขุม\-จตุสจฺจ\-ปฺปฏิเวธํ ปฏิวิชฺฌิตุํ อวิสหนฺตา ชินสาสนา วิคฬิตฺวา ปกฺกมิตฺวา น จิรสฺเสว หีนายาวตฺตนฺติ. กิํ การณํ, ปรมสณฺหสุขุมทุปฺปฏิเวธตาย สจฺจานํ. เอวํ ทุปฺปฏิเวธ\-ภาวํ ทสฺเสนฺติ.}

\prose{กถํ พหุสํวร\-รกฺขิย\-ภาวํ ทสฺเสนฺติ, ยถา มหาราช โกจิเทว ปุริโส มหติ\-มหา\-ยุทฺธภูมิมุ\-ปคโต ปรเสนาย ทิสาวิทิสาหิ สมนฺตา ปริวาริโต สตฺติหตฺถํ ชนมุเปนฺตํ ทิสฺวา ภีโต โอสกฺกติ ปฏินิวตฺตติ ปลายติ. กิํ การณํ, พหุวิธ\-ยุทฺธ\-มุข\-รกฺขณ\-ภยา. เอวเมว โข มหาราช เย เกจิ ปาปกา อสํวุตา อหิริกา อกิริยา อกฺขนฺตี จปลา จลิตา อิตฺตรา พาลชนา ชินสาสเน ปพฺพชนฺติ, เต พหุวิธํ สิกฺขาปทํ ปริรกฺขิตุํ อวิสหนฺตา โอสกฺกิตฺวา ปฏินิวตฺติตฺวา ปลายิตฺวา น จิรสฺเสว หีนายาวตฺตนฺติ. กิํ การณํ, พหุวิธ\-สํวร\-รกฺขิย\-ภาว\-ตฺตา ชินสาสนสฺส. เอวํ พหุวิธสํวรรกฺขิภาวํ ทสฺเสนฺติ.}

\prose{ถลชุตฺตเมปิ มหาราช วสฺสิกาคุมฺเพ กิมิวิทฺธานิ ปุปฺผานิ โหนฺติ, ตานิ องฺกุรานิ สงฺกุฏิตานิ อนฺตรา เยว ปริปตนฺติ, น จ เตสุ ปริปติเตสุ วสฺสิกาคุมฺโพ หีฬิโต นาม โหติ. ยานิ ตตฺถ ฐิตานิ ปุปฺผานิ, ตานิ สมฺมา คนฺเธน ทิสาวิทิสํ อภิพฺยาเปนฺติ. เอวเมว โข มหาราช เย เต ชินสาสเน ปพฺพชิตฺวา หีนายาวตฺตนฺติ, เต ชินสาสเน กิมิวิทฺธานิ วสฺสิ\-กาปุปฺผานิ วิย วณฺณคนฺธ\-รหิตา นิพฺพ\-ณฺณา\-การ\-สีลา อภพฺพา เวปุลฺลาย, น จ เตสํ หีนายาวตฺต\-เนน ชินสาสนํ หีฬิตํ นาม โหติ. เย ตตฺถ ฐิตา ภิกฺขู, เต สเทวกํ โลกํ สีลวร\-คนฺเธน อภิพฺยาเปนฺติ.}

\prose{สาลีนมฺปิ มหาราช นิราตงฺกานํ โลหิตกานํ อนฺตเร กรุมฺภกํ นาม สาลิชาติ อุปฺปชฺชิตฺวา อนฺตรา เยว วินสฺสติ, น จ ตสฺสา วินฏฺฐตฺตา โลหิตกสาลี หีฬิตา นาม โหนฺติ. เย ตตฺถ ฐิตา สาลี, เต ราชูปโภคา โหนฺติ. เอวเมว โข มหาราช เย เต ชินสาสเน \csromanfolio{245}ปพฺพชิตฺวา หีนายาวตฺตนฺติ, เต โลหิตก\-สาลีนม\-นฺตเร กรุมฺภกา วิย ชินสาสเน น วฑฺฒิตฺวา เวปุลฺลตํ น ปาปุณิตฺวา\csromansharedfootnote{799413cdbd2fe3ed}{วฑฺฒิตฺวา เวปุลฺลตํ อปาปุณิตฺวา (สี, ก)} อนฺตรา เยว หีนายา วตฺตนฺติ, น จ เตสํ หีนายาวตฺต\-เนน ชินสาสนํ หีฬิตํ นาม โหติ. เย ตตฺถ ฐิตา ภิกฺขู, เต อรหตฺตสฺส อนุจฺฉวิกา โหนฺติ.}

\proseitem{5}{กามททสฺสาปิ มหาราช มณิรตนสฺส เอกเทสํ\csromansharedfootnote{8da482857bbd4bbc}{เอกเทเส (ก)} กกฺกสํ อุปฺปชฺชติ, น จ ตตฺถ กกฺกสุ\-ปฺปนฺนตฺตา มณิรตนํ หีฬิตํ นาม โหติ. ยํ ตตฺถ ปริสุทฺธํ มณิรตนสฺส, ตํ ชนสฺส หาสกรํ โหติ. เอวเมว โข มหาราช เย เต ชินสาสเน ปพฺพชิตฺวา หีนายาวตฺตนฺติ, กกฺกสา เต ชินสาสเน ปปฏิกา, น จ เตสํ หีนายาวตฺต\-เนน ชินสาสนํ หีฬิตํ นาม โหติ. เย ตตฺถ ฐิตา ภิกฺขู, เต เทวมนุสฺสานํ หาสชนกา โหนฺติ.}

\prose{ชาติสมฺปนฺนสฺสปิ มหาราช โลหิต\-จนฺทนสฺส เอกเทสํ ปูติกํ โหติ อปฺปคนฺธํ. น เตน โลหิตจนฺทนํ หีฬิตํ นาม โหติ. ยํ ตตฺถ อปูติกํ สุคนฺธํ, ตํ สมนฺตา วิธูเปติ อภิพฺยาเปติ. เอวเมว โข มหาราช เย เต ชินสาสเน ปพฺพชิตฺวา หีนายาวตฺตนฺติ, เต โลหิต\-จนฺทนสาร\-นฺตเร ปูติกเทสมิว ฉฑฺฑนียา ชินสาสเน, น จ เตสํ หีนายาวตฺต\-เนน ชินสาสนํ หีฬิตํ นาม โหติ. เย ตตฺถ ฐิตา ภิกฺขู, เต สเทวกํ โลกํ สีลวร\-จนฺทนคนฺเธน อนุลิมฺปยนฺตีติ.}

\prosewithcloserruled{สาธุ ภนฺเต นาคเสน, เตน เตน อนุจฺฉวิเกน เตน เตน สทิเสน การเณน นิรวชฺชม\-นุ\-ปาปิตํ ชินสาสนํ เสฏฺฐภาเวน ปริทีปิตํ, หีนายาวตฺต\-มานาปิ เต ชินสาสนสฺส เสฏฺฐภาวํ เยว ปริทีเปนฺตีติ.}{หีนายา\-วตฺตน\-ปโญฺห ปญฺจโม.}
\csromanmatikaanchor{mtk.169}
\csromantocmarkref{subsubsection}{6. อรหนฺตเวทนาเวทิยนปญฺห}{mtk.169}
\bookmark[dest={mtk.169},level=3]{6. อรหนฺตเวทนาเวทิยนปญฺห}
\csromanheader{3}{6. \textbf{อรหนฺต}\-\textbf{เวทนา}\-\textbf{เวทิย}\-\textbf{น}\-\textbf{ปญฺห}}
\proseitem{5}{ภนฺเต นาคเสน ตุเมฺห ภณถ “อรหา เอก เวทนํ เวทยติ กายิกํ น เจตสิกนฺ”ติ. กินฺนุ โข ภนฺเต นาคเสน อรหโต จิตฺตํ ยํ กายํ นิสฺสาย ปวตฺตติ, ตตฺถ อรหา อนิสฺสโร อสฺสามี \csromanfolio{246}อวสวตฺตีติ. อาม มหาราชาติ. น โข ภนฺเต นาคเสน ยุตฺตเมตํ, ยํ โส สกจิตฺตสฺส ปวตฺตมาเน กาเย อนิสฺสโร โหติ อสฺสามี อวสวตฺตี, สกุโณปิ ตาว ภนฺเต ยสฺมิํ กุลาวเก ปฏิวสติ, ตตฺถ โส อิสฺสโร โหติ สามี วสวตฺตีติ.}

\proseitem{6}{ทส ยิเม มหาราช กายานุคตา ธมฺมา ภเว ภเว กายํ อนุธาวนฺติ อนุปริ\-วตฺตนฺติ. กตเม ทส, สีตํ อุณฺหํ ชิฆจฺฉา ปิปาสา อุจฺจาโร ปสฺสาโว มิทฺธํ\csromansharedfootnote{8feded072d1de0a5}{ถินมิทฺธํ (สพฺพตฺถ)} ชรา พฺยาธิ มรณํ. อิเม โข มหาราช ทส กายานุคตา ธมฺมา ภเว ภเว กายํ อนุธาวนฺติ อนุปริ\-วตฺตนฺติ, ตตฺถ อรหา อนิสฺสโร อสฺสามี อวสวตฺตีติ.}

\prose{ภนฺเต นาคเสน เกน การเณน อรหโต กาเย อาณา นปฺปวตฺตติ อิสฺสริยํ วา, ตตฺถ เม การณํ พฺรูหีติ. ยถา มหาราช เย เกจิ ปถวินิสฺสิตา สตฺตา, สพฺเพ เต ปถวิํ นิสฺสาย จรนฺติ วิหรนฺติ วุตฺติํ กปฺเปนฺติ, อปิ นุ โข มหาราช เตสํ ปถวิยา อาณา ปวตฺตติ อิสฺสริยํ วาติ. น หิ ภนฺเตติ. เอวเมว โข มหาราช อรหโต จิตฺตํ กายํ นิสฺสาย ปวตฺตติ, น จ อรหโต กาเย อาณา ปวตฺตติ อิสฺสริยํ วาติ.}

\prose{ภนฺเต นาคเสน เกน การเณน ปุถุชฺชโน กายิกมฺปิ เจตสิกมฺปิ เวทนํ เวทยตีติ. อภาวิตตฺตา มหาราช จิตฺตสฺส ปุถุชฺชโน กายิกมฺปิ เจตสิกมฺปิ เวทนํ เวทยติ. ยถา มหาราช โคโณ ฉาโต ปริตสิโต อพล\-ทุพฺพล\-ปริตฺตก\-ติเณสุ วา ลตาย วา อุปนิพทฺโธ อสฺส, ยทา โส โคโณ ปริกุปิโต โหติ, ตทา สห อุปนิพนฺธเนน ปกฺกมติ. เอวเมว โข มหาราช อภาวิตจิตฺตสฺส เวทนา อุปฺปชฺชิตฺวา จิตฺตํ ปริโกเปติ, จิตฺตํ ปริกุปิตํ กายํ อาภุชติ นิพฺภุชติ สมฺปริวตฺตกํ กโรติ. อถ โข โส อภาวิตจิตฺโต ตสติ รวติ เภรว\-ราวม\-ภิรวติ, อิทเมตฺถ มหาราช การณํ, เยน การเณน ปุถุชฺชโน กายิกมฺปิ เจตสิกมฺปิ เวทนํ เวทยตีติ.}

\prose{กิํ ปน ตํ การณํ, เยน การเณน อรหา เอกํ เวทนํ เวทยติ กายิกํ น เจตสิกนฺติ. อรหโต มหาราช จิตฺตํ ภาวิตํ โหติ สุภาวิตํ ทนฺตํ สุทนฺตํ อสฺสวํ วจนกรํ, โส ทุกฺขาย เวทนาย ผุฏฺโฐ สมาโน “อนิจฺจนฺ”ติ ทฬฺหํ คณฺหาติ, สมาธิถมฺเภ จิตฺตํ อุปนิพนฺธติ, \csromanfolio{247}ตสฺส ตํ จิตฺตํ สมาธิถมฺเภ อุปนิพนฺธนํ น เวธติ น จลติ, ฐิตํ โหติ อวิกฺขิตฺตํ, ตสฺส เวทนา\-วิการ\-วิปฺผาเรน กาเย อาภุชติ นิพฺภุชติ สมฺปริวตฺตติ, อิทเมตฺถ มหาราช การณํ, เยน การเณน อรหา เอกํ เวทนํ เวทยติ กายิกํ น เจตสิกนฺติ.}

\prosewithcloserruled{ภนฺเต นาคเสน ตํ นาม โลเก อจฺฉริยํ ยํ กาเย จลมาเน จิตฺตํ น จลติ, ตตฺถ เม การณํ พฺรูหีติ. ยถา มหาราช มหติ\-มหารุกฺเข ขนฺธ\-สาขาปลาสสมฺปนฺเน อนิล\-พล\-สมาหเต สาขา จลติ, อปิ นุ ตสฺส ขนฺโธปิ จลตีติ. น หิ ภนฺเตติ. เอวเมว โข มหาราช อรหา ทุกฺขาย เวทนาย ผุฏฺโฐ สมาโน “อนิจฺจนฺ”ติ ทฬฺหํ คณฺหาติ, สมาธิถมฺเภ จิตฺตํ อุปนิพนฺธติ, ตสฺส ตํ จิตฺตํ สมาธิถมฺเภ อุปนิพนฺธนํ น เวธติ น จลติ, ฐิตํ โหติ อวิกฺขิตฺตํ, ตสฺส เวทนา\-วิการ\-วิปฺผาเรน กาโย อาภุชติ นิพฺภุชติ สมฺปริวตฺตติ, จิตฺตํ ปน ตสฺส น เวธติ น จลติ ขนฺโธ วิย มหารุกฺขสฺสาติ. อจฺฉริยํ ภนฺเต นาคเสน, อพฺภุตํ ภนฺเต นาคเสน, น เม เอวรูโป สพฺพกาลิโก ธมฺมปทีโป ทิฏฺฐปุพฺโพติ.}{อรหนฺต\-เวทนา\-เวทิย\-น\-ปโญฺห ฉฏฺโฐ.}
\csromanmatikaanchor{mtk.170}
\csromantocmarkref{subsubsection}{7. อภิสมยนฺตรายกรปญฺห}{mtk.170}
\bookmark[dest={mtk.170},level=3]{7. อภิสมยนฺตรายกรปญฺห}
\csromanheader{3}{7. \textbf{อภิสมย}\-\textbf{นฺตรายกร}\-\textbf{ปญฺห}}
\proseitem{7}{ภนฺเต นาคเสน อิธ โย โกจิ คิหี ปาราชิกํ อชฺฌาปนฺโน ภเวยฺย, โส อปเรน สมเยน ปพฺพาเชยฺย, อตฺตนาปิ โส น ชาเนยฺย “คิหิปาราชิกํ อชฺฌาปนฺโนสฺมี”ติ, นปิ ตสฺส อญฺโญ โกจิ อาจิกฺเขยฺย “คิหิปาราชิกํ อชฺฌาปนฺโนสี”ติ. โส จ ตถตฺตาย ปฏิปชฺเชยฺย, อปิ นุ ตสฺส ธมฺมา\-ภิสมโย ภเวยฺยาติ. น หิ มหาราชาติ. เกน ภนฺเต การเณนาติ. โย ตสฺส เหตุ ธมฺมา\-ภิสมยาย, โส ตสฺส สมุจฺฉินฺโน, ตสฺมา ธมฺมา\-ภิสมโย น ภวตีติ.}

\prose{ภนฺเต นาคเสน ตุเมฺห ภณถ “ชานนฺตสฺส กุกฺกุจฺจํ โหติ, กุกฺกุจฺเจ สติ อาวรณํ โหติ, อาวเฏ จิตฺเต ธมฺมา\-ภิสมโย น โหตี”ติ. อิมสฺส ปน อชานนฺตสฺส อกุกฺกุจฺจ\-ชาตสฺส สนฺตจิตฺตสฺส วิหรโต เกน การเณน ธมฺมา\-ภิสมโย น โหติ, วิสเมน วิสเมเนโส ปโญฺห คจฺฉติ, จินฺเตตฺวา วิสชฺเชถาติ.}

\prose{\csromanfolio{248}รุหติ มหาราช สุกฏฺเฐ สุกลเล มณฺฑเขตฺเต สารทํ สุขสยิตํ พีชนฺติ. อาม ภนฺเตติ. อปิ นุ มหาราช ตญฺเญว พีชํ ฆนเสล\-สิลา\-ตเล รุเหยฺยาติ. น หิ ภนฺเตติ. กิสฺส ปน มหาราช ตญฺเญว พีชํ กลเล รุหติ, กิสฺส ฆนเสเล น รุหตีติ. นตฺถิ ภนฺเต ตสฺส พีชสฺส รูหนาย ฆนเสเล เหตุ, อเหตุนา พีชํ น รุหตีติ. เอวเมว โข มหาราช เยน เหตุนา ตสฺส ธมฺมา\-ภิสมโย ภเวยฺย, โส ตสฺส เหตุ สมุจฺฉินฺโน, อเหตุนา ธมฺมา\-ภิสมโย น โหติ.}

\prose{ยถา วา ปน มหาราช ทณฺฑ\-เลฑฺฑุ\-ลคุฬ\-มุคฺครา ปถวิยา ฐานมุปคจฺฉนฺติ, อปิ นุ มหาราช เต เยว ทณฺฑ\-เลฑฺฑุ\-ลคุฬ\-มุคฺครา คคเน ฐานมุปคจฺฉนฺตีติ. น หิ ภนฺเตติ. กิํ ปเนตฺถ มหาราช การณํ, เยน การเณน เต เยว ทณฺฑ\-เลฑฺฑุ\-ลคุฬ\-มุคฺครา ปถวิยา ฐานมุปคจฺฉนฺติ, เกน การเณน คคเน น ติฏฺฐนฺตีติ. นตฺถิ ภนฺเต เตสํ ทณฺฑ\-เลฑฺฑุ\-ลคุฬ\-มุคฺครานํ ปติฏฺฐานาย อากาเส เหตุ, อเหตุนา น ติฏฺฐนฺตีติ. เอวเมว โข มหาราช ตสฺส เตน โทเสน อภิสมยเหตุ สมุจฺฉินฺโน, เหตุสมุคฺฆาเต อเหตุนา อภิสมโย น โหตีติ.}

\prose{ยถา วา ปน มหาราช ถเล อคฺคิ ชลติ, อปิ นุ โข มหาราช โส เยว\csromansharedfootnote{cb63cef12ad01f03}{โสเยว สนฺธิ \csromandash{} ม.พ.ป.} อคฺคิ อุทเก ชลตีติ. น หิ ภนฺเตติ. กิํ ปเนตฺถ มหาราช การณํ, เยน การเณน โส เยว\csromansharedfootnote{cb63cef12ad01f03}{โสเยว สนฺธิ \csromandash{} ม.พ.ป.} อคฺคิ ถเล ชลติ, เกน การเณน อุทเก น ชลตีติ. นตฺถิ ภนฺเต อคฺคิสฺส ชลนาย อุทเก เหตุ, อเหตุนา นชลตีติ. เอวเมว โข มหาราช ตสฺส เตน โทเสน อภิสมยเหตุ สมุจฺฉินฺโน, เหตุสมุคฺฆาเต อเหตุนา ธมฺมา\-ภิสมโย น โหตีติ.}

\prose{ภนฺเต นาคเสน ปุนเปตํ อตฺถํ จินฺเตหิ, น เม ตตฺถ จิตฺตสญฺญตฺติ ภวติ, อชานนฺตสฺส อสติ กุกฺกุจฺเจ อาวรณํ โหตีติ, การเณน มํ สญฺญาเปหีติ. อปิ นุ มหาราช วิสํ หลาหลํ อชานนฺเตน ขายิตํ ชีวิตํ หรตีติ. อาม ภนฺเตติ. เอวเมว โข มหาราช อชานนฺเตนปิ กตํ ปาปํ อภิสมย\-นฺตรายกรํ โหติ.}

\prose{อปิ นุ มหาราช อคฺคิ อชานิตฺวา อกฺกมนฺตํ ฑหตีติ. อาม ภนฺเตติ. เอวเมว โข มหาราช อชานนฺเตนปิ กตํ ปาปํ อภิสมย\-นฺตรายกรํ โหติ.}

\prose{\csromanfolio{249}อปิ นุ มหาราช อชานนฺตํ อาสีวิโส ฑํสิตฺวา ชีวิตํ หรตีติ. อาม ภนฺเตติ. เอวเมว โข มหาราช อชานนฺเตนปิ กตํ ปาปํ อภิสมย\-นฺตรายกรํ โหติ.}

\prosewithcloserruled{นนุ มหาราช กาลิงฺคราชา สมณโกลญฺโญ สตฺตรตน\-ปริกิณฺโณ หตฺถิรตนม\-ภิรุยฺห กุลทสฺสนาย คจฺฉนฺโต อชานนฺโตปิ นาสกฺขิ โพธิมณฺฑสฺส อุปริโต คนฺตุํ, อิทเมตฺถ มหาราช การณํ, เยน การเณน อชานนฺเตนปิ กตํ ปาปํ อภิสมย\-นฺตรายกรํ โหตีติ. ชินภาสิตํ ภนฺเต นาคเสน การณํ น สกฺกา ปฏิกฺโกสิตุํ, เอโสเวตสฺส อตฺโถ ตถา สมฺปฏิจฺฉามีติ.}{อภิสมย\-นฺตรายกร\-ปโญฺห สตฺตโม.}
\csromanmatikaanchor{mtk.171}
\csromantocmarkref{subsubsection}{8. ทุสฺสีลปญฺห}{mtk.171}
\bookmark[dest={mtk.171},level=3]{8. ทุสฺสีลปญฺห}
\csromanheader{3}{8. \textbf{ทุสฺสีลปญฺห}}
\proseitem{8}{ภนฺเต นาคเสน คิหิทุสฺสีลสฺส จ สมณ\-ทุสฺสีลสฺส จ โก วิเสโส, กิํ นานากรณํ, อุโภเปเต สมสมคติกา, อุภินฺนมฺปิ สมสโม วิปาโก โหติ, อุทาหุ กิญฺจิ นานาการณํ อตฺถีติ.}

\prose{ทส ยิเม มหาราช คุณา สมณ\-ทุสฺสีลสฺส คิหิทุสฺสีลโต วิเสเสน อติเรกา, ทสหิ จ การเณหิ อุตฺตริํ ทกฺขิณํ วิโสเธติ.}

\prose{กตเม ทส คุณา สมณ\-ทุสฺสีลสฺส คิหิทุสฺสีลโต วิเสเสน อติเรกา, อิธ มหาราช สมณทุสฺสีโล พุทฺเธ สคารโว โหติ, ธมฺเม สคารโว โหติ, สํเฆ สคารโว โหติ, สพฺรหฺมจารีสุ สคารโว โหติ, อุทฺเทส\-ปริปุจฺฉาย วายมติ, สวนพหุโล โหติ, ภินฺนสีโลปิ มหาราช ทุสฺสีโล ปริสคโต อากปฺปํ อุปฏฺฐเปติ, ครหภยา กายิกํ วาจสิกํ รกฺขติ, ปธานาภิ\-มุข\-ญฺจสฺส โหติ จิตฺตํ, ภิกฺขุ\-สามญฺญํ อุปคโต โหติ. กโรนฺโตปิ มหาราช สมณทุสฺสีโล ปาปํ ปฏิจฺฉนฺนํ อาจรติ. ยถา มหาราช อิตฺถี สปติกา นิลียิตฺวา รหสฺเสเนว ปาปมาจรติ. เอวเมว โข มหาราช \csromanfolio{250}กโรนฺโตปิ สมณทุสฺสีโล ปาปํ ปฏิจฺฉนฺนํ อาจรติ. อิเม โข มหาราช ทส คุณา สมณ\-ทุสฺสีลสฺส คิหิทุสฺสีลโต วิเสเสน อติเรกา.}

\prose{กตเมหิ ทสหิ การเณหิ อุตฺตริํ ทกฺขิณํ วิโสเธติ, อนวชฺช\-กวจ\-ธารณตายปิ ทกฺขิณํ วิโสเธติ, อิสิสามญฺญภณฺฑุลิงฺคธารณโตปิ ทกฺขิณํ วิโสเธติ, สํฆสมยมนุปฺปวิฏฺฐตายปิ ทกฺขิณํ วิโสเธติ, พุทฺธธมฺมสํฆสรณคตตายปิ ทกฺขิณํ วิโสเธติ, ปธานา\-สย\-นิเกต\-วาสิตายปิ ทกฺขิณํ วิโสเธติ, ชินสาสนธร\csromansharedfootnote{f3be36e3579b8540}{ชินสาสนธน (สี, อิ)} ปริเยสนโตปิ ทกฺขิณํ วิโสเธติ, ปวรธมฺมเทสนโตปิ ทกฺขิณํ วิโสเธติ, ธมฺมทีป\-คติ\-ปรายณ\-ตายปิ ทกฺขิณํ วิโสเธติ, “อคฺโค พุทฺโธ”ติ เอกนฺตอุชุทิฏฺฐิตายปิ ทกฺขิณํ วิโสเธติ, อุโปสถสมาทานโตปิ ทกฺขิณํ วิโสเธติ. อิเมหิ โข มหาราช ทสหิ การเณหิ อุตฺตริํ ทกฺขิณํ วิโสเธติ.}

\prose{สุวิปนฺโนปิ หิ มหาราช สมณทุสฺสีโล ทายกานํ ทกฺขิณํ วิโสเธติ. ยถา มหาราช อุทกํ สุพหลมฺปิ กลล\-กทฺทม\-รโชชลฺลํ อปเนติ. เอวเมว โข มหาราช สุวิปนฺโนปิ สมณทุสฺสีโล ทายกานํ ทกฺขิณํ วิโสเธติ.}

\prose{ยถา วา ปน มหาราช อุโณฺหทกํ สุกุธิตมฺปิ\csromansharedfootnote{7a5f41ddcf774025}{สุกฐิตมฺปิ (สี, อิ), สุขุฐิตมฺปิ (สฺยา)} ปชฺชลนฺตํ มหนฺตํ อคฺคิกฺขนฺธํ นิพฺพาเปติ. เอวเมว โข มหาราช สุวิปนฺโนปิ สมณทุสฺสีโล ทายกานํ ทกฺขิณํ วิโสเธติ.}

\prose{ยถา วา ปน มหาราช โภชนํ วิรสมฺปิ ขุทาทุพฺพลฺยํ อปเนติ. เอวเมว โข มหาราช สุวิปนฺโนปิ สมณทุสฺสีโล ทายกานํ ทกฺขิณํ วิโสเธติ.}

\prose{ภาสิตมฺเปตํ มหาราช ตถาคเตน เทวาติเทเวน มชฺฌิมนิกาย\-วร\-ลญฺฉเก ทกฺขิณ\-วิภงฺเค เวยฺยากรเณ\csromandash{}}

\setlength{\csromangathapagewidth}{0pt}
\setlength{\csromangathawakwidth}{0pt}
\csromangathameasuregroup{}{“โย สีลวา ทุสฺสีเลสุ ททาติ ทานํ, \\ ธมฺเมน ลทฺธํ สุปสนฺนจิตฺโต. \\ อภิสทฺทหํ กมฺมผลํ อุฬารํ, \\ สา ทกฺขิณา ทายกโต วิสุชฺฌตี”ติ\textsuperscript{1}.}
\csromangathameasurewak{“โย สีลวา ทุสฺสีเลสุ ททาติ ทานํ, \\ ธมฺเมน ลทฺธํ สุปสนฺนจิตฺโต. \\ อภิสทฺทหํ กมฺมผลํ อุฬารํ, \\ สา ทกฺขิณา ทายกโต วิสุชฺฌตี”ติ\textsuperscript{1}.}
\setlength{\csromangathaleftcol}{0pt}
\csromangathagroup{\csromangathastanza{\csromangathawak{“โย สีลวา ทุสฺสีเลสุ ททาติ ทานํ, \\ ธมฺเมน ลทฺธํ สุปสนฺนจิตฺโต. \\ อภิสทฺทหํ กมฺมผลํ อุฬารํ, \\ สา ทกฺขิณา ทายกโต วิสุชฺฌตี”ติ\csromansharedfootnote{e2bddd2ca1bb4026}{ม 3. 300 ปิฏฺเฐ.}.}}}

\prosewithcloserruled{\csromanfolio{251}อจฺฉริยํ ภนฺเต นาคเสน, อพฺภุตํ ภนฺเต นาคเสน, ตาวตกํ มยํ ปญฺหํ อปุจฺฉิมฺห, ตํ ตฺวํ โอปมฺเมหิ การเณหิ วิภาเวนฺโต อมตมธุรํ สฺ1อวนูปคํ อกาสิ. ยถา นาม ภนฺเต สูโท วา สูทนฺเตวาสี วา ตาวตกํ มํสํ ลภิตฺวา นานาวิเธหิ สมฺภาเรหิ สมฺปาเทตฺวา ราชูปโภคํ กโรติ. เอวเมว โข ภนฺเต นาคเสน ตาวตกํ มยํ ปญฺหํ อปุจฺฉิมฺห, ตํ ตฺวํ โอปมฺเมหิ การเณหิ วิภาเวตฺวา อมตมธุรํ สวนูปคํ อกาสีติ.}{ทุสฺสีลปโญฺห อฏฺฐโม.}
\csromanmatikaanchor{mtk.172}
\csromantocmarkref{subsubsection}{9. อุทกสตฺตชีวปญฺห}{mtk.172}
\bookmark[dest={mtk.172},level=3]{9. อุทกสตฺตชีวปญฺห}
\csromanheader{3}{9. \textbf{อุทก}\-\textbf{สตฺตชีว}\-\textbf{ปญฺห}}
\proseitem{9}{ภนฺเต นาคเสน อิมํ อุทกํ อคฺคิมฺหิ ตปฺปมานํ จิจฺจิฏายติ จิฏิจิฏายติ สทฺทายติ พหุวิธํ, กินฺนุ โข ภนฺเต นาคเสน อุทกํ ชีวติ, กิํ กีฬมานํ สทฺทายติ, อุทาหุ อญฺเญน ปฏิปีฬิตํ สทฺทายตีติ. น หิ มหาราช อุทกํ ชีวติ, นตฺถิ อุทเก ชีโว วา สตฺโต วา, อปิ จ มหาราช อคฺคิสนฺตาป\-เวคสฺส มหนฺตตาย อุทกํ จิจฺจิฏายติ จิฏิจิฏายติ สทฺทายติ พหุวิธนฺติ.}

\prose{ภนฺเต นาคเสน อิเธกจฺเจ ติตฺถิยา อุทกํ ชีวตีติ สีโตทกํ ปฏิกฺขิปิตฺวา อุทกํ ตาเปตฺวา เวกติกเวกติกํ ปริภุญฺชนฺติ, เต ตุเมฺห ครหนฺติ ปริภวนฺติ “เอกินฺทฺริยํ สมณา สกฺยปุตฺติยา ชีวํ วิเหเฐนฺตี”ติ, ตํ เตสํ ครหํ ปริภวํ วิโนเทหิ อปเนหิ นิจฺฉาเรหีติ. น หิ มหาราช อุทกํ ชีวติ, นตฺถิ มหาราช อุทเก ชีโว วา สตฺโต วา, อปิ จ มหาราช อคฺคิสนฺตาป\-เวคสฺส มหนฺตตาย อุทกํ จิจฺจิฏายติ จิฏิจิฏายติ สทฺทายติ พหุวิธํ.}

\prose{ยถา มหาราช อุทกํ โสพฺภ\-สร\-สริต\-ทห\-ตฬาก\-กนฺทร\-ปทร- อุทปาน\-นินฺน\-โปกฺขรณิคตํ วาตาตป\-เวคสฺส มหนฺตตาย ปริยาทิยติ ปริกฺขยํ คจฺฉติ, อปิ นุ ตตฺถ อุทกํ จิจฺจิฏายติ จิฏิจิฏายติ สทฺทายติ พหุวิธนฺติ. น หิ ภนฺเตติ. ยทิ มหาราช อุทกํ ชีเวยฺย, ตตฺถาปิ อุทกํ สทฺทาเยยฺย, อิมินาปิ มหาราช การเณน ชานาหิ “นตฺถิ อุทเก ชีโว วา สตฺโต วา, อคฺคิสนฺตาป\-เวคสฺส มหนฺตตาย อุทกํ จิจฺจิฏายติ จิฏิจิฏายติ สทฺทายติ พหุวิธนฺ”ติ.}

\prose{\csromanfolio{252}อปรมฺปิ มหาราช อุตฺตริํ การณํ สุโณหิ “นตฺถิ อุทเก ชีโว วา สตฺโต วา, อคฺคิสนฺตาป\-เวคสฺส มหนฺตตาย อุทกํ สทฺทายตี”ติ. ยทา ปน มหาราช อุทกํ ตณฺฑุเลหิ สมฺมิสฺสิตํ ภาชนคตํ โหติ ปิหิตํ อุทฺธเน อฐปิตํ, อปิ นุ ตตฺถ อุทกํ สทฺทายตีติ. น หิ ภนฺเต อจลํ โหติ สนฺตสนฺตนฺติ. ตํ เยว\csromansharedfootnote{cf7dc7e6d0b8f44d}{ตํเยว นิคฺคหีตสนฺธิ \csromandash{} ม.พ.ป.} ปน มหาราช อุทกํ ภาชนคตํ อคฺคิํ อุชฺชาเลตฺวา อุทฺธเน ฐปิตํ โหติ, อปิ นุ ตตฺถ อุทกํ อจลํ โหติ สนฺตสนฺตนฺติ. น หิ ภนฺเต, จลติ ขุพฺภติ ลุฬติ อาวิลติ อูมิชาตํ โหติ, อุทฺธมโธ ทิสาวิทิสํ คจฺฉติ, อุตฺตรติ ปตรติ เผณมาลี โหตีติ. กิสฺส ปน ตํ มหาราช ปากติกํ อุทกํ น จลติ สนฺตสนฺตํ โหติ, กิสฺส ปน อคฺคิคตํ จลติ ขุพฺภติ ลุฬติ อาวิลติ อูมิชาตํ โหติ, อุทฺธมโธ ทิสาวิทิสํ คจฺฉติ, อุตฺตรติ ปตรติ เผณมาลี โหตีติ. ปากติกํ ภนฺเต อุทกํ น จลติ, อคฺคิกตํ ปน อุทกํ อคฺคิสนฺตาป\-เวคสฺส มหนฺตตาย จิจฺจิฏายติ จิฏิจิฏายติ สทฺทายติ พหุวิธนฺติลฺ. อิมินาปิ มหาราช การเณน ชานาหิ “นตฺถิ อุทเก ชีโว วา สตฺโต วา, อคฺคิสนฺตาป\-เวคสฺส มหนฺตตาย อุทกํ สทฺทายตี”ติ.}

\prose{อปรมฺปิ มหาราช อุตฺตริํ การณํ สุโณหิ, นตฺถิ อุทเก ชีโว วา สตฺโต วา, อคฺคิสนฺตาป\-เวคสฺส มหนฺตตาย อุทกํ สทฺทายติ. โหติ ตํ มหาราช อุทกํ ฆเร ฆเร อุทกวารกคตํ ปิหิตนฺติ. อาม ภนฺเตติ. อปิ นุ ตํ มหาราช อุทกํ จลติ ขุพฺภติ ลุฬติ อาวิลติ อูมิชาตํ โหติ, อุทฺธมโธ ทิสาวิทิสํ คจฺฉติ, อุตฺตรติ ปตรติ เผณมาลี โหตีติ. น หิ ภนฺเต, อจลํ ตํ โหติ ปากติกํ อุทกวารกคตํ อุทกนฺติ.}

\prose{สุตปุพฺพํ ปน ตยา มหาราช “มหาสมุทฺเท อุทกํ จลติ ขุพฺภติ ลุฬติ อาวิลติ อูมิชาตํ โหติ, อุทฺธมโธ ทิสาวิทิสํ คจฺฉติ, อุตฺตรติ ปตรติ เผณมาลี โหติ, อุสฺสกฺกิตฺวา โอสฺสกฺกิตฺวา\csromansharedfootnote{951c07a4e4dcdec0}{“โอสฺสกฺกิตฺวา”ติ ปทํ สี- อิโปตฺถเกสุ นตฺถิ.} เวลาย ปหรติ สทฺทายติ พหุวิธนฺ”ติ. อาม ภนฺเต สุตปุพฺพํ เอตํ มยา ทิฏฺฐปุพฺพญฺจ “มหาสมุทฺเท อุทกํ หตฺถสตมฺปิ เทฺวปิ หตฺถสตานิ คคเน อุสฺสกฺกตี”ติ. กิสฺส มหาราช อุทกวารกคตํ อุทกํ น จลติ น สทฺทายติ, กิสฺส ปน มหาสมุทฺเท อุทกํ จลติ สทฺทายตีติ. วาตเวคสฺส มหนฺตตาย ภนฺเต มหาสมุทฺเท อุทกํ จลติ สทฺทายติ, อุทกวารกคตํ อุทกํ อฆฏฺฏิตํ \csromanfolio{253}เกหิจิ น จลติ น สทฺทายตีติ. ยถา มหาราช วาตเวคสฺส มหนฺตตาย มหาสมุทฺเท อุทกํ จลติ สทฺทายติ. เอวเมว อคฺคิสนฺตาป\-เวคสฺส มหนฺตตาย อุทกํ สทฺทายตีติ.}

\prose{นนุ มหาราช เภริโปกฺขรํ สุกฺขํ สุกฺเขน โคจมฺเมน โอนนฺธนฺตีติ. อาม ภนฺเต. อปิ นุ มหาราช เภริยา ชีโว วา สตฺโต วา อตฺถีติ. น หิ ภนฺเตติ. กิสฺส ปน มหาราช เภรี สทฺทายตีติ. อิตฺถิยา วา ภนฺเต ปุริสสฺส วา ตชฺเชน วายาเมนาติ. ยถา มหาราช อิตฺถิยา วา ปุริสสฺส วา ตชฺเชน วายาเมน เภรี สทฺทายติ. เอวเมว อคฺคิสนฺตาป\-เวคสฺส มหนฺตตาย อุทกํ สทฺทายติ, อิมินาปิ มหาราช การเณน ชานาหิ “นตฺถิ อุทเก ชีโว วา สตฺโต วา, อคฺคิสนฺตาป\-เวคสฺส มหนฺตตาย อุทกํ สทฺทายตี”ติ.}

\prose{มยฺหมฺปิ ตาว มหาราช ตว ปุจฺฉิตพฺพํ อตฺถิ, เอวเมโส ปโญฺห สุวินิจฺฉิโต โหติ, กินฺนุ โข มหาราช สพฺเพหิปิ ภาชเนหิ อุทกํ ตปฺปมานํ สทฺทายติ, อุทาหุ เอกจฺเจหิ เยว ภาชเนหิ ตปฺปมานํ สทฺทายตีติ. น หิ ภนฺเต, สพฺเพหิปิ ภาชเนหิ อุทกํ ตปฺปมานํ สทฺทายติ, เอกจฺเจหิ เยว ภาชเนหิ อุทกํ ตปฺปมานํ สทฺทายตีติ. เตน หิ มหาราช ชหิโตสิ สกสมยํ, ปจฺจาคโตสิ มม วิสยํ, นตฺถิ อุทเก ชีโว วา สตฺโต วา. ยทิ มหาราช สพฺเพหิปิ ภาชเนหิ อุทกํ ตปฺปมานํ สทฺทาเยยฺย, ยุตฺตมิทํ “อุทกํ ชีวตี”ติ วตฺตุํ. น หิ มหาราช อุทกํ ทฺวยํ โหติ, ยํ สทฺทายติ, ตํ ชีวติ, ยํ น สทฺทายติ, ตํ น ชีวตีติ. ยทิ มหาราช อุทกํ ชีเวยฺย มหนฺตานํ หตฺถินาคานํ อุสฺสนฺนกายานํ ปภินฺนานํ โสณฺฑาย อุสฺสิญฺจิตฺวา มุเข ปกฺขิปิตฺวา กุจฺฉิํ ปเวสยนฺตานํ, ตมฺปิ อุทกํ เตสํ ทนฺตนฺตเร จิปฺปิยมานํ สทฺทาเยยฺย, หตฺถสติกาปิ มหานาวา ครุกา ภาริกา อเนกสตสหสฺส\-ภาร\-ปริปูรา มหาสมุทฺเท วิจรนฺติ. ตาหิปิ จิปฺปิยมานํ อุทกํ สทฺทาเยยฺย. มหติ\-มหนฺตาปิ มจฺฉา อเนกสต\-โยชนิก\-กายา ติมี ติมิงฺคลา ติมิรปิงฺคลา อพฺภนฺตเร นิมุคฺคา มหาสมุทฺเท นิวาสฏฺฐาน\-ตาย ปฏิวสนฺตา มหาอุทกธารา อาจมนฺติ ธมนฺติ จ, เตสมฺปิ ตํ ทนฺตนฺตเรปิ อุทรนฺตเรปิ จิปฺปิยมานํ อุทกํ สทฺทาเยยฺย. ยสฺมา จ โข มหาราช เอวรูเปหิ เอวรูเปหิ มหนฺเตหิ ปฏิปีฬเนหิ ปฏิปีฬิตํ อุทกํ น สทฺทายติ. ตสฺมาปิ นตฺถิ อุทเก ชีโว วา สตฺโต วาติ, เอวเมตํ มหาราช ธาเรหีติ.}

\proseitem{5}{\csromanfolio{254}สาธุ ภนฺเต นาคเสน โทสาคโต\csromansharedfootnote{1214b8fc76d52658}{เทสาคโต (สี, อิ)} ปโญฺห อนุจฺฉวิกาย วิภตฺติยา วิภตฺโต, ยถา นาม ภนฺเต นาคเสน มหคฺฆํ มณิรตนํ เฉกํ อาจริยํ กุสลํ สิกฺขิตํ มณิการํ ปาปุณิตฺวา กิตฺติํ ลเภยฺย โถมนํ ปสํสํ, มุตฺตารตนํ วา มุตฺติกํ ทุสฺสรตนํ วา ทุสฺสิกํ, โลหิตจนฺทนํ วา คนฺธิกํ ปาปุณิตฺวา กิตฺติํ ลเภยฺย โถมนํ ปสํสํ. เอวเมว โข ภนฺเต นาคเสน โทสาคโต\csromansharedfootnote{1214b8fc76d52658}{เทสาคโต (สี, อิ)} ปโญฺห อนุจฺฉวิกาย วิภตฺติยา วิภตฺโต, เอวเมตํ ตถา สมฺปฏิจฺฉามีติ.}

\vspace{\tipitakaparskip}
\csromancenter{อุทก\-สตฺตชีว\-ปโญฺห นวโม.}
\nitthitammidruled{\textbf{พุทฺธวคฺโค ปฐโม.}}
\csromancenter{อิมสฺมิํ วคฺเค นว ปญฺหา.}
\csromanmatikaanchor{mtk.173}
\csromantocmarknopage{subsection}{2. นิปฺปปญฺจวคฺค}{mtk.173}
\bookmark[dest={mtk.173},level=2]{2. นิปฺปปญฺจวคฺค}
\csromanheader{2}{2. \textbf{นิปฺปปญฺจ}\-\textbf{วคฺค}}
\csromanmatikaanchor{mtk.174}
\csromantocmarkref{subsubsection}{1. นิปฺปปญฺจปญฺห}{mtk.174}
\bookmark[dest={mtk.174},level=3]{1. นิปฺปปญฺจปญฺห}
\csromanheader{3}{1. \textbf{นิปฺปปญฺจ}\-\textbf{ปญฺห}}
\proseitem{1}{\csromanfolio{255}ภนฺเต นาคเสน ภาสิตมฺเปตํ ภควตา “นิปฺปปญฺจารามา ภิกฺขเว วิหรถ นิปฺปปญฺจ\-รติโน”ติ, กตมํ ตํ นิปฺปปญฺจนฺติ. โสตาปตฺติผลํ มหาราช นิปฺปปญฺจํ, สกทาคามิผลํ นิปฺปปญฺจํ, อนาคามิผลํ นิปฺปปญฺจํ, อรหตฺตผลํ นิปฺปปญฺจนฺติ.}

\prose{ยทิ ภนฺเต นาคเสน โสตาปตฺติผลํ นิปฺปปญฺจํ, สกทาคามิ- อนาคามิอรหตฺตผลํ นิปฺปปญฺจํ, กิสฺส ปน อิเม ภิกฺขู อุทฺทิสนฺติ ปริปุจฺฉนฺติ สุตฺตํ เคยฺยํ เวยฺยากรณํ คาถํ อุทานํ อิติวุตฺตกํ ชาตกํ อพฺภุตธมฺมํ เวทลฺลํ, นวกมฺเมน ปลิพุชฺฌนฺติ ทาเนน จ ปูชาย จ, นนุ เต ชินปฺปฏิกฺขิตฺตํ กมฺมํ กโรนฺตีติ.}

\prose{เย เต มหาราช ภิกฺขู อุทฺทิสนฺติ ปริปุจฺฉนฺติ สุตฺตํ เคยฺยํ เวยฺยากรณํ คาถํ อุทานํ อิติวุตฺตกํ ชาตกํ อพฺภุตธมฺมํ เวทลฺลํ, นวกมฺเมน ปลิพุชฺฌนฺติ ทาเนน จ ปูชาย จ, สพฺเพ เต นิปฺปปญฺจสฺส ปตฺติยา กโรนฺติ. เย เต มหาราช สภาว\-ปริสุทฺธา ปุพฺเพ วาสิตวาสนา, เต เอกจิตฺต\-กฺขเณน นิปฺปปญฺจา โหนฺติ. เย ปน เต ภิกฺขู มหารชกฺขา, เต อิเมหิ ปโยเคหิ นิปฺปปญฺจา โหนฺติ.}

\prose{ยถา มหาราช เอโก ปุริโส เขตฺเต พีชํ โรเปตฺวา อตฺตโน ยถา\-พลวีริเยน วินา ปาการวติยา ธญฺญํ อุทฺธเรยฺย, เอโก ปุริโส เขตฺเต พีชํ โรเปตฺวา วนํ ปวิสิตฺวา กฏฺฐญฺจ สาขญฺจ ฉินฺทิตฺวา วติปาการํ กตฺวา ธญฺญํ อุทฺธเรยฺย. ยา ตตฺถ ตสฺส วติปาการ\-ปริเยสนา, สา ธญฺญตฺถาย. เอวเมว โข มหาราช เย เต สภาว\-ปริสุทฺธา ปุพฺเพ วาสิตวาสนา, เต เอกจิตฺต\-กฺขเณน นิปฺปปญฺจา โหนฺติ, วินา วติปาการํ ปุริโส วิย ธญฺญุทฺธาโร. เย ปน เต ภิกฺขู มหารชกฺขา, เต อิเมหิ ปโยเคหิ นิปฺปปญฺจา โหนฺติ, วติปาการํ กตฺวา ปุริโส วิย ธญฺญุทฺธาโร.}

\prose{ยถา วา ปน มหาราช ปุริโส มหติมหนฺเต อมฺพรุกฺข\-มตฺถเก ผลปิณฺฑิ ภเวยฺย, อถ ตตฺถ โย โกจิ อิทฺธิมา อาคนฺตฺวา ตสฺส ผลํ หเรยฺย, โย ปน ตตฺถ อนิทฺธิมา, โส กฏฺฐญฺจ วลฺลิญฺจ ฉินฺทิตฺวา นิสฺเสณิํ พนฺธิตฺวา ตาย ตํ รุกฺขํ อภิรุหิตฺวา ผลํ หเรยฺย. ยา \csromanfolio{256}ตตฺถ ตสฺส นิสฺเสณิ\-ปริเยสนา, สา ผลตฺถาย. เอวเมว โข มหาราช เย เต สภาว\-ปริสุทฺธา ปุพฺเพ วาสิตวาสนา, เต เอกจิตฺต\-กฺขเณน นิปฺปปญฺจา โหนฺติ, อิทฺธิมา วิย รุกฺขผลํ หรนฺโต. เย ปน เต ภิกฺขู มหารชกฺขา, เต อิมินา ปโยเคน สจฺจานิ อภิสเมนฺติ, นิสฺเสณิยา วิย ปุริโส รุกฺขผลํ หรนฺโต.}

\prose{ยถา วา ปน มหาราช เอโก ปุริโส อตฺถกรณิโย เอกโก เยว สามิกํ อุปคนฺตฺวา อตฺถํ สาเธติ. เอโก ธนวา ธนวเสน ปริสํ วฑฺเฒตฺวา ปริสาย อตฺถํ สาเธติ. ยา ตตฺถ ตสฺส ปริส\-ปริเยสนา, สา อตฺถตฺถาย. เอวเมว โข มหาราช เย เต สภาว\-ปริสุทฺธา ปุพฺเพ วาสิตวาสนา, เต เอกจิตฺต\-กฺขเณน ฉสุ อภิญฺญาสุ วสิภาวํ ปาปุณนฺติ, ปุริโส วิย เอกโก อตฺถสิทฺธิํ กโรนฺโต. เย ปน เต ภิกฺขู มหารชกฺขา, เต อิเมหิ ปโยเคหิ สามญฺญ\-ตฺถม\-ภิสา\-เธนฺติ, ปริสาย วิย ปุริโส อตฺถสิทฺธิํ กโรนฺโต.}

\prose{อุทฺเทโสปิ มหาราช พหุกาโร, ปริปุจฺฉาปิ พหุการา, นวกมฺมมฺปิ พหุการํ, ทานมฺปิ พหุการํ, ปูชาปิ พหุการา เตสุ เตสุ กรณีเยสุ. ยถา มหาราช ปุริโส ราชูปเสวี กตาวี อมจฺจภฏพลโทวาริก- อนีกฏฺฐปาริสชฺชชเนหิ, เต ตสฺส กรณีเย อนุปฺปตฺเต สพฺเพปิ อุปการา โหนฺติ. เอวเมว โข มหาราช อุทฺเทโสปิ พหุกาโร, ปริปุจฺฉาปิ พหุกาโร, นวกมฺมมฺปิ พหุการํ, ทานมฺปิ พหุการํ, ปูชาปิ พหุการา เตสุ เตสุ กรณีเยสุ. ยทิ มหาราช สพฺเพปิ อภิชาติ\-ปริสุทฺธา ภเวยฺยุํ, อนุสาสเนน\csromansharedfootnote{8aaebb2dd5000b4c}{อนุสาสเกน (สี, อิ)} กรณียํ น ภเวยฺย. ยสฺมา จ โข มหาราช สวเนน กรณียํ โหติ, เถโร มหาราช สาริปุตฺโต อปริมิตม\-สงฺเขยฺย\-กปฺปํ อุปาทาย อุปจิต\-กุสล\-มูโล ปญฺญาย โกฏิํ คโต, โสปิ วินา สวเนน นาสกฺขิ อาสวกฺขยํ ปาปุณิตุํ. ตสฺมา มหาราช พหุการํ สวนํ ตถา อุทฺเทโสปิ ปริปุจฺฉาปิ. ตสฺมา อุทฺเทส\-ปริปุจฺฉาปิ นิปฺปปญฺจา สงฺขตาติ. สุนิชฺฌาปิโต ภนฺเต นาคเสน ปโญฺห, เอวเมตํ ตถา สมฺปฏิจฺฉามีติ.}

\vspace{\tipitakaparskip}
\csromancenter{นิปฺปปญฺจ\-ปโญฺห ปฐโม.}
\csromanmatikaanchor{mtk.175}
\csromantocmarkref{subsubsection}{2. ขีณาสวภาวปญฺห}{mtk.175}
\bookmark[dest={mtk.175},level=3]{2. ขีณาสวภาวปญฺห}
\csromanmarkh{5. อนุมานปญฺห}\csromanmarkhii{2. ขีณาสว\-ภาว\-ปญฺห}\csromanheadernomark{3}{5. อนุมานปญฺห \textbf{2. ขีณาสว}\-\textbf{ภาว}\-\textbf{ปญฺห}}
\proseitem{2}{\csromanfolio{257}ภนฺเต นาคเสน ตุเมฺห ภณถ\csromansymbolfootnote{*}{เหฏฺฐา 165 ปิฏฺเฐ.}“โย คิหี อรหตฺตํ ปตฺโต, เทฺว วาสฺส คติโย ภวนฺติ อนญฺญา, ตสฺมิํ เยว ทิวเส ปพฺพชติ วา ปรินิพฺพายติ วา. น โส ทิวโส สกฺกา อติกฺกเมตุนฺ”ติ. สเจ โส ภนฺเต นาคเสน ตสฺมิํ ทิวเส อาจริยํ วา อุปชฺฌายํ วา ปตฺตจีวรํ วา น ลเภถ, อปิ นุ โข โส อรหา สยํ วา ปพฺพเชยฺย ทิวสํ วา อติกฺกเมยฺย, อญฺโญ โกจิ อรหา อิทฺธิมา อาคนฺตฺวา ตํ ปพฺพาเชยฺย วา ปรินิพฺพาเยยฺย วาติ. น โส มหาราช อรหา สยํ ปพฺพเชยฺย, สยํ ปพฺพชนฺโต เถยฺยํ อาปชฺชติ, น จ ทิวสํ อติกฺกเมยฺย, อญฺญสฺส อรหนฺตสฺส อาคมนํ ภเวยฺย วา น วา ภเวยฺย, ตสฺมิํ เยว ทิวเส ปรินิพฺพาเยยฺยาติ. เตน หิ ภนฺเต นาคเสน อรหตฺตสฺส สนฺตภาโว วิชหิโต โหติ, เยน อธิคตสฺส ชีวิตหาโร ภวตีติ.}

\prose{วิสมํ มหาราช คิหิลิงฺคํ, วิสเม ลิงฺเค ลิงฺค\-ทุพฺพลตาย อรหตฺตํ ปตฺโต คิหี ตสฺมิํ เยว ทิวเส ปพฺพชติ วา ปรินิพฺพายติ วา, เนโส มหาราช โทโส อรหตฺตสฺส, คิหิลิงฺคสฺเส\-เวโส โทโส ยทิทํ ลิงฺคทุพฺพลตา.}

\prose{ยถา มหาราช โภชนํ สพฺพสตฺตานํ อายุปาลกํ ชีวิต\-รกฺขกํ วิสม\-โกฏฺฐสฺส มนฺท\-ทุพฺพล\-คหณิกสฺส อวิปาเกน ชีวิตํ หรติ, เนโส มหาราช โทโส โภชนสฺส, โกฏฺฐสฺเสเวโส โทโส ยทิทํ อคฺคิทุพฺพลตา. เอวเมว โข มหาราช วิสเม ลิงฺเค ลิงฺค\-ทุพฺพลตาย อรหตฺตํ ปตฺโต คิหี ตสฺมิํ เยว ทิวเส ปพฺพชติ วา ปรินิพฺพายติ วา, เนโส มหาราช โทโส อรหตฺตสฺส, คิหิลิงฺคสฺเส\-เวโส โทโส ยทิทํ ลิงฺคทุพฺพลตา.}

\prose{ยถา วา ปน มหาราช ปริตฺตํ ติณสลากํ อุปริ ครุเก ปาสาเณ ฐปิเต ทุพฺพลตาย ภิชฺชิตฺวา ปตติ. เอวเมว โข มหาราช อรหตฺตํ ปตฺโต คิหี เตน ลิงฺเคน อรหตฺตํ ธเรตุํ อสกฺโกนฺโต ตสฺมิํ เยว ทิวเส ปพฺพชติ วา ปรินิพฺพายติ วา.}

\prosewithcloserruled{ยถา วา ปน มหาราช ปุริโส อพโล ทุพฺพโล นิหีนชจฺโจ ปริตฺตปุญฺโญ มหติ\-มหารชฺชํ ลภิตฺวา ขเณน ปริปตติ ปริธํสติ โอสกฺกติ, น สกฺโกติ อิสฺสริยํ ธาเรตุํ. เอวเมว โข มหาราช \csromanfolio{258}อรหตฺตํ ปตฺโต คิหี เตน ลิงฺเคน อรหตฺตํ ธาเรตุํ น สกฺโกติ, เตน การเณน ตสฺมิํ เยว ทิวเส ปพฺพชติ วา ปรินิพฺพายติ วาติ. สาธุ ภนฺเต นาคเสน, เอวเมตํ ตถา สมฺปฏิจฺฉามีติ.}{ขีณาสว\-ภาว\-ปโญฺห ทุติโย.}
\csromanmatikaanchor{mtk.176}
\csromantocmarkref{subsubsection}{3. ขีณาสวสติสมฺโมสปญฺห}{mtk.176}
\bookmark[dest={mtk.176},level=3]{3. ขีณาสวสติสมฺโมสปญฺห}
\csromanheader{3}{3. \textbf{ขีณาสว}\-\textbf{สติสมฺโมส}\-\textbf{ปญฺห}}
\proseitem{3}{ภนฺเต นาคเสน อตฺถิ อรหโต สติสมฺโมโสติ. วิคต\-สติสมฺโมสา โข มหาราช อรหนฺโต, นตฺถิ อรหนฺตานํ สติสมฺโมโสติ. อาปชฺเชยฺย ปน ภนฺเต อรหา อาปตฺตินฺติ. อาม มหาราชาติ. กิสฺมิํ วตฺถุสฺมินฺติ. กุฏิกาเร, มหาราช สญฺจริตฺเต, วิกาเล กาลสญฺญาย, ปวาริเต อปฺปวาริต\-สญฺญาย, อนติริตฺเต อติริตฺต\-สญฺญายาติ.}

\prose{ภนฺเต นาคเสน ตุเมฺห ภณถ “เย อาปตฺติํ อาปชฺชนฺติ, เต ทฺวีหิ การเณหิ อาปชฺชนฺติ อนาทริเยน วา อชานเนน วา”ติ. อปิ นุ โข ภนฺเต อรหโต อนาทริยํ โหติ, ยํ อรหา อาปตฺติํ อาปชฺชตีติ. น หิ มหาราชาติ.}

\prose{ยทิ ภนฺเต นาคเสน อรหา อาปตฺติํ อาปชฺชติ, นตฺถิ จ อรหโต อนาทริยํ, เตน หิ อตฺถิ อรหโต สติสมฺโมโสติ. นตฺถิ มหาราช อรหโต สติสมฺโมโส, อาปตฺติญฺจ อรหา อาปชฺชตีติ.}

\prose{เตน หิ ภนฺเต การเณน มํ สญฺญาเปหิ, กิํ ตตฺถ การณนฺติ. เทฺวเม มหาราช กิเลสา โลกวชฺชํ ปณฺณตฺติวชฺช\-ญฺจาติ. กตมํ มหาราช โลกวชฺชํ, ทส อกุสล\-กมฺมปถา, อิทํ วุจฺจติ โลกวชฺชํ. กตมํ ปณฺณตฺติวชฺชํ, ยํ โลเก อตฺถิ สมณานํ อนนุจฺฉวิกํ อนนุโลมิกํ, คิหีนํ อนวชฺชํ. ตตฺถ ภควา สาวกานํ สิกฺขาปทํ ปญฺญเปติ “ยาวชีวํ อนติกฺกมนียนฺ”ติ. วิกาลโภชนํ มหาราช โลกสฺส อนวชฺชํ, ตํ ชินสาสเน วชฺชํ. ภูตคาม\-วิโกปนํ มหาราช โลกสฺส อนวชฺชํ, ตํ ชินสาสเน วชฺชํ. อุทเก หสฺสธมฺมํ มหาราช โลกสฺส อนวชฺชํ, ตํ ชินสาสเน วชฺชํ. อิติ เอวรูปานิ เอวรูปานิ \csromanfolio{259}มหาราช ชินสาสเน วชฺชานิ, อิทํ วุจฺจติ ปณฺณตฺติวชฺชํ. โลกวชฺชํ อภพฺโพ ขีณาสโว ตํ อชฺฌาจริตุํ.}

\prosewithcloserruled{ยํ กิเลสํ ปณฺณตฺติวชฺชํ, ตํ อชานนฺโต อาปชฺเชยฺย. อวิสโย มหาราช เอกจฺจสฺส อรหโต สพฺพํ ชานิตุํ, น หิ ตสฺส พลํ อตฺถิ สพฺพํ ชานิตุํ. อนญฺญาตํ มหาราช อรหโต อิตฺถิปุริสานํ นามมฺปิ โคตฺตมฺปิ, มคฺโคปิ ตสฺส มหิยา อนญฺญาโต, วิมุตฺติํ เยว มหาราช เอกจฺโจ อรหา ชาเนยฺย, ฉฬภิญฺโญ อรหา สกวิสยํ ชาเนยฺย, สพฺพญฺญู มหาราช ตถาคโตว สพฺพํ ชานาตีติ. สาธุ ภนฺเต นาคเสน, เอวเมตํ ตถา สมฺปฏิจฺฉามีติ.}{ขีณาสว\-สติสมฺโมส\-ปโญฺห ตติโย.}
\csromanmatikaanchor{mtk.177}
\csromantocmarkref{subsubsection}{4. โลเก นตฺถิภาวปญฺห}{mtk.177}
\bookmark[dest={mtk.177},level=3]{4. โลเก นตฺถิภาวปญฺห}
\csromanheader{3}{4. \textbf{โลเก นตฺถิภาวปญฺห}}
\proseitem{3}{ภนฺเต นาคเสน ทิสฺสนฺติ โลเก พุทฺธา, ทิสฺสนฺติ ปจฺเจกพุทฺธา, ทิสฺสนฺติ ตถาคตสฺส สาวกา, ทิสฺสนฺติ จกฺกวตฺติ\-ราชาโน, ทิสฺสนฺติ ปเทสราชาโน, ทิสฺสนฺติ เทวมนุสฺสา, ทิสฺสนฺติ สธนา, ทิสฺสนฺติ อธนา, ทิสฺสนฺติ สุคตา, ทิสฺสนฺติ ทุคฺคตา, ทิสฺสติ ปุริสสฺส อิตฺถิลิงฺคํ ปาตุภูตํ, ทิสฺสติ อิตฺถิยา ปุริสลิงฺคํ ปาตุภูตํ, ทิสฺสติ สุกตํ ทุกฺกตํ กมฺมํ, ทิสฺสนฺติ กลฺยาณ\-ปาปกานํ กมฺมานํ วิปากู\-ปโภคิโน สตฺตา, อตฺถิ โลเก สตฺตา อณฺฑชาชลาพุชา สํเสทชา โอปปาติกา, อตฺถิ สตฺตา อปทา ทฺวิปทา จตุปฺปทา พหุปฺปทา, อตฺถิ โลเก ยกฺขา รกฺขสา กุมฺภณฺฑา อสุรา ทานวา คนฺธพฺพา เปตา ปิสาจา, อตฺถิ กินฺนรา มโหรคา นาคา สุปณฺณา สิทฺธา วิชฺชาธรา, อตฺถิ หตฺถี อสฺสา คาโว มหิํสา\csromansharedfootnote{3befb57da19b7705}{มหิสา (สี, อิ)} โอฏฺฐา คทฺรภา อชา เอฬกา มิคา สูกรา สีหา พฺยคฺฆา ทีปี อจฺฉา โกกา ตรจฺฉา โสณา สิงฺคาลา, อตฺถิ พหุวิธา สกุณา, อตฺถิ สุวณฺณํ รชตํ มุตฺตา มณิ สงฺโข สิลา ปวาฬํ โลหิตงฺโก มสารคลฺลํ เวฬุริโย วชิรํ ผลิกํ กาฬโลหํ ตมฺพโลหํ วฏฺฏโลหํ กํสโลหํ, อตฺถิ โขมํ โกเสยฺยํ กปฺปาสิกํ สาณํ ภงฺคํ กมฺพลํ, อตฺถิ สาลิ วีหิ ยโว กงฺคุ กุทฺรูโส วรโก \csromanfolio{260}โคธูโม มุคฺโค มาโส ติลํ กุลตฺถํ, อตฺถิ มูลคนฺโธ สารคนฺโธ เผคฺคุคนฺโธ ตจคนฺโธ ปตฺตคนฺโธ ปุปฺผคนฺโธ ผลคนฺโธ สพฺพคนฺโธ, อตฺถิ ติณ ลตา คจฺฉ รุกฺข โอสธิ วนปฺปติ นที ปพฺพต สมุทฺท มจฺฉกจฺฉปา สพฺพํ โลเก อตฺถิ. ยํ ภนฺเต โลเก นตฺถิ, ตํ เม กเถหีติ.}

\proseitemwithcloserruled{4}{ตีณิมานิ มหาราช โลเก นตฺถิ. กตมานิ ตีณิ, สเจตนา วา อเจตนา วา อชรามรา โลเก นตฺถิ, สงฺขารานํ นิจฺจตา นตฺถิ, ปรมตฺเถน สตฺตูปลทฺธิ นตฺถิ, อิมานิ โข มหาราช ตีณิ โลเก นตฺถีติ. สาธุ ภนฺเต นาคเสน, เอวเมตํ ตถา สมฺปฏิจฺฉามีติ.}{โลเก นตฺถิภาว\-ปโญฺห จตุตฺโถ.}
\csromanmatikaanchor{mtk.178}
\csromantocmarkref{subsubsection}{5. อกมฺมชาทิปญฺห}{mtk.178}
\bookmark[dest={mtk.178},level=3]{5. อกมฺมชาทิปญฺห}
\csromanheader{3}{5. \textbf{อกมฺมชา}\-\textbf{ทิ}\-\textbf{ปญฺห}}
\proseitem{4}{ภนฺเต นาคเสน ทิสฺสนฺติ โลเก กมฺมนิพฺพตฺตา, ทิสฺสนฺติ เหตุนิพฺพตฺตา, ทิสฺสนฺติ อุตุนิพฺพตฺตา, ยํ โลเก อกมฺมชํ อเหตุชํ อนุตุชํ, ตํ เม กเถหีติ. เทฺวเม มหาราช โลกสฺมิํ อกมฺมชา อเหตุชา อนุตุชา. กตเม เทฺว, อากาโส มหาราช อกมฺมโช อเหตุโช อนุตุโช, นิพฺพานํ มหาราช อกมฺมชํ อเหตุชํ อนุตุชํ. อิเม โข มหาราช เทฺว อกมฺมชา อเหตุชา อนุตุชาติ.}

\prose{มา ภนฺเต นาคเสน ชินวจนํ มกฺเขหิ, มา อชานิตฺวา ปญฺหํ พฺยากโรหีติ. กิํ โข มหาราช อหํ วทามิ, ยํ มํ ตฺวํ เอวํ วเทสิ “มา ภนฺเต นาคเสน ชินวจนํ มกฺเขหิ, มา อชานิตฺวา ปญฺหํ พฺยากโรหี”ติ. ภนฺเต นาคเสน ยุตฺตมิทํ ตาว วตฺตุํ “อากาโส อกมฺมโช อเหตุโช อนุตุโช”ติ. อเนกสเตหิ ปน ภนฺเต นาคเสน การเณหิ ภควตา สาวกานํ นิพฺพานสฺส สจฺฉิกิริยาย มคฺโค อกฺขาโต, อถ จ ปน ตฺวํ เอวํ วเทสิ “อเหตุชํ นิพฺพานนฺ”ติ. สจฺจํ มหาราช ภควตา อเนกสเตหิ การเณหิ สาวกานํ นิพฺพานสฺส สจฺฉิกิริยาย มคฺโค อกฺขาโต, น จ ปน นิพฺพานสฺส อุปฺปาทาย เหตุ อกฺขาโตติ.}

\prose{เอตฺถ มยํ ภนฺเต นาคเสน อนฺธการโต อนฺธการตรํ ปวิสาม, วนโต วนตรํ ปวิสาม, คหนโต คหนตรํ\csromansharedfootnote{705fb6003121f35d}{คหนนฺตรโต คหนนฺตรํ (ก)} ปวิสาม, ยตฺร หิ \csromanfolio{261}นาม นิพฺพานสฺส สจฺฉิกิริยาย เหตุ อตฺถิ, ตสฺส ปน ธมฺมสฺส อุปฺปาทาย เหตุ นตฺถิ. ยทิ ภนฺเต นาคเสน นิพฺพานสฺส สจฺฉิกิริยาย เหตุ อตฺถิ, เตน หิ นิพฺพานสฺส อุปฺปาทายปิ เหตุ อิจฺฉิตพฺโพ.}

\proseitem{5}{ยถา ปน ภนฺเต นาคเสน ปุตฺตสฺส ปิตา อตฺถิ, เตน การเณน ปิตุโนปิ ปิตา อิจฺฉิตพฺโพ. ยถา อนฺเตวาสิกสฺส อาจริโย อตฺถิ, เตน การเณน อาจริยสฺสปิ อาจริโย อิจฺฉิตพฺโพ. ยถา องฺกุรสฺส พีชํ อตฺถิ, เตน การเณน พีชสฺสปิ พีชํ อิจฺฉิตพฺพํ. เอวเมว โข ภนฺเต นาคเสน ยทิ นิพฺพานสฺส สจฺฉิกิริยาย เหตุ อตฺถิ, เตน การเณน นิพฺพานสฺส อุปฺปาทายปิ เหตุ อิจฺฉิตพฺโพ.}

\prose{ยถา รุกฺขสฺส วา ลตาย วา อคฺเค สติ เตน การเณน มชฺฌมฺปิ อตฺถิ, มูลมฺปิ อตฺถิ. เอวเมว โข ภนฺเต นาคเสน ยทิ นิพฺพานสฺส สจฺฉิกิริยาย เหตุ อตฺถิ, เตน การเณน นิพฺพานสฺส อุปฺปาทายปิ เหตุ อิจฺฉิตพฺโพติ.}

\prose{อนุปฺปาทนียํ มหาราช นิพฺพานํ, ตสฺมา น นิพฺพานสฺส อุปฺปาทาย เหตุ อกฺขาโตติ. อิงฺฆ ภนฺเต นาคเสน การณํ ทสฺเสตฺวา การเณน มํ สญฺญาเปหิ, ยถาหํ ชาเนยฺยํ นิพฺพานสฺส สจฺฉิกิริยาย เหตุ อตฺถิ, นิพฺพานสฺส อุปฺปาทาย เหตุ นตฺถีติ.}

\prose{เตน หิ มหาราช สกฺกจฺจํ โสตํ โอทห, สาธุกํ สุโณหิ, วกฺขามิ ตตฺถ การณํ, สกฺกุเณยฺย มหาราช ปุริโส ปากติเกน พเลน อิโต หิมวนฺตํ ปพฺพตราชํ อุปคนฺตุนฺติ. อาม ภนฺเตติ. สกฺกุเณยฺย ปน โส มหาราช ปุริโส ปากติเกน พเลน หิมวนฺตํ ปพฺพตราชํ อิธ อาหริตุนฺติ. น หิ ภนฺเตติ. เอวเมว โข มหาราช สกฺกา นิพฺพานสฺส สจฺฉิกิริยาย มคฺโค อกฺขาตุํ, น สกฺกา นิพฺพานสฺส อุปฺปาทาย เหตุ ทสฺเสตุํ.}

\prose{สกฺกุเณยฺย มหาราช ปุริโส ปากติเกน พเลน มหาสมุทฺทํ นาวาย อุตฺตริตฺวา ปาริมตีรํ คนฺตุนฺติ. อาม ภนฺเตติ. สกฺกุเณยฺย ปน โส มหาราช ปุริโส ปากติเกน พเลน มหาสมุทฺทสฺส ปาริมตีรํ อิธ อาหริตุนฺติ. น หิ ภนฺเตติ. เอวเมว โข มหาราช สกฺกา นิพฺพานสฺส สจฺฉิกิริยาย มคฺโค อกฺขาตุํ, น สกฺกา นิพฺพานสฺส อุปฺปาทาย เหตุ ทสฺเสตุํ. กิํ การณา, อสงฺขตตฺตา ธมฺมสฺสาติ.}

\prose{\csromanfolio{262}อสงฺขตํ ภนฺเต นาคเสน นิพฺพานนฺติ. อาม มหาราช อสงฺขตํ นิพฺพานํ น เกหิจิ กตํ, นิพฺพานํ มหาราช น วตฺตพฺพํ อุปฺปนฺนนฺติ วา อนุปฺปนฺนนฺติ วา อุปฺปาทนียนฺติ วา อตีตนฺติ วา อนาคตนฺติ วา ปจฺจุปฺปนฺนนฺติ วา จกฺขุ\-วิญฺเญยฺยนฺติ วา โสต\-วิญฺเญยฺยนฺติ วา ฆาน\-วิญฺเญยฺยนฺติ วา ชิวฺหา\-วิญฺเญยฺยนฺติ วา กาย\-วิญฺเญยฺยนฺติ วาติ. ยทิ ภนฺเต นาคเสน นิพฺพานํ น อุปฺปนฺนํ น อนุปฺปนฺนํ น อุปฺปาทนียํ น อตีตํ น อนาคตํ น ปจฺจุปฺปนฺนํ น จกฺขุ\-วิญฺเญยฺยํ น โสตวิญฺเญยฺยํ น ฆานวิญฺเญยฺยํ น ชิวฺหาวิญฺเญยฺยํ น กายวิญฺเญยฺยํ, เตน หิ ภนฺเต นาคเสน ตุเมฺห นตฺถิธมฺมํ นิพฺพานํ อปทิสถ “นตฺถิ นิพฺพานนฺ”ติ. อตฺถิ มหาราช นิพฺพานํ, มโนวิญฺเญยฺยํ นิพฺพานํ, วิสุทฺเธน มานเสน ปณีเตน อุชุเกน อนาวรเณน นิรามิเสน สมฺมาปฏิปนฺโน อริยสาวโก นิพฺพานํ ปสฺสตีติ.}

\prosewithcloserruled{กีทิสํ ปน ตํ ภนฺเต นิพฺพานํ, ยํ ตํ โอปมฺเมหิ อาทีปนียํ การเณหิ มํ สญฺญาเปหิ, ยถา อตฺถิธมฺมํ โอปมฺเมหิ อาทีปนียนฺติ. อตฺถิ มหาราช วาโต นามาติ. อาม ภนฺเตติ. อิงฺฆ มหาราช วาตํ ทสฺเสหิ วณฺณโต วา สณฺฐานโต วา อณุํ วา ถูลํ วา ทีฆํ วา รสฺสํ วาติ. น สกฺกา ภนฺเต นาคเสน วาโต อุปทสฺสยิตุํ, น โส วาโต หตฺถคฺคหณํ วา นิมฺมทฺทนํ วา อุเปติ, อปิ จ อตฺถิ โส วาโตติ. ยทิ มหาราช นสกฺกา วาโต อุปทสฺสยิตุํ, เตน หิ นตฺถิ วาโตติ. ชานามหํ ภนฺเต นาคเสน, วาโต อตฺถีติ เม หทเย อนุปวิฏฺฐํ, น จาหํ สกฺโกมิ วาตํ อุปทสฺสยิ\-ตุนฺติ. เอวเมว โข มหาราช อตฺถิ นิพฺพานํ, น จ สกฺกา นิพฺพานํ อุปทสฺสยิตุํ วณฺเณน วา สณฺฐาเนน วาติ. สาธุ ภนฺเต นาคเสน, สูปทสฺสิตํ โอปมฺมํ, สุนิทฺทิฏฺฐํ การณํ, เอวเมตํ ตถา สมฺปฏิจฺฉามิ “อตฺถิ นิพฺพานนฺ”ติ.}{อกมฺมชา\-ทิ\-ปโญฺห ปญฺจโม.}
\csromanmatikaanchor{mtk.179}
\csromantocmarkref{subsubsection}{6. กมฺมชาทิปญฺห}{mtk.179}
\bookmark[dest={mtk.179},level=3]{6. กมฺมชาทิปญฺห}
\csromanheader{3}{6. \textbf{กมฺมชาทิปญฺห}}
\proseitemwithcloserruled{6}{ภนฺเต นาคเสน กตเม เอตฺถ กมฺมชา, กตเม เหตุชา, กตเม อุตุชา, กตเม น กมฺมชา, น เหตุชา, น อุตุชาติ. เย เกจิ มหาราช สตฺตา สเจตนา, สพฺเพ เต กมฺมชา, อคฺคิ จ สพฺพานิ \csromanfolio{263}จ พีชชาตานิ เหตุชานิ, ปถวี จ ปพฺพตา จ อุทกญฺจ วาโต จ, สพฺเพ เต อุตุชา, อากาโส จ นิพฺพานญฺจ อิเม เทฺว อกมฺมชา อเหตุชา อนุตุชา, นิพฺพานํ ปน มหาราช น วตฺตพฺพํ กมฺมชนฺติ วา เหตุชนฺติ วา อุตุชนฺติ วา อุปฺปนฺนนฺติ วา อนุปฺปนฺนนฺติ วา อุปฺปาทนียนฺติ วา อตีตนฺติ วา อนาคตนฺติ วา ปจฺจุปฺปนฺนนฺติ วา จกฺขุ\-วิญฺเญยฺยนฺติ วา โสต\-วิญฺเญยฺยนฺติ วา ฆาน\-วิญฺเญยฺยนฺติ วา ชิวฺหา\-วิญฺเญยฺยนฺติ วา กาย\-วิญฺเญยฺยนฺติ วา, อปิ จ มหาราช มโนวิญฺเญยฺยํ นิพฺพานํ, ยํ โส สมฺมาปฏิปนฺโน อริยสาวโก วิสุทฺเธน ญาเณน ปสฺสตีติ. รมณีโย ภนฺเต นาคเสน ปโญฺห สุวินิจฺฉิโต นิสฺสํสโย เอกนฺตคโต, วิมติ อุปฺปจฺฉินฺนา, ตฺวํ คณิวร\-ปวรมา\-สชฺชาติ.}{กมฺมชาทิปโญฺห ฉฏฺโฐ.}
\csromanmatikaanchor{mtk.180}
\csromantocmarkref{subsubsection}{7. ยกฺขปญฺห}{mtk.180}
\bookmark[dest={mtk.180},level=3]{7. ยกฺขปญฺห}
\csromanheader{3}{7. \textbf{ยกฺขปญฺห}}
\proseitemwithcloserruled{7}{ภนฺเต นาคเสน อตฺถิ โลเก ยกฺขา นามาติ. อาม มหาราช อตฺถิ โลเก ยกฺขา นามาติ. จวนฺติ ปน เต ภนฺเต ยกฺขา ตมฺหา โยนิยาติ. อาม มหาราช จวนฺติ เต ยกฺขา ตมฺหา โยนิยาติ. กิสฺส ปน ภนฺเต นาคเสน เตสํ มตานํ ยกฺขานํ สรีรํ น ทิสฺสติ, กุณปคนฺโธปิ น วายตีติ. ทิสฺสติ มหาราช มตานํ ยกฺขานํ สรีรํ, กุณปคนฺโธปิ เตสํ วายติ, มตานํ มหาราช ยกฺขานํ สรีรํ กีฏวณฺเณน วา ทิสฺสติ, กิมิวณฺเณน วา ทิสฺสติ, กิปิลฺลิก\-วณฺเณน วา ทิสฺสติ, ปฏงฺควณฺเณน วา ทิสฺสติ, อหิวณฺเณน วา ทิสฺสติ, วิจฺฉิก\-วณฺเณน วา ทิสฺสติ, สตปทิ\-วณฺเณน วา ทิสฺสติ, ทิชวณฺเณน วา ทิสฺสติ, มิควณฺเณน วา ทิสฺสตีติ. โก หิ ภนฺเต นาคเสน อญฺโญ อิทํ ปญฺหํ ปุฏฺโฐ วิสชฺเชยฺย อญฺญตฺร ตวาทิเสน พุทฺธิมตาติ.}{ยกฺขปโญฺห สตฺตโม.}
\csromanmatikaanchor{mtk.181}
\csromantocmarkref{subsubsection}{8. อนวเสสสิกฺขาปทปญฺห}{mtk.181}
\bookmark[dest={mtk.181},level=3]{8. อนวเสสสิกฺขาปทปญฺห}
\csromanheader{3}{8. \textbf{อนวเสส}\-\textbf{สิกฺขาปท}\-\textbf{ปญฺห}}
\proseitem{8}{ภนฺเต นาคเสน เย เต อเหสุํ ติกิจฺฉกานํ ปุพฺพกา อาจริยา. เสยฺยถีทํ, นารโท ธมฺมนฺตรี\csromansharedfootnote{5af40e53edd969a4}{ธนฺวนฺตรี (?)} องฺคิรโส กปิโล กณฺฑรคฺคิ \csromanfolio{264}สาโม อตุโล ปุพฺพกจฺจายโน, สพฺเพเปเต อาจริยา สกิํ เยว โรคุปฺปตฺติญฺจ นิทานญฺจ สภาวญฺจ สมุฏฺฐานญฺจ ติกิจฺฉญฺจ กิริยญฺจ สิทฺธา\-สิทฺธญฺจ สพฺพํ ตํ\csromansharedfootnote{93bf52174ead1d01}{สนฺตํ (ก)} นิรวเสสํ ชานิตฺวา “อิมสฺมิํ กาเย เอตฺตกา โรคา อุปฺปชฺชิสฺสนฺตี”ติ เอกปฺปหาเรน กลาปคฺคาหํ กริตฺวา สุตฺตํ พนฺธิํสุ, อสพฺพญฺญุโน เอเต สพฺเพ, กิสฺส ปน ตถาคโต สพฺพญฺญู สมาโน อนาคตํ กิริยํ พุทฺธญาเณน ชานิตฺวา “เอตฺตเก นาม วตฺถุสฺมิํ เอตฺตกํ นาม สิกฺขาปทํ ปญฺญเปตพฺพํ ภวิสฺสตี”ติ ปริจฺฉินฺทิตฺวา อนวเสสโต สิกฺขาปทํ น ปญฺญเปสิ, อุปฺปนฺนุปฺปนฺเน วตฺถุสฺมิํ อยเส ปากเฏ โทเส วิตฺถาริเก ปุถุคเต อุชฺฌายนฺเตสุ มนุสฺเสสุ ตสฺมิํ ตสฺมิํ กาเล สาวกานํ สิกฺขาปทํ ปญฺญเปสีติ.}

\prosewithcloserruled{ญาตเมตํ มหาราช ตถาคตสฺส “อิมสฺมิํ สมเย อิเมสุ มนุสฺเสสุ สาธิกํ ทิยฑฺฒ\-สิกฺขาปทสตํ ปญฺญเปตพฺพํ ภวิสฺสตี”ติ, อปิ จ ตถาคตสฺส เอวํ อโหสิ “สเจ โข อหํ สาธิกํ ทิยฑฺฒ\-สิกฺขาปทสตํ เอกปฺปหารํ ปญฺญเปสฺสามิ, มหาชโน สนฺตา\-สมาปชฺชิสฺสติ ‘พหุกํ อิธ รกฺขิตพฺพํ, ทุกฺกรํ วต โภ สมณสฺส โคตมสฺส สาสเน ปพฺพชิตุนฺ’ติ, ปพฺพชิตุกามาปิ น ปพฺพชิสฺสนฺติ, วจนญฺจ เม น สทฺทหิสฺสนฺติ, อสทฺทหนฺตา เต มนุสฺสา อปายคามิโน ภวิสฺสนฺติ. อุปฺปนฺนุปฺปนฺเน วตฺถุสฺมิํ ธมฺมเทสนาย วิญฺญาเปตฺวา ปากเฏ โทเส สิกฺขาปทํ ปญฺญเปสฺสามี”ติ. อจฺฉริยํ ภนฺเต นาคเสน พุทฺธานํ, อพฺภุตํ ภนฺเต นาคเสน พุทฺธานํ, ยาว มหนฺตํ ตถาคตสฺส สพฺพญฺญุต\-ญาณํ, เอวเมตํ ภนฺเต นาคเสน, สุนิทฺทิฏฺโฐ เอโส อตฺโถ ตถาคเตน, “พหุกํ อิธ สิกฺขิตพฺพนฺ”ติ สุตฺวา สตฺตานํ สนฺตาโส อุปฺปชฺเชยฺย, เอโกปิ ชินสาสเน น ปพฺพเชยฺย, เอวเมตํ ตถา สมฺปฏิจฺฉามีติ.}{อนวเสส\-สิกฺขาปท\-ปโญฺห อฏฺฐโม.}
\csromanmatikaanchor{mtk.182}
\csromantocmarkref{subsubsection}{9. สูริยตปนปญฺห}{mtk.182}
\bookmark[dest={mtk.182},level=3]{9. สูริยตปนปญฺห}
\csromanheader{3}{9. \textbf{สูริย}\-\textbf{ตปน}\-\textbf{ปญฺห}}
\proseitemwithcloserruled{9}{ภนฺเต นาคเสน อยํ สูริโย สพฺพกาลํ กฐินํ ตปติ, อุทาหุ กิญฺจิกาลํ มนฺทํ ตปตีติ. สพฺพกาลํ มหาราช สูริโย กฐินํ \csromanfolio{265}ตปติ, น กิญฺจิกาลํ มนฺทํ ตปตีติ. ยทิ ภนฺเต นาคเสน สูริโย สพฺพกาลํ กฐินํ ตปติ, กิสฺส ปน อปฺเปกทา สูริโย กฐินํ ตปติ, อปฺเปกทา มนฺทํ ตปตีติ. จตฺตาโรเม มหาราช สูริยสฺส โรคา, เยสํ อญฺญตเรน โรเคน ปฏิปีฬิโต สูริโย มนฺทํ ตปติ. กตเม จตฺตาโร, อพฺภํ มหาราช สูริยสฺส โรโค, เตน โรเคน ปฏิปีฬิโต สูริโย มนฺทํ ตปติ. มหิกา มหาราช สูริยสฺส โรโค, เตน โรเคน ปฏิปีฬิโต สูริโย มนฺทํ ตปติ. เมโฆ มหาราช สูริยสฺส โรโค, เตน โรเคน ปฏิปีฬิโต สูริโย มนฺทํ ตปติ. ราหุ มหาราช สูริยสฺส โรโค, เตน โรเคน ปฏิปีฬิโต สูริโย มนฺทํ ตปติ. อิเม โข มหาราช จตฺตาโร สูริยสฺส โรคา, เยสํ อญฺญตเรน โรเคน ปฏิปีฬิโต สูริโย มนฺทํ ตปตีติ. อจฺฉริยํ ภนฺเต นาคเสน, อพฺภุตํ ภนฺเต นาคเสน, สูริยสฺสปิ ตาว เตโช\-สมฺปนฺนสฺส โรโค อุปฺปชฺชิสฺสติ, กิมงฺคํ ปน อญฺเญสํ สตฺตานํ, นตฺถิ ภนฺเต เอสา วิภตฺติ อญฺญสฺส อญฺญตฺร ตวาทิเสน พุทฺธิมตาติ.}{สูริย\-ตปน\-ปโญฺห นวโม.}
\csromanmatikaanchor{mtk.183}
\csromantocmarkref{subsubsection}{10. กฐินตปนปญฺห}{mtk.183}
\bookmark[dest={mtk.183},level=3]{10. กฐินตปนปญฺห}
\csromanheader{3}{10. \textbf{กฐิน}\-\textbf{ตปน}\-\textbf{ปญฺห}}
\proseitem{10}{ภนฺเต นาคเสน กิสฺส เหมนฺเต สูริโย กฐินํ ตปติ, โน ตถา คิเมฺหติ. คิเมฺห มหาราช อนุปหตํ โหติ รโชชลฺลํ, วาตกฺขุภิตา เรณู คคนานุคตา โหนฺติ, อากาเสปิ อพฺภา สุพหลา โหนฺติ, มหาวาโต จ อธิมตฺตํ วายติ, เต สพฺเพ นานากุลา สมายุตา สูริยรํสิโย ปิทหนฺติ, เตน คิเมฺห สูริโย มนฺทํ ตปติ.}

\prose{เหมนฺเต ปน มหาราช เหฏฺฐา ปถวี นิพฺพุตา โหติ, อุปริ มหาเมโฆ อุปฏฺฐิโต โหติ, อุปสนฺตํ โหติ รโชชลฺลํ, เรณุ จ สนฺตสนฺตํ คคเน จรติ, วิคตวลาหโก จ โหติ อากาโส, \csromanfolio{266}วาโต จ มนฺทมนฺทํ วายติ, เอเตสํ อุปรติยา วิสุทฺธา\csromansharedfootnote{8eae28bb867b90a7}{วิสทา (สี, อิ)} โหนฺติ สูริยรํสิโย, อุปฆาต\-วิมุตฺตสฺส สูริยสฺส ตาโป อติ วิย ตปติ. อิทเมตฺถ มหาราช การณํ, เยน การเณน สูริโย เหมนฺเต กฐินํ ตปติ, โน ตถา คิเมฺหติ. สพฺพีติมุตฺโต ภนฺเต สูริโย กฐินํ ตปติ, เมฆา\-ทิ\-สหคโต กฐินํ น ตปตีติ.}

\vspace{\tipitakaparskip}
\csromancenter{กฐิน\-ตปน\-ปโญฺห ทสโม.}
\nitthitammidruled{\textbf{นิปฺปปญฺจ}\-\textbf{วคฺโค ทุติโย.}}
\csromancenter{อิมสฺมิํ วคฺเค ทส ปญฺหา.}
\csromanmatikaanchor{mtk.184}
\csromantocmarknopage{subsection}{3. เวสฺสนฺตรวคฺค}{mtk.184}
\bookmark[dest={mtk.184},level=2]{3. เวสฺสนฺตรวคฺค}
\csromanheader{2}{3. \textbf{เวสฺสนฺตร}\-\textbf{วคฺค}}
\csromanmatikaanchor{mtk.185}
\csromantocmarkref{subsubsection}{1. เวสฺสนฺตรปญฺห}{mtk.185}
\bookmark[dest={mtk.185},level=3]{1. เวสฺสนฺตรปญฺห}
\csromanheader{3}{1. \textbf{เวสฺสนฺตร}\-\textbf{ปญฺห}}
\proseitem{1}{\csromanfolio{267}ภนฺเต นาคเสน สพฺเพว โพธิสตฺตา ปุตฺตทารํ เทนฺติ, อุทาหุ เวสฺสนฺตเรเนว รญฺญา ปุตฺตทารํ ทินฺนนฺติ. สพฺเพปิ มหาราช โพธิสตฺตา ปุตฺตทารํ เทนฺติ, น เวสฺสนฺตเรเนว รญฺญา ปุตฺตทารํ ทินฺนนฺติ. อปิ จ โข ภนฺเต นาคเสน เตสํ อนุมเตน เทนฺตีติ. ภริยา มหาราช อนุมตา, ทารกา ปน พาลตาย วิลปิํสุ\csromansharedfootnote{338e035521bcca4b}{ลาลปิํสุ (สี, อิ)}, ยทิ เต อตฺถโต ชาเนยฺยุํ, เตปิ อนุโมเทยฺยุํ, น เต วิลเปยฺยุนฺติ.}

\prose{ทุกฺกรํ ภนฺเต นาคเสน โพธิสตฺเตน กตํ, ยํ โส อตฺตโน โอรเส ปิเย ปุตฺเต พฺราหฺมณสฺส ทาสตฺถาย อทาสิ.}

\prose{อิทมฺปิ ทุติยํ ทุกฺกรโต ทุกฺกรตรํ, ยํ โส อตฺตโน โอรเส ปิเย ปุตฺเต พาลเก ตรุณเก ลตาย พนฺธิตฺวา เตน พฺราหฺมเณน ลตาย อนุมชฺชียนฺเต ทิสฺวา อชฺฌุเปกฺขิ.}

\prose{อิทมฺปิ ตติยํ ทุกฺกรโต ทุกฺกรตรํ, ยํ โส สเกน พเลน พนฺธนา มุจฺจิตฺวา อาคเต ทารเก สารชฺชมุ\-ปคเต ปุนเทว ลตาย พนฺธิตฺวา อทาสิ.}

\prose{อิทมฺปิ จตุตฺถํ ทุกฺกรโต ทุกฺกรตรํ, ยํ โส ทารเก “อยํ โข ตาต ยกฺโข ขาทิตุํ เนติ อเมฺห”ติ วิลปนฺเต “มา ภายิตฺถา”ติ น อสฺสาเสสิ.}

\prose{อิทมฺปิ ปญฺจมํ ทุกฺกรโต ทุกฺกรตรํ, ยํ โส ชาลิสฺส กุมารสฺส รุทมานสฺส ปาเทสุ นิปติตฺวา “อลํ ตาต กณฺหาชินํ นิวตฺเตหิ, อหเมว คจฺฉามิ ยกฺเขน สห, ขาทตุ มํ ยกฺโข”ติ ยาจมานสฺส เอวํ น สมฺปฏิจฺฉิ.}

\prose{อิทมฺปิ ฉฏฺฐํ ทุกฺกรโต ทุกฺกรตรํ, ยํ โส ชาลิสฺส กุมารสฺส “ปาสาณสมํ นูน เต ตาต หทยํ, ยํ ตฺวํ อมฺหากํ ทุกฺขิตานํ เปกฺขมาโน นิมฺมนุสฺสเก พฺรหารญฺเญ ยกฺเขน นียมาเน น นิวาเรสี”ติ วิลปมานสฺส การุญฺญํ นากาสิ.}

\prose{\csromanfolio{268}อิทมฺปิ สตฺตมํ ทุกฺกรโต ทุกฺกรตรํ, ยํ ตสฺส รุฬรุฬสฺส ภีมภีมสฺส นีเต ทารเก อทสฺสนํ คมิเต น ผลิ หทยํ สตธา วา สหสฺสธา วา, ปุญฺญกาเมน มนุเชน กิํ ปรทุกฺขาปเนน, นนุ นาม สกทานํ ทาตพฺพํ โหตีติ.}

\prose{ทุกฺกรสฺส มหาราช กตตฺตา โพธิสตฺตสฺส กิตฺติสทฺโท ทสสหสฺสิยา โลกธาตุยา สเทว\-มนุสฺเสสุ อพฺภุคฺคโต, เทวา เทวภวเน ปกิตฺเตนฺติ, อสุรา อสุรภวเน ปกิตฺเตนฺติ, ครุฬา ครุฬภวเน ปกิตฺเตนฺติ, นาคา นาคภวเน ปกิตฺเตนฺติ, ยกฺขา ยกฺขภวเน ปกิตฺเตนฺติ, อนุปุพฺเพน ตสฺส กิตฺติสทฺโท ปรมฺปราย อชฺเชตรหิ อิธ อมฺหากํ สมยํ อนุปฺปตฺโต, ตํ มยํ ทานํ ปกิตฺเตนฺตา วิโกเปนฺตา นิสินฺนา สุทินฺนํ, อุทาหุ ทุทฺทินฺนนฺติ. โส โข ปนายํ มหาราช กิตฺติสทฺโท นิปุณานํ วิญฺญูนํ วิทูนํ วิภาวีนํ โพธิสตฺตานํ ทส คุเณ อนุทสฺสติ. กตเม ทส, อเคธตา นิราลยตา จาโค ปหานํ อปุนราวตฺติตา สุขุมตา มหนฺตตา ทุรนุโพธตา ทุลฺลภตา อสทิสตา พุทฺธธมฺมสฺส, โส โข ปนายํ มหาราช กิตฺติสทฺโท นิปุณานํ วิญฺญูนํ วิทูนํ วิภาวีนํ โพธิสตฺตานํ อิเม ทส คุเณ อนุทสฺสตีติ.}

\prose{ภนฺเต นาคเสน โย ปรํ ทุกฺขาเปตฺวา ทานํ เทติ, อปิ นุ ตํ ทานํ สุขวิปากํ โหติ สคฺค\-สํวตฺตนิกนฺติ. อาม มหาราช กิํ วตฺตพฺพนฺติ. อิงฺฆ ภนฺเต นาคเสน การณํ อุปทสฺเสหีติ. อิธ มหาราช โกจิ สมโณ วา พฺราหฺมโณ วา สีลวา โหติ กลฺยาณธมฺโม, โส ภเวยฺย ปกฺขหโต วา ปีฐสปฺปี วา อญฺญตรํ วา พฺยาธิํ อาปนฺโน, ตเมนํ โย โกจิ ปุญฺญกาโม ยานํ อาโรเปตฺวา ปตฺถิตํ เทสม\-นุปาเปยฺย, อปิ นุ โข มหาราช ตสฺส ปุริสสฺส ตโตนิทานํ กิญฺจิ สุขํ นิพฺพตฺเตยฺย สคฺค\-สํวตฺตนิกํ ตํ กมฺมนฺติ. อาม ภนฺเต กิํ วตฺตพฺพํ, หตฺถิยานํ วา โส ภนฺเต ปุริโส ลเภยฺย อสฺสยานํ วา รถยานํ วา, ถเล ถลยานํ ชเล ชลยานํ เทเวสุ เทวยานํ มนุสฺเสสุ มนุสฺสยานํ, ตทนุจฺฉวิกํ ตทนุโลมิกํ ภเว ภเว นิพฺพตฺเตยฺย, ตทนุจฺฉวิกานิ ตทนุโลมิกานิ จสฺส สุขานิ นิพฺพตฺเตยฺยุํ, สุคติโต สุคติํ คจฺเฉยฺย, เตเนว กมฺมา\-ภิสนฺเทน อิทฺธิยานํ อภิรุยฺห ปตฺถิตํ นิพฺพานนครํ ปาปุเณยฺยาติ. เตน หิ มหาราช ปรทุกฺขาปเนน ทินฺนทานํ สุขวิปากํ \csromanfolio{269}โหติ สคฺค\-สํวตฺตนิกํ, ยํ โส ปุริโส พลีพทฺเท ทุกฺขาเปตฺวา เอวรูปํ สุขํ อนุภวติ.}

\prose{อปรมฺปิ มหาราช อุตฺตริํ การณํ สุโณหิ, ยถา ปรทุกฺขาปเนน ทินฺนทานํ สุขวิปากํ โหติ สคฺค\-สํวตฺตนิกํ, อิธ มหาราช โย โกจิ ราชา ชนปทโต ธมฺมิกํ พลิํ อุทฺธราเปตฺวา อาณาปวตฺตเนน ทานํ ทเทยฺย, อปิ นุ โข โส มหาราช ราชา ตโตนิทานํ กิญฺจิ สุขํ อนุภเวยฺย สคฺค\-สํวตฺตนิกํ ตํ ทานนฺติ. อาม ภนฺเต กิํ วตฺตพฺพํ, ตโตนิทานํ โส ภนฺเต ราชา อุตฺตริํ อเนก\-สตสหสฺสคุณํ ลเภยฺย. ราชูนํ อติราชา ภเวยฺย, เทวานํ อติเทโว ภเวยฺย, พฺรหฺมานํ อติพฺรหฺมา ภเวยฺย, สมณานํ อติสมโณ ภเวยฺย, พฺราหฺมณานํ อติพฺรา- อาหฺมโณ ภเวยฺย, อรหนฺตานํ อติอรหา\csromansharedfootnote{e0aba1ec675f263f}{อติอรหนฺโต (สฺยา, ก)} ภเวยฺยาติ. เตน หิ มหาราช ปรทุกฺขาปเนน ทินฺนทานํ สุขวิปากํ โหติ สคฺค\-สํวตฺตนิกํ, ยํ โส ราชา พลินา ชนํ ปีเฬตฺวา ทินฺนทาเนน เอวรูปํ อุตฺตริํ ยสสุขํ อนุภวตีติ.}

\prose{อติทานํ ภนฺเต นาคเสน เวสฺสนฺตเรน รญฺญา ทินฺนํ, ยํ โส สกํ ภริยํ ปรสฺส ภริยตฺถาย อทาสิ, สเก โอรเส ปุตฺเต พฺราหฺมณสฺส ทาสตฺถาย อทาสิ, อติทานํ นาม ภนฺเต นาคเสน โลเก วิทูหิ นินฺทิตํ ครหิตํ, ยถา นาม ภนฺเต นาคเสน อติภาเรน สกฏสฺส อกฺโข ภิชฺชติ, อติภาเรน นาวา โอสีทติ, อติภุตฺเตน โภชนํ วิสมํ ปริณมติ, อติวสฺเสน ธญฺญํ วินสฺสติ, อติทาเนน โภคกฺขยํ อุเปติ, อติตาเปน ปถวี อุปฑยฺหติ, อติราเคน อุมฺมตฺตโก โหติ, อติโทเสน วชฺโฌ โหติ, อติโมเหน อนยํ อาปชฺชติ, อติโลเภน โจร\-คฺคหณมุ\-ปคจฺฉติ, อติภเยน นิรุชฺฌติ, อติปูเรน นที อุตฺตรติ, อติวาเตน อสนิ ปตติ, อติ- อคฺคินา โอทนํ อุตฺตรติ, อติสญฺจรเณน น จิรํ ชีวติ. เอวเมว โข ภนฺเต นาคเสน อติทานํ นาม โลเก วิทูหิ นินฺทิตํ ครหิตํ, อติทานํ ภนฺเต นาคเสน เวสฺสนฺตเรน รญฺญา ทินฺนํ, น ตตฺถ กิญฺจิ ผลํ อิจฺฉิตพฺพนฺติ.}

\prose{อติทานํ มหาราช โลเก วิทูหิ วณฺณิตํ ถุตํ ปสตฺถํ, เย เกจิ ยาทิสํ กีทิสํ ทานํ เทนฺติ\csromansharedfootnote{525e3afb3ef7dcfa}{เกจิ ยาทิสํ ตาทิสํ ทานํ เทนฺติ (สฺยา)}, อติทานทายี โลเก กิตฺติํ ปาปุณาติ. \csromanfolio{270}ยถา มหาราช อติปวรตาย ทิพฺพํ วนมูลํ คหิตมฺปิ หตฺถปาเส ฐิตานํ ปรชนานํ น ทสฺสยติ, อคโท อติชจฺจตาย\csromansharedfootnote{a12f81eaee9c2b39}{อติอุสภตาย (ก)} ปีฬาย สมุคฺฆาตโก โรคานํ อนฺตกโร, อคฺคิ อติโชติตาย ฑหติ, อุทกํ อติสีตตาย นิพฺพาเปติ, ปทุมํ ปริสุทฺธตาย น อุปลิมฺปติ วาริกทฺทเมน, มณิ อติคุณตาย กามทโท, วชิรํ อติติขิณตาย วิชฺฌติ มณิมุตฺตา\-ผลิกํ, ปถวี อติมหนฺตตาย นโรรคมิคปกฺขิชล เสล ปพฺพต ทุเม ธาเรติ, สมุทฺโท อติมหนฺตตาย อปริปูรโณ, สิเนรุ อติภารตาย อจโล, อากาโส อติวิตฺถาร\-ตาย อนนฺโต, สูริโย อติปฺปภตาย ติมิรํ ฆาเตติ, สีโห อติชาติตาย วิคตภโต, มลฺโล อติพลวตาย ปฏิมลฺลํ ขิปฺปํ อุกฺขิปติ, ราชา อติปุญฺญตาย อธิปติ, ภิกฺขุ อติสีลวนฺต\-ตาย นาค\-ยกฺข\-นร\-มรูหิ นมสฺสนีโย, พุทฺโธ อติอคฺคตาย\csromansharedfootnote{eabbab33d65320ec}{อติวิสิฏฺฐตาย (สฺยา)} อนุปโม. เอวเมว โข มหาราช อติทานํ นาม โลเก วิทูหิ วณฺณิตํ ถุตํ ปสตฺถํ, เย เกจิ ยาทิสํ กีทิสํ ทานํ เทนฺติ, อติทานทายี โลเก กิตฺติํ ปาปุณาติ, อติทาเนน เวสฺสนฺตโร ราชา ทสสหสฺสิยา โลกธาตุยา วณฺณิโต ถุโต ปสตฺโถ มหิโต กิตฺติโต, เตเนว อติทาเนน เวสฺสนฺตโร ราชา อชฺเชตรหิ พุทฺโธ ชาโต อคฺโค สเทวเก โลเก.}

\proseitem{5}{อตฺถิ ปน มหาราช โลเก ฐปนียํ ทานํ, ยํ ทกฺขิเณยฺเย อนุปฺปตฺเต น ทาตพฺพนฺติ. ทส โข ปนิมานิ ภนฺเต นาคเสน ทานานิ, ยานิ โลเก อทานสมฺมตานิ, โย ตานิ ทานานิ เทติ, โส อปายคามี โหติ. กตมานิ ทส, มชฺชทานํ ภนฺเต นาคเสน โลเก อทานสมฺมตํ, โย ตํ ทานํ เทติ, โส อปายคามี โหติ. สมชฺชทานํ ฯเปฯ อิตฺถิทานํ. อุสภทานํ. จิตฺตกมฺม\-ทานํ. สตฺถทานํ. วิสทานํ. สงฺขลิก\-ทานํ. กุกฺกุฏสูกร\-ทานํ. ตุลากูฏ\-มานกูฏ\-ทานํ ภนฺเต นาคเสน โลเก อทานสมฺมตํ โหติ, โย ตํ ทานํ เทติ, โส อปายคามี โหติ. อิมานิ โข ภนฺเต นาคเสน ทส ทานานิ โลเก อทานสมฺมตานิ, โย ตานิ ทานานิ เทติ, โส อปายคามี โหตีติ.}

\prose{นาหํ ตํ มหาราช อทานสมฺมตํ ปุจฺฉามิ, อิมํ ขฺวาหํ มหาราช ตํ ปุจฺฉามิ “อตฺถิ ปน มหาราช โลเก ฐปนียํ ทานํ, ยํ ทกฺขิเณยฺเย อนุปฺปตฺเต น ทาตพฺพนฺ”ติ. นตฺถิ ภนฺเต นาคเสน โลเก ฐปนียํ ทานํ, \csromanfolio{271}ยํ ทกฺขิเณยฺเย อนุปฺปตฺเต น ทาตพฺพํ, จิตฺตปฺปสาเท อุปฺปนฺเน เกจิ ทกฺขิเณยฺยานํ โภชนํ เทนฺติ, เกจิ อจฺฉาทนํ, เกจิ สยนํ, กิจิ อาวสถํ, เกจิ อตฺถรณ\-ปาวุรณํ, เกจิ ทาสิทาสํ, เกจิ เขตฺตวตฺถุํ, เกจิ ทฺวิปท\-จตุปฺปทํ, เกจิ สตํ สหสฺสํ สตสหสฺสํ, เกจิ มหารชฺชํ, เกจิ ชีวิตมฺปิ เทนฺตีติ. ยทิ ปน มหาราช เกจิ ชีวิตมฺปิ เทนฺติ, กิํ การณา เวสฺสนฺตรํ ทานปติํ อติพาฬฺหํ ปริปาเตสิ สุทินฺเน ปุตฺเต จ ทาเร จ.}

\prose{อปิ นุ โข มหาราช อตฺถิ โลกปกติ โลกาจิณฺณํ, ลภติ ปิตา ปุตฺตํ อิณฏฺโฏ วา อาชีวิกปกโต วา อาวปิตุํ วา\csromansharedfootnote{bf23c474a8e40433}{อาธาเปตุํ วา (สฺยา), อาธปิตุํ วา (ก)} วิกฺกิณิตุํ วาติ. อาม ภนฺเต ลภติ ปิตา ปุตฺตํ อิณฏฺโฏ วา อาชีวิกปกโต วา อาวปิตุํ วา วิกฺกิณิตุํ วาติ. ยทิ มหาราช ลภติ ปิตา ปุตฺตํ อิณฏฺโฏ วา อาชีวิกปกโต วา อาวปิตุํ วา วิกฺกิณิตุํ วา, เวสฺสนฺตโรปิ มหาราช ราชา อลภมาโน สพฺพญฺญุต\-ญาณํ อุปทฺทุโต ทุกฺขิโต ตสฺส ธมฺมธนสฺส ปฏิลาภาย ปุตฺตทารํ อาวเปสิ จ วิกฺกิณิ จ, อิติ มหาราช เวสฺสนฺตเรน รญฺญา อญฺเญสํ ทินฺนํ เยว ทินฺนํ, กตํ เยว กตํ. กิสฺส ปน ตฺวํ มหาราช เตน ทาเนน เวสฺสนฺตรํ ทานปติํ อติพาฬฺหํ อปสาเทสีติ.}

\prose{นาหํ ภนฺเต นาคเสน เวสฺสนฺตรสฺส ทานปติโน ทานํ ครหามิ, อปิ จ ปุตฺตทารํ ยาจนฺเต นิมินิตฺวา อตฺตานํ ทาตพฺพนฺติ. เอตํ โข มหาราช อสพฺภิการณํ, ยํ ปุตฺตทารํ ยาจนฺเต อตฺตานํ ทเทยฺย, ยํ ยํ หิ ยาจนฺเต ตํ ตเทว ทาตพฺพํ, เอตํ สปฺปุริสานํ กมฺมํ. ยถา มหาราช โกจิ ปุริโส ปานียํ อาหราเปยฺย, ตสฺส โย โภชนํ ทเทยฺย, อปิ นุ โส มหาราช ปุริโส ตสฺส กิจฺจการี อสฺสาติ. น หิ ภนฺเต, ยํ โส อาหราเปติ, ตเมว ตสฺส เทนฺโต กิจฺจการี อสฺสาติ. เอวเมว โข มหาราช เวสฺสนฺตโร ราชา พฺราหฺมเณ ปุตฺตทารํ ยาจนฺเต ปุตฺตทารํ เยว อทาสิ. สเจ มหาราช พฺราหฺมโณ เวสฺสนฺตรสฺส สรีรํ ยาเจยฺย, น โส มหาราช อตฺตานํ รกฺเขยฺย น กมฺเปยฺย น รชฺเชยฺย, ตสฺส ทินฺนํ ปริจฺจตฺตํ เยว สรีรํ ภเวยฺย. สเจ มหาราช โกจิ เวสฺสนฺตรํ ทานปติํ อุปคนฺตฺวา ยาเจยฺย “ทาสตฺตํ \csromanfolio{272}เม อุเปหี”ติ, ทินฺนํ ปริจฺจตฺตํ เยวสฺส สรีรํ ภเวยฺย, น โส ทตฺวา ตเปยฺย\csromansharedfootnote{30be235847da88de}{ฐเปยฺย (สี)}, รญฺโญ มหาราช เวสฺสนฺตรสฺส กาโย พหุสาธารโณ.}

\prose{ยถา มหาราช ปกฺกา มํสเปสิ พหุสาธารณา, เอวเมว โข มหาราช รญฺโญ เวสฺสนฺตรสฺส กาโย พหุสาธาราโณ. ยถา วา ปน มหาราช ผลิโต รุกฺโข นานา\-ทิชคณ\-สาธารโณ, เอวเมว โข มหาราช รญฺโญ เวสฺสนฺตรสฺส กาโย พหุสาธารโณ. กิํ การณา, เอวาหํ ปฏิปชฺชนฺโต สมฺมาสมฺโพธิํ ปาปุณิสฺสามีติ.}

\prose{ยถา มหาราช ปุริโส อธโน ธนตฺถิโก ธนปริเยสนํ จรมาโน อชปถํ สงฺกุปถํ เวตฺตปถํ คจฺฉติ, ชลถล\-วาณิชฺชํ กโรติ, กาเยน วาจาย มนสา ธนํ อาราเธติ, ธนปฺปฏิลาภาย วายมติ. เอวเมว โข มหาราช เวสฺสนฺตโร ทานปติ อธโน พุทฺธธเนน สพฺพญฺญุตญาณ\-รตน\-ปฺปฏิลาภาย ยาจกานํ ธนธญฺญํ ทาสิทาสํ ยานวาหนํ สกล\-สาปเตยฺยํ สกํ ปุตฺตทารํ อตฺตานญฺจ จชิตฺวา สมฺมาสมฺโพธิํ เยว ปริเยสติ.}

\prose{ยถา วา ปน มหาราช อมจฺโจ มุทฺทกาโม มุทฺทา\-ธิกรณํ ยํ กิญฺจิ เคเห ธนธญฺญํ หิรญฺญ\-สุวณฺณํ, ตํ สพฺพํ ทตฺวาปิ มุทฺท\-ปฺปฏิลาภาย วายมติ. เอวเมว โข มหาราช เวสฺสนฺตโร ทานปติ สพฺพํ ตํ พาหิร\-พฺภนฺตรธนํ ทตฺวา ชีวิตมฺปิ ปเรสํ ทตฺวา สมฺมาสมฺโพธิํ เยว ปริเยสติ.}

\prose{อปิ จ มหาราช เวสฺสนฺตรสฺส ทานปติโน เอวํ อโหสิ “ยํ โส พฺราหฺมโณ ยาจติ, ตเมวาหํ ตสฺส เทนฺโต กิจฺจการี นาม โหมี”ติ. เอวํ โส ตสฺส ปุตฺตทารม\-ทาสิ. น โข มหาราช เวสฺสนฺตโร ทานปติ เทสฺสตาย พฺราหฺมณสฺส ปุตฺตทารม\-ทาสิ, น อทสฺสน\-กามตาย ปุตฺตทารม\-ทาสิ, น อติพหุกา เม ปุตฺตทารา, “น สกฺโกมิ เต โปเสตุนฺ”ติ ปุตฺตทารม\-ทาสิ, น อุกฺกณฺฐิโต “อปฺปิยา เม”ติ นีหริตุกาม\-ตาย ปุตฺตทารม\-ทาสิ. อถ โข สพฺพญฺญุตญาณ\-รตนสฺเสว ปิยตฺตา สพฺพญฺญุต\-ญาณสฺส การณา เวสฺสนฺตโร ราชา เอวรูปํ อตุลํ วิปุลม\-นุตฺตรํ ปิยํ มนาปํ ทยิตํ ปาณสมํ ปุตฺตทาร\-ทานวรํ พฺราหฺมณสฺส อทาสิ.}

\prose{\csromanfolio{273}ภาสิตมฺเปตํ มหาราช ภควตา เทวาติเทวน \textbf{จริยาปิฏเก} *\csromandash{}}

\setlength{\csromangathapagewidth}{0pt}
\setlength{\csromangathawakwidth}{0pt}
\csromangathameasuregroup{“น เม เทสฺสา อุโภ ปุตฺตา, \\ สพฺพญฺญุตํ ปิยํ มยฺหํ,}{\csromangathabat{“น เม เทสฺสา อุโภ ปุตฺตา,}{มทฺที เทวี น เทสฺสิยา.} \\ \csromangathabat{สพฺพญฺญุตํ ปิยํ มยฺหํ,}{ตสฺมา ปิเย อทาสหนฺ”ติ.}}
\csromangathasetleft{“น เม เทสฺสา อุโภ ปุตฺตา, \\ สพฺพญฺญุตํ ปิยํ มยฺหํ,}
\csromangathagroup{\csromangathastanza{\csromangathabat{“น เม เทสฺสา อุโภ ปุตฺตา,}{มทฺที เทวี น เทสฺสิยา.} \\ \csromangathabat{สพฺพญฺญุตํ ปิยํ มยฺหํ,}{ตสฺมา ปิเย อทาสหนฺ”ติ.}}}

\prose{ตสฺมา มหาราช เวสฺสนฺตโร ราชา ปุตฺตทานํ\csromansharedfootnote{d11bdc7916da7d08}{ปุตฺตทารํ (ก)} ทตฺวา ปณฺณสาลํ ปวิสิตฺวา นิปชฺชิ. ตสฺส อติเปเมน ทุกฺขิตสฺส พลวโสโก อุปฺปชฺชิ, หทยวตฺถุ อุณฺหมโหสิ. นาสิกาย อปฺปโหนฺติยา มุเขน อุเณฺห อสฺสาสปสฺสาเส วิสฺสชฺเชสิ, อสฺสูนิ ปริวตฺติตฺวา โลหิตพินฺทูนิ หุตฺวา เนตฺเตหิ นิกฺขมิํสุ. เอวเมว โข มหาราช ทุกฺเขน เวสฺสนฺตโร ราชา พฺราหฺมณสฺส ปุตฺตทารม\-ทาสิ “มา เม ทานปโถ ปริหายี”ติ.}

\prose{อปิ จ มหาราช เวสฺสนฺตโร ราชา เทฺว อตฺถวเส ปฏิจฺจ พฺราหฺมณสฺส เทฺว ทารเก อทาสิ. กตเม เทฺว, ทานปโถ จ เม อปริหีโน ภวิสฺสติ, ทุกฺขิเต จ เม ปุตฺตเก วนมูลผเลหิ อิโตนิทานํ อยฺยโก โมเจสฺสตีติ. ชานาติ หิ มหาราช เวสฺสนฺตโร ราชา “น เม ทารกา สกฺกา เกนจิ ทาสโภเคน ภุญฺชิตุํ, อิเม จ ทารเก อยฺยโก นิกฺกิณิสฺสติ, เอวํ อมฺหากมฺปิ คมนํ ภวิสฺสตี”ติ. อิเม โข มหาราช เทฺว อตฺถวเส ปฏิจฺจ พฺราหฺมณสฺส เทฺว ทารเก อทาสิ.}

\prose{อปิ จ มหาราช เวสฺสนฺตโร ราชา ชานาติ “อยํ โข พฺราหฺมโณ ชิณฺโณ วุฑฺโฒ มหลฺลโก ทุพฺพโล ภคฺโค ทณฺฑปรายโณ ขีณายุโก ปริตฺตปุญฺโญ, เนโส สมตฺโถ อิเม ทารเก ทาสโภเคน ภุญฺชิตุนฺ”ติ. สกฺกุเณยฺย ปน มหาราช ปุริโส ปากติเกน พเลน อิเม จนฺทิมสูริเย เอวํมหิทฺธิเก เอวํ\-มหา\-นุภาเว คเหตฺวา เปฬาย วา สมุคฺเค วา ปกฺขิปิตฺวา นิปฺปเภ กตฺวา ถาลก\-ปริโภเคน\csromansharedfootnote{a0dd97f25429f5e8}{ปทีปปริโภเคน (สฺยา)} ปริภุญฺชิ\-ตุนฺติ. น หิ ภนฺเตติ. เอวเมว โข มหาราช อิมสฺมิํ โลเก จนฺทิมสูริย\-ปฺปฏิภาคสฺส เวสฺสนฺตรสฺส ทารกา น สกฺกา เกนจิ ทาสโภเคน ภุญฺชิตุนฺติ.}

\prose{อปรมฺปิ มหาราช อุตฺตริํ การณํ สุโณหิ, เยน การเณน เวสฺสนฺตรสฺส ทารกา น สกฺกา เกนจิ ทาสโภเคน ภุญฺชิตุํ, ยถา มหาราช รญฺโญ จกฺกวตฺติสฺส มณิรตนํ สุภํ ชาติมนฺตํ อฏฺฐํสํ สุปริกมฺมกตํ จตุหตฺถา\-ยามํ สกฏ\-นาภิ\-ปริณาหํ น สกฺกา เกนจิ}

\prose{\csromanfolio{274}ปิโลติกาย เวเฐตฺวา เปฬาย ปกฺขิปิตฺวา สตฺถก\-นิสา\-น\-ปริโภเคน ปริภุญฺชิตุํ. เอวเมว โข มหาราช โลเก จกฺกวตฺติรญฺโญ มณิรตน\-ปฺปฏิภาคสฺส เวสฺสนฺตรสฺส ทารกา น สกฺกา เกนจิ ทาสโภเคน ภุญฺชิตุํ.}

\prose{อปรมฺปิ มหาราช อุตฺตริํ การณํ สุโณหิ, เยน การเณน เวสฺสนฺตรสฺส ทารกา น สกฺกา เกนจิ ทาสโภเคน ภุญฺชิตุํ, ยถา มหาราช ติธา ปภินฺโน สพฺพเสโต สตฺต\-ปฺปติฏฺฐิโต อฏฺฐรตนุ\-พฺเพโธ นวรต\-นายาม\-ปริณาโห ปาสาทิโก ทสฺสนีโย อุโปสโถ นาคราชา นสกฺกา เกนจิ สุปฺเปน วา สราเวน วา ปิทหิตุํ, โควจฺฉโก วิย วจฺฉกสาลาย ปกฺขิปิตฺวา ปริหริตุํ วา. เอวเมว โข มหาราช โลเก อุโปสถ\-นาคราช\-ปฺปฏิภาคสฺส เวสฺสนฺตรสฺส ทารกา น สกฺกา เกนจิ ทาสโภเคน ภุญฺชิตุํ.}

\prose{อปรมฺปิ มหาราช อุตฺตริํ การณํ สุโณหิ, เยน การเณน เวสฺสนฺตรสฺส ทารกา น สกฺกา เกนจิ ทาสโภเคน ภุญฺชิตุํ, ยถา มหาราช มหาสมุทฺโท ทีฆปุถุล\-วิตฺถิณฺโณ คมฺภีโร อปฺปเมยฺโย ทุรุตฺตโร อปริโยคาโฬฺห อนาวโฏ น สกฺกา เกนจิ สพฺพตฺถ ปิทหิตฺวา เอกติตฺเถน ปริโภคํ กาตุํ, เอวเมว โข มหาราช โลเก มหาสมุทฺท\-ปฺปฏิภาคสฺส เวสฺสนฺตรสฺส ทารกา น สกฺกา เกนจิ ทาสโภเคน ภุญฺชิตุํ.}

\prose{อปรมฺปิ มหาราช อุตฺตริํ การณํ สุโณหิ, เยน การเณน เวสฺสนฺตรสฺส ทารกา น สกฺกา เกนจิ ทาสโภเคน ภุญฺชิตุํ, ยถา มหาราช หิมวนฺโต ปพฺพตราชา ปญฺจ\-โยชนสตํ อจฺจุคฺคโต นเภ ติสหสฺส\-โยชนา\-ยาม\-วิตฺถาโร จตุราสีติ\-กูฏ\-สหสฺส\-ปฺปฏิมณฺฑิโต ปญฺจนฺนํ มหานที\-สตานํ ปภโว มหาภูต\-คณา\-ลโย นานาวิธ\-คนฺธ\-ธโร ทิพฺโพ\-สธ\-สต\-สมลงฺกโต นเภ วลาหโก วิย อจฺจุคฺคโต ทิสฺสติ. เอวเมว โข มหาราช โลเก หิมวนฺต\-ปพฺพตราช\-ปฺปฏิภาคสฺส เวสฺสนฺตรสฺส ทารกา น สกฺกา เกนจิ ทาสโภเคน ภุญฺชิตุํ.}

\prose{อปรมฺปิ มหาราช อุตฺตริํ การณํ สุโณหิ, เยน การเณน เวสฺสนฺตรสฺส ทารกา นสกฺกา เกนจิ ทาสโภเคน ภุญฺชิตุํ, ยถา มหาราช รตฺต\-นฺธการ\-ติมิสายํ อุปริปพฺพต\-คฺเค ชลมาโน มหา- อคฺคิกฺขนฺโธ สุวิทูเรปิ ปญฺญายติ, เอวเมว โข มหาราช เวสฺสนฺตโร ราชา ปพฺพตคฺเค}

\proseitem{5}{\csromanfolio{275}ชลมาโน มหา\-อคฺคิกฺขนฺโธ วิย สุวิทูเรปิ ปากโฏ ปญฺญายติ, ตสฺส ทารกา น สกฺกา เกนจิ ทาสโภเคน ภุญฺชิตุํ.}

\prose{อปรมฺปิ มหาราช อุตฺตริํ การณํ สุโณหิ, เยน การเณน เวสฺสนฺตรสฺส ทารกา น สกฺกา เกนจิ ทาสโภเคน ภุญฺชิตุํ, ยถา มหาราช หิมวนฺเต ปพฺพเต นาค\-ปุปฺผ\-สมเย อุชุวาเต วายนฺเต ทส ทฺวาทส โยชนานิ ปุปฺผคนฺโธ วายติ, เอวเมว โข มหาราช เวสฺสนฺตรสฺส รญฺโญ อปิ โยชน\-สหสฺเสหิปิ ยาว อกนิฏฺฐ\-ภวนํ เอตฺถนฺตเร สุราสุรครุฬคนฺธพฺพยกฺขสมโหรค กินฺนร- อินฺทภวเนสุ กิตฺติสทฺโท อพฺพุคฺคโต, สีลวรคนฺโธ จสฺส สมฺปวายติ, เตน ตสฺส ทารกา น สกฺกา เกนจิ ทาสโภเคน ภุญฺชิตุํ. อนุสิฏฺโฐ มหาราช ชาลี กุมาโร ปิตรา เวสฺสนฺตเรน รญฺญา “อยฺยโก เต ตาต ตุเมฺห พฺราหฺมณสฺส ธนํ ทตฺวา นิกฺกิณนฺโต ตํ นิกฺขสหสฺสํ ทตฺวา นิกฺกิณาตุ, กณฺหาชินํ นิกฺกิณนฺโต ทาสสตํ ทาสิสตํ หตฺถิสตํ อสฺสสตํ เธนุสตํ อุสภสตํ นิกฺขสตนฺติ สพฺพสตํ ทตฺวา นิกฺกิณาตุ, ยทิ เต ตาต อยฺยโก ตุเมฺห พฺราหฺมณสฺส หตฺถโต อาณาย พลสา มุธา คณฺหาติ, มา ตุเมฺห อยฺยกสฺส วจนํ กริตฺถ, พฺราหฺมณสฺเสว อนุยายิโน โหถา”ติ, เอวมนุสาสิตฺวา ปุตฺเต เปเสสิ, ตโต ชาลีกุมาโร คนฺตฺวา อยฺยเกน ปุฏฺโฐ กเถสิ\csromandash{}}

\setlength{\csromangathapagewidth}{0pt}
\setlength{\csromangathawakwidth}{0pt}
\csromangathameasuregroup{\textsuperscript{*}“สหสฺสคฺฆํ หิ มํ ตาต, \\ อโถ กณฺหาชินํ กญฺญํ,}{\csromangathabat{\textsuperscript{*}“สหสฺสคฺฆํ หิ มํ ตาต,}{พฺราหฺมณสฺส ปิตา อทา.} \\ \csromangathabat{อโถ กณฺหาชินํ กญฺญํ,}{หตฺถีนญฺจ สเตน จา”ติ.}}
\csromangathasetleft{\textsuperscript{*}“สหสฺสคฺฆํ หิ มํ ตาต, \\ อโถ กณฺหาชินํ กญฺญํ,}
\csromangathagroup{\csromangathastanza{\csromangathabat{\csromansymbolfootnote{*}{ขุ 6. 368 ปิฏฺเฐ.}“สหสฺสคฺฆํ หิ มํ ตาต,}{พฺราหฺมณสฺส ปิตา อทา.} \\ \csromangathabat{อโถ กณฺหาชินํ กญฺญํ,}{หตฺถีนญฺจ สเตน จา”ติ.}}}

\prosewithcloserruled{สุนิพฺเพฐิโต ภนฺเต นาคเสน ปโญฺห, สุภินฺนํ ทิฏฺฐิชาลํ, สุมทฺทิโต ปรวาโท, สกสมโย สุทีปิโต, พฺยญฺชนํ สุปริโสธิตํ, สุวิภตฺโต อตฺโถ, เอวเมตํ ตถา สมฺปฏิจฺฉามีติ.}{เวสฺสนฺตร\-ปโญฺห ปฐโม.}
\csromanmatikaanchor{mtk.186}
\csromantocmarkref{subsubsection}{2. ทุกฺกรการิกปญฺห}{mtk.186}
\bookmark[dest={mtk.186},level=3]{2. ทุกฺกรการิกปญฺห}
\csromanheader{3}{2. \textbf{ทุกฺกรการิก}\-\textbf{ปญฺห}}
\proseitem{5}{ภนฺเต นาคเสน สพฺเพว โพธิสตฺตา ทุกฺกรการิกํ กโรนฺติ, อุทาหุ โคตเมเนว โพธิสตฺเตน ทุกฺกรการิกา กตาติ. นตฺถิ \csromanfolio{276}มหาราช สพฺเพสํ โพธิสตฺตานํ ทุกฺกรการิกา, โคตเมเนว โพธิสตฺเตน ทุกฺกรการิกา กตาติ.}

\proseitem{2}{ภนฺเต นาคเสน ยทิ เอวํ อยุตฺตํ, ยํ โพธิสตฺตานํ โพธิสตฺเตหิ เวมตฺตตา โหตีติ. จตูหิ มหาราช ฐาเนหิ โพธิสตฺตานํ โพธิสตฺเตหิ เวมตฺตตา โหติ. กตเมหิ จตูหิ, กุลเวมตฺตตา ปธาน\-เวมตฺตตา\csromansharedfootnote{916c939efc389823}{อทฺธานเวมตฺตตา (สี, สฺยา, อิ)} อายุเวมตฺตตา ปมาณ\-เวมตฺตตาติ. อิเมหิ โข มหาราช จตูหิ ฐาเนหิ โพธิสตฺตานํ โพธิสตฺเตหิ เวมตฺตตา โหติ. สพฺเพสมฺปิ มหาราช พุทฺธานํ รูเป สีเล สมาธิมฺหิ ปญฺญาย วิมุตฺติยา วิมุตฺติ\-ญาณ\-ทสฺสเน จตุเวสารชฺเช ทสตถาคตพเล ฉอสาธารณ\-ญาเณ จุทฺทส\-พุทฺธญาเณ อฏฺฐารส\-พุทฺธธมฺเม เกวเล จ พุทฺธคุเณ\csromansharedfootnote{0bf3dd7d50aab561}{พุทฺธธมฺเม (สี, อิ)} นตฺถิ เวมตฺตตา, สพฺเพปิ พุทฺธา พุทฺธธมฺเมหิ สมสมาติ.}

\prose{ยทิ ภนฺเต นาคเสน สพฺเพปิ พุทฺธา พุทฺธธมฺเมหิ สมสมา, เกน การเณน โคตเมเนว โพธิสตฺเตน ทุกฺกรการิกา กตาติ. อปริปกฺเก มหาราช ญาเณ อปริปกฺกาย โพธิยา โคตโม โพธิสตฺโต เนกฺขมฺมม\-ภินิกฺขนฺโต อปริปกฺกํ ญาณํ ปริปาจ\-ยมา\-เนน ทุกฺกรการิกา กตาติ.}

\prose{ภนฺเต นาคเสน เกน การเณน โพธิสตฺโต อปริปกฺเก ญาเณ อปริปกฺกาย โพธิยา มหา\-ภินิกฺขมนํ นิกฺขนฺโต, นนุ นาม ญาณํ ปริปาเจตฺวา ปริปกฺเก ญาเณ นิกฺขมิ\-ตพฺพนฺติ.}

\prose{โพธิสตฺโต มหาราช วิปรีตํ อิตฺถาคารํ ทิสฺวา วิปฺปฏิสารี อโหสิ, ตสฺส วิปฺปฏิสาริสฺส อรติ อุปฺปชฺชิ, อรติจิตฺตํ อุปฺปนฺนํ ทิสฺวา อญฺญตโร มารกายิโก เทวปุตฺโต “อยํ โข กาโล อรติจิตฺตสฺส วิโนทนายา”ติ เวหาเส ฐตฺวา อิทํ วจนมพฺรวิ\csromandash{}}

\prose{“มาริส มา โข ตฺวํ อุกฺกณฺฐิโต อโหสิ, อิโต เต สตฺตเม ทิวเส ทิพฺพํ จกฺกรตนํ ปาตุภวิสฺสติ สหสฺสารํ สเนมิกํ สนาภิกํ สพฺพา\-การ\-ปริปูรํ, ปถวิคตานิ จ เต รตนานิ อากาสฏฺฐานิ จ สยเมว อุปคจฺฉิสฺสนฺติ, ทฺวิสหสฺส\-ปริตฺต\-ทีป\-ปริวาเรสุ จตูสุ มหาทีเปสุ เอกมุเขน อาณา ปวตฺติสฺสติ, ปโรสหสฺสญฺจ เต ปุตฺตา ภวิสฺสนฺติ สูรา วีรงฺครูปา \csromanfolio{277}ปรเสน\-ปฺปมทฺทนา, เตหิ ปุตฺเตหิ ปริกิณฺโณ สตฺต\-รตน\-สมนฺนาคโต จตุทฺทีปมนุสาสิสฺสสี”ติ.}

\prose{ยถา นาม ทิวส\-สนฺตตฺตํ อโยสูลํ สพฺพตฺถ อุปฑหนฺตํ กณฺณโสตํ ปวิเสยฺย, เอวเมว โข มหาราช โพธิสตฺตสฺส ตํ วจนํ กณฺณโสตํ ปวิสิตฺถ, อิติ โส ปกติยาว อุกฺกณฺฐิโต ตสฺสา เทวตาย วจเนน ภิยฺโยโส\-มตฺตาย\csromansharedfootnote{ab67f8da3c3dc98c}{ภิยฺโยโส มตฺตาย เทฺว ปทานิ \csromandash{} ม.พ.ป.} อุพฺพิชฺชิ สํวิชฺชิ สํเวคมาปชฺชิ.}

\prose{ยถา ปน มหาราช มหติ\-มหา\-อคฺคิกฺขนฺโธ ชลมาโน อญฺเญน กฏฺเฐน อุปฑหิโต ภิยฺโยโส\-มตฺตาย\csromansharedfootnote{ab67f8da3c3dc98c}{ภิยฺโยโส มตฺตาย เทฺว ปทานิ \csromandash{} ม.พ.ป.} ชเลยฺย, เอวเมว โข มหาราช โพธิสตฺโต ปกติยาว อุกฺกณฺฐิโต ตสฺสา เทวตาย วเจเนน ภิยฺโยโส\-มตฺตาย\csromansharedfootnote{ab67f8da3c3dc98c}{ภิยฺโยโส มตฺตาย เทฺว ปทานิ \csromandash{} ม.พ.ป.} อุพฺพิชฺชิ สํวิชฺชิ สํเวคมาปชฺชิ.}

\prose{ยถา วา ปน มหาราช มหาปถวี ปกติตินฺตา นิพฺพตฺต\-หริต\-สทฺทลา อาสิตฺโตทกา จิกฺขลฺลชาตา ปุนเทว มหาเมเฆ อภิวุฏฺเฐ ภิยฺโยโส\-มตฺตาย\csromansharedfootnote{ab67f8da3c3dc98c}{ภิยฺโยโส มตฺตาย เทฺว ปทานิ \csromandash{} ม.พ.ป.} จิกฺขลฺลตรา อสฺส, เอวเมว โข มหาราช โพธิสตฺโต ปกติยาว อุกฺกณฺฐิโต ตสฺสา เทวตาย วจเนน ภิยฺโยโส\-มตฺตาย\csromansharedfootnote{ab67f8da3c3dc98c}{ภิยฺโยโส มตฺตาย เทฺว ปทานิ \csromandash{} ม.พ.ป.} อุพฺพิชฺชิ สํวิชฺชิ สํเวคมาปชฺชีติ.}

\prose{อปิ นุ โข ภนฺเต นาคเสน โพธิสตฺตสฺส ยทิ สตฺตเม ทิวเส ทิพฺพํ จกฺกรตนํ นิพฺพตฺเตยฺย, ปฏินิวตฺเตยฺย โพธิสตฺโต ทิพฺเพ จกฺกรตเน นิพฺพตฺเตติ. น หิ มหาราช สตฺตเม ทิวเส โพธิสตฺตสฺส ทิพฺพํ จกฺกรตนํ นิพฺพตฺเตยฺย, อปิ จ ปโลภน\-ตฺถาย ตาย เทวตาย มุสา ภณิตํ, ยทิปิ มหาราช สตฺตเม ทิวเส ทิพฺพํ จกฺกรตนํ นิพฺพตฺเตยฺย, โพธิสตฺโต น นิวตฺเตยฺย, กิํ การณํ “อนิจฺจนฺ”ติ มหาราช โพธิสตฺโต ทฬฺหํ อคฺคเหสิ, “ทุกฺขํ อนตฺตี”ติ ทฬฺหํ อคฺคเหสิ, อุปาทานกฺขยํ ปตฺโต.}

\prose{ยถา มหาราช อโนตตฺตทหโต อุทกํ คงฺคํ นทิํ ปวิสติ, คงฺคาย นทิยา มหาสมุทฺทํ ปวิสติ, มหาสมุทฺทโต ปาตาลมุขํ ปวิสติ, อปิ นุ มหาราช ตํ อุทกํ ปาตาล\-มุขคตํ ปฏินิวตฺติตฺวา มหาสมุทฺทํ ปวิเสยฺย, มหาสมุทฺทโต คงฺคํ นทิํ ปวิเสยฺย, คงฺคาย นทิยา ปุน อโนตตฺตํ ปวิเสยฺยาติ. น หิ ภนฺเตติ. เอวเมว โข มหาราช โพธิสตฺเตน กปฺปานํ สตสหสฺสํ จตุโร จ อสงฺเขฺยยฺเย กุสลํ ปริปาจิตํ อิมสฺส ภวสฺส การณา, โสยํ อนฺติมภโว \csromanfolio{278}อนุปฺปตฺโต ปริปกฺกํ โพธิญาณํ ฉหิ วสฺเสหิ พุทฺโธ ภวิสฺสติ สพฺพญฺญู โลเก อคฺคปุคฺคโล, อปิ นุ โข มหาราชา โพธิสตฺโต จกฺกรตน\-การณา\csromansharedfootnote{4a4e14c1d465ccaa}{จกฺกรตนสฺส การณา (สี, สฺยา, อิ)} ปฏินิวตฺเตยฺยาติ\csromansharedfootnote{a398a2460ba875a1}{ปรินิวตฺเตยฺยาติ (สี, อิ, ก)}. น หิ ภนฺเตติ.}

\proseitem{5}{อปิ จ มหาราช มหาปถวี ปริวตฺเตยฺย สกานนา สปพฺพตา, นเตฺวว โพธิสกฺโก ปฏินิวตฺเตยฺย อปตฺวา สมฺมาสมฺโพธิํ. อาโรเหยฺยปิ เจ มหาราช คงฺคาย อุทกํ ปฏิโสตํ, นเตฺวว โพธิสตฺโต ปฏินิวตฺเตยฺย อปตฺวา สมฺมาสมฺโพธิํ, วิสุสฺเสยฺยปิ เจ มหาราช มหาสมุทฺโท อปริมิต\-ชลธโร โคปเท อุทกํ วิย, นเตฺวว โพธิสตฺโต ปฏินิวตฺเตยฺย อปตฺวา สมฺมาสมฺโพธิํ. ผเลยฺยปิ เจ มหาราช สิเนรุ\-ปพฺพตราชา สตธา วา สหสฺสธา วา, นเตฺวว โพธิสตฺโต ปฏินิวตฺเตยฺย อปตฺวา สมฺมาสมฺโพธิํ. ปเตยฺยุมฺปิ เจ มหาราช จนฺทิมสูริยา สตารกา เลฑฺฑุ วิย ฉมายํ, นเตฺวว โพธิสตฺโต ปฏินิวตฺเตยฺย อปตฺวา สมฺมาสมฺโพธิํ. สํวตฺเตยฺยปิ เจ มหาราช อากาโส กิลญฺชมิว, นเตฺวว โพธิสตฺโต ปฏินิวตฺเตยฺย อปตฺวา สมฺมาสมฺโพธิํ, กิํ การณา, ปทาลิตตฺตา สพฺพพนฺธนานนฺติ.}

\prose{ภนฺเต นาคเสน กติ โลเก พนฺธนานีติ. ทส โข ปนิมานิ มหาราช โลเก พนฺธนานิ, เยหิ พนฺธเนหิ พทฺธา สตฺตา น นิกฺขมนฺติ นิกฺขมิตฺวาปิ ปฏินิวตฺตนฺติ. กตมานิ ทส, มาตา มหาราช โลเก พนฺธนํ, ปิตา มหาราช โลเก พนฺธนํ, ภริยา มหาราช โลเก พนฺธนํ, ปุตฺตา มหาราช โลเก พนฺธนํ, ญาตี มหาราช โลเก พนฺธนํ, มิตฺตํ มหาราช โลเก พนฺธนํ, ธนํ มหาราช โลเก พนฺธนํ, ลาภสกฺกาโร มหาราช โลเก พนฺธนํ, อิสฺสริยํ มหาราช โลเก พนฺธนํ, ปญฺจ กามคุณา มหาราช โลเก พนฺธนํ, อิมานิ โข มหาราช ทส โลเก พนฺธนานิ, เยหิ พนฺธเนหิ พทฺธา สตฺตา น นิกฺขมนฺติ นิกฺขมิตฺวาปิ ปฏินิวตฺตนฺติ, ตานิ ทส พนฺธนานิ โพธิสตฺตสฺส ฉินฺนานิ ปทาลิตานิ, ตสฺมา มหาราช โพธิสตฺโต น ปฏินิวตฺตตีติ.}

\prose{ภนฺเต นาคเสน ยทิ โพธิสตฺโต อุปฺปนฺเน อรติจิตฺเต เทวตาย วจเนน อปริปกฺเก ญาเณ อปริปกฺกาย โพธิยา เนกฺขมฺมม\-ภินิกฺขนฺโต, \csromanfolio{279}กิํ ตสฺส ทุกฺกร\-การิกาย กตาย, นนุ นาม สพฺพภกฺเขน ภวิตพฺพํ ญาณปริปากํ อาคมยมาเนนาติ.}

\prose{ทส โข ปนิเม มหาราช ปุคฺคลา โลกสฺมิํ โอญฺญาตา อวญฺญาตา หีฬิตา ขีฬิตา ครหิตา ปริภูตา อจิตฺตีกตา. กตเม ทส, อิตฺถี มหาราช วิธวา โลกสฺมิํ โอญฺญาตา อวญฺญาตา หีฬิตา ขีฬิตา ครหิตา ปริภูตา อจิตฺตีกตา. ทุพฺพโล มหาราช ปุคฺคโล ฯเปฯ. อมิตฺตญาติ มหาราช ปุคฺคโล. มหคฺฆโส มหาราช ปุคฺคโล. อครุกุลวาสิโก มหาราช ปุคฺคโล. ปาปมิตฺโต มหาราช ปุคฺคโล. ธนหีโน มหาราช ปุคฺคโล. อาจารหีโน มหาราช ปุคฺคโล. กมฺมหีโน มหาราช ปุคฺคโล. ปโยคหีโน มหาราช ปุคฺคโล โลกสฺมิํ โอญฺญาโต อวญฺญาโต หีฬิโต ขีฬิโต ครหิโต ปริภูโต อจิตฺตีกโต. อิเม โข มหาราช ทส ปุคฺคลา โลกสฺมิํ โอญฺญาตา อวญฺญาตา หีฬิตา ขีฬิตา ครหิตา ปริภูตา อจิตฺตีกตา. อิมานิ โข มหาราช ทส ฐานานิ อนุสฺสร\-มานสฺส โพธิสตฺตสฺส เอวํ สญฺญา อุปฺปชฺชิ “มาหํ กมฺมหีโน อสฺสํ ปโยคหีโน ครหีโต เทวมนุสฺสานํ, ยนฺนูนาหํ กมฺมสฺสามี อสฺสํ กมฺมครุ กมฺมา\-ธิปเตยฺโย กมฺมสีโล กมฺมโธรโยฺห กมฺมนิเกตวา อปฺปมตฺโต วิหเรยฺยนฺ”ติ, เอวํ โข มหาราช โพธิสตฺโต ญาณํ ปริปาเจนฺโต ทุกฺกรการิกํ อกาสีติ.}

\prose{ภนฺเต นาคเสน โพธิสตฺโต ทุกฺกรการิกํ กโรนฺโต เอวมาห “น โข ปนาหํ อิมาย กฏุกาย ทุกฺกร\-การิกาย อธิคจฺฉามิ อุตฺตริ\-มนุสฺสธมฺมํ อลม\-ริย\-ญาณ\-ทสฺสน\-วิเสสํ, สิยา นุ โข อญฺโญ มคฺโค โพธายาติ. อปิ นุ ตสฺมิํ สมเย โพธิสตฺตสฺส มคฺคํ อารพฺภ สติสมฺโมโส อโหสี”ติ.}

\prose{ปญฺจวีสติ โข ปนิเม มหาราช จิตฺต\-ทุพฺพลีกรณา ธมฺมา, เยหิ ทุพฺพลีกตํ จิตฺตํ น สมฺมา สมาธิยติ อาสวานํ ขยาย. กตเม ปญฺจวีสติ, โกโธ มหาราช จิตฺต\-ทุพฺพลีกรโณ ธมฺโม, เยน ทุพฺพลีกตํ จิตฺตํ น สมฺมา สมาธิยติ อาสวานํ ขยาย, อุปนาโห ฯเปฯ. มกฺโข. ปฬาโส. อิสฺสา. มจฺฉริยํ. มายา. สาเฐยฺยํ. ถมฺโภ. สารมฺโภ. มาโน. \csromanfolio{280}อติมาโน. มโท. ปมาโท. ถินมิทฺธํ. ตนฺทิ\csromansharedfootnote{7bea62046ab201d8}{นนฺที (อิ, ก)}. อาลสฺยํ. ปาปมิตฺตตา. รูปา. สทฺทา. คนฺธา. รสา. โผฏฺฐพฺพา. ขุทาปิปาสา. อรติ มหาราช จิตฺต\-ทุพฺพลีกรโณ ธมฺโม, เยน ทุพฺพลีกตํ จิตฺตํ น สมฺมา สมาธิยติ อาสวานํ ขยาย, อิเม โข มหาราช ปญฺจวีสติ จิตฺต\-ทุพฺพลีกรณา ธมฺมา, เยหิ ทุพฺพลีกตํ จิตฺตํ น สมฺมา สมาธิยติ อาสวานํ ขยาย.}

\prosewithcloserruled{โพธิสตฺตสฺส โข มหาราช ขุทาปิปาสา\csromansharedfootnote{64daa288f0b2fa79}{ขุปฺปิปาสา (ก)} กายํ ปริยาทิยิํสุ, กาเย ปริยาทินฺเน จิตฺตํ น สมฺมา สมาธิยติ อาสวานํ ขยาย. สตสหสฺสํ มหาราช กปฺปานํ\csromansharedfootnote{02bc66dab254c879}{กปฺเป (ก)} จตุโร จ อสงฺเขฺยยฺเย กปฺเป โพธิสตฺโต จตุนฺนํ เยว อริยสจฺจานํ อภิสมยํ อเนฺวสิ ตาสุ ตาสุ ชาตีสุ, กิํ ปนสฺส ปจฺฉิเม ภเว อภิสมย\-ชาติยํ มคฺคํ อารพฺภ สติสมฺโมโส เหสฺสติ, อปิ จ มหาราช โพธิสตฺตสฺส สญฺญามตฺตํ อุปฺปชฺชิ “สิยา นุ โข อญฺโญ มคฺโค โพธายา”ติ. ปุพฺเพ โข มหาราช โพธสตฺโต เอกมาสิโก สมาโน ปิตุ สกฺกสฺส กมฺมนฺเต สีตาย ชมฺพุจฺฉายาย สิริสยเน ปลฺลงฺกํ อาภุชิตฺวา นิสินฺโน วิวิจฺเจว กาเมหิ วิวิจฺจ อกุสเลหิ ธมฺเมหิ สวิตกฺกํ สวิจารํ วิเวกชํ ปีติสุขํ ปฐมํ ฌานํ อุปสมฺปชฺช วิหาสิ ฯเปฯ จตุตฺถํ ฌานํ อุปสมฺปชฺช วิหาสีติ. สาธุ ภนฺเต นาคเสน, เอวเมตํ ตถา สมฺปฏิจฺฉามิ, ญาณํ ปริปาเจนฺโต โพธิสตฺโต ทุกฺกรการิกํ อกาสีติ.}{ทุกฺกรการิก\-ปโญฺห ทุติโย.}
\csromanmatikaanchor{mtk.187}
\csromantocmarkref{subsubsection}{3. กุสลากุสลพลวตรปญฺห}{mtk.187}
\bookmark[dest={mtk.187},level=3]{3. กุสลากุสลพลวตรปญฺห}
\csromanheader{3}{3. \textbf{กุสลากุสล}\-\textbf{พลวตร}\-\textbf{ปญฺห}}
\proseitem{3}{ภนฺเต นาคเสน กตมํ อธิมตฺตํ พลวตรํ กุสลํ วา อกุสลํ วาติ. กุสลํ มหาราช อธิมตฺตํ พลวตรํ, โน ตถา อกุสลนฺติ. นาหํ ภนฺเต นาคเสน ตํ วจนํ สมฺปฏิจฺฉามิ “กุสลํ อธิมตฺตํ พลวตรํ, โน ตถา อกุสลนฺ”ติ, ทิสฺสนฺติ ภนฺเต นาคเสน อิธ ปาณาติปาติโน อทินฺนาทายิโน กาเมสุ\-มิจฺฉาจาริโน มุสาวาทิโน คามฆาติกา ปนฺถทูสกา เนกติกา วญฺจนิกา, สพฺเพ เต ตาวตเกน ปาเปน ลภนฺติ หตฺถจฺเฉทํ ปาทจฺเฉทํ หตฺถปาท\-จฺเฉทํ กณฺณจฺเฉทํ นาสจฺเฉทํ \csromanfolio{281}กณฺณนาส\-จฺเฉทํ พิลงฺค\-ถาลิกํ สงฺขมุณฺฑิกํ ราหุมุขํ โชติมาลิกํ หตฺถ\-ปชฺโชติกํ เอรกวตฺติกํ จีรกวาสิกํ เอเณยฺยกํ พฬิสมํสิกํ กหาปณิกํ ขารา\-ปตจฺฉิกํ ปลิฆ\-ปริวตฺติกํ ปลาลปีฐกํ ตตฺเตนปิ เตเลน โอสิญฺจนํ สุนเขหิปิ ขาทาปนํ ชีวสูลา\-โรปนํ อสินาปิ สีสจฺเฉทํ, เกจิ รตฺติํ ปาปํ กตฺวา รตฺติํ เยว วิปากํ อนุภวนฺติ, เกจิ รตฺติํ กตฺวา ทิวา เยว อนุภวนฺติ, เกจิ ทิวา กตฺวา ทิวา เยว อนุภวนฺติ, เกจิ ทิวา กตฺวา รตฺติํ เยว อนุภวนฺติ, เกจิ เทฺว ตโย ทิวเส วีติวตฺเต อนุภวนฺติ, สพฺเพปิ เต ทิฏฺเฐว ธมฺเม วิปากํ อนุภวนฺติ. อตฺถิ ปน ภนฺเต นาคเสน โกจิ เอกสฺส วา ทฺวินฺนํ วา ติณฺณํ วา จตุนฺนํ วา ปญฺจนฺนํ วา ทสนฺนํ วา สตสฺส วา สหสฺสสฺส วา สตสหสฺสสฺส วา สปริวารํ ทานํ ทตฺวา ทิฏฺฐ\-ธมฺมิกํ โภคํ วา ยสํ วา สุขํ วา อนุภวิตา สีเลน วา อุโปสถ\-กมฺเมน วาติ.}

\prose{อตฺถิ มหาราช จตฺตาโร ปุริสา ทานํ ทตฺวา สีลํ สมาทิยิตฺวา อุโปสถกมฺมํ กตฺวา ทิฏฺเฐว ธมฺเม เตเนว สรีรเทเหน ติทสปุเร สมนุปฺปตฺตาติ\csromansharedfootnote{f95d0bad603f3648}{ยสมนุปตฺตาติ (สี, อิ)}. โก จ โก จ ภนฺเตติ. มนฺธาตา มหาราช ราชา นิมิ ราชา สาธีโน ราชา คุตฺติโล จ คนฺธพฺโพติ.}

\prose{ภนฺเต นาคเสน อเนเกหิ ตํ ภวสหสฺเสหิ อนฺตริตํ, ทฺวินฺนมฺเปตํ อมฺหากํ\csromansharedfootnote{47e9dda68a9d46b3}{ทีปิตํ, อมฺหากมฺเปตํ (ก)} ปโรกฺขํ, ยทิ สมตฺโถสิ วตฺตมานเก ภเว ภควโต ธรมานกาเล กเถหีติ. วตฺตมานเกปิ มหาราช ภเว ปุณฺณโก ทาโส เถรสฺส สาริปุตฺตสฺส โภชนํ ทตฺวา ตทเหว เสฏฺฐิฏฺฐานํ อชฺฌุปคโต, โส เอตรหิ ปุณฺณโก เสฏฺฐีติ ปญฺญายิ, โคปาลมาตา เทวี อตฺตโน เกเส วิกฺกิณิตฺวา ลทฺเธหิ อฏฺฐหิ กหาปเณหิ เถรสฺส มหา\-กจฺจายนสฺส อตฺต\-ฏฺฐมกสฺส ปิณฺฑปาตํ ทตฺวา ตทเหว รญฺโญ จนฺท\-ปชฺโชตสฺส\csromansharedfootnote{4801fa11412db5ef}{อุเทนสฺส (สี, อิ)} อคฺคมเหสิ\-ฏฺฐานํ ปตฺตา. สุปฺปิยา อุปาสิกา อญฺญตรสฺส คิลาน\-ภิกฺขุโน อตฺตโน อูรุมํเสน ปฏิจฺฉาทนียํ ทตฺวา ทุติยทิวเส เยว รูฬฺหวณา สญฺฉวี\csromansharedfootnote{6388daa8873b9eb1}{สจฺฉวี (สี, อิ)} อโรคา ชาตา. มิลฺลิกา เทวี ภควโต อาภิโทสิกํ กุมฺมาสปิณฺฑํ ทตฺวา ตทเหว รญฺโญ โกสลสฺส อคฺคมเหสี ชาตา. สุมโน มาลากาโร อฏฺฐหิ สุมน\-ปุปฺผมุฏฺฐีหิ ภควนฺตํ ปูเชตฺวา ตํ ทิวสํ เยว มหาสมฺปตฺติํ ปตฺโต. \csromanfolio{282}เอกสาฏโก พฺราหฺมโณ อุตฺตรสาฏเกน ภควนฺตํ ปูเชตฺวา ตํ ทิวสํ เยว สพฺพฏฺฐกํ ลภิ, สพฺเพเปเต มหาราช ทิฏฺฐ\-ธมฺมิกํ โภคญฺจ ยสญฺจ อนุภวิํสูติ.}

\proseitem{5}{ภนฺเต นาคเสน วิจินิตฺวา ปริเยสิตฺวา ฉ ชเน เยว อทฺทสาสีติ. อาม มหาราชาติ. เตน หิ ภนฺเต นาคเสน อกุสลํ เยว อธิมตฺตํ พลวตรํ, โน ตถา กุสลํ, อหญฺหิ ภนฺเต นาคเสน เอกทิวสํ เยว ทสปิ ปุริเส ปสฺสามิ ปาปสฺส กมฺมสฺส วิปาเกน สูเลสุ อาโรเปนฺเต, วีสมฺปิ ติํสมฺปิ จตฺตาลีสมฺปิ ปญฺญาสมฺปิ ปุริสสตมฺปิ ปุริส\-สหสฺสมฺปิ ปสฺสามิ ปาปสฺส กมฺมสฺส วิปาเกน สูเลสุ อาโรเปนฺเต. นนฺทกุลสฺส ภนฺเต นาคเสน ภทฺทสาโล นาม เสนาปติปุตฺโต อโหสิ, เตน จ รญฺญา จนฺทคุตฺเตน สงฺคาโม สมุปพฺยูโฬฺห อโหสิ, ตสฺมิํ โข ปน ภนฺเต นาคเสน สงฺคาเม อุภโต พลกาเย อสีติกพนฺธรูปนิ อเหสุํ, เอกสฺมิํ กิร สีสกพนฺเธ ปริปาเต\csromansharedfootnote{28ca647604d533dc}{ปริปุณฺเณ (สพฺพตฺถ)} เอกํ กพนฺธรูปํ อุฏฺฐหติ, สพฺเพเปเต ปาปสฺเสว กมฺมสฺส วิปาเกน อนยพฺยสนํ อาปนฺนา, อิมินาปิ ภนฺเต นาคเสน การเณน ภณามิ อกุสลํ เยว อธิมตฺตํ พลวตรํ, โน ตถา กุสลนฺติ.}

\prose{สุยฺยติ ภนฺเต นาคเสน อิมสฺมิํ พุทฺธสาสเน โกสเลน รญฺญา อสทิสทานํ ทินฺนนฺติ. อาม มหาราช สุยฺยตีติ. อปิ นุ โข ภนฺเต นาคเสน โกสลราชา ตํ อสทิสํ ทานํ ทตฺวา ตโตนิทานํ กญฺจิ ทิฏฺฐ\-ธมฺมิกํ โภคํ วา ยสํ วา สุขํ วา ปฏิลภีติ\csromansharedfootnote{1a1b074552fe670e}{ปฏิลภตีติ. (ก)}. น หิ มหาราชาติ. ยทิ ภนฺเต นาคเสน โกสลราชา เอวรูปํ อนุตฺตรํ ทานํ ทตฺวาปิ น ลภิ\csromansharedfootnote{3ebe8c216eaffcbe}{น ลภติ (ก)} ตโตนิทานํ กญฺจิ ทิฏฺฐ\-ธมฺมิกํ โภคํ วา ยสํ วา สุขํ วา, เตน หิ ภนฺเต นาคเสน อกุสลํ เยว อธิมตฺตํ พลวตรํ, โน ตถา กุสลนฺติ.}

\prose{ปริตฺตตฺตา มหาราช อกุสลํ ขิปฺปํ ปริณมติ, วิปุลตฺตา กุสลํ ทีเฆน กาเลน ปริณมติ, อุปมายปิ มหาราช เอตํ อุปปริ\-กฺขิ\-ตพฺพํ. ยถา มหาราช อปรนฺเต ชนปเท กุมุทภณฺฑิกา นาม ธญฺญชาติ มาสลูนา\csromansharedfootnote{b86ab0362a9e7d34}{มาสปูรา (ก)} อนฺโตเคหคตา โหติ, สาลโย ฉปฺปญฺจ\-มาเสหิ \csromanfolio{283}ปริณมนฺติ, กิํ ปเนตฺถ มหาราช อนฺตรํ โก วิเสโส กุมุท\-ภณฺฑิกาย จ สาลีนญฺจาติ. ปริตฺตตฺตา ภนฺเต กุมุท\-ภณฺฑิกาย วิปุลตฺตา จ สาลีนํ, สาลโย ภนฺเต นาคเสน ราชารหา ราชโภชนํ, กุมุทภณฺฑิกา ทาส\-กมฺมกรานํ โภชนนฺติ. เอวเมว โข มหาราช ปริตฺตตฺตา อกุสลํ ขิปฺปํ ปริณมติ, วิปุลตฺตา กุสลํ ทีเฆน กาเลน ปริณมตีติ.}

\prose{ยํ ตตฺถ ภนฺเต นาคเสน ขิปฺปํ ปริณมติ, ตํ นาม โลเก อธิมตฺตํ พลวตรํ, ตสฺมา อกุสลํ พลวตรํ, โน ตถา กุสลํ. ยถา นาม ภนฺเต นาคเสน โย โกจิ โยโธ มหติ\-มหา\-ยุทฺธํ ปวิสิตฺวา ปฏิสตฺตุํ อุปกจฺฉเก คเหตฺวา อากฑฺฒิตฺวา ขิปฺปตรํ สามิโน อุปเนยฺย, โส โยโธ โลเก สมตฺโถ สูโร นาม. โย จ ภิสกฺโก ขิปฺปํ สลฺลํ อุทฺธรติ โรคมปเนติ, โส ภิสกฺโก เฉโก นาม. โย คณโก สีฆสีฆํ คเณตฺวา ขิปฺปํ ทสฺสยติ, โส คณโก เฉโก นาม. โย มลฺโล ขิปฺปํ ปฏิมลฺลํ อุกฺขิปิตฺวา อุตฺตานกํ ปาเตติ, โส มลฺโล สมตฺโถ สูโร นาม. เอวเมว โข ภนฺเต นาคเสน ยํ ขิปฺปํ ปริณมติ กุสลํ วา อกุสลํ วา, ตํ โลเก อธิมตฺตํ พลวตรนฺติ.}

\prosewithcloserruled{อุภยมฺปิ ตํ มหาราช กมฺมํ สมฺปรายเวทนียเมว, อปิ จ โข อกุสลํ สาวชฺชตาย ขเณน ทิฏฺฐธมฺม\-เวทนียํ โหติ, ปุพฺพเกหิ มหาราช ขตฺติเยหิ ฐปิโต เอโส นิยโม “โย ปาณํ หนติ, โส ทณฺฑารโห. โย อทินฺนํ อาทิยติ. โย ปรทารํ คจฺฉติ. โย มุสา ภณติ. โย คามํ ฆาเตติ. โย ปนฺถํ ทูเสติ. โย นิกติํ กโรติ. โย วญฺจนํ กโรติ, โส ทณฺฑารโห วธิตพฺโพ เฉตฺตพฺโพ เภตฺตพฺโพ หนฺตพฺโพ”ติ. ตํ เต อุปาทาย วิจินิตฺวา วิจินิตฺวา ทณฺเฑนฺติ วเธนฺติ ฉินฺทนฺติ ภินฺทนฺติ หนนฺติ จ, อปิ นุ มหาราช อตฺถิ เกหิจิ ฐปิโต นิยโม “โย ทานํ วา เทติ, สีลํ วา รกฺขติ, อุโปสถกมฺมํ วา กโรติ, ตสฺส ธนํ วา ยสํ วา ทาตพฺพนฺ”ติ. อปิ นุ ตํ วิจินิตฺวา วิจินิตฺวา ธนํ วา ยสํ วา เทนฺติ, โจรสฺส กตกมฺมสฺส วธพนฺธนํ วิยาติ. น หิ ภนฺเตติ. ยทิ มหาราช ทายกานํ วิจินิตฺวา วิจินิตฺวา ธนํ วา ยสํ วา ทเทยฺยุํ, กุสลมฺปิ ทิฏฺฐธมฺม\-เวทนียํ ภเวยฺย, ยสฺมา จ โข มหาราช ทายเก น วิจินนฺติ “ธนํ วา ยสํ วา ทสฺสามา”ติ, ตสฺมา กุสลํ น ทิฏฺฐธมฺม\-เวทนียํ. อิมินา มหาราช การเณน อกุสลํ ทิฏฺฐธมฺม\-เวทนียํ, \csromanfolio{284}สมฺปรา เยว โส อธิมตฺตํ พลวตรํ เวทนํ เวทยตีติ. สาธุ ภนฺเต นาคเสน, ตวาทิเสน พุทฺธิมนฺเตน วินา เนโส ปโญฺห สุนิพฺเพฐิโย, โลกิกํ ภนฺเต นาคเสน โลกุตฺตเรน วิญฺญาปิตนฺติ.}{กุสลากุสล\-พลวตร\-ปโญฺห ตติโย.}
\csromanmatikaanchor{mtk.188}
\csromantocmarkref{subsubsection}{4. ปุพฺพเปตาทิสปญฺห}{mtk.188}
\bookmark[dest={mtk.188},level=3]{4. ปุพฺพเปตาทิสปญฺห}
\csromanheader{3}{4. \textbf{ปุพฺพเปตา}\-\textbf{ทิส}\-\textbf{ปญฺห}}
\proseitem{5}{ภนฺเต นาคเสน อิเม ทายกา ทานํ ทตฺวา ปุพฺพเปตานํ อาทิสนฺติ\csromansharedfootnote{d77c4b883e6063fa}{อุทฺทิสนฺติ (ก, สี)} “อิทํ เตสํ ปาปุณาตู”ติ, อปิ นุ เต กิญฺจิ ตโตนิทานํ วิปากํ ปฏิลภนฺตีติ. เกจิ มหาราช ปฏิลภนฺติ, เกจิ นปฺปฏิลภนฺตีติ. เก ภนฺเต ปฏิลภนฺติ, เก นปฺปฏิลภนฺตีติ.\csromansymbolfootnote{*}{อํ 3. 478 ปิฏฺเฐ.}นิรยูปปนฺนา มหาราช นปฺปฏิลภนฺติ, สคฺคคตา นปฺปฏิลนฺติ, ติรจฺฉานโยนิ\-คตา นปฺปฏิลภนฺติ, จตุนฺนํ เปตานํ ตโย เปตา นปฺปฏิลภนฺติ วนฺตาสิกา ขุปฺปิปาสิโน นิชฺฌาม\-ตณฺหิกา, ลภนฺติ เปตา ปรทตฺตู\-ปชีวิโน, เตปิ สรมานา เยว ลภนฺตีติ.}

\prose{เตน หิ ภนฺเต นาคเสน ทายกานํ ทานํ วิโสสิตํ\csromansharedfootnote{80838716ba347ae0}{วิโสตํ (สี, อิ)} โหติ อผลํ, เยสํ อุทฺทิสฺส กตํ ยทิ เต นปฺปฏิลภนฺตีติ. น หิ ตํ มหาราช ทานํ อผลํ โหติ อวิปากํ, ทายกา เยว ตสฺส ผลํ อนุภวนฺตีติ. เตน หิ ภนฺเต นาคเสน การเณน มํ สญฺญาเปหีติ. อิธ มหาราช เกจิ มนุสฺสา มจฺฉมํส\-สุรา\-ภตฺต\-ขชฺชกานิ ปฏิยาเทตฺวา ญาติกุลํ คจฺฉนฺติ, ยทิ เต ญาตกา ตํ อุปายนํ น สมฺปฏิจฺเฉยฺยุํ, อปิ นุ ตํ อุปายนํ วิโสสิตํ คจฺเฉยฺย วินสฺเสยฺย วาติ. น หิ ภนฺเต, สามิกานํ เยว ตํ โหตีติ, เอวเมว โข มหาราช ทายกา เยว ตสฺส ผลํ อนุภวนฺติ. ยถา ปน มหาราช ปุริโส คพฺภํ ปวิฏฺโฐ อสติ ปุรโต นิกฺขมน\-มุเข เกน นิกฺขเมยฺยาติ. ปวิฏฺเฐเนว ภนฺเตติ. เอวเมว โข มหาราช ทายกา เยว ตสฺส ผลํ อนุภวนฺตีติ. โหตุ ภนฺเต นาคเสน, เอวเมตํ ตถา สมฺปฏิจฺฉามิ, ทายกา เยว ตสฺส ผลํ อนุภวนฺติ, น มยํ ตํ การณํ วิโลเมมาติ.}

\prose{\csromanfolio{285}ภนฺเต นาคเสน ยทิ อิเมสํ ทายกานํ ทินฺนทานํ ปุพฺพเปตานํ ปาปุณาติ, เต จ ตสฺส วิปากํ อนุภวนฺติ. เตน หิ โย ปาณาติปาตี ลุทฺโท โลหิตปาณี ปทุฏฺฐ\-มนสงฺกปฺโป มนุสฺเส ฆาเตตฺวา ทารุณํ กมฺมํ กตฺวา ปุพฺพเปตานํ อาทิเสยฺย “อิมสฺส เม กมฺมสฺส วิปาโก ปุพฺพเปตานํ ปาปุณาตู”ติ, อปิ นุ ตสฺส วิปาโก ปุพฺพเปตานํ ปาปุณาตีติ. น หิ มหาราชาติ.}

\prose{ภนฺเต นาคเสน โก ตตฺถ เหตุ กิํ การณํ, เยน กุสลํ ปาปุณาติ อกุสลํ น ปาปุณาตีติ. เนโส มหาราช ปโญฺห ปุจฺฉิตพฺโพ, มา จ ตฺวํ มหาราช “วิสชฺชโก อตฺถี”ติ อปุจฺฉิตพฺพํ ปุจฺฉิ, “กิสฺส อากาโส นิราลมฺโพ, กิสฺส คงฺคา อุทฺธมฺมุขา น สนฺทติ, กิสฺส อิเม มนุสฺสา จ ทิชา จ ทฺวิปทา มิคา จตุปฺปทา”ติ ตมฺปิ มํ ตฺวํ ปุจฺฉิสฺสตีติ. นาหํ ตํ ภนฺเต นาคเสน วิเหสาเปกฺโข ปุจฺฉามิ, อปิ จ นิพฺพาหน\-ตฺถาย\csromansharedfootnote{c1b786675a9ebf9e}{นิพฺพานตฺถาย (ก)} สนฺเทหสฺส ปุจฺฉามิ, พหู มนุสฺสา โลเก วามคามิโน\csromansharedfootnote{edc572d4a721d67f}{ปาปคาหิโน (สฺยา)} วิจกฺขุกา, “กินฺติ เต โอตารํ น ลเภยฺยูนฺ”ติ เอวาหํ ตํ ปุจฺฉามีติ. น สกฺกา มหาราช สห อกเตน อนนุมเตน สห ปาปํ กมฺมํ สํวิภชิตุํ.}

\prose{ยถา มหาราช มนุสฺสา อุทก\-นิพฺพาหเนน อุทกํ สุวิทูรมฺปิ หรนฺติ, อปิ นุ มหาราช สกฺกา ฆนมหาเสล\-ปพฺพโต\csromansharedfootnote{1dc42ef557db60fc}{ปพฺพตโต (ก)} นิพฺพาหเนน ยถิจฺฉิตํ หริตุนฺติ. น หิ ภนฺเตติ. เอวเมว โข มหาราช สกฺกา กุสลํ สํวิภชิตุํ, น สกฺกา อกุสลํ สํวิภชิตุํ. ยถา วา ปน มหาราช สกฺกา เตเลน ปทีโป ชาเลตุํ, อปิ นุ มหาราช สกฺกา อุทเกน ปทีโป ชาเลตุนฺติ. น หิ ภนฺเตติ. เอวเมว โข มหาราช สกฺกา กุสลํ สํวิภชิตุํ, น สกฺกา อกุสลํ สํวิภชิตุํ. ยถา วา ปน มหาราช กสฺสกา ตฬากโต อุทกํ นีหริตฺวา ธญฺญํ ปริปาเจนฺติ. อปิ นุ โข มหาราช สกฺกา มหาสมุทฺทโต อุทกํ นีหริตฺวา ธญฺญํ ปริปาเจตุนฺติ. น หิ ภนฺเตติ. เอวเมว โข มหาราช สกฺกา กุสลํ สํวิภชิตุํ, น สกฺโก อกุสลํ สํวิภชิตุนฺติ.}

\prose{ภนฺเต นาคเสน เกน การเณน สกฺกา กุสลํ สํวิภชิตุํ, น สกฺกา อกุสลํ สํวิภชิตุํ. การเณน มํ สญฺญาเปหิ, นาหํ อนฺโธ}

\prose{\csromanfolio{286}อนาโลโก สุตฺวา เวทิสฺสามีติ. อกุสลํ มหาราช โถกํ, กุสลํ พหุกํ, โถกตฺตา อกุสลํ กตฺตารํ เยว ปริยาทิยติ, พหุกตฺตา กุสลํ สเทวกํ โลกํ อชฺโฌตฺถรตีติ. โอปมฺมํ กโรหีติ.}

\prose{ยถา มหาราช ปริตฺตํ เอกํ อุทกพินฺทุ ปถวิยํ นิปเตยฺย, อปิ นุ โข ตํ มหาราช อุทกพินฺทุ ทสปิ ทฺวาทสปิ โยชนานิ อชฺโฌตฺถเรยฺยาติ. น หิ ภนฺเต, ยตฺถ ตํ อุทกพินฺทุ นิปติตํ, ตตฺเถว ปริยาทิยตีติ. เกน การเณน มหาราชาติ. ปริตฺตตฺตา ภนฺเต อุทกพินฺทุสฺสาติ. เอวเมว โข มหาราช ปริตฺตํ อกุสลํ ปริตฺตตฺตา กตฺตารํ เยว ปริยาทิยติ, น สกฺกา สํวิภชิตุํ.}

\prose{ยถา วา ปน มหาราช มหติ\-มหาเมโฆ อภิวสฺเสยฺย ตปฺปยนฺโต ธรณิตลํ, อปิ นุ โข โส มหาราช มหาเมโฆ สมนฺตโต โอตฺถเรยฺยาติ. อาม ภนฺเต ปูรยิตฺวา โส มหาเมโฆ โสพฺภ สร สริตสาขากนฺทรปทร ทห ตฬาก\csromansharedfootnote{9f74b389928f4cdc}{มาติกาตฬาก (ก)} อุทปาน\-โปกฺขรณิโย ทสปิ ทฺวาทสปิ โยชนานิ อชฺโฌตฺถเรยฺยาติ. เกน การเณน มหาราชาติ. มหนฺตตฺตา ภนฺเต เมฆสฺสาติ. เอวเมว โข มหาราช กุสลํ พหุกํ, พหุกตฺตา สกฺกา เทวมนุสฺเสหิปิ สํวิภชิตุนฺติ.}

\prose{ภนฺเต นาคเสน เกน การเณน อกุสลํ โถกํ กุสลํ พหุตรนฺติ. อิธ มหาราช โย โกจิ ทานํ เทติ, สีลํ สมาทิยติ, อุโปสถกมฺมํ กโรติ, โส หฏฺโฐ ปหฏฺโฐ หสิโต ปมุทิโต ปสนฺนมานโส เวทชาโต โหติ, ตสฺส อปราปรํ ปีติ อุปฺปชฺชติ, ปีติมนสฺส ภิยฺโย ภิยฺโย กุสลํ ปวฑฺฒติ.}

\prose{ยถา มหาราช อุทปาเน พหุสลิล\-สมฺปุณฺเณ เอเกน เทเสน อุทกํ ปวิเสยฺย, เอเกน นิกฺขเมยฺย, นิกฺขมนฺเตปิ อปราปรํ อุปฺปชฺชติ, น สกฺกา โหติ ขยํ ปาเปตุํ. เอวเมว โข มหาราช กุสลํ ภิยฺโย ภิยฺโย ปวฑฺฒติ. วสฺสสเตปิ เจ มหาราช ปุริโส กตํ กุสลํ อาวชฺเชยฺย, อาวชฺชิเต อาวชฺชิเต ภิยฺโย ภิยฺโย กุสลํ ปวฑฺฒติ. ตสฺส ตํ กุสลํ สกฺกา โหติ ยถิจฺฉเกหิ สทฺธิํ สํวิภชิตุํ, อิทเมตฺถ มหาราช การณํ, เยน การเณน กุสลํ พหุตรํ.}

\prosewithcloserruled{\csromanfolio{287}อกุสลํ ปน มหาราช กโรนฺโต ปจฺฉา วิปฺปฏิสารี โหติ, วิปฺปฏิสาริโน จิตฺตํ ปฏิลียติ ปฏิกุฏติ ปฏิวตฺตติ น สมฺปสารียติ โสจติ ตปฺปติ หายติ ขียติ น ปริวฑฺฒติ ตตฺเถว ปริยาทิยติ. ยถา มหาราช สุกฺขาย นทิยา มหาปุฬินาย อุนฺนตาวนตาย กุฏิลส\-งฺกุฏิลาย อุปริโต ปริตฺตํ อุทกํ อาคจฺฉนฺตํ หายติ ขียติ น ปริวฑฺฒติ ตตฺเถว ปริยาทิยติ. เอวเมว โข มหาราช อกุสลํ กโรนฺตสฺส จิตฺตํ ปฏิลียติ ปฏิกุฏติ ปฏิวตฺตติ น สมฺปสารียติ โสจติ ตปฺปติ หายติ ขียติ น ปริวฑฺฒติ ตตฺเถว ปริยาทิยติ, อิทเมตฺถ มหาราช การณํ, เยน การเณน อกุสลํ โถกนฺติ. สาธุ ภนฺเต นาคเสน, เอวเมตํ ตถา สมฺปฏิจฺฉามีติ.}{ปุพฺพเปตา\-ทิส\-ปโญฺห จตุตฺโถ.}
\csromanmatikaanchor{mtk.189}
\csromantocmarkref{subsubsection}{5. สุปินปญฺห}{mtk.189}
\bookmark[dest={mtk.189},level=3]{5. สุปินปญฺห}
\csromanheader{3}{5. \textbf{สุปินปญฺห}}
\proseitem{5}{ภนฺเต นาคเสน อิมสฺมิํ โลเก นรนาริโย สุปินํ ปสฺสนฺติ กลฺยาณมฺปิ ปาปกมฺปิ, ทิฏฺฐปุพฺพมฺปิ อทิฏฺฐ\-ปุพฺพมฺปิ, กตปุพฺพมฺปิ อกตปุพฺพมฺปิ, เขมมฺปิ สภยมฺปิ, ทูเรปิ สนฺติเกปิ, พหุวิธานิปิ อเนกวณฺณ\-สหสฺสานิ ทิสฺสนฺติ, กิญฺเจตํ สุปินํ นาม, โก เจตํ ปสฺสตีติ. นิมิตฺตเมตํ มหาราช สุปินํ นาม, ยํ จิตฺตสฺส อาปาต\csromansharedfootnote{d93774ff73259885}{อาปาถ (สี, อิ)} มุปคจฺฉติ. ฉยิเม มหาราช สุปินํ ปสฺสนฺติ, วาติโก สุปินํ ปสฺสติ, ปิตฺติโก สุปินํ ปสฺสติ, เสมฺหิโก สุปินํ ปสฺสติ, เทวตู\-ปสํหารโต สุปินํ ปสฺสติ, สมุทาจิณฺณโต สุปินํ ปสฺสติ, ปุพฺพนิมิตฺตโต สุปินํ ปสฺสติ, ตตฺร มหาราช ยํ ปุพฺพนิมิตฺตโต สุปินํ ปสฺสติ, ตํ เยว\csromansharedfootnote{cf7dc7e6d0b8f44d}{ตํเยว นิคฺคหีตสนฺธิ \csromandash{} ม.พ.ป.} สจฺจํ, อวเสสํ มิจฺฉาติ.}

\prose{ภนฺเต นาคเสน โย ปุพฺพนิมิตฺตโต สุปินํ ปสฺสติ, กิํ ตสฺส จิตฺตํ สยํ คนฺตฺวา ตํ นิมิตฺตํ วิจินาติ, ตํ วา นิมิตฺตํ จิตฺตสฺส อาปาตมุ\-ปคจฺฉติ, อญฺโญ วา อาคนฺตฺวา ตสฺส อาโรเจตีติ. น มหาราช ตสฺส จิตฺตํ สยํ คนฺตฺวา ตํ นิมิตฺตํ วิจินาติ, นาปิ อญฺโญ โกจิ อาคนฺตฺวา ตสฺส อาโรเจติ, อถ โข ตํ เยว\csromansharedfootnote{cf7dc7e6d0b8f44d}{ตํเยว นิคฺคหีตสนฺธิ \csromandash{} ม.พ.ป.} นิมิตฺตํ จิตฺตสฺส อาปาตมุ\-ปคจฺฉติ. ยถา มหาราช อาทาโส น สยํ กุหิญฺจิ คนฺตฺวา ฉายํ วิจินาติ, นาปิ อญฺโญ \csromanfolio{288}โกจิ ฉายํ อาเนตฺวา อาทาสํ อาโรเปติ\csromansharedfootnote{878bec0a1e6ee50c}{อาโรเจติ (ก)}, อถ โข ยโต กุโตจิ ฉายา อาคนฺตฺวา อาทาสสฺส อาปาตมุ\-ปคจฺฉติ, เอวเมว โข มหาราช น ตสฺส จิตฺตํ สยํ คนฺตฺวา ตํ นิมิตฺตํ วิจินาติ, นาปิ อญฺโญ โกจิ อาคนฺตฺวา อาโรเจติ, อถ โข ยโต กุโตจิ นิมิตฺตํ อาคนฺตฺวา จิตฺตสฺส อาปาตมุ\-ปคจฺฉตีติ.}

\prose{ภนฺเต นาคเสน ยํ ตํ จิตฺตํ สุปินํ ปสฺสติ, อปิ นุ ตํ จิตฺตํ ชานาติ “เอวํ นาม วิปาโก ภวิสฺสติ เขมํ วา ภยํ วา”ติ. น หิ มหาราช ตํ จิตฺตํ ชานาติ “เอวํวิปาโก ภวิสฺสติ เขมํ วา ภยํ วา”ติ, นิมิตฺเต ปน อุปฺปนฺเน อญฺเญสํ กเถติ, ตโต เต อตฺถํ กเถนฺตีติ.}

\prose{อิงฺฆ ภนฺเต นาคเสน การณํ เม ทสฺเสหีติ. ยถา มหาราช สรีเร ติลกา ปีฬกา ททฺทูนิ อุฏฺฐหนฺติ ลาภาย วา อลาภาย วา, ยสาย วา อยสาย วา, นินฺทาย วา ปสํสาย วา, สุขาย วา ทุกฺขาย วา, อปิ นุ ตา มหาราช ปีฬกา ชานิตฺวา อุปฺปชฺชนฺติ “อิมํ นาม มยํ อตฺถํ นิปฺผาเทสฺสามี”ติ. น หิ ภนฺเต, ยาทิเส ตา โอกาเส ปีฬกา สมฺภวนฺติ, ตตฺถ ตา ปีฬกา ทิสฺวา เนมิตฺตกา พฺยากโรนฺติ “เอวํ นาม วิปาโก ภวิสฺสตี”ติ. เอวเมว โข มหาราช ยํ ตํ จิตฺตํ สุปินํ ปสฺสติ, น ตํ จิตฺตํ ชานาติ “เอวํ นาม วิปาโก ภวิสฺสติ เขมํ วา ภยํ วา”ติ, นิมิตฺเต ปน อุปฺปนฺเน อญฺเญสํ กเถติ, ตโต เต อตฺถํ กเถนฺตีติ.}

\prose{ภนฺเต นาคเสน โย สุปินํ ปสฺสติ, โส นิทฺทายนฺโต, อุทาหุ ชาครนฺโต\csromansharedfootnote{fde7e1eb2a67f1a6}{ชคฺคนฺโต (สี, อิ)} ปสฺสตีติ. โย โส มหาราช สุปินํ ปสฺสติ, น โส นิทฺทายนฺโต ปสฺสติ, นาปิ ชาครนฺโต ปสฺสติ. อปิ จ โอกฺกนฺเต มิทฺเธ อสมฺปตฺเต ภวงฺเค เอตฺถนฺตเร สุปินํ ปสฺสติ. มิทฺธ\-สมารูฬฺหสฺส มหาราช จิตฺตํ ภวงฺคคตํ โหติ, ภวงฺคคตํ จิตฺตํ นปฺปวตฺตติ, อปฺปวตฺตํ จิตฺตํ สุขทุกฺขํ นปฺปชานาติ, อปฺปฏิวิชานนฺตสฺส สุปิโน น โหติ, ปวตฺตมาเน จิตฺเต สุปินํ ปสฺสติ.}

\prose{ยถา มหาราช ติมิเร อนฺธกาเร อปฺปภาเส สุปริสุทฺเธปิ อาทาเส ฉายา น ทิสฺสติ, เอวเมว โข มหาราช มิทฺธ\-สมารูเฬฺห จิตฺเต ภวงฺคคเต ติฏฺฐมาเนปิ สรีเร จิตฺตํ อปฺปวตฺตํ โหติ, อปฺปวตฺเต จิตฺเต สุปินํ น \csromanfolio{289}ปสฺสติ. ยถา มหาราช อาทาโส, เอวํ สรีรํ ทฏฺฐพฺพํ ยถา อนฺธกาโร, เอวํ มิทฺธํ ทฏฺฐพฺพํ, ยถา อาโลโก, เอวํ จิตฺตํ ทฏฺฐพฺพํ.}

\prose{ยถา วา ปน มหาราช มหิโก\-ตฺถฏสฺส สูริยสฺส ปภา น ทิสฺสติ สนฺตา เยว สูริยรสฺมิ อปฺปวตฺตา โหติ, อปฺปวตฺตาย สูริยรสฺมิยา อาโลโก น โหติ, เอวเมว โข มหาราช มิทฺธ\-สมารูฬฺหสฺส จิตฺตํ ภวงฺคคตํ โหติ, ภวงฺคคตํ จิตฺตํ นปฺปวตฺตติ, อปฺปวตฺเต จิตฺต สุปินํ น ปสฺสติ. ยถา มหาราช สูริโย, เอวํ สรีรํ ทฏฺฐพฺพํ, ยถา มหิโก\-ตฺถรณํ, เอวํ มิทฺธํ ทฏฺฐพฺพํ, ยถา สูริยรสฺมิ, เอวํ จิตฺตํ ทฏฺฐพฺพํ.}

\prose{ทฺวินฺนํ มหาราช สนฺเตปิ สรีเร จิตฺตํ อปฺปวตฺตํ โหติ, มิทฺธ\-สมารูฬฺหสฺส ภวงฺค\-คตสฺส สนฺเตปิ สรีเร จิตฺตํ อปฺปวตฺตํ โหติ, นิโรธ\-สมาปนฺนสฺส สนฺเตปิ สรีเร จิตฺตํ อปฺปวตฺตํ โหติ, ชาครนฺตสฺส มหาราช จิตฺตํ โลลํ โหติ วิวฏํ ปากฏํ อนิพทฺธํ, เอวรูปสฺส จิตฺเต นิมิตฺตํ อาปาตํ น อุเปติ, ยถา มหาราช ปุริสํ วิวฏํ ปากฏํ อกิริยํ อรหสฺสํ รหสฺสกามา ปริวชฺเชนฺติ, เอวเมว โข มหาราช ชาครนฺตสฺส ทิพฺโพ อตฺโถ อาปาตํ น อุเปติ, ตสฺมา ชาครนฺโต สุปินํ น ปสฺสติ. ยถา วา ปน มหาราช ภิกฺขุํ ภินฺนาชีวํ อนาจารํ ปาปมิตฺตํ ทุสฺสีลํ กุสีตํ หีนวีริยํ กุสลา โพธิปกฺขิยา ธมฺมา อาปาตํ น อุเปนฺติ, เอวเมว โข มหาราช ชาครนฺตสฺส ทิพฺโพ อตฺโถ อาปาตํ น อุเปติ, ตสฺมา ชาครนฺโต สุปินํ น ปสฺสตีติ.}

\prosewithcloserruled{ภนฺเต นาคเสน อตฺถิ มิทฺธสฺส อาทิ\-มชฺฌ\-ปริโยสานนฺติ. อาม มหาราช อตฺถิ มิทฺธสฺส อาทิ\-มชฺฌ\-ปริโยสานนฺติ. กตมํ อาทิ, กตมํ มชฺฌํ, กตมํ ปริโยสานนฺติ. โย มหาราช กายสฺส โอนาโห ปริโยนาโห ทุพฺพลฺยํ มนฺทตา อกมฺมญฺญตา กายสฺส, อยํ มิทฺธสฺส อาทิ, โย มหาราช กปินิทฺทา\-ปเรโต โวกิณฺณกํ ชคฺคติ\csromansharedfootnote{e91493c0b295166d}{โวกิณฺณตํ คจฺฉติ (นิสฺย)}, อิทํ มิทฺธสฺส มชฺฌํ. ภวงฺคคติ ปริโยสนํ. มชฺฌูปคโต มหาราช กปินิทฺทา\-ปเรโต สุปินํ ปสฺสติ. ยถา มหาราช โกจิ ยตจารี สมาหิตจิตฺโต ฐิตธมฺโม อจลพุทฺธิ ปหีน\-โกตูหล\-สทฺทํ วนม\-ชฺโฌคาหิตฺวา สุขุมํ อตฺถํ จินฺตยติ, น จ โส ตตฺถ มิทฺธํ โอกฺกมติ, โส ตตฺถ สมาหิโต เอกคฺคจิตฺโต สุขุมํ อตฺถํ ปฏิวิชฺฌติ, เอวเมว โข มหาราช ชาคโร น \csromanfolio{290}มิทฺธ\-สมาปนฺโน มชฺฌูปคโต กปินิทฺทา\-ปเรโต สุปินํ ปสฺสติ. ยถา มหาราช โกตูหลสทฺโท, เอวํ ชาครํ ทฏฺฐพฺพํ, ยถา วิวิตฺตํ วนํ, เอวํ กปินิทฺทา\-ปเรโต ทฏฺฐพฺโพ, ยถา โส โกตูหลสทฺทํ โอหาย มิทฺธํ วิวชฺเชตฺวา มชฺฌตฺตภูโต สุขุมํ อตฺถํ ปฏิวิชฺฌติ, เอวํ ชาคโร น มิทฺธ\-สมาปนฺโน กปินิทฺทา\-ปเรโต สุปินํ ปสฺสตีติ. สาธุ ภนฺเต นาคเสน, เอวเมตํ ตถา สมฺปฏิจฺฉามีติ.}{สุปินปโญฺห ปญฺจโม.}
\csromanmatikaanchor{mtk.190}
\csromantocmarkref{subsubsection}{6. อกาลมรณปญฺห}{mtk.190}
\bookmark[dest={mtk.190},level=3]{6. อกาลมรณปญฺห}
\csromanheader{3}{6. \textbf{อกาลมรณ}\-\textbf{ปญฺห}}
\proseitem{5}{ภนฺเต นาคเสน เย เต สตฺตา มรนฺติ, สพฺเพ เต กาเล เยว มรนฺติ, อุทาหุ อกาเลปิ มรนฺตีติ. อตฺถิ มหาราช กาเลปิ มรณํ, อตฺถิ อกาเลปิ มรณนฺติ.}

\prose{ภนฺเต นาคเสน เก กาเล มรนฺติ, เก อกาเล มรนฺตีติ. ทิฏฺฐปุพฺพา ปน มหาราช ตยา อมฺพรุกฺขา วา ชมฺพุรุกฺขา วา, อญฺญสฺมา วา ปน ผลรุกฺขา ผลานิ ปตนฺตานิ อามานิ จ ปกฺกานิ จาติ. อาม ภนฺเตติ. ยานิ ตานิ มหาราช ผลานิ รุกฺขโต ปตนฺติ, สพฺพานิ ตานิ กาเล เยว ปตนฺติ, อุทาหุ อกาเลปีติ. ยานิ ตานิ ภนฺเต นาคเสน ผลานิ ปริปกฺกานิ วิลีนานิ ปตนฺติ, สพฺพานิ ตานิ กาเล ปตนฺติ. ยานิ ปน ตานิ อวเสสานิ ผลานิ เตสุ กานิจิ กิมิวิทฺธานิ ปตนฺติ, กานิจิ ลคุฬหตานิ\csromansharedfootnote{063cd1cf31cea267}{สกุณปหตา (สฺยา), ลกุฏหตานิ (สี, อิ, ก)} ปตนฺติ, กานิจิ วาตปฺปหตานิ ปตนฺติ, กานิจิ อนฺโตปูติกานิ หุตฺวา ปตนฺติ, สพฺพานิ ตานิ อกาเล ปตนฺตีติ. เอวเมว โข มหาราช เย เต ชราเวคหตา มรนฺติ, เต เยว กาเล มรนฺติ, อวเสสา เกจิ กมฺม\-ปฺปฏิพาฬฺหา มรนฺติ, เกจิ คติปฺปฏิพาฬฺหา มรนฺติ, เกจิ กิริย\-ปฺปฏิพาฬฺหา มรนฺตีติ.}

\prose{ภนฺเต นาคเสน เย เต กมฺม\-ปฺปฏิพาฬฺหา มรนฺติ, เยปิ เต คติปฺปฏิพาฬฺหา มรนฺติ, เยปิ เต กิริย\-ปฺปฏิพาฬฺหา มรนฺติ, เยปิ เต ชรา\-เวค\-ปฺปฏิพาฬฺหา มรนฺติ, สพฺเพ เต กาเล เยว มรนฺติ, โยปิ มาตุกุจฺฉิคโต มรติ, โส ตสฺส กาโล, กาเล เยว โส มรติ. โยปิ วิชาตฆเร มรติ, โส ตสฺส กาโล, โสปิ กาเล เยว มรติ. โยปิ \csromanfolio{291}มาสิโก มรติ ฯเปฯ. โยปิ วสฺสสติโก มรติ, โส ตสฺส กาโล, กาเล เยว โส มรติ, เตน หิ ภนฺเต นาคเสน อกาเล มรณํ นาม น โหติ, เย เกจิ มรนฺติ, สพฺเพ เต กาเล เยว มรนฺตีติ.}

\proseitem{6}{สตฺติเม มหาราช วิชฺชมาเนปิ อุตฺตริํ อายุสฺมิํ อกาเล มรนฺติ. กตเม สตฺต, ชิฆจฺฉิโต มหาราช โภชนํ อลภมาโน อุปหต\-พฺภนฺตโร วิชฺชมาเนปิ อุตฺตริํ อายุสฺมิํ อกาเล มรติ, ปิปาสิโต มหาราช ปานียํ อลภมาโน ปริสุกฺข\-หทโย วิชฺชมาเนปิ อุตฺตริํ อายุสฺมิํ อกาเล มรติ, อหินา ทฏฺโฐ มหาราช วิสเวคา\-ภิหโต ติกิจฺฉกํ อลภมาโน วิชฺชมาเนปิ อุตฺตริํ อายุสฺมิํ อกาเล มรติ, วิสมาสิโต มหาราช ฑยฺหนฺเตสุ องฺคปจฺจงฺเคสุ อคทํ อลภมาโน วิชฺชมาเนปิ อุตฺตริํ อายุสฺมิํ อกาเล มรติ, อคฺคิคโต มหาราช ฌายมาโน นิพฺพาปนํ อลภมาโน วิชฺชมาเนปิ อุตฺตริํ อายุสฺมิํ อกาเล มรติ, อุทกคโต มหาราช ปติฏฺฐํ อลภมาโน วิชฺชมาเนปิ อุตฺตริํ อายุสฺมิํ อกาเล มรติ, สตฺติหโต มหาราช อาพาธิโก ภิสกฺกํ อลภมาโน วิชฺชมาเนปิ อุตฺตริํ อายุสฺมิํ อกาเล มรติ, อิเม โข มหาราช สตฺต วิชฺชมาเนปิ อุตฺตริํ อายุสฺมิํ อกาเล มรนฺติ. ตตฺราปาหํ มหาราช เอกํเสน วทามิ.}

\prose{อฏฺฐวิเธน มหาราช สตฺตานํ กาลํกิริยา โหติ, วาต\-สมุฏฺฐาเนน ปิตฺต\-สมุฏฺฐาเนน เสมฺห\-สมุฏฺฐาเนน สนฺนิปาติเกน อุตุวิปริณาเมน วิสม\-ปริหาเรน โอปกฺกมิเกน กมฺมวิปาเกน มหาราช สตฺตานํ กาลํกิริยา โหติ. ตตฺร มหาราช ยทิทํ กมฺมวิปาเกน กาลํกิริยา, สา เยว ตตฺถ สามยิกา\csromansharedfootnote{8b9f971bb7a7b998}{สามายิกา (ก)} กาลํกิริยา, อวเสสา อสามหิกา กาลํกิริยาติ. ภวติ จ\csromandash{}}

\setlength{\csromangathapagewidth}{0pt}
\setlength{\csromangathawakwidth}{0pt}
\csromangathameasuregroup{“ชิฆจฺฉาย ปิปาสาย, \\ อคฺคิอุทกสตฺตีหิ, \\ วาตปิตฺเตน เสเมฺหน, \\ วิสโมปกฺกมกมฺเมหิ,}{\csromangathabat{“ชิฆจฺฉาย ปิปาสาย,}{อหิทฏฺฐา\textsuperscript{1} วิเสน จ.} \\ \csromangathabat{อคฺคิอุทกสตฺตีหิ,}{อกาเล ตตฺถ มียติ.} \\ \csromangathabat{วาตปิตฺเตน เสเมฺหน,}{สนฺนิปาเตนุตูหิ จ.} \\ \csromangathabat{วิสโมปกฺกมกมฺเมหิ,}{อกาเล ตตฺถ มียตี”ติ.}}
\csromangathasetleft{“ชิฆจฺฉาย ปิปาสาย, \\ อคฺคิอุทกสตฺตีหิ, \\ วาตปิตฺเตน เสเมฺหน, \\ วิสโมปกฺกมกมฺเมหิ,}
\csromangathagroup{\csromangathastanza{\csromangathabat{“ชิฆจฺฉาย ปิปาสาย,}{อหิทฏฺฐา\csromansharedfootnote{8ee412518de62684}{อหิทฏฺโฐ (สี), อหินา ทฏฺโฐ (อิ)} วิเสน จ.} \\ \csromangathabat{อคฺคิอุทกสตฺตีหิ,}{อกาเล ตตฺถ มียติ.}}\csromangathastanza{\csromangathabat{วาตปิตฺเตน เสเมฺหน,}{สนฺนิปาเตนุตูหิ จ.} \\ \csromangathabat{วิสโมปกฺกมกมฺเมหิ,}{อกาเล ตตฺถ มียตี”ติ.}}}

\proseitem{5}{\csromanfolio{292}เกจิ มหาราช สตฺตา ปุพฺเพ กเตน เตน เตน อกุสล\-กมฺมวิปาเกน มรนฺติ. อิธ มหาราช โย ปุพฺเพ ปเร ชิฆจฺฉาย มาเรติ, โส พหูนิ วสฺส\-สตสหสฺสานิ ชิฆจฺฉาย ปริปีฬิโต ฉาโต ปริกิลนฺโต สุกฺข\-มิลาต\-หทโย พุภุกฺขิโต\csromansharedfootnote{e1f9415f69cc774f}{สุกฺขิโต (สี, อิ, ก)} วิสุกฺขิโต ฌายนฺโต อพฺภนฺตรํ ปริฑยฺหนฺโต ชิฆจฺฉาย เยว มรติ ทหโรปิ มชฺฌิโมปิ มหลฺลโกปิ, อิทมฺปิ ตสฺส สามยิก\-มรณํ.}

\prose{โย ปุพฺเพ ปเร ปิปาสาย มาเรติ, โส พหูนิ วสฺส\-สตสหสฺสานิ เปโต หุตฺวา นิชฺฌาม\-ตณฺหิโก สมาโน ลูโข กิโส ปริสุกฺขิตหทโย ปิปาสาย เยว มรติ ทหโรปิ มชฺฌิโมปิ มหลฺลโกปิ, อิทมฺปิ ตสฺส สามยิก\-มรณํ.}

\prose{โย ปุพฺเพ ปเร อหินา ฑํสาเปตฺวา มาเรติ, โส พหูนิ วสฺส\-สตสหสฺสานิ อชคร\-มุเขเนว อชครมุขํ กณฺหสปฺป\-มุเขเนว กณฺหสปฺปมุขํ ปริวตฺติตฺวา เตหิ ขายิตขายิโต อหีหิ ทฏฺโฐ เยว มรติ ทหโรปิ มชฺฌิโมปิ มหลฺลโกปิ, อิทมฺปิ ตสฺส สามยิก\-มรณํ.}

\prose{โย ปุพฺเพ ปเร วิสํ ทตฺวา มาเรติ, โส พหูนิ วสฺส\-สตสหสฺสานิ ฑยฺหนฺเตหิ องฺคปจฺจงฺเคหิ ภิชฺชมาเนน สรีเรน กุณปคนฺธํ วายนฺโต วิเสเนว มรติ ทหโรปิ มชฺฌิโมปิ มหลฺลโกปิ, อิทมฺปิ ตสฺส สามยิก\-มรณํ.}

\prose{โย พุพฺเพ ปเร อคฺคินา มาเรติ, โส พหูนิ วสฺส\-สตสหสฺสานิ องฺคารปพฺพเตเนว องฺคารปพฺพตํ ยมวิสเยเนว ยมวิสยํ ปริวตฺติตฺวา ทฑฺฒ\-วิท\-ฑฺฒ\-คตฺโต อคฺคินา เยว มรติ ทหโรปิ มชฺฌิโมปิ มหลฺลโกปิ, อิทมฺปิ ตสฺส สามยิก\-มรณํ.}

\prose{โย ปุพฺเพ ปเร อุทเกน มาเรติ, โส พหูนิ วสฺส\-สตสหสฺสานิ หตวิลุตฺต\-ภคฺค\-ทุพฺพล\-คตฺโต ขุพฺภิตจิตฺโต\csromansharedfootnote{8d3008c92b67335e}{ขุภิตจิตฺโต (สี, อิ)} อุทเกเนว\csromansharedfootnote{01e20b35cdd35bf9}{อุทเก เยว (พหูสุ)} มรติ ทหโรปิ มชฺฌิโมปิ มหลฺลโกปิ, อิทมฺปิ ตสฺส สามยิก\-มรณํ.}

\prose{โย ปุพฺเพ ปเร สตฺติยา มาเรติ, โส พหูนิ วสฺส\-สตสหสฺสานิ ฉินฺนภินฺนโกฏฺฏิตวิโกฏฺฏิโต สตฺติ\-มุข\-สมาหโต สตฺติยา เยว มรติ ทหโรปิ มชฺฌิโมปิ มหลฺลโกปิ, อิทมฺปิ ตสฺส สามยิก\-มรณํ.}

\prose{\csromanfolio{293}ภนฺเต นาคเสน อกาเล มรณํ อตฺถีติ ยํ วเทสิ, อิงฺฆ เม ตฺวํ ตตฺถ การณํ อติทิสาติ. ยถา มหาราช มหติ\-มหา\-อคฺคิกฺขนฺโธ อาทินฺนติ -ณกฏฺฐสาขาปลาโส ปริยาทินฺน\-ภกฺโข อุปาทาน\-สงฺขยา นิพฺพายติ, โส อคฺคิ วุจฺจติ “อนีติโก อนุปทฺทโว สมเย นิพฺพุโต นามา”ติ, เอวเมว โข มหาราช โย โกจิ พหูนิ ทิวส\-สหสฺสานิ ชีวิตฺวา ชราชิณฺโณ อายุกฺขยา อนีติโก อนุปทฺทโว มรติ, โส วุจฺจติ “สมเย มรณมุปคโต”ติ.}

\prose{ยถา วา ปน มหาราช มหติ\-มหา\-อคฺคิกฺขนฺโธ อาทินฺน\-ติณกฏฺฐ\-สาขา\-ปลาโส อสฺส, ตํ อปริยาทินฺเน เยว ติณกฏฺฐ\-สาขา\-ปลาเส มหติ\-มหาเมโฆ อภิปฺปวสฺสิตฺวา นิพฺพาเปยฺย, อปิ นุ โข มหาราช มหา\-อคฺคิกฺขนฺโธ สมเย นิพฺพุโต นาม โหตีติ. น หิ ภนฺเตติ. กิสฺส ปน โส มหาราช ปจฺฉิโม อคฺคิกฺขนฺโธ ปุริมเกน อคฺคิกฺขนฺเธน สมสมคติโก นาโหสีติ. อาคนฺตุเกน ภนฺเต เมเฆน ปฏิปีฬิโต โส อคฺคิกฺขนฺโธ อสมเย นิพฺพุโตติ. เอวเมว โข มหาราช โย โกจิ อกาเล มรติ, โส อาคนฺตุเกน โรเคน ปฏิปีฬิโต วาต\-สมุฏฺฐาเนน วา ปิตฺต\-สมุฏฺฐาเนน วา เสมฺห\-สมุฏฺฐาเนน วา สนฺนิปาติเกน วา อุตุปริณาม\-เชน วา วิสม\-ปริหาร\-เชน วา โอปกฺกมิเกน วา ชิฆจฺฉาย วา ปิปาสาย วา สปฺปทฏฺเฐน วา วิสมาสิเตน วา อคฺคินา วา อุทเกน วา สตฺติ\-เวค\-ปฺปฏิปีฬิโต วา อกาเล มรติ, อิทเมตฺถ มหาราช การณํ, เยน การเณน อกาเล มรณํ อตฺถิ.}

\prose{ยถา วา ปน มหาราช คคเน มหติ\-มหา\-วลาหโก อุฏฺฐหิตฺวา นินฺนญฺจ ถลญฺจ ปริปูรยนฺโต อภิวสฺสติ, โส วุจฺจติ “เมโฆ อนีติโก อนุปทฺทโว วสฺสตี”ติ. เอวเมว โข มหาราช โย โกจิ จิรํ ชีวิตฺวา ชราชิณฺโณ อายุกฺขยา อนีติโก อนุปทฺทโว มรติ, โส วุจฺจติ “สมเย มรณมุปคโต”ติ.}

\prose{ยถา วา ปน มหาราช คคเน มหติ\-มหา\-วลาหโก อุฏฺฐหิตฺวา อนฺตราเยว มหตา วาเตน อพฺภตฺถํ คจฺเฉยฺย, อปิ นุ โข โส มหาราช มหาวลาหโก สมเย วิคโต นาม โหตีติ. น หิ ภนฺเตติ. กิสฺส ปน โส มหาราช ปจฺฉิโม วลาหโก ปุริเมน วลาหเกน สมสมคติโก นาโหสีติ. อา คนฺตุเกน ภนฺเต วาเตน ปฏิปีฬิโต โส วลาหโก อสมยปฺปตฺโต เยว}

\prose{\csromanfolio{294}วิคโตติ. เอวเมว โข มหาราช โย โกจิ อกาเล มรติ, โส อาคนฺตุเกน โรเคน ปฏิปีฬิโต วาต\-สมุฏฺฐาเนน วา ฯเปฯ สตฺติ\-เวค\-ปฺปฏิปีฬิโต วา อกาเล มรติ, อิทเมตฺถ มหาราช การณํ, เยน การเณน อกาเล มรณํ อตฺถีติ.}

\prose{ยถา วา ปน มหาราช พลวา อาสีวิโส กุปิโต กิญฺจิเทว ปุริสํ ฑํเสยฺย, ตสฺส ตํ วิสํ อนีติกํ อนุปทฺทวํ มรณํ ปาเปยฺย, ตํ วิสํ วุจฺจติ “อนีติกม\-นุปทฺทวํ โกฏิคตนฺ”ติ. เอวเมว โข มหาราช โย โกจิ จิรํ ชีวิตฺวา ชราชิณฺโณ อายุกฺขยา อนีติโก อนุปทฺทโว มรติ, โส วุจฺจติ “อนีติโก อนุปทฺทโว ชีวิตโกฏิคโต สามยิกํ มรณมุปคโต”ติ.}

\prose{ยถา วา ปน มหาราช พลวตา อาสีวิเสน ทฏฺฐสฺส อนฺตราเยว อาหิตุณฺฑิโก อคทํ ทตฺวา อวิสํ กเรยฺย, อปิ นุ โข ตํ มหาราช วิสํ สมเย วิคตํ นาม โหตีติ. น หิ ภนฺเตติ. กิสฺส ปน ตํ มหาราช ปจฺฉิมํ วิสํ ปุริมเกน วิเสน สมสม\-คติกํ นาโหสีติ. อาคนฺตุเกน ภนฺเต อคเทน ปฏิปีฬิตํ วิสํ อโกฏิคตํ เยว วิคตนฺติ. เอวเมว โข มหาราช โย โกจิ อกาเล มรติ, โส อาคนฺตุเกน โรเคน ปฏิปีฬิโต วาต\-สมุฏฺฐาเนน วา ฯเปฯ สตฺติ\-เวค\-ปฺปฏิปีฬิโต วา อกาเล มรติ, อิทเมตฺถ มหาราช การณํ, เยน การเณน อกาเล มรณํ อตฺถีติ.}

\prose{ยถา วา ปน มหาราช อิสฺสาโส สรํ ปาเตยฺย, สเจ โส สโร ยถาคติ\-คมน\-ปถ\-มตฺถกํ คจฺฉติ, โส สโร วุจฺจติ “อนีติโก อนุปทฺทโว ยถาคติมนปถมตฺถกํ คโต นามา”ติ. เอวเมว โข มหาราช โย โกจิ จิรํ ชีวิตฺวา ชราชิณฺโณ อายุกฺขยา อนีติโก อนุปทฺทโว มรติ, โส วุจฺจติ “อนีติโก อนุปทฺทโว สมเย มรณมุปคโต”ติ.}

\prose{ยถา วา ปน มหาราช อิสฺสาโส สรํ ปาเตยฺย, ตสฺส ตํ สรํ ตสฺมิํ เยว ขเณ โกจิ คเณฺหยฺย, อปิ นุ โข โส มหาราช สโร ยถาคติ\-คมน\-ปถ\-มตฺถกํ คโต นาม โหตีติ. น หิ ภนฺเตติ. กิสฺส ปน โส มหาราช ปจฺฉิโม สโร ปุริมเกน สเรน สมสมคติโก นาโหสีติ. อาคนฺตุเกน ภนฺเต คหเณน ตสฺส \csromanfolio{295}สรสฺส คมนํ อุปจฺฉินฺนนฺติ. เอวเมว โข มหาราช โย โกจิ อกาเล มรติ, โส อาคนฺตุเกน โรเคน ปฏิปีฬิโต วาต\-สมุฏฺฐาเนน วา ฯเปฯ สตฺติ\-เวค\-ปฺปฏิปีฬิโต วา อกาเล มรติ, อิทเมตฺถ มหาราช การณํ, เยน การเณน อกาเล มรณํ อตฺถีติ.}

\proseitem{5}{ยถา วา ปน มหาราช โย โกจิ โลหมยํ ภาชนํ อาโกเฏยฺย, ตสฺส อาโกฏเนน สทฺโท นิพฺพตฺติตฺวา ยถาคติ\-คมน\-ปถ\-มตฺถกํ คจฺฉติ, โส สทฺโท วุจฺจติ “อนีติโก อนุปทฺทโว ยถาคติ\-คมน\-ปถ\-มตฺถกํ คโต นามา”ติ. เอวเมว โข มหาราช โย โกจิ พหูนิ ทิวส\-สหสฺสานิ ชีวิตฺวา ชราชิณฺโณ อายุกฺขยา อนีติโก อนุปทฺทโว มรติ, โส วุจฺจติ “อนีติโก อนุปทฺทโว สมเย มรณมุปาคโต”ติ.}

\prose{ยถา วา ปน มหาราช โย โกจิ โลหมยํ ภาชนํ อาโกเฏยฺย, ตสฺส อาโกฏเนน สทฺโท นิพฺพตฺเตยฺย, นิพฺพตฺเต สทฺเท อทูรคเต โกจิ อามเสยฺย, สห อามสเนน สทฺโท นิรุชฺเฌยฺย, อปิ นุ โข โส มหาราช สทฺโท ยถาคติ\-คมน\-ปถ\-มตฺถกํ คโต นาม โหตีติ. น หิ ภนฺเตติ. กิสฺส ปน มหาราช ปจฺฉิโม สทฺโท ปุริมเกน สทฺเทน สมสมคติโก นาโหสีติ. อาคนฺตุเกน ภนฺเต อามสเนน โส สทฺโท อุปรโตติ. เอวเมว โข มหาราช โย โกจิ อกาเล มรติ, โส อาคนฺตุเกน โรเคน ปฏิปีฬิโต วาต\-สมุฏฺฐาเนน วา ฯเปฯ สตฺติ\-เวค\-ปฺปฏิปีฬิโต วา อกาเล มรติ, อิทเมตฺถ มหาราช การณํ, เยน การเณน อกาเล มรณํ อตฺถีติ.}

\prose{ยถา วา ปน มหาราช เขตฺเต สุวิรูฬฺหํ ธญฺญพีชํ สมฺมา ปวตฺตมาเนน วสฺเสน โอตต\-วิตต\-อากิณฺณ\-พหุผลํ หุตฺวา\csromansharedfootnote{f865033bc2bfb885}{อุฏฺฐิตอากิณฺณพหุผลํ ภวิตฺวา (สฺยา)} สสฺสุ\-ฏฺฐาน\-สมยํ ปาปุณาติ, ตํ ธญฺญํ วุจฺจติ “อนีติกม\-นุปทฺทวํ สมย\-สมฺปตฺตํ นาม โหตี”ติ. เอวเมว โข มหาราช โย โกจิ พหูนิ ทิวส\-สหสฺสานิ ชีวิตฺวา ชราชิณฺโณ อายุกฺขยา อนีติโก อนุปทฺทโว มรติ, โส วุจฺจติ “อนีติโก อนุปทฺทโว สมเย มรณมุปคโต”ติ.}

\prose{ยถา วา ปน มหาราช เขตฺต สุวิรูฬํ ธญฺญพีชํ อุทเกน วิกลํ มเรยฺย, อปิ นุ โข ตํ มหาราช ธญฺญํ สมย\-สมฺปตฺตํ นาม โหตีติ. \csromanfolio{296}น หิ ภนฺเตติ. กิสฺส ปน ตํ มหาราช ปจฺฉิมํ ธญฺญํ ปุริมเกน ธญฺเญน สมสม\-คติกํ นาโหสีติ. อาคนฺตุเกน ภนฺเต อุเณฺหน ปฏิปีฬิตํ ตํ ธญฺญํ มตนฺติ. เอวเมว โข มหาราช โย โกจิ อกาเล มรติ, โส อาคนฺตุเกน โรเคน ปฏิปีฬิโต วาต\-สมุฏฺฐาเนน วา ฯเปฯ สตฺติ\-เวค\-ปฺปฏิปีฬิโต วา อกาเล มรติ, อิทเมตฺถ มหาราช การณํ, เยน การเณน อกาเล มรณํ อตฺถีติ.}

\prose{สุตปุพฺพํ ปน ตยา มหาราช “สมฺปนฺน\-ตรุณ\-สสฺสํ กิมโย อุฏฺฐหิตฺวา สมูลํ นาเสนฺตี”ติ. สุตปุพฺพญฺเจว ตํ ภนฺเต อเมฺหหิ ทิฏฺฐปุพฺพ\-ญฺจาติ. กินฺนุ โข ตํ มหาราช สสฺสํ กาเล นฏฺฐํ, อุทาหุ อกาเล นฏฺฐนฺติ. อกาเล ภนฺเต, ยทิ โข ตํ ภนฺเต สสฺสํ กิมโย น ขาเทยฺยุํ, สสฺสุ\-ทฺธรณ\-สมยํ ปาปุเณยฺยาติ. กิํ ปน มหาราช อาคนฺตุเกน อุปฆาเตน สสฺสํ วินสฺสติ, นิรุปฆาตํ สสฺสํ สสฺสุ\-ทฺธรณ\-สมยํ ปาปุณาตีติ. อาม ภนฺเตติ. เอวเมว โข มหาราช โย โกจิ อกาเล มรติ, โส อาคนฺตุเกน โรเคน ปฏิปีฬิโต วาต\-สมุฏฺฐาเนน วา ฯเปฯ สตฺติ\-เวค\-ปฺปฏิปีฬิโต วา มรติ, อิทเมตฺถ มหาราช การณํ, เยน การเณน อกาเล มรณํ อตฺถีติ.}

\prose{สุตปุพฺพํ ปน ตยา มหาราช “สมฺปนฺเน สสฺเส ผลภาร\-น\-มิเต มญฺชริตปตฺเต กรกวสฺสํ นาม วสฺสชาติ นิปติตฺวา วินาเสติ อผลํ กโรตี”ติ. สุตปุพฺพญฺเจว ตํ ภนฺเต อเมฺหหิ ทิฏฺฐปุพฺพ\-ญฺจาติ. อปิ นุ โข ตํ มหาราช สสฺสํ กาเล นฏฺฐํ, อุทาหุ อกาเล นฏฺฐนฺติ. อกาเล ภนฺเต, ยทิ โข ตํ ภนฺเต สสฺสํ กรกวสฺสํ น วสฺเสยฺย สสฺสุ\-ทฺธรณ\-สมยํ อาปุเณยฺยาติ. กิํ ปน มหาราช อาคนฺตุเกน อุปฆาเตน สสฺสํ วินสฺสติ, นิรุปฆาตํ สสฺสํ สสฺสุ\-ทฺธรณ\-สมยํ ปาปุณาตีติ. อาม ภนฺเตติ. เอวเมว โข มหาราช โย โกจิ อกาเล มรติ, โส อาคนฺตุเกน โรเคน ปฏิปีฬิโต วาต\-สมุฏฺฐาเนน วา ปิตฺต\-สมุฏฺฐาเนน วา เสมฺห\-สมุฏฺฐาเนน วา สนฺนิปาติเกน วา อุตุปริณาม\-เชน วา วิสม\-ปริหาร\-เชน วา โอปกฺกมิเกน วา ชิฆจฺฉาย วา ปิปาสาย วา สปฺปทฏฺเฐน วา วิสมาสิเตน วา อคฺคินา วา อุทเกน วา สตฺติ\-เวค\-ปฺปฏิปีฬิโต วา อกาเล มรติ. ยทิ ปน อาคนฺตุเกน โรเคน ปฏิปีฬิโต น ภเวยฺย, สมเยว มรณํ ปาปุเณยฺย, อิทเมตฺถ มหาราช การณํ, เยน การเณน อกาเล มรณํ อตฺถีติ.}

\prosewithcloserruled{\csromanfolio{297}อจฺฉริยํ ภนฺเต นาคเสน, อพฺภุตํ ภนฺเต นาคเสน, สุทสฺสิตํ การณํ, สุทสฺสิตํ โอปมฺมํ อกาเล มรณสฺส ปริทีปนาย, “อตฺถิ อกาเล มรณนฺ”ติ อุตฺตานีกตํ ปากฏํ กตํ วิภูตํ กตํ, อจิตฺตวิกฺขิตฺตโกปิ ภนฺเต นาคเสน มนุโช เอกเมเกนปิ ตาว โอปมฺเมน นิฏฺฐํ คจฺเฉยฺย “อตฺถิ อกาเล มรณนฺ”ติ. กิํ ปน มนุโช สเจตโน, ปฐโม\-ปมฺเมเน\-วาหํ ภนฺเต สญฺญตฺโต “อตฺถิ อกาเล มรณนฺ”ติ, อปิ จ อปราปรํ นิพฺพาหนํ โสตุกาโม น สมฺปฏิจฺฉินฺติ.}{อกาลมรณ\-ปโญฺห ฉฏฺโฐ.}
\csromanmatikaanchor{mtk.191}
\csromantocmarkref{subsubsection}{7. เจติยปาฏิหาริยปญฺห}{mtk.191}
\bookmark[dest={mtk.191},level=3]{7. เจติยปาฏิหาริยปญฺห}
\csromanheader{3}{7. \textbf{เจติย}\-\textbf{ปาฏิหาริย}\-\textbf{ปญฺห}}
\proseitem{7}{ภนฺเต นาคเสน สพฺเพสํ ปรินิพฺพุตานํ เจติเย ปาฏิหีรํ โหติ, อุทาหุ เอกจฺจานํ เยว โหตีติ. เอกจฺจานํ มหาราช โหติ, เอกจฺจานํ น โหตีติ. กตเมสํ ภนฺเต โหติ, กตเมสํ น โหตีติ. ติณฺณนฺนํ มหาราช อญฺญตรสฺส อธิฏฺฐานา ปรินิพฺพุตสฺส เจติเย ปาฏิหีรํ โหติ. กตเมสํ ติณฺณนฺนํ, อิธ มหาราช อรหา เทวมนุสฺสานํ อนุกมฺปาย ติฏฺฐนฺโตว อธิฏฺฐาติ “เอวํนาม เจติเย ปาฏิหีรํ โหตู”ติ, ตสฺส อธิฏฺฐาน\-วเสน เจติเย ปาฏิหีรํ โหติ, เอวํ อรหโต อธิฏฺฐาน\-วเสน ปรินิพฺพุตสฺส เจติเย ปาฏิหีรํ โหติ.}

\prose{ปุน จปรํ มหาราช เทวตา มนุสฺสานํ อนุกมฺปาย ปรินิพฺพุตสฺส เจติเย ปาฏิหีรํ ทสฺเสนฺติ “อิมินา ปาฏิหีเรน สทฺธมฺโม นิจฺจสมฺปคฺคหิโต ภวิสฺสติ, มนุสฺสา จ ปสนฺนา กุสเลน อภิวฑฺฒิสฺสนฺตี”ติ, เอวํ เทวตานํ อธิฏฺฐาน\-วเสน ปรินิพฺพุตสฺส เจติเย ปาฏิหีรํ โหติ.}

\prose{ปุน จปรํ มหาราช อิตฺถี วา ปุริโส วา สทฺโธ ปสนฺโน ปณฺฑิโต พฺยตฺโต เมธาวี พุทฺธิสมฺปนฺโน โยนิโส จินฺตยิตฺวา คนฺธํ วา มาลํ วา ทุสฺสํ วา อญฺญตรํ วา กิญฺจิ อธิฏฺฐหิตฺวา เจติเย อุกฺขิปติ “เอวํนาม โหตู”ติ, ตสฺสปิ อธิฏฺฐาน\-วเสน ปรินิพฺพุตสฺส เจติเย ปาฏิหีรํ โหติ, เอวํ มนุสฺสานํ อธิฏฺฐาน\-วเสน ปรินิพฺพุตสฺส เจติเย ปาฏิหีรํ โหติ.}

\prose{\csromanfolio{298}อิเมสํ โข มหาราช ติณฺณนฺนํ อญฺญตรสฺส อธิฏฺฐาน\-วเสน ปรินิพฺพุตสฺส เจติเย ปาฏิหีรํ โหติ.}

\prosewithcloserruled{ยทิ มหาราช เตสํ อธิฏฺฐานํ น โหติ, ขีณาสวสฺสปิ ฉฬภิญฺญสฺส เจโต\-วสิ\-ปฺปตฺตสฺส เจติเย ปาฏิหีรํ น โหติ, อสติปิ มหาราช ปาฏิหีเร จริตํ ทิสฺวา สุปริสุทฺธํ โอกปฺเปตพฺพํ นิฏฺฐํ คนฺตพฺพํ สทฺทหิตพฺพํ “สุปรินิพฺพุโต อยํ พุทฺธปุตฺโต”ติ. สาธุ ภนฺเต นาคเสน, เอวเมตํ ตถา สมฺปฏิจฺฉามีติ.}{เจติย\-ปาฏิหาริย\-ปโญฺห สตฺตโม.}
\csromanmatikaanchor{mtk.192}
\csromantocmarkref{subsubsection}{8. ธมฺมาภิสมยปญฺห}{mtk.192}
\bookmark[dest={mtk.192},level=3]{8. ธมฺมาภิสมยปญฺห}
\csromanheader{3}{8. \textbf{ธมฺมา}\-\textbf{ภิสมยปญฺห}}
\proseitem{8}{ภนฺเต นาคเสน เย เต สมฺมา ปฏิปชฺชนฺติ, เตสํ สพฺเพสํ เยว ธมฺมา\-ภิสมโย โหติ, อุทาหุ กสฺสจิ น โหตีติ. กสฺสจิ มหาราช โหติ, กสฺสจิ น โหตีติ. กสฺส ภนฺเต โหติ, กสฺส น โหตีติ. ติรจฺฉาน\-คตสฺส มหาราช สุปฺปฏิปนฺนสฺสาปิ ธมฺมา\-ภิสมโย น โหติ, เปตฺติวิสยู\-ปปนฺนสฺส. มิจฺฉา\-ทิฏฺฐิกสฺส. กุหกสฺส. มาตุฆาตกสฺส. ปิตุฆาตกสฺส. อรหนฺต\-ฆาตกสฺส. สํฆเภทกสฺส. โลหิตุ\-ปฺปาทกสฺส. เถยฺย\-สํวาสกสฺส. ติตฺถิย\-ปกฺกนฺตสฺส. ภิกฺขุนิ\-ทูสกสฺส. เตรสนฺนํ ครุกาปตฺตีนํ อญฺญตรํ อาปชฺชิตฺวา อวุฏฺฐิตสฺส. ปณฺฑกสฺส. อุภโต\-พฺยญฺชนกสฺส สุปฺปฏิปนฺนสฺสาปิ ธมฺมา\-ภิสมโย น โหติ. โยปิ มนุสฺส\-ทหรโก อูนก\-สตฺตวสฺสิโก, ตสฺส สุปฺปฏิปนฺนสฺสาปิ ธมฺมา\-ภิสมโย น โหติ. อิเมสํ โข มหาราช โสฬสนฺนํ ปุคฺคลานํ สุปฺปฏิปนฺนานมฺปิ ธมฺมา\-ภิสมโย น โหตีติ.}

\prose{ภนฺเต นาคเสน เย เต ปนฺนรส ปุคฺคลา วิรุทฺธา เยว, เตสํ ธมฺมา\-ภิสมโย โหตุ วา มา วา โหตุ, อถ เกน การเณน มนุสฺส\-ทหรกสฺส อูนกสตฺตวสฺเสกสฺส สุปฺปฏิปนฺนสฺสาปิ ธมฺมา\-ภิสมโย น โหติ, เอตฺถ ตาว ปโญฺห ภวติ “นนุ นาม ทหรกสฺส น ราโค โหติ, น โทโส โหติ, น โมโห โหติ, น มาโน โหติ, น มิจฺฉาทิฏฺฐิ โหติ, น อรติ โหติ, น กามวิตกฺโก โหติ, อมิสฺสิโต กิเลเสหิ, โส นาม ทหรโก ยุตฺโต จ ปตฺโต จ อรหติ จ จตฺตาริ สจฺจานิ เอก\-ปฺปฏิเวเธน ปฏิวิชฺฌิตุนฺ”ติ.}

\prose{\csromanfolio{299}ตญฺเญเวตฺถ มหาราช การณํ, เยนาหํ การเณน ภณามิ “อูนก\-สตฺตวสฺสิกสฺส สุปฺปฏิปนฺนสฺสาปิ ธมฺมา\-ภิสมโย น โหตี”ติ. ยทิ มหาราช อูนก\-สตฺตวสฺสิโก รชนีเย รชฺเชยฺย, ทุสฺสนีเย ทุสฺเสยฺย, โมหนีเย มุเยฺหยฺย, มทนีเย มชฺเชยฺย, ทิฏฺฐิํ วิชาเนยฺย, รติญฺจ อรติญฺจ วิชาเนยฺย, กุสลากุสลํ วิตกฺเกยฺย, ภเวยฺย ตสฺส ธมฺมา\-ภิสมโย, อปิ จ มหาราช อูนก\-สตฺตวสฺสิกสฺส จิตฺตํ อพลํ ทุพฺพลํ ปริตฺตํ อปฺปํ โถกํ มนฺทํ อวิภูตํ, อสงฺขตา นิพฺพานธาตุ ครุกา ภาริกา วิปุลา มหตี. อูนก\-สตฺตวสฺสิโก มหาราช เตน ทุพฺพเลน จิตฺเตน ปริตฺตเกน มนฺเทน อวิภูเตน น สกฺโกติ ครุกํ ภารุกํ วิปุลํ มหติํ อสงฺขตํ นิพฺพานธาตุํ ปฏิวิชฺฌิตุํ.}

\prose{ยถา มหาราช สิเนรุ\-ปพฺพตราชา ครุโก ภาริโก วิปุโล มหนฺโต, อปิ นุ โข ตํ มหาราช ปุริโส อตฺตโน ปากติเกน ถาม\-พลวีริเยน สกฺกุเณยฺย สิเนรุ\-ปพฺพตราชานํ อุทฺธริตุนฺติ. น หิ ภนฺเตติ. เกน การเณน มหาราชาติ. ทุพฺพลตฺตา ภนฺเต ปุริสสฺส, มหนฺตตฺตา สิเนรุ\-ปพฺพตราชาสฺสาติ. เอวเมว โข มหาราช อูนก\-สตฺตวสฺสิกสฺส จิตฺตํ อพลํ ทุพฺพลํ ปริตฺตํ อปฺปํ โถกํ มนฺทํ อวิภูตํ, อสงฺขตา นิพฺพานธาตุ ครุกา ภาริกา วิปุลา มหตี. อูนก\-สตฺตวสฺสิโก เตน ทุพฺพเลน จิตฺเตน ปริตฺเตน มนฺเทน อวิภูเตน น สกฺโกติ ครุกํ ภาริกํ วิปุลํ มหติํ อสงฺขตํ นิพฺพานธาตุํ ปฏิวิชฺฌิตุํ, เตน การเณน อูนก\-สตฺตวสฺสิกสฺส สุปฺปฏิปนฺนสฺสาปิ ธมฺมา\-ภิสมโย น โหติ.}

\prose{ยถา วา ปน มหาราช อยํ มหาปถวี ทีฆา อายตา ปุถุลา วิตฺถตา วิสาลา วิตฺถิณฺณา วิปุลา มหนฺตา, อปิ นุ โข ตํ มหาราช มหาปถวิํ สกฺกา ปริตฺตเกน อุทกพินฺทุเกน เตเมตฺวา อุทก\-จิกฺขลฺลํ กาตุนฺติ. น หิ ภนฺเตติ. เกน การเณน มหาราชาติ. ปริตฺตตฺตา ภนฺเต อุทกพินฺทุสฺส, มหนฺตตฺตา มหาปถวิยาติ. เอวเมว โข มหาราช อูนก\-สตฺตวสฺสิกสฺส จิตฺตํ อพลํ ทุพฺพลํ ปริตฺตํ อปฺปํ โถกํ มนฺทํ อวิภูตํ, อสงฺขตา นิพฺพานธาตุ ทีฆา อายตา ปุถุลา วิตฺถตา วิสาลา วิตฺถิณฺณา วิปุลา มหนฺตา. อูนก\-สตฺตวสฺสิโก เตน ทุพฺพเลน จิตฺเตน ปริตฺตเกน มนฺเทน อวิภูเตน น สกฺโกติ มหติํ อสงฺขตํ}

\prose{\csromanfolio{300}นิพฺพานธาตุํ ปฏิวิชฺฌิตุํ, เตน การเณน อูนก\-สตฺตวสฺสิกสฺส สุปฺปฏิปนฺนสฺสาปิ ธมฺมา\-ภิสมโย น โหติ.}

\prose{ยถา วา ปน มหาราช อพล\-ทุพฺพล\-ปริตฺต\-อปฺป\-โถก\-มนฺทคฺคิ ภเวยฺย, อปิ นุ โข มหาราช ตาวตเกน มนฺเทน อคฺคินา สกฺกา สเทวเก โลเก อนฺธการํ วิธมิตฺวา อาโลกํ ทสฺเสตุนฺติ. น หิ ภนฺเตติ. เกน การเณน มหาราชาติ. มนฺทตฺตา ภนฺเต อคฺคิสฺส, โลกสฺส มหนฺตตฺตาติ. เอวเมว โข มหาราช อูนก\-สตฺตวสฺสิกสฺส จิตฺตํ อพลํ ทุพฺพลํ ปริตฺตํ อปฺปํ โถกํ มนฺทํ อวิภูตํ, มหตา จ อวิชฺช\-นฺธกาเรน ปิหิตํ. ตสฺมา ทุกฺกรํ ญาณาโลกํ ทสฺสยิตุํ, เตน การเณน อูนก\-สตฺตวสฺสิกสฺส สุปฺปฏิปนฺนสฺสาปิ ธมฺมา\-ภิสมโย น โหติ.}

\prosewithcloserruled{ยถา วา ปน มหาราช อาตุโร กิโส อณุปริมิต\-กาโย สาลกกิมิ หตฺถินาคํ ติธา ปภินฺนํ นวายตํ ติวิตฺถตํ ทสปริณาหํ อฏฺฐรตนิกํ สกฏฺฐานมุ\-ปคตํ ทิสฺวา คิลิตุํ ปริกฑฺเฒยฺย, อปิ นุ โข โส มหาราช สาลกกิมิ สกฺกุเณยฺย ตํ หตฺถินาคํ คิลิตุนฺติ. น หิ ภนฺเตติ. เกน การเณน มหาราชาติ. ปริตฺตตฺตา ภนฺเต สาลกกิมิสฺส, มหนฺตตฺตา หตฺถินาคสฺสาติ. เอวเมว โข มหาราช อูนก\-สตฺตวสฺสิกสฺส จิตฺตํ อพลํ ทุพฺพลํ ปริตฺตํ อปฺปํ โถกํ มนฺทํ อวิภูตํ, มหตี อสงฺขตา นิพฺพานธาตุ. โส เตน ทุพฺพเลน จิตฺเตน ปริตฺตเกน มนฺเทน อวิภูเตน น สกฺโกติ มหติํ อสงฺขตํ นิพฺพานธาตุํ ปฏิวิชฺฌิตุํ, เตน การเณน อูนก\-สตฺตวสฺสิกสฺส สุปฺปฏิปนฺนสฺสาปิ ธมฺมา\-ภิสมโย น โหตีติ. สาธุ ภนฺเต นาคเสน, เอวเมตํ ตถา สมฺปฏิจฺฉามีติ.}{ธมฺมาภิสมย\-ปโญฺห อฏฺฐโม.}
\csromanmatikaanchor{mtk.193}
\csromantocmarkref{subsubsection}{9. เอกนฺตสุขนิพฺพานปญฺห}{mtk.193}
\bookmark[dest={mtk.193},level=3]{9. เอกนฺตสุขนิพฺพานปญฺห}
\csromanheader{3}{9. \textbf{เอกนฺตสุข}\-\textbf{นิพฺพาน}\-\textbf{ปญฺห}}
\proseitem{9}{ภนฺเต นาคเสน กิํ เอกนฺตสุขํ นิพฺพานํ, อุทาหุ ทุกฺเขน มิสฺสนฺติ. เอกนฺตสุขํ มหาราช นิพฺพานํ, ทุกฺเขน อมิสฺสนฺติ.}

\prose{น มยํ ตํ ภนฺเต นาคเสน วจนํ สทฺทหาม “เอกนฺตสุขํ นิพฺพานนฺ”ติ, เอวเมตฺถ มยํ ภนฺเต นาคเสน ปจฺเจม “นิพฺพานํ ทุกฺเขน มิสฺสนฺ”ติ, การณญฺเจตฺถ อุปลภาม “นิพฺพานํ ทุกฺเขน มิสฺสนฺ”ติ. กตมํ เอตฺถ \csromanfolio{301}การาณํ, เย เต ภนฺเต นาคเสน นิพฺพานํ ปริเยสนฺติ, เตสํ ทิสฺสติ กายสฺส จ จิตฺตสฺส จ อาตาโป ปริตาโป ฐานจงฺกม\-นิสชฺชา\-สย\-นาหาร\-ปริคฺคโห มิทฺธสฺส จ อุปโรโธ อายตนานญฺจ ปฏิปีฬนํ ธนธญฺญ\-ปิย\-ญาติมิตฺต\-ปฺปชหนํ. เย เกจิ โลเก สุขิตา สุขสมปฺปิตา, เต สพฺเพปิ ปญฺจหิ กามคุเณหิ อายตเน รเมนฺติ พฺรูเหนฺติ, มนาปิก\-มนาปิก\-พหุวิธ\-สุภนิมิตฺเตน รูเปน จกฺขุํ รเมนฺติ พฺรูเหนฺติ, มนาปิก\-มนาปิก\-คีตวาทิต\-พหุวิธ\-สุภนิมิตฺเตน สทฺเทน โสตํ รเมนฺติ พฺรูเหนฺติ, มนาปิก\-มนาปิก\-ปุปฺผ\-ผล\-ปตฺต\-ตจ\-มูล\-สาร\-พหุวิธ\-สุภนิมิตฺเตน คนฺเธน ฆานํ รเมนฺติ พฺรูเหนฺติ, มนาปิกมนาปิกขชฺชโภชฺชเลยฺยเปยฺย- สายนียพหุวิธสุภนิมิตฺเตน รเสน ชิวฺหํ รเมนฺติ พฺรูเหนฺติ, มนาปิก\-มนาปิก\-สณฺหสุขุม\-มุทุ\-มทฺทว\-พหุวิธ\-สุภนิมิตฺเตน ผสฺเสน กายํ รเมนฺติ พฺรูเหนฺติ, มนาปิกมนาปิกกลฺยาณปาปกสุภา- สุภพหุวิธวิตกฺกมนสิกาเรน มนํ รเมนฺติ พฺรูเหนฺติ. ตุเมฺห ตํ จกฺขุโสต\-ฆาน\-ชิวฺหา\-กายมโน\-พฺรูหนํ หนถ อุปหนถ, ฉินฺทถ อุปจฺฉินฺทถ, รุนฺธถ อุปรุนฺธถ. เตน กาโยปิ ปริตปติ, จิตฺตมฺปิ ปริตปติ, กาเย ปริตตฺเต กายิก\-ทุกฺขเวทนํ เวทิยติ, จิตฺเต ปริตตฺเต เจตสิก\-ทุกฺขเวทนํ เวทยติ. นนุ มาคณฺฑิโยปิ\csromansharedfootnote{203dab4a65fd9673}{มาคนฺทิโยปิ (สี, อิ)} ปริพฺพาชโก ภควนฺตํ ครหมาโน เอวมาห\csromansymbolfootnote{*}{ม 2. 169 ปิฏฺเฐ.}“ภูนหุ\csromansharedfootnote{c0a94df100613837}{ภูตหจฺโจ (อิ), ภูนหจฺโจ (ก)} สมโณ โคตโม”ติ. อิทเมตฺถ การณํ, เยนาหํ การเณน พฺรูมิ “นิพฺพานํ ทุกฺเขน มิสฺสนฺ”ติ.}

\prose{น หิ มหาราช นิพฺพานํ ทุกฺเขน มิสฺสํ, เอกนฺตสุขํ นิพฺพานํ. ยํ ปน ตฺวํ มหาราช พฺรูสิ “นิพฺพานํ ทุกฺขนฺ”ติ. เนตํ ทุกฺขํ นิพฺพานํ นาม, นิพฺพานสฺส ปน สจฺฉิกิริยาย ปุพฺพภาโค เอโส, นิพฺพาน\-ปริเยสนํ เอตํ, เอกนฺตสุขํ เยว มหาราช นิพฺพานํ, น ทุกฺเขน มิสฺสํ. เอตฺถ การณํ วทามิ, อตฺถิ มหาราช ราชูนํ รชฺชสุขํ นามาติ. อาม ภนฺเต อตฺถิ ราชูนํ รชฺชสุขนฺติ. อปิ นุ โข ตํ มหาราช รชฺชสุขํ ทุกฺเขน มิสฺสนฺติ. น หิ ภนฺเตติ. กิสฺส ปน เต มหาราช ราชาโน ปจฺจนฺเต กุปิเต เตสํ ปจฺจนฺต\-นิสฺสิตานํ ปฏิเสธาย อมจฺเจหิ ปริณายเกหิ ภเฏหิ พลตฺเถหิ ปริวุตา ปวาสํ คนฺตฺวา ฑํสมกส\-วาตาตป\-ปฏิปีฬิตา สมวิสเม ปริธาวนฺติ, มหายุทฺธญฺจ กโรนฺติ, ชีวิตสํสยญฺจ ปาปุณนฺตีติ. เนตํ ภนฺเต นาคเสน รชฺชสุขํ นาม, รชฺชสุขสฺส ปริเยสนาย ปุพฺพภาโค เอโส, ทุกฺเขน ภนฺเต นาคเสน \csromanfolio{302}ราชาโน รชฺชํ ปริเยสิตฺวา รชฺชสุขํ อนุภวนฺติ, เอวํ ภนฺเต นาคเสน รชฺชสุขํ ทุกฺเขน อมิสฺสํ, อญฺญํ ตํ รชฺชสุขํ, อญฺญํ ทุกฺขนฺติ. เอวเมว โข มหาราช เอกนฺตสุขํ นิพฺพานํ, น ทุกฺเขน มิสฺสํ. เย ปน ตํ นิพฺพานํ ปริเยสนฺติ, เต กายญฺจ จิตฺตญฺจ อาตาเปตฺวา ฐานจงฺกม\-นิสชฺชา\-สย\-นาหารํ ปริคฺคเหตฺวา มิทฺธํ อุปรุนฺธิตฺวา อายตนานิ ปฏิปีเฬตฺวา กายญฺจ ชีวิตญฺจ ปริจฺจชิตฺวา ทุกฺเขน นิพฺพานํ ปริเยสิตฺวา เอกนฺตสุขํ นิพฺพานํ อนุภวนฺติ, นิหต\-ปจฺจามิตฺตา วิย ราชาโน รชฺชสุขํ. เอวํ มหาราช เอกนฺตสุขํ นิพฺพานํ, น ทุกฺเขน มิสฺสํ, อญฺญํ นิพฺพานํ, อญฺญํ ทุกฺขนฺติ.}

\proseitem{5}{อปรมฺปิ มหาราช อุตฺตริํ การณํ สุโณหิ เอกนฺตสุขํ นิพฺพานํ, น ทุกฺเขน มิสฺสํ, อญฺญํ ทุกฺขํ, อญฺญํ นิพฺพานนฺติ, อตฺถิ มหาราช อาจริยานํ สิปฺปวนฺตานํ สิปฺปสุขํ นามาติ. อาม ภนฺเต อตฺถิ อาจริยานํ สิปฺปวนฺตานํ สิปฺปสุขนฺติ. อปิ นุ โข ตํ มหาราช สิปฺปสุขํ ทุกฺเขน มิสฺสนฺติ. น หิ ภนฺเตติ. กิสฺส ปน เต มหาราช อาจริยา\csromansharedfootnote{0dc61a19fd4f3eaf}{อิทํ ปทํ สีอิโปตฺถเกสุ นตฺถิ.} อาจริยานํ อภิวาทน\-ปจฺจุฏฺฐาเนน อุทกา\-หรณ\-ฆร\-สมฺมชฺชน\-ทนฺตกฏฺฐ\-มุโขทกา\-นุปฺปทาเนน อุจฺฉิฏฺฐปฏิคฺคหณอุจฺฉาทนนหาปนปาทปริกมฺเมน สกจิตฺตํ นิกฺขิปิตฺวา ปรจิตฺตา\-นุวตฺตเนน ทุกฺขเสยฺยาย วิสม\-โภชเนน กายํ อาตาเปนฺตีติ. เนตํ ภนฺเต นาคเสน สิปฺปสุขํ นาม, สิปฺป\-ปริเยสนาย ปุพฺพภาโค เอโส, ทุกฺเขน ภนฺเต นาคเสน อาจริยา สิปฺปํ ปริเยสิตฺวา สิปฺปสุขํ อนุภวนฺติ, เอวํ ภนฺเต นาคเสน สิปฺปสุขํ ทุกฺเขน อมิสฺสํ, อญฺญํ ตํ สิปฺปสุขํ, อญฺญํ ทุกฺขนฺติ. เอวเมว โข มหาราช เอกนฺตสุขํ นิพฺพานํ, น ทุกฺเขน มิสฺสํ. เย ปน ตํ นิพฺพานํ ปริเยสนฺติ, เต กายญฺจ จิตฺตญฺจ อาตาเปตฺวาฐานจงฺกมนิสชฺชาสยนาหารํ ปริคฺคเหตฺวา มิทฺธํ อุปรุนฺธิตฺวา อายตนานิ ปฏิปีเฬตฺวา กายญฺจ ชีวิตญฺจ ปริจฺจชิตฺวา ทุกฺเขน นิพฺพานํ ปริเยสิตฺวา เอกนฺตสุขํ นิพฺพานํ อนุภวนฺติ, อาจริยา วิย สิปฺปสุขํ. เอวํ มหาราช เอกนฺตสุขํ นิพฺพานํ, น ทุกฺเขน มิสฺสํ, อญฺญํ ทุกฺขํ, อญฺญํ นิพฺพานนฺติ. สาธุ ภนฺเต นาคเสน, เอวเมตํ ตถา สมฺปฏิจฺฉามีติ.}

\vspace{\tipitakaparskip}
\csromancenter{เอกนฺตสุข\-นิพฺพาน\-ปโญฺห นวโม.}
\csromanmatikaanchor{mtk.194}
\csromantocmarkref{subsubsection}{10. นิพฺพานรูปสณฺฐานปญฺห}{mtk.194}
\bookmark[dest={mtk.194},level=3]{10. นิพฺพานรูปสณฺฐานปญฺห}
\csromanmarkh{5. อนุมานปญฺห}\csromanmarkhii{10. นิพฺพาน\-รูป\-สณฺฐาน\-ปญฺห}\csromanheadernomark{3}{5. อนุมานปญฺห \textbf{10. นิพฺพาน}\-\textbf{รูป}\-\textbf{สณฺฐาน}\-\textbf{ปญฺห}}
\proseitem{10}{\csromanfolio{303}ภนฺเต นาคเสน “นิพฺพานํ นิพฺพานนฺ”ติ ยํ วเทสิ, สกฺกา ปน ตสฺส นิพฺพานสฺส รูปํ วา สณฺฐานํ วา วยํ วา ปมาณํ วา โอปมฺเมน วา การเณน วา เหตุนา วา นเยน วา อุปทสฺสยิ\-ตุนฺติ. อปฺปฏิภาคํ มหาราช นิพฺพานํ, น สกฺกา นิพฺพานสฺส รูปํ วา สณฺฐานํ วา วยํ วา ปมาณํ วา โอปมฺเมน วา การเณน วา เหตุนา วา นเยน วา อุปทสฺสยิ\-ตุนฺติ. เอตมฺปาหํ ภนฺเต นาคเสน น สมฺปฏิจฺฉามิ, ยํ อตฺถิธมฺมสฺส นิพฺพานสฺส รูปํ วา สณฺฐานํ วา วยํ วา ปมาณํ วา โอปมฺเมน วา การเณน วา เหตุนา วา นเยน วา อปญฺญาปนํ, การเณน มํ สญฺญาเปหีติ. โหตุ มหาราช การเณน ตํ สญฺญาเปสฺสามิ, อตฺถิ มหาราช มหาสมุทฺโท นามติ. อาม ภนฺเต อตฺเถโส มหาสมุทฺโทติ. สเจ ตํ มหาราช โกจิ เอวํ ปุจฺเฉยฺย “กิตฺตกํ มหาราช มหาสมุทฺเท อุทกํ, กติ ปน เต สตฺตา, เย มหาสมุทฺเท ปฏิวสนฺตี”ติ, เอวํ ปุฏฺโฐ ตฺวํ มหาราช กินฺติ ตสฺส พฺยากเรยฺยาสีติ. สเจ มํ ภนฺเต โกจิ เอวํ ปุจฺเฉยฺย “กิตฺตกํ มหาราช มหาสมุทฺเท อุทกํ, กติ ปน เต สตฺตา, เย มหาสมุทฺเท ปฏิวสนฺตี”ติ. ตมหํ ภนฺเต เอวํ วเทยฺยํ “อปุจฺฉิตพฺพํ มํ ตฺวํ อมฺโภ ปุริส ปุจฺฉสิ, เนสา ปุจฺฉา เกนจิ ปุจฺฉิตพฺพา, ฐปนีโย เอโส ปโญฺห. อวิภตฺโต โลกกฺขายิเกหิ มหาสมุทฺโท, น สกฺกา มหาสมุทฺเท อุทกํ ปริมินิตุํ สตฺตา วา เย ตตฺถ วาสมุปคตาติ เอวาหํ ภนฺเต ตสฺส ปฏิวจนํ ทเทยฺยนฺ”ติ.}

\prose{กิสฺส ปน ตฺวํ มหาราช อตฺถิธมฺเม มหาสมุทฺเท เอวํ ปฏิวจนํ ทเทยฺยาสิ, นนุ วิคเณตฺวา\csromansharedfootnote{1481ab848104aba5}{มินิตฺวา (ก)} ตสฺส อาจิกฺขิตพฺพํ “เอตฺตกํ มหาสมุทฺเท อุทกํ, เอตฺตกา จ สตฺตา มหาสมุทฺเท ปฏิวสนฺตี”ติ. น สกฺกา ภนฺเต อวิสโย เอโส ปโญฺหติ.}

\prose{ยถา มหาราช อตฺถิธมฺเม เยว มหาสมุทฺเท น สกฺกา อุทกํ ปริคเณตุํ\csromansharedfootnote{b7f4f1fe258f7dfb}{ปริมินิตุํ (ก)} สตฺตา วา เย ตตฺถ วาสมุปคตา. เอวเมว โข มหาราช อตฺถิ\-ธมฺมสฺเสว นิพฺพานสฺส น สกฺกา รูปํ วา สณฺฐานํ วา วยํ วา ปมาณํ วา โอปมฺเมน วา การเณน วา เหตุนา วา นเยน วา อุปทสฺสยิตุํ, วิคเณยฺย มหาราช อิทฺธิมา เจโตวสิปฺปตฺโต มหาสมุทฺเท อุทกํ \csromanfolio{304}ตตฺราสเย จ สตฺเต, น เตฺวว โส อิทฺธิมา เจโตวสิปฺปตฺโต สกฺกุเณยฺย นิพฺพานสฺส รูปํ วา สณฺฐานํ วา วยํ วา ปมาณํ วา โอปมฺเมน วา การเณน วา เหตุนา วา นเยน วา อุปทสฺสยิตุํ.}

\prose{อปรมฺปิ มหาราช อุตฺตริํ การณํ สุโณหิ, อตฺถิ\-ธมฺมสฺเสว นิพฺพานสฺส น สกฺกา รูปํ วา สณฺฐานํ วา วยํ วา ปมาณํ วา โอปมฺเมน วา การเณน วา เหตุนา วา นเยน วา อุปทสฺสยิ\-ตุนฺติ. อตฺถิ มหาราช เทเวสุ อรูปกายิกา นาม เทวาติ. อาม ภนฺเต สุยฺยติ “อตฺถิ เทเวสุ อรูปกายิกา นาม เทวา”ติ. สกฺกา ปน มหาราช เตสํ อรูปกายิกานํ เทวานํ รูปํ วา สณฺฐานํ วา วยํ วา ปมาณํ วา โอปมฺเมน วา การเณน วา เหตุนา วา นเยน วา อุปทสฺสยิ\-ตุนฺติ. น หิ ภนฺเตติ. เตน หิ มหาราช นตฺถิ อรูปกายิกา เทวาติ. อตฺถิ ภนฺเต อรูปกายิกา เทวา, น จ สกฺกา เตสํ รูปํ วา สณฺฐานํ วา วยํ วา ปมาณํ วา โอปมฺเมน วา การเณน วา เหตุนา วา นเยน วา อุปทสฺสยิ\-ตุนฺติ. ยถา มหาราช อตฺถิสตฺตานํ เยว อรูปกายิกานํ เทวานํ น สกฺกา รูปํ วา สณฺฐานํ วา วยํ วา ปมาณํ วา โอปมฺเมน วา การเณน วา เหตุนา วา นเยน วา อุปทสฺสยิตุํ. เอวเมว โข มหาราช อตฺถิ\-ธมฺมสฺเสว นิพฺพานสฺส น สกฺกา รูปํ วา สณฺฐานํ วา วยํ วา ปมาณํ วา โอปมฺเมน วา การเณน วา เหตุนา วา นเยน วา อุปทสฺสยิ\-ตุนฺติ.}

\prose{ภนฺเต นาคเสน โหตุ เอกนฺตสุขํ นิพฺพานํ, น จ สกฺกา ตสฺส รูปํ วา สณฺฐานํ วา วยํ วา ปมาณํ วา โอปมฺเมน วา การเณน วา เหตุนา วา นเยน วา อุปทสฺสยิตุํ. อตฺถิ ปน ภนฺเต นิพฺพานสฺส คุณํ อญฺเญหิ อนุปวิฏฺฐํ กิญฺจิ โอปมฺมนิทสฺสนมตฺตนฺติ. สรูปโต มหาราช นตฺถิ, คุณโต ปน สกฺกา กิญฺจิ โอปมฺม\-นิทสฺสนมตฺตํ อุปทสฺสยิ\-ตุนฺติ. สาธุ ภนฺเต นาคเสน, ยถาหํ ลภามิ นิพฺพานสฺส คุณโตปิ เอกเทส\-ปริทีปน\-มตฺตํ, ตถา สีฆํ พฺรูหิ, นิพฺพาเปหิ เม หทย\-ปริฬาหํ วินย สีตล\-มธุร\-วจน\-มาลุเตนาติ.}

\prose{ปทุมสฺส มหาราช เอโก คุโณ นิพฺพานํ อนุปวิฏฺโฐ, อุทกสฺส เทฺว คุณา, อคทสฺส ตโย คุณา, มหาสมุทฺทสฺส จตฺตาโร คุณา, โภชนสฺส ปญฺจ คุณา, อากาสสฺส ทส คุณา, มณิรตนสฺส ตโย \csromanfolio{305}คุณา, โลหิต\-จนฺทนสฺส ตโย คุณา, สปฺปิมณฺฑสฺส ตโย คุณา, คิริสิขรสฺส ปญฺจ คุณา นิพฺพานํ อนุปวิฏฺฐาติ.}

\prose{ภนฺเต นาคเสน “ปทุมสฺส เอโก คุโณ นิพฺพานํ อนุปวิฏฺโฐ”ติ ยํ วเทสิ, กตโม ปทุมสฺส เอโก คุโณ นิพฺพานํ อนุปวิฏฺโฐติ. ยถา มหาราช ปทุมํ อนุปลิตฺตํ อุทเกน, เอวเมว โข มหาราช นิพฺพานํ สพฺพกิเลเสหิ อนุปลิตฺตํ. อยํ มหาราช ปทุมสฺส เอโก คุโณ นิพฺพานํ อนุปวิฏฺโฐติ.}

\prose{ภนฺเต นาคเสน “อุทกสฺส เทฺว คุณา นิพฺพานํ อนุปวิฏฺฐา”ติ ยํ วเทสิ, กตเม อุทกสฺส เทฺว คุณา นิพฺพานํ อนุปวิฏฺฐาติ. ยถา มหาราช อุทกํ สีตลํ ปริฬาห\-นิพฺพาปนํ, เอวเมว โข มหาราช นิพฺพานํ สีตลํ สพฺพ\-กิเลสปริฬาห\-นิพฺพาปนํ. อยํ มหาราช อุทกสฺส ปฐโม คุโณ นิพฺพานํ อนุปวิฏฺโฐ. ปุน จปรํ มหาราช อุทกํ กิลนฺต\-ตสิต\-ปิปาสิต\-ฆมฺมาภิตตฺตานํ ชนปสุ\-ปชานํ ปิปาสา\-วินยนํ, เอวเมว โข มหาราช นิพฺพานํ กามตณฺหา\-ภวตณฺหา\-วิภวตณฺหา\-ปิปาสา\-วินยนํ. อยํ มหาราช อุทกสฺส ทุติโย คุโณ นิพฺพานํ อนุปวิฏฺโฐ. อิเม โข มหาราช อุทกสฺส เทฺว คุณา นิพฺพานํ อนุปวิฏฺฐาติ.}

\prose{ภนฺเต นาคเสน “อคทสฺส ตโย คุณา นิพฺพานํ อนุปวิฏฺฐา”ติ ยํ วเทสิ, กตเม อคทสฺส ตโย คุณา นิพฺพานํ อนุปวิฏฺฐาติ. ยถา มหาราช อคโท วิสปีฬิตานํ สตฺตานํ ปฏิสรณํ, เอวเมว โข มหาราช นิพฺพานํ กิเลส\-วิส\-ปีฬิตานํ สตฺตานํ ปฏิสรณํ. อยํ มหาราช อคทสฺส ปฐโม คุโณ นิพฺพานํ อนุปวิฏฺโฐ. ปุน จปรํ มหาราช อคโท โรคานํ อนฺตกโร, เอวเมว โข มหาราช นิพฺพานํ สพฺพทุกฺขานํ อนฺตกรํ. อยํ มหาราช อคทสฺส ทุติโย คุโณ นิพฺพานํ อนุปวิฏฺโฐ. ปุน จปรํ มหาราช อคโท อมตํ, เอวเมว โข มหาราช นิพฺพานํ อมตํ. อยํ มหาราช อคทสฺส ตติโย คุโณ นิพฺพานํ อนุปวิฏฺโฐ. อิเม โข มหาราช อคทสฺส ตโย คุณา นิพฺพานํ อนุปวิฏฺฐาติ.}

\prose{ภนฺเต นาคเสน “มหาสมุทฺทสฺส จตฺตาโร คุณา นิพฺพานํ อนุปวิฏฺฐา”ติ ยํ วเทสิ, กตเม มหาสมุทฺทสฺส จตฺตาโร คุณา นิพฺพานํ อนุปวิฏฺฐาติ. ยถา มหาราช มหาสมุทฺโท สุญฺโญ สพฺพกุณเปหิ, เอวเมว โข \csromanfolio{306}มหาราช นิพฺพานํ สุญฺญํ สพฺพกิเลส\-กุณเปหิ. อยํ มหาราช มหาสมุทฺทสฺส ปฐโม คุโณ นิพฺพานํ อนุปวิฏฺโฐ. ปุน จปรํ มหาราช มหาสมุทฺโท มหนฺโต อโนรปาโร, น ปริปูรติ สพฺพสวนฺตีหิ, เอวเมว โข มหาราช นิพฺพานํ มหนฺตํ อโนรปารํ, น ปูรติ สพฺพสตฺเตหิ. อยํ มหาราช มหาสมุทฺทสฺส ทุติโย คุโณ นิพฺพานํ อนุปวิฏฺโฐ. ปุน จปรํ มหาราช มหาสมุทฺโท มหนฺตานํ ภูตานํ อาวาโส, เอวเมว โข มหาราช นิพฺพานํ มหนฺตานํ อรหนฺตานํ วิมล\-ขีณาสวพล\-ปฺปตฺต\-วสีภูต\-มหาภูตานํ อาวาโส. อยํ มหาราช มหาสมุทฺทสฺส ตติโย คุโณ นิพฺพานํ อนุปวิฏฺโฐ. ปุน จปรํ มหาราช มหาสมุทฺโท อปริมิตวิวิธวิปุลวีจิปุปฺผสํกุสุมิโต, เอวเมว โข มหาราช นิพฺพานํ อปริมิตวิวิธวิปุลปริสุทฺธวิชฺชาวิมุตฺติปุปฺผสํกุสุมิตํ. อยํ มหาราช มหาสมุทฺทสฺส จตุตฺโถ คุโณ นิพฺพานํ อนุปวิฏฺโฐ. อิเม โข มหาราช มหาสมุทฺทสฺส จตฺตาโร คุณา นิพฺพานํ อนุปวิฏฺฐาติ.}

\proseitem{5}{ภนฺเต นาคเสน “โภชนสฺส ปญฺจ คุณา นิพฺพานํ อนุปวิฏฺฐา”ติ ยํ วเทสิ, กตเม โภชนสฺส ปญฺจ คุณา นิพฺพานํ อนุปวิฏฺฐาติ. ยถา มหาราช โภชนํ สพฺพสตฺตานํ อายุธารณํ, เอวเมว โข มหาราช นิพฺพานํ สจฺฉิกตํ ชรามรณ\-นาสนโต อายุธารณํ. อยํ มหาราช โภชนสฺส ปฐโม คุโณ นิพฺพานํ อนุปวิฏฺโฐ. ปุน จปรํ มหาราช โภชนํ สพฺพสตฺตานํ พลวฑฺฒนํ, เอวเมว โข มหาราช นิพฺพานํ สจฺฉิกตํ สพฺพสตฺตานํ อิทฺธิพล\-วฑฺฒนํ. อยํ มหาราช โภชนสฺส ทุติโย คุโณ นิพฺพานํ อนุปวิฏฺโฐ. ปุน จปรํ มหาราช โภชนํ สพฺพสตฺตานํ วณฺณชนนํ, เอวเมว โข มหาราช นิพฺพานํ สจฺฉิกตํ สพฺพสตฺตานํ คุณวณฺณ\-ชนนํ. อยํ มหาราช โภชนสฺส ตติโย คุโณ นิพฺพานํ อนุปวิฏฺโฐ. ปุน จปรํ มหาราช โภชนํ สพฺพสตฺตานํ ทรถ\-วูปสมนํ, เอวเมว โข มหาราช นิพฺพานํ สจฺฉิกตํ สพฺพสตฺตานํ สพฺพกิเลส\-ทรถ\-วูปสมนํ. อยํ มหาราช โภชนสฺส จตุตฺโถ คุโณ นิพฺพานํ อนุปวิฏฺโฐ. ปุน จปรํ มหาราช โภชนํ สพฺพสตฺตานํ ชิฆจฺฉาทุพฺพลฺย\-ปฏิวิโนทนํ, เอวเมว โข มหาราช นิพฺพานํ สจฺฉิกตํ สพฺพสตฺตานํ สพฺพทุกฺข\-ชิฆจฺฉาทุพฺพลฺย\-ปฏิวิโนทนํ. อยํ มหาราช โภชนสฺส ปญฺจโม คุโณ นิพฺพานํ อนุปวิฏฺโฐ. อิเม โข มหาราช โภชนสฺส ปญฺจ คุณา นิพฺพานํ อนุปวิฏฺฐาติ.}

\prose{\csromanfolio{307}ภนฺเต นาคเสน “อากาสสฺส ทส คุณา นิพฺพานํ อนุปวิฏฺฐา”ติ ยํ วเทสิ, กตเม อากาสสฺส ทส คุณา นิพฺพานํ อนุปวิฏฺฐาติ. ยถา มหาราช อากาโส น ชายติ, น ชียติ, น มียติ, น จวติ, น อุปฺปชฺชติ, ทุปฺปสโห, อโจราหรโณ, อนิสฺสิโต, วิหคคมโน, นิราวรโณ, อนนฺโต. เอวเมว โข มหาราช นิพฺพานํ น ชายติ, น ชียติ, น มียติ, น จวติ, น อุปฺปชฺชติ, ทุปฺปสหํ, อโจราหรณํ, อนิสฺสิตํ, อริยคมนํ, นิราวรณํ, อนนฺตํ. อิเม โข มหาราช อากาสสฺส ทส คุณา นิพฺพานํ อนุปวิฏฺฐาติ.}

\prose{ภนฺเต นาคเสน “มณิรตนสฺส ตโย คุโณ นิพฺพานํ อนุปวิฏฺฐา”ติ ยํ วเทสิ, กตเม มณิรตนสฺส ตโย คุณา นิพฺพานํ อนุปวิฏฺฐาติ. ยถา มหาราช มณิรตนํ กามททํ, เอวเมว โข มหาราช นิพฺพานํ กามททํ. อยํ มหาราช มณิรตนสฺส ปฐโม คุโณ นิพฺพานํ อนุปวิฏฺโฐ. ปุน จปรํ มหาราช มณิรตนํ หาสกรํ, เอวเมว โข มหาราช นิพฺพานํ หาสกรํ. อยํ มหาราช มณิรตนสฺส ทุติโย คุโณ นิพฺพานํ อนุปวิฏฺโฐ. ปุน จปรํ มหาราช มณิรตนํ อุชฺโชตตฺตกรํ, เอวเมว โข มหาราช นิพฺพานํ อุชฺโชตตฺตกรํ\csromansharedfootnote{a901afe4e5740232}{อุชฺโชตตฺถกรํ (สี, อิ), อุชฺโชติตตฺถกรํ (สฺยา)}. อยํ มหาราช มณิรตนสฺส ตติโย คุโณ นิพฺพานํ อนุปวิฏฺโฐ. อิเม โข มหาราช มณิรตนสฺส ตโย คุณา นิพฺพานํ อนุปวิฏฺฐาติ.}

\prose{ภนฺเต นาคเสน “โลหิต\-จนฺทนสฺส ตโย คุณา นิพฺพานํ อนุปวิฏฺฐา”ติ ยํ วเทสิ, กตเม โลหิต\-จนฺทนสฺส ตโย คุณา นิพฺพานํ อนุปวิฏฺฐาติ. ยถา มหาราช โลหิตจนฺทนํ ทุลฺลภํ, เอวเมว โข มหาราช นิพฺพานํ ทุลฺลภํ. อยํ มหาราช โลหิต\-จนฺทนสฺส ปฐโม คุโณ นิพฺพานํ อนุปวิฏฺโฐ. ปุน จปรํ มหาราช โลหิตจนฺทนํ อสมสุคนฺธํ, เอวเมว โข มหาราช นิพฺพานํ อสมสุคนฺธํ. อยํ มหาราช โลหิต\-จนฺทนสฺส ทุติโย คุโณ นิพฺพานํ, อนุปวิฏฺโฐ. ปุน จปรํ มหาราช โลหิตจนฺทนํ สชฺชน\-ปสตฺถํ\csromansharedfootnote{5dcbfb6fc0f3886d}{สพฺพชนปสตฺถํ (สฺยา)}, เอวเมว โข มหาราช นิพฺพานํ อริย\-สชฺชน\-ปสตฺถํ. อยํ มหาราช โลหิต\-จนฺทนสฺส ตติโย คุโณ นิพฺพานํ อนุปวิฏฺโฐ. อิเม โข มหาราช โลหิต\-จนฺทนสฺส ตโย คุณา นิพฺพานํ อนุปวิฏฺฐาติ.}

\prose{\csromanfolio{308}ภนฺเต นาคเสน “สปฺปิมณฺฑสฺส ตโย คุณา นิพฺพานํ อนุปวิฏฺฐา”ติ ยํ วเทสิ, กตเม สปฺปิมณฺฑสฺส ตโย คุณา นิพฺพานํ อนุปวิฏฺฐาติ. ยถา มหาราช สปฺปิมณฺโฑ วณฺณสมฺปนฺโน, เอวเมว โข มหาราช นิพฺพานํ คุณวณฺณสมฺปนฺนํ. อยํ มหาราช สปฺปิมณฺฑสฺส ปฐโม คุโณ นิพฺพานํ อนุปวิฏฺโฐ. ปุน จปรํ มหาราช สปฺปิมณฺโฑ คนฺธสมฺปนฺโน, เอวเมว โข มหาราช นิพฺพานํ สีลคนฺธ\-สมฺปนฺนํ. อยํ มหาราช สปฺปิมณฺฑสฺส ทุติโย คุโณ นิพฺพานํ อนุปวิฏฺโฐ. ปุน จปรํ มหาราช สปฺปิมณฺโฑ รสสมฺปนฺโน, เอวเมว โข มหาราช นิพฺพานํ รสสมฺปนฺนํ. อยํ มหาราช สปฺปิมณฺฑสฺส ตติโย คุโณ นิพฺพานํ อนุปวิฏฺโฐ. อิเม โข มหาราช สปฺปิมณฺฑสฺส ตโย คุณา นิพฺพานํ อนุปวิฏฺฐาติ.}

\prose{ภนฺเต นาคเสน “คิริสิขรสฺส ปญฺจ คุณา นิพฺพานํ อนุปวิฏฺฐา”ติ ยํ วเทสิ, กตเม คิริสิขรสฺส ปญฺจ คุณา นิพฺพานํ อนุปวิฏฺฐาติ. ยถา มหาราช คิริสิขรํ อจฺจุคฺคตํ, เอวเมว โข มหาราช นิพฺพานํ อจฺจุคฺคตํ. อยํ มหาราช คิริสิขรสฺส ปฐโม คุโณ นิพฺพานํ อนุปวิฏฺโฐ. ปุน จปรํ มหาราช คิริสิขรํ อจลํ, เอวเมว โข มหาราช นิพฺพานํ อจลํ. อยํ มหาราช คิริสิขรสฺส ทุติโย คุโณ นิพฺพานํ อนุปวิฏฺโฐ. ปุน จปรํ มหาราช คิริสิขรํ ทุรธิโรหํ, เอวเมว โข มหาราช นิพฺพานํ ทุรธิโรหํ สพฺพกิเลสานํ. อยํ มหาราช คิริสิขรสฺส ตติโย คุโณ นิพฺพานํ อนุปวิฏฺโฐ. ปุน จปรํ มหาราช คิริสิขรํ สพฺพพีชานํ อวิรูหนํ, เอวเมว โข มหาราช นิพฺพานํ สพฺพกิเลสานํ อวิรูหนํ. อยํ มหาราช คิริสิขรสฺส จตุตฺโถ คุโณ นิพฺพานํ อนุปวิฏฺโฐ. ปุน จปรํ มหาราช คิริสิขรํ อนุนย\-ปฺปฏิฆ\-วิปฺปมุตฺตํ, เอวเมว โข มหาราช นิพฺพานํ อนุนย\-ปฺปฏิฆ\-วิปฺปมุตฺตํ. อยํ มหาราช คิริสิขรสฺส ปญฺจโม คุโณ นิพฺพานํ อนุปวิฏฺโฐ. อิเม โข มหาราช คิริสิขรสฺส ปญฺจ คุณา นิพฺพานํ อนุปวิฏฺฐาติ. สาธุ ภนฺเต นาคเสน, เอวเมตํ ตถา สมฺปฏิจฺฉามีติ.}

\vspace{\tipitakaparskip}
\csromancenter{นิพฺพาน\-รูป\-สณฺฐาน\-ปโญฺห ทสโม.}
\csromanmatikaanchor{mtk.195}
\csromantocmarkref{subsubsection}{11. นิพฺพานสจฺฉิกรณปญฺห}{mtk.195}
\bookmark[dest={mtk.195},level=3]{11. นิพฺพานสจฺฉิกรณปญฺห}
\csromanmarkh{5. อนุมานปญฺห}\csromanmarkhii{11. นิพฺพาน\-สจฺฉิกรณ\-ปญฺห}\csromanheadernomark{3}{5. อนุมานปญฺห \textbf{11. นิพฺพาน}\-\textbf{สจฺฉิกรณ}\-\textbf{ปญฺห}}
\proseitem{11}{\csromanfolio{309}ภนฺเต นาคเสน ตุเมฺห ภณถ “นิพฺพานํ น อตีตํ, น อนาคตํ, น ปจฺจุปฺปนฺนํ, น อุปฺปนฺนํ น อนุปฺปนฺนํ น อุปฺปาทนียนฺ”ติ. อิธ ภนฺเต นาคเสน โย โกจิ สมฺมาปฏิปนฺโน นิพฺพานํ สจฺฉิกโรติ, โส อุปฺปนฺนํ สจฺฉิกโรติ, อุทาหุ อุปฺปาเทตฺวา สจฺฉิกโรตีติ. โย โกจิ มหาราช สมฺมาปฏิปนฺโน นิพฺพานํ สจฺฉิกโรติ, โส น อุปฺปนฺนํ สจฺฉิกโรติ, น อุปฺปาเทตฺวา สจฺฉิกโรติ, อปิ จ มหาราช อตฺเถสา นิพฺพานธาตุ ยํ โส สมฺมาปฏิปนฺโน สจฺฉิกโรตีติ.}

\prose{มา ภนฺเต นาคเสน อิมํ ปญฺหํ ปฏิจฺฉนฺนํ กตฺวา ทีเปติ, วิวฏํ ปากฏํ กตฺวา ทีเปหิ ฉนฺทชาโต อุสฺสาหชาโต, ยํ เต สิกฺขิตํ, ตํ สพฺพํ เอตฺเถวากิราหิ, เอตฺถายํ ชโน สมฺมูโฬฺห วิมติชาโต สํสย\-ปกฺขนฺโท, ภินฺเทตํ อนฺโตโทสสลฺลนฺติ. อตฺเถสา มหาราช นิพฺพานธาตุ สนฺตา สุขา ปณีตา, ตํ สมฺมาปฏิปนฺโน ชินานุสิฏฺฐิยา สงฺขาเร สมฺมสนฺโต ปญฺญาย สจฺฉิกโรติ. ยถา มหาราช อนฺเตวาสิโก อาจริยา\-นุสิฏฺฐิยา วิชฺชํ ปญฺญาย สจฺฉิกโรติ, เอวเมว โข มหาราช สมฺมาปฏิปนฺโน ชินานุสิฏฺฐิยา ปญฺญาย นิพฺพานํ สจฺฉิกโรติ.}

\prose{กถํ ปน ตํ นิพฺพานํ ทฏฺฐพฺพนฺติ. อนีติโต นิรุปทฺทวโต อภยโต เขมโต สนฺตโต สุขโต สาตโต ปณีตโต สุจิโต สีตลโต ทฏฺฐพฺพํ.}

\prose{ยถา มหาราช ปุริโส พหุกฏฺฐปุญฺเชน ชลิต\-กฐิเตน อคฺคินา ทยฺหมาโน วายาเมน ตโต มุญฺจิตฺวา นิรคฺคิโกกาสํ ปวิสิตฺวา ตตฺถ ปรมสุขํ ลเภยฺย, เอวเมว โข มหาราช โย สมฺมาปฏิปนฺโน, โส โยนิโส มนสิกาเรน พฺยปคตติวิธคฺคิสนฺตาปํ ปรมสุขํ นิพฺพานํ สจฺฉิกโรติ. ยถา มหาราช อคฺคิ, เอวํ ติวิธคฺคิ ทฏฺฐพฺโพ, ยถา อคฺคิคโต ปุริโส, เอวํ สมฺมาปฏิปนฺโน ทฏฺฐพฺโพ, ยถา นิรคฺคิโกกาโส, เอวํ นิพฺพานํ ทฏฺฐพฺพํ.}

\prose{ยถา วา ปน มหาราช ปุริโส อหิกุกฺกุร มนุสฺสกุณปสรีร วฬญฺชโกฏฺฐาสราสิคโต กุณป\-ชฏา\-ชฏิ\-ตนฺตรม\-นุปวิฏฺโฐ วายาเมน ตโต มุญฺจิตฺวา นิกฺกุณโปกาสํ ปวิสิตฺวา ตตฺถ ปรมสุขํ เลเภยฺย, \csromanfolio{310}เอวเมว โข มหาราช โย สมฺมาปฏิปนฺโน, โส โยนิโส มนสิกาเรน พฺยปคตกิเลสกุณปํ ปรมสุขํ นิพฺพานํ สจฺฉิกโรติ, ยถา มหาราช กุณปํ, เอวํ ปญฺจ กามคุณา ทฏฺฐพฺพา, ยถา กุณปคโต ปุริโส, เอวํ สมฺมาปฏิปนฺโน ทฏฺฐพฺโพ, ยถา นิกฺกุณโปกาโส, เอวํ นิพฺพานํ ทฏฺฐพฺพํ.}

\prose{ยถา วา ปน มหาราช ปุริโส ภีโต ตสิโต กมฺปิโต วิปรีต\-วิพฺภนฺตจิตฺโต วายาเมน ตโต มุญฺจิตฺวา ถิรํ อจลํ อภยฏฺฐานํ ปวิสิตฺวา ตตฺถ ปรมสุขํ ลเภยฺย, เอวเมว โข มหาราช โย สมฺมาปฏิปนฺโน, โส โยนิโส มนสิกาเยน พฺยปคตภยสนฺตาสํ ปรมสุขํ นิพฺพานํ สจฺฉิกโรติ. ยถา มหาราช ภยํ, เอวํ ชาติ\-ชรา\-พฺยาธิ\-มรณํ ปฏิจฺจ อปราปรํ ปวตฺตภยํ ทฏฺฐพฺพํ, ยถา ภีโต ปุริโส, เอวํ สมฺมาปฏิปนฺโน ทฏฺฐพฺโพ, ยถา อภยฏฺฐานํ, เอวํ นิพฺพานํ ทฏฺฐพฺพํ.}

\prose{ยถา วา ปน มหาราช ปุริโส กิลิฏฺฐ\-มลิน\-กลล\-กทฺทม\-เทเส ปติโต วายาเมน ตํ กลลกทฺทมํ อปวาเหตฺวา ปริสุทฺธ\-วิมล\-เทสมุ\-ปคนฺตฺวา ตตฺถ ปรมสุขํ ลเภยฺย, เอวเมว โข มหาราช โย สมฺมาปฏิปนฺโน, โส โยนิโส มนสิกาเรน พฺยปคตกิเลสมลกทฺทมํ ปรมสุขํ นิพฺพานํ สจฺฉิกโรติ. ยถา มหาราช กลลํ, เอวํ ลาภ\-สกฺการ\-สิโลโก ทฏฺฐพฺโพ, ยถา กลลคโต ปุริโส, เอวํ สมฺมาปฏิปนฺโน ทฏฺฐพฺโพ, ยถา ปริสุทฺธ\-วิมล\-เทโส, เอวํ นิพฺพานํ ทฏฺฐพฺพํ.}

\prose{ตญฺจ ปน นิพฺพานํ สมฺมาปฏิปนฺโน กินฺติ สจฺฉิกโรติ, โย โส มหาราช สมฺมาปฏิปนฺโน, โส สงฺขารานํ ปวตฺตํ สมฺมสติ. ปวตฺตํ สมฺมสมาโน ตตฺถ ชาติํ ปสฺสติ ชรํ ปสฺสติ พฺยาธิํ ปสฺสติ มรณํ ปสฺสติ, น ตตฺถ กิญฺจิ สุขํ สาตํ ปสฺสติ อาทิโตปิ มชฺฌโตปิ ปริโยสานโตปิ. โส ตตฺถ กิญฺจิ น คยฺหูปคํ ปสฺสติ. ยถา มหาราช ปุริโส ทิวสสนฺตตฺเต อโยคุเฬ ชลิเต ตตฺเต กฐิเต อาทิโตปิ มชฺฌโตปิ ปริโยสานโตปิ น กิญฺจิ คยฺหูปคํ ปเทสํ ปสฺสติ, เอวเมว โข มหาราช โย สงฺขารานํ ปวตฺตํ สมฺมสติ, โส ปวตฺตํ สมฺมสมาโน ตตฺถ ชาติํ ปสฺสติ ชรํ ปสฺสติ พฺยาธิํ ปสฺสติ มรณํ ปสฺสติ, น ตตฺถ กิญฺจิ สุขํ สาตํ ปสฺสติ อาทิโตปิ มชฺฌโตปิ ปริโยสานโตปิ. โส ตตฺถ\csromansharedfootnote{b87641f43a4275dd}{อิทํ ปททฺวยํ สีอิโปตฺถเกสุ นตฺถิ.} น กิญฺจิ คยฺหูปคํ ปสฺสติ, ตสฺส \csromanfolio{311}คยฺหูปคํ อปสฺสนฺตสฺส จิตฺเต อรติ สณฺฐาติ, กายสฺมิํ ฑาโห โอกฺกมติ, โส อตาโณ อสรโณ อสรณีภูโต ภเวสุ นิพฺพินฺทติ.}

\prose{ยถา มหาราช ปุริโส ชลิตชาลํ มหนฺตํ อคฺคิกฺขนฺธํ ปวิเสยฺย, โส ตตฺถ อตาโณ อสรโณ อสรณีภูโต อคฺคิมฺหิ นิพฺพินฺเทยฺย, เอวเมว โข มหาราช ตสฺส คยฺหูปคํ อปสฺสนฺตสฺส จิตฺเต อรติ สณฺฐาติ, กายสฺมิํ ฑาโห โอกฺกมติ, โส อตาโณ อสรโณ อสรณีภูโต ภเวสุ นิพฺพินฺทติ.}

\prose{ตสฺส ปวตฺเต ภยทสฺสาวิสฺส เอวํ จิตฺตํ อุปฺปชฺชติ “สนฺตตฺตํ โข ปเนตํ ปวตฺตํ อาทิตฺตํ สมชฺชลิตํ พหุทุกฺขํ พหูปายาสํ, ยทิ โกจิ ลเภถ อปฺปวตฺตํ เอตํ สนฺตํ เอตํ ปณีตํ, ยทิทํ สพฺพสงฺขาราสมโถ สพฺพู\-ปธิ\-ปฏินิสฺสคฺโค ตณฺหกฺขโย วิราโค นิโรโธ นิพฺพานนฺ”ติ. อิติ เหตํ\csromansharedfootnote{2a99c3fad0dd7b26}{หิตํ (สี), หิ อิทํ (อิ)} ตสฺส อปฺปวตฺเต จิตฺตํ ปกฺขนฺทติ ปสีทติ ปหํสยติ ตุสยติ\csromansharedfootnote{66bc5eadb953bf18}{ปหํสียติ กุหียติ (สี, อิ)} “ปฏิลทฺธํ โข เม นิสฺสรณนฺ”ติ.}

\prose{ยถา มหาราช ปุริโส วิปฺปนฏฺโฐ วิเทส\-ปกฺขนฺโท นิพฺพาหน\-มคฺคํ ทิสฺวา ตตฺถ ปกฺขนฺทติ ปสีทติ ปหํสยติ ตุสยติ\csromansharedfootnote{66bc5eadb953bf18}{ปหํสียติ กุหียติ (สี, อิ)} “ปฏิลทฺโธ เม นิพฺพาหนมคฺโค”ติ, เอวเมว โข มหาราช ปวตฺเต ภยทสฺสาวิสฺส อปฺปวตฺเต จิตฺตํ ปกฺขนฺทติ ปสีทติ ปหํสยติ ตุสยติ\csromansharedfootnote{66bc5eadb953bf18}{ปหํสียติ กุหียติ (สี, อิ)} “ปฏิลทฺธํ โข เม นิสฺสรณนฺ”ติ.}

\prose{โส อปฺปวตฺต\-ตฺถาย มคฺคํ อายูหติ คเวสติ ภาเวติ พหุลีกโรติ, ตสฺส ตทตฺถํ สติ สนฺติฏฺฐติ, ตทตฺถํ วีริยํ สนฺติฏฺฐติ, ตทตฺถํ ปีติ สนฺติฏฺฐติ, ตสฺส ตํ จิตฺตํ อปราปรํ มนสิกโรโต ปวตฺตํ สมติกฺกมิตฺวา อปฺปวตฺตํ โอกฺกมติ, อปฺปวตฺตม\-นุปฺปตฺโต มหาราช สมฺมาปฏิปนฺโน “นิพฺพานํ สจฺฉิกโรตี”ติ วุจฺจตีติ. สาธุ ภนฺเต นาคเสน, เอวเมตํ ตถา สมฺปฏิจฺฉามีติ. นิพฺพาน\-สจฺฉิกรณ\-ปโญฺห เอกาทสโม.}

\csromanmatikaanchor{mtk.196}
\csromantocmarkref{subsubsection}{12. นิพฺพานสนฺนิหิตปญฺห}{mtk.196}
\bookmark[dest={mtk.196},level=3]{12. นิพฺพานสนฺนิหิตปญฺห}
\csromanheader{3}{12. \textbf{นิพฺพาน}\-\textbf{สนฺนิหิต}\-\textbf{ปญฺห}}
\proseitem{12}{\csromanfolio{312}ภนฺเต นาคเสน อตฺถิ โส ปเทโส ปุรตฺถิมาย วา ทิสาย ทกฺขิณาย วา ทิสาย ปจฺฉิมาย วา ทิสาย อุตฺตราย วา ทิสาย อุทฺธํ วา อโธ วา ติริยํ วา, ยตฺถ นิพฺพานํ สนฺนิหิตนฺติ. นตฺถิ มหาราช โส ปเทโส ปุรตฺถิมาย วา ทิสาย ทกฺขิณาย วา ทิสาย ปจฺฉิมาย วา ทิสาย อุตฺตราย วา ทิสาย อุทฺธํ วา อโธ วา ติริยํ วา, ยตฺถ นิพฺพานํ สนฺนิหิตนฺติ.}

\prose{ยทิ ภนฺเต นาคเสน นตฺถิ นิพฺพานสฺส สนฺนิหิโตกาโส, เตน หิ นตฺถิ นิพฺพานํ. เยสญฺจ ตํ นิพฺพานํ สจฺฉิกตํ, เตสมฺปิ สจฺฉิกิริยา มิจฺฉา, การณํ ตตฺถ วกฺขามิ, ยถา ภนฺเต นาคเสน มหิยา ธญฺญุฏฺฐานํ เขตฺตํ อตฺถิ, คนฺธุฏฺฐานํ ปุปฺผํ อตฺถิ, ปุปฺผุฏฺฐานํ คุมฺโพ อตฺถิ, ผลุฏฺฐานํ รุกฺโข อตฺถิ, รตนุฏฺฐานํ อากโร อตฺถิ, ตตฺถ โย โกจิ ยํ ยํ อิจฺฉติ, โส ตตฺถ คนฺตฺวา ตํ ตํ หรติ, เอวเมว โข ภนฺเต นาคเสน ยทิ นิพฺพานํ อตฺถิ, ตสฺส นิพฺพานสฺส อุฏฺฐาโนกาโสปิ อิจฺฉิตพฺโพ, ยสฺมา จ โข ภนฺเต นาคเสน นิพฺพานสฺส อุฏฺฐาโนกาโส นตฺถิ, ตสฺมา นตฺถิ นิพฺพานนฺติ พฺรูมิ, เยสญฺจ นิพฺพานํ สจฺฉิกตํ, เตสมฺปิ สจฺฉิกิริยา มิจฺฉาติ.}

\prose{นตฺถิ มหาราช นิพฺพานสฺส สนฺนิหิโตกาโส, อตฺถิ เจตํ นิพฺพานํ, สมฺมาปฏิปนฺโน โยนิโส มนสิกาเรน นิพฺพานํ สจฺฉิกโรติ. ยถา ปน มหาราช อตฺถิ อคฺคิ นาม, นตฺถิ ตสฺส สนฺนิหิโตกาโส, เทฺว กฏฺฐานิ สงฺฆฏฺเฏนฺโต อคฺคิํ อธิคจฺฉติ. เอวเมว โข มหาราช อตฺถิ นิพฺพานํ, นตฺถิ ตสฺส สนฺนิหิโตกาโส, สมฺมาปฏิปนฺโน โยนิโส มนสิกาเรน นิพฺพานํ สจฺฉิกโรติ.}

\prose{ยถา วา ปน มหาราช อตฺถิ สตฺต รตนานิ นาม. เสยฺยถีทํ, จกฺกรตนํ หตฺถิรตนํ อสฺสรตนํ มณิรตนํ อิตฺถิรตนํ คหปติ\-รตนํ ปริณายก\-รตนํ, น จ เตสํ รตนานํ สนฺนิหิโตกาโส อตฺถิ, ขตฺติยสฺส ปน สมฺมา\-ปฏิปนฺนสฺส ปฏิปตฺติ\-พเลน ตานิ รตานานิ อุปคจฺฉนฺติ. เอวเมว โข มหาราช อตฺถิ นิพฺพานํ, นตฺถิ ตสฺส สนฺนิหิโตกาโส, สมฺมาปฏิปนฺโน โยนิโส มนสิกาเรน นิพฺพานํ สจฺฉิกโรตีติ.}

\prose{ภนฺเต นาคเสน นิพฺพานสฺส สนฺนิหิโตกาโส มา โหตุ, อตฺถิ ปน ตํ ฐานํ, ยตฺถ ฐิโต สมฺมาปฏิปนฺโน นิพฺพานํ สจฺฉิกโรตีติ. อาม \csromanfolio{313}มหาราช อตฺถิ ตํ ฐานํ, ยตฺถ ฐิโต สมฺมาปฏิปนฺโน นิพฺพานํ สจฺฉิกโรตีติ.}

\prose{กตมํ ปน ภนฺเต ตํ ฐานํ, ยตฺถ ฐิโต สมฺมาปฏิปนฺโน นิพฺพานํ สจฺฉิกโรตีติ. สีลํ มหาราช ฐานํ, สีเล ปติฏฺฐิโต โยนิโส มนสิกโรนฺโต สกฺกยวเนปิ จีนวิลาเตปิ อลสนฺเทปิ นิคุมฺเพปิ\csromansharedfootnote{54c84562252d42a9}{นิกุมฺเพปิ (สี, สฺยา, อิ)} กาสิโกสเลปิ กสฺมีเรปิ คนฺธาเรปิ นคมุทฺธนิปิ พฺรหฺมโลเกปิ ยตฺถ กตฺถจิปิ ฐิโต สมฺมาปฏิปนฺโน นิพฺพานํ สจฺฉิกโรติ. ยถา มหาราช โย โกจิ จกฺขุมา ปุริโส สกยวเนปิ จีนวิลาเตปิ อลสนฺเทปิ นิคุมฺเพปิ กาสิโกสเลปิ กสฺมีเรปิ คนฺธาเรปิ นคมุทฺธนิปิ พฺรหฺมโลเกปิ ยตฺถ กตฺถจิปิ ฐิโต อากาสํ ปสฺสติ, เอวเมว โข มหาราช สีเล ปติฏฺฐิโต โยนิโส มนสิกโรนฺโต สกยวเนปิ ฯเปฯ ยตฺถ กตฺถจิปิ ฐิโต สมฺมาปฏิปนฺโน นิพฺพานํ สจฺฉิกโรติ.}

\prose{ยถา วา ปน มหาราช สกยวเนปิ ฯเปฯ ยตฺถ กตฺถจิปิ ฐิตสฺส ปุพฺพทิสา อตฺถิ, เอวเมว โข มหาราช สีเล ปติฏฺฐิตสฺส โยนิโส มนสิ\-กโรนฺตสฺส สกฺกยวเนปิ ฯเปฯ ยตฺถ กตฺถจิปิ ฐิตสฺส สมฺมา\-ปฏิปนฺนสฺส อตฺถิ นิพฺพาน\-สจฺฉิกิริยาติ. สาธุ ภนฺเต นาคเสน, เทสิตํ ตยา นิพฺพานํ, เทสิตา นิพฺพาน\-สจฺฉิกิริยา, ปริกฺขตา สีลคุณา, ทสฺสิตา สมฺมาปฏิปตฺติ, อุสฺสาปิโต ธมฺมทฺธโช, สณฺฐปิตา ธมฺมเนตฺติ, อวญฺโฌ สุปฺปยุตฺตานํ สมฺมาปโยโค, เอวเมตํ คณิวรปวร ตถา สมฺปฏิจฺฉามีติ.}

\vspace{\tipitakaparskip}
\csromancenter{นิพฺพาน\-สนฺนิหิต\-ปโญฺห ทฺวาทสโม.}
\nitthitammidruled{\textbf{เวสฺสนฺตร}\-\textbf{วคฺโค ตติโย.}}
\csromancenter{อิมสฺมิํ วคฺเค ทฺวาทส ปญฺหา.}
\csromanmatikaanchor{mtk.197}
\csromantocmarknopage{subsection}{4. อนุมานวคฺค}{mtk.197}
\bookmark[dest={mtk.197},level=2]{4. อนุมานวคฺค}
\csromanheader{2}{4. \textbf{อนุมานวคฺค}}
\csromanmatikaanchor{mtk.198}
\csromantocmarkref{subsubsection}{1. อนุมานปญฺห}{mtk.198}
\bookmark[dest={mtk.198},level=3]{1. อนุมานปญฺห}
\proseitem{1}{\csromanfolio{314}อถ โข มิลินฺโท ราชา เยนายสฺมา นาคเสโน เตนุปสงฺกมิ, อุปสงฺกมิตฺวา อายสฺมนฺตํ นาคเสนํ อภิวาเทตฺวา เอกมนฺตํ นิสีทิ, เอกมนฺตํ นิสินฺโน โข มิลินฺโท ราชา ญาตุกาโม โสตุกาโม ธาเรตุกาโม ญาณาโลกํ ทฏฺฐุกาโม อญฺญาณํ ภินฺทิตุกาโม ญาณาโลกํ อุปฺปาเทตุกาโม อวิชฺช\-นฺธการํ นาเสตุกาโม อธิมตฺตํ ธิติญฺจ อุสฺสาหญฺจ สติญฺจ สมฺปชญฺญญฺจ อุปฏฺฐเปตฺวา อายสฺมนฺตํ นาคเสนํ เอตทโวจ “ภนฺเต นาคเสน กิํ ปน พุทฺโธ ตยา ทิฏฺโฐ”ติ. น หิ มหาราชาติ. กิํ ปน เต อาจริเยหิ พุทฺโธ ทิฏฺโฐติ. น หิ มหาราชาติ. ภนฺเต นาคเสน น กิร ตยา พุทฺโธ ทิฏฺโฐ, นาปิ กิร เต อาจริเยหิ พุทฺโธ ทิฏฺโฐ, เตน หิ ภนฺเต นาคเสน นตฺถิ พุทฺโธ, น เหตฺถ พุทฺโธ ปญฺญายตีติ.}

\prose{อตฺถิ ปน เต มหาราช ปุพฺพกา ขตฺติยา, เย เต ตว ขตฺติย\-วํสสฺส ปุพฺพงฺคมาติ. อาม ภนฺเต. โก สํสโย, อตฺถิ ปุพฺพกา ขตฺติยา, เย มม ขตฺติย\-วํสสฺส ปุพฺพงฺคมาติ. ทิฏฺฐปุพฺพา ตยา มหาราช ปุพฺพกา ขตฺติยาติ. น หิ ภนฺเตติ. เย ปน ตํ มหาราช อนุสาสนฺติ ปุโรหิตา เสนาปติโน อกฺขทสฺสา มหามตฺตา, เตหิ ปุพฺพกา ขตฺติยา ทิฏฺฐปุพฺพาติ. น หิ ภนฺเตติ. ยทิ ปน เต มหาราช ปุพฺพกา ขตฺติยา น ทิฏฺฐา, นาปิ กิร เต อนุสาสเกหิ ปุพฺพกา ขตฺติยา ทิฏฺฐา, เตน หิ นตฺถิ ปุพฺพกา ขตฺติยา, น เหตฺถ ปุพฺพกา ขตฺติยา ปญฺญายนฺตีติ.}

\prose{ทิสฺสนฺติ ภนฺเต นาคเสน ปุพฺพกานํ ขตฺติยานํ อนุภูตานิ ปริโภค\-ภณฺฑานิ. เสยฺยถีทํ, เสตจฺฉตฺตํ อุณฺหีสํ ปาทุกา วาลพีชนี ขคฺครตนํ มหารหานิ จ สยนานิ. เยหิ มยํ ชาเนยฺยาม สทฺทเหยฺยาม “อตฺถิ ปุพฺพกา ขตฺติยา”ติ. เอวเมว โข มหาราช มยมฺเปตํ ภควนฺตํ ชาเนยฺยาม สทฺทเหยฺยาม, อตฺถิ ตํ การณํ, เยน มยํ การเณน ชาเนยฺยาม สทฺทเหยฺยาม “อตฺถิ โส ภควา”ติ. กตมํ ตํ การณํ, อตฺถิ โข มหาราช เตน ภควตา ชานตา ปสฺสตา อรหตา สมฺมา\-สมฺพุทฺเธน อนุภูตานิ ปริโภค\-ภณฺฑานิ. เสยฺยถีทํ, จตฺตาโร \csromanfolio{315}สติปฏฺฐานา จตฺตาโร สมฺมปฺปธานา จตฺตาโร อิทฺธิปาทา ปญฺจินฺทฺริยานิ ปญฺจ พลานิ สตฺต โพชฺฌงฺคา อริโย อฏฺฐงฺคิโต มคฺโค, เยหิ สเทวโก โลโก ชานาติ สทฺทหติ “อตฺถิ โส ภควา”ติ, อิมินา มหาราช การเณน อิมินา เหตุนา อิมินา นเยน อิมินา อนุมาเนน ญาตพฺโพ “อตฺถิ โส ภควา”ติ.}

\setlength{\csromangathapagewidth}{0pt}
\setlength{\csromangathawakwidth}{0pt}
\csromangathameasuregroup{“พหู ชเน ตารยิตฺวา, \\ อนุมาเนน ญาตพฺพํ,}{\csromangathabat{“พหู ชเน ตารยิตฺวา,}{นิพฺพุโต อุปธิกฺขเย.} \\ \csromangathabat{อนุมาเนน ญาตพฺพํ,}{อตฺถิ โส ทฺวิปทุตฺตโม”ติ.}}
\csromangathasetleft{“พหู ชเน ตารยิตฺวา, \\ อนุมาเนน ญาตพฺพํ,}
\csromangathagroup{\csromangathastanza{\csromangathabat{“พหู ชเน ตารยิตฺวา,}{นิพฺพุโต อุปธิกฺขเย.} \\ \csromangathabat{อนุมาเนน ญาตพฺพํ,}{อตฺถิ โส ทฺวิปทุตฺตโม”ติ.}}}

\prose{ภนฺเต นาคเสน โอปมฺมํ กโรหีติ. ยถา มหาราช นครวฑฺฒกี นครํ มาเปตุกาโม ปฐมํ ตาว สมํ อนุนฺนต\-มโนนตํ อสกฺขร\-ปาสาณํ นิรุปทฺทวม\-นวชฺชํ รมณียํ ภูมิภาคํ อนุวิโลเกตฺวา ยํ ตตฺถ วิสมํ, ตํ สมํ การาเปตฺวา ขาณุกณฺฏกํ วิโสธาเปตฺวา ตตฺถ นครํ มาเปยฺย โสภนํ วิภตฺตํ ภาคโส มิตํ อุกฺกิณฺณปริขา\-ปาการํ ทฬฺห\-โคปุร\-ฏฺฏาล\-โกฏฺฏกํ ปุถุจจฺจร\-จตุกฺก\-สนฺธิสิงฺฆาฏกํ สุจิสมตล\-ราช\-มคฺคํ สุวิภตฺต\-อนฺตราปณํ อารามุ\-ยฺยาน\-ตฬาก\-โปกฺขรณิ\-อุทปาน\-สมฺปนฺนํ พหุวิธ\-เทวฏฺฐาน\-ปฺปฏิมณฺฑิตํ สพฺพ\-โทส\-วิรหิตํ, โส ตสฺมิํ นาคเร สพฺพถา เวปุลฺลตฺตํ ปตฺเต อญฺญํ เทสํ อุปคจฺเฉยฺย, อถ ตํ นครํ อปเรน สมเยน อิทฺธํ ภเวยฺย ผีตํ สุภิกฺขํ เขมํ สมิทฺธํ สิวํ อนีติกํ นุรุปทฺทวํ นานาชน\-สมากุลํ, ปุถู ขตฺติยา พฺราหฺมณา เวสฺสา สุทฺทา หตฺถาโรหา อสฺสาโรหา รถิกา ปตฺติกา ธนุคฺคหา ถรุคฺคหา เจลกา จลกา ปิณฺฑทายกา อุคฺคา ราชปุตฺตา ปกฺขนฺทิโน มหานาคา สูรา วมฺมิโน โยธิโน ทาสิกปุตฺตา ภฏิปุตฺตา\csromansharedfootnote{5c103dadb0a6729a}{ทาสปุตฺตา พฏฺฏิปุตฺตา (สี, อิ)} มลฺลกา คณกา อาฬาริกา สูทา กปฺปกา นหาปกา จุนฺทา มาลาการา สุวณฺณการา สชฺฌุการา สีสการา ติปุการา โลหการา วฏฺฏการา อโยการา มณิการา เปสการา กุมฺภการา เวณุการา โลณการา จมฺมการา รถการา ทนฺตการา รชฺชุการา โกจฺฉการา สุตฺตการา วิลีวการา ธนุการา ชิยการา อุสุการา จิตฺตการา รงฺคการา รชกา ตนฺตวายา ตุนฺนวายา เหรญฺญิกา ทุสฺสิกา คนฺธิกา ติณหารกา กฏฺฐหารกา ภถกา ปณฺณิกา ผลิกา\csromansharedfootnote{cd35ce782cf64d3a}{ผลฺลิกา (สี, อิ)} มูลิกา โอกนิกา ปูวิกา มจฺฉิกา มํสิกา มชฺชิกา นฏกา นจฺจกา ลงฺฆกา อินฺทชาลิกา เวตาลิกา \csromanfolio{316}มลฺลา ฉวฑาหกา ปุปฺผฉฑฺฑกา เวนา เนสาทา คณิกา ลาสิกา กุมฺภทาสิโย สกฺกยวนจีนวิลาตา อุชฺเชนกา ภารุกจฺฉกา กาสิโกสลา ปรนฺตกา มาคธกา สาเกตกา โสเรยฺยกา\csromansharedfootnote{0fc6df80729f14e4}{โสรฏฺฐกา (สี, อิ)} ปาเวยฺยกา โกฏุมฺพรมาถุรกา อลสนฺทกสฺมีรคนฺธารา ตํ นครํ วาสาย อุปคตา นานาวิสยิโน ชนา นวํ สุวิภตฺตํ อโทสม\-นวชฺชํ รมณียํ ตํ นครํ ปสฺสิตฺวา อนุมาเนน ชานนฺติ “เฉโก วต โภ โส นครวฑฺฒกี, โย อิมสฺส นครสฺส มาเปตา”ติ. เอวเมว โข มหาราช โส ภควา อสโม อสมสโม อปฺปฏิสโม อสทิโส อตุโล อสงฺเขฺยโย อปฺปเมยฺโย อปริเมยฺโย อมิตคุโณ คุณปารมิปฺปตฺโต อนนฺตธิติ อนนฺตเตโช อนนฺตวีริโย อนนฺตพโล พุทฺธพล\-ปารมิํ คโต สเสนมารํ ปราเชตฺวา ทิฏฺฐิชาลํ ปทาเลตฺวา อวิชฺชํ เขเปตฺวา วิชฺชํ อุปฺปาเทตฺวา ธมฺมุกฺกํ ธารยิตฺวา สพฺพญฺญุตํ ปาปุณิตฺวา วิชิตสงฺคาโม ธมฺมนครํ มาเปสิ.}

\proseitem{5}{ภควโต โข มหาราช ธมฺมนครํ สีลปาการํ หิริปริขํ ญาณ\-ทฺวารโกฏฺฐกํ วีริย\-อฏฺฏาลกํ สทฺธาเอสิกํ สติโทวาริกํ ปญฺญาปาสาทํ สุตฺตนฺต\-จจฺจรํ อภิธมฺม\-สิงฺฆาฏกํ วินย\-วินิจฺฉยํ สติปฏฺฐาน\-วีถิกํ, ตสฺส โข ปน มหาราช สติปฏฺฐาน\-วีถิยํ เอวรูปา อาปณา ปสาริตา โหนฺติ. เสยฺยถีทํ, ปุปฺผาปณํ คนฺธาปณํ ผลาปณํ อคทาปณํ โอสธาปณํ อมตาปณํ รตนาปณํ สพฺพาปณนฺติ.}

\prose{ภนฺเต นาคเสน กตมํ พุทฺธสฺส ภควโต ปุปฺผาปณนฺติ. อตฺถิ โข ปน มหาราช เตน ภควตา ชานตา ปสฺสตา อรหตา สมฺมา\-สมฺพุทฺเธน อารมฺมณ\-วิภตฺติโย อกฺขาตา. เสยฺยถีทํ, อนิจฺจสญฺญา ทุกฺขสญฺญา อนตฺตสญฺญา อสุภสญฺญา อาทีนวสญฺญา ปหานสญฺญา วิราคสญฺญา นิโรธสญฺญา สพฺพโลเก อนภิรติ\-สญฺญา สพฺพ\-สงฺขาเรสุ อนิจฺจสญฺญา อานาปานสฺสติ อุทฺธุมาตก\-สญฺญา วินีลกสญฺญา วิปุพฺพกสญฺญา วิจฺฉิทฺทก\-สญฺญา วิกฺขายิตก\-สญฺญา วิกฺขิตฺตก\-สญฺญา หตวิกฺขิตฺตก\-สญฺญา โลหิตกสญฺญา ปุฬวกสญฺญา อฏฺฐิกสญฺญา เมตฺตาสญฺญา กรุณาสญฺญา มุทิตาสญฺญา อุเปกฺขาสญฺญา มรณนุสฺสติ กายคตาสติ, อิเม โข มหาราช พุทฺเธน ภควตา อารมฺมณ\-วิภตฺติโย อกฺขาตา. ตตฺถ โย โกจิ ชรามรณา มุจฺจิตุกาโม, โส เตสุ \csromanfolio{317}อญฺญตรํ อารมฺมณํ คณฺหาติ, เตน อารมฺมเณน ราคา วิมุจฺจติ, โทสา วิมุจฺจติ, โมหา วิมุจฺจติ, มานโต วิมุจฺจติ, ทิฏฺฐิโต วิมุจฺจติ, สํสารํ ตรติ, ตณฺหาโสตํ นิวาเรติ, ติวิธํ มลํ วิโสเธติ, สพฺพกิเลเส อุปหนฺตฺวา อมลํ วิรชํ สุทฺธํ ปณฺฑรํ อชาติํ อชรํ อมรํ สุขํ สีติภูตํ อภยํ นครุตฺตมํ นิพฺพานนครํ ปวิสิตฺวา อรหตฺเต จิตฺตํ วิโมเจติ, อิทํ วุจฺจติ มหาราช “ภควโต ปุปฺผาปณนฺ”ติ.}

\setlength{\csromangathapagewidth}{0pt}
\setlength{\csromangathawakwidth}{0pt}
\csromangathameasuregroup{กมฺมมูลํ คเหตฺวาน, \\ อารมฺมณํ กิณิตฺวาน,}{\csromangathabat{กมฺมมูลํ คเหตฺวาน,}{อาปณํ อุปคจฺฉถ.} \\ \csromangathabat{อารมฺมณํ กิณิตฺวาน,}{ตโต มุจฺจถ มุตฺติยาติ.}}
\csromangathasetleft{กมฺมมูลํ คเหตฺวาน, \\ อารมฺมณํ กิณิตฺวาน,}
\csromangathagroup{\csromangathastanza{\csromangathabat{กมฺมมูลํ คเหตฺวาน,}{อาปณํ อุปคจฺฉถ.} \\ \csromangathabat{อารมฺมณํ กิณิตฺวาน,}{ตโต มุจฺจถ มุตฺติยาติ.}}}

\prose{ภนฺเต นาคเสน กตมํ พุทฺธสฺส ภควโต คนฺธาปณนฺติ. อตฺถิ โข ปน มหาราช เตน ภควตา สีลวิภตฺติโย อกฺขาตา, เยน สีลคนฺเธน อนุลิตฺตา ภควโต ปุตฺตา สเทวกํ โลกํ สีลคนฺเธน ธูเปนฺติ สมฺปธูเปนฺติ, ทิสมฺปิ อนุทิสมฺปิ อนุวาตมฺปิ ปฏิวาตมฺปิ วายนฺติ อติวายนฺติ, ผริตฺวา ติฏฺฐนฺติ. กตมา ตา สีลวิภตฺติโย, สรณสีลํ ปญฺจงฺคสีลํ อฏฺฐงฺคสีลํ ทสงฺคสีลํ ปญฺจุทฺเทส\-ปริยาปนฺนํ ปาติโมกฺข\-สํวรสีลํ, อิทํ วุจฺจติ มหาราช “ภควโต คนฺธาปณนฺ”ติ. ภาสิตมฺเปตํ มหาราช ภควตา เทวาติเทเวน\csromandash{}}

\setlength{\csromangathapagewidth}{0pt}
\setlength{\csromangathawakwidth}{0pt}
\csromangathameasuregroup{จนฺทนํ ตครํ วาปิ, \\ เอเตสํ คนฺธชาตานํ, \\ อปฺปมตฺโต อยํ คนฺโธ, \\ โย จ สีลวตํ คนฺโธ,}{“น ปุปฺผคนฺโธ ปฏิวาตเมติ, \\ น จนฺทนํ ตคฺครมลฺลิกา วา. \\ สตญฺจ คนฺโธ ปฏิวาตเมติ, \\ สพฺพา ทิสา สปฺปุริโส ปวายติ\textsuperscript{1}. \\ \csromangathabat{จนฺทนํ ตครํ วาปิ,}{อุปฺปลํ อถ วสฺสิกี.} \\ \csromangathabat{เอเตสํ คนฺธชาตานํ,}{สีลคนฺโธ อนุตฺตโร.} \\ \csromangathabat{อปฺปมตฺโต อยํ คนฺโธ,}{ยฺวายํ ตครจนฺทนํ.} \\ \csromangathabat{โย จ สีลวตํ คนฺโธ,}{วาติ เทเวสุ อุตฺตโม”ติ.}}
\csromangathameasurewak{“น ปุปฺผคนฺโธ ปฏิวาตเมติ, \\ น จนฺทนํ ตคฺครมลฺลิกา วา. \\ สตญฺจ คนฺโธ ปฏิวาตเมติ, \\ สพฺพา ทิสา สปฺปุริโส ปวายติ\textsuperscript{1}.}
\csromangathasetleft{จนฺทนํ ตครํ วาปิ, \\ เอเตสํ คนฺธชาตานํ, \\ อปฺปมตฺโต อยํ คนฺโธ, \\ โย จ สีลวตํ คนฺโธ,}
\csromangathagroup{\csromangathastanza{\csromangathawak{“น ปุปฺผคนฺโธ ปฏิวาตเมติ, \\ น จนฺทนํ ตคฺครมลฺลิกา วา. \\ สตญฺจ คนฺโธ ปฏิวาตเมติ, \\ สพฺพา ทิสา สปฺปุริโส ปวายติ\csromansharedfootnote{8447997b3fa6121c}{ปวาติ (สี, อิ) อํ 1. 227; ขุ 1. 21 ปิฏฺเฐ.}.}}\csromangathastanza{\csromangathabat{จนฺทนํ ตครํ วาปิ,}{อุปฺปลํ อถ วสฺสิกี.} \\ \csromangathabat{เอเตสํ คนฺธชาตานํ,}{สีลคนฺโธ อนุตฺตโร.}}\csromangathastanza{\csromangathabat{อปฺปมตฺโต อยํ คนฺโธ,}{ยฺวายํ ตครจนฺทนํ.} \\ \csromangathabat{โย จ สีลวตํ คนฺโธ,}{วาติ เทเวสุ อุตฺตโม”ติ.}}}

\prose{ภนฺเต นาคเสน กตมํ พุทฺธสฺส ภควโต ผลาปณนฺติ. ผลานิ โข มหาราช ภควตา อกฺขาตานิ. เสยฺยถีทํ, โสตาปตฺติผลํ สกทาคามิผลํ อนาคามิผลํ อรหตฺตผลํ สุญฺญต\-ผลสมาปตฺติ อนิมิตฺต\-ผลสมาปตฺติ อปฺปณิหิต\-ผลสมาปตฺติ. ตตฺถ โย โกจิ \csromanfolio{318}ยํ ผลํ อิจฺฉติ, โส กมฺมมูลํ ทตฺวา ปตฺถิตํ ผลํ กิณาติ. ยทิ โสตาปตฺติผลํ, ยทิ สกทาคามิผลํ, ยทิ อนาคามิผลํ, ยทิ อรหตฺตผลํ, ยทิ สุญฺญต\-ผล\-สมาปตฺติํ, ยทิ อนิมิตฺต\-ผล\-สมาปตฺติํ, ยทิ อปฺปณิหิต\-ผล\-สมาปตฺติํ. ยถา มหาราช กสฺสจิ ปุริสสฺส ธุวผโล อมฺโพ ภเวยฺย, โส น ตาว ตโต ผลานิ ปาเตติ, ยาว กยิกา น อาคจฺฉนฺติ, อนุปฺปตฺเต ปน กยิเก มูลํ คเหตฺวา เอวํ อาจิกฺขติ “อมฺโภ ปุริส เอโส โข ธุวผโล อมฺโพ, ตโต ยํ อิจฺฉสิ, เอตฺตกํ ผลํ คณฺหาหิ สลาฏุกํ วา โทวิลํ วา เกสิกํ วา อามํ วา ปกฺกํ วา”ติ, โส เตน อตฺตนา ทินฺนมูเลน ยทิ สลาฏุกํ อิจฺฉติ, สลาฏุกํ คณฺหาติ, ยทิ โทวิลํ อิจฺฉติ, โทวิลํ คณฺหาติ, ยทิ เกสิกํ อิจฺฉติ, เกสิกํ คณฺหาติ, ยทิ อามกํ อิจฺฉติ, อามกํ คณฺหาติ, ยทิ ปกฺกํ อิจฺฉติ, ปกฺกํ คณฺหาติ. เอวเมว โข มหาราช โย ยํ ผลํ อิจฺฉติ, โส กมฺมมูลํ ทตฺวา ปตฺถิตํ ผลํ คณฺหาติ, ยทิ โสตาปตฺติผลํ ฯเปฯ ยทิ อปฺปณิหิต\-ผล\-สมาปตฺติํ, อิทํ วุจฺจติ มหาราช “ภควโต ผลาปณนฺ”ติ.}

\setlength{\csromangathapagewidth}{0pt}
\setlength{\csromangathawakwidth}{0pt}
\csromangathameasuregroup{กมฺมมูลํ ชนา ทตฺวา, \\ เตน เต สุขิตา โหนฺติ,}{\csromangathabat{กมฺมมูลํ ชนา ทตฺวา,}{คณฺหนฺติ อมตปฺผลํ.} \\ \csromangathabat{เตน เต สุขิตา โหนฺติ,}{เย กีตา อมตปฺผลนฺติ.}}
\csromangathasetleft{กมฺมมูลํ ชนา ทตฺวา, \\ เตน เต สุขิตา โหนฺติ,}
\csromangathagroup{\csromangathastanza{\csromangathabat{กมฺมมูลํ ชนา ทตฺวา,}{คณฺหนฺติ อมตปฺผลํ.} \\ \csromangathabat{เตน เต สุขิตา โหนฺติ,}{เย กีตา อมตปฺผลนฺติ.}}}

\prose{ภนฺเต นาคเสน กตมํ พุทฺธสฺส ภควโต อคทาปณนฺติ. อคทานิ โข มหาราช ภควตา อกฺขาตานิ, เยหิ อคเทหิ โส ภควา สเทวกํ โลกํ กิเลสวิสโต ปริโมเจติ. กตมานิ ปน ตานิ อคทานิ, ยานิมานิ มหาราช ภควตา จตฺตาริ อริยสจฺจานิ อกฺขาตานิ. เสยฺยถีทํ, ทุกฺขํ อริยสจฺจํ ทุกฺข\-สมุทยํ อริยสจฺจํ ทุกฺขนิโรธํ อริยสจฺจํ ทุกฺขนิโรธ\-คามินี ปฏิปทา อริยสจฺจํ, ตตฺถ เย เกจิ อญฺญาเปกฺขา จตุสจฺจํ ธมฺมํ สุณนฺติ, เต ชาติยา ปริมุจฺจนฺติ, ชราย ปริมุจฺจนฺติ, มรณา ปริมุจฺจนฺติ, โสก\-ปริเทว\-ทุกฺข\-โทมนสฺสุปายาเสหิ ปริมุจฺจนฺติ, อิทํ วุจฺจติ มหาราช “ภควโต อคทาปณนฺ”ติ.}

\setlength{\csromangathapagewidth}{0pt}
\setlength{\csromangathawakwidth}{0pt}
\csromangathameasuregroup{เย เกจิ อคทา โลเก\textsuperscript{1}, \\ ธมฺมาคทสมํ นตฺถิ,}{\csromangathabat{เย เกจิ อคทา โลเก\textsuperscript{1},}{วิสานํ ปฏิพาหกา.} \\ \csromangathabat{ธมฺมาคทสมํ นตฺถิ,}{เอตํ ปิวถ ภิกฺขโวติ.}}
\csromangathasetleft{เย เกจิ อคทา โลเก\textsuperscript{1}, \\ ธมฺมาคทสมํ นตฺถิ,}
\csromangathagroup{\csromangathastanza{\csromangathabat{เย เกจิ อคทา โลเก\csromansharedfootnote{ecc4a585f4546449}{โลเก อคทา (อิ)},}{วิสานํ ปฏิพาหกา.} \\ \csromangathabat{ธมฺมาคทสมํ นตฺถิ,}{เอตํ ปิวถ ภิกฺขโวติ.}}}

\prose{\csromanfolio{319}ภนฺเต นาคเสน กตมํ พุทฺธสฺส ภควโต โอสธาปณนฺติ. โอสธานิ โข มหาราช ภควตา อกฺขาตานิ, เยหิ โอสเธหิ โส ภควา เทวมนุสฺเส ติกิจฺฉติ. เสยฺยถีทํ, จตฺตาโร สติปฏฺฐานา จตฺตาโร สมฺมปฺปธานา จตฺตาโร อิทฺธิปาทา ปญฺจินฺทฺริยานิ ปญฺจ พลานิ สตฺต โพชฺฌงฺคา อริโย อฏฺฐงฺคิโก มคฺโค, เอเตหิ โอสเธหิ ภควา มิจฺฉาทิฏฺฐิํ วิเรเจติ, มิจฺฉา\-สงฺกปฺปํ วิเรเจติ, มิจฺฉาวาจํ วิเรเจติ, มิจฺฉา\-กมฺมนฺตํ วิเรเจติ, มิจฺฉาอาชีวํ วิเรเจติ, มิจฺฉาวายามํ วิเรเจติ, มิจฺฉาสติํ วิเรเจติ, มิจฺฉาสมาธิํ วิเรเจติ, โลภวมนํ กาเรติ, โทสวมนํ กาเรติ, โมหวมนํ กาเรติ, มานวมนํ กาเรติ, ทิฏฺฐิวมนํ กาเรติ, วิจิกิจฺฉา\-วมนํ กาเรติ, อุทฺธจฺจ\-วมนํ กาเรติ, ถินมิทฺธ\-วมนํ กาเรติ, อหิริกา\-โนตฺตปฺป\-วมนํ กาเรติ, สพฺพกิเลส\-วมนํ กาเรติ, อิทํ วุจฺจติ มหาราช “ภควโต โอสธาปณนฺ”ติ.}

\setlength{\csromangathapagewidth}{0pt}
\setlength{\csromangathawakwidth}{0pt}
\csromangathameasuregroup{เย เกจิ โอสธา โลเก, \\ ธมฺโมสธสมํ นตฺถิ, \\ ธมฺโมสธํ ปิวิตฺวาน, \\ ภาวยิตฺวา จ ปสฺสิตฺวา,}{\csromangathabat{เย เกจิ โอสธา โลเก,}{วิชฺชนฺติ วิวิธา พหู.} \\ \csromangathabat{ธมฺโมสธสมํ นตฺถิ,}{เอตํ ปิวถ ภิกฺขโว.} \\ \csromangathabat{ธมฺโมสธํ ปิวิตฺวาน,}{อชรามรณา สิยุํ.} \\ \csromangathabat{ภาวยิตฺวา จ ปสฺสิตฺวา,}{นิพฺพุตา อุปธิกฺขเยติ.}}
\csromangathasetleft{เย เกจิ โอสธา โลเก, \\ ธมฺโมสธสมํ นตฺถิ, \\ ธมฺโมสธํ ปิวิตฺวาน, \\ ภาวยิตฺวา จ ปสฺสิตฺวา,}
\csromangathagroup{\csromangathastanza{\csromangathabat{เย เกจิ โอสธา โลเก,}{วิชฺชนฺติ วิวิธา พหู.} \\ \csromangathabat{ธมฺโมสธสมํ นตฺถิ,}{เอตํ ปิวถ ภิกฺขโว.}}\csromangathastanza{\csromangathabat{ธมฺโมสธํ ปิวิตฺวาน,}{อชรามรณา สิยุํ.} \\ \csromangathabat{ภาวยิตฺวา จ ปสฺสิตฺวา,}{นิพฺพุตา อุปธิกฺขเยติ.}}}

\prose{ภนฺเต นาคเสน กตมํ พุทฺธสฺส ภควโต อมตาปณนฺติ. อมตํ โข มหาราช ภควตา อกฺขาตํ, เยน อมเตน โส ภควา สเทวกํ โลกํ อภิสิญฺจิ, เยน อมเตน อภิสิตฺตา เทวมนุสฺสา ชาติชราพฺยาธิมรณ\-โสก\-ปริเทว\-ทุกฺข\-โทมนสฺสุปายาเสหิ ปริมุจฺจิํสุ. กตมํ ตํ อมตํ, ยทิทํ กายคตาสติ. ภาสิตมฺเปตํ มหาราช ภควตา เทวาติเทเวน\csromansymbolfootnote{*}{อํ 1. 47 ปิฏฺเฐ.}“อมตํ เต ภิกฺขเว ปริภุญฺชนฺติ, เย กายคตาสติํ ปริภุญฺชนฺตี”ติ, อิทํ วุจฺจติ มหาราช “ภควโต อมตาปณนฺ”ติ.}

\setlength{\csromangathapagewidth}{0pt}
\setlength{\csromangathawakwidth}{0pt}
\csromangathameasuregroup{พฺยาธิตํ ชนตํ ทิสฺวา, \\ กมฺเมน ตํ กิณิตฺวาน,}{\csromangathabat{พฺยาธิตํ ชนตํ ทิสฺวา,}{อมตาปณํ ปสารยิ.} \\ \csromangathabat{กมฺเมน ตํ กิณิตฺวาน,}{อมตํ อาเทถ ภิกฺขโวติ.}}
\csromangathasetleft{พฺยาธิตํ ชนตํ ทิสฺวา, \\ กมฺเมน ตํ กิณิตฺวาน,}
\csromangathagroup{\csromangathastanza{\csromangathabat{พฺยาธิตํ ชนตํ ทิสฺวา,}{อมตาปณํ ปสารยิ.} \\ \csromangathabat{กมฺเมน ตํ กิณิตฺวาน,}{อมตํ อาเทถ ภิกฺขโวติ.}}}

\prose{ภนฺเต นาคเสน กตมํ พุทฺธสฺส ภควโต รตนาปณนฺติ. รตนานิ โข มหาราช ภควตา อกฺขาตานิ, เยหิ รตเนหิ วิภูสิตา ภควโต ปุตฺตา สเทวกํ โลกํ วิโรจนฺติ โอภาเสนฺติ ปภาเสนฺติ ชลนฺติ ปชฺชลนฺติ อุทฺธํ อโธ ติริยํ อาโลกํ ทสฺเสนฺติ. กตมานิ}

\prose{\csromanfolio{320}ตานิ รตนานิ, สีลรตนํ สมาธิรตนํ ปญฺญารตนํ วิมุตฺติรตนํ วิมุตฺติ\-ญาณ\-ทสฺสนรตนํ ปฏิสมฺภิทา\-รตนํ โพชฺฌงฺค\-รตนํ.}

\prose{กตมํ มหาราช ภควโต สีลรตนํ, ปาติโมกฺข\-สํวรสีลํ อินฺทฺริยสํวรสีลํ อาชีว\-ปาริสุทฺธิสีลํ ปจฺจยสนฺนิสฺสิตสีลํ จูฬสีลํ มชฺฌิมสีลํ มหาสีลํ มคฺคสีลํ ผลสีลํ. สีลรตเนน โข มหาราช วิภูสิตสฺส ปุคฺคลสฺส สเทวโก โลโก สมารโก สพฺรหฺมโก สสฺสมณ\-พฺราหฺมณี ปชา ปิหยติ ปตฺเถติ, สีลรตน\-ปิฬนฺโธ โข มหาราช ภิกฺขุ ทิสมฺปิ อนุทิสมฺปิ อุทฺธมฺปิ อโธปิ ติริยมฺปิ วิโรจติ อติวิโรจติ\csromansharedfootnote{d1390573a3cf485f}{อติโรจติ (สี, อิ)}, เหฏฺฐโต อวีจิํ อุปริโต ภวคฺคํ อุปาทาย เอตฺถนฺตเร สพฺพรตนานิ อติกฺกมิตฺวา อภิภวิตฺวา\csromansharedfootnote{77dd5461f29be15c}{อติสยิตฺวา (สี, อิ)} อชฺโฌตฺถริตฺวา ติฏฺฐติ, เอวรูปานิ โข มหาราช สีลรตนานิ ภควโต รตนาปเณ ปสาริตานิ, อิทํ วุจฺจติ มหาราช “ภควโต สีลรตนนฺ”ติ.}

\setlength{\csromangathapagewidth}{0pt}
\setlength{\csromangathawakwidth}{0pt}
\csromangathameasuregroup{เอวรูปานิ สีลานิ, \\ กมฺเมน ตํ กิณิตฺวาน,}{\csromangathabat{เอวรูปานิ สีลานิ,}{สนฺติ พุทฺธสฺส อาปเณ.} \\ \csromangathabat{กมฺเมน ตํ กิณิตฺวาน,}{รตนํ โว ปิฬินฺธถาติ.}}
\csromangathasetleft{เอวรูปานิ สีลานิ, \\ กมฺเมน ตํ กิณิตฺวาน,}
\csromangathagroup{\csromangathastanza{\csromangathabat{เอวรูปานิ สีลานิ,}{สนฺติ พุทฺธสฺส อาปเณ.} \\ \csromangathabat{กมฺเมน ตํ กิณิตฺวาน,}{รตนํ โว ปิฬินฺธถาติ.}}}

\prose{กตมํ มหาราช ภควโต สมาธิรตนํ, สวิตกฺก\-สวิจาโร สมาธิ อวิตกฺก\-วิจารมตฺโต สมาธิ อวิตกฺก\-อวิจาโร สมาธิ สุญฺญโต สมาธิ อนิมิตฺโต สมาธิ อปฺปณิหิโต สมาธิ, สมาธิรตนํ โข มหาราช ปิฬนฺธสฺส ภิกฺขุโน เย เต กามวิตกฺกพฺยาปาทวิตกฺก วิหิํสาวิตกฺกมานุทฺธจฺจทิฏฺฐิวิจิกิจฺฉากิเลสวตฺถูนิ วิวิธานิ จ กุวิตกฺกานิ, เต สพฺเพ สมาธิํ อาสชฺช วิกิรนฺติ วิธมนฺติ วิทฺธํสนฺติ น สณฺฐนฺติ\csromansharedfootnote{9ad3bf8b55469b6f}{น สณฺฐหนฺติ (สี)} น อุปลิมฺปนฺติ\csromansharedfootnote{e8a939d7502058ee}{น อุปลิปฺปนฺติ (สี, อิ)}. ยถา มหาราช วาริ โปกฺขรปตฺเต วิกิรติ วิธมติ วิทฺธํสติ น สณฺฐาติ น อุปลิมฺปติ, ตํ กิสฺส เหตุ, ปริสุทฺธตฺตา ปทุมสฺส. เอวเมว โข มหาราช สมาธิรตนํ ปิฬนฺธสฺส ภิกฺขุโน เย เต กามวิตกฺกพฺยาปาทวิตกฺก วิหิํสาวิตกฺก มานุทฺธจฺจ ทิฏฺฐิ วิจิกิจฺฉา กิเลสวตฺถูนิ วิวิธานิ จ กุวิตกฺกานิ, เต สพฺเพ สมาธิํ อาสชฺช วิกิรนฺติ วิธมนฺติ วิทฺธํสนฺติ น สณฺฐนฺติ น อุปลิมฺปนฺติ, ตํ กิสฺส เหตุ, ปริสุทฺธตฺตา สมาธิสฺส, อิทํ วุจฺจติ มหาราช “ภควโต สมาธิรตนนฺ”ติ, เอวรูปานิ โข มหาราช สมาธิ\-รตนานิ ภควโต รตนาปเณ ปสาริตานิ.}

\setlength{\csromangathapagewidth}{0pt}
\setlength{\csromangathawakwidth}{0pt}
\csromangathameasuregroup{สมาธิรตนมาลสฺส, \\ น จ วิกฺขิปเต จิตฺตํ,}{\csromangathabat{สมาธิรตนมาลสฺส,}{กุวิตกฺกา น ชายเร.} \\ \csromangathabat{น จ วิกฺขิปเต จิตฺตํ,}{เอตํ ตุเมฺห ปิฬนฺธถาติ.}}
\csromangathasetleft{สมาธิรตนมาลสฺส, \\ น จ วิกฺขิปเต จิตฺตํ,}
\csromangathagroup{\csromangathastanza{\csromanfolio{321}\csromangathabat{สมาธิรตนมาลสฺส,}{กุวิตกฺกา น ชายเร.} \\ \csromangathabat{น จ วิกฺขิปเต จิตฺตํ,}{เอตํ ตุเมฺห ปิฬนฺธถาติ.}}}

\prose{กตมํ มหาราช ภควโต ปญฺญารตนํ, ยาย มหาราช ปญฺญาย อริยสาวโก “อิทํ กุสลนฺ”ติ ยถาภูตํ ปชานาติ, “อิทํ อกุสลนฺ”ติ ยถาภูตํ ปชานาติ, “อิทํ สาวชฺชํ, อิทํ อนวชฺชํ, อิทํ เสวิตพฺพํ, อิทํ น เสวิตพฺพํ, อิทํ หีนํ, อิทํ ปณีตํ, อิทํ กณฺหํ, อิทํ สุกฺกํ, อิทํ กณฺหสุกฺกสปฺปฏิภาคนฺ”ติ ยถาภูตํ ปชานาติ, “อิทํ ทุกฺขนฺ”ติ ยถาภูตํ ปชานาติ, “อยํ ทุกฺขสมุทโย”ติ ยถาภูตํ ปชานาติ, “อยํ ทุกฺขนิโรโธ”ติ ยถาภูตํ ปชานาติ, “อยํ ทุกฺขนิโรธ\-คามินี ปฏิปทา”ติ ยถาภูตํ ปชานาติ. อิทํ วุจฺจติ มหาราช “ภควโต ปญฺญารตนนฺ”ติ.}

\setlength{\csromangathapagewidth}{0pt}
\setlength{\csromangathawakwidth}{0pt}
\csromangathameasuregroup{ปญฺญารตนมาลสฺส, \\ ขิปฺปํ ผสฺเสติ\textsuperscript{1} อมตํ,}{\csromangathabat{ปญฺญารตนมาลสฺส,}{น จิรํ วตฺตเต ภโว.} \\ \csromangathabat{ขิปฺปํ ผสฺเสติ\textsuperscript{1} อมตํ,}{น จ โส โรจเต ภเวติ.}}
\csromangathasetleft{ปญฺญารตนมาลสฺส, \\ ขิปฺปํ ผสฺเสติ\textsuperscript{1} อมตํ,}
\csromangathagroup{\csromangathastanza{\csromangathabat{ปญฺญารตนมาลสฺส,}{น จิรํ วตฺตเต ภโว.} \\ \csromangathabat{ขิปฺปํ ผสฺเสติ\csromansharedfootnote{dd58b13cb4208c01}{ผุสฺเสติ (สฺยา), ปสฺสติ (ก)} อมตํ,}{น จ โส โรจเต ภเวติ.}}}

\prose{กตมํ มหาราช ภควโต วิมุตฺติรตนํ, วิมุตฺติรตนํ\csromansharedfootnote{c2baf97e5fc53061}{วิมุตฺติรตนนฺติ (สี, อิ)} โข มหาราช อรหตฺตํ วุจฺจติ, อรหตฺตํ ปตฺโต โข มหาราช ภิกฺขุ “วิมุตฺติรตนํ ปิฬนฺโธ”ติ วุจฺจติ. ยถา มหาราช ปุริโส มุตฺตา\-กลาป\-มณิ\-กลาป\-ปวาฬ\-กลาปา\-ภรณ\-ปฺปฏิมณฺฑิโต\csromansharedfootnote{af40b4a7c60e54c5}{ปวาฬาภรณปฏิปณฺฑิโต (สี, อิ)} อคลุ\-ตครตา\-ลี\-สก\-โลหิตจนฺทนา\-นุลิตฺต\-คตฺโต นาคปุนฺนาค สาล สลฬ จมฺปก ยูถิกาติมุตฺตกปาฏลุปฺปลวสฺสิกมลฺลิกาวิจิตฺโต เสสชเน อติกฺกมิตฺวา วิโรจติ อติวิโรจติ โอภาสติ ปภาสติ สมฺปภาสติ ชลติ ปชฺชลติ อภิภวติ อชฺโฌตฺถรติ มาลาคนฺธ\-รตนา\-ภรเณหิ. เอวเมว โข มหาราช อรหตฺตํ ปตฺโต ขีณาสโว วิมุตฺติ\-รตน\-ปิฬนฺโธ อุปาทายุปาทาย วิมุตฺตานํ ภิกฺขูนํ อติกฺกมิตฺวา สมติกฺกมิตฺวา วิโรจติ อติวิโรจติ โอภาสติ ปภาสติ สมฺปภาสติ ชลติ ปชฺชลติ อภิภวติ อชฺโฌตฺถรติ วิมุตฺติยา, ตํ กิสฺส เหตุ, อคฺคํ มหาราช เอตํ ปิฬนฺธนํ สพฺพ\-ปิฬนฺธนานํ, ยทิทํ วิมุตฺติ\-ปิฬนฺธนํ, อิทํ วุจฺจติ มหาราช “ภควโต วิมุตฺติรตนนฺ”ติ.}

\setlength{\csromangathapagewidth}{0pt}
\setlength{\csromangathawakwidth}{0pt}
\csromangathameasuregroup{มณิมาลาธรํ เคห, \\ วิมุตฺติรตนมาลนฺตุ,}{\csromangathabat{มณิมาลาธรํ เคห,}{ชโน\textsuperscript{1} สามิํ อุทิกฺขติ.} \\ \csromangathabat{วิมุตฺติรตนมาลนฺตุ,}{อุทิกฺขนฺติ สเทวกาติ.}}
\csromangathasetleft{มณิมาลาธรํ เคห, \\ วิมุตฺติรตนมาลนฺตุ,}
\csromangathagroup{\csromangathastanza{\csromangathabat{มณิมาลาธรํ เคห,}{ชโน\csromansharedfootnote{bc78d1a585637ed6}{เคหํ, ชโน (ก)} สามิํ อุทิกฺขติ.} \\ \csromangathabat{วิมุตฺติรตนมาลนฺตุ,}{อุทิกฺขนฺติ สเทวกาติ.}}}

\prose{\csromanfolio{322}กตมํ มหาราช ภควโต วิมุตฺติญณทสฺสนรตนํ, ปจฺจเวกฺขณ\-ญาณํ มหาราช ภควโต วิมุตฺติญาณทสฺสนรตนนฺติ วุจฺจติ, เยน ญาเณน อริยสาวโก มคฺค\-ผลนิพฺพานานิ ปหีนกิเลสา\-วสิฏฺฐ\-กิเลเส จ ปจฺจเวกฺขติ.}

\setlength{\csromangathapagewidth}{0pt}
\setlength{\csromangathawakwidth}{0pt}
\csromangathameasuregroup{เยน ญาเณน พุชฺฌนฺติ, \\ ตํ ญาณรตนํ ลทฺธุํ,}{\csromangathabat{เยน ญาเณน พุชฺฌนฺติ,}{อริยา กตกิจฺจตํ.} \\ \csromangathabat{ตํ ญาณรตนํ ลทฺธุํ,}{วายเมถ ชิโนรสาติ.}}
\csromangathasetleft{เยน ญาเณน พุชฺฌนฺติ, \\ ตํ ญาณรตนํ ลทฺธุํ,}
\csromangathagroup{\csromangathastanza{\csromangathabat{เยน ญาเณน พุชฺฌนฺติ,}{อริยา กตกิจฺจตํ.} \\ \csromangathabat{ตํ ญาณรตนํ ลทฺธุํ,}{วายเมถ ชิโนรสาติ.}}}

\prose{กตมํ มหาราช ภควโต ปฏิสมฺภิทา\-รตนํ, จตสฺโส โข มหาราช ปฏิสมฺภิทาโย อตฺถ\-ปฏิสมฺภิทา ธมฺม\-ปฏิสมฺภิทา นิรุตฺติ\-ปฏิสมฺภิทา ปฏิภาน\-ปฏิสมฺภิทาติ, อิเมหิ โข มหาราช จตูหิ ปฏิสมฺภิทา\-รตเนหิ สมลงฺกโต ภิกฺขุ ยํ ยํ ปริสํ อุปสงฺกมติ, ยทิ ขตฺติย\-ปริสํ, ยทิ พฺราหฺมณ\-ปริสํ, ยทิ คหปติ\-ปริสํ, ยทิ สมณปริสํ, วิสารโท อุปสงฺกมติ อมงฺกุภูโต อภีรุ อจฺฉมฺภี อนุตฺราสี วิคต\-โลมหํโส ปริสํ อุปสงฺกมติ.}

\prose{ยถา มหาราช โยโธ สงฺคามสูโร สนฺนทฺธ\-ปญฺจา\-วุโธ อจฺฉมฺภิโต\csromansharedfootnote{8073eee923ed8de8}{อสมฺภีโต (สี, อิ)} สงฺคามํ โอตรติ, สเจ อมิตฺตา ทูเร ภวิสฺสนฺติ อุสุนา ปาตยิสฺสามิ, ตโต โอรโต ภวิสฺสนฺติ สตฺติยา ปหริสฺสามิ, ตโต โอรโต ภวิสฺสนฺติ กณเยน ปหริสฺสามิ, อุปคตํ สนฺตํ มณฺฑลคฺเคน ทฺวิธา ฉินฺทิสฺสามิ, กายูปคตํ ฉุริกาย วินิวิชฺฌิสฺสามีติ\csromansharedfootnote{d9d577c72f4cf18a}{วิชฺฌิสฺสามีติ (สี)}, เอวเมว โข มหาราช จตุปฏิสมฺภิทา\-รตน\-มณฺฑิโต ภิกฺขุ อจฺฉมฺภิโต ปริสํ อุปสงฺกมติ, โย โกจิ มํ อตฺถ\-ปฏิสมฺภิเท ปญฺหํ ปุจฺฉิสฺสติ, ตสฺส อตฺเถน อตฺถํ กถยิสฺสามิ, การเณน การณํ กถยิสฺสามิ, เหตุนา เหตุํ กถยิสฺสามิ, นเยน นยํ กถยิสฺสามิ, นิสฺสํสยํ กริสฺสามิ, วิมติํ วิเวเจสฺสามิ, โตสยิสฺสามิ ปญฺห\-เวยฺยากรเณน.}

\prose{โย โกจิ มํ ธมฺม\-ปฏิสมฺภิเท ปญฺหํ ปุจฺฉิสฺสติ, ตสฺส ธมฺเมน ธมฺมํ กถยิสฺสามิ, อมเตน อมตํ กถยิสฺสามิ, อสงฺขเตน อสงฺขตํ กถยิสฺสามิ, นิพฺพาเนน นิพฺพานํ กถยิสฺสามิ, สุญฺญเตน สุญฺญตํ กถยิสฺสามิ, อนิมิตฺเตน อนิมิตฺตํ กถยิสฺสามิ, อปฺปณิหิเตน อปฺปณิหิตํ กถยิสฺสามิ, อเนเชน อเนชํ กถยิสฺสามิ, นิสฺสํสยํ กริสฺสามิ, วิมติํ วิเวเจสฺสามิ, โตสยิสฺสามิ ปญฺห\-เวยฺยากรเณน.}

\prose{\csromanfolio{323}โย โกจิ มํ นิรุตฺติ\-ปฏิสมฺภิเท ปญฺหํ ปุจฺฉิสฺสติ, ตสฺส นิรุตฺติยา นิรุตฺติํ กถยิสฺสามิ, ปเทน ปทํ กถยิสฺสามิ, อนุปเทน อนุปทํ กถยิสฺสานิ, อกฺขเรน อกฺขรํ กถยิสฺสามิ, สนฺธิยา สนฺธิํ กถยิสฺสามิ, พฺยญฺชเนน พฺยญฺชนํ กถยิสฺสามิ, อนุพฺยญฺชเนน อนุพฺยญฺชนํ กถยิสฺสามิ, วณฺเณน วณฺณํ กถยิสฺสามิ, สเรน สรํ กถยิสฺสามิ, ปญฺญตฺติยา ปญฺญตฺติํ กถยิสฺสามิ, โวหาเรน โวหารํ กถยิสฺสามิ, นิสฺสํสยํ กริสฺสามิ, วิมติํ วิเวเจสฺสามิ, โตสยิสฺสามิ ปญฺห\-เวยฺยากรเณน.}

\prose{โย โกจิ มํ ปฏิภาน\-ปฏิสมฺภิเท ปญฺหํ ปุจฺฉิสฺสติ, ตสฺส ปฏิภาเนน ปฏิภานํ กถยิสฺสามิ, โอปมฺเมน โอปมฺมํ กถยิสฺสามิ, ลกฺขเณน ลกฺขณํ กถยิสฺสามิ, รเสน รสํ กถยิสฺสามิ, นิสฺสํสยํ กริสฺสามิ, วิมติํ วิเวเจสฺสามิ, โตสยิสฺสามิ ปญฺหเวยฺยากรเณนาติ, อิทํ วุจฺจติ มหาราช “ภควโต ปฏิสมฺภิทารตนนฺ”ติ.}

\setlength{\csromangathapagewidth}{0pt}
\setlength{\csromangathawakwidth}{0pt}
\csromangathameasuregroup{ปฏิสมฺภิทา กิณิตฺวาน, \\ อจฺฉมฺภิโต อนุพฺพิคฺโค,}{\csromangathabat{ปฏิสมฺภิทา กิณิตฺวาน,}{ญาเณน ผสฺสเยยฺย โย.} \\ \csromangathabat{อจฺฉมฺภิโต อนุพฺพิคฺโค,}{อติโรจติ สเทวเกติ.}}
\csromangathasetleft{ปฏิสมฺภิทา กิณิตฺวาน, \\ อจฺฉมฺภิโต อนุพฺพิคฺโค,}
\csromangathagroup{\csromangathastanza{\csromangathabat{ปฏิสมฺภิทา กิณิตฺวาน,}{ญาเณน ผสฺสเยยฺย โย.} \\ \csromangathabat{อจฺฉมฺภิโต อนุพฺพิคฺโค,}{อติโรจติ สเทวเกติ.}}}

\prose{กตมํ มหาราช ภควโต โพชฺฌงฺค\-รตนํ. สตฺติเม มหาราช โพชฺฌงฺคา, สติสมฺโพชฺฌงฺโค ธมฺมวิจย\-สมฺโพชฺฌงฺโค วีริย\-สมฺโพชฺฌงฺโค ปีติสมฺโพชฺฌงฺโค ปสฺสทฺธิ\-สมฺโพชฺฌงฺโค สมาธิ\-สมฺโพชฺฌงฺโค อุเปกฺขา\-สมฺโพชฺฌงฺโค, อิเมหิ โข มหาราช สตฺตหิ โพชฺฌงฺค\-รตเนหิ ปฏิมณฺฑิโต ภิกฺขุ สพฺพํ ตมํ อภิภุยฺย สเทวกํ โลกํ โอภาเสติ ปภาเสติ อาโลกํ ชเนติ, อิทํ วุจฺจติ มหาราช “ภควโต โพชฺฌงฺครตนนฺ”ติ.}

\setlength{\csromangathapagewidth}{0pt}
\setlength{\csromangathawakwidth}{0pt}
\csromangathameasuregroup{โพชฺฌงฺครตนมาลสฺส, \\ กมฺเมน ตํ กิณิตฺวาน,}{\csromangathabat{โพชฺฌงฺครตนมาลสฺส,}{อุฏฺฐหนฺติ\textsuperscript{1} สเทวกา.} \\ \csromangathabat{กมฺเมน ตํ กิณิตฺวาน,}{รตนํ โว ปิฬนฺธถาติ.}}
\csromangathasetleft{โพชฺฌงฺครตนมาลสฺส, \\ กมฺเมน ตํ กิณิตฺวาน,}
\csromangathagroup{\csromangathastanza{\csromangathabat{โพชฺฌงฺครตนมาลสฺส,}{อุฏฺฐหนฺติ\csromansharedfootnote{93fae70cd1b0614f}{อุปฏฺฐหนฺติ (ก), อุทิกฺขนฺติ (สฺยา)} สเทวกา.} \\ \csromangathabat{กมฺเมน ตํ กิณิตฺวาน,}{รตนํ โว ปิฬนฺธถาติ.}}}

\prose{ภนฺเต นาคเสน กตมํ พุทฺธสฺส ภควโต สพฺพาปณนฺติ. สพฺพาปณํ โข มหาราช ภควโต นวงฺคํ พุทฺธวจนํ สารีริกานิ ปาริโภคิกานิ เจติยานิ สํฆรตนญฺจ, สพฺพาปเณ มหาราช ภควตา ชาติสมฺปตฺติ ปสาริตา, โภคสมฺปตฺติ ปสาริตา, อายุสมฺปตฺติ ปสาริตา, อาโรคฺยสมฺปตฺติ ปสาริตา, วณฺณสมฺปตฺติ ปสาริตา, ปญฺญาสมฺปตฺติ ปสาริตา, มานุสิก\-สมฺปตฺติ ปสาริตา, ทิพฺพสมฺปตฺติ ปสาริตา,}

\prose{\csromanfolio{324}นิพฺพานสมฺปตฺติ ปสาริตา. ตตฺถ เย ตํ ตํ สมฺปตฺติํ อิจฺฉนฺติ, เต กมฺมมูลํ ทตฺวา ปตฺถิต\-ปตฺถิตํ สมฺปตฺติํ กิณนฺติ, เกจิ สีลสมาทาเนน กิณนฺติ, เกจิ อุโปสถ\-กมฺเมน กิณนฺติ, อปฺป\-มตฺต\-เกนปิ กมฺมมูเลน อุปาทายุปาทาย สมฺปตฺติโย ปฏิลภนฺติ. ยถา มหาราช อาปณิกสฺส อาปเณ ติลมุคฺคมาเส ปริตฺตเกนปิ ตณฺฑุล\-มุคฺค\-มาเสน อปฺปเกนปิ มูเลน อุปาทายุปาทาย คณฺหนฺติ. เอวเมว โข มหาราช ภควโต สพฺพาปเณ อปฺป\-มตฺต\-เกนปิ กมฺมมูเลน อุปาทายุปาทาย สมฺปตฺติโย ปฏิลภนฺติ. อิทํ วุจฺจติ มหาราช “ภควโต สพฺพปณนฺ”ติ.}

\setlength{\csromangathapagewidth}{0pt}
\setlength{\csromangathawakwidth}{0pt}
\csromangathameasuregroup{อายุ อโรคตา วณฺณํ, \\ อสงฺขตญฺจ อมตํ, \\ อปฺเปน พหุเกนาปิ, \\ กิณิตฺวา สทฺธามูเลน,}{\csromangathabat{อายุ อโรคตา วณฺณํ,}{สคฺคํ อุจฺจากุลีนตา.} \\ \csromangathabat{อสงฺขตญฺจ อมตํ,}{อตฺถิ สพฺพาปเณ ชิเน.} \\ \csromangathabat{อปฺเปน พหุเกนาปิ,}{กมฺมมูเลน คยฺหติ.} \\ \csromangathabat{กิณิตฺวา สทฺธามูเลน,}{สมิทฺธา โหถ ภิกฺขโวติ.}}
\csromangathasetleft{อายุ อโรคตา วณฺณํ, \\ อสงฺขตญฺจ อมตํ, \\ อปฺเปน พหุเกนาปิ, \\ กิณิตฺวา สทฺธามูเลน,}
\csromangathagroup{\csromangathastanza{\csromangathabat{อายุ อโรคตา วณฺณํ,}{สคฺคํ อุจฺจากุลีนตา.} \\ \csromangathabat{อสงฺขตญฺจ อมตํ,}{อตฺถิ สพฺพาปเณ ชิเน.}}\csromangathastanza{\csromangathabat{อปฺเปน พหุเกนาปิ,}{กมฺมมูเลน คยฺหติ.} \\ \csromangathabat{กิณิตฺวา สทฺธามูเลน,}{สมิทฺธา โหถ ภิกฺขโวติ.}}}

\prose{ภควโต โข มหาราช ธมฺมนคเร เอวรูปา ชนา ปฏิวสนฺติ, สุตฺตนฺติกา เวนยิกา อาภิธมฺมิกา ธมฺมกถิกา ชาตกภาณกา ทีฆภาณกา มชฺฌิม\-ภาณกา สํยุตฺต\-ภาณกา องฺคุตฺตร\-ภาณกา ขุทฺทก\-ภาณกา สีลสมฺปนฺนา สมาธิ\-สมฺปนฺนา ปญฺญาสมฺปนฺนา โพชฺฌงฺค\-ภาวนารตา วิปสฺสกา สทตฺถม\-นุยุตฺตา อารญฺญิกา รุกฺขมูลิกา อพฺโภกาสิกา ปลาลปุญฺชิกา โสสานิกา เนสชฺชิกา ปฏิปนฺนกา ผลฏฺฐา เสกฺขา ผลสมงฺคิโน โสตาปนฺนา สกทาคามิโน อนาคามิโน อรหนฺโต เตวิชฺชา ฉฬภิญฺญา อิทฺธิมนฺโต ปญฺญาย ปารมิํคตา สติปฏฺฐาน\-สมฺมปฺปธาน- อิทฺธิปาท- อินฺทฺริยพลโพชฺฌงฺคมคฺควรฌานวิโมกฺขรูปารูปสนฺตสุขสมาปตฺติ กุสลา, เตหิ อรหนฺเตหิ อากุลํ สมากุลํ อากิณฺณํ สมากิณฺณํ นฬวน\-สรวนมิว ธมฺมนครํ อโหสิ. ภวตีห\csromandash{}}

\setlength{\csromangathapagewidth}{0pt}
\setlength{\csromangathawakwidth}{0pt}
\csromangathameasuregroup{วีตราคา วีตโทสา, \\ วีตตณฺหา อนาทานา, \\ อารญฺญิกา ธุตธรา, \\ วิเวกาภิรตา ธีรา, \\ เนสชฺชิกา สนฺถติกา, \\ ปํสุกูลธรา สพฺเพ, \\ ติจีวรธรา สนฺตา, \\ รตา เอกาสเน วิญฺญู, \\ อปฺปิจฺฉา นิปกา ธีรา, \\ ลาภาลาเภน สนฺตุฏฺฐา, \\ ฌายี ฌานรตา ธีรา, \\ อากิญฺจญฺญํ ปตฺถยานา, \\ ปฏิปนฺนา ผลฏฺฐา จ, \\ อาสีสกา\textsuperscript{1} อุตฺตมตฺถํ, \\ โสตาปนฺนา จ วิมลา, \\ อนาคามี จ อรหนฺโต, \\ สติปฏฺฐานกุสลา, \\ วิปสฺสกา ธมฺมธรา, \\ อิทฺธิปาเทสุ กุสลา, \\ สมฺมปฺปธานานุยุตฺตา, \\ อภิญฺญาปารมิปฺปตฺตา, \\ อนฺตลิกฺขมฺหิ จรณา, \\ โอกฺขิตฺตจกฺขู มิตภาณี, \\ สุทนฺโต อุตฺตเม ทมฺเม\textsuperscript{1}, \\ เตวิชฺชา ฉฬภิญฺญา จ, \\ ปญฺญาย ปารมิปฺปตฺตา,}{\csromangathabat{วีตราคา วีตโทสา,}{วีตโมหา อนาสวา.} \\ \csromangathabat{วีตตณฺหา อนาทานา,}{ธมฺมนคเร วสนฺติ เต.} \\ \csromangathabat{อารญฺญิกา ธุตธรา,}{ฌายิโน ลูขจีวรา.} \\ \csromangathabat{วิเวกาภิรตา ธีรา,}{ธมฺมนคเร วสนฺติ เต.} \\ \csromangathabat{เนสชฺชิกา สนฺถติกา,}{อโถปิ ฐานจงฺกมา.} \\ \csromangathabat{ปํสุกูลธรา สพฺเพ,}{ธมฺมนคเร วสนฺติ เต.} \\ \csromangathabat{ติจีวรธรา สนฺตา,}{จมฺมขณฺฑจตุตฺถกา.} \\ \csromangathabat{รตา เอกาสเน วิญฺญู,}{ธมฺมนคเร วสนฺติ เต.} \\ \csromangathabat{อปฺปิจฺฉา นิปกา ธีรา,}{อปฺปาหารา อโลลุปา.} \\ \csromangathabat{ลาภาลาเภน สนฺตุฏฺฐา,}{ธมฺมนคเร วสนฺติ เต.} \\ \csromangathabat{ฌายี ฌานรตา ธีรา,}{สนฺตจิตฺตา สมาหิตา.} \\ \csromangathabat{อากิญฺจญฺญํ ปตฺถยานา,}{ธมฺมนคเร วสนฺติ เต.} \\ \csromangathabat{ปฏิปนฺนา ผลฏฺฐา จ,}{เสกฺขา ผลสมงฺคิโน.} \\ \csromangathabat{อาสีสกา\textsuperscript{1} อุตฺตมตฺถํ,}{ธมฺมนคเร วสนฺติ เต.} \\ \csromangathabat{โสตาปนฺนา จ วิมลา,}{สกทาคามิโน จ เย.} \\ \csromangathabat{อนาคามี จ อรหนฺโต,}{ธมฺมนคเร วสนฺติ เต.} \\ \csromangathabat{สติปฏฺฐานกุสลา,}{โพชฺฌงฺคภาวนารตา.} \\ \csromangathabat{วิปสฺสกา ธมฺมธรา,}{ธมฺมนคเร วสนฺติ เต.} \\ \csromangathabat{อิทฺธิปาเทสุ กุสลา,}{สมาธิภาวนารตา.} \\ \csromangathabat{สมฺมปฺปธานานุยุตฺตา,}{ธมฺมนคเร วสนฺติ เต.} \\ \csromangathabat{อภิญฺญาปารมิปฺปตฺตา,}{เปตฺติเก โคจเร รตา.} \\ \csromangathabat{อนฺตลิกฺขมฺหิ จรณา,}{ธมฺมนคเร วสนฺติ เต.} \\ \csromangathabat{โอกฺขิตฺตจกฺขู มิตภาณี,}{คุตฺตทฺวารา สุสํวุตา.} \\ \csromangathabat{สุทนฺโต อุตฺตเม ทมฺเม\textsuperscript{1},}{ธมฺมนคเร วสนฺติ เต.} \\ \csromangathabat{เตวิชฺชา ฉฬภิญฺญา จ,}{อิทฺธิยา ปารมิํ คตา.} \\ \csromangathabat{ปญฺญาย ปารมิปฺปตฺตา,}{ธมฺมนคเร วสนฺติ เตติ.}}
\csromangathasetleft{วีตราคา วีตโทสา, \\ วีตตณฺหา อนาทานา, \\ อารญฺญิกา ธุตธรา, \\ วิเวกาภิรตา ธีรา, \\ เนสชฺชิกา สนฺถติกา, \\ ปํสุกูลธรา สพฺเพ, \\ ติจีวรธรา สนฺตา, \\ รตา เอกาสเน วิญฺญู, \\ อปฺปิจฺฉา นิปกา ธีรา, \\ ลาภาลาเภน สนฺตุฏฺฐา, \\ ฌายี ฌานรตา ธีรา, \\ อากิญฺจญฺญํ ปตฺถยานา, \\ ปฏิปนฺนา ผลฏฺฐา จ, \\ อาสีสกา\textsuperscript{1} อุตฺตมตฺถํ, \\ โสตาปนฺนา จ วิมลา, \\ อนาคามี จ อรหนฺโต, \\ สติปฏฺฐานกุสลา, \\ วิปสฺสกา ธมฺมธรา, \\ อิทฺธิปาเทสุ กุสลา, \\ สมฺมปฺปธานานุยุตฺตา, \\ อภิญฺญาปารมิปฺปตฺตา, \\ อนฺตลิกฺขมฺหิ จรณา, \\ โอกฺขิตฺตจกฺขู มิตภาณี, \\ สุทนฺโต อุตฺตเม ทมฺเม\textsuperscript{1}, \\ เตวิชฺชา ฉฬภิญฺญา จ, \\ ปญฺญาย ปารมิปฺปตฺตา,}
\csromangathagroup{\csromangathastanza{\csromangathabat{วีตราคา วีตโทสา,}{วีตโมหา อนาสวา.} \\ \csromangathabat{วีตตณฺหา อนาทานา,}{ธมฺมนคเร วสนฺติ เต.}}\csromangathastanza{\csromangathabat{อารญฺญิกา ธุตธรา,}{ฌายิโน ลูขจีวรา.} \\ \csromangathabat{วิเวกาภิรตา ธีรา,}{ธมฺมนคเร วสนฺติ เต.}}\csromangathastanza{\csromangathabat{เนสชฺชิกา สนฺถติกา,}{อโถปิ ฐานจงฺกมา.} \\ \csromangathabat{ปํสุกูลธรา สพฺเพ,}{ธมฺมนคเร วสนฺติ เต.}}\csromangathastanza{\csromanfolio{325}\csromangathabat{ติจีวรธรา สนฺตา,}{จมฺมขณฺฑจตุตฺถกา.} \\ \csromangathabat{รตา เอกาสเน วิญฺญู,}{ธมฺมนคเร วสนฺติ เต.}}\csromangathastanza{\csromangathabat{อปฺปิจฺฉา นิปกา ธีรา,}{อปฺปาหารา อโลลุปา.} \\ \csromangathabat{ลาภาลาเภน สนฺตุฏฺฐา,}{ธมฺมนคเร วสนฺติ เต.}}\csromangathastanza{\csromangathabat{ฌายี ฌานรตา ธีรา,}{สนฺตจิตฺตา สมาหิตา.} \\ \csromangathabat{อากิญฺจญฺญํ ปตฺถยานา,}{ธมฺมนคเร วสนฺติ เต.}}\csromangathastanza{\csromangathabat{ปฏิปนฺนา ผลฏฺฐา จ,}{เสกฺขา ผลสมงฺคิโน.} \\ \csromangathabat{อาสีสกา\csromansharedfootnote{63ba153d8aba474d}{อาสิํสกา (สี, อิ)} อุตฺตมตฺถํ,}{ธมฺมนคเร วสนฺติ เต.}}\csromangathastanza{\csromangathabat{โสตาปนฺนา จ วิมลา,}{สกทาคามิโน จ เย.} \\ \csromangathabat{อนาคามี จ อรหนฺโต,}{ธมฺมนคเร วสนฺติ เต.}}\csromangathastanza{\csromangathabat{สติปฏฺฐาน\-กุสลา,}{โพชฺฌงฺค\-ภาวนารตา.} \\ \csromangathabat{วิปสฺสกา ธมฺมธรา,}{ธมฺมนคเร วสนฺติ เต.}}\csromangathastanza{\csromangathabat{อิทฺธิปาเทสุ กุสลา,}{สมาธิ\-ภาวนารตา.} \\ \csromangathabat{สมฺมปฺปธานานุยุตฺตา,}{ธมฺมนคเร วสนฺติ เต.}}\csromangathastanza{\csromangathabat{อภิญฺญาปารมิ\-ปฺปตฺตา,}{เปตฺติเก โคจเร รตา.} \\ \csromangathabat{อนฺตลิกฺขมฺหิ จรณา,}{ธมฺมนคเร วสนฺติ เต.}}\csromangathastanza{\csromangathabat{โอกฺขิตฺตจกฺขู มิตภาณี,}{คุตฺตทฺวารา สุสํวุตา.} \\ \csromangathabat{สุทนฺโต อุตฺตเม ทมฺเม\csromansharedfootnote{19a4578373188b83}{ธมฺเม (สี, อิ)},}{ธมฺมนคเร วสนฺติ เต.}}\csromangathastanza{\csromangathabat{เตวิชฺชา ฉฬภิญฺญา จ,}{อิทฺธิยา ปารมิํ คตา.} \\ \csromangathabat{ปญฺญาย ปารมิปฺปตฺตา,}{ธมฺมนคเร วสนฺติ เตติ.}}}

\prose{เย โข เต มหาราช ภิกฺขู อปริมิต\-ญาณ\-วรธรา อสงฺคา อตุลคุณา\csromansharedfootnote{11f821419ac3c868}{อตุลิยคุณา (สี, อิ, ก)} อตุลยสา อตุลพลา อตุลเตชา ธมฺมจกฺกา\-นุ\-ปฺปวตฺตกา ปญฺญาปารมิํ คตา, เอวรูปา โข มหาราช ภิกฺขู ภควโต ธมฺมนคเร \textbf{“ธมฺม}\-\textbf{เสนาปติโน”}ติ วุจฺจนฺติ.}

\prose{เย ปน เต มหาราช ภิกฺขู อิทฺธิมนฺโต อธิคต\-ปฺปฏิสมฺภิทาปตฺต\-เวสารชฺชา คคนจรา ทุราสทา ทุปฺปสหา อนาลมฺพจรา สสา\-คร\-มหิ\-ธร\-ปถวิ\-กมฺปกา จนฺทสูริย\-ปริมชฺชกา วิกุพฺพนา\-ธิฏฺฐานา\-ภินีหารกุสลา}

\prose{\csromanfolio{326}อิทฺธิยา ปารมิํ คตา, เอวรูปา โข มหาราช ภิกฺขู ภควโต ธมฺมนคเร \textbf{“ปุโรหิตา”ติ} วุจฺจนฺติ.}

\prose{เย ปน เต มหาราช ภิกฺขู ธุตงฺคม\-นุคตา อปฺปิจฺฉา สนฺตุฏฺฐา วิญฺญตฺติม\-เนสน\-ชิคุจฺฉกา ปิณฺฑาย สปทานจาริโน ภมราว คนฺธม\-นุ\-ฆายิตฺวา ปวิสนฺติ วิวิตฺตกานนํ, กาเย จ ชีวิเต จ นิรเปกฺขา อรหตฺตม\-นุปฺปตฺตา ธุตงฺคคุเณ อคฺคนิกฺขิตฺตา, เอวรูปา โข มหาราช ภิกฺขู ภควโต ธมฺมนคเร \textbf{“อกฺขทสฺสา”ติ} วุจฺจนฺติ.}

\prose{เย ปน เต มหาราช ภิกฺขู ปริสุทฺธา วิมลา นิกฺกิเลสา จุตูปปาต\-กุสลา ทิพฺพจกฺขุมฺหิ ปารมิํ คตา, เอวรูปา โข มหาร ภิกฺขู ภควโต ธมฺมนคเร \textbf{“นครโชตกา”ติ} วุจฺจนฺติ.}

\prose{เย ปน เต มหาราช ภิกฺขู พหุสฺสุตา อาคตาคมา ธมฺมธรา วินยธรา มาติกาธรา สิถิล\-ธนิต\-ทีฆ\-รสฺส\-ครุกลหุก\-กฺขร\-ปริจฺเฉท\-กุสลา นวงฺค\-สาสน\-ธรา, เอวรูปา โข มหาราช ภิกฺขู ภควโต ธมฺมนคเร \textbf{“ธมฺมรกฺขา”ติ} วุจฺจนฺติ.}

\prose{เย ปน เต มหาราช ภิกฺขู วินยญฺญู วินยโกวิทา ฐานาฏฺฐาน\-กุสลา\csromansharedfootnote{ac25714d0f1c8feb}{นิทานปฐนกุสลา (สี, อิ), นิทานวตฺถุกุสลา (สฺยา)} อาปตฺตานาปตฺติครูกลหุกสเตกิจฺฉ- อเตกิจฺฉวุฏฺฐานเทสนานิคฺคหปฏิกมฺม- โอสารณนิสฺสารณปฏิสารณกุสลา วินเย ปารมิํ คตา, เอวรูปา โข มหาราช ภิกฺขู ภควโต ธมฺมนคเร \textbf{“รูปรกฺขา”ติ}\csromansharedfootnote{fc255596a6c88bcc}{รูปทกฺขาติ (สี, สฺยา, อิ)} วุจฺจนฺติ.}

\prose{เย ปน เต มหาราช ภิกฺขู วิมุตฺติ\-วร\-กุสุม\-มาล\-พทฺธา วรปวร\-มหคฺฆ\-เสฏฺฐภาวม\-นุปฺปตฺตา พหุชนกนฺตม\-ภิปตฺถิตา, เอวรูปา โข มหาราช ภิกฺขู ภควโต ธมฺมนคเร \textbf{“ปุปฺผาปณิกา”ติ} วุจฺจนฺติ.}

\prose{เย ปน เต มหาราช ภิกฺขู จตุสจฺจาภิสมย\-ปฺปฏิวิทฺธา ทิฏฺฐสจฺจา วิญฺญาตสาสนา จตูสุ สามญฺญผเลสุ ติณฺณ\-วิจิกิจฺฉา ปฏิลทฺธ\-ผลสุขา อญฺเญสมฺปิ ปฏิปนฺนานํ เต ผเล สํวิภชนฺติ, เอวรูปา โข มหาราช ภิกฺขู ภควโต ธมฺมนคเร \textbf{“ผลาปณิกา”ติ} วุจฺจนฺติ.}

\prose{\csromanfolio{327}เย ปน เต มหาราช ภิกฺขู สีลสํวร\-คนฺธม\-นุลิตฺตา\csromansharedfootnote{5bd5faaad383b9fd}{สีลวรสุคนฺธมนุลิตฺตา (สี, อิ)} อเนกวิธ\-พหุ\-คุณ\-ธรา กิเลสมล\-ทุคฺคนฺธ\-วิธมกา, เอวรูปา โข มหาราช ภิกฺขู ภควโต ธมฺมนคเร \textbf{“คนฺธาปณิกา”ติ} วุจฺจนฺติ.}

\prose{เย ปน เต มหาราช ภิกฺขู ธมฺมกามา ปิยสมุทาหารา อภิธมฺเม อภิวินเย อุฬารปาโมชฺชา อรญฺญคตาปิ รุกฺขมูลคตาปิ สุญฺญาคารคตาปิ ธมฺมวรรสํ ปิวนฺติ, กาเยน วาจาย มนสา ธมฺมวร\-รสโม\-คาฬฺหา อธิมตฺต\-ปฏิภานา ธมฺเมสุ ธมฺเมสน\-ปฺปฏิปนฺนา อิโต วา ตโต วา ยตฺถ ยตฺถ อปฺปิจฺฉกถา สนฺตุฏฺฐิกถา ปวิเวกกถา อสํสคฺคกถา วีริยา\-รมฺภ\-กถา สีลกถา สมาธิกถา ปญฺญากถา วิมุตฺติกถา วิมุตฺติ\-ญาณ\-ทสฺสน\-กถา, ตตฺถ ตตฺถ คนฺตฺวา ตํ ตํ กถารสํ ปิวนฺติ, เอวรูปา โข มหาราช ภิกฺขู ภควโต ธมฺมนคเร \textbf{“โสณฺฑา ปิปาสา”ติ} วุจฺจนฺติ.}

\prose{เย ปน เต มหาราช ภิกฺขู ปุพฺพรตฺตาปรรตฺตํ ชาคริยา\-นุโยคม\-นุยุตฺตา นิสชฺชฏฺฐาน\-จงฺกเมหิ รตฺตินฺทิวํ วีตินาเมนฺติ, ภาวนา\-นุโยคม\-นุยุตฺตา กิเลส\-ปฏิพาหนาย สทตฺถ\-ปฺปสุตา, เอวรูปา โข มหาราช ภิกฺขู ภควโต ธมฺมนคเร \textbf{“นครคุตฺติกา”ติ} วุจฺจนฺติ.}

\prose{เย ปน เต มหาราช ภิกฺขู นวงฺคํ พุทฺธวจนํ อตฺถโต จ พฺยญฺชนโต จ นยโต จ การณโต จ เหตุโต จ อุทาหรณโต จ วาเจนฺติ อนุวาเจนฺติ ภาสนฺติ อนุภาสนฺติ, เอวรูปา โข มหาราช ภิกฺขู ภควโต ธมฺมนคเร \textbf{“ธมฺมาปณิกา”ติ} วุจฺจนฺติ.}

\prose{เย ปน เต มหาราช ภิกฺขู ธมฺมรตน\-โภเคน อาคม\-ปริยตฺติ\-สุต\-โภเคน โภคิโน ธนิโน นิทฺทิฏฺฐ\-สร\-พฺยญฺชน\-ลกฺขณ\-ปฺปฏิเวธา วิญฺญู ผรณา, เอวรูปา โข มหาราช ภิกฺขู ภควโต ธมฺมนคเร \textbf{“ธมฺมเสฏฺฐิโน”ติ} วุจฺจนฺติ.}

\prose{เย ปน เต มหาราช ภิกฺขู อุฬาร\-เทสนา\-ปฏิเวธา ปริจิณฺณา\-รมฺมณ\-วิภตฺติ\-นิทฺเทสา สิกฺขา\-คุณ\-ปารมิปฺปตฺตา, เอวรูปา โข มหาราช ภิกฺขู ภควโต ธมฺมนคเร \textbf{“วิสฺสุต}\-\textbf{ธมฺมิกา”ติ} วุจฺจนฺติ.}

\prose{เอวํ สุวิภตฺตํ โข มหาราช ภควโต ธมฺมนครํ เอวํ สุมาปิตํ เอวํ สุวิหิตํ เอวํ สุปริปูริตํ เอวํ สุววตฺถาปิตํ เอวํ สุรกฺขิตํ เอวํ สุโคปิตํ}

\prose{\csromanfolio{328}เอวํ ทุปฺปสยฺหํ ปจฺจตฺถิเกหิ ปจฺจามิตฺเตหิ, อิมินา มหาราช การเณน อิมินา เหตุนา อิมินา นเยน อิมินา อนุมาเนน ญาตพฺพํ อตฺถิ โส ภควาติ.}

\setlength{\csromangathapagewidth}{0pt}
\setlength{\csromangathawakwidth}{0pt}
\csromangathameasuregroup{“ยถาปิ นครํ ทิสฺวา, \\ อนุมาเนน ชานนฺติ, \\ ตเถว โลกนาถสฺส, \\ อนุมาเนน ชานนฺติ, \\ อนุมาเนน ชานนฺติ, \\ ยถายํ ทิสฺสเต อูมิ, \\ ตถา พุทฺธํ โสกนุทํ, \\ ตณฺหกฺขยมนุปตฺตํ, \\ อนุมาเนน ญาตพฺพํ, \\ ยถา ธมฺมูมิวิปฺผาโร, \\ อนุมาเนน ชานนฺติ, \\ ยถา อจฺจุคฺคโต เอโส, \\ ตถา ทิสฺวา ธมฺมคิริํ, \\ อจฺจุคฺคตํ ภควโต, \\ อนุมาเนน ญาตพฺพํ, \\ ตถา หิ โส มหาวีโร, \\ ยถาปิ คชราชสฺส, \\ อนุมาเนน ชานนฺติ, \\ ตเถว พุทฺธนาคสฺส, \\ อนุมาเนน ชานนฺติ, \\ อนุมาเนน ชานนฺติ, \\ มิคราชสฺส สทฺเทน, \\ ตเถว ติตฺถิเย ทิสฺวา, \\ อนุมาเนน ญาตพฺพํ, \\ นิพฺพุตํ ปถวิํ ทิสฺวา, \\ อนุมาเนน ชานนฺติ, \\ ตเถวิมํ ชนํ ทิสฺวา, \\ อนุมาเนน ญาตพฺพํ, \\ ลคฺคํ ทิสฺวา ภุสํ ปงฺกํ, \\ อนุมาเนน ชานนฺติ, \\ ตเถวิมํ ชนํ ทิสฺวา, \\ วหิตํ ธมฺมนทิยา, \\ ธมฺมามตคตํ ทิสฺวา, \\ อนุมาเนน ญาตพฺพํ, \\ อนุมาเนน ชานนฺติ, \\ ยถายํ วายเต คนฺโธ, \\ ตเถวายํ สีลคนฺโธ, \\ อนุมาเนน ญาตพฺพํ,}{\csromangathabat{“ยถาปิ นครํ ทิสฺวา,}{สุวิภตฺตํ มโนรมํ.} \\ \csromangathabat{อนุมาเนน ชานนฺติ,}{วฑฺฒกิสฺส มหตฺตนํ.} \\ \csromangathabat{ตเถว โลกนาถสฺส,}{ทิสฺวา ธมฺมปุรํ วรํ.} \\ \csromangathabat{อนุมาเนน ชานนฺติ,}{อตฺถิ โส ภควา อิติ.} \\ \csromangathabat{อนุมาเนน ชานนฺติ,}{อูมิํ ทิสฺวาน สาคเร.} \\ \csromangathabat{ยถายํ ทิสฺสเต อูมิ,}{มหนฺโต โส ภวิสฺสติ.} \\ \csromangathabat{ตถา พุทฺธํ โสกนุทํ,}{สพฺพตฺถ มปราชิตํ.} \\ \csromangathabat{ตณฺหกฺขยมนุปตฺตํ,}{ภวสํสารโมจนํ.} \\ \csromangathabat{อนุมาเนน ญาตพฺพํ,}{อูมิํ ทิสฺวา สเทวเก.} \\ \csromangathabat{ยถา ธมฺมูมิวิปฺผาโร,}{อคฺโค พุทฺโธ ภวิสฺสติ.} \\ \csromangathabat{อนุมาเนน ชานนฺติ,}{ทิสฺวา อจฺจุคฺคตํ คิริํ.} \\ \csromangathabat{ยถา อจฺจุคฺคโต เอโส,}{หิมวา โส ภวิสฺสติ.} \\ \csromangathabat{ตถา ทิสฺวา ธมฺมคิริํ,}{สีตีภูตํ นิรูปธิํ.} \\ \csromangathabat{อจฺจุคฺคตํ ภควโต,}{อจลํ สุปฺปติฏฺฐิตํ.} \\ \csromangathabat{อนุมาเนน ญาตพฺพํ,}{ทิสฺวาน ธมฺมปพฺพตํ.} \\ \csromangathabat{ตถา หิ โส มหาวีโร,}{อคฺโค พุทฺโธ ภวิสฺสติ.} \\ \csromangathabat{ยถาปิ คชราชสฺส,}{ปทํ ทิสฺวาน มานุสา.} \\ \csromangathabat{อนุมาเนน ชานนฺติ,}{มหา เอโส คโช อิติ.} \\ \csromangathabat{ตเถว พุทฺธนาคสฺส,}{ปทํ ทิสฺวา วิภาวิโน.} \\ \csromangathabat{อนุมาเนน ชานนฺติ,}{อุฬาโร โส ภวิสฺสติ.} \\ \csromangathabat{อนุมาเนน ชานนฺติ,}{ภีเต ทิสฺวาน กุมฺมิเค.} \\ \csromangathabat{มิคราชสฺส สทฺเทน,}{ภีตาเม กุมฺมิคา อิติ.} \\ \csromangathabat{ตเถว ติตฺถิเย ทิสฺวา,}{วิตฺถเต ภีตมานเส.} \\ \csromangathabat{อนุมาเนน ญาตพฺพํ,}{ธมฺมราเชน คชฺชิตํ.} \\ \csromangathabat{นิพฺพุตํ ปถวิํ ทิสฺวา,}{หริตปตฺตํ มโหทิกํ.} \\ \csromangathabat{อนุมาเนน ชานนฺติ,}{มหาเมเฆน นิพฺพุตํ.} \\ \csromangathabat{ตเถวิมํ ชนํ ทิสฺวา,}{อาโมทิตปโมทิตํ.} \\ \csromangathabat{อนุมาเนน ญาตพฺพํ,}{ธมฺมเมเฆน ตปฺปิตํ.} \\ \csromangathabat{ลคฺคํ ทิสฺวา ภุสํ ปงฺกํ,}{กลลทฺทคตํ มหิํ.} \\ \csromangathabat{อนุมาเนน ชานนฺติ,}{วาริกฺขนฺโธ มหา คโต.} \\ \csromangathabat{ตเถวิมํ ชนํ ทิสฺวา,}{รชปงฺกสโมหิตํ.} \\ \csromangathabat{วหิตํ ธมฺมนทิยา,}{วิสฏฺฐํ ธมฺมสาคเร.} \\ \csromangathabat{ธมฺมามตคตํ ทิสฺวา,}{สเทวกมิมํ มหิํ.} \\ \csromangathabat{อนุมาเนน ญาตพฺพํ,}{ธมฺมกฺขนฺโธ มหา คโต.} \\ \csromangathabat{อนุมาเนน ชานนฺติ,}{ฆายิตฺวา คนฺธมุตฺตมํ.} \\ \csromangathabat{ยถายํ วายเต คนฺโธ,}{เหสฺสนฺติ ปุปฺผิตา ทุมา.} \\ \csromangathabat{ตเถวายํ สีลคนฺโธ,}{ปวายติ สเทวเก.} \\ \csromangathabat{อนุมาเนน ญาตพฺพํ,}{อตฺถิ พุทฺโธ อนุตฺตโร”ติ.}}
\csromangathasetleft{“ยถาปิ นครํ ทิสฺวา, \\ อนุมาเนน ชานนฺติ, \\ ตเถว โลกนาถสฺส, \\ อนุมาเนน ชานนฺติ, \\ อนุมาเนน ชานนฺติ, \\ ยถายํ ทิสฺสเต อูมิ, \\ ตถา พุทฺธํ โสกนุทํ, \\ ตณฺหกฺขยมนุปตฺตํ, \\ อนุมาเนน ญาตพฺพํ, \\ ยถา ธมฺมูมิวิปฺผาโร, \\ อนุมาเนน ชานนฺติ, \\ ยถา อจฺจุคฺคโต เอโส, \\ ตถา ทิสฺวา ธมฺมคิริํ, \\ อจฺจุคฺคตํ ภควโต, \\ อนุมาเนน ญาตพฺพํ, \\ ตถา หิ โส มหาวีโร, \\ ยถาปิ คชราชสฺส, \\ อนุมาเนน ชานนฺติ, \\ ตเถว พุทฺธนาคสฺส, \\ อนุมาเนน ชานนฺติ, \\ อนุมาเนน ชานนฺติ, \\ มิคราชสฺส สทฺเทน, \\ ตเถว ติตฺถิเย ทิสฺวา, \\ อนุมาเนน ญาตพฺพํ, \\ นิพฺพุตํ ปถวิํ ทิสฺวา, \\ อนุมาเนน ชานนฺติ, \\ ตเถวิมํ ชนํ ทิสฺวา, \\ อนุมาเนน ญาตพฺพํ, \\ ลคฺคํ ทิสฺวา ภุสํ ปงฺกํ, \\ อนุมาเนน ชานนฺติ, \\ ตเถวิมํ ชนํ ทิสฺวา, \\ วหิตํ ธมฺมนทิยา, \\ ธมฺมามตคตํ ทิสฺวา, \\ อนุมาเนน ญาตพฺพํ, \\ อนุมาเนน ชานนฺติ, \\ ยถายํ วายเต คนฺโธ, \\ ตเถวายํ สีลคนฺโธ, \\ อนุมาเนน ญาตพฺพํ,}
\csromangathagroup{\csromangathastanza{\csromangathabat{“ยถาปิ นครํ ทิสฺวา,}{สุวิภตฺตํ มโนรมํ.} \\ \csromangathabat{อนุมาเนน ชานนฺติ,}{วฑฺฒกิสฺส มหตฺตนํ.}}\csromangathastanza{\csromangathabat{ตเถว โลกนาถสฺส,}{ทิสฺวา ธมฺมปุรํ วรํ.} \\ \csromangathabat{อนุมาเนน ชานนฺติ,}{อตฺถิ โส ภควา อิติ.}}\csromangathastanza{\csromangathabat{อนุมาเนน ชานนฺติ,}{อูมิํ ทิสฺวาน สาคเร.} \\ \csromangathabat{ยถายํ ทิสฺสเต อูมิ,}{มหนฺโต โส ภวิสฺสติ.}}\csromangathastanza{\csromangathabat{ตถา พุทฺธํ โสกนุทํ,}{สพฺพตฺถ มปราชิตํ.} \\ \csromangathabat{ตณฺหกฺขยมนุปตฺตํ,}{ภวสํสาร\-โมจนํ.}}\csromangathastanza{\csromangathabat{อนุมาเนน ญาตพฺพํ,}{อูมิํ ทิสฺวา สเทวเก.} \\ \csromangathabat{ยถา ธมฺมูมิวิปฺผาโร,}{อคฺโค พุทฺโธ ภวิสฺสติ.}}\csromangathastanza{\csromangathabat{อนุมาเนน ชานนฺติ,}{ทิสฺวา อจฺจุคฺคตํ คิริํ.} \\ \csromangathabat{ยถา อจฺจุคฺคโต เอโส,}{หิมวา โส ภวิสฺสติ.}}\csromangathastanza{\csromangathabat{ตถา ทิสฺวา ธมฺมคิริํ,}{สีตีภูตํ นิรูปธิํ.} \\ \csromangathabat{อจฺจุคฺคตํ ภควโต,}{อจลํ สุปฺปติฏฺฐิตํ.}}\csromangathastanza{\csromangathabat{อนุมาเนน ญาตพฺพํ,}{ทิสฺวาน ธมฺมปพฺพตํ.} \\ \csromangathabat{ตถา หิ โส มหาวีโร,}{อคฺโค พุทฺโธ ภวิสฺสติ.}}\csromangathastanza{\csromangathabat{ยถาปิ คชราชสฺส,}{ปทํ ทิสฺวาน มานุสา.} \\ \csromangathabat{อนุมาเนน ชานนฺติ,}{มหา เอโส คโช อิติ.}}\csromangathastanza{\csromangathabat{ตเถว พุทฺธนาคสฺส,}{ปทํ ทิสฺวา วิภาวิโน.} \\ \csromangathabat{อนุมาเนน ชานนฺติ,}{อุฬาโร โส ภวิสฺสติ.}}\csromangathastanza{\csromangathabat{อนุมาเนน ชานนฺติ,}{ภีเต ทิสฺวาน กุมฺมิเค.} \\ \csromangathabat{มิคราชสฺส สทฺเทน,}{ภีตาเม กุมฺมิคา อิติ.}}\csromangathastanza{\csromangathabat{ตเถว ติตฺถิเย ทิสฺวา,}{วิตฺถเต ภีตมานเส.} \\ \csromangathabat{อนุมาเนน ญาตพฺพํ,}{ธมฺมราเชน คชฺชิตํ.}}\csromangathastanza{\csromanfolio{329}\csromangathabat{นิพฺพุตํ ปถวิํ ทิสฺวา,}{หริตปตฺตํ มโหทิกํ.} \\ \csromangathabat{อนุมาเนน ชานนฺติ,}{มหาเมเฆน นิพฺพุตํ.}}\csromangathastanza{\csromangathabat{ตเถวิมํ ชนํ ทิสฺวา,}{อาโมทิตปโมทิตํ.} \\ \csromangathabat{อนุมาเนน ญาตพฺพํ,}{ธมฺมเมเฆน ตปฺปิตํ.}}\csromangathastanza{\csromangathabat{ลคฺคํ ทิสฺวา ภุสํ ปงฺกํ,}{กลลทฺทคตํ มหิํ.} \\ \csromangathabat{อนุมาเนน ชานนฺติ,}{วาริกฺขนฺโธ มหา คโต.}}\csromangathastanza{\csromangathabat{ตเถวิมํ ชนํ ทิสฺวา,}{รชปงฺกสโมหิตํ.} \\ \csromangathabat{วหิตํ ธมฺมนทิยา,}{วิสฏฺฐํ ธมฺมสาคเร.}}\csromangathastanza{\csromangathabat{ธมฺมามตคตํ ทิสฺวา,}{สเทวกมิมํ มหิํ.} \\ \csromangathabat{อนุมาเนน ญาตพฺพํ,}{ธมฺมกฺขนฺโธ มหา คโต.}}\csromangathastanza{\csromangathabat{อนุมาเนน ชานนฺติ,}{ฆายิตฺวา คนฺธมุตฺตมํ.} \\ \csromangathabat{ยถายํ วายเต คนฺโธ,}{เหสฺสนฺติ ปุปฺผิตา ทุมา.}}\csromangathastanza{\csromangathabat{ตเถวายํ สีลคนฺโธ,}{ปวายติ สเทวเก.} \\ \csromangathabat{อนุมาเนน ญาตพฺพํ,}{อตฺถิ พุทฺโธ อนุตฺตโร”ติ.}}}

\prose{เอวรูเปน โข มหาราช การณสเตน การณ\-สหสฺเสน เหตุสเตน เหตุสหสฺเสน นยสเตน นยสหสฺเสน โอปมฺมสเตน โอปมฺม\-สหสฺเสน สกฺกา พุทฺธพลํ อุปทสฺสยิตุํ, ยถา มหาราช ทกฺโข มาลากาโร นานา\-ปุปฺผราสิมฺหา อาจริยา\-นุสิฏฺฐิยา ปจฺจตฺต\-ปุริสกาเรน วิจิตฺตํ มาลาคุณราสิํ กเรยฺย, เอวเมว โข มหาราช โส ภควา วิจิตฺต\-ปุปฺผราสิ วิย อนนฺตคุโณ อปฺปเมยฺยคุโณ, อหเมตรหิ ชินสาสเน มาลากาโร วิย ปุปฺผคนฺถโก ปุพฺพกานํ อาจริยานํ มคฺเคนปิ มยฺหํ พุทฺธิพเลนปิ อสงฺเขฺยยฺเยนปิ การเณน อนุมาเนน พุทฺธพลํ ทีปยิสฺสามิ, ตฺวํ ปเนตฺถ ฉนฺทํ ชเนหิ สวนายาติ.}

\prose{ทุกฺกรํ ภนฺเต นาคเสน อญฺเญสํ เอวรูเปน การเณน อนุมาเนน พุทฺธพลํ อุปทสฺสยิตุํ, นิพฺพุโตสฺมิ ภนฺเต นาคเสน ตุมฺหากํ ปรม\-วิจิตฺเตน ปญฺหเวยฺยากรเณนาติ.}

\vspace{\tipitakaparskip}
\csromancenter{อนุมานปโญฺห ปฐโม.}
\csromanmatikaanchor{mtk.199}
\csromantocmarkref{subsubsection}{2. ธุตงฺคปญฺห}{mtk.199}
\bookmark[dest={mtk.199},level=3]{2. ธุตงฺคปญฺห}
\csromanheader{3}{2. \textbf{ธุตงฺคปญฺห}}
\setlength{\csromangathapagewidth}{0pt}
\setlength{\csromangathawakwidth}{0pt}
\csromangathameasuregroup{“ปสฺสตารญฺญเก ภิกฺขู, \\ ปุน ปสฺสติ คิหี ราชา, \\ อุโภปิ เต วิโลเกตฺวา, \\ พุชฺเฌยฺย เจ คิหี ธมฺเม, \\ ปรวาทิวาทมถนํ, \\ หนฺท ปุจฺเฉ กถิเสฏฺฐํ,}{\csromangathabat{“ปสฺสตารญฺญเก ภิกฺขู,}{อชฺโฌคาเฬฺห ธุเต คุเณ.} \\ \csromangathabat{ปุน ปสฺสติ คิหี ราชา,}{อนาคามิผเล ฐิเต.} \\ \csromangathabat{อุโภปิ เต วิโลเกตฺวา,}{อุปฺปชฺชิ สํสโย มหา.} \\ \csromangathabat{พุชฺเฌยฺย เจ คิหี ธมฺเม,}{ธุตงฺคํ นิปฺผลํ สิยา.} \\ \csromangathabat{ปรวาทิวาทมถนํ,}{นิปุณํ ปิฏกตฺตเย.} \\ \csromangathabat{หนฺท ปุจฺเฉ กถิเสฏฺฐํ,}{โส เม กงฺขํ วิเนสฺสตี”ติ.}}
\csromangathasetleft{“ปสฺสตารญฺญเก ภิกฺขู, \\ ปุน ปสฺสติ คิหี ราชา, \\ อุโภปิ เต วิโลเกตฺวา, \\ พุชฺเฌยฺย เจ คิหี ธมฺเม, \\ ปรวาทิวาทมถนํ, \\ หนฺท ปุจฺเฉ กถิเสฏฺฐํ,}
\csromangathagroup{\csromangathastanza{\csromanfolio{330}\csromangathabat{“ปสฺสตารญฺญเก ภิกฺขู,}{อชฺโฌคาเฬฺห ธุเต คุเณ.} \\ \csromangathabat{ปุน ปสฺสติ คิหี ราชา,}{อนาคามิผเล ฐิเต.}}\csromangathastanza{\csromangathabat{อุโภปิ เต วิโลเกตฺวา,}{อุปฺปชฺชิ สํสโย มหา.} \\ \csromangathabat{พุชฺเฌยฺย เจ คิหี ธมฺเม,}{ธุตงฺคํ นิปฺผลํ สิยา.}}\csromangathastanza{\csromangathabat{ปรวาทิวาทมถนํ,}{นิปุณํ ปิฏกตฺตเย.} \\ \csromangathabat{หนฺท ปุจฺเฉ กถิเสฏฺฐํ,}{โส เม กงฺขํ วิเนสฺสตี”ติ.}}}

\prose{อถ โข มิลินฺโท ราชา เยนายสฺมา นาคเสโน เตนุปสงฺกมิ, อุปสงฺกมิตฺวา อยสฺมนฺตํ นาคเสนํ อภิวาเทตฺวา เอกมนฺตํ นิสีทิ, เอกมนฺตํ นิสินฺโน โข มิลินฺโท ราชา อายสฺมนฺตํ นาคเสนํ เอตทโวจ “ภนฺเต นาคเสน อตฺถิ โกจิ คิหี อคาริโก กามโภคี ปุตฺตทาร\-สมฺพาธ\-สยนํ อชฺฌาวสนฺโต กาสิกจนฺทนํ ปจฺจนุโภนฺโต มาลา\-คนฺธ\-วิเลปนํ ธารยนฺโต ชาตรูป\-รชตํ สาทิยนฺโต มณิมุตฺตา\-กญฺจน\-วิจิตฺต\-โมฬิพทฺโธ เยน สนฺตํ ปรมตฺถํ นิพฺพานํ สจฺฉิกตนฺ”ติ.}

\prose{น มหาราช เอกญฺเญว สตํ น เทฺว สตานิ น ตีณิ จตฺตาริ ปญฺจ สตานิ น สหสฺสํ น สตสหสฺสํ น โกฏิสตํ น โกฏิสหสฺสํ น โกฏิ\-สต\-สหสฺสํ, ติฏฺฐตุ มหาราช ทสนฺนํ วีสติยา สตสฺส สหสฺสสฺส อภิสมโย, กตเมน เต ปริยาเยน อนุโยคํ ทมฺมีติ.}

\prose{ตฺวเมเวตํ พฺรูหีติ. เตนหิ\csromansharedfootnote{cb56481398405b41}{เตน หิ เทฺว ปทานิ \csromandash{} ม.พ.ป.} เต มหาราช กถยิสฺสามิ สเตน วา สหสฺเสน วา สตสหสฺเสน วา โกฏิยา วา โกฏิสเตน วา โกฏิสหสฺเสน วา โกฏิ\-สต\-สหสฺเสน วา, ยา กาจิ นวงฺเค พุทฺธวจเน สลฺเลขิตา\-จาร\-ปฺปฏิปตฺติ\-ธุต\-วรงฺค\-คุณ\-นิสฺสิตา\csromansharedfootnote{4569ff34712045d2}{ธุตงฺคคุณธรนิสฺสิตา (ก)} กถา, ตา สพฺพา อิธ สโมสริสฺสนฺติ. ยถา มหาราช นินฺนุ\-นฺนต\-สมวิสม\-ถลา\-ถล\-เทส\-ภาเค อภิวุฏฺฐํ อุทกํ, สพฺพํ ตํ ตโต วินิคฬิตฺวา มโหทธิํ สาครํ สโมสรติ. เอวเมว โข มหาราช สมฺปาทเก สติ ยา กาจิ นวงฺเค พุทฺธวจเน สลฺเลขิตา\-จาร\-ปฺปฏิปตฺติ\-ธุตงฺคคุณ\-ธร\-นิสฺสิตา กถา, ตา สพฺพา อิธ สโมสริสฺสนฺติ.}

\prose{\csromanfolio{331}มยฺหมฺเปตฺถ มหาราช ปริพฺยตฺตตาย พุทฺธิยา การณ\-ปริทีปนํ สโมสริสฺสติ, เตเนโส อตฺโถ สุวิภตฺโต วิจิตฺโต ปริปุณฺโณ ปริสุทฺโธ สมานีโต ภวิสฺสติ. ยถา มหาราช กุสโล เลขาจริโย อนุสิฏฺโฐ เลขํ โอสาเรนฺโต อตฺตโน พฺยตฺตตาย พุทฺธิยา การณ\-ปริทีปเนน เลขํ ปริปูเรติ, เอวํ สา เลขา สมตฺตา ปริปุณฺณา อนูนิกา ภวิสฺสติ. เอวเมว มยฺหมฺเปตฺถ ปริพฺยตฺตตาย พุทฺธิยา การณ\-ปริทีปนํ สโมสริสฺสติ, เตเนโส อตฺโถ สุวิภตฺโต วิจิตฺโต ปริปุณฺโณ ปริสุทฺโธ สมานีโต ภวิสฺสติ.}

\prose{นคเร มหาราช สาวตฺถิยา ปญฺจโกฏิมตฺตา อริยสาวกา ภควโต อุปาสกอุปาสิกาโย สตฺต\-ปณฺณาส\-สหสฺสานิ ตีณิ จ สตสหสฺสานิ อนาคามิผเล ปติฏฺฐิตา, เต สพฺเพปิ คิหี เยว น ปพฺพชิตา. ปุน ตตฺเถว กณฺฑมฺพมูเล ยมก\-ปาฏิหาริเย วีสติ ปาณโกฏิโย อภิสมิํสุ, ปน จูฬราหุโลวาเท\csromansharedfootnote{3c4f91313656872d}{มหาราหุโลวาเท (สี, อิ)}, มหามงฺคล\-สุตฺตนฺเต, สมจิตฺต\-ปริยาเย, ปราภว\-สุตฺตนฺเต, ปุราเภท\-สุตฺตนฺเต, กลหวิวาท\-สุตฺตนฺเต, จูฬพฺยูหสุตฺต\-นฺเต, มหาพฺยูห\-สุตฺตนฺเต, ตุวฏก\-สุตฺตนฺเต, สาริปุตฺต\-สุตฺตนฺเต คณนปถม\-ตีตานํ เทวตานํ ธมฺมา\-ภิสมโย อโหสิ.}

\prose{นคเร ราชคเห ปญฺญาส\-สหสฺสานิ ตีณิ จ สตสหสฺสานิ อริยสาวกา ภควโต อุปาสกอุปาสิกาโย, ปุน ตตฺเถว ธนปาล\-หตฺถินาค\-ทมเน นวุติ ปาณโกฏิโย, ปารายน\-สมาคเม ปาสาณกเจติเย จุทฺทส ปาณโกฏิโย, ปุน อินฺทสาล\-คุหายํ อสีติ เทวตาโกฏิโย, ปุน พาราณสิยํ อิสิปตเน มิคทาเย ปฐเม ธมฺมเทสเน อฏฺฐารส พฺรหฺมโกฏิโย อปริมาณา จ เทวตาโย, ปุน ตาวติํส\-ภวเน ปณฺฑุ\-กมฺพล\-สิลายํ อภิธมฺม\-เทสนาย อสีติ เทวตาโกฏิโย, เทโวโรหเณ สงฺกสฺส\-นครทฺวาเร โลกวิวรณ\-ปาฏิหาริเย ปสนฺนานํ นรมรูนํ ติํส โกฏิโย อภิสมิํสุ.}

\prose{ปุน สกฺเกสุ กปิล\-วตฺถุสฺมิํ นิโคฺรธาราเม พุทฺธวํส\-เทสนาย มหาสมย\-สุตฺตนฺต\-เทสนาย จ คณนปถม\-ตีตานํ เทวตานํ ธมฺมา\-ภิสมโย}

\prose{\csromanfolio{332}อโหสิ. ปุน สุมนมาลาการ\-สมาคเม, ครห\-ทินฺน\-สมาคเม, อานนฺท\-เสฏฺฐิ\-สมาคเม, ชมฺพุกา\-ชีวก\-สมาคเม, มณฺฑุก\-เทวปุตฺต\-สมาคเม, มฏฺฐกุณฺฑลิเทวปุตฺต\-สมาคเม, สุลสา\-นครโสภินิ\-สมาคเม, สิริมา\-นครโสภินิ\-สมาคเม, เปสการ\-ธีตุ\-สมาคเม, จูฬสุภทฺทา\-สมาคเม, สาเกต\-พฺราหฺมณสฺส อาฬาหน\-ทสฺสน\-สมาคเม, สูนา\-ปร\-นฺตก\-สมาคเม, สกฺกปญฺห\-สมาคเม, ติโรกุฏฺฏ\-สมาคเม\csromansharedfootnote{8351025be59a17ea}{ติโรกุฑฺฑสมาคเม (สี, อิ)}, รตนสุตฺต\-สมาคเม ปจฺเจกํ จตุราสีติยา ปาณสหสฺสานํ ธมฺมา\-ภิสมโย อโหสิ, ยาวตา มหาราช ภควา โลเก อฏฺฐาสิ, ตาว ตีสุ มณฺฑเลสุ โสฬสสุ มหาชนปเทสุ ยตฺถ ยตฺถ ภควา วิหาสิ, ตตฺถ ตตฺถ เยภุยฺยน เทฺว ตโย จตฺตาโร ปญฺจ สตํ สหสฺสํ สตสหสฺสํ เทวา จ มนุสฺสา จ สนฺตํ ปรมตฺถํ นิพฺพานํ สจฺฉิกริํสุ. เย เต มหาราช เทวา คิหี เยว น เต ปพฺพชิตา, เอตานิ เจว มหาราช อญฺญานิ จ อเนกานิ เทวตา\-โกฏิ\-สต\-สหสฺสานิ คิหี อคาริกา กามโภคิโน สนฺตํ ปรมตฺถํ นิพฺพานํ สจฺฉิกริํสูติ.}

\prose{ยทิ ภนฺเต นาคเสน คิหี อคาริกา กามโภคิโน สนฺตํ ปรมตฺถํ นิพฺพานํ สจฺฉิกโรนฺติ, อถ อิมานิ ธุตงฺคานิ กิมตฺถํ สาเธนฺติ, เตน การเณน ธุตงฺคานิ อกิจฺจกรานิ โหนฺติ. ยทิ ภนฺเต นาคเสน วินา มนฺโตสเธหิ พฺยาธโย วูปสมนฺติ, กิํ วมน\-วิเรจนา\-ทินา สรีร\-ทุพฺพลกรเณน. ยทิ มุฏฺฐีหิ ปฏิสตฺตุ\-นิคฺคโห ภวติ, กิํ อสิสตฺติ\-สร\-ธนุ\-โกทณฺฑ\-ลคุฬ\-มุคฺคเรหิ. ยทิ คณฺฐิกุฏิลสุสิรกณฺฏลตาสาขา อาลมฺพิตฺวา รุกฺขม\-ภิรูหนํ ภวติ, กิํ ทีฆทฬฺห\-นิสฺเสณิ\-ปริเยส\-เนน, ยทิ ถณฺฑิล\-เสยฺยาย ธาตุสมตา ภวติ, กิํ สุขสมฺผสฺส\-มหติ\-มหา\-สิริสยน\-ปริเยส\-เนน. ยทิ เอกโก สาสงฺก\-สปฺปฏิภย\-วิสม\-กนฺตาร\-ตรณ\-สมตฺโถ ภวติ, กิํ สนฺนทฺธ\-สชฺช\-มหติ\-มหา\-สตฺถ\-ปริเยส\-เนน. ยทิ นทิสรํ พาหุนา ตริตุํ สมตฺโถ ภวติ, กิํ ธุวเสตุ\-นาวา\-ปริเยส\-เนน. ยทิ สกสนฺตเกน ฆาสจฺฉาทนํ กาตุํ ปโหติ, กิํ ปรูปเสวน\-ปิย\-สมุลฺลาป\-ปจฺฉาปุเร\-ธาวเนน. ยทิ อขาตตฬาเก อุทกํ ลภติ, กิํ อุทปาน\-ตฬาก\-โปกฺขรณิ\-ขณเนน. เอวเมว โข ภนฺเต นาคเสน ยทิ คิหี อคาริกา กามโภคิโน สนฺตํ ปรมตฺถํ นิพฺพานํ สจฺฉิกโรนฺติ, กิํ ธุตคุณ\-วร\-สมาทิย\-เนนาติ.}

\prose{\csromanfolio{333}อฏฺฐวีสติ โข ปนิเม มหาราช ธุตงฺคคุณา ยถาภุจฺจ\-คุณา, เยหิ คุเณหิ ธุตงฺคานิ สพฺพพุทฺธานํ ปิหยิตานิ ปตฺถิตานิ. กตเม อฏฺฐวีสติ, อิธ มหาราช ธุตงฺคํ สุทฺธาชีวํ สุขผลํ อนวชฺชํ น ปรทุกฺขาปนํ อภยํ อสมฺปีฬนํ เอกนฺต\-วฑฺฒิกํ อปริหานิยํ อมายํ อารกฺขา ปตฺถิตททํ สพฺพสตฺต\-ทมนํ สํวรหิตํ ปติรูปํ อนิสฺสิตํ วิปฺปมุตฺตํ ราคกฺขยํ โทสกฺขยํ โมหกฺขยํ มานปฺปหานํ กุวิตกฺก\-จฺเฉทนํ กงฺขา\-วิตรณํ โกสชฺช\-วิทฺธํสนํ อรติปฺปหานํ ขมนํ อตุลํ อปฺปมาณํ สพฺพทุกฺขกฺขย\-คมนํ, อิเม โข มหาราช อฏฺฐวีสติ ธุตงฺคคุณา ยถาภุจฺจ\-คุณา เยหิ คุเณหิ ธุตงฺคานิ สพฺพพุทฺธานํ ปิหยิตานิ ปตฺถิตานิ.}

\prose{เย โข เต มหาราช ธุตคุเณ สมฺมา อุปเสวนฺติ, เต อฏฺฐารสหิ คุเณหิ สมุเปตา ภวนฺติ. กตเมหิ อฏฺฐารสหิ, อาจาโร เตสํ สุวิสุทฺโธ โภติ, ปฏิปทา สุปูริตา โหติ, กายิกํ วาจสิกํ สุรกฺขิตํ โหติ, มโนสมาจาโร สุวิสุทฺโธ โหติ, วีริยํ สุปคฺคหิตํ โหติ, ภยํ วูปสมฺมติ, อตฺตานุทิฏฺฐิพฺยปคตา โหติ, อาฆาโต อุปรโต โหติ, เมตฺตา อุปฏฺฐิตา โหติ, อาหาโร ปริญฺญาโต โหติ, สพฺพสตฺตานํ ครุกโต โหติ, โภชเน มตฺตญฺญู โหติ, ชาคริยม\-นุยุตฺโต โหติ, อนิเกโต โหติ, ยตฺถ ผาสุ ตตฺถ วิหารี โหติ, ปาปเชคุจฺฉี โหติ, วิเวการาโม โหติ, สตตํ อปฺปมตฺโต โหติ, เย เต มหาราช ธุตคุเณ สมฺมา อุปเสวนฺติ, เต อิเมหิ อฏฺฐารสหิ คุเณหิ สมุเปตา ภวนฺติ.}

\prose{ทส อิเม มหาราช ปุคฺคลา ธุตคุณารหา. กตเม ทส, สทฺโธ โหติ หิริมา ธิติมา อกุโห อตฺถวสี อโลโล สิกฺขากาโม ทฬฺหสมาทาโน อนุชฺฌาน\-พหุโล เมตฺตาวิหารี, อิเม โข มหาราช ทส ปุคฺคลา ธุตคุณารหา.}

\prose{เย เต มหาราช คิหี อคาริกา กามโภคิโน สนฺตํ ปรมตฺถํ นิพฺพานํ สจฺฉิกโรนฺติ, สพฺเพ เต ปุริมาสุ ชาตีสุ เตรสสุ ธุตคุเณสุ กตูปาสนา กตภูมิกมฺมา, เต ตตฺถ จารญฺจ ปฏิปตฺติญฺจ โสธยิตฺวา อชฺเชตรหิ คิหี เยว สนฺตา สนฺตํ ปรมตฺถํ นิพฺพานํ สจฺฉิกโรนฺติ.}

\prose{\csromanfolio{334}ยถา มหาราช กุสโล อิสฺสาโส อนฺเตวาสิเก ปฐมํ ตาว อุปาสนสาลายํ จาปเภทจาปาโรปนคฺคหณมุฏฺฐิปฺปฏิปีฬน- องฺคุลิ\-วินามน\-ปาทฐปน\-สรคฺคหณ\-สนฺนหน- อากฑฺฒนสนฺธารณลกฺขนิยมนขิปเน ติณปุริสก\-ฉกณ\csromansharedfootnote{fc9c5bb312c5488b}{ฉณก (สี, อิ)} ติณปลาล\-มตฺติกาปุญฺช\-ผลก\-ลกฺข\-เวเธ อนุสิกฺขาเปตฺวา รญฺโญ สนฺติเก อุปาสนํ อาราธยิตฺวา อาชญฺญรถคชตุรงฺคธนธญฺญ- หิรญฺญสุวณฺณทาสิทาสภริยคามวรํ ลภติ. เอวเมว โข มหาราช เย เต คิหี อคาริกา กามโภคิโน สนฺตํ ปรมตฺถํ นิพฺพานํ สจฺฉิกโรนฺติ, เต สพฺเพ ปุริมาสุ ชาตีสุ เตรสสุ ธุตคุเณสุ กตูปาสนา กตภูมิกมฺมา, เต ตตฺเถว จารญฺจ ปฏิปตฺติญฺจ โสธยิตฺวา อชฺเชตรหิ คิหี เยว สนฺตา สนฺตํ ปรมตฺถํ นิพฺพานํ สจฺฉิกโรนฺติ. น มหาราช ธุตคุเณสุ ปุพฺพาเสวนํ วินา เอกิสฺสา เยว ชาติยา อรหตฺตํ สจฺฉิกิริยา โหติ, อุตฺตเมน ปน วีริเยน อุตฺตมาย ปฏิปตฺติยา ตถารูเปน อาจริเยน กลฺยาณมิตฺเตน อรหตฺตํ สจฺฉิกิริยา โหติ.}

\prose{ยถา วา ปน มหาราช ภิสกฺโก สลฺลกตฺโต อาจริยํ ธเนน วา วตฺต\-ปฺปฏิปตฺติยา วา อาราเธตฺวา สตฺถคฺคหณเฉทนเลขนเวธนสลฺลุทฺธรณวณ โธวน โสสนเภสชฺชา นุลิมฺปนวมนวิเรจนานุวาสนกิริยมนุสิกฺขิตฺวา วิชฺชาสุ กตสิกฺโข กตูปาสโน กตหตฺโถ อาตุเร อุปสงฺกมติ ติกิจฺฉาย. เอวเมว โข มหาราช เย เต คิหี อคาริกา กามโภคิโน สนฺตํ ปรมตฺถํ นิพฺพานํ สจฺฉิกโรนฺติ, เต สพฺเพ ปุริมาสุ ชาตีสุ เตรสสุ ธุตคุเณสุ กตูปาสนา กตภูมิกมฺมา, เต ตตฺเถว จารญฺจ ปฏิปตฺติญฺจ โสธยิตฺวา อชฺเชตรหิ คิหี เยว สนฺตา สนฺตํ ปรมตฺถํ นิพฺพานํ สจฺฉิกโรนฺติ, น มหาราช ธุตคุเณหิ อวิสุทฺธานํ ธมฺมา\-ภิสมโย โหติ.}

\prose{ยถา มหาราช อุทกสฺส อเสจเนน พีชานํ อวิรูหนํ โหติ, เอวเมว โข มหาราช ธุตคุเณหิ อวิสุทฺธานํ ธมฺมา\-ภิสมโย น โหติ.}

\prose{ยถา วา ปน มหาราช อกตกุสลานํ อกต\-กลฺยาณานํ สุคติคมนํ น โหติ, เอวเมว โข มหาราช ธุตคุเณหิ อวิสุทฺธานํ ธมฺมา\-ภิสมโย น โหติ.}

\prose{\csromanfolio{335}ปถวิสมํ มหาราช ธุตคุณํ วิสุทฺธิ\-กามานํ ปติฏฺฐาน\-ฏฺเฐน.}

\prose{อาโปสมํ มหาราช ธุตคุณํ วิสุทฺธิ\-กามานํ สพฺพ\-กิเลสมล\-โธว\-นฏฺเฐน.}

\prose{เตโชสมํ มหาราช ธุตคุณํ วิสุทฺธิ\-กามานํ สพฺพกิเลส\-วน\-ชฺฌาปน\-ฏฺเฐน.}

\prose{วาโยสมํ มหาราช ธุตคุณํ วิสุทฺธิ\-กามานํ สพฺพ\-กิเลสมล\-รโช\-ปวาห\-นฏฺเฐน.}

\prose{อคทสมํ มหาราช ธุตคุณํ วิสุทฺธิ\-กามานํ สพฺพ\-กิเลสพฺยาธิ\-วูปสม\-นฏฺเฐน.}

\prose{อมตสมํ มหาราช ธุตคุณํ วิสุทฺธิ\-กามานํ สพฺพกิเลสวิสนาสนฏฺฏฺเนน.}

\prose{เขตฺตสมํ มหาราช ธุตคุณํ วิสุทฺธิ\-กามานํ สพฺพ\-สามญฺญ\-คุณ\-สสฺส\-วิรูห\-นฏฺเฐน.}

\prose{มโนหรสมํ มหาราช ธุตคุณํ วิสุทฺธิ\-กามานํ ปตฺถิติ\-จฺฉิต\-สพฺพสมฺปตฺติ\-วรท\-ทฏฺเฐน.}

\prose{นาวาสมํ มหาราช ธุตคุณํ วิสุทฺธิ\-กามานํ สํสาร\-มหณฺณว\-ปารคม\-นฏฺเฐน.}

\prose{ภีรุตฺตาณสมํ มหาราช ธุตคุณํ วิสุทฺธิ\-กามานํ ชรามรณ\-ภีตานํ อสฺสาส\-กรณ\-ฏฺเฐน.}

\prose{มาตุสมํ มหาราช ธุตคุณํ วิสุทฺธิ\-กามานํ กิเลสทุกฺข\-ปฺปฏิปีฬิตานํ อนุคฺคาหก\-ฏฺเฐน.}

\prose{ปิตุสมํ มหาราช ธุตคุณํ วิสุทฺธิ\-กามานํ กุสล\-วฑฺฒิ\-กามานํ สพฺพ\-สามญฺญ\-คุณ\-ชน\-กฏฺเฐน.}

\prose{มิตฺตสมํ มหาราช ธุตคุณํ วิสุทฺธิ\-กามานํ สพฺพ\-สามญฺญ\-คุณ\-ปริเยสน\-อวิสํวาท\-กฏฺเฐน.}

\prose{ปทุมสมํ มหาราช ธุตคุณํ วิสุทฺธิ\-กามานํ สพฺพ\-กิเลสมเลหิ อนุปลิตฺต\-ฏฺเฐน.}

\prose{\csromanfolio{336}จตุชฺชา\-ติย\-วร\-คนฺธสมํ มหาราช ธุตคุณํ วิสุทฺธิ\-กามานํ กิเลส\-ทุคฺคนฺธ\-ปฏิวิโนทน\-ฏฺเฐน.}

\prose{คิริราชวรสมํ มหาราช ธุตคุณํ วิสุทฺธิ\-กามานํ อฏฺฐ\-โลกธมฺม\-วาเตหิ อกมฺปิยฏฺเฐน.}

\prose{อากาสสมํ มหาราช ธุตคุณํ วิสุทฺธิ\-กามานํ สพฺพตฺถ คหณาปคตอุรุวิสฏวิตฺถตมหนฺตฏฺเฐน.}

\prose{นทีสมํ มหาราช ธุตคุณํ วิสุทฺธิ\-กามานํ กิเลสมล\-ปวาห\-นฏฺเฐน.}

\prose{สุเทสกสมํ มหาราช ธุตคุณํ วิสุทฺธิ\-กามานํ ชาติ\-กนฺตาร\-กิเลส\-วนคหน\-นิตฺถรณ\-ฏฺเฐน.}

\prose{มหา\-สตฺถวาหสมํ มหาราช ธุตคุณํ วิสุทฺธิ\-กามานํ สพฺพ\-ภย\-สุญฺญ\-เขม\-อภย\-วร\-ปวร\-นิพฺพาน\-นคร\-สมฺปาปน\-ฏฺเฐน.}

\prose{สุมชฺชิต\-วิมลา\-ทาสสมํ มหาราช ธุตคุณํ วิสุทฺธิ\-กามานํ สงฺขารานํ สภาว\-ทสฺสนฏฺเฐน.}

\prose{ผลกสมํ มหาราช ธุตคุณํ วิสุทฺธิ\-กามานํ กิเลส\-ลคุฬ\-สร\-สตฺติ\-ปฏิพาห\-นฏฺเฐน.}

\prose{ฉตฺตสมํ มหาราช ธุตคุณํ วิสุทฺธิ\-กามานํ กิเลส\-วสฺสติ\-วิธ\-คฺคิสนฺตาปา\-ตป\-ปฏิพาห\-นฏฺเฐน.}

\prose{จนฺทสมํ มหาราช ธุตคุณํ วิสุทฺธิ\-กามานํ ปิหยิ\-ตป\-ตฺถิต\-ฏฺเฐน.}

\prose{สูริยสมํ มหาราช ธุตคุณํ วิสุทฺธิ\-กามานํ โมหตม\-ติมิร\-นาส\-นฏฺเฐน.}

\prose{สาครสมํ มหาราช ธุตคุณํ วิสุทฺธิ\-กามานํ อเนกวิธ\-สามญฺญ\-คุณ\-วร\-รตนุ\-ฏฺฐา\-นฏฺเฐน, อปริมิต\-อสงฺเขฺยยฺย- อปฺปเมยฺย\-ฏฺเฐน จ.}

\prose{เอวํ โข มหาราช ธุตคุณํ วิสุทฺธิ\-กามานํ พหูปการํ สพฺพ\-ทรถ\-ปริฬาหนุทํ อรตินุทํ ภยนุทํ ภวนุทํ ขีลนุทํ มลนุทํ โสกนุทํ ทุกฺขนุทํ ราคนุทํ โทสนุทํ โมหนุทํ มานนุทํ ทิฏฺฐินุทํ สพฺพา\-กุสลธมฺมนุทํ ยสาวหํ หิตาวหํ สุขาวหํ ผาสุกรํ ปีติกรํ โยคกฺเขมกรํ อนวชฺชํ อิฏฺฐ\-สุขวิปากํ คุณราสิคุณปุญฺช อปริมิต\-อสงฺเขฺยยฺย อปฺปเมยฺยคุณํ วรํ ปวรํ อคฺคํ.}

\prose{\csromanfolio{337}ยถา มหาราช มนุสฺสา อุปตฺถมฺภ\-วเสน โภชนํ อุปเสวนฺติ, หิตวเสน เภสชฺชํ อุปเสวนฺติ, อุปการวเสน มิตฺตํ อุปเสวนฺติ, ตารณวเสน นาวํ อุปเสวนฺติ, สุคนฺธวเสน มาลาคนฺธํ อุปเสวนฺติ, อภยวเสน ภีรุตฺตาณํ อุปเสวนฺติ, ปติฏฺฐาวเสน\csromansharedfootnote{c3b07499269f0bd9}{ปติฏฺฐานวเสน (ก)} ปถวิํ อุปเสวนฺติ, สิปฺปวเสน อาจริยํ อุปเสวนฺติ, ยสวเสน ราชานํ อุปเสวนฺติ, กามททวเสน มณิรตนํ อุปเสวนฺติ. เอวเมว โข มหาราช สพฺพ\-สามญฺญ\-คุณ\-ทท\-วเสน อริยา ธุตคุณํ อุปเสวนฺติ.}

\prose{ยถา วา ปน มหาราช อุทกํ พีชวิรูหนาย, อคฺคิ ฌาปนาย, อาหาโร พลาหรณาย, ลตา พนฺธนาย, สตฺถํ เฉทนาย, ปานียํ ปิปาสา\-วินยนาย, นิธิ อสฺสาสกรณาย, นาวา ตีรสมฺปาปนาย, เภสชฺชํ พฺยาธิ\-วูปสมนาย, ยานํ สุขคมนาย, ภีรุตฺตาณํ ภยวิโนทนาย, ราชา อารกฺขตฺถาย, ผลกํ ทณฺฑ\-เลฑฺฑุ\-ลคุฬ\-สร\-สตฺติ\-ปฏิพาหนาย, อาจริโย อนุสาสนาย, มาตา โปสนาย, อาทาโส โอโลกนาย, อลงฺกาโร โสภนาย, วตฺถํ ปฏิจฺฉาทนาย, นิสฺเสณี อาโรหนาย, ตุลา วิสม\-วิกฺเขป\-นาย\csromansharedfootnote{d600b0c47bf7dc16}{นิกฺเขปนาย (สี, อิ)}, มนฺตํ ปริชปฺปนาย, อาวุธํ ตชฺชนีย\-ปฏิพาหนาย, ปทีโป อนฺธการ\-วิธมนาย, วาโต ปริฬาห\-นิพฺพาปนาย, สิปฺปํ วุตฺติ\-นิปฺผาทนาย, อคทํ ชีวิต\-รกฺขณาย, อากโร รตนุ\-ปฺปาทนาย, รตนํ อลงฺกราย, อาณา อนติกฺกมนาย, อิสฺสริยํ วสวตฺตนาย. เอวเมว โข มหาราช ธุตคุณํ สามญฺญ\-พีช\-วิรูหนาย, กิเลสมล\-ฌาปนาย, อิทฺธิพลา\-หรณาย, สติสํวร\-นิพนฺธนาย, วิมติ\-วิจิกิจฺฉา\-สมุจฺเฉทนาย, ตณฺหา\-ปิปาสา\-วินยนาย, อภิสมย- อสฺสาสกรณาย, จตุโรฆนิตฺถรณาย, กิเลสพฺยาธิ\-วูปสมาย, นิพฺพานสุข\-ปฺปฏิลาภาย, ชาติชราพฺยาธิมรณโสกปริเทวทุกฺขโทมนสฺสุปายาสภยวิโน ทนาย, สามญฺญ\-คุณ\-ปริรกฺข\-ณาย, อรติ\-กุวิตกฺก\-ปฏิพาหนาย, สกล\-สามญฺญตฺถา\-นุสาสนาย, สพฺพ\-สามญฺญ\-คุณ\-โปสนาย, สมถ\-วิปสฺสนามคฺค\-ผลนิพฺพาน\-ทสฺสนาย, สกล\-โลก\-ถุต\-โถมิต\-มหติ\-มหา\-โสภน\-กรณาย, สพฺพา\-ปาย\-ปิทหนาย, สามญฺญตฺถ\-เสล\-สิขร\-มุทฺธนิ อภิรูหนาย, วงฺก\-กุฏิล\-วิสม\-จิตฺตวิกฺเขป\-นาย\csromansharedfootnote{6ec08c0faaa86622}{จิตฺตนิกฺเขปนาย (สี, อิ)}, เสวิตพฺพา\-เสวิตพฺพธมฺเม สาธุ\-สชฺฌาย\-กรณาย,}

\prose{\csromanfolio{338}สพฺพกิเลส\-ปฏิสตฺตุ\-ตชฺชนาย, อวิชฺชนฺธการ\-วิธมนาย, ติวิธคฺคิ\-สนฺตาป\-ปริฬาห\-นิพฺพาปนาย, สณฺหสุขุม\-สนฺต\-สมาปตฺติ\-นิปฺผาทนาย, สกล\-สามญฺญ\-คุณ\-ปริรกฺข\-ณาย, โพชฺฌงฺค\-วร\-รตนุ\-ปฺปาทนาย, โยคิ\-ชนา\-ลงฺกรณาย, อนวชฺช\-นิปุณ\-สุขุม\-สนฺติสุ\-ขม\-นติกฺกมนาย, สกล\-สามญฺญ\-อริยธมฺม\-วสวตฺตนาย. อิติ มหาราช อิเมสํ คุณานํ อธิคมาย ยทิทํ เอกเมกํ ธุตคุณํ, เอวํ มหาราช อตุลิยํ ธุตคุณํ อปฺปเมยฺยํ อสมํ อปฺปฏิสมํ อปฺปฏิภาคํ อปฺปฏิเสฏฺฐํ อุตฺตรํ เสฏฺฐํ วิสิฏฺฐํ อธิกํ อายตํ ปุถุลํ วิสฏํ วิตฺถตํ ครุกํ ภาริยํ มหนฺตํ.}

\prose{โย โข มหาราช ปุคฺคโล ปาปิจฺโฉ อิจฺฉาปกโต กุหโก ลุทฺโธ โอทริโก ลาภกาโม ยสกาโม กิตฺติกาโม อยุตฺโต อปฺปตฺโต อนนุจฺฉวิโก อนรโห อปฺปติรูโป ธุตงฺคํ\csromansharedfootnote{a6a370c82a0116eb}{ตธุคุณํ (ก) เอวมุปริปิ.} สมาทิยติ, โส ทิคุณํ ทณฺฑมาปชฺชติ, สพฺพคุณ\-ฆาตมา\-ปชฺชติ, ทิฏฺฐ\-ธมฺมิกํ หีฬนํ ขีฬนํ ครหนํ อุปณฺฑนํ ขิปนํ อสมฺโภคํ นิสฺสารณํ นิจฺฉุภนํ ปวาหนํ ปพฺพาชนํ ปฏิลภติ, สมฺปราเยปิ สตโยชนิเก อวีจิ\-มหานิรเย อุณฺหกฏฺฐิตตตฺตสนฺตตฺตอจฺจิชาลามาลเก อเนก\-วสฺส\-โกฏิสตสหสฺสานิ อุทฺธมโธ ติริยํ เผณุทฺเทหกํ สมฺปริวตฺตกํ ปจฺจติ, ตโต มุจฺจิตฺวา\csromansharedfootnote{977ba4560d3c3007}{มุญฺจิตฺวา (ก)} กิสผรุสกาฬงฺคปจฺจงฺโค สูนุทฺธุมาตสุสิรุตฺตมงฺโค\csromansharedfootnote{b0380d61b367ef45}{สูนุทฺธุมาตสูจิมุขปมาณสุสิรุตฺตมงฺโค (สี, อิ)} ฉาโต ปิปาสิโต วิสม\-ภีมรูป\-วณฺโณ ภคฺค\-กณฺณ\-โสโต อุมฺมีลิตนิมีลิตเนตฺตนยโน อรุคตฺต\-ปกฺกคตฺโต ปุลวา\-กิณฺณ\-สพฺพกาโย วาตมุเข ชลมาโน วิย อคฺคิกฺขนฺโธ อนฺโต ชลมาโน ปชฺชลมาโน อตาโณ อสรโณ อารุณฺณรุณฺณการุญฺญรวํ ปริเทวมาโน นิชฺฌาม\-ตณฺหิโก สมณมหาเปโต หุตฺวา อาหิณฺฑมาโน มหิยา อฏฺฏสฺสรํ กโรติ.}

\prose{ยถา มหาราช โกจิ อยุตฺโต อปฺปตฺโต อนนุจฺฉวิโก อนรโห อปฺปติรูโป หีโน กุชาติโก ขตฺติยา\-ภิเสเกน อภิสิญฺจติ, โส ลภติ หตฺถจฺเฉทํ ปาทจฺเฉทํ หตฺถปาท\-จฺเฉทํ กณฺณจฺเฉทํ นาสจฺเฉทํ กณฺณนาส\-จฺเฉทํ พิลงฺค\-ถาลิกํ สงฺขมุณฺฑิกํ ราหุมุขํ โชติมาลิกํ หตฺถ\-ปชฺโชติกํ เอรกวตฺติกํ จีรกวาสิกํ เอเณยฺยกํ พฬิสมํสิกํ กหาปณกํ ขารา\-ปตจฺฉิกํ ปลิฆ\-ปริวตฺติกํ ปลาลปีฐกํ ตตฺเตน \csromanfolio{339}เกเลน โอสิญฺจนํ สุนเขหิ ขาทาปนํ ชีวสูลา\-โรปนํ อสินา สีสจฺเฉทํ อเนกวิหิตมฺปิ กมฺมการณํ อนุภวติ, กิํ การณา อยุตฺโต อปฺปตฺโต อนนุจฺฉวิโก อนรโห อปฺปติรูโป หีโน กุชาติโก มหนฺเต อิสฺสริเย ฐาเน อตฺตานํ ฐเปสิมฺ, เวลํ ฆาเตสิ, เอวเมว โข มหาราช โย โกจิ ปุคฺคโล ปาปิจฺโฉ ฯเปฯ มหิยา อฏฺฏสฺสรํ กโรติ.}

\proseitem{5}{โย ปน มหาราช ปุคฺคโล ยุตฺโต ปตฺโต อนุจฺฉวิโก อรโห ปติรูโป อปฺปิจฺโฉ สนฺตุฏฺโฐ ปวิวิตฺโต อสํสฏฺโฐ อารทฺธวีริโย ปหิตตฺโต อสโฐ อมาโย อโนทริโก อลาภกาโม อยสกาโม อกิตฺติกาโม สทฺโธ สทฺธา\-ปพฺพชิโต ชรามรณา มุจฺจิตุกาโม “สาสนํ ปคฺคณฺหิสฺสามี”ติ ธุตงฺคํ สมาทิยติ, โส ทิคุณํ ปูชํ อรหติ เทวานญฺจ ปิโย โหติ มนาโป ปิหยิโต ปตฺถิโต, ชาติสุมน\-มลฺลิกาทีนํ วิย ปุปฺผํ นหาตา\-นุลิตฺตสฺส, ชิฆจฺฉิตสฺส วิย ปณีตโภชนํ, ปิปาสิตสฺส วิย สีตล\-วิมล\-สุรภิ\-ปานียํ, วิสคตสฺส วิย โอสธวรํ, สีฆคมน\-กามสฺส วิย อาชญฺญรถ\-วรุตฺตมํ, อตฺถกามสฺส วิย มโนหร\-มณิรตนํ, อภิสิญฺจิตุกามสฺส วิย ปณฺฑร\-วิมล\-เสตจฺฉตฺตํ, ธมฺมกามสฺส วิย อรหตฺตผลา\-ธิคมม\-นุตฺตรํ. ตสฺส จตฺตาโร สติปฏฺฐานา ภาวนา\-ปาริปูริํ คจฺฉนฺติ, จตฺตาโร สมฺมปฺปธานา จตฺตาโร อิทฺธิปาทา ปญฺจินฺทฺริยานิ ปญฺจ พลานิ สตฺต โพชฺฌงฺคา อริโย อฏฺฐงฺคิโก มคฺโค ภาวนา\-ปาริปูริํ คจฺฉติ, สมถวิปสฺสานา อธิคจฺฉติ, อธิคม\-ปฺปฏิปตฺติ ปริณมติ, จตฺตาริ สามญฺญผลานิ จตสฺโส ปฏิสมฺภิทา ติสฺโส วิชฺชา ฉฬภิญฺญา เกวโล จ สมณธมฺโม สพฺเพ ตสฺสาเธยฺยา โหนฺติ, วิมุตฺติ\-ปณฺฑร\-วิมล\-เสตจฺฉตฺเตน อภิสิญฺจติ.}

\prose{ยถา มหาราช รญฺโญ ขตฺติยสฺส อภิชาต\-กุล\-กุลีนสฺส ขตฺติยา\-ภิเสเกน อภิสิตฺตสฺส ปริจรนฺติ สรฏฺฐ\-เนคมชานปท\-ภฏ\-พลา\csromansharedfootnote{23f6a712e3720d8f}{พลตฺถา (สี, อิ)} อฏฺฐตฺติํสา จ ราชปริสา นฏนจฺจกา มุขมงฺคลิกา โสตฺถิวาจกา สมณพฺราหฺมณ\-สพฺพ\-ปาสณฺฑ\-คณา อภิคจฺฉนฺติ, ยํ กิญฺจิ ปถวิยา ปฏฺฏน\-รตนากร\-นคร\-สุงฺกฏฺฐาน\-เวรชฺชก\-เฉชฺชเภชฺช\-ชนม\-นุสาสนํ สพฺพตฺถ สามิโก ภวติ, เอวเมว โข มหาราช โย โกจิ ปุคฺคโล ยุตฺโต ปตฺโต ฯเปฯ วิมุตฺติ\-ปณฺฑร\-วิมล\-เสตจฺฉตฺเตน อภิสิญฺจติ.}

\prose{\csromanfolio{340}เตรสิมานิ มหาราช ธุตงฺคานิ, เยหิ สุทฺธิกโต นิพฺพาน\-มหาสมุทฺทํ ปวิสิตฺวา พหุวิธํ ธมฺมกีฬมภิกีฬติ, รูปารูป\-อฏฺฐ\-สมาปตฺติโย วฬญฺเชติ, อิทฺธิวิธํ ทิพฺพโสตธาตุํ ปรจิตฺต\-วิชานนํ ปุพฺเพ\-นิวาสา\-นุสฺสติํ ทิพฺพจกฺขุํ สพฺพา\-สวกฺขยญฺจ ปาปุณาติ. กเตเม เตรส, ปํสุกูลิ\-กงฺคํ เตจีวริกงฺคํ ปิณฺฑปาติ\-กงฺคํ สปทาน\-จาริกงฺคํ เอกาสนิกงฺคํ ปตฺต\-ปิณฺฑิ\-กงฺคํ ขลุปจฺฉา\-ภตฺติ\-กงฺคํ อารญฺญิกงฺคํ รุกฺข\-มูลิ\-กงฺคํ อพฺโภกาสิกงฺคํ โสสานิกงฺคํ ยถา\-สนฺถติกงฺคํ เนสชฺชิกงฺคํ, อิเมหิ โข มหาราช เตรสหิ ธุตคุเณหิ ปุพฺเพ อาเสวิเตหิ นิเสวิเตหิ จิณฺเณหิ ปริจิณฺเณหิ จริเตหิ อุปจริเตหิ ปริปูริเตหิ เกวลํ สามญฺญํ ปฏิลภติ, ตสฺสาเธยฺยา โหนฺติ เกวลา สนฺตา สุขา สมาปตฺติโย.}

\prose{ยถา มหาราช สธโน นาวิโก ปฏฺฏเน สุฏฺฐุ กตสุงฺโก มหาสมุทฺทํ ปวิสิตฺวา วงฺคํ ตกฺโกลํ จีนํ โสวีรํ สุรฏฺฐํ อลสนฺทํ โกลปฏฺฏนํ สุวณฺณภูมิํ คจฺฉติ อญฺญมฺปิ ยํ กิญฺจิ นาวาสญฺจรณํ. เอวเมว โข มหาราช อิเมหิ เตรสหิ ธุตคุเณหิ ปุพฺเพ อาเสวิเตหิ นิเสวิเตหิ จิณฺเณหิ ปริจิณฺเณหิ จริเตหิ อุปจริเตหิ ปริปูริเตหิ เกวลํ สามญฺญํ ปฏิลภติ, ตสฺสาเธยฺยา โหนฺติ เกวลา สนฺตา สุขา สมาปตฺติโย.}

\prose{ยถา มหาราช กสฺสโก ปฐมํ เขตฺตโทสํ ติณกฏฺฐ\-ปาสาณํ อปเนตฺวา กสิตฺวา วปิตฺวา สมฺมา อุทกํ ปเวเสตฺวา รกฺขิตฺวา โคเปตฺวา ลวน\-มทฺทเนน พหุธญฺญโก โหติ, ตสฺสาเธยฺยา ภวนฺติ เย เกจิ อธนา กปณา ทลิทฺทา ทุคฺคตชนา. เอวเมว โข มหาราช อิเมหิ เตรสหิ ธุตคุเณหิ ปุพฺเพ อาเสวิเตหิ ฯเปฯ เกวลา สนฺตา สุขา สมาปตฺติโย.}

\prose{ยถา วา ปน มหาราช ขตฺติโย มุทฺธาวสิตฺโต อภิชาตกุลีโน เฉชฺชเภชฺช\-ชนม\-นุสาสเน อิสฺสโร โหติ วสวตฺตี สามิโก อิจฺฉากรโณ, เกวลา จ มหาปถวี ตสฺสาเธยฺยา โหติ. เอวเมว โข มหาราช อิเมหิ เตรสหิ ธุตคุเณหิ ปุพฺเพ อาเสวิเตหิ นิเสวิเตหิ จิณฺเณหิ ปริจิณฺเณหิ จริเตหิ อุปจริเตหิ ปริปูริเตหิ ชินสาสนวเร อิสฺสโร โหติ วสวตฺตี สามิโก อิจฺฉากรโณ, เกวลา จ สมณคุณา ตสฺสาเธยฺยา โหนฺติ.}

\prose{\csromanfolio{341}นนุ มหาราช\csromansymbolfootnote{*}{วิ 1. 336 ปิฏฺเฐ.}เถรา อุปเสโน วงฺคนฺตปุตฺโต สลฺเลข\-ธุตคุเณ ปริปูร\-การิ\-ตาย อนาทิยิตฺวา สาวตฺถิยา สํฆสฺส กติกํ สปริโส นรทมฺม\-สารถิํ ปฏิสลฺลานคตํ อุปสงฺกมิตฺวา ภควโต ปาเท สิรสาวนฺทิตฺวา เอกมนฺตํ นิสีทิ, ภควา จ ตํ สุวินีตํ ปริสํ โอโลเกตฺวา หฏฺฐตุฏฺโฐ ปมุทิโต อุทคฺโค ปริสาย สทฺธิํ สลฺลาปํ สลฺลปิตฺวา อสมฺภินฺเนน พฺรหฺมสฺสเรน เอตทโวจ “ปาสาทิกา โข ปน ตฺยายํ อุปเสน ปริสา, กถํ ตฺวํ อุปเสน ปริสํ วิเนสี”ติ. โสปิ สพฺพญฺญุนา ทสพเลน เทวาติเทเวน ปุฏฺโฐ ยถาภูต\-สภาว\-คุณ\-วเสน ภควนฺตํ เอตทโวจ\csromandash{}}

\prose{โย โกจิ มํ ภนฺเต อุปสงฺกมิตฺวา ปพฺพชฺชํ วา นิสฺสยํ วา ยาจติ, ตมหํ เอวํ วทามิ “อหํ โข อาวุโส อารญฺญิโก ปิณฺฑปาติโก ปํสุกูลิโก เตจีวริโก. สเจ ตฺวมฺปิ อารญฺญิโก ภวิสฺสสิ ปิณฺฑปาติโก ปํสุกูลิโก เตจีวริโก, เอวาหํ ตํ ปพฺพาเชสฺสามิ นิสฺสยํ ทสฺสามี”ติ, สเจ โส เม ภนฺเต ปฏิสฺสุณิตฺวา นนฺทติ โอรมติ, เอวาหํ ตํ ปพฺพาเชมิ นิสฺสยํ เทมิ, สเจ น นนฺทติ น โอรมติ, น ตํ ปพฺพาเชมิ, น นิสฺสยํ เทมิ, เอวาหํ ภนฺเต ปริสํ วิเนมีติ. เอวํ โข\csromansharedfootnote{06cc5a50ceee2208}{เอวมฺปิ (ก)} มหาราช ธุตคุณ\-วร\-สมาทิณฺโณ ชินสาสนวเร อิสฺสโร โหติ. วสวตฺตี สามิโก อิจฺฉากรโณ, ตสฺสาเธยฺยา โหนฺติ เกวลา สนฺตา สุขา สมาปตฺติโย.}

\prose{ยถา มหาราช ปทุมํ อภิวุทฺธปริสุทฺธอุทิจฺจชาติปฺปภวํ สินิทฺธํ มุทุ โลภนียํ สุคนฺธํ ปิยํ ปตฺถิตํ ปสตฺถํ ชลกทฺทมม\-นุปลิตฺตํ อณุปตฺต\-เกสร\-กณฺณิกาภิ\-มณฺฑิตํ ภมรคณ\-เสวิตํ สีตล\-สลิล\-สํวทฺธํ, เอวเมว โข มหาราช อิเมหิ เตรสหิ ธุตคุเณหิ ปุพฺเพ อาเสวิเตหิ นิเสวิเตหิ จิณฺเณหิ ปริจิณฺเณหิ จริเตหิ อุปจริเตหิ ปริปูริเตหิ อริยสาวโก ติํสคุณ\-วเรหิ สมุเปโต โหติ.}

\prose{กตเมหิ ติํสคุณ\-วเรหิ, สินิทฺธ\-มุทุ\-มทฺทว\-เมตฺตจิตฺโต โหติ, ฆาติต\-หต\-วิหต\-กิเลโส โหติ, หตนิหตมาน\-ทพฺโพ โหติ, อจล\-ทฬฺห\-นิวิฏฺฐ\-นิพฺเพมติก\-สทฺโธ โหติ, ปริปุณฺณ\-ปีณิต\-ปหฏฺฐ\-โลภนีย\-สนฺต\-สุข\-สมาปตฺติ\-ลาภี โหติ, สีลวร\-ปวร\-อสม\-สุจิ\-คนฺธปริภาวิโต โหติ, เทวมนุสฺสานํ ปิโย โหติ มนาโป, ขีณาสว\-อริย\-วร\-ปุคฺคล\-ปตฺถิโต, เทวมนุสฺสานํ วนฺทิตปูชิโต, พุธวิพุธปณฺฑิตชานานํ}

\prose{\csromanfolio{342}ถุตถวิต\-โถมิต\-ปสตฺโถ, อิธ วา หุรํ วา โลเกน อนุปลิตฺโต, อปฺปโถกวชฺเชปิ\csromansharedfootnote{ad1241dc71886c43}{อปฺปสาวชฺเชปิ (ก)} ภยทสฺสาวี, วิปุล\-วร\-สมฺปตฺติ\-กามานํ มคฺคผล\-วร\-ตฺถ\-สาธโน, อายาจิต\-วิปุล\-ปณีต\-ปจฺจย\-ภาคี, อนิเกตสยโน, ฌาน\-ชฺโฌสิต\-ตปฺป\-วร\-วิหารี, วิชฏิต\-กิเลส\-ชาล\-วตฺถุ, ภินฺน\-ภคฺค\-สงฺกุฏิต\-สญฺฉินฺน\-คติ\-นีวรโณ, อกุปฺปธมฺโม, อภินีตวาโส, อนวชฺชโภคี, คติวิมุตฺโต, อุตฺติณฺณ\-สพฺพ\-วิจิกิจฺโฉ, วิมุตฺติชฺโฌสิตตฺโถ\csromansharedfootnote{e755ee51a3a74461}{วิมุตฺติชฺฌาสิตตฺโต (สี, อิ) 3. สมาจิณฺณกุสลมูโล (ก)}, ทิฏฺฐธมฺโม, อจล\-ทฬฺห\-ภีรุตฺตาณมุ\-ปคโต, สมุจฺฉินฺนา\-นุสโย, สพฺพา\-สวกฺขยํ ปตฺโต, สนฺต\-สุข\-สมาปตฺติ\-วิหาร\-พหุโล, สพฺพ\-สมณ\-คุณ\-สมุเปโต, อิเมหิ ติํสคุณ\-วเรหิ สมุเปโต โหติ.}

\prose{นนุ มหาราช เถโร สาริปุตฺโต ทสสหสฺสิ\-โลกธาตุยา อคฺคปุริโส ฐเปตฺวา ทสพลํ โลกาจริยํ, โสปิ อปริมิตม\-สงฺเขฺยยฺยกปฺเป สมา\-จิต\-กุสล\-มูโล\csromansharedfootnote{cf2cbc004bdd8150}{สมาจิณฺณกุสลมูโล (ก)} พฺราหฺมณกุล\-กุลีโน มนาปิกํ กามรติํ อเนกสต\-สงฺขํ ธนวรญฺจ โอหาย ชินสาสเน ปพฺพชิตฺวา อิเมหิ เตรสหิ ธุตคุเณหิ กายวจีจิตฺตํ ทมยิตฺวา อชฺเชตรหิ อนนฺตคุณ\-สมนฺนาคโต โคตมสฺส ภควโต สาสนวเร ธมฺมจกฺกม\-นุ\-ปฺปวตฺตโก ชาโต. ภาสิตมฺเปตํ มหาราช ภควตา เทวาติเทเวน เอกงฺคุ\-ตฺตร\-นิกาย\-วร\-ลญฺฉเก\csromandash{} * “นาหํ ภิกฺขเว อญฺญํ เอกปุคฺคลมฺปิ สมนุปสฺสามิ, โย เอวํ ตถาคเตน อนุตฺตรํ ธมฺมจกฺกํ ปวตฺติตํ สมฺมเทว อนุปฺปวตฺเตติ, ยถยิทํ ภิกฺขเว สาริปุตฺโต, สาริปุตฺโต ภิกฺขเว ตถาคเตน อนุตฺตรํ ธมฺมจกฺกํ ปวตฺติตํ สมฺมเทว อนุปฺปวตฺเตตี”ติ.}

\prose{สาธุ ภนฺเต นาคเสน, ยํ กิญฺจิ นวงฺคํ พุทฺธวจนํ, ยา จ โลกุตฺตรา กิริยา, ยา จ โลเก อธิคม\-วิปุล\-วร\-สมฺปตฺติโย, สพฺพํ ตํ เตรสสุ ธุตคุเณสุ สโมธาโน\-ปคต\csromansharedfootnote{29c234f30a963c88}{สโมธาเนตพฺพํ (ก) * อํ 1. 23 ปิฏฺเฐ.}นฺติ.}

\vspace{\tipitakaparskip}
\csromancenter{ธุตงฺคปโญฺห ทุติโย.}
\nitthitammidruled{\textbf{อนุมานวคฺโค จตุตฺโถ.}}
\csromanmatikaanchor{mtk.200}
\csromantocmarknopage{subsection}{5. โอปมฺมกถาปญฺห}{mtk.200}
\bookmark[dest={mtk.200},level=2]{5. โอปมฺมกถาปญฺห}
\csromanheader{2}{6. \textbf{โอปมฺม}\-\textbf{กถา}\-\textbf{ปญฺห}}
\csromanheader{2}{\textbf{มาติกา}}
\prose{\csromanfolio{343}ภนฺเต นาคเสน กติหงฺเคหิ สมนฺนาคโต ภิกฺขุ อรหตฺตํ สจฺฉิกโรตีติ.}

\prose{อิธ มหาราช อรหตฺตํ สจฺฉิ\-กาตุกาเมน ภิกฺขุนา\csromandash{}}

\prose{คทฺรภสฺส\csromansharedfootnote{b229418e8b5d176a}{โฆรสฺสรสฺส (สี, สฺยา, อิ)} เอกํ องฺคํ คเหตพฺพํ.}

\prose{กุกฺกุฏสฺส ปญฺจ องฺคานิ คเหตพฺพานิ.}

\prose{กลนฺทกสฺส เอกํ องฺคํ คเหตพฺพํ.}

\prose{ทีปินิยา เอกํ องฺคํ คเหตพฺพํ.}

\prose{ทีปิกสฺส เทฺว องฺคานิ คเหตพฺพานิ.}

\prose{กุมฺมสฺส ปญฺจ องฺคานิ คเหตพฺพานิ.}

\prose{วํสสฺส เอกํ องฺคํ คเหตพฺพํ.}

\prose{จาปสฺส เอกํ องฺคํ คเหตพฺพํ.}

\prose{วายสสฺส เทฺว องฺคานิ คเหตพฺพานิ.}

\prosewithclosermidruled{มกฺกฏสฺส เทฺว องฺคานิ คเหตพฺพานิ.}{คทฺรภวคฺโค ปฐโม.}
\prose{ลาพุลตาย เอกํ องฺคํ คเหตพฺพํ.}

\prose{ปทุมสฺส ตีณิ องฺคานิ คเหตพฺพานิ.}

\prose{พีชสฺส เทฺว องฺคานิ คเหตพฺพานิ.}

\prose{สาล\-กลฺยาณิกาย เอกํ องฺคํ คเหตพฺพํ.}

\prose{นาวาย ตีณิ องฺคานิ คเหตพฺพานิ.}

\prose{นาวา\-ลคฺคน\-กสฺส\csromansharedfootnote{0016d5d34016e36a}{นาวาลคนกสฺส (สี, อิ)} เทฺว องฺคานิ คเหตพฺพานิ.}

\prose{\csromanfolio{344}กูปสฺส\csromansharedfootnote{712cb164837e00c7}{กูปกสฺส (ก)} เอกํ องฺคํ คเหตพฺพํ.}

\prose{นิยามกสฺส ตีณิ องฺคานิ คเหตพฺพานิ.}

\prose{กมฺมการสฺส เอกํ องฺคํ คเหตพฺพํ.}

\prosewithclosermidruled{สมุทฺทสฺส ปญฺจ องฺคานิ คเหตพฺพานิ.}{สมุทฺทวคฺโค ทุติโย.}
\prose{ปถวิยา ปญฺจ องฺคานิ คเหตพฺพานิ.}

\prose{อาปสฺส ปญฺจ องฺคานิ คเหตพฺพานิ.}

\prose{เตชสฺส ปญฺจ องฺคานิ คเหตพฺพานิ.}

\prose{วายุสฺส ปญฺจ องฺคานิ คเหตพฺพานิ.}

\prose{ปพฺพตสฺส ปญฺจ องฺคานิ คเหตพฺพานิ.}

\prose{อากาสสฺส ปญฺจ องฺคานิ คเหตพฺพานิ.}

\prose{จนฺทสฺส ปญฺจ องฺคานิ คเหตพฺพานิ.}

\prose{สูริยสฺส สตฺต องฺคานิ คเหตพฺพานิ.}

\prose{สกฺกสฺส ตีณิ องฺคานิ คเหตพฺพานิ.}

\prosewithclosermidruled{จกฺกวตฺติสฺส จตฺตาริ องฺคานิ คเหตพฺพานิ.}{ปถวีวคฺโค ตติโย.}
\prose{อุปจิกาย เอกํ องฺคํ คเหตพฺพํ.}

\prose{พิฬารสฺส เทฺว องฺคานิ คเหตพฺพานิ.}

\prose{อุนฺทูรสฺส เอกํ องฺคํ คเหตพฺพํ.}

\prose{วิจฺฉิกสฺส เอกํ องฺคํ คเหตพฺพํ.}

\prose{นกุลสฺส เอกํ องฺคํ คเหตพฺพํ.}

\prose{ชรสิงฺคาลสฺส เทฺว องฺคานิ คเหตพฺพานิ.}

\prose{\csromanfolio{345}มิคสฺส ตีณิ องฺคานิ คเหตพฺพานิ.}

\prose{โครูปสฺส จตฺตาริ องฺคานิ คเหตพฺพานิ.}

\prose{วราหสฺส เทฺว องฺคานิ คเหตพฺพานิ.}

\prosewithclosermidruled{หตฺถิสฺส ปญฺจ องฺคานิ คเหตพฺพานิ.}{อุปจิกาวคฺโค จตุตฺโถ.}
\prose{สีหสฺส สตฺต องฺคานิ คเหตพฺพานิ.}

\prose{จกฺกวากสฺส ตีณิ องฺคานิ คเหตพฺพานิ.}

\prose{เปณาหิกาย เทฺว องฺคานิ คเหตพฺพานิ.}

\prose{ฆรกโปตสฺส เอกํ องฺคํ คเหตพฺพํ.}

\prose{อุลูกสฺส เทฺว องฺคานิ คเหตพฺพานิ.}

\prose{สตปตฺตสฺส เอกํ องฺคํ คเหตพฺพํ.}

\prose{วคฺคุลิสฺส เทฺว องฺคานิ คเหตพฺพานิ.}

\prose{ชลูกาย เอกํ องฺคํ คเหตพฺพํ.}

\prose{สปฺปสฺส ตีณิ องฺคานิ คเหตพฺพานิ.}

\prosewithclosermidruled{อชครสฺส เอกํ องฺคํ คเหตพฺพํ.}{สีลวคฺโค ปญฺจโม.}
\csromancenter{ปนฺถ\-มกฺกฏกสฺส เอกํ องฺคํ คเหตพฺพํ.}
\csromancenter{ถนสิ\-ตทา\-รกสฺส\csromansharedfootnote{c4a1cb410dc86000}{ถนปีตทารกสฺส (สฺยา)} เอกํ องฺคํ คเหตพฺพํ.}
\csromancenter{จิตฺตก\-ธร\-กุมฺมสฺส\csromansharedfootnote{d99349c5ae147073}{จิตฺตกถลกุมฺมสฺส (ก)} เอกํ องฺคํ คเหตพฺพํ.}
\prose{ปวนสฺส ปญฺจ องฺคานิ คเหตพฺพานิ.}

\prose{รูกฺขสฺส ตีณิ องฺคานิ คเหตพฺพานิ.}

\prose{เมฆสฺส ปญฺจ องฺคานิ คเหตพฺพานิ.}

\prose{มณิรตนสฺส ตีณิ องฺคานิ คเหตพฺพานิ.}

\prose{\csromanfolio{346}มาควิกสฺส จตฺตาริ องฺคานิ คเหตพฺพานิ.}

\prosewithclosermidruled{พาฬิสิกสฺส เทฺว องฺคานิ คเหตพฺพานิ.}{ตจฺฉกสฺส เทฺว องฺคานิ คเหตพฺพานิ. มกฺกฏวคฺโค ฉฏฺโฐ.}
\prose{กุมฺภสฺส เอกํ องฺคํ คเหตพฺพํ.}

\prose{กาฬายสสฺส\csromansharedfootnote{621e434f4b551ac2}{กาฬหํสสฺส (ก)} เทฺว องฺคานิ คเหตพฺพานิ.}

\prose{ฉตฺตสฺส ตีณิ องฺคานิ คเหตพฺพานิ.}

\prose{เขตฺตสฺส ตีณิ องฺคานิ คเหตพฺพานิ.}

\prose{อคทสฺส เทฺว องฺคานิ คเหตพฺพานิ.}

\prose{โภชนสฺส ตีณิ องฺคานิ คเหตพฺพานิ.}

\prosewithclosermidruled{อิสฺสาสสฺส จตฺตาริ องฺคานิ คเหตพฺพานิ.}{กุมฺภวคฺโค สตฺตโม.}
\prose{รญฺโญ จตฺตาริ องฺคานิ คเหตพฺพานิ.}

\prose{โทวาริกสฺส เทฺว องฺคานิ คเหตพฺพานิ.}

\prose{นิสทาย เอกํ องฺคํ คเหตพฺพํ.}

\setlength{\csromangathapagewidth}{0pt}
\setlength{\csromangathawakwidth}{0pt}
\csromangathameasuregroup{}{ปทีปสฺส เทฺว องฺคานิ คเหตพฺพานิ, \\ มยูรสฺส เทฺว องฺคานิ คเหตพฺพานิ. \\ ตุรงฺคสฺส เทฺว องฺคานิ คเหตพฺพานิ. \\ โสณฺฑิกสฺส เทฺว องฺคานิ คเหตพฺพานิ. \\ อินฺทขีลสฺส เทฺว องฺคานิ คเหตพฺพานิ. \\ ตุลาย เอกํ องฺคํ คเหตพฺพํ. \\ ขคฺคสฺส เทฺว องฺคานิ คเหตพฺพานิ.}
\csromangathameasurewak{อินฺทขีลสฺส เทฺว องฺคานิ คเหตพฺพานิ. \\ ตุลาย เอกํ องฺคํ คเหตพฺพํ. \\ ขคฺคสฺส เทฺว องฺคานิ คเหตพฺพานิ.}
\setlength{\csromangathaleftcol}{0pt}
\csromangathagroup{\csromangathastanza{ปทีปสฺส เทฺว องฺคานิ คเหตพฺพานิ, \\ มยูรสฺส เทฺว องฺคานิ คเหตพฺพานิ. \\ ตุรงฺคสฺส เทฺว องฺคานิ คเหตพฺพานิ. \\ โสณฺฑิกสฺส เทฺว องฺคานิ คเหตพฺพานิ.}\csromangathastanza{\csromangathawak{อินฺทขีลสฺส เทฺว องฺคานิ คเหตพฺพานิ. \\ ตุลาย เอกํ องฺคํ คเหตพฺพํ. \\ ขคฺคสฺส เทฺว องฺคานิ คเหตพฺพานิ.}}}

\prose{\csromanfolio{347}มจฺฉสฺส เทฺว องฺคานิ คเหตพฺพานิ.}

\prose{อิณคฺคาหกสฺส เอกํ องฺคํ คเหตพฺพํ.}

\prose{พฺยาธิตสฺส เทฺว องฺคานิ คเหตพฺพานิ.}

\prose{มตสฺส เทฺว องฺคานิ คเหตพฺพานิ.}

\prose{นทิยา เทฺว องฺคานิ คเหตพฺพานิ.}

\prose{อุสภสฺส เอกํ องฺคํ คเหตพฺพํ.}

\prose{มคฺคสฺส เทฺว องฺคานิ คเหตพฺพานิ.}

\prose{สุงฺกสายิกสฺส เอกํ องฺคํ คเหตพฺพํ.}

\prose{โจรสฺส ตีณิ องฺคานิ คเหตพฺพานิ.}

\prose{สกุณคฺฆิยา เอกํ องฺคํ คเหตพฺพํ.}

\prose{สุนขสฺส เอกํ องฺคํ คเหตพฺพํ.}

\prose{ติกิจฺฉกสฺส ตีณิ องฺคนิ คเหตพฺพานิ.}

\prose{คพฺภินิยา เทฺว องฺคานิ คเหตพฺพานิ.}

\prose{จมริยา เอกํ องฺคํ คเหตพฺพํ.}

\prose{กิกิยา เทฺว องฺคานิ คเหตพฺพานิ.}

\prose{กโปติกาย ตีณิ องฺคานิ คเหตพฺพานิ.}

\prose{เอกนยนสฺส เทฺว องฺคานิ คเหตพฺพานิ.}

\prose{กสฺสกสฺส ตีณิ องฺคานิ คเหตพฺพานิ.}

\prose{ชมฺพุก\-สิงฺคา\-ลิยา เอกํ องฺคํ คเหตพฺพํ.}

\prose{จงฺควารกสฺส เทฺว องฺคานิ คเหตพฺพานิ.}

\prose{ทพฺพิยา เอกํ องฺคํ คเหตพฺพํ.}

\prose{อิณสาธกสฺส ตีณิ องฺคานิ คเหตพฺพานิ.}

\prose{อนุวิจินกสฺส เอกํ องฺคํ คเหตพฺพํ.}

\prose{สารถิสฺส เทฺว องฺคานิ คเหตพฺพานิ.}

\prose{\csromanfolio{348}โภชกสฺส เทฺว องฺคานิ คเหตพฺพานิ.}

\prose{ตุนฺนวายสฺส เอกํ องฺคํ คเหตพฺพํ.}

\prose{นาวิกสฺส เอกํ องฺคํ คเหตพฺพํ.}

\prosewithcloserruled{ภมรสฺส เทฺว องฺคานิ คเหตพฺพานีติ.}{มาติกา นิฏฺฐิตา.}
\csromanmatikaanchor{mtk.201}
\csromanmatikaanchor{mtk.202}
\csromantocmarknopage{paragraph}{1. คทฺรภวคฺค}{mtk.201}
\bookmark[dest={mtk.201},level=4]{1. คทฺรภวคฺค}
\csromantocmarkref{subparagraph}{1. คทฺรภงฺคปญฺห}{mtk.202}
\bookmark[dest={mtk.202},level=5]{1. คทฺรภงฺคปญฺห}
\csromanmarkh{1. คทฺรภวคฺค}\csromanmarkhii{1. คทฺรภ\-งฺค\-ปญฺห}\csromanheadernomark{5}{1. \textbf{คทฺรภวคฺค 1. คทฺรภ}\-\textbf{งฺค}\-\textbf{ปญฺห}}
\proseitem{1}{ภนฺเต นาคเสน “คทฺรภสฺส เอกํ องฺคํ คเหตพฺพนฺ”ติ ยํ วเทสิ, กตมํ ตํ เอกํ องฺคํ คเหตพฺพนฺติ. ยถา มหาราช คทฺรโภ นาม สงฺการกูเฏปิ จรุกฺเกปิ สิงฺฆาฏเกปิ คามทฺวาเรปิ ถุสราสิมฺหิปิ ยตฺถ กตฺถจิ สยติ, น สยนพหุโล โหติ. เอวเมว โข มหาราช โยคินา โยคาวจเรน ติณสนฺถาเรปิ ปณฺณ\-สนฺถาเรปิ กฏฺฐมญฺจเกปิ ฉมายปิ ตตฺถ กตฺถจิ จมฺมขณฺฑํ ปตฺถริตฺวา ยตฺถ กตฺถจิ สยิตพฺพํ, น สยนพหุเลน ภวิตพฺพํ. อิทํ มหาราช คทฺรภสฺส เอกํ องฺคํ คเหตพฺพํ. ภาสิตมฺเปตํ มหาราช ภควตา เทวาติเทเวน\csromansymbolfootnote{*}{สํ 1. 458 ปิฏฺเฐ.}“กลิงฺครู\-ปธานา ภิกฺขเว มหาราช เอตรหิ มม สาวกา วิหรนฺติ อปฺปมตฺตา อาตาปิโน ปธานสฺมินฺ”ติ. ภาสิตมฺเปตํ มหาราช เถเรน สาริปุตฺเตน ธมฺม\-เสนาปตินาปิ\csromandash{}}

\setlength{\csromangathapagewidth}{0pt}
\setlength{\csromangathawakwidth}{0pt}
\csromangathameasuregroup{\textsuperscript{+}“ปลฺลงฺเกน นิสินฺนสฺส, \\ อลํ ผาสุวิหาราย,}{\csromangathabat{\textsuperscript{+}“ปลฺลงฺเกน นิสินฺนสฺส,}{ชณฺณุเกนาภิวสฺสติ.} \\ \csromangathabat{อลํ ผาสุวิหาราย,}{ปหิตตฺตสฺส ภิกฺขุโน”ติ.}}
\csromangathasetleft{\textsuperscript{+}“ปลฺลงฺเกน นิสินฺนสฺส, \\ อลํ ผาสุวิหาราย,}
\csromangathagroupwithcloserruled{\csromangathastanza{\csromangathabat{\csromansymbolfootnote{+}{ขุ 2. 344 ปิฏฺเฐ.}“ปลฺลงฺเกน นิสินฺนสฺส,}{ชณฺณุเกนาภิวสฺสติ.} \\ \csromangathabat{อลํ ผาสุวิหาราย,}{ปหิตตฺตสฺส ภิกฺขุโน”ติ.}}}{คทฺรภ\-งฺค\-ปโญฺห ปฐโม.}
\csromanmatikaanchor{mtk.203}
\csromantocmarkref{subparagraph}{2. กุกฺกุฏงฺคปญฺห}{mtk.203}
\bookmark[dest={mtk.203},level=5]{2. กุกฺกุฏงฺคปญฺห}
\csromanheader{5}{2. \textbf{กุกฺกุฏงฺคปญฺห}}
\proseitem{2}{ภนฺเต นาคเสน “กุกฺกุฏสฺส ปญฺจ องฺคานิ คเหตพฺพานี”ติ ยํ วเทสิ, กตมานิ ตานิ ปญฺจ องฺคานิ คเหตพฺพานีติ. ยถา มหาราช \csromanfolio{349}กุกฺกุโฏ กาเลน สมเยน ปฏิสลฺลียติ. เอวเมว โข มหาราช โยคินา โยคาวจเรน กาเลน สมเยเนว เจติยงฺคณํ สมฺมชฺชิตฺวา ปานียํ ปริโภชนียํ อุปฏฺฐเปตฺวา สรีรํ ปฏิชคฺคิตฺวา นหายิตฺวา เจติยํ วนฺทิตฺวา วุฑฺฒานํ ภิกฺขูนํ ทสฺสนาย คนฺตฺวา กาเลน สมเยน สุญฺญาคารํ ปวิสิตพฺพํ. อิทํ มหาราช กุกฺกุฏสฺส ปฐมํ องฺคํ คเหตพฺพํ.}

\prose{ปุน จปรํ มหาราช กุกฺกุโฏ กาเลน สมเยเนว วุฏฺฐาติ. เอวเมว โข มหาราช โยคินา โยคาวจเรน กาเลน สมเยเนว วุฏฺฐหิตฺวา เจติยงฺคณํ สมฺมชฺชิตฺวา ปานียํ ปริโภชนียํ อุปฏฺฐเปตฺวา สรีรํ ปฏิชคฺคิตฺวา เจติยํ วนฺทิตฺวา ปุนเทว สุญฺญาคารํ ปวิสิตพฺพํ. อิทํ มหาราช กุกฺกุฏสฺส ทุติยํ องฺคํ คเหตพฺพํ.}

\prose{ปุน จปรํ มหาราช กุกฺกุโฏ ปถวิํ ขณิตฺวา ขณิตฺวา อชฺโฌหารํ อชฺโฌหรติ. เอวเมว โข มหาราช โยคินา โยคาวจเรน ปจฺจเวกฺขิตฺวา ปจฺจเวกฺขิตฺวา อชฺโฌหารํ อชฺโฌหริตพฺพํ “เนว ทวาย น มทาย น มณฺฑนาย น วิภูสนาย, ยาวเทว อิมสฺส กายสฺส ฐิติยา ยาปนาย วิหิํสู\-ปรติยา พฺรหฺมจริยา\-นุคฺคหาย, อิติ ปุราณญฺจ เวทนํ ปฏิหงฺขามิ นวญฺจ เวทนํ น อุปฺปาเทสฺสามิ, ยาตฺรา จ เม ภวิสฺสติ อนวชฺชตา จ ผาสุวิหาโร จา”ติ. อิทํ มหาราช กุกฺกุฏสฺส ตติยํ องฺคํ คเหตพฺพํ. ภาสิตมฺเปตํ มหาราช ภควตา เทวาติเทเวน\csromandash{}}

\setlength{\csromangathapagewidth}{0pt}
\setlength{\csromangathawakwidth}{0pt}
\csromangathameasuregroup{\textsuperscript{*}“กนฺตาเร ปุตฺตมํสํว, \\ เอวํ อาหริ อาหารํ,}{\csromangathabat{\textsuperscript{*}“กนฺตาเร ปุตฺตมํสํว,}{อกฺขสฺสพฺภญฺชนํ ยถา.} \\ \csromangathabat{เอวํ อาหริ อาหารํ,}{ยาปนตฺถมมุจฺฉิโต”ติ.}}
\csromangathasetleft{\textsuperscript{*}“กนฺตาเร ปุตฺตมํสํว, \\ เอวํ อาหริ อาหารํ,}
\csromangathagroup{\csromangathastanza{\csromangathabat{\csromansymbolfootnote{*}{วิสุทฺธิ 1. 43 ปิฏฺเฐ.}“กนฺตาเร ปุตฺตมํสํว,}{อกฺขสฺสพฺภญฺชนํ ยถา.} \\ \csromangathabat{เอวํ อาหริ อาหารํ,}{ยาปนตฺถมมุจฺฉิโต”ติ.}}}

\prose{ปุน จปรํ มหาราช กุกฺกุโฏ สจกฺขุโกปิ รตฺติํ อนฺโธ โหติ. เอวเมว โข มหาราช โยคินา โยคาวจเรน อนนฺเธเนว อนฺเธน วิย ภวิตพฺพํ, อรญฺเญปิ โคจรคาเม ปิณฺฑาย จรนฺเตนปิ รชนีเยสุ รูปสทฺทคนฺธรสโผฏฺฐพฺพ\-ธมฺเมสุ อนฺเธน พธิเรน มูเคน วิย ภวิตพฺพํ, น นิมิตฺตํ คเหตพฺพํ, นานุพฺยญฺชนํ คเหตพฺพํ. อิทํ มหาราช กุกฺกุฏสฺส จตุตฺถํ องฺคํ คเหตพฺพํ. ภาสิตมฺเปตํ มหาราช เถเรน มหา\-กจฺจายเนน\csromandash{}}

\setlength{\csromangathapagewidth}{0pt}
\setlength{\csromangathawakwidth}{0pt}
\csromangathameasuregroup{\textsuperscript{+}“จกฺขุมาสฺส ยถา อนฺโธ, \\ ปญฺญวาสฺส\textsuperscript{1} ยถา มูโค, \\ อตฺตอตฺเถ\textsuperscript{1} สมุปฺปนฺเน,}{\csromangathabat{\textsuperscript{+}“จกฺขุมาสฺส ยถา อนฺโธ,}{โสตวา พธิโร ยถา.} \\ \csromangathabat{ปญฺญวาสฺส\textsuperscript{1} ยถา มูโค,}{พลวา ทุพฺพโลริว.} \\ \csromangathabat{อตฺตอตฺเถ\textsuperscript{1} สมุปฺปนฺเน,}{สเยถ มตสายิก นฺ”ติ.}}
\csromangathasetleft{\textsuperscript{+}“จกฺขุมาสฺส ยถา อนฺโธ, \\ ปญฺญวาสฺส\textsuperscript{1} ยถา มูโค, \\ อตฺตอตฺเถ\textsuperscript{1} สมุปฺปนฺเน,}
\csromangathagroup{\csromangathastanza{\csromangathabat{\csromansymbolfootnote{+}{ขุ 2. 295 ปิฏฺเฐ.}“จกฺขุมาสฺส ยถา อนฺโธ,}{โสตวา พธิโร ยถา.} \\ \csromangathabat{ปญฺญวาสฺส\csromansharedfootnote{565992db637c5c64}{ชิวฺหาวสฺส (สี, อิ)} ยถา มูโค,}{พลวา ทุพฺพโลริว.} \\ \csromangathabat{อตฺตอตฺเถ\csromansharedfootnote{565992db637c5c64}{ชิวฺหาวสฺส (สี, อิ)} สมุปฺปนฺเน,}{สเยถ มตสายิก นฺ”ติ.}}}

\proseitem{6}{\csromanfolio{350}ปุน จปรํ มหาราช กุกฺกุโฏ เลฑฺฑุ\-ทณฺฑ\-ลคุฬ\-มุคฺคเรหิ ปริปาติยนฺโตปิ สกํ เคหํ น วิชหติ. เอวเมว โข มหาราช โยคินา โยคาวจเรน จีวรกมฺมํ กโรนฺเตนปิ นวกมฺมํ กโรนฺเตนปิ วตฺต\-ปฺปฏิวตฺตํ กโรนฺเตนปิ อุทฺทิสนฺเตนปิ อุทฺทิสาเป\-นฺเตนปิ โยนิโส มนสิกาโร น วิชหิตพฺโพ, สกํ โข ปเนตํ มหาราช โยคิโน เคหํ, ยทิทํ โยนิโส มนสิกาโร. อิทํ มหาราช กุกฺกุฏสฺส ปญฺจมํ องฺคํ คเหตพฺพํ. ภาสิตมฺเปตํ มหาราช ภควตา เทวาติเทเวน\csromansymbolfootnote{*}{สํ 3. 128 ปิฏฺเฐ.}“โก จ ภิกฺขเว ภิกฺขุโน โคจโร สโก เปตฺติโก วิสโย, ยทิทํ จตฺตาโร สติปฏฺฐานา”ติ. ภาสิตมฺเปตํ มหาราช เถเรน สาริปุตฺเตน ธมฺม\-เสนาปตินาปิ\csromandash{}}

\setlength{\csromangathapagewidth}{0pt}
\setlength{\csromangathawakwidth}{0pt}
\csromangathameasuregroup{“ยถา สุทนฺโต มาตงฺโค, \\ ภกฺขาภกฺขํ วิชานาติ, \\ ตเถว พุทฺธปุตฺเตน, \\ ชินวจนํ น มทฺทิตพฺพํ,}{\csromangathabat{“ยถา สุทนฺโต มาตงฺโค,}{สกํ โสณฺฑํ น มทฺทติ.} \\ \csromangathabat{ภกฺขาภกฺขํ วิชานาติ,}{อตฺตโน วุตฺติกปฺปนํ.} \\ \csromangathabat{ตเถว พุทฺธปุตฺเตน,}{อปฺปมตฺเตน วา ปน.} \\ \csromangathabat{ชินวจนํ น มทฺทิตพฺพํ,}{มนสิการวรุตฺตมนฺ”ติ.}}
\csromangathasetleft{“ยถา สุทนฺโต มาตงฺโค, \\ ภกฺขาภกฺขํ วิชานาติ, \\ ตเถว พุทฺธปุตฺเตน, \\ ชินวจนํ น มทฺทิตพฺพํ,}
\csromangathagroupwithcloserruled{\csromangathastanza{\csromangathabat{“ยถา สุทนฺโต มาตงฺโค,}{สกํ โสณฺฑํ น มทฺทติ.} \\ \csromangathabat{ภกฺขาภกฺขํ วิชานาติ,}{อตฺตโน วุตฺติกปฺปนํ.}}\csromangathastanza{\csromangathabat{ตเถว พุทฺธปุตฺเตน,}{อปฺปมตฺเตน วา ปน.} \\ \csromangathabat{ชินวจนํ น มทฺทิตพฺพํ,}{มนสิการวรุตฺตมนฺ”ติ.}}}{กุกฺกุฏงฺคปโญฺห ทุติโย.}
\csromanmatikaanchor{mtk.204}
\csromantocmarkref{subparagraph}{3. กลนฺทกงฺคปญฺห}{mtk.204}
\bookmark[dest={mtk.204},level=5]{3. กลนฺทกงฺคปญฺห}
\csromanheader{5}{3. \textbf{กลนฺทก}\-\textbf{งฺค}\-\textbf{ปญฺห}}
\proseitem{3}{ภนฺเต นาคเสน “กลนฺทกสฺส เอกํ องฺคํ คเหตพฺพนฺ”ติ ยํ วเทสิ, กตมํ ตํ เอกํ องฺคํ คเหตพฺพนฺติ. ยถา มหาราช กลนฺทโก ปฏิสตฺตุมฺหิ โอปตนฺเต นงฺคุฏฺฐํ ปปฺโผเฏตฺวา มหนฺตํ กตฺวา เตเนว นงฺคุฏฺฐ\-ลคุเฬน ปฏิสตฺตุํ ปฏิพาหติ. เอวเมว โข มหาราช โยคินา โยคาวจเรน กิเลสสตฺตุมฺหิ โอปตนฺเต สติปฏฺฐาน\-ลคุฬํ ปปฺโผเฏตฺวา มหนฺตํ กตฺวา เตเนว สติปฏฺฐาน\-ลคุเฬน สพฺพกิเสสา ปฏิพาหิตพฺพา. อิทํ มหาราช กลนฺทกสฺส เอกํ องฺคํ คเหตพฺพํ. ภาสิตมฺเปตํ มหาราช เถเรน จูฬปนฺถเกน\csromandash{}}

\setlength{\csromangathapagewidth}{0pt}
\setlength{\csromangathawakwidth}{0pt}
\csromangathameasuregroup{“ยทา กิเลสา โอปตนฺติ, \\ สติปฏฺฐานลคุเฬน,}{\csromangathabat{“ยทา กิเลสา โอปตนฺติ,}{สามญฺญคุณธํสนา.} \\ \csromangathabat{สติปฏฺฐานลคุเฬน,}{หนฺตพฺพา เต ปุนปฺปุนนฺ”ติ.}}
\csromangathasetleft{“ยทา กิเลสา โอปตนฺติ, \\ สติปฏฺฐานลคุเฬน,}
\csromangathagroup{\csromangathastanza{\csromangathabat{“ยทา กิเลสา โอปตนฺติ,}{สามญฺญคุณธํสนา.} \\ \csromangathabat{สติปฏฺฐาน\-ลคุเฬน,}{หนฺตพฺพา เต ปุนปฺปุนนฺ”ติ.}}}

\vspace{\tipitakaparskip}
\csromancenter{กลนฺทก\-งฺค\-ปโญฺห ตติโย.}
\csromanmatikaanchor{mtk.205}
\csromantocmarkref{subparagraph}{4. ทีปินิยงฺคปญฺห}{mtk.205}
\bookmark[dest={mtk.205},level=5]{4. ทีปินิยงฺคปญฺห}
\csromanmarkh{6. โอปมฺม\-กถา\-ปญฺห}\csromanmarkhii{4. ทีปินิย\-งฺค\-ปญฺห}\csromanheadernomark{5}{6. โอปมฺม\-กถา\-ปญฺห \textbf{4. ทีปินิย}\-\textbf{งฺค}\-\textbf{ปญฺห}}
\proseitem{4}{\csromanfolio{351}ภนฺเต นาคเสน “ทีปินิยา เอกํ องฺคํ คเหตพฺพนฺ”ติ ยํ วเทสิ, กตมํ ตํ เอกํ องฺคํ คเหตพฺพนฺติ. ยถา มหาราช ทีปินี สกิํ เยว คพฺภํ คณฺหาติ, น ปุนปฺปุนํ ปุริสํ อุเปติ. เอวเมว โข มหาราช โยคินา โยคาวจเรน อายติํ ปฏิสนฺธิํ อุปฺปตฺติํ คพฺภเสยฺยํ จุติํ เภทํ ขยํ วินาสํ สํสารภยํ ทุคฺคติํ วิสมํ สมฺปีฬิตํ ทิสฺวา “ปุนพฺภเว นปฺปฏิสนฺทหิสฺสามี”ติ โยนิโส มนสิกาโร กรณีโย. อิทํ มหาราช ทีปินิยา เอกํ องฺคํ คเหตพฺพํ. ภาสิตมฺเปตํ มหาราช ภควตา เทวาติเทเวน สุตฺตนิปาเต ธนิย\-โคปาลกสุตฺเต*\csromandash{}}

\setlength{\csromangathapagewidth}{0pt}
\setlength{\csromangathawakwidth}{0pt}
\csromangathameasuregroup{}{“อุสโภริว เฉตฺว พนฺธนานิ, \\ นาโค ปูติลตํว ทาลยิตฺวา. \\ นาหํ ปุนุเปสฺสํ คพฺภเสยฺยํ, \\ อถ เจ ปตฺถยสี ปวสฺส เทวา”ติ.}
\csromangathameasurewak{“อุสโภริว เฉตฺว พนฺธนานิ, \\ นาโค ปูติลตํว ทาลยิตฺวา. \\ นาหํ ปุนุเปสฺสํ คพฺภเสยฺยํ, \\ อถ เจ ปตฺถยสี ปวสฺส เทวา”ติ.}
\setlength{\csromangathaleftcol}{0pt}
\csromangathagroupwithcloserruled{\csromangathastanza{\csromangathawak{“อุสโภริว เฉตฺว พนฺธนานิ, \\ นาโค ปูติลตํว ทาลยิตฺวา. \\ นาหํ ปุนุเปสฺสํ คพฺภเสยฺยํ, \\ อถ เจ ปตฺถยสี ปวสฺส เทวา”ติ.}}}{ทีปินิย\-งฺค\-ปโญฺห จตุตฺโถ.}
\csromanmatikaanchor{mtk.206}
\csromantocmarkref{subparagraph}{5. ทีปิกงฺคปญฺห}{mtk.206}
\bookmark[dest={mtk.206},level=5]{5. ทีปิกงฺคปญฺห}
\csromanheader{5}{5. \textbf{ทีปิกงฺคปญฺห}}
\proseitem{5}{ภนฺเต นาคเสน “ทีปิกสฺส เทฺว องฺคานิ คเหตพฺพานี”ติ ยํ วเทสิ, กตมานิ ตานิ เทฺว องฺคานิ คเหตพฺพานีติ. ยถา มหาราช ทีปิโก อรญฺเญ ติณคหนํ วา วนคหนํ วา ปพฺพตคหนํ วา นิสฺสาย นิลียิตฺวา มิเค คณฺหาติ. เอวเมว โข มหาราช โยคินา โยคาวจเรน วิเวกํ เสวิตพฺพํ อรญฺญํ รุกฺขมูลํ ปพฺพตํ กนฺทรํ คิริคุหํ สุสานํ วนปตฺถํ อพฺโภกาสํ ปลาลปุญฺชํ อปฺปสทฺทํ อปฺปนิคฺโฆสํ วิชนวาตํ มนุสฺส\-ราหเสยฺยกํ ปฏิสลฺลาน\-สารุปฺปํ วิเวกํ เสวมาโน หิ มหาราช โยคี โยคาวจโร นจิรสฺเสว ฉฬภิญฺญาสุ จ วสิภาวํ ปาปุณาติ. อิทํ มหาราช ทีปิกสฺส ปฐมํ องฺคํ คเหตพฺพํ. ภาสิตมฺเปตํ มหาราช เถเรหิ ธมฺม\-สงฺคาหเกหิ\csromandash{}}

\setlength{\csromangathapagewidth}{0pt}
\setlength{\csromangathawakwidth}{0pt}
\csromangathameasuregroup{\textsuperscript{*}“ยถาปิ ทีปิโก นาม, \\ ตเถวายํ พุทฺธปุตฺโต, \\ อรญฺญํ ปวิสิตฺวาน,}{\csromangathabat{\textsuperscript{*}“ยถาปิ ทีปิโก นาม,}{นิลียิตฺวา คณฺหเต1 มิเค.} \\ \csromangathabat{ตเถวายํ พุทฺธปุตฺโต,}{ยุตฺตโยโค วิปสฺสโก.} \\ \csromangathabat{อรญฺญํ ปวิสิตฺวาน,}{คณฺหาติ ผลมุตฺตมนฺ”ติ.}}
\csromangathasetleft{\textsuperscript{*}“ยถาปิ ทีปิโก นาม, \\ ตเถวายํ พุทฺธปุตฺโต, \\ อรญฺญํ ปวิสิตฺวาน,}
\csromangathagroup{\csromangathastanza{\csromanfolio{352}\csromangathabat{\csromansymbolfootnote{*}{1. คณฺหตี (สี, อิ), วิสุทฺธิ 1. 262; วิฏฺฐ 2. 12; ทีฏฺฐ 2. 354; มฏฺฐ 1. 253, ปติสํฏฺฐ 2. 89 ปิฏฺฐาทีสุ.}“ยถาปิ ทีปิโก นาม,}{นิลียิตฺวา คณฺหเต1 มิเค.} \\ \csromangathabat{ตเถวายํ พุทฺธปุตฺโต,}{ยุตฺตโยโค วิปสฺสโก.} \\ \csromangathabat{อรญฺญํ ปวิสิตฺวาน,}{คณฺหาติ ผลมุตฺตมนฺ”ติ.}}}

\prose{ปุน จปรํ มหาราช ทีปิโก ยํ กิญฺจิ ปสุํ วธิตฺวา วาเมน ปสฺเสน ปติตํ น ภกฺเขติ. เอวเมว โข มหาราช โยคินา โยคาวจเรน เวฬุทาเนน วา ปตฺตทาเนน วา ปุปฺผทาเนน วา ผลทาเนน วา สินานทาเนน วา มตฺติกาทาเนน วา จุณฺณทาเนน วา ทนฺตกฏฺฐ\-ทาเนน วา มุโขทก\-ทาเนน วา จาตุกมฺยตาย วา มุคฺคสุปฺยตาย\csromansharedfootnote{8d0565006ab4858b}{มุคฺคสุปฺปตาย (สี, อิ); ขุ 7. 290; ขุ 8. 258 ปิฏฺเฐสุ.} วา ปาริภฏฺยตาย\csromansharedfootnote{8dbcb0c52b1eb848}{ปาริภฏฺฏตาย (สี, อิ)} วา ชงฺฆ\-เป\-สนี\-เยน วา เวชฺชกมฺเมน วา ทูตกมฺเมน วา ปหิณคมเนน วา ปิณฺฑปฏิปิณฺเฑน วา ทานา\-นุปฺปทาเนน วา วตฺถุวิชฺชาย วา นกฺขตฺต\-วิชฺชาย วา องฺควิชฺชาย\csromansharedfootnote{78e50c364c3b9fd5}{นควิชฺชาย (ก)} วา อญฺญตร\-ญฺญตเรน วา พุทฺธ\-ปฺปฏิกุฏฺเฐน มิจฺฉาชีเวน นิปฺผาทิตํ โภชนํ น ภุญฺชิตพฺพํ\csromansharedfootnote{c04bd7a2d081a27a}{ปริภุญฺชิตพฺพํ (สี, อิ)} วาเมน ปสฺเสน ปติตํ ปสุํ วิย ทีปิโก. อิทํ มหาราช ทีปิกสฺส ทุติยํ องฺคํ คเหตพฺพํ. ภาสิตมฺเปตํ มหาราช เถเรน สาริปุตฺเตน ธมฺม\-เสนาปตินา\csromandash{}}

\setlength{\csromangathapagewidth}{0pt}
\setlength{\csromangathawakwidth}{0pt}
\csromangathameasuregroup{\textsuperscript{*}“วจีวิญฺญตฺติวิปฺผารา, \\ สเจ ภุตฺโต ภเวยฺยาหํ, \\ ยทิปิ เม อนฺตคุณํ, \\ เนว ภินฺเทยฺยมาชีวํ,}{\csromangathabat{\textsuperscript{*}“วจีวิญฺญตฺติวิปฺผารา,}{อุปฺปนฺนํ มธุปายสํ.} \\ \csromangathabat{สเจ ภุตฺโต ภเวยฺยาหํ,}{สาชีโว ครหิโต มม.} \\ \csromangathabat{ยทิปิ เม อนฺตคุณํ,}{นิกฺขมิตฺวา พหี จเร.} \\ \csromangathabat{เนว ภินฺเทยฺยมาชีวํ,}{จชมาโนปิ ชีวิตนฺ”ติ.}}
\csromangathasetleft{\textsuperscript{*}“วจีวิญฺญตฺติวิปฺผารา, \\ สเจ ภุตฺโต ภเวยฺยาหํ, \\ ยทิปิ เม อนฺตคุณํ, \\ เนว ภินฺเทยฺยมาชีวํ,}
\csromangathagroupwithcloserruled{\csromangathastanza{\csromangathabat{\csromansymbolmark{*}“วจีวิญฺญตฺติวิปฺผารา,}{อุปฺปนฺนํ มธุปายสํ.} \\ \csromangathabat{สเจ ภุตฺโต ภเวยฺยาหํ,}{สาชีโว ครหิโต มม.}}\csromangathastanza{\csromangathabat{ยทิปิ เม อนฺตคุณํ,}{นิกฺขมิตฺวา พหี จเร.} \\ \csromangathabat{เนว ภินฺเทยฺยมาชีวํ,}{จชมาโนปิ ชีวิตนฺ”ติ.}}}{ทีปิกงฺคปโญฺห ปญฺจโม.}
\csromanmatikaanchor{mtk.207}
\csromantocmarkref{subparagraph}{6. กุมฺมงฺคปญฺห}{mtk.207}
\bookmark[dest={mtk.207},level=5]{6. กุมฺมงฺคปญฺห}
\csromanheader{5}{6. \textbf{กุมฺมงฺคปญฺห}}
\proseitem{6}{ภนฺเต นาคเสน “กุมฺมสฺส ปญฺจ องฺคานิ คเหตพฺพานี”ติ ยํ วเทสิ, กตมานิ ตานิ ปญฺจ องฺคานิ คเหตพฺพานีติ. ยถา มหาราช กุมฺโม อุทกจโร อุทเก เยว วาสํ กปฺเปติ. เอวเมว โข มหาราช โยคินา โยคาวจเรน สพฺพปาณภูต\-ปุคฺคลานํ หิตานุกมฺปินา เมตฺตา\-สหคเตน เจตสา วิปุเลน มหคฺคเตน อปฺปมาเณน อเวเรน อพฺยาปชฺเชน สพฺพาวนฺตํ โลกํ ผริตฺวา วิหริตพฺพํ. อิทํ มหาราช กุมฺมสฺส ปฐมํ องฺคํ คเหตพฺพํ.}

\prose{\csromanfolio{353}ปุน จปรํ มหาราช กุมฺโม อุทเก อุปฺปิลวนฺโต สีสํ อุกฺขิปิตฺวา ยทิ โกจิ ปสฺสติ, ตตฺเถว นิมุชฺชติ คาฬฺหโมคาหติ “มา มํ เต ปุน ปสฺเสยฺยุนฺ”ติ, เอวเมว โข มหาราช โยคินา โยคาวจเรน กิเลเสสุ โอปตนฺเตสุ อารมฺมณสเร นิมุชฺชิตพฺพํ คาฬฺหโม\-คาหิตพฺพํ “มา มํ กิเลสา ปุน ปสฺเสยฺยุนฺ”ติ. อิทํ มหาราช กุมฺมสฺส ทุติยํ องฺคํ คเหตพฺพํ.}

\prose{ปุน จปรํ มหาราช กุมฺโม อุทกโต นิกฺขมิตฺวา กายํ โอตาเปติ. เอวเมว โข มหาราช โยคินา โยคาวจเรน นิสชฺชฏฺฐาน\-สยน\-จงฺกมโต มานสํ นีหริตฺวา สมฺมปฺปธาเน มานสํ โอตาเปตพฺพํ. อิทํ มหาราช กุมฺมสฺส ตติยํ องฺคํ คเหตพฺพํ.}

\prose{ปุน จปรํ มหาราช กุมฺโม ปถวิํ ขณิตฺวา วิวิตฺเต วาสํ กปฺเปติ. เอวเมว โข มหาราช โยคินา โยคาวจเรน ลาภ\-สกฺการ\-สิโลกํ ปชหิตฺวา สุญฺญํ วิวิตฺตํ กานนํ วนปตฺถํ ปพฺพตํ กนฺทรํ คิริคุหํ อปฺปสทฺทํ อปฺปนิโฆสํ ปวิวิตฺตโม\-คาหิตฺวา วิวิตฺเต เยว วาจํ อุปคนฺตพฺพํ. อิทํ มหาราช กุมฺมสฺส จตุตฺถํ องฺคํ คเหตพฺพํ. ภาสิตมฺเปตํ มหาราช เถเรน อุปเสเนน วงฺคนฺต\-ปุตฺเตน\csromandash{}}

\setlength{\csromangathapagewidth}{0pt}
\setlength{\csromangathawakwidth}{0pt}
\csromangathameasuregroup{\textsuperscript{*}“วิวิตฺตํ อปฺปนิคฺโฆสํ, \\ เสเว เสนาสนํ ภิกฺขุ,}{\csromangathabat{\textsuperscript{*}“วิวิตฺตํ อปฺปนิคฺโฆสํ,}{วาฬมิคนิเสวิตํ.} \\ \csromangathabat{เสเว เสนาสนํ ภิกฺขุ,}{ปฏิสลฺลานการณา”ติ.}}
\csromangathasetleft{\textsuperscript{*}“วิวิตฺตํ อปฺปนิคฺโฆสํ, \\ เสเว เสนาสนํ ภิกฺขุ,}
\csromangathagroup{\csromangathastanza{\csromangathabat{\csromansymbolfootnote{*}{ขุ 2. 305 ปิฏฺเฐ.}“วิวิตฺตํ อปฺปนิคฺโฆสํ,}{วาฬมิค\-นิเสวิตํ.} \\ \csromangathabat{เสเว เสนาสนํ ภิกฺขุ,}{ปฏิสลฺลานการณา”ติ.}}}

\prose{ปุน จปรํ มหาราช กุมฺโม จาริกํ จรมาโน ยทิ กิญฺจิ ปสฺสติ วา, สทฺทํ สุณาติ วา, โสณฺฑิปญฺจมานิ องฺคานิ สเก กปาเล นิทหิตฺวา อปฺโปสฺสุกฺโก ตุณฺหีภูโต ติฏฺฐติ กายม\-นุรกฺขนฺโต. เอวเมว โข มหาราช โยคินา โยคาวจเรน สพฺพตฺถ รูปสทฺทคนฺธรสโผฏฺฐพฺพ\-ธมฺเมสุ อาปตนฺเตสุ ฉสุ ทฺวาเรสุ สํวรกวาฏํ อนุคฺฆาเฏตฺวา มานสํ สโมทหิตฺวา สํวรํ กตฺวา สเตน สมฺปชาเนน วิหาตพฺพํ สมณธมฺมํ อนุรกฺข\-มาเนน. อิทํ มหาราช กุมฺมสฺส ปญฺจมํ องฺคํ คเหตพฺพํ. ภาสิตมฺเปตํ มหาราช ภควตา เทวาติเทเวน สํยุตฺตนิกาย\-วเร กุมฺมู\-ปม\-สุตฺตนฺเต\csromandash{}}

\setlength{\csromangathapagewidth}{0pt}
\setlength{\csromangathawakwidth}{0pt}
\csromangathameasuregroup{}{\textsuperscript{*}“กุมฺโมว องฺคานิ สเก กปาเล, \\ สเมทหํ ภิกฺขุ มโนวิตกฺเก. \\ อนิสฺสิโต อญฺญมเหฐยาโน, \\ ปรินิพฺพุโตนูปเวเทยฺย กญฺจี”ติ.}
\csromangathameasurewak{\textsuperscript{*}“กุมฺโมว องฺคานิ สเก กปาเล, \\ สเมทหํ ภิกฺขุ มโนวิตกฺเก. \\ อนิสฺสิโต อญฺญมเหฐยาโน, \\ ปรินิพฺพุโตนูปเวเทยฺย กญฺจี”ติ.}
\setlength{\csromangathaleftcol}{0pt}
\csromangathagroupwithcloserruled{\csromangathastanza{\csromangathawak{\csromanfolio{354}\csromansymbolfootnote{*}{สํ 2. 386 ปิฏฺเฐ.}“กุมฺโมว องฺคานิ สเก กปาเล, \\ สเมทหํ ภิกฺขุ มโนวิตกฺเก. \\ อนิสฺสิโต อญฺญมเหฐยาโน, \\ ปรินิพฺพุโตนูปเวเทยฺย กญฺจี”ติ.}}}{กุมฺมงฺคปโญฺห ฉฏฺโฐ.}
\csromanmatikaanchor{mtk.208}
\csromantocmarkref{subparagraph}{7. วํสงฺคปญฺห}{mtk.208}
\bookmark[dest={mtk.208},level=5]{7. วํสงฺคปญฺห}
\csromanheader{5}{7. \textbf{วํสงฺคปญฺห}}
\proseitem{7}{ภนฺเต นาคเสน “วํสสฺส เอกํ องฺคํ คเหตพฺพนฺ”ติ ยํ วเทสิ, กตมํ ตํ เอกํ องฺคํ คเหตพฺพนฺติ. ยถา มหาราช วํโส ยตฺถ วาโต, ตตฺถ อนุโลเมติ, นาญฺญ\-ตฺถม\-นุธาวติ. เอวเมว โข มหาราช โยคินา โยคาวจเรน ยํ พุทฺเธน ภควตา ภาสิตํ นวงฺคํ สตฺถุ สาสนํ, ตํ อนุโลมยิตฺวา กปฺปิเย อนวชฺเช ฐตฺวา สมณธมฺมํ เยว ปริเยสิตพฺพํ. อิทํ มหาราช วํสสฺส เอกํ องฺคํ คเหตพฺพํ. ภาสิตมฺเปตํ มหาราช เถเรน ราหุเลน\_\_}

\setlength{\csromangathapagewidth}{0pt}
\setlength{\csromangathawakwidth}{0pt}
\csromangathameasuregroup{“นวงฺคํ พุทฺธวจนํ, \\ กปฺปิเย อนวชฺชสฺมิํ,}{\csromangathabat{“นวงฺคํ พุทฺธวจนํ,}{อนุโลเมตฺวาน สพฺพทา.} \\ \csromangathabat{กปฺปิเย อนวชฺชสฺมิํ,}{ฐตฺวาปายํ สมุตฺตรินฺ”ติ.}}
\csromangathasetleft{“นวงฺคํ พุทฺธวจนํ, \\ กปฺปิเย อนวชฺชสฺมิํ,}
\csromangathagroupwithcloserruled{\csromangathastanza{\csromangathabat{“นวงฺคํ พุทฺธวจนํ,}{อนุโลเมตฺวาน สพฺพทา.} \\ \csromangathabat{กปฺปิเย อนวชฺชสฺมิํ,}{ฐตฺวาปายํ สมุตฺตรินฺ”ติ.}}}{วํสงฺคปโญฺห สตฺตโม.}
\csromanmatikaanchor{mtk.209}
\csromantocmarkref{subparagraph}{8. จาปงฺคปญฺห}{mtk.209}
\bookmark[dest={mtk.209},level=5]{8. จาปงฺคปญฺห}
\csromanheader{5}{8. \textbf{จาปงฺคปญฺห}}
\proseitem{8}{ภนฺเต นาคเสน “จาปสฺส เอกํ องฺคํ คเหตพฺพนฺ”ติ ยํ วเทสิ, กตมํ ตํ เอกํ องฺคํ คเหตพฺพนฺติ. ยถา มหาราช จาโป สุตจฺฉิโต นมิโต\csromansharedfootnote{aabb03b633f1f261}{มิโต (สี, อิ, ก)} ยาวคฺคมูลํ สมกเมว อนุนมติ นปฺปฏิตฺถมฺภติ. เอวเมว โข มหาราช โยคินา โยคาวจเรน เถรนวมชฺฌิม\-สมเกสุ อนุนมิตพฺพํ นปฺปฏิผริตพฺพํ. อิทํ มหาราช จาปสฺส เอกํ องฺคํ คเหตพฺพํ. ภาสิตมฺเปตํ มหาราช ภควตา เทวาติเทเวน วุธุร(ปุณฺณก) ชาตเก\csromandash{}}

\setlength{\csromangathapagewidth}{0pt}
\setlength{\csromangathawakwidth}{0pt}
\csromangathameasuregroup{“จาโปวูนุทโร ธีโร, \\ ปฏิโลมํ น วตฺเตยฺย,}{\csromangathabat{“จาโปวูนุทโร ธีโร,}{วํโส วาปิ ปกมฺปเย\textsuperscript{1}.} \\ \csromangathabat{ปฏิโลมํ น วตฺเตยฺย,}{ส ราชวสติํ วเส”ติ.}}
\csromangathasetleft{“จาโปวูนุทโร ธีโร, \\ ปฏิโลมํ น วตฺเตยฺย,}
\csromangathagroupwithcloserruled{\csromangathastanza{\csromanfolio{355}\csromangathabat{“จาโปวูนุทโร ธีโร,}{วํโส วาปิ ปกมฺปเย\csromansharedfootnote{31fc77139993bb48}{จาโป วา นุน เม ธีโร, วํโสว อนุโลมยํ. (สี, อิ, ก) ขุ 6. 292 ปิฏฺเฐ.}.} \\ \csromangathabat{ปฏิโลมํ น วตฺเตยฺย,}{ส ราชวสติํ วเส”ติ.}}}{จาปงฺคปโญฺห อฏฺฐโม.}
\csromanmatikaanchor{mtk.210}
\csromantocmarkref{subparagraph}{9. วายสงฺคปญฺห}{mtk.210}
\bookmark[dest={mtk.210},level=5]{9. วายสงฺคปญฺห}
\csromanheader{5}{9. \textbf{วายสงฺคปญฺห}}
\proseitem{9}{ภนฺเต นาคเสน “วายสสฺส เทฺว องฺคานิ คเหตพฺพานี”ติ ยํ วเทสิ, กตมานิ ตานิ เทฺว องฺคานิ คเหตพฺพานีติ. ยถา มหาราช วายโส อาสงฺกิต\-ปริสงฺกิโต ยตฺตปฺปยตฺโต จรติ. เอวเมว โข มหาราช โยคินา โยคาวจเรน อาสงฺกิต\-ปริสงฺกิเตน ยตฺตปยตฺเตน อุปฏฺฐิตาย สติยา สํวุเตหิ อินฺทฺริเยหิ จริตพฺพํ. อิทํ มหาราช วายสสฺส ปฐมํ องฺคํ คเหตพฺพํ.}

\prose{ปุน จปรํ มหาราช วายโส ยํ กิญฺจิ โภชนํ ทิสฺวา ญาตีหิ สํวิภชิตฺวา ภุญฺชติ. เอวเมว โข มหาราช โยคินา โยคาจเรน เย เต ลาภา ธมฺมิกา ธมฺมลทฺธา อนฺตมโส ปตฺต\-ปริยาปนฺน\-มตฺตมฺปิ, ตถารูเปหิ ลาเภหิ ปฏิวิภตฺต\-โภคินา ภวิตพฺพํ สีลวนฺเตหิ สพฺรหฺมจารีหิ. อิทํ มหาราช วายสสฺส ทุติยํ องฺคํ คเหตพฺพํ. ภาสิตมฺเปตํ มหาราช เถเรน สาริปุตฺเตน ธมฺม\-เสนาปตินา\csromandash{}}

\setlength{\csromangathapagewidth}{0pt}
\setlength{\csromangathawakwidth}{0pt}
\csromangathameasuregroup{“สเจ เม อุปนาเมนฺติ, \\ สพฺเพ สํวิภชิตฺวาน,}{\csromangathabat{“สเจ เม อุปนาเมนฺติ,}{ยถาลทฺธํ ตปสฺสิโน.} \\ \csromangathabat{สพฺเพ สํวิภชิตฺวาน,}{ตโต ภุญฺชามิ โภชนนฺ”ติ.}}
\csromangathasetleft{“สเจ เม อุปนาเมนฺติ, \\ สพฺเพ สํวิภชิตฺวาน,}
\csromangathagroupwithcloserruled{\csromangathastanza{\csromangathabat{“สเจ เม อุปนาเมนฺติ,}{ยถาลทฺธํ ตปสฺสิโน.} \\ \csromangathabat{สพฺเพ สํวิภชิตฺวาน,}{ตโต ภุญฺชามิ โภชนนฺ”ติ.}}}{วายสงฺคปโญฺห นวโม.}
\csromanmatikaanchor{mtk.211}
\csromantocmarkref{subparagraph}{10. มกฺกฏงฺคปญฺห}{mtk.211}
\bookmark[dest={mtk.211},level=5]{10. มกฺกฏงฺคปญฺห}
\csromanheader{5}{10. \textbf{มกฺกฏงฺคปญฺห}}
\proseitem{10}{ภนฺเต นาคเสน “มกฺกฏสฺส เทฺว องฺคานิ คเหตพฺพานี”ติ ยํ วเทสิ, กตมานิ ตานิ เทฺว องฺคานิ คเหตพฺพานีติ. ยถา มหาราช มกฺกโฏ วาสมุ\-ปคจฺฉนฺโต ตถารูเป โอกาเส มหติ\-มหารุกฺเข ปวิวิตฺเต สพฺพฏฺฐก\-สาเข\csromansharedfootnote{72d0aee6d8e39c12}{สพฺพตฺถกสาเข (สฺยา), สพฺพฏฺฐสาเข (ก)} ภีรุตฺตาเณ วาสมุ\-ปคจฺฉติ. เอวเมว โข \csromanfolio{356}มหาราช โยคินา โยคาวจเรน ลชฺชิํ เปสลํ สีลวนฺตํ กลฺยาณธมฺมํ พหุสฺสุตํ ธมฺมธรํ วินยธรํ ปิยํ ครุภาวนียํ วตฺตารํ วจนกฺขมํ โอวาทกํ วิญฺญาปกํ สนฺทสฺสกํ สมาทปกํ สมุตฺเตชกํ สมฺปหํสกํ เอวรูปํ กลฺยาณมิตฺตํ อาจริยํ นิสฺสาย วิหริตพฺพํ. อิทํ มหาราช มกฺกฏสฺส ปฐมํ องฺคํ คเหตพฺพํ.}

\prose{ปุน จปรํ มหาราช มกฺกโฏ รุกฺเข เยว จรติ ติฏฺฐติ นิสีทติ, ยทิ นิทฺทํ โอกฺกมติ, ตตฺเถว รตฺติํ วาสม\-นุภวติ. เอวเมว โข มหาราช โยคินา โยคาวจเรน ปวนา\-ภิมุเขน ภวิตพฺพํ, ปวเน เยว ฐานจงฺกม\-นิสชฺชา\-สยนํ นิทฺทํ โอกฺกมิตพฺพํ, ตตฺเถว สติปฏฺฐานม\-นุภวิตพฺพํ. อิทํ มหาราช มกฺกฏสฺส ทุติยํ องฺคํ คเหตพฺพํ. ภาสิตมฺเปตํ มหาราช เถเรน สาริปุตฺเตน ธมฺม\-เสนาปตินา\csromandash{}}

\setlength{\csromangathapagewidth}{0pt}
\setlength{\csromangathawakwidth}{0pt}
\csromangathameasuregroup{“จงฺกมนฺโตปิ ติฏฺฐนฺโต, \\ ปวเน โสภเต ภิกฺขุ,}{\csromangathabat{“จงฺกมนฺโตปิ ติฏฺฐนฺโต,}{นิสชฺชาสยเนน วา.} \\ \csromangathabat{ปวเน โสภเต ภิกฺขุ,}{ปวนนฺตํว วณฺณิตนฺ”ติ.}}
\csromangathasetleft{“จงฺกมนฺโตปิ ติฏฺฐนฺโต, \\ ปวเน โสภเต ภิกฺขุ,}
\csromangathagroup{\csromangathastanza{\csromangathabat{“จงฺกมนฺโตปิ ติฏฺฐนฺโต,}{นิสชฺชาสยเนน วา.} \\ \csromangathabat{ปวเน โสภเต ภิกฺขุ,}{ปวนนฺตํว วณฺณิตนฺ”ติ.}}}

\vspace{\tipitakaparskip}
\csromancenter{มกฺกฏงฺคปโญฺห ทสโม.}
\nitthitammidruled{\textbf{คทฺรภวคฺโค ปฐโม.}}
\csromancenter{\textbf{ตสฺสุทฺทานํ}}
\setlength{\csromangathapagewidth}{0pt}
\setlength{\csromangathawakwidth}{0pt}
\csromangathameasuregroup{คทฺรโภ เจว\textsuperscript{1} กุกฺกุโฏ, \\ กุมฺโม วํโส จ จาโป จ,}{\csromangathabat{คทฺรโภ เจว\textsuperscript{1} กุกฺกุโฏ,}{กลนฺโท ทีปินิ ทีปิโก.} \\ \csromangathabat{กุมฺโม วํโส จ จาโป จ,}{วายโส อถ มกฺกโฏติ.}}
\csromangathasetleft{คทฺรโภ เจว\textsuperscript{1} กุกฺกุโฏ, \\ กุมฺโม วํโส จ จาโป จ,}
\csromangathagroup{\csromangathastanza{\csromangathabat{คทฺรโภ เจว\csromansharedfootnote{90fa0787eea1fc7e}{โฆรสฺสโร จ (สี, สฺยา, อิ)} กุกฺกุโฏ,}{กลนฺโท ทีปินิ ทีปิโก.} \\ \csromangathabat{กุมฺโม วํโส จ จาโป จ,}{วายโส อถ มกฺกโฏติ.}}}

\csromanmatikaanchor{mtk.212}
\csromantocmarknopage{paragraph}{2. สมุทฺทวคฺค}{mtk.212}
\bookmark[dest={mtk.212},level=4]{2. สมุทฺทวคฺค}
\csromanheader{4}{2. \textbf{สมุทฺทวคฺค}}
\csromanmatikaanchor{mtk.213}
\csromantocmarkref{subparagraph}{1. ลาพุลตงฺคปญฺห}{mtk.213}
\bookmark[dest={mtk.213},level=5]{1. ลาพุลตงฺคปญฺห}
\csromanheader{5}{1. \textbf{ลาพุ}\-\textbf{ลต}\-\textbf{งฺค}\-\textbf{ปญฺห}}
\proseitem{1}{\csromanfolio{357}ภนฺเต นาคเสน “ลาพุลตาย เอกํ องฺคํ คเหตพฺพนฺ”ติ ยํ วเทสิ, กตมํ ตํ องฺคํ คเหตพฺพนฺติ. ยถา มหาราช ลาพุลตา ติเณ วา กฏฺเฐ วา ลตาย วา โสณฺฑิกาหิ อาลมฺพิตฺวา ตสฺสูปริ วฑฺฒติ. เอวเมว โข มหาราช โยคินา โยคาวจเรน อรหตฺเต อภิวฑฺฒิตุ\-กาเมน มนสา อารมฺมณํ อาลมฺพิตฺวา อรหตฺเต อภิวฑฺฒิตพฺพํ. อิทํ มหาราช ลาพุลตาย เอกํ องฺคํ คเหตพฺพํ. ภาสิตมฺเปตํ มหาราช เถเรน สาริปุตฺเตน ธมฺม\-เสนาปตินา\csromandash{}}

\setlength{\csromangathapagewidth}{0pt}
\setlength{\csromangathawakwidth}{0pt}
\csromangathameasuregroup{“ยถา ลาพุลตา นาม, \\ อาลมฺพิตฺวา โสณฺฑิกาหิ, \\ ตเถว พุทฺธปุตฺเตน, \\ อารมฺมณํ อาลมฺพิตฺวา,}{\csromangathabat{“ยถา ลาพุลตา นาม,}{ติเณ กฏฺเฐ ลตาย วา.} \\ \csromangathabat{อาลมฺพิตฺวา โสณฺฑิกาหิ,}{ตโต วฑฺฒติ อุปฺปริ.} \\ \csromangathabat{ตเถว พุทฺธปุตฺเตน,}{อรหตฺตผลกามินา.} \\ \csromangathabat{อารมฺมณํ อาลมฺพิตฺวา,}{วฑฺฒิตพฺพํ อเสกฺขผเล”ติ.}}
\csromangathasetleft{“ยถา ลาพุลตา นาม, \\ อาลมฺพิตฺวา โสณฺฑิกาหิ, \\ ตเถว พุทฺธปุตฺเตน, \\ อารมฺมณํ อาลมฺพิตฺวา,}
\csromangathagroupwithcloserruled{\csromangathastanza{\csromangathabat{“ยถา ลาพุลตา นาม,}{ติเณ กฏฺเฐ ลตาย วา.} \\ \csromangathabat{อาลมฺพิตฺวา โสณฺฑิกาหิ,}{ตโต วฑฺฒติ อุปฺปริ.}}\csromangathastanza{\csromangathabat{ตเถว พุทฺธปุตฺเตน,}{อรหตฺตผลกามินา.} \\ \csromangathabat{อารมฺมณํ อาลมฺพิตฺวา,}{วฑฺฒิตพฺพํ อเสกฺขผเล”ติ.}}}{ลาพุ\-ลต\-งฺค\-ปโญฺห ปฐโม.}
\csromanmatikaanchor{mtk.214}
\csromantocmarkref{subparagraph}{2. ปทุมงฺคปญฺห}{mtk.214}
\bookmark[dest={mtk.214},level=5]{2. ปทุมงฺคปญฺห}
\csromanheader{5}{2. \textbf{ปทุมงฺคปญฺห}}
\proseitem{2}{ภนฺเต นาคเสน “ปทุมสฺส ตีณิ องฺคานิ คเหตพฺพานี”ติ ยํ วเทสิ, กตมานิ ตานิ ตีณิ องฺคานิ คเหตพฺพานีติ. ยถา มหาราช ปทุมํ อุทเก ชาตํ อุทเก สํวทฺธํ อนุปลิตฺตํ อุทเกน. เอวเมว โข มหาราช โยคินา โยคาวเจเรน กุเล คเณ ลาเภ ยเส สกฺกาเร สมฺมานนาย ปริโภค\-ปจฺจเยสุ จ สพฺพตฺถ อนุปลิตฺเตน ภวิตพฺพํ. อิทํ มหาราช ปทุมสฺส ปฐมํ องฺคํ คเหตพฺพํ.}

\prose{ปุน จปรํ มหาราช ปทุมํ อุทกา อจฺจุคฺคมฺม ฐาติ. เอวเมว โข มหาราช โยคินา โยคาวจเรน สพฺพโลกํ อภิภวิตฺวา อจฺจุคฺคมฺม โลกุตฺตร\-ธมฺเม ฐาตพฺพํ. อิทํ มหาราช ปทุมสฺส ทุติยํ องฺคํ คเหตพฺพํ.}

\csromanmatikaanchor{mtk.215}
\csromantocmarkref{subparagraph}{3. วีชงฺคปญฺห}{mtk.215}
\bookmark[dest={mtk.215},level=5]{3. วีชงฺคปญฺห}
\prosewithcloserruled{\csromanfolio{358}ปุน จปรํ มหาราช ปทุมํ อปฺป\-มตฺต\-เกนปิ อนิเลน เอริตํ จลติ. เอวเมว โข มหาราช โยคินา โยคาวจเรน อปฺป\-มตฺต\-เกสุปิ กิเลเสสุ สํยโม กรณีโย, ภยทสฺสาวินา วิหริตพฺพํ. อิทํ มหาราช ปทุมสฺส ตติยํ องฺคํ คเหตพฺพํ. ภาสิตมฺเปตํ มหาราช ภควตา เทวาติเทเวน\csromansymbolfootnote{*}{อภิ 2. 253; ที 1. 59; ม 1. 39; สํ 3. 163; อํ 1. 65 ปิฏฺฐาทีสุ.}“อณุมตฺเตสุ วชฺเชสุ ภยทสฺสาวี สมาทาย สิกฺขติ สิกฺขาปเทสู”ติ.}{ปทุมงฺคปโญฺห ทุติโย.}
\csromanheader{2}{3. \textbf{พีชงฺคปญฺห}}
\proseitem{3}{ภนฺเต นาคเสน “พีชสฺส เทฺว องฺคานิ คเหตพฺพานี”ติ ยํ วเทสิ, กตมานิ ตานิ เทฺว องฺคานิ คเหตพฺพานีติ. ยถา มหาราช พีชํ อปฺปกมฺปิ สมานํ ภทฺทเก เขตฺเต วุตฺตํ เทเว สมฺมา ธารํ ปเวจฺฉนฺเต สุพหูนิ ผลานิ อนุทสฺสติ. เอวเมว โข มหาราช โยคินา โยคาวจเรน ยถา ปฏิปาทิตํ สีลํ เกวลํ สามญฺญผลม\-นุ\-ทสฺสติ. เอวํ สมฺมา ปฏิปชฺชิตพฺพํ. อิทํ มหาราช พีชสฺส ปฐมํ องฺคํ คเหตพฺพํ.}

\prose{ปุน จปรํ มหาราช พีชํ สุปริโสธิเต เขตฺเต โรปิตํ ขิปฺปเมว สํวิรูหติ. เอวเมว โข มหาราช โยคินา โยคาวจเรน มานสํ สุปริคฺคหิตํ สุญฺญาคาเร ปริโสธิตํ สติปฏฺฐาน\-เขตฺต\-วเร ขิตฺตํ ขิปฺปเมว วิรูหติ. อิทํ มหาราช พีชสฺส ทุติยํ องฺคํ คเหตพฺพํ. ภาสิตมฺเปตํ มหาราช เถเรน อนุรุทฺเธน\csromandash{}}

\setlength{\csromangathapagewidth}{0pt}
\setlength{\csromangathawakwidth}{0pt}
\csromangathameasuregroup{“ยถาปิ เขตฺเต\textsuperscript{1} ปริสุทฺเธ, \\ วิปุลํ ตสฺส ผลํ โหติ, \\ ตเถว โยคินา จิตฺตํ, \\ สติปฏฺฐานเขตฺตมฺหิ,}{\csromangathabat{“ยถาปิ เขตฺเต\textsuperscript{1} ปริสุทฺเธ,}{พีชญฺจสฺส ปติฏฺฐิตํ.} \\ \csromangathabat{วิปุลํ ตสฺส ผลํ โหติ,}{อปิ โตเสติ กสฺสกํ.} \\ \csromangathabat{ตเถว โยคินา จิตฺตํ,}{สุญฺญาคาเร วิโสธิตํ.} \\ \csromangathabat{สติปฏฺฐานเขตฺตมฺหิ,}{ขิปฺปเมว วิรูหตี”ติ.}}
\csromangathasetleft{“ยถาปิ เขตฺเต\textsuperscript{1} ปริสุทฺเธ, \\ วิปุลํ ตสฺส ผลํ โหติ, \\ ตเถว โยคินา จิตฺตํ, \\ สติปฏฺฐานเขตฺตมฺหิ,}
\csromangathagroup{\csromangathastanza{\csromangathabat{“ยถาปิ เขตฺเต\csromansharedfootnote{3f8b01435c8fb9d3}{ยถา เขตฺเต (สี)} ปริสุทฺเธ,}{พีชญฺจสฺส ปติฏฺฐิตํ.} \\ \csromangathabat{วิปุลํ ตสฺส ผลํ โหติ,}{อปิ โตเสติ กสฺสกํ.}}\csromangathastanza{\csromangathabat{ตเถว โยคินา จิตฺตํ,}{สุญฺญาคาเร วิโสธิตํ.} \\ \csromangathabat{สติปฏฺฐานเขตฺตมฺหิ,}{ขิปฺปเมว วิรูหตี”ติ.}}}

\vspace{\tipitakaparskip}
\csromancenter{พีชงฺคปโญฺห ตติโย.}
\csromanmatikaanchor{mtk.216}
\csromantocmarkref{subparagraph}{4. สาลกลฺยาณิกงฺคปญฺห}{mtk.216}
\bookmark[dest={mtk.216},level=5]{4. สาลกลฺยาณิกงฺคปญฺห}
\csromanmarkh{6. โอปมฺม\-กถา\-ปญฺห}\csromanmarkhii{4. สาล\-กลฺยาณิก\-งฺค\-ปญฺห}\csromanheadernomark{5}{6. โอปมฺม\-กถา\-ปญฺห \textbf{4. สาล}\-\textbf{กลฺยาณิก}\-\textbf{งฺค}\-\textbf{ปญฺห}}
\proseitem{4}{\csromanfolio{359}ภนฺเต นาคเสน “สาล\-กลฺยาณิกาย เอกํ องฺคํ คเหตพฺพนฺ”ติ ยํ วเทสิ, กตมํ ตํ เอกํ องฺคํ คเหตพฺพนฺติ. ยถา มหาราช สาลกลฺยาณิกา นาม อนฺโตปถวิยํ เยว อภิวฑฺฒติ หตฺถสตมฺปิ ภิยฺโยปิ. เอวเมว โข มหาราช โยคินา โยคาวจเรน จตฺตาริ สามญฺญผลานิ จตสฺโส ปฏิสมฺภิทา ฉฬภิญฺญาโย เกวลญฺจ สมณธมฺมํ สุญฺญาคาเร เยว ปริปูร\-ยิตพฺพํ. อิทํ มหาราช สาล\-กลฺยาณิกาย เอกํ องฺคํ คเหตพฺพํ. ภาสิตมฺเปตํ มหาราช เถเรน ราหุเลน\csromandash{}}

\setlength{\csromangathapagewidth}{0pt}
\setlength{\csromangathawakwidth}{0pt}
\csromangathameasuregroup{“สาลกลฺยาณิกา นาม, \\ อนฺโตปถวิยํ เยว, \\ ยถา กาลมฺหิ สมฺปตฺเต, \\ อุคฺคญฺฉิตฺวาน เอกาหํ, \\ เอวเมวาหํ มหาวีร,}{\csromangathabat{“สาลกลฺยาณิกา นาม,}{ปาทโป ธรณีรุโห.} \\ \csromangathabat{อนฺโตปถวิยํ เยว,}{สตหตฺโถปิ วฑฺฒติ.} \\ \csromangathabat{ยถา กาลมฺหิ สมฺปตฺเต,}{ปริปาเกน โส ทุโม.} \\ \csromangathabat{อุคฺคญฺฉิตฺวาน เอกาหํ,}{สตหตฺโถปิ วฑฺฒติ.} \\ \csromangathabat{เอวเมวาหํ มหาวีร,}{สาลกลฺยาณิกา วิย.}}
\csromangathasetleft{“สาลกลฺยาณิกา นาม, \\ อนฺโตปถวิยํ เยว, \\ ยถา กาลมฺหิ สมฺปตฺเต, \\ อุคฺคญฺฉิตฺวาน เอกาหํ, \\ เอวเมวาหํ มหาวีร,}
\csromangathagroup{\csromangathastanza{\csromangathabat{“สาลกลฺยาณิกา นาม,}{ปาทโป ธรณีรุโห.} \\ \csromangathabat{อนฺโตปถวิยํ เยว,}{สตหตฺโถปิ วฑฺฒติ.}}\csromangathastanza{\csromangathabat{ยถา กาลมฺหิ สมฺปตฺเต,}{ปริปาเกน โส ทุโม.} \\ \csromangathabat{อุคฺคญฺฉิตฺวาน เอกาหํ,}{สตหตฺโถปิ วฑฺฒติ.} \\ \csromangathabat{เอวเมวาหํ มหาวีร,}{สาลกลฺยาณิกา วิย.}}}

\prose{อพฺภนฺตเร สุญฺญาคาเร, ธมฺมโต อภิวฑฺฒยินฺ”ติ. สาล\-กลฺยาณิก\-งฺค\-ปโญฺห จตุตฺโถ.}
\csromansectionrule

\csromanmatikaanchor{mtk.217}
\csromantocmarkref{subparagraph}{5. นาวงฺคปญฺห}{mtk.217}
\bookmark[dest={mtk.217},level=5]{5. นาวงฺคปญฺห}
\csromanheader{5}{5. \textbf{นาวงฺคปญฺห}}
\proseitem{5}{ภนฺเต นาคเสน “นาวาย ตีณิ องฺคานิ คเหตพฺพานี”ติ ยํ วเทสิ, กตมานิ ตานิ ตีณิ องฺคานิ คเหตพฺพานีติ. ยถา มหาราช นาวา พหุวิธ\-ทารุสงฺฆาฏ\-สมวาเยน พหุมฺปิ ชนํ ตารยติ. เอวเมว โข มหาราช โยคินา โยคาวจเรน อาจาร\-สีล\-คุณ\-วตฺต\-ปฺปฏิวตฺต\-พหุวิธ\-ธมฺม\-สงฺฆาฏ\-สมวาเยน สเทวโก โลโก ตารยิตพฺโพ. อิทํ มหาราช นาวาย ปฐมํ องฺคํ คเหตพฺพํ.}

\prose{ปุน จปรํ มหาราช นาวา พหุวิธ- อูมิตฺถนิตเวควิสฏมาฏฺฏเวคํ สหติ. เอวเมว โข มหาราช โยคินา โยคาวจเรน พหุวิธ\-กิเลส\-อูมิ\-เวคํ ลาภสกฺการ\-ยส\-สิโลก\-ปูชน\-วนฺทนา ปรกุเลสุ นินฺทาปสํสา\-สุขทุกฺข\-สมฺมานน\-วิมานน\-พหุวิธ\-โทส\-อูมิ\-เวคญฺจ สหิตพฺพํ. อิทํ มหาราช นาวาย ทุติยํ องฺคํ คเหตพฺพํ.}

\prosewithcloserruled{\csromanfolio{360}ปุน จปรํ มหาราช นาวา อปริมิตมนนฺตมปารมกฺโขภิตคมฺภีเร มหติ\-มหา\-โฆเส ติมิติมิงฺคล\-มกร\-มจฺฉคณา\-กุเล มหติ\-มหาสมุทฺเท จรติ. เอวเมว โข โยคินา โยคาวจเรน ติปริวฏฺฏ\-ทฺวาทสาการ จตุสจฺจาภิสมย\-ปฺปฏิเวเธ มานสํ สญฺจาร\-ยิตพฺพํ. อิทํ มหาราช นาวาย ตติยํ องฺคํ คเหตพฺพํ. ภาสิตมฺเปตํ มหาราช ภควตา เทวาติเทเวน สํยุตฺตนิกาย\-วเร สจฺจสํยุตฺเต\csromandash{}\csromansymbolfootnote{*}{สํ 3. 366 ปิฏฺเฐ.}“วิตกฺเกนฺตา จ โข ตุเมฺห ภิกฺขเว “อิทํ ทุกฺขนฺ”ติ วิตกฺเกยฺยาถ, “อยํ ทุกฺขสมุทโย”ติ วิตกฺเกยฺยาถ, “อยํ ทุกฺขนิโรโธ”ติ วิตกฺเกยฺยาถ, “อยํ ทุกฺขนิโรธ\-คามินี ปฏิปที”ติ วิตกฺเกยฺยาถี”ติ.}{นาวงฺคปโญฺห ปญฺจโม.}
\csromanmatikaanchor{mtk.218}
\csromantocmarkref{subparagraph}{6. นาวาลคฺคนกงฺคปญฺห}{mtk.218}
\bookmark[dest={mtk.218},level=5]{6. นาวาลคฺคนกงฺคปญฺห}
\csromanheader{5}{6. \textbf{นาวาลคฺคนกงฺคปญฺห}}
\proseitem{6}{ภนฺเต นาคเสน “นาวา\-ลคฺคน\-กสฺส เทฺว องฺคานิ คเหตพฺพานี”ติ ยํ วเทสิ, กตมานิ ตานิ เทฺว องฺคานิ คเหตพฺพานีติ. ยถา มหาราช นาวาลคฺคนกํ พหุอูมิ\-ชาลากุล\-วิกฺโขภิต\-สลิล\-ตเล มหติ\-มหาสมุทฺเท นาวํ ลคฺเคติ ฐเปติ, น เทติ ทิสาวิทิสํ หริตุํ. เอวเมว โข มหาราช โยคินา โยคาวจเรน ราคโทสโมคูมิชาเล มหติ\-มหาวิตกฺก\-สมฺปหาเร จิตฺตํ ลคฺเคตพฺพํ, น ทาตพฺพํ ทิสาวิทิสํ หริตุํ. อิทํ มหาราช นาวา\-ลคฺคน\-กสฺส ปฐมํ องฺคํ คเหตพฺพํ.}

\prose{ปุน จปรํ มหาราช นาวาลคฺคนกํ น ปฺลวติ\csromansharedfootnote{ee5b79700b6f635c}{น ปิลวติ (สี, อิ)} วิสีทติ, หตฺถสเตปิ อุทเก นาวํ ลคฺเคติ ฐานมุปเนติ. เอวเมว โข มหาราช โยคินา โยคาวจเรน ลาภ\-ยส\-สกฺการ\-มานน\-วนฺทนปูชน\-อปจิตีสุ ลาภคฺคยสคฺเคปิ น ปฺลวิตพฺพํ, สรีร\-ยาปน\-มตฺตเก เยว จิตฺตํ ฐเปตพฺพํ. อิทํ มหาราช นาวา\-ลคฺคน\-กสฺส ทุติยํ องฺคํ คเหตพฺพํ. ภาสิตมฺเปตํ มหาราช เถเรน สาริปุตฺเตน ธมฺม\-เสนาปตินา\csromandash{}}

\setlength{\csromangathapagewidth}{0pt}
\setlength{\csromangathawakwidth}{0pt}
\csromangathameasuregroup{“ยถา สมุทฺเท ลคฺคนกํ, \\ ตเถว ลาภสกฺกาเร,}{\csromangathabat{“ยถา สมุทฺเท ลคฺคนกํ,}{น ปฺลวติ วิสีทติ.} \\ \csromangathabat{ตเถว ลาภสกฺกาเร,}{มา ปฺลวถ วิสีทถา”ติ.}}
\csromangathasetleft{“ยถา สมุทฺเท ลคฺคนกํ, \\ ตเถว ลาภสกฺกาเร,}
\csromangathagroupwithcloserruled{\csromangathastanza{\csromanfolio{361}\csromangathabat{“ยถา สมุทฺเท ลคฺคนกํ,}{น ปฺลวติ วิสีทติ.} \\ \csromangathabat{ตเถว ลาภสกฺกาเร,}{มา ปฺลวถ วิสีทถา”ติ.}}}{นาวาลคฺคนกงฺคปโญฺห ฉฏฺโฐ.}
\csromanmatikaanchor{mtk.219}
\csromantocmarkref{subparagraph}{7. กูปงฺคปญฺห}{mtk.219}
\bookmark[dest={mtk.219},level=5]{7. กูปงฺคปญฺห}
\csromanheader{5}{7. \textbf{กูปงฺคปญฺห}}
\proseitemwithcloserruled{7}{ภนฺเต นาคเสน “กูปสฺส เอกํ องฺคํ คเหตพฺพนฺ”ติ ยํ วเทสิ, กตมํ ตํ เอกํ องฺคํ คเหตพฺพนฺติ. ยถา มหาราช กูโป รชฺชุญฺจ วรตฺตญฺจ ลงฺการญฺจ ธาเรติ. เอวเมว โข มหาราช โยคินา โยคาวจเรน สติสมฺปชญฺญ\-สมนฺนาคเตน ภวิตพฺพํ. อภิกฺกนฺเต ปฏิกฺกนฺเต อาโลกิเต วิโลกิเต สมิญฺชิเต ปสาริเต สงฺฆาฏิ\-ปตฺต\-จีวร\-ธารเณ อสิเต ปีเต ขายิเต สายิเต อุจฺจาร\-ปสฺสาว\-กมฺเม คเต ฐิเต นิสินฺเน สุตฺเต ชาคริเต ภาสิเต ตุณฺหีภาเว สมฺปชาน\-การินา ภวิตพฺพํ. อิทํ มหาราช กูปสฺส เอกํ องฺคํ คเหตพฺพํ. ภาสิตมฺเปตํ มหาราช ภควตา เทวาติเทเวน\csromansymbolfootnote{*}{สํ 3. 123 ปิฏฺเฐ.}“สโต ภิกฺขเว ภิกฺขุ วิหเรยฺย สมฺปชาโน, อยํ โว อมฺหากํ อนุสาสนี”ติ.}{กูปงฺคปโญฺห สตฺตโม.}
\csromanmatikaanchor{mtk.220}
\csromantocmarkref{subparagraph}{8. นิยามกงฺคปญฺห}{mtk.220}
\bookmark[dest={mtk.220},level=5]{8. นิยามกงฺคปญฺห}
\csromanheader{5}{8. \textbf{นิยามก}\-\textbf{งฺค}\-\textbf{ปญฺห}}
\proseitem{8}{ภนฺเต นาคเสน “นิยามกสฺส ตีณิ องฺคานิ คเหตพฺพานี”ติ ยํ วเทสิ, กตมานิ ตานิ ตีณิ องฺคานิ คเหตพฺพานีติ. ยถา มหาราช นิยามโก รตฺตินฺทิวํ สตตํ สมิตํ อปฺปมตฺโต ยตฺตปฺปยตฺโต นาวํ สาเรติ. เอวเมว โข มหาราช โยคินา โยคาวจเรน จิตฺตํ นิยามยมาเนน รตฺตินฺทิวํ สตตํ สมิตํ อปฺปมตฺเตน ยตฺต\-ปฺป\-ยตฺเตน โยนิโส มนสิกาเรน จิตฺตํ นิยาเมตพฺพํ. อิทํ มหาราช นิยามกสฺส ปฐมํ องฺคํ คเหตพฺพํ. ภาสิตมฺเปตํ มหาราช ภควตา เทวาติเทเวน ธมฺมปเท\csromandash{}}

\setlength{\csromangathapagewidth}{0pt}
\setlength{\csromangathawakwidth}{0pt}
\csromangathameasuregroup{\textsuperscript{+}“อปฺปมาทรตา โหถ, \\ ทุคฺคา อุทฺธรถ’ตฺตานํ,}{\csromangathabat{\textsuperscript{+}“อปฺปมาทรตา โหถ,}{สจิตฺตมนุรกฺขถ.} \\ \csromangathabat{ทุคฺคา อุทฺธรถ’ตฺตานํ,}{ปงฺเก สนฺโนว\textsuperscript{1} กุญฺชโร”ติ.}}
\csromangathasetleft{\textsuperscript{+}“อปฺปมาทรตา โหถ, \\ ทุคฺคา อุทฺธรถ’ตฺตานํ,}
\csromangathagroup{\csromangathastanza{\csromangathabat{\csromansymbolfootnote{+}{ขุ 1. 60 ปิฏฺเฐ.}“อปฺปมาทรตา โหถ,}{สจิตฺตม\-นุรกฺขถ.} \\ \csromangathabat{ทุคฺคา อุทฺธรถ’ตฺตานํ,}{ปงฺเก สนฺโนว\csromansharedfootnote{60adf04150c5c155}{สตฺโตว (สี, อิ)} กุญฺชโร”ติ.}}}

\prose{\csromanfolio{362}ปุน จปรํ มหาราช นิยามกสฺส ยํ กิญฺจิ มหาสมุทฺเท กลฺยาณํ วา ปาปกํ วา, สพฺพํ ตํ วิทิตํ โหติ. เอวเมว โข มหาราช โยคินา โยคาวจเรน กุสลากุสลํ สาวชฺชา\-นวชฺชํ หีนปฺปณีตํ กณฺห\-สุกฺก\-สปฺปฏิภาคํ วิชานิตพฺพํ. อิทํ มหาราช นิยามกสฺส ทุติยํ องฺคํ คเหตพฺพํ.}

\prosewithcloserruled{ปุน จปรํ มหาราช นิยามโก ยนฺเต มุทฺทิกํ เทติ “มา โกจิ ยนฺตํ อามสิตฺถา”ติ. เอวเมว โข มหาราช โยคินา โยคาวจเรน จิตฺเต สํวรมุทฺทิกา ทาตพฺพา “มา กิญฺจิ ปาปกํ อกุสล\-วิตกฺกํ วิตกฺเกสี”ติ, อิทํ มหาราช นิยามกสฺส ตติยํ องฺคํ คเหตพฺพํ. ภาสิตมฺเปตํ มหาราช ภควตา เทวาติเทเวน สํยุตฺตนิกาย\-วเร\csromansymbolfootnote{*}{สํ 3. 366 ปิฏฺเฐ.}“มา ภิกฺขเว ปาปเก อกุสเล วิตกฺเก วิตกฺเกยฺยาถ. เสยฺยถีทํ, กามวิตกฺกํ พฺยาปาทตกฺกํ วิหิํสาวิตกฺกนฺ”ติ.}{นิยามก\-งฺค\-ปโญฺห อฏฺฐโม.}
\csromanmatikaanchor{mtk.221}
\csromantocmarkref{subparagraph}{9. กมฺมการงฺคปญฺห}{mtk.221}
\bookmark[dest={mtk.221},level=5]{9. กมฺมการงฺคปญฺห}
\csromanheader{5}{9. \textbf{กมฺมกา}\-\textbf{รงฺค}\-\textbf{ปญฺห}}
\proseitem{9}{ภนฺเต นาคเสน “กมฺมการสฺส เอกํ องฺคํ คเหตพฺพนฺ”ติ ยํ วเทสิ, กตมํ ตํ เอกํ องฺคํ คเหตพฺพนฺติ. ยถา มหาราช กมฺมกาโร เอวํ จินฺตยติ “ภตโก อหํ อิมาย นาวาย กมฺมํ กโรมิ, อิมายาหํ นาวาย วาหสา ภตฺตเวตนํ ลภามิ, น เม ปมาโท กรณีโย, อปฺปมาเทน เม อยํ นาวา วาเหตพฺพา”ติ, เอวเมว โข มหาราช โยคินา โยคาวจเรนเอวํจินฺตยิตพฺพํ “อิมํ โข อหํ จาตุมหาภูติกํ กายํ สมฺมสนฺโต สตตํ สมิตํ อปฺปมตฺโต อุปฏฺฐิตสฺสติ สโต สมฺปชาโน สมาหิโต เอกคฺคจิตฺโต ชาติชราพฺยาธิมรณ\-โสก\-ปริเทว\-ทุกฺข\-โทมนสฺสุปายาเสหิ ปริมุจฺจิสฺสามีติ อปฺปมาโท เม กรณีโย”ติ, อิทํ มหาราช กมฺมการสฺส เอกํ องฺคํ คเหตพฺพํ. ภาสิตมฺเปตํ มหาราช เถเรน สาริปุตฺเตน ธมฺม\-เสนาปตินา\csromandash{}}

\setlength{\csromangathapagewidth}{0pt}
\setlength{\csromangathawakwidth}{0pt}
\csromangathameasuregroup{“กายํ อิมํ สมฺมสถ, \\ กาเย สภาวํ ทิสฺวาน,}{\csromangathabat{“กายํ อิมํ สมฺมสถ,}{ปริชานาถ ปุนปฺปุนํ.} \\ \csromangathabat{กาเย สภาวํ ทิสฺวาน,}{ทุกฺขสฺสนฺตํ กริสฺสถา”ติ.}}
\csromangathasetleft{“กายํ อิมํ สมฺมสถ, \\ กาเย สภาวํ ทิสฺวาน,}
\csromangathagroup{\csromangathastanza{\csromangathabat{“กายํ อิมํ สมฺมสถ,}{ปริชานาถ ปุนปฺปุนํ.} \\ \csromangathabat{กาเย สภาวํ ทิสฺวาน,}{ทุกฺขสฺสนฺตํ กริสฺสถา”ติ.}}}

\vspace{\tipitakaparskip}
\csromancenter{กมฺมกา\-รงฺค\-ปโญฺห นวโม.}
\csromanmatikaanchor{mtk.222}
\csromantocmarkref{subparagraph}{10. สมุทฺทงฺคปญฺห}{mtk.222}
\bookmark[dest={mtk.222},level=5]{10. สมุทฺทงฺคปญฺห}
\csromanmarkh{6. โอปมฺม\-กถา\-ปญฺห}\csromanmarkhii{10. สมุทฺท\-งฺค\-ปญฺห}\csromanheadernomark{5}{6. โอปมฺม\-กถา\-ปญฺห \textbf{10. สมุทฺท}\-\textbf{งฺค}\-\textbf{ปญฺห}}
\proseitem{10}{\csromanfolio{363}ภนฺเต นาคเสน “สมุทฺทสฺส ปญฺจ องฺคานิ คเหตพฺพานี”ติ ยํ วเทสิ, กตมานิ ตานิ ปญฺจ องฺคานิ คเหตพฺพานีติ. ยถา มหาราช มหาสมุทฺโท มเตน กุณเปน สทฺธิํ น สํวสติ, เอวเมว โข มหาราช โยคินา โยคาวจเรน ราคโทสโมหมานทิฏฺฐิมกฺขปฬาส- อิสฺสามจฺฉริยมายาสาเฐยฺยกุฏิลวิสมทุจฺจริตกิเลสมเลหิ สทฺธิํ น สํวสิตพฺพํ, อิทํ มหาราช สมุทฺทสฺส ปฐมํ องฺคํ คเหตพฺพํ.}

\prose{ปุน จปรํ มหาราช มหาสมุทฺโท มุตฺตา\-มณิ\-เวฬุริย\-สงฺข\-สิลา\-ปวาฬ\-ผลิก\-มณิ\-วิวิธ\-รตน\-นิจยํ ธาเรนฺโต ปิทหติ, น พหิ วิกิรติ, เอวเมว โข มหาราช โยคินา โยคาวจเรน มคฺค\-ผล\-ฌาน\-วิโมกฺข\-สมาธิ\-สมาปตฺติ\-วิปสฺสนา\-ภิญฺญา\-วิวิธ\-คุณรตฺ อนานิ อธิคนฺตฺวา ปทหิตพฺพานิ, น พหิ นีหริตพฺพานิ, อิทํ มหาราช สมุทฺทสฺส ทุติยํ องฺคํ คเหตพฺพํ.}

\prose{ปุน จปรํ มหาราช มหาสมุทฺโท มหนฺเตหิ มหาภูเตหิ สทฺธิํ สํวสติ, เอวเมว โข มหาราช โยคินา โยคาวจเรน อปฺปิจฺฉํ สนฺตุฏฺฐํ ธุตวาทํ สลฺเลขวุตฺติํ อาจารสมฺปนฺนํ ลชฺชิํ เปสลํ ครุํ ภาวนียํ วตฺตารํ วจนกฺขมํ โจทกํ ปาปครหิํ โอวาทกํ อนุสาสกํ วิญฺญาปกํ สนฺทสฺสกํ สมาทปกํ สมุตฺเตชกํ สมฺปหํสกํ กลฺยาณมิตฺตํ สพฺรหฺมจาริํ นิสฺสาย วสิตพฺพํ, อิทํ มหาราช มหาสมุทฺทสฺส ตติยํ องฺคํ คเหตพฺพํ.}

\prose{ปุน จปรํ มหาราช มหาสมุทฺโท นวสลิล\-สมฺปุณฺณาหิ คงฺคา\-ยมุนา\-อจิรวตี\-สรภู\-มหี\-อาทีหิ นทีสตสหสฺเสหิ อนฺตลิกฺเข สลิลธาราหิ จ ปูริโตปิ สกํ เวลํ นาติวตฺตติ, เอวเมว โข มหาราช โยคินา โยคาวจเรน ลาภสกฺการสิโลก\-วนฺทน\-มานน\-ปูชน\-การณา ชีวิตเหตุปิ สญฺจิจฺจ สิกฺขาปท\-วีติกฺกโม น กรณิโย, อิทํ มหาราช มหาสมุทฺทสฺส จตุตฺถํ องฺคํ คเหตพฺพํ. ภาสิตมฺเปตํ มหาราช ภควตา เทวาติเทเวน\csromandash{}}

\prose{\csromansymbolfootnote{*}{วิ 4. 421; อํ 3. 41; ขุ 1. 142 ปิฏฺเฐสุ.}“เสยฺยถาปิ มหาราช\csromansharedfootnote{534a6ef0db708649}{ปหาราธ (สี, อิ)} มหาสมุทฺโท ฐิตธมฺโม เวลํ นาติกฺกมติ เอวเมว โข มหาราช ยํ มยา สาวกานํ สิกฺขาปทํ ปญฺญตฺตํ, ตํ มม สาวกา ชีวิตเหตุปิ นาติกฺกมนฺตี”ติ.}

\prose{\csromanfolio{364}ปุน จปรํ มหาราช มหาสมุทฺโท สพฺพสวนฺตีหิ คงฺคายมุนา- อจิรวตีสรภูมหีหิ อนฺตลิกฺเข อุทกราหิปิ น ปริปูรติ, เอวเมว โข มหาราช โยคินา โยคาวจเรน อุทฺเทส\-ปริปุจฺฉา\-สวน\-ธารณ\-วินิจฺฉย- อภิธมฺมวินยคาฬฺหสุตฺตนฺตวิคฺคหปทนิกฺเขปปทสนฺธิ ปทวิภตฺติ นวงฺคชินสาสนวรํ สุณนฺเตนาปิ น ตปฺปิตพฺพํ, อิทํ มหาราช มหาสมุทฺทสฺส ปญฺจมํ องฺคํ คเหตพฺพํ. ภาสิตมฺเปตํ มหาราช ภควตา เทวาติเทเวน สุตโสมชาตเก\csromandash{}}

\setlength{\csromangathapagewidth}{0pt}
\setlength{\csromangathawakwidth}{0pt}
\csromangathameasuregroup{}{\textsuperscript{*}“อคฺคิ ยถา ติณกฏฺฐํ ทหนฺโต, \\ น ตปฺปติ สาคโร วา นทีหิ. \\ เอวมฺปิ เจ1 ปณฺฑิตา ราชเสฏฺฐ, \\ สุตฺวา น ตปฺปนฺติ สุภาสิเตนา”ติ.}
\csromangathameasurewak{\textsuperscript{*}“อคฺคิ ยถา ติณกฏฺฐํ ทหนฺโต, \\ น ตปฺปติ สาคโร วา นทีหิ. \\ เอวมฺปิ เจ1 ปณฺฑิตา ราชเสฏฺฐ, \\ สุตฺวา น ตปฺปนฺติ สุภาสิเตนา”ติ.}
\setlength{\csromangathaleftcol}{0pt}
\csromangathagroup{\csromangathastanza{\csromangathawak{\csromansymbolfootnote{*}{1. เอวมฺปิ เว (สฺยา), เอวํ หิ เม (ก) ขุ 6. 137 ปิฏฺเฐ.}“อคฺคิ ยถา ติณกฏฺฐํ ทหนฺโต, \\ น ตปฺปติ สาคโร วา นทีหิ. \\ เอวมฺปิ เจ1 ปณฺฑิตา ราชเสฏฺฐ, \\ สุตฺวา น ตปฺปนฺติ สุภาสิเตนา”ติ.}}}

\vspace{\tipitakaparskip}
\csromancenter{สมุทฺท\-งฺค\-ปโญฺห ทสโม.}
\nitthitammidruled{\textbf{สมุทฺทวคฺโค ทุติโย.}}
\csromancenter{\textbf{ตสฺสุทฺทานํ}}
\setlength{\csromangathapagewidth}{0pt}
\setlength{\csromangathawakwidth}{0pt}
\csromangathameasuregroup{ลาพุลตา จ ปทุมํ, \\ นาวา จ นาวาลคฺคนํ, \\ กมฺมกาโร สมุทฺโท จ,}{\csromangathabat{ลาพุลตา จ ปทุมํ,}{พีชํ สาลกลฺยาณิกา.} \\ \csromangathabat{นาวา จ นาวาลคฺคนํ,}{กูโป นิยามโก ตถา.} \\ \csromangathabat{กมฺมกาโร สมุทฺโท จ,}{วคฺโค เตน ปวุจฺจตีติ.}}
\csromangathasetleft{ลาพุลตา จ ปทุมํ, \\ นาวา จ นาวาลคฺคนํ, \\ กมฺมกาโร สมุทฺโท จ,}
\csromangathagroup{\csromangathastanza{\csromangathabat{ลาพุลตา จ ปทุมํ,}{พีชํ สาลกลฺยาณิกา.} \\ \csromangathabat{นาวา จ นาวาลคฺคนํ,}{กูโป นิยามโก ตถา.} \\ \csromangathabat{กมฺมกาโร สมุทฺโท จ,}{วคฺโค เตน ปวุจฺจตีติ.}}}

\csromanmatikaanchor{mtk.223}
\csromantocmarknopage{paragraph}{3. ปถวีวคฺค}{mtk.223}
\bookmark[dest={mtk.223},level=4]{3. ปถวีวคฺค}
\csromanheader{4}{3. \textbf{ปถวีวคฺค}}
\csromanmatikaanchor{mtk.224}
\csromantocmarkref{subparagraph}{1. ปถวีองฺคปญฺห}{mtk.224}
\bookmark[dest={mtk.224},level=5]{1. ปถวีองฺคปญฺห}
\csromanheader{5}{1. ปถวี\-องฺค\-ปญฺห}
\proseitem{1}{\csromanfolio{365}ภนฺเต นาคเสน “ปถวิยา ปญฺจ องฺคานิ คเหตพฺพานี”ติ ยํ วเทสิ, กตมานิ ตานิ ปญฺจ องฺคานิ คเหตพฺพานีติ. ยถา มหาราช ปถวี อิฏฺฐานิฏฺฐานิ กปฺปูรา\-ครุ\-ตครจนฺทน\-กุงฺกุมาทีนิ อากิรนฺเตปิ ปิตฺต\-เสมฺห\-ปุพฺพ\-รุหิร\-เสท\-เมท\-เขฬสิงฺฆาณิก\-ลสิก\-มุตฺตกรีสาทีนิ อากิรนฺเตปิ ตาทิสา เยว, เอวเมว โข มหาราช โยคินา โยคาวจเรน อิฏฺฐานิฏฺเฐ ลาภาลาเภ ยสายเส นินฺทาปสํสาย สุขทุกฺเข สพฺพตฺถ ตาทินา เยว ภวิตพฺพํ, อิทํ มหาราช ปถวิยา ปฐมํ องฺคํ คเหตพฺพํ.}

\prose{ปุน จปรํ มหาราช ปถวี มณฺฑน\-วิภูส\-นาปคตา สกคนฺธปริภาวิตา, เอวเมว โข มหาราช โยคินา โยคาวจเรน วิภูส\-นาปคเตน สกสีล\-คนฺธปริภาวิเตน ภวิตพฺพํ, อิทํ มหาราช ปถวิยา ทุติยํ องฺคํ คเหตพฺพํ.}

\prose{ปุน จปรํ มหาราช ปถวี นิรนฺตรา อขณฺฑจฺฉิทฺทา อสุสิรา พหลา ฆนา วิตฺถิณฺณา, เอวเมว โข มหาราช โยคินา โยคาวจเรน นิรนฺต\-รม\-ขณฺฑ\-จฺฉิทฺทม\-สุสิร\-พหล\-ฆน\-วิตฺถิณฺณ\-สีเลน ภวิตพฺพํ, อิทํ มหาราช ปถวิยา ตติยํ องฺคํ คเหตพฺพํ.}

\prose{ปุน จปรํ มหาราช ปถวี คามนิคมนครชนปทรุกฺขปพฺพตน- ทีตฬากโปกฺขรณีมิคปกฺขิมนุชนรนาริคณํ ธาเรนฺตีปิ อกิลาสุ โหติ, เอวเมว โข มหาราช โยคินา โยคาวจเรน โอวทนฺเตนปิ อนุสาสนฺเตนปิ วิญฺญาเปนฺเตนปิ สนฺทสฺเส\-นฺเตนปิ สมา\-ทเป\-นฺเตนปิ สมุตฺเตเชนฺเตนปิ สมฺป\-หํเส\-นฺเตนปิ ธมฺมเทสนาสุ อกิลาสุนา ภวิตพฺพํ, อิทํ มหาราช ปถวิยา จตุตฺถํ องฺคํ คเหตพฺพํ.}

\prose{ปุน จปรํ มหาราช ปถวี อนุนย\-ปฺปฏิฆ\-วิปฺปมุตฺตา, เอวเมว โข มหาราช โยคินา โยคาวจเรน อนุนย\-ปฺปฏิฆ\-วิปฺปมุตฺเตน ปถวิสเมน เจตสา วิหริตพฺพํ, อิทํ มหาราช ปถวิยา ปญฺจมํ องฺคํ คเหตพฺพํ. ภาสิตมฺเปตํ มหาราช อุปาสิกาย จูฬสุภทฺทาย สกสมเณ ปริกิตฺตยมานาย\csromandash{}}

\setlength{\csromangathapagewidth}{0pt}
\setlength{\csromangathawakwidth}{0pt}
\csromangathameasuregroup{“เอกญฺเจ พาหํ วาสิยา, \\ เอกญฺเจ พาหํ คนฺเธน, \\ อมุสฺมิํ ปฏิโฆ นตฺถิ, \\ ปถวีสมจิตฺตา เต,}{\csromangathabat{“เอกญฺเจ พาหํ วาสิยา,}{ตจฺเฉ กุปิตมานสา\textsuperscript{1}.} \\ \csromangathabat{เอกญฺเจ พาหํ คนฺเธน,}{อาลิมฺเปยฺย ปโมทิตา\textsuperscript{1}.} \\ \csromangathabat{อมุสฺมิํ ปฏิโฆ นตฺถิ,}{ราโค อสฺมิํ น วิชฺชติ.} \\ \csromangathabat{ปถวีสมจิตฺตา เต,}{ตาทิสา สมณา มมา”ติ.}}
\csromangathasetleft{“เอกญฺเจ พาหํ วาสิยา, \\ เอกญฺเจ พาหํ คนฺเธน, \\ อมุสฺมิํ ปฏิโฆ นตฺถิ, \\ ปถวีสมจิตฺตา เต,}
\csromangathagroup{\csromangathastanza{\csromanfolio{366}\csromangathabat{“เอกญฺเจ พาหํ วาสิยา,}{ตจฺเฉ กุปิตมานสา\csromansharedfootnote{a55ab0345b0f9f6f}{กุปิตมานโส (ก)}.} \\ \csromangathabat{เอกญฺเจ พาหํ คนฺเธน,}{อาลิมฺเปยฺย ปโมทิตา\csromansharedfootnote{a55ab0345b0f9f6f}{กุปิตมานโส (ก)}.}}\csromangathastanza{\csromangathabat{อมุสฺมิํ ปฏิโฆ นตฺถิ,}{ราโค อสฺมิํ น วิชฺชติ.} \\ \csromangathabat{ปถวีสมจิตฺตา เต,}{ตาทิสา สมณา มมา”ติ.}}}

\prose{ปถวี\-องฺค\-ปโญฺห ปฐโม.}
\csromansectionrule

\csromanmatikaanchor{mtk.225}
\csromantocmarkref{subparagraph}{2. อาปงฺคปญฺห}{mtk.225}
\bookmark[dest={mtk.225},level=5]{2. อาปงฺคปญฺห}
\csromanheader{5}{2. \textbf{อาปงฺคปญฺห}}
\proseitem{2}{ภนฺเต นาคเสน “อาปสฺส ปญฺจ องฺคานิ คเหตพฺพานี”ติ ยํ วเทสิ, กตมานิ ตานิ ปญฺจ องฺคานิ คเหตพฺพานีติ. ยถา มหาราช อาโป สุสณฺฐิตมกมฺปิตมลุฬิตสภาวปริสุทฺโธ, เอวเมว โข มหาราช โยคินา โยคาวจเรน กุหนลปนเนมิตฺตก นิปฺเปสิกตํ อปเนตฺวา สุสณฺฐิตมกมฺปิตมลุฬิตสภาวปริสุทฺธาจาเรน ภวิตพฺพํ, อิทํ มหาราช อาปสฺส ปฐมํ องฺคํ คเหตพฺพํ.}

\prose{ปุน จปรํ มหาราช อาโป สีตล\-สภาว\-สณฺฐิโต, เอวเมว โข มหาราช โยคินา โยคาวจเรน สพฺพสตฺเตสุ ขนฺติเมตฺตา\-นุทฺทย\-สมฺปนฺเนน หิเตสินา อนุกมฺเปเกน ภวิตพฺพํ, อิทํ มหาราช อาปสฺส ทุติยํ องฺคํ คเหตพฺพํ.}

\prose{ปุน จปรํ มหาราช อาโป อสุจิํ สุจิํ กโรติ, เอวเมว โข มหาราช โยคินา โยคาวจเรน คาเม วา อรญฺเญ วา อุปชฺฌาเย อุปชฺฌายมตฺเตสุ อาจริเย อาจริยมตฺเตสุ สพฺพตฺถ อนธิกรเณน ภวิตพฺพํ อนวเสส\-การินา, อิทํ มหาราช อาปสฺส ตติยํ องฺคํ คเหตพฺพํ.}

\prose{ปุน จปรํ มหาราช อาโป พหุชน\-ปตฺถิโต, เอวเมว โข มหาราช โยคินา โยคาวจเรน อปฺปิจฺฉ\-สนฺตุฏฺฐ\-ปวิวิตฺต\-ปฏิสลฺลาเนน สตตํ สพฺพโลกม\-ภิปตฺถิเตน ภวิตพฺพํ, อิทํ มหาราช อาปสฺส จตุตฺถํ องฺคํ คเหตพฺพํ.}

\prose{\csromanfolio{367}ปุน จปรํ มหาราช อาโป น กสฺสจิ อหิตมุ\-ปทหติ, เอวเมว โข มหาราช โยคินา โยคาวจเรน ปรภณฺฑนกลหวิคฺคหวิวาท\-ริตฺต\-ชฺฌาน\-อรติ\-ชนนํ กาย\-วจี\-จิตฺเตหิ ปาปกํ น กรณียํ, อิทํ มหาราช อาปสฺส ปญฺจมํ องฺคํ คเหตพฺพํ. ภาสิตมฺเปตํ มหาราช ภควตา เทวาติเทเวน กณฺหชาตเก\csromandash{}}

\setlength{\csromangathapagewidth}{0pt}
\setlength{\csromangathawakwidth}{0pt}
\csromangathameasuregroup{\textsuperscript{*}“วรญฺเจ เม อโท สกฺก, \\ น มโน วา สรีรํ วา, \\ กทาจิ อุปหญฺเญถ,}{\csromangathabat{\textsuperscript{*}“วรญฺเจ เม อโท สกฺก,}{สพฺพภูตานมิสฺสร.} \\ \csromangathabat{น มโน วา สรีรํ วา,}{มํกเต สกฺก กสฺสจิ.} \\ \csromangathabat{กทาจิ อุปหญฺเญถ,}{เอตํ สกฺก วรํ วเร”ติ.}}
\csromangathasetleft{\textsuperscript{*}“วรญฺเจ เม อโท สกฺก, \\ น มโน วา สรีรํ วา, \\ กทาจิ อุปหญฺเญถ,}
\csromangathagroupwithcloserruled{\csromangathastanza{\csromangathabat{\csromansymbolfootnote{*}{ขุ 5. 203 ปิฏฺเฐ.}“วรญฺเจ เม อโท สกฺก,}{สพฺพภูตานมิ\-สฺสร.} \\ \csromangathabat{น มโน วา สรีรํ วา,}{มํกเต สกฺก กสฺสจิ.} \\ \csromangathabat{กทาจิ อุปหญฺเญถ,}{เอตํ สกฺก วรํ วเร”ติ.}}}{อาปงฺคปโญฺห ทุติโย.}
\csromanmatikaanchor{mtk.226}
\csromantocmarkref{subparagraph}{3. เตชงฺคปญฺห}{mtk.226}
\bookmark[dest={mtk.226},level=5]{3. เตชงฺคปญฺห}
\csromanheader{5}{3. \textbf{เตชงฺคปญฺห}}
\proseitem{3}{ภนฺเต นาคเสน “เตชสฺส ปญฺจ องฺคานิ คเหตพฺพานี”ติ ยํ วเทสิ, กตมานิ ตานิ ปญฺจ องฺคานิ คเหตพฺพานีติ. ยถา มหาราช เตโช ติณกฏฺฐ\-สาขา\-ปลาสํ ฑหติ, เอวเมว โข มหาราช โยคินา โยคาวจเรน เย เต อพฺภนฺตรา วา พาหิรา วา กิเลสา อิฏฺฐา\-นิฏฺฐารมฺมณา\-นุภวนา, สพฺเพ เต ญาณคฺคินา ฑหิตพฺพา, อิทํ มหาราช เตชสฺส ปฐมํ องฺคํ คเหตพฺพํ.}

\prose{ปุน จปรํ มหาราช เตโช นิทฺทโย อการุณิโก, เอวเมว โข มหาราช โยคินา โยคาวจเรน สพฺพกิเลเสสุ การุญฺญานุทฺทยา น กาตพฺพา, อิทํ มหาราช เตชสฺส ทุติยํ องฺคํ คเหตพฺพํ.}

\prose{ปุน จปรํ มหาราช เตโช สีตํ ปฏิหนติ, เอวเมว โข มหาราช โยคินา โยคาวจเรน วีริย\-สนฺตาป\-เตชํ อภิชเนตฺวา กิเลสา ปฏิหนฺตพฺพา, อิทํ มหาราช เตชสฺส ตติยํ องฺคํ คเหตพฺพํ.}

\prose{ปุน จปรํ มหาราช เตโช อนุนย\-ปฺปฏิฆ\-วิปฺปมุตฺโต อุณฺหมภิชเนติ, เอวเมว โข มหาราช โยคินา โยคาวจเรน อนุนย\-ปฺปฏิฆ\-วิปฺปมุตฺเตน เตโชสเมน เจตสา วิหริตพฺพํ, อิทํ มหาราช เตชสฺส จุตุตฺถํ องฺคํ คเหตพฺพํ.}

\prosewithcloserruled{\csromanfolio{368}ปุน จปรํ มหาราช เตโช อนฺธการํ วิธมิตฺวา\csromansharedfootnote{61d7c8d0f08facb6}{วิธมติ (สี, อิ, ก)} อาโลกํ ทสฺสยติ, เอวเมว โข มหาราช โยคินา โยคาวจเรน อวิชฺช\-นฺธการํ วิธมิตฺวา ญาณาโลกํ ทสฺสยิตพฺพํ, อิทํ มหาราช เตชสฺส ปญฺจมํ องฺคํ คเหตพฺพํ. ภาสิตมฺเปตํ มหาราช ภควตา เทวาติเทเวน สกํ ปุตฺตํ ราหุลํ โอวทนฺเตน\csromandash{} “เตโชสมํ\csromansharedfootnote{0bf288a524a99bab}{ม 2. 87 ปิฏฺเฐ.} ราหุล ภาวนํ ภาเวหิ, เตโชสมํ หิ เต ราหุล ภาวนํ ภาวยโต อุปฺปนฺนา มนาปามนาปา ผสฺสา จิตฺตํ น ปริยาทาย ฐสฺสนฺตี”ติ.}{เตชงฺคปโญฺห ตติโย.}
\csromanmatikaanchor{mtk.227}
\csromantocmarkref{subparagraph}{4. วายุงฺคปญฺห}{mtk.227}
\bookmark[dest={mtk.227},level=5]{4. วายุงฺคปญฺห}
\csromanheader{5}{4. \textbf{วายุงฺคปญฺห}}
\proseitem{4}{ภนฺเต นาคเสน “วายุสฺส ปญฺจ องฺคานิ คเหตพฺพานี”ติ ยํ วเทสิ, กตมานิ ตานิ ปญฺจ องฺคานิ คเหตพฺพานีติ. ยถา มหาราช วายุ สุปุปฺผิต\-วนสณฺฑ\-นฺตรํ อภิวายติ, เอวเมว โข มหาราช โยคินา โยคาวจเรน วิมุตฺติ\-วร\-กุสุม\-ปุปฺผิตา\-รมฺมณ\-วนนฺตเร รมิตพฺพํ, อิทํ มหาราช วายุสฺส ปฐมํ องฺคํ คเหตพฺพํ.}

\prose{ปุน จปรํ มหาราช วายุ ธรณีรุห\-ปาทป\-คเณ มถยติ, เอวเมว โข มหาราช โยคินา โยคาวจเรน วนนฺตร\-คเตน สงฺขาเร วิจินนฺเตน กิเลสา มถยิตพฺพา. อิทํ มหาราช วายุสฺส ทุติยํ องฺคํ คเหตพฺพํ.}

\prose{ปุน จปรํ มหาราช วายุ อากาเส จรติ, เอวเมว โข มหาราช โยคินา โยคาวจเรน โลกุตฺตร\-ธมฺเมสุ มานสํ สญฺจาร\-ยิตพฺพํ, อิทํ มหาราช วายุสฺส ตติยํ องฺคํ คเหตพฺพํ.}

\prose{ปุน จปรํ มหาราช วายุ คนฺธํ อนุภวติ, เอวเมว โข มหาราช โยคินา โยคาวจเรน อตฺตโน สีลวร\-สุรภิคนฺโธ\csromansharedfootnote{973d4c379a3df9a4}{สีลสุรภิคนฺโธ (สี, อิ)} อนุภวิตพฺโพ, อิทํ มหาราช วายุสฺส จตุตฺถํ องฺคํ คเหตพฺพํ.}

\prose{ปุน จปรํ มหาราช วายุ นิราลโย อนิเกตวาสี, เอวเมว โข มหาราช โยคินา โยคาวจเรน นิราลยม\-นิเกตม\-สนฺถเวน \csromanfolio{369}สพฺพตฺถ วิมุตฺเตน ภวิตพฺพํ, อิทํ มหาราช วายุสฺส ปญฺจมํ องฺคํ คเหตพฺพํ. ภาสิตมฺเปตํ มหาราช ภควตา เทวาติเทเวน สุตฺตนิปาเต\csromandash{}}

\setlength{\csromangathapagewidth}{0pt}
\setlength{\csromangathawakwidth}{0pt}
\csromangathameasuregroup{\textsuperscript{*}“สนฺถวาโต ภยํ ชาตํ, \\ อนิเกตมสนฺถวํ,}{\csromangathabat{\textsuperscript{*}“สนฺถวาโต ภยํ ชาตํ,}{นิเกตา ชายเต รโช.} \\ \csromangathabat{อนิเกตมสนฺถวํ,}{เอตํ เว มุนิทสฺสนนฺ”ติ.}}
\csromangathasetleft{\textsuperscript{*}“สนฺถวาโต ภยํ ชาตํ, \\ อนิเกตมสนฺถวํ,}
\csromangathagroupwithcloserruled{\csromangathastanza{\csromangathabat{\csromansymbolfootnote{*}{ขุ 1. 309 ปิฏฺเฐ.}“สนฺถวาโต ภยํ ชาตํ,}{นิเกตา ชายเต รโช.} \\ \csromangathabat{อนิเกตม\-สนฺถวํ,}{เอตํ เว มุนิทสฺสนนฺ”ติ.}}}{วายุงฺคปโญฺห จตุตฺโถ.}
\csromanmatikaanchor{mtk.228}
\csromantocmarkref{subparagraph}{5. ปพฺพตงฺคปญฺห}{mtk.228}
\bookmark[dest={mtk.228},level=5]{5. ปพฺพตงฺคปญฺห}
\csromanheader{5}{5. \textbf{ปพฺพต}\-\textbf{งฺค}\-\textbf{ปญฺห}}
\proseitem{4}{ภนฺเต นาคเสน “ปพฺพตสฺส ปญฺจ องฺคานิ คเหตพฺพานี”ติ ยํ วเทสิ, กตมานิ ตานิ ปญฺจ องฺคานิ คเหตพฺพานีติ. ยถา มหาราช ปพฺพโต อจโล อกมฺปิโต\csromansharedfootnote{b57833e614ae5de3}{อกมฺปิโย (สี, อิ)} อสมฺปเวธี, เอวเมว โข มหาราช โยคินา โยคาวจเรน สมฺมานเน วิมานเน สกฺกาเร อสกฺกาเร ครุกาเร อครุกาเร ยเส อยเส นินฺทาย ปสํสาย สุเข ทุกฺเข อิฏฺฐานิฏฺเฐสุ สพฺพตฺถ รูปสทฺทคนฺธรสโผฏฺฐพฺพ\-ธมฺเมสุ รชนีเยสุ น รชฺชิตพฺพํ, ทุสฺสนีเยสุ น ทุสฺสิตพฺพํ, มุยฺหนีเยสุ น มุยฺหิตพฺพํ, น กมฺปิตพฺพํ น จลิตพฺพํ, ปพฺพเตน วิย อจเลน ภวิตพฺพํ, อิทํ มหาราช ปพฺพตสฺส ปฐมํ องฺคํ คเหตพฺพํ. ภาสิตมฺเปตํ มหาราช ภควตา เทวาติเทเวน\csromandash{}}

\setlength{\csromangathapagewidth}{0pt}
\setlength{\csromangathawakwidth}{0pt}
\csromangathameasuregroup{“เสโล ยถา เอกฆโน\textsuperscript{1}, \\ เอวํ นินฺทาปสํสาสุ,}{\csromangathabat{“เสโล ยถา เอกฆโน\textsuperscript{1},}{วาเตน น สมีรติ.} \\ \csromangathabat{เอวํ นินฺทาปสํสาสุ,}{น สมิญฺชนฺติ ปณฺฑิตา”ติ.}}
\csromangathasetleft{“เสโล ยถา เอกฆโน\textsuperscript{1}, \\ เอวํ นินฺทาปสํสาสุ,}
\csromangathagroup{\csromangathastanza{\csromangathabat{“เสโล ยถา เอกฆโน\csromansharedfootnote{32bd57b394d48588}{เอกคฺฆโน (ก) ขุ 1. 25 ปิฏฺเฐ ธมฺมปเท.},}{วาเตน น สมีรติ.} \\ \csromangathabat{เอวํ นินฺทาปสํสาสุ,}{น สมิญฺชนฺติ ปณฺฑิตา”ติ.}}}

\prose{ปุน จปรํ มหาราช ปพฺพโต ถทฺโธ น เกนจิ สํสฏฺโฐ, เอวเมว โข มหาราช โยคินา โยคาวจเรน ถทฺเธน อสํสฏฺเฐน ภวิตพฺพํ, น เกนจิ สํสคฺโค กรณีโย, อิทํ มหาราช ปพฺพตสฺส ทุติยํ องฺคํ คเหตพฺพํ. ภาสิตมฺเปตํ มหาราช ภควตา เทวาติเทเวน\csromandash{}}

\setlength{\csromangathapagewidth}{0pt}
\setlength{\csromangathawakwidth}{0pt}
\csromangathameasuregroup{\textsuperscript{+}“อสํสฏฺฐํ คหฏฺเฐหิ, \\ อโนกสาริมปฺปิจฺฉํ,}{\csromangathabat{\textsuperscript{+}“อสํสฏฺฐํ คหฏฺเฐหิ,}{อนาคาเรหิ จูภยํ.} \\ \csromangathabat{อโนกสาริมปฺปิจฺฉํ,}{ตมหํ พฺรูมิ พฺราหฺมณนฺ”ติ.}}
\csromangathasetleft{\textsuperscript{+}“อสํสฏฺฐํ คหฏฺเฐหิ, \\ อโนกสาริมปฺปิจฺฉํ,}
\csromangathagroup{\csromangathastanza{\csromangathabat{\csromansymbolfootnote{+}{ขุ 1. 71, 376 ปิฏฺเฐสุ.}“อสํสฏฺฐํ คหฏฺเฐหิ,}{อนาคาเรหิ จูภยํ.} \\ \csromangathabat{อโนกสาริม\-ปฺปิจฺฉํ,}{ตมหํ พฺรูมิ พฺราหฺมณนฺ”ติ.}}}

\prose{ปุน จปรํ มหาราช ปพฺพเต พีชํ น วิรูหติ, เอวเมว โข มหาราช โยคินา โยคาวจเรน สกมานเส กิเลสา น วิรูหาเปตพฺพา, อิทํ มหาราช ปพฺพตสฺส ตติยํ องฺคํ คเหตพฺพํ. ภาสิตมฺเปตํ มหาราช เถเรน สุภูตินา\csromandash{}}

\setlength{\csromangathapagewidth}{0pt}
\setlength{\csromangathawakwidth}{0pt}
\csromangathameasuregroup{\textsuperscript{*}“ราคูสํหิตํ จิตฺตํ, \\ สยํว ปจฺจเวกฺขามิ\textsuperscript{1}, \\ รชฺชเส\textsuperscript{1} รชนีเย จ, \\ มุยฺหเส\textsuperscript{1} โมหนีเย จ, \\ วิสุทฺธานํ อยํ วาโส, \\ มา โข วิสุทฺธํ ทูเสสิ,}{\csromangathabat{\textsuperscript{*}“ราคูสํหิตํ จิตฺตํ,}{ยทา อุปฺปชฺชเต มม.} \\ \csromangathabat{สยํว ปจฺจเวกฺขามิ\textsuperscript{1},}{เอกคฺโค\textsuperscript{1} ตํ ทเมมหํ.} \\ \csromangathabat{รชฺชเส\textsuperscript{1} รชนีเย จ,}{ทุสฺสนีเย จ ทุสฺสเส.} \\ \csromangathabat{มุยฺหเส\textsuperscript{1} โมหนีเย จ,}{นิกฺขมสฺสุ วนา ตุวํ.} \\ \csromangathabat{วิสุทฺธานํ อยํ วาโส,}{นิมฺมลานํ ตปสฺสินํ.} \\ \csromangathabat{มา โข วิสุทฺธํ ทูเสสิ,}{นิกฺขมสฺสุ วนา ตุวนฺ”ติ.}}
\csromangathasetleft{\textsuperscript{*}“ราคูสํหิตํ จิตฺตํ, \\ สยํว ปจฺจเวกฺขามิ\textsuperscript{1}, \\ รชฺชเส\textsuperscript{1} รชนีเย จ, \\ มุยฺหเส\textsuperscript{1} โมหนีเย จ, \\ วิสุทฺธานํ อยํ วาโส, \\ มา โข วิสุทฺธํ ทูเสสิ,}
\csromangathagroup{\csromangathastanza{\csromanfolio{370}\csromangathabat{\csromansymbolfootnote{*}{ขุ 3. 74 ปิฏฺเฐ อปทาเน.}“ราคูสํหิตํ จิตฺตํ,}{ยทา อุปฺปชฺชเต มม.} \\ \csromangathabat{สยํว ปจฺจเวกฺขามิ\csromansharedfootnote{ae13fcd3e61d04ee}{ปจฺจเวกฺขิตฺวา (สพฺพตฺถ)},}{เอกคฺโค\csromansharedfootnote{a87fd6d5cf0de505}{เอกโก (สพฺพตฺถ)} ตํ ทเมมหํ.}}\csromangathastanza{\csromangathabat{รชฺชเส\csromansharedfootnote{dadcb7ca2ef54a8c}{รชฺชสิ (สี), รญฺชสิ (อิ)} รชนีเย จ,}{ทุสฺสนีเย จ ทุสฺสเส.} \\ \csromangathabat{มุยฺหเส\csromansharedfootnote{dadcb7ca2ef54a8c}{รชฺชสิ (สี), รญฺชสิ (อิ)} โมหนีเย จ,}{นิกฺขมสฺสุ วนา ตุวํ.}}\csromangathastanza{\csromangathabat{วิสุทฺธานํ อยํ วาโส,}{นิมฺมลานํ ตปสฺสินํ.} \\ \csromangathabat{มา โข วิสุทฺธํ ทูเสสิ,}{นิกฺขมสฺสุ วนา ตุวนฺ”ติ.}}}

\proseitem{5}{ปุน จปรํ มหาราช ปพฺพโต อจฺจุคฺคโต, เอวเมว โข มหาราช โยคินา โยคาวจเรน ญาณจฺจุคฺคเตน ภวิตพฺพํ, อิทํ มหาราช ปพฺพตสฺส จตุตฺถํ องฺคํ คเหตพฺพํ. ภาสิตมฺเปตํ มหาราช ภควตา เทวาติเทเวน\csromandash{}}

\setlength{\csromangathapagewidth}{0pt}
\setlength{\csromangathawakwidth}{0pt}
\csromangathameasuregroup{\textsuperscript{+}“ปมาทํ อปฺปมาเทน, \\ ปญฺญาปาสาทมารุยฺห, \\ ปพฺพตฏฺโฐว ภูมฏฺเฐ\textsuperscript{1},}{\csromangathabat{\textsuperscript{+}“ปมาทํ อปฺปมาเทน,}{ยทา นุทติ ปณฺฑิโต.} \\ \csromangathabat{ปญฺญาปาสาทมารุยฺห,}{อโสโก โสกินิํ ปชํ.} \\ \csromangathabat{ปพฺพตฏฺโฐว ภูมฏฺเฐ\textsuperscript{1},}{ธีโร พาเล อเวกฺขตี”ติ.}}
\csromangathasetleft{\textsuperscript{+}“ปมาทํ อปฺปมาเทน, \\ ปญฺญาปาสาทมารุยฺห, \\ ปพฺพตฏฺโฐว ภูมฏฺเฐ\textsuperscript{1},}
\csromangathagroup{\csromangathastanza{\csromangathabat{\csromansymbolfootnote{+}{ขุ 1. 17 ปิฏฺเฐ ธมฺมปเท.}“ปมาทํ อปฺปมาเทน,}{ยทา นุทติ ปณฺฑิโต.} \\ \csromangathabat{ปญฺญาปาสาทมารุยฺห,}{อโสโก โสกินิํ ปชํ.} \\ \csromangathabat{ปพฺพตฏฺโฐว ภูมฏฺเฐ\csromansharedfootnote{884ef7389ff3bf2f}{ภุมฺมฏฺเฐ (สี, อิ)},}{ธีโร พาเล อเวกฺขตี”ติ.}}}

\prose{ปุน จปรํ มหาราช ปพฺพโต อนุนฺนโต อโนนโต, เอวเมว โข มหาราช โยคินา โยคาวจเรน อุนฺนตาวนติ น กรณียา, อิทํอา มหาราช ปพฺพตสฺส ปญฺจมํ องฺคํ คเหตพฺพํ. ภาสิตมฺปตํ มหาราช อุปาสิกาย จูฬสุภทฺทาย สกสมเณ ปริกิตฺตยมานาย\csromandash{}}

\setlength{\csromangathapagewidth}{0pt}
\setlength{\csromangathawakwidth}{0pt}
\csromangathameasuregroup{“ลาเภน อุนฺนโต โลโก, \\ ลาภาลาเภน เอกฏฺฐา\textsuperscript{1},}{\csromangathabat{“ลาเภน อุนฺนโต โลโก,}{อลาเภน จ โอนโต.} \\ \csromangathabat{ลาภาลาเภน เอกฏฺฐา\textsuperscript{1},}{ตาทิสา สมณา มมา”ติ.}}
\csromangathasetleft{“ลาเภน อุนฺนโต โลโก, \\ ลาภาลาเภน เอกฏฺฐา\textsuperscript{1},}
\csromangathagroupwithcloserruled{\csromangathastanza{\csromangathabat{“ลาเภน อุนฺนโต โลโก,}{อลาเภน จ โอนโต.} \\ \csromangathabat{ลาภาลาเภน เอกฏฺฐา\csromansharedfootnote{97e7858fd2d948f5}{เอกตฺถา (ก) ธมฺมปทฏฺฐ 2. 292 ปิฏฺเฐ.},}{ตาทิสา สมณา มมา”ติ.}}}{ปพฺพต\-งฺค\-ปโญฺห ปญฺจโม.}
\csromanmatikaanchor{mtk.229}
\csromantocmarkref{subparagraph}{6. อากาสงฺคปญฺห}{mtk.229}
\bookmark[dest={mtk.229},level=5]{6. อากาสงฺคปญฺห}
\csromanheader{5}{6. \textbf{อากาสงฺคปญฺห}}
\proseitem{6}{ภนฺเต นาคเสน “อากาสสฺส ปญฺจ องฺคานิ คเหตพฺพานี”ติ ยํ วเทสิ, กตมานิ ตานิ ปญฺจ องฺคานิ คเหตพฺพานีติ. ยถา มหาราช อากาโส สพฺพโส อคโยฺห, เอวเมว โข มหาราช โยคินา \csromanfolio{371}โยคาวจเรน สพฺพโส กิเลเสหิ อคเยฺหน ภวิตพฺพํ, อิทํ มหาราช อากาสสฺส ปฐมํ องฺคํ คเหตพฺพํ.}

\prose{ปุน จปรํ มหาราช อากาโส อิสิตา\-ปส\-ภูต\-ทิชคณา\-นุสญฺจริโต, เอวเมว โข มหาราช โยคินา โยคาวจเรน “อนิจฺจํ ทุกฺขํ อนตฺตา”ติ สงฺขาเรสุ มานสํ สญฺจาร\-ยิตพฺพํ, อิทํ มหาราช อากาสสฺส ทุติยํ องฺคํ คเหตพฺพํ.}

\prose{ปุน จปรํ มหาราช อากาโส สนฺตาสนีโย, เอวเมว โข มหาราช โยคินา โยคาวจเรน สพฺพภว\-ปฏิสนฺธีสุ มานสํ อุพฺเพชยิตพฺพํ, อสฺสาโท น กาตพฺโพ, อิทํ มหาราช อากาสสฺส ตติยํ องฺคํ คเหตพฺพํ.}

\prose{ปุน จปรํ มหาราช อากาโส อนนฺโต อปฺปมาโณ อปริเมยฺโย, เอวเมว โข มหาราช โยคินา โยคาวจเรน อนนฺตสีเลน อปริธิตญาเณน ภวิตพฺพํ, อิทํ มหาราช อากาสสฺส จตุตฺถํ องฺคํ คเหตพฺพํ.}

\prosewithcloserruled{ปุน จปรํ มหาราช อากาโส อลคฺโค อสตฺโต อปฺปติฏฺฐิโต อปลิพุทฺโธ, เอวเมว โข มหาราช โยคินา โยคาวจเรน กุเล คเณ ลาเภ อาวาเส ปลิโพเธ ปจฺจเย สพฺพกิเลเสสุ จ สพฺพตฺถ อลคฺเคน ภวิตพฺพํ, อนาสตฺเตน อปฺปติฏฺฐิเตน อปลิพุทฺเธน ภวิตพฺพํ, อิทํ มหาราช อากาสสฺส ปญฺจมํ องฺคํ คเหตพฺพํ. ภาสิตมฺเปตํ มหาราช ภควตา เทวาติเทเวน สกํ ปุตฺตํ ราหุลํ โอวทนฺเตน\csromandash{}“เสยฺยถาปิ ราหุล\csromansharedfootnote{0bf288a524a99bab}{ม 2. 87 ปิฏฺเฐ.} อากาโส น กตฺถจิ ปติฏฺฐิโต, เอวเมว โข ตฺวํ ราหุล อากาสสมํ ภวนํ ภาเวหิ, อากาสสมํ หิ เต ราหุล ภาวนํ ภาวยโต อุปฺปนฺนา มนาปามนาปา ผสฺสา จิตฺตํ น ปริยาทาย ฐสฺสนฺตี”ติ.}{อากาสงฺคปโญฺห ฉฏฺโฐ.}
\csromanmatikaanchor{mtk.230}
\csromantocmarkref{subparagraph}{7. จนฺทงฺคปญฺห}{mtk.230}
\bookmark[dest={mtk.230},level=5]{7. จนฺทงฺคปญฺห}
\csromanheader{5}{7. \textbf{จนฺทงฺคปญฺห}}
\proseitem{7}{ภนฺเต นาคเสน “จนฺทสฺส ปญฺจ องฺคานิ คเหตพฺพานี”ติ ยํ วเทสิ, กตมานิ ตานิ ปญฺจ องฺคานิ คเหตพฺพานีติ. ยถา มหาราช จนฺโท \csromanfolio{372}สุกฺกปกฺเข อุทยนฺโต อุตฺตรุตฺตริํ วฑฺฒติ, เอวเมว โข มหาราช โยคินา โยคาวจเรน อาจาร\-สีล\-คุณ\-วตฺตปฺปฏิปตฺติยา อาคมาธิคเม ปฏิสลฺลาเน สติปฏฺฐาเน อินฺทฺริเยสุ คุตฺตทฺวารตาย โภชเน มตฺตญฺญุตาย ชาคริยานุโยเค อุตฺตรุตฺตริํ วฑฺฒิตพฺพํ, อิทํ มหาราช จนฺทสฺส ปฐมํ องฺคํ คเหตพฺพํ.}

\prose{ปุน จปรํ มหาราช จนฺโท อุฬาราธิปติ, เอวเมว โข มหาราช โยคินา โยคาวจเรน อุฬาเรน ฉนฺทา\-ธิปตินา ภวิตพฺพํ, อิทํ มหาราช จนฺทสฺส ทุติยํ องฺคํ คเหตพฺพํ.}

\prose{ปุน จปรํ มหาราช จนฺโท นิสาย จรติ, เอวเมว โข มหาราช โยคินา โยคาวจเรน ปวิวิตฺเตน ภวิตพฺพํ, อิทํ มหาราช จนฺทสฺส ตติยํ องฺคํ คเหตพฺพํ.}

\prose{ปุน จปรํ มหาราช จนฺโท วิมานเกตุ, เอวเมว โข มหาราช โยคินา โยคาวจเรน สีลเกตุนา ภวิตพฺพํ, อิทํ มหาราช จนฺทสฺส จตุตฺถํ องฺคํ คเหตพฺพํ.}

\prosewithcloserruled{ปุน จปรํ มหาราช จนฺโท อายาจิต\-ปตฺถิโต อุเทติ, เอวเมว โข มหาราช โยคินา โยคาวจเรน อายาจิต\-ปตฺถิเตน กุลานิ อุปสงฺกมิตพฺพานิ, อิทํ มหาราช จนฺทสฺส ปญฺจมํ องฺคํ คเหตพฺพํ. ภาสิตมฺเปตํ มหาราช ภควตา เทวาติเทเวน สํยุตฺตนิกาย\-วเร\csromandash{}“จนฺทูปมา ภิกฺขเว กุลานิ อุปสงฺกมถ, อปกสฺเสว กายํ อปกสฺส จิตฺตํ นิจฺจนวกา กุเลสุ อปฺปคพฺภา”ติ\csromansharedfootnote{755e1e2715cbde8a}{อปฺปคพฺพาติ (ก) สํ 1. 401 ปิฏฺเฐ.}.}{จนฺทงฺคปโญฺห สตฺตโม.}
\csromanmatikaanchor{mtk.231}
\csromantocmarkref{subparagraph}{8. สูริยงฺคปญฺห}{mtk.231}
\bookmark[dest={mtk.231},level=5]{8. สูริยงฺคปญฺห}
\csromanheader{5}{8. \textbf{สูริยงฺคปญฺห}}
\proseitem{8}{ภนฺเต นาคเสน “สูริยสฺส\csromansharedfootnote{acc421982f842023}{สุริยสฺส (สี, สฺยา, อิ)} สตฺต องฺคานิ คเหตพฺพานี”ติ ยํ วเทสิ, กตมานิ ตานิ สตฺต องฺคานิ คเหตพฺพานีติ. ยถา มหาราช สูริโย สพฺพํ อุทกํ ปริโสเสติ, เอวเมว โข มหาราช โยคินา โยคาวจเรน สพฺพกิเลสา อนวเสสํ ปริโสเสตพฺพา, อิทํ มหาราช สูริยสฺส ปฐมํ องฺคํ คเหตพฺพํ.}

\prose{\csromanfolio{373}ปุน จปรํ มหาราช สูริโย ตมนฺธการํ วิธมติ, เอวเมว โข มหาราช โยคินา โยคาวจเรน สพฺพํ ราคตมํ โทสตมํ โมหตมํ มานตมํ ทิฏฺฐิตมํ กิเลสตมํ สพฺพํ ทุจฺจริตตมํ วิธม\-ยิตพฺพํ, อิทํ มหาราช สูริยสฺส ทุติยํ องฺคํ คเหตพฺพํ.}

\prose{ปุน จปรํ มหาราช สูริโย อภิกฺขณํ จรติ, เอวเมว โข มหาราช โยคินา โยคาวจเรน อภิกฺขณํ โยนิโส มนสิกาโร กาตพฺโพ, อิทํ มหาราช สูริยสฺส ตติยํ องฺคํ คเหตพฺพํ.}

\prose{ปุน จปรํ มหาราช สูริโย รํสิมาลี, เอวเมว โข มหาราช โยคินา โยคาวจเรน อารมฺมณมาลินา ภวิตพฺพํ, อิทํ มหาราช สูริยสฺส จตุตฺถํ องฺคํ คเหตพฺพํ.}

\prose{ปุน จปรํ มหาราช สูริโย มหาชนกายํ สนฺตาเปนฺโต จรติ, เอวเมว โข มหาราช โยคินา โยคาวจเรน อาจาร\-สีล\-คุณ\-วตฺตปฺปฏิปตฺติยา ฌาน\-วิโมกฺข\-สมาธิ\-สมาปตฺติ- อินฺทฺริยพลโพชฺฌงฺคสติปฏฺฐานสมฺมปฺปธานอิทฺธิปาเทหิ สเทวโก โลโก สนฺตาปยิตพฺโพ, อิทํ มหาราช สูริยสฺส ปญฺจมํ องฺคํ คเหตพฺพํ.}

\prose{ปุน จปรํ มหาราช สูริโย ราหุภยา ภีโต จรติ, เอวเมว โข มหาราช โยคินา โยคาวจเรน ทุจฺจริต\-ทุคฺคติ\-วิสม\-กนฺตาร\-วิปาก\-วินิปาต\-กิเลส\-ชาล\-ชฏิเต ทิฏฺฐิ\-สงฺฆาฏ\-ปฏิมุกฺเก กุปถ\-ปกฺขนฺเท กุมฺมคฺค\-ปฏิปนฺเน\csromansharedfootnote{d6b2b84fc2c62c92}{กุมคฺคปฏิปนฺเน (สฺยา, ก)} สตฺเต ทิสฺวา มหตา สํเวคภเยน มานสํ สํเวเชตพฺพํ, อิทํ มหาราช สูริยสฺส ฉฏฺฐํ องฺคํ คเหตพฺพํ.}

\prose{ปุน จปรํ มหาราช สูริโย กลฺยาณปาปเก ทสฺเสติ, เอวเมว โข มหาราช โยคินา โยคาวจเรน อินฺทฺริยพลโพชฺฌงฺคสติปฏฺฐานสมฺมปฺปธาน- อิทฺธิปาทโลกิยโลกุตฺตรธมฺมา ทสฺเสตพฺพา, อิทํ มหาราช สูริยสฺส สตฺตมํ องฺคํ คเหตพฺพํ. ภาสิตมฺเปตํ มหาราช เถเรน วงฺคีเสน\csromandash{}}

\setlength{\csromangathapagewidth}{0pt}
\setlength{\csromangathawakwidth}{0pt}
\csromangathameasuregroup{“ยถาปิ สูริโย อุทยนฺโต, \\ สุจิญฺจ อสุจิญฺจาปิ, \\ ตถา ภิกฺขุ ธมฺมธโร, \\ ปถํ ทสฺเสติ วิวิธํ,}{\csromangathabat{“ยถาปิ สูริโย อุทยนฺโต,}{รูปํ ทสฺเสติ ปาณินํ.} \\ \csromangathabat{สุจิญฺจ อสุจิญฺจาปิ,}{กลฺยาณญฺจาปิ ปาปกํ.} \\ \csromangathabat{ตถา ภิกฺขุ ธมฺมธโร,}{อวิชฺชาปิหิตํ ชนํ.} \\ \csromangathabat{ปถํ ทสฺเสติ วิวิธํ,}{อาทิจฺโจวุทยํ ยถา”ติ.}}
\csromangathasetleft{“ยถาปิ สูริโย อุทยนฺโต, \\ สุจิญฺจ อสุจิญฺจาปิ, \\ ตถา ภิกฺขุ ธมฺมธโร, \\ ปถํ ทสฺเสติ วิวิธํ,}
\csromangathagroupwithcloserruled{\csromangathastanza{\csromangathabat{“ยถาปิ สูริโย อุทยนฺโต,}{รูปํ ทสฺเสติ ปาณินํ.} \\ \csromangathabat{สุจิญฺจ อสุจิญฺจาปิ,}{กลฺยาณญฺจาปิ ปาปกํ.}}\csromangathastanza{\csromanfolio{374}\csromangathabat{ตถา ภิกฺขุ ธมฺมธโร,}{อวิชฺชาปิหิตํ ชนํ.} \\ \csromangathabat{ปถํ ทสฺเสติ วิวิธํ,}{อาทิจฺโจวุทยํ ยถา”ติ.}}}{สูริยงฺคปโญฺห อฏฺฐโม.}
\csromanmatikaanchor{mtk.232}
\csromantocmarkref{subparagraph}{9. สกฺกงฺคปญฺห}{mtk.232}
\bookmark[dest={mtk.232},level=5]{9. สกฺกงฺคปญฺห}
\csromanheader{5}{9. \textbf{สกฺกงฺคปญฺห}}
\proseitem{9}{ภนฺเต นาคเสน “สกฺกสฺส ตีณิ องฺคานิ คเหตพฺพานี”ติ ยํ วเทสิ, กตมานิ ตานิ ตีณิ องฺคานิ คเหตพฺพานีติ. ยถา มหาราช สกฺโก เอกนฺต\-สุขสมปฺปิโต, เอวเมว โข มหาราช โยคินา โยคาวจเรน เอกนฺต\-ปวิเวกสุขาภิ\-รเตน ภวิตพฺพํ, อิทํ มหาราช สกฺกสฺส ปฐมํ องฺคํ คเหตพฺพํ.}

\prose{ปุน จปรํ มหาราช สกฺโก เทเว ทิสฺวา ปคฺคณฺหาติ, หาสม\-ภิ\-ชเนติ, เอวเมว โข มหาราช โยคินา โยคาวจเรน กุสเลสุ ธมฺเมสุ อลีนม\-ตนฺทิตํ สนฺตํ มานสํ ปคฺคเหตพฺพํ, หาสม\-ภิ\-ชเนตพฺพํ, อุฏฺฐหิตพฺพํ ฆฏิตพฺพํ วายมิตพฺพํ, อิทํ มหาราช สกฺกสฺส ทุติยํ องฺคํ คเหตพฺพํ.}

\prose{ปุน จปรํ มหาราช สกฺกสฺส อนภิรติ นุปฺปชฺชติ, เอวเมว โข มหาราช โยคินา โยคาวจเรน สุญฺญาคาเร อนภิรติ น อุปฺปาเทตพฺพา, อิทํ มหาราช สกฺกสฺส ตติยํ องฺคํ คเหตพฺพํ. ภาสิตมฺเปตํ มหาราช เถเรน สุภูตินา\csromandash{}}

\setlength{\csromangathapagewidth}{0pt}
\setlength{\csromangathawakwidth}{0pt}
\csromangathameasuregroup{“สาสเน เต มหาวีร, \\ นาภิชานามิ อุปฺปนฺนํ,}{\csromangathabat{“สาสเน เต มหาวีร,}{ยโต ปพฺพชิโต อหํ.} \\ \csromangathabat{นาภิชานามิ อุปฺปนฺนํ,}{มานสํ กามสํหิตนฺ”ติ.}}
\csromangathasetleft{“สาสเน เต มหาวีร, \\ นาภิชานามิ อุปฺปนฺนํ,}
\csromangathagroupwithcloserruled{\csromangathastanza{\csromangathabat{“สาสเน เต มหาวีร,}{ยโต ปพฺพชิโต อหํ.} \\ \csromangathabat{นาภิชานามิ อุปฺปนฺนํ,}{มานสํ กามสํหิตนฺ”ติ.}}}{สกฺกงฺคปโญฺห นวโม.}
\csromanmatikaanchor{mtk.233}
\csromantocmarkref{subparagraph}{10. จกฺกวตฺติงฺคปญฺห}{mtk.233}
\bookmark[dest={mtk.233},level=5]{10. จกฺกวตฺติงฺคปญฺห}
\csromanheader{5}{10. \textbf{จกฺกวตฺติงฺคปญฺห}}
\proseitem{30}{ภนฺเต นาคเสน “จกฺกวตฺติสฺส จตฺตาริ องฺคานิ คเหตพฺพานี”ติ ยํ วเทสิ, กตมานิ ตานิ จตฺตาริ องฺคานิ คเหตพฺพานีติ. ยถา \csromanfolio{375}มหาราช จกฺกวตฺตี จตูหิ สงฺคห\-วตฺถูหิ ชนํ สงฺคณฺหาติ, เอวเมว โข มหาราช โยคินา โยคาวจเรน จตสฺสนฺนํ ปริสานํ มานสํ สงฺคเหตพฺพํ อนุคฺคเหตพฺพํ สมฺป\-หํเส\-ตพฺพํ, อิทํ มหาราช จกฺกวตฺติสฺส ปฐมํ องฺคํ คเหตพฺพํ.}

\prose{ปุน จปรํ มหาราช จกฺกวตฺติสฺส วิชิเต โจรา น อุฏฺฐหนฺติ, เอวเมว โข มหาราช โยคินา โยคาวจเรน กามราค\-พฺยาปาท\-วิหิํสาวิตกฺกา น อุปฺปาเทตพฺพา, อิทํ มหาราช จกฺกวตฺติสฺส ทุติยํ องฺคํ คเหตพฺพํ. ภาสิตมฺเปตํ มหาราช ภควตา เทวาติเทเวน\csromandash{}}

\setlength{\csromangathapagewidth}{0pt}
\setlength{\csromangathawakwidth}{0pt}
\csromangathameasuregroup{}{“วิตกฺกูปสเม จ โย รโต,}
\csromangathameasurewak{“วิตกฺกูปสเม จ โย รโต,}
\setlength{\csromangathaleftcol}{0pt}
\csromangathagroup{\csromangathastanza{\csromangathawak{“วิตกฺกูปสเม จ โย รโต,}}}

\prose{อสุภํ ภาวยเต\csromansharedfootnote{7b2388b42f0a97fb}{ภาวยตี (สฺยา) ขุ 1. 64 ปิฏฺเฐ ธมฺมปเท.} สทา สโต.}

\prose{เอส โข พฺยนฺติกาหิติ,}

\prose{เอส เฉจฺฉติ มารพนฺธนนฺ”ติ.}

\prose{ปุน จปรํ มหาราช จกฺกวตฺตี ทิวเส ทิวเส สมุทฺท\-ปริยนฺตํ มหาปถวิํ อนุยายติ กลฺยาณ\-ปาปกานิ วิจินมาโน, เอวเมว โข มหาราช โยคินา โยคาวจเรน กายกมฺมํ วจีกมฺมํ มโนกมฺมํ ทิวเส ทิวเส ปจฺจเวกฺขิตพฺพํ “กินฺนุ โข เม อิเมหิ ตีหิ ฐาเนหิ อนุปวชฺชสฺส ทิวโส วีติวตฺตตี”ติ, อิทํ มหาราช จกฺกวตฺติสฺส ตติยํ องฺคํ คเหตพฺพํ. ภาสิตมฺเปตํ มหาราช ภควตา เทวาติเทเวน องฺคุตฺตร\-นิกายวเร\csromandash{}}

\prose{\csromansymbolfootnote{*}{อํ 3. 325 ปิฏฺเฐ.}“กถมฺภูตสฺส เม รตฺตินฺทิวา วีติวตฺตนฺตีติ\csromansharedfootnote{bdaefc0f5d2268a7}{วีติปตนฺตีติ (สี, อิ)} ปพฺพชิเตน อภิณฺหํ ปจฺจเวกฺขิตพฺพนฺ”ติ.}

\prose{ปุน จปรํ มหาราช จกฺกวตฺติสฺส อพฺภนฺตร\-พาหิรา\-รกฺขา สุสํวิหิตา โหติ, เอวเมว โข มหาราช โยคินา โยคาวจเรน อพฺภนฺตรานํ พาหิรานํ กิเลสานํ อารกฺขาย สติโทวาริโก ฐเปตพฺโพ, \csromanfolio{376}อิทํ มหาราช จกฺกวตฺติสฺส จตุตฺถํ องฺคํ คเหตพฺพํ. ภาสิตมฺเปตํ มหาราช ภควตา เทวาติเทเวน\csromandash{}}

\proseitem{6}{\csromansymbolfootnote{*}{อํ 2. 481 ปิฏฺเฐ.}“สติโทวาริโก ภิกฺขเว อริยสาวโก อกุสลํ ปชหติ กุสลํ ภาเวติ, สาวชฺชํ ปชหติ, อนวชฺชํ ภาเวติ, สุทฺธมตฺตานํ ปริหรตี”ติ.}

\vspace{\tipitakaparskip}
\csromancenter{จกฺกวตฺติงฺคปโญฺห ทสโม.}
\nitthitammidruled{\textbf{ปถวีวคฺโค ตติโย.}}
\csromancenter{\textbf{ตสฺสุทฺทานํ}}
\setlength{\csromangathapagewidth}{0pt}
\setlength{\csromangathawakwidth}{0pt}
\csromangathameasuregroup{ปถวี อาโป จ เตโช จ, \\ อากาโส จนฺทสูริโย จ,}{\csromangathabat{ปถวี อาโป จ เตโช จ,}{วาโย จ ปพฺพเตน จ.} \\ \csromangathabat{อากาโส จนฺทสูริโย จ,}{สกฺโก จ จกฺกวตฺตินาติ.}}
\csromangathasetleft{ปถวี อาโป จ เตโช จ, \\ อากาโส จนฺทสูริโย จ,}
\csromangathagroup{\csromangathastanza{\csromangathabat{ปถวี อาโป จ เตโช จ,}{วาโย จ ปพฺพเตน จ.} \\ \csromangathabat{อากาโส จนฺทสูริโย จ,}{สกฺโก จ จกฺกวตฺตินาติ.}}}

\csromanmatikaanchor{mtk.234}
\csromantocmarknopage{paragraph}{4. อุปจิกาวคฺค}{mtk.234}
\bookmark[dest={mtk.234},level=4]{4. อุปจิกาวคฺค}
\csromanheader{4}{4. \textbf{อุปจิกาวคฺค}}
\csromanmatikaanchor{mtk.235}
\csromantocmarkref{subparagraph}{1. อุปจิกงฺคปญฺห}{mtk.235}
\bookmark[dest={mtk.235},level=5]{1. อุปจิกงฺคปญฺห}
\csromanheader{5}{1. \textbf{อุปจิก}\-\textbf{งฺค}\-\textbf{ปญฺห}}
\proseitem{1}{\csromanfolio{377}ภนฺเต นาคเสน “อุปจิกาย เอกํ องฺคํ คเหตพฺพนฺ”ติ ยํ วเทสิ, กตมํ ตํ เอกํ องฺคํ คเหตพฺพนฺติ. ยถา มหาราช อุปจิกา อุปริ ฉทนํ กตฺวา อตฺตานํ ปิทหิตฺวา โคจราย จรติ, เอวเมว โข มหาราช โยคินา โยคาวจเรน สีลสํวร\-ฉทนํ กตฺวา มานสํ ปิทหิตฺวา ปิณฺฑาย จริตพฺพํ, สีลสํวร\-ฉทเนน โข มหาราช โยคี โยคาวจโร สพฺพ\-ภย\-สมติกฺกนฺโต โหติ, อิทํ มหาราช อุปจิกาย เอกํ องฺคํ คเหตพฺพํ, ภาสิตมฺเปตํ มหาราช เถเรน อุปเสเนน วงฺคนฺต\-ปุตฺเตน\csromandash{}}

\setlength{\csromangathapagewidth}{0pt}
\setlength{\csromangathawakwidth}{0pt}
\csromangathameasuregroup{“สีลสํวรฉทนํ, \\ อนุปลิตฺโต โลเกน,}{\csromangathabat{“สีลสํวรฉทนํ,}{โยคี กตฺวาน มานสํ.} \\ \csromangathabat{อนุปลิตฺโต โลเกน,}{ภยา จ ปริมุจฺจตี”ติ.}}
\csromangathasetleft{“สีลสํวรฉทนํ, \\ อนุปลิตฺโต โลเกน,}
\csromangathagroupwithcloserruled{\csromangathastanza{\csromangathabat{“สีลสํวร\-ฉทนํ,}{โยคี กตฺวาน มานสํ.} \\ \csromangathabat{อนุปลิตฺโต โลเกน,}{ภยา จ ปริมุจฺจตี”ติ.}}}{อุปจิก\-งฺค\-ปโญฺห ปฐโม.}
\csromanmatikaanchor{mtk.236}
\csromantocmarkref{subparagraph}{2. พิฬารงฺคปญฺห}{mtk.236}
\bookmark[dest={mtk.236},level=5]{2. พิฬารงฺคปญฺห}
\csromanheader{5}{2. \textbf{พิฬารงฺคปญฺห}}
\proseitem{2}{ภนฺเต นาคเสน “พิฬารสฺส เทฺว องฺคานิ คเหตพฺพานี”ติ ยํ วเทสิ, กตมานิ ตานิ เทฺว องฺคานิ คเหตพฺพานีติ. ยถา มหาราช พิฬาโร คุหาคโตปิ สุสิรคโตปิ หมฺมิยนฺตรคโตปิ อุนฺทูรํ เยว ปริเยสติ, เอวเมว โข มหาราช โยคินา โยคาวจเรน คามคเตนาปิ อรญฺญคเตนาปิ รุกฺขมูลคเตนาปิ สุญฺญาคารคเตนาปิ สตตํ สมิตํ อปฺปมตฺเตน กายคตาสติ\-โภชนํ เยว ปริเยสิตพฺพํ, อิทํ มหาราช พิฬารสฺส ปฐมํ องฺคํ คเหตพฺพํ.}

\prose{ปุน จปรํ มหาราช พิฬาโร อาสนฺเน เยว โคจรํ ปริเยสติ, เอวเมว โข มหาราช โยคินา โยคาวจเรน อิเมสุ เยว ปญฺจสุ อุปาทาน\-กฺขนฺเธสุ อุทยพฺพยา\-นุปสฺสินา วิหริตพฺพํ “อิติ รูปํ อิติ รูปสฺส สมุทโย อิติ รูปสฺส อตฺถงฺคโม, อิติ เวทนา อิติ เวทนาย สมุทโย อิติ เวทนาย อตฺถงฺคโม, อิติ สญฺญา อิติ สญฺญาย \csromanfolio{378}สมุทโย อิติ สญฺญาย อตฺถงฺคโม, อิติ สงฺขารา อิติ สงฺขารานํ สมุทโย อิติ สงฺขารานํ อตฺถงฺคโม, อิติ วิญฺญาณํ อิติ วิญฺญาณสฺส สมุทโย อิติ วิญฺญาณสฺส อตฺถงฺคโม”ติ. อิทํ มหาราช พิฬารสฺส ทุติยํ องฺคํ คเหตพฺพํ. ภาสิตมฺเปตํ มหาราช ภควตา เทวาติเทเวน\csromandash{}}

\setlength{\csromangathapagewidth}{0pt}
\setlength{\csromangathawakwidth}{0pt}
\csromangathameasuregroup{“น อิโต ทูเร ภวิตพฺพํ, \\ ปจฺจุปฺปนฺนมฺหิ โวหาเร,}{\csromangathabat{“น อิโต ทูเร ภวิตพฺพํ,}{ภวคฺคํ กิํ กริสฺสติ.} \\ \csromangathabat{ปจฺจุปฺปนฺนมฺหิ โวหาเร,}{สเก กายมฺหิ วินฺทถา”ติ.}}
\csromangathasetleft{“น อิโต ทูเร ภวิตพฺพํ, \\ ปจฺจุปฺปนฺนมฺหิ โวหาเร,}
\csromangathagroupwithcloserruled{\csromangathastanza{\csromangathabat{“น อิโต ทูเร ภวิตพฺพํ,}{ภวคฺคํ กิํ กริสฺสติ.} \\ \csromangathabat{ปจฺจุปฺปนฺนมฺหิ โวหาเร,}{สเก กายมฺหิ วินฺทถา”ติ.}}}{พิฬารงฺคปโญฺห ทุติโย.}
\csromanmatikaanchor{mtk.237}
\csromantocmarkref{subparagraph}{3. อุนฺทูรงฺคปญฺห}{mtk.237}
\bookmark[dest={mtk.237},level=5]{3. อุนฺทูรงฺคปญฺห}
\csromanheader{5}{3. \textbf{อุนฺทูรงฺคปญฺห}}
\proseitem{3}{ภนฺเต นาคเสน “อุนฺทูรสฺส\csromansharedfootnote{d8cecea569a5b732}{อุนฺทุรสฺส (สฺยา, ก)} เอกํ องฺคํ คเหตพฺพนฺ”ติ ยํ วเทสิ, กตมํ ตํ เอกํ องฺคํ คเหตพฺพนฺติ. ยถา มหาราช อุนฺทูโร อิโตจิโต จ วิจรนฺโต อาหารูปาสีสโก เยว จรติ, เอวเมว โข มหาราช โยคินา โยคาวจเรน อิโตจิโต จ วิจรนฺเตน โยนิโส มนสิ\-กา\-รูปา\-สีสเกเนว ภวิตพฺพํ, อิทํ มหาราช อุนฺทูรสฺส เอกํ องฺคํ คเหตพฺพํ. ภาสิตมฺเปตํ มหาราช เถเรน อุปเสเนน วงฺคนฺต\-ปุตฺเตน\csromandash{}}

\setlength{\csromangathapagewidth}{0pt}
\setlength{\csromangathawakwidth}{0pt}
\csromangathameasuregroup{“ธมฺมาสีสํ\textsuperscript{1} กริตฺวาน, \\ อโนลีโน วิหรติ,}{\csromangathabat{“ธมฺมาสีสํ\textsuperscript{1} กริตฺวาน,}{วิหรนฺโต วิปสฺสโก.} \\ \csromangathabat{อโนลีโน วิหรติ,}{อุปสนฺโต สทา สโต”ติ.}}
\csromangathasetleft{“ธมฺมาสีสํ\textsuperscript{1} กริตฺวาน, \\ อโนลีโน วิหรติ,}
\csromangathagroupwithcloserruled{\csromangathastanza{\csromangathabat{“ธมฺมาสีสํ\csromansharedfootnote{7694a626a9e36285}{ธมฺมสีสํ (สี, อิ)} กริตฺวาน,}{วิหรนฺโต วิปสฺสโก.} \\ \csromangathabat{อโนลีโน วิหรติ,}{อุปสนฺโต สทา สโต”ติ.}}}{อุนฺทูรงฺคปโญฺห ตติโย.}
\csromanmatikaanchor{mtk.238}
\csromantocmarkref{subparagraph}{4. วิจฺฉิกงฺคปญฺห}{mtk.238}
\bookmark[dest={mtk.238},level=5]{4. วิจฺฉิกงฺคปญฺห}
\csromanheader{5}{4. \textbf{วิจฺฉิก}\-\textbf{งฺค}\-\textbf{ปญฺห}}
\proseitem{2}{ภนฺเต นาคเสน “วิจฺฉิกสฺส เอกํ องฺคํ คเหตพฺพนฺ”ติ ยํ วเทสิ, กตมํ ตํ เอกํ องฺคํ คเหตพฺพนฺติ. ยถา มหาราช วิจฺฉิโก นงฺคุลาวุโธ นงฺคุลํ อุสฺสาเปตฺวา จรติ, เอวเมว โข มหาราช โยคินา โยคาวจเรน ญาณาวุเธน ภวิตพฺพํ, ญาณํ อุสฺสาเปตฺวา \csromanfolio{379}วิหริตพฺพํ, อิทํ มหาราช วิจฺฉิกสฺส เอกํ องฺคํ คเหตพฺพํ. ภาสิตมฺเปตํ มหาราช เถเรน อุปเสเนน วงฺคนฺต\-ปุตฺเตน\csromandash{}}

\setlength{\csromangathapagewidth}{0pt}
\setlength{\csromangathawakwidth}{0pt}
\csromangathameasuregroup{“ญาณขคฺคํ คเหตฺวาน, \\ ปริมุจฺจติ สพฺพภยา,}{\csromangathabat{“ญาณขคฺคํ คเหตฺวาน,}{วิหรนฺโต วิปสฺสโก.} \\ \csromangathabat{ปริมุจฺจติ สพฺพภยา,}{ทุปฺปสโห จ โส ภเว”ติ.}}
\csromangathasetleft{“ญาณขคฺคํ คเหตฺวาน, \\ ปริมุจฺจติ สพฺพภยา,}
\csromangathagroupwithcloserruled{\csromangathastanza{\csromangathabat{“ญาณขคฺคํ คเหตฺวาน,}{วิหรนฺโต วิปสฺสโก.} \\ \csromangathabat{ปริมุจฺจติ สพฺพภยา,}{ทุปฺปสโห จ โส ภเว”ติ.}}}{วิจฺฉิก\-งฺค\-ปโญฺห จตุตฺโถ.}
\csromanmatikaanchor{mtk.239}
\csromantocmarkref{subparagraph}{5. นกุลงฺคปญฺห}{mtk.239}
\bookmark[dest={mtk.239},level=5]{5. นกุลงฺคปญฺห}
\csromanheader{5}{5. \textbf{นกุลงฺคปญฺห}}
\proseitem{4}{ภนฺเต นาคเสน “นกุลสฺส เอกํ องฺคํ คเหตพฺพนฺ”ติ ยํ วเทสิ, กตมํ ตํ เอกํ องฺคํ คเหตพฺพนฺติ. ยถา มหาราช นกุโล อุรคมุ\-ปคจฺฉนฺโต เภสชฺเชน กายํ ปริภาเวตฺวา อุรคมุ\-ปคจฺฉติ คเหตุํ, เอวเมว โข มหาราช โยคินา โยคาวจเรน โกธา\-ฆาต\-พหุลํ กลห\-วิคฺคห\-วิวาท\-วิโรธา\-ภิภูตํ โลกมุ\-ปคจฺฉนฺเตน เมตฺตา\-เภสชฺเชน มานสํ อนุลิมฺปิตพฺพํ, อิทํ มหาราช นกุลสฺส เอกํ องฺคํ คเหตพฺพํ. ภาสิตมฺเปตํ มหาราช เถเรน สาริปุตฺเตน ธมฺม\-เสนาปตินา\csromandash{}}

\setlength{\csromangathapagewidth}{0pt}
\setlength{\csromangathawakwidth}{0pt}
\csromangathameasuregroup{“ตสฺมา สกํ ปเรสมฺปิ, \\ เมตฺตจิตฺเตน ผริตพฺพํ,}{\csromangathabat{“ตสฺมา สกํ ปเรสมฺปิ,}{กาตพฺพา เมตฺตภาวนา.} \\ \csromangathabat{เมตฺตจิตฺเตน ผริตพฺพํ,}{เอตํ พุทฺธาน สาสนนฺ”ติ.}}
\csromangathasetleft{“ตสฺมา สกํ ปเรสมฺปิ, \\ เมตฺตจิตฺเตน ผริตพฺพํ,}
\csromangathagroupwithcloserruled{\csromangathastanza{\csromangathabat{“ตสฺมา สกํ ปเรสมฺปิ,}{กาตพฺพา เมตฺตภาวนา.} \\ \csromangathabat{เมตฺตจิตฺเตน ผริตพฺพํ,}{เอตํ พุทฺธาน สาสนนฺ”ติ.}}}{นกุลงฺคปโญฺห ปญฺจโม.}
\csromanmatikaanchor{mtk.240}
\csromantocmarkref{subparagraph}{6. ชรสิงฺคาลงฺคปญฺห}{mtk.240}
\bookmark[dest={mtk.240},level=5]{6. ชรสิงฺคาลงฺคปญฺห}
\csromanheader{5}{6. \textbf{ชรสิงฺคาล}\-\textbf{งฺค}\-\textbf{ปญฺห}}
\proseitem{4}{ภนฺเต นาคเสน “ชรสิงฺคาลสฺส เทฺว องฺคานิ คเหตพฺพานี”ติ ยํ วเทสิ, กตมานิ ตานิ เทฺว องฺคานิ คเหตพฺพานีติ. ยถา มหาราช ชรสิงฺคาโล โภชนํ ปฏิลภิตฺวา อชิคุจฺฉมาโน ยาวทตฺถํ อาหรยติ, เอวเมว โข มหาราช โยคินา โยคาวจเรน โภชนํ ปฏิลภิตฺวา อชิคุจฺฉมาเนน สรีร\-ยาปนมตฺตเมว ปริภุญฺชิ\-ตพฺพํ, อิทํ มหาราช ชรสิงฺคาลสฺส ปฐมํ องฺคํ คเหตพฺพํ. ภาสิตมฺเปตํ มหาราช เถเรน มหากสฺสเปน\csromandash{}}

\setlength{\csromangathapagewidth}{0pt}
\setlength{\csromangathawakwidth}{0pt}
\csromangathameasuregroup{\textsuperscript{*}“เสนาสนมฺหา โอรุยฺห, \\ ภุญฺชนฺตํ ปุริสํ กุฏฺฐํ, \\ โส เม ปกฺเกน หตฺเถน, \\ อาโลปํ ปกฺขิปนฺตสฺส, \\ กุฏฺฏมูลญฺจ นิสฺสาย, \\ ภุญฺชมาเน วา ภุตฺเต วา,}{\csromangathabat{\textsuperscript{*}“เสนาสนมฺหา โอรุยฺห,}{คามํ ปิณฺฑาย ปาวิสิํ.} \\ \csromangathabat{ภุญฺชนฺตํ ปุริสํ กุฏฺฐํ,}{สกฺกจฺจ นํ อุปฏฺฐหิํ.} \\ \csromangathabat{โส เม ปกฺเกน หตฺเถน,}{อาโลปํ อุปนามยิ.} \\ \csromangathabat{อาโลปํ ปกฺขิปนฺตสฺส,}{องฺคุลิเปตฺถ ฉิชฺชถ.} \\ \csromangathabat{กุฏฺฏมูลญฺจ นิสฺสาย,}{อาโลปํ ตํ อภุญฺชิสํ.} \\ \csromangathabat{ภุญฺชมาเน วา ภุตฺเต วา,}{เชคุจฺฉํ เม น วิชฺชตี”ติ.}}
\csromangathasetleft{\textsuperscript{*}“เสนาสนมฺหา โอรุยฺห, \\ ภุญฺชนฺตํ ปุริสํ กุฏฺฐํ, \\ โส เม ปกฺเกน หตฺเถน, \\ อาโลปํ ปกฺขิปนฺตสฺส, \\ กุฏฺฏมูลญฺจ นิสฺสาย, \\ ภุญฺชมาเน วา ภุตฺเต วา,}
\csromangathagroup{\csromangathastanza{\csromangathabat{\csromansymbolfootnote{*}{ขุ 2. 350 ปิฏฺเฐ เถรคาถายํ.}“เสนาสนมฺหา โอรุยฺห,}{คามํ ปิณฺฑาย ปาวิสิํ.} \\ \csromangathabat{ภุญฺชนฺตํ ปุริสํ กุฏฺฐํ,}{สกฺกจฺจ นํ อุปฏฺฐหิํ.}}\csromangathastanza{\csromanfolio{380}\csromangathabat{โส เม ปกฺเกน หตฺเถน,}{อาโลปํ อุปนามยิ.} \\ \csromangathabat{อาโลปํ ปกฺขิปนฺตสฺส,}{องฺคุลิเปตฺถ ฉิชฺชถ.}}\csromangathastanza{\csromangathabat{กุฏฺฏมูลญฺจ นิสฺสาย,}{อาโลปํ ตํ อภุญฺชิสํ.} \\ \csromangathabat{ภุญฺชมาเน วา ภุตฺเต วา,}{เชคุจฺฉํ เม น วิชฺชตี”ติ.}}}

\proseitem{6}{ปุน จปรํ มหาราช ชรสิงฺคาโล โภชนํ ปฏิลภิตฺวา น วิจินาติ ลูขํ วา ปณีตํ วาติ, เอวเมว โข มหาราช โยคินา โยคาวจเรน โภชนํ ปฏิลภิตฺวา น วิจินิตพฺพํ “ลูขํ วา ปณีตํ วา สมฺปนฺนํ วา อสมฺปนฺนํ วา”ติ, ยถา ลทฺเธน สนฺตุสฺสิตพฺพํ, อิทํ มหาราช ชรสิงฺคาลสฺส ทุติยํ องฺคํ คเหตพฺพํ. ภาสิตมฺเปตํ มหาราช เถเรน อุปเสเนน วงฺคนฺต\-ปุตฺเตน\csromandash{}}

\setlength{\csromangathapagewidth}{0pt}
\setlength{\csromangathawakwidth}{0pt}
\csromangathameasuregroup{\textsuperscript{*}“ลูเขนปิ จ สนฺตุสฺเส, \\ รเสสุ อนุคิทฺธสฺส, \\ อิตรีตเรน สนฺตุฏฺโฐ\textsuperscript{1},}{\csromangathabat{\textsuperscript{*}“ลูเขนปิ จ สนฺตุสฺเส,}{นาญฺญํ ปตฺเถ รสํ พหุํ.} \\ \csromangathabat{รเสสุ อนุคิทฺธสฺส,}{ฌาเน น รมเต\textsuperscript{1} มโน.} \\ \csromangathabat{อิตรีตเรน สนฺตุฏฺโฐ\textsuperscript{1},}{สามญฺญํ ปริปูรตี”ติ.}}
\csromangathasetleft{\textsuperscript{*}“ลูเขนปิ จ สนฺตุสฺเส, \\ รเสสุ อนุคิทฺธสฺส, \\ อิตรีตเรน สนฺตุฏฺโฐ\textsuperscript{1},}
\csromangathagroupwithcloserruled{\csromangathastanza{\csromangathabat{\csromansymbolfootnote{*}{ขุ 2. 305 ปิฏฺเฐ.}“ลูเขนปิ จ สนฺตุสฺเส,}{นาญฺญํ ปตฺเถ รสํ พหุํ.} \\ \csromangathabat{รเสสุ อนุคิทฺธสฺส,}{ฌาเน น รมเต\csromansharedfootnote{dbce7c77354f138c}{รมตี (สี, อิ)} มโน.} \\ \csromangathabat{อิตรีตเรน สนฺตุฏฺโฐ\csromansharedfootnote{dbce7c77354f138c}{รมตี (สี, อิ)},}{สามญฺญํ ปริปูรตี”ติ.}}}{ชรสิงฺคาล\-งฺค\-ปโญฺห ฉฏฺโฐ.}
\csromanmatikaanchor{mtk.241}
\csromantocmarkref{subparagraph}{7. มิคงฺคปญฺห}{mtk.241}
\bookmark[dest={mtk.241},level=5]{7. มิคงฺคปญฺห}
\csromanheader{5}{7. \textbf{มิคงฺคปญฺห}}
\proseitem{7}{ภนฺเต นาคเสน “มิคสฺส ตีณิ องฺคานิ คเหตพฺพานี”ติ ยํ วเทสิ, กตมานิ ตานิ ตีณิ องฺคานิ คเหตพฺพานีติ. ยถา มหาราช มิโค ทิวา อรญฺเญ จรติ, รตฺติํ อพฺโภกาเส, เอวเมว โข มหาราช โยคินา โยคาวจเรน ทิวา อรญฺเญ วิหริตพฺพํ, รตฺติํ อพฺโภกาเส, อิทํ มหาราช มิคสฺส ปฐมํ องฺคํ คเหตพฺพํ.}

\prose{ภาสิตมฺเปตํ มหาราช ภควตา เทวาติเทเวน\csromansymbolfootnote{+}{ม 1. 113 ปิฏฺเฐ.}โลมหํสน\-ปริยาเย\csromandash{}}

\prose{“โส โข อหํ สาริปุตฺต ยา ตา รตฺติโย สีตา เหมนฺติกา อนฺตรฏฺฐกา หิมปาตสมยา\csromansharedfootnote{f8a1c4057d776046}{อนฺตรฏฺฐเก หิมปาตสมเย (สี, อิ, ก)}, ตถารูปาสุ รตฺถีสุ รตฺติํ อพฺโภกาเส วิหรามิ, ทิวา วนสณฺเฑ. คิมฺหานํ ปจฺฉิเม มาเส ทิวา อพฺโภกาเส วิหรามิ, รตฺติํ วนสณฺเฑ”ติ.}

\prose{\csromanfolio{381}ปุน จปรํ มหาราช มิโค สตฺติมฺหิ วา สเร วา โอปตนฺเต วญฺเจติ\csromansharedfootnote{eb3414d73468c1a5}{วชฺเชติ (ก)} ปลายติ, น กายมุปเนติ, เอวเมว โข มหาราช โยคินา โยคาวจเรน กิเลเสสุ โอปตนฺเตสุ วญฺจยิตพฺพํ\csromansharedfootnote{4adb465eb6910ae6}{วชฺชยิตพฺพํ (ก)} ปลายิตพฺพํ, น จิตฺตมุ\-ปเนตพฺพํ, อิทํ มหาราช มิคสฺส ทุติยํ องฺคํ คเหตพฺพํ.}

\prose{ปุน จปรํ มหาราช มิโค มนุสฺเส ทิสฺวา เยน วา เตน วา ปลายติ “มา มํ เต อทฺทสํสู”ติ, เอวเมว โข มหาราช โยคินา โยคาวจเรน ภณฺฑนกลหวิคฺคหวิวาท\-สีเล ทุสฺสีเล กุสีเต สงฺคณิการาเม ทิสฺวา เยน วา เตน วา ปลายิตพฺพํ “มา มํ เต อทฺทสํสุ, อหญฺจ เต มา อทฺทสนฺ”ติ. อิทํ มหาราช มิคสฺส ตติยํ องฺคํ คเหตพฺพํ. ภาสิตมฺเปตํ มหาราช เถเรน สาริปุตฺเตน ธมฺม\-เสนาปตินา\csromandash{}}

\setlength{\csromangathapagewidth}{0pt}
\setlength{\csromangathawakwidth}{0pt}
\csromangathameasuregroup{\textsuperscript{*}“มา เม กทาจิ ปาปิจฺโฉ, \\ อปฺปสฺสุโต อนาจาโร,}{\csromangathabat{\textsuperscript{*}“มา เม กทาจิ ปาปิจฺโฉ,}{กุสีโต หีนวีริโย.} \\ \csromangathabat{อปฺปสฺสุโต อนาจาโร,}{สมฺมโต\textsuperscript{1} อหุ กตฺถจี”ติ.}}
\csromangathasetleft{\textsuperscript{*}“มา เม กทาจิ ปาปิจฺโฉ, \\ อปฺปสฺสุโต อนาจาโร,}
\csromangathagroupwithcloserruled{\csromangathastanza{\csromangathabat{\csromansymbolfootnote{*}{ขุ 2. 344 ปิฏฺเฐ.}“มา เม กทาจิ ปาปิจฺโฉ,}{กุสีโต หีนวีริโย.} \\ \csromangathabat{อปฺปสฺสุโต อนาจาโร,}{สมฺมโต\csromansharedfootnote{eeb530451dc0a289}{สเมโต (สี, อิ)} อหุ กตฺถจี”ติ.}}}{มิคงฺคปโญฺห สตฺตโม.}
\csromanmatikaanchor{mtk.242}
\csromantocmarkref{subparagraph}{8. โครูปงฺคปญฺห}{mtk.242}
\bookmark[dest={mtk.242},level=5]{8. โครูปงฺคปญฺห}
\csromanheader{5}{8. \textbf{โครูปงฺคปญฺห}}
\proseitem{8}{ภนฺเต นาคเสน “โครูปสฺส จตฺตาริ องฺคานิ คเหตพฺพานี”ติ ยํ เวเทสิ, กตมานิ ตานิ จตฺตาริ องฺคานิ คเหตพฺพานีติ. ยถา มหาราช โครูโป สกํ เคหํ น วิชหติ, เอวเมว โข มหาราช โยคินา โยคาวจเรน สโก กาโย น วิชหิตพฺโพ “อนิจฺจุจฺฉาทน ปริมทฺทน\-เภทน วิกิรณวิทฺธํสนธมฺโม อยํ กาโย”ติ. อิทํ มหาราช โครูปสฺส ปฐมํ องฺคํ คเหตพฺพํ.}

\prose{ปุน จปรํ มหาราช โครูโป อาทินฺนธุโร สุขทุกฺเขน ธุรํ วหติ, เอวเมว โข มหาราช โยคินา โยคาวจเรน อาทินฺน\-พฺรหฺมจริเยน สุขทุกฺเขน ยาว ชีวตปริยาทานา อาปาณโกฏิกํ พฺรหฺมจริยํ จริตพฺพํ. อิทํ มหาราช โครูปสฺส ทุติยํ องฺคํ คเหตพฺพํ.}

\csromanmatikaanchor{mtk.243}
\csromantocmarkref{subparagraph}{9. วหาหงฺคปญฺห}{mtk.243}
\bookmark[dest={mtk.243},level=5]{9. วหาหงฺคปญฺห}
\prose{ปุน จปรํ มหาราช โครูโป ฉนฺเทน ฆายมาโน ปานียํ ปิวติ, เอวเมว โข มหาราช โยคินา โยคาวจเรน อาจริยุ\-ปชฺฌายานํ \csromanfolio{382}อนุสิฏฺฐิ ฉนฺเทน เปเมน ปสาเทน ฆายมาเนน ปฏิคฺคเหตพฺพา. อิทํ มหาราช โครูปสฺส ตติยํ องฺคํ คเหตพฺพํ.}

\prose{ปุน จปรํ มหาราช โครูโป เยน เกนจิ วาหิยมาโน วหติ, เอวเมว โข มหาราช โยคินา โยคาวจเรน เถรนวมชฺฌิมภิกฺขูนมฺปิ คิหิอุปาสกสฺสาปิ โอวาทานุสาสนี สิรสา สมฺปฏิจฺฉิตพฺพา. อิทํ มหาราช โครูปสฺส จตุตฺถํ องฺคํ คเหตพฺพํ. ภาสิตมฺเปตํ มหาราช เถเรน สาริปุตฺเตน ธมฺม\-เสนาปตินา\csromandash{}}

\setlength{\csromangathapagewidth}{0pt}
\setlength{\csromangathawakwidth}{0pt}
\csromangathameasuregroup{“ตทหุ ปพฺพชิโต สนฺโต, \\ โสปิ มํ อนุสาเสยฺย, \\ ติพฺพํ ฉนฺทญฺจ เปมญฺจ, \\ ฐเปยฺยาจริยฏฺฐาเน,}{\csromangathabat{“ตทหุ ปพฺพชิโต สนฺโต,}{ชาติยา สตฺตวสฺสิโก.} \\ \csromangathabat{โสปิ มํ อนุสาเสยฺย,}{สมฺปฏิจฺฉามิ มตฺถเก\textsuperscript{1}.} \\ \csromangathabat{ติพฺพํ ฉนฺทญฺจ เปมญฺจ,}{ตสฺมิํ ทิสฺวา อุปฏฺฐเป.} \\ \csromangathabat{ฐเปยฺยาจริยฏฺฐาเน,}{สกฺกจฺจ นํ ปุนปฺปุนนฺ”ติ.}}
\csromangathasetleft{“ตทหุ ปพฺพชิโต สนฺโต, \\ โสปิ มํ อนุสาเสยฺย, \\ ติพฺพํ ฉนฺทญฺจ เปมญฺจ, \\ ฐเปยฺยาจริยฏฺฐาเน,}
\csromangathagroupwithcloserruled{\csromangathastanza{\csromangathabat{“ตทหุ ปพฺพชิโต สนฺโต,}{ชาติยา สตฺตวสฺสิโก.} \\ \csromangathabat{โสปิ มํ อนุสาเสยฺย,}{สมฺปฏิจฺฉามิ มตฺถเก\csromansharedfootnote{407b0746f53ed6c5}{มุทฺธนา (สี)}.}}\csromangathastanza{\csromangathabat{ติพฺพํ ฉนฺทญฺจ เปมญฺจ,}{ตสฺมิํ ทิสฺวา อุปฏฺฐเป.} \\ \csromangathabat{ฐเปยฺยาจริยฏฺฐาเน,}{สกฺกจฺจ นํ ปุนปฺปุนนฺ”ติ.}}}{โครูปงฺคปโญฺห อฏฺฐโม.}
\csromanheader{2}{9. \textbf{วราหงฺคปญฺห}}
\proseitem{9}{ภนฺเต นาคเสน “วาราหสฺส เทฺว องฺคานิ คเหตพฺพานี”ติ ยํ วเทสิ, กตมานิ ตานิ เทฺว องฺคานิ คเหตพฺพานีติ. ยถา มหาราช วราโห สนฺตตฺต\-กฐิเต\csromansharedfootnote{29d8a0e99de9c760}{สนฺตตฺตกฐิเน (สี, อิ)} คิมฺหสมเย สมฺปตฺเต อุทกํ อุปคจฺฉติ, เอวเมว โข มหาราช โยคินา โยคาวจเรน โทเสน จิตฺเต อาลุฬิต\-ขลิต\-วิพฺภนฺต\-สนฺตตฺเต สีตลา\-มต\-ปณีต\-เมตฺตาภาวนํ อุปคนฺตพฺพํ. อิทํ มหาราช วราหสฺส ปฐมํ องฺคํ คเหตพฺพํ.}

\prose{ปุน จปรํ มหาราช วราโห จิกฺขลฺลมุ\-ทกมุ\-ปคนฺตฺวา นาสิกาย ปถวิํ ขณิตฺวา โทณํ กตฺวา โทณิกาย สยติ, เอวเมว โข มหาราช โยคินา โยคาวจเรน มานเส กายํ นิกฺขิปิตฺวา อารมฺมณ\-นฺตร\-คเตน สยิตพฺพํ. อิทํ มหาราช วราหสฺส ทุติยํ องฺคํ คเหตพฺพํ. ภาสิตมฺเปตํ มหาราช เถเรน ปิณฺโฑล\-ภารทฺวาเชน\csromandash{}}

\setlength{\csromangathapagewidth}{0pt}
\setlength{\csromangathawakwidth}{0pt}
\csromangathameasuregroup{“กาเย\textsuperscript{1} สภาวํ ทิสฺวาน, \\ เอกากิโย อทุติโย,}{\csromangathabat{“กาเย\textsuperscript{1} สภาวํ ทิสฺวาน,}{วิจินิตฺวา วิปสฺสโก.} \\ \csromangathabat{เอกากิโย อทุติโย,}{เสติ อารมฺมณนฺตเร”ติ.}}
\csromangathasetleft{“กาเย\textsuperscript{1} สภาวํ ทิสฺวาน, \\ เอกากิโย อทุติโย,}
\csromangathagroup{\csromangathastanza{\csromangathabat{“กาเย\csromansharedfootnote{c01764ff5660ddb1}{กาเยน (ก)} สภาวํ ทิสฺวาน,}{วิจินิตฺวา วิปสฺสโก.} \\ \csromangathabat{เอกากิโย อทุติโย,}{เสติ อารมฺมณนฺตเร”ติ.}}}

\vspace{\tipitakaparskip}
\csromancenter{วราหงฺคปโญฺห นวโม.}
\csromanmatikaanchor{mtk.244}
\csromantocmarkref{subparagraph}{10. หตฺถิงฺคปญฺห}{mtk.244}
\bookmark[dest={mtk.244},level=5]{10. หตฺถิงฺคปญฺห}
\csromanmarkh{6. โอปมฺม\-กถา\-ปญฺห}\csromanmarkhii{10. หตฺถิงฺคปญฺห}\csromanheadernomark{5}{6. โอปมฺม\-กถา\-ปญฺห \textbf{10. หตฺถิงฺคปญฺห}}
\proseitem{10}{\csromanfolio{383}ภนฺเต นาคเสน “หตฺถิสฺส ปญฺจ องฺคานิ คเหตพฺพานี”ติ ยํ วเทสิ, กตมานิ ตานิ ปญฺจ องฺคานิ คเหตพฺพานีติ. ยถา มหาราช หตฺถี นาม จรนฺโต เยว ปถวิํ ทาเลติ, เอวเมว โข มหาราช โยคินา โยคาวจเรน กายํ สมฺมาสมาเนเนว สพฺเพ กิเลสา ทาเลตพฺพา, อิทํ มหาราช หตฺถิสฺส ปฐมํ องฺคํ คเหตพฺพํ.}

\prose{ปุน จปรํ มหาราช หตฺถี สพฺพกาเยเนว อปโลเกติ, อุชุกํ เยว เปกฺขติ, น ทิสาวิทิสา วิโลเกติ, เอวเมว โข มหาราช โยคินา โยคาวจเรน สพฺพกาเยน อปโลกินา ภวิตพฺพํ, น ทิสาวิทิสา วิลิเกตพฺพา, น อุทฺธํ อุลฺโลเกตพฺพํ, น อโธ โอโลเกตพฺพํ, ยุคมตฺตปฺเปกฺขินา ภวิตพฺพํ, อิทํ มหาราช หตฺถิสฺส ทุติยํ องฺคํ คเหตพฺพํ.}

\prose{ปุน จปรํ มหาราช หตฺถี อนิพทฺธสยโน โคจรายมนุคนฺตฺวา น ตเมว เทสํ วาส\-ตฺถมุ\-ปคจฺฉติ, น ธุวปฺปติฏฺฐา\-ลโย, เอวเมว โข มหาราช โยคิน โยคาวจเรน อนิพทฺธสยเนน ภวิตพฺพํ, นิราลเยน ปิณฺฑาย คนฺตพฺพํ, ยทิ ปสฺสติ วิปสฺสโก มนุญฺญํ ปติรูปํ รุจิตเทเส ภวํ มณฺฑปํ วา รุกฺขมูลํ วา คุหํ วา ปพฺภารํ วา, ตตฺเถว วาสมุ\-ปคนฺตพฺพํ, ธุวปฺปติฏฺฐา\-ลโย น กาตพฺโพ, อิทํ มหาราช หตฺถิสฺส ตติยํ องฺคํ คเหตพฺพํ.}

\prose{ปุน จปรํ มหาราช หตฺถี อุทกํ โอคาหิตฺวา สุจิวิมล\-สีตล\-สลิล\-ปริปุณฺณํ กุมุทุ\-ปฺปล\-ปทุม\-ปุณฺฑรีก\-สญฺฉนฺนํ มหติมหนฺตํ ปทุมสรํ โอคาหิตฺวา กีฬติ คชวรกีฬํ, เอวเมว โข มหาราช โยคินา โยคาวจเรน สุจิวิมล\-วิปฺปสนฺนมนาวิล\-ธมฺมวร\-วาริ\-ปุณฺณํ วิมุตฺติ\-กุสุม\-สญฺฉนฺนํ มหาสติปฏฺฐานโปกฺขรณํ โอคาหิตฺวา ญาเณน สงฺขารา โอธุนิตพฺพา วิธุนิตพฺพา, โยคาวจรกีฬา กีฬิตพฺพา, อิทํ มหาราช หตฺถิสฺส จตุตฺถํ องฺคํ คเหตพฺพํ.}

\prose{ปุน จปรํ มหาราช หตฺถี สโต ปาทํ อุทฺธรติ, สโต ปาทํ นิกฺขิปติ, เอวเมว โข มหาราช โยคินา โยคาวจเรน สเตน สมฺปชาเนน ปาทํ อุทฺธริตพฺพํ, สเตน สมฺปชาเนน ปาทํ นิกฺขิปิตพฺพํ, อภิกฺกม\-ปฏิกฺกเม สมิญฺชน\-ปสารเณ สพฺพตฺถ สเตน สมฺปชาเนน ภวิตพฺพํ, อิทํ มหาราช \csromanfolio{384}หตฺถิสฺส ปญฺจมํ องฺคํ คเหตพฺพํ. ภาสิตมฺเปตํ มหาราช ภควตา เทวาติเทเวน สํยุตฺตนิกาย\-วเร\csromandash{}}

\setlength{\csromangathapagewidth}{0pt}
\setlength{\csromangathawakwidth}{0pt}
\csromangathameasuregroup{\textsuperscript{*}“กาเยน สํวโร สาธุ, \\ มนสา สํวโร สาธุ, \\ สพฺพตฺถ สํวุโต ลชฺชี,}{\csromangathabat{\textsuperscript{*}“กาเยน สํวโร สาธุ,}{สาธุ วาจาย สํวโร.} \\ \csromangathabat{มนสา สํวโร สาธุ,}{สาธุ สพฺพตฺถ สํวโร.} \\ \csromangathabat{สพฺพตฺถ สํวุโต ลชฺชี,}{รกฺขิโตติ ปวุจฺจตี”ติ.}}
\csromangathasetleft{\textsuperscript{*}“กาเยน สํวโร สาธุ, \\ มนสา สํวโร สาธุ, \\ สพฺพตฺถ สํวุโต ลชฺชี,}
\csromangathagroup{\csromangathastanza{\csromangathabat{\csromansymbolfootnote{*}{สํ 1. 73; ขุ 1. 65 ปิฏฺเฐสุ.}“กาเยน สํวโร สาธุ,}{สาธุ วาจาย สํวโร.} \\ \csromangathabat{มนสา สํวโร สาธุ,}{สาธุ สพฺพตฺถ สํวโร.} \\ \csromangathabat{สพฺพตฺถ สํวุโต ลชฺชี,}{รกฺขิโตติ ปวุจฺจตี”ติ.}}}

\vspace{\tipitakaparskip}
\csromancenter{หตฺถิงฺคปโญฺห ทสโม.}
\nitthitammidruled{\textbf{อุปจิกาวคฺโค จตุตฺโถ.}}
\csromancenter{\textbf{ตสฺสุทฺทานํ}}
\setlength{\csromangathapagewidth}{0pt}
\setlength{\csromangathawakwidth}{0pt}
\csromangathameasuregroup{อุปจิกา พิฬาโร จ,}{\csromangathabat{อุปจิกา พิฬาโร จ,}{อุนฺทูโร วิจฺฉิเกน จ.} \\ นกุโล สิงฺคาโล มิโค. \\ โครูโป วราโห หตฺถินา ทสาติ.}
\csromangathasetleft{อุปจิกา พิฬาโร จ,}
\csromangathagroup{\csromangathastanza{\csromangathabat{อุปจิกา พิฬาโร จ,}{อุนฺทูโร วิจฺฉิเกน จ.} \\ นกุโล สิงฺคาโล มิโค. \\ โครูโป วราโห หตฺถินา ทสาติ.}}

\csromanmatikaanchor{mtk.245}
\csromantocmarknopage{paragraph}{5. สีหวคฺค}{mtk.245}
\bookmark[dest={mtk.245},level=4]{5. สีหวคฺค}
\csromanheader{4}{5. \textbf{สีหวคฺค}}
\csromanmatikaanchor{mtk.246}
\csromantocmarkref{subparagraph}{1. สีหงฺคปญฺห}{mtk.246}
\bookmark[dest={mtk.246},level=5]{1. สีหงฺคปญฺห}
\csromanheader{5}{1. \textbf{สีหงฺคปญฺห}}
\proseitem{1}{\csromanfolio{385}ภนฺเต นาคเสน “สีหสฺส สตฺต องฺคานิ คเหตพฺพานี”ติ ยํ วเทสิ, กตมานิ ตานิ สตฺต องฺคานิ คเหตพฺพานีติ. ยถา มหาราช สีโห นาม เสต\-วิมล\-ปริสุทฺธ\-ปณฺฑโร, เอวเมว โข มหาราช โยคินา โยคาวจเรน เสต\-วิมล\-ปริสุทฺธ\-ปณฺฑร\-จิตฺเตน พฺยปคตกุกฺกุจฺเจน ภวิตพฺพํ. อิทํ มหาราช สีหสฺส ปฐมํ องฺคํ คเหตพฺพํ.}

\prose{ปุน จปรํ มหาราช สีโห จตุจรโณ วิกฺกนฺตจารี, เอวเมว โข มหาราช โยคินา โยคาวจเรน จตุริ\-ทฺธิปาท\-จรเณน ภวิตพฺพํ. อิทํ มหาราช สีหสฺส ทุติยํ องฺคํ คเหตพฺพํ.}

\prose{ปุน จปรํ สีโห อภิรูป\-รุจิร\-เกสรี, เอวเมว โข มหาราช โยคินา โยคาวจเรน อภิรูป\-รุจิร\-สีล\-เกสรินา ภวิตพฺพํ. อิทํ มหาราช สีหสฺส ตติยํ องฺคํ คเหตพฺพํ.}

\prose{ปุน จปรํ มหาราช สีโห ชีวิต\-ปริยาทาเนปิ น กสฺสจิ โอนมติ, เอวเมว โข มหาราช โยคินา โยคาวจเรน จีวร\-ปิณฺฑปาต\-เสนาสน\-คิลานปจฺจย\-เภสชฺช\-ปริกฺขาร\-ปริยาทาเนปิ น กสฺสจิ โอนมิตพฺพํ. อิทํ มหาราช สีหสฺส จตุตฺตํ องฺคํ คเหตพฺพํ.}

\prose{ปุน จปรํ มหาราช สีโห สปทาน\-ภกฺโข ยสฺมิํ โอกาเส นิปตติ, ตตฺเถว ยาวทตฺถํ ภกฺขยติ, น วรมํสํ วิจินาติ, เอวเมว โข มหาราช โยคินา โยคาวจเรน สปทาน\-ภกฺเขน ภวิตพฺพํ, น กุลานิ วิจินิตพฺพานิ, น ปุพฺพเคหํ หิตฺวา กุลานิ อุปสงฺกมิตพฺพานิ, น โภชนํ วิจินิตพฺพํ, ยสฺมิํ โอกาเส กพฬํ อาทียติ, ตสฺมิํ เยว โอกาเส ภุญฺชิตพฺพํ สรีร\-ยาปนตฺถํ\csromansharedfootnote{e009b26a033a7ce9}{สรีรยาปนมตฺตํ (สี, อิ)}, น วรโภชนํ วิจินิตพฺพํ. อิทํ มหาราช สีหสฺส ปญฺจมํ องฺคํ คเหตพฺพํ.}

\prose{ปุน จปรํ มหาราช สีโห อสนฺนิธิ\-ภกฺโข, สกิํ โคจรํ ภกฺขยิตฺวา น ปุน ตํ อุปคจฺฉติ, เอวเมว โข มหาราช โยคินา โยคาวจเรน \csromanfolio{386}อสนฺนิธิการปริโภคินา ภวิตพฺพํ. อิทํ มหาราช สีหสฺส ฉฏฺฐํ องฺคํ คเหตพฺพํ.}

\prosewithcloserruled{ปุน จปรํ มหาราช สีโห โภชนํ อลทฺธา น ปริตสฺสติ. ลทฺธาปิ โภชนํ อคธิโต\csromansharedfootnote{235dafb2128ac3dd}{อคถิโต (สี)} อมุจฺฉิโต อนชฺโฌสนฺโน ปริภุญฺชติ, เอวเมว โข มหาราช โยคินา โยคาวจเรน โภชนํ อลทฺธา น ปริตสฺสิตพฺพํ, ลทฺธาปิ โภชนํ อคธิเตน อมุจฺฉิเตน อนชฺโฌสนฺเนน อาทีนว\-ทสฺสาวินา นิสฺสรณ\-ปญฺเญน ปริภุญฺชิ\-ตพฺพํ. อิทํ มหาราช สีหสฺส สตฺตมํ องฺคํ คเหตพฺพํ. ภาสิตมฺเปตํ มหาราช ภควตา เทวาติเทเวน สํยุตฺตนิกาย\-วเร เถรํ มหากสฺสปํ ปริกิตฺตยมาเนน\csromandash{}\csromansymbolfootnote{*}{สํ 1. 399 ปิฏฺเฐ.}“สนฺตุฏฺโฐยํ ภิกฺขเว กสฺสโป อิตรีตเรน ปิณฺฑปาเตน, อิตรีตร\-ปิณฺฑปาต\-สนฺตุฏฺฐิยา จ วณฺณวาที, น จ ปิณฺฑปาตเหตุ อเนสนํ อปฺปติรูปํ อาปชฺชติ, อลทฺธา จ ปิณฺฑปาตํ น ปริตสฺสติ, ลทฺธา จ ปิณฺฑปาตํ อคธิโต อมุจฺฉิโต อนชฺโฌสนฺโน อาทีนวทสฺสาวี นิสฺสรณปญฺโญ ปริภุญฺชตี”ติ.}{สีหงฺคปโญฺห ปฐโม.}
\csromanmatikaanchor{mtk.247}
\csromantocmarkref{subparagraph}{2. จกฺกวากงฺคปญฺห}{mtk.247}
\bookmark[dest={mtk.247},level=5]{2. จกฺกวากงฺคปญฺห}
\csromanheader{5}{2. \textbf{จกฺกวาก}\-\textbf{งฺค}\-\textbf{ปญฺห}}
\proseitem{1}{ภนฺเต นาคเสน “จกฺกวากสฺส ตีณิ องฺคานิ คเหตพฺพานี”ติ ยํ วเทสิ, กตมานิ ตานิ ตีณิ องฺคานิ คเหตพฺพานีติ. ยถา มหาราช จกฺกวาโก ยาว ชีวิต\-ปริยาทานา ทุติยิกํ น วิชหติ, เอวเมว โข มหาราช โยคินา โยคาวจเรน ยาว ชีวิต\-ปริยาทานา โยนิโส มนสิกาโร น วิชหิตพฺโพ. อิทํ มหาราช จกฺกวากสฺส ปฐมํ องฺคํ คเหตพฺพํ.}

\prose{ปุน จปรํ มหาราช จกฺกวาโก เสวาลปณก\-ภกฺโข, เตน จ สนฺตุฏฺฐิํ อาปชฺชติ, ตาย จ สนฺตุฏฺฐิยา พเลน จ วณฺเณน จ น ปริหายติ, เอวเมว โข มหาราช โยคินา โยคาวจเรน ยถาลาภ\-สนฺโตโส กรณีโย, ยถา\-ลาภ\-สนฺตุฏฺโฐ โข มหาราช โยคี โยคาวจโร น ปริหายติ สีเลน, น ปริหายติ สมาธินา, \csromanfolio{387}น ปริหายติ ปญฺญาย, น ปริหายติ วิมุตฺติยา, น ปริหายติ วิมุตฺติ\-ญาณ\-ทสฺสเนน, น ปริหายติ สพฺเพหิ กุสเลหิ ธมฺเมหิ, อิทํ มหาราช จกฺกวากสฺส ทุติยํ องฺคํ คเหตพฺพํ.}

\proseitem{2}{ปุน จปรํ มหาราช จกฺกวาโก ปาเณ น วิเหฐยติ, เอวเมว โข มหาราช โยคินา โยคาวจเรน นิหิตทณฺเฑน นิหิตสตฺเถน ลชฺชินา ทยาปนฺเนน สพฺพปาณภูต\-หิตา\-นุกมฺปินา ภวิตพฺพํ. อิทํ มหาราช จกฺกวากสฺส ตติยํ องฺคํ คเหตพฺพํ. ภาสิตมฺเปตํ มหาราช ภควตา เทวาติเทเวน จกฺกวากชาตเก*\csromandash{}}

\setlength{\csromangathapagewidth}{0pt}
\setlength{\csromangathawakwidth}{0pt}
\csromangathameasuregroup{“โย น หนฺติ น ฆาเตติ, \\ เมตฺตํโส\textsuperscript{1} สพฺพภูเตสุ,}{\csromangathabat{“โย น หนฺติ น ฆาเตติ,}{น ชินาติ น ชาปเย.} \\ \csromangathabat{เมตฺตํโส\textsuperscript{1} สพฺพภูเตสุ,}{เวรํ ตสฺส น เกนจี”ติ.}}
\csromangathasetleft{“โย น หนฺติ น ฆาเตติ, \\ เมตฺตํโส\textsuperscript{1} สพฺพภูเตสุ,}
\csromangathagroupwithcloserruled{\csromangathastanza{\csromangathabat{“โย น หนฺติ น ฆาเตติ,}{น ชินาติ น ชาปเย.} \\ \csromangathabat{เมตฺตํโส\csromansharedfootnote{a30ed7eff8f270ae}{อหิํสา (สี, อิ)} สพฺพภูเตสุ,}{เวรํ ตสฺส น เกนจี”ติ.}}}{จกฺกวาก\-งฺค\-ปโญฺห ทุติโย.}
\csromanmatikaanchor{mtk.248}
\csromantocmarkref{subparagraph}{3. เปณาหิกงฺคปญฺห}{mtk.248}
\bookmark[dest={mtk.248},level=5]{3. เปณาหิกงฺคปญฺห}
\csromanheader{5}{3. \textbf{เปณาหิกงฺคปญฺห}}
\proseitem{3}{ภนฺเต นาคเสน “เปณาหิกาย เทฺว องฺคานิ คเหตพฺพานี”ติ ยํ วเทสิ, กตมานิ ตานิ เทฺว องฺคานิ คเหตพฺพานีติ. ยถา มหาราช เปณาหิกา สกปติมฺหิ อุสูยาย ฉาปเก น โปสยติ, เอวเมว โข มหาราช โยคินา โยคาวจเรน สกมเน\csromansharedfootnote{0fe938b9b1c1247c}{สกมโน (ก)} กิเลเส อุปฺปนฺเน อุสูยายิตพฺพํ, สติปฏฺฐาเนน สมฺมา\-สํวร\-สุสิเร ปกฺขิปิตฺวา มโนทฺวาเร กายคตาสติ ภาเวตพฺพา. อิทํ มหาราช เปณาหิกาย ปฐมํ องฺคํ คเหตพฺพํ.}

\prose{ปุน จปรํ มหาราช เปณาหิกา ปวเน ทิวสํ โคจรํ จริตฺวา สายํ ปกฺขิคณํ อุเปติ อตฺตโน คุตฺติยา, เอวเมว โข มหาราช โยคินา โยคาวจเรน เอกเกน ปวิเวกํ เสวิตพฺพํ สํโยชน\-ปริมุตฺติยา, ตตฺร รติํ อลภมาเนน อุปวาทภย\-ปริรกฺข\-ณาย สํฆํ โอสริตฺวา สํฆรกฺขิเตน วสิตพฺพํ. อิทํ มหาราช เปณาหิกาย ทุติยํ องฺคํ คเหตพฺพํ. ภาสิตมฺเปตํ มหาราช พฺรหฺมุนา สหมฺปตินา ภควโต สนฺติเก\csromandash{}}

\setlength{\csromangathapagewidth}{0pt}
\setlength{\csromangathawakwidth}{0pt}
\csromangathameasuregroup{\textsuperscript{*}“เสเวถ ปนฺตานิ เสนาสนานิ,}{\csromangathabat{\textsuperscript{*}“เสเวถ ปนฺตานิ เสนาสนานิ,}{จเรยฺย สํโยชนวิปฺปโมกฺขา.} \\ สเจ รติํ นาธิคจฺเฉยฺย ตตฺถ. \\ สํเฆ วเส รกฺขิตตฺโถ สตีมา”ติ.}
\csromangathasetleft{\textsuperscript{*}“เสเวถ ปนฺตานิ เสนาสนานิ,}
\csromangathagroupwithcloserruled{\csromangathastanza{\csromanfolio{388}\csromangathabat{\csromansymbolfootnote{*}{สํ 1. 156 ปิฏฺเฐ.}“เสเวถ ปนฺตานิ เสนาสนานิ,}{จเรยฺย สํโยชนวิปฺปโมกฺขา.} \\ สเจ รติํ นาธิคจฺเฉยฺย ตตฺถ. \\ สํเฆ วเส รกฺขิตตฺโถ สตีมา”ติ.}}{เปณาหิกงฺคปโญฺห ตติโย.}
\csromanmatikaanchor{mtk.249}
\csromantocmarkref{subparagraph}{4. ฆรกโปตงฺคปญฺห}{mtk.249}
\bookmark[dest={mtk.249},level=5]{4. ฆรกโปตงฺคปญฺห}
\csromanheader{5}{4. \textbf{ฆรกโปต}\-\textbf{งฺค}\-\textbf{ปญฺห}}
\proseitem{44}{ภนฺเต นาคเสน “ฆรกโปตสฺส เอกํ องฺคํ คเหตพฺพนฺ”ติ ยํ วเทสิ, กตมํ ตํ เอกํ องฺคํ คเหตพฺพนฺติ. ยถา มหาราช ฆรกโปโต ปรเคเห วสมาโน น เตสํ กิญฺจิ ภณฺฑสฺส นิมิตฺตํ คณฺหาติ. มชฺฌตฺโต วสติ สญฺญาพหุโล, เอวเมว โข มหาราช โยคินา โยคาวจเรน ปรกุลํ อุปคเตน ตสฺมิํ กุเล อิตฺถีนํ วา ปุริสานํ วา มญฺเจ วา ปีเฐ วา วตฺเถ วา อลงฺกาเร วา อุปโภเค วา ปริโภเค วา โภชน\-วิกตีสุ วา น นิมิตฺตํ คเหตพฺพํ, มชฺฌตฺเตน ภวิตพฺพํ, สมณสญฺญา ปจฺจุปฏฺฐเปตพฺพา. อิทํ มหาราช ฆรกโปตสฺส เอกํ องฺคํ คเหตพฺพํ. ภาสิตมฺเปตํ มหาราช ภควตา เทวาติเทเวน จูฬานารทชาตเก\csromandash{}}

\setlength{\csromangathapagewidth}{0pt}
\setlength{\csromangathawakwidth}{0pt}
\csromangathameasuregroup{“ปวิสิตฺวา ปรกุลํ, \\ มิตํ ขาเท มิตํ ภุญฺเช,}{\csromangathabat{“ปวิสิตฺวา ปรกุลํ,}{ปานตฺถํ โภชนาย วา\textsuperscript{1}.} \\ \csromangathabat{มิตํ ขาเท มิตํ ภุญฺเช,}{น จ รูเป มนํ กเร”ติ.}}
\csromangathasetleft{“ปวิสิตฺวา ปรกุลํ, \\ มิตํ ขาเท มิตํ ภุญฺเช,}
\csromangathagroupwithcloserruled{\csromangathastanza{\csromangathabat{“ปวิสิตฺวา ปรกุลํ,}{ปานตฺถํ โภชนาย วา\csromansharedfootnote{53b08fbe0b124c10}{ปาเนสุ โภชเนสุ วา (สี, อิ) ขุ 5. 263 ปิฏฺเฐ.}.} \\ \csromangathabat{มิตํ ขาเท มิตํ ภุญฺเช,}{น จ รูเป มนํ กเร”ติ.}}}{ฆรกโปต\-งฺค\-ปโญฺห จตุตฺโถ.}
\csromanmatikaanchor{mtk.250}
\csromantocmarkref{subparagraph}{5. อุลูกงฺคปญฺห}{mtk.250}
\bookmark[dest={mtk.250},level=5]{5. อุลูกงฺคปญฺห}
\csromanheader{5}{5. \textbf{อุลูกงฺคปญฺห}}
\proseitem{5}{ภนฺเต นาคเสน “อุลูกสฺส เทฺว องฺคานิ คเหตพฺพานี”ติ ยํ วเทสิ, กตมานิ ตานิ เทฺว องฺคานิ คเหตพฺพานีติ. ยถา มหาราช อุลูโก กาเกหิ ปฏิวิรุทฺโธ, รตฺติํ กากสํฆํ คนฺตฺวา พหูปิ กาเก หนติ, เอวเมว โข มหาราช โยคินา โยคาวจเรน อญฺญาเณน ปฏิวิรุทฺโธ กาตพฺโพ, เอเกน รโห นิสีทิตฺวา อญฺญาณํ สมฺปมทฺทิตพฺพํ, มูลโต ฉินฺทิตพฺพํ. อิทํ มหาราช อุลูกสฺส ปฐมํ องฺคํ คเหตพฺพํ.}

\csromanmatikaanchor{mtk.251}
\csromantocmarkref{subparagraph}{6. สตปตฺตงฺคปญฺห}{mtk.251}
\bookmark[dest={mtk.251},level=5]{6. สตปตฺตงฺคปญฺห}
\prosewithcloserruled{\csromanfolio{389}ปุน จปรํ มหาราช อุลูโก สุปฺปฏิสลฺลีโน โหติ, เอวเมว โข มหาราช โยคินา โยคาวจเรน ปฏิสลฺลานา\-ราเมน ภวิตพฺพํ ปฏิสลฺลาน\-รเตน. อิทํ มหาราช อุลูกสฺส ทุติยํ องฺคํ คเหตพฺพํ. ภาสิตมฺเปตํ มหาราช ภควตา เทวาติเทเวน สํยุตฺตนิกาย\-วเร\csromandash{}\csromansymbolfootnote{*}{สํ 3. 363 ปิฏฺเฐ.}“อิธ ภิกฺขเว ภิกฺขุ ปฏิสลฺลานา\-ราโม ปฏิสลฺลาน\-รโต “อิทํ ทุกฺขนฺ”ติ ยถาภูตํ ปชานาติ, “อยํ ทุกฺขสมุทโย”ติ ยถาภูตํ ปชานาติ, “อยํ ทุกฺขนิโรโธ”ติ ยถาภูตํ ปชานาติ, “อยํ ทุกฺขนิโรธ\-คามินี ปฏิปทา”ติ ยถาภูตํ ปชานาตี”ติ.}{อุลูกงฺคปโญฺห ปญฺจโม.}
\csromanheader{2}{6. \textbf{สตปตฺตงฺคปญฺหา}}
\proseitem{6}{ภนฺเต นาคเสน “สตปตฺตสฺส เอกํ องฺคํ คเหตพฺพนฺ”ติ ยํ วเทสิ, กตมํ ตํ เอกํ องฺคํ คเหตพฺพนฺติ. ยถา มหาราช สตปตฺโต รวิตฺวา ปเรสํ เขมํ วา ภยํ วา อาจิกฺขติ, เอวเมว โข มหาราช โยคินา โยคาวจเรน ปเรสํ ธมฺมํ เทสยมาเนน วินิปาตํ ภยโต ทสฺสยิตพฺพํ, นิพฺพานํ เขมโต ทสฺสยิตพฺพํ. อิทํ มหาราช สตปตฺตสฺส เอกํ องฺคํ คเหตพฺพํ. ภาสิตมฺเปตํ มหาราช เถเรน ปิณฺโฑล\-ภารทฺวาเชน\csromandash{}}

\setlength{\csromangathapagewidth}{0pt}
\setlength{\csromangathawakwidth}{0pt}
\csromangathameasuregroup{“นิรเย ภยสนฺตาสํ,}{\csromangathabat{“นิรเย ภยสนฺตาสํ,}{นิพฺพาเน วิปุลํ สุขํ.}}
\csromangathasetleft{“นิรเย ภยสนฺตาสํ,}
\csromangathagroup{\csromangathastanza{\csromangathabat{“นิรเย ภยสนฺตาสํ,}{นิพฺพาเน วิปุลํ สุขํ.}}}

\prosewithcloserruled{อุภยา\-เนตาน\-ตฺถานิ ทสฺเสตพฺพานิ โยคินา”ติ.}{สตปตฺตงฺค\-ปโญฺห ฉฏฺโฐ.}
\csromanmatikaanchor{mtk.252}
\csromantocmarkref{subparagraph}{7. วคฺคุลิงฺคปญฺห}{mtk.252}
\bookmark[dest={mtk.252},level=5]{7. วคฺคุลิงฺคปญฺห}
\csromanheader{5}{7. \textbf{วคฺคุ}\-\textbf{ลิงฺค}\-\textbf{ปญฺห}}
\proseitem{7}{ภนฺเต นาคเสน “วคฺคุลิสฺส เทฺว องฺคานิ คเหตพฺพานี”ติ ยํ วเทสิ, กตมานิ ตานิ เทฺว องฺคานิ คเหตพฺพานีติ. ยถา มหาราช วคฺคุลิ เคหํ ปวิสิตฺวา วิจริตฺวา นิกฺขมติ, น ตตฺถ ปลิพุทฺธติ, เอวเมว โข มหาราช โยคินา โยคาวจเรน คามํ ปิณฺฑาย ปวิสิตฺวา สปทานํ \csromanfolio{390}วิจริตฺวา ปฏิลทฺธ\-ลาเภน ขิปฺปเมว นิกฺขมิตพฺพํ, น ตตฺถ ปลิพุทฺเธน ภวิตพฺพํ, อิทํ มหาราช วคฺคุลิสฺส ปฐมํ องฺคํ คเหตพฺพํ.}

\prose{ปุน จปรํ มหาราช วคฺคุลิ ปรเคเห วสมาโน น เตสํ ปริหานิํ กโรติ, เอวเมว โข มหาราช โยคินา โยคาวจเรน กุลานิ อุปสงฺกมิตฺวา อติยาจนาย วา วิญฺญตฺติ\-พหุลตาย วา กายโทส\-พหุลตาย วา อติภาณิตาย วา สมาน\-สุข\-ทุกฺข\-ตาย วา น เตสํ โกจิ วิปฺปฏิสาโร กรณีโย, นปิ เตสํ มูลกมฺมํ ปริหาเปตพฺพํ, สพฺพถา วฑฺฒิ เยว อิจฺฉิตพฺพา. อิทํ มหาราช วคฺคุลิสฺส ทุติยํ องฺคํ คเหตพฺพํ. ภาสิตมฺเปตํ มหาราช ภควตา เทวาติเทเวน ทีฆนิกาย\-วเร ลกฺขณ\-สุตฺตนฺเต*\csromandash{}}

\setlength{\csromangathapagewidth}{0pt}
\setlength{\csromangathawakwidth}{0pt}
\csromangathameasuregroup{}{“สทฺธาย สีเลน สุเตน พุทฺธิยา, \\ จาเคน ธมฺเมน พหูหิ สาธุหิ. \\ ธเนน ธญฺเญน จ เขตฺตวตฺถุนา, \\ ปุตฺเตหิ ทาเรหิ จตุปฺปเทหิ จ. \\ ญาตีหิ มิตฺเตหิ จ พนฺธเวหิ, \\ พเลน วณฺเณน สุเขน จูภยํ. \\ กถํ น หาเยยฺยุํ ปเรติ อิจฺฉติ, \\ อตฺถสมิทฺธิญฺจ ปนาภิกงฺขตี”ติ.}
\csromangathameasurewak{“สทฺธาย สีเลน สุเตน พุทฺธิยา, \\ จาเคน ธมฺเมน พหูหิ สาธุหิ. \\ ธเนน ธญฺเญน จ เขตฺตวตฺถุนา, \\ ปุตฺเตหิ ทาเรหิ จตุปฺปเทหิ จ. \\ ญาตีหิ มิตฺเตหิ จ พนฺธเวหิ, \\ พเลน วณฺเณน สุเขน จูภยํ. \\ กถํ น หาเยยฺยุํ ปเรติ อิจฺฉติ, \\ อตฺถสมิทฺธิญฺจ ปนาภิกงฺขตี”ติ.}
\setlength{\csromangathaleftcol}{0pt}
\csromangathagroupwithcloserruled{\csromangathastanza{\csromangathawak{“สทฺธาย สีเลน สุเตน พุทฺธิยา, \\ จาเคน ธมฺเมน พหูหิ สาธุหิ. \\ ธเนน ธญฺเญน จ เขตฺตวตฺถุนา, \\ ปุตฺเตหิ ทาเรหิ จตุปฺปเทหิ จ.}}\csromangathastanza{\csromangathawak{ญาตีหิ มิตฺเตหิ จ พนฺธเวหิ, \\ พเลน วณฺเณน สุเขน จูภยํ. \\ กถํ น หาเยยฺยุํ ปเรติ อิจฺฉติ, \\ อตฺถสมิทฺธิญฺจ ปนาภิกงฺขตี”ติ.}}}{วคฺคุ\-ลิงฺค\-ปโญฺห สตฺตโม.}
\csromanmatikaanchor{mtk.253}
\csromantocmarkref{subparagraph}{8. ชลูกงฺคปญฺห}{mtk.253}
\bookmark[dest={mtk.253},level=5]{8. ชลูกงฺคปญฺห}
\csromanheader{5}{8. \textbf{ชลูกงฺคปญฺห}}
\proseitem{8}{ภนฺเต นาคเสน “ชลูกาย เอกํ องฺคํ คเหตพฺพนฺ”ติ ยํ วเทสิ, กตมํ ตํ เอกํ องฺคํ คเหตพฺพนฺติ. ยถา มหาราช ชลูกา ยตฺถ อลฺลียติ, ตตฺเถว ทฬฺหํ อลฺลียิตฺวา รุหิรํ ปิวติ, เอวเมว โข มหาราช โยคินา โยคาวจเรน ยสฺมิํ อารมฺมเณ จิตฺตํ อลฺลียติ, ตํ อารมฺมณํ วณฺณโต จ สณฺฐานโต จ ทิสโต จ โอกาสโต จ ปริจฺเฉทโต จ ลิงฺคโต จ นิมิตฺตโต จ ทฬฺหํ ปติฏฺฐาเปตฺวา เตเน\-วา\-รมฺมเณน วิมุตฺติรสม\-เสจนกํ ปาตพฺพํ. อิทํ มหาราช ชลูกาย เอกํ องฺคํ คเหตพฺพํ. ภาสิตมฺเปตํ มหาราช เถเรน อนุรุทฺเธน\csromandash{}}

\setlength{\csromangathapagewidth}{0pt}
\setlength{\csromangathawakwidth}{0pt}
\csromangathameasuregroup{“ปริสุทฺเธน จิตฺเตน, \\ เตน จิตฺเตน ปาตพฺพํ,}{\csromangathabat{“ปริสุทฺเธน จิตฺเตน,}{อารมฺมเณ ปติฏฺฐาย.} \\ \csromangathabat{เตน จิตฺเตน ปาตพฺพํ,}{วิมุตฺติรสมเสจนนฺ”ติ.}}
\csromangathasetleft{“ปริสุทฺเธน จิตฺเตน, \\ เตน จิตฺเตน ปาตพฺพํ,}
\csromangathagroupwithcloserruled{\csromangathastanza{\csromanfolio{391}\csromangathabat{“ปริสุทฺเธน จิตฺเตน,}{อารมฺมเณ ปติฏฺฐาย.} \\ \csromangathabat{เตน จิตฺเตน ปาตพฺพํ,}{วิมุตฺติรสมเสจนนฺ”ติ.}}}{ชลูกงฺคปโญฺห อฏฺฐโม.}
\csromanmatikaanchor{mtk.254}
\csromantocmarkref{subparagraph}{9. สปฺปงฺคปญฺห}{mtk.254}
\bookmark[dest={mtk.254},level=5]{9. สปฺปงฺคปญฺห}
\csromanheader{5}{9. \textbf{สปฺปงฺคปญฺห}}
\proseitem{9}{ภนฺเต นาคเสน “สปฺปสฺส ตีณิ องฺคานิ คเหตพฺพานี”ติ ยํ วเทสิ, กตมานิ ตานิ ตีณิ องฺคานิ คเหตพฺพานีติ. ยถา มหาราช สปฺโป อุเรน คจฺฉติ, เอวเมว โข มหาราช โยคินา โยคาวจเรน ปญฺญาย จริตพฺพํ, ปญฺญาย จรมานสฺส โข มหาราช โยคิโน จิตฺตํ ญาเย จรติ, วิลกฺขณํ วิวชฺเชติ, สลกฺขณํ ภาเวติ. อิทํ มหาราช สปฺปสฺส ปฐมํ องฺคํ คเหตพฺพํ.}

\prose{ปุน จปรํ มหาราช สปฺโป จรมาโน โอสธํ ปริวชฺเชนฺโต จรติ, เอวเมว โข มหาราช โยคินา โยคาวจเรน ทุจฺจริตํ ปริวชฺเชนฺเตน จริตพฺพํ. อิทํ มหาราช สปฺปสฺส ทุติยํ องฺคํ คเหตพฺพํ.}

\prose{ปุน จปรํ มหาราช สปฺโป มนุสฺเส ทิสฺวา ตปฺปติ\csromansharedfootnote{cd16a6da0ca1f0ff}{มนุสฺสํ ทิสฺวา กมฺปติ (ก)} โสจติ จินฺตยติ, เอวเมว โข มหาราช โยคินา โยคาวจเรน กุวิตกฺเก วิตกฺเกตฺวา อรติํ อุปฺปาทยิตฺวา ตปฺปิตพฺพํ โสจิตพฺพํ จินฺตยิตพฺพํ “ปมาเทน เม ทิวโส วีตินามิโต, น โส ปุน สกฺกา ลทฺธุนฺ”ติ. อิทํ มหาราช สปฺปสฺส ตติยํ องฺคํ คเหตพฺพํ. ภาสิตมฺเปตํ มหาราช ภควตา ภลฺลาฏิยชาตเก ทฺวินฺนํ กินฺนรานํ\csromandash{}}

\setlength{\csromangathapagewidth}{0pt}
\setlength{\csromangathawakwidth}{0pt}
\csromangathameasuregroup{}{“มเยกรตฺตํ\textsuperscript{1} วิปฺปวสิมฺห ลุทฺท, \\ อกามกา อญฺญมญฺญํ สรนฺตา. \\ ตเมกรตฺตํ\textsuperscript{1} อนุตปฺปมานา, \\ โสจาม ‘สา รตฺติ ปน นเหสฺสตี’ติ”.}
\csromangathameasurewak{“มเยกรตฺตํ\textsuperscript{1} วิปฺปวสิมฺห ลุทฺท, \\ อกามกา อญฺญมญฺญํ สรนฺตา. \\ ตเมกรตฺตํ\textsuperscript{1} อนุตปฺปมานา, \\ โสจาม ‘สา รตฺติ ปน นเหสฺสตี’ติ”.}
\setlength{\csromangathaleftcol}{0pt}
\csromangathagroup{\csromangathastanza{\csromangathawak{“มเยกรตฺตํ\csromansharedfootnote{a677fff8a2d5c714}{ยเมกรตฺติํ (สี, อิ) ขุ 5. 331 ปิฏฺเฐ.} วิปฺปวสิมฺห ลุทฺท, \\ อกามกา อญฺญมญฺญํ สรนฺตา. \\ ตเมกรตฺตํ\csromansharedfootnote{47290e3c9cebdd9b}{ตเมกรตฺติํ (สี, อิ)} อนุตปฺปมานา, \\ โสจาม ‘สา รตฺติ ปน นเหสฺสตี’ติ”.}}}

\vspace{\tipitakaparskip}
\csromancenter{สปฺปงฺคปโญฺห นวโม.}
\csromanmatikaanchor{mtk.255}
\csromantocmarkref{subparagraph}{10. อชครงฺคปญฺห}{mtk.255}
\bookmark[dest={mtk.255},level=5]{10. อชครงฺคปญฺห}
\csromanheader{5}{10. \textbf{อชคร}\-\textbf{งฺค}\-\textbf{ปญฺห}}
\proseitem{10}{\csromanfolio{392}ภนฺเต นาคเสน “อชครสฺส เอกํ องฺคํ คเหตพฺพนฺ”ติ ยํ วเทสิ, กตมํ ตํ เอกํ องฺคํ คเหตพฺพนฺติ. ยถา มหาราช อชคโร มหติมหากาโย พหูปิ ทิวเส อูนูทโร ทีนตโร กุจฺฉิปูรํ อาหารํ น ลภติ, อปริปุณฺโณ เยว ยาวเทว สรีร\-ยาปน\-มตฺตเกน ยาเปติ, เอวเมว โข มหาราช โยคิโน โยคาวจรสฺส ภิกฺขาจริย\-ปฺปสุตสฺส ปรปิณฺฑมุ\-ปคตสฺส ปรทินฺน\-ปฺปาฏิกงฺขิสฺส สยํคาห\-ปฺปฏิวิรตสฺส ทุลฺลภํ อุทรปริปูรํ อาหารํ, อปิ จ อตฺถวสิเกน กุลปุตฺเตน จตฺตาโร ปญฺจ อาโลเป อภุญฺชิตฺวา อวเสสํ อุทเกน ปริปูเรตพฺพํ. อิทํ มหาราช อชครสฺส เอกํ องฺคํ คเหตพฺพํ. ภาสิตมฺเปตํ มหาราช เถเรน สาริปุตฺเตน ธมฺม\-เสนาปตินา\csromandash{}}

\setlength{\csromangathapagewidth}{0pt}
\setlength{\csromangathawakwidth}{0pt}
\csromangathameasuregroup{\textsuperscript{*}“อลฺลํ สุกฺขํ วา ภุญฺชนฺโต, \\ อูนูทโร มิตาหาโร, \\ จตฺตาโร ปญฺจ อาโลเป, \\ อลํ ผาสุ วิหาราย,}{\csromangathabat{\textsuperscript{*}“อลฺลํ สุกฺขํ วา ภุญฺชนฺโต,}{น พาฬฺหํ สุหิโต สิยา.} \\ \csromangathabat{อูนูทโร มิตาหาโร,}{สโต ภิกฺขุ ปริพฺพเช.} \\ \csromangathabat{จตฺตาโร ปญฺจ อาโลเป,}{อภุตฺวา อุทกํ ปิเว.} \\ \csromangathabat{อลํ ผาสุ วิหาราย,}{ปหิตตฺตสฺส ภิกฺขุโน”ติ.}}
\csromangathasetleft{\textsuperscript{*}“อลฺลํ สุกฺขํ วา ภุญฺชนฺโต, \\ อูนูทโร มิตาหาโร, \\ จตฺตาโร ปญฺจ อาโลเป, \\ อลํ ผาสุ วิหาราย,}
\csromangathagroup{\csromangathastanza{\csromangathabat{\csromansymbolfootnote{*}{ขุ 2. 343 ปิฏฺเฐ.}“อลฺลํ สุกฺขํ วา ภุญฺชนฺโต,}{น พาฬฺหํ สุหิโต สิยา.} \\ \csromangathabat{อูนูทโร มิตาหาโร,}{สโต ภิกฺขุ ปริพฺพเช.}}\csromangathastanza{\csromangathabat{จตฺตาโร ปญฺจ อาโลเป,}{อภุตฺวา อุทกํ ปิเว.} \\ \csromangathabat{อลํ ผาสุ วิหาราย,}{ปหิตตฺตสฺส ภิกฺขุโน”ติ.}}}

\vspace{\tipitakaparskip}
\csromancenter{อชคร\-งฺค\-ปโญฺห ทสโม.}
\nitthitammidruled{\textbf{สีหวคฺโค ปญฺจโม.}}
\csromancenter{\textbf{ตสฺสุทฺทานํ}}
\setlength{\csromangathapagewidth}{0pt}
\setlength{\csromangathawakwidth}{0pt}
\csromangathameasuregroup{เกสรี จกฺกวาโก จ, \\ อุลูโก สตปตฺโต จ, \\ สปฺโป อชคโร เจว,}{\csromangathabat{เกสรี จกฺกวาโก จ,}{เปณาหิ ฆรกโปตโก.} \\ \csromangathabat{อุลูโก สตปตฺโต จ,}{วคฺคุลิ จ ชลูปิกา.} \\ \csromangathabat{สปฺโป อชคโร เจว,}{วคฺโค เตน ปวุจฺจตีติ.}}
\csromangathasetleft{เกสรี จกฺกวาโก จ, \\ อุลูโก สตปตฺโต จ, \\ สปฺโป อชคโร เจว,}
\csromangathagroup{\csromangathastanza{\csromangathabat{เกสรี จกฺกวาโก จ,}{เปณาหิ ฆรกโปตโก.} \\ \csromangathabat{อุลูโก สตปตฺโต จ,}{วคฺคุลิ จ ชลูปิกา.} \\ \csromangathabat{สปฺโป อชคโร เจว,}{วคฺโค เตน ปวุจฺจตีติ.}}}

\csromanmatikaanchor{mtk.256}
\csromantocmarknopage{paragraph}{6. มกฺกฏกวคฺค}{mtk.256}
\bookmark[dest={mtk.256},level=4]{6. มกฺกฏกวคฺค}
\csromanheader{4}{6. \textbf{มกฺกฏกวคฺค}}
\csromanmatikaanchor{mtk.257}
\csromantocmarkref{subparagraph}{1. ปนฺถมกฺกฏกงฺคปญฺห}{mtk.257}
\bookmark[dest={mtk.257},level=5]{1. ปนฺถมกฺกฏกงฺคปญฺห}
\csromanheader{5}{1. \textbf{ปนฺถ}\-\textbf{มกฺกฏก}\-\textbf{งฺค}\-\textbf{ปญฺห}}
\proseitem{1}{\csromanfolio{393}ภนฺเต นาคเสน “ปนฺถ\-มกฺกฏกสฺส เอกํ องฺคํ คเหตพฺพนฺ”ติ ยํ วเทสิ, กตมํ ตํ เอกํ องฺคํ คเหตพฺพนฺติ. ยถา มหาราช ปนฺถ\-มกฺกฏโก ปนฺเถ มกฺกฏ\-ชาล\-วิตานํ กตฺวา ยทิ ตตฺถ ชาลเก ลคฺคติ กิมิ วา มกฺขิกา วา ปฏงฺโค วา, ตํ คเหตฺวา ภกฺขยติ, เอวเมว โข มหาราช โยคินา โยคาวจเรน ฉสุ ทฺวาเรสุ สติปฏฺฐาน\-ชาล\-วิตานํ กตฺวา ยทิ ตตฺถ กิเลสมกฺขิกา พชฺฌนฺติ, ตตฺเถว ฆาเตตพฺพา. อิทํ มหาราช ปนฺถ\-มกฺกฏกสฺส เอกํ องฺคํ คเหตพฺพํ. ภาสิตมฺเปตํ มหาราช เถเรน อนุรุทฺเธน\csromandash{}}

\setlength{\csromangathapagewidth}{0pt}
\setlength{\csromangathawakwidth}{0pt}
\csromangathameasuregroup{“จิตฺตํ นิยเม ฉสุ ทฺวาเรสุ, \\ กิเลสา ตตฺถ ลคฺคา เจ,}{\csromangathabat{“จิตฺตํ นิยเม ฉสุ ทฺวาเรสุ,}{สติปฏฺฐานวรุตฺตเม.} \\ \csromangathabat{กิเลสา ตตฺถ ลคฺคา เจ,}{หนฺตพฺพา เต วิปสฺสินา”ติ.}}
\csromangathasetleft{“จิตฺตํ นิยเม ฉสุ ทฺวาเรสุ, \\ กิเลสา ตตฺถ ลคฺคา เจ,}
\csromangathagroupwithcloserruled{\csromangathastanza{\csromangathabat{“จิตฺตํ นิยเม ฉสุ ทฺวาเรสุ,}{สติปฏฺฐานวรุตฺตเม.} \\ \csromangathabat{กิเลสา ตตฺถ ลคฺคา เจ,}{หนฺตพฺพา เต วิปสฺสินา”ติ.}}}{ปนฺถ\-มกฺกฏก\-งฺค\-ปโญฺห ปฐโม.}
\csromanmatikaanchor{mtk.258}
\csromantocmarkref{subparagraph}{2. ถนสฺสิตทารกงฺคปญฺห}{mtk.258}
\bookmark[dest={mtk.258},level=5]{2. ถนสฺสิตทารกงฺคปญฺห}
\csromanheader{5}{2. \textbf{ถนสฺสิต}\-\textbf{ทารก}\-\textbf{งฺค}\-\textbf{ปญฺห}}
\proseitem{2}{ภนฺเต นาคเสน “ถนสฺสิต\-ทารกสฺส เอกํ องฺคํ คเหตพฺพนฺ”ติ ยํ วเทสิ, กตมํ ตํ เอกํ องฺคํ คเหตพฺพนฺติ. ยถา มหาราช ถนสฺสิต\-ทารโก สทตฺเถ ลคฺคติ, ขีรตฺถิโก โรทติ, เอวเมว โข มหาราช โยคินา โยคาวจเรน สทตฺเถ ลคฺคิตพฺพํ, สพฺพตฺถ ธมฺมญาเณน ภวิตพฺพํ, อุทฺเทเส ปริปุจฺฉาย สมฺมปฺปโยเค ปวิเวเก ครุสํวาเส กลฺยาณมิตฺต\-เสวเน. อิทํ มหาราช ถนสฺสิต\-ทารกสฺส เอกํ องฺคํ คเหตพฺพํ. ภาสิตมฺเปตํ มหาราช ภควตา เทวาติเทเวน ทีฆนิกาย\-วเร ปรินิพฺพาน\-สุตฺตนฺเต\csromandash{}}

\setlength{\csromangathapagewidth}{0pt}
\setlength{\csromangathawakwidth}{0pt}
\csromangathameasuregroup{}{“อิงฺฆ ตุเมฺห อานนฺท สารตฺเถ\textsuperscript{1} ฆฏถ, สารตฺเถ อนุยุญฺชถ, \\ สารตฺเถ อปฺปมตฺตา อาตาปิโน ปหิตตฺตา วิหรถา”ติ.}
\csromangathameasurewak{“อิงฺฆ ตุเมฺห อานนฺท สารตฺเถ\textsuperscript{1} ฆฏถ, สารตฺเถ อนุยุญฺชถ, \\ สารตฺเถ อปฺปมตฺตา อาตาปิโน ปหิตตฺตา วิหรถา”ติ.}
\setlength{\csromangathaleftcol}{0pt}
\csromangathagroup{\csromangathastanza{\csromangathawak{“อิงฺฆ ตุเมฺห อานนฺท สารตฺเถ\csromansharedfootnote{689cb935299b99a0}{สทตฺเถ (สี, อิ) ที 2. 117 ปิฏฺเฐ.} ฆฏถ, สารตฺเถ อนุยุญฺชถ, \\ สารตฺเถ อปฺปมตฺตา อาตาปิโน ปหิตตฺตา วิหรถา”ติ.}}}

\vspace{\tipitakaparskip}
\csromancenter{ถนสฺสิต\-ทารก\-งฺค\-ปโญฺห ทุติโย.}
\csromanmatikaanchor{mtk.259}
\csromantocmarkref{subparagraph}{3. จิตฺตกธรกุมฺมงฺคปญฺห}{mtk.259}
\bookmark[dest={mtk.259},level=5]{3. จิตฺตกธรกุมฺมงฺคปญฺห}
\csromanheader{5}{3. \textbf{จิตฺตก}\-\textbf{ธร}\-\textbf{กุมฺม}\-\textbf{งฺค}\-\textbf{ปญฺห}}
\proseitem{3}{\csromanfolio{394}ภนฺเต นาคเสน “จิตฺตก\-ธร\-กุมฺมสฺส\csromansharedfootnote{da173f235e9b8d82}{จิตฺตกถลกุมฺมสฺส (สี, อิ)} เอกํ องฺคํ คเหตพฺพนฺ”ติ ยํ วเทสิ, กตมํ ตํ เอกํ องฺคํ คเหตพฺพนฺติ. ยถา มหาราช จิตฺตก\-ธร\-กุมฺโม อุทกภยา อุทกํ ปริวชฺเชตฺวา วิจรติ, ตาย จ ปน อุทกํ ปริวชฺชนาย อายุนา น ปริหายติ, เอวเมว โข มหาราช โยคินา โยคาวจเรน ปมาเท ภยทสฺสาวินา ภวิตพฺพํ, อปฺปมาเท คุณวิเสส\-ทสฺสาวินา. ตาย จ ปน ภยทสฺสาวิ\-ตาย น ปริหายติ สามญฺญา, นิพฺพานสฺส สนฺติเก อุเปติ\csromansharedfootnote{4367f2a5c6c36f0c}{ฐเปติ (ก)}. อิทํ มหาราช จิตฺตก\-ธร\-กุมฺมสฺส เอกํ องฺคํ คเหตพฺพํ. ภาสิตมฺเปตํ มหาราช ภควตา เทวาติเทเวน ธมฺมปเท\csromandash{}}

\setlength{\csromangathapagewidth}{0pt}
\setlength{\csromangathawakwidth}{0pt}
\csromangathameasuregroup{\textsuperscript{*}“อปฺปมาทรโต ภิกฺขุ, \\ อภพฺโพ ปริหานาย,}{\csromangathabat{\textsuperscript{*}“อปฺปมาทรโต ภิกฺขุ,}{ปมาเท ภยทสฺสิ วา.} \\ \csromangathabat{อภพฺโพ ปริหานาย,}{นิพฺพานสฺเสว สนฺติเก”ติ.}}
\csromangathasetleft{\textsuperscript{*}“อปฺปมาทรโต ภิกฺขุ, \\ อภพฺโพ ปริหานาย,}
\csromangathagroupwithcloserruled{\csromangathastanza{\csromangathabat{\csromansymbolfootnote{*}{ขุ 1. 17 ปิฏฺเฐ.}“อปฺปมาทรโต ภิกฺขุ,}{ปมาเท ภยทสฺสิ วา.} \\ \csromangathabat{อภพฺโพ ปริหานาย,}{นิพฺพานสฺเสว สนฺติเก”ติ.}}}{จิตฺตก\-ธร\-กุมฺม\-งฺค\-ปโญฺห จตุตฺโถ.}
\csromanmatikaanchor{mtk.260}
\csromantocmarkref{subparagraph}{4. ปวนงฺคปญฺห}{mtk.260}
\bookmark[dest={mtk.260},level=5]{4. ปวนงฺคปญฺห}
\csromanheader{5}{4. \textbf{ปวนงฺคปญฺห}}
\proseitem{4}{ภนฺเต นาคเสน “ปวนสฺส ปญฺจ องฺคานิ คเหตพฺพานี”ติ ยํ วเทสิ, กตมานิ ตานิ ปญฺจ องฺคานิ คเหตพฺพานีติ. ยถา มหาราช ปวนํ นาม อสุจิชนํ ปฏิจฺฉาเทติ, เอวเมว โข มหาราช โยคินา โยคาวจเรน ปเรสํ อปรทฺธํ ขลิตํ ปฏิจฺฉาเทตพฺพํ น วิวริตพฺพํ. อิทํ มหาราช ปวนสฺส ปฐมํ องฺคํ คเหตพฺพํ.}

\prose{ปุน จปรํ มหาราช ปวนํ สุญฺญํ ปจุรชเนหิ, เอวเมว โข มหาราช โยคินา โยคาวจเรน ราค\-โทสโมหมาน\-ทิฏฺฐิชาเลหิ สพฺเพหิ จ กิเลเสหิ สุญฺเญน ภวิตพฺพํ. อิทํ มหาราช ปวนสฺส ทุติยํ องฺคํ คเหตพฺพํ.}

\prose{ปุน จปรํ มหาราช ปวนํ วิวิตฺตํ ชนสมฺพาธ\-รหิตํ, เอวเมว โข มหาราช โยคินา โยคาวจเรน ปาปเกหิ อกุสเลหิ ธมฺเมหิ อนริเยหิ ปวิวิตฺเตน ภวิตพฺพํ. อิทํ มหาราช ปวนสฺส ตติยํ องฺคํ คเหตพฺพํ.}

\prose{\csromanfolio{395}ปุน จปรํ มหาราช ปวนํ สนฺตํ ปริสุทฺธํ, เอวเมว โข มหาราช โยคินา โยคาวจเรน สนฺเตน ปริสุทฺเธน ภวิตพฺพํ, นิพฺพุเตน ปหีนมาเนน ปหีนมกฺเขน ภวิตพฺพํ. อิทํ มหาราช ปวนสฺส จตุตฺถํ องฺคํ คเหตพฺพํ.}

\prose{ปุน จปรํ มหาราช ปวนํ อริย\-ชน\-สํเสวิตํ, เอวเมว โข มหาราช โยคินา โยคาวจเรน อริย\-ชน\-สํเสวิเตน ภวิตพฺพํ. อิทํ มหาราช ปวนสฺส ปญฺจมํ องฺคํ คเหตพฺพํ. ภาสิตมฺเปตํ มหาราช ภควตา เทวาติเทเวน สํยุตฺตนิกาย\-วเร*\csromandash{}}

\setlength{\csromangathapagewidth}{0pt}
\setlength{\csromangathawakwidth}{0pt}
\csromangathameasuregroup{“ปวิวิตฺเตหิ อริเยหิ, \\ นิจฺจํ อารทฺธวีริเยหิ,}{\csromangathabat{“ปวิวิตฺเตหิ อริเยหิ,}{ปหิตตฺเตหิ ฌายิภิ.} \\ \csromangathabat{นิจฺจํ อารทฺธวีริเยหิ,}{ปณฺฑิเตหิ สหาวเส”ติ.}}
\csromangathasetleft{“ปวิวิตฺเตหิ อริเยหิ, \\ นิจฺจํ อารทฺธวีริเยหิ,}
\csromangathagroupwithcloserruled{\csromangathastanza{\csromangathabat{“ปวิวิตฺเตหิ อริเยหิ,}{ปหิตตฺเตหิ ฌายิภิ.} \\ \csromangathabat{นิจฺจํ อารทฺธ\-วีริเยหิ,}{ปณฺฑิเตหิ สหาวเส”ติ.}}}{ปวนงฺคปโญฺห จตุตฺโถ.}
\csromanmatikaanchor{mtk.261}
\csromantocmarkref{subparagraph}{5. รุกฺขงฺคปญฺห}{mtk.261}
\bookmark[dest={mtk.261},level=5]{5. รุกฺขงฺคปญฺห}
\csromanheader{5}{5. \textbf{รุกฺขงฺคปญฺห}}
\proseitem{5}{ภนฺเต นาคเสน “รุกฺขสฺส ตีณิ องฺคานิ คเหตพฺพานี”ติ ยํ วเทสิ, กตมานิ ตานิ ตีณิ องฺคานิ คเหตพฺพานีติ. ยถา มหาราช รุกฺโข นาม ปุปฺผ\-ผล\-ธโร, เอวเมว โข มหาราช โยคินา โยคาวจเรน วิมุตฺติ\-ปุปฺผ\-สามญฺญผล\-ธารินา ภวิตพฺพํ. อิทํ มหาราช รุกฺขสฺส ปฐมํ องฺคํ คเหตพฺพํ.}

\prose{ปุน จปรํ มหาราช รุกฺโข อุปคตานม\-นุปฺปวิฏฺฐานํ ชนานํ ฉายํ เทติ, เอวเมว โข มหาราช โยคินา โยคาวจเรน อุปคตานม\-นุปฺปวิฏฺฐานํ ปุคฺคลานํ อามิสปฺปฏิสนฺถาเรน วา ธมฺม\-ปฺปฏิสนฺถาเรน วา ปฏิสนฺถริตพฺพํ. อิทํ มหาราช รุกฺขสฺส ทุติยํ องฺคํ คเหตพฺพํ.}

\prose{ปุน จปรํ มหาราช รุกฺโข ฉายาเวมตฺตํ น กโรติ, เอวเมว โข มหาราช โยคินา โยคาวจเรน สพฺพสตฺเตสุ เวมตฺตตา น กาตพฺพา, โจรวธกปจฺจตฺถิเกสุปิ อตฺตนิปิ สมสมา เมตฺตาภาวนา กาตพฺพา “กินฺติ อิเม สตฺตา อเวรา อพฺยาปชฺชา\csromansharedfootnote{c2d66c02fdf53257}{อพฺยาปชฺฌา (สี)} อนีฆา สุขี อตฺตานํ ปริหเรยฺยูนฺ”ติ. อิทํ มหาราช รุกฺขสฺส ตติยํ องฺคํ คเหตพฺพํ. ภาสิตมฺเปตํ มหาราช เถเรน สาริปุตฺเตน ธมฺม\-เสนาปตินา\csromandash{}}

\setlength{\csromangathapagewidth}{0pt}
\setlength{\csromangathawakwidth}{0pt}
\csromangathameasuregroup{“วธเก เทวทตฺตมฺหิ, \\ ธนปาเล ราหุเล จ,}{\csromangathabat{“วธเก เทวทตฺตมฺหิ,}{โจเร องฺคุลิมาลเก.} \\ \csromangathabat{ธนปาเล ราหุเล จ,}{สพฺพตฺถ สมโก มุนี”ติ.}}
\csromangathasetleft{“วธเก เทวทตฺตมฺหิ, \\ ธนปาเล ราหุเล จ,}
\csromangathagroupwithcloserruled{\csromangathastanza{\csromanfolio{396}\csromangathabat{“วธเก เทวทตฺตมฺหิ,}{โจเร องฺคุลิมาลเก.} \\ \csromangathabat{ธนปาเล ราหุเล จ,}{สพฺพตฺถ สมโก มุนี”ติ.}}}{รุกฺขงฺคปโญฺห ปญฺจโม.}
\csromanmatikaanchor{mtk.262}
\csromantocmarkref{subparagraph}{6. เมฆงฺคปญฺห}{mtk.262}
\bookmark[dest={mtk.262},level=5]{6. เมฆงฺคปญฺห}
\csromanheader{5}{6. \textbf{เมฆงฺคปญฺห}}
\proseitem{6}{ภนฺเต นาคเสน “เมฆสฺส ปญฺจ องฺคานิ คเหตพฺพานี”ติ ยํ วเทสิ, กตมานิ ตานิ ปญฺจ องฺคานิ คเหตพฺพานีติ. ยถา มหาราช เมโฆ อุปฺปนฺนํ รโชชลฺลํ วูปสเมติ, เอวเมว โข มหาราช โยคินา โยคาวจเรน อุปฺปนฺนํ กิเลส\-รโชชลฺลํ วูปสเมตพฺพํ. อิทํ มหาราช เมฆสฺส ปฐมํ องฺคํ คเหตพฺพํ.}

\prose{ปุน จปรํ มหาราช เมโฆ ปถวิยา อุณฺหํ นิพฺพาเปติ, เอวเมว โข มหาราช โยคินา โยคาวจเรน เมตฺตาภาวนาย สเทวโก โลโก นิพฺพาเปตพฺโพ. อิทํ มหาราช เมฆสฺส ทุติยํ องฺคํ คเหตพฺพํ.}

\prose{ปุน จปรํ มหาราช เมโฆ สพฺพพีชานิ วิรุหาเปติ, เอวเมว โข มหาราช โยคินา โยคาวจเรน สพฺพสตฺตานํ สทฺธํ อุปฺปาเทตฺวา ตํ สทฺธาพีชํ ตีสุ สมฺปตฺตีสุ โรเปตพฺพํ, ทิพฺพ\-มานุสิกาสุ สุขสมฺปตฺตีสุ ยาว\-ปรมตฺถ\-นิพฺพานสุข\-สมฺปตฺติ. อิทํ มหาราชาช เมฆสฺส ตติยํ องฺคํ คเหตพฺพํ.}

\prose{ปุน จปรํ มหาราช เมโฆ อุตุโต สมุฏฺฐหิตฺวา ธรณิ\-ตล\-รุเห ติณรุกฺขลตาคุมฺพโอสธิวนปฺปตโย ปริรกฺขติ, เอวเมว โข มหาราช โยคินา โยคาวจเรน โยนิโส มนสิการํ นิพฺพตฺเตตฺวา เตน โยนิโส มนสิกาเรน สมณธมฺโม ปริรกฺขิตพฺโพ, โยนิโส มนสิการ\-มูลกา สพฺเพ กุสลา ธมฺมา. อิทํ มหาราช เมฆสฺส จตุตฺถํ องฺคํ คเหตพฺพํ.}

\prose{ปุน จปรํ มหาราช เมโฆ วสฺสมาโน นทิตฬาก\-โปกฺขรณิโย กนฺทรปทรสรโสพฺภอุทปานานิ จ ปริปูเรติ อุทกธาราหิ, เอวเมว โข มหาราช โยคินา โยคาวจเรน อาคม\-ปริยตฺติยา ธมฺมเมฆมภิวสฺสยิตฺวา อธิคม\-กามานํ มานสํ ปริปูร\-ยิตพฺพํ. อิทํ มหาราช \csromanfolio{397}เมฆสฺส ปญฺจมํ องฺคํ คเหตพฺพํ. ภาสิตมฺเปตํ มหาราช เถเรน สาริปุตฺเตน ธมฺม\-เสนาปตินา\csromandash{}}

\setlength{\csromangathapagewidth}{0pt}
\setlength{\csromangathawakwidth}{0pt}
\csromangathameasuregroup{“โภธเนยฺยํ ชนํ ทิสฺวา, \\ ขเณน อุปคนฺตฺวาน,}{\csromangathabat{“โภธเนยฺยํ ชนํ ทิสฺวา,}{สตสหสฺเสปิ โยชเน.} \\ \csromangathabat{ขเณน อุปคนฺตฺวาน,}{โพเธติ ตํ มหามุนี”ติ.}}
\csromangathasetleft{“โภธเนยฺยํ ชนํ ทิสฺวา, \\ ขเณน อุปคนฺตฺวาน,}
\csromangathagroupwithcloserruled{\csromangathastanza{\csromangathabat{“โภธเนยฺยํ ชนํ ทิสฺวา,}{สตสหสฺเสปิ โยชเน.} \\ \csromangathabat{ขเณน อุปคนฺตฺวาน,}{โพเธติ ตํ มหามุนี”ติ.}}}{เมฆงฺคปโญฺห ฉฏฺโฐ.}
\csromanmatikaanchor{mtk.263}
\csromantocmarkref{subparagraph}{7. มณิรตนงฺคปญฺห}{mtk.263}
\bookmark[dest={mtk.263},level=5]{7. มณิรตนงฺคปญฺห}
\csromanheader{5}{7. \textbf{มณิรตน}\-\textbf{งฺค}\-\textbf{ปญฺห}}
\proseitem{7}{ภนฺเต นาคเสน “มณิรตนสฺส ตีณิ องฺคานิ คเหตพฺพานี”ติ ยํ วเทสิ, กตมานิ ตานิ ตีณิ องฺคานิ คเหตพฺพานีติ. ยถา มหาราช มณิรตนํ เอกนฺต\-ปริสุทฺธํ, เอวเมว โข มหราช โยคินา โยคาวจเรน เอกนฺต\-ปริสุทฺธาชีเวน ภวิตพฺพํ. อิทํ มหาราช มณิรตนสฺส ปฐมํ องฺคํ คเหตพฺพํ.}

\prose{ปุน จปรํ มหาราช มณิรตนํ น เกนจิ สทฺธิํ มิสฺสียติ, เอวเมว โข มหาราช โยคินา โยคาวจรน ปาเปหิ ปาปสหาเยหิ สทฺธิํ น มิสฺสิตพฺพํ. อิทํ มหาราช มณิรตนสฺส ทุติยํ องฺคํ คเหตพฺพํ.}

\prose{ปุน จปรํ มหาราช มณิรตนํ ชาติรตเนหิ โยชียติ, เอวเมว โข มหาราช โยคินา โยคาวจเรน อุตฺตม\-วร\-ชาติมนฺเตหิ สทฺธิํ สํวสิตพฺพํ, ปฏิปนฺนกผลฬเสกฺขผลสมงฺคีหิ โสตาปนฺนสกทาคามิอนาคามิ- อรหนฺตเตวิชฺชฉฬภิญฺญสมณมณิรตเนหิ สทฺธิํ สํวสิตพฺพํ. อิทํ มหาราช มณิรตนสฺส ตติยํ องฺคํ คเหตพฺพํ. ภาสิตมฺเปตํ มหาราช ภควตา เทวาติเทเวน สุตฺตนิปาเต\csromandash{}}

\setlength{\csromangathapagewidth}{0pt}
\setlength{\csromangathawakwidth}{0pt}
\csromangathameasuregroup{\textsuperscript{*}“สุทฺธา สุทฺเธหิ สํวาสํ, \\ ตโต สมคฺคา นิปกา,}{\csromangathabat{\textsuperscript{*}“สุทฺธา สุทฺเธหิ สํวาสํ,}{กปฺปยโวฺห ปติสฺสตา.} \\ \csromangathabat{ตโต สมคฺคา นิปกา,}{ทุกฺขสฺสนฺตํ กริสฺสถา”ติ.}}
\csromangathasetleft{\textsuperscript{*}“สุทฺธา สุทฺเธหิ สํวาสํ, \\ ตโต สมคฺคา นิปกา,}
\csromangathagroupwithcloserruled{\csromangathastanza{\csromangathabat{\csromansymbolfootnote{*}{ขุ 1. 322 ปิฏฺเฐ.}“สุทฺธา สุทฺเธหิ สํวาสํ,}{กปฺปยโวฺห ปติสฺสตา.} \\ \csromangathabat{ตโต สมคฺคา นิปกา,}{ทุกฺขสฺสนฺตํ กริสฺสถา”ติ.}}}{มณิรตน\-ปโญฺห สตฺตโม.}
\csromanmatikaanchor{mtk.264}
\csromantocmarkref{subparagraph}{8. มาควิกงฺคปญฺห}{mtk.264}
\bookmark[dest={mtk.264},level=5]{8. มาควิกงฺคปญฺห}
\csromanheader{5}{8. \textbf{มาควิก}\-\textbf{งฺค}\-\textbf{ปญฺห}}
\csromanmatikaanchor{mtk.265}
\csromantocmarkref{subparagraph}{9. ภาฬิสิกงฺคปญฺห}{mtk.265}
\bookmark[dest={mtk.265},level=5]{9. ภาฬิสิกงฺคปญฺห}
\proseitem{8}{ภนฺเต นาคเสน “มาควิกสฺส จตฺตาริ องฺคานิ คเหตพฺพานี”ติ ยํ วเทสิ, กตมานิ ตานิ จตฺตาริ องฺคานิ คเหตพฺพานีติ. ยถา มหาราช \csromanfolio{398}มาควิโก อปฺปมิทฺโธ โหติ, เอวเมว โข มหาราช โยคินา โยคาวจเรน อปฺปมิทฺเธน ภวิตพฺพํ. อิทํ มหาราช มาควิกสฺส ปฐมํ องฺคํ คเหตพฺพํ.}

\prose{ปุน จปรํ มหาราช มาควิโก มิเคสุ เยว จิตฺตํ อุปนิพนฺธติ, เอวเมว โข มหาราช โยคินา โยคาวจเรน อารมฺมเณสุ เยว จิตฺตํ อุปนิพนฺธิตพฺพํ. อิทํ มหาราช มาควิกสฺส ทุติยํ องฺคํ คเหตพฺพํ.}

\prose{ปุน จปรํ มหาราช มาควิโก กาลํ กมฺมสฺส ชานาติ, เอวเมว โข มหาราช โยคินา โยคาวจเรน ปฏิสลฺลานสฺส กาโล ชานิตพฺโพ “อยํ กาโล ปฏิสลฺลานสฺส, อยํ กาโล นิกฺขมนายา”ติ. อิทํ มหาราช มาควิกสฺส ตติยํ องฺคํ คเหตพฺพํ.}

\prose{ปุน จปรํ มหาราช มาควิโก มิคํ ทิสฺวา หาสม\-ภิ\-ชเนติ “อิมํ ลจฺฉามี”ติ, เอวเมว โข มหาราช โยคินา โยคาวจเรน อารมฺมเณ อภิรมิตพฺพํ, หาสม\-ภิ\-ชเนตพฺพํ “อุตฺตริํ วิเสสมธิคจฺฉิสฺสามี”ติ. อิทํ มหาราช มาควิกสฺส จตุตฺถํ องฺคํ คเหตพฺพํ. ภาสิตมฺเปตํ มหาราช เถเรน โมฆราเชน\csromandash{}}

\setlength{\csromangathapagewidth}{0pt}
\setlength{\csromangathawakwidth}{0pt}
\csromangathameasuregroup{“อารมฺมเณ ลภิตฺวาน, \\ ภิยฺโย หาโส ชเนตพฺโพ,}{\csromangathabat{“อารมฺมเณ ลภิตฺวาน,}{ปหิตตฺเตน ภิกฺขุนา.} \\ \csromangathabat{ภิยฺโย หาโส ชเนตพฺโพ,}{อธิคจฺฉิสฺสามิ อุตฺตรินฺ”ติ.}}
\csromangathasetleft{“อารมฺมเณ ลภิตฺวาน, \\ ภิยฺโย หาโส ชเนตพฺโพ,}
\csromangathagroupwithcloserruled{\csromangathastanza{\csromangathabat{“อารมฺมเณ ลภิตฺวาน,}{ปหิตตฺเตน ภิกฺขุนา.} \\ \csromangathabat{ภิยฺโย หาโส ชเนตพฺโพ,}{อธิคจฺฉิสฺสามิ อุตฺตรินฺ”ติ.}}}{มาควิก\-งฺค\-ปโญฺห อฏฺฐโม.}
\csromanheader{2}{9. \textbf{พาฬิสิก}\-\textbf{งฺค}\-\textbf{ปญฺห}}
\proseitem{9}{ภนฺเต นาคเสน “พาฬิสิกสฺส เทฺว องฺคานิ คเหตพฺพานี”ติ ยํ วเทสิ, กตมานิ ตานิ เทฺว องฺคานิ คเหตพฺพานีติ. ยถา มหาราช พาฬิสิโก พฬิเสน มจฺเฉ อุทฺธรติ, เอวเมว โข มหาราช โยคินา โยคาวจเรน ญาเณน อุตฺตริํ สามญฺญผลานิ อุทฺธริตพฺพานิ. อิทํ มหาราช พาฬิสิกสฺส ปฐมํ องฺคํ คเหตพฺพํ.}

\prose{\csromanfolio{399}ปุน จปรํ มหาราช พาฬิสิโก ปริตฺตกํ วธิตฺวา วิปุลํ ลาภม\-ธิคจฺฉติ, เอวเมว โข มหาราช โยคินา โยคาวจเรน ปริตฺต\-โลกามิส\-มตฺตํ ปริจฺจชิตพฺพํ, โลกามิสมตฺตํ มหาราช ปริจฺจชิตฺวา โยคี โยคาวจโร วิปุลํ สามญฺญผลํ อธิคจฺฉติ. อิทํ มหาราช พาฬิสิกสฺส ทุติยํ องฺคํ คเหตพฺพํ. ภาสิตมฺเปตํ มหาราช เถเรน ราหุเลน\csromandash{}}

\setlength{\csromangathapagewidth}{0pt}
\setlength{\csromangathawakwidth}{0pt}
\csromangathameasuregroup{“สุญฺญตญฺจานิมิตฺตญฺจ, \\ จตุโร ผเล ฉฬภิญฺญา,}{\csromangathabat{“สุญฺญตญฺจานิมิตฺตญฺจ,}{วิโมกฺขญฺจาปฺปณิหิตํ.} \\ \csromangathabat{จตุโร ผเล ฉฬภิญฺญา,}{จชิตฺวา โลกามิสํ ลเภ”ติ.}}
\csromangathasetleft{“สุญฺญตญฺจานิมิตฺตญฺจ, \\ จตุโร ผเล ฉฬภิญฺญา,}
\csromangathagroupwithcloserruled{\csromangathastanza{\csromangathabat{“สุญฺญตญฺจานิมิตฺตญฺจ,}{วิโมกฺขญฺจาปฺปณิหิตํ.} \\ \csromangathabat{จตุโร ผเล ฉฬภิญฺญา,}{จชิตฺวา โลกามิสํ ลเภ”ติ.}}}{พาฬิสิก\-งฺค\-ปโญฺห นวโม.}
\csromanmatikaanchor{mtk.266}
\csromantocmarkref{subparagraph}{10. ตจฺฉกงฺคปญฺห}{mtk.266}
\bookmark[dest={mtk.266},level=5]{10. ตจฺฉกงฺคปญฺห}
\csromanheader{5}{10. \textbf{ตจฺฉก}\-\textbf{งฺค}\-\textbf{ปญฺห}}
\proseitem{10}{ภนฺเต นาคเสน “ตจฺฉกสฺส เทฺว องฺคานิ คเหตพฺพานี”ติ ยํ วเทสิ, กตมานิ ตานิ เทฺว องฺคานิ คเหตพฺพานีติ. ยถา มหาราช ตจฺฉโก กาฬสุตฺตํ อนุโลเมตฺวา รุกฺขํ ตจฺฉติ, เอวเมว โข มหาราช โยคินา โยคาวจเรน ชินสาสนมนุโลมยิตฺวา สีลปถวิยํ ปติฏฺฐหิตฺวา สทฺธาหตฺเถน ปญฺญาวาสิํ คเหตฺวา กิเลสา ตจฺเฉตพฺพา. อิทํ มหาราช ตจฺฉกสฺส ปฐมํ องฺคํ คเหตพฺพํ.}

\prose{ปุน จปรํ มหาราช ตจฺฉโก เผคฺคุํ อปหริตฺวา สารมาทิยติ, เอวเมว โข มหาราช โยคินา โยคาวจเรน สสฺสตํ อุจฺเฉทํ ตํ ชีวํ ตํ สรีรํ อญฺญํ ชีวํ อญฺญํ สรีรํ ตทุตฺตมํ อญฺญทุตฺตมํ อกตมภพฺพํ อปุริสการํ อพฺรหฺมจริย\-วาสํ สตฺตวินาสํ นวสตฺต\-ปาตุภาวํ สงฺขาร\-สสฺสต\-ภาวํ โย กโรติ, โส ปฏิสํเวเทติ, อญฺโญ กโรติ, อญฺโญ ปฏิสํเวเทติ, กมฺมผล\-ทสฺสนา จ กิริย\-ผล\-ทิฏฺฐิ จ อิติ เอวรูปานิ เจว อญฺญานิ จ วิวาทปถานิ อปเนตฺวา สงฺขารานํ สภาวํ ปรมสุญฺญตํ นิรีห\-นิชฺชีวตํ\csromansharedfootnote{4b02a00f7ab33ce7}{นิสตฺตนิชีวตํ (ก)} อจฺจนฺตํ สุญฺญตํ อาทิยิตพฺพํ. อิทํ มหาราช ตจฺฉกสฺส \csromanfolio{400}ทุติยํ องฺคํ คเหตพฺพํ. ภาสิตมฺเปตํ มหาราช ภควตา เทวาติเทเวน สุตฺตนิปาเต*\csromandash{}}

\setlength{\csromangathapagewidth}{0pt}
\setlength{\csromangathawakwidth}{0pt}
\csromangathameasuregroup{“การณฺฑวํ นิทฺธมถ, \\ ตโต ปลาเป วาเหถ, \\ นิทฺธมิตฺวาน ปาปิจฺเฉ, \\ สุทฺธา สุทฺเธหิ สํวาสํ, \\ ตโต สมคฺคา นิปกา,}{\csromangathabat{“การณฺฑวํ นิทฺธมถ,}{กสมฺพุํ อปกสฺสถ.} \\ \csromangathabat{ตโต ปลาเป วาเหถ,}{อสฺสมเณ สมณมานิเน.} \\ \csromangathabat{นิทฺธมิตฺวาน ปาปิจฺเฉ,}{ปาปอาจารโคจเร.} \\ \csromangathabat{สุทฺธา สุทฺเธหิ สํวาสํ,}{กปฺปยโวฺห ปติสฺสตา.} \\ \csromangathabat{ตโต สมคฺคา นิปกา,}{ทุกฺขสฺสนฺตํ กริสฺสถา”ติ.}}
\csromangathasetleft{“การณฺฑวํ นิทฺธมถ, \\ ตโต ปลาเป วาเหถ, \\ นิทฺธมิตฺวาน ปาปิจฺเฉ, \\ สุทฺธา สุทฺเธหิ สํวาสํ, \\ ตโต สมคฺคา นิปกา,}
\csromangathagroup{\csromangathastanza{\csromangathabat{“การณฺฑวํ นิทฺธมถ,}{กสมฺพุํ อปกสฺสถ.} \\ \csromangathabat{ตโต ปลาเป วาเหถ,}{อสฺสมเณ สมณมานิเน.}}\csromangathastanza{\csromangathabat{นิทฺธมิตฺวาน ปาปิจฺเฉ,}{ปาป\-อาจาร\-โคจเร.} \\ \csromangathabat{สุทฺธา สุทฺเธหิ สํวาสํ,}{กปฺปยโวฺห ปติสฺสตา.} \\ \csromangathabat{ตโต สมคฺคา นิปกา,}{ทุกฺขสฺสนฺตํ กริสฺสถา”ติ.}}}

\vspace{\tipitakaparskip}
\csromancenter{ตจฺฉก\-งฺค\-ปโญฺห ทสโม.}
\nitthitammidruled{\textbf{มกฺกฏกวคฺโค ฉฏฺโฐ.}}
\csromancenter{\textbf{ตสฺสุทฺทานํ}}
\setlength{\csromangathapagewidth}{0pt}
\setlength{\csromangathawakwidth}{0pt}
\csromangathameasuregroup{มกฺกโฏ ทารโก กุมฺโม, \\ เมโฆ มณิ มาควิโก,}{\csromangathabat{มกฺกโฏ ทารโก กุมฺโม,}{วนํ รุกฺโข จ ปญฺจโม.} \\ \csromangathabat{เมโฆ มณิ มาควิโก,}{พาฬิสี ตจฺฉเกน จาติ.}}
\csromangathasetleft{มกฺกโฏ ทารโก กุมฺโม, \\ เมโฆ มณิ มาควิโก,}
\csromangathagroup{\csromangathastanza{\csromangathabat{มกฺกโฏ ทารโก กุมฺโม,}{วนํ รุกฺโข จ ปญฺจโม.} \\ \csromangathabat{เมโฆ มณิ มาควิโก,}{พาฬิสี ตจฺฉเกน จาติ.}}}

\csromanmatikaanchor{mtk.267}
\csromantocmarknopage{paragraph}{7. กุมฺภวคฺค}{mtk.267}
\bookmark[dest={mtk.267},level=4]{7. กุมฺภวคฺค}
\csromanheader{4}{7. \textbf{กุมฺภวคฺค}}
\csromanmatikaanchor{mtk.268}
\csromantocmarkref{subparagraph}{1. กุมฺภงฺคปญฺห}{mtk.268}
\bookmark[dest={mtk.268},level=5]{1. กุมฺภงฺคปญฺห}
\csromanheader{5}{1. \textbf{กุมฺภงฺคปญฺห}}
\proseitem{1}{\csromanfolio{401}ภนฺเต นาคเสน “กุมฺภสฺส เอกํ องฺคํ คเหตพฺพนฺ”ติ ยํ วเทสิ, กตมํ ตํ เอกํ องฺคํ คเหตพฺพนฺติ. ยถา มหาราช กุมฺโภ สมฺปุณฺโณ น สณติ, เอวเมว โข มหาราช โยคินา โยคาวจเรน อาคเม อธิคเม ปริยตฺติยํ สามญฺเญ ปารมิํ ปตฺวา น สณิตพฺพํ, น เตน มาโน กรณีโย, น ทพฺโพ\csromansharedfootnote{2da831ac26dde604}{ทปฺโป (สี)} ทสฺเสตพฺโพ, นิหตมาเนน นิหตทพฺเพน ภวิตพฺพํ, อุชุเกน อมุขเรน อวิกตฺถินา. อิทํ มหาราช กุมฺภสฺส เอกํ องฺคํ คเหตพฺพํ. ภาสิตมฺเปตํ มหาราช ภควตา เทวาติเทเนน สุตฺตนิปาเต*\csromandash{}}

\setlength{\csromangathapagewidth}{0pt}
\setlength{\csromangathawakwidth}{0pt}
\csromangathameasuregroup{“ยทูนกํ ตํ สณติ, \\ อฑฺฒกุมฺภูปโม\textsuperscript{1} พาโล,}{\csromangathabat{“ยทูนกํ ตํ สณติ,}{ยํ ปูรํ สนฺตเมว ตํ.} \\ \csromangathabat{อฑฺฒกุมฺภูปโม\textsuperscript{1} พาโล,}{รหโท ปูโรว ปณฺฑิโต”ติ.}}
\csromangathasetleft{“ยทูนกํ ตํ สณติ, \\ อฑฺฒกุมฺภูปโม\textsuperscript{1} พาโล,}
\csromangathagroupwithcloserruled{\csromangathastanza{\csromangathabat{“ยทูนกํ ตํ สณติ,}{ยํ ปูรํ สนฺตเมว ตํ.} \\ \csromangathabat{อฑฺฒกุมฺภูปโม\csromansharedfootnote{fff168db9b2a6c94}{ริตฺตกุมฺภูปโม (สี)} พาโล,}{รหโท ปูโรว ปณฺฑิโต”ติ.}}}{กุมฺภงฺคปโญฺห ปฐโม.}
\csromanmatikaanchor{mtk.269}
\csromantocmarkref{subparagraph}{2. กาฬายสงฺคปญฺห}{mtk.269}
\bookmark[dest={mtk.269},level=5]{2. กาฬายสงฺคปญฺห}
\csromanheader{5}{2. \textbf{กาฬาย}\-\textbf{สงฺค}\-\textbf{ปญฺห}}
\proseitem{2}{ภนฺเต นาคเสน “กาฬายสสฺส\csromansharedfootnote{621e434f4b551ac2}{กาฬหํสสฺส (ก)} เทฺว องฺคานิ คเหตพฺพานี”ติ ยํ วเทสิ, กตมานิ ตานิ เทฺว องฺคานิ คเหตพฺพานีติ. ยถา มหาราช กาฬายโส สุปีโต\csromansharedfootnote{a009ca54dc9d83b2}{สุถิโต (ก)} วมติ\csromansharedfootnote{61a079c492501b68}{วหติ (สฺยา, ก)}, เอวเมว โข มหาราช โยคิโน โยคาวจรสฺส มานสํ โยนิโส มนสิกาเรน\csromansharedfootnote{cefceb75e6433db5}{โยนิโส มนสิกาเร (สี, สฺยา, ก)} อปีตํ วมติ. อิทํ มหาราช กาฬายสสฺส ปฐมํ องฺคํ คเหตพฺพํ.}

\prose{ปุน จปรํ มหาราช กาฬายโส สกิํ ปีตํ อุทกํ น วมติ, เอวเมว โข มหาราช โยคินา โยคาวจเรน โย สกิํ อุปฺปนฺโน ปสาโท, น ปุน โส วมิตพฺโพ “อุฬาโร โส ภควา สมฺมาสมฺพุทฺโธ, สฺวากฺขาโต ธมฺโม, สุปฺปฏิปนฺโน สํโฆ”ติ. รูปํ อนิจฺจํ, เวทนา อนิจฺจา, สญฺญา อนิจฺจา, สงฺขารา อนิจฺจา, วิญฺญาณํ อนิจฺจนฺติ ยํ สกิํ อุปฺปนฺนํ ญาณํ, น ปุน ตํ วมิตพฺพํ. \csromanfolio{402}อิทํ มหาราช กาฬายสสฺส ทุติยํ องฺคํ คเหตพฺพํ. ภาสิตมฺเปตํ มหาราช ภควตา เทวาติเทเวน\csromandash{}}

\setlength{\csromangathapagewidth}{0pt}
\setlength{\csromangathawakwidth}{0pt}
\csromangathameasuregroup{}{“ทสฺสนมฺหิ ปริโสธิโต\textsuperscript{1} นโร, \\ อริยธมฺเม นิยโต วิเสสคู. \\ นปฺปเวธติ อเนกภาคโส, \\ สพฺพโส จ มุขภาวเมว โส”ติ.}
\csromangathameasurewak{“ทสฺสนมฺหิ ปริโสธิโต\textsuperscript{1} นโร, \\ อริยธมฺเม นิยโต วิเสสคู. \\ นปฺปเวธติ อเนกภาคโส, \\ สพฺพโส จ มุขภาวเมว โส”ติ.}
\setlength{\csromangathaleftcol}{0pt}
\csromangathagroupwithcloserruled{\csromangathastanza{\csromangathawak{“ทสฺสนมฺหิ ปริโสธิโต\csromansharedfootnote{39e9417b32c6475e}{ปริโสธิเต (สี, ก)} นโร, \\ อริยธมฺเม นิยโต วิเสสคู. \\ นปฺปเวธติ อเนกภาคโส, \\ สพฺพโส จ มุขภาวเมว โส”ติ.}}}{กาฬาย\-สงฺค\-ปโญฺห ทุติโย.}
\csromanmatikaanchor{mtk.270}
\csromantocmarkref{subparagraph}{3. ฉตฺตงฺคปญฺห}{mtk.270}
\bookmark[dest={mtk.270},level=5]{3. ฉตฺตงฺคปญฺห}
\csromanheader{5}{3. \textbf{ฉตฺตงฺคปญฺห}}
\proseitem{3}{ภนฺเต นาคเสน “ฉตฺตสฺส ตีณิ องฺคานิ คเหตพฺพานี”ติ ยํ วเทสิ, กตามานิ ตานิ ตีณิ องฺคานิ คเหตพฺพานีติ. ยถา มหาราช ฉตฺตํ อุปริ มุทฺธนิ จรติ, เอวเมว โข มหาราช โยคินา โยคาวจเรน กิเลสานํ อุปริ มุทฺธนิ จเรน ภวิตพฺพํ. อิทํ มหาราช ฉตฺตสฺส ปฐมํ องฺคํ คเหตพฺพํ.}

\prose{ปุน จปรํ มหาราช ฉตฺตํ มุทฺธนุ\-ปตฺถมฺภํ โหติ, เอวเมว โข มหาราช โยคินา โยคาวจเรน โยนิโส มนสิการุ\-ปตฺถมฺเภน ภวิตพฺพํ. อิทํ มหาราช ฉตฺตสฺส ทุติยํ องฺคํ คเหตพฺพํ.}

\prose{ปุน จปรํ มหาราช ฉตฺตํ วาตาตป\-เมฆ\-วุฏฺฐิโย ปฏิหนติ, เอวเมว โข มหาราช โยคินา โยคาวจเรน นานาวิธ\-ทิฏฺฐิ\-ปุถุสมณพฺราหฺมณานํ มตวาต\csromansharedfootnote{436f36376451cc22}{มหาวาต (ก)} ติวิธคฺคิ\-สนฺตาป\-กิเลส\-วุฏฺฐิโย ปฏิหนฺตพฺพา. อิทํ มหาราช ฉตฺตสฺส ตติยํ องฺคํ คเหตพฺพํ. ภาสิตมฺเปตํ มหาราช เถเรน สาริปุตฺเตน ธมฺม\-เสนาปตินา\csromandash{}}

\setlength{\csromangathapagewidth}{0pt}
\setlength{\csromangathawakwidth}{0pt}
\csromangathameasuregroup{“ยถาปิ ฉตฺตํ วิปุลํ, \\ วาตาตปํ นิวาเรติ, \\ ตเถว พุทฺธปุตฺโตปิ, \\ กิเลสวุฏฺฐิํ วาเรติ,}{\csromangathabat{“ยถาปิ ฉตฺตํ วิปุลํ,}{อจฺฉิทฺทํ ถิรสํหิตํ.} \\ \csromangathabat{วาตาตปํ นิวาเรติ,}{มหตี เมฆวุฏฺฐิโย.} \\ \csromangathabat{ตเถว พุทฺธปุตฺโตปิ,}{สีลฉตฺตธโร สุจิ.} \\ \csromangathabat{กิเลสวุฏฺฐิํ วาเรติ,}{สนฺตาปติวิธคฺคโย”ติ.}}
\csromangathasetleft{“ยถาปิ ฉตฺตํ วิปุลํ, \\ วาตาตปํ นิวาเรติ, \\ ตเถว พุทฺธปุตฺโตปิ, \\ กิเลสวุฏฺฐิํ วาเรติ,}
\csromangathagroup{\csromangathastanza{\csromangathabat{“ยถาปิ ฉตฺตํ วิปุลํ,}{อจฺฉิทฺทํ ถิรสํหิตํ.} \\ \csromangathabat{วาตาตปํ นิวาเรติ,}{มหตี เมฆวุฏฺฐิโย.}}\csromangathastanza{\csromangathabat{ตเถว พุทฺธปุตฺโตปิ,}{สีลฉตฺตธโร สุจิ.} \\ \csromangathabat{กิเลสวุฏฺฐิํ วาเรติ,}{สนฺตาปติวิธคฺคโย”ติ.}}}

\vspace{\tipitakaparskip}
\csromancenter{ฉตฺตงฺคปโญฺห ตติโย.}
\csromanmatikaanchor{mtk.271}
\csromantocmarkref{subparagraph}{4. เขตฺตงฺคปญฺห}{mtk.271}
\bookmark[dest={mtk.271},level=5]{4. เขตฺตงฺคปญฺห}
\csromanmarkh{6. โอปมฺม\-กถา\-ปญฺห}\csromanmarkhii{4. เขตฺตงฺคปญฺห}\csromanheadernomark{5}{6. โอปมฺม\-กถา\-ปญฺห \textbf{4. เขตฺตงฺคปญฺห}}
\proseitem{4}{\csromanfolio{403}ภนฺเต นาคเสน “เขตฺตสฺส ตีณิ องฺคานิ คเหตพฺพานี”ติ ยํ วเทสิ, กตมานิ ตานิ ตีณิ องฺคานิ คเหตพฺพานีติ. ยถา มหาราช เขตฺตํ มาติกา\-สมฺปนฺนํ โหติ, เอวเมว โข มหาราช โยคินา โยคาวจเรน สุจริต\-วตฺต\-ปฺปฏิวตฺต\-มาติกาสมฺปนฺเนน ภวิตพฺพํ. อิทํ มหาราช เขตฺตสฺส ปฐมํ องฺคํ คเหตพฺพํ.}

\prose{ปุน จปรํ มหาราช เขตฺตํ มริยาทา\-สมฺปนฺนํ โหติ, ตาย จ มริยาทาย อุทกํ รกฺขิตฺวา ธญฺญํ ปริปาเจติ, เอวเมว โข มหาราช โยคินา โยคาวจเรน สีลหิริ\-มริยาทา\-สมฺปนฺเนน ภวิตพฺพํ, ตาย จ สีลหิริ\-มริยาทาย สามญฺญํ รกฺขิตฺวา จตฺตาริ สามญฺญผลานิ คเหตพฺพานิ. อิทํ มหาราช เขตฺตสฺส ทุติยํ องฺคํ คเหตพฺพํ.}

\prose{ปุน จปรํ มหาราช เขตฺตํ อุฏฺฐาน\-สมฺปนฺนํ โหติ, กสฺสกสฺส หาสชนกํ อปฺปมฺปิ พีชํ วุตฺตํ พหุ โหติ, พหุ วุตฺตํ พหุตรํ โหติ, เอวเมว โข มหาราช โยคินา โยคาวจเรน อุฏฺฐาน\-สมฺปนฺเนน วิปุล\-ผลทายินา ภวิตพฺพํ, ทายกานํ หาสชนเกน ภวิตพฺพํ, ยถา อปฺปํ ทินฺนํ พหุ โหติ, พหุ ทินฺนํ พหุตรํ โหติ. อิทํ มหาราช เขตฺตสฺส ตติยํ องฺคํ คเหตพฺพํ. ภาสิตมฺเปตํ มหาราช เถเรน อุปาลินา วินยธเรน\csromandash{}}

\setlength{\csromangathapagewidth}{0pt}
\setlength{\csromangathawakwidth}{0pt}
\csromangathameasuregroup{“เขตฺตู ปเมน ภวิตพฺพํ, \\ เอส เขตฺตวโร นาม,}{\csromangathabat{“เขตฺตู ปเมน ภวิตพฺพํ,}{อุฏฺฐานวิปุลทายินา.} \\ \csromangathabat{เอส เขตฺตวโร นาม,}{โย ททาติ วิปุลํ ผลนฺ”ติ.}}
\csromangathasetleft{“เขตฺตู ปเมน ภวิตพฺพํ, \\ เอส เขตฺตวโร นาม,}
\csromangathagroupwithcloserruled{\csromangathastanza{\csromangathabat{“เขตฺตู ปเมน ภวิตพฺพํ,}{อุฏฺฐานวิปุลทายินา.} \\ \csromangathabat{เอส เขตฺตวโร นาม,}{โย ททาติ วิปุลํ ผลนฺ”ติ.}}}{เขตฺตงฺคปโญฺห จตุตฺโถ.}
\csromanmatikaanchor{mtk.272}
\csromantocmarkref{subparagraph}{5. อคทงฺคปญฺห}{mtk.272}
\bookmark[dest={mtk.272},level=5]{5. อคทงฺคปญฺห}
\csromanheader{5}{5. \textbf{อคทงฺคปญฺห}}
\proseitem{5}{ภนฺเต นาคเสน “อคทสฺส เทฺว องฺคานิ คเหตพฺพานี”ติ ยํ วเทสิ, กตมานิ ตานิ เทฺว องฺคานิ คเหตพฺพานีติ. ยถา มหาราช อคเท กีมี น สณฺฐหนฺติ, เอวเมว โข มหาราช โยคินา โยคาวจเรน มานเส กิเลสา น สณฺฐเปตพฺพา. อิทํ มหาราช อคทสฺส ปฐมํ องฺคํ คเหตพฺพํ.}

\prose{\csromanfolio{404}ปุน จปรํ มหาราช อคโท ทฏฺฐผุฏฺฐทิฏฺฐ- อสิต\-ปีต\-ขายิต\-สายิตํ สพฺพํ วิสํ ปฏิหนติ, เอวเมว โข มหาราช โยคินา โยคาวจเรน ราค\-โทสโมหมาน\-ทิฏฺฐิวิสํ สพฺพํ ปฏิหนิตพฺพํ. อิทํ มหาราช อคทสฺส ทุติยํ องฺคํ คเหตพฺพํ. ภาสิตมฺเปตํ มหาราช ภควตา เทวาติเทเวน\csromandash{}}

\setlength{\csromangathapagewidth}{0pt}
\setlength{\csromangathawakwidth}{0pt}
\csromangathameasuregroup{“สงฺขารานํ สภาวตฺถํ, \\ อคเทเนว โหตพฺพํ,}{\csromangathabat{“สงฺขารานํ สภาวตฺถํ,}{ทฏฺฐุกาเมน โยคินา.} \\ \csromangathabat{อคเทเนว โหตพฺพํ,}{กิเลสวิสนาสเน”ติ.}}
\csromangathasetleft{“สงฺขารานํ สภาวตฺถํ, \\ อคเทเนว โหตพฺพํ,}
\csromangathagroupwithcloserruled{\csromangathastanza{\csromangathabat{“สงฺขารานํ สภาวตฺถํ,}{ทฏฺฐุกาเมน โยคินา.} \\ \csromangathabat{อคเทเนว โหตพฺพํ,}{กิเลสวิสนาสเน”ติ.}}}{อคทงฺคปโญฺห ปญฺจโม.}
\csromanmatikaanchor{mtk.273}
\csromantocmarkref{subparagraph}{6. โภชนงฺคปญฺห}{mtk.273}
\bookmark[dest={mtk.273},level=5]{6. โภชนงฺคปญฺห}
\csromanheader{5}{6. \textbf{โภชน}\-\textbf{งฺค}\-\textbf{ปญฺห}}
\proseitem{6}{ภนฺเต นาคเสน “โภชนสฺส ตีณิ องฺคานิ คเหตพฺพานี”ติ ยํ วเทสิ, กตมานิ ตานิ ตีณิ องฺคานิ คเหตพฺพานีติ. ยถา มหาราช โภชนํ สพฺพสตฺตานํ อุปตฺถมฺโภ, เอวเมว โข มหาราช โยคินา โยคาวจเรน สพฺพสตฺตานํ มคฺคุ\-ปตฺถมฺเภน ภวิตพฺพํ. อิทํ มหาราช โภชนสฺส ปฐมํ องฺคํ คเหตพฺพํ.}

\prose{ปุน จปรํ มหาราช โภชนํ สพฺพสตฺตานํ พลํ วฑฺเฒติ, เอวเมว โข มหาราช โยคินา โยคาวจเรน ปุญฺญวฑฺฒิยา วฑฺฒิตพฺพํ. อิทํ มหาราช โภชนสฺส ทุติยํ องฺคํ คเหตพฺพํ.}

\prose{ปุน จปรํ โภชนํ สพฺพสตฺตานํ อภิปตฺถิตํ, เอวเมว โข มหาราช โยคินา โยคาวจเรน สพฺพโลกา\-ภิปตฺถิเตน ภวิตพฺพํ. อิทํ มหาราช โภชนสฺส ตติยํ องฺคํ คเหตพฺพํ. ภาสิตมฺเปตํ มหาราช เถเรน มหา\-โมคฺคลฺลาเนน\csromandash{}}

\setlength{\csromangathapagewidth}{0pt}
\setlength{\csromangathawakwidth}{0pt}
\csromangathameasuregroup{“สํยเมน นิยเมน, \\ ปตฺถิเตน ภวิตพฺพํ,}{\csromangathabat{“สํยเมน นิยเมน,}{สีเลน ปฏิปตฺติยา.} \\ \csromangathabat{ปตฺถิเตน ภวิตพฺพํ,}{สพฺพโลกสฺส โยคินา”ติ.}}
\csromangathasetleft{“สํยเมน นิยเมน, \\ ปตฺถิเตน ภวิตพฺพํ,}
\csromangathagroup{\csromangathastanza{\csromangathabat{“สํยเมน นิยเมน,}{สีเลน ปฏิปตฺติยา.} \\ \csromangathabat{ปตฺถิเตน ภวิตพฺพํ,}{สพฺพโลกสฺส โยคินา”ติ.}}}

\vspace{\tipitakaparskip}
\csromancenter{โภชน\-งฺค\-ปโญฺห ฉฏฺโฐ.}
\csromanmatikaanchor{mtk.274}
\csromantocmarkref{subparagraph}{7. อิสฺสาสงฺคปญฺห}{mtk.274}
\bookmark[dest={mtk.274},level=5]{7. อิสฺสาสงฺคปญฺห}
\csromanmarkh{6. โอปมฺม\-กถา\-ปญฺห}\csromanmarkhii{7. อิสฺสาสงฺคปญฺห}\csromanheadernomark{5}{6. โอปมฺม\-กถา\-ปญฺห \textbf{7. อิสฺสาสงฺคปญฺห}}
\proseitem{7}{\csromanfolio{405}ภนฺเต นาคเสน “อิสฺสาสสฺส จตฺตาริ องฺคานิ คเหตพฺพานี”ติ ยํ วเทสิ, กตมานิ ตานิ จตฺตาริ องฺคานิ คเหตพฺพานีติ. ยถา มหาราช อิสฺสาโส สเร ปาตยนฺโต อุโภ ปาเท ปถวิยํ ทฬฺหํ ปติฏฺฐาเปติ, ชณฺณุอเวกลฺลํ กโรติ, สรกลาปํ กฏิสนฺธิมฺหิ ฐเปติ, กายํ อุปตฺถทฺธํ กโรติ, เทฺว หตฺเถ สนฺธิฏฺฐานํ อาโรเปติ, มุฏฺฐิํ ปีฬยติ, องฺคุลิโย นิรนฺตรํ กโรติ, คีวํ ปคฺคณฺหาติ, จกฺขูนิ มุขญฺจ ปิทหติ, นิมิตฺตํ อุชุํ กโรติ, หาสมุปฺปาเทติ “วิชฺฌิสฺสามี”ติ. เอวเมว โข มหาราช โยคินา โยคาวจเรน สีลปถวิยํ วีริยปาเท ปติฏฺฐาเป\-ตพฺพํ, ขนฺติโสรจฺจํ อเวกลฺลํ กาตพฺพํ, สํวเร จิตฺตํ ฐเปตพฺพํ, สํยมนิยเม อตฺตา อุปเนตพฺโพ, อิจฺฉา มุจฺฉา ปีฬยิตพฺพา, โยนิโส มนสิกาเร จิตฺตํ นิรนฺตรํ กาตพฺพํ, วีริยํ ปคฺคเหตพฺพํ, ฉ ทฺวารา ปิทหิตพฺพา, สติ อุปฏฺฐเปตพฺพา, หาสมุ\-ปฺปาเทตพฺพํ “สพฺพกิเลเส ญาณนาราเจน วิชฺฌิสฺสามี”ติ. อิทํ มหาราช อิสฺสาสสฺส ปฐมํ องฺคํ คเหตพฺพํ.}

\prose{ปุน จปรํ มหาราช อิสฺสาโส อาฬกํ ปริหรติ วงฺก\-ชิมฺห\-กุฏิล\-นาราจสฺส อุชุกรณาย, เอวเมว โข มหาราช โยคินา โยคาวจเรน อิมสฺมิํ กาเย สติปฏฺฐาน\-อาฬกํ ปริหริตพฺพํ วงฺก\-ชิมฺห\-กุฏิล\-จิตฺตสฺส อุชุกรณาย. อิทํ มหาราช อิสฺสาสสฺส ทุติยํ องฺคํ คเหตพฺพํ.}

\prose{ปุน จปรํ มหาราช อิสฺสาโส ลกฺเข อุปาเสติ, เอวเมว โข มหาราช โยคินา โยคาวจเรน อิมสฺมิํ กาเย อุปาสิตพฺพํ. กถํ มหาราช โยคินา โยคาวจเรน อิมสฺมิํ กาเย อุปาสิตพฺพํ,\csromansymbolfootnote{*}{ขุ 7. 40; ขุ 8. 38; ขุ 9. 411 ปิฏฺเฐสุปิ.}อนิจฺจโต อุปาสิตพฺพํ, ทุกฺขโต อุปาสิตพฺพํ, อนตฺตโต อุปาสิตพฺพํ, โรคโต ฯเปฯ คณฺฑโต. สลฺลโต. อฆโต. อาพาธโต. ปรโต. ปโลกโต. ¢ติโต. อุปทฺทวโต. ภยโต. อุปสคฺคโต. จลโต. ปภงฺคุโต. อทฺธุวโต. อตาณโต. อเลณโต. อสรณโต. ริตฺตโต. ตุจฺฉโต. สุญฺญโต. อาทีนวโต. วิปริณาม\-ธมฺมโต. อสารโต. อฆมูลโต. วธกโต. วิภวโต. สาสวโต. สงฺขตโต. มารามิสโต. ชาติธมฺมโต. ชราธมฺมโต. พฺยาธิธมฺมโต. มรณธมฺมโต. โสกธมฺมโต. ปริเทว\-ธมฺมโต. อุปายา\-สธมฺมโต. สํกิเลส\-ธมฺมโต. เอวํ โข มหาราช โยคินา \csromanfolio{406}โยคาวจเรน อิมสฺมิํ กาเย อุปาสิตพฺพํ. อิทํ มหาราช อิสฺสาสสฺส ตติยํ องฺคํ คเหตพฺพํ.}

\prose{ปุน จปรํ มหาราช อิสฺสาโส สายํ ปาตํ อุปาสิติ, เอวเมว โข มหาราช โยคินา โยคาวจเรน สายํ ปาตํ อารมฺมเณ อุปาสิตพฺพํ. อิทํ มหาราช อิสฺสาสสฺส จตุตฺถํ องฺคํ คเหตพฺพํ. ภาสิตมฺเปตํ มหาราช เถเรน สาริปุตฺเตน ธมฺม\-เสนาปตินา\csromandash{}}

\setlength{\csromangathapagewidth}{0pt}
\setlength{\csromangathawakwidth}{0pt}
\csromangathameasuregroup{“ยถา อิสฺสาสโก นาม, \\ อุปาสนํ อริญฺจนฺโต\textsuperscript{1}, \\ ตเถว พุทฺธปุตฺโตปิ, \\ กายุปาสนํ อริญฺจนฺโต,}{\csromangathabat{“ยถา อิสฺสาสโก นาม,}{สายํ ปาตํ อุปาสติ.} \\ \csromangathabat{อุปาสนํ อริญฺจนฺโต\textsuperscript{1},}{ลภเต ภตฺตเวตนํ.} \\ \csromangathabat{ตเถว พุทฺธปุตฺโตปิ,}{กโรติ กายุปาสนํ.} \\ \csromangathabat{กายุปาสนํ อริญฺจนฺโต,}{อรหตฺตมธิคจฺฉตี”ติ.}}
\csromangathasetleft{“ยถา อิสฺสาสโก นาม, \\ อุปาสนํ อริญฺจนฺโต\textsuperscript{1}, \\ ตเถว พุทฺธปุตฺโตปิ, \\ กายุปาสนํ อริญฺจนฺโต,}
\csromangathagroup{\csromangathastanza{\csromangathabat{“ยถา อิสฺสาสโก นาม,}{สายํ ปาตํ อุปาสติ.} \\ \csromangathabat{อุปาสนํ อริญฺจนฺโต\csromansharedfootnote{4aecf1430b278354}{น ริจฺฉนฺโต (สี, ก)},}{ลภเต ภตฺตเวตนํ.}}\csromangathastanza{\csromangathabat{ตเถว พุทฺธปุตฺโตปิ,}{กโรติ กายุปาสนํ.} \\ \csromangathabat{กายุปาสนํ อริญฺจนฺโต,}{อรหตฺตมธิคจฺฉตี”ติ.}}}

\vspace{\tipitakaparskip}
\csromancenter{อิสฺสาสงฺคปโญฺห สตฺตโม.}
\nitthitammidruled{กุมฺภวคฺโค สตฺตโม\csromansharedfootnote{eb66d931022fe143}{อิโต ปรํ ราชงฺคปญฺหาทิกา อฏฺฐติํส ปญฺหา วินฏฺฐา, เยหิ ตา ทิฏฺฐา, เตหิ โน อาโรเจตพฺพา ปุน มุทฺทาปนกาเล ปกฺขิปนตฺถายาติ (น, พุ, ส)}.}
\csromancenter{\textbf{ตสฺสุทฺทานํ}}
\setlength{\csromangathapagewidth}{0pt}
\setlength{\csromangathawakwidth}{0pt}
\csromangathameasuregroup{กุมฺโภ จ กาฬายโส จ, \\ โภชเนน จ อิสฺสาโส,}{\csromangathabat{กุมฺโภ จ กาฬายโส จ,}{ฉตฺตํ เขตฺตญฺจ อคโท.} \\ \csromangathabat{โภชเนน จ อิสฺสาโส,}{วุตฺตํ ทานิ วิทูหีติ.}}
\csromangathasetleft{กุมฺโภ จ กาฬายโส จ, \\ โภชเนน จ อิสฺสาโส,}
\csromangathagroupwithcloser{\csromangathastanza{\csromangathabat{กุมฺโภ จ กาฬายโส จ,}{ฉตฺตํ เขตฺตญฺจ อคโท.} \\ \csromangathabat{โภชเนน จ อิสฺสาโส,}{วุตฺตํ ทานิ วิทูหีติ.}}}{\textbf{โอปมฺม}\-\textbf{กถา}\-\textbf{ปโญฺห นิฏฺฐิโต.}}
\csromanheader{2}{นิคมน \textbf{นิคมน}}
\prose{\csromanfolio{407}อิติ ฉสุ กณฺเฑสุ พาวีสติวคฺค\-ปติ\-มณฺฑิเตสุ ทฺวาสฏฺฐิ\-อธิกา เทฺวสตา อิมสฺมิํ โปตฺถเก อาคตา มิลินฺทปญฺหา สมตฺตา, อนาคตา จ ปน ทฺวาจตฺตาลีสา โหนฺติ, อาคตา จ อนาคตา จ สพฺพา สโมธาเนตฺวา จตูหิ อธิกา ติสตปญฺหา โหนฺติ, สพฺพาว มิลินฺทปญฺหาติ สงฺขํ คจฺฉนฺติ.}

\prose{รญฺโญ จ เถรสฺส จ ปุจฺฉา\-วิสชฺช\-นา\-วสาเน จตุราสีติ\-สตสหสฺส\-โยชน\-พหลา อุทก\-ปริยนฺตํ กตฺวา อยํ มหาปถวี ฉธา กมฺปิตฺถ, วิชฺชุลฺลตา นิจฺฉริํสุ, เทวตา ทิพฺพ\-ปุปฺผ\-วสฺสํ ปวสฺสิํสุ, มหาพฺรหฺมา สาธุการม\-ทาสิ, มหาสมุทฺท\-กุจฺฉิยํ เมฆตฺถนิตนิคฺฆาโส วิย มหาโฆโส อโหสิ, อิติ โส มิลินฺโท ราชา จ โอโรธคณา จ สิรสา อญฺชลิํ ปณาเมตฺวา วนฺทิํสุ.}

\prose{มิลินฺโท ราชา อติวิย ปมุทิตหทโย สุมถิต\-มาน\-หทโย พุทฺธสาสเน สารมติโน รตนตฺตเย สุนิกฺกงฺโข นิคฺคุมฺโพ นิตฺถทฺโธ หุตฺวา เถรสฺส คุเณสุ ปพฺพชฺชาสุ ปฏิปทา\-อิริยาปเถสุ จ อติวิย ปสนฺโน วิสฺสตฺโถ นิราลโย นิหตมาน\-ตฺถมฺโภ อุทฺธฏทาโฐ วิย ภุชคินฺโท เอวมาห “สาธุ ภนฺเต นาคเสน พุทฺธวิสโย ปโญฺห ตยา วิสชฺชิโต, อิมสฺมิํ พุทฺธสาสเน ฐเปตฺวา ธมฺม\-เสนาปติํ สาริปุตฺตตฺเถรํ อญฺโญ ตยา สทิโส ปญฺหวิสชฺชเน นตฺถิ, ขมถ ภนฺเต นาคเสน มม อจฺจยํ, อุปาสกํ มํ ภนฺเต นาคเสน ธาเรถ อชฺชตคฺเค ปาณุเปตํ สรณํ คตนฺ”ติ.}

\prose{ตทา ราชา สห พลกาเยหิ นาคเสน\-ตฺเถรํ ปยิรุปาสิตฺวา มิลินฺทํ นาม วิหารํ กาเรตฺวา เถรสฺส นิยฺยาเตตฺวา จตูหิ ปจฺจเยหิ นาคเสนํ, โกฏิสเตหิ ภิกฺขูหิ สทฺธิํ ปริจริ, ปุนปิ เถรสฺส ปญฺญาย ปสีทิตฺวา ปุตฺตสฺส รชฺชํ นิยฺยาเตตฺวา อคารสฺมา อนคาริยํ ปพฺพชิตฺวา วิปสฺสนํ วฑฺเฒตฺวา อรหตฺตํ ปาปุณิ, เตน วุตฺตํ\csromandash{}}

\setlength{\csromangathapagewidth}{0pt}
\setlength{\csromangathawakwidth}{0pt}
\csromangathameasuregroup{“ปญฺญา ปสตฺถา โลกสฺมิํ, \\ ปญฺญาย วิมติํ หนฺตฺวา, \\ ยสฺมิํ ขนฺเธ ฐิตา ปญฺญา, \\ ปูชา วิเสสสฺสาธาโร, \\ ตสฺมา หิ ปณฺฑิโต โปโส, \\ ปญฺญวนฺตํภิปูเชยฺย, \\ ลงฺกายํ โทณินคเร, \\ มหาเถเรน เลขิตฺวา, \\ มิลินฺทราชปโญฺห จ, \\ มิลินฺโท หิ มหาปญฺโญ, \\ อิมินา ปุญฺญกมฺเมน, \\ เมตฺเตยฺยํ’นาคเต ปสฺเส,}{\csromangathabat{“ปญฺญา ปสตฺถา โลกสฺมิํ,}{กตา สทฺธมฺมฏฺฐิติยา.} \\ \csromangathabat{ปญฺญาย วิมติํ หนฺตฺวา,}{สนฺติํ ปปฺโปนฺติ ปณฺฑิตา.} \\ \csromangathabat{ยสฺมิํ ขนฺเธ ฐิตา ปญฺญา,}{สติ ตตฺถ อนูนกา.} \\ \csromangathabat{ปูชา วิเสสสฺสาธาโร,}{อคฺโค เสฏฺโฐ\textsuperscript{1} อนุตฺตโร.} \\ \csromangathabat{ตสฺมา หิ ปณฺฑิโต โปโส,}{สมฺปสฺสํ หิตมตฺตโน\textsuperscript{1}.} \\ \csromangathabat{ปญฺญวนฺตํภิปูเชยฺย,}{เจติยํ วิย สาทโร”ติ\textsuperscript{1}.} \\ \csromangathabat{ลงฺกายํ โทณินคเร,}{วสตา โทณินามินา.} \\ \csromangathabat{มหาเถเรน เลขิตฺวา,}{สุฏฺฐปิตํ ยถาสุตํ.} \\ \csromangathabat{มิลินฺทราชปโญฺห จ,}{นาคเสนวิสชฺชนํ.} \\ \csromangathabat{มิลินฺโท หิ มหาปญฺโญ,}{นาคเสโน สุปณฺฑิโต.} \\ \csromangathabat{อิมินา ปุญฺญกมฺเมน,}{อิโต คจฺฉามิ ตุสฺสิตํ.} \\ \csromangathabat{เมตฺเตยฺยํ’นาคเต ปสฺเส,}{สุเณยฺยํ ธมฺมมุตฺตมนฺติ.}}
\csromangathasetleft{“ปญฺญา ปสตฺถา โลกสฺมิํ, \\ ปญฺญาย วิมติํ หนฺตฺวา, \\ ยสฺมิํ ขนฺเธ ฐิตา ปญฺญา, \\ ปูชา วิเสสสฺสาธาโร, \\ ตสฺมา หิ ปณฺฑิโต โปโส, \\ ปญฺญวนฺตํภิปูเชยฺย, \\ ลงฺกายํ โทณินคเร, \\ มหาเถเรน เลขิตฺวา, \\ มิลินฺทราชปโญฺห จ, \\ มิลินฺโท หิ มหาปญฺโญ, \\ อิมินา ปุญฺญกมฺเมน, \\ เมตฺเตยฺยํ’นาคเต ปสฺเส,}
\csromangathagroupwithcloser{\csromangathastanza{\csromangathabat{“ปญฺญา ปสตฺถา โลกสฺมิํ,}{กตา สทฺธมฺม\-ฏฺฐิติยา.} \\ \csromangathabat{ปญฺญาย วิมติํ หนฺตฺวา,}{สนฺติํ ปปฺโปนฺติ ปณฺฑิตา.}}\csromangathastanza{\csromanfolio{408}\csromangathabat{ยสฺมิํ ขนฺเธ ฐิตา ปญฺญา,}{สติ ตตฺถ อนูนกา.} \\ \csromangathabat{ปูชา วิเสสสฺสาธาโร,}{อคฺโค เสฏฺโฐ\csromansharedfootnote{0ff480be6c481f01}{โสว (อิ)} อนุตฺตโร.}}\csromangathastanza{\csromangathabat{ตสฺมา หิ ปณฺฑิโต โปโส,}{สมฺปสฺสํ หิตมตฺตโน\csromansharedfootnote{99cb5a4cad9bfd20}{อตฺถมตฺตโน (อิ)}.} \\ \csromangathabat{ปญฺญวนฺตํภิปูเชยฺย,}{เจติยํ วิย สาทโร”ติ\csromansharedfootnote{99cb5a4cad9bfd20}{อตฺถมตฺตโน (อิ)}.}}\csromangathastanza{\csromangathabat{ลงฺกายํ โทณินคเร,}{วสตา โทณินามินา.} \\ \csromangathabat{มหาเถเรน เลขิตฺวา,}{สุฏฺฐปิตํ ยถาสุตํ.}}\csromangathastanza{\csromangathabat{มิลินฺทราชปโญฺห จ,}{นาคเสนวิสชฺชนํ.} \\ \csromangathabat{มิลินฺโท หิ มหาปญฺโญ,}{นาคเสโน สุปณฺฑิโต.}}\csromangathastanza{\csromangathabat{อิมินา ปุญฺญกมฺเมน,}{อิโต คจฺฉามิ ตุสฺสิตํ.} \\ \csromangathabat{เมตฺเตยฺยํ’นาคเต ปสฺเส,}{สุเณยฺยํ ธมฺมมุตฺตมนฺติ.}}}{\textbf{มิลินฺทปโญฺห นิฏฺฐิโต.}}
\orphannote{อํ 2. 481 ปิฏฺเฐ.}

\orphannote{สํ 3. 100 ปิฏฺเฐ.}

\orphannote{สํ 2. 12, 302; สํ 3. 363; อํ 3. 517; ขุ 10. 50 ปิฏฺเฐสุ.}

\orphannote{ม 3. 244 ปิฏฺเฐ.}

\orphannote{ม 3. 221; อํ 1. 140 ปิฏฺเฐสุ.}

\orphannote{ม 3. 221; อํ 1. 140 ปิฏฺเฐสุ.}

\orphannote{ม 1. 16 ปิฏฺเฐ.}

\orphannote{ขุ 4. 395 ปิฏฺเฐ.}

\orphannote{ขุ 1. 283 ปิฏฺเฐ.}

\orphannote{ขุ 5. 220 ปิฏฺเฐ.}

\orphannote{ที 3. 134 ปิฏฺเฐ.}

\orphannote{สํ 1. 371 ปิฏฺเฐ.}

\orphannote{ขุ 1. 322 ปิฏฺเฐ.}

\orphannote{ขุ 1. 390 ปิฏฺเฐ.}

